% Generated mechanically from the frozen algebra.tex target.
% Do not edit this generated file; edit the bound locale source instead.

% OK, start here.
%

\setcounter{chapter}{9}
\chapter{交换代数}



\phantomsection
\label{algebra-section-phantom}





\section{引言}
\label{algebra-section-introduction}

\noindent
本文将介绍交换代数的基础知识。可参阅 \cite{MatCA}。






\section{约定}
\label{algebra-section-conventions}

\noindent
环均指带 $1$ 的交换环。零环也是环；事实上，它是唯一没有素理想的环。
我们将使用 Kronecker 符号 $\delta_{ij}$。若 $R \to S$ 是环同态，且
$\mathfrak q$ 是 $S$ 的一个素理想，则用记号
``$\mathfrak p = R \cap \mathfrak q$''
表示 $\mathfrak q$ 在 $R \to S$ 下的逆像素理想，即使 $R$ 不是 $S$ 的子环，
甚至即使 $R \to S$ 不是单射，也采用这一记号。






\section{基本概念}
\label{algebra-section-rings-basic}

\noindent
下面列出交换代数中的一些基本概念。其中有些概念将在后文作更详细的讨论，
有些在此表中定义，其余则视为基础概念而不再定义。若你不熟悉以下大多数斜体概念，
建议先阅读一本代数入门教材，再继续阅读。

\begin{enumerate}
\item $R$ 是一个{\it 环}，
\label{algebra-item-ring}
\item $x\in R$ 是{\it 幂零元}，
\label{algebra-item-ring-element-nilpotent}
\item $x\in R$ 是{\it 零因子}，
\label{algebra-item-ring-element-zerodivisor}
\item $x\in R$ 是{\it 可逆元}，
\label{algebra-item-ring-element-unit}
\item $e \in R$ 是{\it 幂等元}，
\label{algebra-item-ring-element-idempotent}
\item 若幂等元 $e \in R$ 满足 $e = 1$ 或 $e = 0$，则称其为{\it 平凡的}，
\label{algebra-item-idempotent-trivial}
\item $\varphi : R_1 \to R_2$ 是{\it 环同态}，
\label{algebra-item-ring-homomorphism}
\item
\label{algebra-item-ring-homomorphism-finite-presentation}
$\varphi : R_1 \to R_2$ 是{\it 有限表示的}，亦即
{\it $R_2$ 是有限表示的 $R_1$-代数}，
见定义 \ref{algebra-definition-finite-type}，
\item
\label{algebra-item-ring-homomorphism-finite-type}
$\varphi : R_1 \to R_2$ 是{\it 有限型的}，亦即
{\it $R_2$ 是有限型 $R_1$-代数}，
见定义 \ref{algebra-definition-finite-type}，
\item
\label{algebra-item-ring-homomorphism-finite}
$\varphi : R_1 \to R_2$ 是{\it 有限的}，亦即
{\it $R_2$ 是有限 $R_1$-代数}，
\item $R$ 是{\it （整）环}，
\label{algebra-item-ring-domain}
\item $R$ 是{\it 约化的}，
\label{algebra-item-ring-reduced}
\item $R$ 是{\it Noether 环}，
\label{algebra-item-ring-Noetherian}
\item $R$ 是{\it 主理想整环}，简称 {\it PID}，
\label{algebra-item-ring-PID}
\item $R$ 是{\it 欧几里得整环}，
\label{algebra-item-ring-Euclidean}
\item $R$ 是{\it 唯一分解整环}，简称 {\it UFD}，
\label{algebra-item-ring-UFD}
\item $R$ 是{\it 离散赋值环}，简称 {\it dvr}，
\label{algebra-item-ring-dvr}
\item $K$ 是{\it 域}，
\label{algebra-item-field}
\item $L/K$ 是{\it 域扩张}，
\label{algebra-item-field-extension}
\item $L/K$ 是{\it 代数域扩张}，
\label{algebra-item-field-extension-algebraic}
\item $\{t_i\}_{i\in I}$ 是 $L$ 在 $K$ 上的{\it 超越基}，
\label{algebra-item-transcendence-basis}
\item $L$ 在 $K$ 上的{\it 超越次数} $\text{trdeg}(L/K)$，
\label{algebra-item-transcendence-degree}
\item 域 $k$ 是{\it 代数闭的}，
\label{algebra-item-algebraically-closed}
\item
\label{algebra-item-extend-into-algebraically-closed}
若 $L/K$ 是代数扩张，且 $\Omega/K$ 是扩张而 $\Omega$
为代数闭域，则存在环同态 $L \to \Omega$，它延拓 $K$ 上给定的同态，
\item $I \subset R$ 是{\it 理想}，
\label{algebra-item-ideal}
\item $I \subset R$ 是{\it 根理想}，
\label{algebra-item-ideal-radical}
\item 若 $I$ 是理想，则有其{\it 根} $\sqrt{I}$，
\label{algebra-item-radical-ideal}
\item
\label{algebra-item-ideal-nilpotent}
$I \subset R$ 为{\it 幂零理想}，意指对某个 $n \in \mathbf{N}$ 有 $I^n = 0$，
\item
\label{algebra-item-ideal-locally-nilpotent}
$I \subset R$ 为{\it 局部幂零理想}，意指 $I$ 的每个元素均为幂零元，
\item $\mathfrak p \subset R$ 是{\it 素理想}，
\label{algebra-item-prime-ideal}
\item
\label{algebra-item-prime-product-ideals}
若 $\mathfrak p \subset R$ 是素理想，$I, J \subset R$ 是理想，且
$IJ\subset \mathfrak p$，则 $I \subset \mathfrak p$ 或 $J \subset \mathfrak p$。
\item $\mathfrak m \subset R$ 是{\it 极大理想}，
\label{algebra-item-maximal-ideal}
\item 任一非零环都有极大理想，
\label{algebra-item-exists-maximal-ideal}
\item
\label{algebra-item-jacobson-radical}
$R$ 的{\it Jacobson 根}是 $R$ 的所有极大理想之交
$\text{rad}(R) = \bigcap_{\mathfrak m \subset R} \mathfrak m$，
\item 由子集 $T \subset R$ {\it 生成}的理想 $(T)$，
\label{algebra-item-ideal-generated-by}
\item {\it 商环} $R/I$，
\label{algebra-item-quotient-ring}
\item 环 $R$ 中的理想 $I$ 是素理想，当且仅当 $R/I$ 是整环，
\label{algebra-item-characterize-prime-ideal}
\item
\label{algebra-item-characterize-maximal-ideal}
环 $R$ 中的理想 $I$ 是极大理想，当且仅当环 $R/I$ 是域，
\item
\label{algebra-item-inverse-image-ideal}
若 $\varphi : R_1 \to R_2$ 是环同态，且 $I \subset R_2$ 是理想，
则 $\varphi^{-1}(I)$ 是 $R_1$ 的理想，
\item
\label{algebra-item-image-ideal}
若 $\varphi : R_1 \to R_2$ 是环同态，且 $I \subset R_1$ 是理想，
则 $\varphi(I) \cdot R_2$（有时记作 $I \cdot R_2$ 或 $IR_2$）
是由 $\varphi(I)$ 生成的 $R_2$ 的理想，
\item
\label{algebra-item-inverse-image-prime}
若 $\varphi : R_1 \to R_2$ 是环同态，且 $\mathfrak p \subset R_2$
是素理想，则 $\varphi^{-1}(\mathfrak p)$ 是 $R_1$ 的素理想，
\item $M$ 是{\it $R$-模}，
\label{algebra-item-module}
\item
\label{algebra-item-annihilator}
对 $m \in M$，元素 $m$ 在 $R$ 中的{\it 零化子}为
$I = \{f \in R \mid fm = 0\}$，
\item $N \subset M$ 是{\it $R$-子模}，
\label{algebra-item-submodule}
\item $M$ 是{\it Noether $R$-模}，
\label{algebra-item-Noetherian-module}
\item $M$ 是{\it 有限 $R$-模}，
\label{algebra-item-finite-module}
\item $M$ 是{\it 有限生成 $R$-模}，
\label{algebra-item-finitely-generated-module}
\item $M$ 是{\it 有限表示 $R$-模}，
\label{algebra-item-finitely-presented-module}
\item $M$ 是{\it 自由 $R$-模}，
\label{algebra-item-free-module}
\item
\label{algebra-item-extension-free}
若 $0 \to K \to L \to M \to 0$ 是 $R$-模的短正合列，且 $K$、$M$
均为自由模，则 $L$ 也是自由模，
\item 若 $N \subset M \subset L$ 是 $R$-模，则 $L/M = (L/N)/(M/N)$，
\label{algebra-item-isomorphism-theorem}
\item $S$ 是 {\it $R$ 的乘法子集}，
\label{algebra-item-multiplicative-subset}
\item $R$ 的{\it 局部化} $R \to S^{-1}R$，
\label{algebra-item-localization-ring}
\item
\label{algebra-item-localization-zero}
若 $R$ 是环，$S$ 是 $R$ 的乘法子集，则 $S^{-1}R$ 是零环，
当且仅当 $S$ 含有 $0$，
\item
\label{algebra-item-localize-nonzerodivisors}
若 $R$ 是环，且乘法子集 $S$ 完全由非零因子组成，
则 $R \to S^{-1}R$ 是单射，
\item 若 $\varphi : R_1 \to R_2$ 是环同态，且 $S$ 是 $R_1$ 的乘法子集，
则 $\varphi(S)$ 是 $R_2$ 的乘法子集，
\item
\label{algebra-item-products-multiplicative-subsets}
若 $S$、$S'$ 是 $R$ 的乘法子集，并以 $SS'$ 表示乘积集合
$SS' = \{r \in R \mid \exists s\in S, \exists s' \in S', r = ss'\}$，
则 $SS'$ 是 $R$ 的乘法子集，
\item
\label{algebra-item-localization-localization}
若 $S$、$S'$ 是 $R$ 的乘法子集，且 $\overline{S}$ 表示 $S$ 在
$(S')^{-1}R$ 中的像，则
$(SS')^{-1}R = \overline{S}^{-1}((S')^{-1}R)$，
\item $R$-模 $M$ 的{\it 局部化} $S^{-1}M$，
\label{algebra-item-localization-module}
\item
\label{algebra-item-localization-exact}
函子 $M \mapsto S^{-1}M$ 保持单射、满射和正合性，
\item
\label{algebra-item-localization-localization-module}
若 $S$、$S'$ 是 $R$ 的乘法子集，且 $M$ 是 $R$-模，则
$(SS')^{-1}M = S^{-1}((S')^{-1}M)$，
\item
\label{algebra-item-localize-ideal}
若 $R$ 是环，$I$ 是 $R$ 的理想，$S$ 是 $R$ 的乘法子集，
则 $S^{-1}I$ 是 $S^{-1}R$ 的理想，并且
$S^{-1}R/S^{-1}I = \overline{S}^{-1}(R/I)$，其中 $\overline{S}$
是 $S$ 在 $R/I$ 中的像，
\item
\label{algebra-item-ideal-in-localization}
若 $R$ 是环，$S$ 是 $R$ 的乘法子集，则 $S^{-1}R$ 的任一理想 $I'$
均形如 $S^{-1}I$，其中可取 $I$ 为 $I'$ 在 $R$ 中的逆像，
\item
\label{algebra-item-submodule-in-localization}
若 $R$ 是环，$M$ 是 $R$-模，$S$ 是 $R$ 的乘法子集，则
$S^{-1}M$ 的任一子模 $N'$ 均形如 $S^{-1}N$，其中 $N \subset M$
是某个子模；并且可取 $N$ 为 $N'$ 在 $M$ 中的逆像，
\item 若 $S = \{1, f, f^2, \ldots\}$，则 $R_f = S^{-1}R$ 且 $M_f = S^{-1}M$，
\label{algebra-item-localize-f}
\item
\label{algebra-item-localize-p}
若对某个素理想 $\mathfrak p$ 有
$S = R \setminus \mathfrak p = \{x\in R \mid x\not\in \mathfrak p\}$，
则通常记 $R_{\mathfrak p} = S^{-1}R$ 且
$M_{\mathfrak p} = S^{-1}M$，
\item {\it 局部环}是恰有一个极大理想的环，
\label{algebra-item-local-ring}
\item {\it 半局部环}是只有有限多个极大理想的环，
\label{algebra-item-semi-local-ring}
\item
\label{algebra-item-localize-p-local-ring}
若 $\mathfrak p$ 是 $R$ 中的素理想，则 $R_{\mathfrak p}$ 是局部环，
其极大理想为 $\mathfrak p R_{\mathfrak p}$，
\item
\label{algebra-item-residue-field}
环 $R$ 中素理想 $\mathfrak p$ 的{\it 剩余域}记作 $\kappa(\mathfrak p)$，
它是整环 $R/\mathfrak p$ 的分式域；并且等于
$R_\mathfrak p/\mathfrak pR_\mathfrak p
= (R \setminus \mathfrak p)^{-1}R/\mathfrak p$，
\item 给定 $R$ 以及 $M_1$、$M_2$，有{\it 张量积} $M_1 \otimes_R M_2$，
\label{algebra-item-tensor-product}
\item
\label{algebra-item-cauchy-binet}
给定环 $R$ 中分别为 $m \times n$ 和 $n \times m$ 的矩阵 $A$、$B$，
则在 $R$ 中有下式：
$$
\det(AB) = \sum \det(A_S)\det({}_SB)
$$
此处求和遍历
所有大小为 $m$ 的子集 $S \subset \{1, \ldots, n\}$，$A_S$ 是由 $A$
中对应于 $S$ 的列构成的 $m \times m$ 子矩阵，而 ${}_SB$ 是由 $B$
中对应于 $S$ 的行构成的 $m \times m$ 子矩阵，
\item 等等。
\end{enumerate}




\section{蛇形引理}
\label{algebra-section-snake}

\noindent
蛇形引理及其各种变体在同调一章的 Abel 范畴框架中讨论，
见同调，第 \ref{homology-section-abelian-categories} 节。

\begin{lemma}
\label{algebra-lemma-snake}
\begin{reference}
\cite[III, Lemma 3.3]{Cartan-Eilenberg}
\end{reference}
给定行正合的 Abel 群交换图
$$
\xymatrix{
& X \ar[r] \ar[d]^\alpha &
Y \ar[r] \ar[d]^\beta &
Z \ar[r] \ar[d]^\gamma &
0 \\
0 \ar[r] & U \ar[r] & V \ar[r] & W
}
$$
则存在典范正合列
$$
\Ker(\alpha) \to \Ker(\beta) \to \Ker(\gamma)
\to
\Coker(\alpha) \to \Coker(\beta) \to \Coker(\gamma)
$$
此外，若 $X \to Y$ 是单射，则第一个映射为单射；若 $V \to W$
是满射，则最后一个映射为满射。
\end{lemma}

\begin{proof}
映射 $\partial : \Ker(\gamma) \to \Coker(\alpha)$ 定义如下。
取 $z \in \Ker(\gamma)$，选择一个映到 $z$ 的 $y \in Y$。
于是 $\beta(y) \in V$ 在 $W$ 中映为零，故 $\beta(y)$ 是某个
$u \in U$ 的像。令 $\partial z = \overline{u}$，即 $u$ 在 $\alpha$
的余核中的类。正合性的证明从略。
\end{proof}

\section{有限模与有限表示模}
\label{algebra-section-module-finite-type}

\noindent
下面给出一些基本记号和引理。

\begin{definition}
\label{algebra-definition-module-finite-type}
设 $R$ 为环，$M$ 为 $R$-模。
\begin{enumerate}
\item 若存在 $n \in \mathbf{N}$ 及 $x_1, \ldots, x_n \in M$，
使 $M$ 的每个元素都是诸 $x_i$ 的 $R$-线性组合，则称 $M$ 为
{\it 有限 $R$-模}，亦称{\it 有限生成 $R$-模}。等价地，这意味着对某个
$n \in \mathbf{N}$ 存在满射 $R^{\oplus n} \to M$。
\item 若存在整数 $n, m \in \mathbf{N}$ 以及正合列
$$
R^{\oplus m} \longrightarrow R^{\oplus n} \longrightarrow M \longrightarrow 0
$$
则称 $M$ 为{\it 有限表示 $R$-模}，亦称{\it 有限表示的 $R$-模}。
\end{enumerate}
\end{definition}

\noindent
非正式地说，$M$ 是有限表示 $R$-模，当且仅当它有限生成，且这些生成元之间的
关系所成的模也有限生成。定义中那样一个正合列的选取称为 $M$ 的一个{\it 表示}。

\begin{lemma}
\label{algebra-lemma-lift-map}
设 $R$ 为环。设 $\alpha : R^{\oplus n} \to M$ 和 $\beta : N \to M$
为模同态。若 $\Im(\alpha) \subset \Im(\beta)$，则存在 $R$-模同态
$\gamma : R^{\oplus n} \to N$，使得 $\alpha = \beta \circ \gamma$。
\end{lemma}

\begin{proof}
令 $e_i = (0, \ldots, 0, 1, 0, \ldots, 0)$ 为 $R^{\oplus n}$ 的第 $i$ 个基向量。
由假设可取 $x_i \in N$，使得 $\alpha(e_i) = \beta(x_i)$。令
$\gamma(a_1, \ldots, a_n) = \sum a_i x_i$。由构造即有
$\alpha = \beta \circ \gamma$。
\end{proof}

\begin{lemma}
\label{algebra-lemma-extension}
设 $R$ 为环，并设
$$
0 \to M_1 \to M_2 \to M_3 \to 0
$$
为 $R$-模的短正合列。
\begin{enumerate}
\item 若 $M_1$ 和 $M_3$ 是有限 $R$-模，则 $M_2$ 是有限 $R$-模。
\item 若 $M_1$ 和 $M_3$ 是有限表示 $R$-模，则 $M_2$ 是有限表示 $R$-模。
\item 若 $M_2$ 是有限 $R$-模，则 $M_3$ 是有限 $R$-模。
\item 若 $M_2$ 是有限表示 $R$-模，且 $M_1$ 是有限 $R$-模，则
$M_3$ 是有限表示 $R$-模。
\item 若 $M_3$ 是有限表示 $R$-模，且 $M_2$ 是有限 $R$-模，则
$M_1$ 是有限 $R$-模。
\end{enumerate}
\end{lemma}

\begin{proof}
先证 (1)。若 $x_1, \ldots, x_n$ 是 $M_1$ 的生成元，而
$y_1, \ldots, y_m \in M_2$ 在 $M_3$ 中的像生成 $M_3$，则
$x_1, \ldots, x_n, y_1, \ldots, y_m$ 生成 $M_2$。

\medskip\noindent
(3) 由定义立即得到。

\medskip\noindent
证明 (5)。设 $M_3$ 有限表示且 $M_2$ 有限。选取一个表示
$$
R^{\oplus m} \to R^{\oplus n} \to M_3 \to 0
$$
由引理 \ref{algebra-lemma-lift-map}，存在映射 $R^{\oplus n} \to M_2$，
使下图中的实线部分
$$
\xymatrix{
& R^{\oplus m} \ar[r] \ar@{..>}[d] & R^{\oplus n} \ar[r] \ar[d] &
M_3 \ar[r] \ar[d]^{\text{id}} & 0 \\
0 \ar[r] & M_1 \ar[r] & M_2 \ar[r] & M_3 \ar[r] & 0
}
$$
交换，因而得到虚线箭头。由蛇形引理（引理 \ref{algebra-lemma-snake}）可知存在同构
$$
\Coker(R^{\oplus m} \to M_1)
\cong
\Coker(R^{\oplus n} \to M_2)
$$
特别地，$\Coker(R^{\oplus m} \to M_1)$ 是有限 $R$-模。又因
$\Im(R^{\oplus m} \to M_1)$ 由 (3) 为有限模，所以由 (1) 得 $M_1$ 有限。

\medskip\noindent
证明 (4)。设 $M_2$ 有限表示且 $M_1$ 有限。选取表示
$R^{\oplus m} \to R^{\oplus n} \to M_2 \to 0$，并选取满射
$R^{\oplus k} \to M_1$。由引理 \ref{algebra-lemma-lift-map}，复合映射
$R^{\oplus k} \to M_1 \to M_2$ 可分解为
$R^{\oplus k} \to R^{\oplus n} \to M_2$。于是
$R^{\oplus k + m} \to R^{\oplus n} \to M_3 \to 0$
是一个表示。

\medskip\noindent
证明 (2)。设 $M_1$ 和 $M_3$ 均为有限表示模。(1) 的证明中的论证给出交换图
$$
\xymatrix{
0 \ar[r] & R^{\oplus n} \ar[d] \ar[r] & R^{\oplus n + m} \ar[d] \ar[r] &
R^{\oplus m} \ar[d] \ar[r] & 0 \\
0 \ar[r] & M_1 \ar[r] & M_2 \ar[r] & M_3 \ar[r] & 0
}
$$
其中竖直箭头均为满射。由蛇形引理得到短正合列
$$
0 \to \Ker(R^{\oplus n} \to M_1) \to
\Ker(R^{\oplus n + m} \to M_2) \to
\Ker(R^{\oplus m} \to M_3) \to 0
$$
由 (5)，外侧两个模均有限，故中间的模也有限。再由 (4)，$M_2$ 为有限表示模。
\end{proof}

\begin{lemma}
\label{algebra-lemma-trivial-filter-finite-module}
\begin{slogan}
有限模存在滤过，使其逐次商均为循环模。
\end{slogan}
设 $R$ 为环，$M$ 为有限 $R$-模。存在由有限 $R$-子模组成的滤过
$$
0 = M_0 \subset M_1 \subset \ldots \subset M_n = M
$$
使每个商 $M_i/M_{i - 1}$ 都同构于 $R/I_i$，其中 $I_i$ 是 $R$ 的某个理想。
\end{lemma}

\begin{proof}
对 $M$ 的生成元数作归纳。设 $x_1, \ldots, x_r \in M$ 为一组生成元。
令 $M' = Rx_1 \subset M$。则 $M/M'$ 有 $r - 1$ 个生成元，故可应用归纳假设。
并且显然 $M' \cong R/I_1$，其中 $I_1 = \{f \in R \mid fx_1 = 0\}$。
\end{proof}

\begin{lemma}
\label{algebra-lemma-finite-over-subring}
设 $R \to S$ 为环同态，$M$ 为 $S$-模。若 $M$ 作为 $R$-模是有限的，
则 $M$ 作为 $S$-模也是有限的。
\end{lemma}

\begin{proof}
事实上，$M$ 的任一 $R$-生成集也是 $M$ 的 $S$-生成集，
因为其 $R$-模结构由 $R$ 在 $S$ 中的像诱导。
\end{proof}


\section{有限型与有限表示的环同态}
\label{algebra-section-finite-type}

\begin{definition}
\label{algebra-definition-finite-type}
设 $R \to S$ 为环同态。
\begin{enumerate}
\item 若存在 $n \in \mathbf{N}$ 以及 $R$-代数满射
$R[x_1, \ldots, x_n] \to S$，则称 $R \to S$ 为{\it 有限型的}，
亦称{\it $S$ 是有限型 $R$-代数}。
% P01-E0263: source wording candidate recorded in ERRATA_CANDIDATES.jsonl.
\item 若存在整数 $n, m \in \mathbf{N}$、多项式
$f_1, \ldots, f_m \in R[x_1, \ldots, x_n]$ 以及下列 $R$-代数同构：
$$
R[x_1, \ldots, x_n]/(f_1, \ldots, f_m) \cong S
$$
则称 $R \to S$ 为{\it 有限表示的}。
\end{enumerate}
\end{definition}

\noindent
非正式地说，$R \to S$ 为有限表示，当且仅当 $S$ 作为 $R$-代数有限生成，
且这些生成元之间的关系所成的理想有限生成。定义中的一个满射
$R[x_1, \ldots, x_n] \to S$ 有时称为 $S$ 的一个{\it 表示}。

\begin{lemma}
\label{algebra-lemma-compose-finite-type}
有限型与有限表示具有下列稳定性质。
\begin{enumerate}
\item 有限型环同态的复合仍为有限型。
\item 有限表示环同态的复合仍为有限表示。
\item 给定 $R \to S' \to S$，若 $R \to S$ 为有限型，则
$S' \to S$ 为有限型。
\item 给定 $R \to S' \to S$，若 $R \to S$ 为有限表示，且
$R \to S'$ 为有限型，则 $S' \to S$ 为有限表示。
\end{enumerate}
\end{lemma}

\begin{proof}
只证最后一个断言。写成
$S = R[x_1, \ldots, x_n]/(f_1, \ldots, f_m)$ 及
$S' = R[y_1, \ldots, y_a]/I$。设 $y_i$ 的类 $\bar y_i$ 在 $S$
中映为 $h_i \bmod (f_1, \ldots, f_m)$。于是显然
$S = S'[x_1, \ldots, x_n]/(f_1, \ldots, f_m,
h_1 - \bar y_1, \ldots, h_a - \bar y_a)$。
\end{proof}

\begin{lemma}
\label{algebra-lemma-finite-presentation-independent}
设 $R \to S$ 为有限表示环同态。对任一满射
$\alpha : R[x_1, \ldots, x_n] \to S$，$\alpha$ 的核都是
$R[x_1, \ldots, x_n]$ 中的有限生成理想。
\end{lemma}

\begin{proof}
写成 $S = R[y_1, \ldots, y_m]/(f_1, \ldots, f_k)$。
选取 $g_i \in R[y_1, \ldots, y_m]$ 作为 $\alpha(x_i)$ 的提升。
于是 $S = R[x_i, y_j]/(f_l, x_i - g_i)$。选取
$h_j \in R[x_1, \ldots, x_n]$，使 $\alpha(h_j)$ 对应于
$y_j \bmod (f_1, \ldots, f_k)$。考虑映射
$\psi : R[x_i, y_j] \to R[x_i]$，$x_i \mapsto x_i$，
$y_j \mapsto h_j$。则 $\alpha$ 的核是 $(f_l, x_i - g_i)$
在 $\psi$ 下的像，结论成立。
\end{proof}

\begin{lemma}
\label{algebra-lemma-finitely-presented-over-subring}
设 $R \to S$ 为环同态，$M$ 为 $S$-模。假设 $R \to S$ 为有限型，
且 $M$ 作为 $R$-模有限表示，则 $M$ 作为 $S$-模也有限表示。
\end{lemma}

\begin{proof}
这与引理 \ref{algebra-lemma-compose-finite-type} 中 (4) 的证明类似。
可设 $S = R[x_1, \ldots, x_n]/J$。选取
$y_1, \ldots, y_m \in M$，使其作为 $R$-模生成 $M$；并选取关系
$\sum a_{ij} y_j = 0$，$i = 1, \ldots, t$，使其生成
$R^{\oplus m} \to M$ 的核。对任意 $i = 1, \ldots, n$ 和
$j = 1, \ldots, m$，写成
$$
x_i y_j = \sum a_{ijk} y_k
$$
其中 $a_{ijk} \in R$。考虑由 $y_1, \ldots, y_m$ 生成的 $S$-模 $N$，
其关系为 $\sum a_{ij} y_j = 0$，$i = 1, \ldots, t$，以及
$x_i y_j = \sum a_{ijk} y_k$，$i = 1, \ldots, n$、
$j = 1, \ldots, m$。于是 $N$ 有表示
$$
S^{\oplus nm + t} \longrightarrow S^{\oplus m} \longrightarrow N
\longrightarrow 0
$$
由构造存在满射 $\varphi : N \to M$。为完成证明，只需证明 $\varphi$
为单射。设对某些 $b_j \in S$ 有 $z = \sum b_j y_j \in N$。
可将 $b_j$ 视为以 $R$ 中元素为系数、关于 $x_1, \ldots, x_n$ 的多项式。
应用形如 $x_i y_j = \sum a_{ijk} y_k$ 的关系，可归纳降低这些多项式的次数。
故对某些 $c_j \in R$ 有 $z = \sum c_j y_j$。于是若
$\varphi(z) = 0$，则向量 $(c_1, \ldots, c_m)$ 是诸向量
$(a_{i1}, \ldots, a_{im})$ 的 $R$-线性组合，从而如所需地得到 $z = 0$。
\end{proof}





\section{有限环同态}
\label{algebra-section-finite}

\noindent
定义如下。

\begin{definition}
\label{algebra-definition-finite-ring-map}
设 $\varphi : R \to S$ 为环同态。若 $S$ 作为 $R$-模有限，
则称 $\varphi : R \to S$ {\it 有限}。
\end{definition}

\begin{lemma}
\label{algebra-lemma-finite-module-over-finite-extension}
设 $R \to S$ 为有限环同态，$M$ 为 $S$-模。则 $M$ 作为 $R$-模有限，
当且仅当 $M$ 作为 $S$-模有限。
\end{lemma}

\begin{proof}
其中一个方向由引理 \ref{algebra-lemma-finite-over-subring} 得到。
为证另一个方向，设 $M$ 作为 $S$-模有限。取
$x_1, \ldots, x_n \in S$，使其作为 $R$-模生成 $S$；取
$y_1, \ldots, y_m \in M$，使其作为 $S$-模生成 $M$。
则诸 $x_i y_j$ 作为 $R$-模生成 $M$。
\end{proof}

\begin{lemma}
\label{algebra-lemma-finite-transitive}
设 $R \to S$ 与 $S \to T$ 均为有限环同态，则 $R \to T$ 有限。
\end{lemma}

\begin{proof}
若诸 $t_i$ 作为 $S$-模生成 $T$，诸 $s_j$ 作为 $R$-模生成 $S$，
则诸 $t_i s_j$ 作为 $R$-模生成 $T$。（亦可由引理
\ref{algebra-lemma-finite-module-over-finite-extension} 得到。）
\end{proof}

\begin{lemma}
\label{algebra-lemma-finite-finite-type}
设 $\varphi : R \to S$ 为环同态。
\begin{enumerate}
\item 若 $\varphi$ 有限，则 $\varphi$ 为有限型。
\item 若 $S$ 作为 $R$-模为有限表示，则 $\varphi$ 为有限表示。
\end{enumerate}
\end{lemma}

\begin{proof}
对 (1)，若 $x_1, \ldots, x_n \in S$ 作为 $R$-模生成 $S$，
则 $x_1, \ldots, x_n$ 也作为 $R$-代数生成 $S$。对 (2)，设关系
$\sum r_j^ix_i = 0$，$j = 1, \ldots, m$ 作为 $S$ 的 $R$-模生成元 $x_i$
之间的关系的一组生成元。此外，对某些 $r_i \in R$ 写成
$1 = \sum r_ix_i$，并对某些 $r_{ij}^k \in R$ 写成
$x_ix_j = \sum r_{ij}^k x_k$。于是作为 $R$-代数有
$$
S = R[t_1, \ldots, t_n]/
(\sum r_j^it_i,\ 1 - \sum r_it_i,\ t_it_j - \sum r_{ij}^k t_k)
$$
这证明了 (2)。
\end{proof}

\noindent
关于有限环同态的更多内容，见第 \ref{algebra-section-finite-ring-extensions} 节。



\section{余极限}
\label{algebra-section-colimits}

\noindent
本节部分内容与范畴论中关于余极限的一般讨论重叠，见范畴论第
\ref{categories-section-limits}--\ref{categories-section-posets-limits} 节。
预序集的概念定义于范畴论，定义 \ref{categories-definition-directed-set}；
它比偏序集的概念稍弱。

\begin{definition}
\label{algebra-definition-directed-system}
设 $(I, \leq)$ 为预序集。一个{\it 以 $I$ 为指标的 $R$-模系统
$(M_i, \mu_{ij})$}，由以 $I$ 为指标的一族 $R$-模
$\{M_i\}_{i\in I}$ 以及一族 $R$-模同态
$\{\mu_{ij} : M_i \to M_j\}_{i \leq j}$ 构成，并且对所有
$i \leq j \leq k$ 有
$$
\mu_{ii} = \text{id}_{M_i}\quad
\mu_{ik} = \mu_{jk}\circ \mu_{ij}
$$
若 $I$ 是有向集，则称 $(M_i, \mu_{ij})$ 为{\it 有向系统}。
\end{definition}

\noindent
这与范畴论，定义 \ref{categories-definition-system-over-poset}
及第 \ref{categories-section-posets-limits} 节中定义的概念相同。
任意范畴中一个图式或系统的余极限之定义，见范畴论，定义
\ref{categories-definition-colimit}。

\begin{lemma}
\label{algebra-lemma-colimit}
设 $(M_i, \mu_{ij})$ 是预序集 $I$ 上的 $R$-模系统。
该系统 $(M_i, \mu_{ij})$ 的余极限是商 $R$-模
$(\bigoplus_{i\in I} M_i) /Q$，其中 $Q$ 是由所有形如
$$
\iota_i(x_i) - \iota_j(\mu_{ij}(x_i))
$$
的元素生成的 $R$-子模，而
$\iota_i : M_i \to \bigoplus_{i\in I} M_i$ 是自然嵌入。
将余极限记作 $M = \colim_i M_i$，投影记作
$\pi : \bigoplus_{i\in I} M_i \to M$，并令
$\phi_i = \pi \circ \iota_i : M_i \to M$。
\end{lemma}

\begin{proof}
此引理是范畴论，引理
\ref{categories-lemma-colimits-coproducts-coequalizers} 的特例，
不过我们也直接证明这一情形。注意在上述构造中
$\phi_i = \phi_j\circ \mu_{ij}$。要证明偶对 $(M, \phi_i)$ 是余极限，
须验证其泛性质：对任一其他这样的偶对 $(Y, \psi_i)$，其中
$\psi_i : M_i \to Y$ 且 $\psi_i = \psi_j\circ \mu_{ij}$，
存在唯一的 $R$-模同态 $g : M \to Y$，使下图交换：
$$
\xymatrix{
M_i \ar[rr]^{\mu_{ij}} \ar[dr]^{\phi_i} \ar[ddr]_{\psi_i} & &
M_j\ar[dl]_{\phi_j} \ar[ddl]^{\psi_j} \\
& M \ar[d]^{g}\\
& Y
}
$$
这是显然的，因为可在直和 $\bigoplus M_i$ 的直和项 $M_i$ 上
取映射 $\psi_i$，由此定义 $g$。
\end{proof}

\begin{lemma}
\label{algebra-lemma-directed-colimit}
设 $(M_i, \mu_{ij})$ 是预序集 $I$ 上的 $R$-模系统，并假设 $I$ 有向。
该系统 $(M_i, \mu_{ij})$ 的余极限典范同构于如下定义的模 $M$：
\begin{enumerate}
\item 作为集合，令
$$
M = \left(\coprod\nolimits_{i \in I} M_i\right)/\sim
$$
其中，对 $m \in M_i$ 和 $m' \in M_{i'}$，规定
$$
m \sim m' \Leftrightarrow
\mu_{ij}(m) = \mu_{i'j}(m')\text{ 对某个 }j \geq i, i'
$$
\item 作为 Abel 群，对 $m \in M_i$ 和 $m' \in M_{i'}$，
定义 $m$ 与 $m'$ 在 $M$ 中的类之和为
$\mu_{ij}(m) + \mu_{i'j}(m')$ 的类，其中 $j \in I$ 是任一满足
$i \leq j$ 且 $i' \leq j$ 的指标；
\item 作为 $R$-模，对 $m \in M_i$ 和 $x \in R$，定义 $x$ 与 $m$
在 $M$ 中的类之积为 $xm$ 在 $M$ 中的类。
\end{enumerate}
典范映射 $\phi_i : M_i \to M$ 由典范映射
$M_i \to \coprod_{i \in I} M_i$ 诱导。
\end{lemma}

\begin{proof}
从略。可与范畴论，第 \ref{categories-section-directed-colimits} 节比较。
\end{proof}

\begin{lemma}
\label{algebra-lemma-zero-directed-limit}
设 $(M_i, \mu_{ij})$ 为有向系统，并设 $M = \colim M_i$，其映射为
$\mu_i : M_i \to M$。则对 $x_i \in M_i$，$\mu_i(x_i) = 0$，
当且仅当存在 $j \geq i$ 使得 $\mu_{ij}(x_i) = 0$。
\end{lemma}

\begin{proof}
这由引理 \ref{algebra-lemma-directed-colimit} 中有向余极限的描述立即得到。
\end{proof}

\begin{example}
\label{algebra-example-zero-colimit-different}
考虑偏序集 $I = \{a, b, c\}$，其中 $a < b$、$a < c$，且无其他严格不等式。
$I$ 上的一个系统 $(M_a, M_b, M_c, \mu_{ab}, \mu_{ac})$ 由三个
$R$-模 $M_a, M_b, M_c$ 以及两个 $R$-模同态
$\mu_{ab} : M_a \to M_b$ 和 $\mu_{ac} : M_a \to M_c$ 构成。
该系统的余极限就是
$$
M := \colim_{i \in I} M_i = \Coker(M_a \to M_b \oplus M_c)
$$
其中映射为 $\mu_{ab} \oplus -\mu_{ac}$。因此，典范映射
$M_a \to M$ 的核为 $\Ker(\mu_{ab}) + \Ker(\mu_{ac})$；
而典范映射 $M_b \to M$ 的核是 $\Ker(\mu_{ac})$ 在映射
$\mu_{ab}$ 下的像。故引理 \ref{algebra-lemma-zero-directed-limit}
的结论对一般系统显然不成立。
\end{example}


\begin{definition}
\label{algebra-definition-homomorphism-directed-systems}
设 $(M_i, \mu_{ij})$、$(N_i, \nu_{ij})$ 是同一预序集 $I$ 上的
$R$-模系统。从 $(M_i, \mu_{ij})$ 到 $(N_i, \nu_{ij})$ 的
一个{\it 系统同态} $\Phi$，按定义是由 $R$-模同态
$\phi_i : M_i \to N_i$ 组成的一族，并且对所有 $i \leq j$ 有
$\phi_j \circ \mu_{ij} = \nu_{ij} \circ \phi_i$。
\end{definition}

\noindent
用范畴语言说，这与相应图式 $M : I \to \text{Mod}_R$ 和
$N : I \to \text{Mod}_R$ 之间的函子变换是同一个概念。
下一个引理是范畴论，引理 \ref{categories-lemma-functorial-colimit} 的特例。

\begin{lemma}
\label{algebra-lemma-homomorphism-limit}
设 $(M_i, \mu_{ij})$、$(N_i, \nu_{ij})$ 是同一预序集上的 $R$-模系统。
从 $(M_i, \mu_{ij})$ 到 $(N_i, \nu_{ij})$ 的系统同态
$\Phi = (\phi_i)$ 诱导唯一同态
$$
\colim \phi_i : \colim M_i \longrightarrow \colim N_i
$$
使得对所有 $i \in I$，下图交换：
$$
\xymatrix{
M_i \ar[r] \ar[d]_{\phi_i} & \colim M_i \ar[d]^{\colim \phi_i} \\
N_i \ar[r] & \colim N_i
}
$$
\end{lemma}

\begin{proof}
写成 $M = \colim M_i$、$N = \colim N_i$ 及 $\phi = \colim \phi_i$
（后者尚待构造）。以下将不再另作说明地使用引理 \ref{algebra-lemma-colimit}
中对 $M$ 和 $N$ 的显式描述。引理的条件等价于下图交换：
$$
\xymatrix{
\bigoplus_{i\in I} M_i \ar[r] \ar[d]_{\bigoplus\phi_i} & M \ar[d]^\phi \\
\bigoplus_{i\in I} N_i \ar[r] & N
}
$$
因此，若 $\phi$ 存在，则其唯一性显然。为证明 $\phi$ 存在，
只须说明上方水平箭头的核在 $\bigoplus \phi_i$ 下映入下方水平箭头的核。
为此，取 $j \leq k$ 及 $x_j \in M_j$，则
$$
(\bigoplus \phi_i)(x_j - \mu_{jk}(x_j))
=
\phi_j(x_j) - \phi_k(\mu_{jk}(x_j))
=
\phi_j(x_j) - \nu_{jk}(\phi_j(x_j))
$$
如所需地属于下方水平箭头的核。
\end{proof}

\begin{lemma}
\label{algebra-lemma-directed-colimit-exact}
\begin{slogan}
滤过余极限是正合的。有向余极限是正合的。
\end{slogan}
设 $I$ 为有向集。设 $(L_i, \lambda_{ij})$、$(M_i, \mu_{ij})$ 和
$(N_i, \nu_{ij})$ 为 $I$ 上的 $R$-模系统。设
$\varphi_i : L_i \to M_i$ 与 $\psi_i : M_i \to N_i$ 为 $I$ 上的系统同态。
假设对每个 $i \in I$，$R$-模序列
$$
\xymatrix{
L_i \ar[r]^{\varphi_i} &
M_i \ar[r]^{\psi_i} &
N_i
}
$$
是同调为 $H_i$ 的复形。则诸 $R$-模 $H_i$ 构成 $I$ 上的系统；
$R$-模序列
$$
\xymatrix{
\colim_i L_i \ar[r]^\varphi &
\colim_i M_i \ar[r]^\psi &
\colim_i N_i
}
$$
也是复形；若以 $H$ 表示其同调，则
$$
H = \colim_i H_i.
$$
\end{lemma}

\begin{proof}
显然
$
\xymatrix{
\colim_i L_i \ar[r]^\varphi &
\colim_i M_i \ar[r]^\psi &
\colim_i N_i
}
$
是复形。对每个 $i \in I$，存在典范 $R$-模同态 $H_i \to H$：
它把每个 $[m] \in H_i = \Ker(\psi_i) / \Im(\varphi_i)$ 映到
$m$ 在 $\colim_i M_i$ 中的像于
$H = \Ker(\psi) / \Im(\varphi)$ 中的剩余类。这些同态给出同态
$\colim_i H_i \to H$。只须证明它既满又单。

\medskip\noindent
以下将反复使用引理 \ref{algebra-lemma-directed-colimit} 中关于 $I$ 上余极限的描述，
而不再另作说明。取 $h \in H$。由于
$H = \Ker(\psi)/\Im(\varphi)$，$h$ 是模 $\Im(\varphi)$ 意义下
某个元素 $[m]$ 的类，而该元素属于 $\Ker(\psi) \subset \colim_i M_i$。
选取 $i$，使 $[m]$
来自某个 $m \in M_i$。再选取 $j \geq i$，使
$\nu_{ij}(\psi_i(m)) = 0$；由于 $[m] \in \Ker(\psi)$，这样的指标
确实存在。以 $j$ 代替 $i$，并以 $\mu_{ij}(m)$ 代替 $m$，
即可设 $m \in \Ker(\psi_i)$。这证明 $\colim_i H_i \to H$ 为满射。

\medskip\noindent
设 $h_i \in H_i$ 在 $H$ 中的像为零。由于
$H_i = \Ker(\psi_i)/\Im(\varphi_i)$，可用元素
$m \in \Ker(\psi_i) \subset M_i$ 表示 $h_i$。
$h_i$ 在 $H$ 中消失的假设意味着 $m$ 在 $\colim_i M_i$ 中的类
属于 $\varphi$ 的像。因此，存在 $j \geq i$ 及 $l \in L_j$，使得
$\varphi_j(l) = \mu_{ij}(m)$。这显然表明 $h_i$ 在 $H_j$ 中的像为零，
从而证明 $\colim_i H_i \to H$ 为单射。
\end{proof}

\begin{example}
\label{algebra-example-colimit-not-exact}
一般而言，取余极限并不正合。考虑偏序集 $I = \{a, b, c\}$，
其中 $a < b$、$a < c$，且无其他严格不等式，正如例
\ref{algebra-example-zero-colimit-different} 中那样。考虑系统同态
$(0, \mathbf{Z}, \mathbf{Z}, 0, 0) \to
(\mathbf{Z}, \mathbf{Z}, \mathbf{Z}, 1, 1)$。
由例 \ref{algebra-example-zero-colimit-different} 中对余极限的描述可知，
相应的余极限同态不是单射，尽管该系统同态在每个对象上都是单射。
故引理 \ref{algebra-lemma-directed-colimit-exact} 的结论对一般系统不成立。
\end{example}

\begin{lemma}
\label{algebra-lemma-almost-directed-colimit-exact}
设 $\mathcal{I}$ 为满足范畴论，引理
\ref{categories-lemma-split-into-directed} 假设的指标范畴。
则在 $\mathcal{I}$ 上取 Abel 群图式的余极限是正合的（即引理
\ref{algebra-lemma-directed-colimit-exact} 的类似结论在此情形成立）。
\end{lemma}

\begin{proof}
由范畴论，引理 \ref{categories-lemma-split-into-directed}，
可写 $\mathcal{I} = \coprod_{j \in J} \mathcal{I}_j$，其中每个
$\mathcal{I}_j$ 都是滤过范畴，而 $J$ 可以为空。由范畴论，引理
\ref{categories-lemma-directed-category-system}，在指标范畴
$\mathcal{I}_j$ 上取余极限，等同于在某个有向集上取余极限。
因此，引理 \ref{algebra-lemma-directed-colimit-exact} 可用于这些余极限。
问题由此归结为证明：在 $R$-模范畴中，对集合 $J$ 取余积是正合的。
换言之，给定 $R$ 模的正合列 $L_j \to M_j \to N_j$，须证明
% P01-E0264: source wording candidate recorded in ERRATA_CANDIDATES.jsonl.
$$
\bigoplus\nolimits_{j \in J} L_j
\longrightarrow
\bigoplus\nolimits_{j \in J} M_j
\longrightarrow
\bigoplus\nolimits_{j \in J} N_j
$$
正合。这可直接验证，即使 $J$ 为空也成立。
\end{proof}
\section{局部化}
\label{algebra-section-localization}

\begin{definition}
\label{algebra-definition-multiplicative-subset}
设 $R$ 为环，$S$ 是 $R$ 的子集。若
$1\in S$，且 $S$ 对乘法封闭，即
$s, s' \in S \Rightarrow ss' \in S$，
则称 $S$ 为{\it $R$ 的乘法子集}。
\end{definition}

\noindent
给定环 $A$ 及其乘法子集 $S$，在 $A \times S$ 上定义如下关系：
$$
(x, s) \sim (y, t)
\Leftrightarrow
\exists u \in S \text{ 使得 } (xt-ys)u = 0
$$
容易验证这是等价关系。用 $x/s$（或 $\frac{x}{s}$）表示
$(x, s)$ 的等价类，并用 $S^{-1}A$ 表示所有等价类组成的集合。
在 $S^{-1}A$ 中定义加法与乘法如下：
$$
x/s + y/t = (xt + ys)/st, \quad
x/s \cdot y/t = xy/st
$$
可以验证，在这些运算下 $S^{-1}A$ 成为环。

\begin{definition}
\label{algebra-definition-localization}
此环称为{\it $A$ 关于 $S$ 的局部化}。
\end{definition}

\noindent
从 $A$ 到其局部化 $S^{-1}A$ 有自然环同态
$$
A \longrightarrow S^{-1}A, \quad x \longmapsto x/1
$$
有时称为{\it 局部化同态}。一般而言，局部化同态不一定是单射；
若 $S$ 不含零因子，则它是单射。事实上，若 $x/1 = 0$，
则存在 $u\in S$，使得 $A$ 中 $xu = 0$；由于 $S$ 中没有零因子，
必有 $x = 0$。环的局部化具有下列泛性质。

\begin{proposition}
\label{algebra-proposition-universal-property-localization}
设 $f : A \to B$ 为环同态，并把 $S$ 的每个元素映为 $B$ 中的单位。
则存在唯一同态 $g : S^{-1}A \to B$，使下图交换：
$$
\xymatrix{
A \ar[rr]^{f} \ar[dr] & & B \\
& S^{-1}A \ar[ur]_g
}
$$
\end{proposition}

\begin{proof}
存在性。如下定义映射 $g$。对 $x/s\in S^{-1}A$，令
$g(x/s) = f(x)f(s)^{-1}\in B$。由定义容易验证这是良定义的环同态，
并且显然使上图交换。

\medskip\noindent
唯一性。现设 $g' : S^{-1}A \to B$ 满足 $g'(x/1) = f(x)$，
证明 $g = g'$。由图的交换性，对 $s \in S$ 有
$f(s) = g'(s/1)$。而 $B$ 中 $g'(1/s)f(s) = 1$，
故 $g'(1/s) = f(s)^{-1}$，从而
$g'(x/s) = g'(x/1)g'(1/s) = f(x)f(s)^{-1} = g(x/s)$。
\end{proof}

\begin{lemma}
\label{algebra-lemma-localization-zero}
局部化 $S^{-1}A$ 是零环，当且仅当 $0\in S$。
\end{lemma}

\begin{proof}
若 $0\in S$，则按定义任意有序对 $(a, s)\sim (0, 1)$。
若 $0\not \in S$，则显然在 $S^{-1}A$ 中 $1/1 \neq 0/1$。
\end{proof}

\begin{lemma}
\label{algebra-lemma-localization-and-modules}
设 $R$ 为环，$S \subset R$ 为乘法子集。$S^{-1}R$-模范畴等价于
满足下列性质的 $R$-模 $N$ 所成范畴：每个 $s \in S$ 在 $N$ 上的作用
都是自同构。
\end{lemma}

\begin{proof}
给出该等价的函子把 $S^{-1}R$-模 $M$ 关联到同一个模，但通过局部化同态
$R \to S^{-1}R$ 把它视为 $R$-模。反之，若 $N$ 是 $R$-模，
且每个 $s \in S$ 都通过自同构 $s_N$ 作用，则可令 $x/s$ 通过
$x_N  \circ s_N^{-1}$ 作用，从而把 $N$ 视为 $S^{-1}R$-模。
这两个函子互为拟逆的验证从略。
\end{proof}

\noindent
环的局部化概念可推广为模的局部化。设 $A$ 为环，$S$ 是 $A$ 的乘法
子集，$M$ 是 $A$-模。在 $M \times S$ 上定义如下关系：
$$
(m, s) \sim (n, t)
\Leftrightarrow
\exists u\in S \text{ 使得 } (mt-ns)u = 0
$$
这显然是等价关系。
% P01-E0265: source wording candidate recorded in ERRATA_CANDIDATES.jsonl.
用 $m/s$（或 $\frac{m}{s}$）表示 $(m, s)$ 的等价类，并用
$S^{-1}M$ 表示所有等价类组成的集合。定义加法与标量乘法如下：
% P01-E0266: source formula candidate recorded in ERRATA_CANDIDATES.jsonl.
$$
m/s + n/t = (mt + ns)/st,\quad
m/s\cdot n/t = mn/st
$$
显然，这使 $S^{-1}M$ 成为 $S^{-1}A$-模。

\begin{definition}
\label{algebra-definition-localization-module}
$S^{-1}A$-模 $S^{-1}M$ 称为 $M$ 关于 $S$ 的{\it 局部化}。
\end{definition}

\noindent
注意，有时把 $A$-模同态 $M \to S^{-1}M$，$m \mapsto m/1$
称为{\it 局部化同态}。它满足下列泛性质。

\begin{lemma}
\label{algebra-lemma-universal-property-localization-module}
设 $R$ 为环。
% P01-E0267: source wording candidate recorded in ERRATA_CANDIDATES.jsonl.
设 $S \subset R$ 为乘法子集，$M$、$N$ 为 $R$-模。
假设 $S$ 的所有元素在 $N$ 上的作用都是自同构。则由局部化同态诱导的
典范映射
$$
\Hom_R(S^{-1}M, N) \longrightarrow \Hom_R(M, N)
$$
是同构。
\end{lemma}

\begin{proof}
该映射良定义且为 $R$-线性的事实是显然的。
单射性：设 $\alpha \in \Hom_R(S^{-1}M, N)$，并任取
$m/s \in S^{-1}M$。由于
$s \cdot \alpha(m/s) = \alpha(m/1)$，故
$ \alpha(m/s) =s^{-1}(\alpha (m/1))$；所以 $\alpha$ 完全由它在
$M$ 于 $S^{-1}M$ 中的像上的取值决定。
满射性：给定一个
% P01-E0268: source typography candidate recorded in ERRATA_CANDIDATES.jsonl.
R-线性映射 $\beta : M \rightarrow N$。须证明它可以“延拓”到
$S^{-1}M$。定义集合映射
$$
M \times S \rightarrow N,\quad
(m,s) \mapsto s^{-1}\beta(m)
$$
显然，该映射保持上述等价关系，故下降为良定义映射
$\alpha : S^{-1}M \rightarrow N$。还须证明它为 $R$-线性的。
取 $r, r' \in R$、$s, s' \in S$ 以及 $m, m' \in M$，则
\begin{align*}
\alpha(r \cdot m/s + r' \cdot m' /s')
& =
\alpha((r \cdot s' \cdot m + r' \cdot s \cdot m') /(ss')) \\
& =
(ss')^{-1}\beta(r \cdot s' \cdot m + r' \cdot s \cdot m') \\
& =
(ss')^{-1} (r \cdot s' \beta (m) + r' \cdot s \beta (m')) \\
& =
r \alpha (m/s) + r' \alpha (m' /s')
\end{align*}
所需结论成立。
\end{proof}

\begin{example}
\label{algebra-example-localize-at-prime}
设 $A$ 为环，$M$ 为 $A$-模。以下是几个重要的局部化例子。
\begin{enumerate}
\item 给定 $A$ 的素理想 $\mathfrak p$，考虑
$S = A\setminus\mathfrak p$。立即可知 $S$ 是乘法子集。在此情形，
分别用 $A_\mathfrak p$ 和 $M_\mathfrak p$ 表示 $A$ 和 $M$ 关于 $S$
的局部化，称为{\it $A$、相应地 $M$ 在 $\mathfrak p$ 处的局部化}。
\item 设 $f\in A$。考虑 $S = \{1, f, f^2, \ldots\}$。
这显然是 $A$ 的乘法子集。此时分别用 $A_f$（相应地 $M_f$）表示局部化
$S^{-1}A$（相应地 $S^{-1}M$），称为
{\it $A$、相应地 $M$ 关于 $f$ 的局部化}。
注意，$A_f = 0$ 当且仅当 $f$ 在 $A$ 中是幂零元。
\item 令 $S = \{f \in A \mid f \text{ is not a zerodivisor in }A\}$。
这是 $A$ 的乘法子集。此时，环 $Q(A) = S^{-1}A$ 称为 $A$ 的
{\it 全商环}，也称为{\it 全分式环}。
\item 若 $A$ 是整环，则全商环 $Q(A)$ 就是 $A$ 的分式域。请参见
域，例 \ref{fields-example-quotient-field}。
\end{enumerate}
\end{example}

\begin{lemma}
\label{algebra-lemma-localization-colimit}
设 $R$ 为环，$S \subset R$ 为乘法子集，$M$ 为 $R$-模。则
$$
S^{-1}M = \colim_{f \in S} M_f
$$
其中 $S$ 上的预序由下式给出：
$f \geq f' \Leftrightarrow f = f'f''$，其中某个 $f'' \in R$；
此时映射 $M_{f'} \to M_f$ 由
$m/(f')^e \mapsto m(f'')^e/f^e$ 给出。
\end{lemma}

\begin{proof}
从略。提示：使用引理
\ref{algebra-lemma-universal-property-localization-module} 的泛性质。
\end{proof}

\noindent
以下一段中，设 $A$ 为环，$M, N$ 为 $A$ 上的模。

\medskip\noindent
若 $S$、$S'$ 是 $A$ 的乘法子集，则显然
$$
SS' = \{ss' : s\in S, \ s'\in S'\}
$$
也是 $A$ 的乘法子集。于是有下列结论。

\begin{proposition}
\label{algebra-proposition-localize-twice}
设 $\overline{S}$ 为 $S$ 在 $S'^{-1}A$ 中的像，则
$(SS')^{-1}A$ 同构于 $\overline{S}^{-1}(S'^{-1}A)$。
\end{proposition}

\begin{proof}
把 $x\in A$ 映到 $x/1\in (SS')^{-1}A$ 的映射，由泛性质诱导出
把 $x/s\in S'^{-1}A$ 映到 $x/s \in (SS')^{-1}A$ 的映射。
$\overline{S}$ 中元素的像在 $(SS')^{-1}A$ 中可逆。再次由泛性质，
得到映射
$f : \overline{S}^{-1}(S'^{-1}A) \to (SS')^{-1}A$，
它把 $(x/s')/(s/1)$ 映到 $x/ss'$。

\medskip\noindent
另一方面，把 $x\in A$ 映到 $(x/1)/(1/1)$ 的、从 $A$ 到
$\overline{S}^{-1}(S'^{-1}A)$ 的映射，再次由泛性质诱导出映射
$g : (SS')^{-1}A \to \overline{S}^{-1}(S'^{-1}A)$，
它把 $x/ss'$ 映到 $(x/s')/(s/1)$。
立即可验证 $f$ 与 $g$ 互为逆映射，故二者都是同构。
\end{proof}

\noindent
对于模 $M$，有下列结论。

\begin{proposition}
\label{algebra-proposition-localize-twice-module}
把 $S'^{-1}M$ 视为 $A$-模，则 $S^{-1}(S'^{-1}M)$ 同构于
$(SS')^{-1}M$。
\end{proposition}

\begin{proof}
% P01-E0269: source wording candidate recorded in ERRATA_CANDIDATES.jsonl.
注意，对于给定的 $A$-模 M，我们尚未证明 $S^{-1}M$ 的任何泛性质。
因此不能像前一个证明那样论证，而必须显式构造此同构。

\medskip\noindent
如下定义映射：
\begin{align*}
& f : S^{-1}(S'^{-1}M) \longrightarrow (SS')^{-1}M, \quad \frac{x/s'}{s}\mapsto
x/ss'\\
& g : (SS')^{-1}M \longrightarrow S^{-1}(S'^{-1}M), \quad x/t\mapsto
\frac{x/s'}{s}\ \text{对某个 }s\in S, s'\in S', \text{ 且 }
t = ss'
\end{align*}
必须验证这些同态良定义，即不依赖于分式的选择。
% P01-E0270: source wording candidate recorded in ERRATA_CANDIDATES.jsonl.
这容易验证；证明它们互为逆映射也很直接。
\end{proof}

\noindent
% P01-E0271: source precision candidate recorded in ERRATA_CANDIDATES.jsonl.
若 $u : M \to N$ 是 $A$ 同态，则局部化确实诱导出良定义的
$S^{-1}A$ 同态 $S^{-1}u : S^{-1}M \to S^{-1}N$，把 $x/s$ 映到
$u(x)/s$。立即可验证此构造具有函子性，故 $S^{-1}$ 实际上是从
$A$-模范畴到 $S^{-1}A$-模范畴的函子。此外，该函子是正合的，
如下命题所示。

\begin{proposition}
\label{algebra-proposition-localization-exact}
\begin{slogan}
局部化是正合的。
\end{slogan}
% P01-E0272: source base-ring candidate recorded in ERRATA_CANDIDATES.jsonl.
设 $L\xrightarrow{u} M\xrightarrow{v} N$ 为 $R$-模的正合列。则
$S^{-1}L \to S^{-1}M \to S^{-1}N$ 也正合。
\end{proposition}

\begin{proof}
首先，由于局部化是函子，显然
$S^{-1}L \to S^{-1}M \to S^{-1}N$ 是复形。其次，设 $x/s$
在 $S^{-1}N$ 中的像为零，其中 $x/s \in S^{-1}M$。则按定义，存在
$t\in S$，使得 $N$ 中 $v(xt) = v(x)t = 0$，亦即
$xt \in \Ker(v)$。由 $L \to M \to N$ 的正合性，存在 $L$ 中的某个
$y$，使 $xt = u(y)$。于是 $x/s$ 是 $y/st$ 的像。这就证明了正合性。
\end{proof}

\begin{lemma}
\label{algebra-lemma-localize-quotient-modules}
局部化保持商；即若 $N$ 是 $M$ 的子模，则
$S^{-1}(M/N)\simeq (S^{-1}M)/(S^{-1}N)$。
\end{lemma}

\begin{proof}
由正合列
$$
0 \longrightarrow N \longrightarrow M \longrightarrow M/N \longrightarrow 0
$$
得到
$$
0 \longrightarrow S^{-1}N \longrightarrow S^{-1}M
\longrightarrow S^{-1}(M/N) \longrightarrow 0
$$
% P01-E0273: source result-type candidate recorded in ERRATA_CANDIDATES.jsonl.
引理随即成立。
\end{proof}

\noindent
% P01-E0274: source result-type candidate recorded in ERRATA_CANDIDATES.jsonl.
在前一引理中，对 $A$ 的理想 $I$ 取 $N = I$、$M = A$，
便知作为 $A$-模有
$S^{-1}A/S^{-1}I \simeq S^{-1}(A/I)$。下一个命题表明它们作为环也同构。

\begin{proposition}
\label{algebra-proposition-localize-quotient}
设 $I$ 是 $A$ 的理想，$S$ 是 $A$ 的乘法子集。则
$S^{-1}I$ 是 $S^{-1}A$ 的理想，且
$\overline{S}^{-1}(A/I)$ 同构于 $S^{-1}A/S^{-1}I$，其中
$\overline{S}$ 是 $S$ 在 $A/I$ 中的像。
\end{proposition}

\begin{proof}
由于 $I$ 本身是理想，$S^{-1}I$ 为理想的事实是显然的。定义
$$
f : S^{-1}A\longrightarrow \overline{S}^{-1}(A/I), \quad x/s\mapsto
\overline{x}/\overline{s}
$$
其中 $\overline{x}$ 与 $\overline{s}$ 分别是 $x$ 与 $s$ 在
$A/I$ 中的像。本证明中继续使用类似记号。由 $S^{-1}A$ 的泛性质，
该映射良定义；其核包含 $S^{-1}I$，故它诱导映射
$$
\overline{f} : S^{-1}A/S^{-1}I \longrightarrow \overline{S}^{-1}(A/I), \quad
\overline{x/s}\mapsto \overline{x}/\overline{s}
$$

\medskip\noindent
另一方面，把 $x$ 映到 $\overline{x/1}$ 的映射
$A \to S^{-1}A/S^{-1}I$ 诱导出把 $\overline{x}$ 映到
$\overline{x/1}$ 的映射 $A/I \to S^{-1}A/S^{-1}I$。
$\overline{S}$ 的像在 $S^{-1}A/S^{-1}I$ 中可逆，故由泛性质诱导映射
$$
g : \overline{S}^{-1}(A/I) \longrightarrow S^{-1}A/S^{-1}I, \quad
\frac{\overline{x}}{\overline{s}}\mapsto \overline{x/s}
$$
显然，$\overline{f}$ 与 $g$ 互为逆映射，故二者都是同构。
\end{proof}

\noindent
现在考察子模在局部化下的行为。

\begin{lemma}
\label{algebra-lemma-submodule-localization}
$S^{-1}M$ 的任一子模 $N'$ 都具有形式 $S^{-1}N$，其中
$N\subset M$；事实上，可以取 $N$ 为 $N'$ 在 $M$ 中的逆像。
\end{lemma}

\begin{proof}
设 $N$ 为 $N'$ 在 $M$ 中的逆像。于是可见
% P01-E0275: source inclusion-direction candidate recorded in ERRATA_CANDIDATES.jsonl.
$S^{-1}N\supset N'$。为证明二者相等，取 $S^{-1}N$ 中的
$x/s$，其中 $s\in S$ 且 $x\in N$。于是 $x/1\in N'$。
% P01-E0276: source base-ring candidate recorded in ERRATA_CANDIDATES.jsonl.
由于 $N'$ 是 $S^{-1}R$-子模，有
$x/s = x/1\cdot 1/s\in N'$。证明完毕。
\end{proof}

\noindent
令 $M = A$，并令 $N = I$ 为 $A$ 的理想，得到下列推论；
它可以视为命题 \ref{algebra-proposition-localize-quotient} 第一部分的逆命题。

\begin{lemma}
\label{algebra-lemma-ideal-in-localization}
\begin{slogan}
环之局部化中的理想都是理想的局部化。
\end{slogan}
$S^{-1}A$ 的每个理想 $I'$ 都具有形式 $S^{-1}I$；
可以取 $I$ 为 $I'$ 在 $A$ 中的逆像。
\end{lemma}

\begin{proof}
立即由引理 \ref{algebra-lemma-submodule-localization} 得到。
\end{proof}

\section{内 Hom}
\label{algebra-section-hom}

\noindent
若 $R$ 是环，$M$、$N$ 是 $R$-模，则
$$
\Hom_R(M, N) = \{ \varphi : M \to N\}
$$
是从 $M$ 到 $N$ 的 $R$-线性映射所成的集合。照常令
$(\varphi + \psi)(m) = \varphi(m) + \psi(m)$，便赋予此集合 Abel 群
结构。事实上，$\Hom_R(M, N)$ 还按规则
$(x \varphi)(m) = x \varphi(m) = \varphi(xm)$ 成为 $R$-模。

\medskip\noindent
给定 $R$-模同态 $a : M \to M'$ 与 $b : N \to N'$，
可分别用 $a$ 与 $b$ 对同态作前复合及后复合。由此得到交换图
$$
\xymatrix{
\Hom_R(M', N) \ar[d]_{- \circ a} \ar[r]_{b \circ -} &
\Hom_R(M', N') \ar[d]^{- \circ a} \\
\Hom_R(M, N) \ar[r]^{b \circ -} &
\Hom_R(M, N')
}
$$
事实上，图中的映射都是 $R$-模同态。因此 $\Hom_R$ 定义了加性函子
$$
\text{Mod}_R^{opp} \times \text{Mod}_R \longrightarrow \text{Mod}_R, \quad
(M, N) \longmapsto \Hom_R(M, N)
$$

\begin{lemma}
\label{algebra-lemma-hom-exact}
正合性与 $\Hom_R$。设 $R$ 为环，$M_1$、$M_2$、$M_3$ 为
$R$-模，并设 $M_1 \to M_2$ 与 $M_2 \to M_3$ 为 $R$-模同态。
\begin{enumerate}
\item $M_1 \to M_2 \to M_3 \to 0$ 正合，当且仅当对所有 $R$-模 $N$，
$0 \to \Hom_R(M_3, N) \to \Hom_R(M_2, N) \to \Hom_R(M_1, N)$
正合。
\item $0 \to M_1 \to M_2 \to M_3$ 正合，当且仅当对所有 $R$-模 $N$，
$$
0 \to \Hom_R(N, M_1) \to \Hom_R(N, M_2) \to \Hom_R(N, M_3)
$$
正合。
\end{enumerate}
\end{lemma}

\begin{proof}
从略。
\end{proof}

\begin{lemma}
\label{algebra-lemma-hom-from-finitely-presented}
设 $R$ 为环，$M$ 为有限表示 $R$-模，$N$ 为 $R$-模。
\begin{enumerate}
\item 对 $f \in R$，有
$\Hom_R(M, N)_f = \Hom_{R_f}(M_f, N_f) = \Hom_R(M_f, N_f)$；
\item 对 $R$ 的乘法子集 $S$，有
$$
S^{-1}\Hom_R(M, N) = \Hom_{S^{-1}R}(S^{-1}M, S^{-1}N) =
\Hom_R(S^{-1}M, S^{-1}N).
$$
\end{enumerate}
\end{lemma}

\begin{proof}
(1) 是 (2) 的特殊情形。(2) 中第二个等号由引理
\ref{algebra-lemma-localization-and-modules} 得到。选取一个表示
$$
\bigoplus\nolimits_{j = 1, \ldots, m} R
\longrightarrow
\bigoplus\nolimits_{i = 1, \ldots, n} R
\to M \to 0.
$$
由引理 \ref{algebra-lemma-hom-exact}，得到正合列
$$
0 \to
\Hom_R(M, N) \to
\bigoplus\nolimits_{i = 1, \ldots, n} N
\longrightarrow
\bigoplus\nolimits_{j = 1, \ldots, m} N.
$$
关于 $S$ 作局部化并使用命题 \ref{algebra-proposition-localization-exact}，
得到正合列
$$
0 \to
S^{-1}\Hom_R(M, N) \to
\bigoplus\nolimits_{i = 1, \ldots, n} S^{-1}N
\longrightarrow
\bigoplus\nolimits_{j = 1, \ldots, m} S^{-1}N
$$
又因 $S^{-1}M$ 位于正合列
$$
\bigoplus\nolimits_{j = 1, \ldots, m} S^{-1}R
\longrightarrow
\bigoplus\nolimits_{i = 1, \ldots, n} S^{-1}R \to S^{-1}M \to 0
$$
中，该列由引理 \ref{algebra-lemma-hom-exact} 诱导出正合列
$$
0 \to
\Hom_{S^{-1}R}(S^{-1}M, S^{-1}N) \to
\bigoplus\nolimits_{i = 1, \ldots, n} S^{-1}N
\longrightarrow
\bigoplus\nolimits_{j = 1, \ldots, m} S^{-1}N
$$
而这正是上面的正合列，故结论成立。
\end{proof}

\section{有限模与有限表示模的刻画}
\label{algebra-section-colim-and-hom}

\noindent
给定环 $R$ 上的模 $N$，可以用函子 $\Hom_R(N, -)$ 来刻画
$N$ 是否为有限模或有限表示模。

\begin{lemma}
\label{algebra-lemma-characterize-finite-module-hom}
设 $R$ 为环，$N$ 为 $R$-模。下列条件等价：
\begin{enumerate}
\item $N$ 是有限 $R$-模；
\item 对 $R$-模的任意滤过余极限 $M = \colim M_i$，映射
$\colim \Hom_R(N, M_i) \to \Hom_R(N, M)$ 是单射。
\end{enumerate}
\end{lemma}

\begin{proof}
设 (1) 成立，并为 $N$ 选取生成元 $x_1, \ldots, x_m$。
若 $N \to M_i$ 是模同态，且复合
$N \to M_i \to M$ 为零，则由于
$M = \colim_{i' \geq i} M_{i'}$，对每个
$j \in \{1, \ldots, m\}$，都能找到
% P01-E0277: source article candidate recorded in ERRATA_CANDIDATES.jsonl.
某个 $i' \geq i$，使 $x_j$ 在 $M_{i'}$ 中的像为零。
由于 $x_j$ 只有有限多个，可以找到一个同时适用于所有生成元的 $i'$。
于是复合 $N \to M_i \to M_{i'}$ 为零，从而该映射是单射，即 (2) 成立。

\medskip\noindent
设 (2) 成立。对有限子集 $E \subset N$，用 $N_E \subset N$
表示由 $E$ 中元素生成的 $R$-子模。则
$0 = \colim N/N_E$ 是滤过余极限。因此，
% P01-E0278: source factorization wording candidate recorded in ERRATA_CANDIDATES.jsonl.
$\text{id} : N \to N$ 对某个 $E$ 通过 $N_E$ 分解；换言之，$N$
是有限生成的。
\end{proof}

\noindent
为便于引用，现定义模中元素之间的关系。

\begin{definition}
\label{algebra-definition-relation}
设 $R$ 为环，$M$ 为 $R$-模。设 $n \geq 0$，并对
$i = 1, \ldots, n$ 给定 $x_i \in M$。$M$ 中
$x_1, \ldots, x_n$ 之间的一个{\it 关系}，是满足
$\sum_{i = 1, \ldots, n} f_i x_i = 0$ 的元素列
$f_1, \ldots, f_n \in R$。
\end{definition}

\begin{lemma}
\label{algebra-lemma-module-colimit-fp}
设 $R$ 为环，$M$ 为 $R$-模。则 $M$ 是某个 $R$-模有向系统
$(M_i, \mu_{ij})$ 的余极限，其中每个 $M_i$ 都是有限表示 $R$-模。
\end{lemma}

\begin{proof}
考虑任意有限子集 $S \subset M$，以及 $S$ 中元素之间的任意有限关系集
$E$。于是，每个 $s \in S$ 对应于 $x_s \in M$，而每个
$e \in E$ 由元素向量 $f_{e, s} \in R$ 组成，满足
$\sum f_{e, s} x_s = 0$。令 $M_{S, E}$ 为映射
$$
R^{\# E} \longrightarrow R^{\# S}, \quad
(g_e)_{e\in E} \longmapsto (\sum g_e f_{e, s})_{s\in S}.
$$
的余核。有典范映射 $M_{S, E} \to M$。若 $S \subset S'$，
且 $E$ 中元素通过此映射对应于 $E'$ 中的关系，则有显然的映射
$M_{S, E} \to M_{S', E'}$，并且它与到 $M$ 的映射交换。
令 $I$ 为所有有序对 $(S, E)$ 组成的集合，并按上述包含关系排序。
显然，这个有向系统的余极限是 $M$。
\end{proof}

\begin{lemma}
\label{algebra-lemma-characterize-finitely-presented-module-hom}
设 $R$ 为环，$N$ 为 $R$-模。下列条件等价：
\begin{enumerate}
\item $N$ 是有限表示 $R$-模；
\item 对 $R$-模的任意滤过余极限 $M = \colim M_i$，映射
$\colim \Hom_R(N, M_i) \to \Hom_R(N, M)$ 是双射。
\end{enumerate}
\end{lemma}

\begin{proof}
设 (1) 成立，并选取正合列 $F_{-1} \to F_0 \to N \to 0$，
其中 $F_i$ 是有限自由模。于是有关于 $R$-模 $M$ 函子性的正合列
$$
0 \to \Hom_R(N, M) \to \Hom_R(F_0, M) \to \Hom_R(F_{-1}, M)
$$
。函子 $\Hom_R(F_i, M)$ 与滤过余极限交换，因为
$\Hom_R(R^{\oplus n}, M) = M^{\oplus n}$。
由于滤过余极限是正合的（引理 \ref{algebra-lemma-directed-colimit-exact}），
可知 (2) 成立。

\medskip\noindent
设 (2) 成立。由引理 \ref{algebra-lemma-module-colimit-fp}，
可把 $N = \colim N_i$ 写成滤过余极限，其中对所有 $i$，$N_i$
都是有限表示模。因此，对某个 $i$，$\text{id}_N$ 通过 $N_i$ 分解。
这意味着 $N$ 是某个有限表示 $R$-模（即 $N_i$）的直和因子，
故其也是有限表示模。
\end{proof}

\section{张量积}
\label{algebra-section-tensor-product}

\begin{definition}
\label{algebra-definition-bilinear}
设 $R$ 为环，$M, N, P$ 为三个 $R$-模。称映射
$f : M \times N \to P$（其中 $M \times N$ 仅视为两个 $R$-模的笛卡尔
积）是{\it $R$-双线性的}，如果对每个 $x \in M$，从 $N$ 到 $P$ 的映射
$y\mapsto f(x, y)$ 都是 $R$-线性的，并且对每个 $y\in N$，映射
$x\mapsto f(x, y)$ 也都是 $R$-线性的。
\end{definition}

\begin{lemma}
\label{algebra-lemma-tensor-product}
设 $M, N$ 为 $R$-模。则存在一对 $(T, g)$，其中 $T$ 是 $R$-模，
$g : M \times N \to T$ 是 $R$-双线性映射，并具有下列泛性质：
对任意 $R$-模 $P$ 及任意 $R$-双线性映射
$f : M \times N \to P$，存在唯一的 $R$-线性映射
$\tilde{f} : T \to P$，使得 $f = \tilde{f} \circ g$。
换言之，下图交换：
$$
\xymatrix{
M \times N \ar[rr]^f \ar[dr]_g & & P\\
& T \ar[ur]_{\tilde f}
}
$$
此外，若 $(T, g)$ 与 $(T', g')$ 是两对具有此性质的对象，
则存在唯一同构 $j : T \to T'$，使 $j\circ g = g'$。
\end{lemma}

\noindent
满足上述泛性质的 $R$-模 $T$ 称为 $R$-模 $M$ 与 $N$ 的
{\it 张量积}，记为 $M \otimes_R N$。

\begin{proof}
先证明这样的 $R$-模 $T$ 存在。设 $M, N$ 为 $R$-模。
令 $T$ 为商模 $P/Q$，其中 $P$ 是自由 $R$-模
$R^{(M \times N)}$，而 $Q$ 是由下列各类元素生成的 $R$-模
（$x\in M, y\in N$）：
\begin{align*}
(x + x', y) - (x, y) - (x', y), \\
(x, y + y') - (x, y) - (x, y'), \\
(ax, y) - a(x, y), \\
(x, ay) - a(x, y)
\end{align*}
用 $\pi : M \times N \to T$ 表示自然映射。验证双线性条件时，
上述关系即说明该映射为 $R$-双线性的。
% P01-E0280: source wording candidate recorded in ERRATA_CANDIDATES.jsonl.
把像记为 $\pi(x, y) = x \otimes y$；这些元素生成 $T$。
现设 $f : M \times N \to P$ 为 $R$-双线性映射，则可把映射
$f'(x \otimes y) = f(x, y)$ 线性延拓，定义
$f' : T \to P$。显然 $f = f'\circ \pi$。此外，$f'$ 由其在生成集
$\{x \otimes y : x\in M, y\in N\}$ 上的取值唯一确定。
设另有一对 $(T', g')$ 满足同样性质。则存在唯一的
$j : T \to T'$ 以及 $j' : T' \to T$，使
$g' = j\circ g$、$g = j'\circ g'$。
% P01-E0279: source composition-order defect recorded in ERRATA_CANDIDATES.jsonl.
但映射 $(j\circ j') \circ g$ 与 $g$ 都满足泛性质，故由唯一性二者相等，
从而 $j'\circ j$ 是 $T$ 上的恒等映射。类似地，
$(j'\circ j) \circ g' = g'$，且 $j\circ j'$ 是 $T'$ 上的恒等映射。
所以 $j$ 是同构。
\end{proof}

\begin{lemma}
\label{algebra-lemma-flip-tensor-product}
设 $M, N, P$ 为 $R$-模，则双线性映射
\begin{align*}
(x, y) & \mapsto y \otimes x\\
(x + y, z) & \mapsto x \otimes z + y \otimes z\\
(r, x) & \mapsto rx
\end{align*}
分别诱导唯一同构
\begin{align*}
M \otimes_R N & \to N \otimes_R M, \\
(M\oplus N)\otimes_R P & \to (M \otimes_R P)\oplus(N \otimes_R P),  \\
R \otimes_R M & \to M
\end{align*}
\end{lemma}

\begin{proof}
从略。
\end{proof}

\noindent
可以把两个 $R$-模的张量积推广到有限多个 $R$-模，并在多重张量积与
多线性映射之间建立对应。用几乎相同的构造，可以证明下列结论。

\begin{lemma}
\label{algebra-lemma-multilinear}
设 $M_1, \ldots, M_r$ 为 $R$-模。则存在一对 $(T, g)$，
由一个 $R$-模 T 及一个具有如下泛性质的 $R$-多线性映射
$g : M_1\times \ldots \times M_r \to T$ 组成：
对任意 $R$-多线性映射 $f : M_1\times \ldots \times M_r \to P$，
存在唯一的 $R$-模同态 $f' : T \to P$，使 $f'\circ g = f$。
这样的模 $T$ 在唯一同构意义下唯一。把它记为
$M_1\otimes_R \ldots \otimes_R M_r$，并把泛多线性映射记为
$(m_1, \ldots, m_r) \mapsto m_1 \otimes \ldots \otimes m_r$。
\end{lemma}

\begin{proof}
从略。
\end{proof}

\begin{lemma}
\label{algebra-lemma-transitive}
同态
$$
(M \otimes_R N)\otimes_R P \to
M \otimes_R N \otimes_R P \to
M \otimes_R (N \otimes_R P)
$$
由下式给出：
$f((x \otimes y)\otimes z) = x \otimes y \otimes z$ 以及
$g(x \otimes y \otimes z) = x \otimes (y \otimes z)$，
其中 $x\in M, y\in N, z\in P$。这些同态良定义且为同构。
\end{lemma}

\begin{proof}
证明 $f$ 良定义且为同构；对 $g$ 的证明完全类似。固定任意
$z\in P$，则映射 $(x, y)\mapsto x \otimes y \otimes z$，
$x\in M, y\in N$，关于 $x$、$y$ 是 $R$-双线性的，因而诱导同态
$f_z : M \otimes N \to M \otimes N \otimes P$，满足
$f_z(x \otimes y) = x \otimes y \otimes z$。
再考虑由 $(w, z)\mapsto f_z(w)$ 给出的映射
$(M \otimes N)\times P \to M \otimes N \otimes P$。
此映射为 $R$-双线性的，故诱导
$f : (M \otimes_R N)\otimes_R P \to M \otimes_R N \otimes_R P$，
且 $f((x \otimes y)\otimes z) = x \otimes y \otimes z$。
为构造其逆映射，注意映射
$\pi : M \times N \times P \to (M \otimes N)\otimes P$
是 $R$-三线性的。因此，它诱导一个符合泛性质的 $R$-线性映射
$h : M \otimes N \otimes P \to (M \otimes N)\otimes P$。
在此可见，$h(x \otimes y \otimes z) = (x \otimes y)\otimes z$。
由 $f$ 与 $h$ 的显式表达式，$f\circ h$ 与 $h\circ f$ 分别是
$M \otimes N \otimes P$ 与 $(M \otimes N)\otimes P$ 上的恒等映射，
故 $f$ 即所求同构。
\end{proof}

\noindent
归纳可见，这一结论推广到多重张量积。再结合引理
\ref{algebra-lemma-flip-tensor-product}，可知 $R$-模范畴上的张量积运算
具有结合性、交换性与分配性。

\begin{definition}
\label{algebra-definition-bimodule}
若 Abel 群 $N$ 既是 $A$-模又是 $B$-模，且对所有 $a \in A$、
$b \in B$，乘以 $a$ 与乘以 $b$ 的作用交换，即对所有
$n \in N$ 有 $b(an) = a(bn)$，则称之为
{\it $(A, B)$-双模}。在这种情形下，通常把 $B$-作用写在右边：
对 $b \in B$ 与 $n \in N$，用 $nb$ 表示 $n$ 乘以 $b$ 的结果。
采用此约定，上述相容性即为：对所有 $a\in A, b\in B, x\in N$，
有 $(ax)b = a(xb)$。用简记 $_AN_B$ 表示 $(A, B)$-双模 $N$。
\end{definition}

\begin{lemma}
\label{algebra-lemma-tensor-with-bimodule}
设 $M$ 为 $A$-模，$P$ 为 $B$-模，$N$ 为 $(A, B)$-双模。
则模 $(M \otimes_A N)\otimes_B P$ 与
$M \otimes_A(N \otimes_B P)$ 都可以赋予 $(A, B)$-双模结构，而且
$$
(M \otimes_A N)\otimes_B P \cong M \otimes_A(N \otimes_B P).
$$
\end{lemma}

\begin{proof}
先验地，$M \otimes_A N$ 是 $A$-模；但可以用规则
$$
(x \otimes y)b = x \otimes yb, \quad x\in M, y\in N, b\in B
$$
赋予它 $B$-模结构。于是 $M \otimes_A N$ 成为 $(A, B)$-双模。
$N \otimes_B P$ 的情形类似，进而
$(M \otimes_A N)\otimes_B P$ 与 $M \otimes_A(N \otimes_B P)$
也同样如此。由引理 \ref{algebra-lemma-transitive}，这两个模通过同一个映射，
% P01-E0281: source duplicated-word candidate recorded in ERRATA_CANDIDATES.jsonl.
既作为 $A$-模也作为 $B$-模同构。
\end{proof}

\begin{lemma}
\label{algebra-lemma-hom-from-tensor-product}
对任意三个 $R$-模 $M, N, P$，有
$$
\Hom_R(M \otimes_R N, P) \cong \Hom_R(M, \Hom_R(N, P))
$$
\end{lemma}

\begin{proof}
$R$-线性映射 $\hat{f}\in \Hom_R(M \otimes_R N, P)$ 对应于
$R$-双线性映射 $f : M \times N \to P$。对每个 $x\in M$，
由泛性质，映射 $y\mapsto f(x, y)$ 是 $R$-线性的。因此，$f$
对应于映射 $\phi_f : M \to \Hom_R(N, P)$。该映射为 $R$-线性的，
因为
$$
\phi_f(ax + y)(z) =
f(ax + y, z) = af(x, z)+f(y, z) =
(a\phi_f(x)+\phi_f(y))(z),
$$
其中任意 $a \in R$、$x \in M$、$y \in M$ 及 $z \in N$。
反之，任意 $f \in \Hom_R(M, \Hom_R(N, P))$ 定义一个 $R$-双线性映射
$M \times N \to P$，即 $(x, y)\mapsto f(x)(y)$。
故这给出两个模 $\Hom_R(M \otimes_R N, P)$ 与
$\Hom_R(M, \Hom_R(N, P))$ 之间的自然一一对应。
\end{proof}

\begin{lemma}[张量积与余极限交换]
\label{algebra-lemma-tensor-products-commute-with-limits}
设 $(M_i, \mu_{ij})$ 是预序集 $I$ 上的系统，$N$ 是 $R$-模。则
$$
\colim (M_i \otimes N) \cong (\colim M_i)\otimes N.
$$
此外，此同构由同态
$\mu_i \otimes 1: M_i \otimes N \to M \otimes N$
诱导，其中 $M = \colim_i M_i$，自然映射为 $\mu_i : M_i \to M$。
\end{lemma}

\begin{proof}
第一种证明。由引理 \ref{algebra-lemma-hom-from-tensor-product}，
函子 $M' \mapsto M' \otimes_R N$ 是函子
$N' \mapsto \Hom_R(N, N')$ 的左伴随。因此，
$M' \mapsto M' \otimes_R N$ 与所有余极限交换；见范畴论，引理
\ref{categories-lemma-adjoint-exact}。

\medskip\noindent
第二种直接证明。令 $P = \colim (M_i \otimes N)$，其余投影为
$\lambda_i : M_i \otimes N \to P$；令 $M = \colim M_i$，其余投影为
$\mu_i : M_i \to M$。则对所有 $i\leq j$，下图交换：
$$
\xymatrix{
M_i \otimes N \ar[r]_{\mu_i \otimes 1} \ar[d]_{\mu_{ij} \otimes 1} &
M \otimes N \ar[d]^{\text{id}} \\
M_j \otimes N \ar[r]^{\mu_j \otimes 1} &
M \otimes N
}
$$
由引理 \ref{algebra-lemma-homomorphism-limit}，这些映射诱导唯一同态
$\psi : P \to M \otimes N$，使
$\mu_i \otimes 1 = \psi \circ \lambda_i$。

\medskip\noindent
为构造逆映射，对每个 $i\in I$，都有典范的 $R$-双线性映射
$g_i : M_i \times N \to M_i \otimes N$。这诱导唯一映射
$\widehat{\phi} : M \times N \to P$，使
$\widehat{\phi} \circ (\mu_i \times 1) = \lambda_i \circ g_i$。
它为 $R$-双线性的，故诱导 $R$-线性映射
$\phi : M \otimes N \to P$。由下列交换图
$$
\xymatrix{
M_i \times N \ar[r]^{g_i} \ar[d]^{\mu_i \times \text{id}} &
M_i \otimes N\ar[r]_{\text{id}} \ar[d]_{\lambda_i} &
M_i \otimes N \ar[d]_{\mu_i \otimes \text{id}} \ar[rd]^{\lambda_i} \\
M \times N \ar[r]^{\widehat{\phi}} &
P \ar[r]^{\psi} & M \otimes N \ar[r]^{\phi} & P
}
$$
可知 $\psi\circ\widehat{\phi} = g$，其中 g 是典范 $R$-双线性映射
$g : M \times N \to M \otimes N$。故 $\psi\circ\phi$ 是
$M \otimes N$ 上的恒等映射。由右侧方形与三角形，
$\phi\circ\psi$ 也是 $P$ 上的恒等映射。
\end{proof}

\begin{lemma}
\label{algebra-lemma-tensor-product-exact}
设
\begin{align*}
M_1\xrightarrow{f} M_2\xrightarrow{g} M_3 \to 0
\end{align*}
为 $R$-模及同态的正合列，并设 $N$ 为任意 $R$-模。则序列
\begin{equation}
\label{algebra-equation-2ndex}
M_1\otimes N\xrightarrow{f \otimes 1} M_2\otimes N \xrightarrow{g \otimes 1}
M_3\otimes N \to 0
\end{equation}
正合。换言之，函子 $- \otimes_R N$ 是{\it 右正合的}；也就是说，
把原右正合列中的每一项与之作张量积会保持正合性。
\end{lemma}

\begin{proof}
对每个 $R$-模 $P$，向第一个正合列应用函子
$\Hom(-, \Hom(N, P))$，得到
$$
0 \to
\Hom(M_3, \Hom(N, P)) \to
\Hom(M_2, \Hom(N, P)) \to
\Hom(M_1, \Hom(N, P))
$$
它由引理 \ref{algebra-lemma-hom-exact} (1) 正合。
由引理 \ref{algebra-lemma-hom-from-tensor-product}，这变为序列
$$
0 \to \Hom(M_3 \otimes N, P) \to
\Hom(M_2 \otimes N, P) \to \Hom(M_1 \otimes N, P)
$$
因此也正合。再次使用引理 \ref{algebra-lemma-hom-exact} (1)，
即得到所求正合列。
\end{proof}

\begin{remark}
\label{algebra-remark-tensor-product-not-exact}
然而，张量积一般并不保持正合列。换言之，若
$M_1 \to M_2 \to M_3$ 正合，则对任意 $R$-模 $N$，
$M_1 \otimes N \to M_2 \otimes N \to M_3 \otimes N$
未必正合。
\end{remark}



\begin{example}
\label{algebra-example-tensor-product-not-exact}
考虑单射 $2 : \mathbf{Z}\to \mathbf{Z}$，把它视为
$\mathbf{Z}$-模同态。令 $N = \mathbf{Z}/2$。则诱导映射
$\mathbf{Z} \otimes \mathbf{Z}/2 \to \mathbf{Z} \otimes \mathbf{Z}/2$
不是单射。事实上，对
$x \otimes y\in \mathbf{Z} \otimes \mathbf{Z}/2$，
$$
(2 \otimes 1)(x \otimes y) = 2x \otimes y = x \otimes 2y = x \otimes 0 = 0
$$
因此，诱导映射是零映射，而 $\mathbf{Z} \otimes N\neq 0$。
\end{example}

\begin{remark}
\label{algebra-remark-flat-module}
对 $R$-模 $N$，若函子 $-\otimes_R N$ 正合，即与 $N$ 作张量积保持所有
正合列，则称 $N$ 为
% P01-E0282: source article candidate recorded in ERRATA_CANDIDATES.jsonl.
{\it 平坦} $R$-模。我们将在第 \ref{algebra-section-flat} 节进一步讨论这一概念。
\end{remark}

\begin{lemma}
\label{algebra-lemma-tensor-finiteness}
设 $R$ 为环，$M$、$N$ 为 $R$-模。
\begin{enumerate}
\item 若 $N$ 与 $M$ 都是有限模，则 $M \otimes_R N$ 也是有限模。
\item 若 $N$ 与 $M$ 都是有限表示模，则 $M \otimes_R N$ 也是有限表示模。
\end{enumerate}
\end{lemma}

\begin{proof}
设 $M$ 为有限模。选取表示
$0 \to K \to R^{\oplus n} \to M \to 0$。由引理
\ref{algebra-lemma-tensor-product-exact}，这给出正合列
$K \otimes_R N \to N^{\oplus n} \to M \otimes_R N \to 0$。
所以，若 $N$ 也有限，则 $M \otimes_R N$ 是有限模的商，因而有限；
见引理 \ref{algebra-lemma-extension}。类似地，若 $N$ 与 $M$ 都是有限表示模，
则 $K$ 是有限模，而 $M \otimes_R N$ 是有限表示模
$N^{\oplus n}$ 关于有限模 $K \otimes_R N$ 的商，因而是有限表示模；
见引理 \ref{algebra-lemma-extension}。
\end{proof}

\begin{lemma}
\label{algebra-lemma-tensor-localization}
设 $M$ 为 $R$-模。则 $S^{-1}R$-模 $S^{-1}M$ 与
$S^{-1}R \otimes_R M$ 典范同构，且典范同构
$f : S^{-1}R \otimes_R M \to S^{-1}M$ 由下式给出：
$$
f((a/s) \otimes m) = am/s, \forall a \in R, m \in M, s \in S
$$
\end{lemma}

\begin{proof}
显然，由 $f'(a/s, m) = am/s$ 给出的映射
$f' : S^{-1}R \times M \to S^{-1}M$ 是双线性的；
故由泛性质，该映射诱导引理陈述中的唯一 $S^{-1}R$-模同态
$f : S^{-1}R \otimes_R M \to S^{-1}M$。
事实上，$S^{-1}M$ 的每个元素都具有形式 $m/s$，其中
$m\in M, s\in S$；而 $S^{-1}R \otimes_R M$ 的每个元素都具有形式
$1/s \otimes m$。为看出后一事实，把 $S^{-1}R \otimes_R M$ 中的元素写成
$$
\sum_k \frac{a_k}{s_k} \otimes m_k =
\sum_k \frac{a_k t_k}{s} \otimes m_k =
\frac{1}{s} \otimes \sum_k {a_k t_k}m_k = \frac{1}{s} \otimes m
$$
% P01-E0283: source capitalization candidate recorded in ERRATA_CANDIDATES.jsonl.
其中 $m = \sum_k {a_k t_k}m_k$。于是 $f$ 显然是满射；
若 $f(\frac{1}{s} \otimes m) = m/s = 0$，则存在 $t \in S$，
使得 $M$ 中 $tm = 0$。于是
$$
\frac{1}{s} \otimes m = \frac{1}{st} \otimes tm = \frac{1}{st} \otimes 0 = 0
$$
因此 $f$ 是单射。
\end{proof}

\begin{lemma}
\label{algebra-lemma-tensor-product-localization}
设 $M, N$ 为 $R$-模，则有典范 $S^{-1}R$-模同构
$f : S^{-1}M \otimes_{S^{-1}R}S^{-1}N \to S^{-1}(M \otimes_R N)$，
由下式给出：
$$
f((m/s)\otimes(n/t)) = (m \otimes n)/st
$$
\end{lemma}

\begin{proof}
反复使用引理 \ref{algebra-lemma-tensor-with-bimodule} 与引理
\ref{algebra-lemma-tensor-localization}，并注意 $S^{-1}R$ 是
$(R, S^{-1}R)$-双模，可知这两个 $S^{-1}R$-模同构：
\begin{align*}
S^{-1}(M \otimes_R N) & \cong S^{-1}R \otimes_R (M \otimes_R N)\\
 & \cong S^{-1}M \otimes_R N\\
 & \cong (S^{-1}M \otimes_{S^{-1}R}S^{-1}R)\otimes_R N\\
 & \cong S^{-1}M \otimes_{S^{-1}R}(S^{-1}R \otimes_R N)\\
 & \cong S^{-1}M \otimes_{S^{-1}R}S^{-1}N
\end{align*}
容易看出，这一同构正是引理所述同构。
\end{proof}
\section{张量代数}
\label{algebra-section-tensor-algebra}

\noindent
设 $R$ 为环，$M$ 为 $R$-模。定义{\it $M$ 在 $R$ 上的张量代数}为
非交换 $R$-代数
$$
\text{T}(M) = \text{T}_R(M) =
\bigoplus\nolimits_{n \geq 0} \text{T}^n(M)
$$
其中
$\text{T}^0(M) = R$、
$\text{T}^1(M) = M$、
$\text{T}^2(M) = M \otimes_R M$、
$\text{T}^3(M) = M \otimes_R M \otimes_R M$，依此类推。
乘法由纯张量上的下列规则定义：
$$
(x_1 \otimes x_2 \otimes \ldots \otimes x_n)
\cdot
(y_1 \otimes y_2 \otimes \ldots \otimes y_m)
=
x_1 \otimes x_2 \otimes \ldots \otimes x_n \otimes
y_1 \otimes y_2 \otimes \ldots \otimes y_m
$$
再按线性延拓。

\medskip\noindent
定义{\it $M$ 在 $R$ 上的外代数 $\wedge(M)$}为 $\text{T}(M)$
关于由元素 $x \otimes x \in \text{T}^2(M)$ 生成的双边理想的商。
纯张量 $x_1 \otimes \ldots \otimes x_n$ 在 $\wedge^n(M)$ 中的像记为
$x_1 \wedge \ldots \wedge x_n$。这些元素生成 $\wedge^n(M)$；
它们关于每个 $x_i$ 都是 $R$-线性的，并且当两个 $x_i$ 相等时为零
（即作为 $x_1, x_2, \ldots, x_n$ 的函数是交替的）。
$\wedge(M)$ 上的乘法是分次交换的；即每个 $x \in M$ 与 $y \in M$
都满足 $x \wedge y = - y \wedge x$。

\medskip\noindent
例如，设 $M = Rx_1 \oplus \ldots \oplus Rx_n$ 是有限自由模。
此时 $\wedge(M)$ 是自由 $R$-模，其一组基为元素
$$
x_{i_1} \wedge \ldots \wedge x_{i_r}
$$
其中 $0 \leq r \leq n$ 且 $1 \leq i_1 < i_2 < \ldots < i_r \leq n$。

\medskip\noindent
定义{\it $M$ 在 $R$ 上的对称代数 $\text{Sym}(M)$}为 $\text{T}(M)$
关于由元素 $x \otimes y - y \otimes x \in \text{T}^2(M)$
生成的双边理想的商。纯张量 $x_1 \otimes \ldots \otimes x_n$
在 $\text{Sym}^n(M)$ 中的像简记为 $x_1 \ldots x_n$。
这些元素生成 $\text{Sym}^n(M)$，并关于每个 $x_i$ 都是 $R$-线性的。
若元素列 $x_1, \ldots, x_n$ 是元素列 $x_1', \ldots, x_n'$ 的一个排列，
则 $x_1 \ldots x_n = x_1' \ldots x_n'$。因此，$\text{Sym}(M)$ 是交换的。

\medskip\noindent
仍以有限自由模 $M = Rx_1 \oplus \ldots \oplus Rx_n$ 为例。
此时 $\text{Sym}(M) = R[x_1, \ldots, x_n]$ 是多项式代数。

\begin{lemma}
\label{algebra-lemma-free-tensor-algebra}
设 $R$ 为环，$M$ 为 $R$-模。若 $M$ 是自由 $R$-模，
则其每个对称幂与外幂也是自由模。
\end{lemma}

\begin{proof}
从略；有限自由的情形见上文。
\end{proof}

\begin{lemma}
\label{algebra-lemma-presentation-sym-exterior}
设 $R$ 为环，并设
$M_2 \to M_1 \to M \to 0$ 为 $R$-模的正合列。则有正合列
$$
M_2 \otimes_R \text{Sym}^{n - 1}(M_1)
\to
\text{Sym}^n(M_1)
\to
\text{Sym}^n(M)
\to
0
$$
以及类似的正合列
$$
M_2 \otimes_R \wedge^{n - 1}(M_1)
\to
\wedge^n(M_1)
\to
\wedge^n(M)
\to
0
$$
\end{lemma}

\begin{proof}
从略。
\end{proof}

\begin{lemma}
\label{algebra-lemma-present-sym-wedge}
设 $R$ 为环，$M$ 为 $R$-模。设 $x_i$（$i \in I$）是 $M$
作为 $R$-模的一组给定生成元，并设 $n \geq 2$。存在典范正合列
$$
\bigoplus_{1 \leq j_1 < j_2 \leq n}
\bigoplus_{i_1, i_2 \in I}
\text{T}^{n - 2}(M)
\oplus
\bigoplus_{1 \leq j_1 < j_2 \leq n}
\bigoplus_{i \in I}
\text{T}^{n - 2}(M)
\to
\text{T}^n(M)
\to
\wedge^n(M)
\to
0
$$
其中，第一个直和项中的纯张量
$m_1 \otimes \ldots \otimes m_{n - 2}$ 映为
\begin{align*}
\underbrace{
m_1 \otimes \ldots \otimes x_{i_1} \otimes \ldots
\otimes x_{i_2} \otimes \ldots \otimes m_{n - 2}
}_{\text{其中 } x_{i_1} \text{ 且 } x_{i_2}
\text{ occupying slots } j_1 \text{ 且 } j_2
\text{ in the tensor}} \\
+
\underbrace{
m_1 \otimes \ldots \otimes x_{i_2} \otimes \ldots
\otimes x_{i_1} \otimes \ldots \otimes m_{n - 2}
}_{\text{其中 } x_{i_2} \text{ 且 } x_{i_1}
\text{ occupying slots } j_1 \text{ 且 } j_2
\text{ in the tensor}}
\end{align*}
而第二个直和项中的 $m_1 \otimes \ldots \otimes m_{n - 2}$ 映为
$$
\underbrace{
m_1 \otimes \ldots \otimes x_i \otimes \ldots
\otimes x_i \otimes \ldots \otimes m_{n - 2}
}_{\text{其中 } x_{i} \text{ 且 } x_{i}
\text{ occupying slots } j_1 \text{ 且 } j_2
\text{ in the tensor}}
$$
还有典范正合列
$$
\bigoplus_{1 \leq j_1 < j_2 \leq n}
\bigoplus_{i_1, i_2 \in I}
\text{T}^{n - 2}(M)
\to
\text{T}^n(M)
\to
\text{Sym}^n(M)
\to
0
$$
其中，纯张量 $m_1 \otimes \ldots \otimes m_{n - 2}$ 映为
\begin{align*}
\underbrace{
m_1 \otimes \ldots \otimes x_{i_1} \otimes \ldots
\otimes x_{i_2} \otimes \ldots \otimes m_{n - 2}
}_{\text{其中 } x_{i_1} \text{ 且 } x_{i_2}
\text{ occupying slots } j_1 \text{ 且 } j_2
\text{ in the tensor}} \\
-
\underbrace{
m_1 \otimes \ldots \otimes x_{i_2} \otimes \ldots
\otimes x_{i_1} \otimes \ldots \otimes m_{n - 2}
}_{\text{其中 } x_{i_2} \text{ 且 } x_{i_1}
\text{ occupying slots } j_1 \text{ 且 } j_2
\text{ in the tensor}}
\end{align*}
\end{lemma}

\begin{proof}
从略。
\end{proof}

\begin{lemma}
\label{algebra-lemma-present-wedge}
设 $A \to B$ 为环同态，$M$ 为 $B$-模，并设 $n > 1$。
$A$-线性映射
$M \otimes_A \ldots \otimes_A M \to \wedge^n_B(M)$ 的核，
作为 $A$-模由下列元素生成：满足 $m_i = m_j$ 的
$m_1 \otimes \ldots \otimes m_n$，其中
$i \not = j$、$m_1, \ldots, m_n \in M$；以及元素
$m_1 \otimes \ldots \otimes bm_i \otimes \ldots \otimes m_n -
m_1 \otimes \ldots \otimes bm_j \otimes \ldots \otimes m_n$，
其中 $i \not = j$、$m_1, \ldots, m_n \in M$ 且 $b \in B$。
\end{lemma}

\begin{proof}
从略。
\end{proof}

\begin{lemma}
\label{algebra-lemma-colimit-tensor-algebra}
\begin{slogan}
取张量代数与滤过余极限交换。
\end{slogan}
设 $R$ 为环，$M_i$ 为 $R$-模的有向系统。则
$\colim_i \text{T}(M_i) = \text{T}(\colim_i M_i)$，
对称代数与外代数也有类似结论。
\end{lemma}

\begin{proof}
从略。提示：应用引理 \ref{algebra-lemma-tensor-products-commute-with-limits}。
\end{proof}

\begin{lemma}
\label{algebra-lemma-tensor-algebra-localization}
设 $R$ 为环，$S \subset R$ 为乘法子集。则对任意 $R$-模 $M$，
% P01-E0284: source notation candidate recorded in ERRATA_CANDIDATES.jsonl.
$S^{-1}T_R(M) = T_{S^{-1}R}(S^{-1}M)$。对称代数与外代数也类似。
\end{lemma}

\begin{proof}
从略。提示：应用引理 \ref{algebra-lemma-tensor-product-localization}。
\end{proof}

\section{基变换}
\label{algebra-section-base-change}

\noindent
如下正式引入代数中的基变换。

\begin{definition}
\label{algebra-definition-base-change}
设 $\varphi : R \to S$ 为环同态，$M$ 为 $S$-模，并设
$R \to R'$ 为任意环同态。$\varphi$ 沿 $R \to R'$ 的
{\it 基变换}是环同态 $R' \to S \otimes_R R'$。在这种情形下，
常记 $S' = S \otimes_R R'$。$S$-模 $M$ 的{\it 基变换}是
$S'$-模 $M \otimes_R R'$。
\end{definition}

\noindent
若对某族变量 $x_i$（$i \in I$）及某族多项式
$f_j \in R[x_i]$（$j \in J$），有 $S = R[x_i]/(f_j)$，则
$S \otimes_R R' = R'[x_i]/(f'_j)$，其中 $f'_j \in R'[x_i]$
是 $f_j$ 在由 $R \to R'$ 诱导的映射
$R[x_i] \to R'[x_i]$ 下的像。这个简单观察是理解基变换的关键。

\begin{lemma}
\label{algebra-lemma-base-change-finiteness}
设 $R \to S$ 为环同态，$M$ 为 $S$-模。设 $R \to R'$ 为环同态，
并令 $S' = S \otimes_R R'$ 与 $M' = M \otimes_R R'$ 为相应基变换。
\begin{enumerate}
\item 若 $M$ 是有限 $S$-模，则基变换 $M'$ 是有限 $S'$-模。
\item 若 $M$ 是有限表示 $S$-模，则基变换 $M'$ 是有限表示 $S'$-模。
\item 若 $R \to S$ 为有限型，则基变换 $R' \to S'$ 也为有限型。
\item 若 $R \to S$ 为有限表示，则基变换 $R' \to S'$ 也为有限表示。
\end{enumerate}
\end{lemma}

\begin{proof}
(1) 的证明。取满的 $S$-线性映射
$S^{\oplus n} \to M \to 0$。由引理
\ref{algebra-lemma-flip-tensor-product} 与 \ref{algebra-lemma-tensor-product-exact}，
与 $R^\prime$ 作张量积后得到满射
${S^\prime}^{\oplus n} \to M^\prime \rightarrow 0$，
故 $M^\prime$ 是有限生成 $S^\prime$-模。
(2) 的证明。取表示
$S^{\oplus m} \to S^{\oplus n} \to M \to 0$。
由引理 \ref{algebra-lemma-flip-tensor-product} 与
\ref{algebra-lemma-tensor-product-exact}，与 $R^\prime$ 作张量积后得到
$S^\prime$-模 $M^\prime$ 的有限表示
${S^\prime}^{\oplus m} \to {S^\prime}^{\oplus n} \to M^\prime \to 0$。
(3) 的证明。由引理前的观察即可；按假设可取 $I$ 为有限集。
(4) 的证明。同理由引理前的观察即可；按假设可取 $I$ 与 $J$ 为有限集。
\end{proof}

\noindent
设 $\varphi : R \to S$ 为环同态。给定 $S$-模 $N$，可按规则
$r \cdot n = \varphi(r)n$ 得到 $R$-模 $N_R$。
这有时称为把 $N$ {\it 限制}到 $R$。

\begin{lemma}
\label{algebra-lemma-adjoint-tensor-restrict}
设 $R \to S$ 为环同态。函子
$\text{Mod}_S \to \text{Mod}_R$，$N \mapsto N_R$（限制），
与函子 $\text{Mod}_R \to \text{Mod}_S$，
$M \mapsto M \otimes_R S$（基变换）互为伴随。公式为
$$
\Hom_R(M, N_R) = \Hom_S(M \otimes_R S, N)
$$
\end{lemma}

\begin{proof}
若 $\alpha : M \to N_R$ 是 $R$-模同态，则按规则
$\alpha'(m \otimes s) = s\alpha(m)$ 定义
$\alpha' : M \otimes_R S \to N$。若
$\beta : M \otimes_R S \to N$ 是 $S$-模同态，则按规则
$\beta'(m) = \beta(m \otimes 1)$ 定义 $\beta' : M \to N_R$。
这两个构造互逆的验证从略。
\end{proof}

\noindent
上述引理说明限制函子有左伴随，即基变换。它也有右伴随。

\begin{lemma}
\label{algebra-lemma-adjoint-hom-restrict}
设 $R \to S$ 为环同态。函子
$\text{Mod}_S \to \text{Mod}_R$，$N \mapsto N_R$（限制），
与函子 $\text{Mod}_R \to \text{Mod}_S$，
$M \mapsto \Hom_R(S, M)$ 互为伴随。公式为
$$
\Hom_R(N_R, M) = \Hom_S(N, \Hom_R(S, M))
$$
\end{lemma}

\begin{proof}
若 $\alpha : N_R \to M$ 是 $R$-模同态，则按规则
$\alpha'(n) = (s \mapsto \alpha(sn))$ 定义
$\alpha' : N \to \Hom_R(S, M)$。若
$\beta : N \to \Hom_R(S, M)$ 是 $S$-模同态，则按规则
$\beta'(n) = \beta(n)(1)$ 定义 $\beta' : N_R \to M$。
这两个构造互逆的验证从略。
\end{proof}

\begin{lemma}
\label{algebra-lemma-hom-from-tensor-product-variant}
设 $R \to S$ 为环同态。给定 $S$-模 $M, N$ 及 $R$-模 $P$，有
$$
\Hom_R(M \otimes_S N, P) = \Hom_S(M, \Hom_R(N, P))
$$
\end{lemma}

\begin{proof}
这可以直接证明，也可由引理 \ref{algebra-lemma-adjoint-hom-restrict}
与 \ref{algebra-lemma-hom-from-tensor-product} 推出。事实上，
\begin{align*}
\Hom_R(M \otimes_S N, P)
& =
\Hom_S(M \otimes_S N, \Hom_R(S, P)) \\
& =
\Hom_S(M, \Hom_S(N, \Hom_R(S, P))) \\
& =
\Hom_S(M, \Hom_R(N, P))
\end{align*}
正合所求。
\end{proof}

\section{杂项}
\label{algebra-section-miscellany}

\noindent
本节的证明不应引用基本概念一节，即第 \ref{algebra-section-rings-basic} 节之外
的任何结果。

\begin{lemma}
\label{algebra-lemma-product-ideals-in-prime}
设 $R$ 为环，$I$、$J$ 为两个理想，$\mathfrak p$ 是包含乘积 $IJ$
的素理想。则 $\mathfrak{p}$ 包含 $I$ 或 $J$。
\end{lemma}

\begin{proof}
反设不然，取 $x \in I \setminus \mathfrak p$ 与
$y \in J \setminus \mathfrak p$。其乘积属于
$IJ \subset \mathfrak p$，这与 $\mathfrak p$ 为素理想矛盾。
\end{proof}

\begin{lemma}[素理想回避]
\label{algebra-lemma-silly}
\begin{slogan}
1. 在仿射概形中，若有限多个点包含于一个开子集中，则它们包含于一个
更小的主开子集中。
2. 仿射开集在仿射概形的给定有限集的邻域中共尾。
\end{slogan}
设 $R$ 为环。设 $I_i \subset R$（$i = 1, \ldots, r$）与
$J \subset R$ 为理想。假设
\begin{enumerate}
\item 对 $i = 1, \ldots, r$，有 $J \not\subset I_i$；且
\item 除至多两个外，所有 $I_i$ 都是素理想。
\end{enumerate}
则存在 $x \in J$，使得对所有 $i$ 都有 $x\not\in I_i$。
\end{lemma}

\begin{proof}
当 $r = 1$ 时结论成立。若 $r = 2$，取 $x, y \in J$，
使 $x \not \in I_1$ 且 $y \not \in I_2$。除非
$x \in I_2$ 且 $y \in I_1$，否则结论已得。在后一情形，
$x + y$ 不可能属于 $I_1$（否则 $x + y - y \in I_1$），
也不可能属于 $I_2$。

\medskip\noindent
设 $r \geq 3$，并假定结论对 $r - 1$ 成立。重新编号后，
可设 $I_r$ 为素理想。还可设各 $I_i$ 之间不存在包含关系。
取 $x \in J$，使对所有 $i = 1, \ldots, r - 1$ 都有
$x \not \in I_i$。若 $x \not\in I_r$，结论已得；故设
$x \in I_r$。若 $J I_1 \ldots I_{r - 1} \subset I_r$，
则由引理 \ref{algebra-lemma-product-ideals-in-prime} 有
$J \subset I_r$，矛盾。取
$y \in J I_1 \ldots I_{r - 1}$ 且 $y \not \in I_r$。
则 $x + y$ 满足要求。
\end{proof}

\begin{lemma}
\label{algebra-lemma-silly-silly}
设 $R$ 为环，$x \in R$，$I \subset R$ 为理想，并设
$\mathfrak p_i$（$i = 1, \ldots, r$）为素理想。
假设对 $i = 1, \ldots, r$，有
$x + I \not \subset \mathfrak p_i$。则存在 $y \in I$，
使得对所有 $i$ 都有 $x + y \not \in \mathfrak p_i$。
\end{lemma}

\begin{proof}
可设各 $\mathfrak p_i$ 之间不存在包含关系。重新排序后，可设
当 $i < s$ 时 $x \not \in \mathfrak p_i$，而当 $i \geq s$ 时
$x \in \mathfrak p_i$。若 $s = r + 1$，结论已得。否则，可找到
$y \in I$，使 $y \not \in \mathfrak p_s$。选取
$f \in \bigcap_{i < s} \mathfrak p_i$，使
$f \not \in \mathfrak p_s$。
% P01-E0285: source element-membership candidate recorded in ERRATA_CANDIDATES.jsonl.
于是 $x + fy$ 不属于 $\mathfrak p_1, \ldots, \mathfrak p_s$。
对 $s$ 归纳即得结论。
\end{proof}

\begin{lemma}[中国剩余定理]
\label{algebra-lemma-chinese-remainder}
设 $R$ 为环。
\begin{enumerate}
\item 若理想 $I_1, \ldots, I_r$ 满足：当 $a \not = b$ 时
$I_a + I_b = R$，则
$I_1 \cap \ldots \cap I_r = I_1I_2\ldots I_r$，且
$R/(I_1I_2\ldots I_r) \cong R/I_1 \times \ldots \times R/I_r$。
\item 若 $\mathfrak m_1, \ldots, \mathfrak m_r$ 是两两不同的极大理想，
则当 $a \not = b$ 时 $\mathfrak m_a + \mathfrak m_b = R$，
故可应用上述结论。
\end{enumerate}
\end{lemma}

\begin{proof}
先证明 $I_1 \cap \ldots \cap I_r = I_1 \ldots I_r$；这也会推出诱导环同态
$R/(I_1 \ldots I_r) \rightarrow R/I_1 \times \ldots \times R/I_r$
的单射性。包含关系
$I_1 \cap \ldots \cap I_r \supset I_1 \ldots I_r$ 总成立，
因为理想对乘以任意环元素封闭。为证明反向包含，断言理想
$$
I_1 \ldots \hat I_i \ldots I_r,\quad i = 1, \ldots, r
$$
生成环 $R$。对 $r$ 归纳。当 $r = 2$ 时成立。若 $r > 2$，
则 $R$ 是理想
$I_1 \ldots \hat I_i \ldots I_{r - 1}$（$i = 1, \ldots, r - 1$）
之和。因此，$I_r$ 是理想
$I_1 \ldots \hat I_i \ldots I_r$（$i = 1, \ldots, r - 1$）
之和。对这些理想按相反次序应用同样论证，可知 $I_1$ 是理想
$I_1 \ldots \hat I_i \ldots I_r$（$i = 2, \ldots, r$）之和。
由假设 $R = I_1 + I_r$，可知 $R$ 是上述所列理想之和。
因而可找到元素
$a_i \in I_1 \ldots \hat I_i \ldots I_r$，使其总和为一。
把此等式乘以 $I_1 \cap \ldots \cap I_r$ 的任一元素，即得反向包含。

还须证明典范映射
$R/(I_1 \ldots I_r) \rightarrow R/I_1 \times \ldots \times R/I_r$
为满射。为此，考察它在上面所得等式
$1 = \sum_{i=1}^r a_i$ 上的作用。一方面，环态射把 1 映到 1；
另一方面，当 $j \neq i$ 时，任意 $a_i$ 在 $R/I_j$ 中的像为零。
因此，$a_i$ 在 $R/I_i$ 中的像是单位元。故给定任意元素
$(\bar{b_1}, \ldots, \bar{b_r}) \in R/I_1 \times \ldots \times R/I_r$，
元素 $\sum_{i=1}^r a_i \cdot b_i$ 是它在 $R$ 中的一个原像。

\medskip\noindent
为证明 (2)，由极大理想两两不同的定义，对 $a \neq b$ 有
$\mathfrak{m}_a + \mathfrak{m}_b = R$，故上述结论适用。
\end{proof}

\begin{lemma}
\label{algebra-lemma-matrix-left-inverse}
设 $R$ 为环，$n \geq m$。设 $A$ 是系数在 $R$ 中的
$n \times m$ 矩阵，并设 $J \subset R$ 是由 $A$ 的
$m \times m$ 子式生成的理想。
\begin{enumerate}
\item 对任意 $f \in J$，存在 $m \times n$ 矩阵 $B$，使
$BA = f 1_{m \times m}$。
\item 若 $f \in R$，且对某个 $m \times n$ 矩阵 $B$ 有
$BA = f 1_{m \times m}$，则 $f^m \in J$。
\end{enumerate}
\end{lemma}

\begin{proof}
对满足 $I \subset \{1, \ldots, n\}$ 与 $|I| = m$ 的这类子集，
用 $E_I$ 表示投影
$$
R^{\oplus n} = \bigoplus\nolimits_{i \in \{1, \ldots, n\}} R
\longrightarrow \bigoplus\nolimits_{i \in I} R
$$
的 $m \times n$ 矩阵，并令 $A_I = E_I A$；即 $A_I$ 是由 $A$
中下标属于 $I$ 的行构成的 $m \times m$ 矩阵。
设 $B_I$ 为 $A_I$ 的伴随矩阵（代数余子式矩阵的转置），即满足
$A_I B_I = B_I A_I = \det(A_I) 1_{m \times m}$。
$A$ 的 $m \times m$ 子式正是所有满足
$I \subset \{1, \ldots, n\}$、$|I| = m$ 的这类子集所对应的行列式
$\det A_I$。若 $f \in J$，则对某些 $c_I \in R$，可写成
$f = \sum c_I \det(A_I)$。令 $B = \sum c_I B_I E_I$，即得 (1)。

\medskip\noindent
若 $f 1_{m \times m} = BA$，则由柯西--比内公式
(\ref{algebra-item-cauchy-binet})，
$f^m = \sum b_I \det(A_I)$，其中 $b_I$ 是由 $B$ 中下标属于 $I$
的列构成的 $m \times m$ 矩阵的行列式。
\end{proof}

\begin{lemma}
\label{algebra-lemma-matrix-right-inverse}
设 $R$ 为环，$n \geq m$。设 $A = (a_{ij})$ 是系数在 $R$ 中的
$n \times m$ 矩阵，按分块形式写为
$$
A =
\left(
\begin{matrix}
A_1 \\
A_2
\end{matrix}
\right)
$$
其中 $A_1$ 的大小为 $m \times m$。设 $B$ 为 $A_1$ 的伴随矩阵
（代数余子式矩阵的转置）。则
$$
AB = 
\left(
\begin{matrix}
f 1_{m \times m} \\
C
\end{matrix}
\right)
$$
其中 $f = \det(A_1)$，而 $c_{ij}$ 在相差一个符号的意义下，
是 $A$ 中由行 $1, \ldots, \hat j, \ldots, m, i$ 确定的
$m \times m$ 子式。
\end{lemma}

\begin{proof}
% P01-E0286: source missing-identity candidate recorded in ERRATA_CANDIDATES.jsonl.
由于伴随矩阵满足 $A_1B = B A_1 = f$，故 $AB$ 表达式的第一个分块正确。
注意
$$
c_{ij} = \sum\nolimits_k a_{ik}b_{kj} = \sum (-1)^{j + k}a_{ik} \det(A_1^{jk})
$$
% P01-E0287: source index-name candidate recorded in ERRATA_CANDIDATES.jsonl.
其中 $A_1^{ij}$ 表示从 $A_1$ 中删去第 $j$ 行与第 $k$ 列所得矩阵。
最后这个表达式正是引理陈述中矩阵行列式的按行展开。
\end{proof}

\begin{lemma}
\label{algebra-lemma-map-cannot-be-injective}
\begin{slogan}
若有限自由模的源的秩大于目标的秩，则该映射不可能是单射。
\end{slogan}
设 $R$ 为非零环，$n \geq 1$，$M$ 为由 $< n$ 个元素生成的 $R$-模。
则任意 $R$-模同态 $f : R^{\oplus n} \to M$ 的核都非零。
\end{lemma}

\begin{proof}
选取满射 $R^{\oplus n - 1} \to M$。可把映射 $f$ 提升为映射
$f' : R^{\oplus n} \to R^{\oplus n - 1}$（引理
\ref{algebra-lemma-lift-map}）。只须证明 $f'$ 的核非零。
映射 $f' : R^{\oplus n} \to R^{\oplus n - 1}$ 由矩阵
$A = (a_{ij})$ 给出。若某个 $a_{ij}$ 不是幂零元，比如
$a = a_{ij}$ 不是，则可把 $R$ 替换为局部化 $R_a$，从而可设
$a_{ij}$ 是单位。事实上，若局部化后找到非零核，则原来就有非零核，
因为局部化是正合的；见命题 \ref{algebra-proposition-localization-exact}。
在这种情形下，可以在 $R^{\oplus n}$ 与 $R^{\oplus n - 1}$
上都作
% P01-E0288: source terminology candidate recorded in ERRATA_CANDIDATES.jsonl.
基的变换，归约到
$$
A =
\left(
\begin{matrix}
1 & 0 & 0 & \ldots \\
0 & a_{22} & a_{23} & \ldots \\
0 & a_{32} & \ldots \\
\ldots & \ldots
\end{matrix}
\right)
$$
的情形。于是，此时对 $n$ 归纳即可。否则，每个 $a_{ij}$ 都是幂零元。
令 $I = (a_{ij}) \subset R$。注意，对某个 $m \geq 0$，
有 $I^{m + 1} = 0$。令 $m$ 为使 $I^m \not = 0$ 的最大整数。
于是 $(I^m)^{\oplus n}$ 包含于该映射的核中，结论得证。
\end{proof}

\begin{lemma}
\label{algebra-lemma-rank}
\begin{slogan}
有限自由模的秩是良定义的。
\end{slogan}
设 $R$ 为非零环，$n, m \geq 0$ 为整数。若作为 $R$-模，
$R^{\oplus n}$ 同构于 $R^{\oplus m}$，则 $n = m$。
\end{lemma}

\begin{proof}
立即由引理 \ref{algebra-lemma-map-cannot-be-injective} 得到。
\end{proof}

\section{凯莱--哈密顿}
\label{algebra-section-cayley-hamilton}

\begin{lemma}
\label{algebra-lemma-charpoly}
设 $R$ 为环，$A = (a_{ij})$ 是系数在 $R$ 中的 $n \times n$ 矩阵。
设 $P(x) \in R[x]$ 为 $A$ 的特征多项式（定义为
$\det(x\text{id}_{n \times n} - A)$）。则在
$\text{Mat}(n \times n, R)$ 中有 $P(A) = 0$。
\end{lemma}

\begin{proof}
分几步把问题归约到线性代数中熟知的凯莱--哈密顿定理：
\begin{enumerate}
\item 若 $\phi :S \rightarrow R$ 是环态射，而 $b_{ij}$ 是
$a_{ij}$ 在此映射下的原像，则只须对 $S$ 与 $(b_{ij})$ 证明该陈述，
因为 $\phi$ 是环态射。
\item 若 $\psi :R \hookrightarrow S$ 是单射环态射，则显然只须对
$S$ 以及视为 $S$ 中元素的 $a_{ij}$ 证明结论。
\item 因此，可先归约到多项式环情形
$R = \mathbf{Z}[X_{ij}]$、$a_{ij} = X_{ij}$，再进一步归约到
$R = \mathbf{Q}(X_{ij})$；在后一情形可直接应用凯莱--哈密顿定理。
\end{enumerate}
\end{proof}

\begin{lemma}
\label{algebra-lemma-charpoly-module}
设 $R$ 为环，$M$ 为有限 $R$-模，$\varphi : M \to M$ 为自同态。
则存在首一多项式 $P \in R[T]$，使得作为 $M$ 的自同态有
$P(\varphi) = 0$。
\end{lemma}

\begin{proof}
选取满的 $R$-模同态 $R^{\oplus n} \to M$，它由
$(a_1, \ldots, a_n) \mapsto \sum a_ix_i$ 给出，其中
$x_i \in M$ 是某些生成元。选取
$(a_{i1}, \ldots, a_{in}) \in R^{\oplus n}$，使
$\varphi(x_i) = \sum a_{ij} x_j$。换言之，下图交换：
$$
\xymatrix{
R^{\oplus n} \ar[d]_A \ar[r] & M \ar[d]^\varphi \\
R^{\oplus n} \ar[r] & M
}
$$
其中 $A = (a_{ij})$。由引理 \ref{algebra-lemma-charpoly}，
存在首一多项式 $P$，使 $P(A) = 0$。于是 $P(\varphi) = 0$。
\end{proof}

\begin{lemma}
\label{algebra-lemma-charpoly-module-ideal}
设 $R$ 为环，$I \subset R$ 为理想，$M$ 为有限 $R$-模。
设 $\varphi : M \to M$ 为满足 $\varphi(M) \subset IM$ 的自同态。
则存在首一多项式
% P01-E0289: source polynomial-variable candidate recorded in ERRATA_CANDIDATES.jsonl.
$P = t^n + a_1 t^{n - 1} + \ldots + a_n \in R[T]$，
使得 $a_j \in I^j$，并且作为 $M$ 的自同态有 $P(\varphi) = 0$。
\end{lemma}

\begin{proof}
选取满的 $R$-模同态 $R^{\oplus n} \to M$，它由
$(a_1, \ldots, a_n) \mapsto \sum a_ix_i$ 给出，其中
$x_i \in M$ 是某些生成元。选取
$(a_{i1}, \ldots, a_{in}) \in I^{\oplus n}$，使
$\varphi(x_i) = \sum a_{ij} x_j$。换言之，下图交换：
$$
\xymatrix{
R^{\oplus n} \ar[d]_A \ar[r] & M \ar[d]^\varphi \\
I^{\oplus n} \ar[r] & M
}
$$
其中 $A = (a_{ij})$。由引理 \ref{algebra-lemma-charpoly}，多项式
$P(t) = \det(t\text{id}_{n \times n} - A)$ 具有所有所需性质。
\end{proof}

\noindent
作为一个有趣的应用，证明下列令人惊讶的引理。

\begin{lemma}
\label{algebra-lemma-fun}
设 $R$ 为环，$M$ 为有限 $R$-模，并设
$\varphi : M \to M$ 为满的 $R$-模同态。则 $\varphi$ 是同构。
\end{lemma}

\begin{proof}[第一种证明]
记 $R' = R[x]$，并把 $M$ 视为有限 $R'$-模，其中 $x$ 通过
$\varphi$ 作用。令 $I = (x) \subset R'$。由 $\varphi$ 为满射的假设，
有 $IM = M$。因此，可对作为 $R'$-模的 $M$、理想 $I$ 以及自同态
$\text{id}_M$ 应用引理 \ref{algebra-lemma-charpoly-module-ideal}。由此得到
$(1 + a_1 + \ldots + a_n)\text{id}_M = 0$，其中 $a_j \in I$。
对某个 $b_j(x) \in R[x]$，写 $a_j = b_j(x)x$。换回 $\varphi$，
得到
$\text{id}_M = -(\sum_{j = 1, \ldots, n} b_j(\varphi)) \varphi$，
故 $\varphi$ 可逆。
\end{proof}

\begin{proof}[第二种证明]
对 $M$ 在 $R$ 上的生成元个数作归纳。若 $M$ 由一个元素生成，
则对某个理想 $I \subset R$，有 $M \cong R/I$。
此时可把 $R$ 替换为 $R/I$，从而 $M = R$。
此时 $\varphi : R \to R$ 由 $M$ 上乘以某个元素 $r \in R$ 给出。
$\varphi$ 的满射性迫使 $r$ 可逆，因为 $\varphi$ 必须命中 $1$；
所以 $\varphi$ 可逆。

\medskip\noindent
现假设引理对由 $n - 1$ 个元素生成的模已得证明，并考察由 $n$
个元素生成的模 $M$。令 $A$ 表示环 $R[t]$，并令 $t$ 通过
$\varphi$ 作用，从而把模 $M$ 视为 $A$-模。由于 $M$ 在 $R$ 上有限，
它在 $R[t]$ 上也有限；而我们要证明的 $\varphi$ 的单射性是集合论性质，
故不妨证明 $A$ 上的自同态 $t : M \to M$ 是单射。
问题已归约到自同态由底环中元素乘法给出的情形。
令 $M' \subset M$ 表示由 $M$ 的前 $n - 1$ 个生成元生成的
子-$A$-模，并考虑交换图
$$
\xymatrix{
0 \ar[r] & M' \ar[r]\ar[d]^{\varphi\mid_{M'}} & M\ar[d]^\varphi \ar[r] &
M/M' \ar[d]^{\varphi \bmod M'} \ar[r] & 0 \\
0 \ar[r] & M' \ar[r]                                  & M \ar[r]
           & M/M' \ar[r]                                  & 0,
}
$$
$\varphi$ 在 $M'$ 上的限制以及 $\varphi$ 在商 $M/M'$ 上诱导的映射
均为良定义，因为 $\varphi$ 是乘以底环中的一个元素，而 $M'$ 与
$M/M'$ 各自都是 $A$-模。由 $n = 1$ 的情形，映射
$M/M' \to M/M'$ 是同构。追图表明 $\varphi|_{M'}$ 为满射，
故由归纳假设 $\varphi|_{M'}$ 是同构。再由蛇形引理，
中间的竖直映射必为同构。
\end{proof}
\section{环的谱}
\label{algebra-section-spectrum-ring}

\noindent
我们约定把作为拓扑空间的环谱纳入代数章，而把仿射概形纳入概形章；
这一划分只是任意的编排选择。

\begin{definition}
\label{algebra-definition-spectrum-ring}
设 $R$ 为环。
\begin{enumerate}
\item $R$ 的{\it 谱}是 $R$ 的全体素理想组成的集合，通常记作
$\Spec(R)$。
\item 给定子集 $T \subset R$，令 $V(T) \subset \Spec(R)$ 为包含
$T$ 的素理想之集，即 $V(T) = \{ \mathfrak p \in
\Spec(R) \mid \forall f\in T, f\in \mathfrak p\}$。
\item 给定元素 $f \in R$，令 $D(f) \subset \Spec(R)$ 为不包含
$f$ 的素理想之集。
\end{enumerate}
\end{definition}

\begin{lemma}
\label{algebra-lemma-Zariski-topology}
设 $R$ 为环。
\begin{enumerate}
\item 环 $R$ 的谱为空，当且仅当 $R$ 为零环。
\item 每个非零环都有极大理想。
\item 每个非零环都有极小素理想。
\item 给定理想 $I \subset R$ 以及素理想
$I \subset \mathfrak p$，存在素理想
$I \subset \mathfrak q \subset \mathfrak p$，使得
$\mathfrak q$ 在 $I$ 上极小。
\item 若 $T \subset R$，且 $(T)$ 是 $T$ 在 $R$ 中生成的理想，
则 $V((T)) = V(T)$。
\item 若 $I$ 为理想而 $\sqrt{I}$ 为其根理想，参见基本概念
(\ref{algebra-item-radical-ideal})，则 $V(I) = V(\sqrt{I})$。
\item 给定 $R$ 的理想 $I$，有 $\sqrt{I} =
\bigcap_{I \subset \mathfrak p} \mathfrak p$。
\item 若 $I$ 为理想，则 $V(I) = \emptyset$ 当且仅当 $I$ 为单位理想。
\item 若 $I$, $J$ 是 $R$ 的理想，则 $V(I) \cup V(J) =
V(I \cap J)$。
\item 若 $(I_a)_{a\in A}$ 是 $R$ 的一族理想，则
$\bigcap_{a\in A} V(I_a) = V(\bigcup_{a\in A} I_a)$。
\item 若 $f \in R$，则 $D(f) \amalg V(f) = \Spec(R)$。
\item 若 $f \in R$，则 $D(f) = \emptyset$ 当且仅当 $f$ 幂零。
\item 若对某个单位 $u \in R$ 有 $f = u f'$，则 $D(f) = D(f')$。
\item 若 $I \subset R$ 为理想，而 $\mathfrak p$ 是 $R$ 的素理想且
$\mathfrak p \not\in V(I)$，则存在 $f \in R$，使得
$\mathfrak p \in D(f)$ 且 $D(f) \cap V(I) = \emptyset$。
\item 若 $f, g \in R$，则 $D(fg) = D(f) \cap D(g)$。
\item 若对 $i \in I$ 有 $f_i \in R$，则
$\bigcup_{i\in I} D(f_i)$ 是 $V(\{f_i \}_{i\in I})$ 在
$\Spec(R)$ 中的补集。
\item 若 $f \in R$ 且 $D(f) = \Spec(R)$，则 $f$ 为单位。
\end{enumerate}
\end{lemma}

\begin{proof}
依次证明以上各项。
\begin{enumerate}
\item 这是 (2) 或 (3) 的直接推论。
\item 令 $\mathfrak{A}$ 为 $R$ 的全体真理想组成的集合。它按包含关系
排序，而且非空，因为 $(0) \in \mathfrak{A}$ 是真理想。设 $A$ 是
$\mathfrak A$ 的全序子集。于是 $\bigcup_{I \in A} I$ 确为理想。
由于对所有 $I \in A$ 都有 1 $\notin I$，这个并不包含 1，因而仍为
真理想。因此 $\bigcup_{I \in A} I$ 属于 $\mathfrak{A}$，并且是
$A$ 的一个上界。由佐恩引理，$\mathfrak{A}$ 有极大元，而它正是所求的
极大理想。
\item 因为 $R$ 非零，它含有一个极大理想，而极大理想是素理想。因此
$R$ 的全体素理想组成的集合 $\mathfrak{A}$ 非空。按反向包含关系排列
$\mathfrak{A}$。设 $A$ 是 $\mathfrak{A}$ 的全序子集。显然
$J = \bigcap_{I \in A} I$ 是理想；但它为素理想并不那么显然。设
$xy \in J$，则对所有 $I \in A$ 都有 $xy \in I$。现令
$B = \{I \in A | y \in I\}$，并令 $K = \bigcap_{I \in B} I$。
由于 $A$ 全序，或者 $K = J$（此时 $y \in J$，结论成立），或者
$K \supset J$，并且对每个真包含于 $K$ 的 $I \in A$，都有
$y \notin I$。但这些 $I$ 都是素理想，所以对所有这样的
$I, x \in I$。因此 $x \in J$。两种情形下 $J$ 都是素理想。
故由佐恩引理得到一个极大元；在这里，它就是极小素理想。
\item 论证与 (3) 完全相同，只需仅考虑包含于 $\mathfrak{p}$ 且包含
$I$ 的素理想。
\item $(T)$ 是包含 $T$ 的最小理想。因此若 $T \subset I$，其中右端
为某个理想，则也有 $(T) \subset I$。所以若 $I \in V(T)$，
则也有 $I \in V((T))$。反向包含显然。
\item 由 $I \subset \sqrt{I}, V(\sqrt{I}) \subset V(I)$。现设
$\mathfrak{p} \in V(I)$，并取 $x \in \sqrt{I}$。则对某个 $n$ 有
$x^n \in I$，从而 $x^n \in \mathfrak{p}$。由于 $\mathfrak{p}$
为素理想，一个直接的归纳论证给出 $x \in \mathfrak{p}$。故
$\sqrt{I} \subset \mathfrak{p}$，于是
$\mathfrak{p} \in V(\sqrt{I})$。
\item 设 $f \in R \setminus \sqrt{I}$。则对所有 $n$ 都有
$f^n \notin I$。因此 $S = \{1, f, f^2, \ldots\}$ 是不含 $0$ 的
乘法子集。取包含 $S^{-1}I$ 的素理想
$\bar{\mathfrak{p}} \subset S^{-1}R$。令 $\mathfrak{p}$ 为
$\bar{\mathfrak{p}}$ 在 $R$ 中的原像；它是包含 $I$ 且不与 $S$
相交的素理想。这说明
$\bigcap_{I \subset \mathfrak p} \mathfrak p \subset \sqrt{I}$。
反之，若 $a \in \sqrt{I}$，则对某个 $n$ 有 $a^n \in I$。所以当
$I \subset \mathfrak{p}$ 时，$a^n \in \mathfrak{p}$。由于
$\mathfrak{p}$ 为素理想，故 $a \in \mathfrak{p}$，从而得证等式。
\item $I$ 不是单位理想，当且仅当 $I$ 包含于某个极大理想（将 (2)
用于环 $R/I$ 即可看出），而该极大理想因而是素理想。
% P01-E0290: malformed ideal-containment disjunction recorded in ERRATA_CANDIDATES.jsonl.
\item 若 $\mathfrak{p} \in V(I) \cup V(J)$，则 $I \subset \mathfrak{p}$
或 $J \subset \mathfrak{p}$，从而 $I \cap J \subset \mathfrak{p}$。
反之，若 $I \cap J \subset \mathfrak{p}$，则
$IJ \subset \mathfrak{p}$；由于 $\mathfrak{p}$ 为素理想，故
$I$ 或 $J$ 中有一个包含于 $\mathfrak{p}$。
\item $\mathfrak{p} \in \bigcap_{a \in A} V(I_a) \Leftrightarrow
I_a \subset \mathfrak{p}, \forall a \in A \Leftrightarrow
\mathfrak{p} \in V(\bigcup_{a\in A} I_a)$。
\item 若 $\mathfrak{p}$ 为素理想且 $f \in R$，则严格地说，或者
$f \in \mathfrak{p}$，或者 $f \notin \mathfrak{p}$；这正是该不交并
所表达的内容。
\item 若 $a \in R$ 幂零，则对某个 $n$ 有 $a^n = 0$。于是对每个
素理想都有 $a^n \in \mathfrak{p}$。归纳可得
$a \in \mathfrak{p}$，故 $D(a) = \emptyset$。反之，由 (7)，若
$a \in R$ 不幂零，则存在不包含它的素理想。
\item $f \in \mathfrak{p} \Leftrightarrow uf \in \mathfrak{p}$，
因为 $u$ 可逆。
\item 若 $\mathfrak{p} \notin V(I)$，则
$\exists f \in I \setminus \mathfrak{p}$。于是
$f \notin \mathfrak{p}$，所以 $\mathfrak{p} \in D(f)$。此外，若
$\mathfrak{q} \in D(f)$，则 $f \notin \mathfrak{q}$，从而 $I$
不包含于 $\mathfrak{q}$。故 $D(f) \cap V(I) = \emptyset$。
\item 若 $fg \in \mathfrak{p}$，则 $f \in \mathfrak{p}$ 或
$g \in \mathfrak{p}$。因此，若 $f \notin \mathfrak{p}$ 且
$g \notin \mathfrak{p}$，则 $fg \notin \mathfrak{p}$。又因为
$\mathfrak{p}$ 是理想，若 $fg \notin \mathfrak{p}$，则
$f \notin \mathfrak{p}$ 且 $g \notin \mathfrak{p}$。
\item $\mathfrak{p} \in \bigcup_{i \in I} D(f_i) \Leftrightarrow \exists i \in
I, f_i \notin \mathfrak{p} \Leftrightarrow \mathfrak{p} \in \Spec(R)
\setminus V(\{f_i\}_{i \in I})$。
\item 若 $D(f) = \Spec(R)$，则 $V(f) = \emptyset$，因而
$fR = R$，所以 $f$ 是单位。
\end{enumerate}
\end{proof}

\noindent
该引理说明，定义 \ref{algebra-definition-spectrum-ring} 中的子集 $V(T)$
构成 $\Spec(R)$ 上某个拓扑的闭集。它还说明集合 $D(f)$ 是开集，
并构成该拓扑的一组基。

\begin{definition}
\label{algebra-definition-Zariski-topology}
设 $R$ 为环。以各集合 $V(T)$ 为闭集的 $\Spec(R)$ 上的拓扑称为
{\it 扎里斯基}拓扑。开子集 $D(f)$ 称为 $\Spec(R)$ 的
{\it 标准开集}。
\end{definition}

\noindent
由上下文应当能够看清我们是仅把 $\Spec(R)$ 视为集合，还是把它视为
拓扑空间。

\begin{lemma}
\label{algebra-lemma-spec-functorial}
\begin{slogan}
谱的函子性
\end{slogan}
设 $\varphi : R \to R'$ 为环同态。诱导映射
$$
\Spec(\varphi) : \Spec(R') \longrightarrow \Spec(R),
\quad
\mathfrak p' \longmapsto \varphi^{-1}(\mathfrak p')
$$
关于扎里斯基拓扑连续。事实上，对任意元素 $f \in R$，有
$\Spec(\varphi)^{-1}(D(f)) = D(\varphi(f))$。
\end{lemma}

\begin{proof}
由基本概念 (\ref{algebra-item-inverse-image-prime})，
$\mathfrak p := \varphi^{-1}(\mathfrak p')$ 确为 $R$ 的素理想。
引理的最后一个断言直接来自定义，并蕴含第一个断言。
\end{proof}

\noindent
若 $\varphi' : R' \to R''$ 是第二个环同态，则复合
$$
\Spec(R'')
\longrightarrow
\Spec(R')
\longrightarrow
\Spec(R)
$$
等于 $\Spec(\varphi' \circ \varphi)$。换言之，$\Spec$ 是从环范畴
到拓扑空间范畴的反变函子。

\begin{lemma}
\label{algebra-lemma-spec-localization}
设 $R$ 为环，$S \subset R$ 为乘法子集。由 $\Spec$ 的函子性，
映射 $R \to S^{-1}R$ 诱导同胚
$$
\Spec(S^{-1}R)
\longrightarrow
\{\mathfrak p \in \Spec(R) \mid S \cap \mathfrak p = \emptyset \}
$$
其中右端赋予从 $\Spec(R)$ 的扎里斯基拓扑诱导的拓扑。逆映射由
$\mathfrak p \mapsto S^{-1}\mathfrak p = \mathfrak p(S^{-1}R)$ 给出。
\end{lemma}

\begin{proof}
记引理中箭头的右端为 $D$。取素理想
% P01-E0291: missing copula in source construction recorded in ERRATA_CANDIDATES.jsonl.
$\mathfrak p' \subset S^{-1}R$，并令 $\mathfrak p$ 为
$\mathfrak p'$ 在 $R$ 中的原像。由于 $\mathfrak p'$ 不包含 $1$，
可知 $\mathfrak p$ 不包含 $S$ 的任何元素。因此 $\mathfrak p \in D$，
从而像包含于 $D$。现设 $\mathfrak p \in D$。依假设，像
$\overline{S}$ 不包含 $0$。由基本概念
(\ref{algebra-item-localization-zero})，$\overline{S}^{-1}(R/\mathfrak p)$
不是零环。由基本概念 (\ref{algebra-item-localize-ideal})，可见
$S^{-1}R / S^{-1}\mathfrak p = \overline{S}^{-1}(R/\mathfrak p)$
为整环，故 $S^{-1}\mathfrak p$ 为素理想。这个环等式还表明
$S^{-1}\mathfrak p$ 在 $R$ 中的原像等于 $\mathfrak p$，因为由基本概念
(\ref{algebra-item-localize-nonzerodivisors})，映射
$R/\mathfrak p \to \overline{S}^{-1}(R/\mathfrak p)$ 为单射。
这就证明了映射 $\Spec(S^{-1}R) \to \Spec(R)$ 双射到 $D$，并且其
逆映射如上所述。由引理 \ref{algebra-lemma-spec-functorial}，该映射连续。
最后，设 $D(g) \subset \Spec(S^{-1}R)$ 为标准开集。将其写成
$g = h/s$，其中 $h\in R$ 且 $s\in S$。由于 $g$ 与
$h/1$ 相差一个单位，在 $\Spec(S^{-1}R)$ 中有
$D(g) = D(h/1)$。于是由引理 \ref{algebra-lemma-spec-functorial} 及上述
双射性，$D(g) = D(h/1)$ 的像为 $D \cap D(h)$。这也证明了该映射
为开映射。
\end{proof}

\begin{lemma}
\label{algebra-lemma-standard-open}
设 $R$ 为环，$f \in R$。由 $\Spec$ 的函子性，映射
$R \to R_f$ 诱导同胚
$$
\Spec(R_f) \longrightarrow D(f) \subset \Spec(R).
$$
其逆映射由 $\mathfrak p \mapsto \mathfrak p \cdot R_f$ 给出。
\end{lemma}

\begin{proof}
这是引理 \ref{algebra-lemma-spec-localization} 的特殊情形。
\end{proof}

\noindent
谱的每个“仿射开集”并不都是标准开集。参见例
\ref{algebra-example-affine-open-not-standard}。

\begin{lemma}
\label{algebra-lemma-spec-closed}
设 $R$ 为环，$I \subset R$ 为理想。由 $\Spec$ 的函子性，映射
$R \to R/I$ 诱导同胚
$$
\Spec(R/I) \longrightarrow V(I) \subset \Spec(R).
$$
其逆映射由 $\mathfrak p \mapsto \mathfrak p / I$ 给出。
\end{lemma}

\begin{proof}
像显然包含于 $V(I)$。另一方面，若 $\mathfrak p \in V(I)$，则
$\mathfrak p \supset I$，从而可以考虑理想
$\mathfrak p /I \subset R/I$。由基本概念
(\ref{algebra-item-isomorphism-theorem})，可知
$(R/I)/(\mathfrak p/I) = R/\mathfrak p$ 为整环，所以
$\mathfrak p/I$ 是素理想。由此立即可见 $D(f + I)$ 的像为
$D(f) \cap V(I)$，因而该映射是同胚。
\end{proof}



\begin{lemma}
\label{algebra-lemma-quasi-compact}
\begin{slogan}
环的谱是拟紧的
\end{slogan}
设 $R$ 为环。空间 $\Spec(R)$ 是拟紧的。
\end{lemma}

\begin{proof}
% P01-E0292: refinement/subcover terminology candidate recorded in ERRATA_CANDIDATES.jsonl.
只需证明 $\Spec(R)$ 的任何标准开集覆盖都可由一个有限覆盖加细。设
$\Spec(R) = \cup D(f_i)$，其中 $\{f_i\}_{i\in I}$ 是 $R$ 中的一族
元素。这意味着 $\cap V(f_i) = \emptyset$。由引理
\ref{algebra-lemma-Zariski-topology}，这意味着
$V(\{f_i \}) = \emptyset$。再由同一引理，这意味着各 $f_i$ 生成
$R$ 的单位理想。也就是说，可以把 $1$ 写成一个{\it 有限}和：
$1 = \sum_{i \in J} r_i f_i$，其中 $J \subset I$ 有限。于是
$\Spec(R) = \cup_{i \in J} D(f_i)$。
\end{proof}

\begin{lemma}
\label{algebra-lemma-topology-spec}
设 $R$ 为环。$X = \Spec(R)$ 上的拓扑具有下列性质：
\begin{enumerate}
\item $X$ 拟紧；
\item $X$ 有一组由拟紧开集构成的拓扑基；
\item 任意两个拟紧开集的交仍拟紧。
\end{enumerate}
\end{lemma}

\begin{proof}
环的谱拟紧，参见引理 \ref{algebra-lemma-quasi-compact}。它有一组由标准开集
$D(f) = \Spec(R_f)$ 构成的拓扑基（引理
\ref{algebra-lemma-standard-open}），而由第一条可知这些开集拟紧。
两个标准开集的交拟紧，因为 $D(f) \cap D(g) = D(fg)$。
给定任意两个拟紧开集 $U, V \subset X$，可写成
$U = D(f_1) \cup \ldots \cup D(f_n)$ 以及
$V = D(g_1) \cup \ldots \cup D(g_m)$。于是
% P01-E0293: unindexed union candidate recorded in ERRATA_CANDIDATES.jsonl.
$U \cap V = \bigcup D(f_ig_j)$，它是拟紧的。
\end{proof}
\section{局部环}
\label{algebra-section-local-rings}

\noindent
局部环是代数几何中最基本而常用的工具。

\begin{definition}
\label{algebra-definition-local-ring}
恰有一个极大理想的环称为{\it 局部环}。若 $R$ 为局部环，通常把它的
极大理想记作 $\mathfrak m_R$，并把域 $R/\mathfrak m_R$ 称为局部环
$R$ 的{\it 剩余域}。我们常说“设 $(R, \mathfrak m)$ 为局部环”或
“设 $(R, \mathfrak m, \kappa)$ 为局部环”，以表示 $R$ 是局部环、
$\mathfrak m$ 是其唯一极大理想，而
$\kappa = R/\mathfrak m$ 是其剩余域。局部环之间的{\it 局部同态}
是环映射 $\varphi : R \to S$，其中 $R$ 与 $S$ 均为局部环，并且
$\varphi(\mathfrak m_R) \subset \mathfrak m_S$。若已知 $R$ 与 $S$
都是局部环，则“{\it 局部环映射 $\varphi : R \to S$}”意指
$\varphi$ 是局部环之间的局部同态。
\end{definition}

\noindent
域是局部环。域之间的任意环映射都是局部环之间的局部同态。

\medskip\noindent
环 $R$ 在素理想 $\mathfrak p$ 处的局部化 $R_\mathfrak p$ 是局部环，
其极大理想为 $\mathfrak p R_\mathfrak p$。事实上，由引理
\ref{algebra-lemma-spec-localization}，$R_\mathfrak p$ 的每个素理想都包含于
素理想 $\mathfrak p R_\mathfrak p$（所以后者是极大理想，也是
$R_\mathfrak p$ 唯一的极大理想）。$R_\mathfrak p$ 的剩余域记作
$\kappa(\mathfrak p)$；我们称它为{\it $\mathfrak p$ 的剩余域}。
由命题 \ref{algebra-proposition-localize-quotient}，可将 $\kappa(\mathfrak p)$
等同于整环 $R/\mathfrak p$ 的分式域。经由复合
$$
\Spec(\kappa(\mathfrak p)) \to \Spec(R_\mathfrak p) \to \Spec(R)
$$
源中的唯一一点映到靶中的点 $\mathfrak p$。

\medskip\noindent
设 $\varphi : R \to S$ 为环映射。令 $\mathfrak q \subset S$ 为
素理想，并考虑 $R$ 的素理想
$\mathfrak p = \varphi^{-1}(\mathfrak q)$。由于
$\varphi(\mathfrak p) \subset \mathfrak q$，诱导环映射
$$
R_\mathfrak p \to S_\mathfrak q,\quad r/g \mapsto \varphi(r)/\varphi(g)
$$
是局部环映射，并且得到剩余域间的诱导映射
$\kappa(\mathfrak p) \to \kappa(\mathfrak q)$。


\begin{example}
\label{algebra-example-not-local}
若 $R$ 为局部环，而 $\mathfrak p \subset R$ 是非极大的素理想，
则 $R \to R_\mathfrak p$ 不是局部同态。
\end{example}

\begin{lemma}
\label{algebra-lemma-characterize-local-ring}
设 $R$ 为环。下列条件等价：
\begin{enumerate}
\item $R$ 是局部环；
\item $\Spec(R)$ 恰有一个闭点；
\item $R$ 有一个极大理想 $\mathfrak m$，并且
$R \setminus \mathfrak m$ 的每个元素都是单位；
\item $R$ 不是零环，并且对每个 $x \in R$，$x$ 与 $1 - x$ 中至少
一个可逆，也可能二者均可逆。
\end{enumerate}
\end{lemma}

\begin{proof}
设 $R$ 为环，$\mathfrak m$ 为极大理想。若
$x \in R \setminus \mathfrak m$ 且 $x$ 不是单位，则存在包含
$x$ 的极大理想 $\mathfrak m'$。因此 $R$ 至少有两个极大理想。
反过来，若 $\mathfrak m'$ 是另一个极大理想，则可取
$x \in \mathfrak m'$ 且 $x \not \in \mathfrak m$。显然 $x$
不是单位。这证明了 (1) 与 (3) 等价。(1) 与 (2) 的等价是同义反复。
若 $R$ 局部，则 (4) 成立，因为 $x$ 或者属于 $\mathfrak m$，或者
不属于它。若 (4) 成立，而 $\mathfrak m$, $\mathfrak m'$ 是不同的
极大理想，则由中国剩余定理（引理
\ref{algebra-lemma-chinese-remainder}），可选取 $x \in R$ 使得
$x \bmod \mathfrak m' = 0$ 且 $x \bmod \mathfrak m = 1$。
此元素 $x$ 不可逆，而 $1 - x$ 也不可逆，矛盾。因此 (4) 与 (1)
等价。
\end{proof}

\begin{lemma}
\label{algebra-lemma-characterize-local-ring-map}
设 $\varphi : R \to S$ 为环映射，并假设 $R$ 与 $S$ 是局部环。
下列条件等价：
\begin{enumerate}
\item $\varphi$ 是局部环映射；
\item $\varphi(\mathfrak m_R) \subset \mathfrak m_S$；
\item $\varphi^{-1}(\mathfrak m_S) = \mathfrak m_R$；
\item 对任意 $x \in R$，若 $\varphi(x)$ 在 $S$ 中可逆，则 $x$
在 $R$ 中可逆。
\end{enumerate}
\end{lemma}

\begin{proof}
由定义，条件 (1) 与 (2) 等价。若 (3) 成立，则 (2) 成立。
反过来，若 (2) 成立，则 $\varphi^{-1}(\mathfrak m_S)$ 是包含极大理想
$\mathfrak m_R$ 的素理想，故
$\varphi^{-1}(\mathfrak m_S) = \mathfrak m_R$。最后，由引理
\ref{algebra-lemma-characterize-local-ring}，(4) 是 (2) 的逆否命题。
\end{proof}

\begin{remark}
\label{algebra-remark-fundamental-diagram}
与环映射 $\varphi : R \to S$ 及素理想 $\mathfrak p \subset R$
相关联的一个基本交换图如下：
$$
\xymatrix{
\kappa(\mathfrak p) \otimes_R S =
S_{\mathfrak p}/{\mathfrak p}S_{\mathfrak p}
&
S_{\mathfrak p} \ar[l]
&
S \ar[r] \ar[l]
&
S/\mathfrak pS \ar[r]
&
(R \setminus \mathfrak p)^{-1}S/\mathfrak pS
\\
\kappa(\mathfrak p) =
R_{\mathfrak p}/{\mathfrak p}R_{\mathfrak p} \ar[u]
&
R_{\mathfrak p} \ar[u] \ar[l]
&
R \ar[u] \ar[r] \ar[l]
&
R/\mathfrak p \ar[u] \ar[r]
&
\kappa(\mathfrak p) \ar[u]
}
$$
在此图中，最左列与最右列相同。在谱上，各水平映射诱导到其像的
同胚，而各方块诱导拓扑空间的纤维积方块（参见引理
\ref{algebra-lemma-spec-localization} 与 \ref{algebra-lemma-spec-closed}）。这说明
$\mathfrak p$ 属于谱上该映射的像，当且仅当
$S \otimes_R \kappa(\mathfrak p)$ 不是零环。若确实存在位于
$\mathfrak p$ 上方的素理想 $\mathfrak q \subset S$，即满足
$\mathfrak p = \varphi^{-1}(\mathfrak q)$，则可把该图扩充为下图：
$$
\xymatrix{
\kappa(\mathfrak q) = S_{\mathfrak q}/{\mathfrak q}S_{\mathfrak q}
&
S_{\mathfrak q} \ar[l]
&
S \ar[r] \ar[l]
&
S/\mathfrak q \ar[r]
&
\kappa(\mathfrak q)
\\
\kappa(\mathfrak p) \otimes_R S =
S_{\mathfrak p}/{\mathfrak p}S_{\mathfrak p} \ar[u]
&
S_{\mathfrak p} \ar[u] \ar[l]
&
S \ar[u] \ar[r] \ar[l]
&
S/\mathfrak pS \ar[u] \ar[r]
&
(R \setminus \mathfrak p)^{-1}S/\mathfrak pS \ar[u]
\\
\kappa(\mathfrak p) =
R_{\mathfrak p}/{\mathfrak p}R_{\mathfrak p} \ar[u]
&
R_{\mathfrak p} \ar[u] \ar[l]
&
R \ar[u] \ar[r] \ar[l]
&
R/\mathfrak p \ar[u] \ar[r]
&
\kappa(\mathfrak p) \ar[u]
}
$$
在此图中，最左列与最右列仍然相同；在谱上，各水平映射仍诱导到其像
的同胚。
\end{remark}

\begin{lemma}
\label{algebra-lemma-in-image}
设 $\varphi : R \to S$ 为环映射，$\mathfrak p$ 为 $R$ 的素理想。
下列条件等价：
\begin{enumerate}
\item $\mathfrak p$ 属于映射 $\Spec(S) \to \Spec(R)$ 的像；
\item $S \otimes_R \kappa(\mathfrak p) \not = 0$；
\item $S_{\mathfrak p}/\mathfrak p S_{\mathfrak p} \not = 0$；
\item $(S/\mathfrak pS)_{\mathfrak p} \not = 0$；
\item $\mathfrak p = \varphi^{-1}(\mathfrak pS)$。
\end{enumerate}
\end{lemma}

\begin{proof}
注 \ref{algebra-remark-fundamental-diagram} 已经给出了前两个条件的等价性；
其余条件只是它的不同表述。
\end{proof}

\begin{remark}
\label{algebra-remark-local-ring-fibre}
设 $R \to S$ 为环映射。令 $\mathfrak q$ 为 $S$ 的素理想，并位于
$R$ 的素理想 $\mathfrak p$ 上方。由注
\ref{algebra-remark-fundamental-diagram}，素理想 $\mathfrak q$ 对应于纤维环
$F = S \otimes_R \kappa(\mathfrak p)$ 的唯一素理想
$\overline{\mathfrak q}$。于是
$$
F_{\overline{\mathfrak q}}
\cong S_\mathfrak q \otimes_{R_\mathfrak p} \kappa(\mathfrak p)
\cong S_\mathfrak q/\mathfrak p S_\mathfrak q
$$
事实上，有一个显然的环映射
$F \to S_\mathfrak q \otimes_{R_\mathfrak p} \kappa(\mathfrak p)$，
容易看出它同构于 $F \to F_{\overline{\mathfrak q}}$。第二个同构来自
$\kappa(\mathfrak p)$ 是 $R_\mathfrak p$ 对
$\mathfrak pR_\mathfrak p$ 的商这一事实。
\end{remark}
\section{环的雅可布森根}
\label{algebra-section-radical}

\noindent
回顾环 $R$ 的{\it 雅可布森根} $\text{rad}(R)$ 是 $R$ 的全体极大理想
之交。若 $R$ 局部，则 $\text{rad}(R)$ 就是 $R$ 的极大理想。

\begin{lemma}
\label{algebra-lemma-contained-in-radical}
设 $R$ 为环，其雅可布森根为 $\text{rad}(R)$。令 $I \subset R$
为理想。下列条件等价：
\begin{enumerate}
\item $I \subset \text{rad}(R)$；
\item $1 + I$ 的每个元素都是 $R$ 中的单位。
\end{enumerate}
在这种情形下，$R$ 中每个映为 $R/I$ 中单位的元素本身都是单位。
\end{lemma}

\begin{proof}
若 $f \in \text{rad}(R)$，则对 $R$ 的每个极大理想
$\mathfrak m$ 都有 $f \in \mathfrak m$。所以对 $R$ 的每个极大理想
$\mathfrak m$ 都有 $1 + f \not \in \mathfrak m$。因此
$\Spec(R)$ 的闭子集 $V(1 + f)$ 为空。这说明 $1 + f$ 是单位，
参见引理 \ref{algebra-lemma-Zariski-topology}。

\medskip\noindent
反过来，假设对所有 $f \in I$，$1 + f$ 都是单位。若
$\mathfrak m$ 是极大理想且 $I \not \subset \mathfrak m$，则
$I + \mathfrak m = R$。故对某个 $g \in \mathfrak m$ 与
$f \in I$ 有 $1 = f + g$。于是 $g = 1 + (-f)$ 不是单位，矛盾。

\medskip\noindent
对最后一个断言，设 $f \in R$ 映为 $R/I$ 中的单位。可以找到
$g \in R$，它映为 $f \bmod I$ 的乘法逆元。于是
$fg = 1 \bmod I$。由 (2)，$fg$ 是 $R$ 中的单位，从而 $f$ 是单位。
\end{proof}

\begin{lemma}
\label{algebra-lemma-surjective-on-spec-units}
设 $\varphi : R \to S$ 为环映射，并且诱导映射
$\Spec(S) \to \Spec(R)$ 为满射。则元素 $x \in R$ 是单位，当且仅当
$\varphi(x) \in S$ 是单位。
\end{lemma}

\begin{proof}
若 $x$ 是单位，则 $\varphi(x)$ 也是单位。反过来，若
$\varphi(x)$ 是单位，则对所有 $\mathfrak q \in \Spec(S)$ 都有
$\varphi(x) \not \in \mathfrak q$。因此对所有
$\mathfrak q \in \Spec(S)$ 都有
$x \not \in \varphi^{-1}(\mathfrak q) = \Spec(\varphi)(\mathfrak q)$。
由于 $\Spec(\varphi)$ 为满射，由引理
\ref{algebra-lemma-Zariski-topology} 的第 (17) 项可知 $x$ 是单位。
\end{proof}




\section{中山引理}
\label{algebra-section-nakayama}

\noindent
我们引述 \cite{MatCA}：“这个简单而重要的引理归功于
T.~Nakayama、G.~Azumaya 和 W.~Krull。其优先权并不明确；虽然它通常
称为中山引理，已故的 Nakayama 教授却不喜欢这个名称。”

\begin{lemma}[中山引理]
\label{algebra-lemma-NAK}
\begin{reference}
\cite[1.M Lemma (NAK) page 11]{MatCA}
\end{reference}
\begin{history}
我们引述 \cite{MatCA}：“这个简单而重要的引理归功于
T.~Nakayama、G.~Azumaya 和 W.~Krull。其优先权并不明确；虽然它通常
称为中山引理，已故的 Nakayama 教授却不喜欢这个名称。”
\end{history}
设 $R$ 为环，其雅可布森根为 $\text{rad}(R)$。令 $M$ 为 $R$-模，
$I \subset R$ 为理想。
\begin{enumerate}
\item
\label{algebra-item-nakayama}
若 $IM = M$ 且 $M$ 有限，则存在 $f \in 1 + I$ 使得 $fM = 0$。
\item 若 $IM = M$、$M$ 有限且 $I \subset \text{rad}(R)$，则 $M = 0$。
\item 若 $N, N' \subset M$、$M = N + IN'$ 且 $N'$ 有限，则存在
$f \in 1 + I$，使得 $fM \subset N$ 且 $M_f = N_f$。
\item 若 $N, N' \subset M$、$M = N + IN'$、$N'$ 有限且
$I \subset \text{rad}(R)$，则 $M = N$。
\item 若 $N \to M$ 为模映射、$N/IN \to M/IM$ 为满射且 $M$ 有限，
则存在 $f \in 1 + I$，使得 $N_f \to M_f$ 为满射。
\item 若 $N \to M$ 为模映射、$N/IN \to M/IM$ 为满射、$M$ 有限且
$I \subset \text{rad}(R)$，则 $N \to M$ 为满射。
\item 若 $x_1, \ldots, x_n \in M$ 生成 $M/IM$ 且 $M$ 有限，则存在
$f \in 1 + I$，使得 $x_1, \ldots, x_n$ 在 $R_f$ 上生成 $M_f$。
\item 若 $x_1, \ldots, x_n \in M$ 生成 $M/IM$、$M$ 有限且
$I \subset \text{rad}(R)$，则 $M$ 由 $x_1, \ldots, x_n$ 生成。
\item 若 $IM = M$ 且 $I$ 幂零，则 $M = 0$。
\item 若 $N, N' \subset M$、$M = N + IN'$ 且 $I$ 幂零，则 $M = N$。
\item 若 $N \to M$ 为模映射、$I$ 幂零且
$N/IN \to M/IM$ 为满射，则 $N \to M$ 为满射。
\item 若 $\{x_\alpha\}_{\alpha \in A}$ 是 $M$ 中生成 $M/IM$ 的一族
元素且 $I$ 幂零，则 $M$ 由各 $x_\alpha$ 生成。
\end{enumerate}
\end{lemma}

\begin{proof}
证明 (\ref{algebra-item-nakayama})。选取 $M$ 在 $R$ 上的一组生成元
$y_1, \ldots, y_m$。对每个 $i$，可写成
$y_i = \sum z_{ij} y_j$，其中 $z_{ij} \in I$（因为 $M = IM$）。
换言之，$\sum_j (\delta_{ij} - z_{ij})y_j = 0$。令 $f$ 为
$m \times m$ 矩阵 $A = (\delta_{ij} - z_{ij})$ 的行列式。注意
$f \in 1 + I$（因为矩阵 $A$ 的每个分量模 $I$ 同余于
$m \times m$ 单位矩阵的相应分量）。由引理
\ref{algebra-lemma-matrix-left-inverse} (1)，存在 $m \times m$ 矩阵 $B$，
使得 $BA = f 1_{m \times m}$。展开可见，对所有 $h$ 与 $j$，有
$\sum_{i} b_{hi} a_{ij} = f \delta_{hj}$；因而对每个 $h$，
$\sum_{i, j} b_{hi} a_{ij} y_j
= \sum_{j} f \delta_{hj} y_j = f y_h$。换言之，对每个 $h$ 都有
$0 = f y_h$（因为每个 $i$ 都满足 $\sum_j a_{ij} y_j = 0$）。
这说明 $f$ 零化 $M$。

\medskip\noindent
% P01-E0294: missing article in source unit statement recorded in ERRATA_CANDIDATES.jsonl.
由引理 \ref{algebra-lemma-contained-in-radical}，$1 + \text{rad}(R)$ 的元素是
$R$ 的可逆元。因此 (\ref{algebra-item-nakayama}) 蕴含 (2)。把 (1) 用于
$M/N$ 得到 (3)，这里由于 $N'$ 有限，该商有限。把 (2) 用于
$M/N$ 得到 (4)，这里同样由于 $N'$ 有限，该商有限。把 (3) 用于
$M$ 及其子模 $\Im(N \to M)$ 和 $M$，得到 (5)。把 (4) 用于 $M$
及其子模 $\Im(N \to M)$ 和 $M$，得到 (6)。把 (5) 用于映射
$R^{\oplus n} \to M$，$(a_1, \ldots, a_n) \mapsto
a_1x_1 + \ldots + a_nx_n$，得到 (7)。把 (6) 用于映射
$R^{\oplus n} \to M$，$(a_1, \ldots, a_n) \mapsto
a_1x_1 + \ldots + a_nx_n$，得到 (8)。

\medskip\noindent
第 (9) 项成立，因为若 $M = IM$，则对所有 $n \geq 0$ 都有
$M = I^nM$，而 $I$ 幂零意味着对某个 $n \gg 0$ 有 $I^n = 0$。
第 (10)、(11)、(12) 项由 (9) 及上述论证推出。
\end{proof}

\begin{lemma}
\label{algebra-lemma-NAK-localization}
设 $R$ 为环，$S \subset R$ 为乘法子集，$I \subset R$ 为理想，
$M$ 为有限 $R$-模。若 $x_1, \ldots, x_r \in M$ 作为
$S^{-1}(R/I)$-模生成 $S^{-1}(M/IM)$，则存在 $f \in S + I$，使得
$x_1, \ldots, x_r$ 作为 $R_f$-模生成 $M_f$。\footnote{特殊情形：
(I) $I = 0$。引理说明，若 $x_1, \ldots, x_r$ 生成 $S^{-1}M$，
则对某个 $f \in S$，$x_1, \ldots, x_r$ 生成 $M_f$。
(II) $I = \mathfrak p$ 是素理想且 $S = R \setminus \mathfrak p$。
引理说明，若 $x_1, \ldots, x_r$ 生成
$M \otimes_R \kappa(\mathfrak p)$，则对某个满足 $f \in R$、
$f \not \in \mathfrak p$ 的元素，$x_1, \ldots, x_r$ 生成 $M_f$。}
\end{lemma}

\begin{proof}
先看特殊情形 $I = 0$。设 $y_1, \ldots, y_s$ 是 $M$ 在 $R$ 上的
生成元。由于 $S^{-1}M$ 由 $x_1, \ldots, x_r$ 生成，对每个 $i$，
可在 $S^{-1}M$ 中写成
$y_i = \sum (a_{ij}/s_{ij})x_j$，其中 $a_{ij} \in R$ 且
$s_{ij} \in S$。乘以各 $s_{ij}$ 的乘积 $s \in S$，可见在
$S^{-1}M$ 中 $sy_i = \sum a'_{ij}x_j$，其中
$a'_{ij} \in R$。这又意味着存在 $t_i \in S$，使得在 $M$ 中有
$t_isy_i = \sum t_ia'_{ij}x_j$。因此，若 $t \in S$ 是各 $t_i$ 的
乘积，则在 $M_{st}$ 中，$y_i$ 属于 $x_1, \ldots, x_r$ 生成的
$R_{st}$-子模。故 $x_1, \ldots, x_r$ 生成 $M_{st}$。
% P01-E0295: subject--verb agreement defect recorded in ERRATA_CANDIDATES.jsonl.

\medskip\noindent
一般情形。由特殊情形，可找到 $s \in S$，使得
$x_1, \ldots, x_r$ 在 $(R/I)_s$ 上生成 $(M/IM)_s$。由引理
\ref{algebra-lemma-NAK}，可找到 $g \in 1 + I_s \subset R_s$，使得
$x_1, \ldots, x_r$ 在 $(R_s)_g$ 上生成 $(M_s)_g$。写成
$g = 1 + i/s'$。则 $f = ss' + is$ 满足要求；细节从略。
\end{proof}

\begin{lemma}
\label{algebra-lemma-when-surjective-local}
设 $A \to B$ 为局部环之间的局部同态。假设：
\begin{enumerate}
\item $B$ 作为 $A$-模有限；
\item $\mathfrak m_B$ 是有限生成理想；
\item $A \to B$ 诱导剩余域之间的同构；
\item $\mathfrak m_A/\mathfrak m_A^2 \to
\mathfrak m_B/\mathfrak m_B^2$ 为满射。
\end{enumerate}
则 $A \to B$ 为满射。
\end{lemma}

\begin{proof}
为证明 $A \to B$ 满射，把它视为 $A$-模映射并应用引理
\ref{algebra-lemma-NAK} (6)。由此只需证明
$A/\mathfrak m_A \to B/\mathfrak m_AB$ 为满射。由于
$A/\mathfrak m_A = B/\mathfrak m_B$，只需证明
$\mathfrak m_AB \to \mathfrak m_B$ 为满射。把
$\mathfrak m_AB \to \mathfrak m_B$ 视为 $B$-模映射并应用引理
\ref{algebra-lemma-NAK} (6)。于是只需看出
$\mathfrak m_AB/\mathfrak m_A\mathfrak m_B \to
\mathfrak m_B/\mathfrak m_B^2$ 为满射。这由假设 (4) 推出。
\end{proof}
\section{谱的开闭子集}
\label{algebra-section-open-and-closed}

\noindent
事实表明，谱中既开又闭的子集与环的幂等元相对应。

\begin{lemma}
\label{algebra-lemma-idempotent-spec}
设 $R$ 为环，$e \in R$ 为幂等元。此时
$$
\Spec(R) = D(e) \amalg D(1-e).
$$
\end{lemma}

\begin{proof}
注意，整环中的幂等元 $e$ 只能是 $1$ 或 $0$。因此
\begin{eqnarray*}
D(e)
& = &
\{ \mathfrak p \in \Spec(R)
\mid
e \not\in \mathfrak p \} \\
& = &
\{ \mathfrak p \in \Spec(R)
\mid
e \not = 0\text{ 于 }\kappa(\mathfrak p) \} \\
& = &
\{ \mathfrak p \in \Spec(R)
\mid
e = 1\text{ 于 }\kappa(\mathfrak p) \}
\end{eqnarray*}
类似地，
\begin{eqnarray*}
D(1-e)
& = &
\{ \mathfrak p \in \Spec(R)
\mid
1 - e \not\in \mathfrak p \} \\
& = &
\{ \mathfrak p \in \Spec(R)
\mid
e \not = 1\text{ 于 }\kappa(\mathfrak p) \} \\
& = &
\{ \mathfrak p \in \Spec(R)
\mid
e = 0\text{ 于 }\kappa(\mathfrak p) \}
\end{eqnarray*}
由于 $e$ 在任意剩余域中的像只能是 $1$ 或 $0$，可知 $D(e)$ 与
$D(1-e)$ 覆盖整个 $\Spec(R)$。
\end{proof}

\begin{lemma}
\label{algebra-lemma-spec-product}
设 $R_1$ 与 $R_2$ 为环，令 $R = R_1 \times R_2$。映射
$R \to R_1$、$(x, y) \mapsto x$ 以及 $R \to R_2$、
$(x, y) \mapsto y$ 分别诱导连续映射
$\Spec(R_1) \to \Spec(R)$ 与 $\Spec(R_2) \to \Spec(R)$。
诱导映射
$$
\Spec(R_1) \amalg \Spec(R_2)
\longrightarrow
\Spec(R)
$$
是同胚。换言之，$R = R_1\times R_2$ 的谱是 $R_1$ 的谱与
$R_2$ 的谱的不交并。
\end{lemma}

\begin{proof}
写成 $1 = e_1 + e_2$，其中 $e_1 = (1, 0)$ 且
$e_2 = (0, 1)$。注意，$e_1$ 与 $e_2 = 1 - e_1$ 都是幂等元。
留给读者证明 $R_1 = R_{e_1}$ 是 $R$ 在 $e_1$ 处的局部化。
$e_2$ 的情形类似。因此，引理的断言由引理
\ref{algebra-lemma-idempotent-spec} 与引理 \ref{algebra-lemma-standard-open} 合并推出。
\end{proof}

\noindent
引入函数的粘合引理后，我们还会再次证明下一个引理。参见节
\ref{algebra-section-tilde-module-sheaf}。

\begin{lemma}
\label{algebra-lemma-disjoint-decomposition}
设 $R$ 为环。对每个既开又闭的 $U \subset \Spec(R)$，存在唯一的
幂等元 $e \in R$，使得 $U = D(e)$。这给出既开又闭的子集
$U \subset \Spec(R)$ 与幂等元 $e \in R$ 之间的一一对应。
\end{lemma}

\begin{proof}
设 $U \subset \Spec(R)$ 既开又闭。因为 $U$ 闭，由引理
\ref{algebra-lemma-quasi-compact}，它拟紧；其补集同样拟紧。把它写成
标准开集的有限并 $U = \bigcup_{i = 1}^n D(f_i)$。类似地，写成
$\Spec(R) \setminus U = \bigcup_{j = 1}^m D(g_j)$。由于
$\emptyset = D(f_i) \cap D(g_j) = D(f_i g_j)$，由引理
\ref{algebra-lemma-Zariski-topology}，$f_i g_j$ 幂零。令
$I = (f_1, \ldots, f_n) \subset R$，并令
$J = (g_1, \ldots, g_m) \subset R$。注意，$V(J)$ 等于 $U$，
$V(I)$ 等于 $U$ 的补集，所以
$\Spec(R) = V(I) \amalg V(J)$。由上述幂零性可知，对某个充分大的
整数 $N$，有 $(IJ)^N = (0)$。由于
$\bigcup D(f_i) \cup \bigcup D(g_j) = \Spec(R)$，由引理
\ref{algebra-lemma-Zariski-topology} 可知 $I + J = R$。把这个等式取
$2N$ 次幂，得到 $I^N + J^N = R$。写成 $1 = x + y$，其中
$x \in I^N$ 且 $y \in J^N$。因为 $I^N J^N = (0)$，有
$0 = xy = x(1 - x)$。故 $x = x^2$ 是幂等元，并且属于
$I^N \subset I$。幂等元 $y = 1 - x$ 属于 $J^N \subset J$。
这说明，对每个 $\mathfrak p \in V(J)$，幂等元 $x$ 在剩余域
$\kappa(\mathfrak p)$ 中映为 $1$；而对每个
$\mathfrak p \in V(I)$，$x$ 在 $\kappa(\mathfrak p)$ 中映为 $0$。

\medskip\noindent
为证明唯一性，设 $e_1, e_2$ 是 $R$ 中不同的幂等元。必须证明存在
素理想 $\mathfrak p$，使得 $e_1 \in \mathfrak p$ 而
$e_2 \not \in \mathfrak p$，或反过来。写
$e_i' = 1 - e_i$。若 $e_1 \not = e_2$，则
$0 \not = e_1 - e_2  = e_1(e_2 + e_2') - (e_1 + e_1')e_2
= e_1 e_2' - e_1' e_2$。所以幂等元 $e_1 e_2' \not = 0$ 或
$e_1' e_2 \not = 0$ 至少有一个成立。
% P01-E0296: nonzero qualification on idempotent recorded in ERRATA_CANDIDATES.jsonl.
幂等元不幂零，因而由引理 \ref{algebra-lemma-Zariski-topology} 可找到素理想
$\mathfrak p$，使得 $e_1e_2' \not \in \mathfrak p$ 或
$e_1'e_2 \not \in \mathfrak p$。容易看出这给出了所需的素理想。
\end{proof}

\begin{lemma}
\label{algebra-lemma-characterize-spec-connected}
设 $R$ 为非零环。则 $\Spec(R)$ 连通，当且仅当 $R$ 没有非平凡
幂等元。
\end{lemma}

\begin{proof}
这由引理 \ref{algebra-lemma-disjoint-decomposition} 及连通拓扑空间的定义
立即得到。
\end{proof}

\begin{lemma}
\label{algebra-lemma-ideal-is-squared-union-connected}
设 $I \subset R$ 为环 $R$ 的有限生成理想，并且 $I = I^2$。则：
\begin{enumerate}
\item 存在幂等元 $e \in R$，使得 $I = (e)$；
\item 对幂等元 $e' = 1 - e \in R$，有 $R/I \cong R_{e'}$；
\item $V(I)$ 在 $\Spec(R)$ 中既开又闭。
\end{enumerate}
\end{lemma}

\begin{proof}
由中山引理 \ref{algebra-lemma-NAK}，存在元素 $f = 1 + i$、$i \in I$，
使得 $fI = 0$。于是 $f^2 = f + fi = f$，故它为幂等元。
考虑幂等元 $e = 1 - f = -i \in I$。对 $j \in I$，有
$ej = j - fj = j$，所以 $I = (e)$。这证明了 (1)。

\medskip\noindent
(2) 与 (3) 由 (1) 推出。事实上，
$V(I) = V(e) = \Spec(R) \setminus D(e)$；由引理
\ref{algebra-lemma-idempotent-spec} 或引理
\ref{algebra-lemma-disjoint-decomposition}，它既开又闭。这证明了 (3)。
为证明 (2)，注意映射 $R \to R_{e'}$ 为满射，因为在 $R_{e'}$ 中
$x/(e')^n = x/e' = xe'/(e')^2 = xe'/e' = x/1$。映射
$R \to R_{e'}$ 的核是 $R$ 中被 $e'$ 的某个正整数次幂零化的元素之集。
由于 $e'$ 幂等，这就是被 $e'$ 零化的元素所成的理想；又因为
$e + e' = 1$ 是一对正交幂等元，该理想就是 $I = (e)$。这证明了 (2)。
\end{proof}
\section{谱的连通分支}
\label{algebra-section-connected-components}

\noindent
谱的连通分支并不像初看时可能以为的那样容易理解。这是因为我们习惯于
局部连通空间的拓扑，而环的谱一般并不局部连通。

\begin{lemma}
\label{algebra-lemma-closed-union-connected-components}
设 $R$ 为环，$T \subset \Spec(R)$ 为谱的子集。下列条件等价：
\begin{enumerate}
\item $T$ 闭，并且是 $\Spec(R)$ 的若干连通分支之并；
\item $T$ 是 $\Spec(R)$ 中若干既开又闭子集之交；
\item $T = V(I)$，其中 $I \subset R$ 是由幂等元生成的理想。
\end{enumerate}
此外，若 (3) 中的理想存在，则它是唯一的。
\end{lemma}

\begin{proof}
由引理 \ref{algebra-lemma-topology-spec} 及拓扑，引理
\ref{topology-lemma-closed-union-connected-components}，可知 (1) 与
(2) 等价。假设 (2)，并写成 $T = \bigcap U_\alpha$，其中
$U_\alpha \subset \Spec(R)$ 既开又闭。由引理
\ref{algebra-lemma-disjoint-decomposition}，对某个幂等元
$e_\alpha \in R$ 有 $U_\alpha = D(e_\alpha)$。令
$I = (1 - e_\alpha)$，便有 $T = V(I)$，即 (3) 成立。最后假设 (3)。
写成 $T = V(I)$，并令 $I = (e_\alpha)$，其中各 $e_\alpha$ 为幂等元。
于是显然
$T = \bigcap V(e_\alpha) = \bigcap D(1 - e_\alpha)$。

\medskip\noindent
设 $I$ 是由幂等元生成的理想。令 $e \in R$ 为满足
$V(I) \subset V(e)$ 的幂等元。由引理 \ref{algebra-lemma-Zariski-topology}，
对某个 $n \geq 1$ 有 $e^n \in I$。由于 $e$ 幂等，这意味着
$e \in I$。因此，$I$ 恰由满足 $T \subset V(e)$ 的那些幂等元 $e$
生成。换言之，闭子集 $T$ 完全确定理想 $I$，这证明了唯一性。
\end{proof}

\begin{lemma}
\label{algebra-lemma-connected-component}
设 $R$ 为环。$\Spec(R)$ 的一个连通分支形如 $V(I)$，其中 $I$
是由幂等元生成的理想，并且 $R$ 的每个幂等元在 $R/I$ 中都映为
$0$ 或 $1$。
\end{lemma}

\begin{proof}
设 $\mathfrak p$ 为 $R$ 的素理想。由引理
\ref{algebra-lemma-topology-spec}，可知拓扑，引理
\ref{topology-lemma-connected-component-intersection} 的假设对拓扑空间
% P01-E0297: source verb-form defect recorded in ERRATA_CANDIDATES.jsonl.
$\Spec(R)$ 成立。因此，$\mathfrak p$ 在 $\Spec(R)$ 中的连通分支是
$\Spec(R)$ 中包含 $\mathfrak p$ 的所有既开又闭子集之交。所以它等于
$V(I)$，其中 $I$ 由幂等元 $e \in R$ 生成，而 $e$ 在
$\kappa(\mathfrak p)$ 中映为 $0$，参见引理
\ref{algebra-lemma-disjoint-decomposition}。
显然，不属于这一族的任意幂等元 $e$ 在 $R/I$ 中映为 $1$。
\end{proof}
\section{粘合性质}
\label{algebra-section-more-glueing}

\noindent
本节汇集若干如下形式的标准结果：若某个性质在一个标准开覆盖的每个
成员上成立，则它整体成立。事实上，往往只需在局部环层面检验，正如下
一个引理所示。

\begin{lemma}
\label{algebra-lemma-characterize-zero-local}
设 $R$ 为环。
\begin{enumerate}
\item 对 $R$-模 $M$ 的元素 $x$，下列条件等价：
\begin{enumerate}
\item $x = 0$；
\item 对所有 $\mathfrak p \in \Spec(R)$，$x$ 在 $M_\mathfrak p$ 中
的像为零；
\item 对 $R$ 的所有极大理想 $\mathfrak m$，$x$ 在
$M_{\mathfrak m}$ 中的像为零。
\end{enumerate}
换言之，映射 $M \to \prod_{\mathfrak m} M_{\mathfrak m}$ 为单射。
\item 给定 $R$-模 $M$，下列条件等价：
\begin{enumerate}
\item $M$ 为零；
\item 对所有 $\mathfrak p \in \Spec(R)$，$M_{\mathfrak p}$ 为零；
\item 对 $R$ 的所有极大理想 $\mathfrak m$，$M_{\mathfrak m}$ 为零。
\end{enumerate}
\item 给定 $R$-模的复形 $M_1 \to M_2 \to M_3$，下列条件等价：
\begin{enumerate}
\item $M_1 \to M_2 \to M_3$ 正合；
\item 对 $R$ 的每个素理想 $\mathfrak p$，局部化
$M_{1, \mathfrak p} \to M_{2, \mathfrak p} \to M_{3, \mathfrak p}$
正合；
\item 对 $R$ 的每个极大理想 $\mathfrak m$，局部化
$M_{1, \mathfrak m} \to M_{2, \mathfrak m} \to M_{3, \mathfrak m}$
正合。
\end{enumerate}
\item 给定 $R$-模映射 $f : M \to M'$，下列条件等价：
\begin{enumerate}
\item $f$ 为单射；
\item 对 $R$ 的所有素理想 $\mathfrak p$，
$f_{\mathfrak p} : M_\mathfrak p \to M'_\mathfrak p$ 为单射；
\item 对 $R$ 的所有极大理想 $\mathfrak m$，
$f_{\mathfrak m} : M_\mathfrak m \to M'_\mathfrak m$ 为单射。
\end{enumerate}
\item 给定 $R$-模映射 $f : M \to M'$，下列条件等价：
\begin{enumerate}
\item $f$ 为满射；
\item 对 $R$ 的所有素理想 $\mathfrak p$，
$f_{\mathfrak p} : M_\mathfrak p \to M'_\mathfrak p$ 为满射；
\item 对 $R$ 的所有极大理想 $\mathfrak m$，
$f_{\mathfrak m} : M_\mathfrak m \to M'_\mathfrak m$ 为满射。
\end{enumerate}
\item 给定 $R$-模映射 $f : M \to M'$，下列条件等价：
\begin{enumerate}
\item $f$ 为双射；
\item 对 $R$ 的所有素理想 $\mathfrak p$，
$f_{\mathfrak p} : M_\mathfrak p \to M'_\mathfrak p$ 为双射；
\item 对 $R$ 的所有极大理想 $\mathfrak m$，
$f_{\mathfrak m} : M_\mathfrak m \to M'_\mathfrak m$ 为双射。
\end{enumerate}
\end{enumerate}
\end{lemma}

\begin{proof}
设 $x \in M$ 如 (1) 所述。令 $I = \{f \in R \mid fx = 0\}$。
容易看出 $I$ 是理想（它是 $x$ 的零化子）。条件 (1)(c) 意味着，
对每个极大理想 $\mathfrak m$，存在 $f \in R \setminus \mathfrak m$
使得 $fx =0$。换言之，$V(I)$ 不包含闭点。由引理
\ref{algebra-lemma-Zariski-topology} 可知 $I$ 是单位理想。因此 $x$ 为零，
即 (1)(a) 成立。这证明了 (1)。

\medskip\noindent
把 (1) 同时应用于 $M$ 的所有元素，即得 (2)。

\medskip\noindent
证明 (3)。令 $H$ 为该序列的同调，即
$H = \Ker(M_2 \to M_3)/\Im(M_1 \to M_2)$。由命题
\ref{algebra-proposition-localization-exact}，$H_\mathfrak p$ 是序列
$M_{1, \mathfrak p} \to M_{2, \mathfrak p} \to M_{3, \mathfrak p}$
的同调。因此 (3) 是 (2) 的推论。

\medskip\noindent
(4) 与 (5) 是 (3) 的特殊情形。把 (4) 与 (5) 合并，形式地得到 (6)。
\end{proof}

\begin{lemma}
\label{algebra-lemma-cover}
\begin{slogan}
模与代数的扎里斯基局部性质
\end{slogan}
设 $R$ 为环，$M$ 为 $R$-模，$S$ 为 $R$-代数。假设
$f_1, \ldots, f_n$ 是 $R$ 中的有限元素列，并满足
$\bigcup D(f_i) = \Spec(R)$，换言之，$(f_1, \ldots, f_n) = R$。
\begin{enumerate}
\item 若每个 $M_{f_i} = 0$，则 $M = 0$。
\item 若每个 $M_{f_i}$ 都是有限 $R_{f_i}$-模，则 $M$ 是有限
$R$-模。
\item 若每个 $M_{f_i}$ 都是有限表示 $R_{f_i}$-模，则 $M$ 是有限表示
$R$-模。
\item 设 $M \to N$ 为 $R$-模映射。若对每个 $i$，
$M_{f_i} \to N_{f_i}$ 都是同构，则 $M \to N$ 是同构。
\item 设 $0 \to M'' \to M \to M' \to 0$ 为 $R$-模的复形。
若对每个 $i$，$0 \to M''_{f_i} \to M_{f_i} \to M'_{f_i} \to 0$
都正合，则 $0 \to M'' \to M \to M' \to 0$ 正合。
\item 若每个 $R_{f_i}$ 都是诺特环，则 $R$ 是诺特环。
\item 若每个 $S_{f_i}$ 都是有限型 $R_{f_i}$-代数，则 $S$ 是有限型
$R$-代数。
\item 若每个 $S_{f_i}$ 都在 $R_{f_i}$ 上有限表示，则 $S$ 是有限表示
$R$-代数。
\end{enumerate}
\end{lemma}

\begin{proof}
依次证明各项。
\begin{enumerate}
\item 由命题 \ref{algebra-proposition-localize-twice}，该假设蕴含对所有
$\mathfrak p \in \Spec(R)$ 都有 $M_\mathfrak p = 0$，故由引理
\ref{algebra-lemma-characterize-zero-local} 得出结论。
\item 对每个 $i$，取 $M_{f_i}$ 的有限生成集 $X_i$。不失一般性，
可假设 $X_i$ 的元素都在局部化映射 $M \rightarrow M_{f_i}$ 的像中，
于是取 $M$ 中由 $X_i$ 各元素的原像组成的有限集合 $Y_i$。令 $Y$
为这些集合之并，它仍是有限集。考虑显然的 $R$-线性映射
$R^Y \rightarrow M$，把基元素 $e_y$ 映到 $y$。依假设，在 $R$
的任意素理想 $\mathfrak p$ 处局部化后，此映射都是满射；故由引理
\ref{algebra-lemma-characterize-zero-local}，它本身满射，从而 $M$ 有限生成。
\item 由 (2)，有短正合列
$$
0 \rightarrow K \rightarrow R^n \rightarrow M \rightarrow 0
$$
由于局部化是正合函子，且 $M_{f_i}$ 有限表示，由引理
\ref{algebra-lemma-extension}，对所有 $1 \leq i \leq n$，$K_{f_i}$ 都有限生成。
% P01-E0298: overloaded cover/generator index recorded in ERRATA_CANDIDATES.jsonl.
由 (2)，这说明 $K$ 是有限 $R$-模，因而 $M$ 有限表示。
\item 由命题 \ref{algebra-proposition-localize-twice}，该假设蕴含在所有素理想
处的诱导局部化态射都是同构，所以由引理
\ref{algebra-lemma-characterize-zero-local} 得出结论。
\item 由命题 \ref{algebra-proposition-localize-twice}，该假设蕴含在所有素理想
处的诱导局部化序列都短正合，所以由引理
\ref{algebra-lemma-characterize-zero-local} 得出结论。
\item 我们证明 $R$ 的每个理想都有有限生成集。为此，令
$I \subset R$ 为任意理想。由命题
\ref{algebra-proposition-localization-exact}，每个
$I_{f_i} \subset R_{f_i}$ 都是理想。依假设，它们都有限生成，故由
(2) 得出结论。
\item 对每个 $i$，取 $S_{f_i}$ 的有限生成集 $X_i$。不失一般性，
可假设 $X_i$ 的元素都在局部化映射 $S \rightarrow S_{f_i}$ 的像中，
于是取 $S$ 中由 $X_i$ 各元素的原像组成的有限集合 $Y_i$。令 $Y$
为这些集合之并，它仍是有限集。考虑由 $Y$ 诱导的代数同态
$R[X_y]_{y \in Y} \rightarrow S$。因为它是代数同态，其像 $T$ 是
$R$-模 $S$ 的 $R$-子模，所以可以考虑商模 $S/T$。依假设，在各
$f_i$ 处局部化后它为零，故由 (1) 它为零，因而 $S$ 是有限型
$R$-代数。
\item 由上一项，存在满的 $R$-代数同态
$R[X_1, \ldots, X_n] \rightarrow S$。令 $K$ 为该映射的核。
这是 $R[X_1, \ldots, X_n]$ 中的理想，并且在每个 $f_i$ 处局部化后
有限生成。由于各 $f_i$ 在 $R$ 中生成单位理想，它们在
$R[X_1, \ldots, X_n]$ 中也生成单位理想，所以应用 (2) 即完成证明。
\end{enumerate}
\end{proof}

\begin{lemma}
\label{algebra-lemma-cover-upstairs}
设 $R \to S$ 为环映射。假设 $g_1, \ldots, g_n$ 是 $S$ 中的有限
元素列，并满足 $\bigcup D(g_i) = \Spec(S)$，换言之，
$(g_1, \ldots, g_n) = S$。
\begin{enumerate}
\item 若每个 $S_{g_i}$ 都是有限型 $R$-代数，则 $S$ 是有限型
$R$-代数。
\item 若每个 $S_{g_i}$ 都在 $R$ 上有限表示，则 $S$ 在 $R$ 上
有限表示。
\end{enumerate}
\end{lemma}

\begin{proof}
选取 $h_1, \ldots, h_n \in S$，使得 $\sum h_i g_i = 1$。

\medskip\noindent
证明 (1)。对每个 $i$，选取元素的有限列
$x_{i, j} \in S_{g_i}$、$j = 1, \ldots, m_i$，使其作为
$R$-代数生成 $S_{g_i}$。对某个 $y_{i, j} \in S$ 及某个
$n_{i, j} \ge 0$，写成
$x_{i, j} = y_{i, j}/g_i^{n_{i, j}}$。考虑由
$g_1, \ldots, g_n$、$h_1, \ldots, h_n$ 以及
$y_{i, j}$（$i = 1, \ldots, n$、$j = 1, \ldots, m_i$）生成的
$R$-子代数 $S' \subset S$。由于局部化正合（命题
\ref{algebra-proposition-localization-exact}），可知
$S'_{g_i} \to S_{g_i}$ 为单射。另一方面，由 $y_{i, j}$ 的选取，
它是满射。由于 $h_1, \ldots, h_n \in S'$，各元素
$g_1, \ldots, g_n$ 在 $S'$ 中生成单位理想。因此，把
$S' \to S$ 视为 $S'$-模映射，由引理 \ref{algebra-lemma-cover}，它是同构。

\medskip\noindent
证明 (2)。已经知道 $S$ 是有限型的。对某个理想 $J$，写成
$S = R[x_1, \ldots, x_m]/J$。对每个 $i$，选取 $g_i$ 的提升
$g'_i \in R[x_1, \ldots, x_m]$，并选取 $h_i$ 的提升
$h'_i \in R[x_1, \ldots, x_m]$。于是
$$
S_{g_i} = R[x_1, \ldots, x_m, y_i]/(J_i + (1 - y_ig'_i))
$$
其中 $J_i$ 是 $R[x_1, \ldots, x_m, y_i]$ 中由 $J$ 生成的理想。
略去一个小细节。由引理 \ref{algebra-lemma-finite-presentation-independent}，
可以选取有限元素列 $f_{i, j} \in J$、$j = 1, \ldots, m_i$，
使得 $f_{i, j}$ 在 $J_i$ 中的像与 $1 - y_ig'_i$ 一起生成理想
$J_i + (1 - y_ig'_i)$。令
$$
S' = R[x_1, \ldots, x_m]/\left(\sum h'_ig'_i - 1, f_{i, j}; 
i = 1, \ldots, n, j = 1, \ldots, m_i\right)
$$
存在满的 $R$-代数映射 $S' \to S$。各元素
$g'_1, \ldots, g'_n$ 在 $S'$ 中的类生成单位理想，并且由构造，映射
$S'_{g'_i} \to S_{g_i}$ 都是单射。因此可以像 (1) 中那样得出结论。
\end{proof}
\section{函数的粘合}
\label{algebra-section-tilde-module-sheaf}

\noindent
本节将证明：给定由标准开集组成的开覆盖
$$
\Spec(R) = \bigcup\nolimits_{i = 1}^n D(f_i)
$$
以及对每个 $i$ 给定元素 $h_i \in R_{f_i}$，若作为
$R_{f_i f_j}$ 的元素有 $h_i = h_j$，则存在唯一的 $h \in R$，
使得 $h$ 在 $R_{f_i}$ 中的像为 $h_i$。这个结果可以作两种解释：
\begin{enumerate}
\item 规则 $D(f) \mapsto R_f$ 在标准开集上给出一个环层，参见层，
节 \ref{sheaves-section-bases}。
\item 若把 $R_f$ 的元素看成 $D(f)$ 上的“代数”函数或“正则”函数，
那么它们可以像拓扑流形上的连续函数、相应地像可微流形上的可微函数
那样粘合。
\end{enumerate}

\begin{lemma}
\label{algebra-lemma-cover-module}
设 $R$ 为环，$f_1, \ldots, f_n$ 是 $R$ 中生成单位理想的元素，
$M$ 为 $R$-模。序列
$$
0 \to
M \xrightarrow{\alpha}
\bigoplus\nolimits_{i = 1}^n M_{f_i} \xrightarrow{\beta}
\bigoplus\nolimits_{i, j = 1}^n M_{f_i f_j}
$$
正合，其中 $\alpha(m) = (m/1, \ldots, m/1)$，并且
$\beta(m_1/f_1^{e_1}, \ldots, m_n/f_n^{e_n})
= (m_i/f_i^{e_i} - m_j/f_j^{e_j})_{(i, j)}$。
\end{lemma}

\begin{proof}
只需证明该序列在任意极大理想 $\mathfrak m$ 处的局部化正合，参见
引理 \ref{algebra-lemma-characterize-zero-local}。由于
$f_1, \ldots, f_n$ 生成单位理想，存在 $i$ 使得
$f_i \not \in \mathfrak m$。重新编号后可假设 $i = 1$。注意
$(M_{f_i})_\mathfrak m = (M_\mathfrak m)_{f_i}$ 以及
$(M_{f_if_j})_\mathfrak m = (M_\mathfrak m)_{f_if_j}$，参见命题
\ref{algebra-proposition-localize-twice-module}。特别地，因为 $f_1$ 是单位，
有 $(M_{f_1})_\mathfrak m = M_\mathfrak m$ 以及
$(M_{f_1 f_i})_\mathfrak m = (M_\mathfrak m)_{f_i}$。注意，序列中的
映射是由引理 \ref{algebra-lemma-universal-property-localization-module} 及
$M$ 上的恒等映射给出的典范映射。综上，把 $R$ 替换为
$R_\mathfrak m$、把 $M$ 替换为 $M_\mathfrak m$、把 $f_i$ 替换为其
在 $R_\mathfrak m$ 中的像，并把 $f_1$ 替换为
$1 \in R_\mathfrak m$，便归约到 $f_1 = 1$ 的情形。

\medskip\noindent
假设 $f_1 = 1$。此时 $\alpha$ 的单射性显然。设
$m = (m_i) \in \bigoplus_{i = 1}^n M_{f_i}$ 属于 $\beta$ 的核。
则 $m_1 \in M_{f_1} = M$。此外，$\beta(m) = 0$ 蕴含 $m_1$ 与
$m_i$ 在 $M_{f_1f_i} = M_{f_i}$ 中映为同一元素。因此
$\alpha(m_1) = m$，证明完成。
\end{proof}

\begin{lemma}
\label{algebra-lemma-standard-covering}
设 $R$ 为环，并设 $f_1, f_2, \ldots f_n\in R$ 在 $R$ 中生成
单位理想。则下列序列正合：
% P01-E0299: missing separator in source element list recorded in ERRATA_CANDIDATES.jsonl.
$$
0 \longrightarrow
R \longrightarrow
\bigoplus\nolimits_i R_{f_i} \longrightarrow
\bigoplus\nolimits_{i, j}R_{f_if_j}
$$
其中映射 $\alpha : R \longrightarrow \bigoplus_i R_{f_i}$ 与
$\beta : \bigoplus_i R_{f_i} \longrightarrow \bigoplus_{i, j} R_{f_if_j}$
定义为
$$
\alpha(x) = \left(\frac{x}{1}, \ldots, \frac{x}{1}\right)
\text{ 且 }
\beta\left(\frac{x_1}{f_1^{r_1}}, \ldots, \frac{x_n}{f_n^{r_n}}\right)
=
\left(\frac{x_i}{f_i^{r_i}}-\frac{x_j}{f_j^{r_j}}~\text{于}~R_{f_if_j}\right).
$$
\end{lemma}

\begin{proof}
这是引理 \ref{algebra-lemma-cover-module} 的特殊情形。
\end{proof}

\noindent
下面的结果我们在上文已经见过，但为了方便，现将它明确陈述出来。

\begin{lemma}
\label{algebra-lemma-disjoint-implies-product}
设 $R$ 为环。若 $\Spec(R) = U \amalg V$，且 $U$ 与 $V$ 都开，
则 $R \cong R_1 \times R_2$，并且经由引理
\ref{algebra-lemma-spec-product} 中的映射，有 $U \cong \Spec(R_1)$ 与
$V \cong \Spec(R_2)$。此外，$R_1$ 与 $R_2$ 都既是环 $R$ 的局部化，
又是它的商。
\end{lemma}

\begin{proof}
由引理 \ref{algebra-lemma-disjoint-decomposition}，对某个幂等元 $e$，有
$U = D(e)$ 及 $V = D(1-e)$。由引理
\ref{algebra-lemma-standard-covering} 可知
$R \cong R_e \times R_{1 - e}$（因为显然 $R_{e(1-e)} = 0$，所以
粘合条件是平凡的；当然，在这种情形下直接证明该乘积分解也很容易）。
引理得证。
\end{proof}

\begin{lemma}
\label{algebra-lemma-when-injective-covering}
设 $R$ 为环。设 $f_1, \ldots, f_n \in R$。
设 $M$ 为 $R$-模。则 $M \to \bigoplus M_{f_i}$ 是单射，当且仅当
$$
M \longrightarrow \bigoplus\nolimits_{i = 1, \ldots, n} M, \quad
m \longmapsto (f_1m, \ldots, f_nm)
$$
是单射。
\end{lemma}

\begin{proof}
映射 $M \to \bigoplus M_{f_i}$ 是单射，当且仅当对于所有
$m \in M$ 及 $e_1, \ldots, e_n \geq 1$，若
$f_i^{e_i}m = 0$，$i = 1, \ldots, n$，则 $m = 0$。
这显然推出所显示的映射是单射。反之，假设所显示的映射是单射，且
$m \in M$、$e_1, \ldots, e_n \geq 1$ 满足
$f_i^{e_i}m = 0$，$i = 1, \ldots, n$。若所有 $i$ 都有
$e_i = 1$，则由所显示映射的单射性立即得到 $m = 0$。
接下来对 $e = \sum e_i$ 归纳，证明任意这样的数据都有同一结论。
基始情形是 $e = n$，已在上面处理。若某个 $e_i > 1$，令
$m' = f_im$。由归纳假设可知 $m' = 0$。因此
$f_i m = 0$，即可以取 $e_i = 1$，这使 $e$ 减小，结论得证。
\end{proof}

\noindent
下面的引理最好在平坦下降的更一般语境中陈述和证明。不过它与上文十分契合，
所以在此陈述也很自然。

\begin{lemma}
\label{algebra-lemma-glue-modules}
设 $R$ 为环，$f_1, \ldots, f_n \in R$。假设给定以下数据：
\begin{enumerate}
\item 对每个 $i$，给定一个 $R_{f_i}$-模 $M_i$。
\item 对每一对 $i, j$，给定一个 $R_{f_if_j}$-模同构
$\psi_{ij} : (M_i)_{f_j} \to (M_j)_{f_i}$。
\end{enumerate}
并且这些数据满足“上圈条件”，即所有图表
$$
\xymatrix{
(M_i)_{f_jf_k}
\ar[rd]_{\psi_{ij}}
\ar[rr]^{\psi_{ik}}
& &
(M_k)_{f_if_j} \\
&
(M_j)_{f_if_k} \ar[ru]_{\psi_{jk}}
}
$$
均交换（对所有三元组 $i, j, k$）。由这些数据定义
$$
M = \Ker\left(
\bigoplus\nolimits_{1 \leq i \leq n} M_i
\longrightarrow
\bigoplus\nolimits_{1 \leq i, j \leq n} (M_i)_{f_j}
\right)
$$
其中 $(m_1, \ldots, m_n)$ 映到一个元素，其第 $(i, j)$ 个分量为
$m_i/1 - \psi_{ji}(m_j/1)$。则自然映射 $M \to M_i$ 诱导同构
$M_{f_i} \to M_i$。此外，对所有 $m \in M$，有
$\psi_{ij}(m/1) = m/1$（采用显然的记号）。
\end{lemma}

\begin{proof}
要证明 $M_{f_1} \to M_1$ 是同构，只须证明它在
$R_{f_1}$ 的每个素理想 $\mathfrak p'$ 处的局部化都是同构，见引理
\ref{algebra-lemma-characterize-zero-local}。由引理 \ref{algebra-lemma-standard-open}，
可写成 $\mathfrak p' = \mathfrak p R_{f_1}$，其中某个素理想
$\mathfrak p \subset R$ 满足 $f_1 \not \in \mathfrak p$。
由于局部化是正合的（命题 \ref{algebra-proposition-localization-exact}），有
\begin{align*}
(M_{f_1})_{\mathfrak p'} & =
M_\mathfrak p \\
& =
\Ker\left(
\bigoplus\nolimits_{1 \leq i \leq n} M_{i, \mathfrak p}
\longrightarrow
\bigoplus\nolimits_{1 \leq i, j \leq n} ((M_i)_{f_j})_\mathfrak p
\right) \\
& =
\Ker\left(
\bigoplus\nolimits_{1 \leq i \leq n} M_{i, \mathfrak p}
\longrightarrow
\bigoplus\nolimits_{1 \leq i, j \leq n} (M_{i, \mathfrak p})_{f_j}
\right)
\end{align*}
这里还使用了命题 \ref{algebra-proposition-localize-twice-module}。
由于 $f_1$ 在 $R_\mathfrak p$ 中是单位，将 $R$ 替换为
$R_\mathfrak p$、将 $f_i$ 替换为其在 $R_\mathfrak p$ 中的像 $f_i$、
将 $M$ 替换为 $M_\mathfrak p$，并将 $f_1$ 替换为 $1$，
便约化到 $f_1 = 1$ 的情形。

\medskip\noindent
假设 $f_1 = 1$。则对 $j = 2, \ldots, n$，
$\psi_{1j} : (M_1)_{f_j} \to M_j$ 是同构。
若用这些同构作恒等识别 $M_j = (M_1)_{f_j}$，则
$\psi_{ij} : (M_1)_{f_if_j} \to (M_1)_{f_if_j}$
就是典范恒等识别。因此复形
$$
0 \to M_1 \to
\bigoplus\nolimits_{1 \leq i \leq n} (M_1)_{f_i}
\longrightarrow
\bigoplus\nolimits_{1 \leq i, j \leq n}
(M_1)_{f_if_j}
$$
由引理 \ref{algebra-lemma-cover-module} 是正合的。因此在这种情形下，第一个映射将
$M_1$ 与 $M$ 等同，其余结论也都清楚。
\end{proof}

\section{零因子与全分式环}
\label{algebra-section-total-quotient-ring}

\noindent
极小素理想处的局部环具有下列性质。

\begin{lemma}
\label{algebra-lemma-minimal-prime-reduced-ring}
设 $\mathfrak p$ 为环 $R$ 的极小素理想。$R_{\mathfrak p}$ 的极大理想
中的每个元素都幂零。若 $R$ 既约，则 $R_{\mathfrak p}$ 是域。
\end{lemma}

\begin{proof}
若 ${\mathfrak p}R_{\mathfrak p}$ 的某个元素 $x$ 不幂零，则
$D(x) \not = \emptyset$，参见引理 \ref{algebra-lemma-Zariski-topology}。
这与 $\mathfrak p$ 的极小性矛盾。若 $R$ 既约，则
${\mathfrak p}R_{\mathfrak p} = 0$，因而该局部环是域。
\end{proof}

\begin{lemma}
\label{algebra-lemma-reduced-ring-sub-product-fields}
设 $R$ 为既约环。则：
\begin{enumerate}
\item $R$ 是某个域乘积的子环；
\item $R \to \prod_{\mathfrak p\text{ minimal}} R_{\mathfrak p}$
把它嵌入一个域的乘积；
\item $\bigcup_{\mathfrak p\text{ minimal}} \mathfrak p$ 是 $R$ 的
零因子之集。
\end{enumerate}
\end{lemma}

\begin{proof}
由引理 \ref{algebra-lemma-minimal-prime-reduced-ring}，每个环
$R_\mathfrak p$ 都是域。特别地，环映射 $R \to R_\mathfrak p$
的核是 $\mathfrak p$。由引理 \ref{algebra-lemma-Zariski-topology}，有
$\bigcap_{\mathfrak p} \mathfrak p = (0)$。所以 (2) 与 (1) 成立。
若 $x y = 0$ 且 $y \not = 0$，则存在极小素理想 $\mathfrak p$，
使得 $y \not \in \mathfrak p$。因此 $x \in \mathfrak p$。所以
$R$ 的每个零因子都属于
$\bigcup_{\mathfrak p\text{ minimal}} \mathfrak p$。反过来，假设对某个
极小素理想 $\mathfrak p$ 有 $x \in \mathfrak p$。于是 $x$ 在
$R_\mathfrak p$ 中映为零，故存在 $y \in R$、
$y \not \in \mathfrak p$，使得 $xy = 0$。换言之，$x$ 是零因子。
这完成了 (3) 及整个引理的证明。
\end{proof}

\noindent
环 $R$ 的全分式环 $Q(R)$ 已在例
\ref{algebra-example-localize-at-prime} 中引入。

\begin{lemma}
\label{algebra-lemma-total-ring-fractions}
设 $R$ 为环，$S \subset R$ 是由非零因子组成的乘法子集。则
$Q(R) \cong Q(S^{-1}R)$。特别地，$Q(R) \cong Q(Q(R))$。
\end{lemma}

\begin{proof}
若 $x \in S^{-1}R$ 是非零因子，并且对某个 $r \in R$、$f \in S$
有 $x = r/f$，则 $r$ 是 $R$ 中的非零因子。引理随即得证。
\end{proof}

\noindent
可以应用粘合结果证明关于全分式环 $Q(R)$ 的一个结论；该环已在例
\ref{algebra-example-localize-at-prime} 中引入。

\begin{lemma}
\label{algebra-lemma-total-ring-fractions-no-embedded-points}
设 $R$ 为环。假设 $R$ 有有限多个极小素理想
$\mathfrak q_1, \ldots, \mathfrak q_t$，并且
$\mathfrak q_1 \cup \ldots \cup \mathfrak q_t$ 是 $R$ 的零因子之集。
则全分式环 $Q(R)$ 等于
$R_{\mathfrak q_1} \times \ldots \times R_{\mathfrak q_t}$。
\end{lemma}

\begin{proof}
因为每个非零因子都属于 $R \setminus \mathfrak q_i$，存在自然映射
% P01-E0300: element/subset wording defect recorded in ERRATA_CANDIDATES.jsonl.
$Q(R) \to R_{\mathfrak q_i}$。因而有自然映射
$Q(R) \to R_{\mathfrak q_1} \times \ldots \times R_{\mathfrak q_t}$。
对任意非极小素理想 $\mathfrak p \subset R$，由引理
\ref{algebra-lemma-silly} 可知
$\mathfrak p \not \subset \mathfrak q_1 \cup \ldots \cup \mathfrak q_t$。
因此 $\Spec(Q(R)) = \{\mathfrak q_1, \ldots, \mathfrak q_t\}$
（将二者视为 $\Spec(R)$ 的子集，参见引理
\ref{algebra-lemma-spec-localization}）。所以 $\Spec(Q(R))$ 是有限离散集，
从而 $Q(R) = A_1 \times \ldots \times A_t$，其中
$\Spec(A_i) = \{\mathfrak{q}_i\}$，参见引理
\ref{algebra-lemma-disjoint-implies-product}。此外，$A_i$ 是 $R$ 的一个局部化，
并且是局部环。因此 $A_i \cong R_{\mathfrak q_i}$。
\end{proof}
\section{谱的不可约分支}
\label{algebra-section-irreducible}

\noindent
我们将证明，环谱的不可约分支对应于环中的极小素理想。

\begin{lemma}
\label{algebra-lemma-irreducible}
设 $R$ 为环。
\begin{enumerate}
\item 对素理想 $\mathfrak p \subset R$，单点集 $\{\mathfrak p\}$
在扎里斯基拓扑中的闭包为 $V(\mathfrak p)$。用公式表示即
$\overline{\{\mathfrak p\}} = V(\mathfrak p)$。
\item $\Spec(R)$ 的不可约闭子集恰为各子集 $V(\mathfrak p)$，其中
$\mathfrak p \subset R$ 为素理想。
\item $\Spec(R)$ 的不可约分支（参见拓扑，定义
\ref{topology-definition-irreducible-components}）恰为各子集
$V(\mathfrak p)$，其中 $\mathfrak p \subset R$ 为极小素理想。
\end{enumerate}
\end{lemma}

\begin{proof}
注意，若 $ \mathfrak p \in V(I)$，则 $I \subset \mathfrak p$。
所以显然 $\overline{\{\mathfrak p\}} = V(\mathfrak p)$。特别地，
$V(\mathfrak p)$ 是单点集的闭包，因而不可约。第二个断言蕴含第三个。
为证明第二个断言，设 $V(I) \subset \Spec(R)$，其中 $I$ 是根理想。
若 $I$ 不是素理想，则选取 $a, b\in R$、$a, b\not \in I$，使得
$ab\in I$。此时 $V(I, a) \cup V(I, b) = V(I)$；但由引理
\ref{algebra-lemma-Zariski-topology}，$V(I, b) = V(I)$ 与
$V(I, a) = V(I)$ 都不成立。因此 $V(I)$ 不可约。
\end{proof}

\noindent
换言之，该引理表明，$\Spec(R)$ 的每个不可约闭子集都形如
$V(\mathfrak p)$，其中 $\mathfrak p$ 为某个素理想。由于
$V(\mathfrak p) = \overline{\{\mathfrak p\}}$，每个不可约闭子集都有
唯一泛点，参见拓扑，定义 \ref{topology-definition-generic-point}。
特别地，$\Spec(R)$ 是清醒拓扑空间。把这个事实记录为下列引理。

\begin{lemma}
\label{algebra-lemma-spec-spectral}
环的谱是谱空间；参见拓扑，定义
\ref{topology-definition-spectral-space}。
\end{lemma}

\begin{proof}
形式上，这由引理 \ref{algebra-lemma-irreducible} 与引理
\ref{algebra-lemma-topology-spec} 推出。另见上文的讨论。
\end{proof}

\begin{lemma}
\label{algebra-lemma-irreducible-components-containing-x}
设 $R$ 为环，$\mathfrak p \subset R$ 为素理想。
\begin{enumerate}
\item $\Spec(R)$ 中经过 $\mathfrak p$ 的不可约闭子集之集，与素理想
$\mathfrak q \subset R_{\mathfrak p}$ 一一对应。
\item $\Spec(R)$ 中经过 $\mathfrak p$ 的不可约分支之集，与极小素理想
$\mathfrak q \subset R_{\mathfrak p}$ 一一对应。
\end{enumerate}
\end{lemma}

\begin{proof}
这由引理 \ref{algebra-lemma-irreducible} 及引理
\ref{algebra-lemma-spec-localization} 中对 $\Spec(R_\mathfrak p)$ 的描述推出；
后者说明 $\Spec(R_\mathfrak p)$ 对应于 $R$ 中满足
$\mathfrak q \subset \mathfrak p$ 的素理想 $\mathfrak q$。
\end{proof}

\begin{lemma}
\label{algebra-lemma-standard-open-containing-maximal-point}
设 $R$ 为环，$\mathfrak p$ 为 $R$ 的极小素理想。令
$W \subset \Spec(R)$ 为不包含点 $\mathfrak p$ 的拟紧开集。则存在
$f \in R$、$f \not \in \mathfrak p$，使得
$D(f) \cap W = \emptyset$。
\end{lemma}

\begin{proof}
因为 $W$ 拟紧，可以把它写成标准仿射开集 $D(g_i)$ 的有限并，其中
$i = 1, \ldots, n$。由于 $\mathfrak p \not \in W$，对所有 $i$ 都有
$g_i \in \mathfrak p$。由引理
\ref{algebra-lemma-minimal-prime-reduced-ring}，每个 $g_i$ 在
$R_{\mathfrak p}$ 中都幂零。因此可以找到
$f \in R$、$f \not \in \mathfrak p$，使得对所有 $i$，某个
$n_i > 0$ 满足 $f g_i^{n_i} = 0$。这个 $D(f)$ 即满足要求。
\end{proof}

\begin{lemma}
\label{algebra-lemma-ring-with-only-minimal-primes}
设 $R$ 为环，并把 $X = \Spec(R)$ 视为拓扑空间。下列条件等价：
\begin{enumerate}
\item $X$ 仿有限；
\item $X$ 是 Hausdorff 空间；
\item $X$ 全不连通；
\item $X$ 的每个拟紧开集都闭；
\item 其素理想之间不存在非平凡包含关系；
\item 每个素理想都是极大理想；
\item 每个素理想都是极小素理想；
\item 每个标准开集 $D(f) \subset X$ 都闭；
% P01-E0301: unfinished source placeholder recorded in ERRATA_CANDIDATES.jsonl.
\item 此处再添加条件。
\end{enumerate}
\end{lemma}

\begin{proof}
第一种证明。显然，(5)、(6)、(7) 等价。由于每个拟紧开集都是标准
开集的有限并，(4) 与 (8) 显然等价。蕴含关系 (7) $\Rightarrow$ (4)
由引理 \ref{algebra-lemma-standard-open-containing-maximal-point} 得到。
假设 (4) 成立。令 $\mathfrak p, \mathfrak p'$ 为 $R$ 的两个不同
素理想。选取 $f \in \mathfrak p'$、$f \not \in \mathfrak p$
（必要时交换 $\mathfrak p$ 与 $\mathfrak p'$）。于是
$\mathfrak p' \not \in D(f)$ 而 $\mathfrak p \in D(f)$。由 (4)，
开集 $D(f)$ 也闭。因此 $\mathfrak p$ 与 $\mathfrak p'$ 分别位于两个
不交开邻域中，且这两个邻域之并为 $X$。所以 $X$ 是 Hausdorff 且
全不连通的。因此 (4) $\Rightarrow$ (2) 和 (3)。若 (3) 成立，则
$\Spec(R)$ 的点之间不可能有特化关系，从而 (5) 成立。若 $X$ 是
Hausdorff 空间，则每个点都闭，所以 (2) 蕴含 (6)。因此 (2)、(3)、
(4)、(5)、(6)、(7)、(8) 等价。每个仿有限空间都是 Hausdorff 的，
所以 (1) 蕴含 (2)。若 $X$ 满足 (2) 与 (3)，则由引理
\ref{algebra-lemma-quasi-compact}，$X$ 拟紧；再由拓扑，引理
\ref{topology-lemma-profinite}，可知该空间仿有限。

\medskip\noindent
第二种证明。除 (4) 与 (8) 的等价外，其余结论由引理
\ref{algebra-lemma-spec-spectral} 及纯拓扑事实推出；参见拓扑，引理
\ref{topology-lemma-characterize-profinite-spectral}。
\end{proof}
\section{环谱的例}
\label{algebra-section-examples-spectra}

\noindent
本节给出若干谱的例子。

\begin{example}
\label{algebra-example-spec-Zxmodx2minus4}
本例描述 $X = \Spec(\mathbf{Z}[x]/(x^2 - 4))$。令 $\mathfrak{p}$ 为
$X$ 中的任意素理想，并令
$\phi : \mathbf{Z} \to \mathbf{Z}[x]/(x^2 - 4)$ 为自然环映射。于是
$ \phi^{-1}(\mathfrak p)$ 是 $\mathbf{Z}$ 中的素理想。若
$ \phi^{-1}(\mathfrak p) = (2)$，则因为 $\mathfrak p$ 包含 $2$，
它经由映射
$ \mathbf{Z}[x]/(x^2 - 4) \to  \mathbf{Z}[x]/(x^2 - 4, 2)$ 对应于
$$
\mathbf{Z}[x]/(x^2 - 4, 2) \cong (\mathbf{Z}/2\mathbf{Z})[x]/(x^2)
$$
中的素理想。$(\mathbf{Z}/2\mathbf{Z})[x]/(x^2)$ 中的任意素理想
对应于 $(\mathbf{Z}/2\mathbf{Z})[x]$ 中包含 $(x^2)$ 的素理想，因而
这样的素理想包含 $x$。因为
$(\mathbf{Z}/2\mathbf{Z}) \cong (\mathbf{Z}/2\mathbf{Z})[x]/(x)$ 是域，
$(x)$ 是极大理想。任意素理想都包含 $(x)$，而 $(x)$ 极大，故该环
仅有一个素理想 $(x)$。所以在这种情形下
$\mathfrak p = (2, x)$。现若 $ \phi^{-1}(\mathfrak p) = (q)$，其中
$q > 2$，则因为 $\mathfrak p$ 包含 $q$，它经由映射
$ \mathbf{Z}[x]/(x^2 - 4) \to  \mathbf{Z}[x]/(x^2 - 4, q)$ 对应于
$$
\mathbf{Z}[x]/(x^2 - 4, q) \cong (\mathbf{Z}/q\mathbf{Z})[x]/(x^2 - 4)
$$
中的素理想。$(\mathbf{Z}/q\mathbf{Z})[x]/(x^2 - 4)$ 中的任意素理想
对应于 $(\mathbf{Z}/q\mathbf{Z})[x]$ 中包含
$(x^2 - 4) = (x -2)(x + 2)$ 的素理想。因此，这些素理想必须包含
$x -2$ 或 $x + 2$。因为 $(\mathbf{Z}/q\mathbf{Z})[x]$ 是主理想整环，
所有非零素理想都极大，故 $(\mathbf{Z}/q\mathbf{Z})[x]$ 中恰有两个包含
$(x-2)(x + 2)$ 的素理想，
即 $(x-2)$ 与 $(x + 2)$。总之，存在两个素理想
$(q, x-2)$ 与 $(q, x + 2)$，因为
$2 \neq -2 \in \mathbf{Z}/(q)$。最后处理
$\phi^{-1}(\mathfrak p) = (0)$ 的情形。注意，$\mathfrak p$ 对应于
$\mathbf{Z}[x]$ 中包含
$(x^2 - 4) = (x -2)(x + 2)$ 的素理想。因此 $\mathfrak p$ 包含
$(x-2)$ 或 $(x + 2)$。所以依假设，$\mathfrak p$ 对应于
$\mathbf{Z}[x]/(x - 2)$ 中的一个素理想，或
$\mathbf{Z}[x]/(x + 2)$ 中的一个素理想，并且它与 $\mathbf{Z}$ 仅交于
$0$。由于 $\mathbf{Z}[x]/(x - 2) \cong \mathbf{Z}$ 且
$\mathbf{Z}[x]/(x + 2) \cong \mathbf{Z}$，这意味着 $\mathfrak p$
必须对应于其中一个环中的 $0$。因此在原环中，
$\mathfrak p = (x - 2)$ 或 $\mathfrak p = (x + 2)$。
\end{example}

\begin{example}
\label{algebra-example-spec-Zx}
本例描述 $X = \Spec(\mathbf{Z}[x])$。固定 $\mathfrak p \in X$。
令 $\phi : \mathbf{Z} \to \mathbf{Z}[x]$，并注意
$\phi^{-1}(\mathfrak p) \in \Spec(\mathbf{Z})$。若
$\phi^{-1}(\mathfrak p) = (q)$，其中 $q$ 是满足 $q > 0$ 的素数，
则 $\mathfrak p$ 对应于 $(\mathbf{Z}/(q))[x]$ 中的素理想，而该理想
必须由一个在 $(\mathbf{Z}/(q))[x]$ 中不可约的多项式生成。
% P01-E0302: omitted zero prime in positive-characteristic fibre recorded in ERRATA_CANDIDATES.jsonl.
若选择这个多项式的最小次数代表元，则它在 $\mathbf{Z}[x]$ 中也不可约。
% P01-E0303: primitive/monic lift qualification recorded in ERRATA_CANDIDATES.jsonl.
因此在这种情形下，$\mathfrak p = (q, f_q)$，其中 $f_q$ 是
$\mathbf{Z}[x]$ 中的不可约多项式，并且将其视为
$(\mathbf{Z}/(q) [x])$ 中的多项式时仍不可约。现假设
$\phi^{-1}(\mathfrak p) = (0)$。在这种情形下，$\mathfrak p$ 必须由
非常数多项式生成；由于 $\mathfrak p$ 为素理想，可假设这些多项式在
$\mathbf{Z}[x]$ 中不可约。
% P01-E0304: omitted zero ideal in generic fibre recorded in ERRATA_CANDIDATES.jsonl.
由高斯引理，这些多项式在 $\mathbf{Q}[x]$ 中也不可约。因为
$\mathbf{Q}[x]$ 是欧几里得整环，若至少有两个不同的不可约元 $f, g$
生成 $\mathfrak p$，则对某些 $a, b \in \mathbf{Q}[x]$，有
$1 = af + bg$。乘以一个公分母，可见对
$\bar{a}, \bar{b} \in \mathbf{Z}[x]$ 及某个非零的
$m \in \mathbf{Z}$，有 $m = \bar{a}f + \bar{b} g$，矛盾。因此
$\mathfrak p$ 由 $\mathbf{Z}[x]$ 中的一个不可约多项式生成。
\end{example}

\begin{example}
\label{algebra-example-spec-kxy}
设 $k$ 为任意域。本例描述 $X = \Spec(k[x, y])$。显然 $(0)$ 为
素理想；由于 $k[x, y]$ 是唯一分解整环，由不可约多项式生成的任何
主理想也都是素理想。现假设 $\mathfrak p$ 是 $X$ 中的非主理想。
因为 $k[x, y]$ 是诺特唯一分解整环，素理想 $\mathfrak p$ 可由有限个
不可约多项式 $(f_1, \ldots, f_n)$ 生成。现断言：若 $f, g$ 是
$k[x, y]$ 中不相伴的不可约多项式，则
$(f, g) \cap k[x] \neq 0$。为此，只需证明把 $f$ 与 $g$ 视为
$k(x)[y]$ 中的多项式时互素。在这种情形下，$k(x)[y]$ 是欧几里得整环；
应用欧几里得算法并清除分母，得到
$p = af + bg$，其中 $p, a, b \in k[x]$。
% P01-E0305: coefficient-ring defect after clearing denominators recorded in ERRATA_CANDIDATES.jsonl.
所以假设情况并非如此，即存在某个非单位 $h \in k(x)[y]$ 同时整除
$f$ 与 $g$。于是由高斯引理，对某些 $a, b \in k(x)$，有
$ah | f$ 与 $bh | g$，其中 $ah, bh \in k[x]$。
% P01-E0306: ambient polynomial-ring defect recorded in ERRATA_CANDIDATES.jsonl.
由不可约性，$ah = f$ 且 $bh = g$（因为 $h \notin k(x)$）。所以回到
$k(x)[y]$ 中，$f, g $ 相伴，因为 $\frac{a}{b} g = f$。由于 $k(x)$
是 $k[x]$ 的分式域，可对没有公因子的元素 $r , s \in k[x]$ 写成
$g = \frac{r}{s} f $。这蕴含在 $k[x, y]$ 中有 $sg = rf$；由于
$k[x, y]$ 是唯一分解整环，$s$ 必须整除 $f$。故 $s = 1$ 或
$s = f$。若 $s = f$，则 $r = g$，这意味着 $f, g \in k[x]$，因而
它们在 $k(x)$ 中必为单位，并在 $k(x)[y]$ 中互素，与假设矛盾。
若 $s = 1$，则 $g = rf$，仍为矛盾。因此 $f, g$ 必须在欧几里得整环
$k(x)[y]$ 中互素。由此归约到 $\mathfrak p$ 包含某个不可约多项式
$p \in k[x] \subset k[x, y]$ 的情形。由上文，$\mathfrak p$ 对应于环
$k[x, y]/(p) = k(\alpha)[y]$ 中的素理想，其中 $\alpha$ 是 $k$ 上的
代数元，其极小多项式为 $p$。这是主理想整环，故任一素理想对应于
$(0)$ 或 $k(\alpha)[y]$ 中的一个不可约多项式。因此
$\mathfrak p$ 形如 $(p)$ 或 $(p, f)$，其中 $f$ 是 $k[x, y]$ 中的
多项式，并且在商 $k[x, y]/(p)$ 中不可约。
\end{example}

\begin{example}
\label{algebra-example-affine-open-not-standard}
考虑环
$$
R = \{ f \in \mathbf{Q}[z]\text{ 其中 }f(0) = f(1) \}.
$$
考虑映射
$$
\varphi : \mathbf{Q}[A, B] \to R
$$
其中 $\varphi(A) = z^2-z$ 且 $\varphi(B) = z^3-z^2$。容易验证
$(A^3 - B^2 + AB) \subset \Ker(\varphi)$，并且
$A^3 - B^2 + AB$ 不可约。假设 $\varphi$ 满射；那么，因为 $R$ 是
整环（它是某个整环的子环），$\Ker(\varphi)$ 必须是
$\mathbf{Q}[A, B]$ 的素理想。包含 $(A^3-B^2 + AB)$ 的素理想是
$(A^3-B^2 + AB)$ 本身，以及任意满足如下条件的极大理想 $(f, g)$：
$f, g\in\mathbf{Q}[A, B]$，且 $f$ 模 $g$ 不可约。
% P01-E0307: unsupported maximal-ideal classification recorded in ERRATA_CANDIDATES.jsonl.
但 $R$ 不是域，所以核必为 $(A^3-B^2 + AB)$；因此 $\varphi$ 给出同构
$R \to \mathbf{Q}[A, B]/(A^3-B^2 + AB)$。
% P01-E0308: induced-isomorphism direction recorded in ERRATA_CANDIDATES.jsonl.

\medskip\noindent
为看出 $\varphi$ 满射，必须把任意 $f\in R$ 表成 $A(z) = z^2-z$
与 $B(z) = z^3-z^2$ 的 $\mathbf{Q}$-系数多项式。注意关系
$zA(z) = B(z)$。令 $a = f(0) = f(1)$。于是 $z(z-1)$ 必须整除
$f(z)-a$，所以可写成
$f(z) = z(z-1)g(z)+a = A(z)g(z)+a$。若 $\deg(g) < 2$，则
$h(z) = c_1z + c_0$，且
$f(z) = A(z)(c_1z + c_0)+a = c_1B(z)+c_0A(z)+a$，故结论成立。
% P01-E0309: h/g variable mismatch recorded in ERRATA_CANDIDATES.jsonl.
若 $\deg(g)\geq 2$，则由多项式除法算法，可写成
$g(z) = A(z)h(z)+b_1z + b_0$（$\deg(h)\leq\deg(g)-2$），所以
$f(z) = A(z)^2h(z)+b_1B(z)+b_0A(z)$。
% P01-E0310: omitted additive constant recorded in ERRATA_CANDIDATES.jsonl.
继续对 $h(z)$ 应用除法并迭代，可将 $f(z)$ 表成 $A(z)$ 与 $B(z)$
的多项式；故 $\varphi$ 满射。

\medskip\noindent
现令 $a \in \mathbf{Q}$、$a \neq 0, \frac{1}{2}, 1$，并考虑
$$
R_a = \{ f \in \mathbf{Q}[z, \frac{1}{z-a}]\text{ 其中 }f(0) = f(1)
\}.
$$
这同样是有限生成 $\mathbf{Q}$-代数：容易验证函数
$z^2-z$、$z^3-z$ 与 $\frac{a^2-a}{z-a}+z$ 作为
$\mathbf{Q}$-代数生成 $R_a$。有下列包含关系：
$$
R\subset R_a\subset\mathbf{Q}[z, \frac{1}{z-a}], \quad
R\subset\mathbf{Q}[z]\subset\mathbf{Q}[z, \frac{1}{z-a}].
$$
回顾（引理 \ref{algebra-lemma-spec-localization}），对环 T 与乘法子集
$S\subset T$，环映射 $T \to S^{-1}T$ 诱导谱上的映射
$\Spec(S^{-1}T) \to \Spec(T)$，它同胚到子集
$$
\{\mathfrak p \in \Spec(T) \mid S \cap \mathfrak p = \emptyset\}
\subset \Spec(T).
$$
若对某个 $f\in T$ 有 $S = \{ 1, f, f^2, \ldots\}$，则这是开集
$D(f)\subset T$。
% P01-E0311: standard-open ambient-space type error recorded in ERRATA_CANDIDATES.jsonl.
现验证环映射 $R \to R_a$ 的相应性质：我们将通过验证 $\theta$ 是单射
局部同胚，证明由包含 $R\subset R_a$ 诱导的映射
$\theta : \Spec(R_a) \to \Spec(R)$ 同胚到 $\Spec(R)$ 的一个开子集。
为此，使用下面将描述的两个基本开集所组成的 $\Spec(R_a)$ 开覆盖。
对任意 $r\in\mathbf{Q}$，令 $\text{ev}_r : R \to \mathbf{Q}$ 为在
$r$ 处求值给出的同态。注意，当 $r = 0$ 及 $r = 1-a$ 时，它可以
延拓为同态 $\text{ev}_r' : R_a \to \mathbf{Q}$（后者是因为
$\frac{1}{z-a}$ 在 $z = 1-a$ 处有定义，这是由于
$a\neq\frac{1}{2}$）。然而，$\text{ev}_a$ 不能延拓到 $R_a$。
写 $\mathfrak{m}_r = \Ker(\text{ev}_r)$。有
$$
\mathfrak{m}_0 = (z^2-z, z^3-z),
$$
$$
\mathfrak{m}_a = ((z-1 + a)(z-a), (z^2-1 + a)(z-a)), \text{ 且}
$$
$$
\mathfrak{m}_{1-a} = ((z-1 + a)(z-a), (z-1 + a)(z^2-a)).
$$
为验证这些等式，注意各右端显然包含于左端。然后把生成元用 $A$ 与
$B$ 表示，并将 $R$ 视为 $\mathbf{Q}[A, B]/(A^3-B^2 + AB)$，由此
检查各右端为极大理想。注意，$\mathfrak{m}_a$ 不在 $\theta$ 的像中：
$$
(z^2 - z)^2(z - a)\left(\frac{a^2 - a}{z - a} + z\right) =
(z^2 - z)^2(a^2 - a) + (z^2 - z)^2(z - a)z
$$
左端属于 $\mathfrak m_a R_a$，因为 $(z^2 - z)(z - a)$ 属于
$\mathfrak m_a$，并且
$(z^2 - z)(\frac{a^2 - a}{z - a} + z)$ 属于 $R_a$。类似地，元素
$(z^2 - z)^2(z - a)z$ 属于 $\mathfrak m_a R_a$，因为
$(z^2 - z)$ 属于 $R_a$，而 $(z^2 - z)(z - a)$ 属于
$\mathfrak m_a$。由于 $a \not \in \{0, 1\}$，可得
$(z^2 - z)^2 \in \mathfrak m_a R_a$。因此 $R_a$ 的任何理想 $I$
都不可能满足 $I \cap R = \mathfrak m_a$，因为这样的 $I$ 必须包含
$(z^2 - z)^2$，而它属于 $R$ 却不属于 $\mathfrak m_a$。基本开集
$D((z-1 + a)(z-a))\subset\Spec(R)$ 等于闭集
$\{\mathfrak{m}_a, \mathfrak{m}_{1-a}\}$ 的补集。然后验证
$R_{(z-1 + a)(z-a)} = (R_a)_{(z-1 + a)(z-a)}$；把这个局部化环记作
$R'$，便可知映射 $R \to R'$ 分解为 $R \to R_a \to R'$。于是由
引理 \ref{algebra-lemma-spec-localization}，这些映射把
$\Spec(R') \subset \Spec(R_a)$ 与 $\Spec(R') \subset \Spec(R)$
表示为开子集；所以 $\theta : \Spec(R_a) \to \Spec(R)$ 限制到
$D((z-1 + a)(z-a))$ 时，同胚到一个开子集。类似地，$\theta$ 限制到
$D((z^2 + z + 2a-2)(z-a)) \subset \Spec(R_a)$ 时，同胚到开子集
$D((z^2 + z + 2a-2)(z-a)) \subset \Spec(R)$。依
$z^2 + z + 2a-2$ 在 $\mathbf{Q}$ 上是否不可约，前一个基本开集的
补集由一或两个闭点以及闭点 $\mathfrak{m}_a$ 组成。此外，在 $R_a$
中由元素 $(z^2 + z + 2a-a)(z-a)$ 与 $(z-1 + a)(z-a)$ 生成的理想
是整个 $R_a$，所以这两个基本开集覆盖 $\Spec(R_a)$。
% P01-E0312: inconsistent coefficient in covering generator recorded in ERRATA_CANDIDATES.jsonl.
因此，为证明 $\theta$ 同胚到
$\Spec(R)-\{\mathfrak{m}_a\}$，只需证明这一或两个点绝不可能等于
$\mathfrak{m}_{1-a}$。事实确实如此，因为 $1-a$ 是
$z^2 + z + 2a-2$ 的根，当且仅当 $a = 0$ 或 $a = 1$；二者均不出现。

\medskip\noindent
尽管这个同胚仿效了在 $R$ 的某个元素处局部化的行为，而且
$\mathbf{Q}[z, \frac{1}{z-a}]$ 是 $\mathbf{Q}[z]$ 在极大理想
$(z-a)$ 处的局部化，环 $R_a$ 却{\it 不是} $R$ 的局部化：任何局部化
$S^{-1}R$ 都会产生比原环 $R$ 更多的单位。
% P01-E0313: localization-at-maximal-ideal error recorded in ERRATA_CANDIDATES.jsonl.
% P01-E0314: overbroad units claim recorded in ERRATA_CANDIDATES.jsonl.
$R$ 的单位是 $\mathbf{Q}^\times$，也就是 $\mathbf{Q}$ 的单位。事实上，
容易看出 $R_a$ 的单位为 $\mathbf{Q}^*$。具体而言，
$\mathbf{Q}[z, \frac{1}{z - a}]$ 的单位形如 $c (z - a)^n$，其中
$c \in \mathbf{Q}^*$ 且 $n \in \mathbf{Z}$；显然，只有当 $n = 0$
时，它们才属于 $R_a$。所以 $R_a$ 的单位并不比 $R$ 多，因而不可能是
$R$ 的局部化。

\medskip\noindent
我们使用 $a\neq 0, 1$ 来保证 $\frac{1}{z-a}$ 在 $z = 0, 1$ 处
有意义。还在若干处使用了 $a\neq 1/2$：(1) 为了能够讨论
$R_a$ 上 $\text{ev}_{1-a}$ 的核，这保证 $\mathfrak{m}_{1-a}$ 是
$R_a$ 的一个点（即 $R_a$ 恰缺少 $R$ 的一个点）。(2) 在最后，为了
推出 $(z-a)^{k + \ell}$ 只有在 $k = \ell = 0$ 时才能属于 $R$；
% P01-E0315: unexplained exponent variables recorded in ERRATA_CANDIDATES.jsonl.
事实上，若 $a = 1/2$，则只要 $k + \ell$ 为偶数，它就属于 $R$。
所以 $R_a$ 的单位确实会比 $R$ 更多，并且 $R_a$ 可能是 $R$ 的局部化。
\end{example}

\section{关于素理想的一则元观察}
\label{algebra-section-oka-families}

\noindent
本节取自 CRing 项目。设 $R$ 为环，$S \subset R$ 为乘法子集。
引理 \ref{algebra-lemma-spec-localization} 的一个推论是：关于“不与 $S$ 相交”
这一性质极大的理想 $I \subset R$ 是素理想。理由是，对局部化环
$S^{-1}R$ 的某个极大理想 $\mathfrak m$，有 $I = R \cap \mathfrak m$。
事实表明，对于理想的许多性质，不具有该性质的极大理想都是素理想。
\cite{Lam-Reyes} 发展了一种看出这一点的一般方法。本节暂离主线，
解释这一现象。

\medskip\noindent
设 $R$ 为环。若 $I$ 是 $R$ 的理想且 $a \in R$，定义
$$
(I : a) = \left\{ x \in R \mid xa \in I\right\}.
$$
更一般地，若 $J \subset R$ 是理想，定义
$$
(I : J) = \left\{ x \in R \mid xJ \subset I\right\}.
$$

\begin{lemma}
\label{algebra-lemma-colon}
设 $R$ 为环。若 $J \subset R$ 为主理想，则对任意理想
$I \subset J$，都有 $I = J (I : J)$。
\end{lemma}

\begin{proof}
设 $J = (a)$。则 $(I : J) = (I : a)$。由于 $I \subset J$，
任意 $y \in I$ 都可写成 $y = xa$，其中 $x \in (I : a)$。
故 $I \subset J (I : J)$。反之，若 $x \in (I : a)$，则
$xJ = (xa) \subset I$，这证明了另一包含关系。
\end{proof}

\noindent
设 $\mathcal{F}$ 为 $R$ 的一些理想组成的集合。我们关心怎样的条件
能够保证 $\mathcal{F}$ 的补集中的极大元为素理想。

\begin{definition}
\label{algebra-definition-oka-family}
设 $R$ 为环，$\mathcal{F}$ 为 $R$ 的理想组成的集合。若
$R \in \mathcal{F}$，并且只要 $I \subset R$ 是理想，且对某个
$a \in R$ 有 $(I : a), (I, a) \in \mathcal{F}$，便有
$I \in \mathcal{F}$，则称 $\mathcal{F}$ 为一个 {\it Oka 族}。
\end{definition}

\noindent
下面给出若干 Oka 族的例子。第一个就是本节引言中讨论的基本例子。

\begin{example}
\label{algebra-example-oka-family-not-meet-multiplicative-set}
设 $R$ 为环，$S$ 为 $R$ 的乘法子集。我们断言
$\mathcal{F} = \{I \subset R \mid I \cap S \not = \emptyset\}$
是 Oka 族。事实上，设对某个 $a \in R$ 有
$(I : a), (I, a) \in \mathcal{F}$。取
$s \in (I, a) \cap S$ 与 $s' \in (I : a) \cap S$。则
$ss' \in I \cap S$，因而 $I \in \mathcal{F}$。所以
$\mathcal{F}$ 是 Oka 族。
\end{example}

\begin{example}
\label{algebra-example-oka-family-finitely-generated}
设 $R$ 为环，$I \subset R$ 为理想，$a \in R$。若
$(I : a)$ 由 $a_1, \ldots, a_n$ 生成，而 $(I, a)$ 由
$a, b_1, \ldots, b_m$ 生成，其中 $b_1, \ldots, b_m \in I$，
则 $I$ 由
$aa_1, \ldots, aa_n, b_1, \ldots, b_m$ 生成。为说明这一点，
注意若 $x \in I$，则 $x \in (I, a)$ 是
$a, b_1, \ldots, b_m$ 的线性组合，但 $a$ 的系数必属于
$(I:a)$。由此可知，有限生成理想组成的族是 Oka 族。
\end{example}

\begin{example}
\label{algebra-example-oka-family-principal}
我们证明环 $R$ 的主理想组成的族是 Oka 族。事实上，设
$I \subset R$ 为理想，$a \in R$，并且 $(I, a)$ 与 $(I : a)$
都是主理想。注意 $(I : a) = (I : (I, a))$。令 $J = (I, a)$，
则 $J$ 与 $(I : J)$ 都是主理想。由引理 \ref{algebra-lemma-colon}，
有 $I = J (I : J)$。因此在这里，由于 $J = (I, a)$ 与
$(I : J)$ 都是主理想，$I$ 也是主理想。
\end{example}

\begin{example}
\label{algebra-example-oka-family-bound-cardinality}
设 $R$ 为环，$\kappa$ 为无限基数。至多可由 $\kappa$ 个元素生成的
理想组成的族是 Oka 族。论证与
例 \ref{algebra-example-oka-family-finitely-generated} 中的论证类似，故略去。
\end{example}

\begin{example}
\label{algebra-example-oka-family-property-modules}
设 $A$ 为环，$I \subset A$ 为理想，$a \in A$ 为元素。有短正合列
$0 \to A/(I : a) \to A/I \to A/(I, a) \to 0$，其中第一个箭头由
乘以 $a$ 给出。因此，若 $P$ 是一个对扩张稳定且对 $0$ 成立的
$A$-模性质，则使 $A/I$ 具有 $P$ 的那些理想 $I$ 组成 Oka 族。
\end{example}

\begin{proposition}
\label{algebra-proposition-oka}
若 $\mathcal{F}$ 是理想的 Oka 族，则 $\mathcal{F}$ 的补集中的
任意极大元都是素理想。
\end{proposition}

\begin{proof}
设 $I \not \in \mathcal{F}$ 关于不属于 $\mathcal{F}$ 极大，
但 $I$ 不是素理想。注意 $I \not = R$，因为
$R \in \mathcal{F}$。由于 $I$ 不素，可找到 $a, b \in R - I$，
使 $ab \in I$。于是 $(I, a) \neq I$，且 $(I : a)$ 包含
$b \not \in I$，故也有 $(I : a) \neq I$。所以
$(I : a), (I, a)$ 都严格包含 $I$，从而必属于 $\mathcal{F}$。
由 Oka 条件，$I \in \mathcal{F}$，矛盾。
\end{proof}

\noindent
至此，我们可以把上面大多数例子转化为关于环中素理想的引理。

\begin{lemma}
\label{algebra-lemma-simple}
设 $R$ 为环，$S$ 为 $R$ 的乘法子集。关于性质
$I \cap S = \emptyset$ 极大的理想 $I \subset R$ 是素理想。
\end{lemma}

\begin{proof}
这就是本节引言中讨论的例子。另一种证明是把
例 \ref{algebra-example-oka-family-not-meet-multiplicative-set}
与命题 \ref{algebra-proposition-oka} 结合起来。
\end{proof}

\begin{lemma}
\label{algebra-lemma-cohen}
设 $R$ 为环。
\begin{enumerate}
\item 关于“不有限生成”极大的理想 $I \subset R$ 是素理想。
\item 若 $R$ 的每个素理想都有限生成，则 $R$ 的每个理想都有限生成
\footnote{稍后我们将称 $R$ 为诺特环。}。
\end{enumerate}
\end{lemma}

\begin{proof}
第一条是例 \ref{algebra-example-oka-family-finitely-generated} 与
命题 \ref{algebra-proposition-oka} 的直接推论。对于第二条，设存在一个
不有限生成的理想 $I \subset R$。由不有限生成理想组成的全序链
$\left\{I_\alpha\right\}$ 的并仍不有限生成；事实上，若
$I = \bigcup I_\alpha$ 由 $a_1, \ldots, a_n$ 生成，则所有生成元
都属于某个 $I_\alpha $，因而会生成它。由佐恩引理，存在一个关于
不有限生成极大的理想。由第一条，该理想是素理想。
\end{proof}

\begin{lemma}
\label{algebra-lemma-primes-principal}
设 $R$ 为环。
\begin{enumerate}
\item 关于“不是主理想”极大的理想 $I \subset R$ 是素理想。
\item 若 $R$ 的每个素理想都是主理想，则 $R$ 的每个理想都是主理想。
\end{enumerate}
\end{lemma}

\begin{proof}
第一条由例 \ref{algebra-example-oka-family-principal} 与
命题 \ref{algebra-proposition-oka} 得出。对于第二条，设存在一个不是主理想的
理想 $I \subset R$。由不是主理想的理想组成的全序链
$\left\{I_\alpha\right\}$ 的并不是主理想；
% P01-E0316: missing copula in frozen source recorded in ERRATA_CANDIDATES.jsonl.
事实上，若 $I = \bigcup I_\alpha$ 由 $a$ 生成，则 $a$ 属于某个
$I_\alpha $，并且 $a$ 会生成该理想。由佐恩引理，存在一个关于不是主理想
极大的理想。由第一条，该理想必为素理想。
\end{proof}

\begin{lemma}
\label{algebra-lemma-characterize-domain}
设 $R$ 为环。
\begin{enumerate}
\item 在不含非零因子的理想中极大的理想是素理想。
\item 若 $R$ 非零，且 $R$ 中每个非零素理想都包含一个非零因子，
则 $R$ 是整环。
\end{enumerate}
\end{lemma}

\begin{proof}
考虑由非零因子组成的集合 $S$。它是 $R$ 的乘法子集。因此，
关于不与 $S$ 相交极大的任何理想都是素理想，见
引理 \ref{algebra-lemma-simple}。所以，若每个非零素理想都包含一个非零因子，
则 $(0)$ 是素理想，即 $R$ 是整环。
\end{proof}

\begin{remark}
\label{algebra-remark-cohen-bound-cardinality}
设 $R$ 为环，$\kappa$ 为无限基数。应用
例 \ref{algebra-example-oka-family-bound-cardinality} 与
命题 \ref{algebra-proposition-oka} 可知，关于不能由 $\kappa$ 个元素生成
这一性质极大的任何理想都是素理想。这个结果不太有用，因为存在一个环
$R$，它的每个素理想都可由 $\aleph_0$ 个元素生成，但某个理想却不可以。
具体地，设 $k$ 为域，$T$ 为基数大于 $\aleph_0$ 的集合，并令
$$
R = k[\{x_n\}_{n \geq 1}, \{z_{t, n}\}_{t \in T, n \geq 0}]/
(x_n^2, z_{t, n}^2, x_n z_{t, n} - z_{t, n - 1})
$$
这是一个局部环，其唯一素理想为 $\mathfrak m = (x_n)$。但是理想
$(z_{t, n})$ 不能由可数多个元素生成。
\end{remark}

\begin{example}
\label{algebra-example-noetherian-topology-on-spec}
\begin{reference}
Lukas Heger 于 2020 年 11 月 12 日的评论。
\end{reference}
设 $R$ 为环，$X = \Spec(R)$。由于 $X$ 的闭子集对应于 $R$ 的根理想
（引理 \ref{algebra-lemma-Zariski-topology}），故 $X$ 是诺特拓扑空间当且仅当
根理想满足升链条件。
% P01-E0317: “radical ideas” source typo recorded in ERRATA_CANDIDATES.jsonl.
这又当且仅当每个根理想都是一个有限生成理想的根（细节略去）。令
$$
\mathcal{F} = \{I \subset R \mid
\sqrt{I} = \sqrt{(f_1, \ldots, f_n)}\text{ 对某个 }n
\text{ 且 }f_1, \ldots, f_n \in R\}.
$$
读者可利用恒等式
$$
\sqrt{I} = \sqrt{(I, a)(I : a)}
$$
证明 $\mathcal{F}$ 是 Oka 族；该恒等式对任意理想 $I \subset R$
及任意元素 $a \in R$ 成立。另一方面，若有一条理想全序链
$\{I_\alpha\}$，其中没有理想属于 $\mathcal{F}$，则其并
$I = \bigcup I_\alpha$ 也不可能属于 $\mathcal{F}$。否则，若
$\sqrt{I} = \sqrt{(f_1, \ldots, f_n)}$，则对某个 $e$ 有
$f_i^e \in I$；于是对某个不依赖于 $i$ 的 $\alpha$，有
$f_i^e \in I_\alpha$；继而
$\sqrt{I_\alpha} = \sqrt{(f_1, \ldots, f_n)}$，矛盾。因此，
若不属于 $\mathcal{F}$ 的理想集合非空，它便有极大元；完全仿照
引理 \ref{algebra-lemma-cohen} 可知，$X$ 是诺特拓扑空间当且仅当 $R$ 的
每个素理想都等于某个 $\sqrt{(f_1, \ldots, f_n)}$，其中
$f_1, \ldots, f_n \in R$。若以后需要这个结果，我们会在此处
严谨地陈述并证明它。
\end{example}

\section{有限表示环映射的像}
\label{algebra-section-images-finite-presentation}

\noindent
本节证明由环映射 $R \to S$ 诱导的映射
$\Spec(S) \to \Spec(R)$ 的拓扑性质，主要是切瓦莱定理。为此，
我们将使用可构造集、拟紧集、逆紧集等概念；这些概念定义于
拓扑学，第 \ref{topology-section-constructible} 节。

\begin{lemma}
\label{algebra-lemma-qc-open}
设 $U \subset \Spec(R)$ 为开集。以下条件等价：
\begin{enumerate}
\item $U$ 在 $\Spec(R)$ 中逆紧；
\item $U$ 拟紧；
\item $U$ 是有限个标准开集之并；
\item 存在有限生成理想 $I \subset R$，使得
$X \setminus V(I) = U$。
% P01-E0318: undefined X in frozen source recorded in ERRATA_CANDIDATES.jsonl.
\end{enumerate}
\end{lemma}

\begin{proof}
因为 $\Spec(R)$ 拟紧，所以 (1) $\Rightarrow$ (2)，见
引理 \ref{algebra-lemma-quasi-compact}。因为标准开集构成拓扑的一组基，
所以 (2) $\Rightarrow$ (3)。证明 (3) $\Rightarrow$ (1)。令
$U = \bigcup_{i = 1\ldots n} D(f_i)$。要证明 $U$ 在 $\Spec(R)$ 中
逆紧，只须证明对 $\Spec(R)$ 的任意拟紧开集 $V$，$U \cap V$ 拟紧。
由 (2) $\Rightarrow$ (3)，可写
$V = \bigcup_{j = 1\ldots m} D(g_j)$。每个标准开集都同胚于某个环的谱，
故为拟紧的，见引理 \ref{algebra-lemma-standard-open} 与
\ref{algebra-lemma-quasi-compact}。因此
$U \cap V =
(\bigcup_{i = 1\ldots n} D(f_i)) \cap (\bigcup_{j = 1\ldots m} D(g_j))
= \bigcup_{i, j} D(f_i g_j)$ 是有限个拟紧开集之并，因而拟紧。
最后，由引理 \ref{algebra-lemma-Zariski-topology}，(4) 与 (3) 等价。
\end{proof}

\begin{lemma}
\label{algebra-lemma-affine-map-quasi-compact}
设 $\varphi : R \to S$ 为环映射。诱导的连续映射
$f : \Spec(S) \to \Spec(R)$ 是拟紧的。对任意可构造集
$E \subset \Spec(R)$，逆像 $f^{-1}(E)$ 在 $\Spec(S)$ 中可构造。
\end{lemma}

\begin{proof}
先证明任意拟紧开集 $U \subset \Spec(R)$ 的逆像拟紧。由
引理 \ref{algebra-lemma-qc-open}，可将 $U$ 写成有限个标准开集之并。
% P01-E0319: “finite open” source typo recorded in ERRATA_CANDIDATES.jsonl.
于是由引理 \ref{algebra-lemma-spec-functorial}，$f^{-1}(U)$ 是有限个标准开集
之并。再次应用引理 \ref{algebra-lemma-qc-open}，可知 $f^{-1}(U)$ 拟紧。
第二个断言由拓扑学，引理
\ref{topology-lemma-inverse-images-constructibles} 得出。
\end{proof}

\begin{lemma}
\label{algebra-lemma-constructible}
设 $R$ 为环。$\Spec(R)$ 的一个子集可构造，当且仅当它可写成有限个
形如 $D(f) \cap V(g_1, \ldots, g_m)$ 的子集之并，其中
$f, g_1, \ldots, g_m \in R$。
\end{lemma}

\begin{proof}
由引理 \ref{algebra-lemma-qc-open}，子集 $D(f)$ 与
$V(g_1, \ldots, g_m)$ 的补集都是逆紧开集。因此
$D(f) \cap V(g_1, \ldots, g_m)$ 是可构造子集，其任意有限并亦然。
反之，设 $T \subset \Spec(R)$ 可构造。由拓扑学，定义
\ref{topology-definition-constructible}，可设
$T = U \cap V^c$，其中 $U, V \subset \Spec(R)$ 是逆紧开集。
由引理 \ref{algebra-lemma-qc-open}，可写
$U = \bigcup_{i = 1, \ldots, n} D(f_i)$ 及
$V = \bigcup_{j = 1, \ldots, m} D(g_j)$。于是
$T = \bigcup_{i = 1, \ldots, n} \big(D(f_i) \cap V(g_1, \ldots, g_m)\big)$。
\end{proof}

\begin{lemma}
\label{algebra-lemma-constructible-is-image}
设 $R$ 为环，$T \subset \Spec(R)$ 可构造。则存在有限表示环映射
$R \to S$，使 $T$ 是 $\Spec(S)$ 在 $\Spec(R)$ 中的像。
\end{lemma}

\begin{proof}
有限个环之积的谱是这些谱的不交并，见
引理 \ref{algebra-lemma-spec-product}。因此，若 $T = T_1 \cup T_2$，且结论对
$T_1$ 与 $T_2$ 成立，则结论对 $T$ 也成立。由
引理 \ref{algebra-lemma-constructible}，可设
$T = D(f) \cap V(g_1, \ldots, g_m)$。在这种情形下，$T$ 是映射
$\Spec((R/(g_1, \ldots, g_m))_f) \to \Spec(R)$ 的像，见引理
\ref{algebra-lemma-standard-open} 与 \ref{algebra-lemma-spec-closed}。
\end{proof}

\begin{lemma}
\label{algebra-lemma-open-fp}
设 $R$ 为环，$f$ 为 $R$ 的元素，并令 $S = R_f$。则
$\Spec(S)$ 的可构造子集的像在 $\Spec(R)$ 中可构造。
\end{lemma}

\begin{proof}
以下不再说明地反复使用引理 \ref{algebra-lemma-qc-open}。设 $U, V$ 是
$\Spec(S)$ 中的拟紧开集。我们证明 $U \cap V^c$ 的像可构造。
在引理 \ref{algebra-lemma-standard-open} 给出的识别
$\Spec(S) = D(f)$ 下，集合 $U, V$ 对应于 $\Spec(R)$ 的拟紧开集
$U', V'$。因而只须证明 $U' \cap (V')^c$ 在 $\Spec(R)$ 中可构造，
这是显然的。
\end{proof}

\begin{lemma}
\label{algebra-lemma-closed-fp}
设 $R$ 为环，$I$ 为 $R$ 的有限生成理想，并令 $S = R/I$。
则 $\Spec(S)$ 的可构造子集的像在 $\Spec(R)$ 中可构造。
\end{lemma}

\begin{proof}
若 $I = (f_1, \ldots, f_m)$，则 $V(I)$ 是
$\bigcup D(f_i)$ 的补集，见引理 \ref{algebra-lemma-Zariski-topology}。
因此，由引理 \ref{algebra-lemma-qc-open}，它是可构造的。将映射
$R \to S$ 记作 $f \mapsto \overline{f}$。我们须证明：若
$\overline{U}, \overline{V}$ 是 $\Spec(S)$ 的逆紧开集，则
$\overline{U} \cap \overline{V}^c$ 在 $\Spec(R)$ 中的像可构造。
由引理 \ref{algebra-lemma-qc-open}，可写
$\overline{U} = \bigcup D(\overline{g_i})$。令
$U = \bigcup D({g_i})$，则 $\overline{U}$ 的像是
$U \cap V(I)$，它在 $\Spec(R)$ 中可构造。类似地，对
$\Spec(R)$ 的某个逆紧开集 $V$，$\overline{V}$ 的像等于
$V \cap V(I)$。故 $\overline{U} \cap \overline{V}^c$ 的像等于
$U \cap V(I) \cap V^c$，如所需。
\end{proof}

\begin{lemma}
\label{algebra-lemma-affineline-open}
设 $R$ 为环。映射 $\Spec(R[x]) \to \Spec(R)$ 是开映射，且任意
标准开集的像都是拟紧开集。
\end{lemma}

\begin{proof}
只须证明标准开集 $D(f)$（其中 $f\in R[x]$）的像是拟紧开集。
$D(f)$ 的像就是 $\Spec(R[x]_f) \to \Spec(R)$ 的像。设
$\mathfrak p \subset R$ 为素理想，并令 $\overline{f}$ 为 $f$ 在
$\kappa(\mathfrak p)[x]$ 中的像。回忆引理 \ref{algebra-lemma-in-image}：
$\mathfrak p$ 属于该像，当且仅当
$R[x]_f \otimes_R \kappa(\mathfrak p) =
\kappa(\mathfrak p)[x]_{\overline{f}}$ 不是零环。这恰好等价于 $f$
在 $\kappa(\mathfrak p)[x]$ 中的像非零；换言之，$f$ 的某个系数
不属于 $\mathfrak p$。因此，若
$f = a_d x^d + \ldots + a_0$，则 $D(f)$ 的像为
$D(a_d) \cup \ldots \cup D(a_0)$。
\end{proof}

\noindent
下面证明稍后将使用的特征多项式性质。

\begin{lemma}
\label{algebra-lemma-characteristic-polynomial-prime}
设 $R \to A$ 为环同态，并假设作为 $R$-模有
$A \cong R^{\oplus n}$。设 $f \in A$。乘法映射 $m_f: A
\to A$ 是 $R$-线性的，因而有特征多项式
$P(T) = T^n + r_{n-1}T^{n-1} + \ldots + r_0 \in R[T]$。
对任意素理想 $\mathfrak{p} \in \Spec(R)$，$f$ 在
$A \otimes_R \kappa(\mathfrak{p})$ 上幂零地作用，当且仅当
$\mathfrak p \in V(r_0, \ldots, r_{n-1})$。
\end{lemma}

\begin{proof}
只须证明：乘法映射
$m_{\bar f}: A \otimes_R \kappa(\mathfrak p) \to
A \otimes_R \kappa(\mathfrak p)$ 的特征多项式
$\bar P(T) \in \kappa(\mathfrak p)[T]$，恰为 $P(T)$ 在
$\kappa(\mathfrak p)[T]$ 中的像，结论便容易得出。这里该映射以
$\bar f$ 乘 $A \otimes_R \kappa(\mathfrak p)$ 的元素，而它是视为
$\kappa(\mathfrak p)$ 中元素的 $f$ 的像。
% P01-E0320: incorrect ambient ring for bar f recorded in ERRATA_CANDIDATES.jsonl.
以 $R$-模 $A$ 的一组基 $e_1, \ldots, e_n$ 表示乘法映射 $m_f$，
设其矩阵为 $(a_{ij})$，矩阵元属于 $R$。那么
$A \otimes_R \kappa(\mathfrak p) \cong (R \otimes_R
\kappa(\mathfrak p))^{\oplus n} = \kappa(\mathfrak p)^n$，这是
$\kappa(\mathfrak p)$ 上以
$e_1 \otimes 1, \ldots, e_n \otimes 1$ 为基的 $n$ 维向量空间。
像为 $\bar f = f \otimes 1$，故乘法映射 $m_{\bar f}$ 的矩阵是
$(a_{ij} \otimes 1)$。因此，特征多项式恰为 $P(T)$ 的像。

\medskip\noindent
由线性代数，一个线性变换在 $n$ 维向量空间上幂零地作用，
当且仅当其特征多项式是 $T^n$（因为特征多项式整除极小多项式的
某个幂）。所以，$f$ 在 $A \otimes_R \kappa(\mathfrak p)$ 上
幂零地作用，当且仅当 $\bar P(T) = T^n$。这又当且仅当对所有
$0 \leq i \leq n - 1$ 都有 $r_i \in \mathfrak p$，即
$\mathfrak p \in V(r_0, \ldots, r_{n - 1}).$
\end{proof}

\begin{lemma}
\label{algebra-lemma-affineline-special}
设 $R$ 为环，$f, g \in R[x]$ 为多项式。假设 $g$ 的首项系数是
$R$ 的单位。存在元素 $r_i\in R$（$i = 1\ldots, n$），使
$D(f) \cap V(g)$ 在 $\Spec(R)$ 中的像为
$\bigcup_{i = 1, \ldots, n} D(r_i)$。
% P01-E0321: subject-verb agreement defect recorded in ERRATA_CANDIDATES.jsonl.
\end{lemma}

\begin{proof}
写成 $g = ux^d + a_{d-1}x^{d-1} + \ldots + a_0$，其中 $d$ 是 $g$
的次数，故 $u \in R^*$。考虑环 $A = R[x]/(g)$。作为 $R$-模，
它是有限自由的，以 $1, x, \ldots, x^{d-1}$ 的像为基。考虑 $A$ 上
乘以 $f$ 的像。这是一个 $R$-模映射。因此可令 $P(T) \in R[T]$
为该映射的特征多项式。写成
$P(T) = T^d + r_{d-1} T^{d-1} + \ldots + r_0$。我们断言
$r_0, \ldots, r_{d-1}$ 具有所需性质。下面将使用特征多项式的性质
$$
\mathfrak p \in V(r_0, \ldots, r_{d-1})
\Leftrightarrow
\text{multiplication by }f\text{ is nilpotent on }
A \otimes_R \kappa(\mathfrak p).
$$
该性质已在引理 \ref{algebra-lemma-characteristic-polynomial-prime} 中证明。

\medskip\noindent
设 $\mathfrak q\in D(f) \cap V(g)$，并令
$\mathfrak p = \mathfrak q \cap R$。则有一个与乘以 $f$ 相容的非零映射
$A \otimes_R \kappa(\mathfrak p) \to \kappa(\mathfrak q)$。
而 $f$ 在 $\kappa(\mathfrak q)$ 上作为单位作用。因此可知
$\mathfrak p \not \in  V(r_0, \ldots, r_{d-1})$。

\medskip\noindent
另一方面，设对 $R$ 的某个素理想 $\mathfrak p$ 及某个
$0 \leq i \leq d - 1$，有 $r_i \not\in \mathfrak p$。则乘以 $f$
在代数 $A \otimes_R \kappa(\mathfrak p)$ 上不是幂零的。因此存在
不包含 $f$ 的像的素理想 $\overline{\mathfrak q} \subset
A \otimes_R \kappa(\mathfrak p)$。$\overline{\mathfrak q}$ 在
$R[x]$ 中的逆像是 $D(f) \cap V(g)$ 的一个元素，并映到
$\mathfrak p$。
\end{proof}

\begin{theorem}[切瓦莱定理]
\label{algebra-theorem-chevalley}
设 $R \to S$ 是有限表示的。则 $\Spec(S)$ 的可构造子集在
$\Spec(R)$ 中的像可构造。
\end{theorem}

\begin{proof}
写成 $S = R[x_1, \ldots, x_n]/(f_1, \ldots, f_m)$。可将
$R \to S$ 分解为
$R \to R[x_1] \to R[x_1, x_2]
\to \ldots \to R[x_1, \ldots, x_{n-1}] \to S$。所以可设
$S = R[x]/(f_1, \ldots, f_m)$。此时再将映射分解为
$R \to R[x] \to S$；由引理 \ref{algebra-lemma-closed-fp}，问题归结到
$S = R[x]$ 的情形。由引理 \ref{algebra-lemma-qc-open}，只须证明若
% P01-E0322: missing “it” in frozen source recorded in ERRATA_CANDIDATES.jsonl.
$T = (\bigcup_{i = 1\ldots n} D(f_i)) \cap V(g_1, \ldots, g_m)$，
其中 $f_i , g_j \in R[x]$，则其在 $\Spec(R)$ 中的像可构造。
由于可构造集的有限并可构造，只须处理 $n = 1$，即
$T = D(f) \cap V(g_1, \ldots, g_m)$ 的情形。

\medskip\noindent
注意若 $c \in R$，则
$$
\Spec(R) =
V(c) \amalg D(c) =
\Spec(R/(c)) \amalg \Spec(R_c),
$$
相应地，$\Spec(R[x]) =
V(c) \amalg D(c) = \Spec(R/(c)[x]) \amalg
\Spec(R_c[x])$。$T = D(f) \cap V(g_1, \ldots, g_m)$ 与每一部分的交
仍有相同形式，其中 $f$、$g_i$ 分别由其在 $R/(c)[x]$、$R_c[x]$
中的像代替。注意，$T$ 在 $\Spec(R)$ 中的像是
$T \cap V(c)$ 与 $T \cap D(c)$ 的像之并。利用引理
\ref{algebra-lemma-open-fp} 和 \ref{algebra-lemma-closed-fp}，只须分别证明这两部分
的像在 $\Spec(R/(c))$ 与 $\Spec(R_c)$ 中可构造。

\medskip\noindent
设如上有 $T = D(f) \cap V(g_1, \ldots, g_m)$，且
$\deg(g_1) \leq \deg(g_2) \leq \ldots \leq \deg(g_m)$。
我们将对 $m$ 以及诸 $g_i$ 的次数作归纳。令
$d_1 = \deg(g_1)$，即 $g_1 = c x^{d_1} + l.o.t$，其中非零的
$c \in R$。把 $R$ 分成 $R/(c)$ 与 $R_c$ 两部分后，要么降低
$g_1$ 的次数（由归纳覆盖），要么归结到 $c$ 可逆的情形。若 $c$ 可逆
且 $m > 1$，写成 $g_2 = c' x^{d_2} + l.o.t$。此时考虑
$g_2' = g_2 - (c'/c) x^{d_2 - d_1} g_1$。因为理想
$(g_1, g_2, \ldots, g_m)$ 与 $(g_1, g_2', g_3, \ldots, g_m)$
相等，所以 $T = D(f) \cap V(g_1, g_2', g_3\ldots, g_m)$。
% P01-E0323: missing separator after g_3 recorded in ERRATA_CANDIDATES.jsonl.
但这里 $g_2'$ 的次数严格小于 $g_2$ 的次数，故此情形由归纳覆盖。

\medskip\noindent
上述归纳的基本情形为：(a) $T = D(f) \cap V(g)$，其中 $g$ 的首项系数
可逆；以及 (b) $T = D(f)$。这两种情形分别由引理
\ref{algebra-lemma-affineline-special} 与 \ref{algebra-lemma-affineline-open} 处理。
% P01-E0324: “bases case” source typo recorded in ERRATA_CANDIDATES.jsonl.
\end{proof}

\section{关于像的进一步结果}
\label{algebra-section-more-images}

\noindent
本节汇集几个关于环映射在 $\Spec$ 上之像的补充引理。例如，
另见第 \ref{algebra-section-going-up} 节。

\begin{lemma}
\label{algebra-lemma-generic-finite-presentation}
设 $R \subset S$ 是整环的包含，并假设 $R \to S$ 是有限型的。
则存在非零的 $f \in R$ 与非零的 $g \in S$，使
$R_f \to S_{fg}$ 是有限表示的。
\end{lemma}

\begin{proof}
对 $S$ 在 $R$ 上的生成元个数作归纳。证明过程中，可将 $R$ 换成
$R_f$，将 $S$ 换成 $S_f$，其中 $f \in R$ 非零。

\medskip\noindent
设 $S$ 在 $R$ 上由一个元素生成。则对某个素理想
$\mathfrak q \subset R[x]$，有 $S = R[x]/\mathfrak q$。若
$\mathfrak q = (0)$，无须证明。若 $\mathfrak q \not = (0)$，
取 $h \in \mathfrak q$ 为关于 $x$ 次数最小的非零元素。写成
$h = f x^d + a_{d - 1} x^{d - 1} + \ldots + a_0$，其中
$a_i \in R$ 且 $f \not = 0$。在 $R$ 与 $S$ 中将 $f$ 可逆化后，
可设 $h$ 首一。于是得到满的 $R$-代数映射
$R[x]/(h) \to S$。作为 $R$-模，有
$R[x]/(h) = R \oplus Rx \oplus \ldots \oplus Rx^{d - 1}$；由 $d$
的极小性，$R[x]/(h)$ 到 $S$ 的映射是单射。因此
$R[x]/(h) \cong S$ 在 $R$ 上有限表示。

\medskip\noindent
设 $S$ 在 $R$ 上由 $n > 1$ 个元素生成。设
$x_1, \ldots, x_n \in S$ 生成 $S$，并记 $S' \subset S$ 为由
$x_1, \ldots, x_{n-1}$ 生成的 $R$-子代数。由归纳假设，存在非零的
$f\in R$ 与 $g \in S'$，使 $R_f \to S'_{fg}$ 有限表示。
% P01-E0325: missing article before “induction hypothesis” recorded in ERRATA_CANDIDATES.jsonl.
接着对 $S'_{fg} \to S_{fg}$ 应用归纳假设，可知存在
$f' \in S'_{fg}$ 与 $g' \in S_{fg}$，使
$S'_{fgf'} \to S_{fgf'g'}$ 有限表示。余下结论留给读者。
\end{proof}

\begin{lemma}
\label{algebra-lemma-characterize-image-finite-type}
设 $R \to S$ 为有限型环映射。记 $X = \Spec(R)$、
$Y = \Spec(S)$，并将诱导的谱映射写成 $f : Y \to X$。
% P01-E0326: malformed “Write ... the induced map” recorded in ERRATA_CANDIDATES.jsonl.
设 $E \subset Y = \Spec(S)$ 为可构造集。若点 $\xi \in X$ 属于
$f(E)$，则 $\overline{\{\xi\}} \cap f(E)$ 包含
$\overline{\{\xi\}}$ 的一个稠密开子集。
\end{lemma}

\begin{proof}
设 $\xi \in X$ 是 $f(E)$ 的一点。选取映到 $\xi$ 的点
$\eta \in E$。令 $\mathfrak p \subset R$ 为对应于 $\xi$ 的素理想，
$\mathfrak q \subset S$ 为对应于 $\eta$ 的素理想。考虑图
$$
\xymatrix{
\eta \ar[r] \ar@{|->}[d] & E \cap Y' \ar[r] \ar[d] &
Y' = \Spec(S/\mathfrak q) \ar[r] \ar[d] &
Y \ar[d] \\
\xi \ar[r] & f(E) \cap X' \ar[r] &
X' = \Spec(R/\mathfrak p) \ar[r] &
X
}
$$
由引理 \ref{algebra-lemma-affine-map-quasi-compact}，$E \cap Y'$ 在 $Y'$ 中
可构造。于是可以用 $X'$ 代替 $X$，用 $Y'$ 代替 $Y$。因此可设
$R \subset S$ 是整环的包含，$\xi$ 是 $X$ 的泛点，而 $\eta$ 是
$Y$ 的泛点。结合引理 \ref{algebra-lemma-generic-finite-presentation} 与
切瓦莱定理（定理 \ref{algebra-theorem-chevalley}），可知存在稠密开集
$U \subset X$、$V \subset Y$，使 $f(V) \subset U$，并且
$f : V \to U$ 将可构造集映为可构造集。注意 $E \cap V$ 在 $V$
中可构造，见拓扑学，引理
\ref{topology-lemma-open-immersion-constructible-inverse-image}。
所以 $f(E \cap V)$ 在 $U$ 中可构造并包含 $\xi$。由拓扑学，引理
\ref{topology-lemma-generic-point-in-constructible}，
$f(E \cap V)$ 包含一个稠密开集 $U' \subset U$。
\end{proof}

\noindent
本节最后给出几个与可构造集无关的谱映射之像的结果。

\begin{lemma}
\label{algebra-lemma-surjective-spec-radical-ideal}
设 $\varphi : R \to S$ 为环映射。以下条件等价：
\begin{enumerate}
\item 映射 $\Spec(S) \to \Spec(R)$ 是满射。
\item 对任意理想 $I \subset R$，$\sqrt{IS}$ 在 $R$ 中的逆像等于
$\sqrt{I}$。
\item 对任意根理想 $I \subset R$，$IS$ 在 $R$ 中的逆像等于 $I$。
\item 对 $R$ 的每个素理想 $\mathfrak p$，$\mathfrak p S$ 在 $R$
中的逆像是 $\mathfrak p$。
\end{enumerate}
在这种情形下，任意基变换后仍有同样的结论：给定环映射
$R \to R'$，环映射 $R' \to R' \otimes_R S$ 也具有等价性质
(1)、(2)、(3)。
\end{lemma}

\begin{proof}
若 $J \subset S$ 为理想，则
$\sqrt{\varphi^{-1}(J)} = \varphi^{-1}(\sqrt{J})$。这说明 (2) 与
(3) 等价。(3) $\Rightarrow$ (4) 是直接的。若 $I \subset R$ 是根理想，
则引理 \ref{algebra-lemma-Zariski-topology} 保证
$I = \bigcap_{I \subset \mathfrak p} \mathfrak p$。因此
(4) $\Rightarrow$ (2)。由引理 \ref{algebra-lemma-in-image}，有
$\mathfrak p = \varphi^{-1}(\mathfrak p S)$ 当且仅当
$\mathfrak p$ 属于该像。所以 (1) $\Leftrightarrow$ (4)。因此
(1)、(2)、(3)、(4) 等价。

\medskip\noindent
假设 (1) 成立。设 $R \to R'$ 为环映射，并设
$\mathfrak p' \subset R'$ 是位于 $R$ 的素理想 $\mathfrak p$ 之上的
素理想。要证明 $\mathfrak p'$ 属于
$\Spec(R' \otimes_R S) \to \Spec(R')$ 的像，由引理
\ref{algebra-lemma-in-image}，须证明
$(R' \otimes_R S) \otimes_{R'} \kappa(\mathfrak p')$ 非零。但有
$$
(R' \otimes_R S) \otimes_{R'} \kappa(\mathfrak p') =
S \otimes_R \kappa(\mathfrak p)
\otimes_{\kappa(\mathfrak p)} \kappa(\mathfrak p')
$$
它非零，因为由假设 $S \otimes_R \kappa(\mathfrak p)$ 非零，且
$\kappa(\mathfrak p) \to \kappa(\mathfrak p')$ 是域扩张。
\end{proof}

\begin{lemma}
\label{algebra-lemma-domain-image-dense-set-points-generic-point}
设 $R$ 为整环，$\varphi : R \to S$ 为环映射。以下条件等价：
\begin{enumerate}
\item 环映射 $R \to S$ 是单射。
\item $\Spec(S) \to \Spec(R)$ 的像包含一个稠密点集。
\item 存在素理想 $\mathfrak q \subset S$，它在 $R$ 中的逆像是
$(0)$。
\end{enumerate}
\end{lemma}

\begin{proof}
设 $K$ 为整环 $R$ 的分式域。假设 $R \to S$ 为单射。由于局部化正合，
可知 $K \to S \otimes_R K$ 是单射。因此，由引理
\ref{algebra-lemma-in-image}，存在一个映到 $(0)$ 的素理想。

\medskip\noindent
注意 $(0)$ 在 $\Spec(R)$ 中稠密，所以最后一个条件蕴含第二个条件。

\medskip\noindent
假设第二个条件成立。设 $f \in R$、$f \not = 0$。由于 $R$ 是整环，
可知 $V(f)$ 是 $R$ 的真闭子集。
% P01-E0327: V(f) has the wrong ambient space in frozen source; recorded in ERRATA_CANDIDATES.jsonl.
由假设，存在 $S$ 的素理想 $\mathfrak q$，使
$\varphi(f) \not \in \mathfrak q$。故 $\varphi(f) \not = 0$，
从而 $R \to S$ 是单射。
\end{proof}

\begin{lemma}
\label{algebra-lemma-injective-minimal-primes-in-image}
设 $R \subset S$ 为单射环映射。则 $\Spec(S) \to \Spec(R)$ 命中
所有极小素理想。
\end{lemma}

\begin{proof}
设 $\mathfrak p \subset R$ 为极小素理想。此时 $R_{\mathfrak p}$
有唯一素理想。因此只须证明 $S_{\mathfrak p}$ 非零。这由局部化正合
得出，见命题 \ref{algebra-proposition-localization-exact}。
\end{proof}

\begin{lemma}
\label{algebra-lemma-image-dense-generic-points}
设 $R \to S$ 为环映射。以下条件等价：
\begin{enumerate}
\item $R \to S$ 的核由幂零元组成。
\item $R$ 的极小素理想属于 $\Spec(S) \to \Spec(R)$ 的像。
\item $\Spec(S) \to \Spec(R)$ 的像在 $\Spec(R)$ 中稠密。
\end{enumerate}
\end{lemma}

\begin{proof}
令 $I = \Ker(R \to S)$。注意
$\sqrt{(0)} = \bigcap_{\mathfrak q \subset S} \mathfrak q$，见引理
\ref{algebra-lemma-Zariski-topology}。所以
$\sqrt{I} = \bigcap_{\mathfrak q \subset S} R \cap \mathfrak q$。
于是 $V(I) = V(\sqrt{I})$ 是 $\Spec(S) \to \Spec(R)$ 之像的闭包。
这说明 (1) 与 (3) 等价。显然 (2) 蕴含 (3)。最后假设 (1)。
可以用 $R/I$ 代替 $R$，用 $S/IS$ 代替 $S$，而不影响谱的拓扑和
映射。因此，(1) $\Rightarrow$ (2) 由引理
\ref{algebra-lemma-injective-minimal-primes-in-image} 得出。
\end{proof}

\begin{lemma}
\label{algebra-lemma-minimal-prime-image-minimal-prime}
设 $R \to S$ 为环映射。若极小素理想 $\mathfrak p \subset R$ 属于
$\Spec(S) \to \Spec(R)$ 的像，则它是某个极小素理想的像。
\end{lemma}

\begin{proof}
设 $\mathfrak p = \mathfrak q \cap R$。选取满足
$\mathfrak r \subset \mathfrak q$ 的极小素理想
$\mathfrak r \subset S$，见引理 \ref{algebra-lemma-Zariski-topology}。
由 $\mathfrak p$ 的极小性可知
$\mathfrak p = \mathfrak r \cap R$。
\end{proof}

\begin{lemma}
\label{algebra-lemma-intersection-not-zero}
设 $A \subset B$ 是整环的包含，且诱导分式域的代数扩张。若
$J \subset B$ 为非零理想，则 $A \cap J$ 也非零。因此，
$\Spec(B)$ 的真闭子集的像在 $\Spec(A)$ 中不稠密。
\end{lemma}

\begin{proof}
取非零元素 $x \in J$。由于 $x$ 在 $A$ 的分式域上代数，存在
$d \geq 1$ 及 $a_0, \ldots, a_d \in A$，其中
$a_0, a_d \not = 0$，使得在 $B$ 中
$a_d x^d + a_{d - 1} x^{d - 1} + \ldots + a_0 = 0$。
于是 $a_0 \in A \cap J$。
\end{proof}

\section{诺特环}
\label{algebra-section-Noetherian}

\noindent
若环 $R$ 的任意理想都有限生成，则称 $R$ 是{\it 诺特的}。
这显然等价于 $R$ 的理想满足升链条件。由引理 \ref{algebra-lemma-cohen}，
只须检验 $R$ 的每个素理想都有限生成。

\begin{lemma}
\label{algebra-lemma-Noetherian-permanence}
\begin{slogan}
诺特性质在取有限型扩张与局部化时保持。
\end{slogan}
诺特环上的任意有限生成环都是诺特环。诺特环的任意局部化都是诺特环。
\end{lemma}

\begin{proof}
关于局部化的断言来自如下事实：任意理想 $J \subset S^{-1}R$ 都形如
$I \cdot S^{-1}R$。诺特环 $R$ 的任意商 $R/I$ 都是诺特的，
因为任意理想 $\overline{J} \subset R/I$ 都形如 $J/I$，其中
$I \subset J \subset R$ 为理想。因此只须证明：若 $R$ 诺特，
则 $R[X]$ 也诺特。设
$J_1 \subset J_2 \subset \ldots$ 是 $R[X]$ 中的理想升链。
考虑理想 $I_{i, d}$，其定义为 $J_i$ 中次数为 $d$ 的多项式的
首项系数在 $R$ 中组成的理想。显然，只要 $i \leq i'$ 且
$d \leq d'$，就有 $I_{i, d} \subset I_{i', d'}$。由 $R$ 中的
升链条件，所有 $I_{i, d}$ 中至多只有有限个互异理想。
（提示：$\mathbf{N} \times \mathbf{N}$ 的任意无限子集都包含
一条递增无限序列。）取充分大的 $i_0$，使对所有 $i \geq i_0$ 及
所有 $d$，都有 $I_{i, d} = I_{i_0, d}$。设对某个
$i \geq i_0$ 有 $f \in J_i$。对次数 $d = \deg(f)$ 作归纳，
证明 $f \in J_{i_0}$。事实上，存在 $g\in J_{i_0}$，其次数为 $d$，
且首项系数与 $f$ 相同。由归纳，$f - g \in J_{i_0}$，结论得证。
\end{proof}

\begin{lemma}
\label{algebra-lemma-Noetherian-power-series}
若 $R$ 是诺特环，则形式幂级数环 $R[[x_1, \ldots, x_n]]$ 也是诺特环。
\end{lemma}

\begin{proof}
由于
$R[[x_1, \ldots, x_{n + 1}]] \cong R[[x_1, \ldots, x_n]][[x_{n + 1}]]$，
只须证明：若 $R$ 诺特，则 $R[[x]]$ 诺特。设 $I \subset R[[x]]$
为理想。须证明 $I$ 有限生成。对每个整数 $d$，记
$I_d = \{a \in R \mid ax^d + \text{h.o.t.} \in I\}$。由于 $R$ 诺特，
可知 $I_0 \subset I_1 \subset \ldots$ 稳定。选取 $d_0$，使
$I_{d_0} = I_{d_0 + 1} = \ldots$。对每个 $d \leq d_0$，选取元素
$f_{d, j} \in I \cap (x^d)$（$j = 1, \ldots, n_d$），使得若写成
$f_{d, j} = a_{d, j}x^d + \text{h.o.t}$，则
$I_d = (a_{d, j})$。记
$I' = (\{f_{d, j}\}_{d = 0, \ldots, d_0, j = 1, \ldots, n_d})$。
显然 $I' \subset I$。取 $f \in I$。首先可选取 $c_{d, i} \in R$，使
$$
f - \sum c_{d, i} f_{d, i} \in (x^{d_0 + 1}) \cap I.
$$
接着可选取 $c_{i, 1} \in R$（$i = 1, \ldots, n_{d_0}$），使
$$
f - \sum c_{d, i} f_{d, i} - \sum c_{i, 1}xf_{d_0, i} \in (x^{d_0 + 2}) \cap I.
$$
再选取 $c_{i, 2} \in R$（$i = 1, \ldots, n_{d_0}$），使
$$
f - \sum c_{d, i} f_{d, i} - \sum c_{i, 1}xf_{d_0, i}
- \sum c_{i, 2}x^2f_{d_0, i}
\in (x^{d_0 + 3}) \cap I.
$$
依此类推。最后可知
$$
f = \sum c_{d, i} f_{d, i} +
\sum\nolimits_i (\sum\nolimits_e c_{i, e} x^e)f_{d_0, i}
$$
属于 $I'$，如所需。
% P01-E0328: element/ideal “contained in” wording recorded in ERRATA_CANDIDATES.jsonl.
\end{proof}

\noindent
下面的引理虽然容易，却很有用，因为有限型 $\mathbf{Z}$-代数经常出现于
一种称为“绝对诺特约化”的技巧中。

\begin{lemma}
\label{algebra-lemma-obvious-Noetherian}
域上的任意有限型代数都是诺特的。$\mathbf{Z}$ 上的任意有限型代数
都是诺特的。
\end{lemma}

\begin{proof}
这直接来自引理 \ref{algebra-lemma-Noetherian-permanence}，以及域都是诺特环、
$\mathbf{Z}$ 是诺特环（因为它是主理想整环）这两个事实。
% P01-E0329: missing article in “is Noetherian ring” recorded in ERRATA_CANDIDATES.jsonl.
\end{proof}

\begin{lemma}
\label{algebra-lemma-Noetherian-finite-type-is-finite-presentation}
设 $R$ 为诺特环。
\begin{enumerate}
\item 任意有限 $R$-模都是有限表示的。
\item 有限 $R$-模的任意子模都是有限的。
\item 任意有限型 $R$-代数在 $R$ 上有限表示。
\end{enumerate}
\end{lemma}

\begin{proof}
设 $M$ 为有限 $R$-模。由引理
\ref{algebra-lemma-trivial-filter-finite-module}，可找到 $M$ 的有限滤过，
其逐次商形如 $R/I$。由于任意理想都有限生成，每个商 $R/I$
都有限表示。因此由引理 \ref{algebra-lemma-extension}，$M$ 有限表示。
这证明了 (1)。

\medskip\noindent
设 $N \subset M$ 为子模。由于 $M$ 有限，商 $M/N$ 有限。
由 (1)，$M/N$ 有限表示。因此，由引理 \ref{algebra-lemma-extension} 的 (5)，
$N$ 有限。这证明了 (2)。

\medskip\noindent
为证明 (3)，注意由引理 \ref{algebra-lemma-Noetherian-permanence}，
$R[x_1, \ldots, x_n]$ 的任意理想都有限生成。
\end{proof}

\begin{lemma}
\label{algebra-lemma-Noetherian-topology}
若 $R$ 是诺特环，则 $\Spec(R)$ 是诺特拓扑空间，见拓扑学，定义
\ref{topology-definition-noetherian}。
\end{lemma}

\begin{proof}
这是因为 $\Spec(R)$ 的任意闭子集都唯一地形如 $V(I)$，其中 $I$
是根理想，见引理 \ref{algebra-lemma-Zariski-topology}。并且这一对应反转包含关系。
所以结论由定义得出。
\end{proof}

\begin{lemma}
\label{algebra-lemma-Noetherian-irreducible-components}
\begin{slogan}
诺特仿射概形只有有限多个泛点。
\end{slogan}
若 $R$ 是诺特环，则 $\Spec(R)$ 只有有限多个不可约分支。换言之，
$R$ 只有有限多个极小素理想。
\end{lemma}

\begin{proof}
由引理 \ref{algebra-lemma-Noetherian-topology} 和拓扑学，引理
\ref{topology-lemma-Noetherian}，可知不可约分支只有有限多个。
由引理 \ref{algebra-lemma-irreducible}，它们对应于 $R$ 的极小素理想。
\end{proof}

\begin{lemma}
\label{algebra-lemma-Noetherian-base-change-finite-type}
设 $R \to S$ 为环映射，$R \to R'$ 为有限型映射。若 $S$ 诺特，
则基变换 $S' = R' \otimes_R S$ 诺特。
\end{lemma}

\begin{proof}
由引理 \ref{algebra-lemma-base-change-finiteness}，有限型在基变换下保持。
因此 $S \to S'$ 是有限型的。由于 $S$ 诺特，可以应用引理
\ref{algebra-lemma-Noetherian-permanence}。
\end{proof}

\begin{lemma}
\label{algebra-lemma-Noetherian-field-extension}
设 $k$ 为域，$R$ 为诺特 $k$-代数。若 $K/k$ 是有限生成域扩张，
则 $K \otimes_k R$ 诺特。
\end{lemma}

\begin{proof}
由于 $K/k$ 是有限生成域扩张，存在有限生成 $k$-代数
$B \subset K$，使 $K$ 是 $B$ 的分式域。换言之，
$K = S^{-1}B$，其中 $S = B \setminus \{0\}$。于是
$K \otimes_k R = S^{-1}(B \otimes_k R)$。由引理
\ref{algebra-lemma-Noetherian-base-change-finite-type}，$B \otimes_k R$
诺特。最后，由引理 \ref{algebra-lemma-Noetherian-permanence}，
$K \otimes_k R = S^{-1}(B \otimes_k R)$ 诺特。
\end{proof}

\noindent
下面是几个有时很有用的有趣引理。

\begin{lemma}
\label{algebra-lemma-subring-of-local-ring}
设 $R$ 为环，$\mathfrak p \subset R$ 为素理想。在以下每种情形下，
都存在 $f \in R$、$f \not \in \mathfrak p$，使
$R_f \to R_\mathfrak p$ 是单射：
\begin{enumerate}
\item $R$ 是整环；
\item $R$ 诺特；或
\item $R$ 约化且只有有限多个极小素理想。
\end{enumerate}
\end{lemma}

\begin{proof}
若 $R$ 是整环，则 $R \subset R_\mathfrak p$，故取 $f = 1$ 即可。
若 $R$ 诺特，则 $R \to R_\mathfrak p$ 的核 $I$ 是有限生成理想，
因而可找到 $f \in R$、$f \not \in \mathfrak p$，使 $IR_f = 0$。
对此 $f$，映射 $R_f \to R_\mathfrak p$ 是单射，故 $f$ 可用。
若 $R$ 约化，且极小素理想只有有限多个
$\mathfrak p_1, \ldots, \mathfrak p_n$，则可选择
$f \in \bigcap_{\mathfrak p_i \not \subset \mathfrak p} \mathfrak p_i$、
$f \not \in \mathfrak p$。事实上，若 $\mathfrak{p}_i\not\subset
\mathfrak{p}$，则存在 $f_i \in \mathfrak{p}_i$、
$f_i \not\in \mathfrak{p}$，并且 $f = \prod f_i$ 可用。对此 $f$，
有 $R_f \subset R_\mathfrak p$，因为 $R_f$ 的极小素理想对应于
$R_\mathfrak p$ 的极小素理想，并且可以应用引理
\ref{algebra-lemma-reduced-ring-sub-product-fields}（略去若干细节）。
\end{proof}

\begin{lemma}
\label{algebra-lemma-surjective-endo-noetherian-ring-is-iso}
诺特环的任意满自同态都是同构。
\end{lemma}

\begin{proof}
若 $f : R \to R$ 是这样的自同态但不是单射，则
$$
\Ker(f) \subset \Ker(f \circ f) \subset
\Ker(f \circ f \circ f) \subset \ldots
$$
将是一条严格递增的理想链。
\end{proof}

\section{局部幂零理想}
\label{algebra-section-locally-nilpotent}

\noindent
定义如下。

\begin{definition}
\label{algebra-definition-locally-nilpotent-ideal}
设 $R$ 为环，$I \subset R$ 为理想。若对每个 $x \in I$，都存在
$n \in \mathbf{N}$ 使 $x^n = 0$，则称 $I$ 是{\it 局部幂零的}。
若存在 $n \in \mathbf{N}$ 使 $I^n = 0$，则称 $I$ 是{\it 幂零的}。
\end{definition}

\begin{example}
\label{algebra-example-locally-nilpotent-not-nilpotent}
设 $R = k[x_n | n \in \mathbf{N}]$ 为域 $k$ 上含无穷多个变量的
多项式环。设 $I$ 为由 $x_n^n$（$n \in \mathbf{N}$）生成的理想，
并令 $S = R/I$。则由诸 $x_n$ 的像（$n \in \mathbf{N}$）生成的理想
$J \subset S$ 局部幂零，但并不幂零。事实上，由于幂零元的
$S$-线性组合仍为幂零元，要证明 $J$ 局部幂零，只须注意其所有生成元
都是幂零的（这显然成立）。另一方面，对每个 $n \in \mathbf{N}$，
都有 $x_{n + 1}^n \not \in I$，故 $J^n \not = 0$。所以 $J$ 不幂零。
\end{example}

\begin{lemma}
\label{algebra-lemma-locally-nilpotent}
设 $R \to R'$ 为环映射，$I \subset R$ 为局部幂零理想。
则 $IR'$ 是 $R'$ 的局部幂零理想。
\end{lemma}

\begin{proof}
这来自如下事实：若 $x, y \in R'$ 都幂零，则 $x + y$ 也幂零。
具体地，若 $x^n = 0$ 且 $y^m = 0$，则
$(x + y)^{n + m - 1} = 0$。
\end{proof}

\begin{lemma}
\label{algebra-lemma-locally-nilpotent-unit}
设 $R$ 为环，$I \subset R$ 为局部幂零理想。元素 $x$ 是 $R$ 的单位，
当且仅当 $x$ 在 $R/I$ 中的像是单位。
\end{lemma}

\begin{proof}
若 $x$ 是 $R$ 中的单位，则其像显然是 $R/I$ 中的单位。只须证明反向。
假设 $y \in R$ 的像是 $x$ 的像在 $R/I$ 中的逆元。于是对某个
$z \in I$，有 $xy = 1 - z$。这意味着模 $xR$ 有
$1\equiv z$。由于 $z$ 属于局部幂零理想 $I$，对某个充分大的 $N$，
有 $z^N = 0$。故模 $xR$ 有 $1 = 1^N \equiv z^N = 0$。
换言之，$x$ 整除 $1$，所以是单位。
\end{proof}

\begin{lemma}
\label{algebra-lemma-Noetherian-power}
\begin{slogan}
诺特环的一个理想若逐元素幂零，则该理想幂零。
\end{slogan}
设 $R$ 为诺特环，$I, J$ 为 $R$ 的理想。假设
$J \subset \sqrt{I}$。则对某个 $n$ 有 $J^n \subset I$。
特别地，在诺特环中，“局部幂零理想”与“幂零理想”这两个概念相同。
\end{lemma}

\begin{proof}
写成 $J = (f_1, \ldots, f_s)$。由假设 $f_i^{d_i} \in I$。
取 $n = d_1 + d_2 + \ldots + d_s + 1$。
\end{proof}

\begin{lemma}
\label{algebra-lemma-lift-idempotents}
设 $R$ 为环，$I \subset R$ 为局部幂零理想。则
$R \to R/I$ 在幂等元上诱导双射。
\end{lemma}

\begin{proof}[引理 \ref{algebra-lemma-lift-idempotents} 的第一种证明]
由于 $I$ 局部幂零，它包含于每个素理想。因此
$\Spec(R/I) = V(I) = \Spec(R)$。所以该引理由
引理 \ref{algebra-lemma-disjoint-decomposition} 得出。
\end{proof}

\begin{proof}[引理 \ref{algebra-lemma-lift-idempotents} 的第二种证明]
设 $\overline{e} \in R/I$ 是幂等元。须将 $\overline{e}$ 提升为 $R$
的幂等元。

\medskip\noindent
首先任选 $\overline{e}$ 的一个提升 $f \in R$，并令
$x = f^2 - f$。则 $x \in I$，所以 $x$ 幂零（因为 $I$ 局部幂零）。
现在令 $J$ 为 $R$ 中由 $x$ 生成的理想。由于它由幂零元 $x$ 生成，
$J$ 是幂零的，而不只是局部幂零的。

\medskip\noindent
现在假设已经找到 $\overline{e}$ 的一个提升 $e \in R$，使对某个
$k \geq 1$ 有 $e^2 - e \in J^k$。令
$e' = e - (2e - 1)(e^2 - e) = 3e^2 - 2e^3$；它是
$\overline{e}$ 的另一个提升（因为 $\overline{e}$ 的幂等性给出
$e^2 - e \in I$）。简单计算可得
$$
(e')^2 - e' = (4e^2 - 4e - 3)(e^2 - e)^2 \in J^{2k}
$$

\medskip\noindent
所以，我们从 $\overline{e}$ 的一个满足 $e^2 - e \in J^k$ 的提升
$e$ 出发，得到了 $\overline{e}$ 的一个满足
$(e')^2 - e' \in J^{2k}$ 的提升 $e'$。以此方式可逐次改进近似
（从 $e = f$ 开始，它对 $k = 1$ 符合要求）。最终到达
$J^k = 0$ 的阶段；此时有 $\overline{e}$ 的一个提升 $e$ 满足
$e^2 - e \in J^k = 0$，即该 $e$ 是幂等元。

\medskip\noindent
由此可见，若 $\overline{e} \in R/I$ 是任意幂等元，则存在一个
$\overline{e}$ 的提升，并且该提升是 $R$ 的幂等元。还须证明此提升唯一。设 $e_1$ 与 $e_2$
为两个这样的提升。须证明 $e_1 = e_2$。

\medskip\noindent
由 $e_1$ 与 $e_2$ 的定义，有 $e_1 \equiv e_2
\mod I$，并且 $e_1$ 与 $e_2$ 都幂等。由
$e_1 \equiv e_2 \mod I$ 可知 $e_1 - e_2 \in I$，故
$e_1 - e_2$ 幂零（因为 $I$ 局部幂零）。直接计算（使用 $e_1$ 与
$e_2$ 的幂等性）得到 $(e_1 - e_2)^3 = e_1 - e_2$。由此归纳可知，
对任意正奇数 $k$，都有 $(e_1 - e_2)^k = e_1 - e_2$。由于所有充分大的
$k$ 都满足 $(e_1 - e_2)^k = 0$（因为 $e_1 - e_2$ 幂零），
可知 $e_1 - e_2 = 0$，即 $e_1 = e_2$，证明完成。
\end{proof}

\begin{lemma}
\label{algebra-lemma-lift-idempotents-noncommutative}
设 $A$ 为可能不交换的代数，$e \in A$ 为满足
$x = e^2 - e$ 幂零的元素。则存在形如
$e' = e + x(\sum a_{i, j}e^ix^j) \in A$ 的幂等元，其中
$a_{i, j} \in \mathbf{Z}$。
\end{lemma}

\begin{proof}
考虑环 $R_n = \mathbf{Z}[e]/((e^2 - e)^n)$。显然，若能对每个 $R_n$
证明结论，则引理成立。在 $R_n$ 中考虑理想 $I = (e^2 - e)$，
并应用引理 \ref{algebra-lemma-lift-idempotents}。
\end{proof}

\begin{lemma}
\label{algebra-lemma-lift-nth-roots}
设 $R$ 为环，$I \subset R$ 为局部幂零理想。设 $n \geq 1$ 为在
$R/I$ 中可逆的整数。则
\begin{enumerate}
\item $n$ 次幂映射 $1 + I \to 1 + I$，$1 + x \mapsto (1 + x)^n$
是双射；
\item $R$ 的一个单位是 $n$ 次幂，当且仅当其在 $R/I$ 中的像是
$n$ 次幂。
% P01-E0330: “a nth power” article defect recorded in ERRATA_CANDIDATES.jsonl.
\end{enumerate}
\end{lemma}

\begin{proof}
设 $a \in R$ 为单位，且其在 $R/I$ 中的像与某个 $b \in R$ 的
$b^n$ 的像相同。则 $b$ 是单位（引理
\ref{algebra-lemma-locally-nilpotent-unit}），且对某个 $x \in I$ 有
$ab^{-n} = 1 + x$。由 (1)，$ab^{-n} = c^n$。所以 (2) 由 (1) 得出。

\medskip\noindent
证明 (1)。这是因为映射 $1 + x \mapsto (1 + x)^n$ 有逆映射。
具体地，可以考虑把 $1 + x$ 映为下式的映射：
\begin{align*}
(1 + x)^{1/n}
& =
1 + {1/n \choose 1}x +
{1/n \choose 2}x^2 +
{1/n \choose 3}x^3 + \ldots \\
& =
1 + \frac{1}{n} x + \frac{1 - n}{2n^2}x^2 +
\frac{(1 - n)(1 - 2n)}{6n^3}x^3 + \ldots
\end{align*}
这与初等微积分中一样。该式有意义，因为当 $k \gg 0$ 时
$x^k = 0$，所以级数是有限的；而且每个系数
${1/n \choose k} \in \mathbf{Z}[1/n]$（细节略去；注意由引理
\ref{algebra-lemma-locally-nilpotent-unit}，$n$ 在 $R$ 中可逆）。
\end{proof}



\section{一个有趣的问题}
\label{algebra-section-curiosity}

\noindent
引理 \ref{algebra-lemma-disjoint-implies-product} 说明了当某个理想
$I \subset R$ 的 $V(I)$ 为开集时会发生什么。但若
$\Spec(S^{-1}R)$ 在 $\Spec(R)$ 中是闭的，又会怎样？以下两个引理
给出部分答案。更多信息见第 \ref{algebra-section-pure-ideals} 节。

\begin{lemma}
\label{algebra-lemma-invert-closed-quotient}
设 $R$ 为环，$S \subset R$ 为乘法子集。假设映射
$\Spec(S^{-1}R) \to \Spec(R)$ 的像是闭的。则对某个理想
$I \subset R$，有 $S^{-1}R \cong R/I$。
\end{lemma}

\begin{proof}
令 $I = \Ker(R \to S^{-1}R)$，于是 $V(I)$ 包含该像。设该像为闭子集
$V(I') \subset \Spec(R)$，其中 $I' \subset R$ 为某个理想。因此
$V(I') \subset V(I)$。对 $f \in I'$，元 $f/1 \in S^{-1}R$ 属于每个
素理想。因此对某个 $n \geq 1$，$f^n$ 在 $S^{-1}R$ 中映为零
（引理 \ref{algebra-lemma-Zariski-topology}）。
% P01-E0331: source uses “is contained in every prime ideal” for an element; recorded in ERRATA_CANDIDATES.jsonl.
故 $V(I') = V(I)$。由此可知，每个 $g \in S$ 模 $I$ 都可逆。因此得到
环映射 $R/I \to S^{-1}R$ 与 $S^{-1}R \to R/I$。由 $I$ 的选取，
第一个映射是单射。第二个映射是
$S^{-1}R \to S^{-1}(R/I) = R/I$；由于局部化是正合的，其核为
$S^{-1}I$。因为 $S^{-1}I = 0$，第二个映射也是单射。因此
$S^{-1}R \cong R/I$。
\end{proof}

\begin{lemma}
\label{algebra-lemma-invert-closed-split}
设 $R$ 为环，$S \subset R$ 为乘法子集。假设映射
$\Spec(S^{-1}R) \to \Spec(R)$ 的像是闭的。若 $R$ 是诺特环，或
$\Spec(R)$ 是诺特拓扑空间，或 $S$ 作为幺半群是有限生成的，则对某个环
$R'$，有 $R \cong S^{-1}R \times R'$。
\end{lemma}

\begin{proof}
由引理 \ref{algebra-lemma-invert-closed-quotient}，对某个理想 $I \subset R$，
有 $S^{-1}R \cong R/I$。由引理 \ref{algebra-lemma-disjoint-implies-product}，
只须证明 $V(I)$ 是开集。若 $R$ 是诺特环，则 $\Spec(R)$ 是诺特拓扑空间，
见引理 \ref{algebra-lemma-Noetherian-topology}。若 $\Spec(R)$ 是诺特拓扑空间，
则补集 $\Spec(R) \setminus V(I)$ 是拟紧的，见拓扑学，引理
\ref{topology-lemma-Noetherian-quasi-compact}。因此存在有限多个
$f_1, \ldots, f_n \in I$，使得 $V(I) = V(f_1, \ldots, f_n)$。
由于每个 $f_i$ 在 $S^{-1}R$ 中都映为零，故存在 $g \in S$，使对
$i = 1, \ldots, n$ 都有 $gf_i = 0$。因而按所需有
$D(g) = V(I)$。若 $S$ 作为幺半群是有限生成的，设 $S$ 由
$g_1, \ldots, g_m$ 生成，则
$S^{-1}R \cong R_{g_1 \ldots g_m}$，从而
$V(I) = D(g_1 \ldots g_m)$。
\end{proof}



\section{希尔伯特零点定理}
\label{algebra-section-nullstellensatz}

\begin{theorem}[希尔伯特零点定理]
\label{algebra-theorem-nullstellensatz}
设 $k$ 为域。
\begin{enumerate}
\item
\label{algebra-item-finite-kappa}
对任意极大理想 $\mathfrak m \subset k[x_1, \ldots, x_n]$，域扩张
$\kappa(\mathfrak m)/k$ 是有限的。
\item
\label{algebra-item-polynomial-ring-Jacobson}
任意根理想 $I \subset k[x_1, \ldots, x_n]$ 都是包含它的极大理想之交。
\end{enumerate}
任意有限型 $k$-代数也满足同样的结论。
\end{theorem}

\begin{proof}
只须对多项式代数 $k[x_1, \ldots, x_n]$ 证明定理的
(\ref{algebra-item-finite-kappa})，因为任意有限生成 $k$-代数都是这种多项式代数
的商。对 $n$ 作归纳。$n = 0$ 的情形显然。设 $\mathfrak m$ 是
$k[x_1, \ldots, x_n]$ 中的极大理想。令
$\mathfrak p \subset k[x_n]$ 为 $\mathfrak m$ 与 $k[x_n]$ 的交。

\medskip\noindent
若 $\mathfrak p \not = (0)$，则 $\mathfrak p$ 是极大理想，并由一个
首一不可约多项式 $P$ 生成（因为 $k[x_n]$ 中有欧几里得算法）。于是
$k' = k[x_n]/\mathfrak p$ 是 $k$ 的有限域扩张，并包含于
$\kappa(\mathfrak m)$。在此情形，得到满射
$$
k'[x_1, \ldots, x_{n-1}]
\to
k'[x_1, \ldots, x_n] =
k' \otimes_k k[x_1, \ldots, x_n]
\longrightarrow
\kappa(\mathfrak m)
$$
因此由归纳假设，$\kappa(\mathfrak m)$ 是 $k'$ 的有限扩张。所以
$\kappa(\mathfrak m)$ 也是 $k$ 的有限扩张。

\medskip\noindent
若 $\mathfrak p = (0)$，考虑环扩张
$k[x_n] \subset k[x_1, \ldots, x_n]/\mathfrak m$。这是有限生成环扩张，
故由引理 \ref{algebra-lemma-obvious-Noetherian} 与
\ref{algebra-lemma-Noetherian-finite-type-is-finite-presentation}，它是有限表示的。
因而由定理 \ref{algebra-theorem-chevalley}，
$\Spec(k[x_1, \ldots, x_n]/\mathfrak m)$ 在 $\Spec(k[x_n])$ 中的像
是可构造的。由于该像包含 $(0)$，可知它包含某个非零
$f\in k[x_n]$ 的标准开集 $D(f)$。显然 $D(f)$ 是无限的，这与
$k[x_1, \ldots, x_n]/\mathfrak m$ 为域（因而其谱只有一个点）的
假设矛盾。

\medskip\noindent
证明 (\ref{algebra-item-polynomial-ring-Jacobson})。设
$I \subset R$ 为根理想，其中 $R$ 是 $k$ 上有限型的。取
$f \in R$ 且 $f \not \in I$。须找一个极大理想
$\mathfrak m \subset R$，使得 $I \subset \mathfrak m$ 且
$f \not \in \mathfrak m$。环 $(R/I)_f$ 非零，因为若其中 $1 = 0$，
则 $f^n \in I$；由于 $I$ 是根理想，这将推出 $f \in I$，与对 $f$
的假设矛盾。因此可在 $(R/I)_f$ 中选取一个极大理想
$\mathfrak m'$，见引理 \ref{algebra-lemma-Zariski-topology}。令
$\mathfrak m \subset R$ 为 $\mathfrak m'$ 在 $R$ 中的逆像。可见
$I \subset \mathfrak m$ 且 $f \not \in \mathfrak m$。若能证明
$\mathfrak m$ 是 $R$ 的极大理想，结论即得。显然有
$$
k \subset R/\mathfrak m \subset \kappa(\mathfrak m').
$$
由 (\ref{algebra-item-finite-kappa})，域扩张 $\kappa(\mathfrak m')/k$ 是有限的。
因而由域论，引理
\ref{fields-lemma-subalgebra-algebraic-extension-field}，$R/\mathfrak m$
是域。所以 $\mathfrak m$ 是极大理想，证明完成。
\end{proof}

\begin{lemma}
\label{algebra-lemma-field-finite-type-over-domain}
设 $R$ 为环，$K$ 为域。若 $R \subset K$ 且 $K$ 在 $R$ 上为有限型，
则存在 $f \in R$，使 $R_f$ 为域，并且 $K/R_f$ 是有限域扩张。
\end{lemma}

\begin{proof}
由引理 \ref{algebra-lemma-characterize-image-finite-type}，存在非空开集
$U \subset \Spec(R)$，它包含于 $\Spec(K) \to \Spec(R)$ 的像
$\{(0)\}$ 中。选取 $f \in R$、$f \not = 0$，使
$D(f) \subset U$，即 $D(f) = \{(0)\}$。于是 $R_f$ 是谱恰有一个点
的整环，故 $R_f$ 是域。此时 $K$ 是域 $R_f$ 上的有限生成代数，
因而由希尔伯特零点定理（定理 \ref{algebra-theorem-nullstellensatz}），
它是 $R_f$ 的有限域扩张。
\end{proof}



\section{雅可布森环}
\label{algebra-section-ring-jacobson}

\noindent
设 $R$ 为环。$\Spec(R)$ 的闭点就是 $R$ 的极大理想。代数几何中自然
出现的环往往有许多极大理想，例如域上或 $\mathbf{Z}$ 上的有限型代数。
我们将证明，这些都是雅可布森环的例子。

\begin{definition}
\label{algebra-definition-ring-jacobson}
设 $R$ 为环。若每个根理想 $I$ 都是包含它的极大理想之交，则称 $R$ 为
\textit{雅可布森环}。
\end{definition}

\begin{lemma}
\label{algebra-lemma-finite-type-field-Jacobson}
域上的任意有限型代数都是雅可布森环。
\end{lemma}

\begin{proof}
这由定理 \ref{algebra-theorem-nullstellensatz} 与定义
\ref{algebra-definition-ring-jacobson} 得出。
\end{proof}

\begin{lemma}
\label{algebra-lemma-jacobson-prime}
设 $R$ 为环。若 $R$ 的每个素理想都是包含它的极大理想之交，则 $R$ 是
雅可布森环。
\end{lemma}

\begin{proof}
这立即来自如下事实：每个根理想 $I \subset R$ 都是包含它的素理想之交。
见引理 \ref{algebra-lemma-Zariski-topology}。
\end{proof}

\begin{lemma}
\label{algebra-lemma-jacobson}
环 $R$ 是雅可布森环，当且仅当 $\Spec(R)$ 是雅可布森空间；见拓扑学，
定义 \ref{topology-definition-space-jacobson}。
\end{lemma}

\begin{proof}
先设 $R$ 是雅可布森环。令 $Z \subset \Spec(R)$ 为闭子集。须证明
$Z$ 中的闭点集在 $Z$ 中稠密。令 $U \subset \Spec(R)$ 为开集，且
$U \cap Z$ 非空。须证明 $Z \cap U$ 含有 $\Spec(R)$ 的一个闭点。
由于标准开集构成 $\Spec(R)$ 拓扑的一组基，可以设 $U = D(f)$。
由引理 \ref{algebra-lemma-Zariski-topology}，可以设 $Z = V(I)$，其中 $I$ 是
根理想。又有 $f \not \in I$。按假设，存在极大理想
$\mathfrak m \subset R$，使 $I \subset \mathfrak m$ 而
$f \not\in \mathfrak m$。因此按所需，
$\mathfrak m \in D(f) \cap V(I) = U \cap Z$。

\medskip\noindent
反之，设 $\Spec(R)$ 是雅可布森空间。令 $I \subset R$ 为根理想，并令
$J = \cap_{I \subset \mathfrak m} \mathfrak m$ 为包含 $I$ 的极大理想之交。
显然 $J$ 是根理想，$V(J) \subset V(I)$，并且 $V(J)$ 是 $V(I)$ 中
包含 $V(I)$ 全部闭点的最小闭子集。由假设可知 $V(J) = V(I)$。
但引理 \ref{algebra-lemma-Zariski-topology} 表明，扎里斯基闭集与根理想之间存在
双射，故按所需有 $I = J$。
\end{proof}

\begin{lemma}
\label{algebra-lemma-characterize-jacobson}
设 $R$ 为环。若 $R$ 不是雅可布森环，则存在素理想
$\mathfrak p \subset R$ 及元素 $f \in R$，满足：
\begin{enumerate}
\item $\mathfrak p$ 不是极大理想；
\item $f \not \in \mathfrak p$；
\item $V(\mathfrak p) \cap D(f) = \{\mathfrak p\}$；
\item $(R/\mathfrak p)_f$ 是域。
\end{enumerate}
另一方面，若 $R$ 是雅可布森环，则对任意满足 (1) 与 (2) 的
$(\mathfrak p, f)$，集合 $V(\mathfrak p) \cap D(f)$ 是无限的。
\end{lemma}

\begin{proof}
设 $R$ 不是雅可布森环。由引理 \ref{algebra-lemma-jacobson}，这意味着存在闭子集
$T \subset \Spec(R)$，其闭点集 $T_0 \subset T$ 在 $T$ 中不稠密。
选取 $f \in R$，使 $T_0 \subset V(f)$ 而 $T \not \subset V(f)$。
注意，若 $T = V(I)$，则 $T \cap D(f)$ 同胚于
$\Spec((R/I)_f)$；见引理 \ref{algebra-lemma-spec-closed} 与
\ref{algebra-lemma-standard-open}。由于任意环都有极大理想（引理
\ref{algebra-lemma-Zariski-topology}），可以选取空间 $T \cap D(f)$ 的一个闭点
$t$。于是 $t$ 对应于素理想 $\mathfrak p \subset R$；它不是极大理想
（因为 $t \not \in T_0$）。所以 (1) 成立。由构造，
$f \not \in \mathfrak p$，故 (2) 成立。由于 $t$ 是
$T \cap D(f)$ 的闭点，可知
$V(\mathfrak p) \cap D(f) = \{\mathfrak p\}$，即 (3) 成立。因此
$(R/\mathfrak p)_f$ 是谱只有一个点的整环，故 (4) 成立（例如合用引理
\ref{algebra-lemma-characterize-local-ring} 与
\ref{algebra-lemma-minimal-prime-reduced-ring}）。

\medskip\noindent
反之，设 $R$ 是雅可布森环，且 $(\mathfrak p, f)$ 满足 (1) 与 (2)。若
$V(\mathfrak p) \cap D(f) =
\{\mathfrak p, \mathfrak q_1, \ldots, \mathfrak q_t\}$，则由
$\mathfrak p \not = \mathfrak q_i$ 可知，存在元素 $g \in R$，使得
$g \not \in \mathfrak p$，但对所有 $i$ 都有 $g \in \mathfrak q_i$。
因而 $V(\mathfrak p) \cap D(fg) = \{\mathfrak p\}$。这是不可能的，
因为 $\Spec(R)$ 是雅可布森拓扑空间，所以 $\Spec(R)$ 的每个局部闭子集
至少含有一个闭点。
\end{proof}

\begin{lemma}
\label{algebra-lemma-pid-jacobson}
环 $\mathbf{Z}$ 是雅可布森环。更一般地，设环 $R$ 满足：
\begin{enumerate}
\item $R$ 是整环；
\item $R$ 是诺特环；
\item 任意非零素理想都是极大理想；
\item $R$ 有无限多个极大理想。
\end{enumerate}
则 $R$ 是雅可布森环。
\end{lemma}

\begin{proof}
设 $R$ 满足 (1)、(2)、(3) 与 (4)。该结论意指
$(0) = \bigcap_{\mathfrak m \subset R} \mathfrak m$。由于 $R$ 有无限多个
极大理想，只须证明任意非零 $x \in R$ 至多属于有限多个极大理想，换言之，
$V(x)$ 是有限的。
% P01-E0332: source uses “x is contained in ... maximal ideals” for an element; recorded in ERRATA_CANDIDATES.jsonl.
由引理 \ref{algebra-lemma-spec-closed}，$V(x)$ 同胚于 $\Spec(R/xR)$。由假设 (3)，
$R/xR$ 的每个素理想都是极小的，因而对应于 $\Spec(R/xR)$ 的一个不可约
分支（引理 \ref{algebra-lemma-irreducible}）。由于 $R/xR$ 是诺特环，拓扑空间
$\Spec(R/xR)$ 是诺特的（引理 \ref{algebra-lemma-Noetherian-topology}），且只有
有限多个不可约分支（拓扑学，引理 \ref{topology-lemma-Noetherian}）。
所以按所需，$V(x)$ 是有限的。
\end{proof}

\begin{example}
\label{algebra-example-infinite-product-fields-jacobson}
设 $A$ 为无限集。对每个 $\alpha \in A$，令 $k_\alpha$ 为域。我们断言
$R = \prod_{\alpha\in A} k_\alpha$ 是雅可布森环。首先注意，任意元素
$f \in R$ 都有形式 $f = ue$，其中 $u \in R$ 为单位，$e\in R$ 为
幂等元（留给读者）。因而 $D(f) = D(e)$，并且
$R_f = R_e = R/(1-e)$ 是 $R$ 的商。事实上，任何具有此性质的环都是
雅可布森环。设 $\mathfrak p \subset R$ 为素理想，且
$f \in R$、$f \not \in \mathfrak p$。须找出 $R$ 的一个极大理想
$\mathfrak m$，使 $\mathfrak p \subset \mathfrak m$ 且
$f \not\in \mathfrak m$。因为 $R_f$ 是 $R$ 的商，$R_f$ 的任意极大
理想都对应于 $R$ 的一个不含 $f$ 的极大理想。因此，选取 $R_f$ 的一个
包含 $\mathfrak p R_f$ 的极大理想即可得到结论。
\end{example}

\begin{example}
\label{algebra-example-not-jacobson}
设整环 $R$ 只有有限多个极大理想 $\mathfrak m_i$，
$i = 1, \ldots, n$。除非它是域，否则它不是雅可布森环。事实上，
在此情形，$(0)$ 不是极大理想之交，因为
$(0) \not =
\mathfrak m_1 \cap \mathfrak m_2 \cap \ldots \cap \mathfrak m_n
\supset \mathfrak m_1 \cdot \mathfrak m_2 \cdot \ldots
\cdot \mathfrak m_n \not = 0$。
特别地，离散赋值环，或任何至少有两个素理想的局部环，都不是雅可布森环。
\end{example}

\begin{lemma}
\label{algebra-lemma-finite-residue-extension-closed}
设 $R \to S$ 为环映射。令 $\mathfrak m \subset R$ 为极大理想，
$\mathfrak q \subset S$ 为位于 $\mathfrak m$ 之上的素理想，并且
$\kappa(\mathfrak q)/\kappa(\mathfrak m)$ 是代数域扩张。则
$\mathfrak q$ 是 $S$ 的极大理想。
\end{lemma}

\begin{proof}
考虑图表
$$
\xymatrix{
S \ar[r] & S/\mathfrak q \ar[r] & \kappa(\mathfrak q) \\
R \ar[r] \ar[u] & R/\mathfrak m \ar[u]
}
$$
可见 $\kappa(\mathfrak m) \subset S/\mathfrak q \subset
\kappa(\mathfrak q)$。由于域扩张
$\kappa(\mathfrak m) \subset \kappa(\mathfrak q)$ 是代数的，介于
$\kappa(\mathfrak m)$ 与 $\kappa(\mathfrak q)$ 之间的任意环都是域
（域论，引理 \ref{fields-lemma-subalgebra-algebraic-extension-field}）。
因此 $S/\mathfrak q$ 是域，并且事实上等于 $\kappa(\mathfrak q)$。
\end{proof}

\begin{lemma}
\label{algebra-lemma-dimension}
设 $k$ 为域，$V$ 为 $k$ 上的非零向量空间。假设 $V$ 的维数（这是一个
基数）小于 $k$ 的基数。则对任意线性算子 $T : V \to V$，都存在某个
首一多项式 $P(t) \in k[t]$，使 $P(T)$ 不可逆。
\end{lemma}

\begin{proof}
否则，$V$ 继承域 $k(t)$ 上的向量空间结构。但 $k(t)$ 在 $k$ 上的维数
至少是 $k$ 的基数；例如，这是因为诸元 $\frac{1}{t - \lambda}$ 在
$k$ 上线性无关。
\end{proof}

\noindent
下面是希尔伯特零点定理的另一版本。

\begin{theorem}
\label{algebra-theorem-uncountable-nullstellensatz}
设 $k$ 为域。令 $S$ 为由诸元 $\{x_i\}_{i \in I}$ 在 $k$ 上生成的
$k$-代数。假设 $I$ 的基数小于 $k$ 的基数。则
\begin{enumerate}
\item 对所有极大理想 $\mathfrak m \subset S$，域扩张
$\kappa(\mathfrak m)/k$ 是代数的；
\item $S$ 是雅可布森环。
\end{enumerate}
\end{theorem}

\begin{proof}
若 $I$ 有限，则结论由希尔伯特零点定理，即定理
\ref{algebra-theorem-nullstellensatz} 得出。以下假设 $I$ 无限。只须对以变量
$x_i$ 生成的多项式环中的极大理想
$\mathfrak m \subset k[\{x_i\}_{i \in I}]$ 证明结论，因为 $S$ 是它的
一个商。由于 $I$ 无限，所有单项式
$x_{i_1}^{e_1} \ldots x_{i_r}^{e_r}$ 的集合的基数至多等于 $I$ 的基数，
其中 $i_1, \ldots, i_r \in I$ 且 $e_1, \ldots, e_r \geq 0$。
这是因为 $I \times \ldots \times I$ 的基数等于 $I$ 的基数，并且
$\bigcup_{n \geq 0} I^n$ 的基数也相同。（若 $I$ 有限，这不成立；在该
情形，只有当 $k$ 不可数时，此证明才适用。）

\medskip\noindent
为得到矛盾，取 $T \in \kappa(\mathfrak m)$，使其在 $k$ 上超越。
注意，由乘以 $T$ 给出的 $k$-线性映射
$T : \kappa(\mathfrak m) \to \kappa(\mathfrak m)$ 具有如下性质：对所有
首一多项式 $P(t) \in k[t]$，$P(T)$ 都可逆。此外，
$\kappa(\mathfrak m)$ 在 $k$ 上的维数至多为 $I$ 的基数，因为它是
$k$ 上向量空间 $k[\{x_i\}_{i \in I}]$ 的商（如上所见，后者的维数为
$\# I$）。这与引理 \ref{algebra-lemma-dimension} 矛盾。

\medskip\noindent
为证明 $S$ 是雅可布森环，作如下论证。否则，由引理
\ref{algebra-lemma-characterize-jacobson}，存在素理想
$\mathfrak q \subset S$ 与元素 $f \in S$、$f \not \in \mathfrak q$，
使 $\mathfrak q$ 不是极大理想，而 $(S/\mathfrak q)_f$ 是域。但注意，
$(S/\mathfrak q)_f$ 可由至多 $\# I + 1$ 个元素生成。因此域扩张
$(S/\mathfrak q)_f/k$ 是代数的（由证明的第一部分）。这推出
$\kappa(\mathfrak q)$ 是 $k$ 的代数扩张，故由引理
\ref{algebra-lemma-finite-residue-extension-closed}，$\mathfrak q$ 是极大理想。
此矛盾完成证明。
\end{proof}

\begin{lemma}
\label{algebra-lemma-base-change-Jacobson}
设 $k$ 为域，$S$ 为 $k$-代数。对任意基数大于 $S$ 的基数的域扩张
$K/k$，有：
\begin{enumerate}
\item 对 $S_K$ 的每个极大理想 $\mathfrak m$，域
$\kappa(\mathfrak m)$ 在 $K$ 上是代数的；
\item $S_K$ 是雅可布森环。
\end{enumerate}
\end{lemma}

\begin{proof}
选取 $k \subset K$，使 $K$ 的基数大于 $S$ 的基数。由于 $S$ 的元素
生成 $K$-代数 $S_K$，可应用定理
\ref{algebra-theorem-uncountable-nullstellensatz}。
\end{proof}

\begin{example}
\label{algebra-example-countable-trick-does-not-work}
若 $k$ 是可数域且 $I$ 也可数，定理
\ref{algebra-theorem-uncountable-nullstellensatz} 证明中的技巧确实不适用。
设 $k$ 为可数域，$x$ 为变量，并令 $k(x)$ 为 $x$ 的有理函数域。考虑
多项式代数 $R = k[x, \{x_f\}_{f \in k[x]-\{0\}}]$。令
$I = (\{fx_f - 1\}_{f\in k[x] - \{0\}})$。注意，$I$ 是 $R$ 中的
真理想。选取极大理想 $I \subset \mathfrak m$。则
$k \subset R/\mathfrak m$ 同构于 $k(x)$，并且不在 $k$ 上代数。
\end{example}

\begin{lemma}
\label{algebra-lemma-Jacobson-invert-element}
设 $R$ 为雅可布森环，$f \in R$。环 $R_f$ 是雅可布森环，并且 $R_f$
的极大理想对应于 $R$ 的不含 $f$ 的极大理想。
\end{lemma}

\begin{proof}
由拓扑学，引理 \ref{topology-lemma-jacobson-inherited}，可知
$D(f) = \Spec(R_f)$ 是雅可布森空间，并且 $D(f)$ 的闭点对应于
$\Spec(R)$ 中恰好落在 $D(f)$ 内的闭点。因此由引理
\ref{algebra-lemma-jacobson}，$R_f$ 是雅可布森环。
\end{proof}

\begin{example}
\label{algebra-example-localize-not-preserve-closed-points}
下面这个简单例子表明，当 $R$ 不是雅可布森环时，引理
\ref{algebra-lemma-Jacobson-invert-element} 的结论不成立。考虑环
$R = \mathbf{Z}_{(2)}$，即 $\mathbf{Z}$ 在素理想 $(2)$ 处的局部化。
$R$ 在元素 $2$ 处的局部化同构于 $\mathbf{Q}$；写成公式即
$R_2 \cong \mathbf{Q}$。显然，映射 $R \to R_2$ 把
$\Spec(\mathbf{Q})$ 的闭点映到 $\Spec(R)$ 的泛点。
\end{example}

\begin{example}
\label{algebra-example-infinite-localize-not-preserve-closed-points}
下面这个简单例子表明，即使 $R$ 是雅可布森环，若对无限多个元素作局部化，
引理 \ref{algebra-lemma-Jacobson-invert-element} 的结论也不成立。令
$R = \mathbf{Z}$，并考虑 $R$ 在所有非零元素组成的集合处的局部化
$(R \setminus \{0\})^{-1}R \cong \mathbf{Q}$。显然，映射
$\mathbf{Z} \to \mathbf{Q}$ 把 $\Spec(\mathbf{Q})$ 的闭点映到
$\Spec(\mathbf{Z})$ 的泛点。
\end{example}

\begin{lemma}
\label{algebra-lemma-Jacobson-mod-ideal}
设 $R$ 为雅可布森环，$I \subset R$ 为理想。环 $R/I$ 是雅可布森环，
并且 $R/I$ 的极大理想对应于 $R$ 的包含 $I$ 的极大理想。
\end{lemma}

\begin{proof}
证明与引理 \ref{algebra-lemma-Jacobson-invert-element} 的证明相同。
\end{proof}

\begin{lemma}
\label{algebra-lemma-silly-jacobson}
设 $R$ 为雅可布森环，$K$ 为域。设 $R \subset K$，且 $K$ 在 $R$ 上
为有限型。则 $R$ 是域，并且 $K/R$ 是有限域扩张。
% P01-E0333: source joins “Let R subset K and K is ...” ungrammatically; recorded in ERRATA_CANDIDATES.jsonl.
\end{lemma}

\begin{proof}
首先注意，$R$ 是整环。由引理
\ref{algebra-lemma-field-finite-type-over-domain}，对某个非零 $f \in R$，
$R_f$ 是域且 $K/R_f$ 是有限域扩张。因此 $(0)$ 是 $R_f$ 的极大理想，
再由引理 \ref{algebra-lemma-Jacobson-invert-element} 可知，$(0)$ 是 $R$ 的
极大理想。
\end{proof}

\begin{proposition}
\label{algebra-proposition-Jacobson-permanence}
设 $R$ 为雅可布森环，$R \to S$ 为有限型环映射。则：
\begin{enumerate}
\item 环 $S$ 是雅可布森环；
\item 映射 $\Spec(S) \to \Spec(R)$ 把闭点映为闭点；
\item 对位于 $\mathfrak m \subset R$ 之上的极大理想
$\mathfrak m' \subset S$，域扩张
$\kappa(\mathfrak m')/\kappa(\mathfrak m)$ 是有限的。
\end{enumerate}
\end{proposition}

\begin{proof}
令 $\mathfrak m' \subset S$ 为极大理想，并令
$R \cap \mathfrak m' = \mathfrak m$。映射
$R/\mathfrak m \to S/\mathfrak m'$ 满足引理
\ref{algebra-lemma-silly-jacobson} 的条件；这是由引理
\ref{algebra-lemma-Jacobson-mod-ideal} 得到的。因此 $R/\mathfrak m$ 是域，
$\mathfrak m$ 是极大理想，并且诱导的剩余域扩张是有限的。这证明 (2)
与 (3)。

\medskip\noindent
若 $S$ 不是雅可布森环，则由引理 \ref{algebra-lemma-characterize-jacobson}，
存在 $S$ 的非极大素理想 $\mathfrak q$ 及
$g \in S$、$g \not\in \mathfrak q$，使 $(S/\mathfrak q)_g$ 是域。
为得到矛盾，证明 $\mathfrak q$ 是极大理想。令
$\mathfrak p = \mathfrak q \cap R$。映射
$R/\mathfrak p \to (S/\mathfrak q)_g$ 满足引理
\ref{algebra-lemma-silly-jacobson} 的条件；这是由引理
\ref{algebra-lemma-Jacobson-mod-ideal} 得到的。因此 $R/\mathfrak p$ 是域，且域扩张
$\kappa(\mathfrak p) \to (S/\mathfrak q)_g = \kappa(\mathfrak q)$
是有限的，从而是代数的。于是由引理
\ref{algebra-lemma-finite-residue-extension-closed}，$\mathfrak q$ 是 $S$ 的
极大理想。矛盾。
\end{proof}

\begin{lemma}
\label{algebra-lemma-corollary-jacobson}
$\mathbf{Z}$ 上的任意有限型代数都是雅可布森环。
\end{lemma}

\begin{proof}
合用引理 \ref{algebra-lemma-pid-jacobson} 与命题
\ref{algebra-proposition-Jacobson-permanence}。
\end{proof}

\begin{lemma}
\label{algebra-lemma-image-finite-type-map-Jacobson-rings}
设 $R \to S$ 为雅可布森环之间的有限型环映射。记
$X = \Spec(R)$、$Y = \Spec(S)$，并将诱导的谱映射记为
$f : Y \to X$。令 $E \subset Y = \Spec(S)$ 为可构造集。用下标
${}_0$ 表示拓扑空间的闭点集。
\begin{enumerate}
\item 有 $f(E)_0 = f(E_0) = X_0 \cap f(E)$。
\item 点 $\xi \in X$ 属于 $f(E)$，当且仅当
$\overline{\{\xi\}} \cap f(E_0)$ 在 $\overline{\{\xi\}}$ 中稠密。
\end{enumerate}
\end{lemma}

\begin{proof}
有连续映射的交换图
$$
\xymatrix{
E \ar[r] \ar[d] & Y \ar[d] \\
f(E) \ar[r] & X
}
$$
设 $x \in f(E)$ 在 $f(E)$ 中是闭的。则
$f^{-1}(\{x\})\cap E$ 非空且在 $E$ 中闭。对两个包含关系
$$
f^{-1}(\{x\}) \cap E \subset E \subset Y
$$
应用拓扑学，引理 \ref{topology-lemma-jacobson-inherited}，可找到点
$y \in f^{-1}(\{x\}) \cap E$，它在 $Y$ 中是闭的。换言之，存在
$y \in Y_0$ 且 $y \in E_0$，并映到 $x$。因此 $x \in f(E_0)$。
这证明 $f(E)_0 \subset f(E_0)$。命题
\ref{algebra-proposition-Jacobson-permanence} 蕴含
$f(E_0) \subset X_0 \cap f(E)$。包含关系
$X_0 \cap f(E) \subset f(E)_0$ 显然。这证明第一个断言。

\medskip\noindent
设 $\xi \in f(E)$。由引理 \ref{algebra-lemma-characterize-image-finite-type}，
集合 $f(E) \cap \overline{\{\xi\}}$ 包含
$\overline{\{\xi\}}$ 的一个稠密开子集。由于 $X$ 是雅可布森空间，
可知 $f(E) \cap \overline{\{\xi\}}$ 含有一个稠密闭点集；见拓扑学，
引理 \ref{topology-lemma-jacobson-inherited}。由本引理的 (1) 得出结论。

\medskip\noindent
另一方面，设 $\overline{\{\xi\}} \cap f(E_0)$ 在
$\overline{\{\xi\}}$ 中稠密。由
引理 \ref{algebra-lemma-constructible-is-image}，存在有限表示环映射
$S \to S'$，使 $E$ 是 $Y' := \Spec(S') \to Y$ 的像。对环映射
$S \to S'$ 应用本引理的第一部分，$E_0$ 是 $Y'_0$ 的像。因此以
$S'$ 代替 $S$，可以假设 $E = Y$。设 $\xi$ 对应于
$\mathfrak p \subset R$。考虑图表
$$
\xymatrix{
S \ar[r] & S/\mathfrak p S \\
R \ar[r] \ar[u] & R/\mathfrak p \ar[u]
}
$$
该图表以及 $f(Y_0) \cap V(\mathfrak p)$ 在 $V(\mathfrak p)$ 中的
稠密性表明，态射 $R/\mathfrak p \to S/\mathfrak p S$ 满足引理
\ref{algebra-lemma-domain-image-dense-set-points-generic-point} 的条件 (2)。
因此存在映到 $(0)$ 的素理想
$\overline{\mathfrak q} \subset S/\mathfrak pS$。换言之，
$\overline{\mathfrak q}$ 在 $S$ 中的逆像 $\mathfrak q$ 按所需映到
$\mathfrak p$。
\end{proof}

\noindent
上述引理的结论是：可以从像的闭点集读出 $f$ 的像。当映射是有限表示时，
这一点稍为更好，因为此时知道可构造集的像仍为可构造集。在陈述之前，先引入
一些记号。以 $\text{Constr}(X)$ 表示可构造集的集合。设
$R \to S$ 为环映射。记 $X = \Spec(R)$、$Y = \Spec(S)$，并将诱导的
谱映射记为 $f : Y \to X$。用下标 ${}_0$ 表示拓扑空间的闭点集。


\begin{lemma}
\label{algebra-lemma-conclude-jacobson-Noetherian}
沿用上述记号。假设 $R$ 是诺特雅可布森环，并进一步假设 $R \to S$
为有限型。存在交换图
$$
\xymatrix{
\text{Constr}(Y) \ar[r]^{E \mapsto E_0} \ar[d]^{E \mapsto f(E)} &
\text{Constr}(Y_0) \ar[d]^{E \mapsto f(E)} \\
\text{Constr}(X) \ar[r]^{E \mapsto E_0} &
\text{Constr}(X_0)
}
$$
其中水平箭头是拓扑学，引理
\ref{topology-lemma-jacobson-equivalent-constructible} 中的双射。
\end{lemma}

\begin{proof}
由于 $R \to S$ 为有限型，由引理
\ref{algebra-lemma-Noetherian-finite-type-is-finite-presentation}，它是有限表示的。
因此由切瓦莱定理（定理 \ref{algebra-theorem-chevalley}），$X$ 中可构造集的像
在 $Y$ 中是可构造的。结合引理
\ref{algebra-lemma-image-finite-type-map-Jacobson-rings} 即得本引理。
% P01-E0334: source reverses the domain/codomain of the constructible-image statement; recorded in ERRATA_CANDIDATES.jsonl.
\end{proof}

\noindent
为说明雅可布森环的用法，给出以下两个例子。

\begin{example}
\label{algebra-example-product-matrices-zero}
设 $k$ 为域。空间 $\Spec(k[x, y]/(xy))$ 有两个不可约分支，即 $x$ 轴
与 $y$ 轴。作为推广，令
$$
R = k[x_{11}, x_{12}, x_{21}, x_{22}, y_{11}, y_{12}, y_{21}, y_{22}]/
\mathfrak a,
$$
其中 $\mathfrak a$ 是
$k[x_{11}, x_{12}, x_{21}, x_{22}, y_{11}, y_{12}, y_{21}, y_{22}]$
中的理想，由下列 $2 \times 2$ 乘积矩阵的各个条目生成：
$$
\left(
\begin{matrix}
x_{11} & x_{12}\\
x_{21} & x_{22}
\end{matrix}
\right)
 \left(
\begin{matrix}
y_{11} & y_{12}\\
y_{21} & y_{22}
\end{matrix}
\right).
$$
本例将描述 $\Spec(R)$。

\medskip\noindent
为证明关于 $\Spec(k[x, y]/(xy))$ 的断言，论证如下。若
$\mathfrak p \subset k[x, y]$ 是任意包含 $xy$ 的理想，则 $x$ 或 $y$
应属于 $\mathfrak p$。因此这样的极小素理想恰为 $(x)$ 与 $(y)$。
% P01-E0335: source says arbitrary ideal, but the product implication requires a prime ideal; recorded in ERRATA_CANDIDATES.jsonl.
当 $k$ 代数闭时，这些分支的 $\text{max-Spec}$ 可分别视为 $y$ 轴与
$x$ 轴的点集。

\medskip\noindent
对上述推广，注意，至少在 $k$ 代数闭时，可以把
$k[x_{11}, x_{12}, x_{21}, x_{22}, y_{11}, y_{12}, y_{21}, y_{22}])$
的谱的闭点与如下矩阵空间等同：
% P01-E0336: source polynomial-ring expression has an unmatched closing parenthesis; recorded in ERRATA_CANDIDATES.jsonl.
$$
\left\{ (X, Y) \in \text{Mat}(2, k)\times \text{Mat}(2, k) \mid
X = \left(
\begin{matrix}
x_{11} & x_{12}\\
x_{21} & x_{22}
\end{matrix}
\right),
Y= \left(
\begin{matrix}
y_{11} & y_{12}\\
y_{21} & y_{22}
\end{matrix}
\right)
\right\}
$$
现在定义
$\text{GL}(2, k)\times \text{GL}(2, k)\times \text{GL}(2, k)$
在矩阵空间 $\{(X, Y)\}$ 上的群作用：
$$
(g_1, g_2, g_3) \times (X, Y) \mapsto ((g_1Xg_2^{-1}, g_2Yg_3^{-1})).
$$
这里还要注意，代数集
$$
\text{GL}(2, k)\times \text{GL}(2, k)\times \text{GL}(2, k) \subset
\text{Mat}(2, k)\times \text{Mat}(2, k) \times \text{Mat}(2, k)
$$
不可约，因为它是下列整环的极大谱：
$$
k[x_{11}, x_{12}, \ldots, z_{21}, z_{22}, (x_{11}x_{22}-x_{12}x_{21})^{-1}
, (y_{11}y_{22}-y_{12}y_{21})^{-1}, (z_{11}z_{22}-z_{12}z_{21})^{-1}].
$$
由于不可约代数集的像仍不可约，只须分类集合
$\{(X, Y)\in \text{Mat}(2, k)\times \text{Mat}(2, k)|XY = 0\}$ 的轨道，
再取其闭包。由标准线性代数，归结为以下三种情形：
\begin{enumerate}
\item $\exists (g_1, g_2)$，使 $g_1Xg_2^{-1} = I_{2\times 2}$。
此时 $Y$ 必为 $0$，而这作为代数集在群作用下不变。因此，此轨道包含于
由素理想 $(y_{11}, y_{12}, y_{21}, y_{22})$ 定义的不可约代数集中。
取闭包可知，$(y_{11}, y_{12}, y_{21}, y_{22})$ 确实是一个分支。
\item $\exists (g_1, g_2)$，使
$$
g_1Xg_2^{-1} = \left(
\begin{matrix}
1 & 0 \\
0 & 0
\end{matrix}
\right).
$$
这一情形当且仅当 $X$ 为秩 1 的矩阵；进一步，这样的 $X$ 消去 $Y$
当且仅当
$$
x_{11}y_{11}+x_{12}y_{21} = 0; \quad x_{11}y_{12}+x_{12}y_{22} = 0;
$$
$$
x_{21}y_{11}+x_{22}y_{21} = 0; \quad x_{21}y_{12}+x_{22}y_{22} = 0.
$$
固定一个秩 1 的 $X$。满足上述方程的非零 $Y$ 构成不可约代数集，
理由如下（$Y = 0$ 属于前一种情形）：由
$0 = g_1Xg_2^{-1}g_2Y$ 可知
$$
g_2Y = \left(
\begin{matrix}
0 & 0 \\
y_{21}' & y_{22}'
\end{matrix}
\right).
$$
再在右侧施以 $g_3$ 给出的 $\text{GL}(2, k)$ 作用，可将 $g_2Y$ 化为
$$
g_2Yg_3^{-1} = \left(
\begin{matrix}
0 & 0 \\
0 & 1
\end{matrix}
\right),
$$
所以这样的 $Y$ 构成不可约代数集，同构于 $\text{GL}(2, k)$ 在此作用下
的像。最后注意，$X$ 的“秩 1”条件构成不可约代数集
$\det X = x_{11}x_{22} - x_{12}x_{21} = 0$ 的稠密开子集。于是可知，
全部五个方程在非零矩阵对空间的开子集中定义一个不可约分支
$(x_{11}y_{11}+x_{12}y_{21}, x_{11}y_{12}+x_{12}y_{22}, x_{21}y_{11}
+x_{22}y_{21}, x_{21}y_{12}+x_{22}y_{22}, x_{11}x_{22}-x_{12}x_{21})$。
可以证明，方程对 $\det X = 0$、$\det Y = 0$ 在 $\Spec(R)$ 中截出
一个不可约分支，且上述轨迹是它的稠密开子集。
\item $\exists (g_1, g_2)$，使 $g_1Xg_2^{-1} = 0$；等价地，$X = 0$。
此时 $Y$ 可以任取，因而该分支由
$(x_{11}, x_{12}, x_{21}, x_{22})$ 定义。
\end{enumerate}
\end{example}


\begin{example}
\label{algebra-example-idempotent-matrices}
再举一例，考虑 $R = k[\{t_{ij}\}_{i, j = 1}^{n}]/\mathfrak a$，其中
$\mathfrak a$ 是由乘积矩阵 $T^2-T$ 的各条目生成的理想，
$T = (t_{ij})$。由线性代数可知，在由
$g, T \mapsto gTg^{-1}$ 定义的 $GL(n, k)$-作用下，$T$ 按其秩分类，
且每个 $T$ 都共轭于某个
$\text{diag}(1, \ldots, 1, 0, \ldots, 0)$，后者有 $r$ 个 1 与
$n-r$ 个 0。因此，这种 $\text{diag}(1, \ldots, 1, 0, \ldots, 0)$
在群作用下的每个轨道都构成一个不可约分支，而每个幂等矩阵都包含在某个
这样的轨道中。接下来证明，任意两个不同轨道必不相交。为此只需构造在不同
轨道上取不同值的多项式函数。在特征为 0 的情形，可以取
$f(t_{ij}) = trace(T) = \sum_{i = 1}^nt_{ii}$。在正特征情形，情况稍微
棘手，因为即使 $T \neq 0$，也可能有 $trace(T) = 0$。例如，当
$char = 3$ 时，
% P01-E0337: source has “classified by the its rank”; recorded in ERRATA_CANDIDATES.jsonl.
$$
trace\left(
\begin{matrix}
1 & & \\
& 1 & \\
& & 1
\end{matrix}
\right) = 3 = 0
$$
无论如何，可以用其他函数分离这些分支。例如在特征 3 的情形，
$tr(\wedge^3T)$ 在对应于 $diag(1, 1, 1)$ 的分支上取值 1，而在其他
分支上取值 0。
% P01-E0338: the last claim fails for general n (e.g. rank 4 in characteristic 3); recorded in ERRATA_CANDIDATES.jsonl.
\end{example}

\section{有限环扩张与整环扩张}
\label{algebra-section-finite-ring-extensions}

\noindent
关于有限环映射与整环映射的平凡引理。先回顾定义。

\begin{definition}
\label{algebra-definition-integral-ring-map}
设 $\varphi : R \to S$ 为环映射。
\begin{enumerate}
\item 元素 $s \in S$ 称为{\it 整于 $R$}，如果存在首一多项式
$P(x) \in R[x]$，使得 $P^\varphi(s) = 0$；其中
$P^\varphi(x) \in S[x]$ 是 $P$ 在
$\varphi : R[x] \to S[x]$ 下的像。
\item  若每个 $s \in S$ 都整于 $R$，则称环映射 $\varphi$ 是{\it 整的}。
\end{enumerate}
\end{definition}

\begin{lemma}
\label{algebra-lemma-characterize-integral-element}
设 $\varphi : R \to S$ 为环映射，$y \in S$。若存在 $S$ 的有限生成
$R$-子模 $M$，使 $1 \in M$ 且 $yM \subset M$，则 $y$ 整于 $R$。
\end{lemma}

\begin{proof}
考虑映射 $\varphi : M \to M$，$x \mapsto y \cdot x$。由引理
\ref{algebra-lemma-charpoly-module}，存在首一多项式 $P \in R[T]$，使
$P(\varphi) = 0$。在环 $S$ 中有
$P(y) = P(y) \cdot 1 = P(\varphi)(1) = 0$。
\end{proof}

\begin{lemma}
\label{algebra-lemma-finite-is-integral}
有限环映射是整的。
\end{lemma}

\begin{proof}
设 $R \to S$ 为有限环映射，$y \in S$。对 $M = S$ 应用引理
\ref{algebra-lemma-characterize-integral-element}，可知 $y$ 整于 $R$。
\end{proof}

\begin{lemma}
\label{algebra-lemma-characterize-integral}
设 $\varphi : R \to S$ 为环映射，且 $s_1, \ldots, s_n$ 是 $S$ 的有限个
元素。则下列条件等价：所有 $s_i$（$i = 1, \ldots, n$）都整于 $R$；
存在一个有限于 $R$ 的 $R$-子代数 $S' \subset S$，包含所有 $s_i$。
\end{lemma}

\begin{proof}
若每个 $s_i$ 都整，则由 $\varphi(R)$ 和各 $s_i$ 生成的子代数
在 $R$ 上有限。事实上，若 $s_i$ 在 $R$ 上满足次数为 $d_i$ 的首一方程，
则该子代数作为 $R$-模由满足 $0 \leq e_i \leq d_i - 1$ 的元素
$s_1^{e_1} \ldots s_n^{e_n}$ 生成。反之，若给定一个包含所有 $s_i$、
且为有限的 $R$-子代数 $S'$，则由引理
\ref{algebra-lemma-finite-is-integral}，所有 $s_i$ 都整。
\end{proof}

\begin{lemma}
\label{algebra-lemma-characterize-finite-in-terms-of-integral}
设 $R \to S$ 为环映射。以下条件等价：
\begin{enumerate}
\item $R \to S$ 是有限的；
\item $R \to S$ 是整的且为有限型的；
\item 存在 $x_1, \ldots, x_n \in S$，它们作为 $R$-代数生成 $S$，
且每个 $x_i$ 都整于 $R$。
\end{enumerate}
\end{lemma}

\begin{proof}
由引理 \ref{algebra-lemma-characterize-integral} 立即得到。
\end{proof}

\begin{lemma}
\label{algebra-lemma-integral-transitive}
\begin{slogan}
整环映射的复合仍为整环映射
\end{slogan}
设 $R \to S$ 与 $S \to T$ 是整环映射，则 $R \to T$ 是整环映射。
\end{lemma}

\begin{proof}
取 $t \in T$。令 $P(x) \in S[x]$ 为满足 $P(t) = 0$ 的首一多项式。
对 $P$ 的有限个系数应用引理
\ref{algebra-lemma-characterize-integral}。于是 $t$ 整于某个在 $R$ 上有限的
子代数 $S' \subset S$。再次应用引理
\ref{algebra-lemma-characterize-integral}，得到一个在 $S'$ 上有限且包含 $t$ 的
子代数 $T' \subset T$。对 $R \to S' \to T'$ 应用引理
\ref{algebra-lemma-finite-transitive}，可知 $T'$ 在 $R$ 上有限。于是由引理
\ref{algebra-lemma-finite-is-integral}，$t$ 整于 $R$。
\end{proof}

\begin{lemma}
\label{algebra-lemma-integral-closure-is-ring}
设 $R \to S$ 为环同态。集合
$$
S' = \{s \in S \mid s\text{ is integral over }R\}
$$
是 $S$ 的一个 $R$-子代数。
\end{lemma}

\begin{proof}
由引理 \ref{algebra-lemma-characterize-integral} 与
\ref{algebra-lemma-finite-is-integral}，结论显然。
\end{proof}

\begin{lemma}
\label{algebra-lemma-finite-product-integral}
设 $R_i\to S_i$（$i = 1, \ldots, n$）为环映射。令 $R$ 与 $S$ 分别为
所有 $R_i$ 与所有 $S_i$ 的乘积。则元素
$s = (s_1, \ldots, s_n) \in S$ 整于 $R$，当且仅当每个 $s_i$ 整于
$R_i$。
\end{lemma}

\begin{proof}
略。
\end{proof}

\begin{definition}
\label{algebra-definition-integral-closure}
设 $R \to S$ 为环映射。S 中整于 $R$ 的元素所成的环
$S' \subset S$（见引理 \ref{algebra-lemma-integral-closure-is-ring}）称为
$R$ 在 $S$ 中的{\it 整闭包}。若 $R \subset S$，且 $R = S'$，则称
$R$ 在 $S$ 中{\it 整闭}。
\end{definition}

\noindent
特别地，$R \to S$ 是整的，当且仅当 $R$ 在 $S$ 中的整闭包就是 $S$。

\begin{lemma}
\label{algebra-lemma-finite-product-integral-closure}
设 $R_i\to S_i$（$i = 1, \ldots, n$）为环映射。记 $R_i$ 在 $S_i$ 中的
整闭包为 $S'_i$。进一步令 $R$ 与 $S$ 分别为所有 $R_i$ 与所有 $S_i$
的乘积。则 $R$ 在 $S$ 中的整闭包是各 $S'_i$ 的乘积。特别地，
$R \to S$ 整闭，当且仅当每个 $R_i \to S_i$ 都整闭。
\end{lemma}

\begin{proof}
由引理 \ref{algebra-lemma-finite-product-integral} 立即得到。
\end{proof}

\begin{lemma}
\label{algebra-lemma-integral-closure-localize}
整闭包与局部化交换：若 $A \to B$ 为环映射，且 $S \subset A$ 为乘法子集，
则 $S^{-1}A$ 在 $S^{-1}B$ 中的整闭包是 $S^{-1}B'$，其中
$B' \subset B$ 是 $A$ 在 $B$ 中的整闭包。
\end{lemma}

\begin{proof}
由于局部化是正合的，有 $S^{-1}B' \subset S^{-1}B$。设 $x \in B'$ 且
$f \in S$。于是对某些 $a_i \in A$，在 $B$ 中有
$$
x^d + \sum_{i = 1, \ldots, d} a_i x^{d - i} = 0
$$
因此在 $S^{-1}B$ 中也有
$$
(x/f)^d + \sum\nolimits_{i = 1, \ldots, d} a_i/f^i (x/f)^{d - i} = 0
$$
由此可知 $S^{-1}B'$ 包含于 $S^{-1}A$ 在 $S^{-1}B$ 中的整闭包。
反之，设 $x/f \in S^{-1}B$ 整于 $S^{-1}A$。于是对某些
$a_i \in A$ 与 $f_i \in S$，在 $S^{-1}B$ 中有
$$
(x/f)^d + \sum\nolimits_{i = 1, \ldots, d} (a_i/f_i) (x/f)^{d - i} = 0
$$
这意味着，对某个适当的 $f' \in S$，有
$$
(f'f_1 \ldots f_d x)^d +
\sum\nolimits_{i = 1, \ldots, d}
f^i(f')^if_1^i \ldots f_i^{i - 1} \ldots f_d^i a_i
(f'f_1 \ldots f_dx)^{d - i} = 0
$$
因此 $f'f_1\ldots f_dx \in B'$，从而按所需
$x/f \in S^{-1}B'$。
\end{proof}

\begin{lemma}
\label{algebra-lemma-integral-closure-stalks}
\begin{slogan}
环上代数的元素整于该环，当且仅当它在该环的每个素理想处都局部整。
\end{slogan}
设 $\varphi : R \to S$ 为环映射，$x \in S$。以下条件等价：
\begin{enumerate}
\item $x$ 整于 $R$；
\item 对每个素理想 $\mathfrak p \subset R$，元素
$x \in S_{\mathfrak p}$ 整于 $R_{\mathfrak p}$。
\end{enumerate}
\end{lemma}

\begin{proof}
(1) 蕴含 (2) 显然。设 (2) 成立。考虑由 $\varphi(R)$ 与 $x$ 生成的
$R$-代数 $S' \subset S$。令 $\mathfrak p$ 为 $R$ 的一个素理想。则在
$S_{\mathfrak p}$ 中，对某些 $a_i \in R_{\mathfrak p}$ 有
$$
x^d + \sum_{i = 1, \ldots, d} \varphi(a_i) x^{d - i} = 0
$$
考察其中出现的分母可知，存在 $f \in R$、$f \not \in \mathfrak p$，使
$a_i \in R_f$ 且在 $S_f$ 中有
$$
x^d + \sum_{i = 1, \ldots, d} \varphi(a_i) x^{d - i} = 0
$$
这说明 $S'_f$ 在 $R_f$ 上有限。由于 $\mathfrak p$ 任意，且
$\Spec(R)$ 拟紧（引理 \ref{algebra-lemma-quasi-compact}），可找到有限多个
$f_1, \ldots, f_n \in R$，它们生成 $R$ 的单位理想，并使每个
$S'_{f_i}$ 在 $R_{f_i}$ 上有限。因此由引理 \ref{algebra-lemma-cover}，
$S'$ 在 $R$ 上有限。于是由引理 \ref{algebra-lemma-characterize-integral}，
$x$ 整于 $R$。
\end{proof}

\begin{lemma}
\label{algebra-lemma-base-change-integral}
\begin{slogan}
整性与有限性在基变换下保持
\end{slogan}
设 $R \to S$ 与 $R \to R'$ 为环映射，令 $S' = R' \otimes_R S$。
\begin{enumerate}
\item 若 $R \to S$ 整，则 $R' \to S'$ 整。
\item 若 $R \to S$ 有限，则 $R' \to S'$ 有限。
\end{enumerate}
\end{lemma}

\begin{proof}
证明 (1)。令 $s_i \in S$ 为生成 $S$ 的一组元素。这些元素各自都满足
一个在 $R$ 上的首一多项式方程 $P_i$，其系数属于 $R$。因此元素
$1 \otimes s_i \in S'$ 在 $R'$ 上生成 $S'$，并满足相应的
$R'$ 上的多项式 $P_i'$。由于这些元素在 $R'$ 上生成 $S'$，可知 $S'$ 整于 $R'$。
证明 (2) 略。
\end{proof}

\begin{lemma}
\label{algebra-lemma-integral-local}
设 $R \to S$ 为环映射，且 $f_1, \ldots, f_n \in R$ 生成单位理想。
\begin{enumerate}
\item 若每个 $R_{f_i} \to S_{f_i}$ 都整，则 $R \to S$ 整。
\item 若每个 $R_{f_i} \to S_{f_i}$ 都有限，则 $R \to S$ 有限。
\end{enumerate}
\end{lemma}

\begin{proof}
证明 (1)。

设 $s \in S$。考虑 R[x] 中满足 $P(s) = 0$ 的多项式 $P$ 所成的理想
$I \subset R[x]$。令 $J \subset R$ 表示 $I$ 中元素的首项系数所成的
理想（原文特意标注为 (!)）。由假设并消去分母可知，对所有 $i$，某些
$n_i \geq 0$，所有 $f_i^{n_i} \in J$。因此 $J$ 包含 $1$，从而 $s$
整于 $R$。证明 (2) 略。
\end{proof}

\begin{lemma}
\label{algebra-lemma-integral-permanence}
设 $A \to B \to C$ 为环映射。
\begin{enumerate}
\item 若 $A \to C$ 是整的，则 $B \to C$ 也是整的。
\item 若 $A \to C$ 是有限的，则 $B \to C$ 也是有限的。
\end{enumerate}
\end{lemma}

\begin{proof}
略。
\end{proof}

\begin{lemma}
\label{algebra-lemma-integral-closure-transitive}
设 $A \to B \to C$ 为环映射。令 $B'$ 为 $A$ 在 $B$ 中的整闭包，令
$C'$ 为 $B'$ 在 $C$ 中的整闭包。则 $C'$ 是 $A$ 在 $C$ 中的整闭包。
\end{lemma}

\begin{proof}
略。
\end{proof}

\begin{lemma}
\label{algebra-lemma-integral-overring-surjective}
设 $R \to S$ 为整环扩张且 $R \subset S$。则
$\varphi : \Spec(S) \to \Spec(R)$ 是满射。
\end{lemma}

\begin{proof}
令 $\mathfrak p \subset R$ 为素理想。由引理
\ref{algebra-lemma-in-image}，须证明 $\mathfrak pS_{\mathfrak p} \not = S_{\mathfrak p}$。
局部化 $R_{\mathfrak p} \to S_{\mathfrak p}$ 是单射（因为局部化正合），
并且由引理 \ref{algebra-lemma-integral-closure-localize} 或
\ref{algebra-lemma-base-change-integral} 是整的。因此可用
$R_{\mathfrak p}$、$S_{\mathfrak p}$ 分别代替 $R$、$S$，并设 $R$ 是
带极大理想 $\mathfrak m$ 的局部环；只须证明 $\mathfrak mS \not = S$。
假设 $1 = \sum f_i s_i$，其中 $f_i \in \mathfrak m$、$s_i \in S$，以导出
矛盾。取 $R \subset S' \subset S$，使 $R \to S'$ 有限且
$s_i \in S'$；见引理 \ref{algebra-lemma-characterize-integral}。等式
$1 = \sum f_i s_i$ 意味着有限 $R$-模 $S'$ 满足 $S' = \mathfrak m S'$。
由中山引理 \ref{algebra-lemma-NAK}，得到 $S' = 0$，矛盾。
\end{proof}

\begin{lemma}
\label{algebra-lemma-integral-under-field}
设 $R$ 为环，$K$ 为域。若 $R \subset K$ 且 $K$ 整于 $R$，则 $R$ 是域，
且 $K$ 是代数扩张。若 $R \subset K$ 且 $K$ 在 $R$ 上有限，则 $R$ 是域，
且 $K$ 是有限代数扩张。
\end{lemma}

\begin{proof}
假设 $R \subset K$ 是整扩张。由引理
\ref{algebra-lemma-integral-overring-surjective}，$\Spec(R)$ 只有 $1$ 个点。
显然 $R$ 是整环，故 $R = R_{(0)}$ 是域（引理
\ref{algebra-lemma-minimal-prime-reduced-ring}）。其余断言立即得到。
\end{proof}

\begin{lemma}
\label{algebra-lemma-integral-over-field}
设 $k$ 为域，$S$ 为 $k$ 上的 $k$-代数。
\begin{enumerate}
\item 若 $S$ 是整环且在 $k$ 上有限维，则 $S$ 是域。
\item 若 $S$ 整于 $k$ 且为整环，则 $S$ 是域。
\item 若 $S$ 整于 $k$，则 $S$ 的每个素理想都是极大理想（更多结论见
引理 \ref{algebra-lemma-ring-with-only-minimal-primes}）。
\end{enumerate}
\end{lemma}

\begin{proof}
关于素理想的断言来自“integral $+$ domain $\Rightarrow$ field”这一断言。设 $S$ 整于 $k$
且 $S$ 为整环。取 $s \in S$。由引理
\ref{algebra-lemma-characterize-integral}，可找到包含 $s$ 的有限维
$k$-子代数 $k \subset S' \subset S$。若证明第一条断言，则可知 $S$ 是域。
现在假设 $S$ 在 $k$ 上有限维且为整环。取 $s\in S$。由于 $S$ 是整环，
乘法映射 $s : S \to S$ 由维数相等可知是满射。因此存在 $s_1 \in S$，
使 $ss_1 = 1$。所以 $S$ 是域。
\end{proof}

\begin{lemma}
\label{algebra-lemma-integral-no-inclusion}
设 $R \to S$ 为整环映射。令
$\mathfrak q, \mathfrak q' \in \Spec(S)$ 为不同的素理想，且它们在
$\Spec(R)$ 中的像相同。则既不存在 $\mathfrak q \subset \mathfrak q'$，
也不存在 $\mathfrak q' \subset \mathfrak q$。
\end{lemma}

\begin{proof}
令 $\mathfrak p \subset R$ 为像。由注记
\ref{algebra-remark-fundamental-diagram}，素理想 $\mathfrak q, \mathfrak q'$
对应于 $S \otimes_R \kappa(\mathfrak p)$ 中的理想。因此由引理
\ref{algebra-lemma-integral-over-field} 得到结论。
\end{proof}

\begin{lemma}
\label{algebra-lemma-finite-finite-fibres}
设 $R \to S$ 为有限映射。则
$\Spec(S) \to \Spec(R)$ 的纤维是有限的。
\end{lemma}

\begin{proof}
由注记 \ref{algebra-remark-fundamental-diagram} 的讨论，纤维是环
$S \otimes_R \kappa(\mathfrak p)$ 的谱。由于 $R \to S$ 有限，这些纤维环
在 $\kappa(\mathfrak p)$ 上有限，因而由引理
\ref{algebra-lemma-Noetherian-permanence} 是诺特的。由引理
\ref{algebra-lemma-integral-no-inclusion}，$S \otimes_R \kappa(\mathfrak p)$
的每个素理想都是极小素理想。因此由引理
\ref{algebra-lemma-Noetherian-irreducible-components}，这样的素理想至多有有限多个。
\end{proof}

\begin{lemma}
\label{algebra-lemma-integral-going-up}
设 $R \to S$ 为环映射，且 $S$ 整于 $R$。令
$\mathfrak p \subset \mathfrak p' \subset R$ 为素理想。设
$\mathfrak q$ 是 $S$ 中映到 $\mathfrak p$ 的素理想。则存在素理想
$\mathfrak q'$，满足 $\mathfrak q \subset \mathfrak q'$，并且映到
$\mathfrak p'$。
\end{lemma}

\begin{proof}
可以用 $R$ 替换为 $R/\mathfrak p$，用 $S$ 替换为 $S/\mathfrak q$。
于是归约为整环扩张 $R \subset S$ 以及素理想
$\mathfrak p' \subset R$ 的情形。由引理
\ref{algebra-lemma-integral-overring-surjective} 即得。
\end{proof}

\noindent
上述引理所表达的性质称为环映射 $R \to S$ 的“上升性质”（going up
property），见定义 \ref{algebra-definition-going-up-down}。

\begin{lemma}
\label{algebra-lemma-finite-finitely-presented-extension}
设 $R \to S$ 为有限且有限表示的环映射，$M$ 为 $S$-模。则 $M$ 作为
$R$-模有限表示，当且仅当 $M$ 作为 $S$-模有限表示。
\end{lemma}

\begin{proof}
其中一个方向由引理 \ref{algebra-lemma-finitely-presented-over-subring} 得出。
反之，设 $M$ 作为 $S$-模有限表示。取表示
$$
S^{\oplus m} \longrightarrow
S^{\oplus n} \longrightarrow
M \longrightarrow 0
$$
由于 $S$ 作为 $R$-模有限，$S^{\oplus n} \to M$ 的核是有限
$R$-模。因此由引理 \ref{algebra-lemma-extension}，只须证明 $S$ 作为
$R$-模是有限表示的。

\medskip\noindent
取 $y_1, \ldots, y_n \in S$，使 $y_1, \ldots, y_n$ 作为 $R$-模生成 $S$。由引理
\ref{algebra-lemma-characterize-integral-element}，每个 $y_i$ 都整于 $R$。
选取首一多项式 $P_i(x) \in R[x]$，使 $P_i(y_i) = 0$。考虑环
$$
S' = R[x_1, \ldots, x_n]/(P_1(x_1), \ldots, P_n(x_n))
$$
由引理 \ref{algebra-lemma-compose-finite-type}，$S$ 作为 $S'$-代数是有限表示的。
由于 $S' \to S$ 是满射，由引理
\ref{algebra-lemma-finite-presentation-independent}，核
$J = \Ker(S' \to S)$ 作为理想是有限生成的。因此 $J$ 是有限
$S'$-模（直接由定义得到）。于是 $S = \Coker(J \to S')$ 作为
$S'$-模是有限表示的，这是引理 \ref{algebra-lemma-extension} 的结论。
因此，如第一段所述，只须证明 $S'$ 作为 $R$-模是有限表示的。事实上，
$S'$ 是自由 $R$-模，其基为满足 $0 \leq e_i < \deg(P_i)$ 的单项式
$x_1^{e_1} \ldots x_n^{e_n}$。具体地，将 $R \to S'$ 写成复合
$$
R \to R[x_1]/(P_1(x_1)) \to
R[x_1, x_2]/(P_1(x_1), P_2(x_2)) \to
\ldots \to S'
$$
这表明，此序列中的第 $i$ 个环作为第 $(i - 1)$ 个环的模是自由的，其基为
$1, x_i, \ldots, x_i^{\deg(P_i) - 1}$。由归纳法容易得到结论。细节略。
\end{proof}

\begin{lemma}
\label{algebra-lemma-silly-normal}
设 $R$ 为环，$x, y \in R$ 为非零因子。令
$R[x/y] \subset R_{xy}$ 为由 $x/y$ 生成的 $R$-子代数，并类似地定义
$R[y/x]$ 与 $R[x/y, y/x]$。若 $R$ 在 $R_x$ 或 $R_y$ 中整闭，则序列
$$
0 \to R \xrightarrow{(-1, 1)} R[x/y] \oplus R[y/x] \xrightarrow{(1, 1)}
R[x/y, y/x] \to 0
$$
是 $R$-模的短正合序列。
\end{lemma}

\begin{proof}
由于 $x/y \cdot y/x = 1$，映射
$R[x/y] \oplus R[y/x] \to R[x/y, y/x]$ 是满射。令
$\alpha \in R[x/y] \cap R[y/x]$。为证明中间处正合，须证明
$\alpha \in R$。由假设，可写成
$$
\alpha = a_0 + a_1 x/y + \ldots + a_n (x/y)^n =
b_0 + b_1 y/x + \ldots + b_m(y/x)^m
$$
其中 $n, m \geq 0$ 且 $a_i, b_j \in R$。取 $N > \max(n, m)$。
考虑 $R_{xy}$ 中由下列元素生成的有限 $R$-子模 $M$：
$$
(x/y)^N, (x/y)^{N - 1}, \ldots, x/y, 1, y/x, \ldots,
(y/x)^{N - 1}, (y/x)^N
$$
断言 $\alpha M \subset M$。事实上，对 $0 \leq i \leq N$，显然有
$(x/y)^i (b_0 + b_1 y/x + \ldots + b_m(y/x)^m) \in M$；并且
$(y/x)^i (a_0 + a_1 x/y + \ldots + a_n(x/y)^n) \in M$；对 $0 \leq i \leq N$。因此由引理
\ref{algebra-lemma-characterize-integral-element}，$\alpha$ 整于 $R$。注意
$\alpha \in R_x$，所以若 $R$ 在 $R_x$ 中整闭，则 $\alpha \in R$，结论得证。
\end{proof}

\section{正规环}
\label{algebra-section-normal-rings}

\noindent
我们先介绍正规整环的概念，然后介绍（非常一般的）正规环概念。

\begin{definition}
\label{algebra-definition-domain-normal}
若整环 $R$ 在其分式域中整闭，则称其为{\it 正规}。
\end{definition}

\begin{lemma}
\label{algebra-lemma-integral-closure-in-normal}
设 $R \to S$ 为环映射。若 $S$ 是正规整环，则 $R$
在 $S$ 中的整闭包是正规整环。
\end{lemma}

\begin{proof}
略。
\end{proof}

\noindent
下述概念在研究正规性时偶尔有用。

\begin{definition}
\label{algebra-definition-almost-integral}
设 $R$ 为整环。
\begin{enumerate}
\item 分式域中的元素 $g$ 称为{\it 几乎整于 $R$}，若存在 $r \in R$，
$r\not = 0$，使得对所有 $n \geq 0$ 都有 $rg^n \in R$。
\item 若 $R$ 的分式域中的每个几乎整元素都属于 $R$，则称整环 $R$
为{\it 完全正规}。
\end{enumerate}
\end{definition}

\noindent
下述引理表明，诺特整环是正规整环，当且仅当它完全正规。

\begin{lemma}
\label{algebra-lemma-almost-integral}
设 $R$ 为分式域 $K$ 中的整环。若 $u, v \in K$ 几乎整于 $R$，
则 $u + v$ 与 $uv$ 也几乎整于 $R$。任何整于 $R$ 的元素 $g \in K$
都几乎整于 $R$。若 $R$ 是诺特环，则反之亦然。
\end{lemma}

\begin{proof}
若对所有 $n \geq 0$ 都有 $ru^n \in R$，且
$v^nr' \in R$，则对所有 $n \geq 0$，
$(uv)^nrr'$ 与 $(u + v)^nrr'$ 都属于 $R$，对所有 $n \geq 0$。于是得到第一个断言。
设 $g \in K$ 整于 $R$。此时存在 $d > 0$，使环 $R[g]$ 作为
$R$-模由 $1, g, \ldots, g^d$ 生成。令 $r \in R$ 为
$1, g, \ldots, g^d \in K$ 的公共分母。于是
$rR[g] \subset R$，所以 $g$ 几乎整于 $R$。

\medskip\noindent
设 $R$ 是诺特环，且 $g \in K$ 几乎整于 $R$。取定义中的
$r \in R$，满足 $r\not = 0$。则作为 $R$-模，
$R[g] \subset \frac{1}{r}R$。由于 $R$ 是诺特环，这意味着
$R[g]$ 在 $R$ 上有限。因此 $g$ 整于 $R$，见引理
\ref{algebra-lemma-finite-is-integral}。
\end{proof}

\begin{lemma}
\label{algebra-lemma-localize-normal-domain}
正规整环的任意局部化仍是正规整环。
\end{lemma}

\begin{proof}
设 $R$ 为正规整环，$S \subset R$ 为乘法子集。设 $g$ 是
$R$ 的分式域中的元素，且整于 $S^{-1}R$。令
$P = x^d + \sum_{j < d} a_j x^j$ 为满足 $P(g) = 0$ 的多项式，
其中 $a_i \in S^{-1}R$。选取 $s \in S$，使得对所有 $i$ 都有
$sa_i \in R$。于是 $sg$ 满足首一多项式
$x^d + \sum_{j < d} s^{d-j}a_j x^j$，其系数
$s^{d-j}a_j$ 属于 $R$。由于 $R$ 正规，故 $sg \in R$。
于是 $g \in S^{-1}R$。
\end{proof}

\begin{lemma}
\label{algebra-lemma-PID-normal}
主理想整环是正规环。
\end{lemma}

\begin{proof}
设 $R$ 为主理想整环。设 $g = a/b$ 是 $R$ 的分式域中
整于 $R$ 的元素。由于 $R$ 是主理想整环，我们可以约去 $a$ 与
$b$ 的公因子，并设 $(a, b) = R$。此时，$r_i \in R$ 的任意方程
$(a/b)^n + r_{n-1} (a/b)^{n-1} + \ldots + r_0 = 0$
都会推出 $a^n \in (b)$。除非 $b$ 是 $R$ 中的单位，否则这与
$(a, b) = R$ 矛盾。
\end{proof}

\begin{lemma}
\label{algebra-lemma-prepare-polynomial-ring-normal}
设 $R$ 为分式域为 $K$ 的整环。设
$f = \sum \alpha_i x^i$ 是 $K[x]$ 中的元素。
\begin{enumerate}
\item 若 $f$ 整于 $R[x]$，则所有的 $\alpha_i$ 都整于 $R$；并且
\item 若 $f$ 几乎整于 $R[x]$，则所有的 $\alpha_i$ 都几乎整于 $R$。
\end{enumerate}
\end{lemma}

\begin{proof}
我们先证明第二个断言。写成
$f = \alpha_0 + \alpha_1 x + \ldots + \alpha_r x^r$，其中
$\alpha_r \not = 0$。根据假设，存在
$h = b_0 + b_1 x + \ldots + b_s x^s \in R[x]$，$b_s \not = 0$，
使得对所有 $n \geq 0$ 都有 $f^n h \in R[x]$。这意味着对所有
$n \geq 0$ 都有 $b_s \alpha_r^n \in R$。因此 $\alpha_r$ 几乎整于
$R$。由于几乎整元素的集合构成一个子环（引理
\ref{algebra-lemma-almost-integral}），我们得到
$f - \alpha_r x^r = \alpha_0 + \alpha_1 x + \ldots + \alpha_{r - 1} x^{r - 1}$
几乎整于 $R[x]$。对 $r$ 归纳即得。

\medskip\noindent
为证明第一个断言，我们使用绝对诺特化约。具体地，写
$\alpha_i = a_i / b_i$，并令
$P(t) = t^d + \sum_{j < d} f_j t^j$ 为系数 $f_j \in R[x]$ 且满足
$P(f) = 0$ 的多项式。令 $f_j = \sum f_{ji}x^i$。考虑由
$R$ 中有限个元素 $a_i, b_i, f_{ji}$ 生成的子环
$R_0 \subset R$。它是整环；令 $K_0$ 为其分式域。由于 $R_0$ 是有限型
$\mathbf{Z}$-代数，它是诺特的，见引理
\ref{algebra-lemma-obvious-Noetherian}。此时 $f \in K_0[x]$ 整于 $R_0[x]$ 仍然成立，
因为 $R$ 中关于元素 $a_i, b_i, f_{ji}$ 的所有恒等式在 $R_0$ 中也成立。
由引理 \ref{algebra-lemma-almost-integral}，元素 $f$ 几乎整于 $R_0[x]$。
由该引理的第二个断言，元素 $\alpha_i$ 几乎整于 $R_0$。又因为 $R_0$ 是诺特的，
它们整于 $R_0$，见引理 \ref{algebra-lemma-almost-integral}。当然，它们于是也整于 $R$。
\end{proof}

\begin{lemma}
\label{algebra-lemma-polynomial-domain-normal}
设 $R$ 为正规整环。则 $R[x]$ 是正规整环。
\end{lemma}

\begin{proof}
当 $R$ 是域 $K$ 时结论成立，因为 $K[x]$ 是欧几里得整环，因而是主理想整环，
再由引理 \ref{algebra-lemma-PID-normal} 得到它是正规整环。设 $g$ 是 $R[x]$ 的分式域中
整于 $R[x]$ 的元素。由于 $g$ 整于 $K[x]$（其中 $K$ 是 $R$ 的分式域），我们可以写成
$g = \alpha_d x^d + \alpha_{d-1}x^{d-1} + \ldots + \alpha_0$，其中 $\alpha_i \in K$。
由引理 \ref{algebra-lemma-prepare-polynomial-ring-normal}，元素 $\alpha_i$ 整于 $R$，
因而属于 $R$。
\end{proof}

\begin{lemma}
\label{algebra-lemma-power-series-over-Noetherian-normal-domain}
设 $R$ 为诺特正规整环。则 $R[[x]]$ 是诺特正规整环。
\end{lemma}

\begin{proof}
由引理 \ref{algebra-lemma-Noetherian-power-series}，幂级数环是诺特的。设
$f, g \in R[[x]]$ 为非零元素，且 $w = f/g$ 整于 $R[[x]]$。令 $K$ 为 $R$ 的分式域。
由于劳朗级数环
$K((x)) = K[[x]][1/x]$ 是域，我们可以写成
$w = a_n x^n + a_{n + 1} x^{n + 1} + \ldots$，其中 $n \in \mathbf{Z}$，
$a_i \in K$，且 $a_n \not = 0$。由引理 \ref{algebra-lemma-almost-integral}，存在
$R[[x]]$ 中的非零元素 $h = b_m x^m + b_{m + 1} x^{m + 1} + \ldots$，其中
$b_m \not = 0$，使得对所有 $e \geq 1$ 都有 $w^e h \in R[[x]]$。我们得到
$n \geq 0$，且对所有 $e \geq 1$ 都有 $b_m a_n^e \in R$。
由于 $R$ 是诺特环，由同一引理可知 $a_n \in R$。现在，若
$a_n, a_{n + 1}, \ldots, a_{N - 1} \in R$，则可对
$w - a_n x^n - \ldots - a_{N - 1} x^{N - 1} = a_N x^N + \ldots$
应用同样的论证。如此可知所有 $a_i \in R$，引理得证。
\end{proof}

\begin{lemma}
\label{algebra-lemma-normality-is-local}
设 $R$ 为整环。以下命题等价：
\begin{enumerate}
\item 整环 $R$ 是正规整环；
\item 对每个素理想 $\mathfrak p \subset R$，局部环
$R_{\mathfrak p}$ 是正规整环；并且
\item 对每个极大理想 $\mathfrak m$，环 $R_{\mathfrak m}$
是正规整环。
\end{enumerate}
\end{lemma}

\begin{proof}
由引理 \ref{algebra-lemma-localize-normal-domain}，得到 (1) $\Rightarrow$ (2)。
(2) $\Rightarrow$ (3) 显然。(3) $\Rightarrow$ (1) 则由以下事实得到：对于任意整环 $R$，
有
$$
R = \bigcap\nolimits_{\mathfrak m} R_{\mathfrak m}
$$
（在 $R$ 的分式域内）。事实上，若 $g$ 属于右端，则理想
$I = \{x \in R \mid xg \in R\}$ 不包含于任何极大理想 $\mathfrak m$ 中，
从而 $I = R$。
\end{proof}

\noindent
引理 \ref{algebra-lemma-normality-is-local} 表明，下面的定义与定义
\ref{algebra-definition-domain-normal} 相容。（这是 EGA 中的定义——见
\cite[IV, 5.13.5 and 0, 4.1.4]{EGA}。）

\begin{definition}
\label{algebra-definition-ring-normal}
若对每个素理想 $\mathfrak p \subset R$，局部化 $R_{\mathfrak p}$
是正规整环（见定义 \ref{algebra-definition-domain-normal}），则称环 $R$ 为{\it 正规}。
\end{definition}

\noindent
注意，正规环是约化环，因为 $R$ 是其在所有素理想处局部化所得环的乘积的子环
（例如见引理 \ref{algebra-lemma-characterize-zero-local}）。

\begin{lemma}
\label{algebra-lemma-normal-ring-integrally-closed}
正规环在其全分式环中整闭。
\end{lemma}

\begin{proof}
设 $R$ 为正规环。设 $x \in Q(R)$ 是 $R$ 的全分式环中整于 $R$ 的元素。
令 $I = \{f \in R, fx \in R\}$。设 $\mathfrak p \subset R$ 为素理想。
由于 $R \to R_{\mathfrak p}$ 是平坦的，有
$R_{\mathfrak p} \subset Q(R) \otimes_R R_{\mathfrak p}$。由于
$R_{\mathfrak p}$ 是正规整环，故 $x \otimes 1$ 是 $R_{\mathfrak p}$ 中的元素。
因此可找到 $a, f \in R$，$f \not \in \mathfrak p$，使得
$x \otimes 1 = a \otimes 1/f$。这意味着 $fx - a$ 在
$Q(R) \otimes_R R_{\mathfrak p} = Q(R)_{\mathfrak p}$ 中映为零，而这又意味着
存在 $f' \in R$，$f' \not \in \mathfrak p$，使得在 $R$ 中有
$f'fx = f'a$。换言之，$ff' \in I$。因此 $I$ 是一个不包含于 $R$ 的任何素理想
的理想，即 $I = R$，且 $x \in R$。
\end{proof}

\begin{lemma}
\label{algebra-lemma-localization-normal-ring}
正规环的局部化仍是正规环。
\end{lemma}

\begin{proof}
略。
\end{proof}

\begin{lemma}
\label{algebra-lemma-polynomial-ring-normal}
设 $R$ 为正规环。则 $R[x]$ 是正规环。
\end{lemma}

\begin{proof}
设 $\mathfrak q$ 是 $R[x]$ 中的素理想。令
$\mathfrak p = R \cap \mathfrak q$。由引理
\ref{algebra-lemma-polynomial-domain-normal}，可知 $R_{\mathfrak p}[x]$ 是正规整环。
于是由引理 \ref{algebra-lemma-localize-normal-domain}，$(R[x])_{\mathfrak q}$ 是正规整环。
\end{proof}

\begin{lemma}
\label{algebra-lemma-finite-product-normal}
正规环的有限乘积是正规环。
\end{lemma}

\begin{proof}
只须证明两个正规环（记为 $R$ 和 $S$）的乘积是正规环。由引理
\ref{algebra-lemma-disjoint-decomposition}，$R\times S$ 的素理想具有
$\mathfrak{p}\times S$ 和 $R\times \mathfrak{q}$ 的形式，其中
$\mathfrak{p}$ 和 $\mathfrak{q}$ 分别是 $R$ 和 $S$ 的素理想。局部化给出
$(R\times S)_{\mathfrak{p}\times S}=R_{\mathfrak{p}}$，它由假设是正规整环。
对 $S$ 的情形同理。
\end{proof}

\begin{lemma}
\label{algebra-lemma-characterize-reduced-ring-normal}
设 $R$ 为环。假设 $R$ 是约化的并且只有有限多个极小素理想。
则以下命题等价：
\begin{enumerate}
\item $R$ 是正规环；
\item $R$ 在其全分式环中整闭；并且
\item $R$ 是正规整环的有限乘积。
\end{enumerate}
\end{lemma}

\begin{proof}
(1) $\Rightarrow$ (2) 以及 (3) $\Rightarrow$ (1) 的蕴含一般都成立，
见引理 \ref{algebra-lemma-normal-ring-integrally-closed} 和
\ref{algebra-lemma-finite-product-normal}。

\medskip\noindent
设 $\mathfrak p_1, \ldots, \mathfrak p_n$ 是 $R$ 的极小素理想。
由引理 \ref{algebra-lemma-reduced-ring-sub-product-fields} 和
\ref{algebra-lemma-total-ring-fractions-no-embedded-points}，有
$Q(R) = R_{\mathfrak p_1} \times \ldots \times R_{\mathfrak p_n}$；
并且由引理 \ref{algebra-lemma-minimal-prime-reduced-ring}，每个因子都是域。
记 $e_i = (0, \ldots, 0, 1, 0, \ldots, 0)$ 为 $Q(R)$ 的第 $i$ 个幂等元。

\medskip\noindent
若 $R$ 在 $Q(R)$ 中整闭，则特别地包含幂等元 $e_i$，从而可知 $R$ 是
$n$ 个整环的乘积（见第 \ref{algebra-section-connected-components} 节和
\ref{algebra-section-tilde-module-sheaf} 节）。每个因子具有
$R/\mathfrak p_i$ 的形式，其分式域为 $R_{\mathfrak p_i}$。
由引理 \ref{algebra-lemma-finite-product-integral-closure}，每个映射
$R/\mathfrak p_i \to R_{\mathfrak p_i}$ 都是整闭的。
因此 $R$ 是正规整环的有限乘积。
\end{proof}

\begin{lemma}
\label{algebra-lemma-colimit-normal-ring}
设 $(R_i, \varphi_{ii'})$ 是环的有向系统
（范畴论定义 \ref{algebra-definition-directed-system}）。若每个 $R_i$ 都是正规环，
则 $R = \colim_i R_i$ 也是正规环。
\end{lemma}

\begin{proof}
设 $\mathfrak p \subset R$ 为素理想。令
$\mathfrak p_i = R_i \cap \mathfrak p$（通常的记号滥用）。于是有
$R_{\mathfrak p} = \colim_i (R_i)_{\mathfrak p_i}$。
由于每个 $(R_i)_{\mathfrak p_i}$ 都是正规整环，我们归约为证明正规整环的情形。
若 $a, b \in R$ 且 $a/b$ 满足 $P(T) \in R[T]$ 这一首一多项式，
则可找到充分大的 $i \in I$，使得 $a, b$ 分别来自 $R_i$ 上的元素
$a_i, b_i$，而 $P$ 来自首一多项式 $P_i\in R_i[T]$，并且
$P_i(a_i/b_i)=0$。由于 $R_i$ 正规，有 $a_i/b_i \in R_i$，从而
$a/b \in R$。
\end{proof}

\section{正规环上整扩张的下降}
\label{algebra-section-going-down-integral-over-normal}

\noindent
我们先稍微研究“整于理想”的元素这一概念，然后证明本节标题所指的定理。

\begin{definition}
\label{algebra-definition-integral-over-ideal}
设 $\varphi : R \to S$ 为环映射。设 $I \subset R$ 为理想。
若元素 $g \in S$ 满足下述条件，则称该元素
{\it 整于 $I$}：存在首一多项式
$P = x^d + \sum_{j < d} a_j x^j$，其系数满足
$a_j \in I^{d-j}$，并且
$P^\varphi(g) = 0$ 在 $S$ 中成立。
\end{definition}

\noindent
这一概念主要用于 $\varphi = \text{id}_R : R \to R$ 的情形。
此时，整于 $I$ 的元素所成的集合 $I'$ 称为{\it $I$ 的整闭包}。
我们将看到 $I'$ 是 $R$ 的一个理想（当然 $I \subset I'$）。

\begin{lemma}
\label{algebra-lemma-characterize-integral-ideal}
设 $\varphi : R \to S$ 为环映射。设 $I \subset R$ 为理想。
令 $A = \sum I^nt^n \subset R[t]$ 为多项式环中由
$R \oplus It \subset R[t]$ 生成的子环。元素 $s \in S$ 整于 $I$，
当且仅当元素 $st \in S[t]$ 整于 $A$。
\end{lemma}

\begin{proof}
假设 $st$ 整于 $A$。令
$P = x^d + \sum_{j < d} a_j x^j$ 为系数在 $A$ 中的首一多项式，
满足 $P^\varphi(st) = 0$。令 $a_j' \in A$ 为 $a_j$ 的
$d-j$ 次部分，换言之，$a_j' = a_j'' t^{d-j}$，其中
$a_j'' \in I^{d-j}$。根据次数的原因，仍有
$(st)^d + \sum_{j < d} \varphi(a_j'') t^{d-j} (st)^j = 0$。
因此
$s^d + \sum_{j < d} \varphi(a_j'') s^j = 0$，
从而 $s$ 整于 $I$。

\medskip\noindent
假设 $s$ 整于 $I$。设
$P = x^d + \sum_{j < d} a_j x^j$，其中 $a_j \in I^{d-j}$。
则立即得到多项式
$Q = x^d + \sum_{j < d} (a_j t^{d-j}) x^j$，其系数在 $A$ 中，
这证明了 $st$ 整于 $A$。
\end{proof}

\begin{lemma}
\label{algebra-lemma-integral-over-ideal-is-submodule}
设 $\varphi : R \to S$ 为环映射。设 $I \subset R$ 为理想。
$S$ 中整于 $I$ 的元素所成的集合是 $S$ 的一个 $R$-子模。
此外，若 $s \in S$ 整于 $R$，而 $s'$ 整于 $I$，则
$ss'$ 整于 $I$。
\end{lemma}

\begin{proof}
我们将使用引理 \ref{algebra-lemma-integral-closure-is-ring} 而不再另行说明。
由引理 \ref{algebra-lemma-characterize-integral-ideal} 的刻画，集合对加法封闭是显然的，
下文沿用该引理的记号。任何整于 $R$ 的元素 $s \in S$，都对应于 $S[t]$ 中的
次数 $0$ 元素 $s$，并且该元素整于 $A$（因为 $R \subset A$）。
因此，$S[t]$ 上乘以 $s$ 的运算保持整于 $A$ 的性质，
\end{proof}

\begin{lemma}
\label{algebra-lemma-integral-integral-over-ideal}
设 $\varphi : R \to S$ 为整环映射。设 $I \subset R$ 为理想。
则 $IS$ 中的每个元素都整于 $I$。
\end{lemma}

\begin{proof}
由引理 \ref{algebra-lemma-integral-over-ideal-is-submodule} 立即得到。
\end{proof}

\begin{lemma}
\label{algebra-lemma-polynomials-divide}
设 $K$ 为域。设 $n, m \in \mathbf{N}$，且
$a_0, \ldots, a_{n - 1}, b_0, \ldots, b_{m - 1} \in K$。
若多项式 $x^n + a_{n - 1}x^{n - 1} + \ldots + a_0$
在 $K[x]$ 中整除多项式 $x^m + b_{m - 1} x^{m - 1} + \ldots + b_0$，
则
\begin{enumerate}
\item 对于包含元素 $b_0, \ldots, b_{m - 1}$ 的 $K$ 的任意子环
$R_0$，$a_0, \ldots, a_{n - 1}$ 都整于该子环；并且
\item 对于包含元素
$a_0, \ldots, a_{n - 1}, b_0, \ldots, b_{m - 1}$ 的任意子环
$R \subset K$，每个 $a_i$ 都属于下列根理想：
$$
\sqrt{(b_0, \ldots, b_{m-1})R}
$$
\end{enumerate}
\end{lemma}

\begin{proof}
设 $L/K$ 为域扩张，使得可以写成
$$
x^m + b_{m - 1} x^{m - 1} + \ldots + b_0 =
\prod_{i = 1}^m (x - \beta_i)
$$
其中 $\beta_i \in L$。见域一章的
\ref{fields-section-splitting-fieds} 节。
每个 $\beta_i$ 都整于 $R_0$。由于每个 $a_i$ 都是
$\beta_1, \ldots, \beta_m$ 的齐次多项式，我们得到 $a_i$ 也具有同样性质
（使用引理 \ref{algebra-lemma-integral-closure-is-ring}）。这证明了 (1)。

\medskip\noindent
设 $R$ 如 (2) 中所述。选取 $c_0, \ldots, c_{m - n - 1} \in K$，使得
$$
\begin{matrix}
x^m + b_{m - 1} x^{m - 1} + \ldots + b_0 =  \\
(x^n + a_{n - 1}x^{n - 1} + \ldots + a_0)
(x^{m - n} + c_{m - n - 1}x^{m - n - 1}+ \ldots + c_0).
\end{matrix}
$$
该等式推出
\begin{align*}
c_{m - n - 1} & = b_{m - 1} - a_{n - 1}, \\
c_{m - n - 2} & = b_{m - 2} - a_{n - 2} - a_{n - 1}c_{m - n - 1}, \\
\ldots
\end{align*}
因此对所有 $j$ 都有 $c_j \in R$。除去根式
$\sqrt{(b_0, \ldots, b_{m - 1})}$ 后得到约化环
$\overline{R}$。我们必须证明像
$\overline{a}_i \in \overline{R}$ 都为零。在
$\overline{R}[x]$ 中有关系
$$
\begin{matrix}
x^m = x^m + \overline{b}_{m - 1} x^{m - 1} + \ldots + \overline{b}_0 = \\
(x^n + \overline{a}_{n - 1}x^{n - 1} + \ldots + \overline{a}_0)
(x^{m - n} + \overline{c}_{m - n - 1}x^{m - n - 1}+ \ldots + \overline{c}_0).
\end{matrix}
$$
容易看出这意味着对所有 $i$ 都有 $\overline{a}_i = 0$。事实上，由引理
\ref{algebra-lemma-minimal-prime-reduced-ring}，$\overline{R}$ 在极小素理想
$\mathfrak{p}$ 处的局部化是域，而 $\overline{R}_{\mathfrak p}[x]$ 是唯一分解整环。
因此
$f = x^n + \sum \overline{a}_i x^i$
与 $x^n$ 相伴；又由于 $f$ 是首一的，在
$\overline{R}_{\mathfrak p}[x]$ 中有 $f = x^n$。
于是存在 $s \in \overline{R}$，$s \not\in \mathfrak p$，使得
$s(f - x^n) = 0$。因此所有 $\overline{a}_i$ 都属于
$\mathfrak p$，再由引理 \ref{algebra-lemma-reduced-ring-sub-product-fields} 得到结论。
\end{proof}

\begin{lemma}
\label{algebra-lemma-minimal-polynomial-normal-domain}
设 $R \subset S$ 为整环的包含。假设 $R$ 正规。设 $g \in S$ 整于
$R$。则 $g$ 的最小多项式的系数属于 $R$。
\end{lemma}

\begin{proof}
设 $P = x^m + b_{m-1} x^{m-1} + \ldots + b_0$
为系数在 $R$ 中且满足 $P(g) = 0$ 的多项式。设
$Q = x^n + a_{n-1}x^{n-1} + \ldots + a_0$
为 $g$ 在 $R$ 的分式域 $K$ 上的最小多项式。于是 $Q$ 在 $K[x]$ 中整除
$P$。由引理 \ref{algebra-lemma-polynomials-divide} 可知 $a_i$ 整于 $R$。
由于 $R$ 正规，这意味着它们属于 $R$。
\end{proof}

\begin{proposition}
\label{algebra-proposition-going-down-normal-integral}
设 $R \subset S$ 为整环的包含。假设 $R$ 正规且 $S$ 整于 $R$。
设 $\mathfrak p \subset \mathfrak p' \subset R$ 为素理想。设
$\mathfrak q'$ 是 $S$ 中满足 $\mathfrak p' = R \cap \mathfrak q'$ 的素理想。
则存在素理想 $\mathfrak q$，满足
$\mathfrak q \subset \mathfrak q'$，并且 $\mathfrak p = R \cap \mathfrak q$。
换言之，$R \to S$ 满足下降性质，见定义
\ref{algebra-definition-going-up-down}。
\end{proposition}

\begin{proof}
令 $\mathfrak p$、$\mathfrak p'$ 和 $\mathfrak q'$ 如命题所述。我们要证明存在
素理想 $\mathfrak q$，使得 $\mathfrak q \subset \mathfrak q'$ 且
$R \cap \mathfrak q = \mathfrak p$。这等价于寻找一个映到 $\mathfrak p$ 的
$S_{\mathfrak q'}$ 的素理想。
根据引理 \ref{algebra-lemma-in-image}，我们需要证明
$\mathfrak p S_{\mathfrak q'} \cap R
= \mathfrak p$。取 $z \in \mathfrak p S_{\mathfrak q'} \cap R$。
可写成 $z = y/g$，其中 $y \in \mathfrak pS$ 且
$g \in S$，$g \not\in \mathfrak q'$。换一种写法就是 $zg = y$。

\medskip\noindent
由引理 \ref{algebra-lemma-integral-integral-over-ideal}，存在首一多项式
$P = x^m + b_{m-1} x^{m-1} + \ldots + b_0$，
其中 $b_i \in \mathfrak p$，并且 $P(y) = 0$。

\medskip\noindent
由引理 \ref{algebra-lemma-minimal-polynomial-normal-domain}，$g$ 在 $K$ 上的
最小多项式的系数属于 $R$。将其写成
$Q = x^n + a_{n-1} x^{n-1} + \ldots
+ a_0$。注意，并非所有 $a_i$（$i = n-1, \ldots, 0$）
都属于 $\mathfrak p$，因为这会推出
$g^n = \sum_{j < n} a_j g^j \in \mathfrak pS
\subset \mathfrak p'S
\subset \mathfrak q'$，
从而产生矛盾。

\medskip\noindent
由于 $y = zg$，由上面的论证立即可知
$Q' = x^n + za_{n-1} x^{n-1} + \ldots + z^{n}a_0$
是 $y$ 的最小多项式。因此 $Q'$ 整除 $P$，由引理
\ref{algebra-lemma-polynomials-divide} 可知
$z^ja_{n - j} \in \sqrt{(b_0, \ldots, b_{m-1})}
\subset \mathfrak p$，$j =  1, \ldots, n$。
由于并非所有 $a_i$（$i = n-1, \ldots, 0$）都属于 $\mathfrak p$，
我们得到 $z \in \mathfrak p$，命题得证。
\end{proof}

\section{平坦模与平坦环映射}
\label{algebra-section-flat}

\noindent
一个常用结果是：若 $M = \colim_{i\in \mathcal{I}} M_i$
是 $R$-模的余极限，且 $N$ 是 $R$-模，则
$$
M \otimes N
=
\colim_{i\in \mathcal{I}} M_i \otimes_R N,
$$
见引理 \ref{algebra-lemma-tensor-products-commute-with-limits}。
通常将这一性质表述为{\it $\otimes$ 与余极限可交换}。
另一个常用结果是：若 $0 \to N_1 \to N_2 \to N_3 \to 0$
是正合序列，且 $M$ 是任意 $R$-模，则
$$
M \otimes_R N_1
\to
M \otimes_R N_2
\to
M \otimes_R N_3
\to
0
$$
仍然正合，见引理 \ref{algebra-lemma-tensor-product-exact}。
这两个性质都说明函子
$N \mapsto M \otimes_R N$ {\it 是右正合的}。
见范畴论的 \ref{categories-section-exact-functor} 节和
同调代数的 \ref{homology-section-functors} 节。
若 $N \mapsto N \otimes_R M$ 也是左正合的，即它是正合的，
则称 $R$-模 $M$ 为平坦模。下面给出精确定义。

\begin{definition}
\label{algebra-definition-flat}
设 $R$ 为环。
\begin{enumerate}
\item 若每当
$N_1 \to N_2 \to N_3$ 是 $R$-模的正合序列时，
序列 $M \otimes_R N_1 \to M \otimes_R N_2 \to M \otimes_R N_3$
也正合，则称 $R$-模 $M$ {\it 平坦}。
\item 若 $R$-模的复形
$N_1 \to N_2 \to N_3$ 正合，当且仅当
序列 $M \otimes_R N_1 \to M \otimes_R N_2 \to M \otimes_R N_3$
正合，则称 $R$-模 $M$ {\it 忠实平坦}。
\item 若环映射 $R \to S$ 使得 $S$ 作为 $R$-模平坦，
则称该环映射{\it 平坦}。
\item 若环映射 $R \to S$ 使得 $S$ 作为 $R$-模忠实平坦，
则称该环映射{\it 忠实平坦}。
\end{enumerate}
\end{definition}

\noindent
下面给出一个使用平坦性条件的例子。

\begin{lemma}
\label{algebra-lemma-flat-intersect-ideals}
设 $R$ 为环。设 $I, J \subset R$ 为理想。设 $M$ 为平坦的
$R$-模。则 $IM \cap JM = (I \cap J)M$。
\end{lemma}

\begin{proof}
考虑正合序列 $0 \to I \cap J \to R \to R/I \oplus R/J$。
与平坦模 $M$ 作张量，得到正合序列
$$
0 \to (I \cap J) \otimes_R M \to M \to M/IM \oplus M/JM
$$
由于 $M \to M/IM \oplus M/JM$ 的核等于
$IM \cap JM$，结论成立。
\end{proof}

\begin{lemma}
\label{algebra-lemma-colimit-flat}
设 $R$ 为环。设 $\{M_i, \varphi_{ii'}\}$ 为平坦 $R$-模所成的有向系统。
则 $\colim_i M_i$ 是平坦的 $R$-模。
\end{lemma}

\begin{proof}
这是因为 $\otimes$ 与余极限可交换，且有向余极限是正合的，见引理
\ref{algebra-lemma-directed-colimit-exact}。
\end{proof}

\begin{lemma}
\label{algebra-lemma-composition-flat}
（忠实）平坦环映射的复合仍是（忠实）平坦的。
若 $R \to R'$ 是（忠实）平坦的，且 $M'$ 是
（忠实）平坦的 $R'$-模，则 $M'$ 是（忠实）平坦的 $R$-模。
\end{lemma}

\begin{proof}
引理的第一个断言是第二个断言的特例，因此只须证明后者。设
$R \to R'$ 为平坦环映射，$M'$ 为平坦的 $R'$-模。
我们需要证明 $M'$ 是平坦的 $R$-模。设
$N_1 \to N_2 \to N_3$ 是 $R$-模的正合复形。
于是复形 $R' \otimes_R N_1 \to
R' \otimes_R N_2 \to R' \otimes_R N_3$ 正合（因为 $R'$
作为 $R$-模平坦），从而复形
$M' \otimes_{R'} \left(R' \otimes_R N_1\right)
\to M' \otimes_{R'} \left(R' \otimes_R N_2\right)
\to M' \otimes_{R'} \left(R' \otimes_R N_3\right)$
正合（因为 $M'$ 是平坦的 $R'$-模）。由于对任意 $R$-模 $N$ 都有函子性同构
$M' \otimes_{R'} \left(R' \otimes_R N\right)
\cong \left(M' \otimes_{R'} R'\right) \otimes_R N
\cong M' \otimes_R N$
（由引理 \ref{algebra-lemma-tensor-with-bimodule} 和
\ref{algebra-lemma-flip-tensor-product}），上述复形同构于复形
$M' \otimes_R N_1 \to M' \otimes_R N_2 \to M' \otimes_R N_3$，
因而后者也正合。这说明 $M'$ 是平坦的 $R$-模。沿相反方向追溯这一论证可知：
若 $R \to R'$ 忠实平坦，且 $M'$ 作为 $R'$-模忠实平坦，
则 $M'$ 作为 $R$-模也忠实平坦。
\end{proof}

\begin{lemma}
\label{algebra-lemma-flat}
设 $M$ 为 $R$-模。以下命题等价：
\begin{enumerate}
\item
\label{algebra-item-flat}
$M$ 在 $R$ 上平坦。
\item
\label{algebra-item-injective}
对于每个 $R$-模的单射 $N \subset N'$，
映射 $N \otimes_R M \to N'\otimes_R M$ 是单射。
\item
\label{algebra-item-f-ideal}
对于每个理想 $I \subset R$，映射
$I \otimes_R M \to R \otimes_R M = M$ 是单射。
\item
\label{algebra-item-ffg-ideal}
对于每个有限生成理想 $I \subset R$，
映射 $I \otimes_R M \to R \otimes_R M = M$ 是单射。
\end{enumerate}
\end{lemma}

\begin{proof}
蕴含（\ref{algebra-item-flat}）蕴含（\ref{algebra-item-injective}）
蕴含（\ref{algebra-item-f-ideal}）蕴含（\ref{algebra-item-ffg-ideal}）
都是平凡的。因此我们证明（\ref{algebra-item-ffg-ideal}）蕴含（\ref{algebra-item-flat}）。
设 $N_1 \to N_2 \to N_3$ 正合。令
$K = \Ker(N_2 \to N_3)$，$Q = \Im(N_2 \to N_3)$。
于是得到映射
$$
N_1 \otimes_R M \to
K \otimes_R M \to
N_2 \otimes_R M \to
Q \otimes_R M \to
N_3 \otimes_R M
$$
注意第一个和第三个箭头都是满射。因此只要证明第二个和第四个箭头是单射即可，
从而完成证明\footnote{更详细的论证如下：
假设第二个和第四个箭头是单射。将引理
\ref{algebra-lemma-tensor-product-exact} 应用于正合序列
$K \to N_2 \to Q \to 0$，得到序列
$K \otimes_R M \to N_2 \otimes_R M \to
Q \otimes_R M \to 0$ 正合。因此
$\Ker \left(N_2 \otimes_R M \to Q \otimes_R M\right)
= \Im \left(K \otimes_R M \to N_2 \otimes_R M\right)$。
由于
$\Im \left(K \otimes_R M \to N_2 \otimes_R M\right)
= \Im \left(N_1 \otimes_R M \to N_2 \otimes_R M\right)$
（这是因为 $N_1 \otimes_R M \to
K \otimes_R M$ 满射），并且
$\Ker \left(N_2 \otimes_R M \to Q \otimes_R M\right)
= \Ker \left(N_2 \otimes_R M \to N_3 \otimes_R M\right)$
（这是因为 $Q \otimes_R M \to
N_3 \otimes_R M$ 单射），于是得到
$\Ker \left(N_2 \otimes_R M \to N_3 \otimes_R M\right)
= \Im \left(N_1 \otimes_R M \to N_2 \otimes_R M\right)$，
这说明函子 $- \otimes_R M$ 是正合的，因而 $M$ 平坦。}。
所以只须证明 $- \otimes_R M$ 将单射的 $R$-模映射变为单射的
$R$-模映射。

\medskip\noindent
假设 $K \to N$ 是单射的 $R$-模映射，并令
$x \in \Ker(K \otimes_R M \to N \otimes_R M)$。
我们要证明 $x$ 为零。$R$-模 $K$ 是其有限
$R$-子模的并，因此 $K \otimes_R M$ 是形如
$K_i \otimes_R M$ 的 $R$-模的余极限，其中 $K_i$ 遍历 $K$ 的所有有限
$R$-子模（因为张量积与余极限可交换）。
因此，对某个 $i$，$x$ 来自元素
$x_i \in K_i \otimes_R M$。于是可假设 $K$
是有限 $R$-模。以下均作此假设。我们把单射
$K \to N$ 视为包含，从而 $K \subset N$。

\medskip\noindent
$R$-模 $N$ 是所有包含 $K$ 的有限
$R$-子模的并。因此 $N \otimes_R M$
是形如 $N_i \otimes_R M$ 的 $R$-模的余极限，其中 $N_i$ 遍历 $N$ 中所有
包含 $K$ 的有限 $R$-子模（同样因为张量积与余极限可交换）。
注意这是一个有向系统上的余极限（因为 $N$ 的两个有限子模之和仍然有限）。
所以由引理 \ref{algebra-lemma-zero-directed-limit}，元素
$x \in K \otimes_R M$ 至少在其中一个这样的 $R$-模
$N_i \otimes_R M$ 中映为零（因为 $x$ 在 $N \otimes_R M$ 中映为零）。
因此可假设 $N$ 是有限 $R$-模。

\medskip\noindent
假设 $N$ 是有限 $R$-模。写成 $N = R^{\oplus n}/L$，并令
$K = L'/L$，其中 $L \subset L' \subset R^{\oplus n}$。
对于任意 $R$-子模 $G \subset R^{\oplus n}$，都有一个典范映射
$G \otimes_R M \to M^{\oplus n}$，它由复合映射
$G \otimes_R M \to R^n \otimes_R M = M^{\oplus n}$ 得到。
只须证明 $L \otimes_R M \to M^{\oplus n}$
和 $L' \otimes_R M \to M^{\oplus n}$ 都是单射。
事实上，此时
$K \otimes_R M = L' \otimes_R M/L \otimes_R M \to M^{\oplus n}/L \otimes_R M$
也是单射\footnote{将 $L' \otimes_R M$ 和 $L \otimes_R M$
分别识别为 $M^{\oplus n}$ 的子模后，这一点显然（这是合法的，因为映射
$L \otimes_R M \to M^{\oplus n}$ 和
$L' \otimes_R M \to M^{\oplus n}$ 都是单射，并且与显然的映射
$L' \otimes_R M \to L \otimes_R M$ 交换）。}。

\medskip\noindent
所以当 $L \subset R^{\oplus n}$ 是 $R$-子模时，只须证明
$L \otimes_R M \to M^{\oplus n}$ 是单射。
我们对 $n$ 归纳证明。下面处理归纳基础情形 $n = 1$。
归纳步骤中设 $n > 1$，并令
$L' = L \cap R \oplus 0^{\oplus n - 1}$。于是 $L'' = L/L'$ 是
$R^{\oplus n - 1}$ 的子模。我们得到图表
$$
\xymatrix{
&
L' \otimes_R M \ar[r] \ar[d] &
L \otimes_R M \ar[r] \ar[d] &
L'' \otimes_R M \ar[r] \ar[d] &
0 \\
0 \ar[r] &
M \ar[r] &
M^{\oplus n} \ar[r] &
M^{\oplus n - 1} \ar[r] & 0
}
$$
由归纳假设和基础情形，左右两条竖直箭头都是单射。两行都是正合的，
因此中间的竖直箭头也是单射。

\medskip\noindent
上述归纳的基础情形是 $L \subset R$ 为理想的情形。
换言之，对 $R$ 的任意理想 $I$，我们要证明
$I \otimes_R M \to M$ 是单射。当 $I$ 有限生成时我们已经知道结论成立。
然而，$I = \bigcup I_\alpha$ 是包含于其中的有限生成理想 $I_\alpha$ 的并。
换言之，$I = \colim I_\alpha$。由于 $\otimes$ 与余极限可交换，有
$I \otimes_R M = \colim I_\alpha \otimes_R M$；又因为由假设所有
$I_\alpha \otimes_R M \to M$ 的映射都是单射，
$I \otimes_R M \to M$ 也具有同样性质。
\end{proof}

\begin{lemma}
\label{algebra-lemma-colimit-rings-flat}
设 $\{R_i, \varphi_{ii'}\}$ 为有向集 $I$ 上的环系统。
令 $R = \colim_i R_i$。
\begin{enumerate}
\item 若 $M$ 是 $R$-模，并且对所有 $i$，$M$ 作为 $R_i$-模都是平坦的，
则 $M$ 是平坦的 $R$-模。
\item 对 $i \in I$，令 $M_i$ 为平坦的 $R_i$-模；对 $i' \geq i$，
令 $f_{ii'} : M_i \to M_{i'}$ 为 $\varphi_{ii'}$-线性映射，满足
$f_{i' i''} \circ f_{i i'} = f_{i i''}$。则
$M = \colim_{i \in I} M_i$ 是平坦的 $R$-模。
\end{enumerate}
\end{lemma}

\begin{proof}
第 (1) 部分是第 (2) 部分在对所有 $i$ 取 $M_i = M$ 且
$f_{i i'} = \text{id}_M$ 时的特例。下面证明 (2)。
设 $\mathfrak a \subset R$ 为有限生成理想。由
引理 \ref{algebra-lemma-flat}，只须证明 $\mathfrak a \otimes_R M \to M$ 是单射。
可找到 $i \in I$ 以及有限生成理想 $\mathfrak a' \subset R_i$，使得
$\mathfrak a = \mathfrak a'R$。于是
$\mathfrak a = \colim_{i' \geq i} \mathfrak a'R_{i'}$。
由于 $\otimes$ 与余极限可交换，映射 $\mathfrak a \otimes_R M \to M$
是下列映射的余极限：
$$
\mathfrak a'R_{i'} \otimes_{R_{i'}} M_{i'} \longrightarrow M_{i'}
$$
由假设这些映射都是单射。由于在 $I$ 上的余极限由引理
\ref{algebra-lemma-directed-colimit-exact} 可知是正合的，结论成立。
\end{proof}

\begin{lemma}
\label{algebra-lemma-flat-base-change}
设 $M$ 在 $R$ 上（忠实）平坦，且 $R \to R'$ 为环映射。
则 $M \otimes_R R'$ 在 $R'$ 上（忠实）平坦。
\end{lemma}

\begin{proof}
对于任意 $R'$-模 $N$，都有典范同构
$N \otimes_{R'} (R'\otimes_R M)
= N \otimes_R M$。因此函子
$-\otimes_{R'}(R'\otimes_R M)$ 所需的正合性性质，
由函子 $-\otimes_R M$ 的相应正合性性质推出。
\end{proof}

\begin{lemma}
\label{algebra-lemma-flatness-descends}
设 $R \to R'$ 为忠实平坦的环映射。设 $M$ 为 $R$-模，并令
$M' = R' \otimes_R M$。则 $M$ 在 $R$ 上平坦，当且仅当
$M'$ 在 $R'$ 上平坦。
\end{lemma}

\begin{proof}
由引理 \ref{algebra-lemma-flat-base-change}，若 $M$ 平坦，则 $M'$ 平坦。
反之，假设 $M'$ 平坦。设
$N_1 \to N_2 \to N_3$ 为 $R$-模的正合序列。
我们要证明 $N_1 \otimes_R M \to N_2 \otimes_R M \to N_3 \otimes_R M$
正合。由于 $R \to R'$ 平坦，我们知道
$N_1 \otimes_R R' \to N_2 \otimes_R R' \to N_3 \otimes_R R'$ 正合。
由 $M'$ 的平坦性，
$N_1 \otimes_R R' \otimes_{R'} M'
\to N_2 \otimes_R R' \otimes_{R'} M'
\to N_3 \otimes_R R' \otimes_{R'} M'$ 正合。
可以将其写成
$N_1 \otimes_R M \otimes_R R'
\to N_2 \otimes_R M \otimes_R R'
\to N_3 \otimes_R M \otimes_R R'$。
最后，由忠实平坦性可知
$N_1 \otimes_R M \to N_2 \otimes_R M \to N_3 \otimes_R M$
正合。
\end{proof}

\begin{lemma}
\label{algebra-lemma-flatness-descends-more-general}
设 $R$ 为环。设 $S \to S'$ 为 $R$-代数的平坦映射。
设 $M$ 为 $S$-模，并令 $M' = S' \otimes_S M$。
\begin{enumerate}
\item 若 $M$ 在 $R$ 上平坦，则 $M'$ 在 $R$ 上平坦。
\item 若 $S \to S'$ 忠实平坦，则 $M$ 在 $R$ 上平坦，当且仅当
$M'$ 在 $R$ 上平坦。
\end{enumerate}
\end{lemma}

\begin{proof}
设 $N \to N'$ 为 $R$-模的单射。由 $S \to S'$ 的平坦性，有
$$
\Ker(N \otimes_R M \to N' \otimes_R M) \otimes_S S'
=
\Ker(N \otimes_R M' \to N' \otimes_R M')
$$
若 $M$ 在 $R$ 上平坦，则左端为零，由平坦性的第二个刻画（引理
\ref{algebra-lemma-flat}）可知 $M'$ 在 $R$ 上平坦。
若 $M'$ 在 $R$ 上平坦，则右端为零；若再假设 $S \to S'$ 忠实平坦，
则这推出 $\Ker(N \otimes_R M \to N' \otimes_R M)$ 为零，
从而 $M$ 在 $R$ 上平坦。
\end{proof}

\begin{lemma}
\label{algebra-lemma-flat-permanence}
设 $R \to S$ 为环映射。设 $M$ 为 $S$-模。
若 $M$ 作为 $R$-模平坦且作为 $S$-模忠实平坦，
则 $R \to S$ 平坦。
\end{lemma}

\begin{proof}
设 $N_1 \to N_2 \to N_3$ 为 $R$-模的正合序列。
由假设，$N_1 \otimes_R M \to N_2 \otimes_R M \to N_3 \otimes_R M$
正合。可将其写成
$$
N_1 \otimes_R S \otimes_S M
\to
N_2 \otimes_R S \otimes_S M
\to
N_3 \otimes_R S \otimes_S M.
$$
由 $M$ 在 $S$ 上的忠实平坦性，得到
$N_1 \otimes_R S \to N_2 \otimes_R S \to N_3 \otimes_R S$ 正合。
因此 $R \to S$ 平坦。
\end{proof}

\noindent
设 $R$ 为环。设 $M$ 为 $R$-模。设
$\sum f_i x_i = 0$ 为 $M$ 中的一个关系。
若存在整数 $m \geq 0$、元素 $y_j \in M$（$j = 1, \ldots, m$）以及元素
$a_{ij} \in R$（$i = 1, \ldots, n$，$j = 1, \ldots, m$），使得
$$
x_i = \sum\nolimits_j a_{ij} y_j, \forall i,
\quad\text{且}\quad
0 = \sum\nolimits_i f_ia_{ij}, \forall j.
$$
则称关系 $\sum f_i x_i$ 为{\it 平凡的}。

\begin{lemma}[平坦性的方程判据]
\label{algebra-lemma-flat-eq}
$R$ 上的模 $M$ 平坦，当且仅当 $M$ 中的每个关系都是平凡的。
\end{lemma}

\begin{proof}
假设 $M$ 平坦，设 $\sum f_i x_i = 0$ 为 $M$ 中的关系。
令 $I = (f_1, \ldots, f_n)$，并令
$K = \Ker(R^n \to I, (a_1, \ldots, a_n) \mapsto \sum_i a_i f_i)$。
于是有短正合序列
$0 \to K \to R^n \to I \to 0$。此时
$\sum f_i \otimes x_i$ 是 $I \otimes_R M$ 中的元素，
并映到 $R \otimes_R M = M$ 中的零。由平坦性，
$\sum f_i \otimes x_i$ 在 $I \otimes_R M$ 中为零。
因此存在 $K \otimes_R M$ 中的元素映到
$\sum e_i \otimes x_i \in R^n \otimes_R M$，其中 $e_i$
是 $R^n$ 的第 $i$ 个基元素。将该元素写成
$\sum k_j \otimes y_j$，再将 $k_j$ 在 $R^n$ 中的像写成
$\sum a_{ij} e_i$，即得结论。

\medskip\noindent
假设每个关系都是平凡的。设 $I$ 为有限生成理想，设
$x = \sum f_i \otimes x_i$ 是 $I \otimes_R M$ 中映到
$R \otimes_R M = M$ 中零的元素。这恰好意味着
$\sum f_i x_i$ 是 $M$ 中的关系。由它平凡这一事实容易得到
$x$ 为零，因为
$$
x
=
\sum f_i \otimes x_i
=
\sum f_i \otimes \left(\sum a_{ij}y_j\right)
=
\sum \left(\sum f_i a_{ij}\right) \otimes y_j
=
0
$$
\end{proof}

\begin{lemma}
\label{algebra-lemma-flat-tor-zero}
设 $R$ 为环，$0 \to M'' \to M' \to M \to 0$
为短正合序列，且 $N$ 为 $R$-模。若 $M$ 平坦，则
$N \otimes_R M'' \to N \otimes_R M'$ 是单射，即序列
$$
0 \to N \otimes_R M'' \to N \otimes_R M' \to N \otimes_R M \to 0
$$
是短正合序列。
\end{lemma}

\begin{proof}
设 $R^{(I)} \to N$ 是从自由模到 $N$ 的满射，核为 $K$。
将蛇形引理应用于下图即可得到结论：
$$
\begin{matrix}
 & & 0 & & 0 & & 0 & & \\
 & & \uparrow & & \uparrow & & \uparrow & & \\
 & & M''\otimes_R N & \to & M' \otimes_R N & \to & M \otimes_R N & \to & 0 \\
 & & \uparrow & & \uparrow & & \uparrow & & \\
0 & \to & (M'')^{(I)} & \to & (M')^{(I)} & \to & M^{(I)} & \to & 0 \\
 & & \uparrow & & \uparrow & & \uparrow & & \\
 & & M''\otimes_R K & \to & M' \otimes_R K & \to & M \otimes_R K & \to & 0 \\
 & & & & & & \uparrow & & \\
 & & & & & & 0 & &
\end{matrix}
$$
其中各行各列均正合。中间一行正合，因为与自由模 $R^{(I)}$ 作张量的运算是正合的。
\end{proof}

\begin{lemma}
\label{algebra-lemma-flat-ses}
设 $0 \to M' \to M \to M'' \to 0$ 为 $R$-模的短正合序列。
若 $M'$ 和 $M''$ 平坦，则 $M$ 也平坦。
若 $M$ 和 $M''$ 平坦，则 $M'$ 也平坦。
\end{lemma}

\begin{proof}
我们使用如下判据：若对于每个理想 $I \subset R$，映射
$N \otimes_R I \to N$ 都是单射，则模 $N$ 平坦，
见引理 \ref{algebra-lemma-flat}。取理想 $I \subset R$。
考虑图表
$$
\begin{matrix}
0 & \to & M' & \to & M & \to & M'' & \to & 0 \\
& & \uparrow & & \uparrow & & \uparrow & & \\
& & M'\otimes_R I & \to & M \otimes_R I & \to & M''\otimes_R I & \to & 0
\end{matrix}
$$
其中两行正合。这立即证明第一个断言。
第二个断言则是因为若 $M''$ 平坦，则由引理
\ref{algebra-lemma-flat-tor-zero}，左下方的水平箭头是单射。
\end{proof}

\begin{lemma}
\label{algebra-lemma-easy-ff}
设 $R$ 为环。设 $M$ 为 $R$-模。以下命题等价：
\begin{enumerate}
\item $M$ 忠实平坦；并且
\item $M$ 平坦，且对所有 $R$-模同态 $\alpha : N \to N'$，
有 $\alpha = 0$ 当且仅当 $\alpha \otimes \text{id}_M = 0$。
\end{enumerate}
\end{lemma}

\begin{proof}
若 $M$ 忠实平坦，则
$0 \to \Ker(\alpha) \to N \to N'$ 正合，当且仅当与 $M$ 作张量后仍然正合。
这证明 (1) 蕴含 (2)。
反之，假设 (2)。设 $N_1 \to N_2 \to N_3$ 为复形，并假设
$N_1 \otimes_R M \to N_2 \otimes_R M \to N_3\otimes_R M$
正合。取 $x \in \Ker(N_2 \to N_3)$，考虑映射
$\alpha : R \to N_2/\Im(N_1)$，
$r \mapsto rx + \Im(N_1)$。由复形 $-\otimes_R M$ 的正合性，
$\alpha \otimes \text{id}_M$ 为零。根据假设 $\alpha$ 为零。
因此 $x$ 属于 $N_1 \to N_2$ 的像。
\end{proof}

\begin{lemma}
\label{algebra-lemma-ff}
\begin{slogan}
平坦模忠实平坦，当且仅当它的纤维非零。
\end{slogan}
设 $M$ 为平坦的 $R$-模。以下命题等价：
\begin{enumerate}
\item $M$ 忠实平坦；
\item 对每个非零 $R$-模 $N$，张量积 $M \otimes_R N$
非零；
\item 对所有 $\mathfrak p \in \Spec(R)$，张量积
$M \otimes_R \kappa(\mathfrak p)$ 非零；并且
\item 对所有 $R$ 的极大理想 $\mathfrak m$，张量积
$M \otimes_R \kappa(\mathfrak m) = M/{\mathfrak m}M$
非零。
\end{enumerate}
\end{lemma}

\begin{proof}
假设 $M$ 忠实平坦且 $N \not = 0$。由引理 \ref{algebra-lemma-easy-ff}，
非零映射 $1 : N \to N$ 诱导非零映射
$M \otimes_R N \to M \otimes_R N$，因此 $M \otimes_R N \not = 0$。
故 (1) 蕴含 (2)。(2) $\Rightarrow$ (3) $\Rightarrow$ (4) 立即成立。

\medskip\noindent
假设 (4)。设 $N_1 \to N_2 \to N_3$ 为复形，且
$N_1 \otimes_R M \to N_2\otimes_R M \to
N_3\otimes_R M$ 正合。令 $H$ 为该复形的上同调，即
$H = \Ker(N_2 \to N_3)/\Im(N_1 \to N_2)$。
为完成证明，我们证明 $H = 0$。由平坦性，$H \otimes_R M = 0$。
取 $x \in H$，令 $I = \{f \in R \mid fx = 0 \}$
为其湮没子。由于 $R/I \subset H$，由 $M$ 的平坦性得到
$M/IM \subset H \otimes_R M = 0$。若 $I \not =  R$，可选取
极大理想 $I \subset \mathfrak m \subset R$。这立即产生矛盾。
\end{proof}

\begin{lemma}
\label{algebra-lemma-ff-rings}
设 $R \to S$ 为平坦环映射。以下命题等价：
\begin{enumerate}
\item $R \to S$ 忠实平坦；
\item 诱导的 $\Spec$ 映射是满射；并且
\item 每个闭点 $x \in \Spec(R)$
都在映射 $\Spec(S) \to \Spec(R)$ 的像中。
\end{enumerate}
\end{lemma}

\begin{proof}
由引理 \ref{algebra-lemma-ff} 很快得到结论，因为我们在注记
\ref{algebra-remark-fundamental-diagram} 中看到，$\mathfrak p$ 在像中，当且仅当环
$S \otimes_R \kappa(\mathfrak p)$ 非零。
\end{proof}

\begin{lemma}
\label{algebra-lemma-local-flat-ff}
局部环之间的平坦局部环同态是忠实平坦的。
\end{lemma}

\begin{proof}
由引理 \ref{algebra-lemma-ff-rings} 立即得到。
\end{proof}

\noindent
平坦性与局部化相容良好。

\begin{lemma}
\label{algebra-lemma-flat-localization}
设 $R$ 为环。设 $S \subset R$ 为乘法子集。
\begin{enumerate}
\item 局部化 $S^{-1}R$ 是平坦的 $R$-代数。
\item 若 $M$ 是 $S^{-1}R$-模，则 $M$ 是平坦的 $R$-模，
当且仅当 $M$ 是平坦的 $S^{-1}R$-模。
\item 设 $M$ 为 $R$-模。则 $M$ 是平坦的 $R$-模，当且仅当对 $R$ 的所有素理想
$\mathfrak p$，$M_{\mathfrak p}$ 是平坦的 $R_{\mathfrak p}$-模。
\item 设 $M$ 为 $R$-模。则 $M$ 是平坦的 $R$-模，当且仅当对 $R$ 的所有极大理想
$\mathfrak m$，$M_{\mathfrak m}$ 是平坦的 $R_{\mathfrak m}$-模。
\item 设 $R \to A$ 为环映射，$M$ 为 $A$-模，且
$g_1, \ldots, g_m \in A$ 为生成 $A$ 的单位理想的元素。则 $M$ 在 $R$ 上平坦，
当且仅当每个局部化 $M_{g_i}$ 在 $R$ 上平坦。
\item 设 $R \to A$ 为环映射，且 $M$ 为 $A$-模。则 $M$ 是平坦的 $R$-模，
当且仅当对 $A$ 的所有素理想 $\mathfrak q$，局部化
$M_{\mathfrak q}$ 是平坦的 $R_{\mathfrak p}$-模
（其中 $\mathfrak p$ 是位于 $\mathfrak q$ 下方的 $R$ 的素理想）。
\item 设 $R \to A$ 为环映射，且 $M$ 为 $A$-模。则 $M$ 是平坦的 $R$-模，
当且仅当对 $A$ 的所有极大理想 $\mathfrak m$，局部化
$M_{\mathfrak m}$ 是平坦的 $R_{\mathfrak p}$-模
（其中 $\mathfrak p = R \cap \mathfrak m$）。
\end{enumerate}
\end{lemma}

\begin{proof}
我们证明引理的最后一个断言。证明中将反复使用局部化是正合的并且与张量积可交换这一事实，
见 \ref{algebra-section-localization} 节和 \ref{algebra-section-tensor-product} 节。

\medskip\noindent
设 $R \to A$ 为环映射，$M$ 为 $A$-模。假设对 $A$ 的所有极大理想
$\mathfrak m$，$M_{\mathfrak m}$ 是平坦的 $R_{\mathfrak p}$-模，其中
$\mathfrak p = R \cap \mathfrak m$。设 $I \subset R$ 为理想。
我们要证明映射 $I \otimes_R M \to M$ 是单射。
可以将其视为 $A$-模的映射。由假设，局部化
$(I \otimes_R M)_{\mathfrak m} \to M_{\mathfrak m}$ 是单射，因为
$(I \otimes_R M)_{\mathfrak m} =
I_{\mathfrak p} \otimes_{R_{\mathfrak p}} M_{\mathfrak m}$。
因此由引理 \ref{algebra-lemma-characterize-zero-local}，$I \otimes_R M \to M$ 的核为零。
于是 $M$ 在 $R$ 上平坦。

\medskip\noindent
反之，假设 $M$ 在 $R$ 上平坦。取 $A$ 中位于 $R$ 的素理想
$\mathfrak p$ 上方的素理想 $\mathfrak q$。设
$I \subset R_{\mathfrak p}$ 为理想。我们要证明
$I \otimes_{R_{\mathfrak p}} M_{\mathfrak q} \to M_{\mathfrak q}$
是单射。可写成 $I = J_{\mathfrak p}$，其中 $J \subset R$ 为理想。
于是映射 $I \otimes_{R_{\mathfrak p}} M_{\mathfrak q} \to M_{\mathfrak q}$
正是映射 $J \otimes_R M \to M$ 在 $\mathfrak q$ 处的局部化，而后者是单射。
由于局部化是正合的，可知 $M_{\mathfrak q}$ 是平坦的
$R_{\mathfrak p}$-模。

\medskip\noindent
这证明了 (7) 和 (6)。其他断言可由最后一个断言直接得到（证明略）。
\end{proof}

\begin{lemma}
\label{algebra-lemma-flat-going-down}
设 $R \to S$ 为平坦映射。设 $\mathfrak p \subset \mathfrak p'$ 为 $R$ 的素理想。
设 $\mathfrak q' \subset S$ 为映到 $\mathfrak p'$ 的 $S$ 的素理想。
则存在素理想 $\mathfrak q \subset \mathfrak q'$，映到 $\mathfrak p$。
\end{lemma}

\begin{proof}
由引理 \ref{algebra-lemma-flat-localization}，局部环映射
$R_{\mathfrak p'} \to S_{\mathfrak q'}$ 平坦。
由引理 \ref{algebra-lemma-local-flat-ff}，该局部环映射忠实平坦。
由引理 \ref{algebra-lemma-ff-rings}，存在一个映到
$\mathfrak p R_{\mathfrak p'}$ 的素理想。该素理想在 $S$ 中的逆像即为所求。
\end{proof}

\noindent
引理中描述的 $R \to S$ 的性质称为“下降性质”。
见定义 \ref{algebra-definition-going-up-down}。

\begin{lemma}
\label{algebra-lemma-colimit-faithfully-flat}
设 $R$ 为环。设 $\{S_i, \varphi_{ii'}\}$ 为忠实平坦的 $R$-代数所成的有向系统。
则 $S = \colim_i S_i$ 是忠实平坦的 $R$-代数。
\end{lemma}

\begin{proof}
由引理 \ref{algebra-lemma-colimit-flat}，可知 $S$ 平坦。
设 $\mathfrak m \subset R$ 为极大理想。由引理
\ref{algebra-lemma-ff-rings}，环 $S_i/\mathfrak m S_i$ 均非零。
因此 $S/\mathfrak mS = \colim S_i/\mathfrak mS_i$ 也非零，
因为 $1$ 不等于零。于是映射 $\Spec(S) \to \Spec(R)$ 的像包含
$\mathfrak m$，由引理 \ref{algebra-lemma-ff-rings} 可知 $R \to S$
忠实平坦。
\end{proof}

\section{支集与湮没子}
\label{algebra-section-supp-and-ann}

\noindent
一些非常基本的定义和引理。

\begin{definition}
\label{algebra-definition-support-module}
设 $R$ 为环，$M$ 为 $R$-模。$M$ 的{\it 支集}是集合
$$
\text{Supp}(M)
=
\{
\mathfrak p \in \Spec(R)
\mid
M_{\mathfrak p} \not = 0
\}
$$
\end{definition}

\begin{lemma}
\label{algebra-lemma-support-zero}
\begin{slogan}
环上的模支集为空，当且仅当它是平凡模。
\end{slogan}
设 $R$ 为环，$M$ 为 $R$-模。则
$$
M = (0) \Leftrightarrow \text{Supp}(M) = \emptyset.
$$
\end{lemma}

\begin{proof}
事实上，引理 \ref{algebra-lemma-characterize-zero-local} 甚至表明：
若 $M$ 非零，则 $\text{Supp}(M)$ 总是包含一个极大理想。
\end{proof}

\begin{definition}
\label{algebra-definition-annihilator}
设 $R$ 为环，$M$ 为 $R$-模。
\begin{enumerate}
\item 给定元素 $m \in M$，$m$ 的{\it 湮没子}是理想
$$
\text{Ann}_R(m) = \text{Ann}(m) = \{f \in R \mid fm = 0\}.
$$
\item $M$ 的{\it 湮没子}是理想
$$
\text{Ann}_R(M) = \text{Ann}(M) = \{f \in R \mid fm = 0\ \forall m \in M\}.
$$
\end{enumerate}
\end{definition}

\begin{lemma}
\label{algebra-lemma-annihilator-flat-base-change}
设 $R \to S$ 为平坦环映射。设 $M$ 为 $R$-模且
$m \in M$。则 $\text{Ann}_R(m) S = \text{Ann}_S(m \otimes 1)$。
若 $M$ 是有限 $R$-模，则
$\text{Ann}_R(M) S = \text{Ann}_S(M \otimes_R S)$。
\end{lemma}

\begin{proof}
令 $I = \text{Ann}_R(m)$。按定义有正合序列
$0 \to I \to R \to M$，其中映射 $R \to M$ 将 $f$ 送到 $fm$。
利用平坦性得到正合序列
$0 \to I \otimes_R S \to S \to M \otimes_R S$，这证明了第一个断言。
若 $m_1, \ldots, m_n$ 是 $M$ 的一组生成元，则
$\text{Ann}_R(M) = \bigcap \text{Ann}_R(m_i)$。类似地，
$\text{Ann}_S(M \otimes_R S) = \bigcap \text{Ann}_S(m_i \otimes 1)$。
令 $I_i = \text{Ann}_R(m_i)$。于是只须证明
$\bigcap_{i = 1, \ldots, n} (I_i S) = (\bigcap_{i = 1, \ldots, n} I_i)S$。
这就是引理 \ref{algebra-lemma-flat-intersect-ideals}。
\end{proof}

\begin{lemma}
\label{algebra-lemma-support-closed}
设 $R$ 为环，$M$ 为 $R$-模。若 $M$ 有限，则
$\text{Supp}(M)$ 是闭集。更准确地，若 $I = \text{Ann}(M)$ 是 $M$ 的湮没子，
则 $V(I) = \text{Supp}(M)$。
\end{lemma}

\begin{proof}
我们证明 $V(I) = \text{Supp}(M)$。

\medskip\noindent
假设 $\mathfrak p \in \text{Supp}(M)$。则 $M_{\mathfrak p} \not = 0$。
取 $m \in M$，使其在 $M_\mathfrak p$ 中的像非零。
由局部化 $M_\mathfrak p$ 的构造，$m$ 的湮没子包含于
$\mathfrak p$。因此当然 $I = \text{Ann}(M)$ 也必须包含于 $\mathfrak p$。

\medskip\noindent
反之，假设 $\mathfrak p \not \in \text{Supp}(M)$。则
$M_{\mathfrak p} = 0$。令 $x_1, \ldots, x_r \in M$ 为生成元。
由引理 \ref{algebra-lemma-localization-colimit}，存在
$f \in R$，$f\not\in \mathfrak p$，使得 $x_i/1 = 0$ 在 $M_f$ 中成立。
因此对某个 $n_i \geq 1$ 有 $f^{n_i} x_i = 0$。
于是当 $n = \max\{n_i\}$ 时有 $f^nM = 0$，结论成立。
\end{proof}

\begin{lemma}
\label{algebra-lemma-support-base-change}
设 $R \to R'$ 为环映射，$M$ 为有限 $R$-模。
则 $\text{Supp}(M \otimes_R R')$ 是 $\text{Supp}(M)$ 的逆像。
\end{lemma}

\begin{proof}
设 $\mathfrak p \in \text{Supp}(M)$。由中山引理
（引理 \ref{algebra-lemma-NAK}）可知
$$
M \otimes_R \kappa(\mathfrak p) = M_\mathfrak p/\mathfrak p M_\mathfrak p
$$
是非零的 $\kappa(\mathfrak p)$-向量空间。
因此对每个位于 $\mathfrak p$ 上方的素理想 $\mathfrak p' \subset R'$，有
$$
(M \otimes_R R')_{\mathfrak p'}/\mathfrak p' (M \otimes_R R')_{\mathfrak p'} =
(M \otimes_R R') \otimes_{R'} \kappa(\mathfrak p') =
M \otimes_R \kappa(\mathfrak p) \otimes_{\kappa(\mathfrak p)}
\kappa(\mathfrak p')
$$
非零。这意味着 $\mathfrak p' \in \text{Supp}(M \otimes_R R')$。
反之，若 $\mathfrak p' \subset R'$ 是位于任意素理想 $\mathfrak p \subset R$ 上方的素理想，
则
$$
(M \otimes_R R')_{\mathfrak p'} =
M_\mathfrak p \otimes_{R_\mathfrak p} R'_{\mathfrak p'}.
$$
因此，若 $\mathfrak p' \in \text{Supp}(M \otimes_R R')$
位于素理想 $\mathfrak p \subset R$ 上方，则
$\mathfrak p \in \text{Supp}(M)$。
\end{proof}

\begin{lemma}
\label{algebra-lemma-support-element}
设 $R$ 为环，$M$ 为 $R$-模，且 $m \in M$。
则 $\mathfrak p \in V(\text{Ann}(m))$，当且仅当
$m$ 不映到 $M_\mathfrak p$ 中的零。
\end{lemma}

\begin{proof}
可用 $Rm \subset M$ 替换 $M$。此时 (1)
$\text{Ann}(m) = \text{Ann}(M)$；(2) $m$ 不映到 $M_\mathfrak p$ 中的零，
当且仅当 $\mathfrak p \in \text{Supp}(M)$。
结论现在由引理 \ref{algebra-lemma-support-closed} 得到。
\end{proof}

\begin{lemma}
\label{algebra-lemma-support-finite-presentation-constructible}
设 $R$ 为环，$M$ 为 $R$-模。若 $M$ 是有限表示的 $R$-模，则
$\text{Supp}(M)$ 是 $\Spec(R)$ 的闭子集，且其补集拟紧。
\end{lemma}

\begin{proof}
取表示
$$
R^{\oplus m} \longrightarrow R^{\oplus n} \longrightarrow M \to 0
$$
令 $A \in \text{Mat}(n \times m, R)$ 为第一个映射的矩阵。
由中山引理 \ref{algebra-lemma-NAK}，有
$$
M_{\mathfrak p} \not = 0 \Leftrightarrow
M \otimes \kappa(\mathfrak p) \not = 0 \Leftrightarrow
\text{rank}(A \bmod \mathfrak p) < n.
$$
因此，若 $I$ 是由 $A$ 的 $n \times n$ 子式生成的 $R$ 的理想，
则 $\text{Supp}(M) = V(I)$。由于 $I$
有限生成，设 $I = (f_1, \ldots, f_t)$，则
$\Spec(R) \setminus V(I)$ 是标准开集 $D(f_i)$ 的有限并，因而拟紧。
\end{proof}

\begin{lemma}
\label{algebra-lemma-support-quotient}
设 $R$ 为环，$M$ 为 $R$-模。
\begin{enumerate}
\item 若 $M$ 有限，则 $M/IM$ 的支集为
$\text{Supp}(M) \cap V(I)$。
\item 若 $N \subset M$，则 $\text{Supp}(N) \subset
\text{Supp}(M)$。
\item 若 $Q$ 是 $M$ 的商模，则 $\text{Supp}(Q) \subset
\text{Supp}(M)$。
\item 若 $0 \to N \to M \to Q \to 0$ 是短正合序列，
则 $\text{Supp}(M) = \text{Supp}(Q) \cup \text{Supp}(N)$。
\end{enumerate}
\end{lemma}

\begin{proof}
函子 $M \mapsto M_{\mathfrak p}$ 是正合的。这立即推出除第一个断言外的所有断言。
对于第一个断言，我们需要证明：若 $M_\mathfrak p \not = 0$ 且
$I \subset \mathfrak p$，则 $(M/IM)_{\mathfrak p}
= M_\mathfrak p/IM_\mathfrak p \not = 0$。这由中山引理
\ref{algebra-lemma-NAK} 得到。
\end{proof}

\section{上升与下降}
\label{algebra-section-going-up}

\noindent
设 $\mathfrak p$、$\mathfrak p'$ 为环 $R$ 的素理想。
令 $X = \Spec(R)$，赋予扎里斯基拓扑。记 $x \in X$ 为对应于
$\mathfrak p$ 的点，记 $x' \in X$ 为对应于
$\mathfrak p'$ 的点。则有：
$$
x' \leadsto x \Leftrightarrow \mathfrak p' \subset \mathfrak p.
$$
换言之，$x$ 是 $x'$ 的一个特殊化，当且仅当
$\mathfrak p' \subset \mathfrak p$。
术语和记号见拓扑一章的 \ref{topology-section-specialization} 节。

\begin{definition}
\label{algebra-definition-going-up-down}
设 $\varphi : R \to S$ 为环映射。
\begin{enumerate}
\item 若给定 $R$ 中的素理想 $\mathfrak p \subset \mathfrak p'$，
以及 $S$ 中位于 $\mathfrak p$ 上方的素理想 $\mathfrak q$，存在 $S$ 中的素理想
$\mathfrak q'$，满足 (a) $\mathfrak q \subset \mathfrak q'$，且 (b)
$\mathfrak q'$ 位于 $\mathfrak p'$ 上方，则称
$\varphi : R \to S$ 满足{\it 上升}。
\item 若给定 $R$ 中的素理想 $\mathfrak p \subset \mathfrak p'$，
以及 $S$ 中位于 $\mathfrak p'$ 上方的素理想 $\mathfrak q'$，存在 $S$ 中的素理想
$\mathfrak q$，满足 (a) $\mathfrak q \subset \mathfrak q'$，且 (b)
$\mathfrak q$ 位于 $\mathfrak p$ 上方，则称
$\varphi : R \to S$ 满足{\it 下降}。
\end{enumerate}
\end{definition}

\noindent
到目前为止，我们已经见到以下情形：
\begin{enumerate}
\item 整环映射满足上升，见引理 \ref{algebra-lemma-integral-going-up}。
\item 作为特例，有限环映射满足上升。
\item 作为特例，商映射 $R \to R/I$ 满足上升。
\item 平坦环映射满足下降，见引理 \ref{algebra-lemma-flat-going-down}。
\item 作为特例，任意局部化满足下降。
\item 若 $R \subset S$ 为整环的扩张，$R$ 正规且 $S$ 整于 $R$，
则满足下降，见命题 \ref{algebra-proposition-going-down-normal-integral}。
\end{enumerate}
下面是下降成立的另一种情形。

\begin{lemma}
\label{algebra-lemma-open-going-down}
设 $R \to S$ 为环映射。若诱导映射
$\varphi : \Spec(S) \to \Spec(R)$ 是开映射，则
$R \to S$ 满足下降。
\end{lemma}

\begin{proof}
设 $\mathfrak p \subset \mathfrak p' \subset R$，且
$\mathfrak q' \subset S$ 位于 $\mathfrak p'$ 上方。由于 $\varphi$ 是开映射，
对每个 $g \in S$，$g \not \in \mathfrak q'$，都有
$\mathfrak p$ 属于 $D(g) \subset \Spec(S)$ 的像。换言之，
$S_g \otimes_R \kappa(\mathfrak p)$ 非零。由于 $S_{\mathfrak q'}$
是这些 $S_g$ 的有向余极限，这意味着
$S_{\mathfrak q'} \otimes_R \kappa(\mathfrak p)$ 非零，见
引理 \ref{algebra-lemma-localization-colimit} 和
\ref{algebra-lemma-tensor-products-commute-with-limits}。
因此 $\mathfrak p$ 属于
$\Spec(S_{\mathfrak q'}) \to \Spec(R)$ 的像，正是所需结论。
\end{proof}

\begin{lemma}
\label{algebra-lemma-going-up-down-specialization}
设 $R \to S$ 为环映射。
\begin{enumerate}
\item $R \to S$ 满足下降，当且仅当一般化沿映射
$\Spec(S) \to \Spec(R)$ 提升，见拓扑一章定义
\ref{topology-definition-lift-specializations}。
\item $R \to S$ 满足上升，当且仅当特殊化沿映射
$\Spec(S) \to \Spec(R)$ 提升，见拓扑一章定义
\ref{topology-definition-lift-specializations}。
\end{enumerate}
\end{lemma}

\begin{proof}
略。
\end{proof}

\begin{lemma}
\label{algebra-lemma-going-up-down-composition}
设 $R \to S$ 和 $S \to T$ 为满足下降的环映射。
则 $R \to T$ 也满足下降。上升的情形同样成立。
\end{lemma}

\begin{proof}
根据引理 \ref{algebra-lemma-going-up-down-specialization}，结论由拓扑一章的引理
\ref{topology-lemma-lift-specialization-composition} 得到。
\end{proof}

\begin{lemma}
\label{algebra-lemma-image-stable-specialization-closed}
设 $R \to S$ 为环映射。令 $T \subset \Spec(R)$
为 $\Spec(S)$ 的像。若 $T$ 在特殊化下稳定，则 $T$ 是闭集。
\end{lemma}

\begin{proof}
我们给出两个证明。

\medskip\noindent
第一个证明。设 $\mathfrak p \subset R$ 为素理想，且
$\Spec(R)$ 中对应的点属于 $T$ 的闭包。这意味着对于每个
$f \in R$，$f \not \in \mathfrak p$，都有
$D(f) \cap T \not = \emptyset$。注意 $D(f) \cap T$
是 $\Spec(S_f)$ 在 $\Spec(R)$ 中的像。因此 $S_f \not = 0$。
换言之，在环 $S_f$ 中 $1 \not = 0$。由于 $S_{\mathfrak p}$ 是环 $S_f$ 的有向余极限，
可知在 $S_{\mathfrak p}$ 中 $1 \not = 0$。
换言之，$S_{\mathfrak p} \not = 0$；考察
$\Spec(S_{\mathfrak p}) \to \Spec(S) \to \Spec(R)$ 的像，
可知存在 $\mathfrak p' \in T$，满足 $\mathfrak p' \subset \mathfrak p$。
由于假设 $T$ 在特殊化下闭，得到 $\mathfrak p$ 是 $T$ 中的点，命题得证。

\medskip\noindent
第二个证明。令 $I = \Ker(R \to S)$。可用 $R/I$ 替换 $R$。
此时环映射 $R \to S$ 是单射。由引理
\ref{algebra-lemma-injective-minimal-primes-in-image}，$R$ 的所有极小素理想
都包含于像 $T$ 中。因此若 $T$ 在特殊化下稳定，则它包含所有素理想。
\end{proof}

\begin{lemma}
\label{algebra-lemma-going-up-closed}
设 $R \to S$ 为环映射。以下命题等价：
\begin{enumerate}
\item $R \to S$ 满足上升；并且
\item 映射 $\Spec(S) \to \Spec(R)$ 是闭映射。
\end{enumerate}
\end{lemma}

\begin{proof}
一般而言，拓扑空间的闭映射会使特殊化沿映射提升，见拓扑一章的引理
\ref{topology-lemma-closed-open-map-specialization}。
因此第二个条件推出第一个条件。

\medskip\noindent
假设 $R \to S$ 满足上升。设 $V(I) \subset \Spec(S)$ 为闭集。
我们要证明 $V(I)$ 在 $\Spec(R)$ 中的像是闭集。
环映射 $S \to S/I$ 显然满足上升。因此由引理
\ref{algebra-lemma-going-up-down-composition}，$R \to S \to S/I$ 满足上升。
将 $S$ 替换为 $S/I$ 后，只须证明 $\Spec(S)$ 在 $\Spec(R)$ 中的像
$T$ 是闭集。由拓扑一章引理
\ref{topology-lemma-open-closed-specialization} 和
\ref{topology-lemma-lift-specializations-images}，该像在特殊化下稳定。
于是结论由引理 \ref{algebra-lemma-image-stable-specialization-closed} 得到。
\end{proof}

\begin{lemma}
\label{algebra-lemma-constructible-stable-specialization-closed}
设 $R$ 为环，设 $E \subset \Spec(R)$ 为可构造子集。
\begin{enumerate}
\item 若 $E$ 在特殊化下稳定，则 $E$ 是闭集。
\item 若 $E$ 在一般化下稳定，则 $E$ 是开集。
\end{enumerate}
\end{lemma}

\begin{proof}
第一个证明。第一个断言由引理
\ref{algebra-lemma-image-stable-specialization-closed} 与引理
\ref{algebra-lemma-constructible-is-image} 结合得到。第二个断言成立是因为可构造集的补集
仍然可构造（见拓扑一章引理 \ref{topology-lemma-constructible}），
再应用本引理的第一部分以及拓扑一章引理
\ref{topology-lemma-open-closed-specialization}。

\medskip\noindent
第二个证明。

由于根据引理 \ref{algebra-lemma-spec-spectral}，$\Spec(R)$ 是谱空间，
这只是拓扑一章引理
\ref{topology-lemma-constructible-stable-specialization-closed} 的一个特例。
\end{proof}

\begin{proposition}
\label{algebra-proposition-fppf-open}
设 $R \to S$ 是平坦且有限表示的环映射。
则 $\Spec(S) \to \Spec(R)$ 是开映射。
更一般地，任意有限表示且满足下降性质的环映射
$R \to S$ 都具有此性质。
\end{proposition}

\begin{proof}
若 $R \to S$ 平坦，则由引理
\ref{algebra-lemma-flat-going-down}，$R \to S$ 满足下降性质。
因此，为证明命题，只须假设 $R \to S$ 有限表示且满足下降性质。

\medskip\noindent
由于 $g \in S$ 的标准开集 $D(g) \subset \Spec(S)$
构成拓扑的一个基，只须证明 $D(g)$ 的像是开集。
回忆（引理 \ref{algebra-lemma-standard-open}），
$\Spec(S_g) \to \Spec(S)$ 将 $\Spec(S_g)$ 同胚映到 $D(g)$。
由于 $S \to S_g$ 满足下降性质（见上文），由引理
\ref{algebra-lemma-going-up-down-composition} 可知 $R \to S_g$
满足下降性质。因此将 $S$ 替换为 $S_g$ 后，只须证明该像是开集。
由切瓦莱定理（定理 \ref{algebra-theorem-chevalley}），该像是可构造集 $E$。
又由于 $R \to S$ 满足下降性质，$E$ 在一般化下稳定，见拓扑一章引理
\ref{topology-lemma-open-closed-specialization} 和
\ref{topology-lemma-lift-specializations-images}。
因此由引理 \ref{algebra-lemma-constructible-stable-specialization-closed}，
$E$ 是开集。
\end{proof}

\begin{lemma}
\label{algebra-lemma-same-image}
设 $k$ 为域，$R$、$S$ 为 $k$-代数。
设 $S' \subset S$ 为 $k$-子代数，且 $f \in S' \otimes_k R$。
在交换图表
$$
\xymatrix{
\Spec((S \otimes_k R)_f) \ar[rd] \ar[rr] & &
\Spec((S' \otimes_k R)_f) \ar[ld] \\
& \Spec(R) &
}
$$
中，两条对角箭头的像相同。
\end{lemma}

\begin{proof}
设 $\mathfrak p \subset R$ 属于西南箭头的像。这意味着（引理
\ref{algebra-lemma-in-image}）
$$
(S' \otimes_k R)_f \otimes_R \kappa(\mathfrak p)
=
(S' \otimes_k \kappa(\mathfrak p))_f
$$
不是零环，即 $S' \otimes_k \kappa(\mathfrak p)$ 不是零环，且 $f$
在其中的像不是幂零元。环映射
$$
S' \otimes_k \kappa(\mathfrak p) \to S \otimes_k \kappa(\mathfrak p)
$$
是单射。因此 $S \otimes_k \kappa(\mathfrak p)$ 也不是零环，且 $f$
在其中的像不是幂零元。于是
$$
(S \otimes_k R)_f \otimes_R \kappa(\mathfrak p)
$$
不是零环。因此（引理 \ref{algebra-lemma-in-image}），$\mathfrak p$ 属于东南箭头的像，
这正是所需结论。
\end{proof}

\begin{lemma}
\label{algebra-lemma-map-into-tensor-algebra-open}
设 $k$ 为域。
设 $R$ 和 $S$ 为 $k$-代数。
映射 $\Spec(S \otimes_k R) \to \Spec(R)$
是开映射。
\end{lemma}

\begin{proof}
设 $f \in S \otimes_k R$。只须证明标准开集 $D(f)$ 的像是开集。
取一个有限型 $k$-子代数 $S' \subset S$，使得
$f \in S' \otimes_k R$。映射 $R \to S' \otimes_k R$ 平坦且有限表示，
所以由命题 \ref{algebra-proposition-fppf-open}，
$$
\Spec((S' \otimes_k R)_f) \to \Spec(R)
$$
的像 $U$ 是开集。由引理 \ref{algebra-lemma-same-image}，这也正是 $D(f)$ 的像，
命题得证。
\end{proof}

\noindent
下面是一个有时很有用的棘手引理。

\begin{lemma}
\label{algebra-lemma-unique-prime-over-localize-below}
设 $R \to S$ 为环映射。
设 $\mathfrak p \subset R$ 为素理想。假设
\begin{enumerate}
\item 在 $\mathfrak p$ 上方存在唯一的素理想
$\mathfrak q \subset S$，并且
\item 下列两者之一成立：
\begin{enumerate}
\item $R \to S$ 满足上升性质，或
\item $R \to S$ 满足下降性质，且 $R$ 的每个素理想上方至多有一个
素理想属于 $S$。
\end{enumerate}
\end{enumerate}
则 $S_{\mathfrak p} = S_{\mathfrak q}$。
\end{lemma}

\begin{proof}
考虑任意一个素理想 $\mathfrak q' \subset S$，它对应于
$\Spec(S_{\mathfrak p})$ 中的一个点。这意味着
$\mathfrak p' = R \cap \mathfrak q'$ 包含于 $\mathfrak p$。
图示如下：
$$
\xymatrix{
\mathfrak q' \ar@{-}[d] \ar@{-}[r] & ? \ar@{-}[r] \ar@{-}[d] & S \ar@{-}[d] \\
\mathfrak p' \ar@{-}[r] & \mathfrak p \ar@{-}[r] & R
}
$$
假设 (1) 和 (2)(a)。由上升性质，存在素理想 $\mathfrak q'' \subset S$，
使得 $\mathfrak q' \subset \mathfrak q''$ 且 $\mathfrak q''$ 位于
$\mathfrak p$ 上方。由 $\mathfrak q$ 的唯一性，得到
$\mathfrak q'' = \mathfrak q$。换言之，$\mathfrak q'$ 定义了
$\Spec(S_{\mathfrak q})$ 中的一个点。

\medskip\noindent
假设 (1) 和 (2)(b)。由下降性质，存在素理想
$\mathfrak q'' \subset \mathfrak q$，位于 $\mathfrak p'$ 上方。
由位于 $\mathfrak p'$ 上方的素理想的唯一性，得到
$\mathfrak q' = \mathfrak q''$。换言之，$\mathfrak q'$ 定义了
$\Spec(S_{\mathfrak q})$ 中的一个点。

\medskip\noindent
在两种情形中，我们都得到双射
$$
\Spec(S_{\mathfrak q}) \to \Spec(S_{\mathfrak p})
$$
。
显然，这意味着 $S - \mathfrak q$ 中的所有元素都在 $S_{\mathfrak p}$ 中可逆，
换言之 $S_{\mathfrak p} = S_{\mathfrak q}$。
\end{proof}

\noindent
下面的引理是平坦环映射的下降性质的一个推广。

\begin{lemma}
\label{algebra-lemma-going-down-flat-module}
设 $R \to S$ 为环映射。设 $N$ 为有限生成的 $S$-模，且作为 $R$-模平坦。
在 $\text{Supp}(N) \subset \Spec(S)$ 上赋予诱导拓扑。
则一般化沿映射 $\text{Supp}(N) \to \Spec(R)$ 提升。
\end{lemma}

\begin{proof}
该陈述的含义如下。设 $\mathfrak p \subset \mathfrak p' \subset R$ 为素理想。
设 $\mathfrak q' \subset S$ 为素理想，且 $\mathfrak q' \in \text{Supp}(N)$。
则存在素理想 $\mathfrak q \subset \mathfrak q'$，满足
$\mathfrak q \in \text{Supp}(N)$ 且位于 $\mathfrak p$ 上方。
由于 $N$ 在 $R$ 上平坦，由引理 \ref{algebra-lemma-flat-localization} 可知
$N_{\mathfrak q'}$ 在 $R_{\mathfrak p'}$ 上平坦。
由于 $N_{\mathfrak q'}$ 是 $S_{\mathfrak q'}$ 上的有限模，且因
$\mathfrak q' \in \text{Supp}(N)$ 而非零，由 Nakayama 引理
\ref{algebra-lemma-NAK} 可知
$$
N_{\mathfrak q'} \otimes_{S_{\mathfrak q'}} \kappa(\mathfrak q')
$$
非零。因此
$$
N_{\mathfrak q'} \otimes_{R_{\mathfrak p'}} \kappa(\mathfrak p')
$$
也非零。由引理 \ref{algebra-lemma-ff} 得到
$$
N_{\mathfrak q'} \otimes_{R_{\mathfrak p'}} \kappa(\mathfrak p)
$$
非零。令
$$
J \subset S_{\mathfrak q'} \otimes_{R_{\mathfrak p'}} \kappa(\mathfrak p)
$$
为有限非零模
$$
N_{\mathfrak q'} \otimes_{R_{\mathfrak p'}} \kappa(\mathfrak p)
$$
的湮没子。由于 $J$ 是真理想，我们可以选取素理想 $\mathfrak q \subset S$，
使其对应于
$$
S_{\mathfrak q'} \otimes_{R_{\mathfrak p'}} \kappa(\mathfrak p)/J
$$
的一个素理想。该素理想属于 $N$ 的支集，位于 $\mathfrak p$ 上方，且包含于
$\mathfrak q'$，正是所需结论。
\end{proof}

\section{可分扩张}
\label{algebra-section-separability}

\noindent
本节讨论非代数域扩张的可分性。这与几何约化代数的概念密切相关，见定义
\ref{algebra-definition-geometrically-reduced}。

\begin{definition}
\label{algebra-definition-separable-field-extension}
设 $K/k$ 为域扩张。
\begin{enumerate}
\item 若存在 $K/k$ 的超越基 $\{x_i; i \in I\}$，使得扩张
$K/k(x_i; i \in I)$ 是可分代数扩张，则称 $K$ 在 $k$ 上{\it 可分生成}。
\item 若对于每个子扩张 $k \subset K' \subset K$，其中 $K'$ 在 $k$ 上有限生成，
扩张 $K'/k$ 都是可分生成的，则称 $K$ 在 $k$ 上{\it 可分}。
\end{enumerate}
\end{definition}

\noindent
由这个略显笨拙的定义，尚不清楚一个可分生成的域扩张本身是否可分。
下面将证明确实如此，见引理 \ref{algebra-lemma-separably-generated-separable}。

\begin{lemma}
\label{algebra-lemma-subextensions-are-separable}
设 $K/k$ 为可分域扩张。对于任意子扩张 $K/K'/k$，域扩张
$K'/k$ 都是可分的。
\end{lemma}

\begin{proof}
这直接由定义得到。
\end{proof}

\begin{lemma}
\label{algebra-lemma-generating-finitely-generated-separable-field-extensions}
设 $K/k$ 是一个可分生成且有限生成的域扩张。
令 $r = \text{trdeg}_k(K)$。则存在 $K$ 中的元素
$x_1, \ldots, x_{r + 1}$，满足
\begin{enumerate}
\item $x_1, \ldots, x_r$ 是 $K$ 在 $k$ 上的超越基；
\item $K = k(x_1, \ldots, x_{r + 1})$；
\item $x_{r + 1}$ 在 $k(x_1, \ldots, x_r)$ 上可分。
\end{enumerate}
\end{lemma}

\begin{proof}
将定义与域一章引理 \ref{fields-lemma-primitive-element} 结合即可。
\end{proof}

\begin{lemma}
\label{algebra-lemma-make-separably-generated}
设 $K/k$ 为有限生成域扩张。存在图表
$$
\xymatrix{
K \ar[r] & K' \\
k \ar[u] \ar[r] & k' \ar[u]
}
$$
其中 $k'/k$、$K'/K$ 是有限纯不可分域扩张，且 $K'/k'$ 是可分生成域扩张。
\end{lemma}

\begin{proof}
只有在 $k$ 的特征为 $p > 0$ 时该引理才有意义。
取 $x_1, \ldots, x_r$ 为 $K$ 在 $k$ 上的超越基。
由于 $K$ 在 $k$ 上有限生成，扩张
$k(x_1, \ldots, x_r) \subset K$ 是有限扩张。
令 $K/K_{sep}/k(x_1, \ldots, x_r)$ 为域一章引理
\ref{fields-lemma-separable-first} 中给出的子扩张。
若 $K = K_{sep}$，则结论成立。我们对
$d = [K : K_{sep}]$ 使用归纳法。

\medskip\noindent
假设 $d > 1$。取 $\beta \in K$，使得
$\alpha = \beta^p \in K_{sep}$ 且 $\beta \not \in K_{sep}$。
令 $P = T^n + a_1T^{n - 1} + \ldots + a_n$
为 $\alpha$ 在 $k(x_1, \ldots, x_r)$ 上的最小多项式。
令 $k'/k$ 为通过添加 $p$ 次根得到的有限纯不可分扩张，
使得每个 $a_i$ 都是
$k'(x_1^{1/p}, \ldots, x_r^{1/p})$ 中的 $p$ 次幂。
这样的扩张存在；细节略去。
令 $L$ 为适合下图的域
$$
\xymatrix{
K \ar[r] & L \\
k(x_1, \ldots, x_r) \ar[u] \ar[r] & k'(x_1^{1/p}, \ldots, x_r^{1/p}) \ar[u]
}
$$
我们可以并且确实假设 $L$ 是 $K$ 与
$k'(x_1^{1/p}, \ldots, x_r^{1/p})$ 的复合域。
令 $L/L_{sep}/k'(x_1^{1/p}, \ldots, x_r^{1/p})$
为域一章引理 \ref{fields-lemma-separable-first} 中给出的子扩张。
则 $L_{sep}$ 是 $K_{sep}$ 与
$k'(x_1^{1/p}, \ldots, x_r^{1/p})$ 的复合域。
元素 $\alpha \in L_{sep}$ 是多项式 $P$ 的一个根；该多项式的所有系数
都是 $k'(x_1^{1/p}, \ldots, x_r^{1/p})$ 中的 $p$ 次幂，且其根两两不同。
由域一章引理 \ref{fields-lemma-pth-root}，可知存在某个
$\alpha' \in L_{sep}$，使得 $\alpha = (\alpha')^p$。
显然，这意味着 $\beta$ 映到 $\alpha' \in L_{sep}$。
换言之，我们得到域塔
$$
\xymatrix{
K \ar[r] & L \\
K_{sep}(\beta) \ar[r] \ar[u] & L_{sep} \ar[u] \\
K_{sep} \ar[r] \ar[u] & L_{sep} \ar@{=}[u] \\
k(x_1, \ldots, x_r) \ar[u] \ar[r] & k'(x_1^{1/p}, \ldots, x_r^{1/p}) \ar[u] \\
k \ar[r] \ar[u] & k' \ar[u]
}
$$
因此该构造导致新的情形
$[L : L_{sep}] < [K : K_{sep}]$。由归纳假设，存在
$k' \subset k''$ 和 $L \subset L'$，使其对于扩张 $L/k'$ 满足本引理的结论。
那么，对于扩张 $K/k$，扩张 $k''/k$ 和 $L'/K$ 就满足要求。
引理得证。
\end{proof}

\section{几何约化代数}
\label{algebra-section-geometrically-reduced}

\noindent
关于几何约化代数的主要结果是引理
\ref{algebra-lemma-geometrically-reduced-finite-purely-inseparable-extension}。
建议读者读完定义后跳到该引理。

\begin{definition}
\label{algebra-definition-geometrically-reduced}
设 $k$ 为域，$S$ 为 $k$-代数。
若对于每个域扩张 $K/k$，$K$-代数
$K \otimes_k S$ 都是约化的，则称 $S$ 在 $k$ 上{\it 几何约化}。
\end{definition}

\noindent
设 $k$ 为域，且 $S$ 为约化的 $k$-代数。要检验 $S$ 是几何约化的，
只须检验 $\overline{k} \otimes_k S$ 是约化的（其中 $\overline{k}$ 表示
$k$ 的代数闭包）。事实上，只须对有限纯不可分域扩张 $k'/k$ 检验即可。见
引理 \ref{algebra-lemma-geometrically-reduced-finite-purely-inseparable-extension}。

\begin{lemma}
\label{algebra-lemma-subalgebra-separable}
几何约化代数的初等性质。设 $k$ 为域，$S$ 为 $k$-代数。
\begin{enumerate}
\item 若 $S$ 在 $k$ 上几何约化，则其每个 $k$-子代数也几何约化。
\item 若 $S$ 的所有有限生成 $k$-子代数都是几何约化的，则 $S$ 也是几何约化的。
\item 几何约化 $k$-代数的有向余极限也是几何约化的。
\item 若 $S$ 在 $k$ 上几何约化，则 $S$ 的任意局部化也在 $k$ 上几何约化。
\end{enumerate}
\end{lemma}

\begin{proof}
略。第二和第三个性质由张量积与余极限交换这一事实得到。
\end{proof}

\begin{lemma}
\label{algebra-lemma-geometrically-reduced-permanence}
设 $k$ 为域。若 $R$ 在 $k$ 上几何约化，且 $S \subset R$ 为乘法子集，
则局部化 $S^{-1}R$ 在 $k$ 上几何约化。
若 $R$ 在 $k$ 上几何约化，则 $R[x]$ 在 $k$ 上几何约化。
\end{lemma}

\begin{proof}
略。提示：约化环的局部化仍是约化的，且局部化与张量积交换。
\end{proof}

\noindent
在下面引理的证明中，我们会反复使用以下观察：若
$R' \subset R$ 且 $S' \subset S$ 是 $k$-代数的包含，则映射
$R' \otimes_k S' \to R \otimes_k S$ 是单射。

\begin{lemma}
\label{algebra-lemma-limit-argument}
设 $k$ 为域，$R$、$S$ 为 $k$-代数。
\begin{enumerate}
\item 若 $R \otimes_k S$ 非约化，则存在有限生成子代数 $R' \subset R$、
$S' \subset S$，使得 $R' \otimes_k S'$ 非约化。
\item 若 $R \otimes_k S$ 含有非零零因子，则存在有限生成子代数
$R' \subset R$、$S' \subset S$，使得 $R' \otimes_k S'$ 含有非零零因子。
\item 若 $R \otimes_k S$ 含有非平凡幂等元，则存在有限生成子代数
$R' \subset R$、$S' \subset S$，使得 $R' \otimes_k S'$ 含有非平凡幂等元。
\end{enumerate}
\end{lemma}

\begin{proof}
设 $z \in R \otimes_k S$ 为幂零元。可写成
$z = \sum_{i = 1, \ldots, n} x_i \otimes y_i$。
因此可以取 $R'$ 为 $x_i$ 生成的 $k$-子代数，取 $S'$ 为 $y_i$
生成的 $k$-子代数。第二和第三个断言的证明相同。
\end{proof}

\begin{lemma}
\label{algebra-lemma-geometrically-reduced-any-reduced-base-change}
令 $k$ 为域。
令 $S$ 为几何约化的 $k$-代数。
令 $R$ 为任意约化的 $k$-代数。
则 $R \otimes_k S$ 是约化的。
\end{lemma}

\begin{proof}
由引理 \ref{algebra-lemma-limit-argument}，我们可以假设 $R$ 是 $k$ 上的有限型代数。
于是，作为约化的诺特环，$R$ 嵌入域的有限乘积中
（见引理 \ref{algebra-lemma-total-ring-fractions-no-embedded-points}、
\ref{algebra-lemma-Noetherian-irreducible-components} 和
\ref{algebra-lemma-minimal-prime-reduced-ring}）。
因此我们可以假设 $R$ 是域的有限乘积。在这种情形下，由定义
\ref{algebra-definition-geometrically-reduced} 可知 $R \otimes_k S$ 是约化的。
\end{proof}

\begin{lemma}
\label{algebra-lemma-separable-extension-preserves-reducedness}
令 $k$ 为域。
令 $S$ 为约化的 $k$-代数。
令 $K/k$ 或为可分域扩张，或为可分生成的域扩张。
则 $K \otimes_k S$ 是约化的。
\end{lemma}

\begin{proof}
假设 $k \subset K$ 为可分扩张。
由引理 \ref{algebra-lemma-limit-argument}，我们可以假设 $S$ 是 $k$ 上的有限型代数，
且 $K$ 在 $k$ 上有限生成。
于是 $S$ 嵌入域的有限乘积中，即其全分式环（见引理
\ref{algebra-lemma-minimal-prime-reduced-ring} 和
\ref{algebra-lemma-total-ring-fractions-no-embedded-points}）。
因此实际上可以假设 $S$ 是整环。
按照引理 \ref{algebra-lemma-generating-finitely-generated-separable-field-extensions}，
取 $x_1, \ldots, x_{r + 1} \in K$。
令 $P \in k(x_1, \ldots, x_r)[T]$ 为 $x_{r + 1}$ 的最小多项式。它是可分多项式。
容易看出
$k[x_1, \ldots, x_r] \otimes_k S = S[x_1, \ldots, x_r]$ 是整环。
这推出 $k(x_1, \ldots, x_r) \otimes_k S$ 是整环，
因为它是 $S[x_1, \ldots, x_r]$ 的局部化。
环扩张 $k(x_1, \ldots, x_r) \otimes_k S \subset K \otimes_k S$
由单个元素 $x_{r + 1}$ 以及一个方程（即 $P$）生成。因此 $K \otimes_k S$ 嵌入
$F[T]/(P)$，其中 $F$ 是 $k(x_1, \ldots, x_r) \otimes_k S$ 的分式域。
由于 $P$ 可分，这个环是域的有限乘积，结论成立。

\medskip\noindent
此时我们尚不知道可分生成的域扩张是可分的，
所以还必须在这种情形下证明该引理。
为此，设 $\{x_i\}_{i \in I}$ 是 $K$ 在 $k$ 上的分离超越基。
对于任意有限集的元素 $\lambda_j \in K$，存在 $T \subset I$ 的有限子集，使得
$k(\{x_i\}_{i\in T}) \subset k(\{x_i\}_{i \in T} \cup \{\lambda_j\})$
是有限可分扩张。因此 $K$ 是 $k$ 上有限生成且可分生成扩张的有向余极限。
于是前一段的论证同样适用于此情形。
\end{proof}

\begin{lemma}
\label{algebra-lemma-generic-points-geometrically-reduced}
令 $k$ 为域，$S$ 为 $k$-代数。假设 $S$ 是约化的，且 $S_{\mathfrak p}$ 是几何
约化的，其中 $\mathfrak p$ 遍历 $S$ 的每个极小素理想。
则 $S$ 是几何约化的。
\end{lemma}

\begin{proof}
由于 $S$ 是约化的，映射
$S \to \prod_{\mathfrak p\text{ minimal}} S_{\mathfrak p}$
是单射，见引理 \ref{algebra-lemma-reduced-ring-sub-product-fields}。
若 $K/k$ 是域扩张，则映射
$$
S \otimes_k K \to (\prod S_\mathfrak p) \otimes_k K \to
\prod S_\mathfrak p \otimes_k K
$$
是单射：第一个映射因为 $k \to K$ 是平坦的，第二个直接可见，
因为 $K$ 是自由的 $k$-模。由于 $S_\mathfrak p$ 是几何约化的，
右侧的环是约化的。因此 $S \otimes_k K$ 是约化环的子环，故为约化的。
\end{proof}

\begin{lemma}
\label{algebra-lemma-separable-algebraic-diagonal}
令 $k'/k$ 为可分代数扩张。
则存在乘法子集 $S \subset k' \otimes_k k'$，使得乘法映射
$k' \otimes_k k' \to k'$ 被识别为
$k' \otimes_k k' \to S^{-1}(k' \otimes_k k')$。
\end{lemma}

\begin{proof}
先假设 $k'/k$ 是有限可分扩张。则 $k' = k(\alpha)$，
见《域》，引理 \ref{fields-lemma-primitive-element}。
令 $P \in k[x]$ 为 $\alpha$ 在 $k$ 上的最小多项式。
则 $P$ 是不可约、可分且首一的多项式，见《域》，
小节 \ref{fields-section-separable-extensions}。
于是 $k'[x]/(P) \to k' \otimes_k k'$，
$\sum \alpha_i x^i \mapsto \alpha_i \otimes \alpha^i$ 是同构。
我们可以在 $k'[x]$ 中分解 $P = (x - \alpha) Q$，由于 $P$ 可分，
可知 $Q(\alpha) \not = 0$。
于是显然，乘法子集 $S'$ 由
$Q$ 在 $k'[x]/(P)$ 中生成并起作用，即
$k' = (S')^{-1}(k'[x]/(P))$。
通过结构传递，$S'$ 在 $k' \otimes_k k'$ 中的像 $S$ 也起作用。

\medskip\noindent
在一般情形下，写 $k' = \bigcup k_i$ 为其有限子域扩张（在 $k$ 上）的并。
对每个 $i$，存在乘法子集 $S_i \subset k_i \otimes_k k_i$，使得
$k_i = S_i^{-1}(k_i \otimes_k k_i)$。令
$S$ 为 $\bigcup S_i \subset k' \otimes_k k'$ 的乘法闭包。
显然，乘法映射将 $S$ 的每个元素映为 $k'$ 中的可逆元素。
利用 $k' \otimes_k k'$ 是环 $k_i \otimes_k k_i$ 的并，可知乘法映射核中的
每个元素都被映为 $S^{-1}(k' \otimes_k k')$ 中的零。
利用局部化的正合性，结论得证。
\end{proof}

\begin{lemma}
\label{algebra-lemma-geometrically-reduced-over-separable-algebraic}
令 $k'/k$ 为可分代数域扩张。
令 $A$ 为 $k'$-代数。则 $A$ 在
$k$ 上几何约化，当且仅当它在 $k'$ 上几何约化。
\end{lemma}

\begin{proof}
假设 $A$ 在 $k'$ 上几何约化。
令 $K/k$ 为域扩张。由引理
\ref{algebra-lemma-separable-extension-preserves-reducedness}，$K \otimes_k k'$ 是约化环。
因此由引理 \ref{algebra-lemma-geometrically-reduced-any-reduced-base-change}，可知
$K \otimes_k A = (K \otimes_k k') \otimes_{k'} A$ 是约化的。

\medskip\noindent
假设 $A$ 在 $k$ 上几何约化。令 $K/k'$ 为域扩张。则
$$
K \otimes_{k'} A = (K \otimes_k A) \otimes_{(k' \otimes_k k')} k'
$$
由于由引理 \ref{algebra-lemma-separable-algebraic-diagonal}，
$k' \otimes_k k' \to k'$ 是局部化，故可知 $K \otimes_{k'} A$
是约化代数的局部化，因此是约化的。
\end{proof}

\section{可分扩张（续）}
\label{algebra-section-separability-continued}

\noindent
本节继续讨论小节 \ref{algebra-section-separability} 中开始的内容。

\begin{lemma}
\label{algebra-lemma-mini-separability}
令 $k$ 为特征为 $p > 1$ 的域。令 $K/k$ 为由
$x_1, \ldots, x_{n + 1} \in K$ 生成的域扩张，满足
\begin{enumerate}
\item $\{x_1, \ldots, x_n\}$ 是 $K/k$ 的超越基，
\item 对于每个 $k$-线性无关的子集
$\{a_1, \ldots, a_m\}$（其中 $K$），集合
$\{a^p_1, \ldots, a_m^p\}$ 是 $k$-线性无关的。
\end{enumerate}
则存在 $1 \leq j \leq n+1$，使得
$\{ x_1, \ldots, \widehat{x}_j, \ldots, x_{n+1}\}$
是 $K / k$ 的分离超越基。
\end{lemma}

\begin{proof}
由假设，$x_{n + 1}$ 在 $k(x_1, \ldots, x_n)$ 上代数，
因此存在非零多项式 $F \in k[X_1, \ldots, X_{n + 1}]$，使得
$F(x_1, \ldots, x_{n+1}) = 0$。取总次数最小的 $F$。
则 $F$ 是不可约的，因为至少有一个不可约因子也必须具有相同的性质。

\medskip\noindent
我们断言，对某个 $i$，$X_i$ 在 $F$ 中出现的幂不全是 $p$ 的倍数。
反设 $F$ 中出现的所有指数都是 $p$ 的倍数，则集合
$$
\{x_1^{\alpha_1} \ldots x^{\alpha_{n+1}}_{n+1} \mid \lambda_\alpha \neq 0\}
\subset
K
$$
是 $k$-线性相关的，其中 $\lambda_\alpha$ 是 $F$ 的系数。
由假设 (2)，我们得到集合
$$
\{x_1^{\alpha_1 / p} \ldots x^{\alpha_{n+1} / p}_{n+1}
\mid \lambda_\alpha \neq 0 \}
$$
也是 $k$-线性相关的，这与 $\deg(F)$ 的最小性矛盾。

\medskip\noindent
取某个 $i$，使 $p$ 次幂以外的 $X_i$ 幂出现在 $F$ 中。于是可见 $x_i$ 在下列域上代数：
$$
L = k(x_1, \ldots, x_{i - 1}, x_{i + 1}, \ldots, x_{n+1})
$$
由《域》，引理 \ref{fields-lemma-transcendence-degree}，可知下列元素：
$$
x_1, \ldots, x_{i - 1}, x_{i + 1}, \ldots, x_{n+1}
$$
是 $K/k$ 的超越基。因此 $L$ 是由下列元素：
$$
x_1, \ldots, x_{i - 1}, x_{i + 1}, \ldots, x_{n + 1}
$$
生成的、在 $k$ 上的多项式环的分式域。
由 Gauss 引理，我们得到
$$
P(T) =
F(x_1, \ldots, x_{i - 1}, T, x_{i + 1}, \ldots, x_{n + 1}) \in L[T]
$$
是不可约的。按构造，$P(T)$ 不包含于 $L[T^p]$ 中。
因此 $K/L$ 是可分的，正如所要求的。
\end{proof}

\noindent
令 $p$ 为素数，令 $k$ 为特征为 $p$ 的域。
在此情形下，我们用 $k^{1/p}$ 表示通过向 $k$ 添入
$p$ 次根（对应于 $k$ 中的所有元素）并添入 $k$ 所得的扩张。
（换言之，它是代数闭包中由 $k$ 生成的、由 $p$ 次根（取自 $k$ 的元素）生成的子域。）

\begin{lemma}
\label{algebra-lemma-characterize-separable-field-extensions}
令 $k$ 为特征为 $p > 0$ 的域。
令 $K/k$ 为域扩张。
以下命题等价：
\begin{enumerate}
\item $K$ 在 $k$ 上可分，
\item 对于每个 $k$-线性无关子集 $\{a_1, \ldots, a_m\}$，
在 $K$ 中，其对应的集合 $\{a^p_1, \ldots, a_m^p\}$ 是 $k$-线性无关的，
\item 环 $K \otimes_k k^{1/p}$ 是约化的，并且
\item $K$ 在 $k$ 上几何约化。
\end{enumerate}
\end{lemma}
\begin{proof}
蕴含关系 (1) $\Rightarrow$ (4) 由引理
\ref{algebra-lemma-separable-extension-preserves-reducedness} 得出。
蕴含关系 (4) $\Rightarrow$ (3) 是直接的。

\medskip\noindent
假设 (3)。考虑由
$m : K \otimes_k k^{1/p} \rightarrow K$ 给出的环同态
$$
\lambda \otimes \mu \rightarrow \lambda^p \mu^p
$$
注意，对于所有 $x \in K \otimes_k k^{1/p}$，都有
$x^p = m(x) \otimes 1$。
由于 $K \otimes_k k^{1/p}$ 是约化的，
可见 $m$ 是单射。若 $\{a_1, \ldots, a_m\} \subset K$
是 $k$-线性无关的，则 $\{a_1 \otimes 1, \ldots, a_m \otimes 1\}$
是 $k^{1/p}$-线性无关的。由 $m$ 的单射性可知，
$a_1^p, \ldots, a_m^p$ 的任何非平凡 $k$-线性组合都不为零。
因此 (3) 推出 (2)。

\medskip\noindent
假设 (2)。为证明 (1)，可假设 $K$ 在 $k$ 上有限生成，
并且需要证明 $K$ 在 $k$ 上可分生成。
令 $\{x_1, \ldots, x_d\}$ 为 $K/k$ 的超越基。由
《域》，引理 \ref{fields-lemma-algebraic-finitely-generated}，
有 $[K : K'] < \infty$，其中 $K' = k(x_1, \ldots, x_d)$。
选取超越基，使不可分次数 $[K : K']_i$ 最小。
若 $K / K'$ 可分，则结论成立。
反设并导出矛盾。则存在 $x_{d + 1} \in K$，它在 $K'$ 上不可分，
特别地，$[K'(x_{d+1}) : K']_i > 1$。于是由
引理 \ref{algebra-lemma-mini-separability}，存在 $1 \leq j \leq n + 1$，
使得 $K'' = k(x_1, \ldots, \widehat{x}_j, \ldots, x_{d+1})$
满足 $[K'(x_{d+1}) : K'']_i = 1$。由次数的乘法性，
$[K : K'']_i < [K : K']_i$，矛盾得证。
\end{proof}

\begin{lemma}
\label{algebra-lemma-separably-generated-separable}
可分生成的域扩张是可分的。
\end{lemma}

\begin{proof}
结合引理 \ref{algebra-lemma-separable-extension-preserves-reducedness} 与
引理 \ref{algebra-lemma-characterize-separable-field-extensions}。
\end{proof}

\noindent
下面的引理将使用完美闭包的概念，
其定义见定义 \ref{algebra-definition-perfection}。

\begin{lemma}
\label{algebra-lemma-geometrically-reduced-finite-purely-inseparable-extension}
令 $k$ 为域。令 $S$ 为 $k$-代数。
以下命题等价：
\begin{enumerate}
\item 对每个有限的纯不可分扩张，$k' \otimes_k S$ 是约化的，其中 $k'$ 是 $k$ 的扩张，
\item $k^{1/p} \otimes_k S$ 是约化的，
\item $k^{perf} \otimes_k S$ 是约化的，其中 $k^{perf}$ 是
$k$ 的完美闭包，
\item $\overline{k} \otimes_k S$ 是约化的，其中 $\overline{k}$ 是
$k$ 的代数闭包，并且
\item $S$ 在 $k$ 上几何约化。
\end{enumerate}
\end{lemma}

\begin{proof}
注意，任何有限纯不可分扩张 $k'/k$ 都嵌入完美闭包中。
此外，在 $k^{perf}$ 中，$k^{1/p}$ 嵌入 $k^{perf}$，而后者嵌入
$\overline{k}$。因此显然有
(5) $\Rightarrow$ (4) $\Rightarrow$ (3) $\Rightarrow$ (2)，
并且 (3) $\Rightarrow$ (1)。

\medskip\noindent
我们证明 (1) $\Rightarrow$ (5)。
假设 $k' \otimes_k S$ 对每个有限
纯不可分扩张 $k'$（在 $k$ 上）都是约化的。令 $K/k$ 为域扩张。我们需要证明 $K \otimes_k S$
是约化的。由引理 \ref{algebra-lemma-limit-argument}，可归约到
$K/k$ 是有限生成域扩张的情形。取如下图表：
$$
\xymatrix{
K \ar[r] & K' \\
k \ar[u] \ar[r] & k' \ar[u]
}
$$
如引理 \ref{algebra-lemma-make-separably-generated} 中所述。
由假设，$k' \otimes_k S$ 是约化的。
由引理 \ref{algebra-lemma-separable-extension-preserves-reducedness}，
可知 $K' \otimes_k S$ 是约化的。
因此 $K \otimes_k S$ 是约化的，正如所要求的。

\medskip\noindent
最后证明 (2) $\Rightarrow$ (5)。
假设 $k^{1/p} \otimes_k S$ 是约化的。则 $S$ 是约化的。
此外，对极小素理想处的每个局部化 $S_{\mathfrak p}$，
素理想 $\mathfrak p$，环 $k^{1/p}\otimes_k S_{\mathfrak p}$
是 $k^{1/p} \otimes_k S$ 的局部化，因而是约化的。
但是由引理 \ref{algebra-lemma-minimal-prime-reduced-ring}，$S_{\mathfrak p}$ 是域，
所以由引理 \ref{algebra-lemma-characterize-separable-field-extensions}，
$S_{\mathfrak p}$ 是几何约化的。
由引理 \ref{algebra-lemma-generic-points-geometrically-reduced}，
可知 $S$ 是几何约化的。
\end{proof}

\section{完美域}
\label{algebra-section-perfect-fields}

\noindent
下面给出定义。

\begin{definition}
\label{algebra-definition-perfect}
令 $k$ 为域。若 $k$ 的每个域扩张都在 $k$ 上可分，
则称 $k$ 为{\it 完美的}。
\end{definition}

\begin{lemma}
\label{algebra-lemma-perfect}
域 $k$ 完美，当且仅当它是特征为 $0$ 的域，或是特征为 $p > 0$ 的域，
且每个元素都有一个 $p$ 次根。
\end{lemma}

\begin{proof}
特征为零的情形显然。
假设 $k$ 的特征为 $p > 0$。
若 $k$ 完美，则所有向 $k$ 添入某个元素的
$p$ 次根所得到的域扩张都必须是平凡的，因此 $k$ 的每个元素都有 $p$ 次根。
反之，若每个元素都有 $p$ 次根，则 $k = k^{1/p}$，且 $k$ 的每个域扩张都
由引理 \ref{algebra-lemma-characterize-separable-field-extensions} 可知是可分的。
\end{proof}

\begin{lemma}
\label{algebra-lemma-make-separable}
令 $K/k$ 为有限生成的域扩张。
存在图表
$$
\xymatrix{
K \ar[r] & K' \\
k \ar[u] \ar[r] & k' \ar[u]
}
$$
其中 $k'/k$、$K'/K$ 是有限纯不可分域扩张，且 $K'/k'$ 是可分域扩张。
在此情形下，可以假设 $K' = k'K$ 是复合域，
并且还可以假设 $K' = (k' \otimes_k K)_{red}$。
\end{lemma}

\begin{proof}
由引理 \ref{algebra-lemma-make-separably-generated}，可找到一个图表，使
$K'/k'$ 可分生成。
由引理 \ref{algebra-lemma-separably-generated-separable}，这推出 $K'$ 在 $k'$ 上可分。
复合域 $k'K$ 是 $K'/k'$ 的子扩张，因此
$k' \subset k'K$ 由引理 \ref{algebra-lemma-subextensions-are-separable} 可知是可分的。
环 $(k' \otimes_k K)_{red}$ 是整环，因为对某个
$n \gg 0$，映射 $x \mapsto x^{p^n}$ 将其映入 $K$。
因此由引理 \ref{algebra-lemma-integral-over-field}，它是域。
于是 $(k' \otimes_k K)_{red} \to K'$ 将其同构地映到 $k'K$ 上。
\end{proof}

\begin{lemma}
\label{algebra-lemma-perfection}
\begin{slogan}
每个域都有唯一的完美闭包。
\end{slogan}
对每个域 $k$，存在纯不可分扩张 $k'/k$，使得 $k'$ 是完美的。
域扩张 $k'/k$ 在唯一同构意义下是唯一的。
\end{lemma}

\begin{proof}
若 $k$ 的特征为零，则 $k' = k$ 是唯一选择。假设 $k$ 的特征为 $p > 0$。
对每个 $n > 0$，存在唯一的代数扩张
$k \subset k^{1/p^n}$，使得 (a) 每个元素 $\lambda \in k$
在 $k^{1/p^n}$ 中都有 $p^n$ 次根，并且 (b) 对每个元素
$\mu \in k^{1/p^n}$，都有 $\mu^{p^n} \in k$。
具体地，考虑环映射 $k \to k^{1/p^n} = k$，$x \mapsto x^{p^n}$。
该映射是单射，并满足 (a) 和 (b)。显然，作为经由映射
$y \mapsto y^p$ 的 $k$-扩张，有
$k^{1/p^n} \subset k^{1/p^{n + 1}}$。于是可取
$k' = \bigcup k^{1/p^n}$。
省略一些细节。
\end{proof}

\begin{definition}
\label{algebra-definition-perfection}
令 $k$ 为域。引理 \ref{algebra-lemma-perfection} 中的域扩张 $k'/k$
称为 $k$ 的{\it 完美闭包}。记作 $k^{perf}/k$。
\end{definition}

\noindent
注意，若 $k'/k$ 是任意代数纯不可分扩张，则
$k'$ 是 $k^{perf}$ 的子扩张，即 $k^{perf}/k'/k$。事实上，
由引理 \ref{algebra-lemma-perfection} 的唯一性，$(k')^{perf}$ 与 $k^{perf}$ 同构。

\begin{lemma}
\label{algebra-lemma-perfect-reduced}
令 $k$ 为完美域。
任意约化 $k$-代数在 $k$ 上都是几何约化的。
令 $R$、$S$ 为 $k$-代数。
假设 $R$ 和 $S$ 都是约化的。
则 $k$-代数 $R \otimes_k S$ 是约化的。
\end{lemma}

\begin{proof}
第一个陈述由引理
\ref{algebra-lemma-geometrically-reduced-finite-purely-inseparable-extension} 得出。
第二个陈述应用第一个陈述以及引理
\ref{algebra-lemma-geometrically-reduced-any-reduced-base-change} 即得。
\end{proof}









\section{万有同胚}
\label{algebra-section-universal-homeomorphism}

\noindent
设 $k'/k$ 为代数纯不可分域扩张。那么，对任意 $k$-代数 $R$，
环映射 $R \to k' \otimes_k R$ 都诱导谱之间的同胚。
其原因是下面稍为一般的引理 \ref{algebra-lemma-p-ring-map}。

\begin{lemma}
\label{algebra-lemma-surjective-locally-nilpotent-kernel}
设 $\varphi : R \to S$ 为满映射，且其核局部幂零。
则 $\varphi$ 诱导谱之间的同胚以及剩余域之间的同构。
对任意环映射 $R \to R'$，环映射
$R' \to R' \otimes_R S$ 也是核局部幂零的满映射。
\end{lemma}

\begin{proof}
由引理 \ref{algebra-lemma-spec-closed}，映射 $\Spec(S) \to \Spec(R)$
是到闭子集 $V(\Ker(\varphi))$ 上的同胚。当然，
$V(\Ker(\varphi)) = \Spec(R)$，因为 $R$ 的每个素理想都包含
$R$ 的每个幂零元。这也推出关于剩余域的陈述。
由张量积的右正合性可知，$\Ker(\varphi)R'$ 是满映射
$R' \to R' \otimes_R S$ 的核。最后的陈述于是由
引理 \ref{algebra-lemma-locally-nilpotent} 得出。
\end{proof}

\begin{lemma}
\label{algebra-lemma-powers-field}
\begin{reference}
\cite[Lemma 3.1.6]{Alper-adequate}
\end{reference}
设 $k'/k$ 为域扩张。下列条件等价：
\begin{enumerate}
\item 对每个 $x \in k'$，存在 $n > 0$，使得 $x^n \in k$；
\item $k' = k$，或者 $k$ 和 $k'$ 的特征为 $p > 0$，并且
$k'/k$ 是纯不可分扩张，或 $k$ 和 $k'$ 都是 $\mathbf{F}_p$ 的代数扩张。
\end{enumerate}
\end{lemma}

\begin{proof}
注意，(2) 中列出的每一种情形都满足 (1)。
因此，假设 $k'/k$ 满足 (1)，并证明我们处于 (2) 的某种情形。
排除 $k = k'$ 的情形后，可以假设 $k' \not = k$。
显然，$k'/k$ 是代数扩张。因此可以假设 $k'/k$ 是非平凡有限扩张。
令 $k'/k'_{sep}/k$ 为域，引理 \ref{fields-lemma-separable-first}
中得到的可分子扩张。须证明 $k = k'_{sep}$，或 $k$ 是
$\mathbf{F}_p$ 的代数扩张。因此可以假设 $k'/k$ 是非平凡有限可分扩张，
而我们须证明 $k$ 是 $\mathbf{F}_p$ 的代数扩张。

\medskip\noindent
取 $x \in k'$，$x \not \in k$。取 $n, m > 0$，使得
$x^n \in k$ 且 $(x + 1)^m \in k$。令 $\overline{k}$ 为 $k$ 的代数闭包。
可以选取嵌入 $\sigma, \tau : k' \to \overline{k}$，使得
$\sigma(x) \not = \tau(x)$。这由域，第
\ref{fields-section-separable-extensions} 节中的讨论得出
（更确切地说，将 $k'$ 换成由 $x$ 生成的 $k$-扩张后，
它由域，引理 \ref{fields-lemma-count-embeddings} 得出）。
于是，对 $\overline{k}$ 中某个 $n$ 次单位根 $\zeta$，有
$\sigma(x) = \zeta \tau(x)$。类似地，对某个 $m$ 次单位根
$\zeta' \in \overline{k}$，有
$\sigma(x + 1) = \zeta' \tau(x + 1)$。
由于 $\sigma(x + 1) \not = \tau(x + 1)$，可知 $\zeta' \not = 1$。于是
$$
\zeta' (\tau(x) + 1) =
\zeta' \tau(x + 1) =
\sigma(x + 1) =
\sigma(x) + 1 =
\zeta \tau(x) + 1
$$
推出
$$
\tau(x) (\zeta' - \zeta) = 1 - \zeta'
$$
因此 $\zeta' \not = \zeta$，且
$$
\tau(x) = (1 - \zeta')/(\zeta' - \zeta)
$$
故 $k'$ 中不属于 $k$ 的每个元素都在素子域上代数。
由于 $k'$ 由 $k'$ 中不属于 $k$ 的元素在素子域上生成，
可知 $k'$（从而 $k$）在素子域上代数。

\medskip\noindent
最后，若 $k$ 的特征为 $0$，则上述结论如下导致矛盾
（鼓励读者自行寻找证明）。对每个有理数 $y$，类似地得到单位根
$\zeta_y$，使得 $\sigma(x + y) = \zeta_y\tau(x + y)$。于是有
$$
\zeta \tau(x) + y = \zeta_y(\tau(x) + y)
$$
并且由上面关于 $\tau(x)$ 的公式可知
$\zeta_y \in \mathbf{Q}(\zeta, \zeta')$。由于数域
$\mathbf{Q}(\zeta, \zeta')$ 仅含有限多个单位根，
可找到两个不同的有理数 $y, y'$，使得
$\zeta_y = \zeta_{y'}$。于是
$$
y - y' =
\sigma(x + y) - \sigma(x + y') =
\zeta_y(\tau(x + y)) - \zeta_{y'}\tau(x + y') = \zeta_y(y - y')
$$
这推出 $\zeta_y = 1$，矛盾。
\end{proof}

\begin{lemma}
\label{algebra-lemma-powers}
设 $\varphi : R \to S$ 为环映射。若
\begin{enumerate}
\item 对任意 $x \in S$，存在 $n > 0$，使得
$x^n$ 属于 $\varphi$ 的像；
\item $\Ker(\varphi)$ 局部幂零，
\end{enumerate}
则 $\varphi$ 诱导谱之间的同胚，并诱导满足
引理 \ref{algebra-lemma-powers-field} 中等价条件的剩余域扩张。
\end{lemma}

\begin{proof}
假设 (1) 和 (2)。设 $\mathfrak q, \mathfrak q'$ 为 $S$ 的素理想，
它们位于 $R$ 的同一素理想 $\mathfrak p$ 上。假设 $x \in S$，且
$x \in \mathfrak q$、$x \not \in \mathfrak q'$。则对所有 $n > 0$，
$x^n \in \mathfrak q$ 且 $x^n \not \in \mathfrak q'$。
若对某个 $n > 0$，有 $x^n = \varphi(y)$，其中 $y \in R$，则
$$
x^n \in \mathfrak q \Rightarrow y \in \mathfrak p \Rightarrow
x^n \in \mathfrak q'
$$
这导致矛盾。因此不存在上述 $x$，从而
$\mathfrak q = \mathfrak q'$；也就是说，谱上的映射是单射。
由假设 (2)，核 $I = \Ker(\varphi)$ 包含于每个素理想中，故作为拓扑空间，
$\Spec(R) = \Spec(R/I)$。由假设 (1)，诱导映射 $R/I \to S$ 是整的，
所以引理 \ref{algebra-lemma-integral-overring-surjective} 表明
$\Spec(S) \to \Spec(R/I)$ 是满射。结合上述结论，
$\Spec(S) \to \Spec(R)$ 是双射。
若任取 $x \in S$，并选取 $y \in R$，使得对某个 $n > 0$，
$\varphi(y) = x^n$，则在上述双射下，开集
$D(x) \subset \Spec(S)$ 对应于开集 $D(y) \subset \Spec(R)$。
因此映射 $\Spec(S) \to \Spec(R)$ 是同胚。

\medskip\noindent
为证明关于剩余域的陈述，设 $\mathfrak q \subset S$ 为位于素理想
$\mathfrak p \subset R$ 上的素理想。令 $x \in \kappa(\mathfrak q)$。
将 $\kappa(\mathfrak q)$ 视为局部环 $S_\mathfrak q$ 的剩余域，
则 $x$ 是某个 $y/z \in S_\mathfrak q$ 的像，其中
$y \in S$、$z \in S$、$z \not \in \mathfrak q$。
选取 $n, m > 0$，使得 $y^n, z^m$ 属于 $\varphi$ 的像。
那么 $x^{nm}$ 是 $(y/z)^{nm} = (y^n)^m/(z^m)^n$ 的剩余类，
而后者属于 $R_\mathfrak p \to S_\mathfrak q$ 的像。
因此 $x^{nm}$ 属于 $\kappa(\mathfrak p) \to \kappa(\mathfrak q)$ 的像。
\end{proof}

\begin{lemma}
\label{algebra-lemma-2-3-ring-map}
设 $\varphi : R \to S$ 为环映射。假设
\begin{enumerate}
\item[(a)] 作为 $R$-代数，$S$ 由满足
$x^2, x^3 \in \varphi(R)$ 的元素 $x$ 生成；
\item[(b)] $\Ker(\varphi)$ 局部幂零。
\end{enumerate}
则 $\varphi$ 诱导剩余域之间的同构以及谱之间的同胚。
对任意环映射 $R \to R'$，环映射 $R' \to R' \otimes_R S$
也满足 (a) 和 (b)。
\end{lemma}

\begin{proof}
假设 (a) 和 (b)。由于 $S$ 在 $R$ 上整，谱上的映射是闭映射；
参见引理 \ref{algebra-lemma-going-up-closed} 和
\ref{algebra-lemma-integral-going-up}。由引理
\ref{algebra-lemma-image-dense-generic-points}，其像稠密。
因此 $\Spec(S) \to \Spec(R)$ 是满射。
若 $\mathfrak q \subset S$ 是位于 $\mathfrak p \subset R$ 上的素理想，
则域扩张 $\kappa(\mathfrak q)/\kappa(\mathfrak p)$ 由元素
$\alpha \in \kappa(\mathfrak q)$ 生成，这些元素的平方和立方
均在 $\kappa(\mathfrak p)$ 中。因此显然
$\alpha \in \kappa(\mathfrak p)$，从而
$\kappa(\mathfrak q) = \kappa(\mathfrak p)$。
若 $\mathfrak q, \mathfrak q'$ 是位于 $\mathfrak p$ 上的两个不同素理想，
则 (a) 中作为 $S$ 的生成元的某个 $x$，在
$\kappa(\mathfrak q) = \kappa(\mathfrak p)$ 和
$\kappa(\mathfrak q') = \kappa(\mathfrak p)$ 中的像必然不同。
这与 $x^2$ 和 $x^3$ 的像都相同这一事实矛盾。
这就证明 $\Spec(S) \to \Spec(R)$ 是单射，因而是同胚
（由已经证明的结论）。

\medskip\noindent
由于 $\varphi$ 诱导谱之间的同胚，它在谱上尤其是满射；
此性质在任意基变换下保持，参见引理
\ref{algebra-lemma-surjective-spec-radical-ideal}。
因此，对任意 $R \to R'$，环映射 $R' \to R' \otimes_R S$ 的核
由幂零元组成；参见引理 \ref{algebra-lemma-image-dense-generic-points}。
换言之，$R' \to R' \otimes_R S$ 满足 (b)。
显然，(a) 在基变换下保持。
\end{proof}

\begin{lemma}
\label{algebra-lemma-help-with-powers}
设 $p$ 为素数，$n, m > 0$ 为两个整数。存在整数 $a$，使得
$(x + y)^{p^a}, p^a(x + y) \in \mathbf{Z}[x^{p^n}, p^nx, y^{p^m}, p^my]$。
\end{lemma}

\begin{proof}
只要 $a \geq n, m$，关于 $p^a(x + y)$ 的结论就是显然的。
事实上，取 $a \gg n, m$。写成
$$
(x + y)^{p^a}  = \sum\nolimits_{i, j \geq 0, i + j = p^a}
{p^a \choose i, j} x^iy^j
$$
对每个满足 $i + j = p^a$ 的 $i, j \geq 0$，写
$i = q p^n + r$，其中 $r \in \{0, \ldots, p^n - 1\}$；并写
$j = q' p^m + r'$，其中 $r' \in \{0, \ldots, p^m - 1\}$。
若
$$
p^{nr + mr'} \text{ divides } {p^a \choose i, j}
$$
成立，则条件
$(x + y)^{p^a} \in \mathbf{Z}[x^{p^n}, p^nx, y^{p^m}, p^my]$
成立。若 $r = r' = 0$，则此整除性成立。若 $r \not = 0$，则写成
$$
{p^a \choose i, j} = \frac{p^a}{i} {p^a - 1 \choose i - 1, j}
$$
由于 $r \not = 0$，有理数 $p^a/i$ 的 $p$-进赋值至少为
$a - (n - 1)$（因为 $i$ 不被 $p^n$ 整除）。
所以在此情形下，${p^a \choose i, j}$ 被 $p^{a - n + 1}$ 整除。
类似地，若 $r' \not = 0$，则 ${p^a \choose i, j}$ 被
$p^{a - m + 1}$ 整除。取 $a = np^n + mp^m + n + m$ 即可。
\end{proof}

\begin{lemma}
\label{algebra-lemma-p-ring-map-field}
设 $k'/k$ 为域扩张，$p$ 为素数。下列条件等价：
\begin{enumerate}
\item 作为 $k$ 的域扩张，$k'$ 由满足下述条件的元素 $x$ 生成：
存在 $n > 0$，使得 $x^{p^n} \in k$ 且 $p^nx \in k$；
\item $k = k'$，或者 $k$ 和 $k'$ 的特征均为 $p$，且
$k'/k$ 是纯不可分扩张。
\end{enumerate}
\end{lemma}

\begin{proof}
令 $x \in k'$。若存在 $n > 0$，使得 $x^{p^n} \in k$ 且
$p^nx \in k$，并且特征不是 $p$，则 $x \in k$。
若特征为 $p$，则 $x^{p^n} \in k$，所以 $x$ 在 $k$ 上纯不可分。
\end{proof}

\begin{lemma}
\label{algebra-lemma-p-ring-map}
设 $\varphi : R \to S$ 为环映射，$p$ 为素数。假设
\begin{enumerate}
\item[(a)] 作为 $R$-代数，$S$ 由满足下述条件的元素 $x$ 生成：
存在 $n > 0$，使得 $x^{p^n} \in \varphi(R)$ 且
$p^nx \in \varphi(R)$；
\item[(b)] $\Ker(\varphi)$ 局部幂零。
\end{enumerate}
则 $\varphi$ 诱导谱之间的同胚，并诱导满足引理
\ref{algebra-lemma-p-ring-map-field} 中等价条件的剩余域扩张。
对任意环映射 $R \to R'$，环映射
$R' \to R' \otimes_R S$ 也满足 (a) 和 (b)。
\end{lemma}

\begin{proof}
假设 (a) 和 (b)。注意，(b) 等价于引理 \ref{algebra-lemma-powers} 的条件 (2)。
令 $T \subset S$ 为由下述元素组成的集合：$x \in S$，且存在整数
$n > 0$，使得 $x^{p^n} , p^n x \in \varphi(R)$。
断言 $T = S$。这将证明引理 \ref{algebra-lemma-powers} 的条件 (1) 成立，
从而 $\varphi$ 诱导谱之间的同胚。
由假设 (a)，只须证明 $T \subset S$ 是 $R$-子代数。
若 $x \in T$ 且 $y \in R$，则显然 $yx \in T$。
假设 $x, y \in T$，并且 $n, m > 0$ 满足
$x^{p^n}, y^{p^m}, p^n x, p^m y \in \varphi(R)$。
则 $(xy)^{p^{n + m}}, p^{n + m}xy \in \varphi(R)$，
因而 $xy \in T$。由引理 \ref{algebra-lemma-help-with-powers}，有 $x + y \in T$，
断言得证。

\medskip\noindent
由于 $\varphi$ 诱导谱之间的同胚，它在谱上尤其是满射；
此性质在任意基变换下保持，参见引理
\ref{algebra-lemma-surjective-spec-radical-ideal}。
因此，对任意 $R \to R'$，环映射 $R' \to R' \otimes_R S$ 的核
由幂零元组成；参见引理 \ref{algebra-lemma-image-dense-generic-points}。
换言之，$R' \to R' \otimes_R S$ 满足 (b)。
显然，(a) 在基变换下保持。
最后，关于剩余域的条件由 (a) 得出，因为 $S$ 作为 $R$-代数的生成元
映到各剩余域扩张的生成元。
\end{proof}

\begin{lemma}
\label{algebra-lemma-radicial}
设 $\varphi : R \to S$ 为环映射。假设
\begin{enumerate}
\item $\varphi$ 诱导谱上的单射；
\item $\varphi$ 诱导纯不可分剩余域扩张。
\end{enumerate}
则对任意环映射 $R \to R'$，映射 $R' \to R' \otimes_R S$
也满足性质 (1) 和 (2)。
\end{lemma}

\begin{proof}
置 $S' = R' \otimes_R S$，于是有环谱的连续映射的交换图
$$
\xymatrix{
\Spec(S') \ar[r] \ar[d] & \Spec(S) \ar[d] \\
\Spec(R') \ar[r] & \Spec(R)
}
$$
设 $\mathfrak p' \subset R'$ 为位于 $\mathfrak p \subset R$ 上的素理想。
若 $S$ 中没有位于 $\mathfrak p$ 上的素理想，则 $S'$ 中也没有位于
$\mathfrak p'$ 上的素理想。否则，由注记 \ref{algebra-remark-fundamental-diagram}，
$F = S \otimes_R \kappa(\mathfrak p)$ 有唯一的素理想 $\mathfrak r$，
且其剩余域在 $\kappa(\mathfrak p)$ 上纯不可分。考虑环映射
$$
\kappa(\mathfrak p) \to F \to \kappa(\mathfrak r)
$$
由引理 \ref{algebra-lemma-minimal-prime-reduced-ring}，理想
$\mathfrak r \subset F$ 局部幂零，所以可将引理
\ref{algebra-lemma-surjective-locally-nilpotent-kernel} 应用于环映射
$F \to \kappa(\mathfrak r)$。还可将引理 \ref{algebra-lemma-p-ring-map}
应用于环映射 $\kappa(\mathfrak p) \to \kappa(\mathfrak r)$。
因此，在映射
$$
\kappa(\mathfrak p') \to
\kappa(\mathfrak p') \otimes_{\kappa(\mathfrak p)} F \to
\kappa(\mathfrak p') \otimes_{\kappa(\mathfrak p)} \kappa(\mathfrak r)
$$
中，复合映射和第二个箭头均诱导谱上的双射和纯不可分剩余域扩张。
这推出第一个映射也具有同样的性质。由于
$$
\kappa(\mathfrak p') \otimes_{\kappa(\mathfrak p)} F =
\kappa(\mathfrak p') \otimes_{\kappa(\mathfrak p)}
\kappa(\mathfrak p) \otimes_R S =
\kappa(\mathfrak p') \otimes_R S =
\kappa(\mathfrak p') \otimes_{R'} R' \otimes_R S
$$
由注记 \ref{algebra-remark-fundamental-diagram} 中的讨论可得结论。
\end{proof}

\begin{lemma}
\label{algebra-lemma-radicial-integral}
设 $\varphi : R \to S$ 为环映射。假设
\begin{enumerate}
\item $\varphi$ 是整的；
\item $\varphi$ 诱导谱上的单射；
\item $\varphi$ 诱导纯不可分剩余域扩张。
\end{enumerate}
则 $\varphi$ 诱导从 $\Spec(S)$ 到 $\Spec(R)$ 的某个闭子集上的同胚，
并且对任意环映射 $R \to R'$，映射 $R' \to R' \otimes_R S$
满足性质 (1)、(2)、(3)。
\end{lemma}

\begin{proof}
由引理 \ref{algebra-lemma-going-up-closed} 和
\ref{algebra-lemma-integral-going-up}，谱上的映射是闭映射。
由引理 \ref{algebra-lemma-radicial} 和 \ref{algebra-lemma-base-change-integral}，
这些性质在基变换下保持。
\end{proof}

\begin{lemma}
\label{algebra-lemma-radicial-integral-bijective}
设 $\varphi : R \to S$ 为环映射。假设
\begin{enumerate}
\item $\varphi$ 是整的；
\item $\varphi$ 诱导谱上的双射；
\item $\varphi$ 诱导纯不可分剩余域扩张。
\end{enumerate}
则 $\varphi$ 诱导谱之间的同胚，并且对任意环映射
$R \to R'$，映射 $R' \to R' \otimes_R S$ 满足性质 (1)、(2)、(3)。
\end{lemma}

\begin{proof}
由引理 \ref{algebra-lemma-radicial-integral} 和
\ref{algebra-lemma-surjective-spec-radical-ideal} 得出。
\end{proof}

\begin{lemma}
\label{algebra-lemma-universally-bijective}
设 $\varphi : R \to S$ 为满足下列条件的环映射：
\begin{enumerate}
\item $\varphi$ 的核局部幂零；
\item 作为 $R$-代数，$S$ 由满足下述条件的元素 $x$ 生成：
存在 $n > 0$ 及多项式 $P(T) \in R[T]$，使其在 $S[T]$ 中的像为
$(T - x)^n$。
\end{enumerate}
则 $\Spec(S) \to \Spec(R)$ 是同胚，并且 $R \to S$
诱导纯不可分剩余域扩张。此外，条件 (1) 和 (2) 在任意基变换下保持。
\end{lemma}

\begin{proof}
可以将 $R$ 替换为 $R/\Ker(\varphi)$；参见引理
\ref{algebra-lemma-surjective-locally-nilpotent-kernel}。
假设 (2) 表明，$S$ 由在 $R$ 上整的元素作为 $R$-代数生成。
因此 $R \subset S$ 是整的
（引理 \ref{algebra-lemma-integral-closure-is-ring}）。特别地，
$\Spec(S) \to \Spec(R)$ 是满的闭映射
（引理 \ref{algebra-lemma-integral-overring-surjective}、
\ref{algebra-lemma-going-up-closed} 和 \ref{algebra-lemma-integral-going-up}）。

\medskip\noindent
设 $x \in S$ 为 (2) 中的某个生成元；也就是说，存在
$n > 0$，使得 $(T - x)^n \in R[T]$。设 $\mathfrak p \subset R$ 为素理想。
由上述结论及引理 \ref{algebra-lemma-in-image}，环
$\kappa(\mathfrak p) \otimes_R S$ 非零。
若 $\kappa(\mathfrak p)$ 的特征为零，则 $nx \in R$ 表明
$1 \otimes x$ 属于 $\kappa(\mathfrak p) \to
\kappa(\mathfrak p) \otimes_R S$ 的像。因此
$\kappa(\mathfrak p) \to \kappa(\mathfrak p) \otimes_R S$ 是同构。
若 $\kappa(\mathfrak p)$ 的特征为 $p > 0$，则写
$n = p^k m$，其中 $m$ 与 $p$ 互素。在
$\kappa(\mathfrak p) \otimes_R S[T]$ 中有
$$
(T - 1 \otimes x)^n = ((T - 1 \otimes x)^{p^k})^m =
(T^{p^k} - 1 \otimes x^{p^k})^m
$$
因而可知 $mx^{p^k} \in R$。这表明 $1 \otimes x^{p^k}$ 属于
$\kappa(\mathfrak p) \to \kappa(\mathfrak p) \otimes_R S$ 的像。
所以可将引理 \ref{algebra-lemma-p-ring-map} 应用于
$\kappa(\mathfrak p) \to \kappa(\mathfrak p) \otimes_R S$。
在两种情形下，都可断定 $\kappa(\mathfrak p) \otimes_R S$
有唯一的素理想，且其剩余域在 $\kappa(\mathfrak p)$ 上纯不可分。
由注记 \ref{algebra-remark-fundamental-diagram} 可知 $\varphi$ 在谱上为双射。

\medskip\noindent
关于基变换的陈述是立即的。
\end{proof}











\section{几何不可约代数}
\label{algebra-section-algebras-over-fields}

\noindent
若对每个域扩张 $k'/k$，代数 $S \otimes_k k'$ 都有唯一的极小素理想，
则称域 $k$ 上的代数 $S$ 几何不可约。本节将发展与此概念有关的一些理论。

\begin{lemma}
\label{algebra-lemma-flat-fibres-irreducible}
设 $R \to S$ 为环映射。假设
\begin{enumerate}
\item[(a)] $\Spec(R)$ 不可约；
\item[(b)] $R \to S$ 平坦；
\item[(c)] $R \to S$ 是有限表示的；
\item[(d)] 对 $R$ 中稠密的一族素理想 $\mathfrak p$，纤维环
$S \otimes_R \kappa(\mathfrak p)$ 的谱不可约。
\end{enumerate}
则 $\Spec(S)$ 不可约。更一般地，将 (b) $+$ (c) 替换为
“映射 $\Spec(S) \to \Spec(R)$ 是开映射”，结论仍成立。
\end{lemma}

\begin{proof}
假设 (b) 和 (c) 表明谱上的映射是开映射；参见命题
\ref{algebra-proposition-fppf-open}。因此，本引理由拓扑，引理
\ref{topology-lemma-irreducible-on-top} 得出。
\end{proof}

\begin{lemma}
\label{algebra-lemma-separably-closed-irreducible}
设 $k$ 为可分闭域，$R$、$S$ 为 $k$-代数。
若 $R$、$S$ 各有唯一的极小素理想，则 $R \otimes_k S$ 也如此。
\end{lemma}

\begin{proof}
设 $k \subset \overline{k}$ 为完美闭包；参见定义
\ref{algebra-definition-perfection}。由假设，$\overline{k}$ 是代数闭域。
环映射 $R \to R \otimes_k \overline{k}$、
$S \to S \otimes_k \overline{k}$ 以及
$R \otimes_k S \to (R \otimes_k S) \otimes_k \overline{k}
= (R \otimes_k \overline{k}) \otimes_{\overline{k}} (S \otimes_k \overline{k})$
满足引理 \ref{algebra-lemma-p-ring-map} 的假设。因此可以假设 $k$ 是代数闭域。

\medskip\noindent
可以把 $R$ 和 $S$ 替换为它们的约化。因此可以假设 $R$ 和 $S$ 是整环。
由引理 \ref{algebra-lemma-perfect-reduced}，$R \otimes_k S$ 是约化的。
所以其谱可约，当且仅当它含有非零零因子。
由引理 \ref{algebra-lemma-limit-argument}，可约化到如下情形：
$R$ 和 $S$ 是代数闭域 $k$ 上的有限型整环。

\medskip\noindent
注意，环映射 $R \to R \otimes_k S$ 是有限表示的且平坦。
此外，对 $R$ 的每个极大理想 $\mathfrak m$，有
$(R \otimes_k S) \otimes_R R/\mathfrak m \cong S$，因为由希尔伯特零点定理
\ref{algebra-theorem-nullstellensatz}，$k \cong R/\mathfrak m$。
而且，由于 $\Spec(R)$ 是雅可布森空间，极大理想之集在 $R$ 的谱中稠密；
参见引理 \ref{algebra-lemma-finite-type-field-Jacobson}。
因此，可将引理 \ref{algebra-lemma-flat-fibres-irreducible} 应用于环映射
$R \to R \otimes_k S$，从而如所需，$R \otimes_k S$ 的谱不可约。
\end{proof}

\begin{lemma}
\label{algebra-lemma-geometrically-irreducible}
设 $k$ 为域，$R$ 为 $k$-代数。下列条件等价：
\begin{enumerate}
\item 对每个域扩张 $k'/k$，$R \otimes_k k'$ 的谱不可约；
\item 对每个有限可分域扩张 $k'/k$，$R \otimes_k k'$ 的谱不可约；
\item $R \otimes_k \overline{k}$ 的谱不可约，其中 $\overline{k}$
是 $k$ 的可分代数闭包；
\item $R \otimes_k \overline{k}$ 的谱不可约，其中 $\overline{k}$
是 $k$ 的代数闭包。
\end{enumerate}
\end{lemma}

\begin{proof}
显然，(1) 推出 (2)。

\medskip\noindent
假设 (2)，并令 $\overline{k}$ 为 $k$ 的可分代数闭包。
假设 $\mathfrak q_i \subset R \otimes_k \overline{k}$，$i = 1, 2$
是两个极小素理想。对每个有限子扩张 $\overline{k}/k'/k$，
扩张 $k'/k$ 可分，并且环映射
$R \otimes_k k' \to R \otimes_k \overline{k}$ 平坦。
所以 $\mathfrak p_i = (R \otimes_k k') \cap \mathfrak q_i$
是极小素理想（因为由引理 \ref{algebra-lemma-flat-going-down}，
平坦环映射满足下降）。于是由假设 (2)，
$\mathfrak p_1 = \mathfrak p_2$。由于 $\overline{k} = \bigcup k'$，
可知 $\mathfrak q_1 = \mathfrak q_2$。因此
$\Spec(R \otimes_k \overline{k})$ 不可约。

\medskip\noindent
假设 (3)，并令 $\overline{k}$ 为 $k$ 的代数闭包。
设 $\overline{k}/\overline{k}'/k$ 为相应的 $k$ 的可分代数闭包。
则 $\overline{k}/\overline{k}'$ 是纯不可分扩张（在正特征情形），
或是平凡扩张。因此，例如由引理 \ref{algebra-lemma-p-ring-map}，映射
$R \otimes_k \overline{k}' \to R \otimes_k \overline{k}$
诱导谱之间的同胚。故 (4) 成立。

\medskip\noindent
假设 (4)。设 $k'/k$ 为任意域扩张，并令 $\overline{k}$ 为 $k$ 的代数闭包。
可以选取一个域 $F$，使得 $k'$ 和 $\overline{k}$ 都同构于 $F$ 的子域。于是
$$
R \otimes_k F = (R \otimes_k \overline{k}) \otimes_{\overline{k}} F
$$
故由引理 \ref{algebra-lemma-separably-closed-irreducible}，
$R \otimes_k F$ 有唯一的极小素理想。最后，环映射
$R \otimes_k k' \to R \otimes_k F$ 平坦且单射，因此
$R \otimes_k k'$ 的任一极小素理想都是 $R \otimes_k F$ 的
某个极小素理想的像（由引理
\ref{algebra-lemma-injective-minimal-primes-in-image} 和下降）。
所以这样的极小素理想只有一个，证明完毕。
\end{proof}

\begin{definition}
\label{algebra-definition-geometrically-irreducible}
设 $k$ 为域，$S$ 为 $k$-代数。若对每个域扩张 $k'/k$，
$S \otimes_k k'$ 的谱不可约\footnote{不可约空间非空。}，
则称 $S$ {it 在 $k$ 上几何不可约}。
\end{definition}

\noindent
由引理 \ref{algebra-lemma-geometrically-irreducible}，只须对有限可分域扩张
$k'/k$ 检验这一点，或者取 $k'$ 等于 $k$ 的可分代数闭包加以检验。

\begin{lemma}
\label{algebra-lemma-separably-closed-irreducible-implies-geometric}
设 $k$ 为域，$R$ 为 $k$-代数。若 $k$ 是可分闭域，则
$R$ 在 $k$ 上几何不可约，当且仅当 $R$ 的谱不可约。
\end{lemma}

\begin{proof}
由定义 \ref{algebra-definition-geometrically-irreducible} 后的注记立即得到。
\end{proof}

\begin{lemma}
\label{algebra-lemma-subalgebra-geometrically-irreducible}
设 $k$ 为域，$S$ 为 $k$-代数。
\begin{enumerate}
\item 若 $S$ 在 $k$ 上几何不可约，则每个 $k$-子代数也几何不可约。
\item 若 $S$ 的所有有限生成 $k$-子代数都几何不可约，
则 $S$ 几何不可约。
\item 几何不可约 $k$-代数的有向余极限几何不可约。
\end{enumerate}
\end{lemma}

\begin{proof}
设 $S' \subset S$ 为子代数。那么，对任意扩张 $k'/k$，环映射
$S' \otimes_k k' \to S \otimes_k k'$ 仍是单射。
故 (1) 由引理 \ref{algebra-lemma-injective-minimal-primes-in-image} 得出
（并使用连续映射下不可约空间的像不可约这一事实）。
第二个和第三个性质由张量积与余极限交换这一事实得出。
\end{proof}

\begin{lemma}
\label{algebra-lemma-geometrically-irreducible-any-base-change}
设 $k$ 为域，$S$ 为几何不可约的 $k$-代数，$R$ 为任意 $k$-代数。映射
$$
\Spec(R \otimes_k S) \longrightarrow \Spec(R)
$$
诱导不可约分支之间的双射。
\end{lemma}

\begin{proof}
回忆，不可约分支对应于极小素理想
（引理 \ref{algebra-lemma-irreducible}）。由于 $R \to R \otimes_k S$ 平坦，
由下降（引理 \ref{algebra-lemma-flat-going-down}），
$R \otimes_k S$ 的任一极小素理想都位于 $R$ 的某个极小素理想上。
反之，若 $\mathfrak p \subset R$ 为（极小）素理想，则由
$R \to R \otimes_k S$ 的平坦性，有
$$
R \otimes_k S/\mathfrak p(R \otimes_k S)
=
(R/\mathfrak p) \otimes_k S
\subset
\kappa(\mathfrak p) \otimes_k S
$$
由假设，环 $\kappa(\mathfrak p) \otimes_k S$ 的谱不可约。
因此，$R \otimes_k S/\mathfrak p(R \otimes_k S)$ 有唯一的极小素理想
（引理 \ref{algebra-lemma-injective-minimal-primes-in-image}）。
换言之，不可约集 $V(\mathfrak p)$ 的逆像不可约。本引理因而得证。
\end{proof}

\noindent
下面对几何不可约域扩张的概念作一些注记。

\begin{lemma}
\label{algebra-lemma-field-extension-geometrically-irreducible}
设 $K/k$ 为域扩张。若 $k$ 在 $K$ 中代数闭，则 $K$ 在 $k$ 上几何不可约。
\end{lemma}

\begin{proof}
假设 $k$ 在 $K$ 中代数闭。由定义
\ref{algebra-definition-geometrically-irreducible} 和引理
\ref{algebra-lemma-geometrically-irreducible}，只须证明：
对每个有限可分扩张 $k'/k$，$K \otimes_k k'$ 的谱不可约。
设 $k'$ 由 $\alpha \in k'$ 在 $k$ 上生成；参见域，引理
\ref{fields-lemma-primitive-element}。设
$P = T^d + a_1 T^{d - 1} + \ldots + a_d \in k[T]$ 为 $\alpha$ 的极小多项式。
则 $K \otimes_k k' \cong K[T]/(P)$。
$K[T]/(P)$ 的谱可约的唯一可能，是 $P$ 在 $K[T]$ 中可约。
假设 $P = P_1 P_2$ 是 $K[T]$ 中的非平凡因式分解，以导出矛盾。
由引理 \ref{algebra-lemma-polynomials-divide}，$P_1$ 和 $P_2$ 的系数在 $k$ 上代数。
我们的假设表明 $P_1$ 和 $P_2$ 的系数属于 $k$，
这与 $P$ 在 $k$ 上不可约矛盾。
\end{proof}

\begin{lemma}
\label{algebra-lemma-geometrically-irreducible-transitive}
设 $K/k$ 为几何不可约域扩张，$S$ 为几何不可约的 $K$-代数。
则 $S$ 在 $k$ 上几何不可约。
\end{lemma}

\begin{proof}
由定义 \ref{algebra-definition-geometrically-irreducible} 和引理
\ref{algebra-lemma-geometrically-irreducible}，只须证明：
对每个有限可分扩张 $k'/k$，$S \otimes_k k'$ 的谱不可约。
由于 $K$ 在 $k$ 上几何不可约，可知 $K' = K \otimes_k k'$
是 $K$ 的有限可分域扩张。因为假设 $S$ 在 $K$ 上几何不可约，
所以 $S \otimes_k k' = S \otimes_K K'$ 的谱不可约。
\end{proof}

\begin{lemma}
\label{algebra-lemma-geometrically-irreducible-base-change-transcendental}
设 $K/k$ 为域扩张。下列条件等价：
\begin{enumerate}
\item $K$ 在 $k$ 上几何不可约；
\item 纯超越扩张所诱导的扩张 $K(t)/k(t)$ 几何不可约。
\end{enumerate}
\end{lemma}

\begin{proof}
假设 (1)。以 $\Omega$ 表示 $k(t)$ 的一个代数闭包。
由定义 \ref{algebra-definition-geometrically-irreducible}，可知
$$
K \otimes_k \Omega = K \otimes_k k(t) \otimes_{k(t)} \Omega
$$
的谱不可约。由于 $K(t)$ 是 $K \otimes_k k(T)$ 的局部化，
可知 $K(t) \otimes_{k(t)} \Omega$ 的谱不可约。
所以由引理 \ref{algebra-lemma-geometrically-irreducible}，
$K(t)/k(t)$ 几何不可约。

\medskip\noindent
假设 (2)。设 $k'/k$ 为域扩张。须证明 $K \otimes_k k'$ 有唯一的极小素理想。
已知
$$
K(t) \otimes_{k(t)} k'(t)
$$
的谱不可约，也就是说，它有唯一的极小素理想。
由于存在单射
$K \otimes_k k' \to K(t) \otimes_{k(t)} k'(t)$（省略细节），
结论由引理 \ref{algebra-lemma-injective-minimal-primes-in-image} 和
\ref{algebra-lemma-minimal-prime-image-minimal-prime} 得出。
\end{proof}

\begin{lemma}
\label{algebra-lemma-geometrically-irreducible-add-transcendental}
设 $K/L/M$ 为域塔，且 $L/M$ 几何不可约。
设 $x \in K$ 在 $L$ 上超越。则 $L(x)/M(x)$ 几何不可约。
\end{lemma}

\begin{proof}
这由引理 \ref{algebra-lemma-geometrically-irreducible-base-change-transcendental}
得出，因为域 $L(x)$ 和 $M(x)$ 分别是 $L$ 和 $M$ 的纯超越扩张。
\end{proof}

\begin{lemma}
\label{algebra-lemma-geometrically-irreducible-separable-elements}
设 $K/k$ 为域扩张。下列条件等价：
\begin{enumerate}
\item $K/k$ 几何不可约；
\item 每个在 $k$ 上可分代数的元素 $\alpha \in K$ 都属于 $k$。
\end{enumerate}
\end{lemma}

\begin{proof}
假设 (1)，并设 $\alpha \in K$ 在 $k$ 上可分代数。
则 $k' = k(\alpha)$ 是 $k$ 的有限可分扩张，并且包含在 $K$ 中。
由引理 \ref{algebra-lemma-subalgebra-geometrically-irreducible}，
扩张 $k'/k$ 几何不可约。特别地，$k' \otimes_k \overline{k}$
的谱不可约（因而，若它是域的积，则恰有一个因子）。
由域，引理 \ref{fields-lemma-finite-separable-tensor-alg-closed}，
可知 $\Hom_k(k', \overline{k})$ 有一个元素；进而由域，引理
\ref{fields-lemma-separable-equality}，可知 $k' = k$。故 (2) 成立。

\medskip\noindent
假设 (2)。令 $k' \subset K$ 为由在 $k$ 上代数的元素组成的子域。
由引理 \ref{algebra-lemma-field-extension-geometrically-irreducible}，
扩张 $K/k'$ 几何不可约。由假设，$k'/k$ 是纯不可分扩张。
由引理 \ref{algebra-lemma-p-ring-map}，扩张 $k'/k$ 几何不可约。
所以由引理 \ref{algebra-lemma-geometrically-irreducible-transitive}，
$K/k$ 几何不可约。
\end{proof}

\begin{lemma}
\label{algebra-lemma-make-geometrically-irreducible}
设 $K/k$ 为域扩张。考虑由在 $k$ 上可分代数的元素组成的子扩张
$K/k'/k$。则 $K$ 在 $k'$ 上几何不可约。若 $K/k$ 为有限生成域扩张，
则 $[k' : k] < \infty$。
\end{lemma}

\begin{proof}
第一个陈述立即由引理
\ref{algebra-lemma-geometrically-irreducible-separable-elements} 以及下述事实得出：
由可分代数扩张的传递性，在 $k'$ 上可分代数的元素都属于 $k'$；
参见域，引理 \ref{fields-lemma-separable-permanence}。
若 $K/k$ 有限生成，则由域，引理
\ref{fields-lemma-algebraic-closure-in-finitely-generated}，
$k'$ 在 $k$ 上有限。
\end{proof}

\begin{lemma}
\label{algebra-lemma-Galois-orbit}
设 $K/k$ 为域扩张，$\overline{k}/k$ 为可分代数闭包。
则 $\text{Gal}(\overline{k}/k)$ 在
$\overline{k} \otimes_k K$ 的素理想上可迁地作用。
\end{lemma}

\begin{proof}
设 $K/k'/k$ 为引理 \ref{algebra-lemma-make-geometrically-irreducible}
中得到的子扩张。注意，由于 $k \subset \overline{k}$ 是整的，
$\overline{k} \otimes_k K$ 和 $\overline{k} \otimes_k k'$ 的所有素理想
都是极大理想；参见引理 \ref{algebra-lemma-integral-no-inclusion}。
由引理 \ref{algebra-lemma-geometrically-irreducible-any-base-change}，映射
$$
\Spec(\overline{k} \otimes_k K) \to \Spec(\overline{k} \otimes_k k')
$$
是双射，因为 (1) 所有素理想都是极小素理想，(2)
$\overline{k} \otimes_k K = (\overline{k} \otimes_k k') \otimes_{k'} K$，
并且 (3) $K$ 在 $k'$ 上几何不可约。
因此，只须证明 $\text{Gal}(\overline{k}/k)$ 在
$\overline{k} \otimes_k k'$ 的素理想上的作用可迁。

\medskip\noindent
由于 $\overline{k} \otimes_k k'$ 的每个素理想都是极大理想，
其剩余域同构于 $\overline{k}$。因此，
$\overline{k} \otimes_k k'$ 的素理想与
$\Hom_k(k', \overline{k})$ 的元素一一对应，其中
$\sigma \in \Hom_k(k', \overline{k})$ 对应于映射
$1 \otimes \sigma : \overline{k} \otimes_k k' \to \overline{k}$
的核 $\mathfrak p_\sigma$。特别地，
$\text{Gal}(\overline{k}/k)$ 如所需地在此集合上可迁作用。
\end{proof}





\section{几何连通代数}
\label{algebra-section-geometrically-connected}

\begin{lemma}
\label{algebra-lemma-separably-closed-connected}
设 $k$ 为可分闭域，$R$、$S$ 为 $k$-代数。若
$\Spec(R)$ 和 $\Spec(S)$ 连通，则
$\Spec(R \otimes_k S)$ 也连通。
\end{lemma}

\begin{proof}
回忆，$\Spec(R)$ 连通，当且仅当 $R$ 没有非平凡幂等元；
参见引理 \ref{algebra-lemma-characterize-spec-connected}。
所以由引理 \ref{algebra-lemma-limit-argument}，可以假设 $R$ 和 $S$ 都是
$k$ 上的有限型代数。此时 $R$ 和 $S$ 是诺特环，并且只有有限多个极小素理想；
参见引理 \ref{algebra-lemma-Noetherian-irreducible-components}。
因此，可对 $n + m$ 作归纳，其中 $n$、$m$ 分别是
$\Spec(R)$、$\Spec(S)$ 的不可约分支数。
当然，$n$ 或 $m$ 为零的情形是平凡的。若 $n = m = 1$，
即 $\Spec(R)$ 和 $\Spec(S)$ 都只有一个不可约分支，
则结论由引理 \ref{algebra-lemma-separably-closed-irreducible} 得出。
假设 $n > 1$。设 $\mathfrak p \subset R$ 为对应于不可约闭子集
$T \subset \Spec(R)$ 的极小素理想。令 $T' \subset \Spec(R)$ 为其余
$n - 1$ 个不可约分支之并。选取理想 $I \subset R$，使得
$T' = V(I) = \Spec(R/I)$（引理 \ref{algebra-lemma-spec-closed}）。
仔细选取极小素理想，还可以使 $T'$ 连通；参见拓扑，引理
\ref{topology-lemma-remove-irreducible-connected}。那么
$T \cup T' = \Spec(R)$，并且
$T \cap T' = V(\mathfrak p + I) = \Spec(R/(\mathfrak p + I))$
非空，因为假设 $\Spec(R)$ 连通。
$T$ 在 $\Spec(R \otimes_k S)$ 中的逆像是
$\Spec(R/\mathfrak p \otimes_k S)$，而 $T'$ 在
$\Spec(R \otimes_k S)$ 中的逆像是 $\Spec(R/I \otimes_k S)$。
由归纳假设，这两个空间都连通。$T \cap T'$ 的逆像是
$\Spec(R/(\mathfrak p + I) \otimes_k S)$，它非空。
因此 $\Spec(R \otimes_k S)$ 连通。
\end{proof}

\begin{lemma}
\label{algebra-lemma-geometrically-connected}
设 $k$ 为域，$R$ 为 $k$-代数。下列条件等价：
\begin{enumerate}
\item 对每个域扩张 $k'/k$，$R \otimes_k k'$ 的谱连通；
\item 对每个有限可分域扩张 $k'/k$，$R \otimes_k k'$ 的谱连通。
\end{enumerate}
\end{lemma}

\begin{proof}
对任意域扩张 $k'/k$，$R \otimes_k k'$ 的谱连通，
等价于 $R \otimes_k k'$ 没有非平凡幂等元；参见引理
\ref{algebra-lemma-characterize-spec-connected}。假设 (2)。
设 $k \subset \overline{k}$ 为 $k$ 的可分代数闭包。
由引理 \ref{algebra-lemma-limit-argument}，可知 (2) 等价于
$R \otimes_k \overline{k}$ 没有非平凡幂等元。
对任意域扩张 $k'/k$，存在域扩张 $\overline{k}'/\overline{k}$，
使得 $k' \subset \overline{k}'$。由引理
\ref{algebra-lemma-separably-closed-connected}，
$R \otimes_k \overline{k}'$ 没有非平凡幂等元。
若 $R \otimes_k k'$ 有非平凡幂等元，则
$R \otimes_k \overline{k}'$ 也有，矛盾。
\end{proof}

\begin{definition}
\label{algebra-definition-geometrically-connected}
设 $k$ 为域，$S$ 为 $k$-代数。若对每个域扩张 $k'/k$，
$S \otimes_k k'$ 的谱连通，则称 $S$ {it 在 $k$ 上几何连通}。
\end{definition}

\noindent
由引理 \ref{algebra-lemma-geometrically-connected}，只须对有限可分域扩张
$k'/k$ 检验这一点。

\begin{lemma}
\label{algebra-lemma-separably-closed-connected-implies-geometric}
设 $k$ 为域，$R$ 为 $k$-代数。若 $k$ 是可分闭域，则
$R$ 在 $k$ 上几何连通，当且仅当 $R$ 的谱连通。
\end{lemma}

\begin{proof}
由定义 \ref{algebra-definition-geometrically-connected} 后的注记立即得到。
\end{proof}

\begin{lemma}
\label{algebra-lemma-subalgebra-geometrically-connected}
设 $k$ 为域，$S$ 为 $k$-代数。
\begin{enumerate}
\item 若 $S$ 在 $k$ 上几何连通，则每个 $k$-子代数也几何连通。
\item 若 $S$ 的所有有限生成 $k$-子代数都几何连通，则 $S$ 几何连通。
\item 几何连通 $k$-代数的有向余极限几何连通。
\end{enumerate}
\end{lemma}

\begin{proof}
这由连通性的如下刻画得出：不存在非平凡幂等元。
第二个和第三个性质由张量积与余极限交换这一事实得出。
\end{proof}

\noindent
下面的引理将被更一般的簇，引理
\ref{varieties-lemma-bijection-connected-components} 所取代。

\begin{lemma}
\label{algebra-lemma-geometrically-connected-any-base-change}
设 $k$ 为域，$S$ 为几何连通的 $k$-代数，$R$ 为任意 $k$-代数。映射
$$
R \longrightarrow R \otimes_k S
$$
诱导幂等元之间的双射，而映射
$$
\Spec(R \otimes_k S) \longrightarrow \Spec(R)
$$
诱导连通分支之间的双射。
\end{lemma}

\begin{proof}
第二个断言由第一个断言结合引理 \ref{algebra-lemma-connected-component} 得出。
由引理 \ref{algebra-lemma-subalgebra-geometrically-connected} 和
\ref{algebra-lemma-limit-argument}，可以假设 $R$ 和 $S$ 都是 $k$ 上的有限型代数。
于是 $R \otimes_k S$ 也是 $k$ 上的有限型代数。
注意，在此情形下，所有这些环都是诺特环，因而它们的谱只有有限多个连通分支
（因为只有有限多个不可约分支；参见引理
\ref{algebra-lemma-Noetherian-irreducible-components}）。
特别地，这里涉及的所有连通分支都是开集！因此，由引理
\ref{algebra-lemma-disjoint-implies-product} 可知，在此情形下，
本引理的第一个陈述等价于第二个陈述。下面证明第二个陈述。
由于代数 $S$ 几何连通且非零，映射
$X = \Spec(R \otimes_k S) \to \Spec(R) = Y$ 的所有纤维都连通且非空。
此外，由于 $R \to R \otimes_k S$ 是有限表示的平坦映射，
映射 $X \to Y$ 是开映射
（命题 \ref{algebra-proposition-fppf-open}）。拓扑，引理
\ref{topology-lemma-connected-fibres-connected-components} 表明
$X \to Y$ 诱导连通分支之间的双射。
\end{proof}




\section{几何整代数}
\label{algebra-section-geometrically-integral}

\noindent
定义如下。

\begin{definition}
\label{algebra-definition-geometrically-integral}
设 $k$ 为域，$S$ 为 $k$-代数。若对每个域扩张 $k'/k$，
环 $S \otimes_k k'$ 都是整环，则称 $S$ {it 在 $k$ 上几何整}。
\end{definition}

\noindent
关于几何整代数的任何问题都可以转化为关于几何约化且不可约代数的问题。

\begin{lemma}
\label{algebra-lemma-geometrically-integral}
设 $k$ 为域，$S$ 为 $k$-代数。此时，$S$ 在 $k$ 上几何整，
当且仅当 $S$ 在 $k$ 上既几何不可约又几何约化。
\end{lemma}

\begin{proof}
省略。
\end{proof}

\begin{lemma}
\label{algebra-lemma-characterize-geometrically-integral}
设 $k$ 为域，$S$ 为 $k$-代数。下列条件等价：
\begin{enumerate}
\item $S$ 在 $k$ 上几何整；
\item 对每个有限域扩张 $k'/k$，环 $S \otimes_k k'$ 是整环；
\item $S \otimes_k \overline{k}$ 是整环，其中 $\overline{k}$
是 $k$ 的代数闭包。
\end{enumerate}
\end{lemma}

\begin{proof}
由引理 \ref{algebra-lemma-geometrically-integral}、
\ref{algebra-lemma-geometrically-reduced-finite-purely-inseparable-extension} 和
\ref{algebra-lemma-geometrically-irreducible} 得出。
\end{proof}

\begin{lemma}
\label{algebra-lemma-geometrically-integral-any-integral-base-change}
设 $k$ 为域，$S$ 为几何整的 $k$-代数。
设 $R$ 为 $k$-代数且为整环。则 $R \otimes_k S$ 是整环。
\end{lemma}

\begin{proof}
由引理 \ref{algebra-lemma-geometrically-reduced-any-reduced-base-change}，
环 $R \otimes_k S$ 是约化的；由引理
\ref{algebra-lemma-geometrically-irreducible-any-base-change}，
环 $R \otimes_k S$ 不可约（其谱只有一个不可约分支），
所以 $R \otimes_k S$ 是整环。
\end{proof}





\section{赋值环}
\label{algebra-section-valuation-rings}

\noindent
先给出一些定义。

\begin{definition}
\label{algebra-definition-valuation-ring}
赋值环：
\begin{enumerate}
\item 设 $K$ 为域，$A$、$B$ 为包含于 $K$ 中的局部环。
若 $A \subset B$ 且 $\mathfrak m_A = A \cap \mathfrak m_B$，
则称 $B$ {it 支配} $A$。
\item 设 $A$ 为环。若 $A$ 是局部整环，并且在 $A$ 的分式域所含局部环中，
$A$ 关于支配关系是极大的，则称 $A$ 为{it 赋值环}。
\item 设 $A$ 为分式域为 $K$ 的赋值环。若 $R \subset K$ 是 $K$ 的子环，
且 $R \subset A$，则称 $A$ {it 以 $R$ 为中心}。
\end{enumerate}
\end{definition}

\noindent
依此定义，域是赋值环。

\begin{lemma}
\label{algebra-lemma-dominate}
设 $K$ 为域，$A \subset K$ 为局部子环。则存在分式域为 $K$ 且支配 $A$ 的赋值环。
\end{lemma}

\begin{proof}
用支配关系将 $K$ 的局部子环之集视为偏序集。
假设 $\{A_i\}_{i \in I}$ 是 $K$ 的一族全序局部子环。
则 $B = \bigcup A_i$ 是支配所有 $A_i$ 的局部子环。
因此，由佐恩引理，只须证明：若 $A \subset K$ 是分式域不等于 $K$ 的局部环，
则存在局部环 $B \subset K$，$B \not = A$，且它支配 $A$。

\medskip\noindent
取不属于 $A$ 的分式域的 $t \in K$。
若 $t$ 在 $A$ 上超越，则 $A[t] \subset K$，从而
$A[t]_{(t, \mathfrak m)} \subset K$ 是不同于 $A$ 且支配 $A$ 的局部环。
假设 $t$ 在 $A$ 上代数。则对某个非零 $a \in A$，元素 $at$ 在 $A$ 上整。
此时，由 $A$ 和 $ta$ 生成的子环 $A' \subset K$ 在 $A$ 上有限。
由引理 \ref{algebra-lemma-integral-overring-surjective}，存在位于
$\mathfrak m$ 上的素理想 $\mathfrak m' \subset A'$。
于是 $A'_{\mathfrak m'}$ 支配 $A$。若 $A = A'_{\mathfrak m'}$，
则 $t$ 属于 $A$ 的分式域，这与假设矛盾。因此如所需，
$A \not = A'_{\mathfrak m'}$。
\end{proof}

\begin{lemma}
\label{algebra-lemma-valuation-ring-normal}
设 $A$ 为赋值环。则 $A$ 是正规整环。
\end{lemma}

\begin{proof}
假设 $x$ 属于 $A$ 的分式域，且在 $A$ 上整。
以 $A'$ 表示由 $A$ 和 $x$ 生成的 $K$ 的子环。
由于 $A\subset A'$ 是整扩张，由引理
\ref{algebra-lemma-integral-overring-surjective}，存在位于
$\mathfrak m$ 上的素理想 $\mathfrak m' \subset A'$。
于是 $A'_{\mathfrak m'}$ 支配 $A$。由于 $A$ 是赋值环，
可知 $A=A'_{\mathfrak m'}$，从而 $x\in A$。
\end{proof}


\begin{lemma}
\label{algebra-lemma-valuation-ring-x-or-x-inverse}
设 $A$ 为极大理想是 $\mathfrak m$、分式域是 $K$ 的赋值环。
令 $x \in K$。则 $x \in A$ 或 $x^{-1} \in A$，或二者皆成立。
\end{lemma}

\begin{proof}
假设 $x$ 不属于 $A$。以 $A'$ 表示由 $A$ 和 $x$ 生成的 $K$ 的子环。
由于 $A$ 是赋值环，$A'$ 中没有位于 $\mathfrak m$ 上的素理想。
因为 $\mathfrak m$ 是极大理想，可知 $V(\mathfrak m A') = \emptyset$。
于是由引理 \ref{algebra-lemma-Zariski-topology}，$\mathfrak m A' = A'$。
故可写成 $1 = \sum_{i = 0}^d t_i x^i$，其中 $t_i \in \mathfrak m$。
这推出 $(1 - t_0) (x^{-1})^d - \sum t_i (x^{-1})^{d - i} = 0$。
特别地，$x^{-1}$ 在 $A$ 上整，所以由引理
\ref{algebra-lemma-valuation-ring-normal}，$x^{-1} \in A$。
\end{proof}

\begin{lemma}
\label{algebra-lemma-x-or-x-inverse-valuation-ring}
设 $A \subset K$ 为域 $K$ 的子环，并且对所有 $x \in K$，
$x \in A$ 或 $x^{-1} \in A$，或二者皆成立。
则 $A$ 是分式域为 $K$ 的赋值环。
\end{lemma}

\begin{proof}
若 $A$ 不等于 $K$，则 $A$ 不是域，并且存在非零极大理想
$\mathfrak m$。若 $\mathfrak m'$ 是另一个极大理想，则选取
$x, y \in A$，使得 $x \in \mathfrak m$、$y \not \in \mathfrak m$、
$x \not \in \mathfrak m'$ 且 $y \in \mathfrak m'$。
那么 $x/y \in A$ 和 $y/x \in A$ 都不成立，这与本引理的假设矛盾。
所以 $A$ 是局部环。假设 $A'$ 是包含于 $K$ 中且支配 $A$ 的局部环。
令 $x \in A'$。须证明 $x \in A$。若不然，则 $x^{-1} \in A$，
而且当然 $x^{-1} \in \mathfrak m_A$。但是这样就有
$x^{-1} \in \mathfrak m_{A'}$，与 $x \in A'$ 矛盾。
\end{proof}

\begin{lemma}
\label{algebra-lemma-colimit-valuation-rings}
\begin{slogan}
赋值环在滤过直极限下稳定
\end{slogan}
设 $I$ 为有向集，$(A_i, \varphi_{ij})$ 为 $I$ 上的赋值环系统。
则 $A = \colim A_i$ 是赋值环。
\end{lemma}

\begin{proof}
显然 $A$ 是整环。设 $a, b \in A$。
引理 \ref{algebra-lemma-x-or-x-inverse-valuation-ring} 表明，须证明在 $A$ 中
$a | b$ 或 $b | a$。选取充分大的 $i$，使得存在映到 $a, b$ 的
$a_i, b_i \in A_i$。将引理 \ref{algebra-lemma-valuation-ring-x-or-x-inverse}
应用于 $A_i$ 中的 $a_i, b_i$，即得 $A$ 中 $a, b$ 的结论。
\end{proof}

\begin{lemma}
\label{algebra-lemma-valuation-ring-cap-field}
设 $L/K$ 为域扩张。若 $B \subset L$ 是赋值环，
则 $A = K \cap B$ 是赋值环。
\end{lemma}

\begin{proof}
可以把 $L$ 替换为 $B$ 的分式域 $F$，并把 $K$ 替换为
$K \cap F$。然后本引理由引理
\ref{algebra-lemma-valuation-ring-x-or-x-inverse} 和
\ref{algebra-lemma-x-or-x-inverse-valuation-ring} 结合得出。
\end{proof}

\begin{lemma}
\label{algebra-lemma-valuation-ring-cap-field-finite}
设 $L/K$ 为代数域扩张。若 $B \subset L$ 是分式域为 $L$
且不是域的赋值环，则 $A = K \cap B$ 是赋值环且不是域。
\end{lemma}

\begin{proof}
由引理 \ref{algebra-lemma-valuation-ring-cap-field}，环 $A$ 是赋值环。
若 $A$ 是域，则 $A = K$。于是 $A = K \subset B$ 是整扩张，
故 $B$ 的素理想之间不存在真包含
（引理 \ref{algebra-lemma-integral-no-inclusion}）。
这与 $B$ 是局部整环而不是域的假设矛盾。
\end{proof}

\begin{lemma}
\label{algebra-lemma-make-valuation-rings}
设 $A$ 为赋值环。对任意素理想 $\mathfrak p \subset A$，
商环 $A/\mathfrak p$ 是赋值环。局部化 $A_\mathfrak p$ 也如此；
事实上，$A$ 的任意局部化都是赋值环。
\end{lemma}

\begin{proof}
使用引理 \ref{algebra-lemma-x-or-x-inverse-valuation-ring}
所给出的赋值环刻画。
\end{proof}

\begin{lemma}
\label{algebra-lemma-stack-valuation-rings}
设 $A'$ 为剩余域是 $K$ 的赋值环，$A$ 为分式域是 $K$ 的赋值环。则
$C = \{\lambda \in A' \mid \lambda \bmod \mathfrak m_{A'} \in A\}$
是赋值环。
\end{lemma}

\begin{proof}
注意，$\mathfrak m_{A'} \subset C$ 且 $C/\mathfrak m_{A'} = A$。
特别地，$C$ 的分式域等于 $A'$ 的分式域。
将使用引理 \ref{algebra-lemma-x-or-x-inverse-valuation-ring} 的判据来证明本引理。
令 $x$ 为 $C$ 的分式域中的元素。由该引理，可以假设 $x \in A'$。
若 $x \in \mathfrak m_{A'}$，则 $x \in C$。若非如此，则 $x$ 是 $A'$ 的单位，
并且还有 $x^{-1} \in A'$。因此，再次由该引理，$x$ 或 $x^{-1}$
映到 $A$ 的某个元素。
\end{proof}

\begin{lemma}
\label{algebra-lemma-find-valuation-rings}
设 $A$ 为分式域是 $K$ 的正规整环。
\begin{enumerate}
\item 对每个 $x \in K$，$x \not \in A$，存在分式域是 $K$ 的赋值环
$A \subset V \subset K$，使得 $x \not \in V$。
\item 若 $A$ 是局部环，还可以选取支配 $A$ 的 $V$。
\end{enumerate}
换言之，$A$ 是 $K$ 中所有包含 $A$ 的赋值环之交；若 $A$ 是局部环，
则 $A$ 是 $K$ 中所有支配 $A$ 的赋值环之交。
\end{lemma}

\begin{proof}
假设 $x \in K$，$x \not \in A$。考虑 $B = A[x^{-1}]$。
则 $x \not \in B$。事实上，若
$x = a_0 + a_1x^{-1} + \ldots + a_d x^{-d}$，
则 $x^{d + 1} - a_0x^d - \ldots - a_d = 0$，于是 $x$ 在 $A$ 上整，
这与 $A$ 正规矛盾。因此 $x^{-1}$ 不是 $B$ 的单位。
所以 $V(x^{-1}) \subset \Spec(B)$ 非空
（引理 \ref{algebra-lemma-Zariski-topology}），并且可以选取满足
$x^{-1} \in \mathfrak p$ 的素理想 $\mathfrak p \subset B$。
选取支配 $B_\mathfrak p$ 的赋值环 $V \subset K$
（引理 \ref{algebra-lemma-dominate}）。由于 $x^{-1} \in \mathfrak m_V$，
有 $x \not \in V$。

\medskip\noindent
若 $A$ 是局部环，则断言
$x^{-1} B + \mathfrak m_A B \not = B$。事实上，若
$1 = (a_0 + a_1x^{-1} + \ldots + a_d x^{-d})x^{-1} +
a'_0 + \ldots + a'_d x^{-d}$，其中 $a_i \in A$ 且
$a'_i \in \mathfrak m_A$，则会得到
$$
(1 - a'_0) x^{d + 1} - (a_0 + a'_1) x^d - \ldots - a_d = 0
$$
由于 $a'_0 \in \mathfrak m_A$，可知 $1 - a'_0$ 是 $A$ 的单位，
从而像先前一样，$x$ 会在 $A$ 上整，导致矛盾。
再选取素理想 $\mathfrak p \supset x^{-1} B + \mathfrak m_A B$，
即可得到支配 $A$ 的 $V$。
\end{proof}

\noindent
{it 全序阿贝尔群}是二元组 $(\Gamma, \geq)$，其中 $\Gamma$ 是赋有全序
$\geq$ 的阿贝尔群，并且对所有 $\gamma, \gamma', \gamma'' \in \Gamma$，
$\gamma \geq \gamma' \Rightarrow
\gamma + \gamma'' \geq \gamma' + \gamma''$。

\begin{lemma}
\label{algebra-lemma-valuation-group}
设 $A$ 为分式域是 $K$ 的赋值环。置
$\Gamma = K^*/A^*$（群运算按加法记）。对 $\gamma, \gamma' \in \Gamma$，
定义 $\gamma \geq \gamma'$，当且仅当 $\gamma - \gamma'$ 属于
$A - \{0\} \to \Gamma$ 的像。则 $(\Gamma, \geq)$ 是全序阿贝尔群。
\end{lemma}

\begin{proof}
省略；但它容易由引理 \ref{algebra-lemma-valuation-ring-x-or-x-inverse} 得出。
注意，当 $A = K$ 时，得到赋有唯一全序的零群 $\Gamma = \{0\}$。
\end{proof}

\begin{definition}
\label{algebra-definition-value-group}
设 $A$ 为赋值环。
\begin{enumerate}
\item 引理 \ref{algebra-lemma-valuation-group} 的全序阿贝尔群
$(\Gamma, \geq)$ 称为赋值环 $A$ 的{it 值群}。
\item 映射 $v : A - \{0\} \to \Gamma$ 以及
$v : K^* \to \Gamma$ 都称为与 $A$ 相伴的{it 赋值}。
\item 若 $\Gamma \cong \mathbf{Z}$，则称赋值环 $A$ 为{it 离散赋值环}。
\end{enumerate}
\end{definition}

\noindent
注意，若 $\Gamma \cong \mathbf{Z}$，则满足 $1 \geq 0$ 的这种同构是唯一的。
按此方式选取同构后，该次序即为整数的通常次序。

\begin{lemma}
\label{algebra-lemma-properties-valuation}
设 $A$ 为赋值环。赋值 $v : A -\{0\} \to \Gamma_{\geq 0}$
具有下列性质：
\begin{enumerate}
\item $v(a) = 0 \Leftrightarrow a \in A^*$；
\item $v(ab) = v(a) + v(b)$；
\item 当 $a + b \not = 0$ 时，$v(a + b) \geq \min(v(a), v(b))$。
\end{enumerate}
\end{lemma}

\begin{proof}
省略。
\end{proof}

\begin{lemma}
\label{algebra-lemma-characterize-valuation-ring}
设 $A$ 为环。下列条件等价：
\begin{enumerate}
\item $A$ 是赋值环；
\item $A$ 是局部整环，并且 $A$ 的每个有限生成理想都是主理想。
\end{enumerate}
\end{lemma}

\begin{proof}
假设 $A$ 是赋值环，并设 $f_1, \ldots, f_n \in A$。
选取 $i$，使得 $v(f_i)$ 在所有 $v(f_j)$ 中最小。
则 $(f_i) = (f_1, \ldots, f_n)$。反之，假设 $A$ 是局部整环，
并且 $A$ 的每个有限生成理想都是主理想。
取 $f, g \in A$，并写成 $(f, g) = (h)$。则对某些
$a, b, c, d \in A$，有 $f = ah$、$g = bh$ 且 $h = cf + dg$。
所以 $ac + bd = 1$，从而 $a$ 或 $b$ 是单位；也就是说，
$g/f$ 或 $f/g$ 属于 $A$。由引理
\ref{algebra-lemma-x-or-x-inverse-valuation-ring}，这表明 $A$ 是赋值环。
\end{proof}

\begin{lemma}
\label{algebra-lemma-valuation-valuation-ring}
设 $(\Gamma, \geq)$ 为全序阿贝尔群，$K$ 为域，
$v : K^* \to \Gamma$ 为阿贝尔群同态，并且对满足
$a, b, a + b$ 均非零的 $a, b \in K$，有
$v(a + b) \geq \min(v(a), v(b))$。则
$$
A =
\{
x \in K \mid x = 0 \text{ 或 } v(x) \geq 0
\}
$$
是值群为 $\Im(v) \subset \Gamma$ 的赋值环，其极大理想为
$$
\mathfrak m =
\{
x \in K \mid x = 0 \text{ 或 } v(x) > 0
\}
$$
而其单位群为
$$
A^* =
\{
x \in K^* \mid v(x) = 0
\}.
$$
\end{lemma}

\begin{proof}
省略。
\end{proof}

\noindent
设 $(\Gamma, \geq)$ 为全序阿贝尔群。$\Gamma$ 的{it 理想}是子集
$I \subset \Gamma$，使得 $I$ 的所有元素都 $\geq 0$，并且
$\gamma \in I$、$\gamma' \geq \gamma$ 推出 $\gamma' \in I$。
若这种理想满足 $0 \not \in I$ 以及
$\gamma + \gamma' \in I, \gamma, \gamma' \geq 0
\Rightarrow \gamma \in I \text{ 或 } \gamma' \in I$，
则称它为{it 素理想}。

\begin{lemma}
\label{algebra-lemma-ideals-valuation-ring}
设 $A$ 为赋值环。$A$ 中的理想与 $\Gamma$ 的理想按 $1 - 1$ 对应。
此双射保持包含关系，并将素理想映到素理想。
\end{lemma}

\begin{proof}
省略。
\end{proof}

\begin{lemma}
\label{algebra-lemma-valuation-ring-Noetherian-discrete}
赋值环是诺特环，当且仅当它是离散赋值环或域。
\end{lemma}

\begin{proof}
假设 $A$ 是离散赋值环，其赋值
$v : A \setminus \{0\} \to \mathbf{Z}$ 已正规化，使得
$\Im(v) = \mathbf{Z}_{\geq 0}$。由引理
\ref{algebra-lemma-ideals-valuation-ring}，$A$ 的理想是各子集
$I_n = \{0\} \cup v^{-1}(\mathbf{Z}_{\geq n})$。
显然，满足 $v(x) = n$ 的任意元素 $x \in A$ 都生成 $I_n$。
因此 $A$ 是主理想整环，当然是诺特环。

\medskip\noindent
假设 $A$ 是值群为 $\Gamma$ 的诺特赋值环。
由引理 \ref{algebra-lemma-ideals-valuation-ring}，可知 $\Gamma$ 的理想满足升链条件。
可以假设 $A$ 不是域；也就是说，存在 $\gamma \in \Gamma$，使得
$\gamma > 0$。将升链条件应用于各子集
$\gamma + \Gamma_{\geq 0}$，其中 $\gamma > 0$，可知存在大于 $0$
的最小元素 $\gamma_0$。设 $\gamma \in \Gamma$ 为满足
$\gamma > 0$ 的元素。考虑元素序列 $\gamma$、
$\gamma - \gamma_0$、$\gamma - 2\gamma_0$ 等。
由升链条件，它们不可能全都 $> 0$。设
$\gamma - n \gamma_0$ 是其中最后一个 $\geq 0$ 的元素。
由 $\gamma_0$ 的极小性，$0 = \gamma - n \gamma_0$。
因此如所需，$\Gamma$ 是循环群。
\end{proof}























\section{再论诺特环}
\label{algebra-section-Noetherian-again}


\begin{lemma}
\label{algebra-lemma-Noetherian-basic}
设 $R$ 为诺特环。任意有限 $R$-模都是有限表示的。
有限 $R$-模的任意子模都是有限的。
有限 $R$-模的 $R$-子模满足升链条件。
\end{lemma}

\begin{proof}
先证明有限 $R$-模 $M$ 的任意子模 $N$ 都是有限的。
对 $M$ 的生成元个数作归纳。若该数为 $1$，则对某些理想
$I \subset J \subset R$，有 $N = J/I \subset M = R/I$。
所以诺特性的定义给出结论。若 $M$ 的生成元多于 $1$ 个，
则可找到短正合列 $0 \to M' \to M \to M'' \to 0$，
其中 $M'$ 和 $M''$ 的生成元更少。注意，置
$N' = M' \cap N$ 及 $N'' = \Im(N \to M'')$，可得到 $N$ 的类似短正合列。
因此，结论由归纳假设得出，因为 $N$ 的生成元个数至多为
$N'$ 的生成元个数与 $N''$ 的生成元个数之和。

\medskip\noindent
为证明 $M$ 是有限表示的，只须把前面的结论应用于一个表示
$R^n \to M$ 的核。

\medskip\noindent
众所周知且容易证明，$M$ 的 $R$-子模满足升链条件，
等价于 $M$ 的每个子模都是有限 $R$-模。省略证明。
\end{proof}

\begin{lemma}[阿廷--里斯]
\label{algebra-lemma-Artin-Rees}
假设 $R$ 是诺特环，$I \subset R$ 是理想。
设 $N \subset M$ 为有限 $R$-模。则存在常数 $c > 0$，使得对所有
$n \geq c$，有 $I^n M \cap N  =  I^{n-c}(I^cM \cap N)$。
\end{lemma}

\begin{proof}
考虑环 $S = R \oplus I \oplus I^2 \oplus \ldots
= \bigoplus_{n \geq 0} I^n$。约定 $I^0 = R$。
乘法映射 $I^n \times I^m$ 通过 $R$ 中的乘法映到 $I^{n + m}$。
注意，若 $I = (f_1, \ldots, f_t)$，则 $S$ 是诺特环
$R[X_1, \ldots, X_t]$ 的商。该映射把单项式
$X_1^{e_1}\ldots X_t^{e_t}$ 映到 $f_1^{e_1}\ldots f_t^{e_t}$。
所以 $S$ 是诺特环。类似地，考虑模
$M \oplus IM \oplus I^2M \oplus \ldots
= \bigoplus_{n \geq 0} I^nM$。这是有限生成 $S$-模。
事实上，若 $x_1, \ldots, x_r$ 在 $R$ 上生成 $M$，
则它们也在 $S$ 上生成 $\bigoplus_{n \geq 0} I^nM$。
接着考虑子模 $\bigoplus_{n \geq 0} I^nM \cap N$。
容易验证，这是 $S$-子模。由引理 \ref{algebra-lemma-Noetherian-basic}，
它是有限生成 $S$-模；设其生成元为
$\xi_j \in \bigoplus_{n \geq 0} I^nM \cap N$，$j = 1, \ldots, s$。
将每个 $\xi_j$ 分解为齐次部分后，可以假设对某个 $d_j$，
每个 $\xi_j \in I^{d_j}M \cap N$。置 $c = \max\{d_j\}$。
则对所有 $n \geq c$，$I^nM \cap N$ 中的每个元素都形如
$\sum h_j \xi_j$，其中 $h_j \in I^{n - d_j}$。
本引理于是由此以及平凡观察
$I^{n-d_j}(I^{d_j}M \cap N) \subset I^{n-c}(I^cM \cap N)$ 得出。
\end{proof}

\begin{lemma}
\label{algebra-lemma-map-AR}
假设 $0 \to K \to M \xrightarrow{f} N$ 是诺特环 $R$ 上
有限生成模的正合列。设 $I \subset R$ 为理想。则存在 $c$，使得对所有
$n \geq c$，有
$$
f^{-1}(I^nN) = K + I^{n-c}f^{-1}(I^cN)
\quad\text{且}\quad
f(M) \cap I^nN \subset f(I^{n - c}M)
$$
\end{lemma}

\begin{proof}
将引理 \ref{algebra-lemma-Artin-Rees} 应用于 $\Im(f) \subset N$，并注意映射
$f : I^{n-c}M \to I^{n-c}f(M)$ 是满射。
\end{proof}

\begin{lemma}[克鲁尔交定理]
\label{algebra-lemma-intersect-powers-ideal-module-zero}
设 $R$ 为诺特局部环，$I \subset R$ 为真理想，$M$ 为有限 $R$-模。
则 $\bigcap_{n \geq 0} I^nM = 0$。
\end{lemma}

\begin{proof}
置 $N = \bigcap_{n \geq 0} I^nM$。则对所有 $n \geq 0$，
$N = I^nM \cap N$。由阿廷--里斯引理 \ref{algebra-lemma-Artin-Rees}，
对某个充分大的 $n$，有 $N = I^nM \cap N \subset IN$。
由中山引理 \ref{algebra-lemma-NAK}，$N = 0$。
\end{proof}

\begin{lemma}
\label{algebra-lemma-intersection-powers-ideal-module}
设 $R$ 为诺特环，$I \subset R$ 为理想，$M$ 为有限 $R$-模。
令 $N = \bigcap_n I^n M$。
\begin{enumerate}
\item 对每个满足 $I \subset \mathfrak p$ 的素理想 $\mathfrak p$，
存在 $f \in R$，$f \not \in \mathfrak p$，使得 $N_f = 0$。
\item 若 $I$ 包含于 $R$ 的雅可布森根，则 $N = 0$。
\end{enumerate}
\end{lemma}

\begin{proof}
证明 (1)。设 $x_1, \ldots, x_n$ 为模 $N$ 的生成元；参见引理
\ref{algebra-lemma-Noetherian-basic}。对每个满足 $I \subset \mathfrak p$
的素理想 $\mathfrak p$，由引理
\ref{algebra-lemma-intersect-powers-ideal-module-zero}，$N$ 在局部化
$M_{\mathfrak p}$ 中的像为零。因此，可找到
$g_i \in R$，$g_i \not \in \mathfrak p$，使得 $x_i$ 在 $N_{g_i}$ 中映到零。
于是 $N_{g_1g_2\ldots g_n} = 0$。

\medskip\noindent
(2) 由 (1) 和引理 \ref{algebra-lemma-characterize-zero-local} 得出。
\end{proof}

\begin{remark}
\label{algebra-remark-intersection-powers-ideal}
引理 \ref{algebra-lemma-intersect-powers-ideal-module-zero} 特别表明：
若 $I \subset R$ 是诺特局部环 $R$ 中的非单位理想，则
$\bigcap_n I^n = (0)$。更一般地，设 $R$ 为诺特环，$I \subset R$ 为理想。
假设 $f \in \bigcap_{n \in \mathbf{N}} I^n$。
则引理 \ref{algebra-lemma-intersection-powers-ideal-module} 表明，
对每个素理想 $I \subset \mathfrak p$，存在
$g \in R$，$g \not \in \mathfrak p$，使得 $f$ 在 $R_g$ 中映到零。
在代数几何中，我们把它表述为：“$f$ 在 $\Spec(R)$ 的闭集
$V(I)$ 的某个开邻域中为零”。
\end{remark}

\begin{lemma}[阿廷--泰特]
\label{algebra-lemma-Artin-Tate}
设 $R$ 为诺特环，$S$ 为有限生成 $R$-代数。
若 $T \subset S$ 是 $R$-子代数，且 $S$ 作为 $T$-模有限生成，
则 $T$ 在 $R$ 上是有限型的。
\end{lemma}

\begin{proof}
选取作为 $R$-代数生成 $S$ 的元素 $x_1, \ldots, x_n \in S$。
在 $S$ 中选取作为 $T$-模生成 $S$ 的 $y_1, \ldots, y_m$。
于是存在 $a_{ij} \in T$，使得 $x_i = \sum a_{ij} y_j$。
还存在 $b_{ijk} \in T$，使得 $y_i y_j = \sum b_{ijk} y_k$。
令 $T' \subset T$ 为由 $a_{ij}$ 和 $b_{ijk}$ 生成的 $R$-子代数。
它是有限生成 $R$-代数，因而是诺特环。考虑代数
$$
S' = T'[Y_1, \ldots, Y_m]/(Y_i Y_j - \sum b_{ijk} Y_k).
$$
注意，$S'$ 在 $T'$ 上有限；事实上，作为 $T'$-模，它由
$1, Y_1, \ldots, Y_m$ 的类生成。考虑把 $Y_i$ 映到 $y_i$ 的
$T'$-代数同态 $S' \to S$。由于 $a_{ij} \in T'$，可知 $x_j$
属于此映射的像。因此 $S' \to S$ 是满射。
故 $S$ 在 $T'$ 上也是有限的。由于 $T'$ 是诺特环，
可知 $T \subset S$ 在 $T'$ 上有限，证明完成。
\end{proof}






















\section{长度}
\label{algebra-section-length}

\begin{definition}
\label{algebra-definition-length}
设 $R$ 为环。对任意 $R$-模 $M$，定义 $M$ 在 $R$ 上的{it 长度}为
$$
\text{length}_R(M)
=
\sup
\{
n
\mid
\exists\ 0 = M_0 \subset M_1 \subset \ldots \subset M_n = M,
\text{ }M_i \not = M_{i + 1}
\}.
$$
\end{definition}

\noindent
换言之，它是子模链长度的上确界。一个子模链何时是另一个子模链的加细，
有一个显然的概念。这在所有子模链之集上给出偏序；若 $M$ 非零，
最小链的形状是 $0 = M_0 \subset M_1 = M$。
注意一个显然事实：若 $M$ 的长度有限，则每条链都可加细为极大链。
但是，所有极大链长度相同并不如此显然（稍后将看到）。

\begin{lemma}
\label{algebra-lemma-finite-length-finite}
\begin{slogan}
有限长度模是有限模。
\end{slogan}
设 $R$ 为环，$M$ 为 $R$-模。
若 $\text{length}_R(M) < \infty$，则 $M$ 是有限 $R$-模。
\end{lemma}

\begin{proof}
省略。
\end{proof}

\begin{lemma}
\label{algebra-lemma-length-additive}
\begin{slogan}
长度在短正合列中可加。
\end{slogan}
若 $0 \to M' \to M \to M'' \to 0$ 是 $R$ 上模的短正合列，
则 $M$ 的长度等于 $M'$ 与 $M''$ 的长度之和。
\end{lemma}

\begin{proof}
给定 $M'$ 和 $M''$ 的长度分别为 $n', n''$ 的滤过，
容易构造 $M$ 的相应滤过，其长度为 $n' + n''$。于是
$\text{length}_R M \geq \text{length}_R M' + \text{length}_R M''$。
反之，给定 $M$ 的滤过
$M_0 \subset M_1 \subset \ldots \subset M_n$，考虑诱导滤过
$M_i' = M_i \cap M'$ 和 $M_i'' = \Im(M_i \to M'')$。
令 $n'$（相应地，$n''$）为滤过 $\{M'_i\}$（相应地，$\{M''_i\}$）
的步数。若 $M_i' = M_{i + 1}'$ 且 $M_i'' = M_{i + 1}''$，
则 $M_i = M_{i + 1}$。所以 $n' + n'' \geq n$。
结合先前结论即得证明。
\end{proof}

\begin{lemma}
\label{algebra-lemma-length-infinite}
设 $R$ 为极大理想是 $\mathfrak m$ 的局部环。
若 $M$ 是 $R$-模，且对所有 $n \geq 0$，
$\mathfrak m^n M \not = 0$，则 $\text{length}_R(M) = \infty$。
换言之，若 $M$ 长度有限，则对某个 $n$，$\mathfrak m^nM = 0$。
\end{lemma}

\begin{proof}
假设对所有 $n\geq 0$，$\mathfrak m^n M \not = 0$。
选取 $x \in M$ 及 $f_1, \ldots, f_n \in \mathfrak m$，使得
$f_1f_2 \ldots f_n x \not = 0$。滤过
$$
0 \subset R f_1 \ldots f_n x \subset R f_1 \ldots f_{n - 1} x
\subset \ldots \subset R x \subset M
$$
的前 $n$ 步彼此不同。例如，若
$R f_1 x = R f_1 f_2 x$，则对某个 $g$，
$f_1 x = g f_1 f_2 x$，从而 $(1 - gf_2) f_1 x = 0$；
又因为 $1 - gf_2$ 是单位，所以 $f_1 x = 0$，
这与 $x$ 和 $f_1, \ldots, f_n$ 的选取矛盾。因此长度为无穷。
\end{proof}

\begin{lemma}
\label{algebra-lemma-length-independent}
设 $R \to S$ 为环映射，$M$ 为 $S$-模。总有
$\text{length}_R(M) \geq \text{length}_S(M)$。
若 $R \to S$ 是满射，则等号成立。
\end{lemma}

\begin{proof}
$M$ 的 $S$-子模滤过给出 $M$ 的 $R$-子模滤过。这证明不等式。
若 $R \to S$ 是满射，则 $M$ 的任意 $R$-子模自动也是 $S$-子模。
所以在此情形下等号成立。
\end{proof}

\begin{lemma}
\label{algebra-lemma-dimension-is-length}
设 $R$ 为极大理想是 $\mathfrak m$ 的环。假设 $M$ 是满足
$\mathfrak m M  =  0$ 的 $R$-模。则 $M$ 作为 $R$-模的长度，
等于 $M$ 作为 $R/\mathfrak m$-向量空间的维数。
该长度有限，当且仅当 $M$ 是有限 $R$-模。
\end{lemma}

\begin{proof}
第一部分是引理 \ref{algebra-lemma-length-independent} 的特例。
所以长度有限，当且仅当 $M$ 作为 $R/\mathfrak m$-向量空间有有限基，
当且仅当 $M$ 作为 $R$-模有有限生成集。
\end{proof}

\begin{lemma}
\label{algebra-lemma-length-localize}
设 $R$ 为环，$M$ 为 $R$-模，$S \subset R$ 为乘法子集。则
$\text{length}_R(M) \geq \text{length}_{S^{-1}R}(S^{-1}M)$。
\end{lemma}

\begin{proof}
由引理 \ref{algebra-lemma-submodule-localization}，任意子模
$N' \subset S^{-1}M$ 都形如 $S^{-1}N$，其中 $N \subset M$ 为 $R$-子模。
本引理得证。
\end{proof}

\begin{lemma}
\label{algebra-lemma-length-finite}
设 $R$ 为极大理想 $\mathfrak m$ 有限生成的环（例如 $R$ 为诺特环）。
假设 $M$ 是有限 $R$-模，并且对某个 $n$，
$\mathfrak m^n M  =  0$。则 $\text{length}_R(M) < \infty$。
\end{lemma}

\begin{proof}
考虑滤过
$0 = \mathfrak m^n M \subset
\mathfrak m^{n-1} M \subset
\ldots \subset \mathfrak m M \subset M$。
所有子商都是可应用引理 \ref{algebra-lemma-dimension-is-length} 的有限生成 $R$-模。
由可加性得出结论；参见引理 \ref{algebra-lemma-length-additive}。
\end{proof}

\begin{definition}
\label{algebra-definition-simple-module}
设 $R$ 为环，$M$ 为 $R$-模。若 $M \not = 0$，
且 $M$ 的每个子模等于 $M$ 或 $0$，则称 $M$ 是{it 单模}。
\end{definition}

\begin{lemma}
\label{algebra-lemma-characterize-length-1}
设 $R$ 为环，$M$ 为 $R$-模。下列条件等价：
\begin{enumerate}
\item $M$ 是单模；
\item $\text{length}_R(M) = 1$；
\item 对某个极大理想 $\mathfrak m \subset R$，有
$M \cong R/\mathfrak m$。
\end{enumerate}
\end{lemma}

\begin{proof}
设 $\mathfrak m$ 为 $R$ 的极大理想。由引理
\ref{algebra-lemma-dimension-is-length}，模 $R/\mathfrak m$ 的长度为 $1$。
前两个断言的等价性是重言式。假设 $M$ 是单模。
取 $x \in M$，$x \not = 0$。由于 $M$ 是单模，有 $M = R \cdot x$。
令 $I \subset R$ 为 $x$ 的零化理想，即
$I = \{f \in R \mid fx = 0\}$。映射 $R/I \to M$，
$f \bmod I \mapsto fx$ 是同构，因而 $R/I$ 是单 $R$-模。
由于 $R/I \not = 0$，可知 $I \not = R$。
令 $I \subset \mathfrak m$ 为包含 $I$ 的极大理想。
若 $I \not = \mathfrak m$，则 $\mathfrak m /I \subset R/I$
是非平凡子模，与 $R/I$ 的单性矛盾。因此如所需，$I = \mathfrak m$。
\end{proof}

\begin{lemma}
\label{algebra-lemma-simple-pieces}
设 $R$ 为环，$M$ 为有限长度 $R$-模。任取子模的极大链
$$
0 = M_0 \subset M_1 \subset M_2 \subset \ldots \subset M_n = M
$$
其中 $M_i \not = M_{i-1}$，$i = 1, \ldots, n$。则
\begin{enumerate}
\item $n = \text{length}_R(M)$；
\item 每个 $M_i/M_{i-1}$ 都是单模；
\item 每个 $M_i/M_{i-1}$ 都形如 $R/\mathfrak m_i$，
其中 $\mathfrak m_i$ 为某个极大理想；
\item 给定极大理想 $\mathfrak m \subset R$，有
$$
\# \{i \mid \mathfrak m_i = \mathfrak m\}
=
\text{length}_{R_{\mathfrak m}} (M_{\mathfrak m}).
$$
\end{enumerate}
\end{lemma}

\begin{proof}
若 $M_i/M_{i-1}$ 不是单模，则可以加细该滤过，因而该滤过不是极大的。
所以 $M_i/M_{i-1}$ 是单模。由引理
\ref{algebra-lemma-characterize-length-1}，模 $M_i/M_{i-1}$ 长度为 $1$，
并且形如 $R/\mathfrak m_i$，其中 $\mathfrak m_i$ 为某个极大理想。
由长度的可加性，即引理 \ref{algebra-lemma-length-additive}，可知
$n = \text{length}_R(M)$。由于局部化正合，
$$
0 = (M_0)_{\mathfrak m}
\subset (M_1)_{\mathfrak m}
\subset (M_2)_{\mathfrak m}
\subset \ldots
\subset (M_n)_{\mathfrak m} = M_{\mathfrak m}
$$
是 $M_{\mathfrak m}$ 的滤过，其相继商为
$(M_i/M_{i-1})_{\mathfrak m}$。所以，最后一个陈述直接由下述事实得出：
给定 $R$ 的极大理想 $\mathfrak m$、$\mathfrak m'$，有
$$
(R/\mathfrak m')_{\mathfrak m}
\cong
\left\{
\begin{matrix}
0 &
\text{若 } \mathfrak m \not = \mathfrak m', \\
R_{\mathfrak m}/\mathfrak m R_{\mathfrak m} &
\text{若 } \mathfrak m  = \mathfrak m'
\end{matrix}
\right.
$$
把这一点留给读者。
\end{proof}

\begin{lemma}
\label{algebra-lemma-pushdown-module}
设 $A$ 为极大理想是 $\mathfrak m$ 的局部环。
设 $B$ 为极大理想是 $\mathfrak m_i$，$i = 1, \ldots, n$ 的半局部环。
假设 $A \to B$ 是同态，每个 $\mathfrak m_i$ 都位于 $\mathfrak m$ 上，且
$$
[\kappa(\mathfrak m_i) : \kappa(\mathfrak m)] < \infty.
$$
设 $M$ 为有限长度 $B$-模。则
$$
\text{length}_A(M) = \sum\nolimits_{i = 1, \ldots, n}
[\kappa(\mathfrak m_i) : \kappa(\mathfrak m)]
\text{length}_{B_{\mathfrak m_i}}(M_{\mathfrak m_i}),
$$
特别地，$\text{length}_A(M) < \infty$。
\end{lemma}

\begin{proof}
选取引理 \ref{algebra-lemma-simple-pieces} 中那样的 $B$-子模极大链
$$
0 = M_0
\subset M_1
\subset M_2
\subset \ldots
\subset M_m = M
$$
则对某个 $i(j) \in \{1, \ldots, n\}$，每个商
$M_j/M_{j - 1}$ 都同构于 $\kappa(\mathfrak m_{i(j)})$。
此外，由引理 \ref{algebra-lemma-dimension-is-length}，还有：
$$
\text{length}_A(\kappa(\mathfrak m_i)) =
[\kappa(\mathfrak m_i) : \kappa(\mathfrak m)]
$$
本引理由长度的可加性得出（引理 \ref{algebra-lemma-length-additive}）。
\end{proof}

\begin{lemma}
\label{algebra-lemma-pullback-module}
设 $A \to B$ 为局部环的平坦局部同态。则对任意 $A$-模 $M$，有
$$
\text{length}_A(M) \text{length}_B(B/\mathfrak m_AB)
=
\text{length}_B(M \otimes_A B).
$$
特别地，若 $\text{length}_B(B/\mathfrak m_AB) < \infty$，
则 $M$ 长度有限，当且仅当 $M \otimes_A B$ 长度有限。
\end{lemma}

\begin{proof}
由引理 \ref{algebra-lemma-local-flat-ff}，环映射 $A \to B$ 忠实平坦。
所以，若 $0 = M_0 \subset M_1 \subset \ldots \subset M_n = M$
是 $M$ 中长度为 $n$ 的链，则相应的链
$0 = M_0 \otimes_A B \subset M_1 \otimes_A B \subset
\ldots \subset M_n \otimes_A B = M \otimes_A B$ 的长度也为 $n$。
这证明
$\text{length}_A(M) = \infty \Rightarrow
\text{length}_B(M \otimes_A B) = \infty$。
接着，假设 $\text{length}_A(M) < \infty$。
此时，$M$ 有一个长度为 $\ell = \text{length}_A(M)$ 的滤过，
其各商都是 $A/\mathfrak m_A$。如上论证，$M \otimes_A B$
有一个长度为 $\ell$ 的滤过，其各商同构于
$B \otimes_A A/\mathfrak m_A =  B/\mathfrak m_AB$。本引理得证。
\end{proof}

\begin{lemma}
\label{algebra-lemma-pullback-transitive}
设 $A \to B \to C$ 为局部环的平坦局部同态。则
$$
\text{length}_B(B/\mathfrak m_A B)
\text{length}_C(C/\mathfrak m_B C)
=
\text{length}_C(C/\mathfrak m_A C)
$$
\end{lemma}

\begin{proof}
将引理 \ref{algebra-lemma-pullback-module} 应用于环映射
$B \to C$ 和 $B$-模 $M = B/\mathfrak m_A B$ 即得。
\end{proof}












\section{阿廷环}
\label{algebra-section-artinian}

\noindent
阿廷环，特别是局部阿廷环，在代数几何中起重要作用，例如在形变理论中。

\begin{definition}
\label{algebra-definition-artinian}
若环 $R$ 的理想满足降链条件，则称它是{\it 阿廷的}。
\end{definition}

\begin{lemma}
\label{algebra-lemma-finite-dimensional-algebra}
假设 $R$ 是域上的有限维代数。则 $R$ 是阿廷环。
\end{lemma}

\begin{proof}
理想的降链条件显然成立。
\end{proof}

\begin{lemma}
\label{algebra-lemma-artinian-finite-nr-max}
若 $R$ 是阿廷环，则 $R$ 只有有限多个极大理想。
\end{lemma}

\begin{proof}
假设 $\mathfrak m_i$，$i = 1, 2, 3, \ldots$ 是两两不同的极大理想。
则
$\mathfrak m_1 \supset \mathfrak m_1\cap \mathfrak m_2
\supset \mathfrak m_1 \cap \mathfrak m_2 \cap \mathfrak m_3 \supset \ldots$
是无穷降链（因为由中国剩余定理，所有映射
$R \to \oplus_{i = 1}^n R/\mathfrak m_i$ 都是满射）。
\end{proof}

\begin{lemma}
\label{algebra-lemma-artinian-radical-nilpotent}
设 $R$ 为阿廷环。则 $R$ 的雅可布森根是幂零理想。
\end{lemma}

\begin{proof}
设 $I \subset R$ 为雅可布森根。注意，
$I \supset I^2 \supset I^3 \supset \ldots$ 是降链。
所以对某个 $n$，$I^n = I^{n + 1}$。置
$J = \{ x\in R \mid xI^n = 0\}$。须证明 $J = R$。
若非如此，选取满足 $J \subset J'$ 的理想 $J' \not = J$，
并使它极小（由阿廷性可行）。则 $J'/J$ 是单 $R$-模，
因而同构于某个极大理想 $\mathfrak m$ 所给出的 $R/\mathfrak m$；
参见引理 \ref{algebra-lemma-characterize-length-1}。
于是 $\mathfrak m I^n$ 零化 $J'$。由于 $I \subset \mathfrak m$，
可知 $I^{n + 1} = I^n$ 零化 $J'$。所以 $J' = J$，矛盾。
\end{proof}

\begin{lemma}
\label{algebra-lemma-product-local}
任意只有有限多个极大理想、且雅可布森根局部幂零的环，
都等于它在各极大理想处的局部化之积。此外，所有素理想都是极大理想。
\end{lemma}

\begin{proof}
设 $R$ 为极大理想只有 $\mathfrak m_1, \ldots, \mathfrak m_n$ 的环。
令 $I = \bigcap_{i = 1}^n \mathfrak m_i$ 为 $R$ 的雅可布森根。
假设 $I$ 局部幂零。设 $\mathfrak p$ 为 $R$ 的素理想。
由于每个素理想都包含 $R$ 的每个幂零元，可知
$ \mathfrak p \supset \mathfrak m_1 \cap \ldots \cap \mathfrak m_n$。
又由于 $\mathfrak m_1 \cap \ldots \cap \mathfrak m_n \supset
\mathfrak m_1 \ldots \mathfrak m_n$，可知
$\mathfrak p \supset \mathfrak m_1 \ldots \mathfrak m_n$。
故对某个 $i$，$\mathfrak p \supset \mathfrak m_i$，从而
$\mathfrak p = \mathfrak m_i$。因此 $R$ 的谱是离散拓扑空间
$\{\mathfrak m_1, \ldots, \mathfrak m_n\}$。
将引理 \ref{algebra-lemma-disjoint-implies-product} 应用 $n - 1$ 次，得到
$R = R_1 \times \ldots \times R_n$，其中对每个 $i$，$R_i$ 的谱都是单点集。
所以 $R_i$ 是局部环；又因为由该引理，它是 $R$ 的局部化，
故它如所需地是 $R$ 的这些局部环之一。
\end{proof}

\begin{lemma}
\label{algebra-lemma-artinian-finite-length}
环 $R$ 是阿廷环，当且仅当它作为自身上的模具有有限长度。
任意这样的环 $R$ 既是阿廷环又是诺特环；$R$ 的任意素理想都是极大理想；
并且 $R$ 等于它在各极大理想处的局部化的（有限）乘积。
\end{lemma}

\begin{proof}
若 $R$ 在自身上长度有限，则它的理想既满足升链条件又满足降链条件。
所以它既是诺特环又是阿廷环。由引理
\ref{algebra-lemma-artinian-finite-nr-max}、
\ref{algebra-lemma-artinian-radical-nilpotent} 和 \ref{algebra-lemma-product-local}，
任意阿廷环都等于它在极大理想处的局部化之积。

\medskip\noindent
假设 $R$ 是阿廷环。将证明 $R$ 在自身上长度有限。
只须给出一条子模链，使其相继商长度有限。
由上述结论，可以假设 $R$ 是极大理想为 $\mathfrak m$ 的局部环。
由引理 \ref{algebra-lemma-artinian-radical-nilpotent}，
对某个 $n$，有 $\mathfrak m^n =0$。考虑序列
$0 = \mathfrak m^n \subset \mathfrak m^{n-1} \subset
\ldots \subset \mathfrak m \subset R$。由引理
\ref{algebra-lemma-dimension-is-length}，每个子商
$\mathfrak m^j/\mathfrak m^{j + 1}$ 的长度，等于它作为
$\kappa(\mathfrak m)$-向量空间的维数。该维数必须有限；
否则会有向量子空间的无穷降链，它对应于 $R$ 中理想的无穷降链。
\end{proof}

\section{本质有限型同态}
\label{algebra-section-essentially-of-finite-type}

\noindent
下面给出关于有限型环同态之局部化的一些简单说明。

\begin{definition}
\label{algebra-definition-essentially-finite-p-t}
设 $R \to S$ 为环同态。
\begin{enumerate}
\item 若 $S$ 是某个有限型 $R$-代数的局部化，则称 $R \to S$
是{\it 本质有限型的}。
\item 若 $S$ 是某个有限表示 $R$-代数的局部化，则称 $R \to S$
是{\it 本质有限表示的}。
\end{enumerate}
\end{definition}

\begin{lemma}
\label{algebra-lemma-composition-essentially-of-finite-type}
本质有限型环同态的类在复合下保持不变。本质有限表示的情形亦然。
\end{lemma}

\begin{proof}
略。
\end{proof}

\begin{lemma}
\label{algebra-lemma-base-change-essentially-of-finite-type}
本质有限型环同态的类在基变换下保持不变。本质有限表示的情形亦然。
\end{lemma}

\begin{proof}
略。
\end{proof}

\begin{lemma}
\label{algebra-lemma-essentially-of-finite-type-into-artinian-local}
设 $R \to S$ 为环同态。假设 $S$ 是极大理想为 $\mathfrak m$
的阿廷局部环。则
\begin{enumerate}
\item $R \to S$ 有限，当且仅当 $R \to S/\mathfrak m$ 有限；
\item $R \to S$ 是有限型的，当且仅当 $R \to S/\mathfrak m$
是有限型的；
\item $R \to S$ 本质有限型，当且仅当复合
$R \to S/\mathfrak m$ 本质有限型。
\end{enumerate}
\end{lemma}

\begin{proof}
若 $R \to S$ 有限，则由引理 \ref{algebra-lemma-finite-transitive}，
$R \to S/\mathfrak m$ 有限。反过来，假设 $R \to S/\mathfrak m$ 有限。
由于 $S$ 在自身上具有有限长度（引理
\ref{algebra-lemma-artinian-finite-length}），可以选取由理想组成的滤过
$$
0 \subset I_1 \subset \ldots \subset I_n = S
$$
使得作为 $S$-模，有 $I_i/I_{i - 1} \cong S/\mathfrak m$。
因此，$S$ 有一条由 $R$-子模 $I_i$ 组成的滤过，其中每个相继商
都是有限 $R$-模。故由引理 \ref{algebra-lemma-extension}，$S$ 是有限
$R$-模。

\medskip\noindent
若 $R \to S$ 是有限型的，则由引理
\ref{algebra-lemma-compose-finite-type}，$R \to S/\mathfrak m$ 是有限型的。
反过来，假设 $R \to S/\mathfrak m$ 是有限型的。选取
$f_1, \ldots, f_n \in S$，使它们的像生成 $S/\mathfrak m$。
于是 $A = R[x_1, \ldots, x_n] \to S$、$x_i \mapsto f_i$
是环同态，且 $A \to S/\mathfrak m$ 满射（特别地，它是有限的）。
故由第 (1) 部分，$A \to S$ 有限；再由引理
\ref{algebra-lemma-compose-finite-type}，可知 $R \to S$ 是有限型的。

\medskip\noindent
若 $R \to S$ 本质有限型，则由引理
\ref{algebra-lemma-composition-essentially-of-finite-type}，
$R \to S/\mathfrak m$ 本质有限型。反过来，假设
$R \to S/\mathfrak m$ 本质有限型。设 $S/\mathfrak m$
是 $R[x_1, \ldots, x_n]/I$ 的局部化。选取
$f_1, \ldots, f_n \in S$，使它们模 $\mathfrak m$ 的剩余类
分别对应于 $x_1, \ldots, x_n$ 模 $I$ 的剩余类。考虑同态
$R[x_1, \ldots, x_n] \to S$、$x_i \mapsto f_i$，并记其核为 $J$。
置 $A = R[x_1, \ldots, x_n]/J \subset S$，并置
$\mathfrak p = A \cap \mathfrak m$。注意，
$A/\mathfrak p \subset S/\mathfrak m$ 等于
$R[x_1, \ldots, x_n]/I$ 在 $S/\mathfrak m$ 中的像。
所以 $\kappa(\mathfrak p) = S/\mathfrak m$。由第 (1) 部分，
$A_\mathfrak p \to S$ 有限。最后由引理
\ref{algebra-lemma-composition-essentially-of-finite-type}，
可知 $S$ 本质有限型。
\end{proof}

\noindent
下面的引理也可利用射影空间的固有性来证明，而不必使用这里给出的代数论证。

\begin{lemma}
\label{algebra-lemma-localization-at-closed-point-special-fibre}
设 $\varphi : R \to S$ 本质有限型，且 $R$ 与 $S$ 都是局部环
（但 $\varphi$ 不必是局部同态）。则存在某个 $n$ 和一个位于
$\mathfrak m_R$ 之上的极大理想
$\mathfrak m \subset R[x_1, \ldots, x_n]$，使得 $S$ 是
$R[x_1, \ldots, x_n]_\mathfrak m$ 的某个商环的局部化。
\end{lemma}

\begin{proof}
可以将 $S$ 写成 $R[x_1, \ldots, x_n]$ 的某个商环的局部化。
因此只须证明如下情形：对某个素理想
$\mathfrak q \subset R[x_1, \ldots, x_n]$，有
$S = R[x_1, \ldots, x_n]_\mathfrak q$。若
$\mathfrak q + \mathfrak m_R R[x_1, \ldots, x_n] \not = R[x_1, \ldots, x_n]$，
则可找到断言中的极大理想 $\mathfrak m$，使得
$\mathfrak q \subset \mathfrak m$；此时结论显然。

\medskip\noindent
选取支配 $R \to \kappa(\mathfrak q)$ 之像的赋值环
$A \subset \kappa(\mathfrak q)$（引理 \ref{algebra-lemma-dominate}）。
若 $x_i$ 的像 $\lambda_i \in \kappa(\mathfrak q)$ 属于 $A$，
则 $\mathfrak q$ 包含于 $\mathfrak m_A$ 在同态
$R[x_1, \ldots, x_n] \to A$ 下的逆像中，这就回到前一种情形。
因此存在某个 $i$，使得 $\lambda_i^{-1} \in A$，并且对所有
$j = 1, \ldots, n$，有 $\lambda_j/\lambda_i \in A$
（因为 $A$ 的值群是全序的，参见引理
\ref{algebra-lemma-valuation-group}）。于是考虑同态
$$
R[y_0, y_1, \ldots, \hat{y_i}, \ldots, y_n]
\to R[x_1, \ldots, x_n]_\mathfrak q,\quad
y_0 \mapsto 1/x_i,\quad y_j \mapsto x_j/x_i
$$
设 $\mathfrak q' \subset R[y_0, \ldots, \hat{y_i}, \ldots, y_n]$
为 $\mathfrak q$ 的逆像。由于 $y_0 \not \in \mathfrak q'$，
容易看出上述箭头在局部化后给出一个同构。另一方面，
第一段的结论适用于 $R[y_0, \ldots, \hat{y_i}, \ldots, y_n]$，
因为 $y_j$ 映到 $A$ 的一个元素。证明完毕。
\end{proof}






















\section{K-群}
\label{algebra-section-K-groups}

\noindent
设 $R$ 为环。下面将引入两个与 $R$ 相伴的阿贝尔群。第一个记作
$K'_0(R)$，它具有如下性质\footnote{该定义对任意环都有意义，
但除非 $R$ 是诺特环，否则很少使用。}：
\begin{enumerate}
\item 对每个有限 $R$-模 $M$，给定 $K'_0(R)$ 中的一个元素 $[M]$；
\item 对有限 $R$-模的每个短正合列
$0 \to M' \to M \to M'' \to 0$，都有关系
$[M] = [M'] + [M'']$；
\item 群 $K'_0(R)$ 由诸元素 $[M]$ 生成；
\item $K'_0(R)$ 中生成元 $[M]$ 之间的所有关系，都是由上述正合列
所产生关系的 $\mathbf{Z}$-线性组合。
\end{enumerate}
实际构造略为麻烦，因为必须顾及所有有限生成 $R$-模的汇集是一个真类。
不过，可以用下述方式克服这个问题：取诸元素 $[R^n/K]$ 作为群
$K'_0(R)$ 的生成元集合，其中 $n$ 遍历所有整数，而 $K$ 遍历所有
子模 $K \subset R^n$。施加于这些元素上的关系子群，可取为项均具有
$R^n/K$ 形式的短正合列所给出的关系。元素 $[M]$ 定义如下：
选取 $n$ 与 $K$，使得 $M \cong R^n/K$，并置
$[M] = [R^n/K]$。细节留给读者。

\begin{lemma}
\label{algebra-lemma-length-K0}
若 $R$ 是阿廷局部环，则长度函数定义一个自然的阿贝尔群同态
$\text{length}_R : K'_0(R) \to \mathbf{Z}$。
\end{lemma}

\begin{proof}
任意有限 $R$-模的长度都有限，因为它是 $R^n$ 的商，而由引理
\ref{algebra-lemma-artinian-finite-length}，后者长度有限。
此外，长度函数具有可加性；参见引理 \ref{algebra-lemma-length-additive}。
\end{proof}

\noindent
上述两个群中的第二个记作 $K_0(R)$，它具有如下性质：
\begin{enumerate}
\item 对每个有限投射 $R$-模 $M$，给定 $K_0(R)$ 中的一个元素 $[M]$；
\item 对有限投射 $R$-模的每个短正合列
$0 \to M' \to M \to M'' \to 0$，都有关系
$[M] = [M'] + [M'']$；
\item 群 $K_0(R)$ 由诸元素 $[M]$ 生成；
\item $K_0(R)$ 中的所有关系，都是由上述正合列所产生关系的
$\mathbf{Z}$-线性组合。
\end{enumerate}
这个群的构造与上面相同。

\medskip\noindent
注意，存在显然的映射 $K_0(R) \to K'_0(R)$，它一般不是同构。

\begin{example}
\label{algebra-example-K0-field}
注意，若 $R = k$ 是域，则显然有
$K_0(k) = K'_0(k) \cong \mathbf{Z}$；该同构由维数函数给出
（它也就是长度函数）。
\end{example}

\begin{example}
\label{algebra-example-K0-PID}
设 $R$ 为主理想整环。断言 $K_0(R) = K'_0(R) = \mathbf{Z}$。
事实上，任意有限投射 $R$-模都是有限自由模。由引理
\ref{algebra-lemma-rank}，有限自由模具有良定义的秩。给定有限自由模的短正合列
$$
0 \to M' \to M \to M'' \to 0
$$
有 $\text{rank}(M) = \text{rank}(M') + \text{rank}(M'')$，
因为在此情形下 $M \cong M' \oplus M'$（例如，由引理
\ref{algebra-lemma-lift-map}，该列分裂）。由此得到
$K_0(R) = \mathbf{Z}$。

\medskip\noindent
主理想整环上模的结构定理表明，任意有限生成 $R$-模都具有形式
$M = R^{\oplus r} \oplus R/(d_1) \oplus \ldots \oplus R/(d_k)$。
考虑短正合列
$$
0 \to (d_i) \to R \to R/(d_i) \to 0
$$
由于理想 $(d_i)$ 作为模同构于 $R$（它是由 $d_i$ 生成的自由模），
在 $K'_0(R)$ 中有 $[(d_i)] = [R]$。于是
$[R/(d_i)] = [(d_i)]-[R] = 0$。由此可知，挠模在
$K'_0(R)$ 中的类为零。利用自由部分的秩，可得到等同
$K'_0(R) = \mathbf{Z}$，且典范同态
$K_0(R) \to K'_0(R)$ 是同构。
\end{example}

\begin{example}
\label{algebra-example-K0-polynomial-ring}
设 $k$ 为域。则 $K_0(k[x]) = K'_0(k[x]) = \mathbf{Z}$。
这是例 \ref{algebra-example-K0-PID} 的直接推论，因为
$R = k[x]$ 是主理想整环。
\end{example}

\begin{example}
\label{algebra-example-K0-node}
设 $k$ 为域。令 $R = \{f \in k[x] \mid f(0) = f(1)\}$；
比较例 \ref{algebra-example-affine-open-not-standard}。在此情形下，
$K_0(R) \cong k^* \oplus \mathbf{Z}$，而
$K'_0(R) = \mathbf{Z}$。
\end{example}

\begin{lemma}
\label{algebra-lemma-K0-product}
设 $R = R_1 \times R_2$。则
$K_0(R) = K_0(R_1) \times K_0(R_2)$，且
$K'_0(R) = K'_0(R_1) \times K'_0(R_2)$。
\end{lemma}

\begin{proof}
略。
\end{proof}

\begin{lemma}
\label{algebra-lemma-K0prime-Artinian}
设 $R$ 为阿廷局部环。引理 \ref{algebra-lemma-length-K0} 中的映射
$\text{length}_R : K'_0(R) \to \mathbf{Z}$ 是同构。
\end{lemma}

\begin{proof}
略。
\end{proof}

\begin{lemma}
\label{algebra-lemma-K0-local}
设 $(R, \mathfrak m)$ 为局部环。每个有限投射 $R$-模都是有限自由模。
由 $[M] \to \text{rank}_R(M)$ 定义的映射
$\text{rank}_R : K_0(R) \to \mathbf{Z}$ 是良定义的同构。
\end{lemma}

\begin{proof}
设 $P$ 为有限投射 $R$-模。选取诸元素
$x_1, \ldots, x_n \in P$，使它们的像构成 $P/\mathfrak m P$ 的一组基。
由 Nakayama 引理 \ref{algebra-lemma-NAK}，这些元素生成 $P$。
相应的满射 $u : R^{\oplus n} \to P$ 有分裂，因为 $P$ 是投射模。
因此 $R^{\oplus n} = P \oplus Q$，其中 $Q = \Ker(u)$。
由此 $Q/\mathfrak m Q = 0$，故由 Nakayama 引理，$Q$ 为零。
这样就看出，每个有限投射 $R$-模都是有限自由模。
由引理 \ref{algebra-lemma-rank}，有限自由模具有良定义的秩。
给定有限自由 $R$-模的短正合列
$$
0 \to M' \to M \to M'' \to 0
$$
有 $\text{rank}(M) = \text{rank}(M') + \text{rank}(M'')$，
因为在此情形下 $M \cong M' \oplus M'$（例如，由引理
\ref{algebra-lemma-lift-map}，该列分裂）。由此得到
$K_0(R) = \mathbf{Z}$。
\end{proof}

\begin{lemma}
\label{algebra-lemma-K0-and-K0prime-Artinian-local}
设 $R$ 为阿廷局部环。存在交换图
$$
\xymatrix{
K_0(R) \ar[rr] \ar[d]_{\text{rank}_R} & &
K'_0(R) \ar[d]^{\text{length}_R} \\
\mathbf{Z} \ar[rr]^{\text{length}_R(R)} & &
\mathbf{Z}
}
$$
其中由引理 \ref{algebra-lemma-K0prime-Artinian} 和
\ref{algebra-lemma-K0-local}，两个竖直映射都是同构。
\end{lemma}

\begin{proof}
设 $P$ 为有限投射 $R$-模。须证明
$\text{length}_R(P) = \text{rank}_R(P) \text{length}_R(R)$。
由引理 \ref{algebra-lemma-K0-local}，模 $P$ 是有限自由模。所以对某个
$n \geq 0$，有 $P \cong R^{\oplus n}$。于是
$\text{rank}_R(P) = n$，并且由长度的可加性（引理
\ref{algebra-lemma-length-additive}），
$\text{length}_R(R^{\oplus n}) = n \text{length}_R(R)$。
因此结论成立。
\end{proof}

















\section{分次环}
\label{algebra-section-graded}

\noindent
这里所谓的{\it 分次环}，是指底层阿贝尔群具有直和分解
$S = \bigoplus_{d \geq 0} S_d$ 的环 $S$，并且
$S_d \cdot S_e \subset S_{d + e}$。注意，不允许负次数部分中有非零元。
{\it 无关理想}是理想 $S_{+} = \bigoplus_{d > 0} S_d$。
所谓{\it 分次模}，是指底层阿贝尔群具有直和分解
$M = \bigoplus_{n\in \mathbf{Z}} M_n$ 的 $S$-模 $M$，并且
$S_d \cdot M_e \subset M_{d + e}$。注意，对模允许负次数部分中有非零元。
置 $S_{-k} = (0)$（$k > 0$），便可将 $S$ 看成分次 $S$-模。
$M$（相应地，$S$）的元素 $x$（相应地，$f$）称为{\it 齐次的}，
如果对某个 $d$，有 $x \in M_d$（相应地，$f \in S_d$）。
分次 $S$-模之间的{\it 同态}，是满足
$\varphi(M_d) \subset M'_d$ 的 $S$-模同态 $\varphi : M \to M'$。
这里不允许同态移动次数。用 $\text{GrHom}_0(M, N)$ 表示从 $M$ 到 $N$
的分次模同态所成的 $S_0$-模。

\medskip\noindent
此时自然有分次理想、分次商环、分次子模、分次商模、分次张量积等概念。
相关定义和引理留给读者。例如，分次模的短正合列在每个次数上都是短正合的。

\medskip\noindent
给定分次环 $S$、分次 $S$-模 $M$ 和 $n \in \mathbf{Z}$，
用 $M(n)$ 表示满足 $M(n)_d = M_{n + d}$ 的分次 $S$-模。
它称为{\it $M$ 由 $n$ 的扭转}。特别地，得到诸模
$S(n)$、$n \in \mathbf{Z}$；它们在射影概形的研究中起重要作用。
存在一些显然的函子性同构，例如
$(M \oplus N)(n) = M(n) \oplus N(n)$、
$(M \otimes_S N)(n) = M \otimes_S N(n) = M(n) \otimes_S N$。
此外，还可以在下列 $S_0$-模上定义分次 $S$-模结构：
$$
\text{GrHom}(M, N) =
\bigoplus\nolimits_{n \in \mathbf{Z}} \text{GrHom}_n(M, N),
\quad
\text{GrHom}_n(M, N) = \text{GrHom}_0(M, N(n)).
$$
乘法的定义从略。

\begin{lemma}
\label{algebra-lemma-graded-NAK}
设 $S$ 为分次环，$M$ 为分次 $S$-模。
\begin{enumerate}
\item 若 $S_+M = M$ 且 $M$ 有限，则 $M = 0$。
\item 若 $N, N' \subset M$ 是分次子模，
$M = N + S_+N'$，且 $N'$ 有限，则 $M = N$。
\item 若 $N \to M$ 是分次模同态，$N/S_+N \to M/S_+M$
满射，且 $M$ 有限，则 $N \to M$ 满射。
\item 若齐次元 $x_1, \ldots, x_n \in M$ 生成 $M/S_+M$，
且 $M$ 有限，则 $x_1, \ldots, x_n$ 生成 $M$。
\end{enumerate}
\end{lemma}

\begin{proof}
证明 (1)。选取 $M$ 在 $S$ 上的一组生成元 $y_1, \ldots, y_r$。
可以假设 $y_i$ 是次数为 $d_i$ 的齐次元。重新编号后，可以假设
$d_r = \min(d_i)$。此时条件 $S_+M = M$ 蕴含 $y_r = 0$。
故对 $r$ 归纳即得 $M = 0$。将 (1) 应用于 $M/N$，即得 (2)。
将 (2) 应用于子模 $\Im(N \to M)$ 与 $M$，即得 (3)。
将 (3) 应用于模同态
$\bigoplus S(-d_i) \to M$、$(s_1, \ldots, s_n) \mapsto \sum s_i x_i$，
即得 (4)。
\end{proof}

\noindent
设 $S$ 为分次环，$d \geq 1$ 为整数。置
$S^{(d)} = \bigoplus_{n \geq 0} S_{nd}$。把 $S^{(d)}$ 看成分次环，
其次数为 $n$ 的分量是 $(S^{(d)})_n = S_{nd}$。给定分次
$S$-模 $M$，同样可以考虑
$M^{(d)} = \bigoplus_{n \in \mathbf{Z}} M_{nd}$；
这是一个分次 $S^{(d)}$-模。

\begin{lemma}
\label{algebra-lemma-uple-generated-degree-1}
设 $S$ 为在 $S_0$ 上有限生成的分次环。则对所有充分可除的 $d$，
代数 $S^{(d)}$ 在 $S_0$ 上由次数 $1$ 的部分生成。
\end{lemma}

\begin{proof}
设 $S$ 在 $S_0$ 上由 $f_1, \ldots, f_r \in S$ 生成。
用 $f_i$ 的齐次部分替换它们后，可以假设 $f_i$ 是次数
$d_i > 0$ 的齐次元。于是 $S_n$ 的任意元素都是单项式
$f_1^{e_1} \ldots f_r^{e_r}$ 的 $S_0$-系数线性组合，
其中 $\sum e_i d_i = n$。设 $m$ 是 $\text{lcm}(d_i)$ 的倍数。
对任意 $N \geq r$，若
$$
\sum e_i d_i = N m
$$
则由一个初等论证，对某个 $i$ 有 $e_i \geq m/d_i$。
所以次数为 $N m$ 的每个单项式，都是次数为 $m$ 的单项式
$f_i^{m/d_i}$ 与次数为 $(N - 1)m$ 的单项式之积。
由此，任意次数为 $nrm$、其中 $n \geq 2$ 的单项式，都是次数为
$rm$ 的诸单项式之积。因此 $S^{(rm)}$ 在 $S_0$ 上由次数 $1$ 的部分生成。
\end{proof}

\begin{lemma}
\label{algebra-lemma-integral-closure-graded}
\begin{reference}
\cite[Theorem 2.3.2]{Huneke-Swanson}
\end{reference}
设 $R \to S$ 为分次环同态，$S' \subset S$ 为 $R$ 在 $S$ 中的整闭包。
则
$$
S' = \bigoplus\nolimits_{d \geq 0} S' \cap S_d,
$$
即 $S'$ 是 $S$ 的分次 $R$-子代数。
\end{lemma}

\begin{proof}
须证明：若 $s = s_n + s_{n + 1} + \ldots + s_m \in S'$，
则每个齐次部分 $s_j \in S'$。对所有分次环同态
$R \to S$，关于 $m - n$ 归纳证明。先注意，$s_0$ 显然在 $R_0$
上整（因而在 $R$ 上整），因为存在与环同态 $R \to R_0$ 相容的
环同态 $S \to S_0$。故以 $s - s_0$ 代替 $s$ 后，可以假设
$n > 0$。考虑分次环扩张 $R[t, t^{-1}] \to S[t, t^{-1}]$，
其中 $t$ 的次数为 $0$。存在交换图
$$
\xymatrix{
S[t, t^{-1}] \ar[rr]_{s \mapsto t^{\deg(s)}s} & & S[t, t^{-1}] \\
R[t, t^{-1}] \ar[u] \ar[rr]^{r \mapsto t^{\deg(r)}r} & &  R[t, t^{-1}] \ar[u]
}
$$
其中横向映射都是环自同构。因此，$S[t, t^{-1}]$ 在
$R[t, t^{-1}]$ 上的整闭包 $C$ 被映入自身。于是
$$
t^m(s_n + s_{n + 1} + \ldots + s_m) -
(t^ns_n + t^{n + 1}s_{n + 1} + \ldots + t^ms_m) \in C
$$
由归纳假设可知，对 $i = n, \ldots, m - 1$，每个
$(t^m - t^i)s_i \in C$。注意，对任意环 $A$ 和
$m > i \geq n > 0$，有
$A[t, t^{-1}]/(t^m - t^i - 1) \cong A[t]/(t^m - t^i - 1) \supset A$，
因为在 $A[t]/(t^m - t^i - 1)$ 中，
$t(t^{m - 1} - t^{i - 1}) = 1$。由于 $t^m - t^i$ 映为 $1$，
可知对 $i = n, \ldots, m - 1$，$s_i$ 在环
$S[t]/(t^m - t^i - 1)$ 中的像，在 $R[t]/(t^m - t^i - 1)$ 上整。
由于 $R \to R[t]/(t^m - t^i - 1)$ 有限，由传递性可知
$s_i$ 在 $R$ 上整；参见引理 \ref{algebra-lemma-integral-transitive}。
最后还可得
$s_m = s - \sum_{i = n, \ldots, m - 1} s_i$ 在 $R$ 上整。
\end{proof}







\section{分次环的 Proj}
\label{algebra-section-proj}

\noindent
设 $S$ 为分次环。所谓{\it 齐次理想}，就是既是该环的理想、
又是 $S$ 的分次子模的 $I \subset S$。等价地，它是由齐次元生成的理想。
再等价地，若 $f \in I$ 且
$$
f = f_0 + f_1 + \ldots + f_n
$$
是 $f$ 在 $S$ 中的齐次部分分解，则对每个 $i$ 都有 $f_i \in I$。
要检验齐次理想 $\mathfrak p$ 是否为素理想，只须检验：
若 $ab \in \mathfrak p$，其中 $a, b$ 齐次，则
$a \in \mathfrak p$ 或 $b \in \mathfrak p$。

\begin{definition}
\label{algebra-definition-proj}
设 $S$ 为分次环。定义 $\text{Proj}(S)$ 为 $S$ 的所有满足
$S_{+} \not \subset \mathfrak p$ 的齐次素理想 $\mathfrak p$ 所成的集合。
集合 $\text{Proj}(S)$ 是 $\Spec(S)$ 的子集，赋予它诱导拓扑。
拓扑空间 $\text{Proj}(S)$ 称为分次环 $S$ 的{\it 齐次谱}。
\end{definition}

\noindent
注意，由构造有连续映射
$$
\text{Proj}(S) \longrightarrow \Spec(S_0).
$$

\medskip\noindent
设 $S = \oplus_{d \geq 0} S_d$ 为分次环。取 $f\in S_d$，
并假设 $d \geq 1$。定义 $S_{(f)}$ 为 $S_f$ 的子环，
其元素具有形式 $r/f^n$，其中 $r$ 齐次且 $\deg(r) = nd$。
若 $M$ 是分次 $S$-模，则定义 $S_{(f)}$-模 $M_{(f)}$
为 $M_f$ 的子模，其元素具有形式 $x/f^n$，其中 $x$ 是次数为
$nd$ 的齐次元。

\begin{lemma}
\label{algebra-lemma-Z-graded}
设 $S$ 是含有一个正次数齐次可逆元的 $\mathbf{Z}$-分次环。
则 $S$ 的 $\mathbf{Z}$-分次素理想所成集合
$G \subset \Spec(S)$（赋予诱导拓扑）同胚地映到 $\Spec(S_0)$。
\end{lemma}

\begin{proof}
先构造逆映射以证明该映射是双射。设 $f \in S_d$、$d > 0$
在 $S$ 中可逆。若 $\mathfrak p_0$ 是 $S_0$ 的素理想，则
$\mathfrak p_0S$ 是 $S$ 的 $\mathbf{Z}$-分次理想，且
$\mathfrak p_0S \cap S_0 = \mathfrak p_0$。若
$ab \in \mathfrak p_0S$，其中 $a$、$b$ 齐次，则
$a^db^d/f^{\deg(a) + \deg(b)} \in \mathfrak p_0$。
所以 $a^d/f^{\deg(a)} \in \mathfrak p_0$ 或
$b^d/f^{\deg(b)} \in \mathfrak p_0$；换言之，
$a^d \in \mathfrak p_0S$ 或 $b^d \in \mathfrak p_0S$。
因此 $\sqrt{\mathfrak p_0S}$ 是 $S$ 的 $\mathbf{Z}$-分次素理想，
且它与 $S_0$ 的交为 $\mathfrak p_0$。

\medskip\noindent
为证明该映射是同胚，只须证明 $G \cap D(g)$ 的像是开集。
若 $g = \sum g_i$，其中 $g_i \in S_i$，则由上述结论，
$G \cap D(g)$ 映满开集 $\bigcup D(g_i^d/f^i)$。
\end{proof}

\noindent
对齐次元 $f \in S$，若其次数为 $> 0$，则定义
$$
D_{+}(f) = \{ \mathfrak p \in \text{Proj}(S) \mid f \not\in \mathfrak p \}.
$$
最后，对齐次理想 $I \subset S$，定义
$$
V_{+}(I) = \{ \mathfrak p \in \text{Proj}(S) \mid I \subset \mathfrak p \}.
$$
更一般地，对齐次元的任意集合 $E$、$E \subset S$，也使用记号
$V_{+}(E)$。

\begin{lemma}[Proj 上的拓扑]
\label{algebra-lemma-topology-proj}
设 $S = \oplus_{d \geq 0} S_d$ 为分次环。
\begin{enumerate}
\item 诸集合 $D_{+}(f)$ 是 $\text{Proj}(S)$ 中的开集。
\item 有 $D_{+}(ff') = D_{+}(f) \cap D_{+}(f')$。
\item 设 $g = g_0 + \ldots + g_m$ 是 $S$ 的元素，且
$g_i \in S_i$。则
$$
D(g) \cap \text{Proj}(S) =
(D(g_0) \cap \text{Proj}(S))
\cup
\bigcup\nolimits_{i \geq 1} D_{+}(g_i).
$$
\item
设 $g_0\in S_0$ 是次数为 $0$ 的齐次元。则
$$
D(g_0) \cap \text{Proj}(S)
=
\bigcup\nolimits_{f \in S_d, \ d\geq 1} D_{+}(g_0 f).
$$
\item 开集 $D_{+}(f)$ 构成 $\text{Proj}(S)$ 的拓扑的一组基。
\item 设 $f \in S$ 是正次数齐次元。环 $S_f$ 有自然的
$\mathbf{Z}$-分次。环同态 $S \to S_f \leftarrow S_{(f)}$ 诱导同胚
$$
D_{+}(f)
\leftarrow
\{\mathbf{Z}\text{-分次素理想，位于 }S_f\}
\to
\Spec(S_{(f)}).
$$
\item 存在某个 $S$，使得 $\text{Proj}(S)$ 不是拟紧的。
\item 诸集合 $V_{+}(I)$ 是闭集。
\item 任意闭子集 $T \subset \text{Proj}(S)$ 都具有形式
$V_{+}(I)$，其中 $I \subset S$ 是某个齐次理想。
\item 对任意分次理想 $I \subset S$，有
$V_{+}(I) = \emptyset$，当且仅当 $S_{+} \subset \sqrt{I}$。
\end{enumerate}
\end{lemma}

\begin{proof}
由于 $D_{+}(f) = \text{Proj}(S) \cap D(f)$，这些集合是开集，
从而 (1) 得证。又因为 $D(ff') = D(f) \cap D(f')$，(2) 亦得证。
同理，诸集合 $V_{+}(I) = \text{Proj}(S) \cap V(I)$ 是闭集，
这证明了 (8)。

\medskip\noindent
假设 $T \subset \text{Proj}(S)$ 是闭集。则对某个理想
$J \subset S$，可写成 $T = \text{Proj}(S) \cap V(J)$。
由齐次理想的定义，若 $g \in J$、$g = g_0 + \ldots + g_m$，
其中 $g_d \in S_d$，则对所有 $\mathfrak p \in T$ 都有
$g_d \in \mathfrak p$。因此，令 $I \subset S$ 为由 $J$ 中诸元素的
齐次部分生成的理想，便有 $T = V_{+}(I)$。这证明了 (9)。

\medskip\noindent
由定义立即得到 $\text{Proj}(S) \cap D(g)$ 的公式，其中 $g \in S$。
这证明了 (3)。再考虑 $\text{Proj}(S) \cap D(g_0)$ 的公式。
右端包含于左端是显然的。反过来，假设
$g_0 \not \in \mathfrak p$，其中 $\mathfrak p \in \text{Proj}(S)$。
若对所有正次数齐次元 $f$ 都有 $g_0f \in \mathfrak p$，
则 $S_{+} \subset \mathfrak p$，矛盾。这给出反向包含，
从而证明 (4)。

\medskip\noindent
开集族 $D(g) \cap \text{Proj}(S)$ 构成一组拓扑基，因为标准开集
$D(g) \subset \Spec(S)$ 构成 $\Spec(S)$ 的一组拓扑基。
由上述公式，可将 $D(g) \cap \text{Proj}(S)$ 表成开集
$D_{+}(f)$ 之并。因此，开集族 $D_{+}(f)$ 也构成一组拓扑基。
这证明了 (5)。

\medskip\noindent
证明 (6)。首先，由引理 \ref{algebra-lemma-standard-open}，可将
$D_{+}(f)$ 等同为 $D(f) = \Spec(S_f)$ 的一个子集
（赋予诱导拓扑）。注意，环 $S_f$ 有一个 $\mathbf{Z}$-分次。
其齐次元具有形式 $r/f^n$，其中 $r \in S$ 齐次，且次数为
$\deg(r/f^n) = \deg(r) - n\deg(f)$。子集 $D_{+}(f)$ 恰好对应于
那些为 $\mathbf{Z}$-分次理想（即由齐次元生成）的素理想
$\mathfrak p \subset S_f$。因此只须证明，$S_f$ 的
$\mathbf{Z}$-分次素理想所成的集合同胚地映到
$\Spec(S_{(f)})$。这由引理 \ref{algebra-lemma-Z-graded} 得出。

\medskip\noindent
令 $S = \mathbf{Z}[X_1, X_2, X_3, \ldots]$，并赋予分次，
使每个 $X_i$ 的次数为 $1$。不难看出，覆盖
$$
\text{Proj}(S) = \bigcup\nolimits_{i = 1}^\infty D_{+}(X_i)
$$
不存在有限加细。这证明了 (7)。

\medskip\noindent
设 $I \subset S$ 为分次理想。若 $\sqrt{I} \supset S_{+}$，
则 $V_{+}(I) = \emptyset$，因为按照定义，每个素理想
$\mathfrak p \in \text{Proj}(S)$ 都不包含 $S_{+}$。
反过来，假设 $S_{+} \not \subset \sqrt{I}$。于是可以找到
$f \in S_{+}$，使得 $f$ 模 $I$ 不幂零。显然，这意味着 $f$
的某个齐次部分模 $I$ 不幂零；换言之，可以且确实假设 $f$ 齐次。
这蕴含 $I S_f \not = S_f$，也就是 $(S/I)_f$ 非零。
又因为 $(S/I)_{(f)} \not = 0$，它是映入 $(S/I)_f$ 的环。
选取素理想 $\mathfrak q \subset (S/I)_{(f)}$。它对应于
$S/I$ 的一个不包含无关理想 $(S/I)_{+}$ 的分次素理想，
后者又对应于 $S$ 的一个包含 $I$、但不包含 $S_{+}$ 的分次素理想
$\mathfrak p$，正合所需。这证明了 (10)，并完成证明。
\end{proof}

\begin{example}
\label{algebra-example-proj-polynomial-ring-1-variable}
设 $R$ 为环。若 $S = R[X]$ 且 $\deg(X) = 1$，则自然映射
$\text{Proj}(S) \to \Spec(R)$ 是双射，实际上是同胚。
事实上，假设 $\mathfrak p \in \text{Proj}(S)$。由于
$S_{+} \not \subset \mathfrak p$，可知 $X \not \in \mathfrak p$。
所以，若 $aX^n \in \mathfrak p$，其中 $a \in R$ 且 $n > 0$，
则 $a \in \mathfrak p$。由此
$\mathfrak p = \mathfrak p_0S$，其中
$\mathfrak p_0 = \mathfrak p \cap R$。
\end{example}

\noindent
若 $\mathfrak p \in \text{Proj}(S)$，则定义 $S_{(\mathfrak p)}$
为如下环：其元素是分式 $r/f$，其中 $r, f \in S$ 是同次数的齐次元，
且 $f \not\in \mathfrak p$。与通常一样，规定
$r/f = r'/f'$，当且仅当存在某个齐次元 $f'' \in S$，
满足 $f'' \not \in \mathfrak p$ 且 $f''(rf' - r'f) = 0$。
给定分次 $S$-模 $M$，令 $M_{(\mathfrak p)}$ 为如下
$S_{(\mathfrak p)}$-模：其元素是分式 $x/f$，其中
$x \in M$ 与 $f \in S$ 齐次且次数相同，并且
$f \not \in \mathfrak p$。规定 $x/f = x'/f'$，当且仅当存在某个
齐次元 $f'' \in S$，满足 $f'' \not \in \mathfrak p$ 且
$f''(xf' - x'f) = 0$。

\begin{lemma}
\label{algebra-lemma-proj-prime}
设 $S$ 为分次环，$M$ 为分次 $S$-模。设 $\mathfrak p$ 是
$\text{Proj}(S)$ 的一个元素。设 $f \in S$ 是正次数齐次元，
且 $f \not \in \mathfrak p$，即 $\mathfrak p \in D_{+}(f)$。
设 $\mathfrak p' \subset S_{(f)}$ 是 $\Spec(S_{(f)})$ 中按照
引理 \ref{algebra-lemma-topology-proj} 对应于 $\mathfrak p$ 的元素。则
$S_{(\mathfrak p)} = (S_{(f)})_{\mathfrak p'}$，并且相容地，
$M_{(\mathfrak p)} = (M_{(f)})_{\mathfrak p'}$。
\end{lemma}

\begin{proof}
定义映射 $\psi : M_{(\mathfrak p)} \to (M_{(f)})_{\mathfrak p'}$。
设 $x/g \in M_{(\mathfrak p)}$。置
$$
\psi(x/g) = (x g^{\deg(f) - 1}/f^{\deg(x)})/(g^{\deg(f)}/f^{\deg(g)}).
$$
这是有意义的，因为 $\deg(x) = \deg(g)$，且
$g^{\deg(f)}/f^{\deg(g)} \not \in \mathfrak p'$。
关于 $\psi$ 是良定义的模同态且为同构的验证从略。提示：
其逆映射把 $(x/f^n)/(g/f^m)$ 送到 $(xf^m)/(g f^n)$。
\end{proof}

\noindent
下面是引理 \ref{algebra-lemma-silly} 的一个分次版本。

\begin{lemma}
\label{algebra-lemma-graded-silly}
设 $S$ 为分次环，$\mathfrak p_i$（$i = 1, \ldots, r$）为齐次素理想，
$I \subset S_{+}$ 为分次理想。假设对所有 $i$ 都有
$I \not\subset \mathfrak p_i$。则存在正次数齐次元 $x\in I$，
使得对所有 $i$ 都有 $x\not\in \mathfrak p_i$。
\end{lemma}

\begin{proof}
可以假设诸 $\mathfrak p_i$ 之间不存在包含关系。当 $r = 1$ 时结论成立。
假设对 $r - 1$ 结论成立。选取正次数齐次元 $x \in I$，
使得对所有 $i = 1, \ldots, r - 1$ 都有
$x \not \in \mathfrak p_i$。若 $x \not\in \mathfrak p_r$，证明完毕。
故假设 $x \in \mathfrak p_r$。若
$I \mathfrak p_1 \ldots \mathfrak p_{r-1} \subset \mathfrak p_r$，
则 $I \subset \mathfrak p_r$，矛盾。选取齐次元
$y \in I\mathfrak p_1 \ldots \mathfrak p_{r-1}$，使得
$y \not \in \mathfrak p_r$。此时
$x^{\deg(y)} + y^{\deg(x)}$ 符合要求。
\end{proof}

\begin{lemma}
\label{algebra-lemma-smear-out}
设 $S$ 为分次环，$\mathfrak p \subset S$ 为素理想。
令 $\mathfrak q$ 为由 $\mathfrak p$ 的齐次元生成的 $S$ 的齐次理想。
则 $\mathfrak q$ 是 $S$ 的素理想。
\end{lemma}

\begin{proof}
为证明 $\mathfrak q$ 为素理想，只须检验：若 $f, g \in S$ 齐次且
$fg \in \mathfrak q$，则 $f$ 或 $g$ 属于 $\mathfrak q$。
因为 $\mathfrak p \subset \mathfrak q$，所以
$fg \in \mathfrak p$。由于 $\mathfrak p$ 为素理想，可知
$f \in \mathfrak p$ 或 $g \in \mathfrak p$。
又因为 $f$ 与 $g$ 齐次，显然 $f \in \mathfrak q$ 或
$g \in \mathfrak q$。
\end{proof}

\begin{lemma}
\label{algebra-lemma-graded-ring-minimal-prime}
设 $S$ 为分次环。
\begin{enumerate}
\item $S$ 的任意极小素理想都是 $S$ 的齐次理想。
\item 给定齐次理想 $I \subset S$，$I$ 上的任意极小素理想都是齐次的。
\end{enumerate}
\end{lemma}

\begin{proof}
第一条成立，因为引理 \ref{algebra-lemma-smear-out} 中构造的素理想
$\mathfrak q$ 满足 $\mathfrak q \subset \mathfrak p$。
第二条成立，因为可以考虑 $S/I$ 并应用第一条。
\end{proof}

\begin{lemma}
\label{algebra-lemma-dehomogenize-finite-type}
设 $R$ 为环，$S$ 为分次 $R$-代数，$f \in S_{+}$ 为齐次元。
假设 $S$ 在 $R$ 上是有限型的。则
\begin{enumerate}
\item 环 $S_{(f)}$ 在 $R$ 上是有限型的；
\item 对任意有限分次 $S$-模 $M$，模 $M_{(f)}$ 是有限
$S_{(f)}$-模。
\end{enumerate}
\end{lemma}

\begin{proof}
选取 $f_1, \ldots, f_n \in S$，使它们作为 $R$-代数生成 $S$。
可以假设每个 $f_i$ 都齐次（把每个 $f_i$ 分解成其齐次分量即可）。
$S_{(f)}$ 的一个元素是如下形式的和：
$$
\sum\nolimits_{e\deg(f) =
\sum e_i\deg(f_i)} \lambda_{e_1 \ldots e_n} f_1^{e_1} \ldots f_n^{e_n}/f^e
$$
其中 $\lambda_{e_1 \ldots e_n} \in R$。因此，作为 $R$-代数，
$S_{(f)}$ 由满足 $e\deg(f) = \sum e_i\deg(f_i)$ 的诸元素
$f_1^{e_1} \ldots f_n^{e_n} /f^e$ 生成。若
$e_i \geq \deg(f)$，则可将它写成
$$
f_1^{e_1} \ldots f_n^{e_n}/f^e =
f_i^{\deg(f)}/f^{\deg(f_i)} \cdot
f_1^{e_1} \ldots f_i^{e_i - \deg(f)} \ldots f_n^{e_n}/f^{e - \deg(f_i)}
$$
所以，只需要诸元素 $f_i^{\deg(f)}/f^{\deg(f_i)}$，以及满足
$e \deg(f) = \sum e_i \deg(f_i)$ 和 $e_i < \deg(f)$ 的诸元素
$f_1^{e_1} \ldots f_n^{e_n} /f^e$。这是一个有限表，故 (1) 成立。

\medskip\noindent
为证明 (2)，假设 $M$ 由齐次元 $x_1, \ldots, x_m$ 生成。
与上述论证相同，可知 $M_{(f)}$ 作为 $S_{(f)}$-模由如下有限表中的
元素生成：
$f_1^{e_1} \ldots f_n^{e_n} x_j /f^e$，其中
$e \deg(f) = \sum e_i \deg(f_i) + \deg(x_j)$ 且
$e_i < \deg(f)$。
\end{proof}

\begin{lemma}
\label{algebra-lemma-homogenize}
设 $R$ 为环，$R'$ 为有限型 $R$-代数，$M$ 为有限 $R'$-模。
则存在分次 $R$-代数 $S$、分次 $S$-模 $N$ 和一个次数为 $1$
的齐次元 $f \in S$，使得
\begin{enumerate}
\item $R' \cong S_{(f)}$ 且 $M \cong N_{(f)}$（作为模）；
\item $S_0 = R$，且 $S$ 在 $R$ 上由有限多个次数为 $1$ 的元素生成；
\item $N$ 是有限 $S$-模。
\end{enumerate}
\end{lemma}

\begin{proof}
对某个理想 $I$，可写 $R' = R[x_1, \ldots, x_n]/I$。
对元素 $g \in R[x_1, \ldots, x_n]$，用
$\tilde g \in R[X_0, \ldots, X_n]$ 表示满足
$g = \tilde g(1, x_1, \ldots, x_n)$ 的最小次数齐次元。
令 $\tilde I \subset R[X_0, \ldots, X_n]$ 为由所有元素
$\tilde g$（$g \in I$）生成的理想。置
$S = R[X_0, \ldots, X_n]/\tilde I$，并以 $f$ 表示 $X_0$
在 $S$ 中的像。由构造有同构
$$
S_{(f)} \longrightarrow R', \quad
X_i/X_0 \longmapsto x_i.
$$
为对模 $M$ 作同样的构造，选取表示
$$
M = (R')^{\oplus r}/\sum\nolimits_{j \in J} R'k_j
$$
其中 $k_j = (k_{1j}, \ldots, k_{rj})$。令
$d_{ij} = \deg(\tilde k_{ij})$，置 $d_j = \max\{d_{ij}\}$。
置 $K_{ij} = X_0^{d_j - d_{ij}}\tilde k_{ij}$；
它是次数为 $d_j$ 的齐次元。用这些记号，置
$$
N = \Coker\Big(
\bigoplus\nolimits_{j \in J} S(-d_j) \xrightarrow{(K_{ij})} S^{\oplus r}
\Big)
$$
即可。若干细节从略。
\end{proof}




\section{诺特分次环}
\label{algebra-section-noetherian-graded}

\noindent
本节给出一些诺特分次环的理论，包括关于希尔伯特多项式的若干材料。

\begin{lemma}
\label{algebra-lemma-S-plus-generated}
设 $S$ 为分次环。一族齐次元 $f_i \in S_{+}$ 作为 $S_0$-代数生成
$S$，当且仅当它们作为 $S$ 的理想生成 $S_{+}$。
\end{lemma}

\begin{proof}
若诸 $f_i$ 作为 $S_0$-代数生成 $S$，则 $S_{+}$ 中的每个元素
都是诸 $f_i$ 的无常数项多项式，因此 $S_{+}$ 作为理想由诸
$f_i$ 生成。反过来，假设 $S_{+} = \sum Sf_i$。将证明 $S$ 的任意元素
$f$ 都可以写成以诸 $f_i$ 为变量、系数在 $S_0$ 中的多项式。
只须对齐次元证明。设 $f$ 的次数为 $d$，对 $d$ 归纳。
$d = 0$ 的情形显然。若 $d > 0$，则 $f \in S_{+}$，故对某些
$g_i \in S$，可写成 $f = \sum g_i f_i$。由于 $S$ 分次，
可以用 $g_i$ 的次数为 $d - \deg(f_i)$ 的齐次分量代替它。
由归纳假设，每个 $g_i$ 都是诸 $f_i$ 的多项式，结论随即成立。
\end{proof}

\begin{lemma}
\label{algebra-lemma-graded-Noetherian}
分次环 $S$ 是诺特环，当且仅当 $S_0$ 是诺特环，且
$S_{+}$ 作为 $S$ 的理想有限生成。
\end{lemma}

\begin{proof}
显然，若 $S$ 是诺特环，则 $S_0 = S/S_{+}$ 是诺特环，且
$S_{+}$ 有限生成。反过来，假设 $S_0$ 是诺特环，且 $S_{+}$
作为 $S$ 的理想有限生成。选取生成元
$S_{+} = (f_1, \ldots, f_n)$。将诸 $f_i$ 分解为齐次部分后，
可以假设每个 $f_i$ 都齐次。由引理 \ref{algebra-lemma-S-plus-generated}，
把 $X_i$ 送到 $f_i$ 的映射
$S_0[X_1, \ldots X_n] \to S$ 是满射。因此由引理
\ref{algebra-lemma-Noetherian-permanence}，$S$ 是诺特环。
\end{proof}

\begin{definition}
\label{algebra-definition-numerical-polynomial}
设 $A$ 为阿贝尔群。设函数 $f : n \mapsto f(n) \in A$
对所有充分大的整数 $n$ 有定义。若存在 $r \geq 0$ 及元素
$a_0, \ldots, a_r\in A$，使得对所有 $n \gg 0$ 都有
$$
f(n) = \sum\nolimits_{i = 0}^r \binom{n}{i} a_i
$$
则称该函数为{\it 数值多项式}。
\end{definition}

\noindent
使用二项式系数的理由是如下初等事实：若多项式
$P \in \mathbf{Q}[T]$ 在所有整数点的值都是整数，则它等于某个和
$P(T) = \sum a_i \binom{T}{i}$，其中 $a_i \in \mathbf{Z}$。
特别地，注意表达式 $\binom{T + 1}{i + 1}$ 具有这种形式。

\begin{lemma}
\label{algebra-lemma-numerical-polynomial-functorial}
若 $A \to A'$ 是阿贝尔群同态，且
$f : n \mapsto f(n) \in A$ 是数值多项式，则其复合也是数值多项式。
\end{lemma}

\begin{proof}
这由定义立即得出。
\end{proof}

\begin{lemma}
\label{algebra-lemma-numerical-polynomial}
假设 $f: n \mapsto f(n) \in A$ 对所有充分大的 $n$ 有定义，
并假设 $n \mapsto f(n) - f(n-1)$ 是数值多项式。
则 $f$ 是数值多项式。
\end{lemma}

\begin{proof}
假设对所有 $n \gg 0$ 都有
$f(n) - f(n-1) = \sum\nolimits_{i = 0}^r \binom{n}{i} a_i$。
置 $g(n) = f(n) - \sum\nolimits_{i = 0}^r \binom{n + 1}{i + 1} a_i$。
于是对所有 $n \gg 0$，有 $g(n) - g(n-1) = 0$。
故 $g$ 最终为常值，设其等于 $a_{-1}$。留给读者证明
$a_{-1} + \sum\nolimits_{i = 0}^r \binom{n + 1}{i + 1} a_i$
具有所需形式（参见本引理上方的说明）。
\end{proof}

\begin{lemma}
\label{algebra-lemma-graded-module-fg}
若 $M$ 是有限生成分次 $S$-模，且 $S$ 在 $S_0$ 上有限生成，
则每个 $M_n$ 都是有限 $S_0$-模。
\end{lemma}

\begin{proof}
设 $M$ 的生成元为 $m_i$，$S$ 的生成元为 $f_i$。
取齐次分量后，可以假设诸 $m_i$ 与诸 $f_i$ 都齐次，
并可假设 $f_i \in S_{+}$。此时显然，每个 $M_n$ 作为
$S_0$-模由次数为 $n$ 的“单项式”$\prod f_i^{e_i} m_j$ 生成。
\end{proof}

\begin{proposition}
\label{algebra-proposition-graded-hilbert-polynomial}
假设 $S$ 是诺特分次环，$M$ 是有限分次 $S$-模。考虑函数
$$
\mathbf{Z} \longrightarrow K'_0(S_0), \quad
n \longmapsto [M_n]
$$
参见引理 \ref{algebra-lemma-graded-module-fg}。若 $S_{+}$ 由次数为 $1$
的元素生成，则该函数是数值多项式。
\end{proposition}

\begin{proof}
对 $S_1$ 的最小生成元数归纳证明。若该数为 $0$，则对所有
$n \gg 0$ 都有 $M_n = 0$，结论成立。为证明归纳步骤，
从一个极小生成元集中选取 $x\in S_1$，使归纳假设适用于分次环
$S/(x)$。

\medskip\noindent
先证明当 $x$ 在 $M$ 上幂零时结论成立。对此，再关于满足
$x^r M  = 0$ 的最小整数 $r$ 归纳。若 $r = 1$，则 $M$ 是
$S/xS$ 上的模，结论由另一个归纳假设成立。若 $r > 1$，
则可以找到短正合列 $0 \to M' \to M \to M'' \to 0$，
使整数 $r', r''$ 都严格小于 $r$。于是已知结论对 $M''$ 与
$M'$ 成立。由于 $K'_0(S_0)$ 中有关系
$
[M_d]  = [M'_d] + [M''_d]
$
故也得到 $M$ 的结论。

\medskip\noindent
若 $x$ 在 $M$ 上不幂零，设 $M' \subset M$ 是 $x$
在其上幂零的最大子模。考虑正合列
$0 \to M' \to M \to M/M' \to 0$，再次可见只须对 $M/M'$
证明结论。换言之，可以假设乘以 $x$ 是单射。

\medskip\noindent
令 $\overline{M} = M/xM$。注意，映射 $x : M \to M$
{\it 不是}分次 $S$-模同态，因为它不把 $M_d$ 映入 $M_d$。
具体而言，对每个 $d$ 都有短正合列
$$
0 \to M_d \xrightarrow{x} M_{d + 1} \to \overline{M}_{d + 1} \to 0
$$
这证明 $[M_{d + 1}] - [M_d] = [\overline{M}_{d + 1}]$。
故由引理 \ref{algebra-lemma-numerical-polynomial} 得到结论。
\end{proof}

\begin{remark}
\label{algebra-remark-period-polynomial}
若 $S$ 仍是诺特环，但 $S$ 不是由次数为 $1$ 的元素生成的，
则与分次 $S$-模相伴的函数是周期多项式
（即对某个 $n$，它在整数模 $n$ 的各同余类上都是数值多项式）。
\end{remark}

\begin{example}
\label{algebra-example-hilbert-function}
假设 $S = k[X_1, \ldots, X_d]$。由例 \ref{algebra-example-K0-field}，
可作等同 $K_0(k) = K'_0(k) = \mathbf{Z}$。因此，任意有限生成分次
$k[X_1, \ldots, X_d]$-模都给出数值多项式
$n \mapsto \dim_k(M_n)$。
\end{example}

\begin{lemma}
\label{algebra-lemma-quotient-smaller-d}
设 $k$ 为域。假设 $I \subset k[X_1, \ldots, X_d]$ 是非零分次理想。
令 $M = k[X_1, \ldots, X_d]/I$。则数值多项式
$n \mapsto \dim_k(M_n)$（参见例 \ref{algebra-example-hilbert-function}）
的次数为 $ < d - 1$（若 $d = 1$，则它为零）。
\end{lemma}

\begin{proof}
与分次模 $k[X_1, \ldots, X_d]$ 相伴的数值多项式是
$n \mapsto \binom{n - 1 + d}{d - 1}$。对任意次数为 $e$ 的非零齐次元
$f \in I$ 以及任意次数 $n >> e$，都有
$I_n \supset f \cdot k[X_1, \ldots, X_d]_{n-e}$，因而
$\dim_k(I_n) \geq \binom{n - e - 1 + d}{d - 1}$。所以
$\dim_k(M_n) \leq \binom{n - 1 + d}{d - 1} - \binom{n - e - 1 + d}{d - 1}$。
最后一个表达式的次数为 $ < d - 1$（若 $d = 1$，则它为零），
故结论成立。
\end{proof}








\section{诺特局部环}
\label{algebra-section-Noetherian-local}

\noindent
在本节中，$(R, \mathfrak m, \kappa)$ 始终表示诺特局部环。
本节将发展一些模的希尔伯特函数理论。设 $M$ 为有限 $R$-模。
定义 $M$ 的{\it 希尔伯特函数}为函数
$$
\varphi_M : n
\longmapsto
\text{length}_R(\mathfrak m^nM/{\mathfrak m}^{n + 1}M)
$$
它对所有整数 $n \geq 0$ 有定义。另一个重要不变量是函数
$$
\chi_M : n
\longmapsto
\text{length}_R(M/{\mathfrak m}^{n + 1}M)
$$
它也对所有整数 $n \geq 0$ 有定义。注意，由引理
\ref{algebra-lemma-length-additive} 有
$$
\chi_M(n) = \sum\nolimits_{i = 0}^n \varphi_M(i).
$$
这一构造还有一个使用定义理想的变体。

\begin{definition}
\label{algebra-definition-ideal-definition}
设 $(R, \mathfrak m)$ 为诺特局部环。满足
$\sqrt{I} = \mathfrak m$ 的理想 $I \subset R$ 称为
{\it $R$ 的定义理想}。
\end{definition}

\noindent
设 $I \subset R$ 为定义理想。因为 $R$ 是诺特环，这意味着对某个
$r$ 有 $\mathfrak m^r \subset I$；参见引理
\ref{algebra-lemma-Noetherian-power}。因此，被 $I$ 的某个幂零化的任意有限
$R$-模都有有限长度；参见引理 \ref{algebra-lemma-length-finite}。
所以可以定义
$$
\varphi_{I, M}(n) = \text{length}_R(I^nM/I^{n + 1}M)
\quad\text{且}\quad
\chi_{I, M}(n) = \text{length}_R(M/I^{n + 1}M)
$$
其中 $n \geq 0$。同样有
$$
\chi_{I, M}(n) = \sum\nolimits_{i = 0}^n \varphi_{I, M}(i).
$$

\begin{lemma}
\label{algebra-lemma-differ-finite}
假设 $M' \subset M$ 是有限 $R$-模，且商的长度有限。
则存在常数 $c_1, c_2$，使得对所有 $n \geq c_2$ 都有
$$
c_1 + \chi_{I, M'}(n - c_2) \leq \chi_{I, M}(n) \leq
c_1 + \chi_{I, M'}(n)
$$
\end{lemma}

\begin{proof}
由于 $M/M'$ 长度有限，存在 $c_2 \geq 0$，使得
$I^{c_2}M \subset M'$。令 $c_1 = \text{length}_R(M/M')$。
对 $n \geq c_2$，有
\begin{eqnarray*}
\chi_{I, M}(n)
& = &
\text{length}_R(M/I^{n + 1}M) \\
& = &
c_1 + \text{length}_R(M'/I^{n + 1}M) \\
& \leq &
c_1 + \text{length}_R(M'/I^{n + 1}M') \\
& = &
c_1 + \chi_{I, M'}(n)
\end{eqnarray*}
另一方面，由于 $I^{c_2}M \subset M'$，对 $n \geq c_2$ 有
$I^nM \subset I^{n - c_2}M'$。因此，对 $n \geq c_2$ 得
\begin{eqnarray*}
\chi_{I, M}(n)
& = &
\text{length}_R(M/I^{n + 1}M) \\
& = &
c_1 + \text{length}_R(M'/I^{n + 1}M) \\
& \geq &
c_1 + \text{length}_R(M'/I^{n + 1 - c_2}M') \\
& = &
c_1 + \chi_{I, M'}(n - c_2)
\end{eqnarray*}
证明完毕。
\end{proof}

\begin{lemma}
\label{algebra-lemma-hilbert-ses}
假设 $0 \to M' \to M \to M'' \to 0$ 是有限 $R$-模的短正合列。
则存在余长度为 $l$ 的子模 $N \subset M'$ 及 $c \geq 0$，使得
$$
\chi_{I, M}(n) = \chi_{I, M''}(n) + \chi_{I, N}(n - c) + l
$$
以及
$$
\varphi_{I, M}(n) = \varphi_{I, M''}(n) + \varphi_{I, N}(n - c)
$$
对所有 $n \geq c$ 成立。
\end{lemma}

\begin{proof}
注意，$M/I^nM \to M''/I^nM''$ 是满射，其核为
$M' / M' \cap I^nM$。由阿廷--里斯引理
\ref{algebra-lemma-Artin-Rees}，存在常数 $c$，使得
$M' \cap I^nM = I^{n - c}(M' \cap I^cM)$。记
$N = M' \cap I^cM$。注意，$I^c M' \subset N \subset M'$。
因此，对 $n \geq c$，
$\text{length}_R(M' / M' \cap I^nM)
= \text{length}_R(M'/N) + \text{length}_R(N/I^{n - c}N)$。
由短正合列
$$
0 \to M' / M' \cap I^nM \to M/I^nM \to M''/I^nM'' \to 0
$$
及长度的可加性（引理 \ref{algebra-lemma-length-additive}），得到等式
$$
\chi_{I, M}(n - 1)
=
\chi_{I, M''}(n - 1)
+
\chi_{I, N}(n - c - 1)
+
\text{length}_R(M'/N)
$$
它对 $n \geq c$ 成立。有
$\varphi_{I, M}(n) = \chi_{I, M}(n) - \chi_{I, M}(n - 1)$，
对模 $M''$ 与 $N$ 亦然。因此，对 $n \geq c$ 得
$\varphi_{I, M}(n) = \varphi_{I, M''}(n) + \varphi_{I, N}(n-c)$。
\end{proof}

\begin{lemma}
\label{algebra-lemma-hilbert-change-I}
假设 $I$、$I'$ 是诺特局部环 $R$ 的两个定义理想。
设 $M$ 为有限 $R$-模。则存在常数 $a$，使得对 $n \geq 1$ 有
$\chi_{I, M}(n) \leq \chi_{I', M}(an)$。
\end{lemma}

\begin{proof}
存在整数 $c \geq 1$，使得 $(I')^c \subset I$。
因此有满射 $M/(I')^{c(n + 1)}M \to M/I^{n + 1}M$。
所以取 $a = 2c - 1$ 即得结论。
\end{proof}

\begin{proposition}
\label{algebra-proposition-hilbert-function-polynomial}
设 $R$ 为诺特局部环，$M$ 为有限 $R$-模，$I \subset R$ 为定义理想。
希尔伯特函数 $\varphi_{I, M}$ 与函数 $\chi_{I, M}$ 都是数值多项式。
\end{proposition}

\begin{proof}
考虑分次环
$S = R/I \oplus I/I^2 \oplus I^2/I^3 \oplus
\ldots = \bigoplus_{d \geq 0} I^d/I^{d + 1}$。再考虑分次
$S$-模
$N = M/IM \oplus IM/I^2M \oplus \ldots =
\bigoplus_{d \geq 0} I^dM/I^{d + 1}M$。这一对 $(S, N)$ 满足命题
\ref{algebra-proposition-graded-hilbert-polynomial} 的假设。
因此，关于 $\varphi_{I, M}$ 的结论由该命题和引理
\ref{algebra-lemma-length-K0} 得出；关于 $\chi_{I, M}$ 的结论再由此及
引理 \ref{algebra-lemma-numerical-polynomial} 得出。
\end{proof}

\begin{definition}
\label{algebra-definition-hilbert-polynomial}
设 $R$ 为诺特局部环，$M$ 为有限 $R$-模。$M$ 在 $R$ 上的
{\it 希尔伯特多项式}，是满足对 $n \gg 0$ 有
$P(n) = \varphi_M(n)$ 的元素 $P(t) \in \mathbf{Q}[t]$。
\end{definition}

\noindent
由命题 \ref{algebra-proposition-hilbert-function-polynomial} 可知，
希尔伯特多项式存在。

\begin{lemma}
\label{algebra-lemma-d-independent}
设 $R$ 为诺特局部环，$M$ 为有限 $R$-模。
\begin{enumerate}
\item 数值多项式 $\varphi_{I, M}$ 的次数不依赖于定义理想 $I$。
\item 数值多项式 $\chi_{I, M}$ 的次数不依赖于定义理想 $I$。
\end{enumerate}
\end{lemma}

\begin{proof}
第 (2) 条由引理 \ref{algebra-lemma-hilbert-change-I} 立即得出。
第 (1) 条由第 (2) 条得出，因为对 $n \geq 1$ 有
$\varphi_{I, M}(n) = \chi_{I, M}(n) - \chi_{I, M}(n - 1)$。
\end{proof}

\begin{definition}
\label{algebra-definition-d}
设 $R$ 为诺特局部环，$M$ 为有限 $R$-模。用{\it $d(M)$}
表示 $\{-\infty, 0, 1, 2, \ldots \}$ 中按如下方式定义的元素：
\begin{enumerate}
\item 若 $M = 0$，置 $d(M) = -\infty$；
\item 若 $M \not = 0$，则 $d(M)$ 是数值多项式 $\chi_M$ 的次数。
\end{enumerate}
\end{definition}

\noindent
若对所有 $n$ 都有 $\mathfrak m^nM \not = 0$，则 $d(M)$ 是
$M$ 的希尔伯特多项式的次数 $+1$。

\begin{lemma}
\label{algebra-lemma-differ-finite-chi}
设 $R$ 为诺特局部环，$I \subset R$ 为定义理想，$M$ 为不具有
有限长度的有限 $R$-模。若 $M' \subset M$ 是余长度有限的子模，
则 $\chi_{I, M} - \chi_{I, M'}$ 的次数为
$<$ 这两个多项式中任一个的次数。
\end{lemma}

\begin{proof}
由引理 \ref{algebra-lemma-differ-finite} 及初等计算得出。
\end{proof}

\begin{lemma}
\label{algebra-lemma-hilbert-ses-chi}
设 $R$ 为诺特局部环，$I \subset R$ 为定义理想，并设
$0 \to M' \to M \to M'' \to 0$ 是有限 $R$-模的短正合列。则
\begin{enumerate}
\item 若 $M'$ 不具有有限长度，则
$\chi_{I, M} - \chi_{I, M''} - \chi_{I, M'}$ 是次数为
$<$ $\chi_{I, M'}$ 的次数的数值多项式；
\item $\max\{ \deg(\chi_{I, M'}), \deg(\chi_{I, M''}) \} = \deg(\chi_{I, M})$；
\item $\max\{d(M'), d(M'')\} = d(M)$。
\end{enumerate}
\end{lemma}

\begin{proof}
先证明 (1)。设 $N \subset M'$ 如引理 \ref{algebra-lemma-hilbert-ses} 中所述。
由引理 \ref{algebra-lemma-differ-finite-chi}，数值多项式
$\chi_{I, M'} - \chi_{I, N}$ 的次数为
$<$ $\chi_{I, M'}$ 与 $\chi_{I, N}$ 的共同次数。
由引理 \ref{algebra-lemma-hilbert-ses}，差
$$
\chi_{I, M}(n) - \chi_{I, M''}(n) - \chi_{I, N}(n - c)
$$
对 $n \gg 0$ 为常数。由初等计算，
$\chi_{I, N}(n) - \chi_{I, N}(n - c)$ 的次数为
$<$ $\chi_{I, N}$ 的次数，后者大于零（见上文）。
合并这些结论即得 (1)。

\medskip\noindent
注意，$\chi_{I, M'}$ 与 $\chi_{I, M''}$ 的首项系数非负。
所以 $\chi_{I, M'} + \chi_{I, M''}$ 的次数等于二者次数的最大值。
因此，若 $M'$ 不具有有限长度，则 (2) 由 (1) 得出。
若 $M'$ 长度有限，则由阿廷--里斯定理（引理
\ref{algebra-lemma-Artin-Rees}），对所有 $n \gg 0$，映射
$I^nM \to I^nM''$ 都是同构。所以对所有 $n \gg 0$，
$M/I^nM \to M''/I^nM''$ 是以 $M'$ 为核的满射，因而
$\chi_{I, M}(n) - \chi_{I, M''}(n) = \text{length}(M')$。
该等式对所有 $n \gg 0$ 成立，所以此时 (2) 也成立。

\medskip\noindent
证明 (3)。除非 $M$、$M'$ 或 $M''$ 之一为零，否则它由 (2) 得出。
这些特殊情形的证明从略。
\end{proof}



































\section{维数}
\label{algebra-section-dimension}

\noindent
请与拓扑学，第 \ref{topology-section-krull-dimension} 节比较。

\begin{definition}
\label{algebra-definition-chain-primes}
设 $R$ 为环。所谓一条{\it 素理想链}，是 $R$ 的素理想序列
$\mathfrak p_0 \subset \mathfrak p_1 \subset \ldots \subset \mathfrak p_n$，
满足对 $i = 0, \ldots, n - 1$ 有
$\mathfrak p_i \not = \mathfrak p_{i + 1}$。这条素理想链的
{\it 长度}为 $n$。
\end{definition}

\noindent
回忆，环 $R$ 的素理想与 $\Spec(R)$ 的不可约闭子集之间有一个
反转包含关系的双射；参见引理 \ref{algebra-lemma-irreducible}。

\begin{definition}
\label{algebra-definition-Krull}
环 $R$ 的{\it 克鲁尔维数}，是拓扑空间 $\Spec(R)$ 的克鲁尔维数；
参见拓扑学，定义 \ref{topology-definition-Krull}。
换言之，它是满足如下条件的整数 $n\geq 0$ 的上确界：
$R$ 有一条素理想链
$$
\mathfrak p_0
\subset
\mathfrak p_1
\subset
\ldots
\subset
\mathfrak p_n, \quad
\mathfrak p_i \not = \mathfrak p_{i + 1}.
$$
其长度为 $n$。
\end{definition}

\begin{definition}
\label{algebra-definition-height}
环 $R$ 的素理想 $\mathfrak p$ 的{\it 高度}，是局部环
$R_{\mathfrak p}$ 的维数。
\end{definition}

\begin{lemma}
\label{algebra-lemma-dimension-height}
$R$ 的克鲁尔维数，是它的（极大）素理想之高度的上确界。
\end{lemma}

\begin{proof}
这是因为总可以在一条素理想链的末端添加一个极大理想。
\end{proof}

\begin{lemma}
\label{algebra-lemma-Noetherian-dimension-0}
维数为 $0$ 的诺特环是阿廷环。反过来，任意阿廷环都是维数为零的诺特环。
\end{lemma}

\begin{proof}
假设 $R$ 是维数为 $0$ 的诺特环。由引理
\ref{algebra-lemma-Noetherian-topology}，空间 $\Spec(R)$ 是诺特空间。
由拓扑学，引理 \ref{topology-lemma-Noetherian} 可知，
$\Spec(R)$ 只有有限多个不可约分支，设
$\Spec(R) = Z_1 \cup \ldots \cup Z_r$。按照引理
\ref{algebra-lemma-irreducible}，每个 $Z_i = V(\mathfrak p_i)$，
其中 $\mathfrak p_i$ 是极小理想。由于维数为 $0$，
这些 $\mathfrak p_i$ 也是极大理想。因此 $\Spec(R)$ 是以
$\mathfrak p_i$ 为元素的离散拓扑空间。雅可布森根
$\bigcap \mathfrak p_i$ 的所有元素 $f$ 都是幂零的；
否则 $R_f$ 不是零环，从而会有另一个素理想。
由引理 \ref{algebra-lemma-product-local}，
$R$ 等于 $\prod R_{\mathfrak p_i}$。由于
$R_{\mathfrak p_i}$ 也是维数为 $0$ 的诺特环，上述论证表明，
它的根 $\mathfrak p_iR_{\mathfrak p_i}$ 局部幂零。
由引理 \ref{algebra-lemma-Noetherian-power}，对某个 $n \geq 1$ 有
$\mathfrak p_i^nR_{\mathfrak p_i} = 0$。
由引理 \ref{algebra-lemma-length-finite}，可知 $R_{\mathfrak p_i}$
在 $R$ 上具有有限长度。因此，由引理
\ref{algebra-lemma-artinian-finite-length} 可知 $R$ 是阿廷环。

\medskip\noindent
若 $R$ 是阿廷环，则由引理 \ref{algebra-lemma-artinian-finite-length}，
它是诺特环。由引理 \ref{algebra-lemma-artinian-finite-nr-max}、
\ref{algebra-lemma-artinian-radical-nilpotent} 和
\ref{algebra-lemma-product-local} 合并可知，它的所有素理想都是极大理想。
\end{proof}

\noindent
下文将使用定义 \ref{algebra-definition-d} 中定义的不变量 $d(-)$。
先给出一个预备引理。

\begin{lemma}
\label{algebra-lemma-dimension-0-d-0}
设 $R$ 为诺特局部环。则
$\dim(R) = 0 \Leftrightarrow d(R) = 0$。
\end{lemma}

\begin{proof}
这是因为 $d(R) = 0$，当且仅当 $R$ 作为 $R$-模具有有限长度。
参见引理 \ref{algebra-lemma-artinian-finite-length}。
\end{proof}

\begin{proposition}
\label{algebra-proposition-dimension-zero-ring}
设 $R$ 为环。下列条件等价：
\begin{enumerate}
\item $R$ 是阿廷环；
\item $R$ 是诺特环且 $\dim(R) = 0$；
\item $R$ 作为自身上的模具有有限长度；
\item $R$ 是有限多个阿廷局部环的乘积；
\item $R$ 是诺特环，且 $\Spec(R)$ 是有限离散拓扑空间；
\item $R$ 是有限多个维数为 $0$ 的诺特局部环的乘积；
\item $R$ 是有限多个满足 $d(R_i) = 0$ 的诺特局部环 $R_i$ 的乘积；
\item $R$ 是有限多个极大理想幂零的诺特局部环 $R_i$ 的乘积；
\item $R$ 是诺特环，只有有限多个极大理想，且它的雅可布森根理想幂零；
\item $R$ 是诺特环，且它的素理想之间不存在严格包含关系。
\end{enumerate}
\end{proposition}

\begin{proof}
这是引理 \ref{algebra-lemma-product-local}、
\ref{algebra-lemma-artinian-finite-length}、
\ref{algebra-lemma-Noetherian-dimension-0} 和
\ref{algebra-lemma-dimension-0-d-0} 的综合。
\end{proof}

\begin{lemma}
\label{algebra-lemma-height-1}
设 $R$ 为诺特局部环。下列条件等价：
\begin{enumerate}
\item
\label{algebra-item-dim-1}
$\dim(R) = 1$；
\item
\label{algebra-item-d-1}
$d(R) = 1$；
\item
\label{algebra-item-Vx}
存在 $x \in \mathfrak m$，其中 $x$ 非幂零，并使得
$V(x) = \{\mathfrak m\}$；
\item
\label{algebra-item-x}
存在 $x \in \mathfrak m$，其中 $x$ 非幂零，并使得
$\mathfrak m = \sqrt{(x)}$；
\item
\label{algebra-item-ideal-1}
存在由 $1$ 个元素生成的定义理想，且不存在由 $0$ 个元素生成的定义理想。
\end{enumerate}
\end{lemma}

\begin{proof}
首先，假设 $\dim(R) = 1$。设 $\mathfrak p_i$ 为 $R$ 的极小素理想。
由于维数为 $1$，$R$ 唯一的其他素理想是 $\mathfrak m$。
按照引理 \ref{algebra-lemma-Noetherian-irreducible-components}，
这样的极小素理想只有有限多个。因此，可以找到
$x \in \mathfrak m$、$x \not \in \mathfrak p_i$；
参见引理 \ref{algebra-lemma-silly}。所以唯一包含 $x$ 的素理想是
$\mathfrak m$，从而得到 (\ref{algebra-item-Vx})。

\medskip\noindent
若 (\ref{algebra-item-Vx}) 成立，则由引理
\ref{algebra-lemma-Zariski-topology}，有 $\mathfrak m = \sqrt{(x)}$，
从而得到 (\ref{algebra-item-x})。反命题也显然。
(\ref{algebra-item-x}) 与 (\ref{algebra-item-ideal-1}) 的等价性直接由定义得出。

\medskip\noindent
假设 (\ref{algebra-item-ideal-1}) 成立。设 $I = (x)$ 为定义理想。
注意，乘以 $x^n$ 给出从 $R/I$ 到 $I^n/I^{n + 1}$ 的满射，
所以 $\text{length}_R(I^n/I^{n + 1})$ 有界。
因此 $d(R) = 0$ 或 $d(R) = 1$；但由 $0$ 不是定义理想这一假设，
排除 $d(R) = 0$。

\medskip\noindent
假设 (\ref{algebra-item-d-1}) 成立。为导出矛盾，假设存在素理想
$\mathfrak p \subset \mathfrak q \subset \mathfrak m$，
且两个包含都严格。选取某个定义理想 $I \subset R$。
将反复使用引理 \ref{algebra-lemma-hilbert-ses-chi}。首先，将它应用于正合列
$0 \to \mathfrak p \to R \to R/\mathfrak p \to 0$，
可得 $d(R/\mathfrak p) \leq 1$。但它显然不可能为零。
选取 $x\in \mathfrak q$、$x\not \in \mathfrak p$。
考虑短正合列
$$
0 \to R/\mathfrak p \to R/\mathfrak p \to R/(xR + \mathfrak p) \to 0.
$$
这蕴含
$\chi_{I, R/\mathfrak p} - \chi_{I, R/\mathfrak p}
- \chi_{I, R/(xR + \mathfrak p)} = - \chi_{I, R/(xR + \mathfrak p)}$
的次数为 $ < 1$。换言之，
$d(R/(xR + \mathfrak p)) = 0$，因而由引理
\ref{algebra-lemma-dimension-0-d-0}，
$\dim(R/(xR + \mathfrak p)) = 0$。但
$R/(xR + \mathfrak p)$ 有两个不同的素理想
$\mathfrak q/(xR + \mathfrak p)$ 与
$\mathfrak m/(xR + \mathfrak p)$，这给出所需矛盾。
\end{proof}

\begin{proposition}
\label{algebra-proposition-dimension}
设 $R$ 为诺特局部环，$d \geq 0$ 为整数。下列条件等价：
\begin{enumerate}
\item
\label{algebra-item-dim-d}
$\dim(R) = d$；
\item
\label{algebra-item-d-d}
$d(R) = d$；
\item
\label{algebra-item-ideal-d}
存在由 $d$ 个元素生成的定义理想，且不存在由少于 $d$ 个元素生成的
定义理想。
\end{enumerate}
\end{proposition}

\begin{proof}
这个证明实际上与引理 \ref{algebra-lemma-height-1} 的证明相同。
对 $d$ 归纳证明该命题。由引理
\ref{algebra-lemma-dimension-0-d-0} 和 \ref{algebra-lemma-height-1}，
可以假设 $d > 1$。用 $d'(R)$ 表示 $R$ 的定义理想的最小生成元数。
将证明不等式
$\dim(R) \geq d'(R) \geq d(R) \geq \dim(R)$，
从而它们全都相等。

\medskip\noindent
首先，假设 $\dim(R) = d$。设 $\mathfrak p_i$ 为 $R$ 的极小素理想。
按照引理 \ref{algebra-lemma-Noetherian-irreducible-components}，
它们只有有限多个。因此，可以找到
$x \in \mathfrak m$、$x \not \in \mathfrak p_i$；
参见引理 \ref{algebra-lemma-silly}。注意，每条极大素理想链都从某个
$\mathfrak p_i$ 开始，所以 $R/xR$ 的维数至多为 $d-1$。
由归纳假设，存在 $x_2, \ldots, x_d$，使它们在 $R/xR$ 中
生成一个定义理想。因此，$R$ 有一个由（至多）$d$ 个元素生成的定义理想。

\medskip\noindent
假设 $d'(R) = d$。令 $I = (x_1, \ldots, x_d)$ 为定义理想。
注意，乘以 $x_1, \ldots, x_d$ 中所有次数为 $n$ 的单项式，
给出从 $\binom{d + n - 1}{d - 1}$ 个 $R/I$ 之直和到
$I^n/I^{n + 1}$ 的满射。所以
$\text{length}_R(I^n/I^{n + 1})$ 受一个次数为 $d-1$ 的多项式控制。
因此 $d(R) \leq d$。

\medskip\noindent
假设 $d(R) = d$。考虑素理想链
$\mathfrak p \subset \mathfrak q \subset
\mathfrak q_2 \subset \ldots \subset \mathfrak q_e = \mathfrak m$，
其中所有包含都严格，且 $e \geq 2$。选取某个定义理想
$I \subset R$。将反复使用引理 \ref{algebra-lemma-hilbert-ses-chi}。
首先，把它应用于正合列
$0 \to \mathfrak p \to R \to R/\mathfrak p \to 0$，
可得 $d(R/\mathfrak p) \leq d$；但它显然不可能为零。
选取 $x\in \mathfrak q$、$x\not \in \mathfrak p$。
考虑短正合列
$$
0 \to R/\mathfrak p \to R/\mathfrak p \to R/(xR + \mathfrak p) \to 0.
$$
这蕴含
$\chi_{I, R/\mathfrak p} - \chi_{I, R/\mathfrak p}
- \chi_{I, R/(xR + \mathfrak p)} = - \chi_{I, R/(xR + \mathfrak p)}$
的次数为 $ < d$。换言之，
$d(R/(xR + \mathfrak p)) \leq d - 1$，所以由归纳假设，
$\dim(R/(xR + \mathfrak p)) \leq d - 1$。现在，
$R/(xR + \mathfrak p)$ 有素理想链
$\mathfrak q/(xR + \mathfrak p) \subset \mathfrak q_2/(xR + \mathfrak p)
\subset \ldots \subset \mathfrak q_e/(xR + \mathfrak p)$，
由此 $e - 1 \leq d - 1$。由于开始时的素理想链是任意的，
这证明 $\dim(R) \leq d(R)$。

\medskip\noindent
回顾上述论证可见，我们已如所需地证明了这组循环不等式。
\end{proof}

\noindent
设 $(R, \mathfrak m)$ 为诺特局部环。由上述结果显然可知，
$\mathfrak m$ 不可能由少于 $\dim(R)$ 个变量生成。
由 Nakayama 引理 \ref{algebra-lemma-NAK}，$\mathfrak m$ 的最小生成元数等于
$\dim_{\kappa(\mathfrak m)} \mathfrak m/\mathfrak m^2$。
所以有如下基本不等式：
$$
\dim(R) \leq \dim_{\kappa(\mathfrak m)} \mathfrak m/\mathfrak m^2.
$$
事实证明，等号成立的环具有许多良好性质。它们称为正则局部环。

\begin{definition}
\label{algebra-definition-regular-local}
设 $(R, \mathfrak m)$ 为维数为 $d$ 的诺特局部环。
\begin{enumerate}
\item {\it $R$ 的一组参数系}，是生成 $R$ 的一个定义理想的元素序列
$x_1, \ldots, x_d \in \mathfrak m$；
\item 若存在 $x_1, \ldots, x_d \in \mathfrak m$，使得
$\mathfrak m = (x_1, \ldots, x_d)$，则称 $R$ 为
{\it 正则局部环}，并称 $x_1, \ldots, x_d$ 为一组
{\it 正则参数系}。
\end{enumerate}
\end{definition}

\noindent
下列引理由上面诸引理与命题的证明显然可得；这里明确写出，以便引用。

\begin{lemma}
\label{algebra-lemma-minimal-over-1}
设 $R$ 为诺特环，$x \in R$。
\begin{enumerate}
\item 若 $\mathfrak p$ 是 $(x)$ 上的极小素理想，则
$\mathfrak p$ 的高度为 $0$ 或 $1$。
\item 若 $\mathfrak p, \mathfrak q \in \Spec(R)$，且
$\mathfrak q$ 是 $(\mathfrak p, x)$ 上的极小素理想，则
$\mathfrak p$ 与 $\mathfrak q$ 之间不存在素理想。
\end{enumerate}
\end{lemma}

\begin{proof}
证明 (1)。若 $\mathfrak p$ 是 $x$ 上的极小素理想，则
$R_\mathfrak p$ 中唯一包含 $x$ 的素理想是极大理想
$\mathfrak p R_\mathfrak p$。这是因为 $R_\mathfrak p$ 的素理想
与 $R$ 中包含于 $\mathfrak p$ 的素理想作 $1$-对-$1$ 对应；参见引理
\ref{algebra-lemma-spec-localization}。因此，若 $x$ 在
$R_\mathfrak p$ 中不幂零，则引理 \ref{algebra-lemma-height-1} 表明
$\dim(R_\mathfrak p) = 1$。当然，若 $x$ 在
$R_\mathfrak p$ 中幂零，则论证表明
$\mathfrak pR_\mathfrak p$ 是唯一的素理想，因而高度为 $0$。

\medskip\noindent
证明 (2)。由第 (1) 部分可知，$\mathfrak q/\mathfrak p$
在 $R/\mathfrak p$ 中是高度为 $1$ 或 $0$ 的素理想。
这立即蕴含 $\mathfrak p$ 与 $\mathfrak q$ 之间不可能有素理想。
\end{proof}

\begin{lemma}
\label{algebra-lemma-minimal-over-r}
设 $R$ 为诺特环，$f_1, \ldots, f_r \in R$。
\begin{enumerate}
\item 若 $\mathfrak p$ 是 $(f_1, \ldots, f_r)$ 上的极小素理想，
则 $\mathfrak p$ 的高度为 $\leq r$。
\item 若 $\mathfrak p, \mathfrak q \in \Spec(R)$，且
$\mathfrak q$ 是 $(\mathfrak p, f_1, \ldots, f_r)$ 上的极小素理想，
则 $\mathfrak p$ 与 $\mathfrak q$ 之间的每条素理想链长度至多为 $r$。
\end{enumerate}
\end{lemma}

\begin{proof}
证明 (1)。若 $\mathfrak p$ 是 $f_1, \ldots, f_r$ 上的极小素理想，
则 $R_\mathfrak p$ 中唯一包含 $f_1, \ldots, f_r$ 的素理想是
极大理想 $\mathfrak p R_\mathfrak p$。这是因为
$R_\mathfrak p$ 的素理想与 $R$ 中包含于 $\mathfrak p$ 的素理想
作 $1$-对-$1$ 对应；参见引理 \ref{algebra-lemma-spec-localization}。
所以命题 \ref{algebra-proposition-dimension} 表明
$\dim(R_\mathfrak p) \leq r$。

\medskip\noindent
证明 (2)。由第 (1) 部分可知，$\mathfrak q/\mathfrak p$
是高度为 $\leq r$ 的素理想。这立即蕴含关于
$\mathfrak p$ 与 $\mathfrak q$ 之间素理想链的断言。
\end{proof}

\begin{lemma}
\label{algebra-lemma-one-equation}
假设 $R$ 为诺特局部环，$x\in \mathfrak m$ 是其极大理想中的元素。
则 $\dim R \leq \dim R/xR + 1$。若 $x$ 不属于 $R$ 的任何极小素理想，
则等号成立。（例如，$x$ 是非零因子时。）
\end{lemma}

\begin{proof}
若环中的下列元素：
$$
x_1, \ldots, x_{\dim R/xR} \in R
$$
在 $R/xR$ 中的像生成
$R/xR$ 的一个定义理想，则 $x, x_1, \ldots,
x_{\dim R/xR}$ 生成 $R$ 的一个定义理想。所以不等式由命题
\ref{algebra-proposition-dimension} 得出。另一方面，若 $x$ 不属于 $R$
的任何极小素理想，则 $R/xR$ 中的素理想链在 $R$ 中都产生一条
距离成为极大链至少还差一步的链。
\end{proof}

\begin{lemma}
\label{algebra-lemma-elements-generate-ideal-definition}
设 $(R, \mathfrak m)$ 为诺特局部环。假设
$x_1, \ldots, x_d \in \mathfrak m$ 生成一个定义理想，且
$d = \dim(R)$。则对所有 $i = 1, \ldots, d$，有
$\dim(R/(x_1, \ldots, x_i)) = d - i$。
\end{lemma}

\begin{proof}
或者直接由命题 \ref{algebra-proposition-dimension} 的证明得出，
或者对 $d$ 归纳并使用引理 \ref{algebra-lemma-one-equation}。
\end{proof}







\section{维数理论的应用}
\label{algebra-section-applications-dimension-theory}

\noindent
可以利用维数方面的结果证明某些环有无穷谱，并构造更多雅可布森环。

\begin{lemma}
\label{algebra-lemma-Noetherian-local-domain-dim-2-infinite-opens}
设 $R$ 为维数 $\geq 2$ 的诺特局部整环。
则任意非空开子集 $U \subset \Spec(R)$ 都是无限的。
\end{lemma}

\begin{proof}
为导出矛盾，假设 $U \subset \Spec(R)$ 有限。
此时 $(0) \in U$，且 $\{(0)\}$ 是 $U$ 的开子集
（因为 $\{(0)\}$ 的补集是其他各点之闭包的并）。
因此可以假设 $U = \{(0)\}$。设 $\mathfrak m \subset R$ 为极大理想。
可以找到 $x \in \mathfrak m$、$x \not = 0$，使得
$V(x) \cup U = \Spec(R)$。换言之，$D(x) = \{(0)\}$。
特别地，由引理 \ref{algebra-lemma-one-equation} 可知
$\dim(R/xR) = \dim(R) - 1 \geq 1$。
由命题 \ref{algebra-proposition-dimension}，选取
$\overline{y}_2, \ldots, \overline{y}_{\dim(R)} \in R/xR$，
使它们生成 $R/xR$ 的一个定义理想。选取提升
$y_2, \ldots, y_{\dim(R)} \in R$，从而
$x, y_2, \ldots, y_{\dim(R)}$ 在 $R$ 中生成一个定义理想。
由引理 \ref{algebra-lemma-elements-generate-ideal-definition}，这蕴含
$\dim(R/(y_2)) = \dim(R) - 1$ 且
$\dim(R/(y_2, x)) = \dim(R) - 2$。因此存在一个包含 $y_2$
但不包含 $x$ 的素理想 $\mathfrak p$。这与
$D(x) = \{(0)\}$ 矛盾。
\end{proof}

\noindent
若 $k$ 为域，则环 $k[[t]]$，或者 $p$-进数环，都是维数为 $1$、
且只有有限多个素理想的诺特环。这是可能出现这种现象的最大维数。

\begin{lemma}
\label{algebra-lemma-Noetherian-finite-nr-primes}
只有有限多个素理想的诺特环，其维数为 $\leq 1$。
\end{lemma}

\begin{proof}
设 $R$ 为只有有限多个素理想的诺特环。若 $R$ 是局部整环，
则结论由引理
\ref{algebra-lemma-Noetherian-local-domain-dim-2-infinite-opens} 得出。
若 $R$ 是整环，则由局部情形，对所有极大理想 $\mathfrak m$，
$R_\mathfrak m$ 的维数为 $\leq 1$。所以由引理
\ref{algebra-lemma-dimension-height}，$\dim(R) \leq 1$。
对一般的 $R$，则对 $R$ 的每个极小素理想 $\mathfrak q$，
都有 $\dim(R/\mathfrak q) \leq 1$。由于每个素理想都包含一个极小素理想
（引理 \ref{algebra-lemma-Zariski-topology}），这蕴含
$\dim(R) \leq 1$。
\end{proof}

\begin{lemma}
\label{algebra-lemma-finite-type-algebra-finite-nr-primes}
设 $S$ 是域 $k$ 上的非零有限型代数。下列条件等价：
\begin{enumerate}
\item $\dim(S) = 0$；
\item $S$ 只有有限多个素理想；
\item $S$ 只有有限多个极大理想；
\item $\Spec(S)$ 满足引理
\ref{algebra-lemma-ring-with-only-minimal-primes} 中任一等价条件；
\item $\dim_k(S) < \infty$；
\item $S$ 是阿廷环；
\item $\Spec(S)$ 是离散拓扑空间；
\item 在此添加更多条件。
\end{enumerate}
\end{lemma}

\begin{proof}
查看引理 \ref{algebra-lemma-ring-with-only-minimal-primes} 的第 (5) 部分，
由定义立即可知 (1) 等价于 (4)。回忆，$\Spec(S)$ 是清醒、诺特且
雅可布森的；参见引理 \ref{algebra-lemma-spec-spectral}、
\ref{algebra-lemma-Noetherian-topology}、
\ref{algebra-lemma-finite-type-field-Jacobson} 和 \ref{algebra-lemma-jacobson}。
若 $S$ 的维数为 $0$，则每个点都定义一个不可约分支，而不可约分支
只有有限多个（拓扑学，引理 \ref{topology-lemma-Noetherian}）。
所以 (1) 蕴含 (2)。显然，(2) 蕴含 (3)。
若 (3) 成立，则由拓扑学，引理
\ref{topology-lemma-finite-jacobson}，$\Spec(S)$ 离散，
因而 $S$ 的维数为 $0$。

\medskip\noindent
至此已知 (1)--(4) 等价。蕴含 (5) $\Rightarrow$ (6)
是引理 \ref{algebra-lemma-finite-dimensional-algebra}。
蕴含 (6) $\Rightarrow$ (7) 由命题
\ref{algebra-proposition-dimension-zero-ring} 得出。
蕴含 (7) $\Rightarrow$ (4) 是显然的。
反过来，若 $S$ 满足 (1)--(4)，则 $S$ 只有有限多个素理想
$\mathfrak m_1, \ldots, \mathfrak m_r$，且它们全都极大。
注意，由希尔伯特零点定理（定理
\ref{algebra-theorem-nullstellensatz}），$\kappa(\mathfrak m_i)$ 是
$k$ 的有限扩张。由命题 \ref{algebra-proposition-dimension-zero-ring}
还可知 $S$ 是阿廷环。接着，引理
\ref{algebra-lemma-artinian-finite-length} 告诉我们
$\text{length}_S(S) < \infty$。所以由引理
\ref{algebra-lemma-pushdown-module}，$\dim_k(S) < \infty$。
由此得出 (1)--(7) 等价。（注意，证明
$\dim(S) = 0 \Rightarrow \dim_k(S) < \infty$ 的另一种更标准的方法
是使用诺特正规化，但目前尚未建立它。）
\end{proof}

\begin{lemma}
\label{algebra-lemma-noetherian-dim-1-Jacobson}
诺特雅可布森环。
\begin{enumerate}
\item 任意维数为 $1$ 且有无穷多个素理想的诺特整环 $R$
都是雅可布森环。
\item 任意满足下述条件的诺特环都是雅可布森环：
每个素理想 $\mathfrak p$ 或者极大，或者包含于无穷多个素理想中。
\end{enumerate}
\end{lemma}

\begin{proof}
第 (1) 部分是引理 \ref{algebra-lemma-pid-jacobson} 的重新表述。

\medskip\noindent
设 $R$ 是如下诺特环：每个非极大素理想 $\mathfrak p$
都包含于无穷多个素理想中。为导出矛盾，假设
$\Spec(R)$ 不是雅可布森空间。由引理
\ref{algebra-lemma-irreducible} 和 \ref{algebra-lemma-Noetherian-topology}，
可知 $\Spec(R)$ 是清醒诺特拓扑空间。由拓扑学，引理
\ref{topology-lemma-non-jacobson-Noetherian-characterize}，
存在非极大理想 $\mathfrak p \subset R$，使得
$\{\mathfrak p\}$ 是 $\Spec(R)$ 的局部闭子集。
换言之，$\mathfrak p$ 非极大，且 $\{\mathfrak p\}$
是 $V(\mathfrak p)$ 的开子集。考虑满足
$\mathfrak p \subset \mathfrak q$ 的素理想
$\mathfrak q \subset R$。回忆，
$(R/\mathfrak p)_{\mathfrak q}
= R_{\mathfrak q}/\mathfrak pR_{\mathfrak q}$ 的谱上的拓扑
由 $\Spec(R)$ 的拓扑诱导；参见引理
\ref{algebra-lemma-spec-localization} 和 \ref{algebra-lemma-spec-closed}。
所以 $\{(0)\}$ 是
$\Spec((R/\mathfrak p)_{\mathfrak q})$ 的局部闭子集。
由引理 \ref{algebra-lemma-Noetherian-local-domain-dim-2-infinite-opens}，
可知 $\dim((R/\mathfrak p)_{\mathfrak q}) = 1$。
由于这对每个 $\mathfrak q \supset \mathfrak p$ 都成立，
可得 $\dim(R/\mathfrak p) = 1$。此时利用
$\mathfrak p$ 包含于无穷多个素理想中的假设，可知
$\Spec(R/\mathfrak p)$ 是无限的。因此，由本引理第 (1) 部分，
$V(\mathfrak p) \cong \Spec(R/\mathfrak p)$ 是其闭点的闭包。
这就是所需矛盾，因为它意味着
$\{\mathfrak p\} \subset V(\mathfrak p)$ 不可能是开集。
\end{proof}








\section{模的支集与维数}
\label{algebra-section-support}

\noindent
下面给出关于模的支集与维数的一些基本结果。

\begin{lemma}
\label{algebra-lemma-filter-Noetherian-module}
设 $R$ 为诺特环，$M$ 为有限 $R$-模。则存在一条由 $R$-子模组成的滤过
$$
0 = M_0 \subset M_1 \subset \ldots \subset M_n = M
$$
使得每个商 $M_i/M_{i-1}$ 都同构于某个 $R/\mathfrak p_i$，
其中 $\mathfrak p_i$ 是 $R$ 的素理想。
\end{lemma}

\begin{proof}[第一种证明]
由引理 \ref{algebra-lemma-trivial-filter-finite-module}，只须处理
$M = R/I$ 的情形，其中 $I$ 是理想。考虑如下理想 $J$ 所成的集合
$S$：本引理对模 $R/J$ 不成立；按包含关系排序。
为导出矛盾，假设 $S$ 非空。因为 $R$ 是诺特环，$S$ 有极大元 $J$。
按照 $S$ 的定义，理想 $J$ 不可能是素理想。选取
$a, b\in R$，使得 $ab \in J$，但 $a \in J$ 与 $b\in J$
都不成立。考虑滤过
$0 \subset aR/(J \cap aR) \subset R/J$。
注意，子模 $aR/(J \cap aR)$ 和商模
$(R/J)/(aR/(J \cap aR))$ 都是循环模；将它们写成
$R/J'$ 与 $R/J''$，便有短正合列
$0 \to R/J' \to R/J \to R/J'' \to 0$。
包含 $J \subset J'$ 是严格的，因为 $b \in J'$；
包含 $J \subset J''$ 也是严格的，因为 $a \in J''$。
因此，由 $J$ 的极大性，$R/J'$ 与 $R/J''$ 都有上述滤过，
从而 $R/J$ 也有。矛盾。
\end{proof}

\begin{proof}[第二种证明]
对 $R$-模 $M$，若存在本引理断言中的滤过，则称 $P(M)$ 成立。
注意，$P$ 在扩张下保持不变，且对 $0$ 成立。由引理
\ref{algebra-lemma-trivial-filter-finite-module}，只须证明对每个理想 $I$，
$P(R/I)$ 成立。否则，由于 $R$ 是诺特环，存在极大反例 $J$。
由例 \ref{algebra-example-oka-family-property-modules} 和命题
\ref{algebra-proposition-oka}，理想 $J$ 是素理想，矛盾。
\end{proof}

\begin{lemma}
\label{algebra-lemma-filter-primes-in-support}
设 $R$、$M$、$M_i$、$\mathfrak p_i$ 如引理
\ref{algebra-lemma-filter-Noetherian-module} 中所述。则
$\text{Supp}(M) = \bigcup V(\mathfrak p_i)$；特别地，
$\mathfrak p_i \in \text{Supp}(M)$。
\end{lemma}

\begin{proof}
这由引理 \ref{algebra-lemma-support-closed} 和
\ref{algebra-lemma-support-quotient} 得出。
\end{proof}

\begin{lemma}
\label{algebra-lemma-support-point}
假设 $R$ 是极大理想为 $\mathfrak m$ 的诺特局部环。
设 $M$ 是非零有限 $R$-模。则
$\text{Supp}(M) = \{ \mathfrak m\}$，当且仅当
$M$ 在 $R$ 上具有有限长度。
\end{lemma}

\begin{proof}
假设 $\text{Supp}(M) = \{ \mathfrak m\}$。
只须证明引理 \ref{algebra-lemma-filter-Noetherian-module} 的滤过中所有素理想
$\mathfrak p_i$ 都是极大理想。这由引理
\ref{algebra-lemma-filter-primes-in-support} 显然成立。

\medskip\noindent
假设 $M$ 在 $R$ 上具有有限长度。由引理
\ref{algebra-lemma-length-infinite}，$\mathfrak m^n M = 0$。
对 $R$ 中的素理想，若 $\mathfrak p \not = \mathfrak m$，
则 $\mathfrak m$ 的某个元素在 $R_{\mathfrak p}$
中映为单位，所以 $M_{\mathfrak p} = 0$。
\end{proof}

\begin{lemma}
\label{algebra-lemma-Noetherian-power-ideal-kills-module}
设 $R$ 为诺特环，$I \subset R$ 为理想，$M$ 为有限 $R$-模。
则对某个 $n \geq 0$ 有 $I^nM = 0$，当且仅当
$\text{Supp}(M) \subset V(I)$。
\end{lemma}

\begin{proof}
事实上，$I^nM = 0$ 等价于 $I^n \subset \text{Ann}(M)$。
由于 $R$ 是诺特环，由引理 \ref{algebra-lemma-Noetherian-power}，
这等价于 $I \subset \sqrt{\text{Ann}(M)}$。
再由引理 \ref{algebra-lemma-Zariski-topology}，这等价于
$V(I) \supset V(\text{Ann}(M))$。
由引理 \ref{algebra-lemma-support-closed}，后者等价于
$V(I) \supset \text{Supp}(M)$。
\end{proof}

\begin{lemma}
\label{algebra-lemma-filter-minimal-primes-in-support}
设 $R$、$M$、$M_i$、$\mathfrak p_i$ 如引理
\ref{algebra-lemma-filter-Noetherian-module} 中所述。集合
$\{\mathfrak p_i\}$ 的极小元正是 $\text{Supp}(M)$ 的极小元。
某个极小素理想 $\mathfrak p$ 出现的次数为
$$
\#\{i \mid \mathfrak p_i = \mathfrak p\}
=
\text{length}_{R_\mathfrak p} M_{\mathfrak p}.
$$
\end{lemma}

\begin{proof}
第一条由
$\text{Supp}(M) = \bigcup V(\mathfrak p_i)$ 得出；参见引理
\ref{algebra-lemma-filter-primes-in-support}。设
$\mathfrak p \in \text{Supp}(M)$ 极小。$M_{\mathfrak p}$ 的支集
是只含极大理想 $\mathfrak p R_{\mathfrak p}$ 的集合。
所以由引理 \ref{algebra-lemma-support-point}，$M_{\mathfrak p}$ 的长度有限且
$> 0$。接着注意，$M_{\mathfrak p}$ 有一条滤过，其子商为
$
(R/\mathfrak p_i)_{\mathfrak p}
=
R_{\mathfrak p}/{\mathfrak p_i}R_{\mathfrak p}
$。
当 $\mathfrak p_i \not \subset \mathfrak p$ 时，它们为零；
当 $\mathfrak p_i \subset \mathfrak p$ 时，它们等于
$\kappa(\mathfrak p)$，因为由 $\mathfrak p$ 的极小性，
此时 $\mathfrak p_i = \mathfrak p$。结论由
$\kappa(\mathfrak p)$ 的长度为 $1$ 得出。
\end{proof}

\begin{lemma}
\label{algebra-lemma-support-dimension-d}
设 $R$ 为诺特局部环，$M$ 为有限 $R$-模。则
$d(M) = \dim(\text{Supp}(M))$，其中 $d(M)$ 如定义
\ref{algebra-definition-d} 中所述。
\end{lemma}

\begin{proof}
设 $M_i, \mathfrak p_i$ 如引理
\ref{algebra-lemma-filter-Noetherian-module} 中所述。由引理
\ref{algebra-lemma-hilbert-ses-chi} 得等式
$d(M) = \max \{ d(R/\mathfrak p_i) \}$。
由命题 \ref{algebra-proposition-dimension}，有
$d(R/\mathfrak p_i) = \dim(R/\mathfrak p_i)$。
显然 $\dim(R/\mathfrak p_i) = \dim V(\mathfrak p_i)$。
由于 $\text{Supp}(M)$ 的所有极小素理想都出现在诸
$\mathfrak p_i$ 中（引理
\ref{algebra-lemma-filter-minimal-primes-in-support}），结论成立。
\end{proof}

\begin{lemma}
\label{algebra-lemma-ses-dimension}
设 $R$ 为诺特环，并设
$0 \to M' \to M \to M'' \to 0$ 是有限 $R$-模的短正合列。则
$\max\{\dim(\text{Supp}(M')), \dim(\text{Supp}(M''))\} =
\dim(\text{Supp}(M))$。
\end{lemma}

\begin{proof}
若 $R$ 是局部环，则结论立即由引理
\ref{algebra-lemma-support-dimension-d} 和 \ref{algebra-lemma-hilbert-ses-chi} 得出。
另一个更初等、且在 $R$ 非局部时也适用的论证，是利用
$\text{Supp}(M')$、$\text{Supp}(M'')$ 与 $\text{Supp}(M)$ 都是闭集
（引理 \ref{algebra-lemma-support-closed}），以及
$\text{Supp}(M) = \text{Supp}(M') \cup \text{Supp}(M'')$
（引理 \ref{algebra-lemma-support-quotient}）。
\end{proof}








\section{相伴素理想}
\label{algebra-section-ass}

\noindent
下面是标准定义。对于非诺特环和非有限模，采用
第 \ref{algebra-section-weakly-ass} 节中的定义可能更为合适。

\begin{definition}
\label{algebra-definition-associated}
设 $R$ 为环，$M$ 为 $R$-模。如果存在元素 $m \in M$，
其零化子为 $\mathfrak p$，则称 $R$ 的素理想 $\mathfrak p$
与 $M$ {\it 相伴}。所有这样的素理想组成的集合记作
$\text{Ass}_R(M)$ 或 $\text{Ass}(M)$。
\end{definition}

\begin{lemma}
\label{algebra-lemma-ass-support}
设 $R$ 为环，$M$ 为 $R$-模。则
$\text{Ass}(M) \subset \text{Supp}(M)$。
\end{lemma}

\begin{proof}
如果 $m \in M$ 的零化子为 $\mathfrak p$，那么特别地，
$R \setminus \mathfrak p$ 中没有元素零化 $m$。
因此 $m$ 是 $M_{\mathfrak p}$ 的非零元素，即
$\mathfrak p \in \text{Supp}(M)$。
\end{proof}

\begin{lemma}
\label{algebra-lemma-ass}
设 $R$ 为环，且 $0 \to M' \to M \to M'' \to 0$ 是
$R$-模的短正合列。则 $\text{Ass}(M') \subset \text{Ass}(M)$，且
$\text{Ass}(M) \subset \text{Ass}(M') \cup \text{Ass}(M'')$。
此外，$\text{Ass}(M' \oplus M'') = \text{Ass}(M') \cup \text{Ass}(M'')$。
\end{lemma}

\begin{proof}
若 $m' \in M'$，则 $m'$ 作为 $M'$ 的元素时的零化子，
与 $m'$ 作为 $M$ 的元素时的零化子相同。因此有包含
$\text{Ass}(M') \subset \text{Ass}(M)$。设 $m \in M$ 的零化子
是素理想 $\mathfrak p$。如果存在 $g \in R$，
$g \not \in \mathfrak p$，使得 $m' = gm \in M'$，
那么 $m'$ 的零化子就是 $\mathfrak p$。如果不存在
$g \in R$，$g \not \in \mathfrak p$，使得 $gm \in M'$，
那么像 $m'' \in M''$（即 $m$ 的像）的零化子就是
$\mathfrak p$。这就证明了包含
$\text{Ass}(M) \subset \text{Ass}(M') \cup \text{Ass}(M'')$。
最后一个断言的证明从略。
\end{proof}

\begin{lemma}
\label{algebra-lemma-ass-filter}
设 $R$ 为环，$M$ 为 $R$-模。假设存在一条由 $R$-子模组成的滤过
$$
0 = M_0 \subset M_1 \subset \ldots \subset M_n = M
$$
使得每个商 $M_i/M_{i-1}$ 都同构于 $R/\mathfrak p_i$，
其中 $\mathfrak p_i$ 是 $R$ 的某个素理想。则
$\text{Ass}(M) \subset \{\mathfrak p_1, \ldots, \mathfrak p_n\}$。
\end{lemma}

\begin{proof}
对滤过 $\{ M_i \}$ 的长度 $n$ 作归纳。选取 $m \in M$，
使其零化子为素理想 $\mathfrak p$。如果 $m \in M_{n-1}$，
则由归纳假设即得结论。否则，$m$ 在
$M/M_{n-1} \cong R/\mathfrak p_n$ 中映为非零元素。因此
$\mathfrak p \subset \mathfrak p_n$。如果二者不相等，则可找到
$f \in \mathfrak p_n$，$f \not\in \mathfrak p$。
此时 $fm$ 的零化子仍为 $\mathfrak p$，且 $fm \in M_{n-1}$。
于是归纳假设给出结论。
\end{proof}

\begin{lemma}
\label{algebra-lemma-finite-ass}
设 $R$ 为诺特环，$M$ 为有限 $R$-模。则
$\text{Ass}(M)$ 是有限集。
\end{lemma}

\begin{proof}
这立即由引理 \ref{algebra-lemma-ass-filter} 和
\ref{algebra-lemma-filter-Noetherian-module} 得出。
\end{proof}

\begin{proposition}
\label{algebra-proposition-minimal-primes-associated-primes}
设 $R$ 为诺特环，$M$ 为有限 $R$-模。下列素理想集合相同：
\begin{enumerate}
\item $M$ 的支集中的极小素理想；
\item $\text{Ass}(M)$ 中的极小素理想；
\item 对任意满足 $M_i/M_{i-1} \cong R/\mathfrak p_i$ 的滤过
$0 = M_0 \subset M_1 \subset \ldots
\subset M_{n-1} \subset M_n = M$，集合
$\{\mathfrak p_i\}$ 中的极小素理想。
\end{enumerate}
\end{proposition}

\begin{proof}
选取 (3) 中的滤过。在引理
\ref{algebra-lemma-filter-minimal-primes-in-support} 中，我们已经看到
(1) 与 (3) 中的集合相同。

\medskip\noindent
设 $\mathfrak p$ 是集合 $\{\mathfrak p_i\}$ 的极小元，
并取最小的 $i$ 使得 $\mathfrak p = \mathfrak p_i$。
选取 $m \in M_i$，$m \not \in M_{i-1}$。元素 $m$ 的零化子
包含于 $\mathfrak p_i = \mathfrak p$，并包含
$\mathfrak p_1 \mathfrak p_2 \ldots \mathfrak p_i$。
由 $i$ 与 $\mathfrak p$ 的选取，对 $j < i$ 有
$\mathfrak p_j \not \subset \mathfrak p$，因而
$\mathfrak p_1 \mathfrak p_2 \ldots \mathfrak p_{i - 1}
\not \subset \mathfrak p_i$。选取
$f \in \mathfrak p_1 \mathfrak p_2 \ldots \mathfrak p_{i - 1}$，
$f \not \in \mathfrak p$。则 $fm$ 的零化子为 $\mathfrak p$。
由此可见，$\mathfrak p$ 是 $M$ 的相伴素理想。
由引理 \ref{algebra-lemma-ass-support} 有
$\text{Ass}(M) \subset \text{Supp}(M)$，所以 $\mathfrak p$
在 $\text{Ass}(M)$ 中极小。因此，(1) 中的素理想集合
包含于 (2) 中的素理想集合。

\medskip\noindent
设 $\mathfrak p$ 是 $\text{Ass}(M)$ 的极小元。
由于 $\text{Ass}(M) \subset \text{Supp}(M)$，存在
$\text{Supp}(M)$ 的极小元 $\mathfrak q$，使得
$\mathfrak q \subset \mathfrak p$。我们刚刚已经证明
$\mathfrak q \in \text{Ass}(M)$。由 $\mathfrak p$ 的极小性，
$\mathfrak q = \mathfrak p$。因此，(2) 中的素理想集合
包含于 (1) 中的素理想集合。
\end{proof}

\begin{lemma}
\label{algebra-lemma-ass-zero}
\begin{slogan}
诺特环上的每个非零模都有相伴素理想。
\end{slogan}
设 $R$ 为诺特环，$M$ 为 $R$-模。则
$$
M = (0) \Leftrightarrow \text{Ass}(M) = \emptyset.
$$
\end{lemma}

\begin{proof}
如果 $M = (0)$，那么依定义
$\text{Ass}(M) = \emptyset$。如果 $M \not = 0$，选取任意非零有限生成
子模 $M' \subset M$，例如由单个非零元素生成的子模。由引理
\ref{algebra-lemma-support-zero} 可知 $\text{Supp}(M')$ 非空。
由命题 \ref{algebra-proposition-minimal-primes-associated-primes}，
这意味着 $\text{Ass}(M')$ 非空。再由引理
\ref{algebra-lemma-ass}，可知 $\text{Ass}(M) \not = \emptyset$。
\end{proof}

\begin{lemma}
\label{algebra-lemma-ass-minimal-prime-support}
设 $R$ 为诺特环，$M$ 为 $R$-模。在 $\text{Supp}(M)$
的元素中极小的任意 $\mathfrak p \in \text{Supp}(M)$ 都属于
$\text{Ass}(M)$。
\end{lemma}

\begin{proof}
如果 $M$ 是有限 $R$-模，这就是命题
\ref{algebra-proposition-minimal-primes-associated-primes} 的推论。
一般地，将 $M = \bigcup M_\lambda$ 写成其有限子模之并，并使用
$\text{Supp}(M) = \bigcup \text{Supp}(M_\lambda)$
以及
$\text{Ass}(M) = \bigcup \text{Ass}(M_\lambda)$。
\end{proof}

\begin{lemma}
\label{algebra-lemma-ass-zero-divisors}
设 $R$ 为诺特环，$M$ 为 $R$-模。并集
$\bigcup_{\mathfrak q \in \text{Ass}(M)} \mathfrak q$
正是 $R$ 中在 $M$ 上为零因子的元素所成的集合。
\end{lemma}

\begin{proof}
任意相伴素理想中的每个元素显然都是 $M$ 上的零因子。
反之，假设 $x \in R$ 是 $M$ 上的零因子。考虑子模
$N = \{m \in M \mid xm = 0\}$。由于 $N$ 非零，由引理
\ref{algebra-lemma-ass-zero}，它有相伴素理想 $\mathfrak q$。
于是 $x \in \mathfrak q$；又由引理
\ref{algebra-lemma-ass}，$\mathfrak q$ 是 $M$ 的相伴素理想。
\end{proof}

\begin{lemma}
\label{algebra-lemma-one-equation-module}
设 $R$ 为诺特局部环，$M$ 为有限 $R$-模，并设
$f \in \mathfrak m$ 是 $R$ 的极大理想中的元素。则
$$
\dim(\text{Supp}(M/fM)) \leq
\dim(\text{Supp}(M)) \leq
\dim(\text{Supp}(M/fM)) + 1
$$
如果 $f$ 不属于 $M$ 的支集中的任何极小素理想
（例如，如果 $f$ 是 $M$ 上的非零因子），
那么右侧不等式取等号。
\end{lemma}

\begin{proof}
（括号中的断言由引理 \ref{algebra-lemma-ass-zero-divisors} 得出。）
第一个不等式由
$\text{Supp}(M/fM) \subset \text{Supp}(M)$ 得出，
参见引理 \ref{algebra-lemma-support-quotient}。对于第二个不等式，注意
$\text{Supp}(M/fM) = \text{Supp}(M) \cap V(f)$，参见引理
\ref{algebra-lemma-support-quotient}。例如，由引理
\ref{algebra-lemma-filter-primes-in-support} 和维数的初等性质可知，只须证明
$\dim V(\mathfrak p) \leq \dim (V(\mathfrak p) \cap V(f)) + 1$
对 $R$ 的素理想 $\mathfrak p$ 成立。这是引理
\ref{algebra-lemma-one-equation} 的推论。最后，如果 $f$ 不包含于
$M$ 的支集中的任何极小素理想，那么
$\text{Supp}(M/fM)$ 中的每条素理想链都会产生
$\text{Supp}(M)$ 中的一条素理想链，而后者距极大链至少还差一步。
\end{proof}

\begin{lemma}
\label{algebra-lemma-ass-functorial}
设 $\varphi : R \to S$ 为环映射，$M$ 为 $S$-模。则
$\Spec(\varphi)(\text{Ass}_S(M)) \subset \text{Ass}_R(M)$。
\end{lemma}

\begin{proof}
如果 $\mathfrak q \in \text{Ass}_S(M)$，那么存在 $M$ 中的
$m$，使得 $m$ 在 $S$ 中的零化子为 $\mathfrak q$。
于是 $m$ 在 $R$ 中的零化子为 $\mathfrak q \cap R$。
\end{proof}

\begin{remark}
\label{algebra-remark-ass-reverse-functorial}
设 $\varphi : R \to S$ 为环映射，$M$ 为 $S$-模。
不一定有
$\Spec(\varphi)(\text{Ass}_S(M)) \supset \text{Ass}_R(M)$。
例如，考虑环映射
$R = k \to S = k[x_1, x_2, x_3, \ldots]/(x_i^2)$ 以及 $M = S$。
此时 $\text{Ass}_R(M)$ 非空，但 $\text{Ass}_S(S)$ 为空。
\end{remark}

\begin{lemma}
\label{algebra-lemma-ass-functorial-Noetherian}
设 $\varphi : R \to S$ 为环映射，$M$ 为 $S$-模。
如果 $S$ 是诺特环，则
$\Spec(\varphi)(\text{Ass}_S(M)) = \text{Ass}_R(M)$。
\end{lemma}

\begin{proof}
我们已经在引理
\ref{algebra-lemma-ass-functorial} 中看到
$\Spec(\varphi)(\text{Ass}_S(M)) \subset \text{Ass}_R(M)$。
反过来，选取素理想 $\mathfrak p \in \text{Ass}_R(M)$。
设 $m \in M$ 满足：$m$ 在 $R$ 中的零化子为
$\mathfrak p$。令
$I = \{g \in S \mid gm = 0\}$ 为 $m$ 在 $S$ 中的零化子。
于是 $R/\mathfrak p \subset S/I$ 是单射。
结合引理 \ref{algebra-lemma-injective-minimal-primes-in-image} 和
\ref{algebra-lemma-minimal-prime-image-minimal-prime}，可知存在包含
$I$ 的极小素理想 $\mathfrak q \subset S$，它映到
$\mathfrak p$。由命题
\ref{algebra-proposition-minimal-primes-associated-primes}，
$\mathfrak q$ 是 $S/I$ 的相伴素理想；再由引理
\ref{algebra-lemma-ass}，$\mathfrak q$ 是 $M$ 的相伴素理想，
结论得证。
\end{proof}

\begin{lemma}
\label{algebra-lemma-ass-quotient-ring}
设 $R$ 为环，$I$ 为理想，$M$ 为 $R/I$-模。
通过典范单射
$\Spec(R/I) \to \Spec(R)$，有
$\text{Ass}_{R/I}(M) = \text{Ass}_R(M)$。
\end{lemma}

\begin{proof}
证明从略。
\end{proof}

\begin{lemma}
\label{algebra-lemma-associated-primes-localize}
设 $R$ 为环，$M$ 为 $R$-模，并设
$\mathfrak p \subset R$ 为素理想。
\begin{enumerate}
\item 如果 $\mathfrak p \in \text{Ass}(M)$，则
$\mathfrak pR_{\mathfrak p} \in \text{Ass}(M_{\mathfrak p})$。
\item 如果 $\mathfrak p$ 是有限生成的，则逆命题也成立。
\end{enumerate}
\end{lemma}

\begin{proof}
如果 $\mathfrak p \in \text{Ass}(M)$，则存在元素 $m \in M$，
其零化子为 $\mathfrak p$。由于局部化正合
（命题 \ref{algebra-proposition-localization-exact}），
$m/1$ 在 $M_{\mathfrak p}$ 中的零化子为
$\mathfrak pR_{\mathfrak p}$，故 (1) 成立。
假设 $\mathfrak pR_{\mathfrak p} \in \text{Ass}(M_{\mathfrak p})$，
并且 $\mathfrak p = (f_1, \ldots, f_n)$。设 $m/g$ 是
$M_{\mathfrak p}$ 中零化子为
$\mathfrak pR_{\mathfrak p}$ 的元素。这意味着 $m$ 的零化子
包含于 $\mathfrak p$。由于 $f_i m/g = 0$ 在
$M_{\mathfrak p}$ 中成立，可知存在
$g_i \in R$，$g_i \not \in \mathfrak p$，使得
$g_i f_i m = 0$ 在 $M$ 中成立。合起来可知
$g_1\ldots g_nm$ 的零化子为 $\mathfrak p$。因此
$\mathfrak p \in \text{Ass}(M)$。
\end{proof}

\begin{lemma}
\label{algebra-lemma-localize-ass}
设 $R$ 为环，$M$ 为 $R$-模，并设 $S \subset R$ 为乘法子集。
通过典范单射 $\Spec(S^{-1}R) \to \Spec(R)$，有
\begin{enumerate}
\item $\text{Ass}_R(S^{-1}M) = \text{Ass}_{S^{-1}R}(S^{-1}M)$；
\item
$\text{Ass}_R(M) \cap \Spec(S^{-1}R) \subset \text{Ass}_R(S^{-1}M)$；且
\item 如果 $R$ 是诺特环，则此包含为等式。
\end{enumerate}
\end{lemma}

\begin{proof}
对 $m \in S^{-1}M$，设 $I \subset R$ 和 $J \subset S^{-1}R$
分别为 $m$ 的零化子。则 $I$ 是 $J$ 在映射
$R \to S^{-1}R$ 下的原像，且 $J = S^{-1}I$。
(1) 中的等式由引理 \ref{algebra-lemma-spec-localization} 对映射
$\Spec(S^{-1}R) \to \Spec(R)$ 的描述得出。
对 $m \in M$，设 $I \subset R$ 为 $m$ 在 $R$ 中的零化子，
并设 $J \subset S^{-1}R$ 为
$m/1 \in S^{-1}M$ 的零化子。我们有 $J = S^{-1}I$，这蕴含 (2)。
诺特情形中的等式由引理
\ref{algebra-lemma-associated-primes-localize} 得出，因为对
$\mathfrak p \in R$，若 $S \cap \mathfrak p = \emptyset$，则
$M_{\mathfrak p} = (S^{-1}M)_{S^{-1}\mathfrak p}$。
\end{proof}

\begin{lemma}
\label{algebra-lemma-localize-ass-nonzero-divisors}
设 $R$ 为环，$M$ 为 $R$-模，并设 $S \subset R$ 为乘法子集。
假设每个 $s \in S$ 都是 $M$ 上的非零因子。则
$$
\text{Ass}_R(M) = \text{Ass}_R(S^{-1}M).
$$
\end{lemma}

\begin{proof}
由假设有 $M \subset S^{-1}M$，故引理
\ref{algebra-lemma-ass} 给出包含
$\text{Ass}(M) \subset \text{Ass}(S^{-1}M)$。
反之，假设 $n/s \in S^{-1}M$ 的零化子为素理想
$\mathfrak p$。则 $n \in M$ 的零化子也是 $\mathfrak p$。
\end{proof}

\begin{lemma}
\label{algebra-lemma-ideal-nonzerodivisor}
设 $R$ 为极大理想是 $\mathfrak m$ 的诺特局部环，
$I \subset \mathfrak m$ 为理想，$M$ 为有限 $R$-模。
下列条件等价：
\begin{enumerate}
\item 存在 $x \in I$，它不是 $M$ 上的零因子；
\item 对所有 $\mathfrak q \in \text{Ass}(M)$，均有
$I \not \subset \mathfrak q$。
\end{enumerate}
\end{lemma}

\begin{proof}
如果 $I$ 中存在非零因子 $x$，那么 $x$ 显然不可能属于
$M$ 的任何相伴素理想。反之，假设对所有
$\mathfrak q \in \text{Ass}(M)$ 均有
$I \not \subset \mathfrak q$。在这种情形下，由引理
\ref{algebra-lemma-finite-ass} 和 \ref{algebra-lemma-silly}，可以选取
$x \in I$，使得对所有 $\mathfrak q \in \text{Ass}(M)$，
$x \not \in \mathfrak q$。由引理
\ref{algebra-lemma-ass-zero-divisors}，元素 $x$ 不是 $M$ 上的零因子。
\end{proof}

\begin{lemma}
\label{algebra-lemma-zero-at-ass-zero}
设 $R$ 为环，$M$ 为 $R$-模。如果 $R$ 是诺特环，则映射
$$
M
\longrightarrow
\prod\nolimits_{\mathfrak p \in \text{Ass}(M)} M_{\mathfrak p}
$$
是单射。
\end{lemma}

\begin{proof}
设 $x \in M$ 是该映射的核中的元素。如果 $\mathfrak p$ 是
$Rx \subset M$ 的相伴素理想，那么一方面
$\mathfrak p \in \text{Ass}(M)$（引理 \ref{algebra-lemma-ass}），
另一方面
$(Rx)_{\mathfrak p} \subset M_{\mathfrak p}$ 非零。
此矛盾说明 $\text{Ass}(Rx) = \emptyset$。因此由引理
\ref{algebra-lemma-ass-zero}，$Rx = 0$。
\end{proof}

\noindent
这个引理或许应该放在别处。

\begin{lemma}
\label{algebra-lemma-dim-not-zero-exists-nonzerodivisor-nonunit}
设 $k$ 为域，$S$ 为有限型 $k$-代数。如果 $\dim(S) > 0$，
则存在既是非零因子又是非单位的元素 $f \in S$。
\end{lemma}

\begin{proof}
由引理 \ref{algebra-lemma-finite-ass}，环 $S$ 只有有限多个相伴素理想。
由引理 \ref{algebra-lemma-finite-type-algebra-finite-nr-primes}，
环 $S$ 有无穷多个极大理想。因此，可以选取极大理想
$\mathfrak m \subset S$，它不是 $S$ 的相伴素理想。
由素理想回避（引理 \ref{algebra-lemma-silly}），可以选取非零元素
$f \in \mathfrak m$，它不属于 $S$ 的任何相伴素理想。
由引理 \ref{algebra-lemma-ass-zero-divisors}，元素 $f$ 是非零因子；
又因 $f \in \mathfrak m$，可知 $f$ 不是单位。
\end{proof}



\section{符号幂}
\label{algebra-section-symbolic-power}

\noindent
定义如下。

\begin{definition}
\label{algebra-definition-symbolic-power}
设 $R$ 为环，$\mathfrak p$ 为素理想。对 $n \geq 0$，
$\mathfrak p$ 的第 $n$ 个{\it 符号幂}是理想
$\mathfrak p^{(n)} = \Ker(R \to R_\mathfrak p/\mathfrak p^nR_\mathfrak p)$。
\end{definition}

\noindent
注意 $\mathfrak p^n \subset \mathfrak p^{(n)}$，但等式并非总成立。

\begin{lemma}
\label{algebra-lemma-symbolic-power-associated}
设 $R$ 为诺特环，$\mathfrak p$ 为素理想，且 $n > 0$。则
$\text{Ass}(R/\mathfrak p^{(n)}) = \{\mathfrak p\}$。
\end{lemma}

\begin{proof}
如果 $\mathfrak q$ 是 $R/\mathfrak p^{(n)}$ 的相伴素理想，
那么显然 $\mathfrak p \subset \mathfrak q$。另一方面，任意
$x \in R$，$x \not \in \mathfrak p$，都是
$R/\mathfrak p^{(n)}$ 上的非零因子。事实上，如果 $y \in R$ 且
$xy \in \mathfrak p^{(n)} = R \cap \mathfrak p^nR_{\mathfrak p}$，
那么 $y \in \mathfrak p^nR_{\mathfrak p}$，因而
$y \in \mathfrak p^{(n)}$。引理得证。
\end{proof}

\begin{lemma}
\label{algebra-lemma-symbolic-power-flat-extension}
设 $R \to S$ 为平坦环映射，$\mathfrak p \subset R$ 为素理想，
且 $\mathfrak q = \mathfrak p S$ 是 $S$ 的素理想。则
$\mathfrak p^{(n)} S = \mathfrak q^{(n)}$。
\end{lemma}

\begin{proof}
由于
$\mathfrak p^{(n)} = \Ker(R \to R_\mathfrak p/\mathfrak p^nR_\mathfrak p)$，
利用平坦性可知 $\mathfrak p^{(n)} S$ 是映射
$S \to S_\mathfrak p/\mathfrak p^nS_\mathfrak p$ 的核。
另一方面，$\mathfrak q^{(n)}$ 是映射
$S \to S_\mathfrak q/\mathfrak q^nS_\mathfrak q =
S_\mathfrak q/\mathfrak p^nS_\mathfrak q$ 的核。因此只须证明
$$
S_\mathfrak p/\mathfrak p^nS_\mathfrak p
\longrightarrow
S_\mathfrak q/\mathfrak p^nS_\mathfrak q
$$
是单射。注意，右侧的模是将左侧的模关于诸元素
$f \in S$，$f \not \in \mathfrak q$ 局部化所得。
因此只须证明这些元素是
$S_\mathfrak p/\mathfrak p^nS_\mathfrak p$ 上的非零因子。
由平坦性，模 $S_\mathfrak p/\mathfrak p^nS_\mathfrak p$
具有有限滤过，其子商为
$$
\mathfrak p^iS_\mathfrak p/\mathfrak p^{i + 1}S_\mathfrak p
\cong \mathfrak p^iR_\mathfrak p/\mathfrak p^{i + 1}R_\mathfrak p
\otimes_{R_\mathfrak p} S_\mathfrak p \cong
V \otimes_{\kappa(\mathfrak p)} (S/\mathfrak q)_\mathfrak p
$$
其中 $V$ 是 $\kappa(\mathfrak p)$-向量空间。因此，$f$
按所需方式可逆地作用。
\end{proof}

\section{相对相伴素理想集}
\label{algebra-section-relative-assassin}

\noindent
下面讨论相对相伴素理想集。设 $R \to S$ 为环映射，
$N$ 为 $S$-模。在这种情形下，可以引入如下 $S$ 的素理想
$\mathfrak q$ 的集合：
\begin{enumerate}
\item $A$：令 $\mathfrak p = R \cap \mathfrak q$，有
$\mathfrak q \in \text{Ass}_S(N \otimes_R \kappa(\mathfrak p))$；
\item $A'$：令 $\mathfrak p = R \cap \mathfrak q$，则在典范映射
$\Spec(S \otimes_R \kappa(\mathfrak p)) \to \Spec(S)$ 下，
$\mathfrak q$ 属于
$\text{Ass}_{S \otimes \kappa(\mathfrak p)}(N \otimes_R \kappa(\mathfrak p))$
的像；
\item $A_{fin}$：令 $\mathfrak p = R \cap \mathfrak q$，有
$\mathfrak q \in \text{Ass}_S(N/\mathfrak pN)$；
\item $A'_{fin}$：对某个素理想 $\mathfrak p' \subset R$，有
$\mathfrak q \in \text{Ass}_S(N/\mathfrak p'N)$；
\item $B$：对某个 $R$-模 $M$，有
$\mathfrak q \in \text{Ass}_S(N \otimes_R M)$；以及
\item $B_{fin}$：对某个有限 $R$-模 $M$，有
$\mathfrak q \in \text{Ass}_S(N \otimes_R M)$。
\end{enumerate}
下面确定这些集合之间的一些关系。

\begin{lemma}
\label{algebra-lemma-compare-relative-assassins}
设 $R \to S$ 为环映射，$N$ 为 $S$-模，并设
$A$、$A'$、$A_{fin}$、$B$ 与 $B_{fin}$ 为上面引入的
$\Spec(S)$ 的子集。
\begin{enumerate}
\item 总有 $A = A'$。
\item 总有 $A_{fin} \subset A$、
$B_{fin} \subset B$、$A_{fin} \subset A'_{fin} \subset B_{fin}$
以及 $A \subset B$。
\item 如果 $S$ 是诺特环，则 $A = A_{fin}$ 且 $B = B_{fin}$。
\item 如果 $N$ 在 $R$ 上平坦，则
$A = A_{fin} = A'_{fin}$ 且 $B = B_{fin}$。
\item 如果 $R$ 是诺特环且 $N$ 在 $R$ 上平坦，则所有这些集合相等，
即 $A = A' = A_{fin} = A'_{fin} = B = B_{fin}$。
\end{enumerate}
\end{lemma}

\begin{proof}
本证明中的一些论证会在后面更精确的引理的证明中重复出现
（因为那些引理处理给定的模 $M$ 或给定的素理想
$\mathfrak p$，而不是它们的全体）。

\medskip\noindent
(1) 的证明。设 $\mathfrak p$ 为 $R$ 的素理想。则有
$$
\text{Ass}_S(N \otimes_R \kappa(\mathfrak p)) =
\text{Ass}_{S/\mathfrak pS}(N \otimes_R \kappa(\mathfrak p)) =
\text{Ass}_{S \otimes_R \kappa(\mathfrak p)}(N \otimes_R \kappa(\mathfrak p))
$$
其中第一个等式由引理
\ref{algebra-lemma-ass-quotient-ring} 得出，第二个等式由引理
\ref{algebra-lemma-localize-ass} 的第 (1) 部分得出。这证明了
$A = A'$。包含 $A_{fin} \subset A'_{fin}$ 是清楚的。

\medskip\noindent
(2) 的证明。除 $A_{fin} \subset A$ 外，各个包含关系都立即由定义得出；
而该包含关系由引理 \ref{algebra-lemma-localize-ass} 以及在
$A_{fin}$ 的表述中要求 $\mathfrak p = R \cap \mathfrak q$
这一事实得出。

\medskip\noindent
(3) 的证明。如果 $S$ 是诺特环，则等式 $A = A_{fin}$ 由引理
\ref{algebra-lemma-localize-ass} 的第 (3) 部分得出。设
$\mathfrak q = (g_1, \ldots, g_m)$ 是 $S$ 的有限生成素理想。
设 $z \in N \otimes_R M$ 的零化子为 $\mathfrak q$。
可以选取有限子模 $M' \subset M$，使 $z$ 是某个
$z' \in N \otimes_R M'$ 的像。于是
$\text{Ann}_S(z') \subset \mathfrak q = \text{Ann}_S(z)$。
由于 $N \otimes_R -$ 与余极限交换，且 $M$ 是有限 $R$-模的
有向余极限，可以找到 $M' \subset M'' \subset M$，使得像
$z'' \in N \otimes_R M''$ 被
$g_1, \ldots, g_m$ 零化。因此
$\text{Ann}_S(z'') = \mathfrak q$。这证明了：如果 $S$ 是诺特环，
则 $B = B_{fin}$。

\medskip\noindent
(4) 的证明。如果 $N$ 平坦，则函子 $N \otimes_R -$ 正合。
特别地，如果 $M' \subset M$，则
$N \otimes_R M' \subset N \otimes_R M$。因此，如果
$z \in N \otimes_R M$ 的零化子
$\mathfrak q = \text{Ann}_S(z)$ 是素理想，那么可以选取任意有限
$R$-子模 $M' \subset M$，使得
$z \in N \otimes_R M'$；而将 $z$ 视为
$N \otimes_R M'$ 的元素时，其零化子仍等于
$\mathfrak q$。因此 $B = B_{fin}$。设 $\mathfrak p'$ 为
$R$ 的素理想，并设 $\mathfrak q$ 为 $S$ 的素理想，且它是
$N/\mathfrak p'N$ 的相伴素理想。
这意味着 $\mathfrak p'S \subset \mathfrak q$。由于 $N$
在 $R$ 上平坦，可知 $N/\mathfrak p'N$ 在整环
$R/\mathfrak p'$ 上平坦。因此，$R/\mathfrak p'$ 的每个非零元素
都是 $N/\mathfrak p'$ 上的非零因子。于是这些元素都不可能映到
$\mathfrak q$ 的元素，故 $\mathfrak p' = R \cap \mathfrak q$。
因此 $A_{fin} = A'_{fin}$。最后，由引理
\ref{algebra-lemma-localize-ass-nonzero-divisors} 可知
$\text{Ass}_S(N/\mathfrak p'N) =
\text{Ass}_S(N \otimes_R \kappa(\mathfrak p'))$，即
$A'_{fin} = A$。

\medskip\noindent
(5) 的证明。由于其他等式已在 (4) 中证明，只须证明
$A'_{fin} = B_{fin}$。为此，设 $M$ 为有限 $R$-模。
由引理 \ref{algebra-lemma-filter-Noetherian-module}，存在一条由
$R$-子模组成的滤过
$$
0 = M_0 \subset M_1 \subset \ldots \subset M_n = M
$$
使得每个商 $M_i/M_{i-1}$ 都同构于某个
$R/\mathfrak p_i$，其中 $\mathfrak p_i$ 是 $R$ 的素理想。
由于 $N$ 平坦，得到一条由 $S$-子模组成的滤过
$$
0 = N \otimes_R M_0 \subset N \otimes_R M_1 \subset \ldots \subset
N \otimes_R M_n = N \otimes_R M
$$
其每个子商都同构于 $N/\mathfrak p_iN$。由引理
\ref{algebra-lemma-ass}，可得
$\text{Ass}_S(N \otimes_R M) \subset \bigcup \text{Ass}_S(N/\mathfrak p_iN)$。
因此 $B_{fin} \subset A'_{fin}$。另一个包含关系是 (2) 的一部分，
故结论得证。
\end{proof}

\noindent
我们将 $N$ 在 $S/R$ 上的相对相伴素理想集定义为上面的集合
$A = A'$。作为动机，指出它只依赖于纤维环上的纤维模
$N \otimes_R \kappa(\mathfrak p)$。与模的相伴素理想集一样，
我们提醒读者：当纤维环
$S \otimes_R \kappa(\mathfrak p)$ 为诺特环时，这个概念最为适用；
例如，当 $R \to S$ 是有限型映射时。

\begin{definition}
\label{algebra-definition-relative-assassin}
设 $R \to S$ 为环映射，$N$ 为 $S$-模。
{\it $N$ 在 $S/R$ 上的相对相伴素理想集}是集合
$$
\text{Ass}_{S/R}(N)
=
\{ \mathfrak q \subset S \mid
\mathfrak q \in \text{Ass}_S(N \otimes_R \kappa(\mathfrak p))
\text{ 其中 }\mathfrak p = R \cap \mathfrak q\}.
$$
这就是引理 \ref{algebra-lemma-compare-relative-assassins} 中记作
$A$ 的集合。
\end{definition}

\noindent
下面几个结果的实质都与相对相伴素理想集有关，
尽管这一点未必显而易见。

\begin{lemma}
\label{algebra-lemma-bourbaki}
设 $R \to S$ 为环映射，$M$ 为 $R$-模，$N$ 为 $S$-模。
如果 $N$ 作为 $R$-模是平坦的，则
$$
\text{Ass}_S(M \otimes_R N)
\supset
\bigcup\nolimits_{\mathfrak p \in \text{Ass}_R(M)} \text{Ass}_S(N/\mathfrak pN)
$$
并且如果 $R$ 是诺特环，则等号成立。
\end{lemma}

\begin{proof}
如果 $\mathfrak p \in \text{Ass}_R(M)$，则存在单射
$R/\mathfrak p \to M$。由于 $N$ 在 $R$ 上平坦，得到单射
$R/\mathfrak p \otimes_R N \to M \otimes_R N$。又因为
$R/\mathfrak p \otimes_R N = N/\mathfrak pN$，由引理
\ref{algebra-lemma-ass} 可得
$\text{Ass}_S(N/\mathfrak pN) \subset \text{Ass}_S(M \otimes_R N)$。
因此右侧包含于左侧。

\medskip\noindent
将 $M = \bigcup M_\lambda$ 写成其有限生成
$R$-子模之并。于是也有
$N \otimes_R M = \bigcup N \otimes_R M_\lambda$
（因为 $N$ 在 $R$ 上平坦）。由相伴素理想的定义可知
$\text{Ass}_S(N \otimes_R M) = \bigcup \text{Ass}_S(N \otimes_R M_\lambda)$
且 $\text{Ass}_R(M) = \bigcup \text{Ass}(M_\lambda)$。
因此可以假设 $M$ 是有限生成的。

\medskip\noindent
设 $\mathfrak q \in \text{Ass}_S(M \otimes_R N)$，并假设
$R$ 是诺特环且 $M$ 是有限 $R$-模。为完成证明，须说明
$\mathfrak q$ 属于右侧。首先注意
$\mathfrak qS_{\mathfrak q} \in
\text{Ass}_{S_{\mathfrak q}}((M \otimes_R N)_{\mathfrak q})$，
参见引理 \ref{algebra-lemma-associated-primes-localize}。
设 $\mathfrak p$ 为 $R$ 中对应的素理想。注意
$$
(M \otimes_R N)_{\mathfrak q} = M \otimes_R N_{\mathfrak q}
= M_{\mathfrak p} \otimes_{R_{\mathfrak p}} N_{\mathfrak q}
$$
如果
$\mathfrak pR_{\mathfrak p} \not \in
\text{Ass}_{R_{\mathfrak p}}(M_{\mathfrak p})$，
那么存在 $x \in \mathfrak pR_{\mathfrak p}$，它是
$M_{\mathfrak p}$ 中的非零因子（参见引理
\ref{algebra-lemma-ideal-nonzerodivisor}）。由于
$N_{\mathfrak q}$ 在 $R_{\mathfrak p}$ 上平坦，可知 $x$ 在
$\mathfrak qS_{\mathfrak q}$ 中的像是
$(M \otimes_R N)_{\mathfrak q}$ 上的非零因子。
这与假设
$\mathfrak qS_{\mathfrak q} \in \text{Ass}_S((M \otimes_R N)_{\mathfrak q})$
矛盾。因此，$\mathfrak p$ 是 $M$ 的相伴素理想之一。

\medskip\noindent
继续这一论证，选取滤过
$$
0 = M_0 \subset M_1 \subset \ldots \subset M_n = M
$$
使得每个商 $M_i/M_{i-1}$ 都同构于某个
$R/\mathfrak p_i$，其中 $\mathfrak p_i$ 是 $R$ 的素理想；
参见引理 \ref{algebra-lemma-filter-Noetherian-module}。
（由引理 \ref{algebra-lemma-ass-filter}，至少对一个 $i$ 有
$\mathfrak p_i = \mathfrak p$。）这给出滤过
$$
0 = M_0 \otimes_R N \subset M_1 \otimes_R N \subset \ldots
\subset M_n \otimes_R N = M \otimes_R N
$$
其子商同构于 $N/\mathfrak p_iN$。如果
$\mathfrak p_i \not = \mathfrak p$，那么由前一段的结果，
$\mathfrak q$ 不可能与模 $N/\mathfrak p_iN$ 相伴
（因为 $\text{Ass}_R(R/\mathfrak p_i) = \{\mathfrak p_i\}$）。
因此，$\mathfrak q$ 按所需与 $N/\mathfrak pN$ 相伴。
\end{proof}

\begin{lemma}
\label{algebra-lemma-post-bourbaki}
设 $R \to S$ 为环映射，$N$ 为 $S$-模。假设 $N$
作为 $R$-模是平坦的，并且 $R$ 是分式域为 $K$ 的整环。则
$$
\text{Ass}_S(N) =
\text{Ass}_S(N \otimes_R K) =
\text{Ass}_{S \otimes_R K}(N \otimes_R K)
$$
其中使用典范包含
$\Spec(S \otimes_R K) \subset \Spec(S)$。
\end{lemma}

\begin{proof}
注意 $S \otimes_R K = (R \setminus \{0\})^{-1}S$ 且
$N \otimes_R K = (R \setminus \{0\})^{-1}N$。
对任意非零 $x \in R$，由于 $N$ 在 $R$ 上平坦，
乘以 $x$ 在 $N$ 上是单射。因此，本引理由引理
\ref{algebra-lemma-localize-ass-nonzero-divisors} 与引理
\ref{algebra-lemma-localize-ass} 的第 (1) 部分结合得出。
\end{proof}

\begin{lemma}
\label{algebra-lemma-bourbaki-fibres}
设 $R \to S$ 为环映射，$M$ 为 $R$-模，$N$ 为 $S$-模。
假设 $N$ 作为 $R$-模是平坦的。则
$$
\text{Ass}_S(M \otimes_R N)
\supset
\bigcup\nolimits_{\mathfrak p \in \text{Ass}_R(M)}
\text{Ass}_{S \otimes_R \kappa(\mathfrak p)}(N \otimes_R \kappa(\mathfrak p))
$$
这里使用注 \ref{algebra-remark-fundamental-diagram}，
将纤维环的谱视为 $\Spec(S)$ 的子集。
如果 $R$ 是诺特环，则此包含为等式。
\end{lemma}

\begin{proof}
这等价于引理
\ref{algebra-lemma-bourbaki}；这是由引理
\ref{algebra-lemma-ass-quotient-ring}、
\ref{algebra-lemma-flat-base-change} 及
\ref{algebra-lemma-post-bourbaki} 得出的。
\end{proof}

\begin{remark}
\label{algebra-remark-bourbaki}
设 $R \to S$ 为环映射，$N$ 为 $S$-模，并设
$\mathfrak p$ 为 $R$ 的素理想。则
$$
\text{Ass}_S(N \otimes_R \kappa(\mathfrak p)) =
\text{Ass}_{S/\mathfrak pS}(N \otimes_R \kappa(\mathfrak p)) =
\text{Ass}_{S \otimes_R \kappa(\mathfrak p)}(N \otimes_R \kappa(\mathfrak p)).
$$
第一个等式由引理
\ref{algebra-lemma-ass-quotient-ring} 得出，第二个等式由引理
\ref{algebra-lemma-localize-ass} 的第 (1) 部分得出。
\end{remark}

\section{弱相伴素理想}
\label{algebra-section-weakly-ass}

\noindent
这是相伴素理想概念的一个变体，对非诺特环和非有限模很有用。

\begin{definition}
\label{algebra-definition-weakly-associated}
设 $R$ 为环，$M$ 为 $R$-模。如果存在元素 $m \in M$，
使得 $\mathfrak p$ 在包含零化子
$\text{Ann}(m) = \{f \in R \mid fm = 0\}$ 的素理想中极小，
则称 $R$ 的素理想 $\mathfrak p$ 与 $M$ {\it 弱相伴}。
所有这样的素理想组成的集合记作
$\text{WeakAss}_R(M)$ 或 $\text{WeakAss}(M)$。
\end{definition}

\noindent
因此，相伴素理想必为弱相伴素理想。
下面用在该素理想处的局部化给出一个刻画。

\begin{lemma}
\label{algebra-lemma-weakly-ass-local}
设 $R$ 为环，$M$ 为 $R$-模，$\mathfrak p$ 为 $R$ 的素理想。
下列条件等价：
\begin{enumerate}
\item $\mathfrak p$ 与 $M$ 弱相伴；
\item $\mathfrak pR_{\mathfrak p}$ 与 $M_{\mathfrak p}$ 弱相伴；以及
\item $M_{\mathfrak p}$ 含有一个元素，其零化子的根等于
$\mathfrak pR_{\mathfrak p}$。
\end{enumerate}
\end{lemma}

\begin{proof}
假设 (1) 成立。则存在元素 $m \in M$，使
$\mathfrak p$ 在包含 $m$ 的零化子
$I = \{x \in R \mid xm = 0\}$ 的素理想中极小。
由于局部化正合，$m$ 在 $M_{\mathfrak p}$ 中的零化子为
$I_{\mathfrak p}$。因此，$\mathfrak pR_{\mathfrak p}$ 是
$R_{\mathfrak p}$ 中包含 $m$ 在 $M_{\mathfrak p}$ 中的零化子
$I_{\mathfrak p}$ 的极小素理想。这说明 (2) 成立；又因为它意味着
$\sqrt{I_{\mathfrak p}} = \mathfrak pR_{\mathfrak p}$，
(3) 也成立。

\medskip\noindent
将蕴含 (1) $\Rightarrow$ (3) 应用于 $R_{\mathfrak p}$ 上的
$M_{\mathfrak p}$，可知 (2) $\Rightarrow$ (3)。

\medskip\noindent
最后，假设 (3) 成立。这意味着存在元素
$m/f \in M_{\mathfrak p}$，其零化子的根等于
$\mathfrak pR_{\mathfrak p}$。于是，$m$ 在 $M$ 中的零化子
$I = \{x \in R \mid xm = 0\}$ 满足
$\sqrt{I_{\mathfrak p}} = \mathfrak pR_{\mathfrak p}$。
这显然意味着 $\mathfrak p$ 包含 $I$，且在包含 $I$ 的素理想中极小，
即 (1) 成立。
\end{proof}

\begin{lemma}
\label{algebra-lemma-reduced-weakly-ass-minimal}
对于约化环，该环的弱相伴素理想正是极小素理想。
\end{lemma}

\begin{proof}
设 $(R, \mathfrak m)$ 为约化局部环。假设 $x \in R$ 的零化子的根为
$\mathfrak m$。如果 $\mathfrak m \not = 0$，则 $x$ 不可能是单位，
所以 $x \in \mathfrak m$。特别地，对某个 $n \geq 0$ 有
$x^{1 + n} = 0$。因此 $x = 0$，这与 $x$ 的零化子包含于
$\mathfrak m$ 的假设矛盾。由此可见 $\mathfrak m = 0$，
即 $R$ 是域。由引理 \ref{algebra-lemma-weakly-ass-local}，
这就蕴含本引理的断言。
\end{proof}

\begin{lemma}
\label{algebra-lemma-weakly-ass}
设 $R$ 为环，并设
$0 \to M' \to M \to M'' \to 0$ 是 $R$-模的短正合列。
则 $\text{WeakAss}(M') \subset \text{WeakAss}(M)$ 且
$\text{WeakAss}(M) \subset \text{WeakAss}(M') \cup \text{WeakAss}(M'')$。
\end{lemma}

\begin{proof}
我们将使用引理 \ref{algebra-lemma-weakly-ass-local}
对弱相伴素理想的刻画。设 $\mathfrak p$ 为 $R$ 的素理想。
由于局部化正合，得到短正合列
$0 \to M'_{\mathfrak p} \to M_{\mathfrak p} \to M''_{\mathfrak p} \to 0$。
假设 $m \in M_{\mathfrak p}$ 的零化子的根为
$\mathfrak pR_{\mathfrak p}$。那么，要么 $m$ 在
$M''_{\mathfrak p}$ 中的像 $\overline{m}$ 为零且
$m \in M'_{\mathfrak p}$，要么 $\overline{m}$ 的零化子的根为
$\mathfrak pR_{\mathfrak p}$。这证明了
$\text{WeakAss}(M) \subset \text{WeakAss}(M') \cup \text{WeakAss}(M'')$。
包含 $\text{WeakAss}(M') \subset \text{WeakAss}(M)$
立即由定义得出。
\end{proof}

\begin{lemma}
\label{algebra-lemma-weakly-ass-zero}
\begin{slogan}
每个非零模都有弱相伴素理想。
\end{slogan}
设 $R$ 为环，$M$ 为 $R$-模。则
$$
M = (0) \Leftrightarrow \text{WeakAss}(M) = \emptyset
$$
\end{lemma}

\begin{proof}
如果 $M = (0)$，则依定义
$\text{WeakAss}(M) = \emptyset$。反之，假设 $M \not = 0$。
选取非零元素 $m \in M$。将 $m$ 的零化子写成
$I = \{x \in R \mid xm = 0\}$。于是 $R/I \subset M$。
因此，由引理 \ref{algebra-lemma-weakly-ass}，
$\text{WeakAss}(R/I) \subset \text{WeakAss}(M)$。
但由于 $I \not = R$，$V(I) = \Spec(R/I)$ 含有极小素理想；
参见引理
\ref{algebra-lemma-Zariski-topology} 和
\ref{algebra-lemma-spec-closed}。结论得证。
\end{proof}

\begin{lemma}
\label{algebra-lemma-weakly-ass-support}
设 $R$ 为环，$M$ 为 $R$-模。则
$$
\text{Ass}(M) \subset \text{WeakAss}(M) \subset \text{Supp}(M).
$$
\end{lemma}

\begin{proof}
第一个包含立即由定义得出。如果
$\mathfrak p \in \text{WeakAss}(M)$，则由引理
\ref{algebra-lemma-weakly-ass-local} 有
$M_{\mathfrak p} \not = 0$，故
$\mathfrak p \in \text{Supp}(M)$。
\end{proof}

\begin{lemma}
\label{algebra-lemma-weakly-ass-zero-divisors}
设 $R$ 为环，$M$ 为 $R$-模。并集
$\bigcup_{\mathfrak q \in \text{WeakAss}(M)} \mathfrak q$
正是 $R$ 中在 $M$ 上为零因子的元素所成的集合。
\end{lemma}

\begin{proof}
假设 $f \in \mathfrak q \in \text{WeakAss}(M)$。
则存在元素 $m \in M$，使得 $\mathfrak q$ 在包含
$I = \{x \in R \mid xm = 0\}$ 的素理想中极小。
因此，存在 $g \in R$，$g \not \in \mathfrak q$ 以及 $n > 0$，
使得 $f^ngm = 0$。注意，由于 $g \not \in I$，
$gm \not = 0$。如果取上述 $n$ 为最小，则
$f (f^{n - 1}gm) = 0$ 且 $f^{n - 1}gm \not = 0$，
所以 $f$ 是 $M$ 上的零因子。
反之，假设 $f \in R$ 是 $M$ 上的零因子。考虑子模
$N = \{m \in M \mid fm = 0\}$。由于 $N$ 非零，由引理
\ref{algebra-lemma-weakly-ass-zero}，它有弱相伴素理想 $\mathfrak q$。
显然 $f \in \mathfrak q$；又由引理
\ref{algebra-lemma-weakly-ass}，$\mathfrak q$ 是 $M$ 的弱相伴素理想。
\end{proof}

\begin{lemma}
\label{algebra-lemma-weakly-ass-minimal-prime-support}
设 $R$ 为环，$M$ 为 $R$-模。在 $\text{Supp}(M)$ 的元素中
极小的任意 $\mathfrak p \in \text{Supp}(M)$ 都属于
$\text{WeakAss}(M)$。
\end{lemma}

\begin{proof}
注意，在 $\Spec(R_{\mathfrak p})$ 中有
$\text{Supp}(M_{\mathfrak p}) = \{\mathfrak pR_{\mathfrak p}\}$。
特别地，$M_{\mathfrak p}$ 非零，所以由引理
\ref{algebra-lemma-weakly-ass-zero}，
$\text{WeakAss}(M_{\mathfrak p}) \not = \emptyset$。
又由引理
\ref{algebra-lemma-weakly-ass-support} 有
$\text{WeakAss}(M_{\mathfrak p}) \subset \text{Supp}(M_{\mathfrak p})$，
故
$\text{WeakAss}(M_{\mathfrak p}) = \{\mathfrak pR_{\mathfrak p}\}$。
于是由引理
\ref{algebra-lemma-weakly-ass-local}，
$\mathfrak p \in \text{WeakAss}(M)$。
\end{proof}

\begin{lemma}
\label{algebra-lemma-ass-weakly-ass}
设 $R$ 为环，$M$ 为 $R$-模，并设 $\mathfrak p$ 是 $R$
的有限生成素理想。则
$$
\mathfrak p \in \text{Ass}(M) \Leftrightarrow
\mathfrak p \in \text{WeakAss}(M).
$$
特别地，如果 $R$ 是诺特环，则
$\text{Ass}(M) = \text{WeakAss}(M)$。
\end{lemma}

\begin{proof}
写 $\mathfrak p = (g_1, \ldots, g_n)$，其中 $g_i \in R$。
由于另一个蕴含一般成立，只须证明蕴含“$\Leftarrow$”；
参见引理 \ref{algebra-lemma-weakly-ass-support}。
假设 $\mathfrak p \in \text{WeakAss}(M)$。由引理
\ref{algebra-lemma-weakly-ass-local}，存在元素 $m \in M_{\mathfrak p}$，
使得
$I = \{x \in R_{\mathfrak p} \mid xm = 0\}$ 的根为
$\mathfrak pR_{\mathfrak p}$。因此，对每个 $i$，存在最小的
$e_i > 0$，使得 $g_i^{e_i}m = 0$ 在 $M_{\mathfrak p}$ 中成立。
如果对某个 $i$ 有 $e_i > 1$，则可将 $m$ 替换为
$g_i^{e_i - 1} m \not = 0$，并减小 $\sum e_i$。
因此可以假设 $m \in M_{\mathfrak p}$ 的零化子为
$(g_1, \ldots, g_n)R_{\mathfrak p} = \mathfrak p R_{\mathfrak p}$。
由引理 \ref{algebra-lemma-associated-primes-localize} 可知
$\mathfrak p \in \text{Ass}(M)$。
\end{proof}

\begin{remark}
\label{algebra-remark-weakly-ass-not-functorial}
设 $\varphi : R \to S$ 为环映射，$M$ 为 $S$-模。
与相伴素理想的情形相反，并非总有
$\Spec(\varphi)(\text{WeakAss}_S(M)) \subset \text{WeakAss}_R(M)$
（参见引理 \ref{algebra-lemma-ass-functorial}）。
一个例子是考虑环映射
$$
R = k[x_1, x_2, x_3, \ldots] \to
S = k[x_1, x_2, x_3, \ldots, y_1, y_2, y_3, \ldots]/
(x_1y_1, x_2y_2, x_3y_3, \ldots)
$$
以及 $M = S$。在这种情形下，
$\mathfrak q = \sum x_iS$ 是 $S$ 的极小素理想，因而是
$M = S$ 的弱相伴素理想（参见引理
\ref{algebra-lemma-weakly-ass-minimal-prime-support}）。
但另一方面，$S$ 的任意非零元素在 $R$ 中的零化子都是有限生成的，
所以其根不可能等于
$R \cap \mathfrak q = (x_1, x_2, x_3, \ldots)$
（细节从略）。
\end{remark}

\begin{lemma}
\label{algebra-lemma-weakly-ass-reverse-functorial}
设 $\varphi : R \to S$ 为环映射，$M$ 为 $S$-模。则
$\Spec(\varphi)(\text{WeakAss}_S(M)) \supset \text{WeakAss}_R(M)$。
\end{lemma}

\begin{proof}
设 $\mathfrak p$ 是 $\text{WeakAss}_R(M)$ 的元素。
则存在 $m \in M_{\mathfrak p}$，其零化子
$I = \{x \in R_{\mathfrak p} \mid xm = 0\}$ 的根为
$\mathfrak pR_{\mathfrak p}$。考虑 $m$ 在 $S_{\mathfrak p}$ 中的零化子
$J = \{x \in S_{\mathfrak p} \mid xm = 0 \}$。
由于 $IS_{\mathfrak p} \subset J$，可知包含 $J$ 的任意极小素理想
$\mathfrak q \subset S_{\mathfrak p}$ 都位于 $\mathfrak p$ 之上。
此外，例如由引理
\ref{algebra-lemma-weakly-ass-local}，这样的 $\mathfrak q$
对应于 $M$ 的一个弱相伴素理想。
\end{proof}

\begin{remark}
\label{algebra-remark-ass-functorial}
设 $\varphi : R \to S$ 为环映射，$M$ 为 $S$-模。
将谱上相伴的映射记作
$f : \Spec(S) \to \Spec(R)$。则有
$$
f(\text{Ass}_S(M)) \subset
\text{Ass}_R(M) \subset
\text{WeakAss}_R(M) \subset
f(\text{WeakAss}_S(M))
$$
参见引理
\ref{algebra-lemma-ass-functorial}、
\ref{algebra-lemma-weakly-ass-reverse-functorial} 和
\ref{algebra-lemma-weakly-ass-support}。
一般而言，所有这些包含都可能是严格的；参见注
\ref{algebra-remark-ass-reverse-functorial} 和
\ref{algebra-remark-weakly-ass-not-functorial}。
如果 $S$ 是诺特环，那么所有这些包含都是等式，因为由引理
\ref{algebra-lemma-ass-weakly-ass}，最外侧的两个包含是等式。
\end{remark}

\begin{lemma}
\label{algebra-lemma-weakly-ass-finite-ring-map}
设 $\varphi : R \to S$ 为环映射，$M$ 为 $S$-模。
将谱上相伴的映射记作
$f : \Spec(S) \to \Spec(R)$。如果 $\varphi$ 是有限环映射，则
$$
\text{WeakAss}_R(M) = f(\text{WeakAss}_S(M)).
$$
\end{lemma}

\begin{proof}
其中一个包含已经证明；参见注
\ref{algebra-remark-ass-functorial}。为证明另一个包含，假设
$\mathfrak q \in \text{WeakAss}_S(M)$，并设
$\mathfrak p$ 为 $R$ 中对应的素理想。设 $m \in M$ 满足：
$\mathfrak q$ 是包含
$J = \{g \in S \mid gm = 0\}$ 的极小素理想。因此
$JS_{\mathfrak q}$ 的根为 $\mathfrak qS_{\mathfrak q}$。
由于 $R \to S$ 是有限的，$\mathfrak p$ 上方只有有限多个素理想
$\mathfrak q = \mathfrak q_1, \mathfrak q_2, \ldots, \mathfrak q_l$；
参见引理 \ref{algebra-lemma-finite-finite-fibres}。
选取 $x \in \mathfrak q$，使得对 $i > 1$ 有
$x \not \in \mathfrak q_i$；参见引理
\ref{algebra-lemma-silly}。由上述讨论，存在元素
$y \in S$，$y \not \in \mathfrak q$ 以及整数 $t > 0$，
使得 $y x^t m = 0$。因此，元素 $ym \in M$ 被
$x^t$ 零化，故对 $i = 2, \ldots, l$，$ym$ 在
$M_{\mathfrak q_i}$ 中映为零。须注意，$ym$ 在
$S_{\mathfrak q}$ 中并不映为零。

\medskip\noindent
由有限环映射的上升性质，环 $S_{\mathfrak p}$ 是半局部环，
其极大理想为 $\mathfrak q_i S_{\mathfrak p}$；参见引理
\ref{algebra-lemma-integral-going-up}。如果
$f \in \mathfrak pR_{\mathfrak p}$，则 $f$ 的某个幂落在
$JS_{\mathfrak q}$ 中，故对某个 $t > 0$，可知
$f^t ym$ 在 $M_{\mathfrak q}$ 中映为零。
由于 $ym$ 在 $S_{\mathfrak p}$ 的其他极大理想处为零，
由引理 \ref{algebra-lemma-characterize-zero-local} 可知
$f^t ym$ 在 $M_{\mathfrak p}$ 中为零。
由此可见，$\mathfrak p$ 是包含 $ym$ 在 $R$ 中的零化子的
极小素理想，结论得证。
\end{proof}

\begin{lemma}
\label{algebra-lemma-weakly-ass-quotient-ring}
设 $R$ 为环，$I$ 为理想，$M$ 为 $R/I$-模。
通过典范单射
$\Spec(R/I) \to \Spec(R)$，有
$\text{WeakAss}_{R/I}(M) = \text{WeakAss}_R(M)$。
\end{lemma}

\begin{proof}
这是引理 \ref{algebra-lemma-weakly-ass-finite-ring-map} 的特殊情形。
\end{proof}

\begin{lemma}
\label{algebra-lemma-localize-weakly-ass}
设 $R$ 为环，$M$ 为 $R$-模，并设 $S \subset R$ 为乘法子集。
通过典范单射 $\Spec(S^{-1}R) \to \Spec(R)$，有
$\text{WeakAss}_R(S^{-1}M) = \text{WeakAss}_{S^{-1}R}(S^{-1}M)$
以及
$$
\text{WeakAss}(M) \cap \Spec(S^{-1}R) = \text{WeakAss}(S^{-1}M).
$$
\end{lemma}

\begin{proof}
假设 $m \in S^{-1}M$。令
$I = \{x \in R \mid xm = 0\}$ 且
$I' = \{x' \in S^{-1}R \mid x'm = 0\}$。则
$I' = S^{-1}I$，并且除非 $I = R$，否则
$I \cap S = \emptyset$（验证从略）。
因此，$S^{-1}R$ 中包含 $I'$ 的极小素理想，与 $R$ 中
包含 $I$ 且避开 $S$ 的极小素理想一一对应。这证明了等式
$\text{WeakAss}_R(S^{-1}M) = \text{WeakAss}_{S^{-1}R}(S^{-1}M)$。
第二个等式由引理
\ref{algebra-lemma-weakly-ass-local} 得出，因为对
$\mathfrak p \in R$，若 $S \cap \mathfrak p = \emptyset$，则
$M_{\mathfrak p} = (S^{-1}M)_{S^{-1}\mathfrak p}$。
\end{proof}

\begin{lemma}
\label{algebra-lemma-localize-weakly-ass-nonzero-divisors}
设 $R$ 为环，$M$ 为 $R$-模，并设 $S \subset R$ 为乘法子集。
假设每个 $s \in S$ 都是 $M$ 上的非零因子。则
$$
\text{WeakAss}(M) = \text{WeakAss}(S^{-1}M).
$$
\end{lemma}

\begin{proof}
由假设有 $M \subset S^{-1}M$，故引理
\ref{algebra-lemma-weakly-ass} 给出
$\text{WeakAss}(M) \subset \text{WeakAss}(S^{-1}M)$。
反之，假设 $n/s \in S^{-1}M$ 的零化子为 $I$，且
$\mathfrak p$ 是包含 $I$ 的极小素理想。则
$n \in M$ 的零化子为 $I$，且 $\mathfrak p$ 是包含
$I$ 的极小素理想。
\end{proof}

\begin{lemma}
\label{algebra-lemma-zero-at-weakly-ass-zero}
设 $R$ 为环，$M$ 为 $R$-模。映射
$$
M
\longrightarrow
\prod\nolimits_{\mathfrak p \in \text{WeakAss}(M)} M_{\mathfrak p}
$$
是单射。
\end{lemma}

\begin{proof}
设 $x \in M$ 为该映射的核中的元素。令
$N = Rx \subset M$。如果 $\mathfrak p$ 是 $N$ 的弱相伴素理想，
那么一方面 $\mathfrak p \in \text{WeakAss}(M)$
（引理 \ref{algebra-lemma-weakly-ass}），另一方面
$N_{\mathfrak p} \subset M_{\mathfrak p}$ 非零。
此矛盾说明 $\text{WeakAss}(N) = \emptyset$。
因此，由引理
\ref{algebra-lemma-weakly-ass-zero}，$N = 0$，即 $x = 0$。
\end{proof}

\begin{lemma}
\label{algebra-lemma-weak-post-bourbaki}
设 $R \to S$ 为环映射，$N$ 为 $S$-模。
假设 $N$ 作为 $R$-模是平坦的，并且 $R$ 是分式域为 $K$ 的整环。
则
$$
\text{WeakAss}_S(N) = \text{WeakAss}_{S \otimes_R K}(N \otimes_R K)
$$
其中使用典范包含
$\Spec(S \otimes_R K) \subset \Spec(S)$。
\end{lemma}

\begin{proof}
注意 $S \otimes_R K = (R \setminus \{0\})^{-1}S$ 且
$N \otimes_R K = (R \setminus \{0\})^{-1}N$。
对任意非零 $x \in R$，由于 $N$ 在 $R$ 上平坦，
乘以 $x$ 在 $N$ 上是单射。因此，本引理由引理
\ref{algebra-lemma-localize-weakly-ass-nonzero-divisors} 得出。
\end{proof}

\begin{lemma}
\label{algebra-lemma-weakly-ass-change-fields}
设 $K/k$ 为域扩张，$R$ 为 $k$-代数，$M$ 为 $R$-模。
设 $\mathfrak q \subset R \otimes_k K$ 是位于
$\mathfrak p \subset R$ 上方的素理想。如果 $\mathfrak q$
与 $M \otimes_k K$ 弱相伴，那么 $\mathfrak p$ 与 $M$ 弱相伴。
\end{lemma}

\begin{proof}
设 $z \in M \otimes_k K$ 满足：$\mathfrak q$ 是包含 $z$ 的零化子
$J \subset R \otimes_k K$ 的极小素理想。
选择有限生成的子扩张 $K/L/k$，使
$z \in M \otimes_k L$。由于
$R \otimes_k L \to R \otimes_k K$ 平坦，由引理
\ref{algebra-lemma-annihilator-flat-base-change} 可知
$J = I(R \otimes_k K)$，其中 $I \subset R \otimes_k L$
是 $z$ 在较小环中的零化子。因此，由下降定理
（引理 \ref{algebra-lemma-flat-going-down}），
$\mathfrak q \cap (R \otimes_k L)$ 是包含 $I$ 的极小素理想。
由此归约到下一段所述的情形。

\medskip\noindent
假设 $K/k$ 是有限生成域扩张。设
$x_1, \ldots, x_r \in K$ 是 $K$ 在 $k$ 上的超越基；
参见《域》的第 \ref{fields-section-transcendence} 节。
令 $L = k(x_1, \ldots, x_r)$，并设 $[K : L] = n$。
于是 $R \otimes_k L \to R \otimes_k K$ 是有限环映射。
因此，由引理 \ref{algebra-lemma-weakly-ass-finite-ring-map}，
$\mathfrak q \cap (R \otimes_k L)$ 是 $M \otimes_k K$
作为 $R \otimes_k L$-模时的弱相伴素理想。
由于作为 $R \otimes_k L$-模有
$M \otimes_k K \cong (M \otimes_k L)^{\oplus n}$，
可知 $\mathfrak q \cap (R \otimes_k L)$ 是
$M \otimes_k L$ 的弱相伴素理想
（例如，使用引理 \ref{algebra-lemma-weakly-ass} 并作归纳）。
由此归约到下一段讨论的情形。

\medskip\noindent
假设 $K = k(x_1, \ldots, x_r)$ 是纯超越域扩张。
可以将 $R$ 替换为 $R_\mathfrak p$，将 $M$ 替换为
$M_\mathfrak p$，并将 $\mathfrak q$ 替换为
$\mathfrak q(R_\mathfrak p \otimes_k K)$。
参见引理 \ref{algebra-lemma-localize-weakly-ass}。
由此归约到下一段讨论的情形。

\medskip\noindent
假设 $K = k(x_1, \ldots, x_r)$ 是纯超越域扩张，且 $R$
是极大理想为 $\mathfrak p$ 的局部环。我们断言，任意
$f \in R \otimes_k K$，$f \not \in \mathfrak p(R \otimes_k K)$，
都是 $M \otimes_k K$ 上的非零因子。事实上，设
$z \in M \otimes_k K$。存在有限 $R$-子模
$M' \subset M$，使得 $z \in M' \otimes_k K$，而且 $M'$
在具有此性质的子模中极小：选取 $K$ 作为 $k$-向量空间的一组基
$\{t_\alpha\}$，写
$z = \sum m_\alpha \otimes t_\alpha$，并令 $M'$ 为诸
$m_\alpha$ 生成的 $R$-子模。如果
$z \in \mathfrak p(M' \otimes_k K) = \mathfrak p M' \otimes_k K$，
那么 $\mathfrak pM' = M'$，由引理 \ref{algebra-lemma-NAK}
可得 $M' = 0$，矛盾。
因此，$z$ 在 $M'/\mathfrak p M' \otimes_k K$ 中的像
$\overline{z}$ 非零。但
$R/\mathfrak p \otimes_k K$ 是
$\kappa(\mathfrak p)[x_1, \ldots, x_n]$ 的一个局部化，故为整环；
而 $M'/\mathfrak p M' \otimes_k K$ 是自由模，所以
$f\overline{z} \not = 0$。断言得证。

\medskip\noindent
最后，选取 $z \in M \otimes_k K$，使得 $\mathfrak q$
是包含 $z$ 的零化子 $J \subset R \otimes_k K$ 的极小素理想。
对 $f \in \mathfrak p$，存在 $n \geq 1$ 及
$g \in R \otimes_k K$，$g \not \in \mathfrak q$，使得
$g f^n z \in J$，即 $g f^n z = 0$。
（这是因为 $\mathfrak q$ 位于 $\mathfrak p$ 上方，且
$\mathfrak q$ 是包含 $J$ 的极小素理想。）
上面已经看到 $g$ 是非零因子，所以 $f^n z = 0$。
这意味着，将 $M \otimes_k K$ 视为 $R$-模时，
$\mathfrak p$ 是它的弱相伴素理想。由于
$M \otimes_k K$ 是若干个 $M$ 的直和，和上面一样可知
$\mathfrak p$ 是 $M$ 的弱相伴素理想。
\end{proof}

\section{嵌入素理想}
\label{algebra-section-embedded-primes}

\noindent
定义如下。

\begin{definition}
\label{algebra-definition-embedded-primes}
设 $R$ 为环，$M$ 为 $R$-模。
\begin{enumerate}
\item $M$ 的相伴素理想中，那些在 $M$ 的相伴素理想中不极小的，
称为 $M$ 的{\it 嵌入相伴素理想}。
\item {\it $R$ 的嵌入素理想}是将 $R$ 视为 $R$-模时的
嵌入相伴素理想。
\end{enumerate}
\end{definition}

\noindent
下面给出一种去除这些素理想的方法。

\begin{lemma}
\label{algebra-lemma-remove-embedded-primes}
设 $R$ 为诺特环，$M$ 为有限 $R$-模。考虑如下 $R$-子模集合：
$$
\{
K \subset M
\mid
\text{Supp}(K)
\text{ nowhere dense in }
\text{Supp}(M)
\}.
$$
这个集合有极大元 $K$，且商
$M' = M/K$ 具有下列性质：
\begin{enumerate}
\item $\text{Supp}(M) = \text{Supp}(M')$；
\item $M'$ 没有嵌入相伴素理想；
\item 对任意包含于 $M$ 的所有嵌入相伴素理想中的
$f \in R$，有 $M_f \cong M'_f$。
\end{enumerate}
\end{lemma}

\begin{proof}
下文将直接使用引理 \ref{algebra-lemma-finite-ass} 和命题
\ref{algebra-proposition-minimal-primes-associated-primes}，不再逐次说明。
令 $\mathfrak q_1, \ldots, \mathfrak q_t$ 表示 $M$ 的支集中的
极小素理想，令
$\mathfrak p_1, \ldots, \mathfrak p_s$ 表示 $M$ 的
嵌入相伴素理想。则
$\text{Ass}(M) = \{\mathfrak q_j, \mathfrak p_i\}$。
令
$$
K = \{m \in M \mid \text{Supp}(Rm) \subset \bigcup V(\mathfrak p_i)\}
$$
立即可见它是子模。由于 $M$ 是诺特环上的有限模，
$K$ 也是有限的。因此，$\text{Supp}(K)$ 在
$\text{Supp}(M)$ 中无处稠密。设 $K' \subset M$ 是另一个
支集在 $\text{Supp}(M)$ 中无处稠密的子模。这意味着
$K_{\mathfrak q_j} = 0$。因此，如果 $m \in K'$，
那么 $m$ 在 $M_{\mathfrak q_j}$ 中映为零，进而
$(Rm)_{\mathfrak q_j} = 0$。另一方面，有
$\text{Ass}(Rm) \subset \text{Ass}(M)$。
所以 $Rm$ 的支集包含于 $\bigcup V(\mathfrak p_i)$。
因此 $m \in K$；由于 $m$ 是 $K'$ 的任意元素，故
$K' \subset K$。

\medskip\noindent
令 $M' = M/K$。由于 $K_{\mathfrak q_j}=0$，可知对所有 $j$ 有
$M'_{\mathfrak q_j} = M_{\mathfrak q_j}$。
因此 $M$ 与 $M'$ 具有相同的支集。

\medskip\noindent
假设 $\mathfrak q = \text{Ann}(\overline{m}) \in \text{Ass}(M')$，
其中 $\overline{m} \in M'$ 是 $m \in M$ 的像。
于是 $m \not \in K$，所以 $Rm$ 的支集必含有某个
$\mathfrak q_j$。由于
$M_{\mathfrak q_j} = M'_{\mathfrak q_j}$，可知
$\overline{m}$ 在 $M'_{\mathfrak q_j}$ 中的像非零。
因此 $\mathfrak q \subset \mathfrak q_j$
（事实上二者相等），这意味着 $M'$ 的所有相伴素理想都不是嵌入的。

\medskip\noindent
设 $f$ 是包含于所有 $\mathfrak p_i$ 中的元素。
则 $D(f) \cap \text{supp}(K) = 0$。因此
$M_f = M'_f$，因为 $K_f = 0$。
\end{proof}

\begin{lemma}
\label{algebra-lemma-remove-embedded-primes-localize}
设 $R$ 为诺特环，$M$ 为有限 $R$-模。对任意 $f \in R$，有
$(M')_f = (M_f)'$，其中 $M \to M'$ 与
$M_f \to (M_f)'$ 是引理
\ref{algebra-lemma-remove-embedded-primes} 中构造的商。
\end{lemma}

\begin{proof}
证明从略。
\end{proof}

\begin{lemma}
\label{algebra-lemma-no-embedded-primes-endos}
设 $R$ 为诺特环，$M$ 为没有嵌入相伴素理想的有限 $R$-模。
令 $I = \{x \in R \mid xM = 0\}$。则环 $R/I$
没有嵌入素理想。
\end{lemma}

\begin{proof}
可以将 $R$ 替换为 $R/I$。因此，可以假设 $R$ 的每个非零元素
都在 $M$ 上非平凡地作用。由引理
\ref{algebra-lemma-support-closed}，这意味着 $\Spec(R)$ 等于 $M$ 的支集。
假设 $\mathfrak p$ 是 $R$ 的嵌入素理想。设 $x \in R$
的零化子为 $\mathfrak p$。考虑非零模
$N = xM \subset M$。它被 $\mathfrak p$ 零化。
因此，$N$ 的任意相伴素理想 $\mathfrak q$ 都包含
$\mathfrak p$，并且也是 $M$ 的相伴素理想。
于是 $\mathfrak q$ 将是 $M$ 的嵌入相伴素理想，
这与本引理的假设矛盾。
\end{proof}

\section{正则序列}
\label{algebra-section-regular-sequences}

\noindent
本节研究正则序列的一些基本性质。

\begin{definition}
\label{algebra-definition-regular-sequence}
设 $R$ 为环，$M$ 为 $R$-模。若 $R$ 中的元素序列
$f_1, \ldots, f_r$ 满足下列条件，则称它为
{\it $M$-正则序列}：
\begin{enumerate}
\item 对每个 $i = 1, \ldots, r$，$f_i$ 是
$M/(f_1, \ldots, f_{i - 1})M$ 上的非零因子；并且
\item 模 $M/(f_1, \ldots, f_r)M$ 非零。
\end{enumerate}
如果 $I$ 是 $R$ 的理想且 $f_1, \ldots, f_r \in I$，
则称 $f_1, \ldots, f_r$ 为{\it $I$ 中的 $M$-正则序列}。
如果 $M = R$，则简称 $f_1, \ldots, f_r$ 为
{\it 正则序列}（在 $I$ 中）。
\end{definition}

\noindent
请注意，这个定义依赖于元素 $f_1, \ldots, f_r$ 的次序
（参见下面的例子）。有些论文或书籍去掉了模
$M/(f_1, \ldots, f_r)M$ 非零这一要求。
这样做的优点是，正则序列这一性质在局部化下保持。
然而，我们主要用这个定义来定义 $R$ 为局部环时模的深度；
在这种情形下，要求 $f_i$ 属于极大理想——而从 $R$ 过渡到局部化
$R_\mathfrak p$ 时，这一条件并不保持。

\begin{example}
\label{algebra-example-global-regular}
设 $k$ 为域。在环 $k[x, y, z]$ 中，序列
$x, y(1-x), z(1-x)$ 是正则的，但序列
$y(1-x), z(1-x), x$ 不是。
\end{example}

\begin{example}
\label{algebra-example-local-regular}
设 $k$ 为域。考虑环
$k[x, y, w_0, w_1, w_2, \ldots]/I$，
其中 $I$ 由 $yw_i$（$i = 0, 1, 2, \ldots$）和
$w_i - xw_{i + 1}$（$i = 0, 1, 2, \ldots$）生成。
序列 $x, y$ 是正则的，但 $y$ 是零因子。
此外，还可以在极大理想 $(x, y, w_i)$ 处局部化，
仍得到一个例子。
\end{example}

\begin{lemma}
\label{algebra-lemma-permute-xi}
设 $R$ 为诺特局部环，$M$ 为有限 $R$-模，并设
$x_1, \ldots, x_c$ 为 $M$-正则序列。则 $x_i$ 的任意排列
也为正则序列。
\end{lemma}

\begin{proof}
先处理 $c = 2$ 的情形。考虑
$K \subset M$，即 $x_2 : M \to M$ 的核。
对任意 $z \in K$，由于 $x_2$ 是 $M/x_1M$ 上的非零因子，
存在某个 $z' \in M$，使得 $z = x_1 z'$。
由于 $x_1$ 是 $M$ 上的非零因子，可知 $x_2 z' = 0$。
因此 $x_1 : K \to K$ 是满射。由 Nakayama 引理
\ref{algebra-lemma-NAK}，$K = 0$。
接下来，考虑 $x_1$ 在 $M/x_2M$ 上的乘法。如果 $z \in M$
在此映射的核中映为元素 $\overline{z} \in M/x_2M$，
则对某个 $y \in M$ 有 $x_1 z = x_2 y$。
但由于 $x_1, x_2$ 是正则序列，存在某个 $y' \in M$，
使得 $y = x_1 y'$。所以
$x_1 ( z - x_2 y' ) =0$，进而 $z = x_2 y'$，
故按所需有 $\overline{z} = 0$。

\medskip\noindent
对一般情形，注意任意排列都是相邻指标换位的复合。
因此，只须证明
$$
x_1, \ldots, x_{i-2}, x_i, x_{i-1}, x_{i + 1}, \ldots, x_c
$$
是 $M$-正则序列。这由刚刚处理的情形得出：将它应用于模
$M/(x_1, \ldots, x_{i-2})$ 以及长度为 $2$ 的正则序列
$x_{i-1}, x_i$。
\end{proof}

\begin{lemma}
\label{algebra-lemma-flat-increases-depth}
\begin{slogan}
平坦局部环同态保持并反映正则序列。
\end{slogan}
设 $R, S$ 为局部环，$R \to S$ 为平坦局部环同态。
设 $x_1, \ldots, x_r$ 为 $R$ 中的序列，$M$ 为 $R$-模。
下列条件等价：
\begin{enumerate}
\item $x_1, \ldots, x_r$ 是 $R$ 中的 $M$-正则序列；并且
\item $x_1, \ldots, x_r$ 在 $S$ 中的像构成
$M \otimes_R S$-正则序列。
\end{enumerate}
\end{lemma}

\begin{proof}
这是因为由引理 \ref{algebra-lemma-local-flat-ff}，
$R \to S$ 忠实平坦。
\end{proof}

\begin{lemma}
\label{algebra-lemma-regular-sequence-in-neighbourhood}
设 $R$ 为诺特环，$M$ 为有限 $R$-模，$\mathfrak p$ 为素理想。
设 $x_1, \ldots, x_r$ 为 $R$ 中的序列，它在
$R_{\mathfrak p}$ 中的像构成 $M_{\mathfrak p}$-正则序列。
则存在 $g \in R$，$g \not \in \mathfrak p$，使得
$x_1, \ldots, x_r$ 在 $R_g$ 中的像构成 $M_g$-正则序列。
\end{lemma}

\begin{proof}
令
$$
K_i = \Ker\left(x_i : M/(x_1, \ldots, x_{i - 1})M \to
M/(x_1, \ldots, x_{i - 1})M\right).
$$
这是有限 $R$-模，依假设，它在 $\mathfrak p$ 处的局部化为零。
因此存在 $g \in R$，$g \not \in \mathfrak p$，使得对所有
$i = 1, \ldots, r$ 都有 $(K_i)_g = 0$。这个 $g$ 即满足要求。
\end{proof}

\begin{lemma}
\label{algebra-lemma-join-regular-sequences}
设 $A$ 为环，$I$ 为 $A$ 中由正则序列
$f_1, \ldots, f_n$ 生成的理想。设
$g_1, \ldots, g_m \in A$ 的像
$\overline{g}_1, \ldots, \overline{g}_m$ 在 $A/I$ 中构成正则序列。
则 $f_1, \ldots, f_n, g_1, \ldots, g_m$ 是 $A$ 中的正则序列。
\end{lemma}

\begin{proof}
这立即由定义得出。
\end{proof}

\begin{lemma}
\label{algebra-lemma-regular-sequence-short-exact-sequence}
设 $R$ 为环，$0 \to M_1 \to M_2 \to M_3 \to 0$
为 $R$-模的短正合列，并设 $f_1, \ldots, f_r \in R$。
如果 $f_1, \ldots, f_r$ 既为 $M_1$-正则序列又为
$M_3$-正则序列，则 $f_1, \ldots, f_r$ 也是 $M_2$-正则序列。
\end{lemma}

\begin{proof}
由引理 \ref{algebra-lemma-snake}，如果
$f_1 : M_1 \to M_1$ 和 $f_1 : M_3 \to M_3$ 都是单射，
那么 $f_1 : M_2 \to M_2$ 也是单射，并得到短正合列
$$
0 \to M_1/f_1M_1 \to M_2/f_1M_2 \to M_3/f_1M_3 \to 0
$$
本引理由此及对 $r$ 的归纳得出。一些细节从略。
\end{proof}

\begin{lemma}
\label{algebra-lemma-regular-sequence-powers}
设 $R$ 为环，$M$ 为 $R$-模。设
$f_1, \ldots, f_r \in R$，且 $e_1, \ldots, e_r > 0$ 为整数。
则 $f_1, \ldots, f_r$ 是 $M$-正则序列，当且仅当
$f_1^{e_1}, \ldots, f_r^{e_r}$ 是 $M$-正则序列。
\end{lemma}

\begin{proof}
对 $r$ 作归纳。如果 $r = 1$，结论由下面两个简单事实得出：
(a) $M$ 上非零因子的幂仍为 $M$ 上的非零因子；
(b) $M$ 上非零因子的因子仍为 $M$ 上的非零因子。
如果 $r > 1$，则将归纳假设应用于 $M/f_1M$，可知
$f_1, f_2, \ldots, f_r$ 是 $M$-正则序列，当且仅当
$f_1, f_2^{e_2}, \ldots, f_r^{e_r}$ 是 $M$-正则序列。
所以只须证明：给定 $e > 0$，
$f_1^e, f_2, \ldots, f_r$ 是 $M$-正则序列，当且仅当
$f_1, \ldots, f_r$ 是 $M$-正则序列。
我们对 $e$ 作归纳。$e = 1$ 的情形是平凡的。
在两种假设下，$f_1$ 都是非零因子
（由 $r = 1$ 的情形），所以有短正合列
$$
0 \to M/f_1M \xrightarrow{f_1^{e - 1}} M/f_1^eM \to M/f_1^{e - 1}M \to 0
$$
假设 $f_1, f_2, \ldots, f_r$ 是 $M$-正则序列。
则由归纳假设，元素 $f_2, \ldots, f_r$ 分别构成
$M/f_1M$-正则序列和 $M/f_1^{e - 1}M$-正则序列。
由引理 \ref{algebra-lemma-regular-sequence-short-exact-sequence}，
$f_2, \ldots, f_r$ 是 $M/f_1^eM$-正则序列。因此
$f_1^e, f_2, \ldots, f_r$ 是 $M$-正则序列。
反之，假设 $f_1^e, f_2, \ldots, f_r$ 是 $M$-正则序列。
那么 $f_2 : M/f_1^eM \to M/f_1^eM$ 是单射，故
$f_2 : M/f_1M \to M/f_1M$ 是单射，再由归纳（！），
$f_2 : M/f_1^{e - 1}M \to M/f_1^{e - 1}M$ 是单射，因而
$$
0 \to
M/(f_1, f_2)M \xrightarrow{f_1^{e - 1}}
M/(f_1^e, f_2)M \to
M/(f_1^{e - 1}, f_2)M \to 0
$$
由引理 \ref{algebra-lemma-snake} 是短正合列。这证明了 $r = 2$ 时的逆命题。
如果 $r > 2$，那么
$f_3 : M/(f_1^e, f_2)M \to M/(f_1^e, f_2)M$ 是单射，故
$f_3 : M/(f_1, f_2)M \to M/(f_1, f_2)M$ 是单射，依此类推。
一些细节从略。
\end{proof}

\begin{lemma}
\label{algebra-lemma-regular-sequence-in-polynomial-ring}
设 $R$ 为环，并设不生成单位理想的
$f_1, \ldots, f_r \in R$。下列条件等价：
\begin{enumerate}
\item $f_1, \ldots, f_r$ 的任意排列都是正则序列；
\item $f_1, \ldots, f_r$ 的任意子序列
（按给定次序）都是正则序列；以及
\item 在多项式环 $R[x_1, \ldots, x_r]$ 中，
$f_1x_1, \ldots, f_rx_r$ 是正则序列。
\end{enumerate}
\end{lemma}

\begin{proof}
显然 (1) 蕴含 (2)。下面对 $r$ 作归纳，证明 (2) 蕴含 (1)。
$r = 1$ 的情形是平凡的。$r = 2$ 的情形断言：如果
$a, b \in R$ 是正则序列且 $b$ 是非零因子，那么
$b, a$ 是正则序列。这是清楚的，因为如果 $a$ 与 $b$
都是非零因子，则
$a : R/(b) \to R/(b)$ 的核同构于
$b : R/(a) \to R/(a)$ 的核。
现在考虑 $r > 2$ 的情形。假设 (2) 成立，并设我们要证明
对某个排列 $\sigma$，$f_{\sigma(1)}, \ldots, f_{\sigma(r)}$
是正则序列。由归纳假设，已经知道
$f_{\sigma(1)}, \ldots, f_{\sigma(r - 1)}$ 是正则序列。
因此，只须说明当 $s = \sigma(r)$ 时，$f_s$ 模掉
$f_1, \ldots, \hat f_s, \ldots, f_r$ 后是非零因子。
如果 $s = r$，结论已经成立。如果 $s < r$，则注意
$f_s$ 和 $f_r$ 在环
$R/(f_1, \ldots, \hat f_s, \ldots, f_{r - 1})$
中都是非零因子（再次由归纳假设）。由于已知 $f_s, f_r$
在该环中是正则序列，由长度为 $2$ 的序列的情形可知
$f_r, f_s$ 也是正则序列。

\medskip\noindent
注意，作为 $R$-模，
$R[x_1, \ldots, x_r]/(f_1x_1, \ldots, f_ix_i)$
是如下诸模的直和：
$$
R/I_E \cdot x_1^{e_1} \ldots x_r^{e_r}
$$
其指标为多重指标 $E = (e_1, \ldots, e_r)$，其中
$I_E$ 是由满足 $1 \leq j \leq i$ 且 $e_j > 0$ 的 $f_j$
生成的理想。因此，$f_{i + 1}x_i$ 在此模上是非零因子，
当且仅当对所有 $E$，$f_{i + 1}$ 都是 $R/I_E$ 上的非零因子。
取所有分量均为正的 $E$，可知 $f_{i + 1}$ 是
$R/(f_1, \ldots, f_i)$ 上的非零因子。因此 (3) 蕴含 (2)。
反之，如果 (2) 成立，那么
$f_1, \ldots, f_i, f_{i + 1}$ 的任意子序列都是正则序列；
特别地，$f_{i + 1}$ 是所有 $R/I_E$ 上的非零因子。
由此可见 (2) 蕴含 (3)。
\end{proof}

\section{拟正则序列}
\label{algebra-section-quasi-regular}

\noindent
我们引入拟正则序列的概念；它比正则序列稍弱，但更容易使用。
设 $R$ 为环，且 $f_1, \ldots, f_c \in R$。
令 $J = (f_1, \ldots, f_c)$，并设 $M$ 为 $R$-模。
则存在分次 $R/J[X_1, \ldots, X_c]$-模的典范映射
\begin{equation}
\label{algebra-equation-quasi-regular}
M/JM \otimes_{R/J} R/J[X_1, \ldots, X_c]
\longrightarrow
\bigoplus\nolimits_{n \geq 0} J^nM/J^{n + 1}M
\end{equation}
它由如下规则定义：
$$
\overline{m} \otimes X_1^{e_1} \ldots X_c^{e_c} \longmapsto
f_1^{e_1} \ldots f_c^{e_c} m \bmod J^{e_1 + \ldots + e_c + 1}M.
$$
注意，(\ref{algebra-equation-quasi-regular}) 总是满射。

\begin{definition}
\label{algebra-definition-quasi-regular-sequence}
设 $R$ 为环，$M$ 为 $R$-模。如果
(\ref{algebra-equation-quasi-regular}) 是同构，则称 $R$ 的元素序列
$f_1, \ldots, f_c$ 为{\it $M$-拟正则序列}。
如果 $M = R$，则简称 $f_1, \ldots, f_c$ 为
{\it 拟正则序列}。
\end{definition}

\noindent
因此，如果 $f_1, \ldots, f_c$ 是拟正则序列，那么
$$
R/J[X_1, \ldots, X_c] = \bigoplus\nolimits_{n \geq 0} J^n/J^{n + 1}
$$
其中 $J = (f_1, \ldots, f_c)$。显然，拟正则序列这一性质
与 $f_1, \ldots, f_c$ 的次序无关。

\begin{lemma}
\label{algebra-lemma-regular-quasi-regular}
设 $R$ 为环。
\begin{enumerate}
\item $R$ 的正则序列 $f_1, \ldots, f_c$ 是拟正则序列。
\item 假设 $M$ 为 $R$-模，且 $f_1, \ldots, f_c$
是 $M$-正则序列。则 $f_1, \ldots, f_c$ 是
$M$-拟正则序列。
\end{enumerate}
\end{lemma}

\begin{proof}
令 $J = (f_1, \ldots, f_c)$。
我们对 $c$ 作归纳，证明第一个断言。须说明：给定任意关系
$\sum_{|I| = n} a_I f^I \in J^{n + 1}$，其中 $a_I \in R$，
实际上对所有多重指标 $I$ 都有 $a_I \in J$。
由于 $J^{n + 1}$ 的任意元素都可写成
$\sum_{|I| = n} b_I f^I$，其中 $b_I \in J$，
用 $a_I - b_I$ 替换 $a_I$ 后，可以假设该关系为
$\sum_{|I| = n} a_I f^I = 0$。可将它改写为
$$
\sum\nolimits_{e = 0}^n
\left(
\sum\nolimits_{|I'| = n - e}
a_{I', e} f^{I'}
\right)
f_c^e
=
0
$$
此处及下文中，“带撇”的多重指标 $I'$ 都要求具有形式
$I' = (i_1, \ldots, i_{c - 1}, 0)$。
我们将对 $l \in \{0, \ldots, n\}$ 作归纳，证明：如果有关系
$$
\sum\nolimits_{e = 0}^l
\left(
\sum\nolimits_{|I'| = n - e}
a_{I', e} f^{I'}
\right)
f_c^e
=
0
$$
那么对所有 $I', e$ 都有 $a_{I', e} \in J$。
事实上，令 $J' = (f_1, \ldots, f_{c-1})$。
注意
$\sum\nolimits_{|I'| = n - l} a_{I', l} f^{I'}$
在乘以 $f_c^{l}$ 后映入 $(J')^{n - l + 1}$。
由归纳假设（对 $c$ 的归纳），可知
$f_c^l a_{I', l} \in J'$。
由于 $f_1, \ldots, f_c$ 是正则序列，$f_c$
不是 $R/J'$ 上的零因子，所以
$a_{I', l} \in J'$。这使我们可以将项
$(\sum\nolimits_{|I'| = n - l} a_{I', l} f^{I'})f_c^l$
改写成
$(\sum\nolimits_{|I'| = n - l + 1} f_c b_{I', l - 1}
f^{I'})f_c^{l-1}$。于是得到如下形式的新关系：
$$
\left(\sum\nolimits_{|I'| = n - l + 1}
(a_{I', l-1} + f_c b_{I', l - 1}) f^{I'}\right)f_c^{l-1}
+
\sum\nolimits_{e = 0}^{l - 2}
\left(
\sum\nolimits_{|I'| = n - e}
a_{I', e} f^{I'}
\right)
f_c^e
=
0
$$
现在由归纳假设（这次是对 $l$ 的归纳），可知所有
$a_{I', l-1} + f_c b_{I', l - 1} \in J$，并且当
$e \leq l - 2$ 时所有 $a_{I', e} \in J$。
这与上面得到的 $a_{I', l} \in J' \subset J$ 合起来，
完成归纳步骤的证明。

\medskip\noindent
第二个断言意味着：给定任意形式表达式
$F = \sum_{|I| = n} m_I X^I$，其中 $m_I \in M$，且
$\sum m_I f^I \in J^{n + 1}M$，则所有系数 $m_I$ 都属于
$J$。其证明与上面对第一个断言的相应结果的证明完全相同。
\end{proof}

\begin{lemma}
\label{algebra-lemma-flat-base-change-quasi-regular}
设 $R \to R'$ 为平坦环映射，$M$ 为 $R$-模。
假设 $f_1, \ldots, f_r \in R$ 构成 $M$-拟正则序列。
则 $f_1, \ldots, f_r$ 在 $R'$ 中的像构成
$M \otimes_R R'$-拟正则序列。
\end{lemma}

\begin{proof}
令 $J = (f_1, \ldots, f_r)$、$J' = JR'$ 且
$M' = M \otimes_R R'$。须证明典范映射
$\mu : R'/J'[X_1, \ldots X_r] \otimes_{R'/J'} M'/J'M' \to
\bigoplus (J')^nM'/(J')^{n + 1}M'$ 是同构。
因为 $R \to R'$ 平坦，正合列
$0 \to J^nM \to M$ 和
$0 \to J^{n + 1}M \to J^nM \to J^nM/J^{n + 1}M \to 0$
在与 $R'$ 张量后仍正合。这首先蕴含
$J^nM \otimes_R R' = (J')^nM'$，继而蕴含
$(J')^nM'/(J')^{n + 1}M' = J^nM/J^{n + 1}M \otimes_R R'$。
因此，$\mu$ 是 (\ref{algebra-equation-quasi-regular}) 与
$\text{id}_{R'}$ 的张量积；依假设前者为同构，结论成立。
\end{proof}

\begin{lemma}
\label{algebra-lemma-quasi-regular-sequence-in-neighbourhood}
设 $R$ 为诺特环，$M$ 为有限 $R$-模，$\mathfrak p$ 为素理想。
设 $x_1, \ldots, x_c$ 为 $R$ 中的序列，它在
$R_{\mathfrak p}$ 中的像构成
$M_{\mathfrak p}$-拟正则序列。则存在
$g \in R$，$g \not \in \mathfrak p$，使得
$x_1, \ldots, x_c$ 在 $R_g$ 中的像构成
$M_g$-拟正则序列。
\end{lemma}

\begin{proof}
考虑映射 (\ref{algebra-equation-quasi-regular}) 的核 $K$。
由于 $M/JM \otimes_{R/J} R/J[X_1, \ldots, X_c]$ 是有限
$R/J[X_1, \ldots, X_c]$-模，且
$R/J[X_1, \ldots, X_c]$ 是诺特环，可知 $K$ 也是有限
$R/J[X_1, \ldots, X_c]$-模。选取齐次生成元
$k_1, \ldots, k_t \in K$。依假设，对每个
$i = 1, \ldots, t$，存在 $g_i \in R$，
$g_i \not \in \mathfrak p$，使得 $g_i k_i = 0$。
因此 $g = g_1 \ldots g_t$ 满足要求。
\end{proof}

\begin{lemma}
\label{algebra-lemma-truncate-quasi-regular}
设 $R$ 为环，$M$ 为 $R$-模，并设
$f_1, \ldots, f_c \in R$ 为 $M$-拟正则序列。
对任意 $i$，$\overline{R} = R/(f_1, \ldots, f_i)$ 中的序列
$\overline{f}_{i + 1}, \ldots, \overline{f}_c$
是 $\overline{M} = M/(f_1, \ldots, f_i)M$-拟正则序列。
\end{lemma}

\begin{proof}
只须证明 $i = 1$ 的情形。令
$\overline{J} = (\overline{f}_2, \ldots, \overline{f}_c) \subset \overline{R}$。
则
\begin{align*}
\overline{J}^n\overline{M}/\overline{J}^{n + 1}\overline{M}
& =
(J^nM + f_1M)/(J^{n + 1}M + f_1M) \\
& = J^nM / (J^{n + 1}M + J^nM \cap f_1M).
\end{align*}
因此，为证明本引理，只须说明
$J^{n + 1}M + J^nM \cap f_1M = J^{n + 1}M + f_1J^{n - 1}M$，
因为这将说明
$\bigoplus_{n \geq 0}
\overline{J}^n\overline{M}/\overline{J}^{n + 1}\overline{M}$
是
$\bigoplus_{n \geq 0} J^nM/J^{n + 1}M \cong M/JM[X_1, \ldots, X_c]$
关于 $X_1$ 的商。事实上，有
$J^nM \cap f_1M = f_1J^{n - 1}M$。
确实，如果 $m \not \in J^{n - 1}M$，那么
$f_1m \not \in J^nM$，因为依假设
$\bigoplus J^nM/J^{n + 1}M$ 是多项式代数
$M/J[X_1, \ldots, X_c]$。
\end{proof}

\begin{lemma}
\label{algebra-lemma-quasi-regular-regular}
设 $(R, \mathfrak m)$ 为诺特局部环，$M$ 为非零有限
$R$-模，并设 $f_1, \ldots, f_c \in \mathfrak m$
为 $M$-拟正则序列。则 $f_1, \ldots, f_c$ 是
$M$-正则序列。
\end{lemma}

\begin{proof}
令 $J = (f_1, \ldots, f_c)$。先说明 $f_1$ 是 $M$ 上的非零因子。
假设 $x \in M$ 非零。由 Krull 交定理，存在整数 $r$，
使得 $x \in J^rM$ 但 $x \not \in J^{r + 1}M$；参见引理
\ref{algebra-lemma-intersect-powers-ideal-module-zero}。
于是 $f_1 x \in J^{r + 1}M$，而由所假设的
$\bigoplus J^nM/J^{n + 1}M$ 的结构，它在
$J^{r + 1}M/J^{r + 2}M$ 中的类非零。因此
$f_1x \not = 0$。

\medskip\noindent
现在使用引理
\ref{algebra-lemma-truncate-quasi-regular} 并对 $c$ 作归纳，
即可完成证明。
\end{proof}

\begin{remark}[其他类型的正则序列]
\label{algebra-remark-koszul-regular}
在论文 \cite{Kabele} 中，作者讨论了环 $R$ 的元素序列
$x_1, \ldots, x_r$ 的另外两种正则性条件。具体而言，如果对
$i \geq 1$ 有 $H_i(K_{\bullet}(R, x_{\bullet})) = 0$，
其中 $K_{\bullet}(R, x_{\bullet})$ 是 Koszul 复形，
则称该序列为{\it Koszul-正则的}。如果
$H_1(K_{\bullet}(R, x_{\bullet})) = 0$，则称该序列为
{\it $H_1$-正则的}。有如下蕴含关系：
正则 $\Rightarrow$ Koszul-正则 $\Rightarrow$
$H_1$-正则 $\Rightarrow$ 拟正则。作者通过例子说明，
即使 $R$ 是（非诺特）局部环且该序列生成 $R$ 的极大理想，
这些蕴含一般也不能逆转。我们在《代数进阶》的第
\ref{more-algebra-section-koszul-regular} 节中更详细地引入这些概念。
\end{remark}

\begin{remark}
\label{algebra-remark-join-quasi-regular-sequences}
设 $k$ 为域。考虑环
$$
A = k[x, y, w, z_0, z_1, z_2, \ldots]/
(y^2z_0 - wx, z_0 - yz_1, z_1 - yz_2, \ldots)
$$
在此环中，$x$ 是非零因子，$y$ 在 $A/xA$ 中的像给出拟正则序列。
但 $x, y$ 在 $A$ 中并非拟正则序列，因为
$(x, y)/(x, y)^2$ 不是 $A/(x, y)$ 上秩为二的自由模：
在 $(x, y)/(x, y)^2$ 中有 $wx = 0$，但 $w$ 在
$A/(x, y)$ 中不为零。因此，引理
\ref{algebra-lemma-join-regular-sequences} 的类似断言对拟正则序列不成立。
\end{remark}

\begin{lemma}
\label{algebra-lemma-quasi-regular-on-quotient}
设 $R$ 为环，$J = (f_1, \ldots, f_r)$ 为 $R$ 的理想，
$M$ 为 $R$-模。令
$\overline{R} = R/\bigcap_{n \geq 0} J^n$、
$\overline{M} = M/\bigcap_{n \geq 0} J^nM$，并以
$\overline{f}_i$ 表示 $f_i$ 在 $\overline{R}$ 中的像。
则 $f_1, \ldots, f_r$ 是 $M$-拟正则序列，当且仅当
$\overline{f}_1, \ldots, \overline{f}_r$ 是
$\overline{M}$-拟正则序列。
\end{lemma}

\begin{proof}
这是因为
$J^nM/J^{n + 1}M \cong
\overline{J}^n\overline{M}/\overline{J}^{n + 1}\overline{M}$。
\end{proof}

\section{爆破代数}
\label{algebra-section-blow-up}

\noindent
本节对爆破作一些初等观察。

\begin{definition}
\label{algebra-definition-blow-up}
设 $R$ 为环。
设 $I \subset R$ 为理想。
\begin{enumerate}
\item 与二元组 $(R, I)$ 相关联的{\it 爆破代数}，或称
{\it Rees 代数}，是分次 $R$-代数
$$
\text{Bl}_I(R) =
\bigoplus\nolimits_{n \geq 0} I^n =
R \oplus I \oplus I^2 \oplus \ldots
$$
其中直和项 $I^n$ 置于次数 $n$。
\item 设 $a \in I$。以 $a^{(1)}$ 表示把 $a$ 看作 Rees 代数中
次数为 $1$ 的元素所得的元素。则{\it 仿射爆破代数}
$R[\frac{I}{a}]$ 是第 \ref{algebra-section-proj} 节所构造的代数
$(\text{Bl}_I(R))_{(a^{(1)})}$。
\end{enumerate}
\end{definition}

\noindent
换言之，$R[\frac{I}{a}]$ 的元素可由形如 $x/a^n$ 的表达式表示，
其中 $x \in I^n$。两个代表元 $x/a^n$ 与 $y/a^m$ 定义同一元素，
当且仅当存在某个 $k \geq 0$，使得
$a^k(a^mx - a^ny) = 0$。

\begin{lemma}
\label{algebra-lemma-affine-blowup}
设 $R$ 为环，$I \subset R$ 为理想，且 $a \in I$。
令 $R' = R[\frac{I}{a}]$ 为仿射爆破代数。则
\begin{enumerate}
\item $a$ 在 $R'$ 中的像是非零因子；
\item $IR' = aR'$；并且
\item $(R')_a = R_a$。
\end{enumerate}
\end{lemma}

\begin{proof}
这直接由上面对 $R[\frac{I}{a}]$ 的描述得到。
\end{proof}

\begin{lemma}
\label{algebra-lemma-blowup-base-change}
设 $R \to S$ 为环映射。设 $I \subset R$ 为理想且
$a \in I$。置 $J = IS$，并令 $b \in J$ 为 $a$ 的像。
则 $S[\frac{J}{b}]$ 是 $S \otimes_R R[\frac{I}{a}]$
关于由 $b$ 的某个幂零化的元素所成理想的商。
\end{lemma}

\begin{proof}
令 $S'$ 为 $S \otimes_R R[\frac{I}{a}]$ 关于其
$b$-幂挠元所成理想的商。环映射
$$
S \otimes_R R[\textstyle{\frac{I}{a}}]
\longrightarrow
S[\textstyle{\frac{J}{b}}]
$$
是满射；又因为 $b$ 在 $S[\frac{J}{b}]$ 中是非零因子，
它将 $a$-幂挠元映为零。因此得到满射 $S' \to S[\frac{J}{b}]$。
为说明其核为零，我们构造逆映射。具体地，设
$z = y/b^n$ 是 $S[\frac{J}{b}]$ 的元素，即 $y \in J^n$。
写成 $y = \sum x_is_i$，其中 $x_i \in I^n$ 且 $s_i \in S$。
把 $z$ 映到 $S'$ 中 $\sum s_i \otimes x_i/a^n$ 的类。
这是良定义的，因为映射 $S \otimes_R I^n \to J^n$ 的核中元素
由 $a^n$ 零化，因而在 $S'$ 中映为零。
\end{proof}

\begin{example}
\label{algebra-example-rees-algebra-polynomial}
设 $R$ 为环。设 $P = R[t_1, \ldots, t_n]$ 为多项式代数。
令 $I = (t_1, \ldots, t_n) \subset P$。沿用定义
\ref{algebra-definition-blow-up} 的记号，有同构
$$
P[T_1, \ldots, T_n]/(t_iT_j - t_jT_i) \longrightarrow \text{Bl}_I(P)
$$
把 $T_i$ 送到 $t_i^{(1)}$。由读者验证此映射是良定义的。
由于 $I$ 由 $t_1, \ldots, t_n$ 生成，可知该映射是满射。
为说明该映射是单射，需要证明：对每个 $e \geq 1$，
$P$-模 $I^e$ 由多重指标 $E = (e_1, \ldots, e_n)$ 中
次数 $|E| = e$ 的那些指标所对应的单项式
$t^E = t_1^{e_1} \ldots x_n^{e_n}$ 生成，而且它们只满足关系
$t_i t^E = t_j t^{E'}$，其中 $|E| = |E'| = e$ 且
$e_a + \delta_{a i} = e'_a + \delta_{a j},\ a = 1, \ldots, n$
（Kronecker delta）。细节从略。
\end{example}

\begin{example}
\label{algebra-example-affine-blowup-algebra-polynomial}
设 $R$ 为环。设 $P = R[t_1, \ldots, t_n]$ 为多项式代数。
令 $I = (t_1, \ldots, t_n) \subset P$，并令 $a = t_1$。
沿用定义 \ref{algebra-definition-blow-up} 的记号，有同构
$$
P[x_2, \ldots, x_n]/(t_1x_2 - t_2, \ldots, t_1x_n - t_n)
\longrightarrow
\textstyle{P[\frac{I}{a}] = P[\frac{I}{t_1}]}
$$
把 $x_i$ 送到 $t_i/t_1$。由读者验证此映射是良定义的。
由于 $I$ 由 $t_1, \ldots, t_n$ 生成，可知该映射是满射。
为说明该映射是单射，读者可以论证其源与靶都没有
$t_1$-挠元，并且在逆转 $t_1$ 后该映射成为同构；参见引理
\ref{algebra-lemma-affine-blowup}。或者，读者可使用例
\ref{algebra-example-rees-algebra-polynomial} 中对 Rees 代数的描述。
细节从略。
\end{example}

\begin{lemma}
\label{algebra-lemma-affine-blowup-quotient-description}
设 $R$ 为环，$I = (a_1, \ldots, a_n)$ 为 $R$ 的理想。
令 $a = a_1$。则存在满射
$$
R[x_2, \ldots, x_n]/(a x_2 - a_2, \ldots, a x_n - a_n)
\longrightarrow
\textstyle{R[\frac{I}{a}]}
$$
其核是源中的 $a$-幂挠元之集。
\end{lemma}

\begin{proof}
考虑环映射 $P = \mathbf{Z}[t_1, \ldots, t_n] \to R$，
它把 $t_i$ 送到 $a_i$。置 $J = (t_1, \ldots, t_n)$。
由例 \ref{algebra-example-affine-blowup-algebra-polynomial}，有
$P[\frac{J}{t_1}] =
P[x_2, \ldots, x_n]/(t_1 x_2 - t_2, \ldots, t_1 x_n - t_n)$。
将引理 \ref{algebra-lemma-blowup-base-change} 应用于映射 $P \to A$ 即得结论。
\end{proof}

\begin{lemma}
\label{algebra-lemma-blowup-in-principal}
设 $R$ 为环，$I \subset R$ 为理想，且 $a \in I$。
置 $R' = R[\frac{I}{a}]$。若 $f \in R$ 满足 $V(f) = V(I)$，
则 $f$ 在 $R'$ 中的像是非零因子，并且
$R'_f = R'_a = R_a$。
\end{lemma}

\begin{proof}
下文不再特别说明地使用引理 \ref{algebra-lemma-affine-blowup} 的结论。
假设 $V(f) = V(I)$ 蕴含 $V(fR') = V(IR') = V(aR')$。
因此存在 $b, c \in R'$ 和某些幂，使得
$a^n = fb$ 且 $f^m = ac$。引理随即成立。
\end{proof}

\begin{lemma}
\label{algebra-lemma-blowup-add-principal}
设 $R$ 为环，$I \subset R$ 为理想，$a \in I$，且 $f \in R$。
置 $R' = R[\frac{I}{a}]$ 与 $R'' = R[\frac{fI}{fa}]$。
则存在满的 $R$-代数映射 $R' \to R''$，其核为 $R'$ 中
$f$-幂挠元之集。
\end{lemma}

\begin{proof}
该映射把 $x \in I^n$ 所对应的 $x/a^n$ 送到
$f^nx/(fa)^n$。直接验证可知此映射良定义且为满射。
由于 $af$ 在 $R''$ 中是非零因子
（引理 \ref{algebra-lemma-affine-blowup}），可知 $f$-幂挠元之集映为零。
反之，若 $x \in R'$ 且对所有 $n > 0$ 都有
$f^n x \not = 0$，则由于 $a$ 在 $R'$ 中是非零因子，
对所有 $n$ 都有 $(af)^n x \not = 0$。根据定义
\ref{algebra-definition-blow-up} 之后对 $R''$ 的描述，$x$ 在 $R''$
中的像不为零。
\end{proof}

\begin{lemma}
\label{algebra-lemma-blowup-reduced}
\begin{slogan}
约化性在爆破下不变
\end{slogan}
若 $R$ 约化，则 $R$ 的每个（仿射）爆破代数都是约化的。
\end{lemma}

\begin{proof}
设 $I \subset R$ 为理想且 $a \in I$。假设 $x \in I^n$
所对应的 $x/a^n$ 是 $R[\frac{I}{a}]$ 的幂零元。
则 $(x/a^n)^m = 0$。故对某个 $N \geq 0$，在 $R$ 中有
$a^N x^m = 0$。必要时增大 $N$，可设 $N = me$，其中
$e \geq 0$。于是 $(a^e x)^m = 0$；由于 $R$ 约化，得到
$a^e x = 0$。这说明 $R[\frac{I}{a}]$ 中 $x/a^n = 0$。
\end{proof}

\begin{lemma}
\label{algebra-lemma-blowup-domain}
设 $R$ 为整环，$I \subset R$ 为理想，且 $a \in I$ 为非零元。
则仿射爆破代数 $R[\frac{I}{a}]$ 是整环。
\end{lemma}

\begin{proof}
假设 $x \in I^n$、$y \in I^m$ 所对应的
$x/a^n$、$y/a^m$ 是 $R[\frac{I}{a}]$ 中乘积为零的元素。
则在 $R$ 中有 $a^N x y = 0$。由于 $R$ 是整环，可知
$x = 0$ 或 $y = 0$。
\end{proof}

\begin{lemma}
\label{algebra-lemma-blowup-dominant}
设 $R$ 为环，$I \subset R$ 为理想，且 $a \in I$。
若 $a$ 不属于 $R$ 的任何极小素理想，则
$\Spec(R[\frac{I}{a}]) \to \Spec(R)$ 的像稠密。
\end{lemma}

\begin{proof}
若对 $x \in R$ 有 $a^k x = 0$，则 $x$ 属于 $R$ 的所有
极小素理想，因而是幂零元；参见引理
\ref{algebra-lemma-Zariski-topology}。所以 $R \to R[\frac{I}{a}]$
的核由幂零元组成。于是结论由引理
\ref{algebra-lemma-image-dense-generic-points} 得到。
\end{proof}

\begin{lemma}
\label{algebra-lemma-valuation-ring-colimit-affine-blowups}
设 $(R, \mathfrak m)$ 为分式域是 $K$ 的局部整环。
设 $R \subset A \subset K$ 为支配 $R$ 的赋值环。则
$$
A = \colim R[\textstyle{\frac{I}{a}}]
$$
是仿射爆破 $R \to R[\frac{I}{a}]$ 的有向余极限，
其中各仿射爆破满足下列性质：
\begin{enumerate}
\item $a \in I \subset \mathfrak m$；
\item $I$ 有限生成；并且
\item $R \to R[\frac{I}{a}]$ 在 $\mathfrak m$ 处的纤维环非零。
\end{enumerate}
\end{lemma}

\begin{proof}
由引理 \ref{algebra-lemma-blowup-domain}，任意爆破代数
$R[\frac{I}{a}]$ 都是包含于 $K$ 的整环。本引理就是说，
$A$ 是其中满足 $a \in I$ 及性质 (1)、(2)、(3) 的那些环的有向并。
若 $R[\frac{I}{a}] \subset A$ 且 $R[\frac{J}{b}] \subset A$，则
$$
R[\textstyle{\frac{I}{a}}] \cup R[\textstyle{\frac{J}{b}}]  \subset
R[\textstyle{\frac{IJ}{ab}}] \subset A
$$
第一个包含成立，是因为 $x/a^n = b^nx/(ab)^n$；第二个包含成立，
是因为若 $z \in (IJ)^n$，则可写成 $z = \sum x_iy_i$，其中
$x_i \in I^n$ 且 $y_i \in J^n$，从而
$z/(ab)^n = \sum (x_i/a^n)(y_i/b^n)$ 属于 $A$。

\medskip\noindent
考虑有限子集 $E \subset A$。设 $E = \{e_1, \ldots, e_n\}$。
选取非零的 $a \in R$，使得对所有 $i = 1, \ldots, n$，
均可写成 $e_i = f_i/a$。置 $I = (f_1, \ldots, f_n, a)$。
我们断言 $R[\frac{I}{a}] \subset A$。这是显然的，因为
$R[\frac{I}{a}]$ 的元素可表示为 $e_i$ 的多项式。
本引理立即由此观察得到。
\end{proof}

\section{Ext 群}
\label{algebra-section-ext}

\noindent
本节将使用一点同调代数，以便建立诺特局部环上深度的一些基本性质。

\begin{lemma}
\label{algebra-lemma-resolution-by-finite-free}
设 $R$ 为环，$M$ 为 $R$-模。
\begin{enumerate}
\item 存在正合复形
$$
\ldots \to F_2 \to F_1 \to F_0 \to M \to 0.
$$
其中 $F_i$ 都是自由 $R$-模。
\item 若 $R$ 是诺特环且 $M$ 是有限 $R$-模，则可以选取这个复形，
使得 $F_i$ 都是有限自由模。换言之，可以找到正合复形
$$
\ldots \to R^{\oplus n_2} \to R^{\oplus n_1} \to R^{\oplus n_0} \to M \to 0.
$$
\end{enumerate}
\end{lemma}

\begin{proof}
只说明诺特情形。第一步选取满射 $R^{n_0} \to M$。
若已构造出长度为 $e$ 的正合复形，只须选取满射
$R^{n_{e + 1}} \to
\Ker(R^{n_e} \to R^{n_{e-1}})$；由于 $R$ 是诺特环，这是可行的。
\end{proof}

\begin{definition}
\label{algebra-definition-finite-free-resolution}
设 $R$ 为环，$M$ 为 $R$-模。
\begin{enumerate}
\item $M$ 的一个（左）{\it 消解} $F_\bullet \to M$ 是 $R$-模的正合复形
$$
\ldots \to F_2 \to F_1 \to F_0 \to M \to 0
$$
\item $M$ 的一个{\it 自由 $R$-模消解}是消解
$F_\bullet \to M$，其中每个 $F_i$ 都是自由 $R$-模。
\item $M$ 的一个{\it 有限自由 $R$-模消解}是消解
$F_\bullet \to M$，其中每个 $F_i$ 都是有限自由 $R$-模。
\end{enumerate}
\end{definition}

\noindent
我们常用记号 $F_{\bullet}$ 表示形如
$$
\ldots \to F_i \to F_{i-1} \to \ldots
$$
的 $R$-模复形。此时常用 $d_i$ 或 $d_{F, i}$ 表示映射
$F_i \to F_{i-1}$。本节始终假设 $F_0$ 是复形中最后一个非零项。
复形 $F_{\bullet}$ 的{\it 第 $i$ 个同调群}是群
$H_i = \Ker(d_{F, i})/\Im(d_{F, i + 1})$。
一个{\it 复形映射 $\alpha : F_{\bullet} \to G_{\bullet}$}
由映射 $\alpha_i : F_i \to G_i$ 给出，并且满足
$\alpha_{i-1} \circ d_{F, i} = d_{G, i-1} \circ \alpha_i$。
这样的映射在同调上诱导映射 $H_i(\alpha) :
H_i(F_{\bullet}) \to H_i(G_{\bullet})$。若 $\alpha, \beta
:  F_{\bullet} \to G_{\bullet}$ 是复形映射，则 $\alpha$ 与
$\beta$ 之间的一个{\it 同伦}由一族映射
$h_i : F_i \to G_{i + 1}$ 给出，并且满足
$\alpha_i - \beta_i = d_{G, i + 1} \circ h_i +
h_{i-1} \circ d_{F, i}$。若 $\alpha$ 与 $\beta$ 之间存在同伦，
则称两个映射 $\alpha, \beta : F_{\bullet} \to G_{\bullet}$
{\it 同伦}。

\medskip\noindent
对于形如 $F^{\bullet}$ 的复形，我们将使用非常类似的记号；它形如
$$
\ldots \to F^i \xrightarrow{d^i} F^{i + 1} \to \ldots
$$
同样有复形映射、同伦等概念。在这种情形下，置
$H^i(F^{\bullet}) =
\Ker(d^i)/\Im(d^{i - 1})$，并称之为{\it 第 $i$ 个上同调群}。

\begin{lemma}
\label{algebra-lemma-homotopic-equal-homology}
任意两个同伦的复形映射在（上）同调群上诱导相同的映射。
\end{lemma}

\begin{proof}
从略。
\end{proof}

\begin{lemma}
\label{algebra-lemma-compare-resolutions}
设 $R$ 为环，$M \to N$ 为 $R$-模映射。
设 $N_\bullet \to N$ 为任意消解。设
$$
\ldots \to F_2 \to F_1 \to F_0 \to M
$$
为 $R$-模复形，其中每个 $F_i$ 都是自由 $R$-模。则
\begin{enumerate}
\item 存在复形映射 $F_\bullet \to N_\bullet$，使得
$$
\xymatrix{
F_0 \ar[r] \ar[d] & M \ar[d] \\
N_0 \ar[r] & N
}
$$
交换；并且
\item 任意两个如 (1) 所述的映射
$\alpha, \beta : F_\bullet \to N_\bullet$ 都同伦。
\end{enumerate}
\end{lemma}

\begin{proof}
证明 (1)。由于 $F_0$ 自由，可以找到提升映射
$F_0 \to M \to N$ 的映射 $F_0 \to N_0$。于是得到诱导映射
$F_1 \to F_0 \to N_0$，其像落在 $N_1 \to N_0$ 的像中。
由于 $F_1$ 自由，可以把它提升为映射 $F_1 \to N_1$。
这又诱导映射 $F_2 \to F_1 \to N_1$，它在 $N_0$ 中映为零。
由于 $N_\bullet$ 正合，该映射的像包含于 $N_2 \to N_1$ 的像中。
所以可以提升得到映射 $F_2 \to N_2$。如此重复。

\medskip\noindent
证明 (2)。为说明 $\alpha, \beta$ 同伦，只须说明差
$\gamma = \alpha - \beta$ 与零同伦。注意到
$\gamma_0 : F_0 \to N_0$ 的像包含于 $N_1 \to N_0$ 的像中。
因此可将 $\gamma_0$ 提升为映射 $h_0 : F_0 \to N_1$。
考虑映射 $\gamma_1' = \gamma_1 - h_0 \circ d_{F, 1}$。
由 $h_0$ 的选取，$\gamma_1'$ 的像包含于
$N_1 \to N_0$ 的核中。由于 $N_\bullet$ 正合，可以把
$\gamma_1'$ 提升为映射 $h_1 : F_1 \to N_2$。
至此有 $\gamma_1 = h_0 \circ d_{F, 1}
+ d_{N, 2} \circ h_1$。如此重复。
\end{proof}

\noindent
现在可以定义群 $\Ext^i_R(M, N)$。具体而言，选取 $M$ 的一个
自由 $R$-模消解 $F_{\bullet}$；参见引理
\ref{algebra-lemma-resolution-by-finite-free}。考虑（上同调）复形
$$
\Hom_R(F_\bullet, N) :
\Hom_R(F_0, N) \to
\Hom_R(F_1, N) \to
\Hom_R(F_2, N) \to \ldots
$$
对 $i \geq 0$，定义 $\Ext^i_R(M, N)$ 为该复形的第 $i$ 个
上同调群\footnote{此时也许说“一个”而不是“这个”Ext 群更为恰当。}。
对 $i < 0$，置 $\Ext^i_R(M, N) = 0$。在继续之前，先指出
$$
\Ext^0_R(M, N) = \Ker(\Hom_R(F_0, N) \to \Hom_R(F_1, N)) =
\Hom_R(M, N)
$$
因为可以把引理 \ref{algebra-lemma-hom-exact} 的 (1) 应用于正合序列
$F_1 \to F_0 \to M \to 0$。下一个引理说明此定义在何种意义下
是良定义的。

\begin{lemma}
\label{algebra-lemma-ext-welldefined}
设 $R$ 为环，$M_1, M_2, N$ 为 $R$-模。
假设 $F_{\bullet}$ 是模 $M_1$ 的自由消解，
$G_{\bullet}$ 是模 $M_2$ 的自由消解。
设 $\varphi : M_1 \to M_2$ 为模映射。
设 $\alpha : F_{\bullet} \to G_{\bullet}$ 为诱导
$M_1 = \Coker(d_{F, 1}) \to M_2 = \Coker(d_{G, 1})$ 上
$\varphi$ 的复形映射；参见引理 \ref{algebra-lemma-compare-resolutions}。
则诱导映射
$$
H^i(\alpha) :
H^i(\Hom_R(F_{\bullet}, N))
\longrightarrow
H^i(\Hom_R(G_{\bullet}, N))
$$
与 $\alpha$ 的选择无关。若 $\varphi$ 是同构，则所有映射
$H^i(\alpha)$ 也都是同构。若 $M_1 = M_2$、
$F_\bullet = G_\bullet$，且 $\varphi$ 是恒等映射，
则所有映射 $H_i(\alpha)$ 也都是恒等映射。
\end{lemma}

\begin{proof}
由引理 \ref{algebra-lemma-compare-resolutions}，另一个诱导 $\varphi$ 的映射
$\beta : F_{\bullet} \to G_{\bullet}$ 与 $\alpha$ 同伦。
因此映射 $\Hom_R(F_\bullet, N) \to
\Hom_R(G_\bullet, N)$ 同伦。于是独立性结论由引理
\ref{algebra-lemma-homotopic-equal-homology} 得到。

\medskip\noindent
假设 $\varphi$ 是同构。设 $\psi : M_2 \to M_1$ 为其逆。
选取诱导
$\psi : M_2 = \Coker(d_{G, 1}) \to M_1 = \Coker(d_{F, 1})$
的映射 $\beta : G_{\bullet} \to F_{\bullet}$；参见引理
\ref{algebra-lemma-compare-resolutions}。好，现在考虑映射
$H^i(\alpha) \circ H^i(\beta) =
H^i(\alpha \circ \beta)$。由上文，映射
$H^i(\alpha \circ \beta)$ 与映射
$H^i(\text{id}_{G_{\bullet}}) = \text{id}$ {\it 相同}。
复合 $H^i(\beta) \circ H^i(\alpha)$ 的情形也类似。
因此 $H^i(\alpha)$ 与 $H^i(\beta)$ 互为逆映射。
\end{proof}

\begin{lemma}
\label{algebra-lemma-long-exact-seq-ext}
设 $R$ 为环，$M$ 为 $R$-模。
设 $0 \to N' \to N \to N'' \to 0$ 为短正合序列。
则得到长正合序列
$$
\begin{matrix}
0
\to \Hom_R(M, N')
\to \Hom_R(M, N)
\to \Hom_R(M, N'')
\\
\phantom{0\ }
\to \Ext^1_R(M, N')
\to \Ext^1_R(M, N)
\to \Ext^1_R(M, N'')
\to \ldots
\end{matrix}
$$
\end{lemma}

\begin{proof}
选取自由消解 $F_{\bullet} \to M$。由于每个 $F_i$ 都自由，
得到复形的短正合序列
$$
0 \to
\Hom_R(F_{\bullet}, N') \to
\Hom_R(F_{\bullet}, N) \to
\Hom_R(F_{\bullet}, N'') \to
0
$$
将蛇引理应用于此序列，即得上述长正合序列。
\end{proof}

\begin{lemma}
\label{algebra-lemma-reverse-long-exact-seq-ext}
设 $R$ 为环，$N$ 为 $R$-模。
设 $0 \to M' \to M \to M'' \to 0$ 为短正合序列。
则得到长正合序列
$$
\begin{matrix}
0
\to \Hom_R(M'', N)
\to \Hom_R(M, N)
\to \Hom_R(M', N)
\\
\phantom{0\ }
\to \Ext^1_R(M'', N)
\to \Ext^1_R(M, N)
\to \Ext^1_R(M', N)
\to \ldots
\end{matrix}
$$
\end{lemma}

\begin{proof}
分别选取 $M'$ 与 $M''$ 的生成元集
$\{m'_{i'}\}_{i' \in I'}$ 和 $\{m''_{i''}\}_{i'' \in I''}$。
对每个 $i'' \in I''$，选取元素 $m''_{i''} \in M''$ 的一个提升
$\tilde m''_{i''} \in M$。置 $F' = \bigoplus_{i' \in I'} R$、
$F'' = \bigoplus_{i'' \in I''} R$ 以及 $F = F' \oplus F''$。
把这些自由模的生成元映到相应选定的生成元，得到满的
$R$-模映射 $F' \to M'$、$F'' \to M''$ 和 $F \to M$。
于是得到短正合序列之间的映射
$$
\begin{matrix}
0 & \to & M' & \to & M & \to & M'' & \to & 0 \\
& & \uparrow & & \uparrow & & \uparrow \\
0 & \to & F' & \to & F & \to & F'' & \to & 0 \\
\end{matrix}
$$
由蛇引理，核的序列 $0 \to K' \to K \to K'' \to 0$
是 $R$-模的短正合序列。因此可以无限地继续这一过程。
换言之，我们得到消解的短正合序列，它嵌入图表
$$
\begin{matrix}
0 & \to & M' & \to & M & \to & M'' & \to & 0 \\
& & \uparrow & & \uparrow & & \uparrow \\
0 & \to & F_\bullet' & \to & F_\bullet & \to & F_\bullet'' & \to & 0 \\
\end{matrix}
$$
由于每个序列 $0 \to F'_n \to F_n \to F''_n \to 0$
依构造都是分裂正合的，应用函子 $\Hom_R(-, N)$ 后，
得到复形的短正合序列
$$
0 \to
\Hom_R(F''_{\bullet}, N) \to
\Hom_R(F_{\bullet}, N) \to
\Hom_R(F'_{\bullet}, N) \to
0
$$
将蛇引理应用于此序列，即得上述长正合序列。
\end{proof}

\begin{lemma}
\label{algebra-lemma-annihilate-ext}
设 $R$ 为环，$M$, $N$ 为 $R$-模。
若 $x\in R$ 满足 $xN = 0$ 或 $xM = 0$，
则该元素零化每个模 $\Ext^i_R(M, N)$。
\end{lemma}

\begin{proof}
选取 $M$ 的自由消解 $F_{\bullet}$。由于 $\Ext^i_R(M, N)$
定义为复形 $\Hom_R(F_{\bullet}, N)$ 的上同调，
当 $xN = 0$ 时，引理显然成立。若 $xM = 0$，则
$F_{\bullet}$ 上的乘以 $x$ 提升了 $M$ 上的零映射。
故由引理 \ref{algebra-lemma-ext-welldefined}，它在 Ext 群上诱导的映射
与零映射诱导的映射相同。
\end{proof}

\begin{lemma}
\label{algebra-lemma-ext-noetherian}
设 $R$ 为诺特环，$M$, $N$ 为有限 $R$-模。
则对所有 $i$，$\Ext^i_R(M, N)$ 都是有限 $R$-模。
\end{lemma}

\begin{proof}
这是因为 $\Ext^i_R(M, N)$ 可计算为复形
$\Hom_R(F_\bullet, N)$ 的上同调群，其中每个 $F_n$
都是有限自由 $R$-模；参见引理
\ref{algebra-lemma-resolution-by-finite-free}。
\end{proof}

\section{深度}
\label{algebra-section-depth}

\noindent
定义如下。

\begin{definition}
\label{algebra-definition-depth}
设 $R$ 为环，$I \subset R$ 为理想，$M$ 为有限 $R$-模。
$M$ 的{\it $I$-深度}，记作 $\text{depth}_I(M)$，定义如下：
\begin{enumerate}
\item 若 $IM \not = M$，则 $\text{depth}_I(M)$ 是
$\{0, 1, 2, \ldots, \infty\}$ 中 $I$ 内 $M$-正则序列之长度的上确界；
\item 若 $IM = M$，则置 $\text{depth}_I(M) = \infty$。
\end{enumerate}
若 $(R, \mathfrak m)$ 是局部环，则简称
$\text{depth}_{\mathfrak m}(M)$ 为 $M$ 的{\it 深度}。
\end{definition}

\noindent
说明。依定义 \ref{algebra-definition-regular-sequence}，空序列不是零模上的
正则序列；但就实际应用而言，把 $0$ 模的深度置为 $+\infty$
是方便的。注意，若 $I = R$，则对所有有限 $R$-模 $M$，
都有 $\text{depth}_I(M) = \infty$。若 $I$ 包含于 $R$ 的
Jacobson 根中（例如，若 $R$ 是局部环且
$I \subset \mathfrak m_R$），则由 Nakayama 引理，
$M \not = 0 \Rightarrow IM \not = M$。
模 $M$ 的 $I$-深度为 $0$，当且仅当 $M$ 非零且
$I$ 不含 $M$ 上的非零因子。

\medskip\noindent
例 \ref{algebra-example-global-regular} 表明，即使环是诺特的，深度的表现
也不良好；例 \ref{algebra-example-local-regular} 表明，当环是局部但非诺特时，
其表现也不良好。我们将看到，对诺特局部环，深度的表现良好。

\begin{lemma}
\label{algebra-lemma-depth-weak-sequence}
设 $R$ 为环，$I \subset R$ 为理想，$M$ 为有限 $R$-模。
则 $\text{depth}_I(M)$ 等于满足下列条件的序列
$f_1, \ldots, f_r \in I$ 之长度的上确界：$f_i$ 是
$M/(f_1, \ldots, f_{i - 1})M$ 上的非零因子。
\end{lemma}

\begin{proof}
假设 $IM = M$。则引理 \ref{algebra-lemma-NAK} 表明，存在
$f \in I$，使得 $f : M \to M$ 是 $\text{id}_M$。
所以 $f, 0, 0, 0, \ldots$ 是由相继非零因子组成的无限序列，
由此可见此情形下二者相等。若 $IM \not =  M$，则本引理所述的
序列是 $M$-正则序列，故此情形下二者也相等。
\end{proof}

\begin{lemma}
\label{algebra-lemma-bound-depth}
设 $(R, \mathfrak m)$ 为诺特局部环，$M$ 为非零有限 $R$-模。
则 $\dim(\text{Supp}(M)) \geq \text{depth}(M)$。
\end{lemma}

\begin{proof}
对 $\dim(\text{Supp}(M))$ 作归纳。若
$\dim(\text{Supp}(M)) = 0$，则
$\text{Supp}(M) = \{\mathfrak m\}$，从而
$\text{Ass}(M) = \{\mathfrak m\}$（由引理
\ref{algebra-lemma-ass-support} 与 \ref{algebra-lemma-ass-zero}），
所以例如由引理 \ref{algebra-lemma-ideal-nonzerodivisor}，$M$ 的深度为零。
在归纳步骤中，假设 $\dim(\text{Supp}(M)) > 0$。
设 $f_1, \ldots, f_d$ 为 $\mathfrak m$ 中的元素序列，
且 $f_i$ 是 $M/(f_1, \ldots, f_{i - 1})M$ 上的非零因子。
依引理 \ref{algebra-lemma-depth-weak-sequence}，只须证明
$\dim(\text{Supp}(M)) \geq d$。不妨设 $d > 0$，否则引理显然成立。
由引理 \ref{algebra-lemma-one-equation-module}，有
$\dim(\text{Supp}(M/f_1M)) = \dim(\text{Supp}(M)) - 1$。
由归纳假设得到 $\dim(\text{Supp}(M/f_1M)) \geq d - 1$，正合所需。
\end{proof}

\begin{lemma}
\label{algebra-lemma-depth-finite-noetherian}
设 $R$ 为诺特环，$I \subset R$ 为理想，$M$ 为非零有限
$R$-模，且 $IM \not = M$。则
$\text{depth}_I(M) < \infty$。
\end{lemma}

\begin{proof}
由于 $M/IM$ 非零，由引理 \ref{algebra-lemma-support-zero}，
可以选取 $\mathfrak p \in \text{Supp}(M/IM)$。
于是 $(M/IM)_\mathfrak p \not = 0$，这蕴含
$I \subset \mathfrak p$；又因为局部化正合，还蕴含
$M_\mathfrak p \not = IM_\mathfrak p$。
设 $f_1, \ldots, f_r \in I$ 为 $M$-正则序列。
由于 $(f_1, \ldots, f_r) \subset I$，
$M_\mathfrak p/(f_1, \ldots, f_r)M_\mathfrak p$ 非零。
由于局部化平坦，$f_1, \ldots, f_r$ 的像在
$I_\mathfrak p$ 中构成 $M_\mathfrak p$-正则序列。
这对 $I$ 中每个 $M$-正则序列都成立，故
$\text{depth}_I(M) \leq \text{depth}_{I_\mathfrak p}(M_\mathfrak p)$。
后者的值 $\leq \text{depth}(M_\mathfrak p)$，而由引理
\ref{algebra-lemma-bound-depth}，该值 $< \infty$。
\end{proof}

\begin{lemma}
\label{algebra-lemma-depth-ext}
设 $R$ 为极大理想是 $\mathfrak m$ 的诺特局部环，
$M$ 为非零有限 $R$-模。则 $\text{depth}(M)$ 等于使得
$\Ext^i_R(R/\mathfrak m, M)$ 非零的最小整数 $i$。
\end{lemma}

\begin{proof}
以 $\delta(M)$ 表示 $M$ 的深度，并以 $i(M)$ 表示使得
$\Ext^i_R(R/\mathfrak m, M)$ 非零的最小整数 $i$。
马上会看到 $i(M) < \infty$。由引理
\ref{algebra-lemma-ideal-nonzerodivisor}，有 $\delta(M) = 0$
当且仅当 $i(M) = 0$，因为
$\mathfrak m \in \text{Ass}(M)$ 恰好意味着 $i(M) = 0$。
因此，若 $\delta(M)$ 或 $i(M)$ 的值 $> 0$，则可以选取
$x \in \mathfrak m$，使得 (a) $x$ 是 $M$ 上的非零因子，
且 (b) $\text{depth}(M/xM) = \delta(M) - 1$。
考虑由引理 \ref{algebra-lemma-long-exact-seq-ext} 从短正合序列
$0 \to M \to M \to M/xM \to 0$ 得到的 Ext 群长正合序列：
$$
\begin{matrix}
0
\to \Hom_R(\kappa, M)
\to \Hom_R(\kappa, M)
\to \Hom_R(\kappa, M/xM)
\\
\phantom{0\ }
\to \Ext^1_R(\kappa, M)
\to \Ext^1_R(\kappa, M)
\to \Ext^1_R(\kappa, M/xM)
\to \ldots
\end{matrix}
$$
由于 $x \in \mathfrak m$，所有映射 $\Ext^i_R(\kappa, M)
\to \Ext^i_R(\kappa, M)$ 都为零；参见引理
\ref{algebra-lemma-annihilate-ext}。因此显然有
$i(M/xM) = i(M) - 1$。对 $\delta(M)$ 归纳即完成证明。
\end{proof}

\begin{lemma}
\label{algebra-lemma-depth-in-ses}
设 $R$ 为诺特局部环，$0 \to N' \to N \to N'' \to 0$
为非零有限 $R$-模的短正合序列。
\begin{enumerate}
\item
$\text{depth}(N) \geq \min\{\text{depth}(N'), \text{depth}(N'')\}$
\item
$\text{depth}(N'') \geq \min\{\text{depth}(N), \text{depth}(N') - 1\}$
\item
$\text{depth}(N') \geq \min\{\text{depth}(N), \text{depth}(N'') + 1\}$
\end{enumerate}
\end{lemma}

\begin{proof}
使用引理 \ref{algebra-lemma-depth-ext} 中由 Ext 群
$\Ext^i(\kappa, N)$ 给出的深度刻画，并使用引理
\ref{algebra-lemma-long-exact-seq-ext} 给出的上同调长正合序列
$$
\begin{matrix}
0
\to \Hom_R(\kappa, N')
\to \Hom_R(\kappa, N)
\to \Hom_R(\kappa, N'')
\\
\phantom{0\ }
\to \Ext^1_R(\kappa, N')
\to \Ext^1_R(\kappa, N)
\to \Ext^1_R(\kappa, N'')
\to \ldots
\end{matrix}
$$
\end{proof}

\begin{lemma}
\label{algebra-lemma-depth-drops-by-one}
设 $R$ 为诺特局部环，$M$ 为非零有限 $R$-模。
\begin{enumerate}
\item 若 $x \in \mathfrak m$ 是 $M$ 上的非零因子，则
$\text{depth}(M/xM) = \text{depth}(M) - 1$。
\item 任意 $M$-正则序列 $x_1, \ldots, x_r$ 都可以扩充为
长度为 $\text{depth}(M)$ 的 $M$-正则序列。
\end{enumerate}
\end{lemma}

\begin{proof}
(2) 是 (1) 的形式推论。设 $x \in R$ 如 (1) 所述。
由短正合序列 $0 \to M \to M \to M/xM \to 0$
与引理 \ref{algebra-lemma-depth-in-ses}，可知深度至多下降 1。
另一方面，若 $x_1, \ldots, x_r \in \mathfrak m$
是 $M/xM$ 的正则序列，则 $x, x_1, \ldots, x_r$
是 $M$ 的正则序列。因此深度至少下降 1。
\end{proof}

\begin{lemma}
\label{algebra-lemma-inherit-minimal-primes}
设 $(R, \mathfrak m)$ 为诺特局部环，$M$ 为有限 $R$-模。
设 $x \in \mathfrak m$、$\mathfrak p \in \text{Ass}(M)$，并设
$\mathfrak q$ 是 $\mathfrak p + (x)$ 上的极小素理想。
则对某个 $n \geq 1$，有 $\mathfrak q \in \text{Ass}(M/x^nM)$。
\end{lemma}

\begin{proof}
选取子模 $N \subset M$，使得 $N \cong R/\mathfrak p$。
由 Artin--Rees 引理（引理 \ref{algebra-lemma-Artin-Rees}），
可以选取 $n > 0$，使得 $N \cap x^nM \subset xN$。
令 $\overline{N} \subset M/x^nM$ 为 $N \to M \to M/x^nM$ 的像。
由引理 \ref{algebra-lemma-ass}，只须证明
$\mathfrak q \in \text{Ass}(\overline{N})$。
根据 $n$ 的选择，存在满射
$\overline{N} \to N/xN = R/\mathfrak p + (x)$，
所以 $\mathfrak q$ 属于 $\overline{N}$ 的支撑。
由于 $\overline{N}$ 由 $x^n$ 和 $\mathfrak p$ 零化，
可知 $\mathfrak q$ 在 $\overline{N}$ 的支撑中的素理想里极小。
因此由引理 \ref{algebra-lemma-ass-minimal-prime-support}，
$\mathfrak q$ 是 $\overline{N}$ 的相伴素理想。
\end{proof}

\begin{lemma}
\label{algebra-lemma-depth-dim-associated-primes}
设 $(R, \mathfrak m)$ 为诺特局部环，$M$ 为有限 $R$-模。
对 $\mathfrak p \in \text{Ass}(M)$，有
$\dim(R/\mathfrak p) \geq \text{depth}(M)$。
\end{lemma}

\begin{proof}
若 $\mathfrak m \in \text{Ass}(M)$，则存在非零元
$x \in M$，它由 $\mathfrak m$ 的所有元素零化。
于是 $\text{depth}(M) = 0$。特别地，此时引理成立。

\medskip\noindent
若 $\text{depth}(M) = 1$，则由上一段可知
$\mathfrak m \not \in \text{Ass}(M)$。故对所有
$\mathfrak p \in \text{Ass}(M)$，有
$\dim(R/\mathfrak p) \geq 1$，此时引理也成立。

\medskip\noindent
一般情形下，对 $\text{depth}(M)$ 作归纳；可以且确实假设其值
$> 1$。选取 $M$ 上的非零因子 $x \in \mathfrak m$。
注意 $x \not \in \mathfrak p$（引理
\ref{algebra-lemma-ass-zero-divisors}）。由引理
\ref{algebra-lemma-one-equation}，有
$\dim(R/\mathfrak p + (x)) = \dim(R/\mathfrak p) - 1$。
因此存在 $\mathfrak p + (x)$ 上的极小素理想 $\mathfrak q$，使得
$\dim(R/\mathfrak q) = \dim(R/\mathfrak p) - 1$
（略去一个小论证；提示：诺特局部环 $A$ 的维数等于
$A/\mathfrak r$ 的维数在 $A$ 的极小素理想 $\mathfrak r$
上所取的最大值）。依引理 \ref{algebra-lemma-inherit-minimal-primes}
选取 $n$，使得 $\mathfrak q$ 是 $M/x^nM$ 的相伴素理想。
可以把归纳假设应用于 $M/x^nM$ 与 $\mathfrak q$，因为由引理
\ref{algebra-lemma-depth-drops-by-one}，
$\text{depth}(M/x^nM) = \text{depth}(M) - 1$。
于是得到 $\dim(R/\mathfrak q) \geq \text{depth}(M/x^nM)$，
结论得证。
\end{proof}

\begin{lemma}
\label{algebra-lemma-depth-localization}
设 $R$ 为诺特局部环，$M$ 为有限 $R$-模。
对素理想 $\mathfrak p \subset R$，有
$\text{depth}(M_\mathfrak p) + \dim(R/\mathfrak p) \geq \text{depth}(M)$。
\end{lemma}

\begin{proof}
若 $M_\mathfrak p = 0$，则
$\text{depth}(M_\mathfrak p) = \infty$，引理成立。
若 $\text{depth}(M) \leq \dim(R/\mathfrak p)$，则引理成立。
若 $\text{depth}(M) > \dim(R/\mathfrak p)$，则由引理
\ref{algebra-lemma-depth-dim-associated-primes}，$\mathfrak p$
不包含于 $M$ 的任何相伴素理想 $\mathfrak q$。
因此由引理 \ref{algebra-lemma-silly} 与
\ref{algebra-lemma-finite-ass}，可以找到 $x \in \mathfrak p$，
它不属于 $M$ 的任何相伴素理想。于是 $x$ 是 $M$ 上的非零因子；
参见引理 \ref{algebra-lemma-ass-zero-divisors}。
故 $\text{depth}(M/xM) = \text{depth}(M) - 1$，并且只要
$M_\mathfrak p$ 非零，就有
$\text{depth}(M_\mathfrak p / x M_\mathfrak p) =
\text{depth}(M_\mathfrak p) - 1$；参见引理
\ref{algebra-lemma-depth-drops-by-one}。因此，对
$\text{depth}(M)$ 作归纳即得结论。
\end{proof}

\begin{lemma}
\label{algebra-lemma-depth-goes-down-finite}
设 $(R, \mathfrak m)$ 为诺特局部环，$R \to S$ 为有限环映射。
设 $\mathfrak m_1, \ldots, \mathfrak m_n$ 为 $S$ 的极大理想，
$N$ 为有限 $S$-模。则
$$
\min\nolimits_{i = 1, \ldots, n} \text{depth}(N_{\mathfrak m_i}) =
\text{depth}_\mathfrak m(N)
$$
\end{lemma}

\begin{proof}
由引理 \ref{algebra-lemma-integral-no-inclusion}、
\ref{algebra-lemma-integral-going-up} 和
\ref{algebra-lemma-finite-finite-fibres}，$S$ 的极大理想恰为
$S$ 中位于 $\mathfrak m$ 上方的素理想，并且它们只有有限多个。
因此本引理的陈述有意义。对
$k = \min\nolimits_{i = 1, \ldots, n} \text{depth}(N_{\mathfrak m_i})$
作归纳。若 $k = 0$，则对某个 $i$ 有
$\text{depth}(N_{\mathfrak m_i}) = 0$。由引理
\ref{algebra-lemma-depth-ext}，这意味着
$\mathfrak m_i S_{\mathfrak m_i}$ 是
$N_{\mathfrak m_i}$ 的相伴素理想，因而
$\mathfrak m_i$ 是 $N$ 的相伴素理想（引理
\ref{algebra-lemma-localize-ass}）。由引理
\ref{algebra-lemma-ass-functorial-Noetherian}，把 $N$ 看作 $R$-模时，
$\mathfrak m$ 是它的相伴素理想。因此
$\text{depth}_\mathfrak m(N) = 0$。这就证明了初始情形。
若 $k > 0$，则 $\mathfrak m_i \not \in \text{Ass}_S(N)$。
再次由引理 \ref{algebra-lemma-ass-functorial-Noetherian}，
$\mathfrak m \not \in \text{Ass}_R(N)$。
因此可以找到 $f \in \mathfrak m$，它不是 $N$ 上的零因子；
参见引理 \ref{algebra-lemma-ideal-nonzerodivisor}。由引理
\ref{algebra-lemma-depth-drops-by-one}，从 $N$ 过渡到 $N/fN$ 时，
所有深度都恰好下降 $1$，其余由归纳假设得到。
\end{proof}

\section{Ext 的函子性}
\label{algebra-section-functoriality-ext}

\noindent
本节简要讨论 $\Ext$ 关于环变换等操作的函子性。
下面列出有待完成的事项。
\begin{enumerate}
\item 给定 $R \to R'$、一个 $R$-模 $M$ 和一个 $R'$-模 $N'$，
$R$-模 $\Ext^i_R(M, N')$ 具有自然的 $R'$-模结构。
此外，存在典范的 $R'$-线性映射
$\Ext^i_{R'}(M \otimes_R R', N') \to
\Ext^i_R(M, N')$。
\item 给定 $R \to R'$ 和 $R$-模 $M$, $N$，存在自然的
$R$-模映射
$\Ext^i_R(M, N) \to \text{Ext}^i_R(M, N \otimes_R R')$。
\end{enumerate}

\begin{lemma}
\label{algebra-lemma-flat-base-change-ext}
给定平坦环映射 $R \to R'$、一个 $R$-模 $M$ 和一个
$R'$-模 $N'$，自然映射
$$
\Ext^i_{R'}(M \otimes_R R', N') \to \text{Ext}^i_R(M, N')
$$
对 $i \geq 0$ 是同构。
\end{lemma}

\begin{proof}
选取 $M$ 的自由消解 $F_\bullet$。由于 $R \to R'$ 平坦，
$F_\bullet \otimes_R R'$ 是 $R'$ 上 $M \otimes_R R'$ 的自由消解。
本陈述即映射
$$
\Hom_{R'}(F_\bullet \otimes_R R', N') \to
\Hom_R(F_\bullet, N')
$$
在同调群上诱导同构；这是真的，因为由引理
\ref{algebra-lemma-adjoint-tensor-restrict}，它是复形的同构。
\end{proof}

\section{Ext 群的一个应用}
\label{algebra-section-ext-application}

\noindent
如下。

\begin{lemma}
\label{algebra-lemma-split-injection-after-completion}
设 $R$ 为诺特环，$I \subset R$ 为包含于 $R$ 的 Jacobson 根中的理想。
设 $N \to M$ 为有限 $R$-模之间的同态。
假设存在任意大的 $n$，使得
$N/I^nN \to M/I^nM$ 是分裂单射。
则 $N \to M$ 是分裂单射。
\end{lemma}

\begin{proof}
假设 $\varphi : N \to M$ 满足本引理的假设。
注意，这蕴含对任意大的 $n$ 都有 $\Ker(\varphi) \subset I^nN$。
因此由引理 \ref{algebra-lemma-intersection-powers-ideal-module}，
$\varphi$ 是单射。令 $Q = M/N$，于是有短正合序列
$$
0 \to N \to M \to Q \to 0.
$$
设
$$
F_2 \xrightarrow{d_2} F_1 \xrightarrow{d_1} F_0 \to Q \to 0
$$
为 $Q$ 的有限自由消解。可以选取提升映射 $F_0 \to Q$ 的映射
$\alpha : F_0 \to M$。这诱导映射
$\beta : F_1 \to N$，并且 $\beta \circ d_2 = 0$。
上面的扩张分裂，当且仅当存在映射 $\gamma : F_0 \to N$，
使得 $\beta = \gamma \circ d_1$。换言之，$\beta$ 在
$\Ext^1_R(Q, N)$ 中的类是上面的短正合序列分裂的障碍。

\medskip\noindent
假设 $n$ 是足够大的整数，使得
$N/I^nN \to M/I^nM$ 是分裂单射。这蕴含
$$
0 \to N/I^nN \to M/I^nM \to Q/I^nQ \to 0.
$$
仍是短正合的。此外，序列
$$
F_1/I^nF_1 \xrightarrow{d_1} F_0/I^nF_0 \to Q/I^nQ \to 0
$$
仍正合。与上面同样论证可知，由 $\beta$ 诱导的映射
$\overline{\beta} : F_1/I^nF_1 \to N/I^nN$
等于某个映射
$\overline{\gamma_n} : F_0/I^nF_0 \to N/I^nN$
所给出的复合 $\overline{\gamma_n} \circ d_1$。
由于 $F_0$ 自由，可以把 $\overline{\gamma_n}$ 提升为映射
$\gamma_n : F_0 \to N$；于是
$\beta - \gamma_n \circ d_1$ 是从 $F_1$ 到 $I^nN$ 的映射。
换言之，对这个 $n$，有
$$
\beta \in
\Im\Big(\Hom_R(F_0, N) \to \Hom_R(F_1, N)\Big) + I^n\Hom_R(F_1, N).
$$

\medskip\noindent
由于依假设，对任意大的 $n$ 都有这一性质，
由引理 \ref{algebra-lemma-intersection-powers-ideal-module}，
$\beta$ 在 $\Hom_R(F_0, N) \to \Hom_R(F_1, N)$ 的余核中的像为零。
因此按所需，$\beta$ 属于映射
$\Hom_R(F_0, N) \to \Hom_R(F_1, N)$ 的像。
\end{proof}

\section{Tor 群与平坦性}
\label{algebra-section-tor}

\noindent
本节使用上一节发展的一些同调代数来说明 Tor 群是什么。
具体地，设 $R$ 为环，$M$、$N$ 为两个 $R$-模。
选取 $M$ 的一个由自由 $R$-模组成的消解 $F_\bullet$，
参见引理 \ref{algebra-lemma-resolution-by-finite-free}。考虑同调复形
$$
F_\bullet \otimes_R N
:
\ldots
\to F_2 \otimes_R N
\to F_1 \otimes_R N
\to F_0 \otimes_R N
$$
我们把 $\text{Tor}^R_i(M, N)$ 定义为这个复形的第 $i$ 个同调群。
下面的引理说明这一对象在何种意义下是良定义的。

\begin{lemma}
\label{algebra-lemma-tor-welldefined}
设 $R$ 为环，$M_1, M_2, N$ 为 $R$-模。
假设 $F_\bullet$ 是模 $M_1$ 的自由消解，
$G_\bullet$ 是模 $M_2$ 的自由消解。
设 $\varphi : M_1 \to M_2$ 为模映射。
设 $\alpha : F_\bullet \to G_\bullet$ 为一个复形映射，
它在
$M_1 = \Coker(d_{F, 1}) \to M_2 = \Coker(d_{G, 1})$
上诱导 $\varphi$；参见引理 \ref{algebra-lemma-compare-resolutions}。
则诱导映射
$$
H_i(\alpha) :
H_i(F_\bullet \otimes_R N)
\longrightarrow
H_i(G_\bullet \otimes_R N)
$$
与 $\alpha$ 的选取无关。若 $\varphi$ 是同构，则所有映射
$H_i(\alpha)$ 都是同构。若 $M_1 = M_2$、
$F_\bullet = G_\bullet$ 且 $\varphi$ 是恒等映射，
则所有映射 $H_i(\alpha)$ 也都是恒等映射。
\end{lemma}

\begin{proof}
本引理的证明与引理 \ref{algebra-lemma-ext-welldefined} 的证明相同。
\end{proof}

\noindent
这个引理不仅蕴含 Tor 模是良定义的，还给出了构造
$(M, N) \mapsto \text{Tor}_i^R(M, N)$
关于第一个变量的函子性。关于第二个变量的函子性当然是显然的。
我们留给读者验证，每个 $\text{Tor}_i^R$ 实际上都是函子
$$
\text{Mod}_R \times \text{Mod}_R \to \text{Mod}_R.
$$
这里 $\text{Mod}_R$ 表示 $R$-模范畴；积范畴的定义见
《范畴论》，定义 \ref{categories-definition-product-category}。
具体地，给定 $R$-模态射 $M_1 \to M_2$ 和 $N_1 \to N_2$，
我们得到交换图
$$
\xymatrix{
\text{Tor}_i^R(M_1, N_1) \ar[r] \ar[d] &
\text{Tor}_i^R(M_1, N_2) \ar[d] \\
\text{Tor}_i^R(M_2, N_1) \ar[r] &
\text{Tor}_i^R(M_2, N_2) \\
}
$$

\begin{lemma}
\label{algebra-lemma-long-exact-sequence-tor}
设 $R$ 为环，$M$ 为 $R$-模。
假设 $0 \to N' \to N \to N'' \to 0$ 是 $R$-模的短正合序列。
存在长正合序列
$$
\text{Tor}_1^R(M, N')
\to \text{Tor}_1^R(M, N)
\to \text{Tor}_1^R(M, N'')
\to
M \otimes_R N'
\to M \otimes_R N
\to M \otimes_R N''
\to 0
$$
\end{lemma}

\begin{proof}
这个证明与引理 \ref{algebra-lemma-long-exact-seq-ext} 的证明相同。
\end{proof}

\noindent
考虑一个由 $R$-模组成的同调双复形
$$
\xymatrix{
\ldots \ar[r]^d &
A_{2, 0} \ar[r]^d &
A_{1, 0} \ar[r]^d &
A_{0, 0} \\
\ldots \ar[r]^d &
A_{2, 1} \ar[r]^d \ar[u]^\delta &
A_{1, 1} \ar[r]^d \ar[u]^\delta &
A_{0, 1} \ar[u]^\delta \\
\ldots \ar[r]^d &
A_{2, 2} \ar[r]^d \ar[u]^\delta &
A_{1, 2} \ar[r]^d \ar[u]^\delta &
A_{0, 2} \ar[u]^\delta \\
&
\ldots \ar[u]^\delta &
\ldots \ar[u]^\delta &
\ldots \ar[u]^\delta \\
}
$$
这意味着 $d_{i, j} : A_{i, j} \to A_{i-1, j}$ 和
$\delta_{i, j} : A_{i, j} \to A_{i, j-1}$ 具有下列性质：
\begin{enumerate}
\item 任意两个 $d_{i, j}$ 的复合为零。换言之，双复形的各行都是复形。
\item 任意两个 $\delta_{i, j}$ 的复合为零。换言之，双复形的各列都是复形。
\item 对任意指标对 $(i, j)$，有 $\delta_{i-1, j} \circ d_{i, j}
= d_{i, j-1} \circ \delta_{i, j}$。换言之，所有方块都交换。
\end{enumerate}
正确的做法是给每个这样的双复形配上一个谱序列。
不过目前我们用稍微容易一些的办法就足够了。

\medskip\noindent
具体地，仅为本节之目的，给定双复形
$(A_{\bullet, \bullet}, d, \delta)$，令
$R(A)_j = \Coker(A_{1, j} \to A_{0, j})$，并令
$U(A)_i = \Coker(A_{i, 1} \to A_{i, 0})$。（字母
$R$ 和 $U$ 分别提示 Right（右）与 Up（上）。）
我们用映射 $\delta$ 赋予 $R(A)_\bullet$ 复形结构；
类似地，用映射 $d$ 赋予 $U(A)_\bullet$ 复形结构。
换言之，我们得到如下巨大的交换图：
$$
\xymatrix{
\ldots \ar[r]^d &
U(A)_2 \ar[r]^d &
U(A)_1 \ar[r]^d &
U(A)_0 &
\\
\ldots \ar[r]^d &
A_{2, 0} \ar[r]^d \ar[u] &
A_{1, 0} \ar[r]^d \ar[u] &
A_{0, 0} \ar[r] \ar[u] &
R(A)_0 \\
\ldots \ar[r]^d &
A_{2, 1} \ar[r]^d \ar[u]^\delta &
A_{1, 1} \ar[r]^d \ar[u]^\delta &
A_{0, 1} \ar[r] \ar[u]^\delta &
R(A)_1 \ar[u]^\delta \\
\ldots \ar[r]^d &
A_{2, 2} \ar[r]^d \ar[u]^\delta &
A_{1, 2} \ar[r]^d \ar[u]^\delta &
A_{0, 2} \ar[r] \ar[u]^\delta &
R(A)_2 \ar[u]^\delta \\
&
\ldots \ar[u]^\delta &
\ldots \ar[u]^\delta &
\ldots \ar[u]^\delta &
\ldots \ar[u]^\delta \\
}
$$
（当然，这已不再是双复形。）双复形的态射
$\Phi : (A_{\bullet, \bullet}, d, \delta)
\to (B_{\bullet, \bullet}, d, \delta)$
是什么意思是清楚的；它会诱导复形态射
$R(\Phi) : R(A)_\bullet \to R(B)_\bullet$ 和
$U(\Phi) : U(A)_\bullet \to U(B)_\bullet$。

\begin{lemma}
\label{algebra-lemma-no-spectral-sequence}
设 $(A_{\bullet, \bullet}, d, \delta)$ 是满足下列条件的双复形：
\begin{enumerate}
\item 每一行 $A_{\bullet, j}$ 都是 $R(A)_j$ 的一个消解。
\item 每一列 $A_{i, \bullet}$ 都是 $U(A)_i$ 的一个消解。
\end{enumerate}
则存在典范同构
$$
H_i(R(A)_\bullet)
\cong
H_i(U(A)_\bullet).
$$
这些同构关于满足上述性质的双复形态射是函子性的。
\end{lemma}

\begin{proof}
我们将证明 $H_i(R(A)_\bullet))$ 与
$H_i(U(A)_\bullet)$ 都典范同构于第三个群，即
$$
\mathbf{H}_i(A) :=
\frac{
\{
(a_{i, 0}, a_{i-1, 1}, \ldots, a_{0, i})
\mid
d(a_{i, 0}) = \delta(a_{i-1, 1}), \ldots,
d(a_{1, i-1}) = \delta(a_{0, i})
\}}
{
\{
d(a_{i + 1, 0}) + \delta(a_{i, 1}),
d(a_{i, 1}) + \delta(a_{i-1, 2}),
\ldots,
d(a_{1, i}) + \delta(a_{0, i + 1})
\}
}
$$
这里采用记号约定：$a_{i, j}$ 表示 $A_{i, j}$ 的一个元素。
换言之，$\mathbf{H}_i$ 的一个元素由一条之字形表示；
当 $i = 2$ 时可画成
$$
\xymatrix{
a_{2, 0} \ar@{|->}[r] & d(a_{2, 0}) = \delta(a_{1, 1}) & \\
& a_{1, 1} \ar@{|->}[u] \ar@{|->}[r] & d(a_{1, 1}) = \delta(a_{0, 2}) \\
& & a_{0, 2} \ar@{|->}[u] \\
}
$$
自然地，我们要除去“平凡”之字形，也就是由形如
$(0, \ldots, 0, -\delta(a_{t + 1, t-i}),
d(a_{t + 1, t-i}), 0, \ldots, 0)$ 的元素生成的子模。
注意存在典范同态
$$
\mathbf{H}_i(A) \to H_i(R(A)_\bullet), \quad
(a_{i, 0}, a_{i-1, 1}, \ldots, a_{0, i}) \mapsto
\text{class of image of }a_{0, i}
$$
以及
$$
\mathbf{H}_i(A) \to H_i(U(A)_\bullet), \quad
(a_{i, 0}, a_{i-1, 1}, \ldots, a_{0, i}) \mapsto
\text{class of image of }a_{i, 0}
$$

\medskip\noindent
先证明这些映射是满射。设
$\overline{r} \in H_i(R(A)_\bullet)$。
令 $r \in R(A)_i$ 为表示 $\overline{r}$ 之类的一个余循环。
取映到 $r$ 的元素 $a_{0, i} \in A_{0, i}$。
因为 $\delta(r) = 0$，所以 $\delta(a_{0, i})$ 属于 $d$ 的像。
故存在 $a_{1, i-1} \in A_{1, i-1}$，使得
$d(a_{1, i-1}) = \delta(a_{0, i})$。
这又蕴含 $\delta(a_{1, i-1})$ 属于 $d$ 的核，因为
$d(\delta(a_{1, i-1})) = \delta(d(a_{1, i-1}))
= \delta(\delta(a_{0, i})) = 0$。由各行的正合性，
可找到元素 $a_{2, i-2}$，使得
$d(a_{2, i-2}) = \delta(a_{1, i-1})$。如此继续，
直到得到一条完整的之字形。当然，
$\mathbf{H}_i \to H_i(U(A))$ 的满射性也可同样证明。

\medskip\noindent
为证明单射性，我们采用完全相同的论证。
具体地，设给定一条映到 $H_i(R(A)_\bullet)$ 中零元的之字形
$(a_{i, 0}, a_{i-1, 1}, \ldots, a_{0, i})$。
这意味着 $a_{0, i}$ 映到
$\Coker(A_{i, 1} \to A_{i, 0})$ 的一个元素，
而此元素属于
$\delta : \Coker(A_{i + 1, 1} \to A_{i + 1, 0}) \to
\Coker(A_{i, 1} \to A_{i, 0})$ 的像。
换言之，$a_{0, i}$ 属于
$\delta \oplus d : A_{0, i + 1} \oplus A_{1, i} \to A_{0, i}$ 的像。
由平凡之字形的定义可知，我们可以用一条平凡之字形修改原之字形，
并假设 $a_{0, i} = 0$。这立刻蕴含
$d(a_{1, i-1}) = 0$。由于各行正合，
$a_{1, i-1}$ 属于 $d : A_{2, i-1} \to A_{1, i-1}$ 的像。
于是可以再次用一条平凡之字形修改，并假设原之字形形如
$(a_{i, 0}, a_{i-1, 1}, \ldots, a_{2, i-2}, 0, 0)$。
如此继续便得到所需的单射性。

\medskip\noindent
若 $\Phi : (A_{\bullet, \bullet}, d, \delta)
\to (B_{\bullet, \bullet}, d, \delta)$
是两个都满足本引理条件的双复形之间的态射，
则显然得到交换图
$$
\xymatrix{
H_i(U(A)_\bullet) \ar[d] &
\mathbf{H}_i(A) \ar[r] \ar[l] \ar[d] &
H_i(R(A)_\bullet) \ar[d] \\
H_i(U(B)_\bullet) &
\mathbf{H}_i(B) \ar[r] \ar[l] &
H_i(R(B)_\bullet) \\
}
$$
这证明了函子性。
\end{proof}

\begin{remark}
\label{algebra-remark-signs-double-complex}
上面构造的同构仅在相差符号的意义下才是“正确”的。
同调代数中很大一部分工作在于为各种映射选取符号，
并说明在插入适当符号后相应图表交换。
目前我们就使用上述证明给出的同构，稍后再处理符号问题。
\end{remark}

\begin{lemma}
\label{algebra-lemma-tor-left-right}
设 $R$ 为环。对任意 $i \geq 0$，函子
$\text{Mod}_R \times \text{Mod}_R \to \text{Mod}_R$，
$(M, N) \mapsto \text{Tor}_i^R(M, N)$ 与
$(M, N) \mapsto \text{Tor}_i^R(N, M)$ 典范同构。
\end{lemma}

\begin{proof}
设 $F_\bullet$ 是模 $M$ 的自由消解，
$G_\bullet$ 是模 $N$ 的自由消解。
考虑如下定义的双复形 $(A_{i, j}, d, \delta)$：
\begin{enumerate}
\item 令 $A_{i, j} = F_i \otimes_R G_j$；
\item 令 $d_{i, j} : F_i \otimes_R G_j \to F_{i-1} \otimes G_j$
等于 $d_{F, i} \otimes \text{id}$；
\item 令 $\delta_{i, j} : F_i \otimes_R G_j \to F_i \otimes G_{j-1}$
等于 $\text{id} \otimes d_{G, j}$。
\end{enumerate}
这个双复形通常简记为 $F_\bullet \otimes_R G_\bullet$。

\medskip\noindent
因为每个 $G_j$ 都自由，因而平坦，所以双复形的每一行除同调次数
$0$ 外都正合。因为每个 $F_i$ 都自由，因而平坦，所以双复形的每一列除同调次数
$0$ 外都正合。因此这个双复形满足引理
\ref{algebra-lemma-no-spectral-sequence} 的条件。

\medskip\noindent
为看清该引理给出什么，我们计算 $R(A)_\bullet$ 与 $U(A)_\bullet$。
具体地，
\begin{eqnarray*}
R(A)_i & = & \Coker(A_{1, i} \to A_{0, i}) \\
& = & \Coker(F_1 \otimes_R G_i \to F_0 \otimes_R G_i) \\
& = & \Coker(F_1 \to F_0) \otimes_R G_i \\
& = & M \otimes_R G_i
\end{eqnarray*}
事实上，这些同构与微分 $\delta$ 相容，因而作为同调复形有
$R(A)_\bullet = M \otimes_R G_\bullet$。
完全类似地，有 $U(A)_\bullet = F_\bullet \otimes_R N$。于是
\begin{eqnarray*}
\text{Tor}_i^R(M, N)
& = & H_i(F_\bullet \otimes_R N) \\
& = & H_i(U(A)_\bullet) \\
& = & H_i(R(A)_\bullet) \\
& = & H_i(M \otimes_R G_\bullet) \\
& = & H_i(G_\bullet \otimes_R M) \\
& = & \text{Tor}_i^R(N, M)
\end{eqnarray*}
这里第三个等号来自引理 \ref{algebra-lemma-no-spectral-sequence}，
第五个等号使用张量积的同构 $V \otimes W = W \otimes V$。

\medskip\noindent
函子性。假设有 $R$-模 $M_\nu$、$N_\nu$，其中
$\nu = 1, 2$。设 $\varphi : M_1 \to M_2$ 和
$\psi : N_1 \to N_2$ 为 $R$-模态射。
假设对 $M_\nu$ 有自由消解 $F_{\nu, \bullet}$，
对 $N_\nu$ 有自由消解 $G_{\nu, \bullet}$。
由引理 \ref{algebra-lemma-compare-resolutions}，可选取与
$\varphi$ 和 $\psi$ 相容的复形映射
$\alpha : F_{1, \bullet} \to F_{2, \bullet}$ 与
$\beta : G_{1, \bullet} \to G_{2, \bullet}$。
我们断言，映射对 $(\alpha, \beta)$ 诱导双复形态射
$$
\alpha \otimes \beta :
F_{1, \bullet} \otimes_R G_{1, \bullet}
\longrightarrow
F_{2, \bullet} \otimes_R G_{2, \bullet}
$$
这确实只需直接验证：映射
$F_{1, i} \otimes_R G_{1, j} \to F_{2, i} \otimes_R G_{2, j}$
由 $\alpha_i \otimes \beta_j$ 给出；其中 $\alpha_i$ 与
$\beta_j$ 分别是 $\alpha$ 的次数 $i$ 分量与 $\beta$ 的次数 $j$ 分量。
读者也容易验证，诱导映射
$R(F_{1, \bullet} \otimes_R G_{1, \bullet})_\bullet
\to R(F_{2, \bullet} \otimes_R G_{2, \bullet})_\bullet$
与由 $\varphi \otimes \beta$ 诱导的映射
$M_1 \otimes_R G_{1, \bullet}
\to M_2 \otimes_R G_{2, \bullet}$ 相同。
在 $U(-)_\bullet$ 复形上诱导的映射也类似。
因此，函子性断言来自引理
\ref{algebra-lemma-no-spectral-sequence} 中的函子性断言。
\end{proof}

\begin{remark}
\label{algebra-remark-curiosity-signs-swap}
上面的 $M = N$ 情形很有意思。此时得到典范映射
$\text{Tor}_i^R(M, M) \to \text{Tor}_i^R(M, M)$。
注意，这个映射不是恒等映射；即使 $i = 0$，它也不是恒等映射！
例如，若 $V$ 是域上的 $n$ 维向量空间，则交换映射
$V \otimes_k V \to V \otimes_k V$ 有
$(n^2 + n)/2$ 个特征值 $+1$ 和
$(n^2-n)/2$ 个特征值 $-1$。在特征 $2$ 时，它甚至不可对角化。
注意，即使改变映射的符号也无法消除这一现象。
\end{remark}

\begin{lemma}
\label{algebra-lemma-tor-noetherian}
设 $R$ 为诺特环，$M$、$N$ 为有限 $R$-模。
则对所有 $p$，$\text{Tor}_p^R(M, N)$ 都是有限 $R$-模。
\end{lemma}

\begin{proof}
这是因为 $\text{Tor}_p^R(M, N)$ 可计算为复形
$F_\bullet \otimes_R N$ 的上同调群，其中每个 $F_n$ 都是有限自由
$R$-模；参见引理 \ref{algebra-lemma-resolution-by-finite-free}。
\end{proof}

\begin{lemma}
\label{algebra-lemma-characterize-flat}
设 $R$ 为环，$M$ 为 $R$-模。下列条件等价：
\begin{enumerate}
\item 模 $M$ 在 $R$ 上平坦。
\item 对所有 $i > 0$，函子 $\text{Tor}_i^R(M, -)$ 为零。
\item 函子 $\text{Tor}_1^R(M, -)$ 为零。
\item 对所有理想 $I \subset R$，有 $\text{Tor}_1^R(M, R/I) = 0$。
\item 对所有有限生成理想 $I \subset R$，有
$\text{Tor}_1^R(M, R/I) = 0$。
\end{enumerate}
\end{lemma}

\begin{proof}
假设 $M$ 平坦。设 $N$ 为 $R$-模，并设 $F_\bullet$ 是
$N$ 的自由消解。由 $M$ 的平坦性，
$F_\bullet \otimes_R M$ 是 $N \otimes_R M$ 的一个消解。
故所有高阶 Tor 群都消失。

\medskip\noindent
现在只需证明最后一个条件蕴含 $M$ 平坦。
设 $I \subset R$ 为理想。考虑短正合序列
$0 \to I \to R \to R/I \to 0$。
应用引理 \ref{algebra-lemma-long-exact-sequence-tor}，得到正合序列
$$
\text{Tor}_1^R(M, R/I) \to
M \otimes_R I \to
M \otimes_R R \to
M \otimes_R R/I \to
0
$$
由于显然 $M \otimes_R R = M$，可知最后一个假设蕴含：
对每个有限生成理想 $I$，映射 $M \otimes_R I \to M$ 都是单射。
因此由引理 \ref{algebra-lemma-flat}，$M$ 平坦。
\end{proof}

\begin{remark}
\label{algebra-remark-Tor-ring-mod-ideal}
引理 \ref{algebra-lemma-characterize-flat} 的证明实际上表明
$$
\text{Tor}_1^R(M, R/I)
=
\Ker(I \otimes_R M \to M).
$$
\end{remark}

\section{Tor 的函子性}
\label{algebra-section-functoriality-tor}

\noindent
本节简要讨论 $\text{Tor}$ 关于环变换等操作的函子性。
下面列出若干需要完成的事项。
\begin{enumerate}
\item 给定环映射 $R \to R'$、一个 $R$-模 $M$ 和一个
$R'$-模 $N'$，$R$-模 $\text{Tor}_i^R(M, N')$ 具有自然的
$R'$-模结构。
\item 给定环映射 $R \to R'$ 以及 $R$-模 $M$、$N$，
存在自然的 $R$-模映射
$\text{Tor}_i^R(M, N) \to \text{Tor}_i^{R'}(M \otimes_R R', N \otimes_R R')$。
\item 给定环映射 $R \to R'$、一个 $R$-模 $M$ 和一个
$R'$-模 $N'$，存在自然的 $R'$-模映射
$\text{Tor}_i^R(M, N') \to \text{Tor}_i^{R'}(M \otimes_R R', N')$。
\end{enumerate}

\begin{lemma}
\label{algebra-lemma-flat-base-change-tor}
给定平坦环映射 $R \to R'$ 以及 $R$-模 $M$、$N$，
对所有 $i$，自然的 $R$-模映射
$\text{Tor}_i^R(M, N)\otimes_R R'
\to \text{Tor}_i^{R'}(M \otimes_R R', N \otimes_R R')$
都是同构。
\end{lemma}

\begin{proof}
略。这是因为 $M$ 在 $R$ 上的自由消解 $F_\bullet$ 在与
$R'$ 关于 $R$ 作张量积后仍正合，因而
$(F_\bullet \otimes_R N)\otimes_R R'$ 计算 $R'$ 上的 Tor 群。
\end{proof}

\noindent
下面的引理似乎放在其他地方都不合适。

\begin{lemma}
\label{algebra-lemma-tor-commutes-filtered-colimits}
设 $R$ 为环，$M = \colim M_i$ 是 $R$-模的滤过余极限，
$N$ 是 $R$-模。则对所有 $n$，
$\text{Tor}_n^R(M, N) = \colim \text{Tor}_n^R(M_i, N)$。
\end{lemma}

\begin{proof}
选取 $N$ 的自由消解 $F_\bullet$。由引理
\ref{algebra-lemma-tensor-products-commute-with-limits}，作为复形有
$F_\bullet \otimes_R M = \colim F_\bullet \otimes_R M_i$。
于是由引理 \ref{algebra-lemma-directed-colimit-exact} 得到结论。
\end{proof}

\section{投射模}
\label{algebra-section-projective}

\noindent
下面给出关于投射模的一些引理。

\begin{definition}
\label{algebra-definition-projective}
设 $R$ 为环。一个 $R$-模 $P$ 称为\textit{投射的}，
当且仅当函子
$\Hom_R(P, -) : \text{Mod}_R \to \text{Mod}_R$ 是正合函子。
\end{definition}

\noindent
对任意 $R$-模 $M$，函子 $\Hom_R(M, - )$ 都左正合，
参见引理 \ref{algebra-lemma-hom-exact}。
因此，$P$ 为投射模这一条件真正表达的是：
给定 $R$-模满射 $N \to N'$，映射
$\Hom_R(P, N) \to \Hom_R(P, N')$ 是满射。

\begin{lemma}
\label{algebra-lemma-characterize-projective}
设 $R$ 为环，$P$ 为 $R$-模。下列条件等价：
\begin{enumerate}
\item $P$ 是投射模；
\item $P$ 是某个自由 $R$-模的直和因子；
\item 对每个 $R$-模 $M$，都有 $\Ext^1_R(P, M) = 0$。
\end{enumerate}
\end{lemma}

\begin{proof}
假设 $P$ 是投射模。选取满射 $\pi : F \to P$，其中 $F$
是自由 $R$-模。因为 $P$ 投射，存在
$i \in \Hom_R(P, F)$，使得 $\pi \circ i = \text{id}_P$。
换言之，$F \cong \Ker(\pi) \oplus i(P)$，从而 $P$ 是 $F$ 的直和因子。

\medskip\noindent
反过来，假设 $P \oplus Q = F$ 是自由 $R$-模。
注意，自由模 $F = \bigoplus_{i \in I} R$ 是投射模，
因为有等式：
$$
\Hom_R(F, M) = \prod_{i \in I} M
$$
而函子 $M \mapsto \prod_{i \in I} M$ 正合。
又因为作为函子有下式：
$$
\Hom_R(F, -) = \Hom_R(P, -) \times \Hom_R(Q, -)
$$
所以 $P$ 与 $Q$ 都是投射模。

\medskip\noindent
假设 $P \oplus Q = F$ 是自由 $R$-模。则有形如
$$
\ldots F \xrightarrow{a} F \xrightarrow{b} F \to P \to 0
$$
的自由消解 $F_\bullet$，其中映射 $a, b$ 交替出现，
并且分别等于到 $P$ 和 $Q$ 上的投影算子。
因此复形 $\Hom_R(F_\bullet, M)$ 在次数 $\geq 1$ 上分裂正合，
从而得到条件 (3) 中的消失性。

\medskip\noindent
假设对每个 $R$-模 $M$ 都有 $\Ext^1_R(P, M) = 0$。
取一个自由消解 $F_\bullet \to P$，并令
$M = \Im(F_1 \to F_0) = \Ker(F_0 \to P)$。
考虑由商映射 $\pi : F_1 \to M$ 的类给出的元素
$\xi \in \Ext^1_R(P, M)$。由于 $\xi$ 为零，存在映射
$s : F_0 \to M$，使得 $\pi = s \circ (F_1 \to F_0)$。
显然，这意味着
$$
F_0 = \Ker(s) \oplus \Ker(F_0 \to P) =
P \oplus \Ker(F_0 \to P)
$$
结论得证。
\end{proof}

\begin{lemma}
\label{algebra-lemma-characterize-finite-projective-noetherian}
设 $R$ 为诺特环，$P$ 为有限 $R$-模。
如果对每个有限 $R$-模 $M$ 都有 $\Ext^1_R(P, M) = 0$，
则 $P$ 是投射模。
\end{lemma}

\noindent
这个引理可以加强：当不假设 $R$ 为诺特环时，
存在一个适用于有限表现 $R$-模的版本；在诺特情形，
也有一个让 $M$ 遍历所有有限长度模的版本。

\begin{proof}
选取满射 $R^{\oplus n} \to P$，其核为 $M$。
因为 $\Ext^1_R(P, M) = 0$，这个满射分裂；
于是由引理 \ref{algebra-lemma-characterize-projective} 得到结论。
\end{proof}

\begin{lemma}
\label{algebra-lemma-direct-sum-projective}
投射模的直和是投射模。
\end{lemma}

\begin{proof}
这由引理 \ref{algebra-lemma-characterize-projective} 中
“投射模是自由模的直和因子”这一刻画得到。
\end{proof}

\begin{lemma}
\label{algebra-lemma-lift-projective-module}
设 $R$ 为环，$I \subset R$ 为幂零理想，
$\overline{P}$ 为投射 $R/I$-模。
则存在投射 $R$-模 $P$，使得 $P/IP \cong \overline{P}$。
\end{lemma}

\begin{proof}
由引理 \ref{algebra-lemma-characterize-projective}，可以选取集合 $A$ 以及直和分解
$\bigoplus_{\alpha \in A} R/I = \overline{P} \oplus \overline{K}$，
其中 $\overline{K}$ 是某个 $R/I$-模。
记 $F = \bigoplus_{\alpha \in A} R$ 为以 $A$ 为基的自由 $R$-模。
设 $\overline{p}$ 是直和因子
$\overline{P}$ 在 $\bigoplus_{\alpha \in A} R/I$ 中对应的投影算子，
选取它的一个提升 $p : F \to F$。
注意，$p^2 - p \in \text{End}_R(F)$ 是 $F$ 的幂零自同态
（因为 $I$ 幂零且 $p^2 - p$ 的矩阵元素属于 $I$；更精确地，
若 $I^n = 0$，则 $(p^2 - p)^n = 0$）。
因此由引理 \ref{algebra-lemma-lift-idempotents-noncommutative}，
可以修改 $p$ 的选取并假设 $p$ 是投影算子。令 $P = \Im(p)$。
\end{proof}

\begin{lemma}
\label{algebra-lemma-lift-finite-projective-module}
设 $R$ 为环，$I \subset R$ 为局部幂零理想，
$\overline{P}$ 为有限投射 $R/I$-模。
则存在有限投射 $R$-模 $P$，使得 $P/IP \cong \overline{P}$。
\end{lemma}

\begin{proof}
回顾一下，由引理 \ref{algebra-lemma-characterize-projective}，
$\overline{P}$ 是自由 $R/I$-模
$\bigoplus_{\alpha \in A} R/I$ 的直和因子。
因为 $\overline{P}$ 有限，所以存在有限子集 $A' \subset A$，
使得 $\overline{P}$ 包含于 $\bigoplus_{\alpha \in A'} R/I$。
因此可以假设存在直和分解
$(R/I)^{\oplus n} = \overline{P} \oplus \overline{K}$，
其中 $n$ 是某个整数，$\overline{K}$ 是某个 $R/I$-模。
设 $\overline{p}$ 是直和因子 $\overline{P}$ 在
$(R/I)^{\oplus n}$ 中对应的投影算子，选取其提升
$p \in \text{Mat}(n \times n, R)$。
注意 $p^2 - p \in \text{Mat}(n \times n, R)$ 是幂零的：
因为 $I$ 局部幂零，且 $p^2 - p$ 的矩阵元素 $c_{ij}$ 属于 $I$，
所以对某个 $t > 0$ 有 $c_{ij}^t = 0$，进而
$(p^2 - p)^{tn^2} = 0$（考察矩阵系数即可）。
因此由引理 \ref{algebra-lemma-lift-idempotents-noncommutative}，
可以修改 $p$ 的选取并假设 $p$ 是投影算子。令 $P = \Im(p)$。
\end{proof}

\begin{lemma}
\label{algebra-lemma-lift-projective}
设 $R$ 为环，$I \subset R$ 为理想，$M$ 为 $R$-模。假设
\begin{enumerate}
\item $I$ 是幂零的；
\item $M/IM$ 是投射 $R/I$-模；
\item $M$ 是平坦 $R$-模。
\end{enumerate}
则 $M$ 是投射 $R$-模。
\end{lemma}

\begin{proof}
由引理 \ref{algebra-lemma-lift-projective-module}，
可以找到投射 $R$-模 $P$ 和同构 $P/IP \to M/IM$。
我们将证明 $M$ 同构于 $P$，从而完成证明。
因为 $P$ 投射，可以把映射 $P \to P/IP \to M/IM$
提升为 $R$-模映射 $P \to M$；它模 $I$ 后是同构。
对某个 $n$ 有 $I^n = 0$，所以可利用滤过
\begin{align*}
0 = I^nM \subset I^{n - 1}M \subset \ldots \subset IM \subset M \\
0 = I^nP \subset I^{n - 1}P \subset \ldots \subset IP \subset P
\end{align*}
看出只需证明诱导映射
$I^aP/I^{a + 1}P \to I^aM/I^{a + 1}M$ 都是双射。
因为 $P$ 与 $M$ 都是平坦 $R$-模，可以把这个映射等同于由
$P \to M$ 诱导的映射
$$
I^a/I^{a + 1} \otimes_{R/I} P/IP
\longrightarrow
I^a/I^{a + 1} \otimes_{R/I} M/IM
$$
我们选取 $P \to M$ 时已使诱导映射
$P/IP \to M/IM$ 成为同构，故结论成立。
\end{proof}

\begin{lemma}
\label{algebra-lemma-projective-modulo-two-ideals}
设 $R$ 为环，$I, J \subset R$ 是满足 $I \cap J = 0$ 的理想。
设 $P$ 为 $R$-模，并且 $P/IP$ 是投射 $R/I$-模，
$P/JP$ 是投射 $R/J$-模。则 $P$ 是投射 $R$-模。
\end{lemma}

\begin{proof}
选取满射 $p : F \to P$，其中 $F$ 是自由 $R$-模。
因为 $P/IP$ 是投射 $R/I$-模，可以选取映射
$f : P/IP \to F/IF$，它是 $p \bmod I$ 的右逆。
考虑由商映射 $F/JF \to F/(I + J)F$ 和 $p$ 诱导的映射
$q : F/JF \to F/(I + J)F \times_{P/(I + J)P} P/JP$。
注意，$q$ 是 $R/J$-模的满射（略去一个小细节）。
再考虑由 $f$ 与恒等映射诱导的映射
$f' : P/JP \to F/(I + J)F \times_{P/(I + J)P} P/JP$。
因为 $P/JP$ 是投射 $R/J$-模且 $q$ 为满射，
可以选取映射 $g : P/JP \to F/JF$，使得 $q \circ g = f'$。
于是 $f \bmod I + J = g \bmod I + J$。
又因为 $I \cap J = 0$，有
$F = F/JF \times_{F/(I + J)F} F/IF$；
所以得到良定义的 $R$-模同态 $h = (f, g) : P \to F$。
映射 $a = p \circ h : P \to P$ 模 $I$ 和模 $J$ 后都化为恒等映射。
于是 $a$ 在子模 $IP$ 上也诱导恒等映射：若
$x = \sum i_\alpha x_\alpha$，其中 $i_\alpha \in I$、
$x_\alpha \in P$，则
$a(x) = \sum i_\alpha a(x_\alpha) = \sum i_\alpha x_\alpha = x$，
因为 $a(x_\alpha) = x_\alpha + j_\alpha$，其中
$j_\alpha \in J$，且 $i_\alpha j_\alpha = 0$。
因此 $a : P \to P$ 在 $IP$ 以及 $P/IP$ 上都作为恒等映射，
从而 $a$ 是 $P$ 的自同构（略去一个小细节）。
故 $P$ 是 $F$ 的直和因子，因而是投射模。
\end{proof}

\section{有限投射模}
\label{algebra-section-finite-projective-modules}

\begin{definition}
\label{algebra-definition-locally-free}
设 $R$ 为环，$M$ 为 $R$-模。
\begin{enumerate}
\item 若可用标准开集 $D(f_i)$（$i \in I$）覆盖 $\Spec(R)$，
使得对所有 $i \in I$，$M_{f_i}$ 都是自由 $R_{f_i}$-模，
则称 $M$ \textit{局部自由}。
\item 若可选取上述覆盖，使每个 $M_{f_i}$ 都有限自由，
则称 $M$ \textit{有限局部自由}。
\item 若可选取上述覆盖，使每个 $M_{f_i}$ 都同构于
$R_{f_i}^{\oplus r}$，则称 $M$ 是\textit{秩为 $r$ 的有限局部自由模}。
\end{enumerate}
\end{definition}

\noindent
注意，由引理 \ref{algebra-lemma-cover}，有限局部自由 $R$-模自动是有限表现的。
此外，若 $M$ 是环 $R$ 上秩为 $r$ 的有限局部自由模且 $R$ 非零，
则由引理 \ref{algebra-lemma-rank}，$r$ 唯一确定
（因为局部化 $R_{f_i}$ 中至少有一个是非零环）。

\begin{lemma}
\label{algebra-lemma-finite-projective}
设 $R$ 为环，$M$ 为 $R$-模。下列条件等价：
\begin{enumerate}
\item $M$ 是有限表现且在 $R$ 上平坦；
\item $M$ 是有限投射模；
\item $M$ 是某个有限自由 $R$-模的直和因子；
\item $M$ 是有限表现的，并且对所有
$\mathfrak p \in \Spec(R)$，局部化 $M_{\mathfrak p}$ 都自由；
\item $M$ 是有限表现的，并且对所有极大理想
$\mathfrak m \subset R$，局部化 $M_{\mathfrak m}$ 都自由；
\item $M$ 有限且局部自由；
\item $M$ 有限局部自由；
\item $M$ 有限，对每个素理想 $\mathfrak p$，模
$M_{\mathfrak p}$ 都自由，并且函数
$$
\rho_M : \Spec(R) \to \mathbf{Z}, \quad
\mathfrak p
\longmapsto
\dim_{\kappa(\mathfrak p)} M \otimes_R \kappa(\mathfrak p)
$$
在 Zariski 拓扑下局部常值。
\end{enumerate}
\end{lemma}

\begin{proof}
先假设 $M$ 是有限投射模，即条件 (2) 成立。
取满射 $R^n \to M$，并令 $K$ 为其核。
因为 $M$ 投射，短正合序列
$0 \to K \to R^n \to M \to 0$ 分裂。
所以 (2) $\Rightarrow$ (3)。由投射模的直和因子仍投射这一事实，
(3) $\Rightarrow$ (2)；参见引理 \ref{algebra-lemma-characterize-projective}。

\medskip\noindent
假设 (3) 成立，因而可写 $K \oplus M \cong R^{\oplus n}$。
于是 $K$ 是 $R^n$ 的直和因子，从而有限生成。
这说明 $M = R^{\oplus n}/K$ 是有限表现的。
换言之，(3) $\Rightarrow$ (1)。

\medskip\noindent
假设 $M$ 是有限表现且平坦的，即 (1) 成立。
我们证明 (7)。任选素理想 $\mathfrak p$，并取
$x_1, \ldots, x_r \in M$，使它们的像构成
$M \otimes_R \kappa(\mathfrak p)$ 的一组基。
由 Nakayama 引理（采用引理 \ref{algebra-lemma-NAK-localization} 的形式），
存在 $g \in R$、$g \not \in \mathfrak p$，使这些元素生成 $M_g$。
相应满射 $\varphi : R_g^{\oplus r} \to M_g$ 有下列两个性质：
(a) $\Ker(\varphi)$ 是有限 $R_g$-模（参见引理 \ref{algebra-lemma-extension}）；
(b) 由 $M_g$ 在 $R_g$ 上的平坦性，
$\Ker(\varphi) \otimes \kappa(\mathfrak p) = 0$
（参见引理 \ref{algebra-lemma-flat-tor-zero}）。
再次使用 Nakayama 引理，存在 $g' \in R_g$，使得
$\Ker(\varphi)_{g'} = 0$。换言之，$M_{gg'}$ 自由。

\medskip\noindent
有限局部自由模是有限模，参见引理 \ref{algebra-lemma-cover}，
所以 (7) $\Rightarrow$ (6)。显然 (6) $\Rightarrow$ (7)，
并且 (7) $\Rightarrow$ (8)。

\medskip\noindent
有限局部自由模是有限表现模，参见引理 \ref{algebra-lemma-cover}，
所以 (7) $\Rightarrow$ (4)。当然 (4) 蕴含 (5)。
由于平坦性可以局部检验（参见引理 \ref{algebra-lemma-flat-localization}），
可知 (5) 蕴含 (1)。此时我们已有
$$
\xymatrix{
(2) \ar@{<=>}[r] & (3) \ar@{=>}[r] & (1) \ar@{=>}[r] &
(7)  \ar@{<=>}[r] \ar@{=>}[rd] \ar@{=>}[d] & (6) \\
& & (5) \ar@{=>}[u] & (4) \ar@{=>}[l] & (8)
}
$$

\medskip\noindent
假设 $M$ 满足 (1)、(4)、(5)、(6) 和 (7)。
我们证明 (3)。只需说明 $M$ 投射，即说明 $\Hom_R(M, -)$ 正合。
设 $0 \to N'' \to N \to N'\to 0$ 是 $R$-模的短正合序列。
我们必须证明
$0 \to \Hom_R(M, N'') \to \Hom_R(M, N) \to
\Hom_R(M, N') \to 0$ 正合。
因为 $M$ 有限局部自由，存在覆盖
$\Spec(R) = \bigcup D(f_i)$，使 $M_{f_i}$ 有限自由。
由引理 \ref{algebra-lemma-hom-from-finitely-presented}，序列
$$
0 \to \Hom_R(M, N'')_{f_i} \to \Hom_R(M, N)_{f_i} \to
\Hom_R(M, N')_{f_i} \to 0
$$
等于
$0 \to \Hom_{R_{f_i}}(M_{f_i}, N''_{f_i}) \to
\Hom_{R_{f_i}}(M_{f_i}, N_{f_i}) \to
\Hom_{R_{f_i}}(M_{f_i}, N'_{f_i}) \to 0$。
该序列正合，因为 $M_{f_i}$ 自由，并且局部化序列
$0 \to N''_{f_i} \to N_{f_i} \to N'_{f_i} \to 0$
正合（局部化正合）。于是由引理 \ref{algebra-lemma-cover}，
$0 \to \Hom_R(M, N'') \to \Hom_R(M, N) \to
\Hom_R(M, N') \to 0$ 正合。

\medskip\noindent
最后，假设 (8) 成立。选取极大理想 $\mathfrak m \subset R$，
并取 $x_1, \ldots, x_r \in M$，使它们的像构成
$M \otimes_R \kappa(\mathfrak m) = M/\mathfrak mM$ 的一组
$\kappa(\mathfrak m)$-基。特别地，$\rho_M(\mathfrak m) = r$。
由 Nakayama 引理 \ref{algebra-lemma-NAK}，存在
$f \in R$、$f \not \in \mathfrak m$，使
$x_1, \ldots, x_r$ 在 $R_f$ 上生成 $M_f$。
由 $\rho_M$ 局部常值的假设，存在
$g \in R$、$g \not \in \mathfrak m$，使
$\rho_M$ 在 $D(g)$ 上恒等于 $r$。我们断言
$$
\Psi : R_{fg}^{\oplus r} \longrightarrow M_{fg}, \quad
(a_1, \ldots, a_r) \longmapsto \sum a_i x_i
$$
是同构。这个断言将说明 $M$ 有限局部自由，即 (7) 成立。
为证明该断言，由引理 \ref{algebra-lemma-characterize-zero-local}，
只需证明对所有 $\mathfrak p \in D(fg)$，诱导的局部化映射
$\Psi_{\mathfrak p} : R_{\mathfrak p}^{\oplus r} \to M_{\mathfrak p}$
都是同构。由 $f$ 的选取，$\Psi_{\mathfrak p}$ 是满射。
由假设 (8)，有
$M_{\mathfrak p} \cong R_{\mathfrak p}^{\oplus \rho_M(\mathfrak p)}$，
而由 $g$ 的选取，有 $\rho_M(\mathfrak p) = r$。
因此 $\Psi_{\mathfrak p}$ 给出满射
$R_{\mathfrak p}^{\oplus r} \to
M_{\mathfrak p} \cong R_{\mathfrak p}^{\oplus r}$，
从而由引理 \ref{algebra-lemma-fun}，它是同构。
（当然，最后这个事实也可由简单的矩阵论证得到。）
\end{proof}

\begin{lemma}
\label{algebra-lemma-finite-projective-reduced}
设 $R$ 为约化环，$M$ 为 $R$-模。
则引理 \ref{algebra-lemma-finite-projective} 中的等价条件还等价于
\begin{enumerate}
\item[(9)] $M$ 有限，并且函数
$\rho_M : \Spec(R) \to \mathbf{Z}$，$\mathfrak p \mapsto
\dim_{\kappa(\mathfrak p)} M \otimes_R \kappa(\mathfrak p)$
在 Zariski 拓扑下局部常值。
\end{enumerate}
\end{lemma}

\begin{proof}
选取极大理想 $\mathfrak m \subset R$，并取
$x_1, \ldots, x_r \in M$，使它们的像构成
$M \otimes_R \kappa(\mathfrak m) = M/\mathfrak mM$ 的一组
$\kappa(\mathfrak m)$-基。特别地，$\rho_M(\mathfrak m) = r$。
由 Nakayama 引理 \ref{algebra-lemma-NAK}，存在
$f \in R$、$f \not \in \mathfrak m$，使
$x_1, \ldots, x_r$ 在 $R_f$ 上生成 $M_f$。
由 $\rho_M$ 局部常值的假设，存在
$g \in R$、$g \not \in \mathfrak m$，使
$\rho_M$ 在 $D(g)$ 上恒等于 $r$。我们断言
$$
\Psi : R_{fg}^{\oplus r} \longrightarrow M_{fg}, \quad
(a_1, \ldots, a_r) \longmapsto \sum a_i x_i
$$
是同构。这个断言将说明 $M$ 有限局部自由，即 (7) 成立。
因为 $\Psi$ 是满射，只需证明 $\Psi$ 是单射。
由于 $R_{fg}$ 约化，由引理
\ref{algebra-lemma-reduced-ring-sub-product-fields}，
只需证明在 $R_{fg}$ 的所有极小素理想 $\mathfrak p$ 处局部化后，
$\Psi$ 都是单射。然而，由引理
\ref{algebra-lemma-minimal-prime-reduced-ring}，有
$R_\mathfrak p = \kappa(\mathfrak p)$；又有
$\rho_M(\mathfrak p) = r$，所以
$\Psi_\mathfrak p :
R_\mathfrak p^{\oplus r} \to M \otimes_R \kappa(\mathfrak p)$
是相同维数的有限维向量空间之间的满射，因而是同构。
\end{proof}

\begin{remark}
\label{algebra-remark-warning}
有限且在 $R$ 上平坦的 $R$-模并不自动投射。
一个反例是：取 $R = \mathcal{C}^\infty(\mathbf{R})$ 为
$\mathbf{R}$ 上无穷次可微函数构成的环，并取
$M = R_{\mathfrak m} = R/I$，其中
$\mathfrak m = \{f \in R \mid f(0) = 0\}$，且
$I = \{f \in R \mid \exists \epsilon, \epsilon > 0 :
f(x) = 0\ \forall x, |x| < \epsilon\}$。
\end{remark}

\begin{lemma}
\label{algebra-lemma-finite-flat-local}
（警告：参见注 \ref{algebra-remark-warning}。）
假设 $R$ 是局部环，$M$ 是有限平坦 $R$-模。
则 $M$ 有限自由。
\end{lemma}

\begin{proof}
这由平坦性的方程判据得到；参见引理 \ref{algebra-lemma-flat-eq}。
具体地，假设 $x_1, \ldots, x_r \in M$ 的像构成
$M/\mathfrak mM$ 的一组基。由 Nakayama 引理
\ref{algebra-lemma-NAK}，这些元素生成 $M$。
我们要证明 $x_i$ 之间不存在关系。为此，改为对 $n$ 归纳证明：
若 $x_1, \ldots, x_n \in M$ 在向量空间
$M/\mathfrak mM$ 中线性无关，则它们在 $R$ 上线性无关。

\medskip\noindent
归纳起点是：有 $x \in M$、$x \not\in \mathfrak mM$，
以及关系 $fx = 0$。由方程判据，存在
$y_j \in M$ 和 $a_j \in R$，使得
$x = \sum a_j y_j$，并且对所有 $j$ 都有 $fa_j = 0$。
因为 $x \not\in \mathfrak mM$，至少有一个 $a_j$ 是单位，
故 $f = 0 $。

\medskip\noindent
假设 $\sum f_i x_i$ 是 $x_1, \ldots, x_n$ 之间的一个关系。
由 $x_i$ 的选取，有 $f_i \in \mathfrak m$。
根据平坦性的方程判据，存在 $a_{ij} \in R$ 与
$y_j \in M$，使得 $x_i = \sum a_{ij} y_j$，且
$\sum f_i a_{ij} = 0$。由于 $x_n \not \in \mathfrak mM$，
至少对一个 $j$ 有 $a_{nj}\not\in \mathfrak m$。
由 $\sum f_i a_{ij} = 0$，得到
$f_n = \sum_{i = 1}^{n-1} (-a_{ij}/a_{nj}) f_i$。
现在关系 $\sum f_i x_i = 0$ 可改写为
$\sum_{i = 1}^{n-1} f_i( x_i + (-a_{ij}/a_{nj}) x_n) = 0$。
注意，元素 $x_i + (-a_{ij}/a_{nj}) x_n$ 在
$M/\mathfrak mM$ 中的像是 $n-1$ 个线性无关元素。
由归纳假设，所有 $f_i$（$i \leq n-1$）都必须为零；
而且 $f_n = \sum_{i = 1}^{n-1} (-a_{ij}/a_{nj}) f_i$。
这完成归纳步骤。
\end{proof}

\begin{lemma}
\label{algebra-lemma-finite-projective-descends}
设 $R \to S$ 是局部环之间的平坦局部同态，
$M$ 是有限 $R$-模。则 $M$ 在 $R$ 上有限投射，
当且仅当 $M \otimes_R S$ 在 $S$ 上有限投射。
\end{lemma}

\begin{proof}
由引理 \ref{algebra-lemma-finite-projective}，
在局部环上有限投射等价于有限自由。
假设 $M \otimes_R S$ 是有限自由 $S$-模。
选取 $x_1, \ldots, x_r \in M$，使它们在
$M/\mathfrak m_RM$ 中的像构成一组
$\kappa(\mathfrak m)$-基。于是
$x_1 \otimes 1, \ldots, x_r \otimes 1$
构成 $M \otimes_R S$ 的一组基。
这蕴含映射 $R^{\oplus r} \to M, (a_i) \mapsto \sum a_i x_i$
与 $S$ 作张量积后成为同构。
由 $R \to S$ 的忠实平坦性（参见引理 \ref{algebra-lemma-local-flat-ff}），
该映射本身是同构。
\end{proof}

\begin{lemma}
\label{algebra-lemma-locally-free-semi-local-free}
设 $R$ 为半局部环，$M$ 为有限局部自由模。
若 $M$ 的秩为常数，则 $M$ 自由。特别地，若
$R$ 的谱连通，则 $M$ 自由。
\end{lemma}

\begin{proof}
略。提示：先利用秩为常数，证明对 $R$ 的所有极大理想
$\mathfrak m_1, \ldots, \mathfrak m_n$，
$M/\mathfrak m_iM$ 都有相同的维数 $d$。
然后利用中国剩余定理，证明存在元素
$x_1, \ldots, x_d \in M$，它们对每个
$M/\mathfrak m_iM$ 都构成一组基。最后证明
$x_1, \ldots, x_d$ 是 $M$ 的一组基。
\end{proof}

\noindent
下面这个技术引理将在群胚一章中使用。

\begin{lemma}
\label{algebra-lemma-semi-local-module-basis-in-submodule}
设 $R$ 为极大理想为 $\mathfrak m$、剩余域无限的局部环。
设 $R \to S$ 为环映射，$M$ 为 $S$-模，
$N \subset M$ 为 $R$-子模。假设
\begin{enumerate}
\item $S$ 是半局部的，并且 $\mathfrak mS$ 包含于 $S$ 的
Jacobson 根中；
\item $M$ 是有限自由 $S$-模；
\item $N$ 作为 $S$-模生成 $M$。
\end{enumerate}
则 $N$ 包含 $M$ 的一组 $S$-基。
\end{lemma}

\begin{proof}
假设 $M$ 是秩为 $n$ 的自由模，并设 $I \subset S$ 为 Jacobson 根。
由 Nakayama 引理 \ref{algebra-lemma-NAK}，元素序列
$m_1, \ldots, m_n$ 是 $M$ 的一组基，当且仅当
$\overline{m}_i \in M/IM$ 生成 $M/IM$。
因此可以把 $M$ 替换为 $M/IM$，把 $N$ 替换为
$N/(N \cap IM)$，把 $R$ 替换为 $R/\mathfrak m$，
并把 $S$ 替换为 $S/IS$。在此情形，$S$ 是域的有限积
$S = k_1 \times \ldots \times k_r$，且
$M = k_1^{\oplus n} \times \ldots \times k_r^{\oplus n}$。
$N \subset M$ 作为 $S$-模生成 $M$ 这一事实意味着：
存在 $x_j \in N$，使得某个线性组合
$\sum a_j x_j$（$a_j \in S$）在每个因子
$k_i^{\oplus n}$ 中都有非零分量。
因为 $R = k$ 是无限域，这又意味着存在某个线性组合
$y = \sum c_j x_j$（$c_j \in k$），
它在每个因子中都有非零分量。
因此 $y \in N$ 生成自由直和因子 $Sy \subset M$。
对 $n$ 作归纳，可知结论对 $M/Sy$ 及其子模
$\overline{N} = N/(N \cap Sy)$ 成立。
换言之，存在 $\overline{N}$ 中的
$\overline{y}_2, \ldots, \overline{y}_n$，
它们（自由地）生成 $M/Sy$。
于是 $y, y_2, \ldots, y_n$（自由地）生成 $M$，结论得证。
\end{proof}

\begin{lemma}
\label{algebra-lemma-evaluation-map-iso-finite-projective}
设 $R$ 为环，$L$、$M$、$N$ 为 $R$-模。若 $M$ 有限投射，
则典范映射
$$
\Hom_R(M, N) \otimes_R L \to \Hom_R(M, N \otimes_R L)
$$
是同构。
\end{lemma}

\begin{proof}
由引理 \ref{algebra-lemma-finite-projective}，$M$ 既有限表现又有限局部自由。
由引理 \ref{algebra-lemma-hom-from-finitely-presented} 与
\ref{algebra-lemma-tensor-product-localization}，箭头两边的构造都与局部化交换。
由引理 \ref{algebra-lemma-cover}，可在局部化后检验该映射是否为同构。
因此可以假设 $M$ 有限自由；此时引理立即成立。
\end{proof}

\section{由模映射定义的开集}
\label{algebra-section-loci-maps}

\noindent
给定模映射为满射或同构的素理想所成集合有时是开的。
在有限投射模的情形，可以考察映射的秩。

\begin{lemma}
\label{algebra-lemma-map-between-finite}
设 $R$ 为环，$\varphi : M \to N$ 为 $R$-模映射，
其中 $N$ 是有限 $R$-模。则有等式
\begin{align*}
U & = \{\mathfrak p \subset R \mid
\varphi_{\mathfrak p} : M_{\mathfrak p} \to N_{\mathfrak p}
\text{ 为满射}\} \\
& = \{\mathfrak p \subset R \mid
\varphi \otimes \kappa(\mathfrak p) :
M \otimes \kappa(\mathfrak p) \to N \otimes \kappa(\mathfrak p)
\text{ 为满射}\}
\end{align*}
并且 $U$ 是 $\Spec(R)$ 的开子集。此外，对任意满足
$D(f) \subset U$ 的 $f \in R$，映射 $M_f \to N_f$ 都是满射。
\end{lemma}

\begin{proof}
显示公式中的等式来自 Nakayama 引理。Nakayama 引理也蕴含
$U$ 是开的；参见引理 \ref{algebra-lemma-NAK}，特别是其中 (3)。
若 $D(f) \subset U$，则 $M_f \to N_f$ 在 $R_f$ 的所有素理想处
局部化后都是满射，因而由引理
\ref{algebra-lemma-characterize-zero-local}，它本身是满射。
\end{proof}

\begin{lemma}
\label{algebra-lemma-map-between-finitely-presented}
设 $R$ 为环，$\varphi : M \to N$ 为 $R$-模映射，
其中 $M$ 有限而 $N$ 有限表现。则
$$
U = \{\mathfrak p \subset R \mid
\varphi_{\mathfrak p} : M_{\mathfrak p} \to N_{\mathfrak p}
\text{ 为同构}\}
$$
是 $\Spec(R)$ 的开子集。
\end{lemma}

\begin{proof}
设 $\mathfrak p \in U$。选取表现
$N = R^{\oplus n}/\sum_{j = 1, \ldots, m} R k_j$。
以 $e_i$ 表示 $R^{\oplus n}$ 的第 $i$ 个基向量在 $N$ 中的像。
对每个 $i \in \{1, \ldots, n\}$，选取元素
$m_i \in M_{\mathfrak p}$，使得对某个
$f_i \in R$、$f_i \not \in \mathfrak p$ 有
$\varphi(m_i) = f_i e_i$。这可以做到，因为
$\varphi_{\mathfrak p}$ 是同构。令 $f = f_1 \ldots f_n$，
并令 $\psi : R_f^{\oplus n} \to M_f$ 为把第 $i$ 个基向量
映到 $m_i/f_i$ 的映射。注意，$\varphi_f \circ \psi$
是给定映射 $R^{\oplus n} \to N$ 在 $f$ 处的局部化。
因为 $\varphi_{\mathfrak p}$ 是同构，$\psi(k_j)$ 是 $M$ 中
映到 $M_{\mathfrak p}$ 中零元的元素。因此存在
$g_j \in R$、$g_j \not \in \mathfrak p$，使得
$g_j \psi(k_j) = 0$。令 $g = g_1 \ldots g_m$，
则 $\psi_g$ 经 $N_{fg}$ 分解，给出映射
$\chi : N_{fg} \to M_{fg}$。按构造，$\chi$ 是
$\varphi_{fg}$ 的右逆。于是 $\chi_\mathfrak p$ 是同构。
由引理 \ref{algebra-lemma-map-between-finite}，存在
$h \in R$、$h \not \in \mathfrak p$，使映射
$\chi_h : N_{fgh} \to M_{fgh}$ 为满射。
故 $\varphi_{fgh}$ 与 $\chi_h$ 互为逆映射，
这蕴含 $D(fgh) \subset U$，如所需。
\end{proof}

\begin{lemma}
\label{algebra-lemma-finitely-presented-localization-free}
设 $R$ 为环，$\mathfrak p \subset R$ 为素理想，
$M$ 为有限表现 $R$-模。若 $M_\mathfrak p$ 自由，
则存在 $f \in R$、$f \not \in \mathfrak p$，
使得 $M_f$ 是自由 $R_f$-模。
\end{lemma}

\begin{proof}
选取一组基 $x_1, \ldots, x_n \in M_\mathfrak p$。
可以选取 $f \in R$、$f \not \in \mathfrak p$，
使每个 $x_i$ 都是某个 $y_i \in M_f$ 的像。
用 $f^m y_i$（$m \gg 0$）替换 $y_i$ 后，可以假设
$y_i \in M$。事实上，这会把 $x_1, \ldots, x_n$ 替换为
$f^mx_1, \ldots, f^mx_n$；因为 $f$ 在 $R_\mathfrak p$
中的像是单位，它们仍是一组基。因此得到 $R$-模同态
$\varphi = (y_1, \ldots, y_n) : R^{\oplus n} \to M$，
其在 $\mathfrak p$ 处的局部化是同构。
由引理 \ref{algebra-lemma-map-between-finitely-presented}，
可以找到 $f \in R$、$f \not \in \mathfrak p$，使得对所有满足
$f \not \in \mathfrak q$ 的素理想 $\mathfrak q \subset R$，
$\varphi_\mathfrak q$ 都是同构。
于是由引理 \ref{algebra-lemma-characterize-zero-local}，
$\varphi_f$ 是同构，证明完成。
\end{proof}

\begin{lemma}
\label{algebra-lemma-cokernel-flat}
设 $R$ 为环，$\varphi : P_1 \to P_2$ 为有限投射模之间的映射。则
\begin{enumerate}
\item 使 $\varphi \otimes \kappa(\mathfrak p)$ 为单射的素理想
$\mathfrak p \in \Spec(R)$ 所成集合 $U$ 是开的；并且对任意满足
$D(f) \subset U$ 的 $f\in R$，有
\begin{enumerate}
\item $P_{1, f} \to P_{2, f}$ 是单射；
\item 模 $\Coker(\varphi)_f$ 在 $R_f$ 上有限投射。
\end{enumerate}
\item 使 $\varphi \otimes \kappa(\mathfrak p)$ 为满射的素理想
$\mathfrak p \in \Spec(R)$ 所成集合 $W$ 是开的；并且对任意满足
$D(f) \subset W$ 的 $f\in R$，有
\begin{enumerate}
\item $P_{1, f} \to P_{2, f}$ 是满射；
\item 模 $\Ker(\varphi)_f$ 在 $R_f$ 上有限投射。
\end{enumerate}
\item 使 $\varphi \otimes \kappa(\mathfrak p)$ 为同构的素理想
$\mathfrak p \in \Spec(R)$ 所成集合 $V$ 是开的；并且对任意满足
$D(f) \subset V$ 的 $f\in R$，映射
$\varphi : P_{1, f} \to P_{2, f}$ 是 $R_f$-模的同构。
\end{enumerate}
\end{lemma}

\begin{proof}
为证明集合 $U$ 是开的，可以在 $\Spec(R)$ 上局部工作。
因此可把 $R$ 替换为适当的局部化，并假设
$P_1 = R^{n_1}$ 且 $P_2 = R^{n_2}$；参见引理
\ref{algebra-lemma-finite-projective}。在此情形，
$\varphi \otimes \kappa(\mathfrak p)$ 为单射，当且仅当
$n_1 \leq n_2$，并且 $\varphi$ 的矩阵的某个
$n_1 \times n_1$ 子式 $f$ 在 $\kappa(\mathfrak p)$ 中可逆。
因此 $D(f) \subset U$。这个论证也说明，对
$\mathfrak p \in U$，映射
$P_{1, \mathfrak p} \to P_{2, \mathfrak p}$ 是单射。

\medskip\noindent
现在假设 $D(f) \subset U$。由上一段的说明以及引理
\ref{algebra-lemma-characterize-zero-local}，
$P_{1, f} \to P_{2, f}$ 是单射，即 (1)(a) 成立。
由引理 \ref{algebra-lemma-finite-projective}，为证明 (1)(b)，
只需证明 $\Coker(\varphi)$ 在 $D(f)$ 上局部有限投射。
因此如上所见，可以假设 $P_1 = R^{n_1}$、$P_2 = R^{n_2}$，
并且 $\varphi$ 的矩阵的某个子式在 $R$ 中可逆。
若该子式对应于 $R^{n_2}$ 的前 $n_1$ 个基向量，
则利用其余 $n_2 - n_1$ 个基向量得到映射
$R^{n_2 - n_1} \to R^{n_2} \to \Coker(\varphi)$，
容易看出它是同构。

\medskip\noindent
$W$ 的开性以及在 $D(f) \subset W$ 时的 (2)(a)
来自引理 \ref{algebra-lemma-map-between-finite}。
因为 $P_{2, f}$ 在 $R_f$ 上投射，映射
$\varphi_f : P_{1, f} \to P_{2, f}$ 有一个截面，
从而 $\Ker(\varphi)_f$ 是 $P_{2, f}$ 的直和因子。
因此 $\Ker(\varphi)_f$ 有限投射，故 (2)(b) 也成立。

\medskip\noindent
显然 $V = U \cap W$ 是开的，而 (3) 中的另一个断言
来自 (1)(a) 和 (2)(a)。
\end{proof}

\section{模投射性的忠实平坦下降}
\label{algebra-section-ffdescent-projectivity-introduction}

\medskip\noindent
接下来的几节依照 Raynaud 和 Gruson \cite{GruRay}，
证明模的投射性沿忠实平坦环映射下降。
证明思路是采用 Kaplansky 式的 d\'evissage（逐层剖析）\cite{Kaplansky}，
把问题约化到可数生成模的情形。给定模 $M$ 的一个性质良好的滤过，
d\'evissage 使我们可以把 $M$ 表为各滤过子模的相继商之直和
（参见第 \ref{algebra-section-transfinite-devissage} 节）。
利用这一技巧，我们证明投射模是可数生成模的直和
（定理 \ref{algebra-theorem-projective-direct-sum}）。
为证明可数生成模投射性的下降，我们引入模上的一个
“Mittag-Leffler”条件，证明可数生成模投射当且仅当它既平坦又是
Mittag-Leffler 模（定理 \ref{algebra-theorem-projectivity-characterization}），
然后证明成为 Mittag-Leffler 模这一性质可以下降
（引理 \ref{algebra-lemma-ffdescent-ML}）。最后，给定任意模 $M$，
若它沿某个忠实平坦环映射的基变换是投射模，
我们就用相继商为可数生成投射模的子模来滤过 $M$，
再由 d\'evissage 断定 $M$ 是投射模的直和，
因而自身也是投射模（定理 \ref{algebra-theorem-ffdescent-projectivity}）。

\medskip\noindent
我们指出，\cite{GruRay} 中投射性的忠实平坦下降之证明有一个错误。
该文由投射性沿一类更一般的环映射的下降，推出其沿忠实平坦环映射的下降
（\cite[Example 3.1.4(1) of Part II]{GruRay}）。
然而，沿这类更一般映射下降的证明并不正确。
在 \cite{G} 中，Gruson 解释了问题所在，但没有为这里关心的情形给出修补。
填补投射性忠实平坦下降证明中的这个缺口，归结为证明
Mittag-Leffler 模这一性质沿忠实平坦环映射下降。
我们将在引理 \ref{algebra-lemma-ffdescent-ML} 中完成这一点。

\section{平坦性的刻画}
\label{algebra-section-characterize-flatness}

\noindent
本节讨论平坦性的判据。本节的主要结果是下面的 Lazard 定理
（定理 \ref{algebra-theorem-lazard}）：平坦模是有限自由模所成有向系的余极限。
我们提醒读者回顾“平坦性的方程判据”，参见引理
\ref{algebra-lemma-flat-eq}。事实表明，可以把它整理成一个看似强得多的性质。

\begin{lemma}
\label{algebra-lemma-flat-factors-free}
设 $M$ 为 $R$-模。下列条件等价：
\begin{enumerate}
\item $M$ 平坦。
\item 若 $f: R^n \to M$ 是模映射且 $x \in \Ker(f)$，
则存在模映射 $h: R^n \to R^m$ 和 $g: R^m \to M$，
使得 $f = g \circ h$ 且 $x \in \Ker(h)$。
\item 假设 $f: R^n \to M$ 是模映射，
$N \subset \Ker(f)$ 是任意子模，而 $h: R^n \to R^{m}$
是满足 $N \subset \Ker(h)$ 且 $f$ 经 $h$ 分解的映射。
则给定任意 $x \in \Ker(f)$，都可找到映射
$h': R^n \to R^{m'}$，使得
$N + Rx \subset \Ker(h')$ 且 $f$ 经 $h'$ 分解。
\item 若 $f: R^n \to M$ 是模映射，而
$N \subset \Ker(f)$ 是有限生成子模，则存在模映射
$h: R^n \to R^m$ 和 $g: R^m \to M$，
使得 $f = g \circ h$ 且 $N \subset \Ker(h)$。
\end{enumerate}
\end{lemma}

\begin{proof}
(1) 与 (2) 的等价只是平坦性方程判据的另一种表述\footnote{
事实上，模映射 $f : R^n \to M$ 对应于选取 $M$ 中的元素
$x_1, x_2, \ldots, x_n$（即标准基向量
$e_1, e_2, \ldots, e_n$ 的像）；此外，元素
$x \in \Ker(f)$ 对应于这些 $x_1, x_2, \ldots, x_n$
之间的一个关系（即关系 $\sum_i f_i x_i = 0$，其中
$f_i$ 是 $x$ 的坐标）。模映射 $h$（表示成一个
$m \times n$ 矩阵）对应于引理 \ref{algebra-lemma-flat-eq} 中的矩阵
$(a_{ij})$，而引理 \ref{algebra-lemma-flat-eq} 中的 $y_j$
是 $R^m$ 的标准基向量在 $g$ 下的像。}。
为证明 (2) 蕴含 (3)，设 $g: R^m \to M$ 是使
$f$ 分解为 $f = g \circ h$ 的映射。
由 (2)，可找到 $h'': R^m \to R^{m'}$，使得
$h''$ 消去 $h(x)$，且 $g: R^m \to M$ 经 $h''$ 分解。
此时取 $h' = h'' \circ h$ 即可。
(3) 蕴含 (4) 可对 (4) 中 $N \subset \Ker(f)$ 的生成元个数作归纳得到。
显然 (4) 蕴含 (2)。
\end{proof}

\begin{lemma}
\label{algebra-lemma-flat-factors-fp}
设 $M$ 为 $R$-模。则 $M$ 平坦，当且仅当下列条件成立：
若 $P$ 是有限表现 $R$-模，$f: P \to M$ 是模映射，
则存在有限自由 $R$-模 $F$ 以及模映射
$h: P \to F$ 和 $g: F \to M$，使得 $f = g \circ h$。
\end{lemma}

\begin{proof}
这只是引理 \ref{algebra-lemma-flat-factors-free} 中条件 (4) 的另一种表述。
\end{proof}

\begin{lemma}
\label{algebra-lemma-flat-surjective-hom}
设 $M$ 为 $R$-模。则 $M$ 平坦，当且仅当下列条件成立：
对每个有限表现 $R$-模 $P$，若 $N \to M$ 是满的 $R$-模映射，
则诱导映射 $\Hom_R(P, N) \to \Hom_R(P, M)$ 是满射。
\end{lemma}

\begin{proof}
先假设 $M$ 平坦。我们必须证明：若 $P$ 有限表现，
则给定映射 $f: P \to M$，它经映射 $N \to M$ 分解。
由引理 \ref{algebra-lemma-flat-factors-fp}，$f$ 经某个映射
$F \to M$ 分解，其中 $F$ 有限自由。
因为 $F$ 自由，这个映射经 $N \to M$ 分解。
所以 $f$ 经 $N \to M$ 分解。

\medskip\noindent
反过来，假设本引理的条件成立。
设 $f: P \to M$ 是从有限表现模 $P$ 出发的映射。
选取自由模 $N$ 以及到 $M$ 上的满射 $N \to M$。
则 $f$ 经 $N \to M$ 分解；又因为 $P$ 有限生成，
$f$ 经 $N$ 的某个有限自由子模分解。
因此 $M$ 满足引理 \ref{algebra-lemma-flat-factors-fp} 的条件，故平坦。
\end{proof}

\begin{theorem}[Lazard 定理]
\label{algebra-theorem-lazard}
设 $M$ 为 $R$-模。则 $M$ 平坦，当且仅当它是有限自由
$R$-模所成有向系的余极限。
\end{theorem}

\begin{proof}
平坦模所成有向系的余极限平坦，因为取有向余极限正合且与张量积交换。
因此若 $M$ 是有限自由模所成有向系的余极限，则 $M$ 平坦。

\medskip\noindent
为证明反向蕴含，先回顾任意模 $M$ 都可以如下写成有限表现模所成
有向系的余极限。对某个集合 $I$ 选取满射
$f: R^I \to M$，并令 $K$ 为其核。
令 $E$ 为有序对 $(J, N)$ 所成集合，其中 $J$ 是 $I$ 的有限子集，
$N$ 是 $R^J \cap K$ 的有限生成子模。
规定 $(J, N) \leq (J', N')$ 当且仅当
$J \subset J'$ 且 $N \subset N'$，从而把 $E$ 变成有向偏序集。
对 $e = (J, N)$ 定义 $M_e = R^J/N$；当 $e \leq e'$ 时，
定义 $f_{ee'}: M_e \to M_{e'}$ 为自然映射。
于是 $(M_e, f_{ee'})$ 是有向系，并且自然映射
$f_e: M_e \to M$ 诱导同构
$\colim_{e
\in E} M_e \xrightarrow{\cong} M$。

\medskip\noindent
现在假设 $M$ 平坦。令 $I = M \times \mathbf{Z}$，
以 $(x_i)$ 表示 $R^{I}$ 的典范基；在上面的讨论中，
取 $f: R^I \to M$ 为把 $x_i$ 映到 $i$ 在 $M$ 上的投影的映射。
为证明定理，只需说明满足 $M_e$ 自由的 $e \in E$
构成 $E$ 的共尾子集。任取 $e = (J, N) \in E$。
由引理 \ref{algebra-lemma-flat-factors-fp}，存在有限自由模 $F$ 以及映射
$h: R^J/N \to F$ 和 $g: F \to M$，使自然映射
$f_e: R^J/N \to M$ 分解为
$R^J/N \xrightarrow{h} F \xrightarrow{g} M$。
我们将把 $F$ 实现为某个 $e' \geq e$ 的 $M_{e'}$。

\medskip\noindent
设 $\{ b_1, \ldots, b_n \}$ 是 $F$ 的一组有限基。
选取 $n$ 个互异元素 $i_1, \ldots, i_n \in I$，使得对所有
$\ell$ 都有 $i_{\ell} \notin J$，并使 $x_{i_{\ell}}$
在 $f: R^I \to M$ 下的像等于 $b_{\ell}$ 在
$g: F \to M$ 下的像。这可以做到，因为依我们对 $I$ 的选取，
$M$ 的每个元素都可对无穷多个互异的 $i \in I$ 写成 $f(x_i)$。
现在令 $J' = J \cup \{i_1, \ldots , i_n \}$，并如下定义
$R^{J'} \to F$：当 $i \in J$ 时令 $x_i \mapsto h(x_i)$，
而当 $\ell = 1, \ldots, n$ 时令
$x_{i_{\ell}} \mapsto b_{\ell}$。令
$N' = \Ker(R^{J'} \to F)$。注意：
\begin{enumerate}
\item 方块
$$
\xymatrix{
R^{J'} \ar[r] \ar@{^{(}->}[d] & F \ar[d]^{g} \\
R^{I} \ar[r]_{f} & M
}
$$
交换，
%$R^{J'} \to F$ factors $f: R^I \to M$,
因此 $N' \subset K = \Ker(f)$；
\item $R^{J'} \to F$ 是到有限自由模上的满射，因而分裂，
所以 $N'$ 有限生成；
\item $J \subset J'$ 且 $N \subset N'$。
\end{enumerate}
由 (1) 和 (2)，$e' = (J', N')$ 属于 $E$；
由 (3)，$e' \geq e$；而按构造
$M_{e'} = R^{J'}/N' \cong F$ 自由。
\end{proof}

\section{普遍单射模映射}
\label{algebra-section-universally-injective}

\noindent
下面讨论普遍单射模映射；在某种意义上，它们与平坦模相互补充
（见引理 \ref{algebra-lemma-flat-universally-injective}）。
我们遵循 Lazard 的学位论文 \cite{Autour}；另见 \cite{Lam}。

\begin{definition}
\label{algebra-definition-universally-injective}
设 $f: M \to N$ 是 $R$-模映射。若对每个 $R$-模 $Q$，映射
$f \otimes_R \text{id}_Q: M \otimes_R Q \to N \otimes_R Q$
都是单射，则称 $f$ 为{\it 普遍单射}。若 $R$-模序列
$0 \to M_1 \to M_2 \to M_3 \to 0$ 正合，且
$M_1 \to M_2$ 普遍单射，则称该序列{\it 普遍正合}。
\end{definition}

\begin{example}
\label{algebra-example-universally-exact}
普遍正合序列的例子如下。
\begin{enumerate}
\item 分裂短正合序列是普遍正合的，因为张量运算与取直和可交换。
\item 普遍正合序列的有向系之余极限仍然普遍正合。这是因为取有向余极限
是正合的，而张量运算与取余极限可交换。特别地，分裂正合序列的有向系
之余极限是普遍正合的。下面将看到，反过来，每个普遍正合序列都以这种
方式产生。
\end{enumerate}
\end{example}

\noindent
下面给出短正合序列普遍正合的若干判据。它们类似于上面给出的平坦性判据。
以下 (3)--(6) 分别对应引理 \ref{algebra-lemma-flat-eq}、
\ref{algebra-lemma-flat-factors-free}、\ref{algebra-lemma-flat-surjective-hom}
和定理 \ref{algebra-theorem-lazard} 中的平坦性判据。

\begin{theorem}
\label{algebra-theorem-universally-exact-criteria}
设
$$
0 \to M_1 \xrightarrow{f_1} M_2 \xrightarrow{f_2} M_3 \to 0
$$
是 $R$-模的正合序列。以下条件等价：
\begin{enumerate}
\item 序列 $0 \to M_1 \to M_2 \to M_3
\to 0$ 普遍正合。
\item 对每个有限表示的 $R$-模 $Q$，序列
$$
0 \to M_1 \otimes_R Q \to M_2 \otimes_R Q \to
M_3 \otimes_R Q \to 0
$$
正合。
\item 给定元素 $x_i \in M_1$ $(i = 1, \ldots, n)$、
$y_j \in M_2$ $(j = 1, \ldots, m)$ 和
$a_{ij} \in R$ $(i = 1, \ldots, n, j = 1, \ldots, m)$，并且对所有 $i$ 都有
$$
f_1(x_i) = \sum\nolimits_j a_{ij} y_j,
$$
则存在 $z_j \in M_1$ $(j =1, \ldots, m)$，使得对所有 $i$ 都有
$$
x_i = \sum\nolimits_j a_{ij} z_j .
$$
\item 给定 $R$-模映射的交换图
$$
\xymatrix{
R^n \ar[r] \ar[d] &  R^m \ar[d] \\
M_1 \ar[r]^{f_1}        &  M_2
}
$$
其中 $m$ 和 $n$ 是整数，则存在映射 $R^m \to M_1$，使上方三角形交换。
\item 对每个有限表示的 $R$-模 $P$，$R$-模映射
$\Hom_R(P, M_2) \to \Hom_R(P, M_3)$ 都是满射。
\item 序列 $0 \to M_1 \to M_2 \to M_3
\to 0$ 是如下形式的分裂正合序列之某个有向系的余极限：
$$
0 \to M_{1} \to M_{2, i} \to M_{3, i} \to 0
$$
其中各 $M_{3, i}$ 都是有限表示的。
\end{enumerate}
\end{theorem}

\begin{proof}
显然 (1) 蕴含 (2)。

\medskip\noindent
下面证明 (2) 蕴含 (3)。设 $f_1(x_i) = \sum_j a_{ij} y_j$ 是 (3)
中的关系。设 $(d_j)$ 是 $R^m$ 的一组基，$(e_i)$ 是 $R^n$ 的一组基，
并令映射 $R^m \to R^n$ 由
$d_j \mapsto \sum_i a_{ij} e_i$ 给出。令 $Q$ 为
$R^m \to R^n$ 的余核。将 $R^m \to R^n \to Q \to 0$
分别张量成上下两行，并用映射 $f_1: M_1 \to M_2$
连接所得序列，得到交换图
$$
\xymatrix{
M_1^{\oplus m} \ar[r] \ar[d] & M_1^{\oplus n} \ar[r] \ar[d] & M_1 \otimes_R Q
\ar[r] \ar[d] & 0 \\
M_2^{\oplus m} \ar[r] & M_2^{\oplus n} \ar[r] & M_2 \otimes_R Q \ar[r] & 0
}
$$
其中 $M_1^{\oplus m} \to M_1^{\oplus n}$ 由
$$
(z_1, \ldots, z_m) \mapsto
(\sum\nolimits_j a_{1j} z_j, \ldots, \sum\nolimits_j a_{nj} z_j),
$$
给出，而 $M_2^{\oplus m} \to M_2^{\oplus n}$ 的定义相同。
我们要证明 $x = (x_1, \ldots, x_n) \in M_1^{\oplus n}$
属于 $M_1^{\oplus m} \to M_1^{\oplus n}$ 的像。由 (2)，映射
$M_1 \otimes Q \to M_2 \otimes Q$ 是单射，故由上行的正合性，
只须证明 $x$ 在 $M_2 \otimes Q$ 中的像为 $0$；再由下行的正合性，
只须证明 $x$ 在 $M_2^{\oplus n}$ 中的像属于
$M_2^{\oplus m} \to M_2^{\oplus n}$ 的像。这由假设成立。

\medskip\noindent
条件 (4) 只是把 (3) 改写成图的形式。

\medskip\noindent
下面证明 (4) 蕴含 (5)。设 $\varphi : P \to M_3$ 是从有限表示的
$R$-模 $P$ 出发的映射。我们必须证明 $\varphi$ 可提升为映射
$P \to M_2$。选取 $P$ 的一个表示
$$
R^n \xrightarrow{g_1} R^m \xrightarrow{g_2} P \to 0.
$$
利用 $R^n$ 和 $R^m$ 的自由性，可以先构造 $h_2: R^m \to M_2$，
再构造 $h_1: R^n \to M_1$，使下图交换：
$$
\xymatrix{
         & R^n \ar[r]^{g_1} \ar[d]^{h_1} & R^m \ar[r]^{g_2} \ar[d]^{h_2} & P
\ar[r] \ar[d]^{\varphi} & 0 \\
0 \ar[r] & M_1 \ar[r]^{f_1} & M_2 \ar[r]^{f_2} & M_3 \ar[r] & 0 .
}
$$
由 (4)，存在映射 $k_1: R^m \to M_1$，使得
$k_1 \circ g_1 = h_1$。现在定义 $h'_2: R^m \to M_2$ 为
$h_2' = h_2 - f_1 \circ k_1$。于是
$$
h'_2 \circ g_1 = h_2 \circ g_1 - f_1 \circ k_1 \circ g_1 =
h_2 \circ g_1 - f_1 \circ h_1 = 0 .
$$
因此取商后，$h'_2$ 定义映射 $\varphi': P \to M_2$，满足
$\varphi' \circ g_2 = h_2'$。用图表示即为
$$
\xymatrix{
R^m \ar[r]^{g_2} \ar[d]_{h'_2} & P \ar[d]^{\varphi} \ar[dl]_{\varphi'} \\
M_2 \ar[r]^{f_2} & M_3.
}
$$
其中上方三角形交换。我们断言 $\varphi'$ 就是所求提升，即
$f_2 \circ \varphi' = \varphi$。由定义，
$$
f_2 \circ \varphi' \circ g_2 = f_2 \circ h'_2 =
f_2 \circ h_2 - f_2 \circ f_1 \circ k_1 = f_2 \circ h_2 =
\varphi \circ g_2.
$$
由于 $g_2$ 是满射，证明完成。

\medskip\noindent
现在证明 (5) 蕴含 (6)。把 $M_{3}$ 写成有限表示模
$M_{3, i}$ 的一个有向系的余极限，见引理
\ref{algebra-lemma-module-colimit-fp}。令 $M_{2, i}$ 为 $M_{3, i}$ 与
$M_{2}$ 在 $M_{3}$ 上的纤维积；按定义，它是 $M_2 \times M_{3, i}$
中由两个投影在 $M_3$ 中的像相等的元素组成的子模。令
$M_{1, i}$ 为投影 $M_{2, i} \to M_{3, i}$ 的核。
于是得到正合序列的有向系
$$
0 \to M_{1, i} \to M_{2, i} \to M_{3, i} \to 0,
$$
并且对每个 $i$ 都有正合序列的映射
$$
\xymatrix{
0  \ar[r] & M_{1, i} \ar[d] \ar[r]  &  M_{2, i}  \ar[r] \ar[d] & M_{3, i} \ar[d]
\ar[r] & 0 \\
0  \ar[r] & M_{1} \ar[r]  &  M_{2}  \ar[r] & M_{3} \ar[r] & 0
}
$$
且这些映射与有向系相容。由纤维积 $M_{2, i}$ 的定义可知，
映射 $M_{1, i} \to M_1$ 是同构。由 (5)，存在提升
$M_{3, i} \to M_{2}$ 的映射 $M_{3, i} \to M_3$；再由纤维积的
泛性质，这给出 $M_{2, i} \to M_{3, i}$ 的一个截面。因此序列
$$
0 \to M_{1, i} \to M_{2, i} \to M_{3, i} \to 0
$$
分裂。取余极限，得到交换图
$$
\xymatrix{
0  \ar[r] & \colim M_{1, i} \ar[d]^{\cong} \ar[r]  &  \colim M_{2, i}
 \ar[r]
\ar[d] & \colim M_{3, i} \ar[d]^{\cong} \ar[r] & 0 \\
0  \ar[r] & M_{1} \ar[r]  &  M_{2}  \ar[r] & M_{3} \ar[r] & 0
}
$$
其两行正合，且外侧竖直映射均为同构。因此
$\colim M_{2, i} \to M_2$ 也是同构，故 (6) 成立。

\medskip\noindent
由例 \ref{algebra-example-universally-exact} (2)，条件 (6) 蕴含 (1)。
\end{proof}

\noindent
上一条定理表明，普遍正合序列总是分裂短正合序列的余极限。
若一个普遍单射映射的余核是有限表示的，则该映射本身事实上分裂：

\begin{lemma}
\label{algebra-lemma-universally-exact-split}
设
$$
0 \to M_1 \to M_2 \to M_3 \to 0
$$
是 $R$-模的正合序列。假设 $M_3$ 是有限表示的。那么
$$
0 \to M_1 \to M_2 \to M_3 \to 0
$$
普遍正合，当且仅当它分裂。
\end{lemma}

\begin{proof}
分裂短正合序列总是普遍正合的，见例
\ref{algebra-example-universally-exact}。
反过来，若该序列普遍正合，则把定理
\ref{algebra-theorem-universally-exact-criteria} (5)
应用于 $P = M_3$，可知映射 $M_2 \to M_3$ 有截面。
\end{proof}

\noindent
下一个引理说明普遍单射映射如何与平坦模相互补充。

\begin{lemma}
\label{algebra-lemma-flat-universally-injective}
设 $M$ 是 $R$-模。那么 $M$ 平坦，当且仅当每个 $R$-模正合序列
$$
0 \to M_1 \to M_2 \to M \to 0
$$
都普遍正合。
\end{lemma}

\begin{proof}
这由引理 \ref{algebra-lemma-flat-surjective-hom} 和定理
\ref{algebra-theorem-universally-exact-criteria} (5) 得出。
\end{proof}

\begin{example}
\label{algebra-example-universally-exact-non-split-non-flat}
非分裂及非平坦的普遍正合序列。
\begin{enumerate}
\item 尽管有引理 \ref{algebra-lemma-universally-exact-split}，仍可能存在
$R$-模的短正合序列
$$
0 \to M_1 \to M_2 \to M_3 \to 0
$$
是普遍正合但非分裂的。例如，取 $R = \mathbf{Z}$，
令 $M_1 = \bigoplus_{n=1}^{\infty} \mathbf{Z}$，
令 $M_{2} = \prod_{n = 1}^{\infty} \mathbf{Z}$，并令
$M_{3}$ 为嵌入 $M_1 \to M_2$ 的余核。由于它们都无挠，
$M_1, M_2, M_3$ 都平坦（《代数进阶》，引理
\ref{more-algebra-lemma-dedekind-torsion-free-flat}），
故由引理 \ref{algebra-lemma-flat-universally-injective}，序列
$$
0 \to M_1 \to M_2 \to M_3 \to 0
$$
普遍正合。然而不存在截面 $s: M_3 \to M_2$。
事实上，若 $x$ 是 $(2, 2^2, 2^3, \ldots) \in M_2$
在 $M_3$ 中的像，则每个模映射 $s: M_3 \to M_2$ 都必须把
$x$ 映为零。这是因为对每个 $n \geq 1$ 都有
$x \in 2^n M_3$，所以对所有 $n \geq 1$，$s(x)$ 都可被
$2^n$ 整除，因而只能是 $0$。
\item 尽管有引理 \ref{algebra-lemma-flat-universally-injective}，仍可能存在
$R$-模的短正合序列
$$
0 \to M_1 \to M_2 \to M_3 \to 0
$$
是普遍正合的，但 $M_1, M_2, M_3$ 全都不平坦。事实上，若 $M$
是任意非平坦模，只须取分裂正合序列
$$
0 \to M \to M \oplus M \to M \to 0.
$$
例如在 $R = \mathbf{Z}$ 上，可取 $M$ 为任意挠模。
\item 将 (1) 中的一个正合序列与 (2) 中的一个正合序列取直和，
便得到 $R$-模的短正合序列
$$
0 \to M_1 \to M_2 \to M_3 \to 0
$$
它普遍正合、非分裂，并且 $M_1, M_2, M_3$ 全都不平坦。
\end{enumerate}
\end{example}

\begin{lemma}
\label{algebra-lemma-ui-flat-domain}
设 $0 \to M_1 \to M_2 \to M_3 \to 0$ 是 $R$-模的普遍正合序列，
并假设 $M_2$ 平坦。那么 $M_1$ 和 $M_3$ 都平坦。
\end{lemma}

\begin{proof}
设 $0 \to N \to N' \to N'' \to 0$ 是 $R$-模的短正合序列。
考虑交换图
$$
\xymatrix{
M_1 \otimes_R N \ar[r] \ar[d] &
M_2 \otimes_R N \ar[r] \ar[d] &
M_3 \otimes_R N \ar[d] \\
M_1 \otimes_R N' \ar[r] \ar[d] &
M_2 \otimes_R N' \ar[r] \ar[d] &
M_3 \otimes_R N' \ar[d] \\
M_1 \otimes_R N'' \ar[r] &
M_2 \otimes_R N'' \ar[r] &
M_3 \otimes_R N''
}
$$
（我们省略了边界上的 $0$）。由假设，各行都是短正合序列，并且箭头
$M_2 \otimes N \to M_2 \otimes N'$ 是单射。这显然蕴含
$M_1 \otimes N \to M_1 \otimes N'$ 是单射，故 $M_1$ 平坦。
特别地，左列和中列都给出短正合序列。作一次追图可知，箭头
$M_3 \otimes N \to M_3 \otimes N'$ 是单射。因此 $M_3$ 平坦。
\end{proof}

\begin{lemma}
\label{algebra-lemma-universally-injective-tensor}
设 $R$ 是环，$M \to M'$ 是普遍单射的 $R$-模映射。
那么对任意 $R$-模 $N$，映射
$M \otimes_R N \to M' \otimes_R N$ 都是普遍单射。
\end{lemma}

\begin{proof}
略。
\end{proof}

\begin{lemma}
\label{algebra-lemma-composition-universally-injective}
设 $R$ 是环。普遍单射的 $R$-模映射的复合仍然普遍单射。
\end{lemma}

\begin{proof}
略。
\end{proof}

\begin{lemma}
\label{algebra-lemma-universally-injective-permanence}
设 $R$ 是环，$M \to M'$ 和 $M' \to M''$ 是 $R$-模映射。
若其复合 $M \to M''$ 普遍单射，则 $M \to M'$ 普遍单射。
\end{lemma}

\begin{proof}
略。
\end{proof}

\begin{lemma}
\label{algebra-lemma-universally-injective-check-stalks}
设 $R \to S$ 是环映射，$M \to M'$ 是 $S$-模映射。
以下条件等价：
\begin{enumerate}
\item $M \to M'$ 作为 $R$-模映射普遍单射；
\item 对 $S$ 的每个素理想 $\mathfrak q$，映射
$M_{\mathfrak q} \to M'_{\mathfrak q}$
作为 $R$-模映射普遍单射；
\item 对 $S$ 的每个极大理想 $\mathfrak m$，映射
$M_{\mathfrak m} \to M'_{\mathfrak m}$
作为 $R$-模映射普遍单射；
\item 对 $S$ 的每个素理想 $\mathfrak q$，映射
$M_{\mathfrak q} \to M'_{\mathfrak q}$
作为 $R_{\mathfrak p}$-模映射普遍单射，其中 $\mathfrak p$ 是
$\mathfrak q$ 在 $R$ 中的逆像；
\item 对 $S$ 的每个极大理想 $\mathfrak m$，映射
$M_{\mathfrak m} \to M'_{\mathfrak m}$
作为 $R_{\mathfrak p}$-模映射普遍单射，其中 $\mathfrak p$ 是
$\mathfrak m$ 在 $R$ 中的逆像。
\end{enumerate}
\end{lemma}

\begin{proof}
设 $N$ 是 $R$-模。设 $\mathfrak q$ 是 $S$ 的素理想，位于
$R$ 的素理想 $\mathfrak p$ 之上。则
$$
(M \otimes_R N)_{\mathfrak q} =
M_{\mathfrak q} \otimes_R N =
M_{\mathfrak q} \otimes_{R_{\mathfrak p}} N_{\mathfrak p}.
$$
此外，对 $M'$ 也有同样的等式，而局部化是正合的。再者，若 $N$
是 $R_{\mathfrak p}$-模，则 $N_{\mathfrak p} = N$。利用这些事实，
可以直接证明上述等价性。

\medskip\noindent
例如，假设 (5) 成立。令
$K = \Ker(M \otimes_R N \to M' \otimes_R N)$。由上述说明，
对 $S$ 的每个极大理想 $\mathfrak m$ 都有
$K_{\mathfrak m} = 0$。因此由引理
\ref{algebra-lemma-characterize-zero-local}，$K = 0$，故 (1) 成立。
反过来，假设 (1) 成立。任取位于
$\mathfrak p \subset R$ 之上的 $\mathfrak q \subset S$，
并任取 $R_{\mathfrak p}$ 上的模 $N$。由假设，
$\Ker(M \otimes_R N \to M' \otimes_R N) = 0$。
于是由上述公式以及 $N = N_{\mathfrak p}$ 可知
$$
\Ker(M_{\mathfrak q} \otimes_{R_{\mathfrak p}} N \to
M'_{\mathfrak q} \otimes_{R_{\mathfrak p}} N) = 0
$$
换言之，(4) 成立。当然，(4) $\Rightarrow$ (5) 是立即的。
因此 (1)、(4) 和 (5) 彼此等价。其余等价性的证明从略。
\end{proof}

\begin{lemma}
\label{algebra-lemma-universally-injective-localize}
设 $\varphi : A \to B$ 是环映射。设 $S \subset A$ 和
$S' \subset B$ 是乘法子集，且 $\varphi(S) \subset S'$。
设 $M \to M'$ 是 $B$-模映射。
\begin{enumerate}
\item 若 $M \to M'$ 作为 $A$-模映射普遍单射，则
$(S')^{-1}M \to (S')^{-1}M'$ 作为 $A$-模映射和作为
$S^{-1}A$-模映射都是普遍单射。
\item 若 $M$ 和 $M'$ 是 $(S')^{-1}B$-模，则 $M \to M'$
作为 $A$-模映射普遍单射，当且仅当它作为 $S^{-1}A$-模映射
普遍单射。
\end{enumerate}
\end{lemma}

\begin{proof}
可以用引理 \ref{algebra-lemma-universally-injective-check-stalks} 证明这一点，
也可以如下直接证明。假设 $M \to M'$ 作为 $A$-模映射普遍单射。
设 $Q$ 是 $A$-模。则 $Q \otimes_A M \to Q \otimes_A M'$ 是单射。
因为局部化正合，所以
$(S')^{-1}(Q \otimes_A M) \to (S')^{-1}(Q \otimes_A M')$
是单射。又因为
$(S')^{-1}(Q \otimes_A M) = Q \otimes_A (S')^{-1}M$，
而对 $M'$ 也同样，所以
$Q \otimes_A (S')^{-1}M \to Q \otimes_A (S')^{-1}M'$ 是单射。
因此 $(S')^{-1}M \to (S')^{-1}M'$ 作为 $A$-模映射普遍单射。
这证明了 (1) 的第一部分。
为证明 (2)，可使用以下两个事实：(a) 若 $Q$ 是
$S^{-1}A$-模，则 $Q \otimes_A S^{-1}A = Q$，即在 $A$ 上与
$Q$ 张量和在 $S^{-1}A$ 上与 $Q$ 张量是一回事；(b) 若 $M$
是任意 $A$-模且 $S$ 中的元素在其上都可逆，则
$M \otimes_A Q = M \otimes_{S^{-1}A} S^{-1}Q$。
(2) 立即由此得出。
\end{proof}

\begin{lemma}
\label{algebra-lemma-check-universally-injective-into-flat}
设 $R$ 是环，$M \to M'$ 是 $R$-模映射。若 $M'$ 平坦，则
$M \to M'$ 普遍单射，当且仅当对 $R$ 的每个有限生成理想 $I$，
$M/IM \to M'/IM'$ 都是单射。
\end{lemma}

\begin{proof}
只须证明对每个有限 $R$-模 $Q$，映射
$M \otimes_R Q \to M' \otimes_R Q$ 都是单射，见定理
\ref{algebra-theorem-universally-exact-criteria}。
此时 $Q$ 有一个由子模组成的有限滤过
$0 = Q_0 \subset Q_1 \subset \ldots \subset Q_n = Q$，
其子商同构于循环模 $R/I_i$，见引理
\ref{algebra-lemma-trivial-filter-finite-module}。由于 $M'$ 平坦，
得到 $M' \otimes_R Q$ 的子模 $M' \otimes_R Q_i$ 所成的滤过
$$
\xymatrix{
M \otimes Q_1 \ar[r] \ar[d] &
M \otimes Q_2 \ar[r] \ar[d] &
\ldots \ar[r] &
M \otimes Q \ar[d] \\
M' \otimes Q_1 \ar@{^{(}->}[r] &
M' \otimes Q_2 \ar@{^{(}->}[r] &
\ldots \ar@{^{(}->}[r] &
M' \otimes Q
}
$$
其相继商为 $M' \otimes_R R/I_i = M'/I_iM'$。
简单的归纳表明，只须检验 $M/I_i M \to M'/I_i M'$ 是单射。
注意，有限生成理想 $I'_i \subset I_i$ 的全体是一个有向集。
因此 $M/I_iM = \colim M/I'_iM$ 是滤过余极限，$M'$ 的情形类似；
各映射 $M/I'_iM \to M'/I'_i M'$ 由假设都是单射，而滤过余极限
正合（引理 \ref{algebra-lemma-directed-colimit-exact}），故结论成立。
\end{proof}

\begin{lemma}
\label{algebra-lemma-base-change-along-universally-injective-ring-map}
设 $R \to S$ 是环映射，并且作为 $R$-模映射普遍单射。
那么 $R$-模上的函子 $M \mapsto M \otimes_R S$
反映单射、满射和同构。
\end{lemma}

\begin{proof}
设 $M \to N$ 是 $R$-模映射，其核为 $K$、余核为 $Q$。
若 $M \otimes_R S \to N \otimes_R S$ 是单射，则
$K \otimes_R S \to M \otimes_R S$ 的像为零。由于
$K \subset K \otimes_R S$ 且 $M \subset M \otimes_R S$，
可知 $K = 0$。另一方面，若
$M \otimes_R S \to N \otimes_R S$ 是满射，则由张量积的右正合性，
$Q \otimes_R S$ 为零。因此 $Q \subset Q \otimes_R S$ 也为零。
\end{proof}

\begin{lemma}
\label{algebra-lemma-faithfully-flat-universally-injective}
设 $R \to S$ 是忠实平坦环映射。那么作为 $R$-模映射，
$R \to S$ 普遍单射。特别地，对任意理想 $I \subset R$ 都有
$R \cap IS = I$。
\end{lemma}

\begin{proof}
设 $N$ 是 $R$-模。我们要证明 $N \to N \otimes_R S$ 是单射。
由于 $S$ 作为 $R$-模忠实平坦，只须在与 $S$ 张量后证明这一点。
因此只须证明映射
$N \otimes_R S \to N \otimes_R S \otimes_R S$，
$n \otimes s \mapsto n \otimes 1 \otimes s$ 是单射。
这是成立的，因为它有收缩映射
$n \otimes s \otimes s' \mapsto n \otimes ss'$。
\end{proof}

\section{有限投射模的下降}
\label{algebra-section-finite-projective}

\noindent
本节给出一个初等证明：作为{\it 有限}投射模的性质沿忠实平坦环映射
下降。去掉有限性条件后，这个证明便不再适用。不过，该方法预示了我们
将用来证明{\it 可数生成}投射模之下降的方法；见本节末尾的说明。

\begin{lemma}
\label{algebra-lemma-finite-projective-again}
设 $M$ 是 $R$-模。那么 $M$ 有限投射，当且仅当 $M$ 有限表示且平坦。
\end{lemma}

\begin{proof}
这是引理 \ref{algebra-lemma-finite-projective} 的一部分。不过在此处，我们可以
如下对该引理中的蕴含 (1) $\Rightarrow$ (2) 给出一个更优雅的证明。
若 $M$ 有限表示且平坦，取一个满射 $R^n \to M$。
把引理 \ref{algebra-lemma-flat-surjective-hom} 应用于 $P = M$，
可知映射 $R^n \to M$ 有截面。因此 $M$ 是自由模的直和因子，
从而是投射模。
\end{proof}

\noindent
以下是模的一些可下降性质。

\begin{lemma}
\label{algebra-lemma-descend-properties-modules}
设 $R \to S$ 是忠实平坦环映射，$M$ 是 $R$-模。那么
\begin{enumerate}
\item 若 $S$-模 $M \otimes_R S$ 是有限型的，则 $M$ 是有限型的；
\item 若 $S$-模 $M \otimes_R S$ 是有限表示的，则 $M$ 是有限表示的；
\item 若 $S$-模 $M \otimes_R S$ 平坦，则 $M$ 平坦；
\item 按需在此添加更多性质。
\end{enumerate}
\end{lemma}

\begin{proof}
假设 $M \otimes_R S$ 是有限型的。设 $y_1, \ldots, y_m$ 是
$M \otimes_R S$ 在 $S$ 上的一组生成元。对某些
$x_1, \ldots, x_n \in M$，写
$y_j = \sum x_i \otimes f_i$。于是映射
$\varphi : R^{\oplus n} \to M$ 具有如下性质：
$\varphi \otimes \text{id}_S : S^{\oplus n} \to M \otimes_R S$
是满射。由于 $R \to S$ 忠实平坦，可知 $\varphi$ 是满射，
故 $M$ 有限生成。

\medskip\noindent
假设 $M \otimes_R S$ 是有限表示的。由 (1)，$M$ 是有限型的。
选取满射 $R^{\oplus n} \to M$，并将其核记为 $K$。
由于 $R \to S$ 平坦，$K \otimes_R S$ 是基变换
$S^{\oplus n} \to M \otimes_R S$ 的核。
由于 $M \otimes_R S$ 是有限表示的，可知 $K \otimes_R S$
是有限型的。因此由 (1)，$K$ 是有限型的，从而 $M$ 是有限表示的。

\medskip\noindent
(3) 就是引理 \ref{algebra-lemma-flatness-descends}。
\end{proof}

\begin{proposition}
\label{algebra-proposition-ffdescent-finite-projectivity}
设 $R \to S$ 是忠实平坦环映射，$M$ 是 $R$-模。
若 $S$-模 $M \otimes_R S$ 有限投射，则 $M$ 有限投射。
\end{proposition}

\begin{proof}
这由引理 \ref{algebra-lemma-finite-projective-again} 和
\ref{algebra-lemma-descend-properties-modules} 得出。
\end{proof}

\noindent
下面数节将利用 d\'evissage（逐层剖析）把问题约化到可数生成情形，
以去掉有限性假设。在可数生成情形中，我们的策略是寻找一个类似于
引理 \ref{algebra-lemma-finite-projective-again} 的可数生成投射模刻画，
然后直接证明该刻画可下降。为此，我们引入 Mittag-Leffler 模的概念，
并证明：若模 $M$ 可数生成，则它是投射模，当且仅当它既平坦又是
Mittag-Leffler 模（定理 \ref{algebra-theorem-projectivity-characterization}）。
当 $M$ 有限生成时，该陈述约化为引理
\ref{algebra-lemma-finite-projective-again}，因为根据例
\ref{algebra-example-ML} (1)，有限生成模是 Mittag-Leffler 模，
当且仅当它有限表示。

\section{模的超限 d\'evissage}
\label{algebra-section-transfinite-devissage}

\noindent
本节引入一种把模分解为直和的 d\'evissage（逐层剖析）技术。
主要结果是投射模可以写成可数生成模的直和（见下文定理
\ref{algebra-theorem-projective-direct-sum}）。我们遵循 \cite{Kaplansky}。


\begin{definition}
\label{algebra-definition-devissage}
设 $M$ 是 $R$-模。$M$ 的一个{\it 直和 d\'evissage} 是一族子模
$(M_{\alpha})_{\alpha \in S}$，它由序数 $S$ 标号，并且按包含关系
递增，且满足：
\begin{enumerate}
\item[(0)] $M_0 = 0$；
\item[(1)] $M = \bigcup_{\alpha} M_{\alpha}$；
\item[(2)] 若 $\alpha \in S$ 是极限序数，则
$M_{\alpha} = \bigcup_{\beta < \alpha} M_{\beta}$；
\item[(3)] 若 $\alpha + 1 \in S$，则 $M_{\alpha}$ 是
$M_{\alpha + 1}$ 的直和因子。
\end{enumerate}
若还满足
\begin{enumerate}
\item[(4)] 对 $\alpha + 1 \in S$，$M_{\alpha + 1}/M_{\alpha}$
可数生成，
\end{enumerate}
则称 $(M_{\alpha})_{\alpha \in S}$ 为 $M$ 的一个
{\it Kaplansky d\'evissage}。
\end{definition}

\noindent
下一个引理说明了上述术语的理由。

\begin{lemma}
\label{algebra-lemma-direct-sum-devissage}
设 $M$ 是 $R$-模。若 $(M_{\alpha})_{\alpha \in S}$ 是 $M$ 的
一个直和 d\'evissage，则
$M \cong \bigoplus_{\alpha + 1 \in S} M_{\alpha + 1}/M_{\alpha}$。
\end{lemma}

\begin{proof}
由直和 d\'evissage 的性质 (3)，对每个 $\alpha \in S$，
都有嵌入 $M_{\alpha + 1}/M_{\alpha} \to M$。考虑由这些嵌入之和给出的映射
$$
f : \bigoplus\nolimits_{\alpha + 1\in S} M_{\alpha + 1}/M_{\alpha} \to M
$$
再对 $\beta\in S$ 考虑限制
$$
f_{\beta} :
\bigoplus\nolimits_{\alpha + 1 \leq \beta} M_{\alpha + 1}/M_{\alpha}
\longrightarrow
M
$$
对 $S$ 作超限归纳可知，$f_{\beta}$ 的像为 $M_{\beta}$。
当 $\beta=0$ 时，这由 $(0)$ 成立。若 $\beta+1$ 是后继序数，
且结论对 $\beta$ 成立，则由 (3)，结论对 $\beta + 1$ 成立。
若 $\beta$ 是极限序数，且结论对 $\alpha < \beta$ 成立，
则由 (2)，结论对 $\beta$ 成立。因此由 (1)，$f$ 是满射。

\medskip\noindent
对 $S$ 作超限归纳还可证明各限制 $f_{\beta}$ 都是单射。
当 $\beta = 0$ 时结论成立。若 $\beta+1$ 是后继序数，
且 $f_{\beta}$ 是单射，任取核中的 $x$，并按其分量写成
$x = (x_{\alpha + 1})_{\alpha + 1 \leq \beta + 1}$，其中
$x_{\alpha + 1} \in M_{\alpha + 1}/M_{\alpha}$。
由性质 (3) 以及 $f_{\beta}$ 的像为 $M_{\beta}$ 这一事实，
$(x_{\alpha + 1})_{\alpha + 1 \leq \beta}$ 和
$x_{\beta + 1}$ 都映为 $0$。因此 $x_{\beta+1} = 0$；
再由限制 $f_{\beta}$ 单射的假设，对每个
$\alpha + 1 \leq \beta$ 都有 $x_{\alpha + 1} = 0$。
故 $x = 0$，而 $f_{\beta+1}$ 是单射。
若 $\beta$ 是极限序数，考虑核中的元素 $x$。则 $x$ 已包含在某个
$\alpha < \beta$ 的 $f_{\alpha}$ 的定义域中，故 $x = 0$，
归纳完成。由于对每个 $\beta \in S$，$f_{\beta}$ 都单射，
可知 $f$ 单射。
\end{proof}

\begin{lemma}
\label{algebra-lemma-Kaplansky-devissage}
设 $M$ 是 $R$-模。那么 $M$ 是可数生成 $R$-模的直和，
当且仅当它有 Kaplansky d\'evissage。
\end{lemma}

\begin{proof}
上一引理给出“若”的方向。反过来，假设
$M = \bigoplus_{i \in I} N_i$，其中每个 $N_i$ 都是可数生成
$R$-模。给 $I$ 配备良序，从而可把它视为序数。令
$M_i = \bigoplus_{j < i} N_j$，便得到 $M$ 的一个
Kaplansky d\'evissage $(M_i)_{i \in I}$。
\end{proof}

\begin{theorem}
\label{algebra-theorem-kaplansky-direct-sum}
假设 $M$ 是可数生成 $R$-模的直和。若 $P$ 是 $M$ 的直和因子，
则 $P$ 也是可数生成 $R$-模的直和。
\end{theorem}

\begin{proof}
写成 $M = P \oplus Q$。我们将构造 $M$ 的一个 Kaplansky
d\'evissage $(M_{\alpha})_{\alpha \in S}$；除定义性质 (0)--(4)
外，它还满足：
\begin{enumerate}
\item[(5)] 每个 $M_{\alpha}$ 都是 $M$ 的直和因子；
\item[(6)] $M_{\alpha} = P_{\alpha} \oplus Q_{\alpha}$，其中
$P_{\alpha} =P \cap M_{\alpha}$ 且
$Q_\alpha = Q \cap M_{\alpha}$。
\end{enumerate}
（注意：若性质 (0)--(2) 成立，则性质 (3) 事实上等价于性质 (5)。）

\medskip\noindent
为说明这如何推出定理，只须证明 $(P_{\alpha})_{\alpha \in S}$
构成 $P$ 的一个 Kaplansky d\'evissage。性质 (0)、(1) 和 (2)
是显然的。由 $(M_{\alpha})$ 的 (5) 和 (6)，每个
$P_{\alpha}$ 都是 $M$ 的直和因子。由于
$P_{\alpha} \subset P_{\alpha + 1}$，这蕴含 $P_{\alpha}$ 是
$P_{\alpha + 1}$ 的直和因子；故 $(P_{\alpha})$ 满足 (3)。
关于 (4)，注意
$$
M_{\alpha + 1}/M_{\alpha} \cong P_{\alpha + 1}/P_{\alpha} \oplus
Q_{\alpha + 1}/Q_{\alpha},
$$
因此 $P_{\alpha + 1}/P_{\alpha}$ 可数生成，因为
$M_{\alpha + 1}/M_{\alpha}$ 可数生成。

\medskip\noindent
还需构造诸 $M_{\alpha}$。写成
$M = \bigoplus_{i \in I} N_i$，其中每个 $N_i$ 都是可数生成
$R$-模。选取 $I$ 的一个良序。我们将通过超限递归定义 $M$
的一族递增子模 $M_{\alpha}$，每个序数 $\alpha$ 对应一个，
并使 $M_{\alpha}$ 是某些 $N_i$ 的直和。

\medskip\noindent
当 $\alpha = 0$ 时，令 $M_{0} = 0$。若 $\alpha$ 是极限序数，
且对所有 $\beta < \alpha$ 已定义 $M_{\beta}$，则定义
$M_{\alpha} = \bigcup_{\beta < \alpha} M_{\beta}$。
由于每个 $\beta < \alpha$ 的 $M_{\beta}$ 都是某些 $N_i$ 的直和，
$M_{\alpha}$ 也有同样的性质。若 $\alpha + 1$ 是后继序数，
且已定义 $M_{\alpha}$，则如下定义 $M_{\alpha + 1}$。
若 $M_{\alpha} = M$，令 $M_{\alpha + 1} = M$。
否则，选取使 $N_j$ 不包含于 $M_{\alpha}$ 的最小 $j \in I$。
我们将构造无穷矩阵 $(x_{mn}), m, n = 1, 2, 3, \ldots$，使得：
\begin{enumerate}
\item $N_j$ 包含于由矩阵条目 $x_{mn}$ 生成的 $M$ 的子模；
\item 若把任意条目 $x_{k\ell}$ 按其 $P$-分量和 $Q$-分量写成
$x_{k\ell} = y_{k\ell} + z_{k\ell}$，则对每个 $N_i$，只要
$y_{k\ell}$ 或 $z_{k\ell}$ 在其中的分量非零，矩阵
$(x_{mn})$ 都含有它的一组生成元。
\end{enumerate}
然后定义 $M_{\alpha + 1}$ 为由 $M_{\alpha}$ 和所有 $x_{mn}$
生成的 $M$ 的子模；由矩阵 $(x_{mn})$ 的性质 (2)，
$M_{\alpha + 1}$ 将是某些 $N_i$ 的直和。
为构造矩阵 $(x_{mn})$，令 $x_{11}, x_{12}, x_{13}, \ldots$
为 $N_j$ 的一组可数生成元。然后，若
$x_{11} = y_{11} + z_{11}$ 是按 $P$-分量和 $Q$-分量的分解，
令 $x_{21}, x_{22}, x_{23}, \ldots$ 为如下诸 $N_i$ 之和的一组
可数生成元：$y_{11}$ 或 $z_{11}$ 在其中的分量非零。
对 $x_{12}$ 重复此过程，得到元素
$x_{31}, x_{32}, \ldots$，即矩阵的第三行。
再对 $x_{21}$ 重复，得到第四行；对 $x_{13}$ 重复，得到第五行；
依此类推，按下图所示沿相继的反对角线向下进行：
$$
\left(
\vcenter{
\xymatrix@R=2mm@C=2mm{
x_{11} & x_{12} \ar[dl] & x_{13} \ar[dl] & x_{14} \ar[dl] & \ldots  \\
x_{21} & x_{22} \ar[dl] & x_{23} \ar[dl] & \ldots  \\
x_{31} & x_{32} \ar[dl] & \ldots \\
x_{41} & \ldots \\
\ldots
}
}
\right).
$$

\medskip\noindent
对 $I$ 作超限归纳（使用如下事实：我们构造 $M_{\alpha + 1}$ 时，
令其包含使 $N_j$ 不包含于 $M_{\alpha}$ 的最小 $j$ 所对应的
$N_j$），可知对每个 $i \in I$，$N_i$ 都包含于某个
$M_{\alpha}$。因此存在充分大的序数 $S$，满足：对每个
$i \in I$，都存在 $\alpha \in S$，使 $N_i$ 包含于
$M_{\alpha}$。这意味着 $(M_{\alpha})_{\alpha \in S}$ 满足
$M$ 的 Kaplansky d\'evissage 的性质 (1)。
此外，族 $(M_{\alpha})_{\alpha \in S}$ 还满足其余定义性质以及
上述 (5) 和 (6)：性质 (0)、(2)、(4) 和 (6) 由构造显然成立；
性质 (5) 成立，因为按构造，每个 $M_{\alpha}$ 都是某些 $N_i$
的直和；而由 (5) 以及 $M_{\alpha} \subset M_{\alpha + 1}$，
可推出 (3)。
\end{proof}

\noindent
作为推论，得到本节开头所述的投射模结果。

\begin{theorem}
\label{algebra-theorem-projective-direct-sum}
\begin{slogan}
任意投射模都是可数生成投射模的直和。
\end{slogan}
若 $P$ 是投射 $R$-模，则 $P$ 是可数生成投射 $R$-模的直和。
\end{theorem}

\begin{proof}
一个模是投射模，当且仅当它是自由模的直和因子，
故结论由定理 \ref{algebra-theorem-kaplansky-direct-sum} 得出。
\end{proof}

\section{局部环上的投射模}
\label{algebra-section-projective-local-ring}

\noindent
本节证明一个非常漂亮的结果：局部环上的投射模 $M$ 是自由模
（见下文定理 \ref{algebra-theorem-projective-free-over-local-ring}）。
注意，在附加 $M$ 有限这一假设后，该结果就是引理
\ref{algebra-lemma-finite-flat-local}。一般地，有：

\begin{lemma}
\label{algebra-lemma-projective-free}
设 $R$ 是环。那么每个投射 $R$-模都是自由模，当且仅当每个
可数生成投射 $R$-模都是自由模。
\end{lemma}

\begin{proof}
立即由定理 \ref{algebra-theorem-projective-direct-sum} 得出。
\end{proof}

\noindent
以下是可数生成模为自由模的一个判据。

\begin{lemma}
\label{algebra-lemma-freeness-criteria}
设 $M$ 是可数生成 $R$-模，并具有如下性质：若
$M = N \oplus N'$，其中 $N'$ 是有限自由 $R$-模，则 $N$
的任意元素都包含在 $N$ 的某个自由直和因子中。那么 $M$ 是自由模。
\end{lemma}

\begin{proof}
设 $x_1, x_2, \ldots$ 是 $M$ 的一组可数生成元。
我们归纳构造 $M$ 的有限自由直和因子 $F_1, F_2, \ldots$，
使得对所有 $n$，$F_1 \oplus \ldots \oplus F_n$ 都是 $M$
的直和因子，并包含 $x_1, \ldots, x_n$。具体地，给定具有所需性质的
$F_1, \ldots, F_n$，写成
$$
M = F_1 \oplus \ldots \oplus F_n \oplus N
$$
并令 $x \in N$ 为 $x_{n + 1}$ 的像。由引理陈述中的假设，
可找到包含 $x$ 的自由直和因子 $F_{n + 1} \subset N$。
当然，可以把 $F_{n + 1}$ 换成 $F_{n + 1}$ 的一个有限自由直和因子，
归纳步骤遂告完成。于是
$M = \bigoplus_{i = 1}^{\infty} F_i$ 是自由模。
\end{proof}

\begin{lemma}
\label{algebra-lemma-projective-freeness-criteria}
设 $P$ 是局部环 $R$ 上的投射模。那么 $P$ 的任意元素都包含在
$P$ 的某个自由直和因子中。
\end{lemma}

\begin{proof}
因为 $P$ 投射，所以它是某个自由 $R$-模 $F$ 的直和因子；写成
$F = P \oplus Q$。设 $x \in P$ 是我们要证明包含于 $P$ 的某个
自由直和因子中的元素。选取 $F$ 的一组基 $B$，使表示 $x$ 所需的
基元素个数最少；设
$x = \sum_{i=1}^n a_i e_i$，其中 $e_i \in B$ 且 $a_i \in R$。
那么任何 $a_j$ 都不能写成其余 $a_i$ 的线性组合；否则，若对某些
$b_i \in R$ 有 $a_j = \sum_{i \neq  j} a_i b_i$，则对
$i \neq j$ 把 $e_i$ 换成 $e_i + b_ie_j$，并保持 $B$ 的其余
元素不变，便得到 $F$ 的一组新基，而 $x$ 关于这组基的表示更短。

\medskip\noindent
把 $e_i$ 按其 $P$-分量和 $Q$-分量分解为
$e_i = y_i + z_i, y_i \in P, z_i \in Q$。写成
$y_i = \sum_{j=1}^{n} b_{ij} e_j + t_i$，其中 $t_i$ 是
$B$ 中除 $e_1, \ldots, e_n$ 外的元素的线性组合。
为完成证明，只须说明矩阵 $(b_{ij})$ 可逆。因为在这种情况下，
对 $i=1, \ldots, n$ 按 $e_i \mapsto y_i$ 定义，并固定
$B \setminus \{e_1, \ldots, e_n\}$ 的映射 $F \to F$ 是同构，
所以 $y_1, \ldots, y_n$ 与
$B \setminus \{e_1, \ldots, e_n\}$ 一起构成 $F$ 的一组基。
于是由 $y_1, \ldots, y_n$ 张成的子模 $N$ 是 $P$ 的自由子模；
因为 $N \subset P$，并且 $N$ 和 $P$ 都是 $F$ 的直和因子，
$N$ 是 $P$ 的直和因子；又因为 $x \in P$ 蕴含
$x = \sum_{i=1}^n a_i e_i = \sum_{i=1}^n a_i y_i$，
所以 $x \in N$。

\medskip\noindent
现在证明 $(b_{ij})$ 可逆。把
$y_i = \sum_{j=1}^{n} b_{ij} e_j + t_i$ 代入
$\sum_{i=1}^n a_i e_i = \sum_{i=1}^n a_i y_i$，比较 $e_j$
的系数，得到 $a_j = \sum_{i=1}^n a_i b_{ij}$。
但如上所述，基 $B$ 的选取保证任何 $a_j$ 都不能写成其余
$a_i$ 的线性组合。因此，当 $i \neq j$ 时，$b_{ij}$ 是非单位；
而对所有 $i$，$1-b_{ii}$ 是非单位，因而特别地 $b_{ii}$ 是单位。
局部环上的矩阵若对角线元素均为单位而其余元素均为非单位，
则其行列式是单位，故矩阵可逆。
\end{proof}

\begin{theorem}
\label{algebra-theorem-projective-free-over-local-ring}
\begin{slogan}
局部环上的投射模是自由模。
\end{slogan}
若 $P$ 是局部环 $R$ 上的投射模，则 $P$ 是自由模。
\end{theorem}

\begin{proof}
结论由引理 \ref{algebra-lemma-projective-free}、
\ref{algebra-lemma-freeness-criteria} 和
\ref{algebra-lemma-projective-freeness-criteria} 得出。
\end{proof}

\section{Mittag-Leffler 系统}
\label{algebra-section-mittag-leffler}

\noindent
本节旨在定义 Mittag-Leffler 系统，并说明这一概念为何有用。

\medskip\noindent
下文中，$I$ 表示一个有向集，见《范畴》，定义
\ref{categories-definition-directed-set}。设
$(A_i, \varphi_{ji}: A_j \to A_i)$ 是以 $I$ 为指标的集合或模的逆系，
见《范畴》，定义 \ref{categories-definition-directed-system}。
由于假设 $I$ 有向，这是一个有向逆系（《范畴》，定义
\ref{categories-definition-directed-system}）。对每个 $i \in I$，当
$j \geq i$ 时，诸像 $\varphi_{ji}(A_j) \subset A_i$ 构成 $A_i$
的子集（或子模）的递减有向族。令
$A'_i = \bigcap_{j \geq i} \varphi_{ji}(A_j)$。
当 $j \geq i$ 时，$\varphi_{ji}(A'_j) \subset A'_i$，故限制映射给出
有向逆系 $(A'_i, \varphi_{ji}|_{A'_j})$。由集合范畴或模范畴中
逆系极限的构造，有 $\lim A_i = \lim A'_i$。逆系
$(A_i, \varphi_{ji})$ 的 Mittag-Leffler 条件是：$A'_i$ 等于某个
$j \geq i$ 所对应的 $\varphi_{ji}(A_j)$（因而对所有 $k \geq j$，
也等于 $\varphi_{ki}(A_k)$）：

\begin{definition}
\label{algebra-definition-ML-system}
设 $(A_i, \varphi_{ji})$ 是 $I$ 上集合的有向逆系。若对每个
$i \in I$，当 $j \geq i$ 时，族 $\varphi_{ji}(A_j) \subset A_i$
稳定，则称 $(A_i, \varphi_{ji})$ 是 {\it Mittag-Leffler} 的。
明确地说，这意味着对每个 $i \in I$，存在 $j \geq i$，使得对
$k \geq j$ 有 $\varphi_{ki}(A_k) = \varphi_{ji}( A_j)$。
若 $(A_i, \varphi_{ji})$ 是环 $R$ 上模的有向逆系，则当其底层集合逆系
是 Mittag-Leffler 的时候，称它是 Mittag-Leffler 的。
\end{definition}

\begin{example}
\label{algebra-example-ML-surjective-maps}
若 $(A_i, \varphi_{ji})$ 是集合或模的有向逆系，且映射
$\varphi_{ji}$ 都是满射，则该逆系显然是 Mittag-Leffler 的。
反之，设 $(A_i, \varphi_{ji})$ 是 Mittag-Leffler 的。令
$A'_i \subset A_i$ 为当 $j \geq i$ 时 $\varphi_{ji}(A_j)$ 的稳定像。
则当 $j \geq i$ 时，$\varphi_{ji}|_{A'_j}: A'_j \to A'_i$ 是满射，
且 $\lim A_i = \lim A'_i$。因此，Mittag-Leffler 逆系
$(A_i, \varphi_{ji})$ 的极限也可写成 $I$ 上一个转移映射均为满射的
有向逆系的极限。
\end{example}

\begin{lemma}
\label{algebra-lemma-ML-limit-nonempty}
设 $(A_i, \varphi_{ji})$ 是 $I$ 上的有向逆系。假设 $I$ 可数。
若 $(A_i, \varphi_{ji})$ 是 Mittag-Leffler 的，且各 $A_i$ 非空，
则 $\lim A_i$ 非空。
\end{lemma}

\begin{proof}
设 $i_1, i_2, i_3, \ldots$ 是 $I$ 的元素的一个枚举。按如下条件归纳定义
元素序列 $j_n \in I$，其中 $n = 1, 2, 3, \ldots$：$j_1 = i_1$，
并且 $j_n \geq i_n$，且对 $m < n$ 有 $j_n \geq j_m$。
于是序列 $j_n$ 递增，并构成 $I$ 的共尾子集。因此可假设
$I =\{1, 2, 3, \ldots \}$。由例
\ref{algebra-example-ML-surjective-maps}，问题归结为证明：由正整数指标、
转移映射均为满射的非空集合逆系，其极限非空。这由选择公理得出。
\end{proof}

\noindent
Mittag-Leffler 条件对我们很重要，因为它具有如下正合性性质。

\begin{lemma}
\label{algebra-lemma-ML-exact-sequence}
设
$$
0 \to A_i \xrightarrow{f_i} B_i \xrightarrow{g_i} C_i \to 0
$$
是 $I$ 上阿贝尔群的有向逆系的一个正合列。假设 $I$ 可数。
若 $(A_i)$ 是 Mittag-Leffler 的，则
$$
0 \to \lim A_i \to \lim B_i \to \lim C_i\to 0
$$
正合。
\end{lemma}

\begin{proof}
对有向逆系取极限是左正合的，故只须证明
$\lim B_i \to \lim C_i$ 满射。取 $(c_i) \in \lim C_i$。
对每个 $i \in I$，令 $E_i = g_i^{-1}(c_i)$；由于
$g_i: B_i \to C_i$ 满射，所得集合非空。逆系 $(B_i)$ 的映射系
$\varphi_{ji}: B_j \to B_i$ 限制为映射 $E_j \to E_i$，从而使
$(E_i)$ 成为非空集合的逆系。只须证明 $(E_i)$ 是 Mittag-Leffler 的。
因为这样一来，引理 \ref{algebra-lemma-ML-limit-nonempty} 就说明
$\lim E_i$ 非空，而取 $\lim E_i$ 的任意元素，便得到 $\lim B_i$
中映到 $(c_i)$ 的一个元素。

\medskip\noindent
借助单射 $f_i: A_i \to B_i$，我们把 $A_i$ 视为 $B_i$ 的子集。
由于 $(A_i)$ 是 Mittag-Leffler 的，若 $i \in I$，则存在
$j \geq i$，使得对 $k \geq j$ 有
$\varphi_{ki}(A_k) = \varphi_{ji}(A_j)$。我们断言，对 $k \geq j$
也有 $\varphi_{ki}(E_k) = \varphi_{ji}(E_j)$。对 $k \geq j$，总有
$\varphi_{ki}(E_k) \subset \varphi_{ji}(E_j)$。反过来，取
$e_j \in E_j$；需要找到 $x_k \in E_k$，使
$\varphi_{ki}(x_k) = \varphi_{ji}(e_j)$。任取 $e'_k \in E_k$，
并令 $e'_j = \varphi_{kj}(e'_k)$。于是
$g_j(e_j - e'_j) = c_j - c_j = 0$，故
$e_j - e'_j = a_j \in A_j$。由于
$\varphi_{ki}(A_k) = \varphi_{ji}(A_j)$，存在 $a_k \in A_k$，使
$\varphi_{ki}(a_k) = \varphi_{ji}(a_j)$。因此
$$
\varphi_{ki}(e'_k + a_k) = \varphi_{ji}(e'_j) + \varphi_{ji}(a_j) =
\varphi_{ji}(e_j),
$$
故可取 $x_k = e'_k + a_k$。
\end{proof}

\section{逆系}
\label{algebra-section-inverse-systems}

\noindent
在许多论文中（以及在本节中），术语{\it 逆系}专指偏序集
$(\mathbf{N}, \geq)$ 上的逆系。本节简要讨论这类逆系。
《同调代数》，第 \ref{homology-section-inverse-systems} 节将更广泛地讨论
这些材料。设给定一个环 $R$ 和一列 $R$-模
$$
M_1 \xleftarrow{\varphi_2} M_2 \xleftarrow{\varphi_3} M_3 \leftarrow \ldots
$$
以及图示的映射。复合相继映射后，只要 $i \geq i'$，就得到映射
$\varphi_{ii'} : M_i \to M_{i'}$；而且只要
$i \geq i' \geq i''$，就有
$\varphi_{ii''} = \varphi_{i'i''} \circ \varphi_{i i'}$。
反之，给定映射族 $\varphi_{ii'}$，令
$\varphi_i = \varphi_{i(i-1)}$，便恢复上面所示的映射。在这种情况下，
$$
\lim M_i
=
\{(x_i) \in \prod M_i \mid \varphi_i(x_i) = x_{i - 1}, \ i = 2, 3, \ldots\}
$$
参见《范畴》，第 \ref{categories-section-limit-sets} 节。
如《同调代数》，第 \ref{homology-section-inverse-systems} 节所说明的，
这实际上是 $R$-模范畴中的极限，其定义见《范畴》，第
\ref{categories-section-limits} 节。

\begin{lemma}
\label{algebra-lemma-Mittag-Leffler}
设 $R$ 是环。设 $0 \to K_i \to L_i \to M_i \to 0$ 是 $R$-模的短正合列，
$i \geq 1$，并且它们嵌入如下短正合列之间的映射：
$$
\xymatrix{
0 \ar[r] &
K_i \ar[r] &
L_i \ar[r] &
M_i \ar[r] &
0 \\
0 \ar[r] &
K_{i + 1} \ar[r] \ar[u] &
L_{i + 1} \ar[r] \ar[u] &
M_{i + 1} \ar[r] \ar[u] &
0}
$$
若对每个 $i$，都存在 $c = c(i) \geq i$，使得对所有 $j \geq c$ 有
$\Im(K_c \to K_i) = \Im(K_j \to K_i)$，则序列
$$
0 \to \lim K_i \to \lim L_i \to \lim M_i \to 0
$$
正合。
\end{lemma}

\begin{proof}
这是更一般的引理 \ref{algebra-lemma-ML-exact-sequence} 的特殊情形。
\end{proof}

\section{Mittag-Leffler 模}
\label{algebra-section-mittag-leffler-modules}

\noindent
粗略地说，Mittag-Leffler 模是一个可写成有向余极限、且其对偶是
Mittag-Leffler 系统的模。为了给出精确定义，还需要做一些准备。

\begin{definition}
\label{algebra-definition-ML-inductive-system}
设 $(M_i, f_{ij})$ 是 $R$-模的有向系。若每个 $M_i$ 都是有限表示
$R$-模，并且对每个 $R$-模 $N$，逆系
$$
(\Hom_R(M_i, N), \Hom_R(f_{ij}, N))
$$
都是 Mittag-Leffler 的，则称 $(M_i, f_{ij})$ 是一个
{\it 模的 Mittag-Leffler 有向系}。
\end{definition}

\noindent
下面将刻画哪些 $R$-模是模的 Mittag-Leffler 有向系的余极限。

\begin{definition}
\label{algebra-definition-domination}
设 $f: M \to N$ 和 $g: M \to M'$ 是 $R$-模映射。若对任意
$R$-模 $Q$ 都有 $\Ker(f
\otimes_R \text{id}_Q) \subset \Ker(g \otimes_R \text{id}_Q)$，
则称 $g$ {\it 支配} $f$。
\end{definition}

\noindent
只需对有限表示模检验这一条件。

\begin{lemma}
\label{algebra-lemma-domination-fp}
设 $f: M \to N$ 和 $g: M \to M'$ 是 $R$-模映射。那么 $g$ 支配 $f$，
当且仅当对任意有限表示 $R$-模 $Q$，都有
$\Ker(f \otimes_R \text{id}_Q) \subset \Ker(g \otimes_R
\text{id}_Q)$。
\end{lemma}

\begin{proof}
假设对所有有限表示模 $Q$，都有
$\Ker(f \otimes_R \text{id}_Q) \subset \Ker(g \otimes_R
\text{id}_Q)$。若 $Q$ 是任意模，把
$Q = \colim_{i \in I} Q_i$ 写成有限表示模 $Q_i$ 的一个有向系的余极限。
于是对所有 $i$，都有 $\Ker(f \otimes_R
\text{id}_{Q_i}) \subset \Ker(g \otimes_R \text{id}_{Q_i})$。
由于取有向余极限是正合的且与张量积交换，遂有
$\Ker(f \otimes_R \text{id}_Q) \subset \Ker(g
\otimes_R \text{id}_Q)$。
\end{proof}

\begin{lemma}
\label{algebra-lemma-domination-universally-injective}
设 $f : M \to N$ 和 $g : M \to M'$ 是 $R$-模映射。考虑 $f$ 与 $g$
的推出
$$
\xymatrix{
M  \ar[r]_f \ar[d]_g & N \ar[d]^{g'} \\
M' \ar[r]^{f'} & N'
}
$$
那么 $g$ 支配 $f$，当且仅当 $f'$ 普遍单射。
\end{lemma}

\begin{proof}
回顾 $N'$ 是 $M' \oplus N$ 对由所有 $x \in M$ 所对应的元素
$(g(x), -f(x))$ 组成的子模的商。由 $N'$ 的构造，有短正合列
$$
0 \to \Ker(f) \cap \Ker(g) \to \Ker(f) \to \Ker(f')
\to 0.
$$
因为张量运算与取推出交换，所以对每个 $R$-模 $Q$，都有这样的短正合列
$$
0 \to \Ker(f \otimes \text{id}_Q ) \cap \Ker(g \otimes
\text{id}_Q) \to \Ker(f \otimes \text{id}_Q)
\to \Ker(f' \otimes \text{id}_Q) \to 0
$$
因此 $f'$ 普遍单射，当且仅当对每个 $Q$ 都有
$\Ker(f \otimes \text{id}_Q ) \subset \Ker(g \otimes
\text{id}_Q)$，而这又当且仅当 $g$ 支配 $f$。
\end{proof}

\noindent
上述支配定义有时与通常的映射支配概念相关，如下引理所示。

\begin{lemma}
\label{algebra-lemma-domination}
设 $f: M \to N$ 和 $g: M \to M'$ 是 $R$-模映射。假设
$\Coker(f)$ 是有限表示的。那么 $g$ 支配 $f$，当且仅当 $g$ 经由
$f$ 分解，即存在模映射 $h: N
\to M'$，使得 $g = h \circ f$。
\end{lemma}

\begin{proof}
考虑引理 \ref{algebra-lemma-domination-universally-injective} 陈述中 $f$ 与 $g$
的推出。由推出的构造可知 $\Coker(f') = \Coker(f)$，故
$\Coker(f')$ 是有限表示的。于是由引理
\ref{algebra-lemma-universally-exact-split}，$f'$ 普遍单射，当且仅当
$$
0 \to M' \xrightarrow{f'} N' \to \Coker(f') \to 0
$$
分裂。这当且仅当存在映射 $h' : N' \to M'$，使
$h' \circ f' = \text{id}_{M'}$。由推出的泛性质，这样的 $h'$ 存在，
当且仅当 $g$ 经由 $f$ 分解。
\end{proof}


\begin{proposition}
\label{algebra-proposition-ML-characterization}
设 $M$ 是 $R$-模。设 $(M_i, f_{ij})$ 是以 $I$ 为指标的有限表示
$R$-模有向系，且 $M = \colim M_i$。令 $f_i:
M_i \to M$ 为典范映射。下列条件等价：
\begin{enumerate}
\item 对每个有限表示 $R$-模 $P$ 和模映射 $f: P
\to M$，都存在有限表示 $R$-模 $Q$ 和模映射 $g: P \to Q$，
使得 $g$ 与 $f$ 相互支配，即对每个 $R$-模 $N$ 都有
$\Ker(f \otimes_R \text{id}_N) = \Ker(g \otimes_R \text{id}_N)$。
\item 对每个 $i \in I$，都存在 $j \geq i$，使得
$f_{ij}: M_i \to M_j$ 支配 $f_i: M_i \to M$。
\item 对每个 $i \in I$，都存在 $j \geq i$，使得对所有
$k \geq i$，$f_{ij}: M_i \to M_j$ 都经由
$f_{ik}: M_i \to M_k$ 分解。
\item 对每个 $R$-模 $N$，逆系
$(\Hom_R(M_i, N), \Hom_R(f_{ij}, N))$ 都是 Mittag-Leffler 的。
\item 对 $N = \prod_{s \in I} M_s$，逆系
$(\Hom_R(M_i, N), \Hom_R(f_{ij}, N))$ 是 Mittag-Leffler 的。
\end{enumerate}
\end{proposition}

\begin{proof}
先证明 (1) 与 (2) 等价。假设 (1) 成立，并取 $i
\in I$。对应于映射 $f_i: M_i \to M$，可如 (1) 选取
$g: M_i \to Q$。由于 $M_i$ 和 $Q$ 都是有限表示的，
$\Coker(g)$ 也是有限表示的。于是由引理 \ref{algebra-lemma-domination}，
$f_i : M_i \to M$ 经由 $g: M_i \to Q$ 分解；设对某个
$h: Q \to M$ 有 $f_i = h \circ g$。又因为 $Q$ 有限表示，
$h$ 经由某个 $j \geq i$ 所对应的 $M_j \to M$ 分解；设对某个
$h': Q \to M_j$ 有 $h = f_j \circ h'$。综上得到交换图表
$$
\xymatrix{
  &  M  & \\
M_i \ar[dr]_g \ar[ur]^{f_i} \ar[rr]^{f_{ij}} &
& M_j \ar[ul]_{f_j} \\
  & Q  \ar[ur]_{h'} &
}
$$
因此 $f_{ij}$ 支配 $g$。而 $g$ 支配 $f_i$，所以 $f_{ij}$ 支配
$f_i$。

\medskip\noindent
反之，假设 (2) 成立。设 $P$ 是有限表示的，并设 $f: P
\to M$ 是模映射。那么对某个 $i \in I$，$f$ 经由
$f_i: M_i \to M$ 分解；设对某个 $g': P \to M_i$ 有
$f = f_i \circ g'$。由 (2) 选取 $j \geq i$，使 $f_{ij}$ 支配
$f_i$。我们有交换图表
$$
\xymatrix{
P \ar[d]_{g'} \ar[r]^{f}           & M  \\
M_i \ar[ur]^{f_i} \ar[r]_{f_{ij}} & M_j \ar[u]_{f_j}
}
$$
由该图表以及 $f_{ij}$ 支配 $f_i$ 这一事实可知，$f$ 与
$f_{ij} \circ g'$ 相互支配。因此取
$g = f_{ij} \circ g' : P \to M_j$ 即可。

\medskip\noindent
下面证明 (2) 与 (3) 等价。取 $i \in I$。对 $k \geq i$，
$f_i$ 总是支配 $f_{ik}$，因为 $f_i$ 经由 $f_{ik}$ 分解。
若 (2) 成立，选取 $j \geq i$，使得 $f_{ij}$ 支配 $f_i$。
由于支配关系具有传递性，对 $k \geq i$，$f_{ij}$ 支配 $f_{ik}$。
所有 $M_i$ 都是有限表示的，故对 $k \geq i$，
$\Coker(f_{ik})$ 是有限表示的。由引理 \ref{algebra-lemma-domination}，
对所有 $k \geq i$，$f_{ij}$ 经由 $f_{ik}$ 分解。因此 (2) 蕴含 (3)。
另一方面，若 (3) 成立，则对任意 $R$-模 $N$ 和 $k \geq i$，
$f_{ij} \otimes_R \text{id}_N$ 经由 $f_{ik}
\otimes_R \text{id}_N$ 分解。因此对 $k
\geq i$，有 $\Ker(f_{ik} \otimes_R
\text{id}_N) \subset \Ker(f_{ij} \otimes_R \text{id}_N)$。
而 $\Ker(f_i \otimes_R \text{id}_N: M_i \otimes_R N
\to M \otimes_R N)$ 是对 $k \geq i$ 的诸
$\Ker(f_{ik} \otimes_R
\text{id}_N)$ 之并。因此对任意 $R$-模 $N$，都有
$\Ker(f_i \otimes_R
\text{id}_N) \subset \Ker(f_{ij} \otimes_R \text{id}_N)$；
按定义，这意味着 $f_{ij}$ 支配 $f_i$。

\medskip\noindent
(3) 蕴含 (4)、(4) 蕴含 (5) 是显然的。下面证明 (5) 蕴含 (3)。令
$N = \prod_{s \in I} M_s$。若 (5) 成立，则给定 $i \in I$，选取
$j \geq i$，使得
$$
\Im( \Hom(M_j, N) \to  \Hom(M_i, N)) =
\Im( \Hom(M_k, N) \to  \Hom(M_i, N))
$$
对所有 $k \geq j$ 成立。把对 $s \in I$ 的乘积移到诸 $\Hom$ 外，
并考察乘积每个分量上的映射，这就是说
$$
\Im( \Hom(M_j, M_s) \to  \Hom(M_i, M_s)) =
\Im( \Hom(M_k, M_s) \to   \Hom(M_i, M_s))
$$
对所有 $k \geq j$ 和 $s \in I$ 成立。取 $s = j$，得到
$$
\Im( \Hom(M_j, M_j) \to  \Hom(M_i, M_j)) =
\Im( \Hom(M_k, M_j) \to   \Hom(M_i, M_j))
$$
对所有 $k \geq j$ 成立。由于 $f_{ij}$ 是
$\text{id} \in  \Hom(M_j, M_j)$ 在映射
$\Hom(M_j, M_j) \to  \Hom(M_i, M_j)$ 下的像，这说明对任意
$k \geq j$，都存在 $h \in \Hom(M_k, M_j)$，使
$f_{ij} = h \circ f_{ik}$。若 $j \geq k$，则可取
$h = f_{kj}$。因此 (3) 成立。
\end{proof}

\begin{definition}
\label{algebra-definition-mittag-leffler-module}
设 $M$ 是 $R$-模。若命题 \ref{algebra-proposition-ML-characterization}
中的诸等价条件成立，则称 $M$ 是 {\it Mittag-Leffler} 的。
\end{definition}

\noindent
特别地，有限表示模是 Mittag-Leffler 的。

\begin{remark}
\label{algebra-remark-flat-ML}
设 $M$ 是平坦 $R$-模。由 Lazard 定理（定理
\ref{algebra-theorem-lazard}），可把 $M = \colim M_i$ 写成有向系
$(M_i, f_{ij})$ 的余极限，其中各 $M_i$ 是有限自由 $R$-模。
要使 $M$ 成为 Mittag-Leffler 模，只需对偶逆系
$(\Hom_R(M_i, R), \Hom_R(f_{ij}, R))$ 是 Mittag-Leffler 的。
这由命题 \ref{algebra-proposition-ML-characterization} 的判据 (4)，以及如下事实
得出：对有限自由 $R$-模 $F$ 和任意 $R$-模 $N$，存在函子性同构
$\Hom_R(F, R) \otimes_R N \cong \Hom_R(F, N)$。
\end{remark}

\begin{lemma}
\label{algebra-lemma-tensor-ML-modules}
若 $R$ 是环，且 $M$、$N$ 是 $R$ 上的 Mittag-Leffler 模，
则 $M \otimes_R N$ 是 Mittag-Leffler 模。
\end{lemma}

\begin{proof}
把 $M = \colim_{i \in I} M_i$ 和 $N = \colim_{j \in J} N_j$
写成有限表示 $R$-模的有向余极限。分别以
$f_{ii'} : M_i \to M_{i'}$ 和 $g_{jj'} : N_j \to N_{j'}$ 表示
转移映射。于是 $M_i \otimes_R N_j$ 是有限表示 $R$-模（见引理
\ref{algebra-lemma-tensor-finiteness}），并且
% P01-E0452: frozen target preserves the source's M_j; see ERRATA_CANDIDATES.jsonl.
$M \otimes_R N = \colim_{(i, j) \in I \times J} M_i \otimes_R M_j$。
取 $(i, j) \in I \times J$。根据 Mittag-Leffler 模的定义，
两个系都满足命题 \ref{algebra-proposition-ML-characterization} (3)。
换言之，存在 $i' \geq i$ 和 $j' \geq j$，使得对每一组
$i'' \geq i$ 和 $j'' \geq j$，都存在映射
% P01-E0453: frozen target preserves the source's M-indexed b; see ERRATA_CANDIDATES.jsonl.
$a : M_{i''} \to M_{i'}$ 和 $b : M_{j''} \to M_{j'}$，满足
$f_{ii'} = a \circ f_{ii''}$ 和 $g_{jj'} = b \circ g_{jj''}$。
显然，
$a \otimes b : M_{i''} \otimes_R N_{j''} \to M_{i'} \otimes_R N_{j'}$
对系 $(M_i \otimes_R N_j, f_{ii'} \otimes g_{jj'})$ 起同样作用。
因此由刻画命题 \ref{algebra-proposition-ML-characterization} (3)，可知
$M \otimes_R N$ 是 Mittag-Leffler 的。
\end{proof}

\begin{lemma}
\label{algebra-lemma-ML-also}
设 $R$ 是环，$M$ 是 $R$-模。那么 $M$ 是 Mittag-Leffler 的，
当且仅当对每个有限自由 $R$-模 $F$ 和模映射 $f: F \to M$，
都存在有限表示 $R$-模 $Q$ 和模映射 $g : F \to Q$，使 $g$ 与
$f$ 相互支配，即对每个 $R$-模 $N$ 都有
$\Ker(f \otimes_R \text{id}_N) = \Ker(g \otimes_R \text{id}_N)$。
\end{lemma}

\begin{proof}
由于该条件显然弱于命题 \ref{algebra-proposition-ML-characterization} 的条件 (1)，
可见 Mittag-Leffler 模满足该条件。反之，假设 $M$ 满足该条件，且
$f : P \to M$ 是从有限表示 $R$-模 $P$ 到 $M$ 的 $R$-模映射。
选取满射 $F \to P$，其中 $F$ 是有限自由 $R$-模。由假设，
可找到映射 $F \to Q$，其中 $Q$ 是有限表示 $R$-模，并且
$F \to Q$ 与 $F \to M$ 相互支配。特别地，$F \to Q$ 的核包含
$F \to P$ 的核，因而得到 $R$-模映射 $g : P \to Q$，使
$F \to Q$ 等于复合 $F \to P \to Q$。设 $N$ 是任意 $R$-模，
并考虑交换图表
$$
\xymatrix{
F \otimes_R N \ar[d] \ar[r] & Q \otimes_R N \\
P \otimes_R N \ar[ru] \ar[r] & M \otimes_R N
}
$$
由假设，$F \otimes_R N \to Q \otimes_R N$ 与
$F \otimes_R N \to M \otimes_R N$ 的核相等。因此，因为
$F \otimes_R N \to P \otimes_R N$ 满射，
$P \otimes_R N \to Q \otimes_R N$ 与
$P \otimes_R N \to M \otimes_R N$ 的核也相等。
\end{proof}

\begin{lemma}
\label{algebra-lemma-restrict-ML-modules}
设 $R \to S$ 是有限且有限表示的环映射。设 $M$ 是 $S$-模。
若 $M$ 作为 $S$-模是 Mittag-Leffler 的，则 $M$ 作为 $R$-模也是
Mittag-Leffler 的。
\end{lemma}

\begin{proof}
假设 $M$ 作为 $S$-模是 Mittag-Leffler 的。把 $M = \colim M_i$
写成有限表示 $S$-模 $M_i$ 的一个有向余极限。由于 $M$ 在 $S$ 上
是 Mittag-Leffler 的，对每个 $i$，都存在指标 $j \geq i$，使得
对所有 $k \geq j$，都有分解 $f_{ij} = h \circ f_{ik}$（其中 $h$
依赖于 $i$、$j$ 的选取和 $k$）。注意，由引理
\ref{algebra-lemma-finite-finitely-presented-extension}，诸模 $M_i$ 作为
$R$-模也是有限表示的。此外，所有映射 $f_{ij}, f_{ik}, h$ 都是
$R$-模映射。因此，把系 $(M_i, f_{ij})$ 视为 $R$-模的系时，
它满足同一条件。所以 $M$ 作为 $R$-模是 Mittag-Leffler 的。
\end{proof}

\begin{lemma}
\label{algebra-lemma-mod-ideal-ML-modules}
设 $R$ 是环。设 $S = R/I$，其中 $I$ 是有限生成理想。
设 $M$ 是 $S$-模。那么 $M$ 作为 $R$-模是 Mittag-Leffler 的，
当且仅当 $M$ 作为 $S$-模是 Mittag-Leffler 的。
\end{lemma}

\begin{proof}
一个方向由引理 \ref{algebra-lemma-restrict-ML-modules} 得出。
为证明另一个方向，假设 $M$ 作为 $R$-模是 Mittag-Leffler 的。
把 $M = \colim M_i$ 写成有限表示 $S$-模的一个有向余极限。
由于 $I$ 有限生成，环 $S$ 作为 $R$-代数是有限且有限表示的，
故诸模 $M_i$ 作为 $R$-模是有限表示的，见引理
\ref{algebra-lemma-finite-finitely-presented-extension}。接着，设 $N$ 是任意
$S$-模。注意，对每个 $i$，都有
$\Hom_R(M_i, N) = \Hom_S(M_i, N)$，因为 $R \to S$ 满射。
因此，逆系 $(\Hom_R(M_i, N))_i$ 满足 Mittag-Leffler 条件，蕴含
逆系 $(\Hom_S(M_i, N))_i$ 满足 Mittag-Leffler 条件。故按定义，
$M$ 在 $S$ 上是 Mittag-Leffler 的。
\end{proof}

\begin{remark}
\label{algebra-remark-go-up-ML-modules}
设 $R \to S$ 是有限且有限表示的环映射。设 $M$ 是一个 $S$-模，
它作为 $R$-模是 Mittag-Leffler 的。一般而言，$M$ 作为 $S$-模未必
是 Mittag-Leffler 的。例如，假设 $S$ 是 $R$ 上的对偶数环，即
$S = R \oplus R\epsilon$，其中 $\epsilon^2 = 0$。那么一个 $S$-模
由一个 $R$-模 $M$ 连同一个平方为零的 $R$-线性自同态
$\epsilon : M \to M$ 构成。现在假设 $M_0$ 是一个非 Mittag-Leffler
的 $R$-模。选取表示 $F_1 \xrightarrow{u} F_0 \to M_0 \to 0$，
其中 $F_1$ 和 $F_0$ 是自由 $R$-模。令 $M = F_1 \oplus F_0$，并取
$$
\epsilon =
\left(
\begin{matrix}
0 & 0 \\
u & 0
\end{matrix}
\right) : M \longrightarrow M.
$$
那么 $M/\epsilon M \cong F_1 \oplus M_0$ 在
$R = S/\epsilon S$ 上不是 Mittag-Leffler 的，因而在 $S$ 上也不是
Mittag-Leffler 的（见引理 \ref{algebra-lemma-mod-ideal-ML-modules}）。
另一方面，$M/\epsilon M = M \otimes_S S/\epsilon S$；若 $M$ 是
Mittag-Leffler 的，则由引理 \ref{algebra-lemma-tensor-ML-modules}，它在
$S$ 上也会是 Mittag-Leffler 的。
\end{remark}

\section{直积与张量积的交换}
\label{algebra-section-products-tensor}

\noindent
设 $M$ 是 $R$-模，且 $(Q_{\alpha})_{\alpha \in A}$ 是一族 $R$-模。
那么有典范映射 $M \otimes_R \left( \prod_{\alpha
\in A} Q_{\alpha} \right) \to \prod_{\alpha \in A} ( M \otimes_R
Q_{\alpha})$，它在纯张量上由 $x \otimes (q_{\alpha}) \mapsto (x
\otimes q_{\alpha})$ 给出。如下例所示，该映射不一定单射，也不一定满射。

\begin{example}
\label{algebra-example-Q-not-ML}
取 $R = \mathbf{Z}$、$M = \mathbf{Q}$，并考虑当 $n \geq 1$ 时的族
$Q_n = \mathbf{Z}/n$。那么 $\prod_n (M \otimes Q_n) = 0$。
然而，将单射 $\mathbf{Z} \to \prod_n Q_n$ 与 $M$ 作张量，得到单射
$\mathbf{Q} \to M \otimes (\prod_n Q_n)$，所以
$M \otimes (\prod_n Q_n)$ 非零。因此
$M \otimes (\prod_n Q_n) \to \prod_n (M \otimes Q_n)$ 不是单射。

\medskip\noindent
另一方面，再取 $R = \mathbf{Z}$、$M = \mathbf{Q}$，并令当
$n \geq 1$ 时 $Q_n = \mathbf{Z}$。映射
$M \otimes (\prod_n Q_n)
\to \prod_n (M \otimes Q_n) = \prod_n M$ 的像恰由形如
$(a_n/m)_{n \geq 1}$ 的序列组成，其中 $a_n \in \mathbf{Z}$，而 $m$
是某个非零整数。因此该映射不是满射。
\end{example}

\noindent
下面确定使映射 $M
\otimes_R \left( \prod_{\alpha} Q_{\alpha} \right) \to \prod_{\alpha}
(M \otimes_R Q_{\alpha})$ 对 $(Q_{\alpha})_{\alpha \in A}$ 的所有选择
分别成为满射、双射或单射时，$M$ 所需满足的精确条件。这一点很重要，
因为使该映射单射的模恰好就是 Mittag-Leffler 模（命题
\ref{algebra-proposition-ML-tensor}）。下文中，若 $M$ 是 $R$-模且 $A$ 是集合，
我们用 $M^A$ 表示直积 $\prod_{\alpha \in A} M$。

\begin{proposition}
\label{algebra-proposition-fg-tensor}
设 $M$ 是 $R$-模。下列条件等价：
\begin{enumerate}
\item $M$ 有限生成。
\item 对每族 $R$-模 $(Q_{\alpha})_{\alpha \in A}$，典范映射
$M \otimes_R \left( \prod_{\alpha} Q_{\alpha} \right)
\to \prod_{\alpha} (M \otimes_R Q_{\alpha})$ 都是满射。
\item 对每个 $R$-模 $Q$ 和每个集合 $A$，典范映射
$M \otimes_R Q^{A} \to (M \otimes_R Q)^{A}$ 都是满射。
\item 对每个集合 $A$，典范映射 $M \otimes_R R^{A} \to
M^{A}$ 都是满射。
\end{enumerate}
\end{proposition}

\begin{proof}
先证明 (1) 蕴含 (2)。选取满射 $R^n \to M$，并考虑交换图表
$$
\xymatrix{
R^n \otimes_R (\prod_{\alpha} Q_{\alpha})  \ar[r]^{\cong} \ar[d] &
\prod_{\alpha} (R^n \otimes_R Q_{\alpha}) \ar[d] \\
M \otimes_R (\prod_{\alpha} Q_{\alpha})  \ar[r] & \prod_{\alpha} ( M
\otimes_R Q_{\alpha}).
}
$$
上方箭头是同构，两个竖直箭头是满射。因此下方箭头是满射。

\medskip\noindent
显然 (2) 蕴含 (3)，(3) 蕴含 (4)，所以只须证明 (4) 蕴含 (1)。
事实上，要使 (1) 成立，只须 $M^M$ 的元素
$d = (x)_{x \in M}$ 属于映射 $f: M \otimes_R R^{M} \to M^M$ 的像。
在这种情况下，对某些 $x_i \in M$ 和 $a_i \in R^M$，有
$d = \sum_{i = 1}^{n} f(x_i \otimes a_i)$。若对 $x \in M$，用
$p_x: M^M \to M$ 表示到第 $x$ 个因子的投影，则
% P01-E0455: p_x was typed only on M^M but is also applied to a_i in R^M; see ERRATA_CANDIDATES.jsonl.
$$
x = p_x(d) = \sum\nolimits_{i = 1}^{n} p_x(f(x_i \otimes a_i)) =
\sum\nolimits_{i=1}^{n} p_x(a_i) x_i.
$$
因此 $x_1, \ldots, x_n$ 生成 $M$。
\end{proof}

\begin{proposition}
\label{algebra-proposition-fp-tensor}
设 $M$ 是 $R$-模。下列条件等价：
\begin{enumerate}
\item $M$ 是有限表示的。
\item 对每族 $R$-模 $(Q_{\alpha})_{\alpha \in A}$，典范映射
$M \otimes_R \left( \prod_{\alpha} Q_{\alpha} \right)
\to \prod_{\alpha} (M \otimes_R Q_{\alpha})$ 都是双射。
\item 对每个 $R$-模 $Q$ 和每个集合 $A$，典范映射
$M \otimes_R Q^{A} \to (M \otimes_R Q)^{A}$ 都是双射。
\item 对每个集合 $A$，典范映射 $M \otimes_R R^{A} \to
M^{A}$ 都是双射。
\end{enumerate}
\end{proposition}

\begin{proof}
先证明 (1) 蕴含 (2)。选取表示 $R^m \to R^n
\to M$，并考虑交换图表
$$
\xymatrix{
R^m \otimes_R (\prod_{\alpha} Q_{\alpha}) \ar[r] \ar[d]^{\cong} & R^n
\otimes_R (\prod_{\alpha} Q_{\alpha}) \ar[r] \ar[d]^{\cong} & M \otimes_R
(\prod_{\alpha} Q_{\alpha}) \ar[r] \ar[d] & 0 \\
\prod_{\alpha} (R^m \otimes_R Q_{\alpha}) \ar[r] & \prod_{\alpha} (R^n
\otimes_R Q_{\alpha}) \ar[r] & \prod_{\alpha} (M \otimes_R Q_{\alpha})
\ar[r] & 0.
}
$$
前两个竖直箭头是同构，两行都正合。这说明映射
$M \otimes_R (\prod_{\alpha} Q_{\alpha}) \to
\prod_{\alpha} ( M \otimes_R Q_{\alpha})$
是满射；图追踪还说明它是单射。因此 (2) 成立。

\medskip\noindent
显然 (2) 蕴含 (3)，(3) 蕴含 (4)，所以只须证明 (4) 蕴含 (1)。
由命题 \ref{algebra-proposition-fg-tensor}，若 (4) 成立，我们已经知道 $M$
有限生成。因此可选取满射 $F \to M$，其中 $F$ 是有限自由模。
令 $K$ 为其核。必须证明 $K$ 有限生成。对任意集合 $A$，有交换图表
$$
\xymatrix{
& K \otimes_R R^A \ar[r] \ar[d]_{f_3} & F \otimes_R R^A \ar[r]
\ar[d]_{f_2}^{\cong} & M \otimes_R R^A \ar[r] \ar[d]_{f_1}^{\cong} & 0 \\
0 \ar[r] & K^A \ar[r] & F^A \ar[r] & M^A \ar[r] & 0 .
}
$$
由假设，映射 $f_1$ 是同构；因为 $F$ 是有限自由模，映射 $f_2$
是同构；两行都正合。图追踪说明 $f_3$ 满射，故由命题
\ref{algebra-proposition-fg-tensor}，$K$ 有限生成。
\end{proof}

\noindent
下一个命题需要如下引理。

\begin{lemma}
\label{algebra-lemma-kernel-tensored-fp}
设 $M$ 是 $R$-模，$P$ 是有限表示 $R$-模，且 $f: P
\to M$ 是映射。设 $Q$ 是 $R$-模，并假设 $x \in \Ker(P
\otimes Q \to M \otimes Q)$。那么存在有限表示 $R$-模 $P'$ 和映射
$f': P \to P'$，使得 $f$ 经由 $f'$ 分解，并且
$x \in \Ker(P \otimes Q \to P' \otimes Q)$。
\end{lemma}

\begin{proof}
把 $M$ 写成有限表示模 $M_i$ 的一个有向系的余极限
$M = \colim_{i \in I} M_i$。因为 $P$ 有限表示，映射 $f:
P \to M$ 经由某个 $j \in I$ 所对应的 $M_j \to M$ 分解。
与 $Q$ 作张量后，得到交换图表
$$
\xymatrix{
& M_j \otimes Q \ar[dr] & \\
P \otimes Q \ar[ur] \ar[rr] & & M \otimes Q .
}
$$
$x$ 在 $M_j \otimes Q$ 中的像 $y$ 属于
$M_j \otimes Q \to M \otimes Q$ 的核。由于
$M \otimes Q = \colim_{i \in I} (M_i
\otimes
Q)$，这意味着对某个 $j' \geq j$，$y$ 在 $M_{j'} \otimes Q$ 中
映为 $0$。因此可取 $P' = M_{j'}$，并取 $f'$ 为复合
$P \to M_j \to M_{j'}$。
\end{proof}

\begin{proposition}
\label{algebra-proposition-ML-tensor}
设 $M$ 是 $R$-模。下列条件等价：
\begin{enumerate}
\item $M$ 是 Mittag-Leffler 的。
\item 对每族 $R$-模 $(Q_{\alpha})_{\alpha \in A}$，典范映射
$M \otimes_R \left( \prod_{\alpha} Q_{\alpha} \right)
\to \prod_{\alpha} (M \otimes_R Q_{\alpha})$ 都是单射。
\end{enumerate}
\end{proposition}

\begin{proof}
先证明 (1) 蕴含 (2)。假设 $M$ 是 Mittag-Leffler 的，并设 $x$ 属于
$M \otimes_R (\prod_{\alpha} Q_{\alpha}) \to
\prod_{\alpha} (M \otimes_R Q_{\alpha})$ 的核。把 $M$ 写成有限表示模
$M_i$ 的一个有向系的余极限 $M =
\colim_{i \in I} M_i$。那么
$M \otimes_R (\prod_{\alpha} Q_{\alpha})$ 是诸
$M_i \otimes_R (\prod_{\alpha} Q_{\alpha})$ 的余极限。因此 $x$
是某个元素 $x_i \in M_i \otimes_R (\prod_{\alpha} Q_{\alpha})$ 的像。
必须说明对某个 $j \geq i$，$x_i$ 在
$M_j \otimes_R (\prod_{\alpha} Q_{\alpha})$ 中映为 $0$。
因为 $M$ 是 Mittag-Leffler 的，可选取 $j \geq i$，使
$M_i \to M_j$ 与 $M_i \to M$ 相互支配。然后考虑交换图表
$$
\xymatrix{
M \otimes_R (\prod_{\alpha} Q_{\alpha}) \ar[r] & \prod_{\alpha} (M
\otimes_R Q_{\alpha}) \\
M_i \otimes_R (\prod_{\alpha} Q_{\alpha}) \ar[r]^{\cong} \ar[d] \ar[u] &
\prod_{\alpha} (M_i \otimes_R Q_{\alpha}) \ar[d] \ar[u] \\
M_j \otimes_R (\prod_{\alpha} Q_{\alpha}) \ar[r]^{\cong} & \prod_{\alpha}
(M_j \otimes_R Q_{\alpha})
}
$$
根据命题 \ref{algebra-proposition-fp-tensor}，下方两个水平映射都是同构。
由于 $x_i$ 在 $\prod_{\alpha} (M
\otimes_R Q_{\alpha})$ 中映为 $0$，它在
$\prod_{\alpha} (M_i \otimes_R
Q_{\alpha})$ 中的像属于映射
$\prod_{\alpha} (M_i \otimes_R
Q_{\alpha}) \to \prod_{\alpha} (M \otimes_R Q_{\alpha})$ 的核。
由 $j$ 的选择，该核等于映射
$\prod_{\alpha} (M_i \otimes_R Q_{\alpha})
\to \prod_{\alpha} (M_j \otimes_R Q_{\alpha})$ 的核。
因此 $x_i$ 在 $\prod_{\alpha} (M_j \otimes_R
Q_{\alpha})$ 中映为 $0$，从而在
$M_j \otimes_R (\prod_{\alpha} Q_{\alpha})$ 中也映为 $0$。

\medskip\noindent
现在假设 (2) 成立。我们证明 $M$ 满足命题
\ref{algebra-proposition-ML-characterization} 中 Mittag-Leffler 性的表述 (1)。
设 $f: P \to M$ 是从有限表示模 $P$ 到 $M$ 的映射。
选取集合 $B$，其中含有有限表示 $R$-模的每个同构类的一个代表。
令 $A$ 为诸偶对 $(Q, x)$ 的集合，其中 $Q \in B$ 且
$x \in \Ker(P \otimes Q \to M \otimes Q)$。对
$\alpha = (Q, x) \in A$，记 $Q_{\alpha}$ 为 $Q$，记 $x_{\alpha}$ 为
$x$。考虑交换图表
$$
\xymatrix{
M \otimes_R (\prod_{\alpha} Q_{\alpha}) \ar[r] &
\prod_{\alpha} (M \otimes_R Q_{\alpha}) \\
P \otimes_R (\prod_{\alpha} Q_{\alpha}) \ar[r]^{\cong} \ar[u] &
\prod_{\alpha} (P \otimes_R Q_{\alpha}) \ar[u] .
}
$$
由假设，上方箭头是单射；由命题 \ref{algebra-proposition-fp-tensor}，
下方箭头是同构。令 $x \in P
\otimes_R (\prod_{\alpha} Q_{\alpha})$ 为该同构下对应于
$(x_{\alpha}) \in \prod_{\alpha} (P \otimes_R Q_{\alpha})$ 的元素。
由于图表中的上方箭头单射，有
$x \in \Ker( P \otimes_R (\prod_{\alpha} Q_{\alpha})
\to M \otimes_R (\prod_{\alpha} Q_{\alpha}))$。由引理
\ref{algebra-lemma-kernel-tensored-fp}，得到有限表示模 $P'$ 和映射
$f': P \to P'$，使 $f: P
\to M$ 经由 $f'$ 分解，并且
$x \in \Ker(P \otimes_R
(\prod_{\alpha} Q_{\alpha}) \to P' \otimes_R (\prod_{\alpha}
Q_{\alpha}))$。我们有交换图表
$$
\xymatrix{
P' \otimes_R (\prod_{\alpha} Q_{\alpha}) \ar[r]^{\cong} &
\prod_{\alpha} (P' \otimes_R Q_{\alpha}) \\
P \otimes_R (\prod_{\alpha} Q_{\alpha}) \ar[r]^{\cong} \ar[u] &
\prod_{\alpha} (P \otimes_R Q_{\alpha}) \ar[u] .
}
$$
其中上下两个箭头由命题 \ref{algebra-proposition-fp-tensor} 都是同构。
因此，因为 $x$ 属于左侧竖直映射的核，$(x_{\alpha})$ 属于右侧竖直
映射的核。这意味着对每个 $\alpha \in A$，都有
$x_{\alpha} \in \Ker(P \otimes_R Q_{\alpha} \to P' \otimes_R
Q_{\alpha})$。由 $A$ 的定义，这意味着对所有有限表示 $Q$，都有
$\Ker(P \otimes_R Q \to P' \otimes_R Q) \supset \Ker(P \otimes_R
Q \to M \otimes_R Q)$；而因为 $f: P
\to M$ 经由 $f': P \to P'$ 分解，实际上等号成立。
由引理 \ref{algebra-lemma-domination-fp}，$f$ 与 $f'$ 相互支配。
\end{proof}

\begin{lemma}
\label{algebra-lemma-minimal-contains}
设 $M$ 是 $R$ 上的平坦 Mittag-Leffler 模。设 $F$ 是 $R$-模，且
$x \in F \otimes_R M$。那么存在使 $x \in F' \otimes_R M$
成立的最小子模 $F' \subset F$。此外，$F'$ 是有限 $R$-模。
\end{lemma}

\begin{proof}
因为 $M$ 平坦，若 $F' \subset F$ 是子模，就有
$F' \otimes_R M \subset F \otimes_R M$，故该陈述有意义。
令 $I = \{F' \subset F \mid x \in F' \otimes_R M\}$；对
$i \in I$，用 $F_i \subset F$ 表示相应的子模。那么 $x$ 在映射
$$
F \otimes_R M \longrightarrow \prod (F/F_i \otimes_R M)
$$
下映为零，因而由命题 \ref{algebra-proposition-ML-tensor}，$x$ 在映射
$$
F \otimes_R M \longrightarrow \left(\prod F/F_i\right) \otimes_R M
$$
下映为零。由于 $M$ 平坦，该箭头的核是
$(\bigcap F_i) \otimes_R M$，这就证明 $F' = \bigcap F_i$。
为说明 $F'$ 是有限模，设
$x = \sum_{j = 1, \ldots, m} f_j \otimes m_j$，其中 $f_j \in F'$
且 $m_j \in M$。那么 $x \in F'' \otimes_R M$，其中
$F'' \subset F'$ 是由 $f_1, \ldots, f_m$ 生成的子模。
于是当然有 $F'' = F'$，最终陈述得证。
\end{proof}

\begin{lemma}
\label{algebra-lemma-pure-submodule-ML}
设 $0 \to M_1 \to M_2 \to M_3 \to 0$ 是 $R$-模的一个普遍正合列。
那么：
\begin{enumerate}
\item 若 $M_2$ 是 Mittag-Leffler 的，则 $M_1$ 是 Mittag-Leffler 的。
\item 若 $M_1$ 和 $M_3$ 是 Mittag-Leffler 的，则 $M_2$ 是
Mittag-Leffler 的。
\end{enumerate}
\end{lemma}

\begin{proof}
对任意一族 $R$-模 $(Q_{\alpha})_{\alpha \in A}$，都有交换图表
$$
\xymatrix{
0 \ar[r] & M_1 \otimes_R (\prod_{\alpha} Q_{\alpha}) \ar[r] \ar[d] & M_2
\otimes_R (\prod_{\alpha} Q_{\alpha}) \ar[r] \ar[d] & M_3 \otimes_R
(\prod_{\alpha} Q_{\alpha}) \ar[r] \ar[d] & 0 \\
0 \ar[r] & \prod_{\alpha}(M_1 \otimes Q_{\alpha}) \ar[r] & \prod_{\alpha}(M_2
\otimes Q_{\alpha}) \ar[r] & \prod_{\alpha}(M_3 \otimes Q_{\alpha})\ar[r] & 0
}
$$
两行都正合。因此 (1) 和 (2) 由命题 \ref{algebra-proposition-ML-tensor} 得出。
\end{proof}

\begin{lemma}
\label{algebra-lemma-quotient-module-ML}
设 $M_1 \to M_2 \to M_3 \to 0$ 是 $R$-模的正合列。
若 $M_1$ 有限生成且 $M_2$ 是 Mittag-Leffler 的，则 $M_3$ 是
Mittag-Leffler 的。
\end{lemma}

\begin{proof}
对任意一族 $R$-模 $(Q_{\alpha})_{\alpha \in A}$，由于张量积右正合，
有交换图表
$$
\xymatrix{
M_1 \otimes_R (\prod_{\alpha} Q_{\alpha}) \ar[r] \ar[d] & M_2
\otimes_R (\prod_{\alpha} Q_{\alpha}) \ar[r] \ar[d] & M_3 \otimes_R
(\prod_{\alpha} Q_{\alpha}) \ar[r] \ar[d] & 0 \\
\prod_{\alpha}(M_1 \otimes Q_{\alpha}) \ar[r] & \prod_{\alpha}(M_2
\otimes Q_{\alpha}) \ar[r] & \prod_{\alpha}(M_3 \otimes Q_{\alpha})\ar[r] & 0
}
$$
两行都正合。由命题 \ref{algebra-proposition-fg-tensor}，左侧竖直箭头满射。
由命题 \ref{algebra-proposition-ML-tensor}，中间竖直箭头单射。
图追踪说明右侧竖直箭头单射。故由命题
\ref{algebra-proposition-ML-tensor}，$M_3$ 是 Mittag-Leffler 的。
\end{proof}


\begin{lemma}
\label{algebra-lemma-colimit-universally-injective-ML}
若 $M = \colim M_i$ 是由 Mittag-Leffler $R$-模 $M_i$ 组成、且转移映射
普遍单射的有向系的余极限，则 $M$ 是 Mittag-Leffler 的。
\end{lemma}

\begin{proof}
设 $(Q_{\alpha})_{\alpha \in A}$ 是一族 $R$-模。必须说明
$M \otimes_R (\prod Q_\alpha) \to \prod M \otimes_R Q_\alpha$
单射；而对每个 $i$，已知
$M_i \otimes_R (\prod Q_\alpha) \to \prod M_i \otimes_R Q_\alpha$
单射，见命题 \ref{algebra-proposition-ML-tensor}。由于 $\otimes$ 与滤过余极限
交换，只须说明 $\prod M_i \otimes_R Q_\alpha \to \prod M \otimes_R Q_\alpha$
单射。这是显然的，因为根据转移映射普遍单射这一假设，每个映射
$M_i \otimes_R Q_\alpha \to M \otimes_R Q_\alpha$ 都是单射。
\end{proof}

\begin{lemma}
\label{algebra-lemma-direct-sum-ML}
若 $M = \bigoplus_{i \in I} M_i$ 是 $R$-模的直和，则 $M$ 是
Mittag-Leffler 的，当且仅当每个 $M_i$ 都是 Mittag-Leffler 的。
\end{lemma}

\begin{proof}
“仅当”方向由引理 \ref{algebra-lemma-pure-submodule-ML} (1) 以及分裂短正合列
普遍正合这一事实得出。反方向由引理
\ref{algebra-lemma-colimit-universally-injective-ML} 得出，但也可如下直接论证。
先注意，若 $I$ 有限，则结论由引理
\ref{algebra-lemma-pure-submodule-ML} (2) 得出。对一般的 $I$，若所有 $M_i$
都是 Mittag-Leffler 的，则通过验证命题
\ref{algebra-proposition-ML-characterization} 的条件 (1)，证明 $M$ 也如此。
设 $f: P \to M$ 是从有限表示模 $P$ 出发的映射。那么对 $I$ 的某个
有限子集 $I'$，$f$ 分解为
$P \xrightarrow{f'} \bigoplus_{i' \in I'} M_{i'} \hookrightarrow
\bigoplus_{i \in I} M_i$。由有限情形，
$\bigoplus_{i' \in I'} M_{i'}$ 是 Mittag-Leffler 的，故存在有限表示模
$Q$ 和映射 $g: P \to Q$，使 $g$ 与 $f'$ 相互支配。
于是 $g$ 与 $f$ 也相互支配。
\end{proof}

\begin{lemma}
\label{algebra-lemma-flat-ML-over-ML-ring}
设 $R \to S$ 是环映射。设 $M$ 是 $S$-模。若 $S$ 作为 $R$-模是
Mittag-Leffler 的，且 $M$ 作为 $S$-模是平坦并且 Mittag-Leffler 的，
则 $M$ 作为 $R$-模是 Mittag-Leffler 的。
\end{lemma}

\begin{proof}
我们由命题 \ref{algebra-proposition-ML-tensor} 的刻画推出这一点。
具体地，假设 $Q_\alpha$ 是一族 $R$-模。考虑复合
$$
\xymatrix{
M \otimes_R \prod_\alpha Q_\alpha =
M \otimes_S S \otimes_R \prod_\alpha Q_\alpha \ar[d] \\
M \otimes_S \prod_\alpha (S \otimes_R Q_\alpha) \ar[d] \\
\prod_\alpha (M \otimes_S S \otimes_R Q_\alpha) =
\prod_\alpha (M \otimes_R Q_\alpha)
}
$$
第一个箭头是单射，因为 $M$ 在 $S$ 上平坦，且 $S$ 在 $R$ 上是
Mittag-Leffler 的；第二个箭头是单射，因为 $M$ 在 $S$ 上是
Mittag-Leffler 的。因此 $M$ 在 $R$ 上是 Mittag-Leffler 的。
\end{proof}

\section{凝聚环}
\label{algebra-section-coherent}

\noindent
我们利用关于交换 $\prod$ 与 $\otimes$ 的讨论，确定哪些环具有如下性质：
平坦模的直积仍然平坦。结果表明，这些恰是所谓的凝聚环。
你或许更熟悉环化空间上的凝聚 $\mathcal{O}_X$-模这一概念，见
Modules，第 \ref{modules-section-coherent} 节。

\begin{definition}
\label{algebra-definition-coherent}
设 $R$ 是环。设 $M$ 是 $R$-模。
\begin{enumerate}
\item 若 $M$ 有限生成，且 $M$ 的每个有限生成子模在 $R$ 上都是
有限表示的，则称其为{\it 凝聚模}。
\item 若 $R$ 作为自身上的模是凝聚的，则称其为{\it 凝聚环}。
\end{enumerate}
\end{definition}

\noindent
因此，一个环是凝聚的，当且仅当每个有限生成理想作为模都是有限表示的。

\begin{example}
\label{algebra-example-valuation-ring-coherent}
赋值环是凝聚环。事实上，每个非零有限生成理想都是主理想（引理
\ref{algebra-lemma-characterize-valuation-ring}）；又因为赋值环是整环，
该理想是自由的，因而是有限表示的。
\end{example}

\noindent
凝聚模范畴是阿贝尔范畴。

\begin{lemma}
\label{algebra-lemma-coherent}
设 $R$ 是环。
\begin{enumerate}
\item 凝聚模的有限子模是凝聚的。
\item 设 $\varphi : N \to M$ 是从有限模到凝聚模的同态。那么
$\Ker(\varphi)$ 有限，$\Im(\varphi)$ 凝聚，并且
$\Coker(\varphi)$ 凝聚。
\item 设 $\varphi : N \to M$ 是凝聚模之间的同态。那么
$\Ker(\varphi)$ 和 $\Coker(\varphi)$ 都是凝聚模。
\item 给定 $R$-模的短正合列
$0 \to M_1 \to M_2 \to M_3 \to 0$，若其中两个模凝聚，
则第三个也凝聚。
\end{enumerate}
\end{lemma}

\begin{proof}
第一条陈述由定义立即得出。

\medskip\noindent
设 $\varphi : N \to M$ 满足 (2) 的假设。首先，
$\Im(\varphi)$ 有限，因而由 (1) 是凝聚的。特别地，
$\Im(\varphi)$ 有限表示；于是把引理 \ref{algebra-lemma-extension} 应用于
正合列 $0 \to \Ker(\varphi) \to N \to \Im(\varphi) \to 0$，
可知 $\Ker(\varphi)$ 有限。为证明 $\Coker(\varphi)$ 凝聚，
设 $E \subset \Coker(\varphi)$ 是有限子模，并令 $E'$ 为它在 $M$
中的逆像。由正合列
$0 \to \Im(\varphi) \to E' \to E \to 0$，并且由于
% P01-E0456: the source invokes finite Ker(phi), but this exact sequence requires finite Im(phi); see ERRATA_CANDIDATES.jsonl.
$\Ker(\varphi)$ 有限，由引理 \ref{algebra-lemma-extension} 可知
$E' \subset M$ 有限，因而因为 $M$ 凝聚，它是有限表示的。
同一正合列继而说明 $E$ 有限表示，断言得证。

\medskip\noindent
(3) 立即由 (1) 和 (2) 得出。

\medskip\noindent
设 $0 \to M_1 \xrightarrow{i} M_2 \xrightarrow{p} M_3 \to 0$
是 (4) 中的 $R$-模短正合列。只须证明：若 $M_1$ 和 $M_3$ 凝聚，
则 $M_2$ 也凝聚。由引理 \ref{algebra-lemma-extension}，$M_2$ 有限。
设 $N_2 \subset M_2$ 是有限子模。令
$N_3 = p(N_2) \subset M_3$，并令
$N_1 = i^{-1}(N_2) \subset M_1$。我们有正合列
$0 \to N_1 \to N_2 \to N_3 \to 0$。显然 $N_3$ 有限（因为它是
$N_2$ 的商），因而是有限表示的（因为它是 $M_3$ 的有限子模）。
由引理 \ref{algebra-lemma-extension} (5)，$N_1$ 有限，因而是有限表示的
（因为它是 $M_1$ 的有限子模）。再由引理
\ref{algebra-lemma-extension} (2)，$N_2$ 有限表示。
\end{proof}

\begin{lemma}
\label{algebra-lemma-coherent-ring}
设 $R$ 是环。若 $R$ 凝聚，则一个模凝聚，当且仅当它有限表示。
\end{lemma}

\begin{proof}
显然，凝聚模是有限表示的（在任意环上都是如此）。反之，若 $R$ 凝聚，
则 $R^{\oplus n}$ 凝聚，且任意映射
$R^{\oplus m} \to R^{\oplus n}$ 的余核也凝聚，见引理
\ref{algebra-lemma-coherent}。
\end{proof}

\begin{lemma}
\label{algebra-lemma-Noetherian-coherent}
Noether 环是凝聚环。
\end{lemma}

\begin{proof}
由引理 \ref{algebra-lemma-Noetherian-finite-type-is-finite-presentation}，
任意有限 $R$-模都是有限表示的。特别地，$R$ 的任意理想都是有限表示的。
\end{proof}

\begin{proposition}
\label{algebra-proposition-characterize-coherent}
\begin{reference}
这是 \cite[Theorem 2.1]{Chase}。
\end{reference}
设 $R$ 是环。下列条件等价：
\begin{enumerate}
\item $R$ 凝聚；
\item 平坦 $R$-模的任意直积都平坦；
\item 对每个集合 $A$，模 $R^A$ 都平坦。
\end{enumerate}
\end{proposition}

\begin{proof}
假设 $R$ 凝聚，并设 $Q_\alpha$（$\alpha \in A$）是一族平坦
$R$-模。必须证明，对 $R$ 的每个有限生成理想 $I$，映射
$I \otimes_R \prod_\alpha Q_\alpha \to \prod Q_\alpha$ 都单射，
见引理 \ref{algebra-lemma-flat}。由于 $R$ 凝聚，$I$ 是有限表示 $R$-模。
因此由命题 \ref{algebra-proposition-fp-tensor}，有
$I \otimes_R \prod_\alpha Q_\alpha = \prod I \otimes_R Q_\alpha$。
所需的单射性由此得出，因为根据 $Q_\alpha$ 的平坦性，
$I \otimes_R Q_\alpha \to Q_\alpha$ 单射。

\medskip\noindent
(2) $\Rightarrow$ (3) 是显然的。

\medskip\noindent
假设对每个集合 $A$，$R$-模 $R^A$ 都平坦。设 $I$ 是 $R$ 的有限生成
理想。由假设，$I \otimes_R R^A \to R^A$ 单射。由命题
\ref{algebra-proposition-fg-tensor} 以及 $I$ 的有限性，其像等于 $I^A$。
因此对每个集合 $A$，都有 $I \otimes_R R^A = I^A$；再由命题
\ref{algebra-proposition-fp-tensor} 可知 $I$ 有限表示。
\end{proof}

\section{Mittag-Leffler 模的例与反例}
\label{algebra-section-examples}

\noindent
本节最后给出一些 Mittag-Leffler 模的例与反例。

\begin{example}
\label{algebra-example-ML}
Mittag-Leffler 模。
\begin{enumerate}
\item 任意有限表示模都是 Mittag-Leffler 的。例如，这由命题
\ref{algebra-proposition-ML-characterization} (1) 得出。一般而言，有限生成模
是 Mittag-Leffler 的，当且仅当它是有限表示的。这由命题
\ref{algebra-proposition-fg-tensor}、\ref{algebra-proposition-fp-tensor} 和
\ref{algebra-proposition-ML-tensor} 得出。
\item 自由模是 Mittag-Leffler 的，因为它满足命题
\ref{algebra-proposition-ML-characterization} 的条件 (1)。
\item 由上一个例子和引理 \ref{algebra-lemma-direct-sum-ML}，投射模是
Mittag-Leffler 的。
\end{enumerate}
\end{example}

\noindent
我们还想把 Noether 环 $R$ 上的幂级数环加入例子列表。这将由如下引理推出。

\begin{lemma}
\label{algebra-lemma-flat-ML-criterion}
设 $M$ 是平坦 $R$-模。下列条件等价：
\begin{enumerate}
\item $M$ 是 Mittag-Leffler 的；
\item 若 $F$ 是有限自由 $R$-模，且 $x \in F \otimes_R M$，则存在
$F$ 的最小子模 $F'$，使得 $x \in F' \otimes_R M$。
\end{enumerate}
\end{lemma}

\begin{proof}
(1) $\Rightarrow$ (2) 是引理 \ref{algebra-lemma-minimal-contains} 的特殊情形。
假设 (2)。由定理 \ref{algebra-theorem-lazard}，可把 $M$ 写成有限自由
$R$-模有向系 $(M_i, f_{ij})$ 的余极限
$M = \colim_{i \in I} M_i$。由评注 \ref{algebra-remark-flat-ML}，只须说明
对偶逆系
$$
(\Hom_R(M_i, R), \Hom_R(f_{ij}, R))
$$
是 Mittag-Leffler 的。
换言之，固定 $i \in I$，并对 $j \geq i$，令 $Q_j$ 为
$\Hom_R(M_j, R)
\to \Hom_R(M_i, R)$ 的像；必须证明诸 $Q_j$ 稳定。

\medskip\noindent
因为 $M_i$ 有限自由，对所有 $j$，都有等同
$\Hom_R(M_i, M_j) = \Hom_R(M_i, R) \otimes_R  M_j$。
利用诸 $M_j$ 自由这一事实可知，对 $j \geq i$，$Q_j$ 是满足
$f_{ij} \in Q_j
\otimes_R
M_j$ 的 $\Hom_R(M_i, R)$ 的最小子模。在等同
$\Hom_R(M_i, M) = \Hom_R(M_i, R)
\otimes_R
M$ 下，典范映射 $f_i: M_i \to M$ 属于
$\Hom_R(M_i, R)
\otimes_R M$。由关于 $M$ 的假设，存在 $\Hom_R(M_i, R)$ 的最小子模
$Q$，使得 $f_i \in Q \otimes_R M$。下面证明诸 $Q_j$ 稳定为 $Q$。

\medskip\noindent
对 $j \geq i$，有交换图表
$$
\xymatrix{
Q_j \otimes_R M_j \ar[r] \ar[d] & \Hom_R(M_i, R) \otimes_R M_j
\ar[d] \\
Q_j \otimes_R M \ar[r] & \Hom_R(M_i, R) \otimes_R M.
}
$$
由于 $f_{ij} \in Q_j \otimes_R M_j$ 映到
$f_i \in \Hom_R(M_i, R)
\otimes_R M$，可知 $f_i \in Q_j \otimes_R M$。因此，由 $Q$ 的选择，
对所有 $j \geq i$，都有 $Q \subset Q_j$。

\medskip\noindent
因为诸 $Q_j$ 递减，且对所有 $j \geq i$ 都有 $Q \subset Q_j$，
要说明诸 $Q_j$ 稳定为 $Q$，只须找到 $j \geq i$，使
$Q_j \subset Q$。把这个典范映射看作下式的元素：
% P01-E0459: the source changes the established index set I to undefined J in both colimits; see ERRATA_CANDIDATES.jsonl.
$$
\Hom_R(M_i, R) \otimes_R M = \colim_{j \in J}
(\Hom_R(M_i, R) \otimes_R
M_j),
$$
则 $f_i$ 是当 $j \geq i$ 时诸 $f_{ij}$ 的余极限，并且 $f_i$ 也属于子模
$$
\colim_{j \in J} (Q \otimes_R M_j) \subset \colim_{j \in J}
(\Hom_R(M_i, R)
\otimes_R M_j).
$$
因此，对某个 $j \geq i$，$f_{ij}$ 属于 $Q \otimes_R M_j$。
由于 $Q_j$ 是满足 $f_{ij}
\in
Q_j \otimes_R M_j$ 的 $\Hom_R(M_i, R)$ 的最小子模，故
$Q_j\subset Q$。
\end{proof}

\begin{lemma}
\label{algebra-lemma-product-over-Noetherian-ring}
设 $R$ 是 Noether 环，$A$ 是集合。那么
$M = R^A$ 是平坦的 Mittag-Leffler $R$-模。
\end{lemma}

\begin{proof}
结合引理 \ref{algebra-lemma-Noetherian-coherent} 与命题
\ref{algebra-proposition-characterize-coherent}，可知 $M$ 在 $R$ 上平坦。
下面说明 $M$ 满足引理 \ref{algebra-lemma-flat-ML-criterion} 的条件。
设 $F$ 是有限自由 $R$-模。若 $F'$ 是 $F$ 的任意子模，则由于 $R$
是 Noether 环，它是有限表示的。因此由命题
\ref{algebra-proposition-fp-tensor}，有交换图表
$$
\xymatrix{
F' \otimes_R M \ar[r] \ar[d]^{\cong} & F \otimes_R M \ar[d]^{\cong} \\
(F')^A \ar[r] & F^A
}
$$
借此可把映射 $F' \otimes_R M \to F \otimes_R M$ 与
$(F')^A \to F^A$ 等同。因此，若 $x \in F \otimes_R M$ 对应于
$(x_\alpha) \in F^A$，则由诸 $x_\alpha$ 生成的 $F$ 的子模 $F'$
是 $F$ 的满足 $x \in F' \otimes_R M$ 的最小子模。
% P01-E0460: the translated sentence renders the clear intended meaning of the malformed source phrase; see ERRATA_CANDIDATES.jsonl.
\end{proof}

\begin{lemma}
\label{algebra-lemma-power-series-ML}
设 $R$ 是 Noether 环，$n$ 是正整数。那么 $R$-模
$M = R[[t_1, \ldots, t_n]]$ 平坦且是 Mittag-Leffler 的。
\end{lemma}

\begin{proof}
作为 $R$-模，对某个（可数）集合 $A$，有 $M = R^A$。
因此该引理是引理 \ref{algebra-lemma-product-over-Noetherian-ring} 的特殊情形。
\end{proof}

\begin{example}
\label{algebra-example-not-ML}
非 Mittag-Leffler 模。
\begin{enumerate}
\item 由例 \ref{algebra-example-Q-not-ML} 和命题
\ref{algebra-proposition-ML-tensor}，$\mathbf{Q}$ 不是 Mittag-Leffler
$\mathbf{Z}$-模。
\item 下文将证明（定理 \ref{algebra-theorem-projectivity-characterization}）：
对平坦可数生成模，投射性等价于 Mittag-Leffler 性。因此，任意平坦、
可数生成且非投射的模 $M$ 都是非 Mittag-Leffler 模的例子。
这样的例子见评注 \ref{algebra-remark-warning}。
\item 设 $k$ 是域。令 $R = k[[x]]$。$R$-模
$M = \prod_{n \in \mathbf{N}} R/(x^n)$ 不是 Mittag-Leffler 的。
具体地，考虑元素 $\xi = (\xi_1, \xi_2, \xi_3, \ldots)$，其定义为
$\xi_{2^m} = x^{2^{m - 1}}$，并在其余情形取 $\xi_n = 0$，所以
$$
\xi = (0, x, 0, x^2, 0, 0, 0, x^4, 0, 0, 0, 0, 0, 0, 0, x^8, \ldots)
$$
当 $m \gg 0$ 时，$\xi$ 在 $M/x^{2^m}M$ 中的零化子由
$x^{2^{m - 1}}$ 生成。但是，若 $M$ 是 Mittag-Leffler 的，
则存在有限 $R$-模 $Q$ 和元素 $\xi' \in Q$，使得对所有
$l \geq 1$，$\xi'$ 在 $Q/x^l Q$ 中的零化子与 $\xi$ 在
$M/x^l M$ 中的零化子一致，见命题
\ref{algebra-proposition-ML-characterization} (1)。现在可以证明，存在整数
$a \geq 0$，使得对所有 $l \gg 0$，$\xi'$ 在 $Q/x^l Q$ 中的零化子
由 $x^a$ 或 $x^{l - a}$ 生成（取决于 $\xi' \in Q$ 是否为挠元）。
结合以上结论，会得到对所有 $l = 2^m >> 0$，都有
% P01-E0462: frozen target preserves >>, which renders as two greater-than signs rather than \gg; see ERRATA_CANDIDATES.jsonl.
$a = l/2$ 或 $l - a = l/2$，这是荒谬的。
\item 同一论证说明，$\bigoplus_{n \in \mathbf{N}} R/(x^n)$ 的
$(x)$-进完备化在 $R = k[[x]]$ 上不是 Mittag-Leffler 的
（提示：$\xi$ 实际上是该完备化的一个元素）。
\item 令 $R = k[a, b]/(a^2, ab, b^2)$。令 $S$ 是具有表示
$S = R[t]/(at - b)$ 的有限表示 $R$-代数。那么作为 $R$-模，
$S$ 可数生成且不可分解（细节略）。另一方面，$R$ 是 Artin 局部环，
因而是完备局部环，进而是 Hensel 局部环，见引理
\ref{algebra-lemma-complete-henselian}。若 $S$ 作为 $R$-模是 Mittag-Leffler 的，
则由引理 \ref{algebra-lemma-split-ML-henselian}，它会是有限 $R$-模的直和。
因此，$S$ 作为 $R$-模不是 Mittag-Leffler 的。
\end{enumerate}
\end{example}

\section{可数生成的 Mittag-Leffler 模}
\label{algebra-section-ML-countable}

\noindent
事实证明，可数生成的 Mittag-Leffler 模具有特别简单的结构。

\begin{lemma}
\label{algebra-lemma-ML-countable-colimit}
设 $M$ 是 $R$-模。写成 $M = \colim_{i \in I} M_i$，其中
$(M_i, f_{ij})$ 是由有限表示 $R$-模组成的有向系。若 $M$ 是
Mittag-Leffler 且可数生成，则存在有向可数子集 $I' \subset I$，
使得 $M \cong \colim_{i \in I'} M_i$。
\end{lemma}

\begin{proof}
设 $x_1, x_2, \ldots$ 是 $M$ 的一个可数生成集。对每个 $x_n$，
选取 $i \in I$，使 $x_n$ 属于典范映射 $f_i:
M_i \to M$ 的像；令 $I'_{0} \subset I$ 为所有这些 $i$ 的集合。
现在，由于 $M$ 是 Mittag-Leffler 的，对每个 $i \in I'_{0}$，
可以选取 $j \in I$，使 $j \geq i$，并且对所有 $k \geq i$，
$f_{ij}: M_i \to M_j$ 都经由 $f_{ik}: M_i \to M_k$ 分解
（命题 \ref{algebra-proposition-ML-characterization} 的条件 (3)）；
令 $I'_1$ 为 $I'_0$ 与所有这些 $j$ 的并。由于 $I'_1$ 可数，
可以把它扩大为可数有向集 $I'_{2} \subset I$。现在对 $I'_{2}$
施行与对 $I'_{0}$ 相同的步骤，得到新的可数集
$I'_{3} \subset I$。再把 $I'_{3}$ 扩大为可数有向集 $I'_{4}$。
% P01-E0464: the frozen source reverses the even/odd stages relative to its explicit I'_0 -> I'_1 -> I'_2 construction; see ERRATA_CANDIDATES.jsonl.
以这种方式继续下去——若 $\ell$ 为奇数，则对每个
$ i \in I'_{\ell}$ 加入命题
\ref{algebra-proposition-ML-characterization} (3) 中的一个 $j$；若 $\ell$
为偶数，则把 $I'_{\ell}$ 扩大为有向集——便得到一列子集
$I'_{\ell} \subset I$，其中 $\ell \geq 0$。并集 $I' =
\bigcup I'_{\ell}$ 满足：
\begin{enumerate}
\item $I'$ 可数且有向；
\item 每个 $x_n$ 都属于某个 $i \in I'$ 所对应的
$f_i: M_i \to M$ 的像；
\item 若 $i \in I'$，则存在 $j \in I'$，使得 $j \geq i$，并且对所有
满足 $k \geq i$ 的 $k \in I$，$f_{ij}:
M_i \to M_j$ 都经由 $f_{ik}: M_i \to M_k$ 分解。特别地，对
$k \geq i$ 有 $\Ker(f_{ik}) \subset \Ker(f_{ij})$。
\end{enumerate}
我们断言典范映射 $\colim_{i \in I'} M_i \to
\colim_{i
\in I} M_i = M$ 是同构。由 (2)，它是满射。为证其为单射，
设 $x \in \colim_{i \in I'} M_i$ 在 $\colim_{i \in
I} M_i$ 中映为 $0$。
用某个 $i \in I'$ 的元素 $\tilde{x} \in M_i$ 表示 $x$，
这意味着对某个 $k \in I, k \geq i$，有
$f_{ik}(\tilde{x}) = 0$。但是，由 (3)，存在
$j \in I', j \geq i,$，使 $f_{ij}(\tilde{x}) = 0$。因此，
$x = 0$ 在 $\colim_{i \in I'} M_i$ 中成立。
\end{proof}

\noindent
引理 \ref{algebra-lemma-ML-countable-colimit}
表明，环 $R$ 上可数生成的 Mittag-Leffler 模 $M$ 是如下系统的余极限：
$$
M_1 \to M_2 \to M_3 \to M_4 \to \ldots
$$
其中每个 $M_n$ 都是有限表示 $R$-模。为说明这一点，仿照引理
\ref{algebra-lemma-ML-limit-nonempty}
% P01-E0465: the frozen source gives the direct-system cofinal order as (N, >=), conflicting with the displayed forward system and the referenced increasing construction; see ERRATA_CANDIDATES.jsonl.
的证明，可知可数有向集具有一个与 $(\mathbf{N}, \geq)$
同构的共尾子集。设
$R = k[x_1, x_2, x_3, \ldots]$ 且 $M = R/(x_i)$。则 $M$
有限生成但非有限表示，因而不是 Mittag-Leffler 的（见例
\ref{algebra-example-ML} 的 (1)）。
不过，当然可以取
$M_n = R/(x_1, \ldots, x_n)$，从而写成 $M = \colim_n M_n$；
因此，能把 $M$ 写成这样的极限并不蕴含 $M$ 是 Mittag-Leffler 的。

\begin{lemma}
\label{algebra-lemma-ML-countable}
设 $R$ 是环。
设 $M$ 是 $R$-模。
假设 $M$ 是 Mittag-Leffler 且可数生成。
对任意 $R$-模映射 $f : P \to M$，若 $P$ 有限生成，则存在
自同态 $\alpha : M \to M$，使得
\begin{enumerate}
\item $\alpha : M \to M$ 经由一个有限表示 $R$-模分解，并且
\item $\alpha \circ f = f$。
\end{enumerate}
\end{lemma}

\begin{proof}
把 $M = \colim_{i \in I} M_i$ 写成由有限表示 $R$-模组成的
有向余极限，其中 $I$ 可数，见引理
\ref{algebra-lemma-ML-countable-colimit}。
转移映射记为 $f_{ij}$，并用 $f_i : M_i \to M$
表示映入 $M$ 的典范映射。置 $N = \prod_{s \in I} M_s$。记
$$
M_i^* = \Hom_R(M_i, N) = \prod\nolimits_{s \in I} \Hom_R(M_i, M_s)
$$
于是 $(M_i^*)$ 是指标集为 $I$ 的 $R$-模逆系。
注意 $\Hom_R(M, N) = \lim M_i^*$。
由于 $M$ 是 Mittag-Leffler 的，对每个
$i \in I$，可以找到指标 $k(i) \geq i$，使得
$$
E_i := \bigcap\nolimits_{i' \geq i} \Im(M_{i'}^* \to M_i^*)
=
\Im(M_{k(i)}^* \to M_i^*)
$$
选定 $j \in I$，使 $\Im(P \to M) \subset \Im(M_j \to M)$。
这是可能的，因为 $P$ 有限生成。置 $k = k(j)$。令
$x = (0, \ldots, 0, \text{id}_{M_k}, 0, \ldots, 0) \in M_k^*$，并注意
它映到 $y = (0, \ldots, 0, f_{jk}, 0, \ldots, 0) \in M_j^*$。
由 $k$ 的选取可知 $y \in E_j$。由例
\ref{algebra-example-ML-surjective-maps}，
对每个 $i \geq j$，转移映射 $E_i \to E_j$ 都是满射，并且
$\lim E_i = \lim M_i^* = \Hom_R(M, N)$。因此，引理
\ref{algebra-lemma-ML-limit-nonempty}
保证存在元素 $z \in \Hom_R(M, N)$，
它在 $E_j \subset M_j^*$ 中映到 $y$。令 $z_k$ 为 $z$ 的第 $k$ 个分量。
那么 $z_k : M \to M_k$ 是同态，并且下图
$$
\xymatrix{
M \ar[r]_{z_k} & M_k \\
M_j \ar[ru]_{f_{jk}} \ar[u]^{f_j}
}
$$
交换。令 $\alpha : M \to M$ 为复合
$f_k \circ z_k : M \to M_k \to M$。
按构造，$\alpha$ 经由有限表示模分解，并且
$\alpha \circ f_j = f_j$。由于 $f$ 的像包含在 $f_j$ 的像中，
这也蕴含 $\alpha \circ f = f$。
\end{proof}

\noindent
稍后将看到（见引理
\ref{algebra-lemma-split-ML-henselian}），
引理 \ref{algebra-lemma-ML-countable}
意味着 Hensel 局部环上的可数生成 Mittag-Leffler 模是有限表示模的直和。

\section{投射模的刻画}
\label{algebra-section-characterize-projective}

\noindent
本节的目标是证明，一个模是投射的，当且仅当它平坦、
是 Mittag-Leffler 的，并且是可数生成模的直和
（见下文定理 \ref{algebra-theorem-projectivity-characterization}）。

\begin{lemma}
\label{algebra-lemma-countgen-projective}
设 $M$ 是 $R$-模。若 $M$ 平坦、是 Mittag-Leffler 的且可数生成，
则 $M$ 是投射的。
\end{lemma}

\begin{proof}
由 Lazard 定理（定理 \ref{algebra-theorem-lazard}），可以把 $M =
\colim_{i
\in I} M_i$ 写成由有限自由 $R$-模构成、指标集为集合 $I$ 的有向系
$(M_i, f_{ij})$。由引理 \ref{algebra-lemma-ML-countable-colimit}，
可以假设 $I$ 可数。现在令
$$
0 \to N_1 \to N_2 \to N_3 \to 0
$$
是 $R$-模的正合列。我们必须证明，施用
$\Hom_R(M, -)$
后仍保持正合性。由于 $M_i$ 是有限自由的，
$$
0 \to \Hom_R(M_i, N_1) \to \Hom_R(M_i, N_2) \to
\Hom_R(M_i, N_3) \to 0
$$
对每个 $i$ 都是正合的。由于 $M$ 是 Mittag-Leffler 的，
$(\Hom_R(M_i, N_{1}))$ 是 Mittag-Leffler 逆系。因此由引理
\ref{algebra-lemma-ML-exact-sequence}，
$$
0 \to \lim_{i \in I} \Hom_R(M_i, N_1) \to
\lim_{i \in I} \Hom_R(M_i, N_2) \to
\lim_{i \in I} \Hom_R(M_i, N_3) \to 0
$$
正合。但对任意 $R$-模 $N$，都有函子性的同构
$\Hom_R(M, N) \cong \lim_{i \in I} \Hom_R(M_i, N)$，所以
$$
0 \to \Hom_R(M, N_1) \to \Hom_R(M, N_2) \to
\Hom_R(M, N_3) \to 0
$$
正合。
\end{proof}

\begin{remark}
\label{algebra-remark-characterize-projective}
若没有可数生成这一假设，引理 \ref{algebra-lemma-countgen-projective}
不成立。例如，$\mathbf Z$-模 $M =
\mathbf{Z}[[x]]$ 平坦且是 Mittag-Leffler 的，但并不投射。
由引理 \ref{algebra-lemma-power-series-ML}，它是 Mittag-Leffler 的。
自由 Abel 群的子群是自由的，因而投射 $\mathbf Z$-模事实上是自由的，
其子模也自由。因此，为证明 $M$ 不投射，只须构造一个非自由子模。
固定一个素数 $p$，并考虑由如下幂级数组成的子模 $N$：
$f(x) = \sum a_i x^i$，使得对每个整数 $m
\geq 1$，除有限多个 $i$ 外，$p^m$ 都整除 $a_i$。于是对所有
$a_i \in \mathbf{Z}$，$\sum a_i p^i x^i$ 都属于 $N$，所以 $N$ 不可数。
因此，若 $N$ 自由，则它的秩不可数，并且 $N/pN$ 在
$\mathbf{Z}/p$ 上的维数不可数。这并不成立，因为诸元素
$x^i \in N/pN$（其中 $i \geq 0$）张成 $N/pN$。
\end{remark}

\begin{theorem}
\label{algebra-theorem-projectivity-characterization}
设 $M$ 是 $R$-模。那么 $M$ 投射，当且仅当
\begin{enumerate}
\item $M$ 平坦；
\item $M$ 是 Mittag-Leffler 的；
\item $M$ 是可数生成 $R$-模的直和。
\end{enumerate}
\end{theorem}

\begin{proof}
先设 $M$ 投射。则 $M$ 是自由模的直和因子，因此 $M$ 平坦且是
Mittag-Leffler 的，因为这些性质传递到直和因子。由 Kaplansky 定理
（定理 \ref{algebra-theorem-projective-direct-sum}），$M$ 满足 (3)。

\medskip\noindent
反过来，设 $M$ 满足 (1)--(3)。由于平坦性和 Mittag-Leffler 性传递到
直和因子，$M$ 是平坦、Mittag-Leffler、可数生成 $R$-模的直和。
引理 \ref{algebra-lemma-countgen-projective}
表明 $M$ 是投射模的直和。因此 $M$ 是投射的。
\end{proof}

\begin{lemma}
\label{algebra-lemma-ML-ui-descent}
设 $f: M \to N$ 是 $R$-模的普遍单射映射。假设
$M$ 是可数生成 $R$-模的直和，并假设 $N$ 平坦且是 Mittag-Leffler 的。
则 $M$ 是投射的。
\end{lemma}

\begin{proof}
由引理 \ref{algebra-lemma-ui-flat-domain} 和
\ref{algebra-lemma-pure-submodule-ML}，
$M$ 平坦且是 Mittag-Leffler 的，故结论由定理
\ref{algebra-theorem-projectivity-characterization} 得出。
\end{proof}

\begin{lemma}
\label{algebra-lemma-universally-injective-submodule-powerseries}
设 $R$ 是 Noether 环，$M$ 是 $R$-模。假设 $M$ 是可数生成
$R$-模的直和，并假设对某个 $n$ 存在普遍单射映射
$M \to R[[t_1, \ldots, t_n]]$。则 $M$ 是投射的。
\end{lemma}

\begin{proof}
这由引理 \ref{algebra-lemma-ML-ui-descent} 和
\ref{algebra-lemma-power-series-ML} 得出。
\end{proof}

\section{模的性质的上升}
\label{algebra-section-ascent}

\noindent
定理 \ref{algebra-theorem-projectivity-characterization}
中模的所有性质都沿任意环映射上升：

\begin{lemma}
\label{algebra-lemma-ascend-properties-modules}
设 $R \to S$ 是环映射。设 $M$ 是 $R$-模。则：
\begin{enumerate}
\item 若 $M$ 平坦，则 $S$-模 $M \otimes_R S$ 平坦。
\item 若 $M$ 是 Mittag-Leffler 的，则 $S$-模 $M \otimes_R S$
是 Mittag-Leffler 的。
\item 若 $M$ 是可数生成 $R$-模的直和，则 $S$-模
$M \otimes_R S$ 是可数生成 $S$-模的直和。
\item 若 $M$ 投射，则 $S$-模 $M \otimes_R S$ 投射。
\end{enumerate}
\end{lemma}

\begin{proof}
除 (2) 外都很显然。对 (2)，使用命题
\ref{algebra-proposition-ML-characterization} 中 Mittag-Leffler 性的表述 (3)，
以及张量运算与取余极限可交换这一事实。也可以使用命题
\ref{algebra-proposition-ML-tensor} 的刻画。
\end{proof}

\section{模的性质的下降}
\label{algebra-section-descent}

\noindent
我们讨论定理 \ref{algebra-theorem-projectivity-characterization}
中刻画投射性的诸性质的忠实平坦下降。
在平坦性成立时，Mittag-Leffler 模这一性质可以下降：

\begin{lemma}
\label{algebra-lemma-ffdescent-ML}
\begin{reference}
Juan Pablo Acosta Lopez 于 2014 年 12 月 20 日发送的电子邮件。
\end{reference}
设 $R \to S$ 是忠实平坦环映射。设 $M$ 是 $R$-模。若
$S$-模 $M \otimes_R S$ 是 Mittag-Leffler 的，则 $M$
是 Mittag-Leffler 的。
\end{lemma}

\begin{proof}
把 $M = \colim_{i\in I} M_i$ 写成有限表示
$R$-模 $M_i$ 的有向余极限。
由命题 \ref{algebra-proposition-ML-characterization} 可知，我们必须证明：
对每个 $i \in I$，存在 $i \leq j$、$j\in I$，使得
$M_i\rightarrow M_j$ 支配 $M_i\rightarrow M$。

\medskip\noindent
取 $N$ 为推出
$$
\xymatrix{
M_i  \ar[r] \ar[d] & M_j \ar[d] \\
M \ar[r] & N
}
$$
于是由引理 \ref{algebra-lemma-domination-universally-injective}，
该引理等价于存在 $j$，使 $M_j\rightarrow N$ 普遍单射。
注意，与 $S$ 张量所得的图表
$$
\xymatrix{
M_i\otimes_R S  \ar[r] \ar[d] & M_j\otimes_R S \ar[d] \\
M\otimes_R S \ar[r] & N\otimes_R S
}
$$
是推出图表。由于
$M \otimes_R S = \colim_{i\in I} M_i \otimes_R S$
把 $M\otimes_R S$ 表为有限表示 $S$-模的余极限，并且
$M\otimes_R S$ 是 Mittag-Leffler 的，所以存在 $j \geq i$，
使 $M_j\otimes_R S\rightarrow N\otimes_R S$ 普遍单射。
再利用 $R\rightarrow S$ 忠实平坦，即知
$M_j\rightarrow N$ 也普遍单射。
\end{proof}

\begin{lemma}
\label{algebra-lemma-ffdescent-countable}
设 $R \to S$ 是忠实平坦环映射。设 $M$ 是 $R$-模。若
$S$-模 $M \otimes_R S$ 可数生成，则 $M$ 可数生成。
\end{lemma}

\begin{proof}
设 $M \otimes_R S$ 由诸元素
$y_i$（$i = 1, 2, 3, \ldots$）生成。写成
$y_i = \sum_{j = 1, \ldots, n_i}
x_{ij} \otimes s_{ij}$，其中 $n_i \geq 0$、$x_{ij} \in M$
且 $s_{ij} \in S$。记 $M' \subset M$ 为由可数个元素 $x_{ij}$
生成的子模。于是 $M' \otimes_R S \to M \otimes_R S$ 是满射，
因为它的像包含诸生成元 $y_i$。由于 $S$ 在 $R$ 上忠实平坦，
可得 $M' = M$，如所欲。
\end{proof}

\noindent
至此，可数生成投射模的忠实平坦下降很容易得出。

\begin{lemma}
\label{algebra-lemma-ffdescent-countable-projectivity}
设 $R \to S$ 是忠实平坦环映射。设 $M$ 是 $R$-模。
若 $S$-模 $M \otimes_R S$ 可数生成且投射，则 $M$
可数生成且投射。
\end{lemma}

\begin{proof}
这由引理 \ref{algebra-lemma-descend-properties-modules}、
\ref{algebra-lemma-ffdescent-ML}、\ref{algebra-lemma-ffdescent-countable}
和定理 \ref{algebra-theorem-projectivity-characterization} 得出。
\end{proof}

\noindent
剩下的只是利用 d\'evissage，把一般情形下投射性的下降约化到
可数生成情形。先给出两个简单引理。

\begin{lemma}
\label{algebra-lemma-lift-countably-generated-submodule}
设 $R \to S$ 是环映射，$M$ 是 $R$-模，并设 $Q$ 是
$M \otimes_R S$ 的可数生成 $S$-子模。则存在 $M$ 的可数生成
$R$-子模 $P$，使得
$\Im(P \otimes_R S \to M \otimes_R S)$ 包含 $Q$。
\end{lemma}

\begin{proof}
设 $y_1, y_2, \ldots$ 是 $Q$ 的生成元，并写成
$y_j = \sum_k x_{jk} \otimes s_{jk}$，其中 $x_{jk} \in M$
且 $s_{jk} \in S$。取 $P$ 为 $M$ 中由诸 $x_{jk}$ 生成的子模。
\end{proof}

\begin{lemma}
\label{algebra-lemma-adapted-submodule}
设 $R \to S$ 是环映射，$M$ 是 $R$-模。假设
$M
\otimes_R S = \bigoplus_{i \in I} Q_i$ 是可数生成 $S$-模 $Q_i$
的直和。若 $N$ 是 $M$ 的可数生成子模，则存在 $M$ 的可数生成子模
$N'$，使得 $N' \supset N$，且对某个子集 $I' \subset I$，
有 $\Im(N' \otimes_R S \to M \otimes_R S) =
\bigoplus_{i \in I'} Q_i$。
\end{lemma}

\begin{proof}
令 $N'_0 = N$。归纳构造可数生成子模的递增列
$N'_{\ell} \subset M$（$\ell = 0, 1, 2, \ldots$），满足：
若 $I'_{\ell}$ 是这样的 $i \in I$ 的集合，使得
$\Im(N'_{\ell} \otimes_R S \to M \otimes_R S)$
在 $Q_i$ 上的投影非零，则
$\Im(N'_{\ell + 1} \otimes_R S \to M \otimes_R
S)$ 对所有 $i \in I'_{\ell}$ 都包含 $Q_i$。为从
$N'_\ell$ 构造 $N'_{\ell + 1}$，令 $Q$ 为
$i \in I'_{\ell}$ 时诸（可数多个）$Q_i$ 之和，按引理
\ref{algebra-lemma-lift-countably-generated-submodule} 选取 $P$，再令
$N'_{\ell + 1} = N'_{\ell} + P$。构造出诸 $N'_{\ell}$ 后，只需取
$N' = \bigcup_{\ell} N'_{\ell}$ 和 $I' = \bigcup_{\ell} I'_{\ell}$。
\end{proof}

\begin{theorem}
\label{algebra-theorem-ffdescent-projectivity}
设 $R \to S$ 是忠实平坦环映射。设 $M$ 是 $R$-模。
若 $S$-模 $M \otimes_R S$ 投射，则 $M$ 投射。
\end{theorem}

\begin{proof}
我们将构造 $M$ 的 Kaplansky d\'evissage，以证明它是投射模的直和，
因而投射。由定理 \ref{algebra-theorem-projective-direct-sum}，可以把
$M \otimes_R S =
\bigoplus_{i \in I} Q_i$ 写成可数生成 $S$-模 $Q_i$ 的直和。
在 $M$ 上选取一个良序。我们将利用超限递归定义 $M$ 的子模递增族
$M_{\alpha}$，每个序数 $\alpha$ 对应一个子模，使得
$M_{\alpha} \otimes_R S$ 是某些 $Q_i$ 的直和。

\medskip\noindent
当 $\alpha = 0$ 时，令 $M_0 = 0$。若 $\alpha$ 是极限序数，并且
对所有 $\beta < \alpha$ 已定义 $M_{\beta}$，则定义
$M_\alpha = \bigcup_{\beta < \alpha} M_{\beta}$。由于每个
$\beta < \alpha$ 所对应的 $M_{\beta} \otimes_R S$
都是某些 $Q_i$ 的直和，$M_{\alpha} \otimes_R S$ 也如此。
若 $\alpha + 1$ 是后继序数且 $M_{\alpha}$ 已定义，则如下定义
$M_{\alpha + 1}$。若 $M_{\alpha} = M$，则令
$M_{\alpha +1} = M$。否则选取满足 $x \notin
M_{\alpha}$ 的最小 $x \in M$（相对于固定的良序）。由于 $S$
在 $R$ 上平坦，$(M/M_{\alpha}) \otimes_R S = M
\otimes_R S/M_{\alpha} \otimes_R S$；又因为
$M_{\alpha} \otimes_R S$ 是某些 $Q_i$ 的直和，
$(M/M_{\alpha}) \otimes_R S$ 也如此。由引理
\ref{algebra-lemma-adapted-submodule}，可以在 $M/M_{\alpha}$ 中找到
可数生成 $R$-子模 $P$，它包含 $x$ 在 $M/M_{\alpha}$ 中的像，
并且 $P \otimes_R S$（因为 $S$ 在 $R$ 上平坦，它等于
% P01-E0471: the frozen source maps P tensor S into M tensor S although P is a submodule of M/M_alpha; see ERRATA_CANDIDATES.jsonl.
$\Im(P
\otimes_R S \to M \otimes_R S)$）是某些 $Q_i$ 的直和。由于
$M \otimes_R S = \bigoplus_{i \in I} Q_i$ 投射，且投射性传递到
直和因子，$P \otimes_R S$ 也投射。因此，由引理
\ref{algebra-lemma-ffdescent-countable-projectivity}，$P$ 投射。最后，把
$M_{\alpha + 1}$ 定义为 $P$ 在 $M$ 中的逆像，于是
$M_{\alpha + 1}/M_{\alpha} = P$ 可数生成且投射。特别地，
$M_{\alpha}$ 是 $M_{\alpha + 1}$ 的直和因子，因为
$M_{\alpha + 1}/M_{\alpha}$ 的投射性意味着正合列
$0 \to M_{\alpha} \to M_{\alpha + 1} \to
M_{\alpha + 1}/M_{\alpha} \to 0$ 分裂。

\medskip\noindent
对 $M$ 作超限归纳（利用如下事实：在构造 $M_{\alpha + 1}$ 时，
我们使其包含不属于 $M_{\alpha}$ 的最小元素 $x \in M$），
可知每个 $x \in M$ 都属于某个 $M_{\alpha}$。因此，存在足够大的序数
% P01-E0473: the frozen source reuses S, already the faithfully flat target ring, as an ordinal; see ERRATA_CANDIDATES.jsonl.
$S$，满足：对每个 $x \in M$，存在 $\alpha \in S$，使得
$x \in M_{\alpha}$。这意味着 $(M_{\alpha})_{\alpha \in S}$
满足 $M$ 的 Kaplansky d\'evissage 的性质 (1)。其余性质由构造即明。
我们得到
$M = \bigoplus_{\alpha + 1 \in S} M_{\alpha + 1}/M_{\alpha}$。
由于按构造每个 $M_{\alpha + 1}/M_{\alpha}$ 都投射，$M$ 投射。
\end{proof}

\section{完备化}
\label{algebra-section-completion}

\noindent
设 $R$ 是环，$I$ 是理想。定义{\it $R$ 关于 $I$ 的完备化}为极限
$$
R^\wedge = \lim_n R/I^n.
$$
$R^\wedge$ 的一个元素由一列元素 $f_n \in R/I^n$ 给出，
它们对所有 $n$ 满足 $f_n \equiv f_{n + 1} \bmod I^n$。
我们把 $R^\wedge$ 看作 $R$-代数。类似地，若 $M$ 是 $R$-模，
则定义{\it $M$ 关于 $I$ 的完备化}为极限
$$
M^\wedge = \lim_n M/I^nM.
$$
$M^\wedge$ 的一个元素由一列元素 $m_n \in M/I^nM$ 给出，
它们对所有 $n$ 满足 $m_n \equiv m_{n + 1} \bmod I^nM$。
我们把 $M^\wedge$ 看作 $R^\wedge$-模。由这一描述可立即看出，
总有典范映射
$$
M \longrightarrow M^\wedge
\quad\text{且}\quad
M \otimes_R R^\wedge \longrightarrow M^\wedge.
$$
此外，给定模映射 $\varphi : M \to N$，会在完备化上诱导映射
$\varphi^\wedge : M^\wedge \to N^\wedge$，使图表
$$
\xymatrix{
M \ar[r] \ar[d] & N \ar[d] \\
M^\wedge \ar[r] & N^\wedge
}
$$
交换。一般而言，完备化不是正合函子，见
《例》，第 \ref{examples-section-completion-not-exact} 节。
下面给出一些初步的肯定结果。

\begin{lemma}
\label{algebra-lemma-completion-generalities}
设 $R$ 是环。设 $I \subset R$ 是理想。
设 $\varphi : M \to N$ 是 $R$-模映射。
\begin{enumerate}
\item 若 $M/IM \to N/IN$ 满射，则 $M^\wedge \to N^\wedge$ 满射。
\item 若 $M \to N$ 满射，则 $M^\wedge \to N^\wedge$ 满射。
\item 若 $0 \to K \to M \to N \to 0$ 是 $R$-模的短正合列且
$N$ 平坦，则
$0 \to K^\wedge \to M^\wedge \to N^\wedge \to 0$ 是短正合列。
\item 对任意有限 $R$-模 $M$，映射
$M \otimes_R R^\wedge \to M^\wedge$ 都满射。
\end{enumerate}
\end{lemma}

\begin{proof}
假设 $M/IM \to N/IN$ 满射。由 Nakayama 引理，对每个
$n \geq 1$，映射 $M/I^nM \to N/I^nN$ 都满射。更确切地说，
在环 $R/I^n$ 及其幂零理想 $I/I^n$ 上，对映射
$M/I^nM \to N/I^nN$ 应用引理 \ref{algebra-lemma-NAK} 的 (11)
即可看出这一点。置
$K_n = \{x \in M \mid \varphi(x) \in I^nN\}$。于是得到短正合列
$$
0 \to K_n/I^nM \to M/I^nM \to N/I^nN \to 0
$$
我们断言典范映射 $K_{n + 1}/I^{n + 1}M \to K_n/I^nM$
是满射。事实上，若 $x \in K_n$，写成
$\varphi(x) = \sum z_j n_j$，其中 $z_j \in I^n$、$n_j \in N$。
由假设，可写成 $n_j = \varphi(m_j) + \sum z_{jk}n_{jk}$，
其中 $m_j \in M$、$z_{jk} \in I$ 且 $n_{jk} \in N$。因此
$$
\varphi(x - \sum z_j m_j) = \sum z_jz_{jk} n_{jk}.
$$
这意味着 $x' = x - \sum z_j m_j \in K_{n + 1}$ 映到
$x \bmod I^nM$，从而证明了断言。现在可对上述短正合列逆系应用
引理 \ref{algebra-lemma-Mittag-Leffler}，得出 (1)。
(2) 是 (1) 的特殊情形。
若 (3) 的假设成立，则对每个 $n$，序列
$$
0 \to K/I^nK \to M/I^nM \to N/I^nN \to 0
$$
由引理 \ref{algebra-lemma-flat-tor-zero} 是短正合的。
因此可以直接应用引理 \ref{algebra-lemma-Mittag-Leffler}，得出 (3)。
为证 (4)，选取生成元 $x_i \in M$，$i = 1, \ldots, n$。
于是映射 $R^{\oplus n} \to M$，
$(a_1, \ldots, a_n) \mapsto \sum a_ix_i$ 满射。因此由 (2)，
映射 $(R^\wedge)^{\oplus n} \to M^\wedge$，
$(a_1, \ldots, a_n) \mapsto \sum a_ix_i$ 满射。(4) 由此得出。
\end{proof}

\begin{definition}
\label{algebra-definition-complete}
设 $R$ 是环。设 $I \subset R$ 是理想。
设 $M$ 是 $R$-模。若映射
$$
M \longrightarrow M^\wedge = \lim_n M/I^nM
$$
是同构\footnote{这包括条件 $\bigcap I^nM = 0$。}，则称 $M$
是{\it $I$-进完备的}。若 $R$ 作为 $R$-模是 $I$-进完备的，
则称 $R$ 是{\it $I$-进完备的}。
\end{definition}

\noindent
$R$-模 $M$ 关于 $I$ 的完备化未必是 $I$-进完备的。例子见
《例》，第 \ref{examples-section-noncomplete-completion} 节。
若理想有限生成，则完备化是完备的。

\begin{lemma}
\label{algebra-lemma-hathat-finitely-generated}
\begin{reference}
\cite[定理 15]{Matlis}。这里给出的简洁证明来自 Bjorn Poonen
于 2016 年 11 月 5 日发送的一封电子邮件。
\end{reference}
设 $R$ 是环。设 $I$ 是 $R$ 的有限生成理想。
设 $M$ 是 $R$-模。则
\begin{enumerate}
\item 完备化 $M^\wedge$ 是 $I$-进完备的，并且
\item 对所有 $n \geq 1$，
$I^nM^\wedge = \Ker(M^\wedge \to M/I^nM) = (I^nM)^\wedge$。
\end{enumerate}
特别地，$R^\wedge$ 是 $I$-进完备的，
$I^nR^\wedge = (I^n)^\wedge$，并且
$R^\wedge/I^nR^\wedge = R/I^n$。
\end{lemma}

\begin{proof}
由于 $I$ 有限生成，$I^n$ 也有限生成；设其生成元为
$f_1, \ldots, f_r$。把引理 \ref{algebra-lemma-completion-generalities} 的
(2) 应用于满射 $(f_1, \ldots, f_r) : M^{\oplus r} \to I^n M$，
得到满射
$$
(M^\wedge)^{\oplus r} \xrightarrow{(f_1, \ldots, f_r)} (I^n M)^\wedge =
\lim_{m \geq n} I^n M/I^m M = \Ker(M^\wedge \to M/I^n M).
$$
另一方面，映射
$(f_1, \ldots, f_r) : (M^\wedge)^{\oplus r} \to M^\wedge$
的像是 $I^n M^\wedge$。
因此 $M^\wedge / I^n M^\wedge \simeq M/I^n M$。
取逆极限得到 $(M^\wedge)^\wedge \simeq M^\wedge$；
也就是说，$M^\wedge$ 是 $I$-进完备的。
\end{proof}

\begin{lemma}
\label{algebra-lemma-completion-differ-by-torsion}
设 $R$ 是环。设 $I \subset R$ 是理想。设
$0 \to M \to N \to Q \to 0$ 是 $R$-模的正合列，
并且 $Q$ 被 $I$ 的某个幂零化。则完备化产生正合列
$0 \to M^\wedge \to N^\wedge \to Q \to 0$。
\end{lemma}

\begin{proof}
设 $I^c Q = 0$。则当 $n \geq c$ 时，$Q/I^nQ = Q$。
另一方面，当 $n \geq c$ 时，显然
$I^nM \subset M \cap I^nN \subset I^{n - c}M$。
因此 $M^\wedge = \lim M/(M \cap I^n N)$。对 $n \geq c$ 时的
正合列逆系
$$
0 \to M/(M \cap I^n N) \to N/I^n N \to Q \to 0
$$
应用引理 \ref{algebra-lemma-Mittag-Leffler} 即得结论。
\end{proof}

\begin{lemma}
\label{algebra-lemma-hathat}
\begin{reference}
取自 Lenstra 和 de Smit 的一份未发表笔记。
\end{reference}
设 $R$ 是环。设 $I \subset R$ 是理想。设 $M$ 是 $R$-模。
记 $K_n = \Ker(M^\wedge \to M/I^nM)$。则 $M^\wedge$
是 $I$-进完备的，当且仅当对所有 $n \geq 1$，
$K_n$ 等于 $I^nM^\wedge$。
\end{lemma}

\begin{proof}
模 $I^n M^\wedge$ 包含在 $K_n$ 中。
因此，对每个 $n \geq 1$，都有典范正合列
$$
0 \to K_n/I^nM^\wedge \to M^\wedge/I^nM^\wedge \to M/I^nM \to 0.
$$
由于 $I^nM^\wedge$ 满射到 $I^nM/I^{n + 1}M$，可知
$K_{n + 1} + I^n M^\wedge = K_n$。因此逆系
$\{K_n/I^n M^\wedge\}_{n \geq 1}$ 的转移映射都是满射。
由引理 \ref{algebra-lemma-Mittag-Leffler}，得到短正合列
$$
0 \to
\lim_n K_n/I^n M^\wedge \to
(M^\wedge)^\wedge \to
M^\wedge \to 0
$$
因此，$M^\wedge$ 完备，当且仅当对所有 $n \geq 1$，
$K_n/I^n M^\wedge = 0$。
\end{proof}

\begin{lemma}
\label{algebra-lemma-radical-completion}
设 $R$ 是环，$I \subset R$ 是理想，并令
$R^\wedge = \lim R/I^n$。
\begin{enumerate}
\item $R^\wedge$ 中任何映到 $R/I$ 的单位的元素都是单位；
\item $1 + I$ 中任何元素都映到 $R^\wedge$ 的可逆元素；
\item $1 + IR^\wedge$ 中任何元素在 $R^\wedge$ 中都可逆；并且
\item 理想 $IR^\wedge$ 和 $\Ker(R^\wedge \to R/I)$
都包含在 $R^\wedge$ 的 Jacobson 根中。
\end{enumerate}
\end{lemma}

\begin{proof}
设 $x \in R^\wedge$ 映到 $R/I$ 中的单位 $x_1$。
由引理 \ref{algebra-lemma-locally-nilpotent-unit}，对每个 $n$，
$x$ 都映到 $R/I^n$ 中的单位 $x_n$。
因此 $y = (x_n^{-1}) \in \lim R/I^n = R^\wedge$ 是 $x$ 的逆元。
(2) 和 (3) 立即由 (1) 得出。
(4) 由 (1) 和引理 \ref{algebra-lemma-contained-in-radical} 得出。
\end{proof}

\begin{lemma}
\label{algebra-lemma-when-surjective-to-completion}
设 $A$ 是环。设 $I = (f_1, \ldots, f_r)$ 是有限生成理想。
% P01-E0474: the frozen lemma uses M throughout without stating that M is an A-module; see ERRATA_CANDIDATES.jsonl.
若对每个 $i$，$M \to \lim M/f_i^nM$ 都满射，则
$M \to \lim M/I^nM$ 满射。
\end{lemma}

\begin{proof}
注意有等式：
$$
\lim M/I^nM = \lim M/(f_1^n, \ldots, f_r^n)M
$$
这是因为：
$$
I^n \supset (f_1^n, \ldots, f_r^n) \supset I^{rn}
$$
于是 $\lim M/(f_1^n, \ldots, f_r^n)M$ 的元素 $\xi$
可以符号化地写成
$$
\xi = \sum\nolimits_{n \geq 0} \sum\nolimits_i f_i^n x_{n, i}
$$
其中 $x_{n, i} \in M$。若 $M \to \lim M/f_i^nM$ 满射，
则存在 $x_i \in M$，它在 $\lim M/f_i^nM$ 中映到
$\sum x_{n, i} f_i^n$。于是 $x = \sum x_i$ 在
$\lim M/I^nM$ 中映到 $\xi$。
\end{proof}

\begin{lemma}
\label{algebra-lemma-complete-by-sub}
设 $A$ 是环。设 $I \subset J \subset A$ 是理想。
% P01-E0475: the frozen lemma uses M without stating that M is an A-module; see ERRATA_CANDIDATES.jsonl.
若 $M$ 是 $J$-进完备的且 $I$ 有限生成，则 $M$ 是 $I$-进完备的。
\end{lemma}

\begin{proof}
假设 $M$ 是 $J$-进完备的且 $I$ 有限生成。
因为 $\bigcap J^nM = 0$，所以 $\bigcap I^nM = 0$。由引理
\ref{algebra-lemma-when-surjective-to-completion}，只须在 $I$
由单个元素生成时证明 $M \to \lim M/I^nM$ 满射。设 $I = (f)$。
令 $x_n \in M$ 满足 $x_{n + 1} - x_n \in f^nM$。我们必须证明，
存在 $x \in M$，使得对所有 $n$，都有 $x_n - x \in f^nM$。
由于 $x_{n + 1} - x_n \in J^nM$，且 $M$ 是 $J$-进完备的，
存在元素 $x \in M$，使 $x_n - x \in J^nM$。
用 $x_n - x$ 代替 $x_n$，可以假设 $x_n \in J^nM$。
为完成证明，我们将说明这蕴含 $x_n \in I^nM$。
事实上，写成 $x_n - x_{n + 1} = f^nz_n$。则
$$
x_n = f^n(z_n + fz_{n + 1} + f^2z_{n + 2} + \ldots)
$$
由于 $f^c \in J^c$，和
$z_n + fz_{n + 1} + f^2z_{n + 2} + \ldots$ 在 $M$ 中收敛。
和 $f^n(z_n + fz_{n + 1} + f^2z_{n + 2} + \ldots)$
在 $M$ 中收敛到 $x_n$，因为其部分和等于
$x_n - x_{n + c}$，而 $x_{n + c} \in J^{n + c}M$。
\end{proof}

\begin{lemma}
\label{algebra-lemma-change-ideal-completion}
设 $R$ 是环。
设 $I$、$J$ 是 $R$ 的理想。
假设存在整数 $c, d > 0$，使
$I^c \subset J$ 且 $J^d \subset I$。
则对任意 $R$-模，关于 $I$ 的完备化与关于 $J$ 的完备化一致。
特别地，$R$-模 $M$ 是 $I$-进完备的，当且仅当它是
$J$-进完备的。
\end{lemma}

\begin{proof}
考虑映射系
$M/I^nM \to M/J^{\lfloor n/c \rfloor}M$ 与映射系
$M/J^mM \to M/I^{\lfloor m/d \rfloor}M$，
得到两个完备化之间互为逆的映射。
\end{proof}

\begin{lemma}
\label{algebra-lemma-quotient-complete}
设 $R$ 是环。设 $I$ 是 $R$ 的理想。
设 $M$ 是 $I$-进完备 $R$-模，并设 $K \subset M$ 是 $R$-子模。
以下条件等价：
\begin{enumerate}
\item $K = \bigcap (K + I^nM)$；并且
\item $M/K$ 是 $I$-进完备的。
\end{enumerate}
\end{lemma}

\begin{proof}
置 $N = M/K$。由引理 \ref{algebra-lemma-completion-generalities}，
映射 $M = M^\wedge \to N^\wedge$ 满射。
因此 $N \to N^\wedge$ 满射。不难看出，映射
$N \to N^\wedge$ 的核是模 $\bigcap (K + I^nM) / K$。
\end{proof}

\begin{lemma}
\label{algebra-lemma-when-finite-module-complete-over-complete-ring}
设 $R$ 是环。设 $I$ 是 $R$ 的理想。
设 $M$ 是 $R$-模。
若 (a) $R$ 是 $I$-进完备的，(b) $M$ 是有限 $R$-模，
且 (c) $\bigcap I^nM = (0)$，则 $M$ 是 $I$-进完备的。
\end{lemma}

\begin{proof}
由引理 \ref{algebra-lemma-completion-generalities}，映射
$M = M \otimes_R R = M \otimes_R R^\wedge \to M^\wedge$
满射。该映射的核是 $\bigcap I^nM$，因而由假设为零。
所以 $M \cong M^\wedge$，且 $M$ 完备。
\end{proof}

\begin{lemma}
\label{algebra-lemma-finite-over-complete-ring}
设 $R$ 是环。设 $I \subset R$ 是理想。设 $M$ 是 $R$-模。
假设
\begin{enumerate}
\item $R$ 是 $I$-进完备的；
\item $\bigcap_{n \geq 1} I^nM = (0)$；并且
\item $M/IM$ 是有限 $R/I$-模。
\end{enumerate}
则 $M$ 是有限 $R$-模。
\end{lemma}

\begin{proof}
设 $x_1, \ldots, x_n \in M$ 是这样的元素：它们在 $M/IM$
中的像作为 $R/I$-模生成 $M/IM$。记 $M' \subset M$ 为由
$x_1, \ldots, x_n$ 生成的 $R$-子模。由引理
\ref{algebra-lemma-completion-generalities}，映射
$(M')^\wedge \to M^\wedge$ 满射。
由于 $\bigcap I^nM = 0$，特别有 $\bigcap I^nM' = (0)$。
因此由引理 \ref{algebra-lemma-when-finite-module-complete-over-complete-ring}，
$M'$ 完备，从而 $M' \to M^\wedge$ 满射。最后，
$M \to M^\wedge$ 的核为零，因为它等于 $\bigcap I^nM = (0)$。
所以 $M \cong M' \cong M^\wedge$ 有限生成。
\end{proof}

\section{Noether 环的完备化}
\label{algebra-section-completion-noetherian}

\noindent
本节讨论 Noether 环中关于理想的完备化。

\begin{lemma}
\label{algebra-lemma-completion-tensor}
设 $I$ 是 Noether 环 $R$ 的理想。
用 ${}^\wedge$ 表示关于 $I$ 的完备化。
\begin{enumerate}
\item 若 $K \to N$ 是有限 $R$-模之间的单射映射，
则完备化上的映射 $K^\wedge \to N^\wedge$ 是单射。
\item 若 $0 \to K \to N \to M \to 0$ 是有限 $R$-模的短正合列，
则 $0 \to K^\wedge \to N^\wedge \to M^\wedge \to 0$ 是短正合列。
\item 若 $M$ 是有限 $R$-模，则 $M^\wedge = M \otimes_R R^\wedge$。
\end{enumerate}
\end{lemma}

\begin{proof}
置 $M = N/K$，可知 (1) 由 (2) 得出。
设 $0 \to K \to N \to M \to 0$ 如 (2) 所述。
对每个 $n$，得到短正合列
$$
0 \to K/(I^nN \cap K) \to N/I^nN \to M/I^nM \to 0.
$$
由引理 \ref{algebra-lemma-Mittag-Leffler}，得到正合列
$$
0 \to \lim K/(I^nN \cap K) \to N^\wedge \to M^\wedge \to 0.
$$
由 Artin--Rees 引理 \ref{algebra-lemma-Artin-Rees}，可以选取 $c$，
使得当 $n \geq c$ 时，
$I^nK \subset I^n N \cap K \subset I^{n-c} K$。
因此 $K^\wedge = \lim K/I^nK = \lim K/(I^nN \cap K)$，
从而 (2) 成立。

\medskip\noindent
设 $M$ 如 (3) 所述，并设
$0 \to K \to R^{\oplus t} \to M \to 0$ 是 $M$ 的一个表示。
得到交换图表
$$
\xymatrix{
&
K \otimes_R R^\wedge \ar[r] \ar[d] &
R^{\oplus t} \otimes_R R^\wedge \ar[r] \ar[d] &
M \otimes_R R^\wedge \ar[r] \ar[d] &
0 \\
0 \ar[r] &
K^\wedge \ar[r] &
(R^{\oplus t})^\wedge \ar[r] &
M^\wedge \ar[r] & 0
}
$$
上行正合，见第 \ref{algebra-section-flat} 节。
下行由 (2) 正合。由引理
\ref{algebra-lemma-completion-generalities}，竖直箭头都满射。
中间的竖直箭头是同构。由蛇引理 \ref{algebra-lemma-snake} 得出 (3)。
\end{proof}

\begin{lemma}
\label{algebra-lemma-completion-flat}
设 $I$ 是 Noether 环 $R$ 的理想。
用 ${}^\wedge$ 表示关于 $I$ 的完备化。
\begin{enumerate}
\item 环映射 $R \to R^\wedge$ 平坦。
\item 函子 $M \mapsto M^\wedge$ 在有限生成 $R$-模范畴上正合。
\end{enumerate}
\end{lemma}

\begin{proof}
考虑 $J \otimes_R R^\wedge \to R \otimes_R R^\wedge = R^\wedge$，
其中 $J$ 是 $R$ 的任意理想。
依照引理 \ref{algebra-lemma-completion-tensor}，该映射等同于
$J^\wedge \to R^\wedge$，而 $J^\wedge \to R^\wedge$ 是单射。
(1) 由引理 \ref{algebra-lemma-flat} 得出。
(2) 是引理 \ref{algebra-lemma-completion-tensor} 的 (2) 的重新表述。
\end{proof}

\begin{lemma}
\label{algebra-lemma-completion-faithfully-flat}
设 $I$ 是 Noether 环 $R$ 的理想。用 $R^\wedge$ 表示
$R$ 关于 $I$ 的完备化。若 $I$ 包含在 $R$ 的 Jacobson 根中，
则环映射 $R \to R^\wedge$ 忠实平坦。特别地，若
$(R, \mathfrak m)$ 是 Noether 局部环，则完备化
$\lim_n R/\mathfrak m^n$ 忠实平坦。
\end{lemma}

\begin{proof}
由引理 \ref{algebra-lemma-completion-flat}，它平坦。
考虑复合 $R \to R^\wedge \to R/I$，其中后一个映射是投影映射
$R^\wedge \to R/I$；这表明 $R$ 的任意极大理想都属于
$\Spec(R^\wedge) \to \Spec(R)$ 的像。因此，由引理
\ref{algebra-lemma-ff}，该映射忠实平坦。
\end{proof}

\begin{lemma}
\label{algebra-lemma-completion-complete}
设 $R$ 是 Noether 环。
设 $I$ 是 $R$ 的理想。
设 $M$ 是 $R$-模。
则 $M$ 关于 $I$ 的完备化 $M^\wedge$ 是 $I$-进完备的，
$I^n M^\wedge = (I^nM)^\wedge$，并且
$M^\wedge/I^nM^\wedge = M/I^nM$。
\end{lemma}

\begin{proof}
这是引理 \ref{algebra-lemma-hathat-finitely-generated} 的特殊情形，
因为 $I$ 是有限生成理想。
\end{proof}

\begin{lemma}
\label{algebra-lemma-completion-Noetherian}
设 $I$ 是环 $R$ 的理想。假设
\begin{enumerate}
\item $R/I$ 是 Noether 环；
\item $I$ 有限生成。
\end{enumerate}
则 $R$ 关于 $I$ 的完备化 $R^\wedge$ 是关于
$IR^\wedge$ 完备的 Noether 环。
\end{lemma}

\begin{proof}
由引理 \ref{algebra-lemma-hathat-finitely-generated} 可知，$R^\wedge$
是 $I$-进完备的。因此它也是 $IR^\wedge$-进完备的。
由于 $R^\wedge/IR^\wedge = R/I$ 是 Noether 的，把 $R$
替换为 $R^\wedge$ 后，除假设 (1) 和 (2) 外，还可以假设
$R$ 是 $I$-进完备的。

\medskip\noindent
设 $f_1, \ldots, f_t$ 是 $I$ 的生成元。
则有环的满射
$R/I[T_1, \ldots, T_t] \to \bigoplus I^n/I^{n + 1}$，
它把 $T_i$ 映到元素 $\overline{f}_i \in I/I^2$。
因此 $\bigoplus I^n/I^{n + 1}$ 是 Noether 环。
设 $J \subset R$ 是理想。考虑理想
$$
\bigoplus J \cap I^n/J \cap I^{n + 1} \subset \bigoplus I^n/I^{n + 1}.
$$
设 $\overline{g}_1, \ldots, \overline{g}_m$ 是该理想的生成元。
可以选取 $\overline{g}_j$ 为次数 $d_j$ 的齐次元素，并可选取
$g_j \in J \cap I^{d_j}$，使其映到
$\overline{g}_j \in J \cap I^{d_j}/J \cap I^{d_j + 1}$。
我们断言 $g_1, \ldots, g_m$ 生成 $J$。

\medskip\noindent
设 $x \in J \cap I^n$。存在
$a_j \in I^{\max(0, n - d_j)}$，使
$x - \sum a_j g_j \in J \cap I^{n + 1}$。
其理由是，按 $g_1, \ldots, g_m$ 的选取，
$J \cap I^n/J \cap I^{n + 1}$ 等于
% P01-E0480: the frozen source writes I^(n-d_j) for every j without restricting to d_j <= n, leaving negative ideal powers undefined; see ERRATA_CANDIDATES.jsonl.
$\sum \overline{g}_j I^{n - d_j}/I^{n - d_j + 1}$。
因此，从 $x \in J$ 出发，可以找到一列向量
$(a_{1, n}, \ldots, a_{m, n})_{n \geq 0}$，其中
$a_{j, n} \in I^{\max(0, n - d_j)}$，使得
$$
x =
\sum\nolimits_{n = 0, \ldots, N}
\sum\nolimits_{j = 1, \ldots, m} a_{j, n} g_j \bmod I^{N + 1}
$$
置 $A_j = \sum_{n \geq 0} a_{j, n}$。由于 $R$ 完备，可知
$x = \sum A_j g_j$。所以 $J$ 有限生成，结论得证。
\end{proof}

\begin{lemma}
\label{algebra-lemma-completion-Noetherian-Noetherian}
设 $R$ 是 Noether 环。
设 $I$ 是 $R$ 的理想。
$R$ 关于 $I$ 的完备化 $R^\wedge$ 是 Noether 的。
\end{lemma}

\begin{proof}
这是引理 \ref{algebra-lemma-completion-Noetherian} 的推论。
也可以如下直接看出。选取 $I$ 的生成元
$f_1, \ldots, f_n$。考虑映射
$$
R[[x_1, \ldots, x_n]] \longrightarrow R^\wedge,
\quad
x_i \longmapsto f_i.
$$
这是定义良好的满射环映射（细节略）。
由于 $R[[x_1, \ldots, x_n]]$ 是 Noether 的（见引理
\ref{algebra-lemma-Noetherian-power-series}），结论得证。
\end{proof}

\noindent
设 $R \to S$ 是局部环 $(R, \mathfrak m)$ 与
$(S, \mathfrak n)$ 之间的局部同态。令 $S^\wedge$ 为
$S$ 关于 $\mathfrak n$ 的完备化。一般而言，$S^\wedge$
不是 $S$ 的 $\mathfrak m$-进完备化。若对某个 $t \geq 1$，
有 $\mathfrak n^t \subset \mathfrak mS$，则由引理
\ref{algebra-lemma-change-ideal-completion}，确有
$S^\wedge = \lim S/\mathfrak m^nS$。在某些情形下，这甚至蕴含
$S^\wedge$ 在 $R^\wedge$ 上有限。

\begin{lemma}
\label{algebra-lemma-finite-after-completion}
设 $R \to S$ 是局部环 $(R, \mathfrak m)$ 与
$(S, \mathfrak n)$ 之间的局部同态。分别令
$R^\wedge$、$S^\wedge$ 为 $R$、$S$ 关于
$\mathfrak m$、$\mathfrak n$ 的完备化。若
$\mathfrak m$ 和 $\mathfrak n$ 有限生成，且
$\dim_{\kappa(\mathfrak m)} S/\mathfrak mS < \infty$，则
\begin{enumerate}
\item $S^\wedge$ 等于 $S$ 的 $\mathfrak m$-进完备化；并且
\item $S^\wedge$ 是有限 $R^\wedge$-模。
\end{enumerate}
\end{lemma}

\begin{proof}
因为 $R \to S$ 是局部环映射，所以
$\mathfrak mS \subset \mathfrak n$。
假设 $\dim_{\kappa(\mathfrak m)} S/\mathfrak mS < \infty$
蕴含 $S/\mathfrak mS$ 是 Artin 环，见引理
\ref{algebra-lemma-finite-dimensional-algebra}。
因此它的维数为 $0$，见引理
\ref{algebra-lemma-Noetherian-dimension-0}，从而
$\mathfrak n = \sqrt{\mathfrak mS}$。
这与 $\mathfrak n$ 有限生成这一事实蕴含：对某个
$t \geq 1$，有 $\mathfrak n^t \subset \mathfrak mS$。
由引理 \ref{algebra-lemma-change-ideal-completion}，可把 $S^\wedge$
等同于 $S$ 的 $\mathfrak m$-进完备化。由于 $\mathfrak m$
有限生成，由引理 \ref{algebra-lemma-hathat-finitely-generated} 可知，
$S^\wedge$ 与 $R^\wedge$ 都是 $\mathfrak m$-进完备的。
此时可以把引理 \ref{algebra-lemma-finite-over-complete-ring}
应用于作为 $R^\wedge$-模的 $S^\wedge$，得出结论。
\end{proof}

\begin{lemma}
\label{algebra-lemma-completion-finite-extension}
设 $R$ 是 Noether 环，$R \to S$ 是有限环映射。
设 $\mathfrak p \subset R$ 是素理想，并设
$\mathfrak q_1, \ldots, \mathfrak q_m$ 是 $S$ 中位于
$\mathfrak p$ 上方的素理想
（引理 \ref{algebra-lemma-finite-finite-fibres}）。则
$$
R_\mathfrak p^\wedge \otimes_R S =
(S_\mathfrak p)^\wedge =
S_{\mathfrak q_1}^\wedge \times \ldots \times S_{\mathfrak q_m}^\wedge
$$
其中 $(S_\mathfrak p)^\wedge$ 是关于 $\mathfrak p$ 的完备化，
而局部环 $R_\mathfrak p$ 和 $S_{\mathfrak q_i}$
都关于各自的极大理想完备化。
\end{lemma}

\begin{proof}
可以把 $R$ 替换为局部化 $R_\mathfrak p$，并把 $S$ 替换为
$S_\mathfrak p = S \otimes_R R_\mathfrak p$。
因此可以假设 $R$ 是 Noether 局部环，并且
$\mathfrak p = \mathfrak m$ 是其极大理想。
由引理 \ref{algebra-lemma-finite-after-completion}，
$\mathfrak q_iS_{\mathfrak q_i}$-进完备化
$S_{\mathfrak q_i}^\wedge$ 等于 $\mathfrak m$-进完备化。
对每个 $n \geq 1$，$S/\mathfrak m^nS$ 的素理想与
$S$ 的极大理想 $\mathfrak q_1, \ldots, \mathfrak q_m$
一一对应（由 $S$ 在 $R$ 上的上升定理，见引理
\ref{algebra-lemma-integral-going-up}）。因此，由引理
\ref{algebra-lemma-artinian-finite-length}，
$$
S/\mathfrak m^nS = \prod S_{\mathfrak q_i}/\mathfrak m^nS_{\mathfrak q_i}
$$
（例如使用命题 \ref{algebra-proposition-dimension-zero-ring}
可知 $S/\mathfrak m^nS$ 是 Artin 的）。
所以 $S$ 的 $\mathfrak m$-进完备化 $S^\wedge$ 等于
$\prod S_{\mathfrak q_i}^\wedge$。最后，由引理
\ref{algebra-lemma-completion-tensor}，
$R^\wedge \otimes_R S = S^\wedge$。
\end{proof}

\begin{lemma}
\label{algebra-lemma-split-completed-sequence}
设 $R$ 是环，$I \subset R$ 是理想。
设 $0 \to K \to P \to M \to 0$ 是 $R$-模的短正合列。
若 $M$ 在 $R$ 上平坦，且 $M/IM$ 是投射 $R/I$-模，
则 $I$-进完备化序列
$$
0 \to K^\wedge \to P^\wedge \to M^\wedge \to 0
$$
是分裂正合列。
\end{lemma}

\begin{proof}
由于 $M$ 平坦，每个序列
$$
0 \to K/I^nK \to P/I^nP \to M/I^nM \to 0
$$
都是短正合的，见引理 \ref{algebra-lemma-flat-tor-zero}；
而序列 $0 \to K^\wedge \to P^\wedge \to M^\wedge \to 0$
是短正合的，见引理 \ref{algebra-lemma-completion-generalities}。
只须证明，可以找到分裂
$s_n : M/I^nM \to P/I^nP$，使
$s_{n + 1} \bmod I^n = s_n$。
我们对 $n$ 归纳构造这些 $s_n$。
任取一个分裂 $s_1$；它存在，因为 $M/IM$ 是投射 $R/I$-模。
假设对某个 $n > 0$ 已给定 $s_n$。置
$P_{n + 1} = \{x \in P \mid x \bmod I^nP \in \Im(s_n)\}$。
映射
$\pi : P_{n + 1}/I^{n + 1}P_{n + 1} \to M/I^{n + 1}M$
满射（细节略）。由引理 \ref{algebra-lemma-lift-projective}，
$M/I^{n + 1}M$ 作为 $R/I^{n + 1}$-模是投射的，因而可以选取
$\pi$ 的截面
$t : M/I^{n + 1}M \to P_{n + 1}/I^{n + 1}P_{n + 1}$。
令 $s_{n + 1}$ 等于 $t$ 与典范映射
$P_{n + 1}/I^{n + 1}P_{n + 1} \to P/I^{n + 1}P$
的复合即可。
\end{proof}

\begin{lemma}
\label{algebra-lemma-complete-modulo-nilpotent}
设 $A$ 是 Noether 环。设 $I, J \subset A$ 是理想。
若 $A$ 是 $I$-进完备的且 $A/I$ 是 $J$-进完备的，
则 $A$ 是 $J$-进完备的。
\end{lemma}

\begin{proof}
设 $B$ 是 $A$ 的 $(I + J)$-进完备化。由引理
\ref{algebra-lemma-completion-flat}，$B/IB$ 是 $A/I$ 的 $J$-进完备化，
故由假设同构于 $A/I$。此外，由引理
\ref{algebra-lemma-complete-by-sub}，$B$ 是 $I$-进完备的。因此，由引理
\ref{algebra-lemma-finite-over-complete-ring}，$B$ 是有限 $A$-模。
由 Nakayama 引理（引理 \ref{algebra-lemma-NAK}；这里使用引理
\ref{algebra-lemma-radical-completion} 所给出的 $I$ 包含在 $A$ 的
Jacobson 根中这一事实），可知 $A \to B$ 满射。
由引理 \ref{algebra-lemma-completion-flat}，映射 $A \to B$ 平坦。
$\Spec(B) \to \Spec(A)$ 的像包含 $V(I)$；又由于 $I$
包含在 $A$ 的 Jacobson 根中，由引理
\ref{algebra-lemma-ff-rings} 可知 $A \to B$ 忠实平坦。因此
$A \to B$ 单射。于是 $A$ 关于 $I + J$ 完备，因而更是关于
$J$ 完备的。
\end{proof}

\section{模的极限}
\label{algebra-section-limits}

\noindent
本节讨论对模取极限时会发生什么。

\begin{lemma}
\label{algebra-lemma-limit-complete-pre}
设 $I \subset A$ 是一个环的有限生成理想。
设 $(M_n)$ 是满足 $I^n M_n = 0$ 的 $A$-模逆系统。
则 $M = \lim M_n$ 是 $I$-进完备的。
\end{lemma}

\begin{proof}
有 $M \to M/I^nM \to M_n$。取极限得到
$M \to M^\wedge \to M$。因此 $M$ 是 $M^\wedge$ 的直和因子。
由引理 \ref{algebra-lemma-hathat-finitely-generated}，$M^\wedge$
是 $I$-进完备的，故 $M$ 也是如此。
\end{proof}

\begin{lemma}
\label{algebra-lemma-limit-complete}
设 $I \subset A$ 是一个环的有限生成理想。
设 $(M_n)$ 是满足
$M_n = M_{n + 1}/I^nM_{n + 1}$ 的 $A$-模逆系统。置 $M = \lim M_n$。
则 $M/I^nM = M_n$，且 $M$ 是 $I$-进完备的。
\end{lemma}

\begin{proof}
由引理 \ref{algebra-lemma-limit-complete-pre} 可知 $M$ 是 $I$-进完备的。
由于转移映射都是满射，映射 $M \to M_n$ 也是满射。
考虑定义 $N_n$ 的短正合列逆系统
$$
0 \to N_n \to M \to M_n \to 0
$$
由于 $M_n = M_{n + 1}/I^nM_{n + 1}$，映射
$N_{n + 1} + I^nM \to N_n$ 是满射。因此
$N_{n + 1}/(N_{n + 1} \cap I^{n + 1}M) \to N_n/(N_n \cap I^nM)$
是满射。对短正合列
$$
0 \to N_n/(N_n \cap I^nM) \to M/I^nM \to M_n \to 0
$$
取逆极限，得到正合列
$$
0 \to \lim N_n/(N_n \cap I^nM) \to M^\wedge \to M
$$
由于 $M$ 是 $I$-进完备的，可知
$\lim N_n/(N_n \cap I^nM) = 0$；再由转移映射的满射性，
对所有 $n$ 都有 $N_n/(N_n \cap I^nM) = 0$。
故 $M_n = M/I^nM$，如所求。
\end{proof}

\begin{lemma}
\label{algebra-lemma-finiteness-graded}
设 $A$ 是 Noether 分次环。设 $I \subset A_+$ 是齐次理想。
设 $(N_n)$ 是有限分次 $A$-模的逆系统，并满足
$N_n = N_{n + 1}/I^n N_{n + 1}$。则存在有限分次 $A$-模
$N$，使得对所有 $n$，作为分次模都有 $N_n = N/I^nN$。
\end{lemma}

\begin{proof}
选取 $r$ 以及次数分别为 $d_1, \ldots, d_r$、
生成 $N_1$ 的齐次元素 $x_{1, 1}, \ldots, x_{1, r} \in N_1$。
由于转移映射都是满射，可以选取提升 $x_{1, i}$ 的相容齐次元素系
$x_{n, i} \in N_n$。由分次 Nakayama 引理
（引理 \ref{algebra-lemma-graded-NAK}）可知，$N_n$ 由次数分别为
$d_1, \ldots, d_r$ 的元素 $x_{n, 1}, \ldots, x_{n, r}$ 生成。
因此，当 $m \leq n$ 时，映射 $N_n \to N_n/I^m N_n$
在次数 $< \min(d_i) + m$ 上是同构（因为 $I^mN_n$
在这些次数上为零）。从而次数为 $d$ 的各部分所成的逆系统
% P01-E0483: frozen source stabilization indices are retained; see part-local ledger.
$$
\ldots
= N_{2 + d - \min(d_i), d}
= N_{1 + d - \min(d_i), d}
= N_{d - \min(d_i), d} \to N_{-1 + d - \min(d_i), d} \to \ldots
$$
按所示方式稳定。令 $N$ 为如下分次 $A$-模：其 $d$ 次齐次部分
就是这个稳定值。特别地，在 $N$ 中有元素
$x_i = \lim x_{n, i}$。我们断言诸 $x_i$ 生成 $N$：
任意 $x \in N_d$ 都是 $x_1, \ldots, x_r$ 的线性组合，
因为可以在 $N_{d - \min(d_i), d}$ 中检验这一点；在那里结论成立，
这是由于诸 $x_{d - \min(d_i), i}$ 生成
$N_{d - \min(d_i)}$。最后，读者可逐次检查，验证满射
$N/I^nN \to N_n$ 是同构。细节从略。
\end{proof}

\begin{lemma}
\label{algebra-lemma-daniel-litt}
设 $A$ 是分次环。设 $I \subset A_+$ 是齐次理想。
记 $A' = \lim A/I^n$。设 $(G_n)$ 是分次 $A$-模的逆系统，
其中 $G_n$ 被 $I^n$ 零化。设 $M$ 是分次 $A$-模，并设
$\varphi_n : M \to G_n$ 是相容的分次 $A$-模映射系。若诱导映射
$$
\varphi : M \otimes_A A' \longrightarrow \lim G_n
$$
是同构，则对每个 $d \in \mathbf{Z}$，映射
$M_d \to \lim G_{n, d}$ 都是同构。
\end{lemma}

\begin{proof}
按约定，分次环集中在次数 $\geq 0$ 上，而分次模可以在任意次数上
具有非零部分，见第 \ref{algebra-section-graded} 节。映射 $\varphi$ 存在，
因为 $G_n$ 被 $I^n$ 零化，所以 $\lim G_n$ 是 $A'$-模。
另一个值得记住的事实是
$$
\bigoplus\nolimits_{d \in \mathbf{Z}} \lim G_{n, d} \subset
\lim G_n \subset
\prod\nolimits_{d \in \mathbf{Z}} \lim G_{n, d}
$$
其中下标 ${\ }_d$ 表示 $d$ 次齐次部分。

\medskip\noindent
单射性。设 $x \in M_d$。若在 $\lim G_{n, d}$ 中有 $x \mapsto 0$，
则在 $M \otimes_A A'$ 中有 $x \otimes 1 = 0$。于是可以找到
有限生成子模 $M' \subset M$，且 $x \in M'$，使得
$x \otimes 1$ 在 $M' \otimes_A A'$ 中为零。
设 $M'$ 由次数为 $d_1, \ldots, d_r$ 的齐次元素生成。
令 $n = d - \min(d_i) + 1$。由于 $A'$ 有到 $A/I^n$ 的映射，
且 $A \to A/I^n$ 在次数 $\leq n - 1$ 上是同构，故可知
% P01-E0484: frozen source degree bound is retained; see part-local ledger.
$M' \to M' \otimes_A A'$ 在次数 $\leq n - 1$ 上是单射。
因此 $x = 0$，如所求。

\medskip\noindent
满射性。设 $y \in \lim G_{n, d}$。在 $M \otimes_A A'$ 中选取
一个映到 $y$ 的有限和 $\sum x_i \otimes f'_i$。
可以假设 $x_i$ 是齐次的，设其次数为 $d_i$。
注意，虽然 $A'$ 不是分次环，但它是分次环 $A/I^nA$ 的极限；
而且在任一给定次数上，转移映射最终成为同构（见上文）。这给出
$$
A = \bigoplus\nolimits_{d \geq 0} A_d \subset A' \subset
\prod\nolimits_{d \geq 0} A_d
$$
因此可以写成
$$
f'_i = \sum\nolimits_{j = 0, \ldots, d - d_i - 1} f_{i, j} + f_i + g'_i
$$
其中 $f_{i, j} \in A_j$、$f_i \in A_{d - d_i}$，且
$g'_i \in A'$ 在 $\prod_{j \leq d - d_i} A_j$ 中的像为零。
现在计算 $\varphi_n(\sum_{i, j} f_{i, j}x_i) \in G_n$，
得到的是次数 $< d$ 的齐次元素之和。因此
$\varphi(\sum x_i \otimes f_{i, j})$ 在 $\lim G_{n, d}$ 中的像为零。
类似地，计算表明元素 $\varphi(\sum x_i \otimes g'_i)$
在 $\prod_{d' \leq d} \lim G_{n, d'}$ 中的像为零。
由于已知 $\varphi(\sum x_i \otimes f'_i)$ 等于 $y$，
可知 $\sum f_ix_i \in M_d$ 映到 $y$，如所求。
\end{proof}

\section{平坦性判据}
\label{algebra-section-criteria-flatness}

\noindent
本节证明 Noether 情形下的一些重要技术引理。我们将在第
\ref{algebra-section-more-flatness-criteria} 节把它们（部分地）推广到
非 Noether 情形。

\begin{lemma}
\label{algebra-lemma-mod-injective}
设 $R \to S$ 是局部环的局部同态，且 $S$ 是 Noether 的。
% P01-E0485: frozen source omits "by" in three notation declarations; see ledger.
记 $\mathfrak m$ 为 $R$ 的极大理想。设 $M$ 是平坦 $R$-模，
$N$ 是有限 $S$-模。设 $u : N \to M$ 是 $R$-模映射。
若 $\overline{u} : N/\mathfrak m N \to M/\mathfrak m M$ 是单射，
则 $u$ 是单射。在这种情形下，$M/u(N)$ 在 $R$ 上平坦。
\end{lemma}

\begin{proof}
首先断言，对所有 $n \geq 1$，
$u_n : N/{\mathfrak m}^nN \to M/{\mathfrak m}^nM$ 都是单射。
我们作归纳；基础情形即 $\overline{u} = u_1$ 是单射。
由 $M$ 在 $R$ 上平坦这一假设，有短正合列
$0 \to M \otimes_R {\mathfrak m}^n/{\mathfrak m}^{n + 1}
\to M/{\mathfrak m}^{n + 1}M \to M/{\mathfrak m}^n M \to 0$。
此外，
$M \otimes_R {\mathfrak m}^n/{\mathfrak m}^{n + 1}
= M/{\mathfrak m}M \otimes_{R/{\mathfrak m}}
{\mathfrak m}^n/{\mathfrak m}^{n + 1}$。对 $N$ 也有类似的正合列
$N \otimes_R {\mathfrak m}^n/{\mathfrak m}^{n + 1}
\to N/{\mathfrak m}^{n + 1}N \to N/{\mathfrak m}^n N \to 0$，
只是左端没有零。还有
$N \otimes_R {\mathfrak m}^n/{\mathfrak m}^{n + 1}
= N/{\mathfrak m}N \otimes_{R/{\mathfrak m}}
{\mathfrak m}^n/{\mathfrak m}^{n + 1}$。因此，
$u_n$ 与映射
$\overline{u} \otimes \text{id}_{{\mathfrak m}^n/{\mathfrak m}^{n + 1}}$
都是单射，从而 $u_{n + 1}$ 也是单射。

\medskip\noindent
把 Krull 交定理（引理 \ref{algebra-lemma-intersect-powers-ideal-module-zero}）
应用于环 $S$ 上的 $N$ 及理想 $\mathfrak mS$，得到
$\bigcap \mathfrak m^nN = 0$。因此，对所有 $n$，$u_n$ 的单射性
蕴含 $u$ 是单射。

\medskip\noindent
为证明 $M/u(N)$ 在 $R$ 上平坦，只须证明对每个理想
$I \subset R$ 都有 $\text{Tor}_1^R(M/u(N), R/I) = 0$，
见引理 \ref{algebra-lemma-characterize-flat}。由短正合列
$$
0 \to N \xrightarrow{u} M \to M/u(N) \to 0
$$
以及 $M$ 的平坦性，得到 Tor 的正合列
$$
0 \to \text{Tor}_1^R(M/u(N), R/I) \to N/IN \to M/IM
$$
见引理 \ref{algebra-lemma-long-exact-sequence-tor}。因此，只须证明
$N/IN$ 单射到 $M/IM$ 中。注意，$R/I \to S/IS$ 是局部环的
局部同态，且 $S/IS$ 是 Noether 的；$N/IN \to M/IM$ 是
$R/I$-模映射；$N/IN$ 在 $S/IS$ 上有限；$M/IM$ 在 $R/I$ 上
平坦；而 $u \bmod I : N/IN \to M/IM$ 模 $\mathfrak m$ 是单射。
于是可以把证明的第一部分应用于 $u \bmod I$，结论即得。
\end{proof}

\begin{lemma}
\label{algebra-lemma-grothendieck}
设 $R \to S$ 是 Noether 局部环之间的平坦局部环同态。
记 $\mathfrak m$ 为 $R$ 的极大理想。设 $f \in S$ 在
$S/{\mathfrak m}S$ 中是非零因子。则 $S/fS$ 在 $R$ 上平坦，
且 $f$ 在 $S$ 中是非零因子。
\end{lemma}

\begin{proof}
直接由引理 \ref{algebra-lemma-mod-injective} 得出。
\end{proof}

\begin{lemma}
\label{algebra-lemma-grothendieck-regular-sequence}
设 $R \to S$ 是 Noether 局部环之间的平坦局部环同态。
记 $\mathfrak m$ 为 $R$ 的极大理想。设 $f_1, \ldots, f_c$
是 $S$ 中的元素序列，使其像
$\overline{f}_1, \ldots, \overline{f}_c$ 在
$S/{\mathfrak m}S$ 中构成正则序列。则 $f_1, \ldots, f_c$
在 $S$ 中构成正则序列，且每个商 $S/(f_1, \ldots, f_i)$
都在 $R$ 上平坦。
\end{lemma}

\begin{proof}
由归纳法和引理 \ref{algebra-lemma-grothendieck} 得出。
\end{proof}

\begin{lemma}
\label{algebra-lemma-free-fibre-flat-free}
设 $R \to S$ 是 Noether 局部环之间的局部同态。
设 $\mathfrak m$ 是 $R$ 的极大理想。设 $M$ 是非零有限 $S$-模。
假设 (a) $M/\mathfrak mM$ 是自由 $S/\mathfrak mS$-模，且
(b) $M$ 在 $R$ 上平坦。则 $M$ 是自由的，且 $S$ 在 $R$ 上平坦。
\end{lemma}

\begin{proof}
设 $\overline{x}_1, \ldots, \overline{x}_n$ 是自由模
$M/\mathfrak mM$ 的一组基。选取 $x_1, \ldots, x_n \in M$，
使 $x_i$ 映到 $\overline{x}_i$。令
$u : S^{\oplus n} \to M$ 为把第 $i$ 个标准基向量映到 $x_i$ 的映射。
由引理 \ref{algebra-lemma-mod-injective}，$u$ 是单射。另一方面，
由 Nakayama 引理 \ref{algebra-lemma-NAK}，该映射是满射。结论随即得出。
\end{proof}

\begin{lemma}
\label{algebra-lemma-complex-exact-mod}
设 $R \to S$ 是局部 Noether 环之间的局部同态。设
$\mathfrak m$ 是 $R$ 的极大理想。设
$0 \to F_e \to F_{e-1} \to \ldots \to F_0$
是有限 $S$-模的有限复形。假设每个 $F_i$ 都在 $R$ 上平坦，
且复形
$0 \to F_e/\mathfrak m F_e \to F_{e-1}/\mathfrak m F_{e-1}
\to \ldots \to F_0 / \mathfrak m F_0$ 正合。则
$0 \to F_e \to F_{e-1} \to \ldots \to F_0$ 正合，
而且模 $\Coker(F_1 \to F_0)$ 在 $R$ 上平坦。
\end{lemma}

\begin{proof}
对 $e$ 作归纳。若 $e = 1$，这正是引理
\ref{algebra-lemma-mod-injective}。若 $e > 1$，则由引理
\ref{algebra-lemma-mod-injective} 可知 $F_e \to F_{e-1}$ 是单射，且
$C = \Coker(F_e \to F_{e-1})$ 是在 $R$ 上平坦的有限 $S$-模。
因此可以把归纳假设应用于复形
$0 \to C \to F_{e-2} \to \ldots \to F_0$。
由此得出 $C \to F_{e-2}$ 是单射，进而得到该复形的正合性，
以及 $F_1 \to F_0$ 的余核的平坦性。
\end{proof}

\noindent
本节余下部分证明所谓“{\it 平坦性的局部判据}”的两个版本。
还请注意下文颇有意思的引理 \ref{algebra-lemma-CM-over-regular-flat}。

\begin{lemma}
\label{algebra-lemma-prepare-local-criterion-flatness}
设 $R$ 是极大理想为 $\mathfrak m$、剩余域为
$\kappa = R/\mathfrak m$ 的局部环。设 $M$ 是 $R$-模。
若 $\text{Tor}_1^R(\kappa, M) = 0$，则对每个有限长度
$R$-模 $N$，都有 $\text{Tor}_1^R(N, M) = 0$。
\end{lemma}

\begin{proof}
% P01-E0486: frozen source says descending induction although the proof is ordinary strong induction.
对 $N$ 的长度作归纳。若 $N$ 的长度为 $1$，则
$N \cong \kappa$，结论成立。若 $N$ 的长度大于 $1$，
则可将 $N$ 置于短正合列
$0 \to N' \to N \to N'' \to 0$ 中，其中 $N'$、$N''$
都是长度更小的有限长度 $R$-模。由
$\text{Tor}_1^R(N', M)$ 和 $\text{Tor}_1^R(N'', M)$ 的消没
（归纳假设），再结合 Tor 群的长正合列，可得
$\text{Tor}_1^R(N, M)$ 消没，见引理
\ref{algebra-lemma-long-exact-sequence-tor}。
\end{proof}

\begin{lemma}[平坦性的局部判据]
\label{algebra-lemma-local-criterion-flatness}
设 $R \to S$ 是局部 Noether 环之间的局部同态。
设 $\mathfrak m$ 是 $R$ 的极大理想，并设
$\kappa = R/\mathfrak m$。设 $M$ 是有限 $S$-模。
若 $\text{Tor}_1^R(\kappa, M) = 0$，则 $M$ 在 $R$ 上平坦。
\end{lemma}

\begin{proof}
设 $I \subset R$ 是理想。由引理 \ref{algebra-lemma-flat}，只须证明
$I \otimes_R M \to M$ 是单射。由注记
\ref{algebra-remark-Tor-ring-mod-ideal}，这个核等于
$\text{Tor}_1^R(M, R/I)$。由引理
\ref{algebra-lemma-prepare-local-criterion-flatness} 可知，
% P01-E0487: frozen source uses J without quantifying it; see ledger.
对所有有限余长度的理想，$J \otimes_R M \to M$ 都是单射。

\medskip\noindent
选取 $n >> 0$，并考虑短正合列
$$
0
\to I \cap \mathfrak m^n
\to I \oplus \mathfrak m^n
\to I + \mathfrak m^n
\to 0
$$
这是短正合列 $0 \to R \to R^{\oplus 2} \to R \to 0$ 的子序列。
因此得到交换图
$$
\xymatrix{
(I\cap \mathfrak m^n) \otimes_R M \ar[r] \ar[d] &
I \otimes_R M \oplus \mathfrak m^n \otimes_R M \ar[r] \ar[d] &
(I + \mathfrak m^n) \otimes_R M \ar[d] \\
M \ar[r] &
M \oplus M \ar[r] &
M
}
$$
注意，$I + \mathfrak m^n$ 和 $\mathfrak m^n$ 都是有限余长度的理想。
因此，追图表明
$\Ker((I \cap \mathfrak m^n)\otimes_R M \to M)
\to \Ker(I \otimes_R M \to M)$ 是满射。特别地，由此可知
$K = \Ker(I \otimes_R M \to M)$ 包含在
$(I \cap \mathfrak m^n) \otimes_R M$ 于 $I \otimes_R M$ 中的像内。
由 Artin--Rees 引理 \ref{algebra-lemma-Artin-Rees}，存在某个 $c > 0$，
使得对所有 $n >> 0$，$K$ 包含在
$\mathfrak m^{n-c}(I \otimes_R M)$ 中。由于
$I \otimes_R M$ 是有限 $S$-模（！），而 $S$ 是 Noether 的，
可知这蕴含 $K = 0$。具体而言，上述结论蕴含 $K$ 在
$I \otimes_R M$ 的 $\mathfrak mS$-进完备化中的像为零。
但由引理 \ref{algebra-lemma-completion-faithfully-flat}，从 $S$ 到其
$\mathfrak mS$-进完备化的映射是忠实平坦的。因此 $K = 0$，如所求。
\end{proof}

\noindent
下文经常遇到条件“$M/IM$ 在 $R/I$ 上平坦且
$\text{Tor}_1^R(R/I, M) = 0$”。下面的引理给出这些条件的一些推论
（它是引理 \ref{algebra-lemma-prepare-local-criterion-flatness} 的推广）。

\begin{lemma}
\label{algebra-lemma-what-does-it-mean}
设 $R$ 是环。设 $I \subset R$ 是理想。设 $M$ 是 $R$-模。
若 $M/IM$ 在 $R/I$ 上平坦，且
$\text{Tor}_1^R(R/I, M) = 0$，则
\begin{enumerate}
\item 对所有 $n \geq 1$，$M/I^nM$ 在 $R/I^n$ 上平坦；并且
\item 对任意被某个 $m \geq 0$ 的 $I^m$ 零化的模 $N$，都有
$\text{Tor}_1^R(N, M) = 0$。
\end{enumerate}
特别地，若 $I$ 是幂零理想，则 $M$ 在 $R$ 上平坦。
\end{lemma}

\begin{proof}
假设 $M/IM$ 在 $R/I$ 上平坦且 $\text{Tor}_1^R(R/I, M) = 0$。
设 $N$ 是 $R/I$-模。选取短正合列
% P01-E0488: frozen source incorrectly uses the ideal I as an arbitrary free-module index set.
$$
0 \to K \to \bigoplus\nolimits_{i \in I} R/I \to N \to 0
$$
由 $\text{Tor}$ 的长正合列和 $\text{Tor}_1^R(R/I, M)$ 的消没，得到
$$
0 \to \text{Tor}_1^R(N, M) \to K \otimes_R M \to
(\bigoplus\nolimits_{i \in I} R/I) \otimes_R M \to N \otimes_R M \to 0
$$
但由于 $K$、$\bigoplus_{i \in I} R/I$ 和 $N$ 都被 $I$ 零化，
可知
\begin{align*}
K \otimes_R M & = K \otimes_{R/I} M/IM, \\
(\bigoplus\nolimits_{i \in I} R/I) \otimes_R M & =
(\bigoplus\nolimits_{i \in I} R/I) \otimes_{R/I} M/IM, \\
N \otimes_R M & = N \otimes_{R/I} M/IM.
\end{align*}
因为 $M/IM$ 在 $R/I$ 上平坦，所以
% P01-E0489: frozen source has malformed tensor base R/; retained exactly.
$$
0 \to K \otimes_{R/I} M/IM \to
(\bigoplus\nolimits_{i \in I} R/I) \otimes_{R/I} M/IM \to
N \otimes_{R/} M/IM \to 0
$$
正合。结合上文，可知对任意被 $I$ 零化的 $R$-模 $N$，都有
$\text{Tor}_1^R(N, M) = 0$。

\medskip\noindent
下面对 $m$ 归纳证明 (2)。$m = 1$ 的情形已在上一段证明。
若 $N$ 被 $I^m$ 零化且 $m > 1$，可以选取正合列
$0 \to N' \to N \to N'' \to 0$，其中 $N'$ 和 $N''$
都被 $I^{m - 1}$ 零化。例如可以取 $N' = IN$ 和 $N'' = N/IN$。
此时，正合列
$$
\text{Tor}_1^R(N', M) \to
\text{Tor}_1^R(N, M) \to
\text{Tor}_1^R(N'', M)
$$
与归纳假设证明了所需的消没。

\medskip\noindent
最后证明 (1)。给定 $n \geq 1$，须证明 $M/I^nM$ 在 $R/I^n$ 上平坦。
换言之，须证明函子
$N \mapsto N \otimes_{R/I^n} M/I^nM$ 在由被 $I^n$ 零化的
$R$-模 $N$ 构成的范畴上正合。但对这样的 $N$，有
$N \otimes_{R/I^n} M/I^nM = N \otimes_R M$。
由 (2) 中 $\text{Tor}_1$ 的消没可知，函子
$N \mapsto N \otimes_R M$ 在被 $I$ 的某个幂零化的 $N$
所构成的范畴上正合，结论即得。
\end{proof}

\begin{lemma}
\label{algebra-lemma-what-does-it-mean-again}
设 $R$ 是环。设 $I \subset R$ 是理想。设 $M$ 是 $R$-模。
\begin{enumerate}
\item 若 $M/IM$ 在 $R/I$ 上平坦，且
$M \otimes_R I/I^2 \to IM/I^2M$ 是单射，则
$M/I^2M$ 在 $R/I^2$ 上平坦。
\item 若 $M/IM$ 在 $R/I$ 上平坦，且对
$n = 1, \ldots, k$，映射
$M \otimes_R I^n/I^{n + 1} \to I^nM/I^{n + 1}M$ 都是单射，
则 $M/I^{k + 1}M$ 在 $R/I^{k + 1}$ 上平坦。
\end{enumerate}
\end{lemma}

\begin{proof}
把引理 \ref{algebra-lemma-what-does-it-mean} 应用于以 $R/I^2$ 代替 $R$、
以 $M/I^2M$ 代替 $M$ 的情形，并使用
$$
\text{Tor}_1^{R/I^2}(M/I^2M, R/I) =
\Ker(M \otimes_R I/I^2 \to IM/I^2M),
$$
即可得到第一条，见注记 \ref{algebra-remark-Tor-ring-mod-ideal}。
第二条以同样方式得出：对 $n$ 作归纳，证明当
$n = 1, \ldots, k$ 时，$M/I^{n + 1}M$ 在 $R/I^{n + 1}$ 上平坦。
这里使用对每个 $n$ 都成立的等式
$$
\text{Tor}_1^{R/I^{n + 1}}(M/I^{n + 1}M, R/I^n) =
\Ker(M \otimes_R I^n/I^{n + 1} \to I^nM/I^{n + 1}M)
$$
\end{proof}

\begin{lemma}[局部判据的变式]
\label{algebra-lemma-variant-local-criterion-flatness}
设 $R \to S$ 是 Noether 局部环之间的局部同态。
设 $I \not = R$ 是 $R$ 中的理想。设 $M$ 是有限 $S$-模。
若 $\text{Tor}_1^R(M, R/I) = 0$，且 $M/IM$ 在 $R/I$ 上平坦，
则 $M$ 在 $R$ 上平坦。
\end{lemma}

\begin{proof}
第一种证明：由引理 \ref{algebra-lemma-what-does-it-mean}，
$\text{Tor}_1^R(\kappa, M)$ 为零，其中 $\kappa$ 是 $R$ 的剩余域。
于是由引理 \ref{algebra-lemma-local-criterion-flatness}，$M$ 在 $R$ 上平坦。

\medskip\noindent
第二种证明：设 $\mathfrak m$ 是 $R$ 的极大理想。我们将证明
$\mathfrak m \otimes_R M \to M$ 是单射，然后应用引理
\ref{algebra-lemma-local-criterion-flatness}。设
$\sum f_i \otimes x_i \in \mathfrak m \otimes_R M$，且在 $M$ 中
$\sum f_i x_i = 0$。把平坦性的等式判据——引理
\ref{algebra-lemma-flat-eq}——应用于 $R/I$ 上的 $M/IM$，可知存在
$\overline{a}_{ij} \in R/I$ 和 $\overline{y}_j \in M/IM$，使得
$x_i \bmod IM = \sum_j \overline{a}_{ij} \overline{y}_j $，且
$0 = \sum_i (f_i \bmod I) \overline{a}_{ij}$。
设 $a_{ij} \in R$ 是 $\overline{a}_{ij}$ 的一个提升；同样地，
设 $y_j \in M$ 是 $\overline{y}_j$ 的一个提升。于是
\begin{eqnarray*}
\sum f_i \otimes x_i
& = &
\sum f_i \otimes x_i +
\sum f_ia_{ij} \otimes y_j -
\sum f_i \otimes a_{ij} y_j
\\
& = &
\sum f_i \otimes (x_i - \sum a_{ij} y_j) +
\sum (\sum f_i a_{ij}) \otimes y_j
\end{eqnarray*}
由于 $x_i - \sum a_{ij} y_j \in IM$，且
$\sum f_i a_{ij} \in I$，可知存在 $I \otimes_R M$ 中的元素，
其在 $\mathfrak m \otimes_R M$ 中的像等于给定元素
$\sum f_i \otimes x_i$。但由假设，$I \otimes_R M \to M$ 是单射
（见注记 \ref{algebra-remark-Tor-ring-mod-ideal}），故结论成立。
\end{proof}

\noindent
特别地，在引理 \ref{algebra-lemma-variant-local-criterion-flatness} 的情形下，
假设 $I = (x)$ 由单个元素 $x$ 生成，且该元素在 $R$ 中是非零因子。
那么，$\text{Tor}_1^R(M, R/(x)) = (0)$ 当且仅当 $x$ 在 $M$ 上是非零因子。

\begin{lemma}
\label{algebra-lemma-flat-module-powers}
设 $R \to S$ 是环映射。设 $I \subset R$ 是理想。
设 $M$ 是 $S$-模。假设
\begin{enumerate}
\item $R$ 是 Noether 环；
\item $S$ 是 Noether 环；
\item $M$ 是有限 $S$-模；并且
\item 对每个 $n \geq 1$，模 $M/I^n M$ 在 $R/I^n$ 上平坦。
\end{enumerate}
则对每个 $\mathfrak q \in V(IS)$，局部化 $M_{\mathfrak q}$
在 $R$ 上平坦。特别地，若 $S$ 是局部环且 $IS$ 包含在其极大理想中，
则 $M$ 在 $R$ 上平坦。
\end{lemma}

\begin{proof}
我们将使用引理 \ref{algebra-lemma-variant-local-criterion-flatness}。
由假设，$M/IM$ 在 $R/I$ 上平坦。因此，只须检验
$\text{Tor}_1^R(M, R/I)$ 在 $\mathfrak q$ 处局部化后为零。
由注记 \ref{algebra-remark-Tor-ring-mod-ideal}，这个 Tor 群等于
$K = \Ker(I \otimes_R M \to M)$。已知对所有 $n \geq 1$，映射
$I/I^n \otimes_{R/I^n} M/I^nM \to M/I^nM$ 的核为零。
因此，$K$ 中任意元素在 $I/I^n \otimes_{R/I^n} M/I^nM$ 中的像为零。
由于
$$
I/I^n \otimes_{R/I^n} M/I^nM = I/I^n \otimes_R M =
(I \otimes_R M)/I^{n - 1}(I \otimes_R M)
$$
可知对所有 $n \geq 1$，都有
$K \subset I^{n - 1}(I \otimes_R M)$。由 Artin--Rees 引理，
更确切地说，由引理 \ref{algebra-lemma-intersection-powers-ideal-module}，
可知 $K_{\mathfrak q} = 0$，如所求。
\end{proof}

\begin{lemma}
\label{algebra-lemma-surjective-on-tor-one}
设 $R \to R' \to R''$ 是环映射。设 $M$ 是 $R$-模。
假设 $M \otimes_R R'$ 在 $R'$ 上平坦。则自然映射
$\text{Tor}_1^R(M, R') \otimes_{R'} R'' \to
\text{Tor}_1^R(M, R'')$ 是满射。
\end{lemma}

\begin{proof}
设 $F_\bullet$ 是 $M$ 在 $R$ 上的自由分解。复形
$F_2 \otimes_R R' \to F_1\otimes_R R' \to F_0 \otimes_R R'$
计算 $\text{Tor}_1^R(M, R')$，而复形
$F_2 \otimes_R R'' \to F_1\otimes_R R'' \to F_0 \otimes_R R''$
计算 $\text{Tor}_1^R(M, R'')$。注意，
$F_i \otimes_R R' \otimes_{R'} R'' = F_i \otimes_R R''$。
令 $K' = \Ker(F_1\otimes_R R' \to F_0 \otimes_R R')$，并类似地令
$K'' = \Ker(F_1\otimes_R R'' \to F_0 \otimes_R R'')$。
于是有正合列
$$
0 \to K' \to F_1\otimes_R R' \to F_0 \otimes_R R' \to M \otimes_R R' \to 0.
$$
由 $M \otimes_R R'$ 在 $R'$ 上平坦这一假设，序列
$$
K' \otimes_{R'} R'' \to
F_1 \otimes_R R'' \to
F_0 \otimes_R R'' \to
M \otimes_R R'' \to 0
$$
仍然正合。这意味着 $K' \otimes_{R'} R'' \to K''$ 是满射。
由于 $\text{Tor}_1^R(M, R')$ 是 $K'$ 的商，而
$\text{Tor}_1^R(M, R'')$ 是 $K''$ 的商，结论即得。
\end{proof}

\begin{lemma}
\label{algebra-lemma-surjective-on-tor-one-trivial}
设 $R \to R'$ 是环映射。设 $I \subset R$ 是理想，并令
$I' = IR'$。设 $M$ 是 $R$-模，并置 $M' = M \otimes_R R'$。
自然映射
$\text{Tor}_1^R(R'/I', M) \to \text{Tor}_1^{R'}(R'/I', M')$
是满射。
\end{lemma}

\begin{proof}
设 $F_2 \to F_1 \to F_0 \to M \to 0$ 是 $M$ 在 $R$ 上的自由分解。
置 $F_i' = F_i \otimes_R R'$。序列
$F_2' \to F_1' \to F_0' \to M' \to 0$ 在 $F_1'$ 处可能不再正合。
因此，$M'$ 在 $R'$ 上的一个自由分解形如
$$
F_2' \oplus F_2'' \to F_1' \to F_0' \to M' \to 0
$$
其中 $F_2''$ 是适当的自由 $R'$-模。其次注意，
$F_i \otimes_R R'/I' = F_i' / IF_i' = F_i'/I'F_i'$。
所以复形 $F_2'/I'F_2' \to F_1'/I'F_1' \to F_0'/I'F_0'$
计算 $\text{Tor}_1^R(M, R'/I')$。另一方面，
$F_i' \otimes_{R'} R'/I' = F_i'/I'F_i'$，而 $F_2''$ 也类似。
因此，复形
$F_2'/I'F_2' \oplus F_2''/I'F_2'' \to F_1'/I'F_1' \to F_0'/I'F_0'$
计算 $\text{Tor}_1^{R'}(M', R'/I')$。由于复形间的竖直映射
$$
\xymatrix{
F_2'/I'F_2' \ar[r] \ar[d] &
F_1'/I'F_1' \ar[r] \ar[d] &
F_0'/I'F_0' \ar[d] \\
F_2'/I'F_2' \oplus F_2''/I'F_2'' \ar[r] &
F_1'/I'F_1' \ar[r] &
F_0'/I'F_0'
}
$$
% P01-E0490: frozen source says cohomology for these homological Tor complexes.
显然在同调上诱导满射，结论即得。
\end{proof}

\begin{lemma}
\label{algebra-lemma-another-variant-local-criterion-flatness}
设
$$
\xymatrix{
S \ar[r] & S' \\
R \ar[r] \ar[u] & R' \ar[u]
}
$$
是局部 Noether 环的局部同态所成的交换图。设 $I \subset R$
是正常理想。设 $M$ 是有限 $S$-模。
% P01-E0491: frozen source omits "by" in the notation declaration; see ledger.
记 $I' = IR'$ 且 $M' = M \otimes_S S'$。假设
\begin{enumerate}
\item $S'$ 是张量积 $S \otimes_R R'$ 的一个局部化；
\item $M/IM$ 在 $R/I$ 上平坦；
\item 映射
$\text{Tor}_1^R(M, R/I) \to \text{Tor}_1^{R'}(M', R'/I')$ 为零。
\end{enumerate}
则 $M'$ 在 $R'$ 上平坦。
\end{lemma}

\begin{proof}
由于 $S'$ 是 $S \otimes_R R'$ 的局部化，$M'$ 是
$M \otimes_R R'$ 的局部化。注意，由引理
\ref{algebra-lemma-flat-base-change}，模
$M/IM \otimes_{R/I} R'/I'
= M \otimes_R R' /I'(M \otimes_R R')$ 在 $R'/I'$ 上平坦。
因此，$M'/I'M'$ 作为平坦模的局部化，在 $R'/I'$ 上也是平坦的。
由引理 \ref{algebra-lemma-variant-local-criterion-flatness}，只须证明
$\text{Tor}_1^{R'}(M', R'/I')$ 为零。由于 $M'$ 是
$M \otimes_R R'$ 的局部化，最后一个假设蕴含只须证明
$\text{Tor}_1^R(M, R/I) \otimes_R R'
\to
\text{Tor}_1^{R'}(M \otimes_R R', R'/I')$
是满射。

\medskip\noindent
由引理 \ref{algebra-lemma-surjective-on-tor-one-trivial}，映射
$\text{Tor}_1^R(M, R'/I') \to \text{Tor}_1^{R'}(M \otimes_R R', R'/I')$
是满射。因此现在只须证明
$\text{Tor}_1^R(M, R/I) \otimes_R R'
\to
\text{Tor}_1^R(M, R'/I')$
是满射。这由引理 \ref{algebra-lemma-surjective-on-tor-one} 得出；
只需考察环映射 $R \to R/I \to R'/I'$ 和模 $M$。
\end{proof}

\noindent
请把下面的引理与引理
\ref{algebra-lemma-criterion-flatness-fibre-nilpotent}（幂零理想的情形）
及引理 \ref{algebra-lemma-criterion-flatness-fibre}（有限表示代数的情形）比较。

\begin{lemma}[按纤维的平坦性判据；Noether 情形]
\label{algebra-lemma-criterion-flatness-fibre-Noetherian}
设 $R$、$S$、$S'$ 是 Noether 局部环，且 $R \to S \to S'$
是局部环同态。设 $\mathfrak m \subset R$ 是极大理想。
设 $M$ 是 $S'$-模。假设
\begin{enumerate}
\item 模 $M$ 在 $S'$ 上有限。
\item 模 $M$ 非零。
\item 模 $M/\mathfrak m M$ 在 $S/\mathfrak m S$ 上平坦。
\item 模 $M$ 是平坦 $R$-模。
\end{enumerate}
则 $S$ 在 $R$ 上平坦，且 $M$ 是平坦 $S$-模。
\end{lemma}

\begin{proof}
置 $I = \mathfrak mS \subset S$。由 (3)，$M/IM$ 是平坦
$S/I$-模。由于 $\mathfrak m \otimes_R S' \to I \otimes_S S'$
是满射，$\mathfrak m \otimes_R M \to I \otimes_S M$ 也是满射。
考虑
$$
\mathfrak m \otimes_R M \to I \otimes_S M \to M.
$$
因为 $M$ 在 $R$ 上平坦，复合映射是单射，故两个箭头都是单射。
特别地，$\text{Tor}_1^S(S/I, M) = 0$，见注记
\ref{algebra-remark-Tor-ring-mod-ideal}。由引理
\ref{algebra-lemma-variant-local-criterion-flatness}，可知 $M$ 在 $S$ 上平坦。
注意，由 Nakayama 引理 \ref{algebra-lemma-NAK}，
$M/\mathfrak m_{S'}M$ 非零，故由引理 \ref{algebra-lemma-ff}，$M$
实际上在 $S$ 上忠实平坦（因为这迫使
$M/\mathfrak m_SM \not = 0$）。

\medskip\noindent
考虑正合列 $0 \to \mathfrak m \to R \to \kappa \to 0$。
它给出正合列
$0 \to \text{Tor}_1^R(\kappa, S) \to \mathfrak m \otimes_R S \to I \to 0$。
由于 $M$ 在 $S$ 上平坦，又得到正合列
$0 \to \text{Tor}_1^R(\kappa, S)\otimes_S M \to
\mathfrak m \otimes_R M \to I \otimes_S M \to 0$。
由上文，这蕴含 $\text{Tor}_1^R(\kappa, S)\otimes_S M = 0$。
由于 $M$ 在 $S$ 上忠实平坦，可知
$\text{Tor}_1^R(\kappa, S) = 0$；再由引理
\ref{algebra-lemma-local-criterion-flatness}，可知 $S$ 在 $R$ 上平坦。
\end{proof}

\begin{lemma}
\label{algebra-lemma-flatness-scallop}
设 $A$ 是环，$M$ 是 $A$-模，且 $f \in A$。若
\begin{enumerate}
\item $f$ 在 $A$ 和 $M$ 上都是非零因子；
\item $M_f$ 是平坦 $A_f$-模；并且
\item $M/fM$ 是平坦 $A/fA$-模，
\end{enumerate}
则 $M$ 是平坦 $A$-模。把“平坦”换成“忠实平坦”时结论仍成立。
\end{lemma}

\begin{proof}
由于 $0 \to A \to A \to A/fA \to 0$ 和
$0 \to M \to M \to M/fM \to 0$ 正合，可知对 $i = 1$ 和 $i = 2$，
都有 $\text{Tor}_i^A(M, A/fA) = 0$。由引理
\ref{algebra-lemma-what-does-it-mean}，对所有被 $f$ 零化的 $A$-模 $N$，
都有 $\text{Tor}_1^A(M, N) = 0$（这里使用 $M/fM$ 在 $A/fA$
上的平坦性）。给定被 $f$ 零化的 $A$-模 $N$，可以选取
$A$-模的短正合列 $0 \to N' \to F \to N \to 0$，其中 $F$
是若干份 $A/fA$ 的直和。由正合列
$$
\text{Tor}_2^A(M, F) \to
\text{Tor}_2^A(M, N) \to
\text{Tor}_1^A(M, N')
$$
可知，对所有被 $f$ 零化的 $A$-模 $N$，都有
$\text{Tor}_2^A(M, N) = 0$。其次，设 $K$ 是任意 $A$-模。
可以把映射 $f : K \to K$ 分解成两个短正合列
$$
0 \to K[f] \to K \to K' \to 0
\quad\text{且}\quad
0 \to K' \to K \to K/fK \to 0
$$
其中 $K' = K/K[f] \cong fK$。应用 Tor 的正合列，得到正合列
$$
\text{Tor}_1^A(K[f], M) \to
\text{Tor}_1^A(K, M) \to
\text{Tor}_1^A(K', M)
$$
以及
$$
\text{Tor}_2^A(K/fK, M) \to
\text{Tor}_1^A(K', M) \to
\text{Tor}_1^A(K, M)
$$
利用 $\text{Tor}_1^A(K[f], M)$ 和
$\text{Tor}_2^A(K/fK, M)$ 的消没，可知 $f : K \to K$
在 $\text{Tor}^A_1(M, K)$ 上诱导单射。换言之，
$\text{Tor}^A_1(M, K)$ 为零，当且仅当
$\text{Tor}^A_1(M, K) \otimes_A A_f$ 为零。
由引理 \ref{algebra-lemma-flat-base-change-tor}，有
$$
\text{Tor}^A_1(M, K) \otimes_A A_f = \text{Tor}^{A_f}_1(M_f, K_f)  = 0
$$
% P01-E0492: frozen source has this sentence fragment twice; Chinese supplies the predicate.
最后一个消没来自 $M_f$ 在 $A_f$ 上的平坦性
（引理 \ref{algebra-lemma-characterize-flat}）。可知对所有 $A$-模 $K$，
都有 $\text{Tor}_1^A(M, K) = 0$。因此，由引理
\ref{algebra-lemma-characterize-flat}，$M$ 在 $A$ 上平坦。

\medskip\noindent
忠实平坦模情形的论证从略。
\end{proof}

\begin{lemma}
\label{algebra-lemma-flatness-scallop-pre}
设 $A$ 是环，$M$ 是 $A$-模。设 $I = (f_1, \ldots, f_r)$
是 $A$ 中由 $r \geq 1$ 个元素生成的理想。若
\begin{enumerate}
\item 对 $i = 1, \ldots, r$，$M_{f_i}$ 是平坦 $A_{f_i}$-模；
\item $M/IM$ 是平坦 $A/I$-模；
\item 对 $i = 1, \ldots, r + 1$，有
$\text{Tor}_i^A(M, A/I) = 0$。
\end{enumerate}
则 $M$ 是平坦 $A$-模。把“平坦”换成“忠实平坦”时结论仍成立。
\end{lemma}

\begin{proof}
由引理 \ref{algebra-lemma-what-does-it-mean}，对所有被 $I$ 零化的
$A$-模 $K$，都有 $\text{Tor}_1^A(M, K) = 0$。
给定被 $I$ 零化的 $A$-模 $K$，可以选取 $A$-模的短正合列
$0 \to N \to F \to K \to 0$，其中 $F$ 是若干份 $A/I$ 的直和。
得到正合列
$$
\text{Tor}_i^A(M, F) \to
\text{Tor}_i^A(M, K) \to
\text{Tor}_{i - 1}^A(M, N)
$$
因此，利用假设 (3) 并对 $i$ 归纳，可知对所有被 $I$ 零化的
$A$-模 $K$ 以及 $i = 1, \ldots, r + 1$，都有
$\text{Tor}_i^A(M, K) = 0$。

\medskip\noindent
假设对于某个 $1 \leq j \leq r$，已经证明：对所有被
$f_1, \ldots, f_j$ 零化的 $A$-模 $K$ 以及
$i = 1, \ldots, j + 1$，都有 $\text{Tor}_i^A(M, K) = 0$。
设 $K$ 是被 $f_1, \ldots, f_{j - 1}$ 零化的 $A$-模。
可以把映射 $f_j : K \to K$ 分解成两个短正合列
$$
0 \to K[f_j] \to K \to K' \to 0
\quad\text{且}\quad
0 \to K' \to K \to K/f_jK \to 0
$$
其中 $K' = K/K[f_j] \cong f_jK$。设 $1 \leq i \leq j$。
应用 Tor 的正合列，得到正合列
$$
\text{Tor}_i^A(K[f_j], M) \to
\text{Tor}_i^A(K, M) \to
\text{Tor}_i^A(K', M)
$$
以及
$$
\text{Tor}_{i + 1}^A(K/f_jK, M) \to
\text{Tor}_i^A(K', M) \to
\text{Tor}_i^A(K, M)
$$
利用 $\text{Tor}_i^A(K[f_j], M)$ 和
$\text{Tor}_{i + 1}^A(K/f_jK, M)$ 的消没，可知
$f_j : K \to K$ 在 $\text{Tor}^A_i(M, K)$ 上诱导单射。
换言之，$\text{Tor}^A_i(M, K)$ 为零，当且仅当
$\text{Tor}^A_i(M, K) \otimes_A A_{f_j}$ 为零。
由引理 \ref{algebra-lemma-flat-base-change-tor}，有
$$
\text{Tor}^A_i(M, K) \otimes_A A_{f_j} =
\text{Tor}^{A_{f_j}}_i(M_{f_j}, K_{f_j})  = 0
$$
最后一个消没来自 $M_{f_j}$ 在 $A_{f_j}$ 上的平坦性
（引理 \ref{algebra-lemma-characterize-flat}）。可知对所有被
$f_1, \ldots, f_{j - 1}$ 零化的 $A$-模 $K$ 以及
$i = 1, \ldots, j$，都有 $\text{Tor}_i^A(M, K) = 0$。
对 $j$ 作降归纳，得到 $j = 0$ 时的结论；也就是说，
对所有 $A$-模 $K$，都有 $\text{Tor}_1^A(M, K) = 0$。
因此，由引理 \ref{algebra-lemma-characterize-flat}，$M$ 在 $A$ 上平坦。

\medskip\noindent
忠实平坦模情形的论证从略。
\end{proof}

\section{基变换与平坦性}
\label{algebra-section-base-change-flat}

\noindent
下面给出若干处理基变换下平坦性变化的引理。

\begin{lemma}
\label{algebra-lemma-base-change-flat-up-down}
设
$$
\xymatrix{
S \ar[r] & S' \\
R \ar[r] \ar[u] & R' \ar[u]
}
$$
是局部环的局部同态所成的交换图。假设 $S'$ 是张量积
$S \otimes_R R'$ 的一个局部化。设 $M$ 是 $S$-模，并置
$M' = S' \otimes_S M$。
\begin{enumerate}
\item 若 $M$ 在 $R$ 上平坦，则 $M'$ 在 $R'$ 上平坦。
\item 若 $M'$ 在 $R'$ 上平坦，且 $R \to R'$ 平坦，则
$M$ 在 $R$ 上平坦。
\end{enumerate}
特别地，有
\begin{enumerate}
\item[(3)] 若 $S$ 在 $R$ 上平坦，则 $S'$ 在 $R'$ 上平坦。
\item[(4)] 若 $R' \to S'$ 和 $R \to R'$ 平坦，则 $S$ 在 $R$ 上平坦。
\end{enumerate}
\end{lemma}

\begin{proof}
(1) 的证明。若 $M$ 在 $R$ 上平坦，则由引理
\ref{algebra-lemma-flat-base-change}，$M \otimes_R R'$ 在 $R'$ 上平坦。
若 $W \subset S \otimes_R R'$ 是满足
$W^{-1}(S \otimes_R R') = S'$ 的乘法子集，则
$M' = W^{-1}(M \otimes_R R')$。因此，$M'$ 作为平坦模的局部化，
在 $R'$ 上平坦，见引理 \ref{algebra-lemma-flat-localization} 的 (5)。
这证明了 (1)，并且特别证明了 (3)。

\medskip\noindent
(2) 的证明。假设 $M'$ 在 $R'$ 上平坦，且 $R \to R'$ 平坦。
把 (3) 应用于沿西北对角线反射后的图，可知 $S \to S'$ 平坦。
于是，由引理 \ref{algebra-lemma-local-flat-ff}，$S \to S'$ 忠实平坦。
我们将用引理 \ref{algebra-lemma-flat} 的判据（\ref{algebra-item-f-ideal}）
证明 $M$ 平坦。设 $I \subset R$ 是理想。
% P01-E0493: frozen source says "has a kernel" although every map has one; intended nonzero kernel.
若 $I \otimes_R M \to M$ 的核非零，则
$(I \otimes_R M) \otimes_S S' \to M \otimes_S S' = M'$
的核也非零。注意，由于 $R \to R'$ 平坦，
$I \otimes_R R' = IR'$，并且
$$
(I \otimes_R M) \otimes_S S' =
(I \otimes_R R') \otimes_{R'} (M \otimes_S S') =
IR' \otimes_{R'} M'.
$$
由 $M'$ 在 $R'$ 上的平坦性，可知它单射到 $M'$ 中。
这完成了 (2) 的证明，因而 (4) 也成立。
\end{proof}

\noindent
下面是平坦性局部判据的又一个应用。

\begin{lemma}
\label{algebra-lemma-yet-another-variant-local-criterion-flatness}
考虑局部环和局部同态所成的交换图
$$
\xymatrix{
S \ar[r] & S' \\
R \ar[r] \ar[u] & R' \ar[u]
}
$$
设 $M$ 是有限 $S$-模。假设
\begin{enumerate}
\item 水平箭头是平坦环映射；
\item $M$ 在 $R$ 上平坦；
\item $\mathfrak m_R R' = \mathfrak m_{R'}$；
\item $R'$ 和 $S'$ 都是 Noether 环。
\end{enumerate}
则 $M' = M \otimes_S S'$ 在 $R'$ 上平坦。
\end{lemma}

\begin{proof}
由于 $\mathfrak m_R \subset R$ 且 $R \to R'$ 平坦，由假设 (3) 得到
$\mathfrak m_R \otimes_R R' = \mathfrak m_R R' = \mathfrak m_{R'}$。
注意，由引理 \ref{algebra-lemma-flatness-descends-more-general}，$M'$ 是有限
$S'$-模，且在 $R$ 上平坦。因此
$\mathfrak m_R \otimes_R M' \to M'$ 是单射。又有
$$
\mathfrak m_R \otimes_R M' =
\mathfrak m_R \otimes_R R' \otimes_{R'} M' =
\mathfrak m_{R'} \otimes_{R'} M'
$$
故 $\mathfrak m_{R'} \otimes_{R'} M' \to M'$ 是单射。
这表明 $\text{Tor}_1^{R'}(\kappa_{R'}, M') = 0$
（注记 \ref{algebra-remark-Tor-ring-mod-ideal}）。因此，由引理
\ref{algebra-lemma-local-criterion-flatness}，$M'$ 在 $R'$ 上平坦。
\end{proof}

\section{Artin 环上的平坦性判据}
\label{algebra-section-flatness-artinian}

\noindent
我们讨论 Artin 环上模的一些平坦性判据。注意，Artin 局部环的
极大理想是幂零的，所以以下两个引理适用于 Artin 局部环。

\begin{lemma}
\label{algebra-lemma-local-artinian-basis-when-flat}
设 $(R, \mathfrak m)$ 是极大理想 $\mathfrak m$ 幂零的局部环。
设 $M$ 是平坦 $R$-模。若 $A$ 是集合，而
$x_\alpha \in M$（$\alpha \in A$）是 $M$ 中的一族元素，
则以下条件等价：
\begin{enumerate}
\item $\{\overline{x}_\alpha\}_{\alpha \in A}$ 构成
$R/\mathfrak m$ 上向量空间 $M/\mathfrak mM$ 的一组基；
\item $\{x_\alpha\}_{\alpha \in A}$ 构成 $M$ 在 $R$ 上的一组基。
\end{enumerate}
\end{lemma}

\begin{proof}
(2) $\Rightarrow$ (1) 是立即的。假设 (1)。由 Nakayama 引理
\ref{algebra-lemma-NAK}，诸 $x_\alpha$ 生成 $M$。于是得到短正合列
$$
0 \to K \to \bigoplus\nolimits_{\alpha \in A} R \to M \to 0
$$
与 $R/\mathfrak m$ 张量，并使用引理 \ref{algebra-lemma-flat-tor-zero}，
得到 $K/\mathfrak mK = 0$。由 Nakayama 引理 \ref{algebra-lemma-NAK}，
可知 $K = 0$。
\end{proof}

\begin{lemma}
\label{algebra-lemma-local-artinian-characterize-flat}
设 $R$ 是极大理想幂零的局部环。设 $M$ 是 $R$-模。
以下条件等价：
\begin{enumerate}
\item $M$ 在 $R$ 上平坦；
\item $M$ 是自由 $R$-模；
\item $M$ 是投射 $R$-模。
\end{enumerate}
\end{lemma}

\begin{proof}
由于任意投射模都是平坦的（它是自由模的直和因子），而每个自由模
都是投射的，只须证明平坦模是自由的。设 $M$ 是平坦模。
设 $A$ 是集合，并设 $x_\alpha \in M$（$\alpha \in A$）是使
$\overline{x_\alpha} \in M/\mathfrak m M$ 构成 $R$ 的剩余域上的
一组基的元素。由引理 \ref{algebra-lemma-local-artinian-basis-when-flat}，
诸 $x_\alpha$ 构成 $M$ 在 $R$ 上的一组基，结论即得。
\end{proof}

\begin{lemma}
\label{algebra-lemma-lift-basis}
设 $R$ 是环。设 $I \subset R$ 是理想。设 $M$ 是 $R$-模。
设 $A$ 是集合，并设 $x_\alpha \in M$（$\alpha \in A$）是
$M$ 中的一族元素。假设
\begin{enumerate}
\item $I$ 是幂零理想；
\item $\{\overline{x}_\alpha\}_{\alpha \in A}$ 构成
$M/IM$ 在 $R/I$ 上的一组基；
\item $\text{Tor}_1^R(R/I, M) = 0$。
\end{enumerate}
则 $M$ 是以 $\{x_\alpha\}_{\alpha \in A}$ 为基的自由 $R$-模。
\end{lemma}

\begin{proof}
设 $R$、$I$、$M$、$\{x_\alpha\}_{\alpha \in A}$ 如引理中所述，
并满足假设 (1)、(2)、(3)。由 Nakayama 引理 \ref{algebra-lemma-NAK}，
诸 $x_\alpha$ 在 $R$ 上生成 $M$。假设
$\text{Tor}_1^R(R/I, M) = 0$ 蕴含有短正合列
$$
0 \to I \otimes_R M \to M \to M/IM \to 0.
$$
设 $\sum f_\alpha x_\alpha = 0$ 是 $M$ 中的一个关系。
由 $x_\alpha$ 的选取可知 $f_\alpha \in I$。因此，在
$I \otimes_R M$ 中有 $\sum f_\alpha \otimes x_\alpha = 0$。
映射 $I \otimes_R M \to I/I^2 \otimes_{R/I} M/IM$ 与
$\{x_\alpha\}_{\alpha \in A}$ 构成 $M/IM$ 的一组基这一事实
蕴含 $f_\alpha \in I^2$！因此，诸 $x_\alpha$ 在
$M/I^2M$ 中的像之间没有关系。换言之，$M/I^2M$ 是自由的，
并以诸 $x_\alpha$ 的像为基。然后使用映射
$I \otimes_R M \to I/I^3 \otimes_{R/I^2} M/I^2M$，
可知 $f_\alpha \in I^3$！如此继续。由假设 (1)，对某个 $n$
有 $I^n = 0$，结论即得。
\end{proof}

\begin{lemma}
\label{algebra-lemma-prepare-lift-flatness}
设 $\varphi : R \to R'$ 是环映射。设 $I \subset R$ 是理想。
设 $M$ 是 $R$-模。假设
\begin{enumerate}
\item $M/IM$ 在 $R/I$ 上平坦；
\item $R' \otimes_R M$ 在 $R'$ 上平坦。
\end{enumerate}
置 $I_2 = \varphi^{-1}(\varphi(I^2)R')$。则 $M/I_2M$
在 $R/I_2$ 上平坦。
\end{lemma}

\begin{proof}
可以分别用 $R/I_2$、$M/I_2M$、$R'/\varphi(I)^2R'$ 代替
$R$、$M$、$R'$。此时 $I^2 = 0$，且 $\varphi$ 是单射。
由引理 \ref{algebra-lemma-what-does-it-mean} 及 $I^2 = 0$ 这一事实，
只须证明
$\text{Tor}^R_1(R/I, M) = K = \Ker(I \otimes_R M \to M)$ 为零。
置 $M' = M \otimes_R R'$ 且 $I' = IR'$。由假设，映射
$I' \otimes_{R'} M' \to M'$ 是单射。因此，$K$ 在
$$
I' \otimes_{R'} M' = I' \otimes_R M = I' \otimes_{R/I} M/IM.
$$
中的像为零。此时 $I \to I'$ 是 $R/I$-模的单射。
由于 $M/IM$ 在 $R/I$ 上平坦，映射
$$
I \otimes_{R/I} M/IM \longrightarrow I' \otimes_{R/I} M/IM
$$
是单射。这蕴含 $K$ 在 $I \otimes_R M = I \otimes_{R/I} M/IM$
中为零，如所求。
\end{proof}

\begin{lemma}
\label{algebra-lemma-lift-flatness}
设 $\varphi : R \to R'$ 是环映射。设 $I \subset R$ 是理想。
设 $M$ 是 $R$-模。假设
\begin{enumerate}
\item $I$ 是幂零理想；
\item $R \to R'$ 是单射；
\item $M/IM$ 在 $R/I$ 上平坦；
\item $R' \otimes_R M$ 在 $R'$ 上平坦。
\end{enumerate}
则 $M$ 在 $R$ 上平坦。
\end{lemma}

\begin{proof}
归纳定义 $I_1 = I$ 以及
$I_{n + 1} = \varphi^{-1}(\varphi(I_n)^2R')$（$n \geq 1$）。
注意，由引理 \ref{algebra-lemma-prepare-lift-flatness}，对每个 $n \geq 1$，
$M/I_nM$ 都在 $R/I_n$ 上平坦。显然，
$\varphi(I_{n + 1}) \subset \varphi(I)^{2^n}R'$。由于 $I$ 幂零，
对某个 $n$ 有 $\varphi(I_n) = 0$。又因 $\varphi$ 是单射，
可知对某个 $n$ 有 $I_n = 0$，结论即得。
\end{proof}

\noindent
下面是平坦性局部判据的局部 Artin 版本。

\begin{lemma}
\label{algebra-lemma-artinian-variant-local-criterion-flatness}
设 $R$ 是 Artin 局部环。设 $M$ 是 $R$-模。
设 $I \subset R$ 是正常理想。以下条件等价：
\begin{enumerate}
\item $M$ 在 $R$ 上平坦；
\item $M/IM$ 在 $R/I$ 上平坦，且 $\text{Tor}_1^R(R/I, M) = 0$。
\end{enumerate}
\end{lemma}

\begin{proof}
(1) $\Rightarrow$ (2) 直接由定义得出。假设 $M/IM$ 在 $R/I$
上平坦，且 $\text{Tor}_1^R(R/I, M) = 0$。由引理
\ref{algebra-lemma-local-artinian-characterize-flat}，这蕴含 $M/IM$
在 $R/I$ 上自由。选取集合 $A$ 和元素 $x_\alpha \in M$，
使它们在 $M/IM$ 中的像构成一组基。由引理
\ref{algebra-lemma-lift-basis}，可知 $M$ 是自由的，因而特别是平坦的。
\end{proof}

\noindent
事实上，平坦性可沿源为 Artin 环的单射同态下降。

\begin{lemma}
\label{algebra-lemma-descent-flatness-injective-map-artinian-rings}
设 $R \to S$ 是环映射。设 $M$ 是 $R$-模。假设
\begin{enumerate}
\item $R$ 是 Artin 环；
\item $R \to S$ 是单射；
\item $M \otimes_R S$ 是平坦 $S$-模。
\end{enumerate}
则 $M$ 是平坦 $R$-模。
\end{lemma}

\begin{proof}
第一种证明：设 $I \subset R$ 是 $R$ 的 Jacobson 根。
则 $I$ 幂零；而 $M/IM$ 在 $R/I$ 上平坦，因为 $R/I$
是域的乘积，见第 \ref{algebra-section-artinian} 节。因此，应用引理
\ref{algebra-lemma-lift-flatness}，可知 $M$ 平坦。

\medskip\noindent
第二种证明：由引理 \ref{algebra-lemma-artinian-finite-length}，可将
$R = \prod R_i$ 写成有限个局部 Artin 环的乘积。这对 $R$ 和 $S$
都诱导出类似的乘积分解。因此可约化到 $R$ 是局部 Artin 环的情形
（细节略）。

\medskip\noindent
假设 $R \to S$、$M$ 如引理中所述，满足 (1)、(2)、(3)，并且
$R$ 还是极大理想为 $\mathfrak m$ 的局部环。设 $A$ 是集合，并设
% P01-E0494: frozen source types x_alpha as an element of the index set A, not of M.
$x_\alpha \in A$ 是使 $\overline{x}_\alpha$ 构成
$M/\mathfrak mM$ 在 $R/\mathfrak m$ 上的一组基的元素。
由 Nakayama 引理 \ref{algebra-lemma-NAK}，诸 $x_\alpha$ 作为 $R$-模生成
$M$。置 $N = S \otimes_R M$ 且 $I = \mathfrak mS$。
于是 $\{1 \otimes x_\alpha\}_{\alpha \in A}$ 是 $N$ 中的一族元素，
并构成 $N/IN$ 的一组基。此外，由于 $N$ 在 $S$ 上平坦，
有 $\text{Tor}_1^S(S/I, N) = 0$。因此，由引理
\ref{algebra-lemma-lift-basis}，$N$ 是以
$\{1 \otimes x_\alpha\}_{\alpha \in A}$ 为基的自由模。
$R \to S$ 的单射性于是保证，诸 $x_\alpha$ 之间不可能有系数在
$R$ 中的非平凡关系。
\end{proof}

\noindent
请把下面的引理与引理
\ref{algebra-lemma-criterion-flatness-fibre-Noetherian}（Noether 局部环的情形）、
引理 \ref{algebra-lemma-criterion-flatness-fibre}（有限表示代数的情形），
以及引理 \ref{algebra-lemma-criterion-flatness-fibre-locally-nilpotent}
（局部幂零理想的情形）比较。

\begin{lemma}[按纤维的平坦性判据：幂零情形]
\label{algebra-lemma-criterion-flatness-fibre-nilpotent}
设
$$
\xymatrix{
S \ar[rr] & & S' \\
& R \ar[lu] \ar[ru]
}
$$
是环范畴中的交换图。设 $I \subset R$ 是幂零理想，
$M$ 是 $S'$-模。假设
\begin{enumerate}
\item 模 $M/IM$ 是平坦 $S/IS$-模。
\item 模 $M$ 是平坦 $R$-模。
\end{enumerate}
则 $M$ 是平坦 $S$-模；而且对每个满足
$M \otimes_S \kappa(\mathfrak q)$ 非零的 $\mathfrak q \subset S$，
$S_{\mathfrak q}$ 在 $R$ 上平坦。
\end{lemma}

\begin{proof}
由于 $M$ 在 $R$ 上平坦，把短正合列
$0 \to I \to R \to R/I \to 0$ 与之张量，得到短正合列
$$
0 \to I \otimes_R M \to M \to M/IM \to 0.
$$
注意，$I \otimes_R M \to IS \otimes_S M$ 是满射。结合上式，
这意味着复合映射中的两个映射
$$
I \otimes_R M \to IS \otimes_S M \to M
$$
都是单射。因此
% P01-E0495: frozen source has Tor_1^S(IS,M); the ideal-kernel criterion requires S/IS.
$\text{Tor}_1^S(IS, M) = 0$（见注记
\ref{algebra-remark-Tor-ring-mod-ideal}）；再由引理
\ref{algebra-lemma-what-does-it-mean}，可知 $M$ 是平坦 $S$-模。
最后，须证明对任何满足 $M \otimes_S \kappa(\mathfrak q)$ 非零的
素理想 $\mathfrak q \subset S$，$S_{\mathfrak q}$ 在 $R$ 上平坦。
这由引理 \ref{algebra-lemma-ff} 和 \ref{algebra-lemma-flat-permanence} 得出。
\end{proof}

\section{复形何时正合？}
\label{algebra-section-complex-exact}

\noindent
本节的部分材料可见 Buchsbaum 和 Eisenbud 的论文
\cite{WhatExact}。

\begin{situation}
\label{algebra-situation-complex}
这里 $R$ 是环，并且有复形
$$
0
\to
R^{n_e}
\xrightarrow{\varphi_e}
R^{n_{e-1}}
\xrightarrow{\varphi_{e-1}}
\ldots
\xrightarrow{\varphi_{i + 1}}
R^{n_i}
\xrightarrow{\varphi_i}
R^{n_{i-1}}
\xrightarrow{\varphi_{i-1}}
\ldots
\xrightarrow{\varphi_1}
R^{n_0}
$$
换言之，对 $i = 1, \ldots, e - 1$，要求
$\varphi_i \circ \varphi_{i + 1} = 0$。
\end{situation}

\begin{lemma}
\label{algebra-lemma-add-trivial-complex}
设 $R$ 是环。设
$$
\ldots
\xrightarrow{\varphi_{i + 1}}
R^{n_i}
\xrightarrow{\varphi_i}
R^{n_{i-1}}
\xrightarrow{\varphi_{i-1}}
\ldots
$$
是有限自由 $R$-模的复形。假设对某个 $i$，映射 $\varphi_i$
的某个矩阵系数可逆。则所示复形同构于一个复形
$$
\ldots \to
R^{n_{i + 2}} \xrightarrow{\varphi_{i + 2}}
R^{n_{i + 1}} \to
R^{n_i - 1} \to
R^{n_{i - 1} - 1} \to
R^{n_{i - 2}} \xrightarrow{\varphi_{i - 2}}
R^{n_{i - 3}} \to
\ldots
$$
与复形 $\ldots \to 0 \to R \to R \to 0 \to \ldots$ 的直和，
其中映射 $R \to R$ 是恒等映射。
\end{lemma}

\begin{proof}
该假设意味着，在对 $R^{n_i}$ 和 $R^{n_{i-1}}$ 换基之后，
$R^{n_i}$ 的第一个基向量经 $\varphi_i$ 映到
$R^{n_{i-1}}$ 的第一个基向量。用 $e_j$ 表示 $R^{n_i}$
的第 $j$ 个基向量，用 $f_k$ 表示 $R^{n_{i-1}}$ 的第 $k$
个基向量。写成 $\varphi_i(e_j) = \sum a_{jk} f_k$。
于是除 $k = 1$ 外均有 $a_{1k} = 0$，且 $a_{11} = 1$。
再次更换 $R^{n_i}$ 的基：对 $j > 1$ 置
$e'_j = e_j - a_{j1} e_1$。换坐标后，对 $j > 1$ 有
$a_{j1} = 0$。注意，映射 $R^{n_{i + 1}} \to R^{n_i}$
的像包含于由 $e_j$（$j > 1$）张成的
% P01-E0497: frozen source says "subspace" although R is an arbitrary ring.
子模。还要注意，因为 $f_1$ 在像中，映射
$R^{n_{i-1}} \to R^{n_{i-2}}$ 必须将它零化。
这些条件以及 $\varphi_i$ 的矩阵 $(a_{jk})$ 的形状即蕴含本引理。
\end{proof}

\noindent
在情形 \ref{algebra-situation-complex} 中，我们称形如
$$
0 \to \ldots \to 0 \to R \xrightarrow{1} R \to 0 \to \ldots \to 0
$$
或形如
$$
0 \to \ldots \to 0 \to R
$$
的复形为{\it 平凡的}。更准确地说，若存在
$e \geq i \geq 1$，使得 $n_i = n_{i - 1} = 1$、
$\varphi_i = \text{id}_R$，且对
$j \not \in \{i, i - 1\}$ 有 $n_j = 0$；或者若
$n_0 = 1$ 且对 $i > 0$ 有 $n_i = 0$，则称
$0 \to R^{n_e} \to R^{n_{e-1}} \to \ldots \to R^{n_0}$
平凡。上面的引理显然说明，局部环上的任意有限自由模有限复形，
在与平凡复形作直和的意义下，都等同于一个所有映射的全部矩阵系数
均属于极大理想的复形。

\begin{lemma}
\label{algebra-lemma-exact-depth-zero-local}
在情形 \ref{algebra-situation-complex} 中，假设 $R$ 是极大理想为
$\mathfrak m$ 的 Noether 局部环。假设
$\mathfrak m \in \text{Ass}(R)$，换言之，$R$ 的深度为 $0$。
假设
$0 \to R^{n_e} \to R^{n_{e-1}} \to \ldots \to R^{n_0}$
在 $R^{n_e}, \ldots, R^{n_1}$ 处正合。则该复形同构于平凡复形
的一个直和。
\end{lemma}

\begin{proof}
选取满足 $x \not = 0$ 且 $\mathfrak m x = 0$ 的 $x \in R$。
设 $i$ 是满足 $n_i > 0$ 的最大指标。若 $i = 0$，则命题成立。
若 $i > 0$，用 $f_1$ 表示 $R^{n_i}$ 的第一个基向量。
由复形的正合性，$xf_1$ 不会映到零。因此，映射
$R^{n_i} \to R^{n_{i - 1}}$ 的某个矩阵系数不属于
$\mathfrak m$。于是，引理 \ref{algebra-lemma-add-trivial-complex}
使我们可以减小 $n_e + \ldots + n_1$。归纳即完成证明。
\end{proof}

\begin{lemma}
\label{algebra-lemma-exact-artinian-local}
在情形 \ref{algebra-situation-complex} 中，设 $R$ 是 Artin 局部环。
假设 $0 \to R^{n_e} \to R^{n_{e-1}} \to \ldots \to R^{n_0}$
在 $R^{n_e}, \ldots, R^{n_1}$ 处正合。则该复形同构于平凡复形
的一个直和。
\end{lemma}

\begin{proof}
这是引理 \ref{algebra-lemma-exact-depth-zero-local} 的特殊情形，因为
Artin 局部环的深度为 $0$。
\end{proof}

\noindent
下面定义有限自由模之间映射的秩。这只是秩的一种可能定义，
也只是适合本节使用的定义；在其他场合，别的定义或许更方便。

\begin{definition}
\label{algebra-definition-rank}
设 $R$ 是环。设 $\varphi : R^m \to R^n$ 是有限自由模之间的映射。
\begin{enumerate}
\item $\varphi$ 的{\it 秩}是使
$\wedge^r \varphi : \wedge^r R^m \to \wedge^r R^n$
非零的最大 $r$。
\item 令 $I(\varphi) \subset R$ 为 $\varphi$ 的矩阵的所有
$r \times r$ 子式生成的理想，其中 $r$ 是上面定义的秩。
\end{enumerate}
\end{definition}

\noindent
映射 $\varphi : R^m \to R^n$ 的秩为 $0$ 当且仅当
$\varphi = 0$；在这种情形下，$I(\varphi) = R$。

\begin{lemma}
\label{algebra-lemma-trivial-case-exact}
在情形 \ref{algebra-situation-complex} 中，假设该复形同构于平凡复形
的一个直和。则
\begin{enumerate}
\item 映射 $\varphi_i$ 的秩为
$r_i = n_i - n_{i + 1} + \ldots + (-1)^{e-i-1} n_{e-1} + (-1)^{e-i} n_e$；
\item 对所有 $i$，$1 \leq i \leq e - 1$，有
$\text{rank}(\varphi_{i + 1}) + \text{rank}(\varphi_i) = n_i$；
\item 每个 $I(\varphi_i) = R$。
\end{enumerate}
\end{lemma}

\begin{proof}
可以假设该复形就是平凡复形的直和。于是，对每个 $i$，
可以把 $R^{n_i}$ 的标准基元素分成两类：一类映到
$R^{n_{i-1}}$ 的基元素，另一类映到零（并且当 $i > 0$ 时，
后一类元素是 $R^{n_{i + 1}}$ 中一些基元素的像）。
从 $i = e$ 开始作降归纳，很容易证明：$R^{n_i}$ 中有
$r_{i + 1}$ 个基元素映到零，并有 $r_i$ 个基元素映到
$R^{n_{i-1}}$ 的基元素。结论由此得出。
\end{proof}

\begin{lemma}
\label{algebra-lemma-div-x-exact-one-less}
在情形 \ref{algebra-situation-complex} 中，假设 $R$ 是极大理想为
$\mathfrak m$ 的局部环。假设
$0 \to R^{n_e} \to R^{n_{e-1}} \to \ldots \to R^{n_0}$
在 $R^{n_e}, \ldots, R^{n_1}$ 处正合。设
$x \in \mathfrak m$ 是非零因子。则复形
$0 \to (R/xR)^{n_e} \to \ldots \to (R/xR)^{n_1}$
在 $(R/xR)^{n_e}, \ldots, (R/xR)^{n_2}$ 处正合。
\end{lemma}

\begin{proof}
用 $F_\bullet$ 表示项为 $F_i = R^{n_i}$、微分由
$\varphi_i$ 给出的复形。于是有复形的短正合列
$$
0 \to F_\bullet \xrightarrow{x} F_\bullet \to F_\bullet/xF_\bullet \to 0
$$
应用蛇引理，得到长正合列
$$
H_i(F_\bullet) \xrightarrow{x} H_i(F_\bullet) \to
H_i(F_\bullet/xF_\bullet) \to H_{i - 1}(F_\bullet)
\xrightarrow{x} H_{i - 1}(F_\bullet)
$$
引理由此得出。
\end{proof}

\begin{lemma}[无环性引理]
\label{algebra-lemma-acyclic}
\begin{reference}
\cite[Lemma 1.8]{Peskine-Szpiro}
\end{reference}
设 $R$ 是 Noether 局部环。设
$0 \to M_e \to M_{e-1} \to \ldots \to M_0$
是有限 $R$-模的复形。假设 $\text{depth}(M_i) \geq i$。
设 $i$ 是使复形在 $M_i$ 处不正合的最大指标。若 $i > 0$，则
$$
\Ker(M_i \to M_{i-1})/\Im(M_{i + 1} \to M_i)
$$
的深度 $\geq 1$。
\end{lemma}

\begin{proof}
令 $H = \Ker(M_i \to M_{i-1})/\Im(M_{i + 1} \to M_i)$ 为所讨论的
% P01-E0501: frozen source calls H a cohomology group although this is a homological complex.
同调群。可以把该复形拆成正合列
% P01-E0502: frozen source calls every listed sequence short exact, but one ends at M_{i-1}.
$0 \to M_e \to M_{e-1} \to K_{e-2} \to 0$、
$0 \to K_j \to M_j \to K_{j-1} \to 0$（$i + 2 \leq j \leq e-2 $）、
$0 \to K_{i + 1} \to M_{i + 1} \to B_i \to 0$、
$0 \to K_i \to M_i \to M_{i-1}$，以及
$0 \to B_i \to K_i \to H \to 0$。
我们沿这些序列逐步向上，反复使用引理
\ref{algebra-lemma-depth-in-ses} 来证明关于深度的断言。
首先，由于 $\text{depth}(M_e) \geq e$ 且
$\text{depth}(M_{e-1}) \geq e-1$，可知
$\text{depth}(K_{e-2}) \geq e - 1$。此时，对
$i + 2 \leq j \leq e-2$，序列
$0 \to K_j \to M_j \to K_{j-1} \to 0$ 同样蕴含
$\text{depth}(K_{j-1}) \geq j$（$i + 2 \leq j \leq e-2$）。
序列
$0 \to K_{i + 1} \to M_{i + 1} \to B_i \to 0$
于是表明 $\text{depth}(B_i) \geq i + 1$。序列
$0 \to K_i \to M_i \to M_{i-1}$ 表明
$\text{depth}(K_i) \geq 1$，因为由假设
$M_i$ 的深度 $\geq i \geq 1$。最后，序列
$0 \to B_i \to K_i \to H \to 0$ 蕴含结论。
\end{proof}

\begin{proposition}
\label{algebra-proposition-what-exact}
\begin{reference}
\cite[Corollary 1]{WhatExact}
\end{reference}
在情形 \ref{algebra-situation-complex} 中，假设 $R$ 是 Noether 局部环。
以下条件等价：
\begin{enumerate}
\item $0 \to R^{n_e} \to R^{n_{e-1}} \to \ldots \to R^{n_0}$
在 $R^{n_e}, \ldots, R^{n_1}$ 处正合；
\item 对所有 $i$，$1 \leq i \leq e$，以下两个条件成立：
\begin{enumerate}
\item $\text{rank}(\varphi_i) = r_i$，其中
$r_i = n_i - n_{i + 1} + \ldots + (-1)^{e-i-1} n_{e-1} + (-1)^{e-i} n_e$；
\item $I(\varphi_i) = R$，或 $I(\varphi_i)$ 包含一个长度为 $i$
的正则序列。
\end{enumerate}
\end{enumerate}
\end{proposition}

\begin{proof}
若对某个 $i$，$\varphi_i$ 的某个矩阵系数不属于
$\mathfrak m$，则应用引理 \ref{algebra-lemma-add-trivial-complex}。
不难看出，一个复形的命题与给该复形加上一个平凡复形后所得复形
的命题等价。因此，可以假设每个 $\varphi_i$ 的所有矩阵元素
都属于极大理想。还可以假设 $e \geq 1$。

\medskip\noindent
假设该复形在 $R^{n_e}, \ldots, R^{n_1}$ 处正合。
设 $\mathfrak q \in \text{Ass}(R)$。注意，环
$R_{\mathfrak q}$ 的深度为 $0$，而该复形在
$\mathfrak q$ 处局部化后仍然正合。对 $R_{\mathfrak q}$
上的局部化复形应用引理 \ref{algebra-lemma-exact-depth-zero-local}
和引理 \ref{algebra-lemma-trivial-case-exact}，可知对所有 $i$，
$\varphi_{i, \mathfrak q}$ 的秩为 $r_i$。由于
$R \to \bigoplus_{\mathfrak q \in \text{Ass}(R)} R_\mathfrak q$
是单射（引理 \ref{algebra-lemma-zero-at-ass-zero}），由定义
\ref{algebra-definition-rank} 中给出的秩的定义可知，$\varphi_i$
在 $R$ 上的秩为 $r_i$。因此，由于秩不变，可见
$I(\varphi_i)_\mathfrak q = I(\varphi_{i, \mathfrak q})$。
上面引用的引理说明，对所有 $e \geq i \geq 1$，理想
$I(\varphi_i)_{\mathfrak q}$ 都等于 $R_{\mathfrak q}$；
故没有任何 $I(\varphi_i)$ 包含于 $\mathfrak q$。这蕴含
$I(\varphi_e)I(\varphi_{e-1})\ldots I(\varphi_1)$ 不包含于
$R$ 的任何伴随素理想。由引理 \ref{algebra-lemma-silly}，可以选取
$$
x \in I(\varphi_e)I(\varphi_{e - 1})\ldots I(\varphi_1)
$$
使得对所有 $\mathfrak q \in \text{Ass}(R)$，
$x \not \in \mathfrak q$。注意，$x$ 是非零因子
（引理 \ref{algebra-lemma-ass-zero-divisors}）。由引理
\ref{algebra-lemma-div-x-exact-one-less}，复形
$0 \to (R/xR)^{n_e} \to \ldots \to (R/xR)^{n_1}$
在 $(R/xR)^{n_e}, \ldots, (R/xR)^{n_2}$ 处正合。
对 $e$ 归纳可知，所有理想 $I(\varphi_i)/xR$ 都含有一个
长度为 $i - 1$ 的正则序列。这证明了 $I(\varphi_i)$
含有一个长度为 $i$ 的正则序列。

\medskip\noindent
假设 (2)(a) 和 (2)(b) 成立。我们对 $\dim(R)$ 归纳证明 (1)。
若 $\dim(R) = 0$，则由 (2)(b)，对 $1 \leq i \leq e$
必有 $I(\varphi_i) = R$。由于 $\varphi_i$ 的系数属于极大理想，
这只有在对所有 $i$ 都有 $r_i = 0$ 时才可能。由 (2)(a)，
可知 $e = 0$，故 (1) 成立。假设 $\dim(R) > 0$。
我们断言，对任何素理想 $\mathfrak p \subset R$，条件
(2)(a) 和 (2)(b) 对 $R_\mathfrak p$ 上的复形
$$
0 \to R_\mathfrak p^{n_e} \to R_\mathfrak p^{n_{e - 1}} \to
\ldots \to R_\mathfrak p^{n_0}
$$
以及映射 $\varphi_{i, \mathfrak p}$ 仍然成立。事实上，由于
$I(\varphi_i)$ 含有非零因子，$I(\varphi_i)$ 在
$R_\mathfrak p$ 中的像非零。
这蕴含 $\varphi_{i, \mathfrak p}$ 的秩与 $\varphi_i$ 的秩相同：
上面定义的环 $R$ 上矩阵 $\varphi$ 的秩，在传递到 $R$-代数
$R'$ 时只能下降，而且下降当且仅当 $I(\varphi)$ 在 $R'$
中的像为零。因此 (2)(a) 成立。由此可知
$I(\varphi_{i, \mathfrak p}) = I(\varphi_i)_\mathfrak p$，
并且 (2)(b) 在局部化下也保持成立。
由对 $R$ 的维数的归纳，可以假设该复形在 $R$ 的任何非极大
素理想 $\mathfrak p$ 处局部化后都正合。因此
$\Ker(\varphi_i)/\Im(\varphi_{i + 1})$ 的支撑包含于
$\{\mathfrak m\}$，故它若非零，深度便为 $0$。
由于第一段所述的理由，对所有 $i$ 都有
% P01-E0503: frozen source overlooks rank-zero maps, for which I(varphi_i)=R.
$I(\varphi_i) \subset \mathfrak m$，所以 (2)(b) 蕴含
$\text{depth}(R) \geq e$。由引理 \ref{algebra-lemma-acyclic}，
该复形在 $R^{n_e}, \ldots, R^{n_1}$ 处正合，证明完成。
\end{proof}

\begin{remark}
\label{algebra-remark-what-exact}
若命题 \ref{algebra-proposition-what-exact} 中的等价条件 (1) 和 (2)
成立，则存在 $j$，使得 $I(\varphi_i) = R$ 当且仅当
$i \geq j$。与命题的证明一样，只须在所有矩阵的系数均属于
$R$ 的极大理想 $\mathfrak m$ 时验证这一点。在这种情形下，
$I(\varphi_j) = R$ 当且仅当 $\varphi_j = 0$。但若
$\varphi_j = 0$，则得到任意长的正合复形
$0 \to R^{n_e} \to R^{n_{e-1}} \to \ldots
\to R^{n_j} \to 0 \to 0 \to \ldots \to 0$，
因而由该命题可知，对 $i > j$，$I(\varphi_i)$ 必须为 $R$
（否则它是 Noether 局部环的一个真理想，却包含任意长的正则序列，
这是不可能的）。
\end{remark}
\section{Cohen--Macaulay 模}
\label{algebra-section-CM}

\noindent
这里将说明 Cohen--Macaulay 模具有良好的性质。我们暂不使用
Ext 群来建立它们与对偶性等概念的联系。

\begin{definition}
\label{algebra-definition-CM}
设 $R$ 是 Noether 局部环。设 $M$ 是有限 $R$-模。若
$\dim(\text{Supp}(M)) = \text{depth}(M)$，则称 $M$ 是
{\it Cohen--Macaulay 模}。
\end{definition}

\noindent
我们的第一个目标是证明命题 \ref{algebra-proposition-CM-module}。
为此，将给出一系列（或许并不标准的）初等引理，其中几乎不使用
前面关于深度的结果。先引入一些记号。

\medskip\noindent
设 $R$ 是 Noether 局部环。设 $M$ 是 Cohen--Macaulay 模，
$f_1, \ldots, f_d$ 是一个 $M$-正则序列，其中
$d = \dim(\text{Supp}(M))$。若对所有
$i = 0, 1, \ldots, d-1$ 都有
$\dim (\text{Supp}(M) \cap V(g, f_1, \ldots, f_i))
= d - i - 1$，则称 $g \in \mathfrak m$ {\it 关于
$(M, f_1, \ldots, f_d)$ 良好}。这等价于：对
$i = 0, 1, \ldots, d - 1$，有
$\dim(\text{Supp}(M/(f_1, \ldots, f_i)M) \cap V(g)) =
d - i - 1$。

\begin{lemma}
\label{algebra-lemma-good-element}
记号和假设同上。若 $g$ 关于 $(M, f_1, \ldots, f_d)$ 良好，
则 (a) $g$ 是 $M$ 上的非零因子，并且 (b) $M/gM$ 是
Cohen--Macaulay 模，其极大正则序列为
$f_1, \ldots, f_{d - 1}$。
\end{lemma}

\begin{proof}
对 $d$ 归纳证明该引理。
% P01-E0504: for d=0 the definition of good is vacuous, contrary to the source's claim that no case applies.
若 $d = 0$，则 $M$ 有限，并且不存在本引理适用的情形。
若 $d = 1$，则须证明 $g : M \to M$ 是单射。由于假设
$\dim \text{Supp}(M) \cap V(g) = 0$，其核 $K$ 的支撑为
$\{\mathfrak m\}$，故 $K$ 具有有限长度。因此，
$f_1 : K \to K$ 的单射性蕴含其像与 $K$ 长度相同，从而
$f_1 K = K$；再由 Nakayama 引理 \ref{algebra-lemma-NAK} 可知
$K = 0$。此外，$\dim \text{Supp}(M/gM) = 0$，所以
$M/gM$ 是深度为 $0$ 的 Cohen--Macaulay 模。

\medskip\noindent
假设 $d > 1$。由定义容易看出，$g$ 关于
$(M/f_1M, f_2, \ldots, f_d)$ 良好。由归纳假设，
(a) $g$ 是 $M/f_1M$ 上的非零因子，并且
(b) $M/(g, f_1)M$ 是 Cohen--Macaulay 模，其极大正则序列为
$f_2, \ldots, f_{d - 1}$。由引理 \ref{algebra-lemma-permute-xi}，
可知 $g, f_1$ 是一个 $M$-正则序列。因此，$g$ 是 $M$ 上的
非零因子，而 $f_1, \ldots, f_{d - 1}$ 是一个
$M/gM$-正则序列。
\end{proof}

\begin{lemma}
\label{algebra-lemma-CM-one-g}
设 $R$ 是 Noether 局部环。设 $M$ 是 $R$ 上的
Cohen--Macaulay 模。假设 $g \in \mathfrak m$ 满足
$\dim(\text{Supp}(M) \cap V(g))
= \dim(\text{Supp}(M)) - 1$。则 (a) $g$ 是 $M$ 上的非零因子，
并且 (b) $M/gM$ 是深度少一的 Cohen--Macaulay 模。
\end{lemma}

\begin{proof}
选取一个 $M$-正则序列 $f_1, \ldots, f_d$，其中
$d = \dim(\text{Supp}(M))$。若 $g$ 关于
$(M, f_1, \ldots, f_d)$ 良好，则由引理
\ref{algebra-lemma-good-element} 即得结论。特别地，当 $d = 1$ 时
引理成立。（$d = 0$ 的情形不会出现。）假设 $d > 1$。
选取元素 $h \in R$，使得
(\romannumeral1) $h$ 关于 $(M, f_1, \ldots, f_d)$ 良好，并且
(\romannumeral2) $\dim(\text{Supp}(M) \cap V(h, g)) = d - 2$。
为说明这样的 $h$ 存在，令 $\{\mathfrak q_j\}$ 为以下闭集的
极小素理想所成的（有限）集合：
$\text{Supp}(M)$，
$\text{Supp}(M)\cap V(f_1, \ldots, f_i)$
（$i = 1, \ldots, d - 1$），以及
$\text{Supp}(M) \cap V(g)$。这些 $\mathfrak q_j$ 都不等于
$\mathfrak m$，故由引理 \ref{algebra-lemma-silly}，可以找到
$h \in \mathfrak m$ 使 $h \not \in \mathfrak q_j$。
显然，$h$ 满足 (\romannumeral1) 和 (\romannumeral2)。
由引理 \ref{algebra-lemma-good-element}，$M/hM$ 是 Cohen--Macaulay 模。
由 (\romannumeral2)，二元组 $(M/hM, g)$ 满足归纳假设。
因此，$M/(h, g)M$ 是 Cohen--Macaulay 模，且
$g : M/hM \to M/hM$ 是单射。由引理
\ref{algebra-lemma-permute-xi} 可知，$g : M \to M$ 和
$h : M/gM \to M/gM$ 都是单射。结合 $M/(g, h)M$
是 Cohen--Macaulay 模这一事实，证明完成。
\end{proof}

\begin{proposition}
\label{algebra-proposition-CM-module}
设 $R$ 是极大理想为 $\mathfrak m$ 的 Noether 局部环。
设 $M$ 是 $R$ 上的 Cohen--Macaulay 模，其支撑的维数为 $d$。
假设 $g_1, \ldots, g_c$ 是 $\mathfrak m$ 中的元素，并且
$\dim(\text{Supp}(M/(g_1, \ldots, g_c)M))
= d - c$。则 $g_1, \ldots, g_c$ 是一个 $M$-正则序列，
并且可以延拓为一个极大 $M$-正则序列。
\end{proposition}

\begin{proof}
令 $Z = \text{Supp}(M) \subset \Spec(R)$。由引理
\ref{algebra-lemma-one-equation}，在链
$Z \supset Z \cap V(g_1) \supset \ldots \supset
Z \cap V(g_1, \ldots, g_c)$ 中，每一步至多使维数下降 $1$。
因此，由假设，每一步都恰好使维数下降 $1$。于是，可以依次把
引理 \ref{algebra-lemma-CM-one-g} 应用于模
% P01-E0505: frozen source omits the terminal M in this module quotient.
$M/(g_1, \ldots, g_i)$ 和元素 $g_{i + 1}$。

\medskip\noindent
若 $c < d$，为了用一个元素延拓 $g_1, \ldots, g_c$，只须使用
引理 \ref{algebra-lemma-silly}，选取一个不属于
$Z \cap V(g_1, \ldots, g_c)$ 的有限多个极小素理想中任何一个的
元素 $g_{c + 1} \in \mathfrak m$。
\end{proof}

\noindent
命题 \ref{algebra-proposition-CM-module} 已得证；下面继续发展标准理论。

\begin{lemma}
\label{algebra-lemma-nonzerodivisor-on-CM}
设 $R$ 是极大理想为 $\mathfrak m$ 的 Noether 局部环。
设 $M$ 是有限 $R$-模。设 $x \in \mathfrak m$ 是 $M$ 上的
非零因子。则 $M$ 是 Cohen--Macaulay 模，当且仅当
$M/xM$ 是 Cohen--Macaulay 模。
\end{lemma}

\begin{proof}
由引理 \ref{algebra-lemma-depth-drops-by-one}，有
$\text{depth}(M/xM) = \text{depth}(M)-1$。
由引理 \ref{algebra-lemma-one-equation-module}，有
$\dim(\text{Supp}(M/xM)) = \dim(\text{Supp}(M)) - 1$。
\end{proof}

\begin{lemma}
\label{algebra-lemma-CM-over-quotient}
设 $R \to S$ 是 Noether 局部环的满同态。设 $N$ 是有限
$S$-模。则 $N$ 作为 $S$-模是 Cohen--Macaulay 模，当且仅当
$N$ 作为 $R$-模是 Cohen--Macaulay 模。
\end{lemma}

\begin{proof}
略。
\end{proof}

\begin{lemma}
\label{algebra-lemma-CM-ass-minimal-support}
\begin{reference}
\cite[Chapter 0, Proposition 16.5.4]{EGA}
\end{reference}
设 $R$ 是 Noether 局部环。设 $M$ 是有限 Cohen--Macaulay
$R$-模。若 $\mathfrak p \in \text{Ass}(M)$，则
$\dim(R/\mathfrak p) = \dim(\text{Supp}(M))$，并且
$\mathfrak p$ 是 $M$ 的支撑中的极小素理想。特别地，
$M$ 没有嵌入伴随素理想。
\end{lemma}

\begin{proof}
由引理 \ref{algebra-lemma-depth-dim-associated-primes}，有
$\text{depth}(M) \leq \dim(R/\mathfrak p)$。
当然，由于 $\mathfrak p \in \text{Supp}(M)$
（引理 \ref{algebra-lemma-ass-support}），有
$\dim(R/\mathfrak p) \leq \dim(\text{Supp}(M))$。
因为 $M$ 是 Cohen--Macaulay 模，这两个不等式都取等号。
于是 $\mathfrak p$ 必须在 $\text{Supp}(M)$ 中极小，否则会有
$\dim(R/\mathfrak p) < \dim(\text{Supp}(M))$。
最后，$M$ 的支撑中的极小素理想等于 $\text{Ass}(M)$ 的极小元素
（命题 \ref{algebra-proposition-minimal-primes-associated-primes}），
所以 $M$ 没有嵌入伴随素理想
（定义 \ref{algebra-definition-embedded-primes}）。
\end{proof}

\begin{definition}
\label{algebra-definition-maximal-CM}
设 $R$ 是 Noether 局部环。若有限 $R$-模 $M$ 满足
$\text{depth}(M) = \dim(R)$，则称它为{\it 极大
Cohen--Macaulay 模}。
\end{definition}

\noindent
换言之，Noether 局部环上的极大 Cohen--Macaulay 模就是该环上
深度尽可能大的有限模。等价地，Noether 局部环 $R$ 上的极大
Cohen--Macaulay 模就是维数等于环的维数的 Cohen--Macaulay 模。
特别地，若 $M$ 是满足 $\Spec(R) = \text{Supp}(M)$ 的
Cohen--Macaulay $R$-模，则 $M$ 是极大 Cohen--Macaulay 模。
所以下面的两个引理讨论极大 Cohen--Macaulay 模。

\begin{lemma}
\label{algebra-lemma-maximal-chain-maximal-CM}
\begin{slogan}
在 Cohen--Macaulay 局部环中，任意极大素理想链的长度都等于
环的维数。
\end{slogan}
设 $R$ 是 Noether 局部环。假设存在满足
$\Spec(R) = \text{Supp}(M)$ 的 Cohen--Macaulay 模 $M$。
则任意极大素理想链
$\mathfrak p_0 \subset \mathfrak p_1 \subset \ldots
\subset \mathfrak p_n$ 的长度均为 $n = \dim(R)$。
\end{lemma}

\begin{proof}
对 $\dim(R)$ 归纳证明。若 $\dim(R) = 0$，结论显然。
假设 $\dim(R) > 0$。于是 $n > 0$。选取
$x \in \mathfrak p_1$，使 $x$ 不属于 $R$ 的任何极小素理想，
特别地，$x \not \in \mathfrak p_0$。（见引理
\ref{algebra-lemma-silly}。）由引理 \ref{algebra-lemma-one-equation}，
$\dim(R/xR) = \dim(R) - 1$。由命题
\ref{algebra-proposition-CM-module} 和引理
\ref{algebra-lemma-CM-over-quotient}，$M/xM$ 是 $R/xR$ 上的
Cohen--Macaulay 模。由引理 \ref{algebra-lemma-support-quotient}，
$M/xM$ 的支撑为 $\Spec(R/xR)$。把 $x$ 替换为
$x^n$（对某个 $n$）后，可以假设 $\mathfrak p_1$ 是
$M/xM$ 的伴随素理想，
见引理 \ref{algebra-lemma-inherit-minimal-primes}。由引理
\ref{algebra-lemma-CM-ass-minimal-support}，可知
$\mathfrak p_1/(x)$ 是 $R/xR$ 的极小素理想。因此，链
$\mathfrak p_1/(x) \subset \ldots \subset \mathfrak p_n/(x)$
是 $R/xR$ 中的极大素理想链。由归纳假设，该链的长度为
$\dim(R/xR) = \dim(R) - 1$，正如所需。
\end{proof}

\begin{lemma}
\label{algebra-lemma-dim-formula-maximal-CM}
设 $R$ 是 Noether 局部环。假设存在满足
$\Spec(R) = \text{Supp}(M)$ 的 Cohen--Macaulay 模 $M$。
则对素理想 $\mathfrak p \subset R$，有
$$
\dim(R) = \dim(R_{\mathfrak p}) + \dim(R/\mathfrak p).
$$
\end{lemma}

\begin{proof}
这立即由引理 \ref{algebra-lemma-maximal-chain-maximal-CM} 得出。
\end{proof}

\begin{lemma}
\label{algebra-lemma-localize-CM-module}
设 $R$ 是 Noether 局部环。设 $M$ 是 $R$ 上的
Cohen--Macaulay 模。对任意素理想 $\mathfrak p \subset R$，
模 $M_{\mathfrak p}$ 都是 $R_\mathfrak p$ 上的
Cohen--Macaulay 模。
\end{lemma}

\begin{proof}
可以且确实假设 $\mathfrak p \not = \mathfrak m$ 且 $M$ 非零。
选取极大素理想链
$\mathfrak p = \mathfrak p_c \subset
\mathfrak p_{c - 1} \subset \ldots \subset \mathfrak p_1
\subset \mathfrak m$。若能证明 $M_{\mathfrak p_1}$ 在
$R_{\mathfrak p_1}$ 上的结论，则对 $c$ 归纳即可得到本引理。
因此，可以假设 $\mathfrak p$ 与 $\mathfrak m$ 之间不存在
严格位于二者之间的素理想。注意，
$\dim(\text{Supp}(M_\mathfrak p))
\leq \dim(\text{Supp}(M)) - 1$，因为 $M_\mathfrak p$
的支撑中的任意素理想链都可以在 $M$ 的支撑中再延长一个素理想
（即 $\mathfrak m$）。另一方面，由引理
\ref{algebra-lemma-depth-localization} 以及对 $\mathfrak p$ 的选择，有
$\text{depth}(M_\mathfrak p) \geq
\text{depth}(M) - \dim(R/\mathfrak p) =
\text{depth}(M) - 1$。因此，正如所需，
$\text{depth}(M_\mathfrak p) \geq
\dim(\text{Supp}(M_\mathfrak p))$
（另一个不等式是引理 \ref{algebra-lemma-bound-depth}）。
\end{proof}

\begin{definition}
\label{algebra-definition-module-CM}
设 $R$ 是 Noether 环。设 $M$ 是有限 $R$-模。若对 $R$ 的所有
素理想 $\mathfrak p$，$M_\mathfrak p$ 都是
$R_\mathfrak p$ 上的 Cohen--Macaulay 模，则称 $M$ 是
{\it Cohen--Macaulay 模}。
\end{definition}

\noindent
由引理 \ref{algebra-lemma-localize-CM-module}，只须在 $R$ 的极大理想
处检验这一条件。

\begin{lemma}
\label{algebra-lemma-maximal-CM-polynomial-algebra}
设 $R$ 是 Noether 环。设 $M$ 是 $R$ 上的 Cohen--Macaulay 模。
则 $M \otimes_R R[x_1, \ldots, x_n]$ 是
$R[x_1, \ldots, x_n]$ 上的 Cohen--Macaulay 模。
\end{lemma}

\begin{proof}
对变量个数归纳，只须证明 $R[x]$ 上的
$M[x] = M \otimes_R R[x]$ 的情形。设
$\mathfrak m \subset R[x]$ 是极大理想，并令
$\mathfrak p = R \cap \mathfrak m$。设
$f_1, \ldots, f_d$ 是 $R_{\mathfrak p}$ 的极大理想中的一个
$M_\mathfrak p$-正则序列，其长度为
$d = \dim(\text{Supp}(M_{\mathfrak p}))$。注意，因为 $R[x]$
在 $R$ 上平坦，局部化 $R[x]_{\mathfrak m}$ 在
$R_{\mathfrak p}$ 上平坦。因此，由引理
\ref{algebra-lemma-flat-increases-depth}，序列 $f_1, \ldots, f_d$
是 $R[x]_{\mathfrak m}$ 中一个长度为 $d$ 的
$M[x]_{\mathfrak m}$-正则序列。商
$$
Q = M[x]_{\mathfrak m}/(f_1, \ldots, f_d)M[x]_{\mathfrak m} =
M_{\mathfrak p}/(f_1, \ldots, f_d)M_{\mathfrak p}
\otimes_{R_\mathfrak p} R[x]_{\mathfrak m}
$$
的支撑等于位于 $\mathfrak p$ 上方的素理想，因为
$R_\mathfrak p \to R[x]_\mathfrak m$ 平坦，而
$M_{\mathfrak p}/(f_1, \ldots, f_d)M_{\mathfrak p}$
的支撑等于 $\{\mathfrak p\}$（细节略；提示：由引理
\ref{algebra-lemma-annihilator-flat-base-change} 和
\ref{algebra-lemma-support-closed} 得出）。所以其维数为 $1$。
为完成证明，只须找到 $f \in \mathfrak m$，使它是 $Q$ 上的
非零因子。由于 $\mathfrak m$ 是极大理想，域扩张
$\kappa(\mathfrak m)/\kappa(\mathfrak p)$ 是有限的
（定理 \ref{algebra-theorem-nullstellensatz}）。因此，可以找到
$f \in \mathfrak m$，使它被视为 $x$ 的多项式时，首项系数
不属于 $\mathfrak p$。这样的 $f$ 作用在
$$
M_{\mathfrak p}/(f_1, \ldots, f_d)M_{\mathfrak p} \otimes_R R[x] =
\bigoplus\nolimits_{n \geq 0}
M_{\mathfrak p}/(f_1, \ldots, f_d)M_{\mathfrak p} \cdot x^n
$$
上是非零因子，因而作用在 $Q$ 上也是非零因子。
\end{proof}

\section{Cohen--Macaulay 环}
\label{algebra-section-CM-ring}

\noindent
本节的大多数结果都是第 \ref{algebra-section-CM} 节中结果的特殊情形。

\begin{definition}
\label{algebra-definition-local-ring-CM}
若 Noether 局部环 $R$ 作为自身上的模是 Cohen--Macaulay 模，
则称它为{\it Cohen--Macaulay 环}。
\end{definition}

\noindent
注意，这等价于要求极大理想中存在一个 $R$-正则序列
$x_1, \ldots, x_d$，使得 $R/(x_1, \ldots, x_d)$ 的维数为
$0$。通常只说“正则序列”，而不说“$R$-正则序列”。

\begin{lemma}
\label{algebra-lemma-reformulate-CM}
\begin{slogan}
Cohen--Macaulay 局部环中的正则序列，可以由它截出的对象是否具有
正确的维数来刻画。
\end{slogan}
设 $R$ 是极大理想为 $\mathfrak m $ 的 Noether
Cohen--Macaulay 局部环。设 $x_1, \ldots, x_c \in \mathfrak m$
为若干元素。则
$$
x_1, \ldots, x_c
\text{ is a regular sequence }
\Leftrightarrow
\dim(R/(x_1, \ldots, x_c)) = \dim(R) - c
$$
若这些条件成立，则 $x_1, \ldots, x_c$ 可以延拓为一个长度为
$\dim(R)$ 的正则序列，并且每个商
$R/(x_1, \ldots, x_i)$ 都是维数为 $\dim(R) - i$ 的
Cohen--Macaulay 环。
\end{lemma}

\begin{proof}
这是命题 \ref{algebra-proposition-CM-module} 的特殊情形。
\end{proof}

\begin{lemma}
\label{algebra-lemma-maximal-chain-CM}
设 $R$ 是 Noether 局部环。假设 $R$ 是维数为 $d$ 的
Cohen--Macaulay 环。任何极大
% P01-E0510: frozen source says a maximal chain of ideals; the dimension theorem requires prime ideals.
理想链 $\mathfrak p_0 \subset
\mathfrak p_1 \subset \ldots \subset \mathfrak p_n$
的长度都为 $n = d$。
\end{lemma}

\begin{proof}
这是引理 \ref{algebra-lemma-maximal-chain-maximal-CM} 的特殊情形。
\end{proof}

\begin{lemma}
\label{algebra-lemma-CM-dim-formula}
假设 $R$ 是维数为 $d$ 的 Noether Cohen--Macaulay 局部环。
对任何素理想 $\mathfrak p \subset R$，有
$$
\dim(R) = \dim(R_{\mathfrak p}) + \dim(R/\mathfrak p).
$$
\end{lemma}

\begin{proof}
这立即由引理 \ref{algebra-lemma-maximal-chain-CM} 得出。
（这也是引理 \ref{algebra-lemma-dim-formula-maximal-CM} 的特殊情形。）
\end{proof}

\begin{lemma}
\label{algebra-lemma-localize-CM}
假设 $R$ 是 Cohen--Macaulay 局部环。对任何素理想
$\mathfrak p \subset R$，环 $R_{\mathfrak p}$ 也是
Cohen--Macaulay 环。
\end{lemma}

\begin{proof}
这是引理 \ref{algebra-lemma-localize-CM-module} 的特殊情形。
\end{proof}

\begin{definition}
\label{algebra-definition-ring-CM}
若 Noether 环 $R$ 的所有局部环都是 Cohen--Macaulay 环，
则称它为{\it Cohen--Macaulay 环}。
\end{definition}

\begin{lemma}
\label{algebra-lemma-CM-polynomial-algebra}
假设 $R$ 是 Noether Cohen--Macaulay 环。则 $R$ 上的任意
多项式代数都是 Cohen--Macaulay 环。
\end{lemma}

\begin{proof}
这是引理 \ref{algebra-lemma-maximal-CM-polynomial-algebra} 的特殊情形。
\end{proof}

\begin{lemma}
\label{algebra-lemma-dimension-shift}
设 $R$ 是维数为 $d$ 的 Noether Cohen--Macaulay 局部环。
设 $0 \to K \to R^{\oplus n} \to M \to 0$ 是 $R$-模的正合列。
则或者 $M = 0$，或者 $\text{depth}(K) > \text{depth}(M)$，
或者 $\text{depth}(K) = \text{depth}(M) = d$。
\end{lemma}

\begin{proof}
这是引理 \ref{algebra-lemma-depth-in-ses} 的特殊情形。
\end{proof}

\begin{lemma}
\label{algebra-lemma-mcm-resolution}
设 $R$ 是维数为 $d$ 的 Noether Cohen--Macaulay 局部环。
设 $M$ 是深度为 $e$ 的有限 $R$-模。存在正合复形
% P01-E0511: for e=d the displayed index d-e-1 is -1; the source needs a separate maximal-CM case.
$$
0 \to K \to F_{d-e-1} \to \ldots \to F_0 \to M \to 0
$$
其中每个 $F_i$ 都是有限自由模，而 $K$ 是极大
Cohen--Macaulay 模。
\end{lemma}

\begin{proof}
这立即由定义和引理 \ref{algebra-lemma-dimension-shift} 得出。
\end{proof}

\begin{lemma}
\label{algebra-lemma-find-sequence-image-regular}
设 $\varphi : A \to B$ 是局部环映射。假设 $B$ 是 Noether
Cohen--Macaulay 环，并且
$\mathfrak m_B = \sqrt{\varphi(\mathfrak m_A) B}$。
则 $A$ 中存在元素序列 $f_1, \ldots, f_{\dim(B)}$，使得
$\varphi(f_1), \ldots, \varphi(f_{\dim(B)})$ 是 $B$ 中的
正则序列。
\end{lemma}

\begin{proof}
对 $\dim(B)$ 归纳，只须证明：若 $\dim(B) \geq 1$，则可以找到
$A$ 的一个元素 $f$，使它在 $B$ 中的像为非零因子。由引理
\ref{algebra-lemma-reformulate-CM}，只须找到 $f \in A$，使其在 $B$
中的像不属于 $B$ 的有限多个极小素理想
$\mathfrak q_1, \ldots, \mathfrak q_r$ 中的任何一个。
由假设 $\mathfrak m_B = \sqrt{\varphi(\mathfrak m_A) B}$，
可知 $\mathfrak m_A \not \subset \varphi^{-1}(\mathfrak q_i)$。
因此，由引理 \ref{algebra-lemma-silly} 可以找到 $f$。
\end{proof}

\section{链状环}
\label{algebra-section-catenary}

\noindent
与拓扑学第 \ref{topology-section-catenary-spaces} 节比较。

\begin{definition}
\label{algebra-definition-catenary}
若对任意一对素理想
$\mathfrak p \subset \mathfrak q$，都存在一个整数作为所有有限素理想链
$\mathfrak p = \mathfrak p_0 \subset \mathfrak p_1 \subset \ldots \subset
\mathfrak p_e = \mathfrak q$ 的长度上界，并且所有这样的极大链都具有相同
长度，则称环 $R$ 是{\it 链状的}。
\end{definition}

\begin{lemma}
\label{algebra-lemma-catenary}
环 $R$ 是链状的，当且仅当拓扑空间 $\Spec(R)$ 是链状的（参见
拓扑学定义 \ref{topology-definition-catenary}）。
\end{lemma}

\begin{proof}
这立即由定义以及引理 \ref{algebra-lemma-irreducible} 中对不可约闭子集的刻画得出。
\end{proof}

\noindent
一般而言，即使 $R$ 是链状的，有限生成 $R$-代数也未必是链状的。
因此我们作如下定义。

\begin{definition}
\label{algebra-definition-universally-catenary}
若 Noether 环 $R$ 的每个有限型 $R$-代数都是链状的，则称该环
是{\it 普遍链状的}。
\end{definition}

\noindent
我们将范围限制在 Noether 环上，因为尚不清楚对非 Noether 环而言，
这个定义是否恰当。由引理 \ref{algebra-lemma-quotient-catenary}，为了检验
Noether 环 $R$ 是否普遍链状，只须检验每个多项式代数
$R[x_1, \ldots, x_n]$ 是否链状。

\begin{lemma}
\label{algebra-lemma-localization-catenary}
链状环的任意局部化都是链状环。Noether 普遍链状环的任意局部化
都是普遍链状环。
\end{lemma}

\begin{proof}
设 $A$ 是环，$S \subset A$ 是乘法子集。引理
\ref{algebra-lemma-spec-localization} 中对 $\Spec(S^{-1}A)$ 的描述表明，
若 $A$ 是链状的，则 $S^{-1}A$ 也是链状的。若
$S^{-1}A \to C$ 是有限型的，则存在有限型环映射 $A \to B$，使得
$C = S^{-1}B$。因此，若 $A$ 是 Noether 且普遍链状的，则 $B$
是链状的，从而 $C$ 也是链状的。这就证明了该引理。
\end{proof}

\begin{lemma}
\label{algebra-lemma-universally-catenary}
设 $A$ 是 Noether 普遍链状环。任意在 $A$ 上本质有限型的
$A$-代数都是普遍链状环。
\end{lemma}

\begin{proof}
若 $B$ 是有限型 $A$-代数，则由引理 \ref{algebra-lemma-Noetherian-permanence}，
$B$ 是 Noether 环。任意有限型 $B$-代数都是有限型 $A$-代数，因而由
$A$ 普遍链状这一假设，它是链状的。因此 $B$ 是普遍链状的。由引理
\ref{algebra-lemma-localization-catenary}，$B$ 的任意局部化都是普遍链状的，
这就完成了证明。
\end{proof}

\begin{lemma}
\label{algebra-lemma-catenary-check-local}
设 $R$ 是环。以下条件等价：
\begin{enumerate}
\item $R$ 是链状的；
\item 对每个素理想 $\mathfrak p$，$R_\mathfrak p$ 都是链状的；
\item 对每个极大理想 $\mathfrak m$，$R_\mathfrak m$ 都是链状的。
\end{enumerate}
假设 $R$ 是 Noether 环。以下条件等价：
\begin{enumerate}
\item $R$ 是普遍链状的；
\item 对每个素理想 $\mathfrak p$，$R_\mathfrak p$ 都是普遍链状的；
\item 对每个极大理想 $\mathfrak m$，$R_\mathfrak m$ 都是普遍链状的。
\end{enumerate}
\end{lemma}

\begin{proof}
在两种情形中，(1) $\Rightarrow$ (2) 都由引理
\ref{algebra-lemma-localization-catenary} 得出，而 (2) $\Rightarrow$ (3)
都是立即的。假设对 $R$ 的每个极大理想 $\mathfrak m$，
$R_\mathfrak m$ 都是链状的。若 $\mathfrak p \subset \mathfrak q$
是 $R$ 的素理想，则选择一个满足 $\mathfrak q \subset \mathfrak m$
的极大理想。$\mathfrak p$ 与 $\mathfrak q$ 之间的素理想链，和
% P01-E0512: the source says “Chains of primes ideals”; Chinese uses the unambiguous intended “prime ideals”.
$\mathfrak pR_\mathfrak m$ 与 $\mathfrak qR_\mathfrak m$ 之间的素理想链
一一对应，故 $R$ 是链状的。现在假设 $R$ 是 Noether 环，并且对
$R$ 的每个极大理想 $\mathfrak m$，$R_\mathfrak m$ 都是普遍链状的。
设 $R \to S$ 是有限型环映射。设 $\mathfrak q$ 是 $S$ 的素理想，
它位于素理想 $\mathfrak p \subset R$ 之上。在 $R$ 中选择一个满足
$\mathfrak p \subset \mathfrak m$ 的极大理想。于是
$R_\mathfrak p$ 是 $R_\mathfrak m$ 的局部化，因而由引理
\ref{algebra-lemma-localization-catenary}，它是普遍链状的。于是
$S_\mathfrak p$ 作为 $R_\mathfrak p$ 上的有限型环是链状的，
而 $S_\mathfrak q$ 作为它的局部化也是链状的。因此，由上面处理的
第一种情形，$S$ 是链状的。
\end{proof}

\begin{lemma}
\label{algebra-lemma-quotient-catenary}
链状环的任意商环都是链状环。Noether 普遍链状环的任意商环都是
普遍链状环。
\end{lemma}

\begin{proof}
设 $A$ 是环，$I \subset A$ 是理想。引理 \ref{algebra-lemma-spec-closed}
中对 $\Spec(A/I)$ 的描述表明，若 $A$ 是链状的，则 $A/I$
也是链状的。第二个断言是引理 \ref{algebra-lemma-universally-catenary}
的特殊情形。
\end{proof}

\begin{lemma}
\label{algebra-lemma-catenary-check-irreducible}
设 $R$ 是 Noether 环。
\begin{enumerate}
\item $R$ 是链状的，当且仅当对每个极小素理想 $\mathfrak p$，
$R/\mathfrak p$ 都是链状的。
\item $R$ 是普遍链状的，当且仅当对每个极小素理想 $\mathfrak p$，
$R/\mathfrak p$ 都是普遍链状的。
\end{enumerate}
\end{lemma}

\begin{proof}
若 $\mathfrak a \subset \mathfrak b$ 是 $R$ 的素理想包含关系，
则可以找到一个满足 $\mathfrak p \subset \mathfrak a$ 的极小素理想，
第一个断言于是显然。第二个断言的证明从略。
\end{proof}

\begin{lemma}
\label{algebra-lemma-CM-ring-catenary}
Noether Cohen--Macaulay 环是普遍链状的。更一般地，若 $R$ 是
Noether 环，而 $M$ 是满足 $\text{Supp}(M) = \Spec(R)$ 的
Cohen--Macaulay $R$-模，则 $R$ 是普遍链状的。
\end{lemma}

\begin{proof}
由引理 \ref{algebra-lemma-CM-polynomial-algebra}，$R$ 上的多项式代数是
Cohen--Macaulay 环，所以只须证明 Cohen--Macaulay 环是链状的。
设 $R$ 是 Cohen--Macaulay 环，$\mathfrak p \subset \mathfrak q$
是 $R$ 的素理想。由定义，$R_{\mathfrak q}$ 和 $R_{\mathfrak p}$
都是 Cohen--Macaulay 环。取一条极大素理想链
$\mathfrak p = \mathfrak p_0 \subset \mathfrak p_1 \subset
\ldots \subset \mathfrak p_n = \mathfrak q$。再选择一条极大素理想链
$\mathfrak q_0 \subset \mathfrak q_1 \subset \ldots \subset
\mathfrak q_m = \mathfrak p$。由引理 \ref{algebra-lemma-maximal-chain-CM}，
有 $n + m = \dim(R_{\mathfrak q})$。由同一引理，还有
$m = \dim(R_{\mathfrak p})$。因此
$n = \dim(R_{\mathfrak q}) - \dim(R_{\mathfrak p})$ 与选择无关。

\medskip\noindent
为了证明更一般的断言，完全照上面论证，但改用引理
\ref{algebra-lemma-maximal-CM-polynomial-algebra} 和
\ref{algebra-lemma-maximal-chain-maximal-CM}。
\end{proof}

\begin{lemma}
\label{algebra-lemma-catenary-Noetherian-local}
设 $(A, \mathfrak m)$ 是 Noether 局部环。以下条件等价：
\begin{enumerate}
\item $A$ 是链状的；
\item $\mathfrak p \mapsto \dim(A/\mathfrak p)$ 是 $\Spec(A)$
上的维数函数。
\end{enumerate}
\end{lemma}

\begin{proof}
若 $A$ 是链状的，则由拓扑学引理
\ref{topology-lemma-locally-dimension-function}（以及引理
\ref{algebra-lemma-catenary}），$\Spec(A)$ 有一个维数函数 $\delta$。
可以假设 $\delta(\mathfrak m) = 0$。于是由拓扑学引理
\ref{topology-lemma-dimension-function-catenary}，可知
$$
\delta(\mathfrak p) = \text{codim}(V(\mathfrak m), V(\mathfrak p)) =
\dim(A/\mathfrak p)
$$
由此可见 (1) 推出 (2)。反向蕴含也由拓扑学引理
\ref{topology-lemma-dimension-function-catenary} 得出。
\end{proof}

\section{正则局部环}
\label{algebra-section-regular}

\noindent
正则局部环定义于定义 \ref{algebra-definition-regular-local}。要证明正则局部环的
所有素局部化都是正则的并不容易。事实上，迄今发展的大量材料都是为
证明这一事实而准备的。参见命题
\ref{algebra-proposition-finite-gl-dim-regular}，并沿其中的引用向前追溯。

\begin{lemma}
\label{algebra-lemma-regular-graded}
设 $(R, \mathfrak m, \kappa)$ 是维数为 $d$ 的正则局部环。
分次环 $\bigoplus \mathfrak m^n / \mathfrak m^{n + 1}$ 同构于
分次多项式代数 $\kappa[X_1, \ldots, X_d]$。
\end{lemma}

\begin{proof}
设 $x_1, \ldots, x_d$ 是极大理想 $\mathfrak m$ 的极小生成集，
参见定义 \ref{algebra-definition-regular-local}。存在满射
$\kappa[X_1, \ldots, X_d]
\to \bigoplus \mathfrak m^n/\mathfrak m^{n + 1}$，它将 $X_i$
映到 $x_i$ 在 $\mathfrak m/\mathfrak m^2$ 中的类。由命题
\ref{algebra-proposition-dimension} 有 $d(R) = d$，故数值多项式
$n \mapsto \dim_\kappa \mathfrak m^n/\mathfrak m^{n + 1}$
的次数为 $d - 1$。由引理 \ref{algebra-lemma-quotient-smaller-d} 可知，
该满射 $\kappa[X_1, \ldots, X_d]
\to \bigoplus \mathfrak m^n/\mathfrak m^{n + 1}$ 是同构。
\end{proof}

\begin{lemma}
\label{algebra-lemma-regular-domain}
任意正则局部环都是整环。
\end{lemma}

\begin{proof}
我们将使用引理 \ref{algebra-lemma-intersect-powers-ideal-module-zero} 给出的
$\bigcap \mathfrak m^n = 0$。设 $f, g \in R$ 且 $fg = 0$。
% P01-E0513: the source implicitly needs f and g nonzero before choosing maximal m-adic orders.
只须考虑二者都非零的情形。假设 $f \in \mathfrak m^a$ 且
$g \in \mathfrak m^b$，其中 $a, b$ 均为最大的。由于
$fg = 0 \in \mathfrak m^{a + b + 1}$，由引理
\ref{algebra-lemma-regular-graded} 的结论可知，或者
$f \in \mathfrak m^{a + 1}$，或者 $g \in \mathfrak m^{b + 1}$。
矛盾。
\end{proof}

\begin{lemma}
\label{algebra-lemma-regular-ring-CM}
设 $R$ 是正则局部环，$x_1, \ldots, x_d$ 是极大理想
$\mathfrak m$ 的极小生成集。则 $x_1, \ldots, x_d$ 是正则序列，
并且每个 $R/(x_1, \ldots, x_c)$ 都是维数为 $d - c$ 的正则局部环。
特别地，$R$ 是 Cohen--Macaulay 环。
\end{lemma}

\begin{proof}
注意，由引理 \ref{algebra-lemma-one-equation}，$R/x_1R$ 是维数
$\geq d - 1$ 的 Noether 局部环，而 $x_2, \ldots, x_d$ 生成其
极大理想。因此，按定义它是正则局部环。由引理
\ref{algebra-lemma-regular-domain}，$R$ 是整环，故 $x_1$ 是非零因子。
\end{proof}

\begin{lemma}
\label{algebra-lemma-regular-quotient-regular}
设 $R$ 是正则局部环。设 $I \subset R$ 是理想，并且 $R/I$
也是正则局部环。则存在 $R$ 的极大理想 $\mathfrak m$ 的极小生成集
$x_1, \ldots, x_d$，使得对某个 $0 \leq c \leq d$，有
$I = (x_1, \ldots, x_c)$。
\end{lemma}

\begin{proof}
记 $\dim(R) = d$ 且 $\dim(R/I) = d - c$。用
$\overline{\mathfrak m} = \mathfrak m/I$ 表示 $R/I$ 的极大理想。
设 $\kappa = R/\mathfrak m$。由正则局部环的定义，有
% P01-E0514: the source omits the residue-field subscript on the third vector-space dimension.
$$
\dim_\kappa((I + \mathfrak m^2)/\mathfrak m^2) =
\dim_\kappa(\mathfrak m/\mathfrak m^2)
- \dim(\overline{\mathfrak m}/\overline{\mathfrak m}^2) = d - (d - c) = c
$$
因此可以选择 $x_1, \ldots, x_c \in I$，使它们在
$\mathfrak m/\mathfrak m^2$ 中的像线性无关，再补上
$x_{c + 1}, \ldots, x_d$，得到 $\mathfrak m$ 的极小生成系。
诱导映射 $R/(x_1, \ldots, x_c) \to R/I$ 是相同维数的正则局部环
之间的满射（引理 \ref{algebra-lemma-regular-ring-CM}）。因此其核为零，
也就是说，$I = (x_1, \ldots, x_c)$。具体地，若非如此，则由引理
\ref{algebra-lemma-regular-domain} 和 \ref{algebra-lemma-one-equation}，将有
$\dim(R/I) < \dim(R/(x_1, \ldots, x_c))$。
\end{proof}

\begin{lemma}
\label{algebra-lemma-free-mod-x}
设 $R$ 是 Noether 局部环，$x \in \mathfrak m$。设 $M$ 是有限
$R$-模，使得 $x$ 在 $M$ 上是非零因子，并且 $M/xM$ 作为
$R/xR$-模是自由的。则 $M$ 作为 $R$-模是自由的。
\end{lemma}

\begin{proof}
设 $m_1, \ldots, m_r$ 是 $M$ 的元素，它们的像构成 $M/xM$ 的一组
$R/xR$-基。由 Nakayama 引理 \ref{algebra-lemma-NAK}，
$m_1, \ldots, m_r$ 生成 $M$。若 $\sum a_i m_i = 0$ 是一个关系，
则对所有 $i$ 都有 $a_i \in xR$。因此，对某些 $b_i \in R$，
有 $a_i = b_i x$。于是 $R^r \to M$ 的核 $K$ 满足 $xK = K$，
所以由 Nakayama 引理，它为零。
\end{proof}

\begin{lemma}
\label{algebra-lemma-regular-mcm-free}
设 $R$ 是正则局部环。任意极大 Cohen--Macaulay $R$-模都是自由的。
\end{lemma}

\begin{proof}
设 $M$ 是极大 Cohen--Macaulay $R$-模。设 $x \in \mathfrak m$
是生成 $\mathfrak m$ 的一个正则序列的一部分。由命题
\ref{algebra-proposition-CM-module}，$x$ 在 $M$ 上是非零因子，而
$M/xM$ 是 $R/xR$ 上的极大 Cohen--Macaulay 模。对
$\dim(R)$ 归纳可知 $M/xM$ 是自由的。由引理
\ref{algebra-lemma-free-mod-x} 即得结论。
\end{proof}

\begin{lemma}
\label{algebra-lemma-regular-mod-x}
假设 $R$ 是 Noether 局部环。设 $x \in \mathfrak m$ 是非零因子，
并且 $R/xR$ 是正则局部环。则 $R$ 是正则局部环。更一般地，若
$x_1, \ldots, x_r$ 是 $R$ 中的正则序列，并且
$R/(x_1, \ldots, x_r)$ 是正则局部环，则 $R$ 是正则局部环。
\end{lemma}

\begin{proof}
这是因为把 $x$ 与 $R/xR$ 的极大理想的一组极小生成元的提升放在一起，
就会得到 $\mathfrak m$ 的 $\dim(R)$ 个生成元。使用引理
\ref{algebra-lemma-one-equation}。最后一个断言由第一个断言和归纳法得出。
\end{proof}

\begin{lemma}
\label{algebra-lemma-colimit-regular}
设 $(R_i, \varphi_{ii'})$ 是局部环的有向系，其过渡映射都是局部环映射。
若每个 $R_i$ 都是正则局部环，并且 $R = \colim R_i$ 是 Noether 环，
则 $R$ 是正则局部环。
\end{lemma}

\begin{proof}
设 $\mathfrak m \subset R$ 是极大理想；它是各极大理想
$\mathfrak m_i \subset R_i$ 的余极限。我们对
% P01-E0516: the source omits the residue-field subscript and grouping in this tangent-space dimension.
$d = \dim \mathfrak m/\mathfrak m^2$ 归纳证明该引理。若 $d = 0$，
则 $R = R/\mathfrak m$ 是域，且 $R$ 是正则局部环。若 $d > 0$，
选择 $x \in \mathfrak m$ 且 $x \not \in \mathfrak m^2$。对某个 $i$，
可以找到映到 $x$ 的 $x_i \in \mathfrak m_i$。注意
$R/xR = \colim_{i' \geq i} R_{i'}/x_iR_{i'}$ 是 Noether 局部环。
由引理 \ref{algebra-lemma-regular-ring-CM}，可知
$R_{i'}/x_iR_{i'}$ 是正则局部环。因此由归纳可知 $R/xR$
是正则局部环。由于每个 $R_i$ 都是整环
% P01-E0515: the cited lemma proves the graded-ring statement; regular-domain is the direct domain lemma.
（引理 \ref{algebra-lemma-regular-graded}），可知 $R$ 是整环。因此 $x$
是非零因子，由引理 \ref{algebra-lemma-regular-mod-x} 可知 $R$ 是正则局部环。
\end{proof}

\section{环的满态射}
\label{algebra-section-epimorphism}

\noindent
在任意范畴中都有{\it 满态射}的概念。本节部分材料取自
\cite{Autour} 和 \cite{Mazet}。

\begin{lemma}
\label{algebra-lemma-epimorphism}
设 $R \to S$ 是环映射。以下条件等价：
\begin{enumerate}
\item $R \to S$ 是满态射；
\item 两个环映射 $S \to S \otimes_R S$ 相等；
\item 任一个环映射 $S \to S \otimes_R S$ 都是同构；
\item 环映射 $S \otimes_R S \to S$ 是同构。
\end{enumerate}
\end{lemma}

\begin{proof}
从略。
\end{proof}

\begin{lemma}
\label{algebra-lemma-composition-epimorphism}
两个环满态射的复合仍是满态射。
\end{lemma}

\begin{proof}
从略。提示：这在任意范畴中都成立。
\end{proof}

\begin{lemma}
\label{algebra-lemma-base-change-epimorphism}
若 $R \to S$ 是环满态射，而 $R \to R'$ 是任意环映射，则
$R' \to R' \otimes_R S$ 是满态射。
\end{lemma}

\begin{proof}
从略。提示：这在任意具有推出的范畴中都成立。
\end{proof}

\begin{lemma}
\label{algebra-lemma-permanence-epimorphism}
若 $A \to B \to C$ 是环映射，而 $A \to C$ 是满态射，则
$B \to C$ 也是满态射。
\end{lemma}

\begin{proof}
从略。提示：这在任意范畴中都成立。
\end{proof}

\noindent
特别地，这意味着：若 $R \to S$ 是满态射，其像为
$\overline{R} \subset S$，则 $\overline{R} \to S$ 是满态射。
% P01-E0517: the source has a spurious comma in “in particular, that”.
因此在证明关于满态射的结果时，常可假设映射是单射。下面的引理特别
蕴含每个局部化都是满态射。

\begin{lemma}
\label{algebra-lemma-epimorphism-local}
设 $R \to S$ 是环映射。以下条件等价：
\begin{enumerate}
\item $R \to S$ 是满态射；
\item 对 $R$ 的每个素理想 $\mathfrak p$，
$R_{\mathfrak p} \to S_{\mathfrak p}$ 都是满态射。
\end{enumerate}
\end{lemma}

\begin{proof}
由于 $S_{\mathfrak p} = R_{\mathfrak p} \otimes_R S$（参见引理
\ref{algebra-lemma-tensor-localization}），由引理
\ref{algebra-lemma-base-change-epimorphism} 可知 (1) 推出 (2)。反过来，
假设 (2) 成立。设 $a, b : S \to A$ 是从 $S$ 到环 $A$ 的两个环映射，
并且它们与 $R \to S$ 复合后相等。由假设可知，对 $R$ 的每个素理想
% P01-E0518: the source phrase “two ring maps ... equalizing the map” is category-theoretically malformed.
$\mathfrak p$，诱导映射
$a_{\mathfrak p}, b_{\mathfrak p} : S_{\mathfrak p} \to A_{\mathfrak p}$
相同。由于 $A \subset \prod_{\mathfrak p} A_{\mathfrak p}$，故
$a = b$，参见引理 \ref{algebra-lemma-characterize-zero-local}。
\end{proof}

\begin{lemma}
\label{algebra-lemma-finite-epimorphism-surjective}
\begin{slogan}
一个环映射是满射，当且仅当它是有限满态射。
\end{slogan}
设 $R \to S$ 是环映射。以下条件等价：
\begin{enumerate}
\item $R \to S$ 是满态射且有限；
\item $R \to S$ 是满射。
\end{enumerate}
\end{lemma}

\begin{proof}
（这个引理似乎已在文献中被多次重新证明，并且有许多不同的证明。）
满射环映射显然是满态射。假设 $R \to S$ 是有限环映射，并且
$S \otimes_R S \to S$ 是同构。我们的目标是证明 $R \to S$ 是满射。
假设 $S/R$ 非零。正合列 $R \to S \to S/R \to 0$ 导出正合列
$$
R \otimes_R S \to S \otimes_R S \to S/R \otimes_R S \to 0.
$$
我们的假设蕴含第一个箭头是同构，故 $S/R \otimes_R S = 0$。
从而还有 $S/R \otimes_R S/R = 0$。由引理
\ref{algebra-lemma-trivial-filter-finite-module}，对某个真理想 $I \subset R$，
存在 $R$-模满射 $S/R \to R/I$。于是存在满射
$S/R \otimes_R S/R \to R/I \otimes_R R/I = R/I \not = 0$，
矛盾。
\end{proof}

\begin{lemma}
\label{algebra-lemma-faithfully-flat-epimorphism}
忠实平坦满态射是同构。
\end{lemma}

\begin{proof}
这由引理 \ref{algebra-lemma-epimorphism} 的第 (3) 部分立即得出，因为映射
$S \to S \otimes_R S$ 是映射 $R \to S$ 与 $S$ 作张量积所得的映射。
\end{proof}

\begin{lemma}
\label{algebra-lemma-epimorphism-over-field}
若 $k \to S$ 是满态射，且 $k$ 是域，则 $S = k$ 或 $S = 0$。
\end{lemma}

\begin{proof}
这由引理 \ref{algebra-lemma-faithfully-flat-epimorphism} 的结论立即得出
（因为 $k$ 上任意非零代数都是忠实平坦的）；也可以直接论证：
% P01-E0519: the source switches to undefined R; the intended map is S -> S tensor_k S.
$R \to R \otimes_k R$ 不可能是满射，除非 $\dim_k(R) \leq 1$。
\end{proof}

\begin{lemma}
\label{algebra-lemma-epimorphism-injective-spec}
设 $R \to S$ 是环满态射。则
\begin{enumerate}
\item $\Spec(S) \to \Spec(R)$ 是单射；
\item 对位于 $\mathfrak p \subset R$ 之上的 $\mathfrak q \subset S$，
有 $\kappa(\mathfrak p) = \kappa(\mathfrak q)$。
\end{enumerate}
\end{lemma}

\begin{proof}
设 $\mathfrak p$ 是 $R$ 的素理想。该映射的纤维是纤维环
$S \otimes_R \kappa(\mathfrak p)$ 的谱。由引理
\ref{algebra-lemma-base-change-epimorphism}，映射
$\kappa(\mathfrak p) \to S \otimes_R \kappa(\mathfrak p)$ 是满态射，
所以由引理 \ref{algebra-lemma-epimorphism-over-field}，或者
$S \otimes_R \kappa(\mathfrak p) = 0$，或者
$S \otimes_R \kappa(\mathfrak p) = \kappa(\mathfrak p)$。
这证明了 (1) 和 (2)。
\end{proof}

\begin{lemma}
\label{algebra-lemma-relations}
设 $R$ 是环。设 $M$、$N$ 是 $R$-模。设
$\{x_i\}_{i \in I}$ 是 $M$ 的一组生成元，
$\{y_j\}_{j \in J}$ 是 $N$ 的一组生成元。设
$\{m_j\}_{j \in J}$ 是 $M$ 中的一族元素，其中除有限多个 $j$ 外，
均有 $m_j = 0$。则
$$
\sum\nolimits_{j \in J} m_j \otimes y_j = 0 \text{ 于 } M \otimes_R N
$$
等价于下述条件：存在 $a_{i, j} \in R$，并且除有限多对
$(i, j)$ 外均有 $a_{i, j} = 0$，使得
\begin{align*}
m_j & = \sum\nolimits_{i \in I} a_{i, j} x_i \quad\text{对所有 } j \in J, \\
0 & = \sum\nolimits_{j \in J} a_{i, j} y_j \quad\text{对所有 } i \in I.
\end{align*}
\end{lemma}

\begin{proof}
充分性是立即的。假设 $\sum_{j \in J} m_j \otimes y_j = 0$。
考虑短正合列
$$
0 \to K \to \bigoplus\nolimits_{j \in J} R \to N \to 0
$$
其中 $\bigoplus\nolimits_{j \in J} R$ 的第 $j$ 个基向量映到 $y_j$。
将它与 $M$ 作张量积，得到正合列
$$
K \otimes_R M \to \bigoplus\nolimits_{j \in J} M \to N \otimes_R M \to 0.
$$
假设蕴含存在元素 $k_i \in K$，使得 $\sum k_i \otimes x_i$ 映到
中间项的元素 $(m_j)_{j \in J}$。写成
% P01-E0520: the source says “Writing ... and we obtain”, leaving the sentence grammatically broken.
$k_i = (a_{i, j})_{j \in J}$，即得所需结论。
\end{proof}

\begin{lemma}
\label{algebra-lemma-kernel-difference-projections}
设 $\varphi : R \to S$ 是环映射，$g \in S$。以下条件等价：
\begin{enumerate}
\item 在 $S \otimes_R S$ 中有 $g \otimes 1 = 1 \otimes g$；
\item 存在 $n \geq 0$、元素 $y_i, z_j \in S$ 以及
$x_{i, j} \in R$（其中 $1 \leq i, j \leq n$），使得
\begin{enumerate}
\item $g = \sum_{i, j \leq n} x_{i, j} y_i z_j$；
\item 对每个 $j$，有 $\sum x_{i, j}y_i \in \varphi(R)$；
\item 对每个 $i$，有 $\sum x_{i, j}z_j \in \varphi(R)$。
\end{enumerate}
\end{enumerate}
\end{lemma}

\begin{proof}
(2) 推出 (1) 是显然的。反过来，假设
$g \otimes 1 = 1 \otimes g$。选择 $S$ 作为 $R$-模的一组生成元
$\{s_i\}_{i \in I}$，使得 $0, 1 \in I$，且 $s_0 = 1$、$s_1 = g$。
将引理 \ref{algebra-lemma-relations} 应用于关系
$g \otimes s_0 + (-1) \otimes s_1 = 0$。可知存在
$a_{i, j} \in R$，使得 $g = \sum_i a_{i, 0} s_i$、
$-1 = \sum_i a_{i, 1} s_i$；当 $j \not = 0, 1$ 时，有
$0 = \sum_i a_{i, j} s_i$；并且对所有 $i$，有
$\sum_j a_{i, j}s_j = 0$。于是
$$
\sum\nolimits_{i, j \not = 0} a_{i, j} s_i s_j = -g + a_{0, 0}
$$
且对每个 $j \not = 0$，有
% P01-E0521: the source writes R here although the lemma consistently uses the image varphi(R).
$\sum_{i \not = 0} a_{i, j}s_i \in R$。这证明了
$-g + a_{0, 0}$ 可写成 (2) 中的形式，从而 $g$ 也可写成 (2)
中的形式。细节从略。提示：证明 $S$ 中具有 (2) 所述表达式的元素
构成 $S$ 的一个 $R$-子代数。
\end{proof}

\begin{remark}
\label{algebra-remark-matrices-associated-to-elements-epicenter}
设 $R \to S$ 是环映射。有时把满足
$g \in S$ 且 $g \otimes 1 = 1 \otimes g$ 的元素集合称为 $S$ 的
{\it 满态中心（epicenter）}。它是一个 $R$-代数。由引理
\ref{algebra-lemma-kernel-difference-projections} 的构造，对满态中心中的每个
$g$，得到矩阵分解
$$
(g) = Y X Z
$$
其中 $X \in \text{Mat}(n \times n, R)$、
$Y \in \text{Mat}(1 \times n, S)$，且
$Z \in \text{Mat}(n \times 1, S)$。具体地，设 $x_{i, j}, y_i, z_j$
如该引理第 (2) 部分所述。令 $X = (x_{i, j})$，令 $y$ 为以 $y_i$
为分量的行向量，令 $z$ 为以 $z_j$ 为分量的列向量。
% P01-E0522: the source defines lowercase y,z but immediately uses uppercase Y,Z in the factorization.
在此记号下，引理 \ref{algebra-lemma-kernel-difference-projections} 的条件
(b) 和 (c) 恰好表示
$Y X \in \text{Mat}(1 \times n, R)$、
$X Z \in \text{Mat}(n \times 1, R)$。考虑矩阵三元组
$(X, YX, XZ)$ 十分方便。给定 $n \in \mathbf{N}$ 和三元组
$(P, U, V)$，若存在如上的矩阵分解，使得
$P = X$、$U = YX$ 且 $V = XZ$，则称 $(P, U, V)$ 是
{\it 与 $g$ 关联的 $n$-三元组}。
\end{remark}

\begin{lemma}
\label{algebra-lemma-epimorphism-cardinality}
设 $R \to S$ 是环满态射。则 $S$ 的基数至多为 $R$ 的基数。
用公式表示：$|S| \leq |R|$。
\end{lemma}

\begin{proof}
$R \to S$ 是满态射这一条件意味着每个 $g \in S$ 都满足
$g \otimes 1 = 1 \otimes g$，参见引理 \ref{algebra-lemma-epimorphism}。
我们将使用评注 \ref{algebra-remark-matrices-associated-to-elements-epicenter}
中引入的记号。假设 $g, g' \in S$，并且 $(P, U, V)$ 是同时与
$g$ 和 $g'$ 关联的 $n$-三元组。我们断言 $g = g'$。事实上，
把与 $g$ 的矩阵分解 $(g) = YXZ$ 写成
$(P, U, V) = (X, YX, XZ)$，把与 $g'$ 的矩阵分解
$(g') = Y'X'Z'$ 写成 $(P, U, V) = (X', Y'X', X'Z')$。
于是
$$
(g) = YXZ = UZ = Y'X'Z = Y'PZ = Y'XZ = Y'V = Y'X'Z' = (g')
$$
故 $g = g'$。这说明 $S$ 的基数被所有可能三元组的数目所控制，
后者的基数至多为 $\sup_{n \in \mathbf{N}} |R|^n$。若 $R$ 是无限的，
则它至多为 $|R|$，参见 \cite[Ch. I, 10.13]{Kunen}。

\medskip\noindent
若 $R$ 是有限环，则上述论证只证明 $S$ 至多可数。事实上，在这种情形
$R$ 是 Artin 环，并且映射 $R \to S$ 是满射。我们略去这一情形的证明。
\end{proof}

\begin{lemma}
\label{algebra-lemma-epimorphism-modules}
对环同态 $R \to S$，以下条件等价：
\begin{enumerate}
\item $R \to S$ 是环满态射；
\item 对任意 $S$-模 $N_1, N_2$，都有
$\Hom_S(N_1, N_2) = \Hom_R(N_1, N_2)$；
\item 限制函子 $\text{Mod}_S \to \text{Mod}_R$ 是全忠实的。
\end{enumerate}
\end{lemma}

\begin{proof}
按定义，(2) 与 (3) 等价。

\medskip\noindent
假设 (1) 成立，设 $N_1, N_2$ 是 $S$-模，并设
$\varphi : N_1 \to N_2$ 是 $R$-线性映射。对任意 $x \in N_1$，
考虑由规则 $g \otimes g' \mapsto g\varphi(g'x)$ 定义的映射
$S \otimes_R S \to N_2$。由于两个映射 $S \to S \otimes_R S$
都是同构（引理 \ref{algebra-lemma-epimorphism}），可得
$g \varphi(g'x) = gg'\varphi(x) = \varphi(gg' x)$。因此 $\varphi$
是 $S$-线性的。

\medskip\noindent
假设 (2) 成立。令 $N_1 = S \otimes_R S$，通过左 $S$-模作用把它
看作 $S$-模；令 $N_2 = S \otimes_R S$，赋予右 $S$-模作用。
由于作为 $R$-模有 $N_1 = N_2$，由 (2) 可知
$S \otimes_R S$ 上的两个 $S$-模结构相同。由引理
\ref{algebra-lemma-epimorphism}，这蕴含 (1)。
\end{proof}

\section{纯理想}
\label{algebra-section-pure-ideals}

\noindent
许多论文讨论了本节材料，例如参见 \cite{Lazard}、\cite{Bkouche}
和 \cite{DeMarco}。

\begin{definition}
\label{algebra-definition-pure-ideal}
设 $R$ 是环。若商环 $R/I$ 作为 $R$-模是平坦的，则称
$I \subset R$ 是{\it 纯的}。
\end{definition}

\begin{lemma}
\label{algebra-lemma-pure}
设 $R$ 是环，$I \subset R$ 是理想。以下条件等价：
\begin{enumerate}
\item $I$ 是纯的；
\item 对每个理想 $J \subset R$，都有 $J \cap I = IJ$；
\item 对每个有限生成理想 $J \subset R$，都有 $J \cap I = JI$；
\item 对每个 $x \in R$，都有 $(x) \cap I = xI$；
\item 对每个 $x \in I$，存在某个 $y \in I$，使得 $x = yx$；
\item 对每个 $x_1, \ldots, x_n \in I$，存在 $y \in I$，使得对所有
$i = 1, \ldots, n$ 都有 $x_i = yx_i$；
\item 对 $R$ 的每个素理想 $\mathfrak p$，有
$IR_{\mathfrak p} = 0$ 或 $IR_{\mathfrak p} = R_{\mathfrak p}$；
\item $\text{Supp}(I) = \Spec(R) \setminus V(I)$；
\item $I$ 是映射 $R \to (1 + I)^{-1}R$ 的核；
\item 对 $R$ 的某个乘法子集 $S$，作为 $R$-代数有
$R/I \cong S^{-1}R$；
\item 作为 $R$-代数有 $R/I \cong (1 + I)^{-1}R$。
\end{enumerate}
\end{lemma}

\begin{proof}
对 $R$ 的任意理想 $J$，都有短正合列
$0 \to J \to R \to R/J \to 0$。与 $R/I$ 作张量积，得到正合列
% P01-E0523: the source's R/I + J term should be the quotient R/(I + J).
$J \otimes_R R/I \to R/I \to R/I + J \to 0$，并且
$J \otimes_R R/I = J/JI$。因此，由引理 \ref{algebra-lemma-flat}，
(1)、(2) 和 (3) 等价。此外，它们蕴含 (4)。

\medskip\noindent
(4) $\Rightarrow$ (5) 是平凡的。假设 (5) 成立，并设
$x_1, \ldots, x_n \in I$。选择 $y_i \in I$，使得
$x_i = y_ix_i$。设 $y \in I$ 为满足
$1 - y = \prod_{i = 1, \ldots, n} (1 - y_i)$ 的元素。则对所有
$i = 1, \ldots, n$ 都有 $x_i = yx_i$。故 (6) 成立，从而
(5) $\Leftrightarrow$ (6)。

\medskip\noindent
假设 (5) 成立。设 $x \in I$。则对某个 $y \in I$ 有 $x = yx$。
所以 $x(1 - y) = 0$，这表明 $x$ 在 $(1 + I)^{-1}R$ 中的像为零。
映射 $R \to (1 + I)^{-1}R$ 的核当然总包含于 $I$。因此 (5)
蕴含 (9)。假设 (9) 成立。则对任意 $x \in I$，存在 $y \in I$
使得 $x(1 - y) = 0$。换言之，$x = yx$。故 (5) 与 (9) 等价。

\medskip\noindent
假设 (5) 成立。设 $\mathfrak p$ 是 $R$ 的素理想。若
$\mathfrak p \not \in V(I)$，则 $IR_{\mathfrak p} = R_{\mathfrak p}$。
若 $\mathfrak p \in V(I)$，换言之，若 $I \subset \mathfrak p$，
则对 $x \in I$，存在 $y \in I$ 使得 $x(1 - y) = 0$，因而
% P01-E0524: the source repeats “implies” and leaves this conditional sentence malformed.
$x$ 在 $R_{\mathfrak p}$ 中的像为零，即 $IR_{\mathfrak p} = 0$。
由此可见 (7) 成立。

\medskip\noindent
假设 (7) 成立。则对 $R$ 的任意素理想 $\mathfrak p$，
$(R/I)_{\mathfrak p}$ 或者是 $0$，或者是 $R_{\mathfrak p}$。
因此，由引理 \ref{algebra-lemma-flat-localization} 可知 (1) 成立。
至此可见 (1) -- (7) 和 (9) 全部等价。

\medskip\noindent
由于 $IR_{\mathfrak p} = I_{\mathfrak p}$，可知 (7) 蕴含 (8)。
最后，若 (8) 成立，则它恰好意味着：$I_{\mathfrak p}$ 是零模，
当且仅当 $\mathfrak p \in V(I)$；这显然就是 (7)。现在
(1) -- (9) 等价。

\medskip\noindent
假设 (1) -- (9) 成立。则由 (9) 有
$R/I \subset (1 + I)^{-1}R$；由 (7) 给出的素理想处局部化描述，映射
$R/I \to (1 + I)^{-1}R$ 也是满射。因此 (11) 成立。

\medskip\noindent
(11) $\Rightarrow$ (10) 是平凡的。而 (10) 蕴含 (1)，因为 $R$
的局部化作为 $R$-模是平坦的，参见引理
\ref{algebra-lemma-flat-localization}。
\end{proof}

\begin{lemma}
\label{algebra-lemma-pure-ideal-determined-by-zero-set}
\begin{slogan}
纯理想由其零点集决定。
\end{slogan}
设 $R$ 是环。若 $I, J \subset R$ 是纯理想，则
$V(I) = V(J)$ 蕴含 $I = J$。
\end{lemma}

\begin{proof}
例如，由引理 \ref{algebra-lemma-pure} 的性质 (7)，可知
$I = \Ker(R \to \prod_{\mathfrak p \in V(I)} R_{\mathfrak p})$，
所以可以从与它关联的闭子集恢复该理想。
\end{proof}

\begin{lemma}
\label{algebra-lemma-pure-open-closed-specializations}
设 $R$ 是环。规则 $I \mapsto V(I)$ 确定双射
$$
\{I \subset R \text{ pure}\}
\leftrightarrow
\{Z \subset \Spec(R)\text{ closed and closed under generalizations}\}
$$
\end{lemma}

\begin{proof}
设 $I$ 是纯理想。由于 $R \to R/I$ 平坦，由降链定理，一般化可以沿
映射 $\Spec(R/I) \to \Spec(R)$ 提升。因此 $V(I)$ 对一般化封闭，
这说明该映射良定义。由引理
\ref{algebra-lemma-pure-ideal-determined-by-zero-set}，该映射是单射。假设
$Z \subset \Spec(R)$ 是闭集且对一般化封闭。设 $J \subset R$
是满足 $Z = V(J)$ 的根理想。令
$I = \{x \in R : x \in xJ\}$。注意 $I$ 是理想：若
$x, y \in I$，则存在 $f, g \in J$，使得 $x = xf$ 且 $y = yg$。
于是
$$
x + y = (x + y)(f + g - fg)
$$
验证留给读者。我们断言 $I$ 是纯的，且 $V(I) = V(J)$。若断言成立，
则引理中的映射是满射，从而引理成立。

\medskip\noindent
注意 $I \subset J$，所以 $V(J) \subset V(I)$。设
$I \subset \mathfrak p$ 是素理想。考虑乘法子集
$S = (R \setminus \mathfrak p)(1 + J)$。由 $I$ 的定义以及
$I \subset \mathfrak p$，可知 $0 \not \in S$。因此可以找到一个与
$S$ 不交的 $R$ 的素理想 $\mathfrak q$，参见引理
\ref{algebra-lemma-localization-zero} 和 \ref{algebra-lemma-spec-localization}。
于是 $\mathfrak q \subset \mathfrak p$ 且
$\mathfrak q \cap (1 + J) = \emptyset$。这蕴含
$\mathfrak q + J$ 是 $R$ 的真理想。设 $\mathfrak m$ 是包含
$\mathfrak q + J$ 的极大理想。则 $\mathfrak m \in V(J)$；由于假设
$Z$ 对一般化封闭，故 $\mathfrak q \in V(J) = Z$。进而，由
$\mathfrak q \subset \mathfrak p$ 可知 $\mathfrak p \in V(J)$。
因此 $V(I) = V(J)$。

\medskip\noindent
最后，由于 $V(I) = V(J)$（且 $J$ 是根理想），可知
$J = \sqrt{I}$。取 $x \in I$；按定义，对某个 $y \in J$ 有
$x = xy$。于是 $x = xy = xy^2 = \ldots = xy^n$。因为对某个
$n > 0$ 有 $y^n \in I$，故引理 \ref{algebra-lemma-pure} 的性质 (5)
成立，从而 $I$ 的确是纯的。
\end{proof}

\begin{lemma}
\label{algebra-lemma-finitely-generated-pure-ideal}
设 $R$ 是环，$I \subset R$ 是理想。以下条件等价：
\begin{enumerate}
\item $I$ 是纯的且有限生成；
\item $I$ 由一个幂等元生成；
\item $I$ 是纯的且 $V(I)$ 是开集；
\item $R/I$ 是投射 $R$-模。
\end{enumerate}
\end{lemma}

\begin{proof}
若 (1) 成立，则由引理 \ref{algebra-lemma-pure} 有
$I = I \cap I = I^2$。于是由引理
\ref{algebra-lemma-ideal-is-squared-union-connected}，$I$ 由一个幂等元生成。
故 (1) $\Rightarrow$ (2)。若 (2) 成立，则 $I = (e)$，并且作为
$R$-模有 $R = (1 - e) \oplus (e)$；因此 $R/I$ 是平坦的，$I$
是纯的，且 $V(I) = D(1 - e)$ 是开集。故 (2) $\Rightarrow$
(1) $+$ (3)。最后，假设 (3) 成立。则 $V(I)$ 既开又闭，所以对
$R$ 的某个幂等元 $e$，有 $V(I) = D(1 - e)$，参见引理
\ref{algebra-lemma-disjoint-decomposition}。理想 $J = (e)$ 是满足
$V(J) = V(I)$ 的纯理想，故由引理
\ref{algebra-lemma-pure-ideal-determined-by-zero-set} 有 $I = J$。由此可见
(3) $\Rightarrow$ (2)。由引理 \ref{algebra-lemma-finite-projective}，
(4) 等价于 $I$ 是纯的且 $R/I$ 有限表现。又由引理
\ref{algebra-lemma-extension}，$R/I$ 有限表现当且仅当 $I$ 有限生成。
因此 (4) 等价于 (1)。
\end{proof}

\noindent
可以用上面的结果刻画使每个有限平坦模都有限表现的那些环。

\begin{lemma}
\label{algebra-lemma-finite-flat-module-finitely-presented}
设 $R$ 是环。以下条件等价：
\begin{enumerate}
\item 每个闭且对一般化封闭的 $Z \subset \Spec(R)$ 也是开集；
\item 任意有限平坦 $R$-模都是有限局部自由的。
\end{enumerate}
\end{lemma}

\begin{proof}
若任意有限平坦 $R$-模都有限局部自由，则对纯理想 $I$，$R/I$
的支撑是开集。因此，(2) $\Rightarrow$ (1) 由引理
% P01-E0525: the cited uniqueness lemma does not supply the pure-ideal/closed-set correspondence used here.
\ref{algebra-lemma-pure-ideal-determined-by-zero-set} 得出。

\medskip\noindent
反过来，假设 $R$ 满足 (1)。设 $M$ 是有限平坦 $R$-模。
$M$ 的支撑 $Z = \text{Supp}(M)$ 是闭集，参见引理
\ref{algebra-lemma-support-closed}。另一方面，若
$\mathfrak p \subset \mathfrak p'$，则由引理
\ref{algebra-lemma-finite-flat-local}，$M_{\mathfrak p'}$ 是自由的，并且
% P01-E0526: the source sentence lacks terminal punctuation after this displayed inline equality.
$M_{\mathfrak p} = M_{\mathfrak p'} \otimes_{R_{\mathfrak p'}} R_{\mathfrak p}$。
因此
$\mathfrak p' \in \text{Supp}(M) \Rightarrow \mathfrak p \in \text{Supp}(M)$；
换言之，该支撑对一般化封闭。由于 $R$ 满足 (1)，可知 $M$ 的支撑
既开又闭。假设 $M$ 由 $r$ 个元素 $m_1, \ldots, m_r$ 生成。
各模 $\wedge^i(M)$（$i = 1, \ldots, r$）也是有限平坦 $R$-模，
因为 $\wedge^i(M)_{\mathfrak p} = \wedge^i(M_{\mathfrak p})$
在 $R_{\mathfrak p}$ 上是自由的。注意
$\text{Supp}(\wedge^{i + 1}(M)) \subset \text{Supp}(\wedge^i(M))$。
因此存在由开闭子集给出的分解
$$
\Spec(R) = U_0 \amalg U_1 \amalg \ldots \amalg U_r
$$
使得对所有 $i = 0, \ldots, r$，$\wedge^i(M)$ 的支撑为
$U_r \cup \ldots \cup U_i$。设 $\mathfrak p$ 是 $R$ 的素理想，
并设 $\mathfrak p \in U_i$。注意
$\wedge^i(M) \otimes_R \kappa(\mathfrak p) =
\wedge^i(M \otimes_R \kappa(\mathfrak p))$。必要时重排
$m_1, \ldots, m_r$，可以假设 $m_1, \ldots, m_i$ 生成
$M \otimes_R \kappa(\mathfrak p)$。由 Nakayama 引理
\ref{algebra-lemma-NAK}，对某个 $f \in R$、$f \not \in \mathfrak p$，得到满射
$$
R_f^{\oplus i} \longrightarrow M_f, \quad
(a_1, \ldots, a_i) \longmapsto \sum a_im_i
$$
还可以假设 $D(f) \subset U_i$。这意味着
$\wedge^i(M_f) = \wedge^i(M)_f$ 是平坦 $R_f$-模，其支撑是整个
$\Spec(R_f)$。由上可知，它由单个元素
$m_1 \wedge \ldots \wedge m_i$ 生成。因此，对某个纯理想
$J \subset R_f$，有 $\wedge^i(M)_f \cong R_f/J$，并且
$V(J) = \Spec(R_f)$。显然这意味着 $J = (0)$，参见引理
\ref{algebra-lemma-pure-ideal-determined-by-zero-set}。所以
$m_1 \wedge \ldots \wedge m_i$ 是 $\wedge^i(M_f)$ 的一组基，
从而所显示的映射既是单射也是满射。这按要求证明了 $M$ 有限局部自由。
\end{proof}

\section{有限整体维数的环}
\label{algebra-section-ring-finite-gl-dim}

\noindent
下面的引理常用于比较给定模的不同投射消解。

\begin{lemma}[Schanuel 引理]
\label{algebra-lemma-Schanuel}
设 $R$ 是环，$M$ 是 $R$-模。假设
$$
0 \to K \xrightarrow{c_1} P_1 \xrightarrow{p_1} M \to 0
\quad\text{且}\quad
0 \to L \xrightarrow{c_2} P_2 \xrightarrow{p_2} M \to 0
$$
是两个短正合列，其中 $P_i$ 是投射模。则
$K \oplus P_2 \cong L \oplus P_1$。更精确地，存在交换图
$$
\xymatrix{
0 \ar[r] &
K \oplus P_2 \ar[r]_{(c_1, \text{id})} \ar[d] &
P_1 \oplus P_2 \ar[r]_{(p_1, 0)} \ar[d] &
M \ar[r] \ar@{=}[d] &
0 \\
0 \ar[r] &
P_1 \oplus L \ar[r]^{(\text{id}, c_2)} &
P_1 \oplus P_2 \ar[r]^{(0, p_2)} &
M \ar[r] &
0
}
$$
使得各竖直箭头都是同构。
\end{lemma}

\begin{proof}
考虑由短正合列
$0 \to N \to P_1 \oplus P_2 \to M \to 0$
定义的模 $N$，其中最后一个映射是两个映射 $P_i \to M$ 之和。
容易看出，投影 $N \to P_1$ 是满射，其核为 $L$；而
$N \to P_2$ 是满射，其核为 $K$。由于 $P_i$ 是投射模，故
$N \cong K \oplus P_2 \cong L \oplus P_1$。这证明了第一个断言。

\medskip\noindent
为证明第二个断言（同时再次证明第一个断言），选择
$a : P_1 \to P_2$ 和 $b : P_2 \to P_1$，使得
$p_1 = p_2 \circ a$ 且 $p_2 = p_1 \circ b$。这是可行的，因为
$P_1$ 和 $P_2$ 都是投射模。于是得到交换图
$$
\xymatrix{
0 \ar[r] &
K \oplus P_2 \ar[r]_{(c_1, \text{id})} &
P_1 \oplus P_2 \ar[r]_{(p_1, 0)} &
M \ar[r] &
0 \\
0 \ar[r] &
N \ar[r] \ar[d] \ar[u] &
P_1 \oplus P_2 \ar[r]_{(p_1, p_2)}
\ar[d]_S \ar[u]^T &
M \ar[r] \ar@{=}[d] \ar@{=}[u] &
0 \\
0 \ar[r] &
P_1 \oplus L \ar[r]^{(\text{id}, c_2)} &
P_1 \oplus P_2 \ar[r]^{(0, p_2)} &
M \ar[r] &
0
}
$$
其中 $T$ 和 $S$ 由矩阵
$$
S = \left(
\begin{matrix}
\text{id} & 0 \\
a & \text{id}
\end{matrix}
\right)
\quad\text{且}\quad
T = \left(
\begin{matrix}
\text{id} & b \\
0 & \text{id}
\end{matrix}
\right)
$$
给出。于是 $S$、$T$ 以及映射 $N \to P_1 \oplus L$ 和
$N \to K \oplus P_2$ 都是所需的同构。
\end{proof}

\begin{definition}
\label{algebra-definition-finite-proj-dim}
设 $R$ 是环，$M$ 是 $R$-模。若 $M$ 有一个由投射 $R$-模组成的
有限长度消解，则称其具有{\it 有限投射维数}。这种消解的最小长度
称为 $M$ 的{\it 投射维数}。
\end{definition}

\noindent
显然，$M$ 的投射维数为 $0$ 当且仅当 $M$ 是投射模。下面的引理说明，
投射维数在何种程度上不依赖于投射消解的选择。

\begin{lemma}
\label{algebra-lemma-independent-resolution}
设 $R$ 是环。假设 $M$ 是投射维数为 $d$ 的 $R$-模。假设
$F_e \to F_{e-1} \to \ldots \to F_0 \to M \to 0$
正合，其中 $F_i$ 是投射模，并且 $e \geq d - 1$。则
$F_e \to F_{e-1}$ 的核是投射模（当 $e = 0$ 时，则指
$F_0 \to M$ 的核是投射模）。
\end{lemma}

\begin{proof}
对 $d$ 归纳。若 $d = 0$，则 $M$ 是投射模。在这种情形下有分裂
$F_0 = \Ker(F_0 \to M) \oplus M$，所以
$\Ker(F_0 \to M)$ 是投射模。若 $e = 0$，证明至此完成；若
$e > 0$，则以 $\Ker(F_0 \to M)$ 代替 $M$，便使 $e$ 减少。

\medskip\noindent
下面假设 $d > 0$。设
$0 \to P_d \to P_{d-1} \to \ldots \to P_0 \to M \to 0$
是由投射模 $P_i$ 组成的最短有限消解。由 Schanuel 引理
\ref{algebra-lemma-Schanuel}，
$P_0 \oplus \Ker(F_0 \to M) \cong F_0 \oplus \Ker(P_0 \to M)$。
当 $d = 1$、$e = 0$ 时，这证明了结论，因为此时右端是投射模
$F_0 \oplus P_1$。因此现在可以假设 $e > 0$。模
$F_0 \oplus \Ker(P_0 \to M)$ 有长度为 $d - 1$ 的有限投射消解
$$
0 \to P_d \to P_{d-1} \to \ldots \to
P_2 \to P_1 \oplus F_0 \to \Ker(P_0 \to M) \oplus F_0 \to 0
$$
把归纳假设应用于长度为 $e - 1$ 的正合列
$$
F_e \to F_{e-1} \to \ldots \to F_2 \to P_0 \oplus F_1 \to
P_0 \oplus \Ker(F_0 \to M) \to 0
$$
可知，当 $e \geq 2$ 时 $\Ker(F_e \to F_{e - 1})$ 是投射模；
或者，$\Ker(F_1 \oplus P_0 \to F_0 \oplus P_0)$ 是投射模。
这蕴含引理。
\end{proof}

\begin{lemma}
\label{algebra-lemma-what-kind-of-resolutions}
设 $R$ 是环，$M$ 是 $R$-模，并设 $d \geq 0$。以下条件等价：
\begin{enumerate}
\item $M$ 的投射维数 $\leq d$；
\item 存在由投射模 $P_i$ 组成的消解
$0 \to P_d \to P_{d - 1} \to \ldots \to P_0 \to M \to 0$；
\item 对某个由投射模 $P_i$ 组成的消解
$\ldots \to P_2 \to P_1 \to P_0 \to M \to 0$，当
$d \geq 2$ 时 $\Ker(P_{d - 1} \to P_{d - 2})$ 是投射模；当
$d = 1$ 时 $\Ker(P_0 \to M)$ 是投射模；当 $d = 0$ 时
$M$ 是投射模；
\item 对任意由投射模 $P_i$ 组成的消解
$\ldots \to P_2 \to P_1 \to P_0 \to M \to 0$，当
$d \geq 2$ 时 $\Ker(P_{d - 1} \to P_{d - 2})$ 是投射模；当
$d = 1$ 时 $\Ker(P_0 \to M)$ 是投射模；当 $d = 0$ 时
$M$ 是投射模。
\end{enumerate}
\end{lemma}

\begin{proof}
(1) 与 (2) 的等价性就是投射维数的定义，参见定义
\ref{algebra-definition-finite-proj-dim}。由引理
\ref{algebra-lemma-independent-resolution} 有 (2) $\Rightarrow$ (4)。
(4) $\Rightarrow$ (3) 和 (3) $\Rightarrow$ (2) 都是直接的。
\end{proof}

\begin{lemma}
\label{algebra-lemma-what-kind-of-resolutions-local}
设 $R$ 是局部环，$M$ 是 $R$-模，并设 $d \geq 0$。引理
\ref{algebra-lemma-what-kind-of-resolutions} 中的等价条件 (1) -- (4)
还等价于
\begin{enumerate}
\item[(5)] 存在由自由模 $P_i$ 组成的消解
$0 \to P_d \to P_{d - 1} \to \ldots \to P_0 \to M \to 0$。
\end{enumerate}
\end{lemma}

\begin{proof}
由引理 \ref{algebra-lemma-what-kind-of-resolutions} 和定理
\ref{algebra-theorem-projective-free-over-local-ring} 得出。
\end{proof}

\begin{lemma}
\label{algebra-lemma-what-kind-of-resolutions-Noetherian}
设 $R$ 是诺特环，$M$ 是有限 $R$-模，并设 $d \geq 0$。引理
\ref{algebra-lemma-what-kind-of-resolutions} 中的等价条件 (1) -- (4)
还等价于
\begin{enumerate}
\item[(6)] 存在由有限投射模 $P_i$ 组成的消解
$0 \to P_d \to P_{d - 1} \to \ldots \to P_0 \to M \to 0$。
\end{enumerate}
\end{lemma}

\begin{proof}
选择由有限自由模 $F_i$ 组成的消解
$\ldots \to F_2 \to F_1 \to F_0 \to M \to 0$
（引理 \ref{algebra-lemma-resolution-by-finite-free}）。由引理
\ref{algebra-lemma-what-kind-of-resolutions} 可知，至少当 $d \geq 2$ 时，
% P01-E0527: the source proof does not address the d = 0 and d = 1 cases included in the statement.
$P_d = \Ker(F_{d - 1} \to F_{d - 2})$ 是投射模。由于 $R$ 是诺特环，
并且 $P_d \subset F_{d - 1}$、后者有限自由，故 $P_d$
是有限 $R$-模。因此
$0 \to P_d \to F_{d - 1} \to \ldots \to F_1 \to F_0 \to M \to 0$
就是所需的消解。
\end{proof}

\begin{lemma}
\label{algebra-lemma-what-kind-of-resolutions-Noetherian-local}
设 $R$ 是诺特局部环，$M$ 是有限 $R$-模，并设 $d \geq 0$。引理
\ref{algebra-lemma-what-kind-of-resolutions} 中的等价条件 (1) -- (4)、
引理 \ref{algebra-lemma-what-kind-of-resolutions-local} 的条件 (5)，以及
引理 \ref{algebra-lemma-what-kind-of-resolutions-Noetherian} 的条件 (6)
还等价于
\begin{enumerate}
\item[(7)] 存在由有限自由模 $F_i$ 组成的消解
$0 \to F_d \to F_{d - 1} \to \ldots \to F_0 \to M \to 0$。
\end{enumerate}
\end{lemma}

\begin{proof}
这由引理 \ref{algebra-lemma-what-kind-of-resolutions}、
\ref{algebra-lemma-what-kind-of-resolutions-local} 和
\ref{algebra-lemma-what-kind-of-resolutions-Noetherian} 得出；还要用到局部环上的
有限投射模是有限自由模，参见引理 \ref{algebra-lemma-finite-projective}。
\end{proof}

\begin{lemma}
\label{algebra-lemma-projective-dimension-ext}
设 $R$ 是环，$M$ 是 $R$-模，并设 $n \geq 0$。以下条件等价：
\begin{enumerate}
\item $M$ 的投射维数 $\leq n$；
\item 对所有 $R$-模 $N$ 和所有 $i \geq n + 1$，都有
$\Ext^i_R(M, N) = 0$；
\item 对所有 $R$-模 $N$，都有 $\Ext^{n + 1}_R(M, N) = 0$。
\end{enumerate}
\end{lemma}

\begin{proof}
假设 (1) 成立。选择 $M$ 的自由消解 $F_\bullet \to M$。记
$d_e : F_e \to F_{e - 1}$。由引理
\ref{algebra-lemma-independent-resolution} 可知，对 $e \geq n - 1$，
$P_e = \Ker(d_e)$ 是投射模。于是对 $e \geq n$ 有
% P01-E0528: for n = 0 the source uses P_{e - 1} at e = 0 without defining P_{-1}.
$F_e \cong P_e \oplus P_{e - 1}$，其中 $d_e$ 把直和项
$P_{e - 1}$ 同构地映到 $F_{e - 1}$ 中的 $P_{e - 1}$。
因此，对任意 $R$-模 $N$，复形 $\Hom_R(F_\bullet, N)$
在次数 $\geq n + 1$ 上分裂正合。故 (2) 成立。
(2) $\Rightarrow$ (3) 是平凡的。

\medskip\noindent
假设 (3) 成立。若 $n = 0$，则由引理
\ref{algebra-lemma-characterize-projective}，$M$ 是投射模，所以 (1) 成立。
若 $n > 0$，选择自由 $R$-模 $F$ 和核为 $K$ 的满射 $F \to M$。
由引理 \ref{algebra-lemma-reverse-long-exact-seq-ext}，以及根据 (1)
对所有 $i > 0$ 都有 $\Ext_R^i(F, N)$ 消失，可知对所有 $R$-模
$N$ 都有 $\Ext_R^n(K, N) = 0$。故由归纳假设，$K$ 的投射维数
$\leq n - 1$。在 $K$ 的任意有限投射消解前添加 $F$，便得到长度
多一的 $M$ 的投射消解，所以 $M$ 的投射维数 $\leq n$。
\end{proof}

\begin{lemma}
\label{algebra-lemma-exact-sequence-projective-dimension}
设 $R$ 是环，并设 $0 \to M' \to M \to M'' \to 0$ 是 $R$-模的短正合列。
\begin{enumerate}
\item 若 $M$ 的投射维数 $\leq n$，且 $M''$ 的投射维数
$\leq n + 1$，则 $M'$ 的投射维数 $\leq n$。
\item 若 $M'$ 和 $M''$ 的投射维数都 $\leq n$，则 $M$
的投射维数 $\leq n$。
\item 若 $M'$ 的投射维数 $\leq n$，且 $M$ 的投射维数
$\leq n + 1$，则 $M''$ 的投射维数 $\leq n + 1$。
\end{enumerate}
\end{lemma}

\begin{proof}
把引理 \ref{algebra-lemma-projective-dimension-ext} 中对投射维数的刻画与引理
\ref{algebra-lemma-reverse-long-exact-seq-ext} 中的 Ext 群长正合列结合起来。
\end{proof}

\begin{definition}
\label{algebra-definition-finite-gl-dim}
设 $R$ 是环。若存在整数 $n$，使得每个 $R$-模都有一个长度至多为
$n$ 的投射 $R$-模消解，则称环 $R$ 具有{\it 有限整体维数}。
满足此条件的最小 $n$ 称为 $R$ 的{\it 整体维数}。
\end{definition}

\noindent
下面引理的证明论证可见 Auslander 的论文 \cite{Auslander}。

\begin{lemma}
\label{algebra-lemma-colimit-projective-dimension}
设 $R$ 是环。假设有模 $M = \bigcup_{e \in E} M_e$，其中各子模
$M_e$ 按包含关系良序排列。假设各商模
$M_e/\bigcup\nolimits_{e' < e} M_{e'}$ 的投射维数 $\leq n$。
则 $M$ 的投射维数 $\leq n$。
\end{lemma}

\begin{proof}
对 $n$ 归纳。

\medskip\noindent
基本情形：$n = 0$。此时
$P_e = M_e/\bigcup_{e' < e} M_{e'}$ 是投射模。因此可以选择投影
$M_e \to P_e$ 的截面 $P_e \to M_e$。我们断言，诱导映射
$\psi : \bigoplus_{e \in E} P_e \to M$ 是同构。事实上，若
$x = \sum x_e \in \bigoplus P_e$ 非零，则令 $e_{max}$ 为使
$x_{e_{max}}$ 非零的最大指标，并令
$y = \psi(x) = \psi(\sum x_e)$。由于
$y \in M_{e_{max}}$ 在 $P_{e_{max}}$ 中的像 $x_{e_{max}}$ 非零，
故该元素非零。
另一方面，设 $y \in M$。则对某个 $e$ 有 $y \in M_e$。
我们对 $e$ 作超限归纳，证明 $y \in \Im(\psi)$。设
$x_e \in P_e$ 是 $y$ 的像。则
$y - \psi(x_e) \in \bigcup_{e' < e} M_{e'}$。由归纳假设，
$y - \psi(x_e) \in \Im(\psi)$，所以 $y \in \Im(\psi)$。
因此断言成立，$\psi$ 是同构。由引理
\ref{algebra-lemma-direct-sum-projective}，$M$ 作为投射模的直和是投射模。

\medskip\noindent
若 $n > 0$，则对 $e \in E$，令 $F_e$ 为以 $M_e$ 的元素集为基的
自由 $R$-模。于是，在良序集 $E$ 上有一族短正合列
$$
0 \to K_e \to F_e \to M_e \to 0
$$
注意，转移映射 $F_{e'} \to F_e$ 和 $K_{e'} \to K_e$ 也都是单射。
令 $F = \bigcup F_e$、$K = \bigcup K_e$。则
$$
0 \to
K_e/\bigcup\nolimits_{e' < e} K_{e'} \to
F_e/\bigcup\nolimits_{e' < e} F_{e'} \to
M_e/\bigcup\nolimits_{e' < e} M_{e'} \to 0
$$
也是 $R$-模的短正合列，并且
$F_e/\bigcup_{e' < e} F_{e'}$ 是以 $M_e$ 中不属于
$\bigcup_{e' < e} M_{e'}$ 的元素集为基的自由 $R$-模。因此，由引理
\ref{algebra-lemma-exact-sequence-projective-dimension} 可知
$K_e/\bigcup_{e' < e} K_{e'}$ 的投射维数至多为 $n - 1$。
由归纳假设，$K$ 的投射维数至多为 $n - 1$。因而 $M$
的投射维数至多为 $n$，得证。
\end{proof}

\begin{lemma}
\label{algebra-lemma-finite-gl-dim}
设 $R$ 是环。以下条件等价：
\begin{enumerate}
\item $R$ 的整体维数有限且 $\leq n$；
\item 每个有限 $R$-模的投射维数 $\leq n$；
\item 每个循环 $R$-模 $R/I$ 的投射维数 $\leq n$。
\end{enumerate}
\end{lemma}

\begin{proof}
显然 (1) $\Rightarrow$ (2)，且 (2) $\Rightarrow$ (3)。
假设 (3) 成立。选择 $M$ 的生成元集 $E \subset M$，并在 $E$
上选择一个良序。对 $e \in E$，记 $M_e$ 为由满足 $e' \leq e$
的元素 $e' \in E$ 生成的 $M$ 的子模。则
$M = \bigcup_{e \in E} M_e$。注意，对每个 $e \in E$，商模
$$
M_e/\bigcup\nolimits_{e' < e} M_{e'}
$$
或者为零，或者由一个元素生成，故由 (3)，其投射维数 $\leq n$。
由引理 \ref{algebra-lemma-colimit-projective-dimension}，这意味着 $M$
的投射维数 $\leq n$。
\end{proof}

\begin{lemma}
\label{algebra-lemma-localize-finite-gl-dim}
设 $R$ 是环，$M$ 是 $R$-模，并设 $S \subset R$ 是乘法子集。
\begin{enumerate}
\item 若 $M$ 的投射维数 $\leq n$，则 $S^{-1}M$ 作为
$S^{-1}R$-模的投射维数 $\leq n$。
\item 若 $R$ 的整体维数有限且 $\leq n$，则 $S^{-1}R$
的整体维数有限且 $\leq n$。
\end{enumerate}
\end{lemma}

\begin{proof}
设
$0 \to P_n \to P_{n - 1} \to \ldots \to P_0 \to M \to 0$
是投射消解。局部化是正合的，参见命题
\ref{algebra-proposition-localization-exact}；且每个 $S^{-1}P_i$ 都是投射
$S^{-1}R$-模，参见引理 \ref{algebra-lemma-ascend-properties-modules}。
因此
$0 \to S^{-1}P_n \to \ldots \to S^{-1}P_0 \to S^{-1}M \to 0$
是 $S^{-1}M$ 的投射消解。这证明了 (1)。设 $M'$ 是
$S^{-1}R$-模。注意 $M' = S^{-1}M'$。因此 (2) 由 (1) 得出。
\end{proof}

\section{正则环与整体维数}
\label{algebra-section-regular-finite-gl-dim}

\noindent
可以用有限整体维数环的有关材料给出正则局部环的另一个刻画。

\begin{proposition}
\label{algebra-proposition-regular-finite-gl-dim}
设 $R$ 是维数为 $d$ 的正则局部环。每个深度为 $e$ 的有限
$R$-模 $M$ 都有有限自由消解
$$
0 \to F_{d-e} \to \ldots \to F_0 \to M \to 0.
$$
特别地，正则局部环的整体维数 $\leq d$。
\end{proposition}

\begin{proof}
第一部分由引理 \ref{algebra-lemma-regular-mcm-free} 和引理
\ref{algebra-lemma-mcm-resolution} 得出。最后一部分由此及引理
\ref{algebra-lemma-finite-gl-dim} 得出。
\end{proof}

\begin{lemma}
\label{algebra-lemma-finite-gl-dim-primes}
设 $R$ 是诺特环，$n \geq 0$ 是整数。则 $R$ 的整体维数有限且
$\leq n$，当且仅当对 $R$ 的所有极大理想 $\mathfrak m$，环
$R_{\mathfrak m}$ 的整体维数 $\leq n$。
\end{lemma}

\begin{proof}
由引理 \ref{algebra-lemma-localize-finite-gl-dim} 已知，若 $R$
的有限整体维数为 $n$，则所有局部化 $R_{\mathfrak m}$
的有限整体维数至多为 $n$。
% P01-E0529: the source reads “We saw, Lemma ... that”, with a malformed citation clause.
反过来，假设所有 $R_{\mathfrak m}$ 的整体维数都 $\leq n$。
设 $M$ 是有限 $R$-模，并设
% P01-E0530: for n = 0 the source resolution uses F_{n-1}=F_{-1} without a convention.
$0 \to K_n \to F_{n-1} \to \ldots \to F_0 \to M \to 0$
是由有限自由模 $F_i$ 组成的消解。则 $K_n$ 是有限 $R$-模。
由引理 \ref{algebra-lemma-independent-resolution} 和假设，所有模
$K_n \otimes_R R_{\mathfrak m}$ 都是投射模。因此，由引理
\ref{algebra-lemma-finite-projective}，模 $K_n$ 是有限投射模。
\end{proof}

\begin{lemma}
\label{algebra-lemma-length-resolution-residue-field}
假设 $R$ 是极大理想为 $\mathfrak m$、剩余域为 $\kappa$
的诺特局部环。在这种情形下，$\kappa$ 的投射维数
$\geq \dim_\kappa \mathfrak m / \mathfrak m^2$。
\end{lemma}

\begin{proof}
设 $x_1 , \ldots, x_n$ 是 $\mathfrak m$ 中的元素，其在
$\mathfrak m / \mathfrak m^2$ 中的像构成一组基。考虑
$x_1, \ldots, x_n$ 上的{\it Koszul 复形}。这是复形
$$
0 \to \wedge^n R^n \to \wedge^{n-1} R^n \to \wedge^{n-2} R^n \to
\ldots \to \wedge^i R^n \to \ldots \to R^n \to R
$$
其映射由
$$
e_{j_1} \wedge \ldots \wedge e_{j_i}
\longmapsto
\sum_{a = 1}^i (-1)^{a + 1} x_{j_a} e_{j_1} \wedge \ldots
\wedge \hat e_{j_a} \wedge \ldots \wedge e_{j_i}
$$
给出。容易看出这是一个复形 $K_{\bullet}(R, x_{\bullet})$。
注意，由引理 \ref{algebra-lemma-NAK} 第 (8) 部分，
$K_{\bullet}(R, x_{\bullet})$ 的最后一个映射的余核是 $\kappa$。

\medskip\noindent
若 $\kappa$ 的投射维数有限且为 $d$，则可以找到由有限自由
$R$-模组成的长度为 $d$ 的消解 $F_{\bullet} \to \kappa$
（引理 \ref{algebra-lemma-what-kind-of-resolutions-Noetherian-local}）。
由引理 \ref{algebra-lemma-add-trivial-complex}，可以假设复形
$F_{\bullet}$ 中的所有映射都满足
$\Im(F_i \to F_{i-1}) \subset \mathfrak m F_{i-1}$，因为从消解中
移除平凡直和项至多只会缩短消解。由引理
\ref{algebra-lemma-compare-resolutions}，可以找到诱导 $\kappa$ 上恒等映射的
复形映射 $\alpha : K_{\bullet}(R, x_{\bullet}) \to F_{\bullet}$。
我们将归纳证明各映射
$\alpha_i : \wedge^i R^n = K_i(R, x_{\bullet}) \to F_i$ 都满足
$\alpha_i \otimes \kappa : \wedge^i \kappa^n \to F_i \otimes \kappa$
是单射。这表明 $F_n \not = 0$，因而按要求有 $d \geq n$。

\medskip\noindent
当 $i = 0$ 时，结论是清楚的，因为复合
$R \xrightarrow{\alpha_0} F_0 \to \kappa$ 非零。注意，$F_0$
的秩必为 $1$，否则余核为单个 $\kappa$ 的映射
$F_1 \to F_0$ 不可能具有包含在 $\mathfrak m F_0$ 中的像。

\medskip\noindent
接着检验 $i = 1$ 的情形，因为我们认为这颇具启发性；读者可以跳过，
因为归纳步骤会由 $i = 0$ 的情形推出 $i = 1$ 的情形。上面已经看到
$F_0 = R$，并且 $F_1 \to F_0 = R$ 的像是 $\mathfrak m$。
有交换图
$$
\begin{matrix}
R^n & = & K_1(R, x_{\bullet}) & \to & K_0(R, x_{\bullet}) & = & R \\
& & \downarrow & & \downarrow & & \downarrow \\
& & F_1 & \to & F_0 & = & R
\end{matrix}
$$
其中最右边的竖直箭头由乘以一个单位给出。因此，复合
$R^n \to F_1 \to F_0 = R$ 的像也等于 $\mathfrak m$。
由于 $\dim_\kappa (\mathfrak m / \mathfrak m^2) = n$，映射
$R^n \otimes \kappa \to F_1 \otimes \kappa$ 必为单射。

\medskip\noindent
设 $i \geq 1$，并假设已经对所有 $j \leq i - 1$ 证明了
$\alpha_j \otimes \kappa$ 的单射性。考虑交换图
$$
\begin{matrix}
\wedge^i R^n & = & K_i(R, x_{\bullet}) & \to & K_{i-1}(R, x_{\bullet})
& = & \wedge^{i-1} R^n \\
& & \downarrow & & \downarrow & & \\
& & F_i & \to & F_{i-1} & &
\end{matrix}
$$
已知 $\wedge^{i-1} \kappa^n \to F_{i-1} \otimes \kappa$
是单射。这证明了
$\wedge^{i-1} \kappa^n \otimes_{\kappa} \mathfrak m/\mathfrak m^2
\to F_{i-1} \otimes \mathfrak m/\mathfrak m^2$ 是单射。
此外，由复形的选择，$F_i$ 映入 $\mathfrak mF_{i-1}$；
Koszul 复形也同样如此。因此得到交换图
$$
\begin{matrix}
\wedge^i \kappa^n & \to &
\wedge^{i-1} \kappa^n \otimes \mathfrak m/\mathfrak m^2 \\
\downarrow & & \downarrow \\
F_i \otimes \kappa & \to & F_{i-1} \otimes \mathfrak m/\mathfrak m^2
\end{matrix}
$$
此时，只须验证映射
$\wedge^i \kappa^n \to
\wedge^{i-1} \kappa^n \otimes \mathfrak m/\mathfrak m^2$
是单射，而这可以直接检验。
\end{proof}

\begin{lemma}
\label{algebra-lemma-dim-gl-dim}
设 $R$ 是诺特局部环。假设剩余域 $\kappa$ 在 $R$
上的投射维数有限且为 $n$。在这种情形下，$\dim(R) \geq n$。
\end{lemma}

\begin{proof}
设 $F_{\bullet}$ 是由有限自由 $R$-模组成的 $\kappa$ 的有限消解
（引理 \ref{algebra-lemma-what-kind-of-resolutions-Noetherian-local}）。
由引理 \ref{algebra-lemma-add-trivial-complex}，可以假设复形
$F_{\bullet}$ 中的所有映射都具有性质
% P01-E0531: the source says “have to property” instead of “have the property”.
$\Im(F_i \to F_{i-1}) \subset \mathfrak m F_{i-1}$，因为从消解中移除
平凡直和项至多只会缩短消解。设 $F_n \not = 0$，且对 $i > n$
有 $F_i = 0$，从而 $\kappa$ 的投射维数为 $n$。由命题
\ref{algebra-proposition-what-exact}，可知
$\text{depth}_{I(\varphi_n)}(R) \geq n$，因为按复形的选择，
$I(\varphi_n)$ 不可能等于 $R$。故由引理
\ref{algebra-lemma-bound-depth} 还有 $\dim(R) \geq n$。
\end{proof}

\begin{proposition}
\label{algebra-proposition-finite-gl-dim-regular}
设 $(R, \mathfrak m, \kappa)$ 是诺特局部环。以下条件等价：
\begin{enumerate}
\item $\kappa$ 作为 $R$-模具有有限投射维数；
\item $R$ 具有有限整体维数；
\item $R$ 是正则局部环。
\end{enumerate}
此外，在这种情形下，$R$ 的整体维数等于
$\dim(R) = \dim_\kappa(\mathfrak m/\mathfrak m^2)$。
\end{proposition}

\begin{proof}
由命题 \ref{algebra-proposition-regular-finite-gl-dim} 有
(3) $\Rightarrow$ (2)。(2) $\Rightarrow$ (1) 是平凡的。
假设 (1) 成立。由引理 \ref{algebra-lemma-length-resolution-residue-field}
和 \ref{algebra-lemma-dim-gl-dim}，可知
$\dim(R) \geq \dim_\kappa(\mathfrak m /\mathfrak m^2)$。
因此 $R$ 是正则的，参见定义 \ref{algebra-definition-regular-local}
及其前面的讨论。假设等价条件 (1) -- (3) 成立。由命题
\ref{algebra-proposition-regular-finite-gl-dim}，$R$ 的整体维数至多为
$\dim(R)$；由引理 \ref{algebra-lemma-length-resolution-residue-field}，
它至少为 $\dim_\kappa(\mathfrak m/\mathfrak m^2)$。故所述等式成立。
\end{proof}

\begin{lemma}
\label{algebra-lemma-localization-of-regular-local-is-regular}
诺特局部环 $R$ 是正则局部环，当且仅当它具有有限整体维数。
在这种情形下，对所有素理想 $\mathfrak p$，$R_{\mathfrak p}$
都是正则局部环。
\end{lemma}

\begin{proof}
由命题 \ref{algebra-proposition-finite-gl-dim-regular} 和
\ref{algebra-proposition-regular-finite-gl-dim} 可知，诺特局部环是正则局部环，
当且仅当它具有有限整体维数。此外，任意局部化
$R_{\mathfrak p}$ 都具有有限整体维数，参见引理
\ref{algebra-lemma-localize-finite-gl-dim}，因而是正则局部环。
\end{proof}

\noindent
由引理 \ref{algebra-lemma-localization-of-regular-local-is-regular}，
作如下定义是合理的，因为它不与此前的正则局部环定义冲突。

\begin{definition}
\label{algebra-definition-regular}
若诺特环 $R$ 在所有素理想处的局部化 $R_{\mathfrak p}$
都是正则局部环，则称其是{\it 正则的}。
\end{definition}

\noindent
只要求极大理想处的局部环正则就已足够。注意，即使假设 $R$
是诺特环，这也不同于要求 $R$ 具有有限整体维数。这是因为存在
不具有有限整体维数的正则诺特环；其原因正是它不具有有限维数。

\begin{lemma}
\label{algebra-lemma-finite-gl-dim-finite-dim-regular}
设 $R$ 是诺特环。以下条件等价：
\begin{enumerate}
\item $R$ 的有限整体维数为 $n$；
\item $R$ 是维数为 $n$ 的正则环；
\item 存在整数 $n$，使得极大理想处的所有局部化
$R_{\mathfrak m}$ 都是维数 $\leq n$ 的正则环，并且至少对一个
$\mathfrak m$ 取到等号；
\item 存在整数 $n$，使得素理想处的所有局部化
$R_{\mathfrak p}$ 都是维数 $\leq n$ 的正则环，并且至少对一个
$\mathfrak p$ 取到等号。
\end{enumerate}
\end{lemma}

\begin{proof}
这由上面的讨论得出。更精确地说，这是把定义
\ref{algebra-definition-regular}、引理 \ref{algebra-lemma-finite-gl-dim-primes}
和命题 \ref{algebra-proposition-finite-gl-dim-regular} 结合起来的结果。
\end{proof}

\begin{lemma}
\label{algebra-lemma-flat-under-regular}
设 $R \to S$ 是诺特局部环的局部同态。假设 $R \to S$
平坦，且 $S$ 正则。则 $R$ 正则。
\end{lemma}

\begin{proof}
设 $\mathfrak m \subset R$ 是极大理想，令
$\kappa = R/\mathfrak m$ 为剩余域，并令 $d = \dim S$。
选择任意消解 $F_\bullet \to \kappa$，其中每个 $F_i$
都是有限自由 $R$-模。令
% P01-E0532: when d = 0 the source defines K_d using F_{-1} and F_{-2}.
$K_d = \Ker(F_{d - 1} \to F_{d - 2})$。由 $R \to S$ 的平坦性，复形
$0 \to K_d \otimes_R S \to F_{d - 1} \otimes_R S \to \ldots
\to F_0 \otimes_R S \to \kappa \otimes_R S \to 0$
仍然正合。因为 $S$ 的整体维数为 $d$，
% P01-E0533: the cited proposition gives only the upper bound; the equality is established by the preceding proposition-finite-gl-dim-regular.
参见命题 \ref{algebra-proposition-regular-finite-gl-dim}，可知
$K_d \otimes_R S$ 是有限自由 $S$-模（另见引理
\ref{algebra-lemma-independent-resolution}）。由引理
\ref{algebra-lemma-finite-projective-descends}，可知 $K_d$
是有限自由 $R$-模。因此 $\kappa$ 具有有限投射维数；由命题
\ref{algebra-proposition-finite-gl-dim-regular}，$R$ 正则。
\end{proof}

\section{Auslander--Buchsbaum 公式}
\label{algebra-section-Auslander-Buchsbaum}

\noindent
下述结果见 \cite{Auslander-Buchsbaum}。

\begin{proposition}
\label{algebra-proposition-Auslander-Buchsbaum}
设 $R$ 是诺特局部环。设 $M$ 是非零有限 $R$-模，且具有有限投射维数
$\text{pd}_R(M)$。则
$$
\text{depth}(R) = \text{pd}_R(M) + \text{depth}(M)
$$
成立。
\end{proposition}

\begin{proof}
我们对 $\text{depth}(M)$ 作归纳。最有意思的情形是
$\text{depth}(M) = 0$。在此情形下，设
$$
0 \to R^{n_e} \to R^{n_{e-1}} \to \ldots \to R^{n_0} \to M \to 0
$$
是一个极小有限自由消解，于是 $e = \text{pd}_R(M)$。
由引理 \ref{algebra-lemma-add-trivial-complex}，可以假设该复形中各映射的
所有矩阵系数都包含在 $R$ 的极大理想中。一方面，由命题
\ref{algebra-proposition-what-exact} 可知 $\text{depth}(R) \geq e$。
另一方面，把这个长正合列分解成短正合列
% P01-E0534: for e = 0 or e = 1, the displayed decomposition uses negative-indexed K terms and needs separate low-dimensional cases.
\begin{align*}
0 \to R^{n_e} \to R^{n_{e - 1}} \to K_{e - 2} \to 0,\\
0 \to K_{e - 2} \to R^{n_{e - 2}} \to K_{e - 3} \to 0,\\
\ldots,\\
0 \to K_0 \to R^{n_0} \to M \to 0
\end{align*}
并使用引理 \ref{algebra-lemma-depth-in-ses}，可得
\begin{align*}
\text{depth}(K_{e - 2}) \geq \text{depth}(R) - 1,\\
\text{depth}(K_{e - 3}) \geq \text{depth}(R) - 2,\\
\ldots,\\
\text{depth}(K_0) \geq \text{depth}(R) - (e - 1),\\
\text{depth}(M) \geq \text{depth}(R) - e
\end{align*}
又因为 $\text{depth}(M) = 0$，所以 $\text{depth}(R) \leq e$。
这就完成了 $\text{depth}(M) = 0$ 情形的证明。

\medskip\noindent
归纳步骤。若 $\text{depth}(M) > 0$，则取 $x \in \mathfrak m$，
使其同时是 $M$ 和 $R$ 上的非零因子。这是可能的，因为要么
$\text{pd}_R(M) > 0$，并且由前述命题
\ref{algebra-proposition-what-exact} 有 $\text{depth}(R) > 0$；要么
$\text{pd}_R(M) = 0$，此时 $M$ 是有限自由模，因而同样有
$\text{depth}(R) = \text{depth}(M) > 0$。
% P01-E0535: the source immediately chooses the same x a second time after already choosing it above.
于是由引理 \ref{algebra-lemma-depth-in-ses}（例如）可知
$\text{depth}(R \oplus M) > 0$，从而可以找到 $x \in \mathfrak m$，
使其同时是 $R$ 和 $M$ 上的非零因子。设
$$
0 \to R^{n_e} \to R^{n_{e-1}} \to \ldots \to R^{n_0} \to M \to 0
$$
是如上的极小消解。应用蛇引理可知
$$
0 \to (R/xR)^{n_e} \to (R/xR)^{n_{e-1}} \to \ldots \to (R/xR)^{n_0} \to
M/xM \to 0
$$
也是极小消解。因此 $\text{pd}_R(M) = \text{pd}_{R/xR}(M/xM)$。
由引理 \ref{algebra-lemma-depth-drops-by-one}，有
$\text{depth}(R/xR) = \text{depth}(R) - 1$ 和
$\text{depth}(M/xM) = \text{depth}(M) - 1$。
到此为止，所有深度都指作为 $R$-模的深度；但我们注意到
$\text{depth}_R(M/xM) = \text{depth}_{R/xR}(M/xM)$，对 $R/xR$ 也同样如此。
由归纳假设，$M/xM$ 作为 $R/xR$-模满足 Auslander--Buchsbaum 公式。
由于 $R/xR$ 和 $M/xM$ 的深度都减一，而投射维数没有改变，
结论随即得出。
\end{proof}

\section{同态与维数}
\label{algebra-section-homomorphism-dimension}

\noindent
本节汇集若干简单结果，用以联系存在环映射时各环的维数。

\begin{lemma}
\label{algebra-lemma-dimension-going-up}
设环映射 $R \to S$ 满足上升（见定义
\ref{algebra-definition-going-up-down}）或满足下降（见定义
\ref{algebra-definition-going-up-down}）。再假设
$\Spec(S) \to \Spec(R)$ 是满射。则 $\dim(R) \leq \dim(S)$。
\end{lemma}

\begin{proof}
先假设满足上升。任取 $R$ 中的素理想链
$\mathfrak p_0 \subset \mathfrak p_1 \subset \ldots
\subset \mathfrak p_e$。由满射性，可以选取映到
$\mathfrak p_0$ 的素理想 $\mathfrak q_0$。由上升，可以把它扩充为
长度为 $e$ 的素理想链，其中 $\mathfrak q_i$ 位于
$\mathfrak p_i$ 上方。因此 $\dim(S) \geq \dim(R)$。
下降情形完全相同。纯拓扑版本另见拓扑一章的引理
\ref{topology-lemma-dimension-specializations-lift}。
\end{proof}

\begin{lemma}
\label{algebra-lemma-going-up-maximal-on-top}
设 $R \to S$ 是满足上升性质的环映射，见定义
\ref{algebra-definition-going-up-down}。若
% P01-E0536: the source ends the conditional clause as a sentence before “Then”.
$\mathfrak q \subset S$ 是极大理想，则 $\mathfrak q$ 在 $R$
中的逆像也是极大理想。
\end{lemma}

\begin{proof}
显然。
\end{proof}

\begin{lemma}
\label{algebra-lemma-integral-dim-up}
设 $R \to S$ 是环映射，且 $S$ 在 $R$ 上整。则
$\dim (R) \geq \dim(S)$，并且 $\Spec(S)$ 的每个闭点都映到
$\Spec(R)$ 的闭点。
\end{lemma}

\begin{proof}
这由引理 \ref{algebra-lemma-integral-no-inclusion}、
\ref{algebra-lemma-going-up-maximal-on-top} 和各定义立即得出。
\end{proof}

\begin{lemma}
\label{algebra-lemma-integral-sub-dim-equal}
设 $R \subset S$，且 $S$ 在 $R$ 上整。则 $\dim(R) = \dim(S)$。
\end{lemma}

\begin{proof}
这是引理 \ref{algebra-lemma-integral-going-up}、
\ref{algebra-lemma-integral-overring-surjective}、
\ref{algebra-lemma-dimension-going-up} 和
\ref{algebra-lemma-integral-dim-up} 的组合。
\end{proof}

\begin{definition}
\label{algebra-definition-fibre}
设 $R \to S$ 是环映射。设 $\mathfrak q \subset S$ 是位于
$R$ 的素理想 $\mathfrak p$ 上方的素理想。所谓
{\it $\mathfrak q$ 处纤维的局部环}，是指局部环
$$
S_{\mathfrak q}/\mathfrak pS_{\mathfrak q}
=
(S/\mathfrak pS)_{\mathfrak q}
=
(S \otimes_R \kappa(\mathfrak p))_{\mathfrak q}
$$
\end{definition}

\begin{lemma}
\label{algebra-lemma-dimension-base-fibre-total}
设 $R \to S$ 是诺特环的同态。设 $\mathfrak q \subset S$
是位于素理想 $\mathfrak p$ 上方的素理想。则
$$
\dim(S_{\mathfrak q})
\leq
\dim(R_{\mathfrak p})
+
\dim(S_{\mathfrak q}/\mathfrak pS_{\mathfrak q}).
$$
\end{lemma}

\begin{proof}
我们使用命题 \ref{algebra-proposition-dimension} 对维数的刻画。
取 $\mathfrak p$ 中的元素 $x_1, \ldots, x_d$，使它们生成
$R_{\mathfrak p}$ 的一个定义理想，其中
$d = \dim(R_{\mathfrak p})$。取 $\mathfrak q$ 中的元素
$y_1, \ldots, y_e$，使它们生成
$S_{\mathfrak q}/\mathfrak pS_{\mathfrak q}$ 的一个定义理想，其中
$e = \dim(S_{\mathfrak q}/\mathfrak pS_{\mathfrak q})$。
显然，$S_{\mathfrak q}/(x_1, \ldots, x_d, y_1, \ldots, y_e)$
具有幂零极大理想。因此
$x_1, \ldots, x_d, y_1, \ldots, y_e$ 生成
$S_{\mathfrak q}$ 的一个定义理想。
\end{proof}

\begin{lemma}
\label{algebra-lemma-dimension-base-fibre-equals-total}
设 $R \to S$ 是诺特环的同态。设 $\mathfrak q \subset S$
是位于素理想 $\mathfrak p$ 上方的素理想。假设 $R \to S$
满足下降性质（例如，当 $R \to S$ 平坦时，引理
\ref{algebra-lemma-flat-going-down} 表明这一点）。则
$$
\dim(S_{\mathfrak q})
=
\dim(R_{\mathfrak p})
+
\dim(S_{\mathfrak q}/\mathfrak pS_{\mathfrak q}).
$$
\end{lemma}

\begin{proof}
由引理 \ref{algebra-lemma-dimension-base-fibre-total}，有不等式
$\dim(S_{\mathfrak q}) \leq
\dim(R_{\mathfrak p}) + \dim(S_{\mathfrak q}/\mathfrak pS_{\mathfrak q})$。
为了得到等式，选取素理想链
$\mathfrak pS \subset \mathfrak q_0 \subset \mathfrak q_1 \subset \ldots
\subset \mathfrak q_d = \mathfrak q$，其中
$d = \dim(S_{\mathfrak q}/\mathfrak pS_{\mathfrak q})$。
另一方面，选取素理想链
$\mathfrak p_0 \subset \mathfrak p_1 \subset \ldots \subset \mathfrak p_e
= \mathfrak p$，其中 $e = \dim(R_{\mathfrak p})$。
由下降定理，可以选取位于 $\mathfrak p_{e-1}$ 上方的
$\mathfrak q_{-1} \subset \mathfrak q_0$。继而可以选取位于
$\mathfrak p_{e-2}$ 上方的
$\mathfrak q_{-2} \subset \mathfrak q_{-1}$。依此归纳，最终得到长度为
$e + d$ 的链
$\mathfrak q_{-e} \subset \ldots \subset \mathfrak q_d$。
\end{proof}

\begin{lemma}
\label{algebra-lemma-flat-over-regular-with-regular-fibre}
设 $R \to S$ 是诺特局部环的局部同态。假设
\begin{enumerate}
\item $R$ 正则；
\item $S/\mathfrak m_RS$ 正则；
\item $R \to S$ 平坦。
\end{enumerate}
则 $S$ 正则。
\end{lemma}

\begin{proof}
由引理 \ref{algebra-lemma-dimension-base-fibre-equals-total}，有
$\dim(S) = \dim(R) + \dim(S/\mathfrak m_RS)$。
选取生成元 $x_1, \ldots, x_d \in \mathfrak m_R$，其中
$d = \dim(R)$；再选取 $y_1, \ldots, y_e \in \mathfrak m_S$，
使其在 $S/\mathfrak m_RS$ 中生成该环的极大理想，其中
$e = \dim(S/\mathfrak m_RS)$。于是
$x_1, \ldots, x_d, y_1, \ldots, y_e$ 生成 $S$ 的极大理想，
并且 $e + d = \dim(S)$。
\end{proof}

\noindent
下面的引理稍后将用于证明：域上的有限型环是 Cohen--Macaulay 环，
当且仅当它在某个多项式环上拟有限且平坦。它是引理
\ref{algebra-lemma-CM-over-regular-flat} 的一个部分逆命题。

\begin{lemma}
\label{algebra-lemma-finite-flat-over-regular-CM}
设 $R \to S$ 是诺特局部环的局部同态。假设 $R$ 是
Cohen--Macaulay 环。若 $S$ 在 $R$ 上有限平坦，或者 $S$ 在
$R$ 上平坦且 $\dim(S) \leq \dim(R)$，则 $S$ 是 Cohen--Macaulay 环，
并且 $\dim(R) = \dim(S)$。
\end{lemma}

\begin{proof}
设 $x_1, \ldots, x_d \in \mathfrak m_R$ 是长度为
$d = \dim(R)$ 的正则序列。由引理 \ref{algebra-lemma-flat-increases-depth}，
它在 $S$ 中的像是正则序列。因此，只要 $\dim(S) \leq d$，
$S$ 就是 Cohen--Macaulay 环。若 $S$ 在 $R$ 上有限平坦，
则由引理 \ref{algebra-lemma-integral-sub-dim-equal} 可知这一不等式成立；
第二种情形中，该不等式就是假设。
\end{proof}

\section{维数公式}
\label{algebra-section-dimension-formula}

\noindent
回顾链状（定义 \ref{algebra-definition-catenary}）和普遍链状
（定义 \ref{algebra-definition-universally-catenary}）的定义。

\begin{lemma}
\label{algebra-lemma-dimension-formula}
设 $R \to S$ 是环映射。设 $\mathfrak q$ 是 $S$ 的素理想，
位于 $R$ 的素理想 $\mathfrak p$ 上方。假设
\begin{enumerate}
\item $R$ 是诺特环；
\item $R \to S$ 是有限型的；
\item $R$、$S$ 都是整环；
\item $R \subset S$。
\end{enumerate}
则
$$
\text{高度}(\mathfrak q)
\leq
\text{高度}(\mathfrak p) + \text{trdeg}_R(S)
- \text{trdeg}_{\kappa(\mathfrak p)} \kappa(\mathfrak q)
$$
成立；若 $R$ 普遍链状，则等号成立。
\end{lemma}

\begin{proof}
设 $R \subset S' \subset S$，其中 $S'$ 是 $S$ 的有限生成
$R$-子代数。此时令 $\mathfrak q' = S' \cap \mathfrak q$。
对环映射 $R \to S'$ 和 $S' \to S$ 应用本引理，再使用域塔中
超越次数的可加性（域一章的引理
\ref{fields-lemma-transcendence-degree-tower}），即可推出
$R \to S$ 的情形。因此，可以对 $S$ 在 $R$ 上的生成元数目作归纳，
把问题约化到 $S$ 由 $R$ 上一个元素生成的情形。

\medskip\noindent
情形 I：$S = R[x]$ 是 $R$ 上的多项式代数。
此时 $\text{trdeg}_R(S) = 1$。此外，$R \to S$ 平坦，因而
$$
\dim(S_{\mathfrak q}) =
\dim(R_{\mathfrak p}) + \dim(S_{\mathfrak q}/\mathfrak pS_{\mathfrak q})
$$
，见引理 \ref{algebra-lemma-dimension-base-fibre-equals-total}。
令 $\mathfrak r = \mathfrak pS$。那么
$\text{trdeg}_{\kappa(\mathfrak p)} \kappa(\mathfrak q) = 1$
等价于 $\mathfrak q = \mathfrak r$，并且蕴含
$\dim(S_{\mathfrak q}/\mathfrak pS_{\mathfrak q}) = 0$。
同理，$\text{trdeg}_{\kappa(\mathfrak p)} \kappa(\mathfrak q) = 0$
等价于存在严格包含
$\mathfrak r \subset \mathfrak q$，而这蕴含
$\dim(S_{\mathfrak q}/\mathfrak pS_{\mathfrak q}) = 1$。
至此情形 I 得证，并且每种情况下等号都成立。

\medskip\noindent
情形 II：$S = R[x]/\mathfrak n$，其中 $\mathfrak n \not = 0$。
此时 $\text{trdeg}_R(S) = 0$。
% P01-E0537: the source omits “by” in “Denote q' ... the prime corresponding to q”.
以 $\mathfrak q' \subset R[x]$ 表示与 $\mathfrak q$ 对应的素理想。于是
$$
S_{\mathfrak q} = (R[x])_{\mathfrak q'}/\mathfrak n(R[x])_{\mathfrak q'}
$$
由上一情形，有
$\dim((R[x])_{\mathfrak q'}) =
\dim(R_{\mathfrak p}) + 1
- \text{trdeg}_{\kappa(\mathfrak p)} \kappa(\mathfrak q)$。
由于 $\mathfrak n \not = 0$，可知 $S_{\mathfrak q}$ 的维数至少减一，
见引理 \ref{algebra-lemma-one-equation}；这证明了引理中的不等式。
% P01-E0538: the source phrase “in case R is universally catenary” is missing the article “the”.
为了在 $R$ 普遍链状时证明等号，注意到
$\mathfrak n \subset R[x]$ 是高度为一的素理想，因为它对应于
$F[x]$ 中的非零素理想，其中 $F$ 是 $R$ 的分式域。因此，
$S_\mathfrak q = R[x]_{\mathfrak q'}/\mathfrak nR[x]_{\mathfrak q'}$
中的任意极大素理想链，都对应于 $R[x]$ 中 $\mathfrak q'$ 与 $(0)$
之间长度大一的极大素理想链。
% P01-E0539: the source's “with length 1 greater between q' and (0)” is grammatically misordered.
若 $R$ 普遍链状，则这些链的长度全都相同，并等于
$\mathfrak q'$ 的高度。这证明
$\dim(S_\mathfrak q) = \dim(R[x]_{\mathfrak q'}) - 1$，
进而按所需得到等号。
\end{proof}

\noindent
下面的引理说明，一般有限映射往往在余维 $1$ 处拟有限。

\begin{lemma}
\label{algebra-lemma-finite-in-codim-1}
设 $A \to B$ 是环映射。假设
\begin{enumerate}
\item $A \subset B$ 是整环的扩张；
\item 所诱导的分式域扩张是有限的；
\item $A$ 是诺特环；
\item $A \to B$ 是有限型的。
\end{enumerate}
设 $\mathfrak p \subset A$ 是高度为 $1$ 的素理想。则 $B$ 中位于
$\mathfrak p$ 上方的素理想至多只有有限多个，并且它们的高度全都为 $1$。
\end{lemma}

\begin{proof}
由维数公式（引理 \ref{algebra-lemma-dimension-formula}），对位于
$\mathfrak p$ 上方的任意素理想 $\mathfrak q$，有
$$
\dim(B_{\mathfrak q}) \leq
\dim(A_{\mathfrak p}) - \text{trdeg}_{\kappa(\mathfrak p)} \kappa(\mathfrak q).
$$
由于整环 $B_\mathfrak q$ 至少有 $2$ 个素理想，可知
$\dim(B_{\mathfrak q}) \geq 1$。因此
$\dim(B_{\mathfrak q}) = 1$，并且扩张
$\kappa(\mathfrak p) \subset \kappa(\mathfrak q)$ 是代数的。
所以，$\mathfrak q$ 给出其纤维
$\Spec(B \otimes_A \kappa(\mathfrak p))$ 的一个闭点，见引理
\ref{algebra-lemma-finite-residue-extension-closed}。
由于 $B \otimes_A \kappa(\mathfrak p)$ 是诺特环，该纤维
$\Spec(B \otimes_A \kappa(\mathfrak p))$ 是诺特拓扑空间，见引理
\ref{algebra-lemma-Noetherian-topology}。只由闭点组成的诺特拓扑空间是有限的，
例如见拓扑一章的引理 \ref{topology-lemma-Noetherian}。
\end{proof}

\section{域上有限型代数的维数}
\label{algebra-section-dimension-finite-type-algebras}

\noindent
本节计算域上多项式环的维数。我们还将证明，域上的有限型整环的维数
等于它在极大理想处各局部环的维数。它与基域上超越次数的联系将在
第 \ref{algebra-section-dimension-finite-type-algebras-reprise} 节建立。

\begin{lemma}
\label{algebra-lemma-dim-affine-space}
设 $\mathfrak m$ 是 $k[x_1, \ldots, x_n]$ 中的极大理想。
理想 $\mathfrak m$ 可由 $n$ 个元素生成。
$k[x_1, \ldots, x_n]_{\mathfrak m}$ 的维数为 $n$。
因此，$k[x_1, \ldots, x_n]_{\mathfrak m}$ 是维数为 $n$ 的正则局部环。
\end{lemma}

\begin{proof}
由 Hilbert 零点定理（定理 \ref{algebra-theorem-nullstellensatz}），
剩余域 $\kappa = \kappa(\mathfrak m)$ 是 $k$ 的有限扩张。
% P01-E0540: the source omits “by” in “Denote alpha_i ... the image of x_i”.
以 $\alpha_i \in \kappa$ 表示 $x_i$ 的像。
% P01-E0541: the source uses the ungrammatical “Denote kappa_i = ...” construction.
记 $\kappa_i = k(\alpha_1, \ldots, \alpha_i)
\subset \kappa$，其中 $i = 1, \ldots, n$，并令 $\kappa_0 = k$。
由域论可知 $\kappa_i = k[\alpha_1, \ldots, \alpha_i]$。
按如下方式归纳定义元素
$f_i \in \mathfrak m \cap k[x_1, \ldots, x_i]$：
设 $P_i(T) \in \kappa_{i-1}[T]$ 是 $\alpha_i $ 在
$\kappa_{i-1}$ 上的首一极小多项式。设
$Q_i(T) \in k[x_1, \ldots, x_{i-1}][T]$ 是 $P_i(T)$
的一个首一提升（次数相同）。令 $f_i = Q_i(x_i)$。
注意，若
$d_i = \deg_T(P_i) = \deg_T(Q_i) = \deg_{x_i}(f_i)$，则由域一章的引理
\ref{fields-lemma-multiplicativity-degrees} 和
\ref{fields-lemma-degree-minimal-polynomial}，
有 $d_1d_2\ldots d_i = [\kappa_i : k]$。

\medskip\noindent
我们断言，对所有 $i = 0, 1, \ldots, n$，都存在同构
$$
\psi_i : k[x_1, \ldots, x_i] /(f_1, \ldots, f_i) \cong \kappa_i.
$$
按构造，复合映射
$k[x_1, \ldots, x_i] \to k[x_1, \ldots, x_n] \to \kappa$
满射到 $\kappa_i$，并且 $f_1, \ldots, f_i$ 都在核中。
这给出一个满同态。我们归纳证明 $\psi_i$ 是单射。
$i = 0$ 时显然成立。假设结论对 $i$ 成立，我们来证明
$i + 1$ 的情形。环扩张
$k[x_1, \ldots, x_i]/(f_1, \ldots, f_i) \to
k[x_1, \ldots, x_{i + 1}]/(f_1, \ldots, f_{i + 1})$
是在一个域上添加 $1$ 个元素并施加一个不可约方程而生成的。
由初等域论，
$k[x_1, \ldots, x_{i + 1}]/(f_1, \ldots, f_{i + 1})$
是域，因而
% P01-E0542: the induction step proves injectivity of psi_{i+1}, but the source says psi_i.
$\psi_i$ 是单射。

\medskip\noindent
这蕴含 $\mathfrak m = (f_1, \ldots, f_n)$。此外，还可推出
$$
k[x_1, \ldots, x_n]/(f_1, \ldots, f_i)
\cong
\kappa_i[x_{i + 1}, \ldots, x_n].
$$
因此 $(f_1, \ldots, f_i)$ 是素理想。于是
$$
(0) \subset (f_1) \subset (f_1, f_2) \subset \ldots \subset
(f_1, \ldots, f_n) = \mathfrak m
$$
是长度为 $n$ 的素理想链。引理得证。
\end{proof}

\begin{proposition}
\label{algebra-proposition-finite-gl-dim-polynomial-ring}
域上的 $n$ 元多项式代数是正则环。它的整体维数为 $n$。
所有在极大理想处的局部化都是维数为 $n$ 的正则局部环。
\end{proposition}

\begin{proof}
由引理 \ref{algebra-lemma-dim-affine-space}，在极大理想处的所有局部化
$k[x_1, \ldots, x_n]_{\mathfrak m}$ 都是维数为 $n$ 的正则局部环。
因此，由引理 \ref{algebra-lemma-finite-gl-dim-finite-dim-regular} 得出结论。
\end{proof}

\begin{lemma}
\label{algebra-lemma-dimension-height-polynomial-ring}
设 $k$ 是域。设
$\mathfrak p \subset \mathfrak q \subset k[x_1, \ldots, x_n]$
是一对素理想。$\mathfrak p$ 与 $\mathfrak q$ 之间的任意极大素理想链
长度均为 $\text{高度}(\mathfrak q) - \text{高度}(\mathfrak p)$。
\end{lemma}

\begin{proof}
由命题 \ref{algebra-proposition-finite-gl-dim-polynomial-ring}，
$k[x_1, \ldots, x_n]$ 的任意局部环都是正则的。
因此，所有局部环都是 Cohen--Macaulay 环，见引理
\ref{algebra-lemma-regular-ring-CM}。极大理想处的局部环维数为 $n$，
所以 $k[x_1, \ldots, x_n]$ 中的每条极大素理想链长度均为 $n$，
见引理 \ref{algebra-lemma-maximal-chain-CM}。于是，$(0)$ 与 $\mathfrak p$
之间的每条极大素理想链长度均为 $\text{高度}(\mathfrak p)$，
例如见引理 \ref{algebra-lemma-CM-dim-formula}。把这些结论合在一起，
即可得到本引理的断言。
\end{proof}

\begin{lemma}
\label{algebra-lemma-dimension-spell-it-out}
设 $k$ 是域。设 $S$ 是有限型 $k$-代数，并且是整环。
则对 $S$ 的任意极大理想 $\mathfrak m$，都有
$\dim(S) = \dim(S_{\mathfrak m})$。换言之，每条极大素理想链的长度
都等于 $S$ 的维数。
\end{lemma}

\begin{proof}
写成 $S = k[x_1, \ldots, x_n]/\mathfrak p$。
由命题 \ref{algebra-proposition-finite-gl-dim-polynomial-ring} 和引理
\ref{algebra-lemma-dimension-height-polynomial-ring}，$S$ 中的所有极大素理想链
（它们必然终止于极大理想）的长度都是
$n - \text{高度}(\mathfrak p)$。因此，这个数既是 $S$ 的维数，
也是 $S$ 的任意极大理想 $\mathfrak m$ 所对应的
$S_{\mathfrak m}$ 的维数。
\end{proof}

\noindent
回顾我们在拓扑一章的定义 \ref{topology-definition-Krull} 中定义了
拓扑空间 $X$ 在点 $x$ 处的维数 $\dim_x(X)$。

\begin{lemma}
\label{algebra-lemma-dimension-at-a-point-finite-type-over-field}
设 $k$ 是域。设 $S$ 是有限型 $k$-代数。令 $X = \Spec(S)$。
设 $\mathfrak p \subset S$ 是素理想，并令 $x \in X$ 为对应的点。
以下各数相等：
\begin{enumerate}
\item $\dim_x(X)$；
\item $\max \dim(Z)$，其中最大值取遍所有经过 $x$ 的
$X$ 的不可约分支 $Z$；
\item $\min \dim(S_{\mathfrak m})$，其中最小值取遍所有满足
$\mathfrak p \subset \mathfrak m$ 的极大理想 $\mathfrak m$。
\end{enumerate}
\end{lemma}

\begin{proof}
设 $X = \bigcup_{i \in I} Z_i$ 是 $X$ 的不可约分支分解。
这些分支只有有限多个（见引理 \ref{algebra-lemma-obvious-Noetherian} 和
\ref{algebra-lemma-Noetherian-topology}）。令
$I' = \{i \mid x \in Z_i\}$，并令
$T = \bigcup_{i \not \in I'} Z_i$。于是 $U = X \setminus T$
是 $X$ 中包含点 $x$ 的开子集。数 (2) 是
$\max_{i \in I'} \dim(Z_i)$。对任意满足 $x \in W$ 的开集
$W \subset U$，$W$ 的不可约分支是不可约集
$W_i = Z_i \cap W$，其中 $i \in I'$，并且 $x$ 属于其中每一个。
注意，每个 $W_i$（$i \in I'$）都包含闭点，因为 $X$ 是 Jacobson 的，
见第 \ref{algebra-section-ring-jacobson} 节。由于 $W_i \subset Z_i$，
有 $\dim(W_i) \leq \dim(Z_i)$。闭点的存在结合引理
\ref{algebra-lemma-dimension-spell-it-out} 表明，开集 $W_i$ 中存在长度为
$\dim(Z_i)$ 的不可约闭子集链。因此，对任意 $i \in I'$，
都有 $\dim(W_i) = \dim(Z_i)$。故 $\dim(W)$ 等于数 (2)。
这证明 (1) $ = $ (2)。

\medskip\noindent
设 $\mathfrak m \supset \mathfrak p$ 是包含 $\mathfrak p$
的任意极大理想。令 $x_0 \in X$ 为对应的点。首先，$x_0$
包含在所有不可约分支 $Z_i$（$i \in I'$）中。
令 $\mathfrak q_i$ 表示 $S$ 中对应于不可约分支 $Z_i$ 的极小素理想。
对每个满足 $x_0 \in Z_i$ 的 $i$（这等价于
$\mathfrak m \supset \mathfrak q_i$），都有满射
$$
S_{\mathfrak m} \longrightarrow
S_\mathfrak m/\mathfrak q_i S_\mathfrak m =(S/\mathfrak q_i)_{\mathfrak m}
$$
此外，这样得到的素理想 $\mathfrak q_i S_\mathfrak m$
穷尽诺特局部环 $S_{\mathfrak m}$ 的极小素理想，见引理
\ref{algebra-lemma-irreducible-components-containing-x}。
由引理 \ref{algebra-lemma-dimension-spell-it-out}，可知
$S_{\mathfrak m}$ 的维数是经过 $x_0$ 的各 $Z_i$ 的维数中的最大值。
为完成引理的证明，只需说明可以选取 $x_0$，使
$x_0 \in Z_i \Rightarrow i \in I'$。由于 $S$ 是 Jacobson 环
（如上所见），只需说明 $V(\mathfrak p) \setminus T$
（其中 $T$ 同上）非空。这是显然的，因为它包含点 $x$
（即 $\mathfrak p$）。
\end{proof}

\begin{lemma}
\label{algebra-lemma-dimension-closed-point-finite-type-field}
设 $k$ 是域。设 $S$ 是有限型 $k$-代数。令 $X = \Spec(S)$。
设 $\mathfrak m \subset S$ 是极大理想，并令 $x \in X$
为相应的闭点。则 $\dim_x(X) = \dim(S_{\mathfrak m})$。
\end{lemma}

\begin{proof}
这是引理 \ref{algebra-lemma-dimension-at-a-point-finite-type-over-field}
的一个特例。
\end{proof}

\begin{lemma}
\label{algebra-lemma-disjoint-decomposition-CM-algebra}
设 $k$ 是域。设 $S$ 是有限型
% P01-E0543: the source omits the hyphen in “finite type k algebra”.
$k$-代数。假设 $S$ 是 Cohen--Macaulay 环。则
$\Spec(S) = \coprod T_d$ 是开闭子集 $T_d$ 的有限不交并，
其中 $T_d$ 是维数为 $d$ 的等维空间（见拓扑一章的定义
\ref{topology-definition-equidimensional}）。等价地，$S$ 是环
$S_d$（$d = 0, \ldots, \dim(S)$）的乘积，并且 $S_d$
的每个极大理想 $\mathfrak m$ 都具有高度 $d$。
\end{lemma}

\begin{proof}
两个陈述的等价性由引理 \ref{algebra-lemma-disjoint-implies-product} 得出。
设 $\mathfrak m \subset S$ 是极大理想。$S_{\mathfrak m}$
中的每条极大素理想链长度都相同，并等于
$\dim(S_{\mathfrak m})$，见引理 \ref{algebra-lemma-maximal-chain-CM}。
因此，经过与 $\mathfrak m$ 对应之点的不可约分支
% P01-E0544: the source has singular “the dimension ... all have”.
维数全都等于 $\dim(S_{\mathfrak m})$，见引理
\ref{algebra-lemma-dimension-spell-it-out}。由于 $\Spec(S)$ 是
Jacobson 拓扑空间，它的任意两个不可约分支若相交，其交集便包含闭点，
见引理 \ref{algebra-lemma-finite-type-field-Jacobson} 和
\ref{algebra-lemma-jacobson}。所以，我们已经证明任意两个相交的不可约分支
具有相同维数。再结合 $\Spec(S)$ 只有有限多个不可约分支这一事实，
即可容易地推出本引理；后者见引理 \ref{algebra-lemma-obvious-Noetherian}
和 \ref{algebra-lemma-Noetherian-topology}。
\end{proof}

\section{Noether 正规化}
\label{algebra-section-Noether-normalization}

\noindent
本节证明 Noether 正规化引理的若干变体。我们将使用的关键成分包含在
下面两个引理中。

\begin{lemma}
\label{algebra-lemma-helper}
设 $n \in \mathbf{N}$。设 $N$ 是由多重指标
$\nu = (\nu_1, \ldots, \nu_n)$ 组成的有限非空集合。
给定 $e = (e_1, \ldots, e_n)$，令
$e \cdot \nu = \sum e_i\nu_i$。则当
$e_1 \gg e_2 \gg \ldots \gg e_{n-1} \gg e_n$ 时，有：
若 $\nu, \nu' \in N$，则
$$
(e \cdot \nu = e \cdot \nu')
\Leftrightarrow
(\nu = \nu')
$$
\end{lemma}

\begin{proof}
设 $N = \{\nu_j\}$，其中
$\nu_j = (\nu_{j1}, \ldots, \nu_{jn})$。令
$A_i = \max_j \nu_{ji} - \min_j \nu_{ji}$。
% P01-E0545: “for each i” makes the displayed sufficient condition use the undefined e_0 at i = 1; the intended range starts at i = 2.
若对每个 $i$ 都有
$e_{i - 1} > A_ie_i + A_{i + 1}e_{i + 1} + \ldots + A_ne_n$，
则引理成立。事实上，假设 $e \cdot (\nu - \nu') = 0$。当
$n \ge 2$ 时，
$$
e_1(\nu_1 - \nu'_1) = \sum\nolimits_{i = 2}^n e_i(\nu'_i - \nu_i).
$$
可以假设 $(\nu_1 - \nu'_1) \ge 0$。若
$(\nu_1 - \nu'_1) > 0$，则
$$
e_1(\nu_1 - \nu'_1) \ge e_1 >
A_2e_2 + \ldots + A_ne_n \ge
\sum\nolimits_{i = 2}^n e_i|\nu'_i - \nu_i| \ge
\sum\nolimits_{i = 2}^n e_i(\nu'_i - \nu_i).
$$
这个矛盾说明 $\nu'_1 = \nu_1$。归纳可得，对
$2 \le i \le n$，都有 $\nu'_i = \nu_i$。
\end{proof}

\begin{lemma}
\label{algebra-lemma-helper-polynomial}
设 $R$ 是环。设 $g \in R[x_1, \ldots, x_n]$ 是非常数元素，
即 $g \not \in R$。当
$e_1 \gg e_2 \gg \ldots \gg e_{n-1} \gg e_n = 1$ 时，多项式
$$
g(x_1 + x_n^{e_1}, x_2 + x_n^{e_2}, \ldots, x_{n - 1} + x_n^{e_{n - 1}}, x_n)
=
ax_n^d + \text{lower order terms in }x_n
$$
，其中 $d > 0$，而 $a \in R$ 是 $g$ 的一个非零系数。
\end{lemma}

\begin{proof}
写成 $g = \sum_{\nu \in N} a_\nu x^\nu$，其中 $a_\nu \in R$
% P01-E0546: the source uses the nonstandard phrase “not zero” for a coefficient.
非零。这里 $N$ 是引理 \ref{algebra-lemma-helper} 中那样的有限多重指标集，
并且 $x^\nu = x_1^{\nu_1} \ldots x_n^{\nu_n}$。注意，
$$
(x_1 + x_n^{e_1})^{\nu_1} \ldots (x_{n-1} + x_n^{e_{n-1}})^{\nu_{n-1}}
x_n^{\nu_n}
\quad\text{is}\quad
x_n^{e_1\nu_1 + \ldots + e_{n-1}\nu_{n-1} + \nu_n}.
$$
中的首项如上所示。因此，引理由引理 \ref{algebra-lemma-helper} 得出：
它保证 $g$ 中恰有一个非零项 $a_\nu x^\nu$ 产生
$g(x_1 + x_n^{e_1}, x_2 + x_n^{e_2}, \ldots, x_{n - 1} + x_n^{e_{n - 1}},
x_n)$ 的首项；也就是说，对唯一满足 $e \cdot \nu$ 取最大值的
$\nu \in N$，有 $a = a_\nu$。
\end{proof}

\begin{lemma}
\label{algebra-lemma-one-relation}
设 $k$ 是域。设 $S = k[x_1, \ldots, x_n]/I$，其中 $I$ 是真理想。
若 $I \not = 0$，则存在
$y_1, \ldots, y_{n-1} \in k[x_1, \ldots, x_n]$，使得 $S$ 在
$k[y_1, \ldots, y_{n-1}]$ 上有限。此外，可以选取 $y_i$ 属于由
$x_1, \ldots, x_n$ 生成的 $k[x_1, \ldots, x_n]$ 的
$\mathbf{Z}$-子代数。
\end{lemma}

\begin{proof}
选取 $f \in I$，$f\not = 0$。因为 $S$ 是
$k[x_1, \ldots, x_n]/(f)$ 的商环，只需对后一个环证明本引理。
对适当的整数 $e_i$，我们将取
$y_i = x_i - x_n^{e_i}$，$i = 1, \ldots, n-1$。
何时这会奏效？只需证明
$\overline{x_n} \in k[x_1, \ldots, x_n]/(f)$ 在环
$k[y_1, \ldots, y_{n-1}]$ 上整。$\overline{x_n}$ 在此环上满足的方程是
$$
f(y_1 + x_n^{e_1}, \ldots, y_{n-1} + x_n^{e_{n-1}}, x_n) = 0.
$$
所以，只要能证明
% P01-E0547: the source has the subject–verb mismatch “there exists integers”.
存在整数 $e_i$，使上述方程关于 $x_n$ 的首项系数是 $k$ 的非零元素，
就完成了证明。例如，取
$e_1 \gg e_2 \gg \ldots \gg e_{n-1}$ 即可，见引理
\ref{algebra-lemma-helper-polynomial}。
\end{proof}

\begin{lemma}
\label{algebra-lemma-Noether-normalization}
\begin{slogan}
Noether 正规化
\end{slogan}
设 $k$ 是域。设 $S = k[x_1, \ldots, x_n]/I$，其中 $I$ 是理想。
若 $I \neq (1)$，则存在 $r\geq 0$ 和
$y_1, \ldots, y_r \in k[x_1, \ldots, x_n]$，使得：
(a) 映射 $k[y_1, \ldots, y_r] \to S$ 是单射，其中源是以
$y_1, \ldots, y_r$ 为变量的多项式环；(b) 映射
$k[y_1, \ldots, y_r] \to S$ 是有限的。在此情形下，整数 $r$
就是 $S$ 的维数。此外，可以选取 $y_i$ 属于由
$x_1, \ldots, x_n$ 生成的 $k[x_1, \ldots, x_n]$ 的
$\mathbf{Z}$-子代数。
\end{lemma}

\begin{proof}
对 $n$ 作归纳，$n = 0$ 时显然。若 $I = 0$，则取
$r = n$ 和 $y_i = x_i$。若 $I \not = 0$，则按引理
\ref{algebra-lemma-one-relation} 选取 $y_1, \ldots, y_{n-1}$。令
$S' \subset S$ 为由各 $y_i$ 的像生成的子环。由归纳假设，
可以选取 $r$ 和
$z_1, \ldots, z_r \in k[y_1, \ldots, y_{n-1}]$，使 (a)、(b)
对 $k[z_1, \ldots, z_r] \to S'$ 成立。由于 $S' \to S$
是单射且有限，可知 (a)、(b) 对
$k[z_1, \ldots, z_r] \to S$ 也成立。
% P01-E0548: “The last assertion” literally points to the Z-subalgebra choice, but the cited dimension lemma proves the preceding assertion that r = dim(S).
关于维数的断言由引理 \ref{algebra-lemma-integral-sub-dim-equal} 得出。
\end{proof}

\begin{lemma}
\label{algebra-lemma-Noether-normalization-at-point}
设 $k$ 是域。设 $S$ 是有限型
% P01-E0549: the source omits the hyphen in “finite type k algebra”.
$k$-代数，并记 $X = \Spec(S)$。设 $\mathfrak q$ 是 $S$ 的素理想，
并令 $x \in X$ 为对应的点。存在 $g \in S$、$g \not \in \mathfrak q$，
使得 $\dim(S_g) = \dim_x(X) =: d$，并且存在有限单射环映射
$k[y_1, \ldots, y_d] \to S_g$。
\end{lemma}

\begin{proof}
注意，按定义，$\dim_x(X)$ 是所有 $S_g$ 的维数中的最小值，
其中 $g \in S$、$g \not \in \mathfrak q$；也就是说，该最小值能够取到。
因此，引理由引理 \ref{algebra-lemma-Noether-normalization} 得出。
\end{proof}

\begin{lemma}
\label{algebra-lemma-refined-Noether-normalization}
设 $k$ 是域。设 $\mathfrak q \subset k[x_1, \ldots, x_n]$
是素理想。令 $r = \text{trdeg}_k\ \kappa(\mathfrak q)$。
则存在有限环映射
$\varphi : k[y_1, \ldots, y_n] \to k[x_1, \ldots, x_n]$，使得
$\varphi^{-1}(\mathfrak q) = (y_{r + 1}, \ldots, y_n)$。
\end{lemma}

\begin{proof}
对 $n$ 作归纳。$n = 0$ 时显然。假设 $n > 0$。
若 $r = n$，则 $\mathfrak q = (0)$，结论显然。
选取非零的 $f \in \mathfrak q$。当然，$f$ 不是常数。
应用如下形式的自同构后
$$
k[x_1, \ldots, x_n] \longrightarrow k[x_1, \ldots, x_n],
\quad
x_n \mapsto x_n,
\quad
x_i \mapsto x_i + x_n^{e_i}\ (i < n)
$$
，可以假设
% P01-E0550: the source says f is monic in x_n over a coefficient ring that itself contains x_n; the intended coefficient ring is k[x_1, ..., x_{n-1}].
$f$ 在 $k[x_1, \ldots, x_n]$ 上关于 $x_n$ 首一，见引理
\ref{algebra-lemma-helper-polynomial}。因此，环映射
$$
k[y_1, \ldots, y_n] \longrightarrow k[x_1, \ldots, x_n],
\quad
y_n \mapsto f,
\quad
y_i \mapsto x_i\ (i < n)
$$
是有限的。此外，按构造，
$y_n \in \mathfrak q \cap k[y_1, \ldots, y_n]$。因此
$\mathfrak q \cap k[y_1, \ldots, y_n] = \mathfrak pk[y_1, \ldots, y_n] + (y_n)$，
其中 $\mathfrak p \subset k[y_1, \ldots, y_{n - 1}]$ 是素理想。
注意，$\kappa(\mathfrak p) \subset \kappa(\mathfrak q)$ 是有限扩张，
因而 $r = \text{trdeg}_k\ \kappa(\mathfrak p)$。
对二元组 $(k[y_1, \ldots, y_{n - 1}], \mathfrak p)$ 应用归纳假设，
得到有限环映射
$k[z_1, \ldots, z_{n - 1}] \to k[y_1, \ldots, y_{n - 1}]$，使得
$\mathfrak p \cap k[z_1, \ldots, z_{n - 1}] = (z_{r + 1}, \ldots, z_{n - 1})$。
把环映射
$k[z_1, \ldots, z_{n - 1}] \to k[y_1, \ldots, y_{n - 1}]$
扩张为环映射
$k[z_1, \ldots, z_n] \to k[y_1, \ldots, y_n]$，方法是把
$z_n$ 映到 $y_n$。环映射的复合
$$
k[z_1, \ldots, z_n] \to k[y_1, \ldots, y_n] \to k[x_1, \ldots, x_n]
$$
解决了问题。
\end{proof}

\begin{lemma}
\label{algebra-lemma-Noether-normalization-over-a-domain}
设 $R \to S$ 是单射有限型环映射，并假设 $R$ 是整环。
则存在整数 $d$ 以及由单射组成的分解
$$
R \to R[y_1, \ldots, y_d] \to S' \to S
$$
，使得 $S'$ 在 $R[y_1, \ldots, y_d]$ 上有限，并且对某个
非零的 $f \in R$，有 $S'_f \cong S_f$。
\end{lemma}

\begin{proof}
选取 $x_1, \ldots, x_n \in S$，使其在 $R$ 上生成 $S$。
令 $K$ 为 $R$ 的分式域，并令 $S_K = S \otimes_R K$。
由引理 \ref{algebra-lemma-Noether-normalization}，可以找到
$y_1, \ldots, y_d \in S$，使
$K[y_1, \ldots, y_d] \to S_K$ 是有限单射环映射。
注意，$y_i \in S$，因为可以选取
% P01-E0551: the source switches from y_i to y_j in the explanation without a change of quantified index.
$y_j$ 属于由 $x_1, \ldots, x_n$ 生成的 $\mathbf{Z}$-代数。
由于有限环映射是整的（见引理 \ref{algebra-lemma-finite-is-integral}），
可以找到首一多项式 $P_i \in K[y_1, \ldots, y_d][T]$，使得
$P_i(x_i) = 0$ 在 $S_K$ 中成立。取非零元素 $f \in R$，使得对所有
$i$ 都有 $fP_i \in R[y_1, \ldots, y_d][T]$。于是
$fP_i(x_i)$ 在 $S_K$ 中映到零。因此，把 $f$ 换成 $R$
的另一个非零元素后，还可以假设 $fP_i(x_i)$ 在 $S$ 中为零。
令 $x_i' = fx_i$，并令 $S' \subset S$ 为由
$y_1, \ldots, y_d$ 和 $x'_1, \ldots, x'_n$ 生成的
$R$-子代数。注意，$x'_i$ 在 $R[y_1, \ldots, y_d]$ 上整，
因为 $Q_i(x_i') = 0$，其中
$Q_i = f^{\deg_T(P_i)}P_i(T/f)$；按 $f$ 的选取，这是系数属于
$R[y_1, \ldots, y_d]$ 的关于 $T$ 的首一多项式。因此，由引理
\ref{algebra-lemma-characterize-finite-in-terms-of-integral}，
$R[y_1, \ldots, y_d] \subset S'$ 是有限的。
因为 $S' \subset S$，所以 $S'_f \subset S_f$（局部化是正合的）。
另一方面，$S'_f$ 中的元素 $x_i = x'_i/f$ 在 $R_f$ 上生成
$S_f$，故 $S'_f \to S_f$ 是满射。因此
$S'_f \cong S_f$，证明完成。
\end{proof}

\section{域上有限型代数的维数（再论）}
\label{algebra-section-dimension-finite-type-algebras-reprise}

\noindent
本节是第 \ref{algebra-section-dimension-finite-type-algebras} 节的延续。
我们将在本节建立域上的有限型整环之维数与其在基域上的超越次数之间的联系。

\begin{lemma}
\label{algebra-lemma-dimension-prime-polynomial-ring}
设 $k$ 是域。设 $S$ 是有限型
% P01-E0552: the source omits the hyphen in “finite type k algebra”.
$k$-代数，并且是整环。设 $K$ 是 $S$ 的分式域。
设 $r = \text{trdeg}(K/k)$ 是 $K$ 在 $k$ 上的超越次数。
则 $\dim(S) = r$。此外，$S$ 在每个极大理想处的局部环维数均为 $r$。
\end{lemma}

\begin{proof}
可以写成 $S = k[x_1, \ldots, x_n]/\mathfrak p$。
由引理 \ref{algebra-lemma-dimension-height-polynomial-ring}，$S$
在各极大理想处的所有局部环维数相同。应用引理
\ref{algebra-lemma-Noether-normalization}，得到有限单射环映射
$$
k[y_1, \ldots, y_d] \to S
$$
，其中 $d = \dim(S)$。显然，
$k(y_1, \ldots, y_d) \subset K$ 是有限扩张，结论得证。
\end{proof}

\begin{lemma}
\label{algebra-lemma-tr-deg-specialization}
设 $k$ 是域。设 $S$ 是有限型 $k$-代数。设
$\mathfrak q \subset \mathfrak q' \subset S$ 是不同的素理想。则
$\text{trdeg}_k\ \kappa(\mathfrak q') < \text{trdeg}_k\ \kappa(\mathfrak q)$。
\end{lemma}

\begin{proof}
由引理 \ref{algebra-lemma-dimension-prime-polynomial-ring}，有
$\dim V(\mathfrak q) = \text{trdeg}_k\ \kappa(\mathfrak q)$，
对 $\mathfrak q'$ 也同样成立。因此，严格包含
$V(\mathfrak q') \subset V(\mathfrak q)$ 蕴含维数的严格不等式，
从而得到结论。
\end{proof}

\noindent
下面的引理推广了引理
\ref{algebra-lemma-dimension-closed-point-finite-type-field}。

\begin{lemma}
\label{algebra-lemma-dimension-at-a-point-finite-type-field}
设 $k$ 是域。设 $S$ 是有限型
% P01-E0553: the source omits the hyphen in “finite type k algebra”.
$k$-代数。令 $X = \Spec(S)$。设 $\mathfrak p \subset S$ 是素理想，
并令 $x \in X$ 为对应的点。则
$$
\dim_x(X) = \dim(S_{\mathfrak p}) + \text{trdeg}_k\ \kappa(\mathfrak p).
$$
\end{lemma}

\begin{proof}
由引理 \ref{algebra-lemma-dimension-prime-polynomial-ring}，可知
$r = \text{trdeg}_k\ \kappa(\mathfrak p)$ 等于
$V(\mathfrak p)$ 的维数。在 $S$ 中任取一条以 $\mathfrak p$ 开始的
极大素理想链
$\mathfrak p \subset \mathfrak p_1 \subset \ldots \subset \mathfrak p_r$。
由引理 \ref{algebra-lemma-dimension-spell-it-out}，它的长度为 $r$。
设 $\mathfrak q_j$（$j \in J$）为 $S$ 中包含在 $\mathfrak p$
内的极小素理想。按规则
$\mathfrak q_j \mapsto \mathfrak q_jS_{\mathfrak p}$，它们与
$S_{\mathfrak p}$ 中的极小素理想成 $1-1$ 对应。
由引理 \ref{algebra-lemma-dimension-at-a-point-finite-type-over-field}，
可知 $\dim_x(X)$ 等于各环 $S/\mathfrak q_j$ 的维数中的最大值。
对每个 $j$，选取极大素理想链
$\mathfrak q_j \subset \mathfrak p'_1 \subset \ldots \subset \mathfrak p'_{s(j)}
= \mathfrak p$。于是
$\dim(S_{\mathfrak p}) = \max_{j \in J} s(j)$。
现在，每条链
$$
\mathfrak q_j \subset \mathfrak p'_1 \subset \ldots \subset
\mathfrak p'_{s(j)} = \mathfrak p \subset
\mathfrak p_1 \subset \ldots \subset \mathfrak p_r
$$
都是 $S/\mathfrak q_j$ 中的极大链，并且由前述讨论有
$\dim_x(X) = \max_{j \in J} r + s(j)$。引理得证。
\end{proof}

\noindent
下面的引理说明，一个有限型谱在另一个有限型谱中的余维等于高度之差。

\begin{lemma}
\label{algebra-lemma-codimension}
设 $k$ 是域。设 $S' \to S$ 是有限型
% P01-E0554: the source omits the hyphen in “finite type k algebras”.
$k$-代数之间的满射。设 $\mathfrak p \subset S$ 是素理想，
并设 $\mathfrak p'$ 是 $S'$ 中对应的素理想。令
$X = \Spec(S)$、相应地 $X' = \Spec(S')$，并令
$x \in X$、相应地 $x'\in X'$ 为分别对应于
$\mathfrak p$、相应地 $\mathfrak p'$ 的点。则
$$
\dim_{x'} X' - \dim_x X =
\text{高度}(\mathfrak p') - \text{高度}(\mathfrak p).
$$
\end{lemma}

\begin{proof}
这由引理 \ref{algebra-lemma-dimension-at-a-point-finite-type-field} 立即得出。
\end{proof}

\begin{lemma}
\label{algebra-lemma-dimension-preserved-field-extension}
设 $k$ 是域。设 $S$ 是有限型 $k$-代数。设 $K/k$ 是域扩张。
则 $\dim(S) = \dim(K \otimes_k S)$。
\end{lemma}

\begin{proof}
由引理 \ref{algebra-lemma-Noether-normalization}，存在有限单射映射
$k[y_1, \ldots, y_d] \to S$，其中 $d = \dim(S)$。
由于 $K$ 在 $k$ 上平坦，还得到有限单射映射
$K[y_1, \ldots, y_d] \to K \otimes_k S$。
结论由引理 \ref{algebra-lemma-integral-sub-dim-equal} 得出。
\end{proof}

\begin{lemma}
\label{algebra-lemma-dimension-at-a-point-preserved-field-extension}
设 $k$ 是域。设 $S$ 是有限型 $k$-代数。令 $X = \Spec(S)$。
设 $K/k$ 是域扩张。令 $S_K = K \otimes_k S$，并令
$X_K = \Spec(S_K)$。设 $\mathfrak q \subset S$ 是对应于
$x \in X$ 的素理想，并设 $\mathfrak q_K \subset S_K$ 是位于
$\mathfrak q$ 上方且对应于 $x_K \in X_K$ 的素理想。则
$\dim_x X = \dim_{x_K} X_K$。
\end{lemma}

\begin{proof}
选取表示 $S = k[x_1, \ldots, x_n]/I$。这给出表示
$K \otimes_k S = K[x_1, \ldots, x_n]/(K \otimes_k I)$。
设 $\mathfrak q_K' \subset K[x_1, \ldots, x_n]$、相应地
$\mathfrak q' \subset k[x_1, \ldots, x_n]$ 为对应的素理想。
考虑下列诺特局部环的交换图
$$
\xymatrix{
K[x_1, \ldots, x_n]_{\mathfrak q_K'} \ar[r] &
(K \otimes_k S)_{\mathfrak q_K} \\
k[x_1, \ldots, x_n]_{\mathfrak q'} \ar[r] \ar[u] &
S_{\mathfrak q} \ar[u]
}
$$
两个竖直箭头都是平坦的，因为它们分别是平坦环映射
$S \to S_K$ 和
$k[x_1, \ldots, x_n] \to K[x_1, \ldots, x_n]$ 的局部化。
此外，两个竖直箭头具有相同的纤维环。因此，由引理
\ref{algebra-lemma-dimension-base-fibre-equals-total} 可知
$\text{高度}(\mathfrak q') - \text{高度}(\mathfrak q)
= \text{高度}(\mathfrak q_K') - \text{高度}(\mathfrak q_K)$。
% P01-E0555: the source omits “by” in the two-object “Denote ... the points” construction.
以 $x' \in X' = \Spec(k[x_1, \ldots, x_n])$
和 $x'_K \in X'_K = \Spec(K[x_1, \ldots, x_n])$
表示分别对应于 $\mathfrak q'$ 和 $\mathfrak q_K'$ 的点。
由引理 \ref{algebra-lemma-codimension} 和上面证明的结论，有
\begin{eqnarray*}
n - \dim_x X & = & \dim_{x'} X' - \dim_x X \\
& = & \text{高度}(\mathfrak q') - \text{高度}(\mathfrak q) \\
& = & \text{高度}(\mathfrak q_K') - \text{高度}(\mathfrak q_K) \\
& = & \dim_{x'_K} X'_K - \dim_{x_K} X_K \\
& = & n - \dim_{x_K} X_K
\end{eqnarray*}
，引理随即得证。
\end{proof}

\begin{lemma}
\label{algebra-lemma-inequalities-under-field-extension}
设 $k$ 是域。设 $S$ 是有限型 $k$-代数。设 $K/k$ 是域扩张。
令 $S_K = K \otimes_k S$。设 $\mathfrak q \subset S$ 是素理想，
并设 $\mathfrak q_K \subset S_K$ 是位于 $\mathfrak q$ 上方的素理想。
则
$$
\dim (S_K \otimes_S \kappa(\mathfrak q))_{\mathfrak q_K} =
\dim (S_K)_{\mathfrak q_K} - \dim S_\mathfrak q =
\text{trdeg}_k \kappa(\mathfrak q) - \text{trdeg}_K \kappa(\mathfrak q_K)
$$
此外，给定 $\mathfrak q$，总可以选取 $\mathfrak q_K$，使得上面的数为零。
\end{lemma}

\begin{proof}
注意，$S_\mathfrak q \to (S_K)_{\mathfrak q_K}$ 是诺特局部环的平坦局部同态，
其特殊纤维为
$(S_K \otimes_S \kappa(\mathfrak q))_{\mathfrak q_K}$。因此，
% P01-E0556: the source sentence “Hence the first equality by Lemma ...” is missing its finite verb.
由引理 \ref{algebra-lemma-dimension-base-fibre-equals-total} 得到第一个等式。
第二个等式来自以下事实：采用引理
\ref{algebra-lemma-dimension-at-a-point-preserved-field-extension} 中的记号，有
$\dim_x X = \dim_{x_K} X_K$；并且由引理
\ref{algebra-lemma-dimension-at-a-point-finite-type-field}，有
$\dim_x X = \dim S_\mathfrak q + \text{trdeg}_k \kappa(\mathfrak q)$，
对 $\dim_{x_K} X_K$ 也同样成立。
若选取 $\mathfrak q_K$ 为 $\mathfrak q S_K$ 上的极小素理想，
则纤维环的维数为零。
\end{proof}

\section{域上分次代数的维数}
\label{algebra-section-dimension-graded}

\noindent
下面给出一个基本结果。

\begin{lemma}
\label{algebra-lemma-dimension-graded}
设 $k$ 是域。设 $S$ 是分次 $k$-代数，并且在 $k$ 上由有限多个
次数为 $1$ 的元素生成。假设 $S_0 = k$。设
$P(T) \in \mathbf{Q}[T]$ 是满足如下性质的多项式：
对所有 $d \gg 0$，都有 $\dim(S_d) = P(d)$。见命题
\ref{algebra-proposition-graded-hilbert-polynomial}。则
\begin{enumerate}
\item 无关理想 $S_{+}$ 是一个极大理想 $\mathfrak m$。
\item $S$ 的任意极小素理想都是齐次理想，并包含于
$S_{+} = \mathfrak m$。
\item 有 $\dim(S) = \deg(P) + 1 = \dim_x\Spec(S)$
（约定 $\deg(0) = -1$），其中 $x$ 是对应于极大理想
$S_{+} = \mathfrak m$ 的点。
\item 局部环 $R = S_{\mathfrak m}$ 的 Hilbert 函数等于 $S$
的 Hilbert 函数。
\end{enumerate}
\end{lemma}

\begin{proof}
第一个陈述显然。第二个陈述由引理
\ref{algebra-lemma-graded-ring-minimal-prime} 得出。由 (2)，每个不可约分支
都经过 $x$。因此，由引理
\ref{algebra-lemma-dimension-at-a-point-finite-type-over-field}，有
$\dim(S) = \dim_x\Spec(S) = \dim(S_\mathfrak m)$。由于
$\mathfrak m^d/\mathfrak m^{d + 1} \cong
\mathfrak m^dS_\mathfrak m/\mathfrak m^{d + 1}S_\mathfrak m$，
可知局部环 $S_\mathfrak m$ 的 Hilbert 函数等于 $S$ 的 Hilbert 函数，
这就是 (4)。最后，由命题 \ref{algebra-proposition-dimension}
得到 (3) 中的最后一个等式。
\end{proof}

\section{一般平坦性}
\label{algebra-section-generic-flatness}

\noindent
这个结果大致是说，整环上的有限型代数在把该整环的一个元素可逆化后会变得平坦。
这个结果有若干版本，强度依次递增。

\begin{lemma}
\label{algebra-lemma-generic-flatness-Noetherian}
设 $R \to S$ 是环映射。
设 $M$ 是 $S$-模。假设
\begin{enumerate}
\item $R$ 是诺特环，
\item $R$ 是整环，
\item $R \to S$ 是有限型的，并且
\item $M$ 是有限型 $S$-模。
\end{enumerate}
则存在非零元 $f \in R$，使得
$M_f$ 是自由 $R_f$-模。
\end{lemma}

\begin{proof}
设 $K$ 为 $R$ 的分式域。置 $S_K = K \otimes_R S$。这是 $K$ 上的
有限型代数。我们对 $d = \dim(S_K)$ 作归纳（这是有限的，例如可由
Noether 正规化得到，见第 \ref{algebra-section-Noether-normalization} 节）。
固定 $d \geq 0$。假设在所有满足 $\dim(S_K) < d$ 的情形中引理均已成立。

\medskip\noindent
设给定引理中的 $R \to S$ 和 $M$，且 $\dim(S_K) = d$。由引理
\ref{algebra-lemma-filter-Noetherian-module}，存在滤过
$0 \subset M_1 \subset M_2 \subset \ldots \subset M_n = M$，
使得 $M_i/M_{i - 1}$ 同构于 $S/\mathfrak q$，其中 $\mathfrak q$
是 $S$ 的某个素理想。注意
$\dim((S/\mathfrak q)_K) \leq \dim(S_K)$。还要注意，自由模的扩张仍是自由模
（见基本概念 \ref{algebra-item-extension-free}）。因此可以假设 $M = S$，并且
$S$ 是 $R$ 上的有限型整环。

\medskip\noindent
若 $R \to S$ 有非平凡的核，则在该核中取非零元 $f \in R$。
此时 $S_f = 0$，引理成立。（这实际上是 $d = -1$ 的情形，也是归纳的起点。）
所以可以假设 $R \to S$ 是诺特整环之间的有限型扩张。

\medskip\noindent
应用引理 \ref{algebra-lemma-Noether-normalization-over-a-domain}，并把 $R$
替换为 $R_f$（其中 $f$ 如该引理所述），得到分解
$$
R \subset R[y_1, \ldots, y_d] \subset S
$$
其中第二个扩张是有限的。选取 $z_1, \ldots, z_r \in S$，使其构成
$S$ 的分式域在 $R[y_1, \ldots, y_d]$ 的分式域上的一组基。这给出短正合列
$$
0 \to
R[y_1, \ldots, y_d]^{\oplus r} \xrightarrow{(z_1, \ldots, z_r)}
S \to N \to 0
$$
由构造，$N$ 是有限 $R[y_1, \ldots, y_d]$-模，其支集不包含
$\Spec(R[y_1, \ldots, y_d])$ 的泛点 $(0)$。由引理
\ref{algebra-lemma-support-closed}，存在非零元 $g \in R[y_1, \ldots, y_d]$
消去 $N$，所以可以把 $N$ 视为
$S' = R[y_1, \ldots, y_d]/(g)$ 上的有限模。因为 $\dim(S'_K) < d$，
归纳假设给出非零元 $f \in R$，使得 $N_f$ 是自由 $R_f$-模。又因为
% P01-E0557: the source clause “is free also we conclude” is grammatically malformed.
$(R[y_1, \ldots, y_d])_f \cong R_f[y_1, \ldots, y_d]$ 是自由的，
由前述自由模的扩张仍自由这一事实可知结论成立。
\end{proof}

\begin{lemma}
\label{algebra-lemma-generic-flatness-finitely-presented}
\begin{slogan}
一般自由性。
\end{slogan}
设 $R \to S$ 是环映射。
设 $M$ 是 $S$-模。假设
\begin{enumerate}
\item $R$ 是整环，
\item $R \to S$ 是有限表示的，并且
\item $M$ 是有限表示 $S$-模。
\end{enumerate}
则存在非零元 $f \in R$，使得
$M_f$ 是自由 $R_f$-模。
\end{lemma}

\begin{proof}
写成 $S = R[x_1, \ldots, x_n]/(g_1, \ldots, g_m)$。
% P01-E0558: the source designation “denote ... its image” omits “by”.
对 $g \in R[x_1, \ldots, x_n]$，以 $\overline{g}$ 表示它在 $S$ 中的像。
可以写成 $M = S^{\oplus t}/\sum Sn_i$，其中
$n_i \in S^{\oplus t}$。把它们写成
$n_i = (\overline{g}_{i1}, \ldots, \overline{g}_{it})$，其中
$g_{ij} \in R[x_1, \ldots, x_n]$。设 $R_0 \subset R$ 是由所有元素
$g_i, g_{ij} \in R[x_1, \ldots, x_n]$ 的全部系数生成的子环。定义
$S_0 = R_0[x_1, \ldots, x_n]/(g_1, \ldots, g_m)$，并定义
$M_0 = S_0^{\oplus t}/\sum S_0n_i$。于是 $R_0$ 是
$\mathbf{Z}$ 上的有限型整环，因而是诺特环（见引理
\ref{algebra-lemma-Noetherian-permanence}）。此外，通过单射 $R_0 \to R$，有
$S \cong R \otimes_{R_0} S_0$ 和 $M \cong R \otimes_{R_0} M_0$。
应用引理 \ref{algebra-lemma-generic-flatness-Noetherian}，得到非零元
$f \in R_0$，使得 $(M_0)_f$ 是自由 $(R_0)_f$-模。因此
$M_f = R_f \otimes_{(R_0)_f} (M_0)_f$ 是自由 $R_f$-模。
\end{proof}

\begin{lemma}
\label{algebra-lemma-generic-flatness}
设 $R \to S$ 是环映射。
设 $M$ 是 $S$-模。假设
\begin{enumerate}
\item $R$ 是整环，
\item $R \to S$ 是有限型的，并且
\item $M$ 是有限型 $S$-模。
\end{enumerate}
则存在非零元 $f \in R$，使得
\begin{enumerate}
\item[(a)] $M_f$ 和 $S_f$ 作为 $R_f$-模都是自由的，并且
\item[(b)] $S_f$ 是有限表示 $R_f$-代数，而 $M_f$ 是有限表示
$S_f$-模。
\end{enumerate}
\end{lemma}

\begin{proof}
我们先对 $S = R[x_1, \ldots, x_n]$ 证明引理，再由此推出一般情形。

\medskip\noindent
假设 $S = R[x_1, \ldots, x_n]$。
选取生成 $M$ 的元素 $m_1, \ldots, m_t$。这给出短正合列
$$
0 \to N \to S^{\oplus t} \xrightarrow{(m_1, \ldots, m_t)} M \to 0.
$$
% P01-E0559: the source twice uses “Denote” without the required preposition “by”.
以 $K$ 表示 $R$ 的分式域。记
$S_K = K \otimes_R S = K[x_1, \ldots, x_n]$，并类似地记
$N_K = K \otimes_R N$、$M_K = K \otimes_R M$。
由于 $R \to K$ 平坦，该序列在与 $K$ 张量后仍正合。由于
$S_K = K[x_1, \ldots, x_n]$ 是诺特环（见引理
\ref{algebra-lemma-Noetherian-permanence}），可以找到有限多个生成它的元素
$n'_1, \ldots, n'_s \in N_K$。选取 $n_1, \ldots, n_r \in N$，
使得对某些 $a_{ij} \in K$ 有 $n'_i = \sum a_{ij}n_j$。置
$$
M' = S^{\oplus t}/\sum\nolimits_{i = 1, \ldots, r} Sn_i
$$
由构造，$M'$ 是有限表示 $S$-模，并且存在满射 $M' \to M$，
它诱导同构 $M'_K \cong M_K$。可以对 $R \to S$ 和 $M'$
应用引理 \ref{algebra-lemma-generic-flatness-finitely-presented}，得到
$f \in R$，使得 $M'_f$ 是自由 $R_f$-模。因此 $M'_f \to M_f$
是整环 $R_f$ 上的模满射，其源是自由模，并且与 $K$ 张量后成为同构。
由于 $M'_f$ 自由，并且 $R_f$ 是分式域为 $K$ 的整环，故
$M'_f \subset M'_K$，从而该映射是单射。因此 $M'_f \to M_f$
是同构，结论得证。

\medskip\noindent
在一般情形中，选取满射 $R[x_1, \ldots, x_n] \to S$。
把 $S$ 和 $M$ 都看作有限 $R[x_1, \ldots, x_n]$-模。
由上面证明的特殊情形，存在非零元 $f \in R$，使得 $S_f$ 和 $M_f$
作为 $R_f$-模都是自由的，并且作为 $R_f[x_1, \ldots, x_n]$-模
都是有限表示的。这显然蕴含 $S_f$ 是有限表示 $R_f$-代数，且
$M_f$ 是有限表示 $S_f$-模。
\end{proof}

\noindent
设 $R \to S$ 是环映射。设 $M$ 是 $S$-模。考虑元素 $f \in R$
所满足的下列条件：
\begin{equation}
\label{algebra-equation-flat-and-finitely-presented}
\left\{
\begin{matrix}
S_f & \text{is of finite presentation over }R_f\\
M_f & \text{is of finite presentation as }S_f\text{-模}\\
S_f, M_f & \text{are free as }R_f\text{-模}
\end{matrix}
\right.
\end{equation}
我们定义
\begin{equation}
\label{algebra-equation-good-locus}
U(R \to S, M)
=
\bigcup\nolimits_{f \in
R\text{ 其中 }(\ref{algebra-equation-flat-and-finitely-presented})}
D(f)
\end{equation}
这是 $\Spec(R)$ 的一个开子集。

\begin{lemma}
\label{algebra-lemma-generic-flatness-locus-extension}
设 $R \to S$ 是环映射。
设 $0 \to M_1 \to M_2 \to M_3 \to 0$ 是 $S$-模的短正合列。
则
$$
U(R \to S, M_1) \cap U(R \to S, M_3) \subset U(R \to S, M_2).
$$
\end{lemma}

\begin{proof}
设 $u \in U(R \to S, M_1) \cap U(R \to S, M_3)$。选取
$f_1, f_3 \in R$，使得 $u \in D(f_1)$、$u \in D(f_3)$，并且
(\ref{algebra-equation-flat-and-finitely-presented}) 对 $f_1$ 和 $M_1$
以及对 $f_3$ 和 $M_3$ 都成立。令 $f = f_1f_3$。于是
$u \in D(f)$，并且 (\ref{algebra-equation-flat-and-finitely-presented})
对 $f$ 以及 $M_1$、$M_3$ 都成立。自由模的扩张仍是自由模，
而有限表示模的扩张仍是有限表示模（引理 \ref{algebra-lemma-extension}）。
因此 (\ref{algebra-equation-flat-and-finitely-presented}) 对 $f$ 和 $M_2$
成立。故 $u \in U(R \to S, M_2)$，结论得证。
\end{proof}

\begin{lemma}
\label{algebra-lemma-generic-flatness-locus-localize}
设 $R \to S$ 是环映射。
设 $M$ 是 $S$-模。设 $f \in R$。
在等同 $\Spec(R_f) = D(f)$ 之下，有
$U(R_f \to S_f, M_f) = D(f) \cap U(R \to S, M)$。
\end{lemma}

\begin{proof}
假设 $u \in U(R_f \to S_f, M_f)$。则存在元素 $g \in R_f$，
使得 $u \in D(g)$，并且 (\ref{algebra-equation-flat-and-finitely-presented})
对二元组 $((R_f)_g \to (S_f)_g, (M_f)_g)$ 成立。把
$g$ 写成 $g = a/f^n$，其中 $a \in R$。令 $h = af$。则
$R_h = (R_f)_g$、$S_h = (S_f)_g$ 且 $M_h = (M_f)_g$。
此外 $u \in D(h)$。故 $u \in U(R \to S, M)$。
反过来，假设 $u \in D(f) \cap U(R \to S, M)$。则存在元素
$g \in R$，使得 $u \in D(g)$，并且
(\ref{algebra-equation-flat-and-finitely-presented}) 对二元组
$(R_g \to S_g, M_g)$ 成立。于是显然，
(\ref{algebra-equation-flat-and-finitely-presented}) 对二元组
$(R_{fg} \to S_{fg}, M_{fg}) = ((R_f)_g \to (S_f)_g, (M_f)_g)$
也成立。故 $u \in U(R_f \to S_f, M_f)$，结论得证。
\end{proof}

\begin{lemma}
\label{algebra-lemma-generic-flatness-locus-reduce}
设 $R \to S$ 是环映射。
设 $M$ 是 $S$-模。设 $U \subset \Spec(R)$ 是稠密开集。
假设存在由开集组成的覆盖 $U = \bigcup_{i \in I} D(f_i)$，
使得对每个 $i \in I$，$U(R_{f_i} \to S_{f_i}, M_{f_i})$
在 $D(f_i)$ 中稠密。则 $U(R \to S, M)$ 在 $\Spec(R)$ 中稠密。
\end{lemma}

\begin{proof}
由引理 \ref{algebra-lemma-generic-flatness-locus-localize}，这纯粹是一个拓扑陈述。
事实上，由该引理可知，对每个 $i \in I$，
$U(R \to S, M) \cap D(f_i)$ 在 $D(f_i)$ 中稠密。
由拓扑学，引理 \ref{topology-lemma-nowhere-dense-local}，可知
$U(R \to S, M) \cap U$ 在 $U$ 中稠密。又因为 $U$ 在
$\Spec(R)$ 中稠密，所以 $U(R \to S, M)$ 在 $\Spec(R)$ 中稠密。
\end{proof}

\begin{lemma}
\label{algebra-lemma-generic-flatness-reduced}
设 $R \to S$ 是环映射。设 $M$ 是 $S$-模。假设
\begin{enumerate}
\item $R \to S$ 是有限型的，
\item $M$ 是有限 $S$-模，并且
\item $R$ 是约化环。
\end{enumerate}
则存在子集 $U \subset \Spec(R)$，使得
\begin{enumerate}
\item $U$ 是 $\Spec(R)$ 中的稠密开集；
\item 对每个 $u \in U$，存在 $f \in R$，使得
$u \in D(f) \subset U$，并且
\begin{enumerate}
\item $M_f$ 和 $S_f$ 在 $R_f$ 上自由，
\item $S_f$ 是有限表示 $R_f$-代数，并且
\item $M_f$ 是有限表示 $S_f$-模。
\end{enumerate}
\end{enumerate}
\end{lemma}

\begin{proof}
注意，该引理等价于开集 $U(R \to S, M)$ 在 $\Spec(R)$ 中稠密；
见方程 (\ref{algebra-equation-good-locus})。我们先对
$S = R[x_1, \ldots, x_n]$ 证明引理，再由此推出一般情形。

\medskip\noindent
先证明 $S = R[x_1, \ldots, x_n]$ 且 $M$ 是 $S$ 上任意有限模的情形。
注意，此时 $S_f = R_f[x_1, \ldots, x_n]$ 在 $R_f$ 上自由且有限表示，
所以无须考虑关于 $S$ 的条件，只须考虑关于 $M$ 的条件。
我们对 $n$ 作归纳。

\medskip\noindent
存在有限滤过
$$
0 \subset M_1 \subset M_2 \subset \ldots \subset M_t = M
$$
使得 $M_i/M_{i - 1} \cong S/J_i$，其中 $J_i \subset S$ 是某个理想，见引理
\ref{algebra-lemma-trivial-filter-finite-module}。由于有限多个稠密开集的交仍是稠密开集，
由引理 \ref{algebra-lemma-generic-flatness-locus-extension} 可知，只需对每个模
% P01-E0560: the frozen source says R/J_i although J_i is an ideal of S and the filtration factors are S/J_i.
$R/J_i$ 证明引理。因此可以假设 $M = S/J$，其中 $J$ 是
$S = R[x_1, \ldots, x_n]$ 的一个理想。

\medskip\noindent
设 $I \subset R$ 是由 $J$ 中元素的系数生成的理想。令
$U_1 = \Spec(R) \setminus V(I)$，并令
$$
U_2 = \Spec(R) \setminus \overline{U_1}.
$$
显然 $U = U_1 \cup U_2$ 在 $\Spec(R)$ 中稠密。设 $f \in R$
满足 (a) $D(f) \subset U_1$ 或 (b) $D(f) \subset U_2$。
若对每个这样的 $f$，引理都对二元组
$(R_f \to R_f[x_1, \ldots, x_n], M_f)$ 成立，则由引理
\ref{algebra-lemma-generic-flatness-locus-reduce} 可知
$U(R \to S, M)$ 在 $\Spec(R)$ 中稠密。因此可以假设
(a) $I = R$，或 (b) $V(I) = \Spec(R)$。

\medskip\noindent
在情形 (b) 中，由于 $R$ 约化，实际上有 $I = 0$！因此
$J = 0$ 且 $M = S$，引理在此情形中成立。

\medskip\noindent
在情形 (a) 中还需多做一点工作。注意，$I$ 的每个元素实际上都是
$J$ 中某个元素的某个单项式的系数，因为 $J$ 中元素的系数之集构成一个理想
（细节略）。所以可以找到元素
$$
g = \sum\nolimits_{K \in E} a_K x^K \in J
$$
其中 $E$ 是多重指标 $K = (k_1, \ldots, k_n)$ 的有限集合，且至少有一个系数
$a_{K_0}$ 是 $R$ 中的单位。事实上，由于在情形 (a) 中 $1 \in I$，
可以找到一个系数等于 $1$ 的这种元素。令
$m = \#\{K \in E \mid a_K \text{ is not a unit}\}$。
注意 $0 \leq m \leq \# E - 1$。我们对 $m$ 作归纳。

\medskip\noindent
情形 $m = 0$。此时 $g$ 的所有系数 $a_K$（$K \in E$）都是单位，且
$E \not = \emptyset$。若 $E = \{K_0\}$ 是单元素集，并且
$K_0 = (0, \ldots, 0)$，则 $g$ 是单位且 $J = S$，所以结论当然成立。
（特别地，当 $n = 0$ 时会出现这种情况，它给出了关于 $n$ 的归纳起点。）
若并非 $E = \{(0, \ldots, 0)\}$，则至少有一个 $K$ 不等于
$(0, \ldots, 0)$，亦即 $g \not \in R$。此时采用 Noether 正规化的惯用技巧。
具体地，考虑
$$
G(y_1, \ldots, y_n) =
g(y_1 + y_n^{e_1}, y_2 + y_n^{e_2}, \ldots, y_{n - 1} + y_n^{e_{n - 1}}, y_n)
$$
其中 $0 \ll e_{n -1} \ll e_{n - 2} \ll \ldots \ll e_1$。由引理
\ref{algebra-lemma-helper-polynomial}，把 $G(y_1, \ldots, y_n)$ 视为
$y_n$ 的多项式时，它形如
$$
a_K y_n^{k_n + \sum_{i = 1, \ldots, n - 1} e_i k_i} +
\text{lower order terms in }y_n
$$
由于 $a_K$ 是单位，可知 $M = R[x_1, \ldots, x_n]/J$ 是
$R[y_1, \ldots, y_{n - 1}]$ 上的有限模。因此
$U(R \to R[x_1, \ldots, x_n], M) = U(R \to R[y_1, \ldots, y_{n - 1}], M)$，
由对 $n$ 的归纳即得结论。

\medskip\noindent
情形 $m > 0$。选取多重指标 $K \in E$，使得 $a_K$ 不是单位。
同前，令
$U_1 = \Spec(R_{a_K}) = \Spec(R) \setminus V(a_K)$，
并令
$$
U_2 = \Spec(R) \setminus \overline{U_1}.
$$
显然 $U = U_1 \cup U_2$ 在 $\Spec(R)$ 中稠密。设 $f \in R$
满足 (a) $D(f) \subset U_1$ 或 (b) $D(f) \subset U_2$。
若对每个这样的 $f$，引理都对二元组
$(R_f \to R_f[x_1, \ldots, x_n], M_f)$ 成立，则由引理
\ref{algebra-lemma-generic-flatness-locus-reduce} 可知
$U(R \to S, M)$ 在 $\Spec(R)$ 中稠密。因此可以假设
(a) $a_KR = R$，或 (b) $V(a_K) = \Spec(R)$。
在情形 (a) 中，$a_K$ 已成为单位，所以数 $m$ 减小。
在情形 (b) 中，由于 $R$ 约化，可知 $a_K = 0$。
于是集合 $E$ 缩小，数 $m$ 也随之减小。在两种情形中都由对 $m$ 的归纳得到结论。

\medskip\noindent
至此已经证明了 $S = R[x_1, \ldots, x_n]$ 的情形。假设
$(R \to S, M)$ 是满足引理条件的任意二元组。选取满射
$R[x_1, \ldots, x_n] \to S$。采用 (\ref{algebra-equation-good-locus})
中引入的记号，有
$$
U(R \to S, M) =
U(R \to R[x_1, \ldots, x_n], S)
\cap
U(R \to R[x_1, \ldots, x_n], M)
$$
% P01-E0561: the source final sentence is grammatically malformed.
右侧两个开集由刚证明的结果都是稠密的，故左侧的开集也是稠密的。
\end{proof}

\section{Krull--Akizuki 定理及相关结果}
\label{algebra-section-krull-akizuki}

\noindent
Krull--Akizuki 定理的一个应用是说明离散赋值环非常丰富。
本节还将更一般地说明如何构造支配诺特局部环的离散赋值环。

\medskip\noindent
首先说明如何通过对极大理想作爆破，用一个 $1$ 维诺特局部整环
支配给定的诺特局部整环。

\begin{lemma}
\label{algebra-lemma-dominate-by-dimension-1}
设 $R$ 是分式域为 $K$ 的诺特局部整环。假设 $R$ 不是域。
则存在 $R \subset R' \subset K$，满足
\begin{enumerate}
\item $R'$ 是维数为 $1$ 的诺特局部环；
\item $R \to R'$ 是局部环映射，亦即 $R'$ 支配 $R$；并且
\item $R \to R'$ 本质有限型。
\end{enumerate}
\end{lemma}

\begin{proof}
选取支配 $R$ 的任意赋值环 $A \subset K$（其存在性由引理
\ref{algebra-lemma-dominate} 保证）。
% P01-E0562: the source designation “Denote v ...” omits the required preposition “by”.
以 $v$ 表示相应的赋值。设 $x_1, \ldots, x_r$ 是 $R$ 的极大理想
$\mathfrak m$ 的一组极小生成元。可以并且确实假设
$v(x_r) = \min\{v(x_1), \ldots, v(x_r)\}$。考虑环
$$
S = R[x_1/x_r, x_2/x_r, \ldots, x_{r - 1}/x_r] \subset K.
$$
注意 $\mathfrak mS = x_rS$ 是主理想。又有 $S \subset A$ 且
$v(x_r) > 0$，因此 $x_rS \not = S$。选取 $x_rS$ 上方的一个极小素理想
$\mathfrak q$。由引理 \ref{algebra-lemma-minimal-over-1}，
$\text{高度}(\mathfrak q) = 1$，并且 $\mathfrak q$ 位于
$\mathfrak m$ 上方。因此 $R' = S_{\mathfrak q}$ 即为所求。
\end{proof}

\begin{lemma}[Koll\'ar]
\label{algebra-lemma-hart-serre-loc-thm}
\begin{reference}
这取自 J\'anos Koll\'ar 即将发表的论文
“Variants of normality for Noetherian schemes”。
\end{reference}
设 $(R, \mathfrak m)$ 是诺特局部环。则下列情形恰有一个成立：
\begin{enumerate}
\item $(R, \mathfrak m)$ 是阿廷环；
\item $(R, \mathfrak m)$ 是维数为 $1$ 的正则环；
\item $\text{depth}(R) \geq 2$；或
\item 存在有限环映射 $R \to R'$，它不是同构，其核与余核都被
$\mathfrak m$ 的某个幂消去，并且 $\mathfrak m$ 不是 $R'$ 的相关素理想，
且 $R' \not = 0$。
\end{enumerate}
\end{lemma}

\begin{proof}
注意，$(R, \mathfrak m)$ 不是阿廷环，当且仅当
$V(\mathfrak m) \subset \Spec(R)$ 无处稠密。见命题
\ref{algebra-proposition-dimension-zero-ring}。以下始终作此假设。

\medskip\noindent
设 $J \subset R$ 是被 $\mathfrak m$ 的某个幂消去的最大理想。
若 $J \not = 0$，则 $R \to R/J$ 表明 $(R, \mathfrak m)$ 属于情形 (4)。

\medskip\noindent
否则 $J = 0$。特别地，$\mathfrak m$ 不是 $R$ 的相关素理想，
由引理 \ref{algebra-lemma-ideal-nonzerodivisor}，存在非零因子
$x \in \mathfrak m$。若 $\mathfrak m$ 不是 $R/xR$ 的相关素理想，
则由同一引理有 $\text{depth}(R) \geq 2$。所以只剩下如下情形：
存在 $y \in R$、$y \not \in xR$，使得 $y \mathfrak m \subset xR$。

\medskip\noindent
若 $y \mathfrak m \subset x \mathfrak m$，则可以考虑映射
$\varphi : \mathfrak m \to \mathfrak m$、$f \mapsto yf/x$
（由于 $x$ 是非零因子，该映射定义良好）。由引理
\ref{algebra-lemma-charpoly-module} 的行列式技巧，存在系数在 $R$ 中的首一多项式
$P$，使得 $P(\varphi) = 0$。于是 $P(y/x) = 0$ 在 $R_x$ 中成立。
设 $R' \subset R_x$ 是由 $R$ 和 $y/x$ 生成的环。则
$R \subset R'$，并且 $R'/R$ 是被 $\mathfrak m$ 的某个幂消去的有限
$R$-模。因此 $R$ 属于情形 (4)。

\medskip\noindent
否则，存在 $t \in \mathfrak m$，使得对 $R$ 的某个单位 $u$ 有
$y t = u x$。把 $t$ 替换为 $u^{-1}t$ 后，得到 $yt = x$。
特别地，$y$ 是非零因子。对任意 $t' \in \mathfrak m$，
都有 $y t' = x s$，其中 $s \in R$。所以
$y (t' - s t ) = x s - x s = 0$。由于 $y$ 不是零因子，
这蕴含 $t' = ts$，因而 $\mathfrak m = (t)$。所以
$(R, \mathfrak m)$ 是维数为 1 的正则环。

\medskip\noindent
到目前为止，论证表明每个 $R$ 都属于上述四种情形之一。
为完成证明，必须说明任意两个性质不能同时成立，也就是检验六种性质组合。
这里只说明 (2) 和 (3) 都排除 (4)；其余情形留给读者。

\medskip\noindent
假设 $R$ 是维数为 $1$ 的正则环，并且 $R \to R'$ 是有限环映射，
其核与余核被 $\mathfrak m$ 的某个幂消去，而 $\mathfrak m$ 不是
$R'$ 的相关素理想。于是 $R'$ 是深度至少为 $1$ 的有限 $R$-模，
因而由引理 \ref{algebra-lemma-regular-mcm-free} 可知它是自由的。
% P01-E0563: the source says “in the generic point” instead of “at the generic point”.
由于 $R \to R'$ 在泛点处是同构，$R'$ 作为 $R$-模必为秩 $1$ 的自由模。
于是 $R' \to \text{End}_R(R') \cong R$ 是映射 $R \to R'$ 的逆。
故 (4) 不可能成立。

\medskip\noindent
假设 $R$ 的深度 $\geq 2$，并且 $R \to R'$ 是有限环映射，
其核与余核被 $\mathfrak m$ 的某个幂消去，而 $\mathfrak m$ 不是
$R'$ 的相关素理想。则 $R \to R'$ 必为单射（否则其核的深度会同时为
$0$ 和 $\geq 1$），并且其余核的深度至少为 $1$（由引理
\ref{algebra-lemma-depth-in-ses} 以及假设所给的 $R'$ 深度 $\geq 1$）。
因此，除非该余核也是 $0$，否则它不可能支撑在
$\{\mathfrak m\}$ 上。故 (4) 不可能成立。
\end{proof}

\begin{lemma}
\label{algebra-lemma-nonregular-dimension-one}
设 $R$ 是极大理想为 $\mathfrak m$ 的局部环。假设 $R$ 是诺特环、
维数为 $1$，并且 $\dim(\mathfrak m/\mathfrak m^2) > 1$。
则存在环映射 $R \to R'$，使得
\begin{enumerate}
\item $R \to R'$ 是有限的；
\item $R \to R'$ 不是同构；
\item $R \to R'$ 的核与余核被 $\mathfrak m$ 的某个幂消去；并且
\item $\mathfrak m$ 不是 $R'$ 的相关素理想。
\end{enumerate}
\end{lemma}

\begin{proof}
这由引理 \ref{algebra-lemma-hart-serre-loc-thm} 以及下列事实得出：
$R$ 不是阿廷环，不正则，并且深度不 $\geq 2$（最后一点是因为，
由引理 \ref{algebra-lemma-bound-depth}，深度不超过维数）。
\end{proof}

\begin{example}
\label{algebra-example-nonreduced}
考虑诺特局部环
$$
R = k[[x, y]]/(y^2)
$$
它的维数为 1，并且是 Cohen--Macaulay 环。
引理 \ref{algebra-lemma-nonregular-dimension-one} 中扩张的一个例子是
$$
k[[x, y]]/(y^2) \subset k[[x, z]]/(z^2), \ \ y \mapsto xz
$$
换言之，这是通过向 $R$ 添加 $y/x$ 得到的扩张。把该构造重复
$n > 1$ 次，其效果是添加元素 $y/x^n$。
\end{example}

\begin{example}
\label{algebra-example-bad-dvr-char-p}
设 $k$ 是特征为 $p > 0$ 的域，并且 $k$ 在其 $p$ 次幂子域
$k^p$ 上的次数为无穷。例如，$k = \mathbf{F}_p(t_1, t_2, t_3, \ldots)$。
考虑环
$$
A =
\left\{
\sum a_i x^i \in k[[x]] \text{ 使得 }
[k^p(a_0, a_1, a_2, \ldots) : k^p] < \infty
\right\}
$$
则 $A$ 是离散赋值环，其完备化为 $A^\wedge = k[[x]]$。
注意，$A \subset k[[x]]$ 所诱导的分式域扩张是无限纯不可分的。
任取 $f \in k[[x]]$、$f \not \in A$。令 $R = A[f] \subset k[[x]]$。
则 $R$ 是维数为 $1$ 的诺特局部整环，而其完备化 $R^\wedge$
非约化（请思考！）。
\end{example}

\begin{remark}
\label{algebra-remark-resolution-dim-1}
假设 $R$ 是 $1$ 维半局部诺特整环。若存在极大理想
$\mathfrak m \subset R$，使得 $R_{\mathfrak m}$ 不正则，
则可以对 $(R, \mathfrak m)$ 应用引理
\ref{algebra-lemma-nonregular-dimension-one}，得到有限环扩张
$R \subset R_1$。（例如，可以这样进行构造，使得
$\Spec(R_1) \to \Spec(R)$ 是 $\Spec(R)$ 沿理想 $\mathfrak m$
的爆破。）当然，$R_1$ 是 $1$ 维半局部诺特整环，并且与 $R$
有相同的分式域。若 $R_1$ 不是正则半局部环，则可以重复该构造，
得到 $R_1 \subset R_2$。这样得到有限环扩张序列
$$
R \subset R_1 \subset R_2 \subset R_3 \subset \ldots
$$
若对某个 $n$，$R_n$ 正则，则该序列可以停止。奇点消解所断言的是，
$R_n$ 最终确实正则。事实上，情况并非如此。具体而言，存在特征为
$0$、维数为 $1$ 的诺特局部整环 $A$，其完备化非约化；见
\cite[Proposition 3.1]{Ferrand-Raynaud}，或参见我们的例子，第
\ref{examples-section-local-completion-nonreduced} 节。
特征 $p > 0$ 的例子见例 \ref{algebra-example-bad-dvr-char-p}。
由于爆破构造与完备化交换，很容易看出该序列永不稳定。
相关讨论见 \cite{Bennett}（主要针对正特征情形）。另一方面，
如果 $R$ 在其所有极大理想处的完备化都是约化的，则该过程会停止
% P01-E0564: the source contains the unresolved editorial placeholder “insert future reference here”.
（在此插入未来参考文献）。
\end{remark}

\begin{lemma}
\label{algebra-lemma-characterize-dvr}
设 $A$ 是环。下列条件等价。
\begin{enumerate}
\item 环 $A$ 是离散赋值环。
\item 环 $A$ 是诺特赋值环，但不是域。
\item 环 $A$ 是维数为 $1$ 的正则局部环。
\item 环 $A$ 是诺特局部整环，其极大理想
$\mathfrak m$ 由一个非零元生成。
\item 环 $A$ 是维数为 $1$ 的诺特正规局部整环。
\end{enumerate}
在这种情形中，若 $\pi$ 是 $A$ 的极大理想的生成元，则 $A$
的每个非零元都能唯一地写成 $u\pi^n$，其中 $u \in A$ 是单位。
\end{lemma}

\begin{proof}
(1) 与 (2) 的等价性就是引理
\ref{algebra-lemma-valuation-ring-Noetherian-discrete}。此外，在引理
\ref{algebra-lemma-valuation-ring-Noetherian-discrete} 的证明中已经看到：
若 $A$ 是离散赋值环，则 $A$ 是 PID，因而满足 (3)。
注意正则局部环是整环（见引理 \ref{algebra-lemma-regular-domain}）。
利用这一点，(3) 与 (4) 的等价性由维数理论得到，见第
\ref{algebra-section-dimension} 节。

\medskip\noindent
假设 (3) 成立，并设 $\pi$ 是极大理想 $\mathfrak m$ 的生成元。
对所有 $n \geq 0$，都有
$\dim_{A/\mathfrak m} \mathfrak m^n/\mathfrak m^{n + 1} = 1$，
因为该空间由 $\pi^n$ 生成（且不可能为零）。特别地，
$\mathfrak m^n = (\pi^n)$，并且分次环
$\bigoplus \mathfrak m^n/\mathfrak m^{n + 1}$ 同构于多项式环
% P01-E0565: the source omits parentheses around the residue-field quotient in A/m[T].
$A/\mathfrak m[T]$。对 $x \in A \setminus \{0\}$，定义
$v(x) = \max\{n \mid x \in \mathfrak m^n\}$。
换言之，$x = u \pi^{v(x)}$，其中 $u \in A^*$。
由上述说明，对所有 $x, y \in A \setminus \{0\}$，都有
$v(xy) = v(x) + v(y)$。通过规定
$v(a/b) = v(a) - v(b)$，把它扩张到 $A$ 的分式域 $K$
（由上面证明的乘法性可知其定义良好）。于是显然，$A$ 恰好由
$K$ 中赋值 $\geq 0$ 的元素组成。因此，由引理
\ref{algebra-lemma-valuation-valuation-ring} 可知 $A$ 是赋值环。

\medskip\noindent
由引理 \ref{algebra-lemma-valuation-ring-normal}，赋值环是正规整环。
所以等价条件 (1) -- (3) 蕴含 (5)。假设 (5) 成立。
为导出矛盾，假设 $\mathfrak m$ 不能由 $1$ 个元素生成。
由引理 \ref{algebra-lemma-nonregular-dimension-one}，存在有限环映射
$A \to A'$，它在把 $\mathfrak m$ 的任意非零元可逆化后成为同构，
但本身不是同构。特别地，可以把 $A'$ 视为 $A$ 的分式域的子集。
由于 $A \to A'$ 有限，它是整的（见引理
\ref{algebra-lemma-finite-is-integral}）。因为 $A$ 正规，得到
$A = A'$，矛盾。
\end{proof}

\begin{definition}
\label{algebra-definition-uniformizer}
设 $A$ 是离散赋值环。若元素 $\pi \in A$ 生成 $A$ 的极大理想，
则称它为{\it 一致化元}。
\end{definition}

\noindent
由引理 \ref{algebra-lemma-characterize-dvr}，离散赋值环的任意两个一致化元都相伴。

\begin{lemma}
\label{algebra-lemma-finite-length}
设 $R$ 是分式域为 $K$ 的整环。设 $M$ 是 $K^{\oplus r}$
的 $R$-子模。假设 $R$ 是维数为 $1$ 的诺特局部环。
对任意非零元 $x \in R$，有 $\text{length}_R(R/xR) < \infty$，并且
$$
\text{length}_R(M/xM) \leq r \cdot \text{length}_R(R/xR).
$$
\end{lemma}

\begin{proof}
若 $x$ 是单位，则结论成立。所以可以假设 $x \in \mathfrak m$，
其中后者是 $R$ 的极大理想。由于 $x$ 非零且 $R$ 是整环，
有 $\dim(R/xR) = 0$，因而 $R/xR$ 的长度有限。
考虑引理中的 $M \subset K^{\oplus r}$。必要时把
$K^{\oplus r}$ 替换为较小的子空间后，可以假设 $M$ 中的元素
生成 $K^{\oplus r}$ 这个 $K$-向量空间。

\medskip\noindent
先假设 $M$ 是有限 $R$-模。在此情形中，可以清除分母并假设
$M \subset R^{\oplus r}$。因为 $M$ 生成 $K^{\oplus r}$ 这个
% P01-E0566: the source has the number-agreement typo “a vectors space”.
向量空间，所以 $R^{\oplus r}/M$ 的长度有限。特别地，存在整数
$c \geq 0$，使得 $x^cR^{\oplus r} \subset M$。注意
$M \supset xM \supset x^2M \supset \ldots$ 是一个模序列，
其逐次商均同构于 $M/xM$。因此
$$
n \text{length}_R(M/xM) = \text{length}_R(M/x^nM).
$$
对 $M = R^{\oplus r}$ 作同样的论证，得到
$$
n \text{length}_R(R^{\oplus r}/xR^{\oplus r}) =
\text{length}_R(R^{\oplus r}/x^nR^{\oplus r}).
$$
由上面对 $c$ 的选取，$x^nM$ 夹在 $x^n R^{\oplus r}$ 与
$x^{n + c}R^{\oplus r}$ 之间。由此容易得到
$$
r(n + c) \text{length}_R(R/xR)
\geq
n \text{length}_R(M/xM)
\geq
r (n - c) \text{length}_R(R/xR)
$$
因此，在有限情形中，引理的结论事实上取等号。

\medskip\noindent
现在假设 $M$ 不是有限模。假设对某个自然数 $k$，
$M/xM$ 的长度 $\geq k$。则可以找到
% P01-E0567: the source display omits the final subset relation before N_k.
$$
0 \subset N_0 \subset N_1 \subset N_2 \subset \ldots N_k \subset M/xM
$$
使得对 $i = 0, \ldots k - 1$，有 $N_i \not = N_{i + 1}$。
对 $i = 1, \ldots, k$，选取元素 $m_i \in M$，使其模 $xM$
的同余类落在 $N_i$ 中但不落在 $N_{i - 1}$ 中。
考虑有限 $R$-模 $M' = Rm_1 + \ldots + Rm_k \subset M$。
设 $N'_i \subset M'/xM'$ 是 $N_i$ 的逆像。
由 $m_i$ 的选取，显然 $N'_i \not =N'_{i + 1}$。
所以 $\text{length}_R(M'/xM') \geq k$。由有限情形可知
$k \leq r\text{length}_R(R/xR)$，如欲得。
\end{proof}

\noindent
下面是第一个应用。

\begin{lemma}
\label{algebra-lemma-finite-extension-residue-fields-dimension-1}
设 $R \to S$ 是整环同态，并诱导分式域的嵌入
$K \subset L$。若 $R$ 是维数为 $1$ 的诺特局部环，并且
$[L : K] < \infty$，则
\begin{enumerate}
\item $S$ 中位于 $R$ 的极大理想 $\mathfrak m$ 上方的每个素理想
$\mathfrak n_i$ 都是极大理想；
\item 这样的素理想只有有限多个；并且
\item 对每个 $i$，都有
$[\kappa(\mathfrak n_i) : \kappa(\mathfrak m)] < \infty$。
\end{enumerate}
\end{lemma}

\begin{proof}
选取非零元 $x \in \mathfrak m$。对
$S \subset L \cong K^{\oplus n}$ 应用引理
\ref{algebra-lemma-finite-length}，其中 $n = [L : K]$。
于是环 $S/xS$ 作为 $R$-模的长度有限。因此
$S/\mathfrak m S$ 作为 $\kappa(\mathfrak m)$-向量空间的长度有限。
换言之，$\dim_{\kappa(\mathfrak m)} S/\mathfrak m S$ 有限
（引理 \ref{algebra-lemma-dimension-is-length}）。所以
$S/\mathfrak mS$ 是阿廷环（引理
\ref{algebra-lemma-finite-dimensional-algebra}）。
% P01-E0568: the source has subject–verb disagreement in “results ... implies”.
阿廷环的结构结果蕴含 (1) 与 (2)，例如可见引理
\ref{algebra-lemma-artinian-finite-length}。上面已经证明的有限性蕴含 (3)。
\end{proof}

\begin{lemma}
\label{algebra-lemma-finite-length-global}
设 $R$ 是分式域为 $K$ 的整环。设 $M$ 是 $K^{\oplus r}$
的 $R$-子模。假设 $R$ 是维数为 $1$ 的诺特环。
对任意非零元 $x \in R$，都有
$\text{length}_R(M/xM) < \infty$。
\end{lemma}

\begin{proof}
由于 $R$ 的维数为 $1$，$x$ 包含于有限多个素理想
$\mathfrak m_i$ 中，其中 $i = 1, \ldots, n$，并且它们都是极大理想。
由于 $R$ 是诺特环，$R/xR$ 是阿廷环，并且由命题
\ref{algebra-proposition-dimension-zero-ring} 和引理
\ref{algebra-lemma-artinian-finite-length}，有
$R/xR = \prod_{i = 1, \ldots, n} (R/xR)_{\mathfrak m_i}$。
因此 $M/xM$ 类似地分解为它在各 $\mathfrak m_i$ 处的局部化之积
$M/xM = \prod (M/xM)_{\mathfrak m_i}$。
对 $R_{\mathfrak m_i}$ 上的 $M_{\mathfrak m_i}$ 应用引理
\ref{algebra-lemma-finite-length}，可知每个
$M_{\mathfrak m_i}/xM_{\mathfrak m_i} = (M/xM)_{\mathfrak m_i}$
作为 $R_{\mathfrak m_i}$-模的长度有限。因此 $M/xM$ 作为 $R$-模
的长度有限，因为上述结论蕴含 $M/xM$ 有一个由 $R$-子模组成的有限滤过，
其逐次商同构于剩余域 $\kappa(\mathfrak m_i)$。
\end{proof}

\begin{lemma}[Krull--Akizuki]
\label{algebra-lemma-krull-akizuki}
设 $R$ 是分式域为 $K$ 的整环。设 $L/K$ 是有限域扩张。
假设 $R$ 是诺特环且 $\dim(R) = 1$。则任意满足
$R \subset A \subset L$ 的环 $A$ 都是诺特环。
\end{lemma}

\begin{proof}
设 $I \subset A$ 是非零理想。由引理
\ref{algebra-lemma-intersection-not-zero}，可以找到非零元
$x \in I \cap R$。于是有 $I/xA \subset A/xA$。
由引理 \ref{algebra-lemma-finite-length-global}，$R$-模 $A/xA$
作为 $R$-模的长度有限。因此 $I/xA$ 作为 $R$-模的长度有限。
所以 $I$ 作为 $A$ 的理想是有限生成的。
\end{proof}

\begin{lemma}
\label{algebra-lemma-exists-dvr}
设 $R$ 是分式域为 $K$ 的诺特局部整环。假设 $R$ 不是域。
设 $L/K$ 是有限生成域扩张。则存在分式域为 $L$ 的离散赋值环
% P01-E0569: the source statement omits the indefinite article before “discrete valuation ring”.
$A$，它支配 $R$。
\end{lemma}

\begin{proof}
若 $L$ 在 $K$ 上不是有限的，则选取 $L$ 在 $K$ 上的一组超越基
$x_1, \ldots, x_r$，并把 $R$ 替换为 $R[x_1, \ldots, x_r]$
在由 $\mathfrak m_R$ 和 $x_1, \ldots, x_r$ 生成的极大理想处的局部化。
于是可以假设 $K \subset L$ 有限。

\medskip\noindent
由引理 \ref{algebra-lemma-dominate-by-dimension-1}，可以假设 $\dim(R) = 1$。

\medskip\noindent
设 $A \subset L$ 是 $R$ 在 $L$ 中的整闭包。由引理
\ref{algebra-lemma-krull-akizuki}，它是诺特环。由引理
\ref{algebra-lemma-integral-overring-surjective}，存在素理想
$\mathfrak q \subset A$ 位于 $R$ 的极大理想上方。
由引理 \ref{algebra-lemma-characterize-dvr}，环 $A_{\mathfrak q}$
是支配 $R$ 的离散赋值环，如欲得。
\end{proof}

\section{因子分解}
\label{algebra-section-factoring}

\noindent
下面给出一些通常在本科一年级代数课程中讲授的概念及其相互关系。

\begin{definition}
\label{algebra-definition-irreducible-prime-element}
设 $R$ 是整环。
\begin{enumerate}
\item 若存在单位 $u \in R^*$，使得 $x = uy$，则称元素
$x, y \in R$ {\it 相伴}。
\item 若元素 $x \in R$ 非零、不是单位，并且每当
$x = yz$、$y, z \in R$ 时，$y$ 是单位或与 $x$ 相伴，
则称 $x$ 为{\it 不可约元}。
\item 若由元素 $x \in R$ 生成的理想是素理想，则称它为{\it 素元}。
\end{enumerate}
\end{definition}

\begin{lemma}
\label{algebra-lemma-easy-divisibility}
设 $R$ 是整环。设 $x, y \in R$。则 $x$、$y$ 相伴，当且仅当
$(x) = (y)$。
\end{lemma}

\begin{proof}
若对某个单位 $u \in R$ 有 $x = uy$，则 $(x) \subset (y)$；
又有 $y = u^{-1}x$，所以也有 $(y) \subset (x)$。
反过来，假设 $(x) = (y)$。则对某些
% P01-E0570: the frozen source uses the undefined ring A instead of R.
$f, g \in A$，有 $x = fy$ 且 $y = gx$。于是 $x = fg x$；
由于 $R$ 是整环，$fg = 1$。所以 $x$ 与 $y$ 相伴。
\end{proof}

\begin{lemma}
\label{algebra-lemma-factorization-exists}
设 $R$ 是整环。考虑下列条件：
\begin{enumerate}
\item 环 $R$ 满足主理想升链条件。
\item 每个非零非单位元 $a \in R$ 都有分解
$a = b_1 \ldots b_k$，其中每个 $b_i$ 都是 $R$ 的不可约元。
\end{enumerate}
则 (1) 蕴含 (2)。
\end{lemma}

\begin{proof}
设 $x$ 是一个不能分解为不可约元之积的非零非单位元。置
$x_1 = x$。可以写成 $x = yz$，其中
% P01-E0571: the source incorrectly requires both factors to be nonirreducible, rather than only nonunits.
$y$ 和 $z$ 都既不是不可约元，也不是单位。
于是，或者 $y$ 没有不可约分解，此时置 $x_2 = y$；
或者 $z$ 没有不可约分解，此时置 $x_2 = z$。
如此继续，得到 $R$ 中的元素序列
$$
\ldots | x_3 | x_2 | x_1
$$
其中 $x_n/x_{n + 1}$ 不是单位。这给出严格递增的主理想序列
$(x_1) \subset (x_2) \subset (x_3) \subset \ldots$，证明完成。
\end{proof}

\begin{definition}
\label{algebra-definition-UFD}
若整环 $R$ 满足如下性质，则称它为{\it 唯一分解整环}，简称 {\it UFD}：
若 $x \in R$ 是非零非单位元，则 $x$ 有不可约分解；并且，若
$$
x = a_1 \ldots a_m = b_1 \ldots b_n
$$
是两个不可约分解，则 $n = m$，并且存在置换
$\sigma : \{1, \ldots, n\} \to \{1, \ldots, n\}$，
使得 $a_i$ 与 $b_{\sigma(i)}$ 相伴。
\end{definition}

\begin{lemma}
\label{algebra-lemma-characterize-UFD}
设 $R$ 是整环。假设每个非零非单位元都能分解为不可约元之积。
则 $R$ 是 UFD，当且仅当每个不可约元都是素元。
\end{lemma}

\begin{proof}
假设 $R$ 是 UFD，并设 $x \in R$ 是不可约元。
设 $ab \in (x)$，亦即 $ab = cx$。选取分解
$a = a_1 \ldots a_n$、$b = b_1 \ldots b_m$ 和
$c = c_1 \ldots c_r$。由分解
$$
a_1 \ldots a_n b_1 \ldots b_m = c_1 \ldots c_r x
$$
的唯一性，可知 $x$ 与元素
% P01-E0572: the source range “a_1, ..., b_m” omits the intermediate endpoint and start.
$a_1, \ldots, b_m$ 中的一个相伴。换言之，
$a \in (x)$ 或 $b \in (x)$，所以 $x$ 是素元。

\medskip\noindent
假设每个不可约元都是素元。需要证明不可约分解在相差置换及取相伴元的意义下唯一。
设 $a_1 \ldots a_m = b_1 \ldots b_n$，其中 $a_i$ 和 $b_j$
均不可约。由于 $a_1$ 是素元，对某个 $j$ 有 $b_j \in (a_1)$。
重新编号后，可以假设 $b_1 \in (a_1)$。于是 $b_1 = a_1 u$；
由于 $b_1$ 不可约，$u$ 是单位。所以 $a_1$ 与 $b_1$ 相伴，并且
% P01-E0573: the frozen source swaps the a- and b-factor endpoints in the cancelled equality.
$a_2 \ldots a_n = ub_2\ldots b_m$。对 $n + m$ 作归纳，
可知 $n = m$，并且对 $i = 2, \ldots, n$，$a_i$ 与
$b_{\sigma(i)}$ 相伴，如欲得。
\end{proof}

\begin{lemma}
\label{algebra-lemma-characterize-UFD-height-1}
设 $R$ 是诺特整环。则 $R$ 是 UFD，当且仅当每个高度为 $1$
的素理想都是主理想。
\end{lemma}

\begin{proof}
假设 $R$ 是 UFD，并设 $\mathfrak p$ 是高度为 1 的素理想。
取非零元 $x \in \mathfrak p$，并设 $x = a_1 \ldots a_n$
是不可约分解。由于 $\mathfrak p$ 是素理想，对某个 $i$ 有
$a_i \in \mathfrak p$。由引理 \ref{algebra-lemma-characterize-UFD}，
理想 $(a_i)$ 是素理想。因为 $\mathfrak p$ 的高度为 $1$，
可知 $(a_i) = \mathfrak p$。

\medskip\noindent
假设每个高度为 $1$ 的素理想都是主理想。由于 $R$ 是诺特环，
每个非零非单位元 $x$ 都有不可约分解，见引理
\ref{algebra-lemma-factorization-exists}。由引理
\ref{algebra-lemma-characterize-UFD}，只须证明不可约元 $x$ 是素元。
设 $(x) \subset \mathfrak p$，其中后者是
$(x)$ 上方的极小素理想。由引理 \ref{algebra-lemma-minimal-over-1}，
$\mathfrak p$ 的高度为 $1$。根据假设，$\mathfrak p = (y)$。
于是 $x = yz$；由于 $x$ 不可约，$z$ 是单位。
所以 $(x) = (y)$，从而 $x$ 是素元。
\end{proof}

\begin{lemma}[Nagata 唯一分解性判据]
\label{algebra-lemma-invert-prime-elements}
\begin{reference}
\cite[Lemma 2]{Nagata-UFD}
\end{reference}
设 $A$ 是整环。设 $S \subset A$ 是由素元生成的乘法子集。
设 $x \in A$ 不可约。则
\begin{enumerate}
\item $x$ 在 $S^{-1}A$ 中的像是不可约元或单位；并且
\item $x$ 是素元，当且仅当 $x$ 在 $S^{-1}A$ 中的像是
$S^{-1}A$ 中的素元或单位。
\end{enumerate}
此外，$A$ 是 UFD，当且仅当 $A$ 的每个非零非单位元都有不可约分解，
并且 $S^{-1}A$ 是 UFD。
\end{lemma}

\begin{proof}
设对 $\alpha, \beta \in S^{-1}A$ 有 $x = \alpha \beta$。
则对某些 $a, b \in A$、$s, s' \in S$，有
$\alpha = a/s$ 且 $\beta = b/s'$。所以 $ss'x = ab$。
根据假设，可以对某些素元 $p_i$ 写成
$ss' = p_1 \ldots p_r$。对每个 $i$，元素 $p_i$ 整除
$a$ 或 $b$。作除法后，得到分解 $x = a' b'$，并且对某些
$s'', s''' \in S$，有 $a = s'' a'$、$b = s''' b'$。
由于 $x$ 不可约，$a'$ 或 $b'$ 是单位。回溯上述步骤可知
$\alpha$ 或 $\beta$ 是单位。这证明了 (1)。

\medskip\noindent
假设 $x$ 是素元。则 $A/(x)$ 是整环。因此
$S^{-1}A/xS^{-1}A = S^{-1}(A/(x))$ 是整环或零环。
所以 $x$ 映到素元或单位。

\medskip\noindent
假设 $x$ 在 $S^{-1}A$ 中的像是单位。则对某些
$s \in S$ 和 $y \in A$，有 $y x = s$。根据假设，
$s = p_1 \ldots p_r$，其中 $p_i$ 是素元。对每个 $i$，
$p_i$ 整除 $y$ 或 $p_i$ 整除 $x$。在后一种情形中，
$p_i$ 与 $x$ 相伴（因为 $x$ 不可约），结论得证。
但若对所有 $i = 1, \ldots, r$ 都出现前一种情形，则 $x$
是单位，矛盾。

\medskip\noindent
假设 $x$ 在 $S^{-1}A$ 中的像是素元。设 $a, b \in A$ 且
$ab \in (x)$。则对某些 $s \in S$ 和 $y \in A$，有
$sa = xy$ 或 $sb = xy$。设出现第一种情形。根据假设，
$s = p_1 \ldots p_r$，其中 $p_i$ 是素元。对每个 $i$，
$p_i$ 整除 $y$ 或 $p_i$ 整除 $x$。在后一种情形中，
$p_i$ 与 $x$ 相伴（因为 $x$ 不可约），结论得证。
若对所有 $i = 1, \ldots, r$ 都出现前一种情形，则
$a \in (x)$，如欲得。这完成了 (2) 的证明。

\medskip\noindent
引理的最后一个陈述由 (1)、(2) 以及引理
\ref{algebra-lemma-characterize-UFD} 得出。
\end{proof}

\begin{lemma}
\label{algebra-lemma-UFD-ascending-chain-condition-principal-ideals}
UFD 满足主理想升链条件。
\end{lemma}

\begin{proof}
考虑 $R$ 中主理想的升链
$(a_1) \subset (a_2) \subset (a_3) \subset \ldots$。
写成 $a_1 = p_1^{e_1} \ldots p_r^{e_r}$，其中 $p_i$ 是素元。
于是对某些 $0 \leq c_i \leq e_i$，$a_n$ 与
$p_1^{c_1} \ldots p_r^{c_r}$ 相伴。因为只有有限种可能，结论成立。
\end{proof}

\begin{lemma}
\label{algebra-lemma-factoring-in-polynomial}
设 $R$ 是整环。假设 $R$ 满足主理想升链条件。
则 $R$ 上的多项式环也满足这一性质。
\end{lemma}

\begin{proof}
考虑 $R[x]$ 中主理想的升链
$(f_1) \subset (f_2) \subset (f_3) \subset \ldots$。
由于 $f_{n + 1}$ 整除 $f_n$，该序列中的次数递减。
所以对所有 $n \gg 0$，$f_n$ 都有固定次数 $d \geq 0$。
设 $a_n$ 是 $f_n$ 的首项系数。条件 $f_n \in (f_{n + 1})$
蕴含对所有 $n$，$a_{n + 1}$ 整除 $a_n$。由对 $R$ 的假设，
对所有充分大的 $n$，$a_{n + 1}$ 与 $a_n$ 相伴（引理
\ref{algebra-lemma-easy-divisibility}）。所以对充分大的 $n$，
有 $f_n = u f_{n + 1}$，其中 $u \in R$（由次数原因）是单位
（因为 $a_n$ 与 $a_{n + 1}$ 相伴）。
\end{proof}

\begin{lemma}
\label{algebra-lemma-polynomial-ring-UFD}
UFD 上的多项式环是 UFD。特别地，若 $k$ 是域，则
$k[x_1, \ldots, x_n]$ 是 UFD。
\end{lemma}

\begin{proof}
设 $R$ 是 UFD。则 $R$ 满足主理想升链条件
（引理 \ref{algebra-lemma-UFD-ascending-chain-condition-principal-ideals}），
所以 $R[x]$ 满足主理想升链条件（引理
\ref{algebra-lemma-factoring-in-polynomial}），从而 $R[x]$ 的每个元素
都有不可约分解（引理 \ref{algebra-lemma-factorization-exists}）。
设 $S \subset R$ 是由素元生成的乘法子集。因为 $R$ 的每个非单位元
都是素元之积，可知 $K = S^{-1}R$ 是 $R$ 的分式域。
注意，$R$ 的每个素元都映到 $R[x]$ 的素元，并且
$S^{-1}(R[x]) = S^{-1}R[x] = K[x]$ 是 UFD（甚至是 PID）。
所以可以应用引理 \ref{algebra-lemma-invert-prime-elements} 得出结论。
\end{proof}

\begin{lemma}
\label{algebra-lemma-UFD-normal}
唯一分解整环是正规的。
\end{lemma}

\begin{proof}
设 $R$ 是 UFD。设 $x$ 是 $R$ 的分式域中整于 $R$ 的元素。
设 $x^d - a_1 x^{d - 1} - \ldots - a_d = 0$，其中
$a_i \in R$。可以写成
$x = u p_1^{e_1} \ldots p_r^{e_r}$，其中 $u$ 是单位，
$e_i \in \mathbf{Z}$，而 $p_1, \ldots, p_r$ 是两两不相伴的不可约元。
为证明引理，需要说明 $e_i \geq 0$。否则，不妨设 $e_1 < 0$；
则对 $N \gg 0$，有
$$
u^d p_2^{de_2 + N} \ldots p_r^{de_r + N} =
p_1^{-de_1}p_2^N \ldots p_r^N(
\sum\nolimits_{i = 1, \ldots, d} a_i x^{d - i} ) \in (p_1)
$$
这与 $R$ 中分解的唯一性矛盾。
\end{proof}

\begin{definition}
\label{algebra-definition-PID}
若整环 $R$ 的每个理想都是主理想，则称它为{\it 主理想整环}，
简称 {\it PID}。
\end{definition}

\begin{lemma}
\label{algebra-lemma-PID-UFD}
主理想整环是唯一分解整环。
\end{lemma}

\begin{proof}
因为 PID 是诺特环，这由引理
\ref{algebra-lemma-characterize-UFD-height-1} 得出。
\end{proof}

\begin{definition}
\label{algebra-definition-dedekind-domain}
若整环 $R$ 的每个非零理想 $I \subset R$ 都能写成非零素理想的乘积
$$
I = \mathfrak p_1 \ldots \mathfrak p_r
$$
并且在相差 $\mathfrak p_i$ 的置换的意义下唯一，
则称它为 {\it Dedekind 整环}。
\end{definition}

\begin{lemma}
\label{algebra-lemma-PID-dedekind}
PID 是 Dedekind 整环。
\end{lemma}

\begin{proof}
设 $R$ 是 PID。由于 $R$ 的每个非零理想都是主理想，并且 $R$
是 UFD（引理 \ref{algebra-lemma-PID-UFD}），结论由如下事实得出：
$R$ 中的每个不可约元都是素元（引理
\ref{algebra-lemma-characterize-UFD}），所以元素的分解转化为素理想分解。
\end{proof}

\begin{lemma}
\label{algebra-lemma-product-ideals-principal}
\begin{slogan}
理想的乘积作为模可逆，当且仅当两个因子都可逆。
\end{slogan}
设 $A$ 是环。设 $I$ 和 $J$ 是 $A$ 的非零理想，并且对某个
非零因子 $f \in A$ 有 $IJ = (f)$。则 $I$ 和 $J$ 都是有限生成理想，
并且作为 $A$-模都是秩为 $1$ 的有限局部自由模。
\end{lemma}

\begin{proof}
由引理 \ref{algebra-lemma-finite-projective}，只须证明 $I$ 和 $J$ 是秩为 $1$
的有限局部自由 $A$-模。为此，写成
$f = \sum_{i = 1, \ldots, n} x_i y_i$，其中 $x_i \in I$ 且
$y_i \in J$。还可以对某些 $a_i \in A$ 写成
$x_i y_i = a_i f$。因为 $f$ 是非零因子，可知 $\sum a_i = 1$。
所以只须说明每个 $I_{a_i}$ 和 $J_{a_i}$ 都是 $A_{a_i}$ 上秩为 $1$
的自由模。把 $A$ 替换为 $A_{a_i}$ 后，可知对某些
$x \in I$ 和 $y \in J$ 有 $f = xy$。注意 $x$ 和 $y$ 都是非零因子。
我们断言 $I = (x)$ 且 $J = (y)$，这就完成了证明。事实上，
若 $x' \in I$，则对某个 $a \in A$ 有 $x'y = af = axy$。
所以 $x' = ax$，结论得证。
\end{proof}

\begin{lemma}
\label{algebra-lemma-characterize-Dedekind}
设 $R$ 是环。下列条件等价：
% P01-E0574: the frozen source omits punctuation after “equivalent”.
\begin{enumerate}
\item $R$ 是 Dedekind 整环；
\item $R$ 是诺特整环，并且对每个非零极大理想 $\mathfrak m$，
局部环 $R_{\mathfrak m}$ 都是离散赋值环；并且
\item $R$ 是诺特正规整环，并且 $\dim(R) \leq 1$。
\end{enumerate}
\end{lemma}

\begin{proof}
假设 (1) 成立。由于在 Dedekind 整环的定义中并未假设 $R$ 是诺特环，
所以论证并非平凡。设 $\mathfrak p \subset R$ 是素理想。
由定义中分解的唯一性，注意到
$\mathfrak p \not = \mathfrak p^2$。选取
$x \in \mathfrak p$，使得 $x \not \in \mathfrak p^2$。
设 $y \in \mathfrak p$ 是第二个元素（例如 $y = 0$）。
写成 $(x, y) = \mathfrak p_1 \ldots \mathfrak p_r$。
由于 $(x, y) \subset \mathfrak p$，至少有一个素理想
$\mathfrak p_i$ 包含于 $\mathfrak p$。但因为
$x \not \in \mathfrak p^2$，至多只有一个。因此
$\mathfrak p_1, \ldots, \mathfrak p_r$ 中恰有一个包含于
$\mathfrak p$，不妨设 $\mathfrak p_1 \subset \mathfrak p$。
可知对 $y$ 的每种选取，
$(x, y)R_\mathfrak p = \mathfrak p_1R_\mathfrak p$ 都是素理想。
我们断言 $(x)R_\mathfrak p = \mathfrak pR_\mathfrak p$。
事实上，选取 $y \in \mathfrak p$。把上述论证应用于 $y^2$，
可知 $(x, y^2)R_\mathfrak p$ 是素理想。因此
$y \in (x, y^2)R_\mathfrak p$，亦即在 $R_\mathfrak p$ 中
$y = ax + by^2$。所以
$(1 - by)y = ax \in (x)R_\mathfrak p$，亦即
$y \in (x)R_\mathfrak p$，如欲得。

\medskip\noindent
重新写成 $(x) = \mathfrak p_1 \ldots \mathfrak p_r$，其中
$\mathfrak p_1 \subset \mathfrak p$。可知
$\mathfrak p_1 R_\mathfrak p = \mathfrak p R_\mathfrak p$，
亦即 $\mathfrak p_1 = \mathfrak p$。此外，由引理
\ref{algebra-lemma-product-ideals-principal}，
$\mathfrak p_1 = \mathfrak p$ 是 $R$ 的有限生成理想。
由引理 \ref{algebra-lemma-cohen}，可知 $R$ 是诺特环。
% P01-E0575: the source mismatches m/p and overstates “every prime ideal”, including the zero prime.
此外，对每个素理想 $\mathfrak p$，$R_\mathfrak m$ 都是离散赋值环；
见引理 \ref{algebra-lemma-characterize-dvr}。

\medskip\noindent
(2) 与 (3) 的等价性由引理
\ref{algebra-lemma-normality-is-local} 和
\ref{algebra-lemma-characterize-dvr} 得出。假设 (2) 和 (3) 成立。
设 $I \subset R$ 是理想。我们构造 $I$ 的分解。
若 $I$ 是素理想，则无须证明。否则，选取
$I \subset \mathfrak p$，其中 $\mathfrak p \subset R$ 是极大理想。
令 $J = \{x \in R \mid x \mathfrak p \subset I\}$。
我们断言 $J \mathfrak p = I$。只须在 $R$ 的各极大理想
$\mathfrak m$ 处局部化后检验这一点（$J$ 的形成与局部化交换，
并且使用引理 \ref{algebra-lemma-characterize-zero-local}）。
于是，或者 $\mathfrak p R_\mathfrak m = R_\mathfrak m$，
此时结论显然；或者
$\mathfrak p R_\mathfrak m = \mathfrak m R_\mathfrak m$。
在后一种情形中，$\mathfrak p R_\mathfrak m = (\pi)$，
而 $\mathfrak p$ 为主理想时结论立即成立。
由诺特归纳，理想 $J$ 有分解，从而得到所需的 $I$ 的分解。
我们略去分解唯一性的证明。
\end{proof}

\noindent
下面是 Krull--Akizuki 引理的一个变体。

\begin{lemma}
\label{algebra-lemma-integral-closure-Dedekind}
设 $A$ 是维数为 $1$、分式域为 $K$ 的诺特整环。
设 $L/K$ 是有限扩张。设 $B$ 是 $A$ 在 $L$ 中的整闭包。
则 $B$ 是 Dedekind 整环，并且
$\Spec(B) \to \Spec(A)$ 是满射、具有有限纤维，并诱导有限剩余域扩张。
\end{lemma}

\begin{proof}
由 Krull--Akizuki 定理（引理 \ref{algebra-lemma-krull-akizuki}），
环 $B$ 是诺特环。由引理
\ref{algebra-lemma-integral-sub-dim-equal}，$\dim(B) = 1$。
所以由引理 \ref{algebra-lemma-characterize-Dedekind}，$B$ 是 Dedekind 整环。
谱上映射的满性由引理
\ref{algebra-lemma-integral-overring-surjective} 得出。
最后两个陈述由引理
\ref{algebra-lemma-finite-extension-residue-fields-dimension-1} 得出。
\end{proof}

\section{消失阶}
\label{algebra-section-orders-of-vanishing}

\begin{lemma}
\label{algebra-lemma-ord-additive}
设 $R$ 是维数为 $1$ 的半局部诺特环。
若 $a, b \in R$ 是非零因子，则
$$
\text{length}_R(R/(ab)) =
\text{length}_R(R/(a)) +
\text{length}_R(R/(b))
$$
且这些长度都是有限的。
\end{lemma}

\begin{proof}
我们已在引理 \ref{algebra-lemma-finite-length-global} 中证明了有限性。
% P01-E0576: the frozen source says “Use length is additive” instead of “Use that length is additive”.
可加性来自短正合列
$0 \to R/(a) \to R/(ab) \to R/(b) \to 0$，其中第一个映射
由乘以 $b$ 给出。（使用长度的可加性；
见引理 \ref{algebra-lemma-length-additive}。）
\end{proof}

\begin{definition}
\label{algebra-definition-ord}
设 $K$ 是域，且 $R \subset K$ 是一个
局部的\footnote{当 $R$ 仅为半局部环时，我们也可以作此定义，
但这大概绝不会真正是你想要的！}
诺特子环，其维数为 $1$，分式域为 $K$。
在这种情形下，我们定义沿 $R$ 的\textit{消失阶}
$$
\text{ord}_R : K^* \longrightarrow \mathbf{Z}
$$
如下：若 $x \in R$，则
$$
\text{ord}_R(x) = \text{length}_R(R/(x))
$$
并且对两个非零元 $x, y \in R$，置
$\text{ord}_R(x/y) = \text{ord}_R(x) - \text{ord}_R(y)$。
\end{definition}

\noindent
我们可以用消失阶比较向量空间中的格。定义如下。

\begin{definition}
\label{algebra-definition-lattice}
设 $R$ 是维数为 $1$、分式域为 $K$ 的诺特局部整环。
设 $V$ 是有限维 $K$-向量空间。
$V$ 中的一个\textit{格}是有限 $R$-子模 $M \subset V$，
满足 $V = K \otimes_R M$。
\end{definition}

\noindent
条件 $V = K \otimes_R M$ 表示 $M$ 包含向量空间 $V$ 的一组基。
我们指出，文献中许多地方仅在环 $R$ 是离散赋值环时定义格。
若 $R$ 是离散赋值环，则任一格都是自由 $R$-模；一般情形下
则未必如此。

\begin{lemma}
\label{algebra-lemma-compare-lattices}
设 $R$ 是维数为 $1$、分式域为 $K$ 的诺特局部整环。
设 $V$ 是有限维 $K$-向量空间。
\begin{enumerate}
\item 若 $M$ 是 $V$ 中的格，且 $M \subset M' \subset V$ 是
包含 $M$ 的 $V$ 的一个 $R$-子模，则下列条件等价：
\begin{enumerate}
\item $M'$ 是格；
\item $\text{length}_R(M'/M)$ 有限；
\item $M'$ 是有限生成的。
\end{enumerate}
\item 若 $M$ 是 $V$ 中的格，而 $M' \subset M$ 是 $M$ 的一个
$R$-子模，则 $M'$ 是格当且仅当
$\text{length}_R(M/M')$ 有限。
\item 若 $M$、$M'$ 是 $V$ 中的格，则
$M \cap M'$ 和 $M + M'$ 也是格。
\item 若 $M \subset M' \subset M'' \subset V$ 是 $V$ 中的格，则
$$
\text{length}_R(M''/M) =
\text{length}_R(M'/M) +
\text{length}_R(M''/M').
$$
\item 若 $M$、$M'$、$N$、$N'$ 是 $V$ 中的格，并且
$N \subset M \cap M'$、$M + M' \subset N'$，则
\begin{eqnarray*}
& & \text{length}_R(M/M \cap M') - \text{length}_R(M'/M \cap M')\\
& = &
\text{length}_R(M/N) - \text{length}_R(M'/N) \\
& = &
\text{length}_R(M + M' / M') - \text{length}_R(M + M'/M) \\
& = &
\text{length}_R(N' / M') - \text{length}_R(N'/M)
\end{eqnarray*}
\end{enumerate}
\end{lemma}

\begin{proof}
证明 (1)。假设 (1)(a) 成立。设 $y_1, \ldots, y_m$ 生成 $M'$。
则每个 $y_i = x_i/f_i$，其中 $x_i \in M$，
而 $f_i \in R$ 非零。
因此 $f_1 \ldots f_m M' \subset M$。
由于 $R$ 是维数为 $1$ 的诺特局部环，
对某个 $n$ 有 $\mathfrak m^n \subset (f_1 \ldots f_m)$
（例如可合用引理 \ref{algebra-lemma-one-equation} 与命题
\ref{algebra-proposition-dimension-zero-ring}，或合用引理
\ref{algebra-lemma-finite-length} 与 \ref{algebra-lemma-length-infinite}）。
% P01-E0577: the frozen source omits sentence punctuation after “for some n”.
换言之，对某个 $n$ 有 $\mathfrak m^nM' \subset M$。
所以由引理 \ref{algebra-lemma-length-finite}，
$\text{length}(M'/M) < \infty$；也就是说 (1)(b) 成立。
假设 (1)(b) 成立。则 $M'/M$ 是有限 $R$-模
（见引理 \ref{algebra-lemma-finite-length-finite}）。
所以 $M'$ 作为两个有限 $R$-模的扩张也是有限 $R$-模，
从而 (1)(c) 成立。蕴含
(1)(c) $\Rightarrow$ (1)(a) 来自定义 \ref{algebra-definition-lattice}
之后的说明。

\medskip\noindent
证明 (2)。设
$M$ 是 $V$ 中的格，而 $M' \subset M$ 是一个 $R$-子模。
我们已在 (1) 中看到，若 $M'$ 是格，则
$\text{length}_R(M/M') < \infty$。反之，假设
$\text{length}_R(M/M') < \infty$。由于 $R$ 是诺特环，
$M'$ 是有限生成的；并且由引理 \ref{algebra-lemma-length-infinite}，
对某个 $n$ 有 $\mathfrak m^n M \subset M'$。
因此 $M'$ 包含 $V$ 的一组基，故 $M'$ 是格。

\medskip\noindent
证明 (3)。设 $M$、$M'$ 是 $V$ 中的格。
由于 $R$ 是诺特环，$M$ 的子模 $M \cap M'$ 是有限的。
因为 $M$ 是格，我们可以取
$x_1, \ldots, x_n \in M$，使其构成 $V$ 的一组 $K$-基。
因为 $M'$ 是格，可以写成 $x_i = y_i/f_i$，其中
$y_i \in M'$ 且 $f_i \in R$。于是 $f_ix_i \in M \cap M'$。
所以 $M \cap M'$ 也是格。
$M + M'$ 是格这一事实由 (1) 得出。

\medskip\noindent
(4) 来自长度的可加性（引理 \ref{algebra-lemma-length-additive}）以及如下正合列。
% P01-E0578: the frozen source omits punctuation after this displayed exact sequence.
$$
0 \to M'/M \to M''/M \to M''/M' \to 0
$$
(5) 由反复应用 (4) 得出。
\end{proof}

\begin{definition}
\label{algebra-definition-distance}
设 $R$ 是维数为 $1$、分式域为 $K$ 的诺特局部整环。
设 $V$ 是有限维 $K$-向量空间。
设 $M$、$M'$ 是 $V$ 中的两个格。$M$ 与 $M'$ 之间的\textit{距离}
是引理 \ref{algebra-lemma-compare-lattices} 第 (5) 部分中的整数
$$
d(M, M') = \text{length}_R(M/M \cap M') - \text{length}_R(M'/M \cap M')
$$
\end{definition}

\noindent
特别地，若 $M' \subset M$，则
$d(M, M') = \text{length}_R(M/M')$。

\begin{lemma}
\label{algebra-lemma-properties-distance-function}
设 $R$ 是维数为 $1$、分式域为 $K$ 的诺特局部整环。
设 $V$ 是有限维 $K$-向量空间。
这个距离函数具有如下性质：给定 $V$ 中任意三个格
$M$、$M'$、$M''$，都有
$$
d(M, M'') = d(M, M') + d(M', M'')
$$
特别地，$d(M, M') = - d(M', M)$。
\end{lemma}

\begin{proof}
略。
\end{proof}

\begin{lemma}
\label{algebra-lemma-order-vanishing-determinant}
设 $R$ 是维数为 $1$、分式域为 $K$ 的诺特局部整环。
设 $V$ 是有限维 $K$-向量空间。
设 $\varphi : V \to V$ 是 $K$-线性同构。
对任意格 $M \subset V$，有
$$
d(M, \varphi(M)) = \text{ord}_R(\det(\varphi))
$$
\end{lemma}

\begin{proof}
如下可见，整数 $d(M, \varphi(M))$ 与格 $M$ 的选择无关。
设 $M'$ 是另一个这样的格。则
\begin{eqnarray*}
d(M, \varphi(M)) & = & d(M, M') + d(M', \varphi(M)) \\
& = & d(M, M') + d(\varphi(M'), \varphi(M)) + d(M', \varphi(M'))
\end{eqnarray*}
由于 $\varphi$ 是同构，有
$d(\varphi(M'), \varphi(M)) = d(M', M) = -d(M, M')$，因而
$d(M, \varphi(M)) = d(M', \varphi(M'))$。此外，等式（即本引理的等式）
两边关于 $\varphi$ 都是可加的，也就是说
$$
\text{ord}_R(\det(\varphi \circ \psi))
=
\text{ord}_R(\det(\varphi))
+
\text{ord}_R(\det(\psi))
$$
以及
% P01-E0579: the frozen display is missing the final parenthesis closing d(...).
\begin{eqnarray*}
d(M, \varphi(\psi((M))) & = &
d(M, \psi(M)) + d(\psi(M), \varphi(\psi(M))) \\
& = & d(M, \psi(M)) + d(M, \varphi(M))
\end{eqnarray*}
其中第二个等式使用了上面证明的独立性。因此，只须对
$\text{GL}(V)$ 的生成元证明本引理。选取同构
$K^{\oplus n} \cong V$。则 $\text{GL}(V) = \text{GL}_n(K)$
由初等矩阵 $E$ 生成。
当 $E$ 为单位矩阵时，结论显然。
若 $E = E_{ij}(\lambda)$，其中 $i \not = j$、$\lambda \in K$、
$\lambda \not = 0$，例如
$$
E_{12}(\lambda) =
\left(
\begin{matrix}
1 & \lambda & \ldots \\
0 & 1 & \ldots \\
\ldots & \ldots & \ldots
\end{matrix}
\right)
$$
则相对于另一组基，它变为 $E_{12}(1)$。
取格 $R^{\oplus n} \subset K^{\oplus n}$，结论对
$E = E_{12}(1)$ 是显然的。最后，若 $E = E_i(a)$，
其中 $a \in K^*$，例如
$$
E_1(a) =
\left(
\begin{matrix}
a & 0 & \ldots \\
0 & 1 & \ldots \\
\ldots & \ldots & \ldots
\end{matrix}
\right)
$$
% P01-E0580: the frozen source uses undefined b in R^{\oplus b}; the ambient rank is n.
则 $E_1(a)(R^{\oplus b}) = aR \oplus R^{\oplus n - 1}$，并且显然
$d(R^{\oplus n}, aR \oplus R^{\oplus n - 1})
= \text{ord}_R(a)$，如所需。
\end{proof}

\begin{lemma}
\label{algebra-lemma-finite-extension-dim-1}
设 $A \to B$ 是环映射。假设
\begin{enumerate}
\item $A$ 是维数为 $1$ 的诺特局部整环；
\item $A \subset B$ 是整环的有限扩张。
\end{enumerate}
设 $L/K$ 是相应的分式域有限扩张。
设 $y \in L^*$，并令 $x = \text{Nm}_{L/K}(y)$。
在这种情形下，$B$ 是半局部环。
设 $\mathfrak m_i$（$i = 1, \ldots, n$）是 $B$ 的极大理想。
则
$$
\text{ord}_A(x) =
\sum\nolimits_i
[\kappa(\mathfrak m_i) : \kappa(\mathfrak m_A)]
\text{ord}_{B_{\mathfrak m_i}}(y)
$$
其中 $\text{ord}$ 如定义 \ref{algebra-definition-ord} 所定义。
\end{lemma}

\begin{proof}
由引理 \ref{algebra-lemma-finite-in-codim-1}，环 $B$ 是半局部的。
将其写成 $y = b/b'$，其中 $b, b' \in B$。
由 $\text{ord}$ 的可加性和 $\text{Nm}$ 的乘法性，只须对
$y = b$ 或 $y = b'$ 证明本引理。换言之，可以假设 $y \in B$。
此时公式右边为
$$
\sum [\kappa(\mathfrak m_i) : \kappa(\mathfrak m_A)]
\text{length}_{B_{\mathfrak m_i}}((B/yB)_{\mathfrak m_i})
$$
由引理 \ref{algebra-lemma-pushdown-module}，它等于
$\text{length}_A(B/yB)$。由引理 \ref{algebra-lemma-order-vanishing-determinant}，
有
$$
\text{length}_A(B/yB) = d(B, yB) =
\text{ord}_A(\det\nolimits_K(L \xrightarrow{y} L)).
$$
根据定义，$x = \text{Nm}_{L/K}(y) = \det\nolimits_K(L \xrightarrow{y} L)$，
故引理得证。
\end{proof}

\section{拟有限映射}
\label{algebra-section-quasi-finite}

\noindent
考虑一个有限型环映射 $R \to S$。
若某点在其纤维中是孤立点，就称映射
$\Spec(S) \to \Spec(R)$ 在该点拟有限。
这表示该纤维在该点处为零维。
本节研究这一重要但技术性的概念的基本性质。
更深入的内容可见下一节。

\begin{lemma}
\label{algebra-lemma-isolated-point}
设 $k$ 是域。
设 $S$ 是有限型 $k$-代数。
设 $\mathfrak q$ 是 $S$ 的素理想。
下列条件等价：
\begin{enumerate}
\item $\mathfrak q$ 是 $\Spec(S)$ 的孤立点；
\item $S_{\mathfrak q}$ 在 $k$ 上有限；
\item 存在 $g \in S$、$g \not\in \mathfrak q$，使得
$D(g) = \{ \mathfrak q \}$；
\item $\dim_{\mathfrak q} \Spec(S) = 0$；
\item $\mathfrak q$ 是 $\Spec(S)$ 的闭点，并且
$\dim(S_{\mathfrak q}) = 0$；
\item 域扩张 $\kappa(\mathfrak q)/k$ 是有限的，并且
$\dim(S_{\mathfrak q}) = 0$。
\end{enumerate}
在这种情形下，对某个有限型 $k$-代数 $S'$，有
$S = S_{\mathfrak q} \times S'$。此外，(3) 中的元素 $g$
还满足 $S_{\mathfrak q} = S_g$。
\end{lemma}

\begin{proof}
假设 $\mathfrak q$ 是 $\Spec(S)$ 的孤立点，即
$\{\mathfrak q\}$ 在 $\Spec(S)$ 中是开集。
因为 $\Spec(S)$ 是 Jacobson 空间（见引理
\ref{algebra-lemma-finite-type-field-Jacobson} 和
\ref{algebra-lemma-jacobson}），所以 $\mathfrak q$ 是闭点。
因此 $\{\mathfrak q\}$ 在 $\Spec(S)$ 中既开又闭。
由引理 \ref{algebra-lemma-disjoint-decomposition} 和
\ref{algebra-lemma-disjoint-implies-product}，可以写成
$S = S_1 \times S_2$，其中 $\mathfrak q$ 对应于
$\Spec(S_1)$ 的唯一一点。
所以 $S_1 = S_{\mathfrak q}$ 是有限型 $k$-代数，且维数为零。
因此它在 $k$ 上有限，例如可由引理
\ref{algebra-lemma-Noether-normalization} 得出。
我们已经证明 (1) 蕴含 (2)。

\medskip\noindent
假设 $S_{\mathfrak q}$ 在 $k$ 上有限。
则由引理 \ref{algebra-lemma-finite-dimensional-algebra}，
$S_{\mathfrak q}$ 是阿廷局部环。因此由引理
\ref{algebra-lemma-artinian-finite-length}，
$\Spec(S_{\mathfrak q}) = \{\mathfrak qS_{\mathfrak q}\}$。
考虑正合列 $0 \to K \to S \to S_{\mathfrak q}
\to Q \to 0$。显然 $K_{\mathfrak q} = Q_{\mathfrak q} = 0$。
又因为 $S$ 是诺特环，$K$ 是有限 $S$-模；而由于
$S_{\mathfrak q}$ 在 $k$ 上有限，$Q$ 是有限 $S$-模。
所以存在 $g \in S$、$g \not \in \mathfrak q$，使得
$K_g = Q_g = 0$。于是 $S_{\mathfrak q} = S_g$，且
$D(g) = \{ \mathfrak q \}$。我们已经证明 (2) 蕴含 (3)。

\medskip\noindent
假设 $D(g) =  \{ \mathfrak q \}$。按照 $\Spec(S)$ 上拓扑的构造，
$D(g)$ 是开集，因而 $\mathfrak q$ 是 $\Spec(S)$ 的孤立点。
我们已经证明 (3) 蕴含 (1)。
换言之，(1)、(2)、(3) 等价。

\medskip\noindent
假设 $\dim_{\mathfrak q} \Spec(S) = 0$。这表示
$\mathfrak q$ 在 $\Spec(S)$ 中有某个零维开邻域。
于是它有一个形如 $D(g)$ 的零维开邻域。因为 $S_g$ 是诺特环，
可知 $S_g$ 是阿廷环，且
$D(g) = \Spec(S_g)$ 是有限离散集；见命题
\ref{algebra-proposition-dimension-zero-ring}。
所以 $\mathfrak q$ 是 $D(g)$ 的孤立点；把上面 (1) 与 (2) 的等价性
用于 $\mathfrak qS_g \subset S_g$，便得到
$S_{\mathfrak q} = (S_g)_{\mathfrak qS_g}$ 在 $k$ 上有限。
所以 (4) 蕴含 (2)。显然 (1) 蕴含 (4)。
故 (1)--(4) 全部等价。

\medskip\noindent
引理 \ref{algebra-lemma-dimension-closed-point-finite-type-field}
给出蕴含 (5) $\Rightarrow$ (4)。
蕴含 (4) $\Rightarrow$ (6) 来自引理
\ref{algebra-lemma-dimension-at-a-point-finite-type-field}。
蕴含 (6) $\Rightarrow$ (5) 来自引理
\ref{algebra-lemma-finite-residue-extension-closed}。
至此，我们知道 (1)--(6) 等价。

\medskip\noindent
% P01-E0581: the frozen source has the object before the finite verb (“The two statements ... we saw”).
在上面证明 (1)、(2)、(3) 等价的过程中，
我们已经看到引理末尾的两个陈述。
\end{proof}

\begin{lemma}
\label{algebra-lemma-isolated-point-fibre}
\begin{slogan}
纤维中孤立点的等价条件
\end{slogan}
设 $R \to S$ 是有限型环映射。
设 $\mathfrak q \subset S$ 是位于
$\mathfrak p \subset R$ 之上的素理想。令
$F = \Spec(S \otimes_R \kappa(\mathfrak p))$
为 $\Spec(S) \to \Spec(R)$ 的纤维；见注
\ref{algebra-remark-fundamental-diagram}。
% P01-E0582: the frozen source says “Denote qbar the point” rather than “Denote by qbar ...”.
以 $\overline{\mathfrak q} \in F$ 表示对应于
$\mathfrak q$ 的点。下列条件等价：
% P01-E0583: the frozen source omits punctuation after “equivalent”.
\begin{enumerate}
\item $\overline{\mathfrak q}$ 是 $F$ 的孤立点；
\item $S_{\mathfrak q}/\mathfrak pS_{\mathfrak q}$ 在
$\kappa(\mathfrak p)$ 上有限；
\item 存在 $g \in S$、$g \not \in \mathfrak q$，使得
$D(g)$ 中映到 $\mathfrak p$ 的素理想只有 $\mathfrak q$；
\item $\dim_{\overline{\mathfrak q}}(F) = 0$；
\item $\overline{\mathfrak q}$ 是 $F$ 的闭点，并且
$\dim(S_{\mathfrak q}/\mathfrak pS_{\mathfrak q}) = 0$；
\item 域扩张 $\kappa(\mathfrak q)/\kappa(\mathfrak p)$
是有限的，并且 $\dim(S_{\mathfrak q}/\mathfrak pS_{\mathfrak q}) = 0$。
\end{enumerate}
\end{lemma}

\begin{proof}
注意，$S_{\mathfrak q}/\mathfrak pS_{\mathfrak q} =
(S \otimes_R \kappa(\mathfrak p))_{\overline{\mathfrak q}}$。
此外，$S \otimes_R \kappa(\mathfrak p)$ 是
$\kappa(\mathfrak p)$ 上的有限型代数。
这些条件恰好对应于引理 \ref{algebra-lemma-isolated-point} 中
$\kappa(\mathfrak p)$-代数 $S \otimes_R \kappa(\mathfrak p)$
及其素理想 $\overline{\mathfrak q}$ 的各个条件，所以它们等价。
\end{proof}

\begin{definition}
\label{algebra-definition-quasi-finite}
设 $R \to S$ 是有限型环映射。
设 $\mathfrak q \subset S$ 是素理想。
\begin{enumerate}
\item 若引理 \ref{algebra-lemma-isolated-point-fibre} 中的等价条件成立，
就称 $R \to S$ \textit{在 $\mathfrak q$ 处拟有限}。
\item 若环映射 $A \to B$ 是有限型的，并且在 $B$ 的每个素理想处
都拟有限，就称它是\textit{拟有限}的。
\end{enumerate}
\end{definition}

\begin{lemma}
\label{algebra-lemma-quasi-finite-above-p}
设 $R \to S$ 是有限型环映射，
$\mathfrak p$ 是 $R$ 的素理想。
则下列条件等价：
\begin{enumerate}
\item $R \to S$ 在 $S$ 中位于 $\mathfrak p$ 之上的所有素理想处拟有限；
\item $S \otimes_R \kappa(\mathfrak p)$ 是有限
$\kappa(\mathfrak p)$-代数；
\item $\Spec(S \otimes_R \kappa(\mathfrak p))$ 是有限集。
\end{enumerate}
\end{lemma}

\begin{proof}
条件 (1) 表示 $\Spec(S \otimes_R \kappa(\mathfrak p))$
上的拓扑是离散的（因为每一点都是开点）。
所以 (1)、(2)、(3) 的等价性来自引理
\ref{algebra-lemma-finite-type-algebra-finite-nr-primes}。
\end{proof}

\begin{lemma}
\label{algebra-lemma-quasi-finite}
设 $R \to S$ 是有限型环映射。则 $R \to S$ 拟有限，当且仅当
对所有素理想 $\mathfrak p \subset R$，环
$S \otimes_R \kappa(\mathfrak p)$ 在 $\kappa(\mathfrak p)$ 上有限。
\end{lemma}

\begin{proof}
这立即由更一般的引理 \ref{algebra-lemma-quasi-finite-above-p} 得出。
\end{proof}

\begin{lemma}
\label{algebra-lemma-quasi-finite-local}
设 $R \to S$ 是有限型环映射。设 $\mathfrak q \subset S$
是位于 $\mathfrak p \subset R$ 之上的素理想。设
$f \in R$、$f \not \in \mathfrak p$，且 $g \in S$、$g \not \in \mathfrak q$。
则 $R \to S$ 在 $\mathfrak q$ 处拟有限，当且仅当
$R_f \to S_{fg}$ 在 $\mathfrak qS_{fg}$ 处拟有限。
\end{lemma}

\begin{proof}
$\Spec(S_{fg}) \to \Spec(R_f)$ 的纤维同胚于
$\Spec(S) \to \Spec(R)$ 的纤维的一个开子集。
因此，本引理由引理 \ref{algebra-lemma-isolated-point-fibre}
等价条件中的 (1) 得出。
\end{proof}

\begin{lemma}
\label{algebra-lemma-four-rings}
设
$$
\xymatrix{
S \ar[r] & S' & &
\mathfrak q \ar@{-}[r] & \mathfrak q' \\
R \ar[u] \ar[r] &  R' \ar[u] & &
\mathfrak p \ar@{-}[r] \ar@{-}[u] & \mathfrak p' \ar@{-}[u]
}
$$
是带有所示素理想的环的交换图。
% P01-E0584: the frozen source omits “is” in “Assume R to S of finite type”.
假设 $R \to S$ 是有限型的，并且
$S \otimes_R R' \to S'$ 是满射。
若 $R \to S$ 在 $\mathfrak q$ 处拟有限，则
$R' \to S'$ 在 $\mathfrak q'$ 处拟有限。
\end{lemma}

\begin{proof}
写成 $S \otimes_R \kappa(\mathfrak p) = S_1 \times S_2$，
其中 $S_1$ 在 $\kappa(\mathfrak p)$ 上有限，且按照引理
\ref{algebra-lemma-isolated-point}，$\mathfrak q$ 对应于 $S_1$ 的一点。
这个乘积分解在任意 $S \otimes_R \kappa(\mathfrak p)$-代数上
诱导相应的乘积分解。特别地，我们得到
$S' \otimes_{R'} \kappa(\mathfrak p') = S'_1 \times S'_2$。
因为 $S \otimes_R R' \to S'$ 是满射，典范映射
$(S \otimes_R \kappa(\mathfrak p)) \otimes_{\kappa(\mathfrak p)}
\kappa(\mathfrak p') \to S' \otimes_{R'} \kappa(\mathfrak p')$
是满射，因而
$S_i \otimes_{\kappa(\mathfrak p)} \kappa(\mathfrak p') \to S'_i$
是满射。所以 $S'_1$ 在 $\kappa(\mathfrak p')$ 上有限。
映射 $S' \otimes_{R'} \kappa(\mathfrak p') \to
\kappa(\mathfrak q')$ 通过 $S_1'$ 分解
（即它在因子 $S_2'$ 上为零），
因为映射 $S \otimes_R \kappa(\mathfrak p) \to
\kappa(\mathfrak q)$ 通过 $S_1$ 分解
（即它在因子 $S_2$ 上为零）。所以在纤维的无交并分解
$\Spec(S' \otimes_{R'} \kappa(\mathfrak p'))
= \Spec(S_1') \amalg \Spec(S_2')$ 中，$\mathfrak q'$ 对应于
$\Spec(S_1')$ 的一点；见引理 \ref{algebra-lemma-spec-product}。
% P01-E0585: the frozen source omits the article in “it is Artinian ring”.
由于 $S_1'$ 在域上有限，它是阿廷环，
从而 $\Spec(S_1')$ 是有限离散集。
（见命题 \ref{algebra-proposition-dimension-zero-ring}。）
因此 $\mathfrak q'$ 在其纤维中是孤立的，如所需。
\end{proof}

\begin{lemma}
\label{algebra-lemma-quasi-finite-composition}
拟有限环映射的复合是拟有限的。
\end{lemma}

\begin{proof}
假设 $A \to B$ 与 $B \to C$ 是拟有限环映射。由引理
\ref{algebra-lemma-compose-finite-type}，$A \to C$ 是有限型的。
设 $\mathfrak r \subset C$ 是 $C$ 的素理想，位于
$\mathfrak q \subset B$ 和 $\mathfrak p \subset A$ 之上。
% P01-E0586: the frozen “Since ..., then ...” construction has a redundant then.
由于 $A \to B$ 与 $B \to C$ 分别在 $\mathfrak q$ 和
$\mathfrak r$ 处拟有限，存在 $b \in B$ 与 $c \in C$，使得
$\mathfrak q$ 是 $D(b)$ 中映到 $\mathfrak p$ 的唯一素理想，且类似地，
$\mathfrak r$ 是 $D(c)$ 中映到 $\mathfrak q$ 的唯一素理想。
若 $c' \in C$ 是 $b \in B$ 的像，则 $\mathfrak r$ 是
$D(cc')$ 中映到 $\mathfrak p$ 的唯一素理想。
所以 $A \to C$ 在 $\mathfrak r$ 处拟有限。
\end{proof}

\begin{lemma}
\label{algebra-lemma-quasi-finite-base-change}
设 $R \to S$ 是有限型环映射。
设 $R \to R'$ 是任意环映射。置 $S' = R' \otimes_R S$。
\begin{enumerate}
\item 集合
$\{\mathfrak q' \mid R' \to S' \text{ quasi-finite at }\mathfrak q'\}$
是 $\Spec(S)$ 中相应集合在典范映射
$\Spec(S') \to \Spec(S)$ 下的逆像。
\item 若 $\Spec(R') \to \Spec(R)$ 是满射，则
$R \to S$ 拟有限当且仅当 $R' \to S'$ 拟有限。
\item 拟有限环映射的任意基变换都是拟有限的。
\end{enumerate}
\end{lemma}

\begin{proof}
设 $\mathfrak p' \subset R'$ 是位于 $\mathfrak p \subset R$
之上的素理想。则纤维环 $S' \otimes_{R'} \kappa(\mathfrak p')$
是纤维环 $S \otimes_R \kappa(\mathfrak p)$
沿域扩张 $\kappa(\mathfrak p) \to \kappa(\mathfrak p')$ 的基变换。
因此，第一个断言来自维数在域扩张下的不变性
（引理 \ref{algebra-lemma-dimension-at-a-point-preserved-field-extension}）
以及引理 \ref{algebra-lemma-isolated-point}。
拟有限映射在基变换下的稳定性，由此和有限型性质在基变换下的稳定性得出。
第二个断言成立，是因为该假设蕴含：给定素理想
$\mathfrak q \subset S$，总能找到位于它之上的素理想
$\mathfrak q' \subset S'$。
\end{proof}

\begin{lemma}
\label{algebra-lemma-quasi-finite-permanence}
设 $A \to B$ 与 $B \to C$ 是环同态，且 $A \to C$
是有限型的。设 $\mathfrak r$ 是 $C$ 的素理想，位于
$\mathfrak q \subset B$ 和 $\mathfrak p \subset A$ 之上。
若 $A \to C$ 在 $\mathfrak r$ 处拟有限，则
$B \to C$ 在 $\mathfrak r$ 处拟有限。
\end{lemma}

\begin{proof}
注意，$B \to C$ 是有限型的
（引理 \ref{algebra-lemma-compose-finite-type}），所以此陈述有意义。
我们使用引理 \ref{algebra-lemma-isolated-point-fibre} 的刻画 (3)。
若 $A \to C$ 在 $\mathfrak r$ 处拟有限，则存在某个
$c \in C$，使得
$$
\{\mathfrak r' \subset C \text{ 位于下列对象之上： }\mathfrak p\} \cap D(c) =
\{\mathfrak{r}\}.
$$
由于位于 $\mathfrak q$ 之上的素理想
$\mathfrak r' \subset C$ 构成位于 $\mathfrak p$ 之上的素理想
$\mathfrak r' \subset C$ 的一个子集，所以 $B \to C$ 在
$\mathfrak r$ 处拟有限。
\end{proof}

\begin{lemma}
\label{algebra-lemma-generically-finite}
设 $R \to S$ 是有限型环映射。
设 $\mathfrak p \subset R$ 是极小素理想。
假设 $S$ 中位于 $\mathfrak p$ 之上的素理想至多有限多个。
则存在 $g \in R$、$g \not \in \mathfrak p$，使得环映射
$R_g \to S_g$ 是有限的。
\end{lemma}

\begin{proof}
设 $x_1, \ldots, x_n$ 是 $S$ 在 $R$ 上的一组生成元。
由引理 \ref{algebra-lemma-quasi-finite-above-p}，该假设表示
$S \otimes_R \kappa(\mathfrak p) = S_{\mathfrak p}/\mathfrak pS_{\mathfrak p}$
是有限 $\kappa(\mathfrak p)$-代数。
因此可以找到首一多项式
$P_i \in R_{\mathfrak p}[X]$，使得 $P_i(x_i)$ 在
$S_{\mathfrak p}/\mathfrak pS_{\mathfrak p}$ 中的像为零。
由于 $\mathfrak p$ 是极小素理想，
$\mathfrak pR_{\mathfrak p}$ 是局部幂零理想；见引理
\ref{algebra-lemma-minimal-prime-reduced-ring}。
所以由引理 \ref{algebra-lemma-locally-nilpotent}，
$\mathfrak pS_{\mathfrak p}$ 是局部幂零理想。
% P01-E0587: the frozen source twice uses undefined P where the indexed polynomial is P_i.
因此存在 $e_i \geq 1$，使得 $P(x_i)^{e_i} = 0$
在 $S_{\mathfrak p}$ 中成立。取 $g_1 \in R$、
$g_1 \not \in \mathfrak p$，使 $P_i$ 对所有 $i$ 都有系数属于
$R[1/g_1]$。再取 $g_2 \in R$、$g_2 \not \in \mathfrak p$，
使得 $P(x_i)^{e_i} = 0$ 在 $S_{g_1g_2}$ 中成立。
置 $g = g_1g_2$ 即得结论。
\end{proof}

\section{Zariski 主定理}
\label{algebra-section-Zariski}

\noindent
本节的目标是证明 Zariski 主定理的代数版本。
在 Stacks project 后续关于概形及概形态射的理论发展中，
这个定理将成为许多结果的基础。

\medskip\noindent
设 $R \to S$ 是有限型环映射。
本节的目标是证明，$\Spec(S)$ 中映射拟有限的点所组成的集合
是\textit{开集}（定理 \ref{algebra-theorem-main-theorem}）。
事实上，将会看到存在有限环映射 $R \to S'$，使得在某种意义下，
$S/R$ 的拟有限点集在 $\Spec(S')$ 中是开的
（但我们不会在代数章中证明这一点，因为这里不发展概形语言；
当 $R \to S$ 拟有限时，见引理
\ref{algebra-lemma-quasi-finite-open-integral-closure}）。
这些陈述的证明颇为棘手，我们通过一长串关于环的整扩张与有限扩张的引理
来完成证明。这些材料可见 \cite{Henselian} 和 \cite{Peskine}。
Thierry Coquand 的笔记对我们也很有帮助。

\begin{lemma}
\label{algebra-lemma-make-integral-trivial}
设 $\varphi : R \to S$ 是环映射。
假设 $t \in S$ 满足关系
$\varphi(a_0) + \varphi(a_1)t + \ldots + \varphi(a_n) t^n = 0$。
则 $\varphi(a_n)t$ 在 $R$ 上是整的。
\end{lemma}

\begin{proof}
事实上，将等式
$\varphi(a_0) + \varphi(a_1)t + \ldots + \varphi(a_n) t^n = 0$
乘以 $\varphi(a_n)^{n-1}$，并将其写成
$\varphi(a_0 a_n^{n-1}) +
\varphi(a_1 a_n^{n-2}) (\varphi(a_n)t) +
\ldots +
(\varphi(a_n) t)^n = 0$。
\end{proof}

\noindent
下面的引理在某种意义下是本节的关键引理。

\begin{lemma}
\label{algebra-lemma-make-integral-trick}
设 $R$ 是环。设 $\varphi : R[x] \to S$ 是环映射。
设 $t \in S$。
假设 (a) $t$ 在 $R[x]$ 上是整的，且
(b) 存在首一多项式 $p \in R[x]$，使得
$t \varphi(p) \in \Im(\varphi)$。则存在 $q \in R[x]$，
使得 $t - \varphi(q)$ 在 $R$ 上是整的。
\end{lemma}

\begin{proof}
对某个 $r \in R[x]$，写成 $t \varphi(p) = \varphi(r)$。
利用欧几里得除法，写成 $r = qp + r'$，其中
$q, r' \in R[x]$ 且 $\deg(r') < \deg(p)$。
可以用 $t - \varphi(q)$ 代替 $t$；它在 $R[x]$ 上仍是整的，
并且这样便有 $t \varphi(p) = \varphi(r')$。
在环 $S_t$ 中，可将它写成
$\varphi(p) - (1/t) \varphi(r') = 0$。
这说明 $\varphi(x)$ 给出局部化 $S_t$ 的一个元素，
它在 $\varphi(R)[1/t] \subset S_t$ 上是整的。
另一方面，$t$ 在子环 $\varphi(R)[\varphi(x)] \subset S$ 上是整的。
合起来，由引理 \ref{algebra-lemma-integral-transitive}，可知 $t$ 在子环
$\varphi(R)[1/t] \subset S_t$ 上是整的。换言之，在 $S_t$ 中
存在形如
$$
t^d + \sum\nolimits_{i < d}
\left(\sum\nolimits_{j = 0, \ldots, n_i} \varphi(r_{i, j})/t^j\right) t^i = 0
$$
的等式，其中 $r_{i, j} \in R$。这表示，对某个充分大的 $N$，
$S$ 中有
$t^{d + N} +
\sum_{i < d} \sum_{j = 0, \ldots, n_i} \varphi(r_{i, j}) t^{i + N - j} = 0$。
换言之，$t$ 在 $R$ 上是整的。
\end{proof}

\begin{lemma}
\label{algebra-lemma-combine-lemmas}
设 $R$ 是环。设 $\varphi : R[x] \to S$ 是环映射。
设 $t \in S$。假设 $t$ 在 $R[x]$ 上是整的。
设 $p \in R[x]$、$p = a_0 + a_1x + \ldots +
a_k x^k$，并且 $t \varphi(p) \in \Im(\varphi)$。
则存在 $q \in R[x]$ 和 $n \geq 0$，使得
$\varphi(a_k)^n t - \varphi(q)$ 在 $R$ 上是整的。
\end{lemma}

\begin{proof}
设 $R'$ 和 $S'$ 分别是 $R$ 和 $S$ 在元素 $a_k$ 处的局部化。
设 $\varphi' : R'[x] \to S'$ 是 $\varphi$ 的局部化。
设 $t' \in S'$ 是 $t$ 的像。置 $p' = p/a_k \in R'[x]$。
由于 $t \varphi(p) \in \Im(\varphi)$，有
$t' \varphi'(p') \in \Im(\varphi')$。
因为 $p'$ 是首一的，由引理 \ref{algebra-lemma-make-integral-trick}，
存在 $q' \in R'[x]$，使得 $t' - \varphi'(q')$ 在 $R'$ 上是整的。
可以选取 $n \geq 0$ 和元素 $q \in R[x]$，使得
$a_k^n q'$ 是 $q$ 的像。
于是 $\varphi(a_k)^n t - \varphi(q)$ 是 $S$ 的一个元素，
它在 $S'$ 中的像在 $R'$ 上是整的。
由引理 \ref{algebra-lemma-integral-closure-localize}，
存在 $m \geq 0$，使得
$\varphi(a_k)^m(\varphi(a_k)^n t - \varphi(q))$ 在 $R$ 上是整的。
因此 $\varphi(a_k)^{m + n}t - \varphi(a_k^m q)$
在 $R$ 上是整的，如所需。
\end{proof}

\begin{situation}
\label{algebra-situation-one-transcendental-element}
设 $R$ 是环。
设 $\varphi : R[x] \to S$ 是有限的。令
$$
J = \{ g \in S \mid gS \subset \Im(\varphi)\}
$$
为 $\varphi$ 的“导子理想”。
% P01-E0588: the frozen source omits “is” in the integrally-closed assumption.
假设 $\varphi(R) \subset S$ 在 $S$ 中是整闭的。
\end{situation}

\begin{lemma}
\label{algebra-lemma-leading-coefficient-in-J}
% P01-E0589: the frozen source makes “In Situation ...” a sentence fragment.
在情形 \ref{algebra-situation-one-transcendental-element} 中，
假设 $u \in S$、$a_0, \ldots, a_k \in R$，且
$u \varphi(a_0 + a_1x + \ldots + a_k x^k) \in J$。
则存在 $m \geq 0$，使得 $u \varphi(a_k)^m \in J$。
\end{lemma}

\begin{proof}
假设 $S$ 作为 $R[x]$-模由 $t_1, \ldots, t_n$ 生成。
在这种情形下，
$J = \{ g \in S \mid gt_i \in \Im(\varphi)\text{ 对所有 }i\}$。
注意，每个元素 $u t_i$ 都在 $R[x]$ 上整；见引理
\ref{algebra-lemma-finite-is-integral}。
有 $\varphi(a_0 + a_1x + \ldots + a_k x^k) u t_i \in
\Im(\varphi)$。由引理 \ref{algebra-lemma-combine-lemmas}，
对每个 $i$，存在整数 $n_i$ 和元素 $q_i \in R[x]$，使得
$\varphi(a_k^{n_i}) u t_i - \varphi(q_i)$ 在 $R$ 上是整的。
由假设，此元素属于 $\varphi(R)$，所以
$\varphi(a_k^{n_i}) u t_i \in \Im(\varphi)$。
因此 $m = \max\{n_1, \ldots, n_n\}$ 可用。
\end{proof}

\begin{lemma}
\label{algebra-lemma-all-coefficients-in-J}
在情形 \ref{algebra-situation-one-transcendental-element} 中，
假设 $u \in S$、$a_0, \ldots, a_k \in R$，且
$u \varphi(a_0 + a_1x + \ldots + a_k x^k) \in \sqrt{J}$。
则对所有 $i$，有 $u \varphi(a_i) \in \sqrt{J}$。
\end{lemma}

\begin{proof}
在本引理的假设下，对某个 $n \geq 1$，有
$u^n \varphi(a_0 + a_1x + \ldots + a_k x^k)^n \in J$。
由引理 \ref{algebra-lemma-leading-coefficient-in-J}，可知对某个
$m \geq 1$，有 $u^n \varphi(a_k^{nm}) \in J$。
所以 $u \varphi(a_k) \in \sqrt{J}$，从而
$u \varphi(a_0 + a_1x + \ldots + a_k x^k) - u \varphi(a_k x^k) =
u \varphi(a_0 + a_1x + \ldots + a_{k-1} x^{k-1}) \in \sqrt{J}$。
对 $k$ 作归纳即得结论。
\end{proof}

\noindent
这个引理启发了如下定义。

\begin{definition}
\label{algebra-definition-strongly-transcendental}
给定环的包含 $R \subset S$ 和元素 $x \in S$，
如果每当
$u(a_0 + a_1 x + \ldots + a_k x^k) = 0$，其中
$u \in S$ 且 $a_i \in R$，就有对所有 $i$ 都成立的
$ua_i = 0$，那么称 $x$ \textit{在 $R$ 上强超越}。
\end{definition}

\noindent
注意，若 $S$ 是整环，则这等价于说，$x$ 作为 $S$ 的分式域的元素，
在 $R$ 的分式域上超越。

\begin{lemma}
\label{algebra-lemma-reduced-strongly-transcendental-minimal-prime}
假设 $R \subset S$ 是约化环的包含，并且
$x \in S$ 在 $R$ 上强超越。
设 $\mathfrak q \subset S$ 是极小素理想，并令
$\mathfrak p = R \cap \mathfrak q$。
则 $x$ 在 $S/\mathfrak q$ 中的像，在子环
$R/\mathfrak p$ 上强超越。
\end{lemma}

\begin{proof}
假设 $u(a_0 + a_1x + \ldots + a_k x^k) \in \mathfrak q$。
由引理 \ref{algebra-lemma-minimal-prime-reduced-ring}，
局部环 $S_{\mathfrak q}$ 是域，所以
$u(a_0 + a_1x + \ldots + a_k x^k) $ 在 $S_{\mathfrak q}$ 中为零。
因此对某个 $u' \in S$、$u' \not\in \mathfrak q$，有
$uu'(a_0 + a_1x + \ldots + a_k x^k) = 0$。
因为 $x$ 在 $R$ 上强超越，所以对所有 $i$ 都有
$uu'a_i = 0$。这又蕴含 $ua_i \in \mathfrak q$。
\end{proof}

\begin{lemma}
\label{algebra-lemma-domains-transcendental-not-quasi-finite}
假设 $R\subset S$ 是整环的包含，并设 $x \in S$。
假设 $x$ 在 $R$ 上（强）超越，并且 $S$ 在 $R[x]$ 上有限。
则 $R \to S$ 在 $S$ 的任何素理想处都不拟有限。
\end{lemma}

\begin{proof}
先假设 $R$ 正规；见定义 \ref{algebra-definition-ring-normal}。
由引理 \ref{algebra-lemma-polynomial-ring-normal}，$R[x]$ 正规。
取素理想 $\mathfrak q \subset S$，并置
$\mathfrak p = R \cap \mathfrak q$。
假设扩张 $\kappa(\mathfrak p)
\subset \kappa(\mathfrak q)$ 是有限的。
若 $R \to S$ 在 $\mathfrak q$ 处拟有限，就会是这种情形。
令 $\mathfrak r = R[x] \cap \mathfrak q$。
由于 $\kappa(\mathfrak p)
\subset \kappa(\mathfrak r) \subset \kappa(\mathfrak q)$，
扩张 $\kappa(\mathfrak p)
\subset \kappa(\mathfrak r)$ 也是有限的。
所以包含 $\mathfrak r \supset \mathfrak p R[x]$ 是严格的。
由 $R[x] \subset S$ 的下降性质（见命题
\ref{algebra-proposition-going-down-normal-integral}），
可找到位于素理想 $\mathfrak pR[x]$ 之上的素理想
$\mathfrak q' \subset \mathfrak q$。
因此，纤维 $\Spec(S \otimes_R \kappa(\mathfrak p))$
包含一个不等于 $\mathfrak q$ 的点，即 $\mathfrak q'$；
它的闭包包含 $\mathfrak q$，故 $\mathfrak q$ 在其纤维中不是孤立点。

\medskip\noindent
若 $R$ 不正规，设 $R \subset R' \subset K$，其中 $R'$ 是
$R$ 在其分式域 $K$ 中的整闭包。设 $S \subset S' \subset L$，
其中 $S'$ 是在 $S$ 的分式域 $L$ 中由 $R'$ 与 $S$ 生成的子环。
注意，根据构造，映射 $S \otimes_R R' \to S'$ 是满射。
这蕴含 $R'[x] \subset S'$ 是有限的。此外，映射
$S \subset S'$ 在 $\Spec$ 上诱导满射；见引理
\ref{algebra-lemma-integral-overring-surjective}。
由引理 \ref{algebra-lemma-four-rings} 和刚讨论的正规情形即得结论。
\end{proof}

\begin{lemma}
\label{algebra-lemma-reduced-strongly-transcendental-not-quasi-finite}
假设 $R \subset S$ 是约化环的包含。
% P01-E0590: the frozen source says “Assume x in S be strongly transcendental”.
假设 $x \in S$ 在 $R$ 上强超越，并且 $S$ 在 $R[x]$ 上有限。
则 $R \to S$ 在 $S$ 的任何素理想处都不拟有限。
\end{lemma}

\begin{proof}
任取素理想 $\mathfrak q \subset S$。
选取极小素理想 $\mathfrak q' \subset \mathfrak q$。
根据引理
\ref{algebra-lemma-reduced-strongly-transcendental-minimal-prime} 和
\ref{algebra-lemma-domains-transcendental-not-quasi-finite}，扩张
$R/(R \cap \mathfrak q') \subset
S/\mathfrak q'$ 在对应于 $\mathfrak q$ 的素理想处不拟有限。
由引理 \ref{algebra-lemma-four-rings}，扩张 $R \to S$
在 $\mathfrak q$ 处不拟有限。
\end{proof}

\begin{lemma}
\label{algebra-lemma-quasi-finite-monogenic}
设 $R$ 是环。设 $S = R[x]/I$。
设 $\mathfrak q \subset S$ 是素理想。
假设 $R \to S$ 在 $\mathfrak q$ 处拟有限。
设 $S' \subset S$ 是 $R$ 在 $S$ 中的整闭包。
则存在元素 $g \in S'$、$g \not\in \mathfrak q$，使得
$S'_g \cong S_g$。
\end{lemma}

\begin{proof}
设 $\mathfrak p$ 是 $\mathfrak q$ 在 $\Spec(R)$ 中的像。
存在 $f \in I$、$f = a_nx^n + \ldots + a_0$，使得
对某个 $i$ 有 $a_i \not \in \mathfrak p$。事实上，否则纤维环
$S \otimes_R \kappa(\mathfrak p)$ 将是 $\kappa(\mathfrak p)[x]$，
而映射在位于 $\mathfrak p$ 之上的任何素理想处都不会拟有限。
由此可知，存在关系
$b_m x^m + \ldots + b_0 = 0$，其中 $b_j \in S'$、
$j = 0, \ldots, m$，并且对某个 $j$ 有
$b_j \not \in \mathfrak q \cap S'$。
我们对 $m$ 作归纳来证明本引理。
% P01-E0591: the frozen source has the duplicated copula “The base case is m = 0 is vacuous”.
基始情形 $m = 0$ 是空的（因为陈述 $b_0 = 0$ 与
$b_0 \not \in \mathfrak q$ 相矛盾）。

\medskip\noindent
考虑 $b_m \not \in \mathfrak q$ 的情形。此时 $x$ 在
$S'_{b_m}$ 上是整的；事实上，由引理
\ref{algebra-lemma-make-integral-trivial}，有 $b_mx \in S'$。
因此，单射 $S'_{b_m} \to S_{b_m}$ 也是满射，
即它是所需的同构。

\medskip\noindent
考虑 $b_m \in \mathfrak q$ 的情形。此时由引理
\ref{algebra-lemma-make-integral-trivial}，有 $b_mx \in S'$。
置 $b'_{m - 1} = b_mx + b_{m - 1}$。则
% P01-E0592: the frozen source omits punctuation after this displayed relation.
$$
b'_{m - 1}x^{m - 1} + b_{m - 2}x^{m - 2} + \ldots + b_0 = 0
$$
由于 $b'_{m - 1}$ 模 $S' \cap \mathfrak q$ 同余于 $b_{m - 1}$，
仍有 $b'_{m - 1}, b_{m - 2}, \ldots, b_0$ 中的某个元素
不属于 $S' \cap \mathfrak q$。
所以对 $m$ 应用归纳即可。
\end{proof}

\begin{theorem}[Zariski 主定理]
\label{algebra-theorem-main-theorem}
设 $R$ 是环。
% P01-E0593: the frozen source calls the ring map itself “a finite type R-algebra”.
设 $R \to S$ 是有限型 $R$-代数。
设 $S' \subset S$ 是 $R$ 在 $S$ 中的整闭包。
设 $\mathfrak q \subset S$ 是 $S$ 的素理想。
若 $R \to S$ 在 $\mathfrak q$ 处拟有限，则存在
$g \in S'$、$g \not \in \mathfrak q$，使得 $S'_g \cong S_g$。
\end{theorem}

\begin{proof}
存在有限多个元素 $x_1, \ldots, x_n \in S$，使得 $S$
在由 $x_1, \ldots, x_n$ 生成的 $R$-子代数上有限
% P01-E0594: the frozen source writes the standard compound “subalgebra” as “sub algebra”.
（例如可取 $S$ 在 $R$ 上的生成元）。
% P01-E0595: the frozen proof calls this theorem “the proposition”.
我们对这样的最小数目 $n$ 作归纳来证明该命题。

\medskip\noindent
情形 $n = 0$ 是平凡的，因为此时 $S' = S$；
见引理 \ref{algebra-lemma-finite-is-integral}。

\medskip\noindent
考虑 $n = 1$ 的情形。可以把 $R$ 替换为其在 $S$ 中的整闭包
（引理 \ref{algebra-lemma-quasi-finite-permanence} 保证 $R \to S$
在 $\mathfrak q$ 处仍拟有限）。所以可以假设
$R \subset S$ 在 $S$ 中整闭，换言之，$R = S'$。
考虑映射 $\varphi : R[x] \to S$、$x \mapsto x_1$。
（下面将看到 $\varphi$ 不是单射。）
由假设，$\varphi$ 是有限的。因此我们处于情形
\ref{algebra-situation-one-transcendental-element}。
设 $J \subset S$ 为情形
\ref{algebra-situation-one-transcendental-element} 中定义的“导子理想”。
考虑图
$$
\xymatrix{
R[x] \ar[r] & S \ar[r] & S/\sqrt{J} & R/(R \cap \sqrt{J})[x] \ar[l]
\\
& R \ar[lu] \ar[r] \ar[u] & R/(R \cap \sqrt{J}) \ar[u] \ar[ru] &
}
$$
根据引理 \ref{algebra-lemma-all-coefficients-in-J}，
$x$ 在商 $S/\sqrt{J}$ 中的像在
$R/ (R \cap \sqrt{J})$ 上强超越。
因此，由引理 \ref{algebra-lemma-reduced-strongly-transcendental-not-quasi-finite}，
环映射 $R/ (R \cap \sqrt{J}) \to S/\sqrt{J}$
在 $S/\sqrt{J}$ 的任何素理想处都不拟有限。
由引理 \ref{algebra-lemma-four-rings}，可知 $\mathfrak q$
不属于 $V(J) \subset \Spec(S)$。
所以存在元素 $s \in J$、$s \not\in \mathfrak q$。
由 $J$ 的定义，可以对某个多项式 $f \in R[x]$ 写成
$s = \varphi(f)$。令 $I = \Ker(\varphi : R[x] \to S)$。
由于 $\varphi(f) \in J$，有
$(R[x]/I)_f \cong S_{\varphi(f)}$。又因为
$s \not \in \mathfrak q$，有 $f \not \in \varphi^{-1}(\mathfrak q)$。
因此 $\varphi^{-1}(\mathfrak q)/I$ 是 $R[x]/I$ 的一个素理想，
$R \to R[x]/I$ 在该素理想处拟有限；见引理
\ref{algebra-lemma-quasi-finite-local}。注意，$R$ 在 $R[x]/I$ 中整闭，
因为 $R$ 在 $S$ 中整闭。由引理
\ref{algebra-lemma-quasi-finite-monogenic}，存在元素
$h \in R$、$h \not \in R \cap \mathfrak q$，使得
$R_h \cong (R[x]/I)_h$。所以
$(R[x]/I)_{fh} = S_{\varphi(fh)}$ 同构于 $R$ 的某个主局部化
$R_{h'}$，其中 $h' \in R$、$h' \not \in \mathfrak q$。

\medskip\noindent
考虑 $n > 1$ 的情形。令子环 $R' \subset S$ 为
$R[x_1, \ldots, x_{n-1}]$ 在 $S$ 中的整闭包。
由引理 \ref{algebra-lemma-quasi-finite-permanence}，扩张
$S/R'$ 在 $\mathfrak q$ 处拟有限。还要注意，$S$ 在
$R'[x_n]$ 上有限。由上面的 $n = 1$ 情形，存在
$g' \in R'$、$g' \not \in \mathfrak q$，使得
$(R')_{g'} \cong S_{g'}$。此时不能对 $R \to R'$ 应用归纳，
因为 $R'$ 在 $R$ 上可能不是有限型的。
由于 $S$ 在 $R$ 上有限生成，特别可知
$(R')_{g'}$ 在 $R$ 上有限生成。设元素
$g'$ 以及 $y_1/(g')^{n_1}, \ldots, y_N/(g')^{n_N}$，
其中 $y_i \in R'$，生成 $(R')_{g'}$ 作为 $R$-代数。
令 $R''$ 为 $R'$ 的 $R$-子代数，由
$x_1, \ldots, x_{n-1}, y_1, \ldots, y_N, g'$ 生成。
它具有性质 $(R'')_{g'} \cong S_{g'}$。
% P01-E0596: the frozen source uses two sentence fragments for surjectivity and injectivity.
满性来自 $y_i$ 的选取方式；单性来自
$R'' \subset R'$ 以及局部化的正合性。
注意，由 $R'$ 的选取及引理 \ref{algebra-lemma-characterize-integral}，
$R''$ 在 $R[x_1, \ldots, x_{n-1}]$ 上有限。
令 $\mathfrak q'' = R'' \cap \mathfrak q$。由于
$(R'')_{\mathfrak q''} = S_{\mathfrak q}$，由引理
\ref{algebra-lemma-isolated-point-fibre}，$R \to R''$
在 $\mathfrak q''$ 处拟有限。
对 $R \to R''$、$\mathfrak q''$ 以及
$x_1, \ldots, x_{n-1} \in R''$ 应用归纳假设，便得到
在 $R$ 上整的子环 $R''' \subset R''$ 和元素
$g'' \in R'''$、$g'' \not \in \mathfrak q''$，使得
$(R''')_{g''} \cong (R'')_{g''}$。
把 $g'$ 在 $(R'')_{g''}$ 中的像写成
$g'''/(g'')^n$，其中 $g''' \in  R'''$。
置 $g = g''g''' \in R'''$。则显然
$g \not\in \mathfrak q$，且 $(R''')_g \cong S_g$。
根据构造，$R''' \subset S'$，所以也有 $S'_g \cong S_g$，如所需。
\end{proof}

\begin{lemma}
\label{algebra-lemma-quasi-finite-open}
设 $R \to S$ 是有限型环映射。
$\Spec(S)$ 中 $S/R$ 拟有限的点 $\mathfrak q$ 所组成的集合
在 $\Spec(S)$ 中是开的。
\end{lemma}

\begin{proof}
设 $\mathfrak q \subset S$ 是环映射拟有限的一个点。
由定理 \ref{algebra-theorem-main-theorem}，存在整环映射
$R \to S'$、$S' \subset S$ 以及元素
$g \in S'$、$g\not \in \mathfrak q$，使得
$S'_g \cong S_g$。由于 $S$、从而 $S_g$ 在 $R$ 上是有限型的，
可以找到 $S'$ 的有限多个元素 $y_1, \ldots, y_N$，使得
$S''_g \cong S_g$，其中 $S'' \subset S'$ 是由
$g, y_1, \ldots, y_N$ 生成的 $R$-子代数。
由于 $S''$ 在 $R$ 上有限
（见引理 \ref{algebra-lemma-characterize-integral}），由引理
\ref{algebra-lemma-quasi-finite}，$S''$ 在 $R$ 上拟有限。
容易看出，这蕴含 $S''_g$ 在 $R$ 上拟有限；
例如，因为在一个素理想处拟有限这一性质仅依赖于该素理想处的局部环。
所以 $S_g$ 在 $R$ 上拟有限。同理，这蕴含
$R \to S$ 在 $S$ 中属于 $D(g)$ 的每个素理想处拟有限。
\end{proof}

\begin{lemma}
\label{algebra-lemma-quasi-finite-open-integral-closure}
设 $R \to S$ 是有限型环映射。
假设 $S$ 在 $R$ 上拟有限。
设 $S' \subset S$ 是 $R$ 在 $S$ 中的整闭包。则
\begin{enumerate}
\item $\Spec(S) \to \Spec(S')$ 是到一个开子集上的同胚；
\item 若 $g \in S'$ 且 $D(g)$ 包含于该映射的像中，
则 $S'_g \cong S_g$；
\item 存在有限 $R$-代数 $S'' \subset S'$，使得环映射
$S'' \to S$ 满足 (1) 和 (2)。
\end{enumerate}
\end{lemma}

\begin{proof}
因为 $S/R$ 拟有限，可以对 $\Spec(S)$ 的每个点
$\mathfrak q$ 应用定理 \ref{algebra-theorem-main-theorem}。
由于 $\Spec(S)$ 拟紧（见引理 \ref{algebra-lemma-quasi-compact}），
可以选取有限多个 $g_i \in S'$、$i = 1, \ldots, n$，使得
$S'_{g_i} = S_{g_i}$，并且 $g_1, \ldots, g_n$ 在 $S$
中生成单位理想（换言之，$\Spec(S)$ 中与
$g_1, \ldots, g_n$ 相应的标准开集覆盖整个 $\Spec(S)$）。

\medskip\noindent
假设 $D(g) \subset \Spec(S')$ 包含于该映射的像中。
则 $D(g) \subset \bigcup D(g_i)$。换言之，
$g_1, \ldots, g_n$ 在 $S'_g$ 中生成单位理想。
注意，根据 $g_i$ 的选取，有 $S'_{gg_i} \cong S_{gg_i}$。
所以由引理 \ref{algebra-lemma-cover}，$S'_g \cong S_g$。

\medskip\noindent
我们构造 (3) 中的有限代数 $S'' \subset S'$。
为此，注意每个 $S'_{g_i} \cong S_{g_i}$ 都是有限型
$R$-代数。对每个 $i$，选取一些元素 $y_{ij} \in S'$，
使得每个 $S'_{g_i}$ 作为 $R$-代数由 $1/g_i$ 与各个
$y_{ij}$ 生成。然后令 $S''$ 等于 $S'$ 的 $R$-子代数，
由所有 $g_i$ 和所有 $y_{ij}$ 生成。细节略去。
\end{proof}

\section{Zariski 主定理的应用}
\label{algebra-section-apply-main-theorem}

\noindent
下面是一个直接应用：在 $1$ 维半局部环的拟有限映射中，
有限映射恰好是引理 \ref{algebra-lemma-finite-extension-dim-1}
的公式总取等号的那些映射。

\begin{lemma}
\label{algebra-lemma-quasi-finite-extension-dim-1}
设 $A \subset B$ 是整环的扩张。假设
\begin{enumerate}
\item $A$ 是维数为 $1$ 的诺特局部环；
\item $A \to B$ 是有限型的；
\item 所诱导的分式域扩张 $L/K$ 是有限的。
\end{enumerate}
则 $B$ 是半局部环。
设 $x \in \mathfrak m_A$、$x \not = 0$。
设 $\mathfrak m_i$（$i = 1, \ldots, n$）是 $B$ 的极大理想。
则
$$
[L : K]\text{ord}_A(x)
\geq
\sum\nolimits_i
[\kappa(\mathfrak m_i) : \kappa(\mathfrak m_A)]
\text{ord}_{B_{\mathfrak m_i}}(x)
$$
其中 $\text{ord}$ 如定义 \ref{algebra-definition-ord} 所定义。
等号成立当且仅当 $A \to B$ 是有限的。
\end{lemma}

\begin{proof}
由引理 \ref{algebra-lemma-finite-in-codim-1}，环 $B$ 是半局部的。
设 $B'$ 为 $A$ 在 $B$ 中的整闭包。由引理
\ref{algebra-lemma-quasi-finite-open-integral-closure}，
可以找到有限 $A$-子代数 $C \subset B'$，使得置
$\mathfrak n_i = C \cap \mathfrak m_i$ 后，有
$C_{\mathfrak n_i} \cong B_{\mathfrak m_i}$，并且各素理想
$\mathfrak n_1, \ldots, \mathfrak n_n$ 两两不同。
由引理 \ref{algebra-lemma-finite-in-codim-1}，环 $C$ 是半局部的。
设 $\mathfrak p_j$（$j = 1, \ldots, m$）是 $C$ 的其他极大理想
（即“缺失点”）。由引理 \ref{algebra-lemma-finite-extension-dim-1}，有
$$
\text{ord}_A(x^{[L : K]}) =
\sum\nolimits_i
[\kappa(\mathfrak n_i) : \kappa(\mathfrak m_A)]
\text{ord}_{C_{\mathfrak n_i}}(x)
+
\sum\nolimits_j
[\kappa(\mathfrak p_j) : \kappa(\mathfrak m_A)]
\text{ord}_{C_{\mathfrak p_j}}(x)
$$
故不等式成立。在等号成立的情形下，可知 $m = 0$
（没有“缺失点”）。因此 $C \subset B$ 是半局部环的包含，
它在极大理想上诱导双射，并在所有极大理想处的局部化上诱导同构。
所以若 $b \in B$，则
$I = \{x \in C \mid xb \in C\}$ 是 $C$ 的一个理想，
不包含于 $C$ 的任何极大理想，因而 $I = C$，从而 $b \in C$。
所以 $B = C$，且 $B$ 在 $A$ 上有限。
\end{proof}

\noindent
下面是 Zariski 主定理对局部环之间局部同态结构的一个更标准的应用。

\begin{lemma}
\label{algebra-lemma-essentially-finite-type-fibre-dim-zero}
设 $(R, \mathfrak m_R) \to (S, \mathfrak m_S)$ 是局部环的局部同态。
假设
\begin{enumerate}
\item $R \to S$ 本质上是有限型的；
\item $\kappa(\mathfrak m_R) \subset \kappa(\mathfrak m_S)$ 是有限扩张；
\item $\dim(S/\mathfrak m_RS) = 0$。
\end{enumerate}
则 $S$ 是某个有限 $R$-代数的局部化。
\end{lemma}

\begin{proof}
设 $S'$ 是有限型 $R$-代数，使得对 $S'$ 的某个素理想
$\mathfrak q'$ 有 $S = S'_{\mathfrak q'}$。
由定义 \ref{algebra-definition-quasi-finite}，可知
$R \to S'$ 在 $\mathfrak q'$ 处拟有限。
用某个 $g' \in S'$、$g' \not \in \mathfrak q'$ 的
$S'_{g'}$ 代替 $S'$ 后，可以假设 $R \to S'$ 拟有限；
见引理 \ref{algebra-lemma-quasi-finite-open}。
于是由引理 \ref{algebra-lemma-quasi-finite-open-integral-closure}，
存在有限 $R$-代数 $S''$ 以及元素
$g' \in S'$、$g' \not \in \mathfrak q'$ 和 $g'' \in S''$，
使得 $S'_{g'} \cong S''_{g''}$ 为 $R$-代数。
这就证明了本引理。
\end{proof}

\begin{lemma}
\label{algebra-lemma-completion-at-quasi-finite-prime}
设 $R \to S$ 是环映射，$\mathfrak q$ 是 $S$ 的素理想，
位于 $R$ 的 $\mathfrak p$ 之上。若
\begin{enumerate}
\item $R$ 是诺特环；
\item $R \to S$ 是有限型的；
\item $R \to S$ 在 $\mathfrak q$ 处拟有限，
\end{enumerate}
则对某个 $R_\mathfrak p^\wedge$-代数 $B$，有
$R_\mathfrak p^\wedge \otimes_R S = S_\mathfrak q^\wedge \times B$。
\end{lemma}

\begin{proof}
存在有限 $R$-代数 $S' \subset S$ 和元素
$g \in S'$、$g \not \in \mathfrak q' = S' \cap \mathfrak q$，
使得 $S'_g = S_g$，特别地，
$S'_{\mathfrak q'} = S_\mathfrak q$；见引理
\ref{algebra-lemma-quasi-finite-open-integral-closure}。
由引理 \ref{algebra-lemma-completion-finite-extension}，有
$$
R_\mathfrak p^\wedge \otimes_R S' = (S'_{\mathfrak q'})^\wedge \times B'
$$
注意，在这个乘积分解下，$g$ 映到一对 $(u, b')$，其中
$u \in (S'_{\mathfrak q'})^\wedge$ 是单位，
因为 $g \not \in \mathfrak q'$。
$R_\mathfrak p^\wedge \otimes_R S'$ 的乘积分解诱导乘积分解
% P01-E0597: the frozen source omits punctuation after this displayed product decomposition.
$$
R_\mathfrak p^\wedge \otimes_R S = A \times B
$$
由于 $S'_g = S_g$，还有
$(R_\mathfrak p^\wedge \otimes_R S')_g = (R_\mathfrak p^\wedge \otimes_R S)_g$；
又因为 $g \mapsto (u, b')$ 且 $u$ 是单位，所以
$(S'_{\mathfrak q'})^\wedge = A$。
同构 $S'_{\mathfrak q'} = S_\mathfrak q$ 在完备化上确定一个同构，
故这也说明 $A = S_\mathfrak q^\wedge$。
证明至此完成，不过还应作一个合理性检验：诱导映射
$\phi : R_\mathfrak p^\wedge \otimes_R S \to A = S_\mathfrak q^\wedge$
确实是自然映射。对于形如 $x \otimes 1$、
$x \in R_\mathfrak p^\wedge$ 的元素，这是显然的，
因为自然映射 $R_\mathfrak p^\wedge \to S_\mathfrak q^\wedge$
通过 $(S'_{\mathfrak q'})^\wedge$ 分解。
对于形如 $1 \otimes y$、$y \in S$ 的元素，可以这样论证：
对某个 $n \geq 1$，元素 $g^ny$ 是某个 $y' \in S'$ 的像。
% P01-E0598: the frozen source uses “image ... by the composition” rather than “under”.
因此，$\phi(1 \otimes g^ny)$ 是 $y'$ 在复合映射
$S' \to (S'_{\mathfrak q'})^\wedge \to S_\mathfrak q^\wedge$ 下的像，
它等于 $g^ny$ 在映射 $S \to S_\mathfrak q^\wedge$ 下的像。
由于 $g$ 映到单位，这也蕴含 $\phi(1 \otimes y)$ 具有正确的值，
即 $y$ 在 $S \to S_\mathfrak q^\wedge$ 下的像。
\end{proof}

\section{纤维的维数}
\label{algebra-section-dimension-fibres}

\noindent
我们利用 Zariski 主定理研究纤维维数的行为。
回顾一下，我们已在拓扑学定义 \ref{topology-definition-Krull} 中
定义了拓扑空间 $X$ 在点 $x$ 处的维数 $\dim_x(X)$。

\begin{definition}
\label{algebra-definition-relative-dimension}
假设 $R \to S$ 是有限型的，并设
$\mathfrak q \subset S$ 是位于 $R$ 的素理想
$\mathfrak p$ 之上的素理想。
我们把 $\Spec(S \otimes_R \kappa(\mathfrak p))$
在对应于 $\mathfrak q$ 的点处的维数，定义为
$S/R$ \textit{在 $\mathfrak q$ 处的相对维数}，记作
$\dim_{\mathfrak q}(S/R)$。令 $\dim(S/R)$ 为所有
$\dim_{\mathfrak q}(S/R)$ 的上确界，其中 $\mathfrak q$ 取遍所有素理想。
它称为 $S/R$ 的\textit{相对维数}。
\end{definition}

\noindent
特别地，$R \to S$ 在 $\mathfrak q$ 处拟有限，当且仅当
$\dim_{\mathfrak q}(S/R) = 0$。下面的引理大致是
Zariski 主定理的改述。

\begin{lemma}
\label{algebra-lemma-quasi-finite-over-polynomial-algebra}
设 $R \to S$ 是有限型环映射。
设 $\mathfrak q \subset S$ 是素理想。
假设 $\dim_{\mathfrak q}(S/R) = n$。
则存在 $g \in S$、$g \not\in \mathfrak q$，使得
$S_g$ 在多项式代数 $R[t_1, \ldots, t_n]$ 上拟有限。
\end{lemma}

\begin{proof}
环 $\overline{S} = S \otimes_R \kappa(\mathfrak p)$ 是
$\kappa(\mathfrak p)$ 上的有限型代数。
设 $\overline{\mathfrak q}$ 为 $\overline{S}$ 中对应于
$\mathfrak q$ 的素理想。
根据拓扑空间在一点处的维数的定义，存在开集
$U \subset \Spec(\overline{S})$，满足
% P01-E0599: the frozen source drops mathfrak from the barred prime symbol.
$\overline{q} \in U$ 且 $\dim(U) = n$。
由于 $\Spec(\overline{S})$ 上的拓扑由 $\Spec(S)$ 上的拓扑诱导
（见注 \ref{algebra-remark-fundamental-diagram}），可以找到
$g \in S$、$g \not \in \mathfrak q$，其像为
$\overline{g} \in \overline{S}$，并且
$D(\overline{g}) \subset U$。
因此，用 $S_g$ 代替 $S$ 后，有 $\dim(\overline{S}) = n$。

\medskip\noindent
接着，选取 $S$ 作为 $R$-代数的一组生成元
$x_1, \ldots, x_N$。由引理 \ref{algebra-lemma-Noether-normalization}，
在由 $x_1, \ldots, x_N$ 生成的 $S$ 的 $\mathbf{Z}$-子代数中，
存在元素 $y_1, \ldots, y_n$，使得映射
$R[t_1, \ldots, t_n] \to S$、$t_i \mapsto y_i$ 具有如下性质：
% P01-E0600: the frozen polynomial algebra omits the comma after t_1.
$\kappa(\mathfrak p)[t_1\ldots, t_n] \to \overline{S}$ 是有限的。
特别地，$S$ 在 $R[t_1, \ldots, t_n]$ 上于
$\mathfrak q$ 处拟有限。因此，由引理 \ref{algebra-lemma-quasi-finite-open}，
可以用某个 $g\in S$、$g \not \in \mathfrak q$ 的 $S_g$ 代替 $S$，
使得 $R[t_1, \ldots, t_n] \to S$ 拟有限。
\end{proof}

\begin{lemma}
\label{algebra-lemma-refined-quasi-finite-over-polynomial-algebra}
设 $R \to S$ 是环映射。设 $\mathfrak q \subset S$
是位于 $R$ 的素理想 $\mathfrak p$ 之上的素理想。
假设
\begin{enumerate}
\item $R \to S$ 是有限型的；
\item $\dim_{\mathfrak q}(S/R) = n$；
\item $\text{trdeg}_{\kappa(\mathfrak p)}\kappa(\mathfrak q) = r$。
\end{enumerate}
则存在 $f \in R$、$f \not \in \mathfrak p$，
$g \in S$、$g \not\in \mathfrak q$，以及拟有限环映射
$$
\varphi : R_f[x_1, \ldots, x_n] \longrightarrow S_g
$$
使得
% P01-E0601: the frozen lemma statement omits punctuation after this displayed equality.
$\varphi^{-1}(\mathfrak qS_g) =
(\mathfrak p, x_{r + 1}, \ldots, x_n)R_f[x_1, \ldots, x_n]$。
\end{lemma}

\begin{proof}
用一个主局部化代替 $S$ 后，可以假设存在拟有限环映射
$\varphi : R[t_1, \ldots, t_n] \to S$；见引理
\ref{algebra-lemma-quasi-finite-over-polynomial-algebra}。
置 $\mathfrak q' = \varphi^{-1}(\mathfrak q)$。
设 $\overline{\mathfrak q}' \subset \kappa(\mathfrak p)[t_1, \ldots, t_n]$
为对应于 $\mathfrak q'$ 的素理想。
由引理 \ref{algebra-lemma-refined-Noether-normalization}，存在有限环映射
$\kappa(\mathfrak p)[x_1, \ldots, x_n] \to
\kappa(\mathfrak p)[t_1, \ldots, t_n]$，使得
$\overline{\mathfrak q}'$ 的逆像为
$(x_{r + 1}, \ldots, x_n)$。
设 $\overline{h}_i \in \kappa(\mathfrak p)[t_1, \ldots, t_n]$
为 $x_i$ 的像。可以找到元素
$f \in R$、$f \not \in \mathfrak p$ 以及
$h_i \in R_f[t_1, \ldots, t_n]$，它们在
$\kappa(\mathfrak p)[t_1, \ldots, t_n]$ 中映到 $\overline{h}_i$。
于是环映射
$$
R_f[x_1, \ldots, x_n] \longrightarrow R_f[t_1, \ldots, t_n]
$$
与 $\kappa(\mathfrak p)$ 张量后成为有限映射。
特别地，$R_f[t_1, \ldots, t_n]$ 在
$R_f[x_1, \ldots, x_n]$ 上于素理想
$\mathfrak q'R_f[t_1, \ldots, t_n]$ 处拟有限。
因此，由引理 \ref{algebra-lemma-quasi-finite-open}，存在
$g \in R_f[t_1, \ldots, t_n]$、
$g \not \in \mathfrak q'R_f[t_1, \ldots, t_n]$，使得
$R_f[x_1, \ldots, x_n] \to R_f[t_1, \ldots, t_n, 1/g]$
拟有限。所以复合
$$
R_f[x_1, \ldots, x_n] \longrightarrow
R_f[t_1, \ldots, t_n, 1/g] \longrightarrow S_{\varphi(g)}
$$
拟有限，结论得证。
\end{proof}

\begin{lemma}
\label{algebra-lemma-dimension-inequality-quasi-finite}
设 $R \to S$ 是有限型环映射。
设 $\mathfrak q \subset S$ 是位于 $\mathfrak p \subset R$
之上的素理想。若 $R \to S$ 在 $\mathfrak q$ 处拟有限，则
$\dim(S_{\mathfrak q}) \leq \dim(R_{\mathfrak p})$。
\end{lemma}

\begin{proof}
若 $R_{\mathfrak p}$ 是诺特环
（从而 $S_{\mathfrak q}$ 也是诺特环，因为它在
$R_{\mathfrak p}$ 上本质上为有限型），
则这立即由引理 \ref{algebra-lemma-dimension-base-fibre-total} 和定义得出。
在一般情形下，设 $S'$ 是 $R_\mathfrak p$ 在 $S_\mathfrak p$
中的整闭包。由 Zariski 主定理 \ref{algebra-theorem-main-theorem}，
% P01-E0602: the frozen source reverses the direction by saying q' in S' lies over q in S.
对某个位于 $\mathfrak q$ 之上的 $\mathfrak q' \subset S'$，有
$S_{\mathfrak q} = S'_{\mathfrak q'}$。
由引理 \ref{algebra-lemma-integral-dim-up}，有
$\dim(S') \leq \dim(R_\mathfrak p)$，从而更有
$\dim(S_\mathfrak q) = \dim(S'_{\mathfrak q'}) \leq \dim(R_\mathfrak p)$。
\end{proof}

\begin{lemma}
\label{algebra-lemma-dimension-quasi-finite-over-polynomial-algebra}
\begin{slogan}
仿射 n 维空间的拟有限覆盖的维数至多为 n。
\end{slogan}
设 $k$ 是域。设 $S$ 是有限型 $k$-代数。
假设存在拟有限 $k$-代数映射
$k[t_1, \ldots, t_n] \subset S$。则 $\dim(S) \leq n$。
\end{lemma}

\begin{proof}
由引理 \ref{algebra-lemma-dim-affine-space}，
$k[t_1, \ldots, t_n]$ 的任意局部环的维数至多为 $n$。
所以结论由引理 \ref{algebra-lemma-dimension-inequality-quasi-finite} 得出。
\end{proof}

\begin{lemma}
\label{algebra-lemma-dimension-fibres-bounded-open-upstairs}
设 $R \to S$ 是有限型环映射。
设 $\mathfrak q \subset S$ 是素理想。
假设 $\dim_{\mathfrak q}(S/R) = n$。
则 $\mathfrak q$ 在 $\Spec(S)$ 中有一个开邻域 $V$，
使得对所有 $\mathfrak q' \in V$，都有
$\dim_{\mathfrak q'}(S/R) \leq n$。
\end{lemma}

\begin{proof}
由引理 \ref{algebra-lemma-quasi-finite-over-polynomial-algebra}，
可知可以假设 $S$ 在多项式代数
$R[t_1, \ldots, t_n]$ 上拟有限。
考察纤维，便归结为引理
\ref{algebra-lemma-dimension-quasi-finite-over-polynomial-algebra}。
\end{proof}

\noindent
换言之，本引理说明，纤维维数 $\leq n$ 的点所组成的集合
在 $\Spec(S)$ 中是开的。下一个引理说明，这个开集的形成与基变换交换。
若环映射是有限表示的，则此集合是拟紧开集（见下文）。

\begin{lemma}
\label{algebra-lemma-dimension-fibres-bounded-open-upstairs-base-change}
设 $R \to S$ 是有限型环映射。
设 $R \to R'$ 是任意环映射。
置 $S' = R' \otimes_R S$，并以
$f : \Spec(S') \to \Spec(S)$ 表示相应的谱上映射。
设 $n \geq 0$。
逆像
$f^{-1}(\{\mathfrak q \in \Spec(S) \mid
\dim_{\mathfrak q}(S/R) \leq n\})$
等于
$\{\mathfrak q' \in \Spec(S') \mid
\dim_{\mathfrak q'}(S'/R') \leq n\}$。
\end{lemma}

\begin{proof}
该条件用纤维环的维数表述，而这些纤维环是域上的有限型代数。
合用引理 \ref{algebra-lemma-dimension-at-a-point-preserved-field-extension}
即可得到本引理。
\end{proof}

\begin{lemma}
\label{algebra-lemma-dimension-fibres-bounded-quasi-compact-open-upstairs}
设 $R \to S$ 是有限表示环同态。
设 $n \geq 0$。集合
$$
V_n = \{\mathfrak q \in \Spec(S) \mid \dim_{\mathfrak q}(S/R) \leq n\}
$$
是 $\Spec(S)$ 的拟紧开子集。
\end{lemma}

\begin{proof}
由引理 \ref{algebra-lemma-dimension-fibres-bounded-open-upstairs}，它是开的。
% P01-E0603: the frozen proof reuses n both for the fibre-dimension bound and for the number of presentation variables.
设 $S = R[x_1, \ldots, x_n]/(f_1, \ldots, f_m)$ 是 $S$ 的一个表示。
设 $R_0$ 是 $R$ 的 $\mathbf{Z}$-子代数，由各多项式
$f_i$ 的系数生成。令
$S_0 = R_0[x_1, \ldots, x_n]/(f_1, \ldots, f_m)$。
则 $S = R \otimes_{R_0} S_0$。由引理
\ref{algebra-lemma-dimension-fibres-bounded-open-upstairs-base-change}，
$V_n$ 是开集 $V_{0, n}$ 在拟紧连续映射
$\Spec(S) \to \Spec(S_0)$ 下的逆像。
由于 $S_0$ 是诺特环，$V_{0, n}$ 是拟紧的。
\end{proof}

\begin{lemma}
\label{algebra-lemma-finite-type-domain-over-valuation-ring-dim-fibres}
设 $R$ 是剩余域为 $k$、分式域为 $K$ 的赋值环。
设 $S$ 是包含 $R$ 的整环，并且 $S$ 在 $R$ 上为有限型。
若 $S \otimes_R k$ 不是零环，则
$$
\dim(S \otimes_R k) = \dim(S \otimes_R K)
$$
% P01-E0604: the frozen source omits punctuation after the displayed equality.
事实上，$\Spec(S \otimes_R k)$ 是等维的。
\end{lemma}

\begin{proof}
只须证明，对 $S$ 中包含 $\mathfrak m_RS$ 的每个素理想
$\mathfrak q$，$\dim_{\mathfrak q}(S/k)$ 都等于
$\dim(S \otimes_R K)$。取这样的素理想。
由引理 \ref{algebra-lemma-dimension-fibres-bounded-open-upstairs}，
不等式 $\dim_{\mathfrak q}(S/k) \geq \dim(S \otimes_R K)$ 成立。
置 $n = \dim_{\mathfrak q}(S/k)$。
由引理 \ref{algebra-lemma-quasi-finite-over-polynomial-algebra}，
用某个 $g \in S$、$g \not \in \mathfrak q$ 的 $S_g$ 代替 $S$ 后，
存在拟有限环映射 $R[t_1, \ldots, t_n] \to S$。
若 $\dim(S \otimes_R K) < n$，则
$K[t_1, \ldots, t_n] \to S \otimes_R K$ 有非零核。
设 $f = \sum a_I t_1^{i_1}\ldots t_n^{i_n}$。
用 $f$ 的一个具有最小赋值的非零系数除 $f$，可以假设
$f\in R[t_1, \ldots, t_n]$，并且某个 $a_I$ 在 $k$ 中的像非零。
所以环映射 $k[t_1, \ldots, t_n] \to S \otimes_R k$
有非零核，这蕴含 $\dim(S \otimes_R k) < n$，矛盾。
\end{proof}

\section{有限表示的代数与模}
\label{algebra-section-finite-presentation}

\noindent
本节讨论一些标准结果，其中的关键特征是：所用假设包含有限型或有限表示条件。

\begin{lemma}
\label{algebra-lemma-finite-type-descends}
设 $R \to S$ 是环映射。
设 $R \to R'$ 是忠实平坦环映射。
置 $S' = R'\otimes_R S$。
则 $R \to S$ 为有限型的，当且仅当 $R' \to S'$ 为有限型的。
\end{lemma}

\begin{proof}
若 $R \to S$ 为有限型的，则 $R' \to S'$ 显然为有限型的。
假设 $R' \to S'$ 为有限型的。
设 $y_1, \ldots, y_m$ 在 $R'$ 上生成 $S'$。
对某些 $a_{ij} \in R'$ 和 $x_{ji} \in S$，写成
$y_j = \sum_i a_{ij} \otimes x_{ji}$。
% P01-E0605: the frozen proof switches the two indices from x_{ji} to x_{ij}.
设 $A \subset S$ 是由这些 $x_{ij}$ 生成的 $R$-子代数。
由平坦性，我们有 $A' := R' \otimes_R A \subset S'$，而且
按构造 $y_j \in A'$。因此 $A' = S'$。
由忠实平坦性，$A = S$。
\end{proof}

\begin{lemma}
\label{algebra-lemma-finite-presentation-descends}
设 $R \to S$ 是环映射。
设 $R \to R'$ 是忠实平坦环映射。
置 $S' = R'\otimes_R S$。
则 $R \to S$ 为有限表示的，当且仅当 $R' \to S'$ 为有限表示的。
\end{lemma}

\begin{proof}
若 $R \to S$ 为有限表示的，则 $R' \to S'$ 显然为有限表示的。
假设 $R' \to S'$ 为有限表示的。
由引理 \ref{algebra-lemma-finite-type-descends}，$R \to S$ 为有限型的。
写成 $S = R[x_1, \ldots, x_n]/I$。
由平坦性，$S' = R'[x_1, \ldots, x_n]/R'\otimes I$。
设 $g_1, \ldots, g_m$ 在 $R'[x_1, \ldots, x_n]$ 上生成
$R'\otimes I$。对某些 $a_{ij} \in R'$ 和 $f_{ji} \in I$，写成
$g_j = \sum_i a_{ij} \otimes f_{ji}$。
% P01-E0606: the frozen proof switches the two indices from f_{ji} to f_{ij}.
设 $J \subset I$ 是由这些 $f_{ij}$ 生成的理想。
由平坦性，$R' \otimes_R J \subset R'\otimes_R I$，而且两者都是
$R'[x_1, \ldots, x_n]$ 上的理想。
按构造 $g_j \in R' \otimes_R J$。因此
$R' \otimes_R J = R'\otimes_R I$。
由忠实平坦性，$J = I$。
\end{proof}

\begin{lemma}
\label{algebra-lemma-construct-fp-module}
设 $R$ 是环。
设 $I \subset R$ 是理想。
设 $S \subset R$ 是乘法子集。
置 $R' = S^{-1}(R/I) = S^{-1}R/S^{-1}I$。
\begin{enumerate}
\item 对任意有限 $R'$-模 $M'$，存在有限 $R$-模 $M$，使得
$S^{-1}(M/IM) \cong M'$。
\item 对任意有限表示 $R'$-模 $M'$，存在有限表示 $R$-模 $M$，
使得 $S^{-1}(M/IM) \cong M'$。
\end{enumerate}
\end{lemma}

\begin{proof}
证明 (1)。选取短正合列
$0 \to K' \to (R')^{\oplus n} \to M' \to 0$。
设 $K \subset R^{\oplus n}$ 是 $K'$ 在映射
$R^{\oplus n} \to (R')^{\oplus n}$ 下的逆像。
则 $M = R^{\oplus n}/K$ 符合要求。

\medskip\noindent
证明 (2)。
选取表示 $(R')^{\oplus m} \to (R')^{\oplus n} \to M' \to 0$。
假设第一个映射由矩阵 $A' = (a'_{ij})$ 给出，第二个映射由生成元
$x'_i \in M'$、$i = 1, \ldots, n$ 确定。由于 $R' = S^{-1}(R/I)$，
可以选取 $s \in S$ 以及一个系数在 $R$ 中的矩阵 $A = (a_{ij})$，使得
$a'_{ij} = a_{ij} / s \bmod S^{-1}I$。设 $M$ 是具有表示
$R^{\oplus m} \to R^{\oplus n} \to M \to 0$
的有限表示 $R$-模，其中第一个映射由矩阵 $A$ 给出，
第二个映射由生成元 $x_i \in M$、$i = 1, \ldots, n$ 确定。
则映射 $M \to M'$、$x_i \mapsto x'_i$ 诱导同构
$S^{-1}(M/IM) \cong M'$。
\end{proof}

\begin{lemma}
\label{algebra-lemma-construct-fp-module-from-localization}
设 $R$ 是环。
设 $S \subset R$ 是乘法子集。
设 $M$ 是 $R$-模。
\begin{enumerate}
\item 若 $S^{-1}M$ 是有限 $S^{-1}R$-模，则存在有限 $R$-模 $M'$ 以及映射
$M' \to M$，它诱导同构 $S^{-1}M' \to S^{-1}M$。
\item 若 $S^{-1}M$ 是有限表示 $S^{-1}R$-模，则存在有限表示
$R$-模 $M'$ 以及映射 $M' \to M$，它诱导同构
$S^{-1}M' \to S^{-1}M$。
\end{enumerate}
\end{lemma}

\begin{proof}
证明 (1)。设 $x_1, \ldots, x_n \in M$ 是作为 $S^{-1}R$-模生成
$S^{-1}M$ 的元。设 $M'$ 是由 $x_1, \ldots, x_n$ 生成的 $M$ 的
$R$-子模。

\medskip\noindent
证明 (2)。设 $x_1, \ldots, x_n \in M$ 是作为 $S^{-1}R$-模生成
$S^{-1}M$ 的元。设
$K = \Ker(R^{\oplus n} \to M)$，其中该映射由规则
$(a_1, \ldots, a_n) \mapsto \sum a_i x_i$ 给出。由引理 \ref{algebra-lemma-extension}，
$S^{-1}K$ 是有限 $S^{-1}R$-模。
由 (1)，可以找到有限子模 $K' \subset K$，使得
$S^{-1}K' = S^{-1}K$。取
$M' = \Coker(K' \to R^{\oplus n})$。
\end{proof}

\begin{lemma}
\label{algebra-lemma-construct-fp-module-from-stalk}
设 $R$ 是环。
设 $\mathfrak p \subset R$ 是素理想。
设 $M$ 是 $R$-模。
\begin{enumerate}
\item 若 $M_{\mathfrak p}$ 是有限 $R_{\mathfrak p}$-模，则存在有限 $R$-模 $M'$ 以及映射
$M' \to M$，它诱导同构 $M'_{\mathfrak p} \to M_{\mathfrak p}$。
\item 若 $M_{\mathfrak p}$ 是有限表示 $R_{\mathfrak p}$-模，则存在有限表示
$R$-模 $M'$ 以及映射 $M' \to M$，它诱导同构
$M'_{\mathfrak p} \to M_{\mathfrak p}$。
\end{enumerate}
\end{lemma}

\begin{proof}
% P01-E0607: the frozen proof omits terminal punctuation.
这是引理 \ref{algebra-lemma-construct-fp-module-from-localization} 的特殊情形。
\end{proof}

\begin{lemma}
\label{algebra-lemma-local-isomorphism}
设 $\varphi : R \to S$ 是环映射。设 $\mathfrak q \subset S$ 是位于
$\mathfrak p \subset R$ 之上的素理想。假设
\begin{enumerate}
\item $S$ 在 $R$ 上为有限表示的，
\item $\varphi$ 诱导同构 $R_\mathfrak p \cong S_\mathfrak q$。
\end{enumerate}
则存在 $f \in R$、$f \not \in \mathfrak p$ 以及一个 $R_f$-代数 $C$，使得
$S_f \cong R_f \times C$ 为 $R_f$-代数同构。
\end{lemma}

\begin{proof}
写成 $S = R[x_1, \ldots, x_n]/(g_1, \ldots, g_m)$。设 $a_i \in R_\mathfrak p$
是映到 $x_i$ 在 $S_\mathfrak q$ 中之像的元。
对某个 $f \in R$、$f \not \in \mathfrak p$，写成 $a_i = b_i/f$。
用 $R_f$ 代替 $R$，并用 $x_i - a_i$ 代替 $x_i$ 后，可以假设
$S = R[x_1, \ldots, x_n]/(g_1, \ldots, g_m)$，且 $x_i$ 在 $S_\mathfrak q$ 中映到零。
若 $c_j$ 表示 $g_j$ 的常数项，则 $c_j$ 在 $R_\mathfrak p$ 中映到零。
再一次代替 $R$ 后，可以假设这些 $g_j$ 的常数系数 $c_j$ 为零。
因而得到 $R$-代数映射 $S \to R$、$x_i \mapsto 0$，它的核为理想
$(x_1, \ldots, x_n)$。

\medskip\noindent
我们有同构 $R_\mathfrak p \to S_\mathfrak q \to R_\mathfrak p$，而映射
$S \to R$ 把 $x_i$ 送到零。因此有：
$$
S_\mathfrak q = R_\mathfrak p[x_1, \ldots, x_n]/(x_1, \ldots, x_n)
$$
从而有：
$$
S_\mathfrak q = S_\mathfrak p/(x_1, \ldots, x_n)S_\mathfrak p
$$
这意味着有限生成理想 $(x_1, \ldots, x_n)S_\mathfrak p$ 在
$S_\mathfrak p$ 中是纯的，见定义 \ref{algebra-definition-pure-ideal}。
因此，由引理 \ref{algebra-lemma-finitely-generated-pure-ideal}，
$(x_1, \ldots, x_n)S_\mathfrak p$ 由 $S_\mathfrak p$ 中的一个幂等元 $e$ 生成。
对某个 $f \in R$、$f \not \in \mathfrak p$，用 $R_f \to S_f$ 代替
$R \to S$ 后，可以找到一个映到 $e$ 的幂等元 $e' \in S$。
此时 $e'S$ 和 $(x_1, \ldots, x_n)S$ 是有限生成理想，它们在
$S_\mathfrak p$ 中变得相等。因此，对某个 $f \in R$、$f \not \in \mathfrak p$，
用 $R_f \to S_f$ 代替 $R \to S$ 后，可以假设
$e'S = (x_1, \ldots, x_n)S$。置 $C = e'S$ 即完成证明。
\end{proof}

\begin{lemma}
\label{algebra-lemma-isomorphic-local-rings}
设 $R$ 是环。
设 $S$、$S'$ 在 $R$ 上为有限表示的。
设 $\mathfrak q \subset S$ 和 $\mathfrak q' \subset S'$ 是素理想。
若 $S_{\mathfrak q} \cong S'_{\mathfrak q'}$ 为 $R$-代数同构，则存在
$g \in S$、$g \not \in \mathfrak q$ 和 $g' \in S'$、$g' \not \in \mathfrak q'$，
使得 $S_g \cong S'_{g'}$ 为 $R$-代数同构。
\end{lemma}

\begin{proof}
设 $\psi : S_{\mathfrak q} \to S'_{\mathfrak q'}$ 是引理假设中的同构。
写成 $S = R[x_1, \ldots, x_n]/(f_1, \ldots, f_r)$ 和
$S' = R[y_1, \ldots, y_m]/J$。
对每个 $i = 1, \ldots, n$，选取分式 $h_i/g_i$，其中
$h_i, g_i \in R[y_1, \ldots, y_m]$，且 $g_i \bmod J$ 不属于
$\mathfrak q'$，使它表示 $x_i$ 在 $\psi$ 下的像。
% P01-E0608: the frozen replacement sentence omits the second “replacing ... by ...” construction.
用 $S'_{g_1 \ldots g_n}$ 代替 $S'$，并换用
$R[y_1, \ldots, y_m, y_{m + 1}]$（把 $y_{m + 1}$ 映到 $1/(g_1\ldots g_n)$）后，
可以假设 $\psi(x_i)$ 是某个 $h_i \in R[y_1, \ldots, y_m]$ 的像。
考虑这些元
$f_j(h_1, \ldots, h_n) \in R[y_1, \ldots, y_m]$。
由于 $\psi$ 消去每个 $f_j$，存在
$g \in R[y_1, \ldots, y_m]$、$g \bmod J \not \in \mathfrak q'$，使得对每个
$j = 1, \ldots, r$ 都有 $g f_j(h_1, \ldots, h_n) \in J$。
像前面一样用 $S'_g$ 代替 $S'$，并换用
$R[y_1, \ldots, y_m, y_{m + 1}]$ 后，可以假设
$f_j(h_1, \ldots, h_n) \in J$。因而得到环映射
$S \to S'$、$x_i \mapsto h_i$，它在局部环上诱导 $\psi$。
由引理 \ref{algebra-lemma-compose-finite-type}，映射 $S \to S'$ 为有限表示的。
由引理 \ref{algebra-lemma-local-isomorphism}，可以假设 $S' = S \times C$。
% P01-E0609: localizing at the idempotent of the C factor yields C, not the required S factor.
因此，在对应于因子 $C$ 的幂等元处局部化 $S'$，即得结论。
\end{proof}

\begin{lemma}
\label{algebra-lemma-finite-type-mod-nilpotent}
设 $R$ 是环。设 $I \subset R$ 是幂零理想。
设 $S$ 是 $R$-代数，且 $R/I \to S/IS$ 为有限型的。
则 $R \to S$ 为有限型的。
\end{lemma}

\begin{proof}
选取 $s_1, \ldots, s_n \in S$，使它们在 $S/IS$ 中的像作为
$R/I$-代数生成 $S/IS$。由引理 \ref{algebra-lemma-NAK} 的 (11)，
$R$-代数映射 $R[x_1, \ldots, x_n] \to S$、$x_i \mapsto s_i$ 为满射，
结论随之得证。
\end{proof}

\begin{lemma}
\label{algebra-lemma-surjective-mod-locally-nilpotent}
设 $R$ 是环。设 $I \subset R$ 是局部幂零理想。
设 $S \to S'$ 是 $R$-代数映射，且 $S \to S'/IS'$ 为满射，
并且 $S'$ 在 $R$ 上为有限型的。则 $S \to S'$ 为满射。
\end{lemma}

\begin{proof}
对某个理想 $K$，写成 $S' = R[x_1, \ldots, x_m]/K$。
% P01-E0610: the frozen proof uses m variables, switches to an undefined n in three polynomial rings, and then returns to m.
由假设，存在 $g_j = x_j + \sum \delta_{j, J} x^J \in R[x_1, \ldots, x_n]$，
其中 $\delta_{j, J} \in I$，且 $g_j \bmod K \in \Im(S \to S')$。
因此，只须证明 $g_1, \ldots, g_m$ 生成
$R[x_1, \ldots, x_n]$。设 $R_0 \subset R$ 是 $R$ 的一个有限生成
$\mathbf{Z}$-子代数，它至少包含这些 $\delta_{j, J}$。
则 $R_0 \cap I$ 是幂零理想（由引理 \ref{algebra-lemma-Noetherian-power}）。
从而 $R_0[x_1, \ldots, x_n]$ 由 $g_1, \ldots, g_m$ 生成
（因为 $x_j \mapsto g_j$ 定义 $R_0[x_1, \ldots, x_m]$ 的一个自同构；
细节从略）。由于 $R$ 是这些子环 $R_0$ 的并，结论得证。
\end{proof}

\begin{lemma}
\label{algebra-lemma-isomorphism-modulo-ideal}
设 $R$ 是环。设 $I \subset R$ 是理想。设 $S \to S'$ 是
$R$-代数映射。设 $IS \subset \mathfrak q \subset S$ 是素理想。
假设
\begin{enumerate}
\item $S \to S'$ 为满射，
\item $S_\mathfrak q/IS_\mathfrak q \to S'_\mathfrak q/IS'_\mathfrak q$ 是同构，
\item $S$ 在 $R$ 上为有限型的，
% P01-E0611: the frozen fourth item omits the finite verb “is”.
\item $S'$ 在 $R$ 上为有限表示的，以及
\item $S'_\mathfrak q$ 在 $R$ 上平坦。
\end{enumerate}
则对某个 $g \in S$、$g \not \in \mathfrak q$，映射
$S_g \to S'_g$ 是同构。
\end{lemma}

\begin{proof}
设 $J = \Ker(S \to S')$。由引理 \ref{algebra-lemma-compose-finite-type}，
$J$ 是有限生成理想。由于 $S'_\mathfrak q$ 在 $R$ 上平坦，我们有
$J_\mathfrak q/IJ_\mathfrak q \subset S_\mathfrak q/IS_{\mathfrak q}$
（对 $0 \to J \to S \to S' \to 0$ 应用引理 \ref{algebra-lemma-flat-tor-zero}）。
由假设 (2)，$J_\mathfrak q/IJ_\mathfrak q$ 为零。
由 Nakayama 引理（引理 \ref{algebra-lemma-NAK}），存在
$g \in S$、$g \not \in \mathfrak q$，使得 $J_g = 0$。
因此 $S_g \cong S'_g$，正如所求。
\end{proof}

\begin{lemma}
\label{algebra-lemma-isomorphism-modulo-locally-nilpotent}
设 $R$ 是环。设 $I \subset R$ 是理想。设 $S \to S'$ 是
$R$-代数映射。假设
\begin{enumerate}
\item $I$ 是局部幂零的，
\item $S/IS \to S'/IS'$ 是同构，
\item $S$ 在 $R$ 上为有限型的，
\item $S'$ 在 $R$ 上为有限表示的，以及
\item $S'$ 在 $R$ 上平坦。
\end{enumerate}
则 $S \to S'$ 是同构。
\end{lemma}

\begin{proof}
由引理 \ref{algebra-lemma-surjective-mod-locally-nilpotent}，映射
$S \to S'$ 是满射。由于 $I$ 是局部幂零的，理想
$IS$ 和 $IS'$ 也是局部幂零的（引理 \ref{algebra-lemma-locally-nilpotent}）。
因此，$S$ 的每个素理想 $\mathfrak q$ 都包含 $IS$，并且（显然地）
$S_\mathfrak q/IS_\mathfrak q \cong S'_\mathfrak q/IS'_\mathfrak q$。
于是引理 \ref{algebra-lemma-isomorphism-modulo-ideal} 可用，并且对每个素理想
$\mathfrak q \subset S$，映射 $S_\mathfrak q \to S'_\mathfrak q$ 都是同构。
% P01-E0612: the frozen conclusion has the awkward word order “injective for example by”.
例如由引理 \ref{algebra-lemma-characterize-zero-local}，可知 $S \to S'$ 为单射。
\end{proof}

\section{余极限与有限表示映射}
\label{algebra-section-colimits-flat}

\noindent
% P01-E0613: the frozen introduction says “prove result” without an article or plural ending.
本节证明一些预备引理，它们最终将帮助我们使用绝对 Noether 化约
来证明结果。在范畴篇第 \ref{categories-section-directed-colimits} 节中，
我们一般地讨论了滤过余极限。下面是这一十分一般的概念的一个例子。

\begin{lemma}
\label{algebra-lemma-ring-colimit-fp-category}
设 $R \to A$ 是环映射。考虑由所有 $R$-代数映射图式
$A' \to A$ 组成的范畴 $\mathcal{I}$，其中 $A'$ 在 $R$ 上为有限表示的。
则 $\mathcal{I}$ 是滤过的，而且这些 $A'$ 在 $\mathcal{I}$ 上的余极限同构于 $A$。
\end{lemma}

\begin{proof}
范畴\footnote{为避免集合论上的困难，我们只考虑这样的 $A' \to A$：
$A'$ 是 $R[x_1, x_2, x_3, \ldots]$ 的商。}
$\mathcal{I}$ 非空，因为 $R \to R$ 是它的对象。
考虑 $\mathcal{I}$ 的一对对象 $A' \to A$、$A'' \to A$。
则 $A' \otimes_R A'' \to A$ 属于 $\mathcal{I}$（使用引理
\ref{algebra-lemma-compose-finite-type} 和 \ref{algebra-lemma-base-change-finiteness}）。
环映射 $A' \to A' \otimes_R A''$ 和 $A'' \to A' \otimes_R A''$
定义了 $\mathcal{I}$ 中的态射，从而证明了滤过范畴的第二个定义性质；
见范畴篇定义 \ref{categories-definition-directed}。
最后，假设有 $\mathcal{I}$ 中的两个态射
$\sigma, \tau : A' \to A''$。若 $x_1, \ldots, x_r \in A'$ 作为
$R$-代数生成 $A'$，则可以考虑
$A''' = A''/(\sigma(x_i) - \tau(x_i))$。
这是有限表示 $R$-代数，而给定的 $R$-代数映射
$A'' \to A$ 通过满射 $\nu : A'' \to A'''$ 分解。
因此 $\nu$ 是 $\mathcal{I}$ 中使 $\sigma$ 和 $\tau$ 相等的态射，正如所求。

\medskip\noindent
% P01-E0614: the frozen proof calls the filtered index category cofiltered.
我们的指标范畴是余滤过的，这一事实意味着，可以在集合范畴中计算
$B = \colim_{A' \to A} A'$ 的值（略去一些细节；可与范畴篇第
\ref{categories-section-directed-colimits} 节的讨论比较）。
要看出 $B \to A$ 为满射，对每个 $a \in A$，可以使用
$R[x] \to A$、$x \mapsto a$ 来看出 $a$ 在 $B \to A$ 的像中。
反之，若 $b \in B$ 在 $A$ 中映到零，则可以找到
$\mathcal{I}$ 中的 $A' \to A$ 以及 $a' \in A'$，使得该元映到 $b$。
那么 $A'/(a') \to A$ 也属于 $\mathcal{I}$，而映射 $A' \to B$
分解为 $A' \to A'/(a') \to B$，这说明 $b = 0$，正如所求。
\end{proof}

\noindent
通常，在预序集上思考余极限更容易。
% P01-E0615: the frozen sentence “Let ... a preordered set” omits “be”.
设 $(\Lambda, \geq)$ 是预序集。
$\Lambda$ 上的一个环系统由以下数据给出：对每个
$\lambda \in \Lambda$ 给出一个环 $R_\lambda$，并且当
$\lambda \leq \mu$ 时给出一个态射 $R_\lambda \to R_\mu$。
这些态射必须满足如下规则：对所有 $\lambda \leq \mu \leq \nu$，
$R_\lambda \to R_\mu \to R_\nu$ 等于映射 $R_\lambda \to R_\nu$。
见范畴篇第 \ref{categories-section-posets-limits} 节。
% P01-E0616: the frozen text switches from Lambda to an undefined I in the directedness assumption.
我们将经常假设 $(I, \leq)$ 是 {\it 有向的}，这意味着 $\Lambda$ 非空，
并且对给定的 $\lambda, \mu \in \Lambda$，存在 $\nu \in \Lambda$，使得
$\lambda \leq \nu$ 和 $\mu \leq \nu$。
回忆，在这种情形下，余极限 $\colim_\lambda R_\lambda$
有时称为“直极限”（但我们不使用这一术语）。

\medskip\noindent
注意，范畴篇引理 \ref{categories-lemma-directed-category-system} 告诉我们，
滤过指标范畴上的余极限与有向集上的余极限是同一回事。

\begin{lemma}
\label{algebra-lemma-ring-colimit-fp}
设 $R \to A$ 是环映射。存在一个由有限表示 $R$-代数 $A_\lambda$
组成的有向系统，使得 $A = \colim_\lambda A_\lambda$。
若 $A$ 在 $R$ 上为有限型的，则可以把它安排成：
$A_\lambda$ 的该系统中所有转移映射都是满射。
\end{lemma}

\begin{proof}
第一种证明是：这由引理 \ref{algebra-lemma-ring-colimit-fp-category} 和
范畴篇引理 \ref{categories-lemma-directed-category-system} 得出。

\medskip\noindent
第二种证明。
与引理 \ref{algebra-lemma-module-colimit-fp} 的证明比较。
考虑任意有限子集 $S \subset A$，以及 $S$ 中元之间的任意有限个
多项式关系所组成的集合 $E$。
于是每个 $s \in S$ 对应于 $x_s \in A$，而每个 $e \in E$
由一个多项式 $f_e \in R[X_s; s\in S]$ 组成，它满足 $f_e(x_s) = 0$。
设 $A_{S, E} = R[X_s; s\in S]/(f_e; e\in E)$，
它是有限表示 $R$-代数。
存在典范映射 $A_{S, E} \to A$。
若 $S \subset S'$，并且 $E$ 的各元经映射
$R[X_s; s \in S] \to R[X_s; s\in S']$ 对应于 $E'$ 的一个子集，
则有一个明显的映射 $A_{S, E} \to A_{S', E'}$，它与到 $A$ 的映射可换。
% P01-E0617: the frozen phrase “setting Lambda equal the set” omits “to”.
因此，令 $\Lambda$ 等于按上述包含关系排序的对 $(S, E)$ 所组成的集合，
就得到一个有向偏序集。
显然，这一有向系统的余极限是 $A$。

\medskip\noindent
对最后一个断言，假设 $A = R[x_1, \ldots, x_n]/I$。
在这种情形下，考虑子集 $\Lambda' \subset \Lambda$，它由上述那些满足
$S = \{x_1, \ldots, x_n\}$ 的系统 $(S, E)$ 组成。
容易看出，仍然有 $A = \colim_{\lambda' \in \Lambda'} A_{\lambda'}$。
此外，转移映射显然都是满射。
\end{proof}

\noindent
事实证明，可以如下刻画有限表示环映射。
在某种意义下，这说有限表示代数是 $R$-代数范畴中的“紧”对象。

\begin{lemma}
\label{algebra-lemma-characterize-finite-presentation}
设 $\varphi : R \to S$ 是环映射。下列条件等价：
\begin{enumerate}
\item $\varphi$ 为有限表示的；
\item 对每个 $R$-代数的有向系统 $A_\lambda$，映射
$$
\colim_\lambda \Hom_R(S, A_\lambda) \longrightarrow
\Hom_R(S, \colim_\lambda A_\lambda)
$$
是双射；以及
\item 对每个 $R$-代数的有向系统 $A_\lambda$，映射
$$
\colim_\lambda \Hom_R(S, A_\lambda) \longrightarrow
\Hom_R(S, \colim_\lambda A_\lambda)
$$
是满射。
\end{enumerate}
\end{lemma}

\begin{proof}
假设 (1)，并写成 $S = R[x_1, \ldots, x_n] / (f_1, \ldots, f_m)$。
设 $A = \colim A_\lambda$。注意，$R$-代数同态 $S \to A$ 或
$S \to A_\lambda$ 由 $x_1, \ldots, x_n$ 的像确定。
因此，$\colim_\lambda \Hom_R(S, A_\lambda) \to \Hom_R(S, A)$ 显然为单射。
要看出它为满射，设 $\chi : S \to A$ 是 $R$-代数同态。
则每个 $x_i$ 都映到某个 $A_{\lambda_i}$ 的像中的一个元。
可以选取 $\mu \geq \lambda_i$、$i = 1, \ldots, n$，并假设对
$i = 1, \ldots, n$，$\chi(x_i)$ 是 $y_i \in A_\mu$ 的像。
考虑 $z_j = f_j(y_1, \ldots, y_n) \in A_\mu$。
因为 $\chi$ 是同态，$z_j$ 在 $A = \colim_\lambda A_\lambda$ 中的像为零。
因此存在 $\mu_j \geq \mu$，使得 $z_j$ 在 $A_{\mu_j}$ 中映到零。
选取 $\nu \geq \mu_j$、$j = 1, \ldots, m$。则 $z_1, \ldots, z_m$
在 $A_\nu$ 中的像都是零。这恰好意味着，这些 $y_i$ 映到满足关系
$f_j(y'_1, \ldots, y'_n) = 0$ 的元 $y'_i \in A_\nu$。
从而得到环映射 $S \to A_\nu$。这证明 (1) 蕴含 (2)。

\medskip\noindent
(2) 显然蕴含 (3)。假设 (3)。
由引理 \ref{algebra-lemma-ring-colimit-fp}，可以写成
$S = \colim_\lambda S_\lambda$，其中 $S_\lambda$ 在 $R$ 上为有限表示的。
那么对某个 $\lambda$，恒等映射分解为
$$
S \to S_\lambda \to S
$$
对 $S \to S_\lambda \to S$ 应用引理 \ref{algebra-lemma-compose-finite-type} 的 (4)，
这蕴含 $S$ 在 $S_\lambda$ 上为有限表示的。
对 $R \to S_\lambda \to S$ 应用同一引理的 (2)，得出
$S$ 在 $R$ 上为有限表示的。
\end{proof}

\noindent
使用上面的基本材料，可以如下给出判别一个代数 $A$
何时是某一给定类型代数的滤过余极限的准则。

\begin{lemma}
\label{algebra-lemma-when-colimit}
设 $R \to \Lambda$ 是环映射。设 $\mathcal{E}$ 是一个 $R$-代数集合，且每个
$A \in \mathcal{E}$ 在 $R$ 上为有限表示的。则下列两个断言等价：
\begin{enumerate}
\item $\Lambda$ 是 $\mathcal{E}$ 中元的滤过余极限；以及
% P01-E0619: the frozen phrase “R algebra map” omits the compound hyphen.
\item 对任意 $R$-代数映射 $A \to \Lambda$，若 $A$ 在 $R$ 上为有限表示的，
则可以找到分解 $A \to B \to \Lambda$，其中 $B \in \mathcal{E}$。
\end{enumerate}
\end{lemma}

\begin{proof}
假设 $\mathcal{I} \to \mathcal{E}$、$i \mapsto A_i$ 是滤过图式，且
$\Lambda = \colim_i A_i$。设 $A \to \Lambda$ 是 $R$-代数映射，
且 $A$ 在 $R$ 上为有限表示的。
应用引理 \ref{algebra-lemma-characterize-finite-presentation}，得到分解
$A \to A_i \to \Lambda$。因此 (1) 蕴含 (2)。

\medskip\noindent
考虑引理 \ref{algebra-lemma-ring-colimit-fp-category} 中的范畴 $\mathcal{I}$。
由范畴篇引理 \ref{categories-lemma-cofinal-in-filtered}，
由这些满足 $A \in \mathcal{E}$ 的 $A \to \Lambda$ 组成的充满子范畴
$\mathcal{J}$ 在 $\mathcal{I}$ 中是共终的，并且它是滤过范畴。
再由范畴篇引理 \ref{categories-lemma-cofinal}，$\Lambda$ 也是
$\mathcal{J}$ 上的余极限。
\end{proof}

\noindent
但是还有更强的结果。也就是说，给定 $R = \colim_\lambda R_\lambda$，
% P01-E0620: the frozen text says “the limit of the category” although the following lemma states a colimit of categories.
有限表示 $R$-模的范畴等价于有限表示 $R_\lambda$-模范畴的极限。
对有限表示 $R$-代数的范畴也同样。

\begin{lemma}
\label{algebra-lemma-module-map-property-in-colimit}
设 $A$ 是环，设 $M, N$ 是 $A$-模。
假设 $R = \colim_{i \in I} R_i$ 是 $A$-代数的有向余极限。
\begin{enumerate}
\item 若 $M$ 是有限 $A$-模，且 $u, u' : M \to N$ 是 $A$-模映射，并且
$u \otimes 1 = u' \otimes 1 : M \otimes_A R \to N \otimes_A R$，
则对某个 $i$，有
$u \otimes 1 = u' \otimes 1 : M \otimes_A R_i \to N \otimes_A R_i$。
\item 若 $N$ 是有限 $A$-模，且 $u : M \to N$ 是 $A$-模映射，并且
$u \otimes 1 : M \otimes_A R \to N \otimes_A R$ 为满射，
则对某个 $i$，映射
$u \otimes 1 : M \otimes_A R_i \to N \otimes_A R_i$ 为满射。
\item 若 $N$ 是有限表示 $A$-模，且
$v : N \otimes_A R \to M \otimes_A R$ 是 $R$-模映射，则存在 $i$
和 $R_i$-模映射 $v_i : N \otimes_A R_i \to M \otimes_A R_i$，
使得 $v = v_i \otimes 1$。
\item 若 $M$ 是有限 $A$-模，$N$ 是有限表示 $A$-模，而且
$u : M \to N$ 是 $A$-模映射，使得
$u \otimes 1 : M \otimes_A R \to N \otimes_A R$ 是同构，则对某个 $i$，映射
$u \otimes 1 : M \otimes_A R_i \to N \otimes_A R_i$ 是同构。
\end{enumerate}
\end{lemma}

\begin{proof}
为证明 (1)，假设 $u$ 如 (1) 所述，并设 $x_1, \ldots, x_m \in M$
是生成元。由于 $N \otimes_A R = \colim_i N \otimes_A R_i$，
% P01-E0621: u(x_j) and u'(x_j) belong to N, but the frozen proof puts their equality in M tensor R_i.
可以选取 $i \in I$，使得对 $j = 1, \ldots, m$，
$u(x_j) \otimes 1 = u'(x_j) \otimes 1$ 在 $M \otimes_A R_i$ 中成立。
对这样的 $i$，我们有
$u \otimes 1 = u' \otimes 1 : M \otimes_A R_i \to N \otimes_A R_i$。

\medskip\noindent
为证明 (2)，假设 $u \otimes 1$ 为满射，并设
$y_1, \ldots, y_m \in N$ 是生成元。由于
$N \otimes_A R = \colim_i N \otimes_A R_i$，
% P01-E0622: the frozen proof says elements z_j of M tensor R_i have images in N tensor R, omitting the map u tensor 1.
可以选取 $i \in I$ 和 $z_j \in M \otimes_A R_i$、$j = 1, \ldots, m$，
使得它们在 $N \otimes_A R$ 中的像等于 $y_j \otimes 1$。
对这样的 $i$，映射 $u \otimes 1 : M \otimes_A R_i \to N \otimes_A R_i$
为满射。

\medskip\noindent
为证明 (3)，设 $y_1, \ldots, y_m \in N$ 是生成元。设
$K = \Ker(A^{\oplus m} \to N)$，其中该映射由规则
% P01-E0623: the frozen rule uses undefined x_j although the generators were named y_j.
$(a_1, \ldots, a_m) \mapsto \sum a_j x_j$ 给出。设 $k_1, \ldots, k_t$
是 $K$ 的生成元。写成 $k_s = (k_{s1}, \ldots, k_{sm})$。
由于 $M \otimes_A R = \colim_i M \otimes_A R_i$，可以选取
$i \in I$ 和 $z_j \in M \otimes_A R_i$、$j = 1, \ldots, m$，使得它们在
$M \otimes_A R$ 中的像等于 $v(y_j \otimes 1)$。
我们想用这些 $z_j$ 定义映射
$v_i : N \otimes_A R_i \to M \otimes_A R_i$。
由于 $K \otimes_A R_i \to R_i^{\oplus m} \to N \otimes_A R_i \to 0$
是一个表示，只须验证对每个 $s = 1, \ldots, t$，
$\xi_s = \sum_j k_{sj}z_j$ 在 $M \otimes_A R_i$ 中为零。
这不一定已经成立，但由于 $\xi_s$ 在 $M \otimes_A R$ 中的像为零，
把 $i$ 适当增大后它将成立。

\medskip\noindent
为证明 (4)，假设 $u \otimes 1$ 是同构、$M$ 是有限的，
且 $N$ 是有限表示的。设
$v : N \otimes_A R \to M \otimes_A R$ 是 $u \otimes 1$ 的逆。
应用 (3)，对某个 $i$ 得到映射
$v_i : N \otimes_A R_i \to M \otimes_A R_i$。
% P01-E0624: the frozen identities change the tensor base from A to R, making M tensor_R R_i and N tensor_R R_i generally undefined.
应用 (1)可知，增大 $i$ 后，有
$v_i \circ (u \otimes 1) = \text{id}_{M \otimes_R R_i}$ 和
$(u \otimes 1) \circ v_i = \text{id}_{N \otimes_R R_i}$。
\end{proof}

\begin{lemma}
\label{algebra-lemma-colimit-category-fp-modules}
假设 $R = \colim_{\lambda \in \Lambda} R_\lambda$ 是环的有向余极限。
则有限表示 $R$-模的范畴是有限表示 $R_\lambda$-模范畴的余极限。
更精确地说：
\begin{enumerate}
\item 给定有限表示 $R$-模 $M$，存在 $\lambda \in \Lambda$ 和有限表示
$R_\lambda$-模 $M_\lambda$，使得 $M \cong M_\lambda \otimes_{R_\lambda} R$。
\item 给定 $\lambda \in \Lambda$、有限表示 $R_\lambda$-模
$M_\lambda, N_\lambda$ 以及 $R$-模映射
$\varphi : M_\lambda \otimes_{R_\lambda} R \to N_\lambda \otimes_{R_\lambda} R$，
存在 $\mu \geq \lambda$ 和 $R_\mu$-模映射
$\varphi_\mu : M_\lambda \otimes_{R_\lambda} R_\mu \to
N_\lambda \otimes_{R_\lambda} R_\mu$，使得 $\varphi = \varphi_\mu \otimes 1_R$。
\item 给定 $\lambda \in \Lambda$、有限表示 $R_\lambda$-模
$M_\lambda, N_\lambda$ 以及 $R_\lambda$-模映射
$\varphi, \psi : M_\lambda \to N_\lambda$，若
$\varphi \otimes 1_R = \psi \otimes 1_R$，则对某个 $\mu \geq \lambda$，
$\varphi \otimes 1_{R_\mu} = \psi \otimes 1_{R_\mu}$。
\end{enumerate}
\end{lemma}

\begin{proof}
为证明 (1)，选取表示
$R^{\oplus m} \to R^{\oplus n} \to M \to 0$。
假设第一个映射由矩阵 $A = (a_{ij})$ 给出。
可以选取 $\lambda \in \Lambda$ 以及系数在 $R_\lambda$ 中的矩阵
$A_\lambda = (a_{\lambda, ij})$，使它在 $R$ 中映到 $A$。
然后，只需令 $M_\lambda$ 是具有表示
$R_\lambda^{\oplus m} \to R_\lambda^{\oplus n} \to M_\lambda \to 0$
的 $R_\lambda$-模，其中第一个箭头由 $A_\lambda$ 给出。

\medskip\noindent
(2) 和 (3) 由引理 \ref{algebra-lemma-module-map-property-in-colimit} 得出。
\end{proof}

\begin{lemma}
\label{algebra-lemma-algebra-map-property-in-colimit}
设 $A$ 是环，设 $B, C$ 是 $A$-代数。
假设 $R = \colim_{i \in I} R_i$ 是 $A$-代数的有向余极限。
\begin{enumerate}
\item 若 $B$ 是有限型 $A$-代数，且 $u, u' : B \to C$ 是 $A$-代数映射，并且
$u \otimes 1 = u' \otimes 1 : B \otimes_A R \to C \otimes_A R$，
则对某个 $i$，有
$u \otimes 1 = u' \otimes 1 : B \otimes_A R_i \to C \otimes_A R_i$。
\item 若 $C$ 是有限型 $A$-代数，且 $u : B \to C$ 是 $A$-代数映射，并且
$u \otimes 1 : B \otimes_A R \to C \otimes_A R$ 为满射，
则对某个 $i$，映射
$u \otimes 1 : B \otimes_A R_i \to C \otimes_A R_i$ 为满射。
\item 若 $C$ 在 $A$ 上为有限表示的，且
$v : C \otimes_A R \to B \otimes_A R$ 是 $R$-代数映射，则存在 $i$
和 $R_i$-代数映射 $v_i : C \otimes_A R_i \to B \otimes_A R_i$，
使得 $v = v_i \otimes 1$。
\item 若 $B$ 是有限型 $A$-代数，$C$ 是有限表示 $A$-代数，而且
$u \otimes 1 : B \otimes_A R \to C \otimes_A R$ 是同构，则对某个 $i$，映射
$u \otimes 1 : B \otimes_A R_i \to C \otimes_A R_i$ 是同构。
\end{enumerate}
\end{lemma}

\begin{proof}
为证明 (1)，假设 $u$ 如 (1) 所述，并设
$x_1, \ldots, x_m \in B$ 是生成元。
% P01-E0625: u(x_j) and u'(x_j) belong to C, but the frozen proof invokes the colimit for B and places their equality in B tensor R_i.
由于 $B \otimes_A R = \colim_i B \otimes_A R_i$，可以选取 $i \in I$，
使得对 $j = 1, \ldots, m$，$u(x_j) \otimes 1 = u'(x_j) \otimes 1$
在 $B \otimes_A R_i$ 中成立。对这样的 $i$，我们有
$u \otimes 1 = u' \otimes 1 : B \otimes_A R_i \to C \otimes_A R_i$。

\medskip\noindent
为证明 (2)，假设 $u \otimes 1$ 为满射，并设
$y_1, \ldots, y_m \in C$ 是生成元。由于
$B \otimes_A R = \colim_i B \otimes_A R_i$，
% P01-E0626: the frozen proof omits u tensor 1 when describing the images of z_j in C tensor R.
可以选取 $i \in I$ 和 $z_j \in B \otimes_A R_i$、$j = 1, \ldots, m$，
使得它们在 $C \otimes_A R$ 中的像等于 $y_j \otimes 1$。
对这样的 $i$，映射 $u \otimes 1 : B \otimes_A R_i \to C \otimes_A R_i$
为满射。

\medskip\noindent
为证明 (3)，设 $c_1, \ldots, c_m \in C$ 是生成元。设
% P01-E0627: the frozen presentation maps the polynomial ring to an undefined N instead of C.
$K = \Ker(A[x_1, \ldots, x_m] \to N)$，其中该映射由规则
% P01-E0628: the frozen rule sends every x_j to the index-free sum of c_j rather than to c_j.
$x_j \mapsto \sum c_j$ 给出。设 $f_1, \ldots, f_t$
是 $K$ 作为 $A[x_1, \ldots, x_m]$ 中理想的生成元。
我们把 $f_j = f_j(x_1, \ldots, x_m)$ 看作多项式。
由于 $B \otimes_A R = \colim_i B \otimes_A R_i$，可以选取
$i \in I$ 和 $z_j \in B \otimes_A R_i$、$j = 1, \ldots, m$，使得它们在
$B \otimes_A R$ 中的像等于 $v(c_j \otimes 1)$。
我们想用这些 $z_j$ 定义映射
$v_i : C \otimes_A R_i \to B \otimes_A R_i$。
由于 $K \otimes_A R_i \to R_i[x_1, \ldots, x_m] \to C \otimes_A R_i \to 0$
是一个表示，只须验证：
% P01-E0629: the frozen definition uses f_j while the element and quantifier are indexed by s.
$\xi_s = f_j(z_1, \ldots, z_m)$ 对每个 $s = 1, \ldots, t$
都在 $B \otimes_A R_i$ 中为零。这不一定已经成立，但由于
$\xi_s$ 在 $B \otimes_A R$ 中的像为零，把 $i$ 适当增大后它将成立。

\medskip\noindent
为证明 (4)，假设 $u \otimes 1$ 是同构、$B$ 是有限型
$A$-代数，并且 $C$ 是有限表示 $A$-代数。
% P01-E0630: the frozen inverse v is typed B tensor R to C tensor R, the same direction as u, while the next sentence correctly types v_i in the opposite direction.
设 $v : B \otimes_A R \to C \otimes_A R$ 是 $u \otimes 1$ 的逆。
设 $v_i : C \otimes_A R_i \to B \otimes_A R_i$ 如 (3) 中所述。
% P01-E0631: the frozen identities change the tensor base from A to R.
应用 (1) 可知，增大 $i$ 后，有
$v_i \circ (u \otimes 1) = \text{id}_{B \otimes_R R_i}$ 和
$(u \otimes 1) \circ v_i = \text{id}_{C \otimes_R R_i}$。
\end{proof}

\begin{lemma}
\label{algebra-lemma-colimit-category-fp-algebras}
假设 $R = \colim_{\lambda \in \Lambda} R_\lambda$ 是环的有向余极限。
则有限表示 $R$-代数的范畴是有限表示 $R_\lambda$-代数范畴的余极限。
更精确地说：
\begin{enumerate}
\item 给定有限表示 $R$-代数 $A$，存在 $\lambda \in \Lambda$ 和有限表示
$R_\lambda$-代数 $A_\lambda$，使得 $A \cong A_\lambda \otimes_{R_\lambda} R$。
\item 给定 $\lambda \in \Lambda$、有限表示 $R_\lambda$-代数
$A_\lambda, B_\lambda$ 以及 $R$-代数映射
$\varphi : A_\lambda \otimes_{R_\lambda} R \to B_\lambda \otimes_{R_\lambda} R$，
存在 $\mu \geq \lambda$ 和 $R_\mu$-代数映射
$\varphi_\mu : A_\lambda \otimes_{R_\lambda} R_\mu \to
B_\lambda \otimes_{R_\lambda} R_\mu$，使得 $\varphi = \varphi_\mu \otimes 1_R$。
\item 给定 $\lambda \in \Lambda$、有限表示 $R_\lambda$-代数
$A_\lambda, B_\lambda$ 以及 $R_\lambda$-代数映射
$\varphi_\lambda, \psi_\lambda : A_\lambda \to B_\lambda$，
% P01-E0632: the frozen equalities drop the lambda subscripts from the maps just introduced.
若 $\varphi \otimes 1_R = \psi \otimes 1_R$，则对某个 $\mu \geq \lambda$，
$\varphi \otimes 1_{R_\mu} = \psi \otimes 1_{R_\mu}$。
\end{enumerate}
\end{lemma}

\begin{proof}
为证明 (1)，选取表示
$A = R[x_1, \ldots, x_n]/(f_1, \ldots, f_m)$。
可以选取 $\lambda \in \Lambda$ 和映到
$f_j \in R[x_1, \ldots, x_n]$ 的元
$f_{\lambda, j} \in R_\lambda[x_1, \ldots, x_n]$。
然后只需令
$A_\lambda =
R_\lambda[x_1, \ldots, x_n]/(f_{\lambda, 1}, \ldots, f_{\lambda, m})$。

\medskip\noindent
(2) 和 (3) 由引理 \ref{algebra-lemma-algebra-map-property-in-colimit} 得出。
\end{proof}

\begin{lemma}
\label{algebra-lemma-limit-no-condition-local}
假设 $R \to S$ 是局部环的局部同态。
存在有向集 $(\Lambda, \leq)$ 以及局部环的局部同态
$R_\lambda \to S_\lambda$ 所组成的系统，使得
\begin{enumerate}
\item 系统 $R_\lambda \to S_\lambda$ 的余极限等于 $R \to S$。
\item 每个 $R_\lambda$ 在 $\mathbf{Z}$ 上本质有限型。
\item 每个 $S_\lambda$ 在 $R_\lambda$ 上本质有限型。
\end{enumerate}
\end{lemma}

\begin{proof}
% P01-E0633: the frozen phrase “Denote phi ... the ring map” omits “by”.
以 $\varphi : R  \to S$ 表示该环映射。
设 $\mathfrak m \subset R$ 是 $R$ 的极大理想，设
$\mathfrak n \subset S$ 是 $S$ 的极大理想。令
$$
\Lambda = \{
(A, B)
\mid
A \subset R, B \subset S, \# A < \infty, \# B < \infty, \varphi(A) \subset B
\}.
$$
我们取包含关系作为偏序。对每个
$\lambda = (A, B) \in \Lambda$，设 $R'_\lambda$ 是由 $a \in A$ 生成的
$\mathbf{Z}$-子代数，并设 $S'_\lambda$ 是由 $b$、$b \in B$ 生成的
$\mathbf{Z}$-子代数。设 $R_\lambda$ 是 $R'_\lambda$ 在素理想
$R'_\lambda \cap \mathfrak m$ 处的局部化，并设 $S_\lambda$ 是
$S'_\lambda$ 在素理想 $S'_\lambda \cap \mathfrak n$ 处的局部化。
用图表示为
$$
\xymatrix{
B \ar[r] &
S'_\lambda \ar[r] &
S_\lambda \ar[r] &
S \\
A \ar[r] \ar[u] &
R'_\lambda \ar[r] \ar[u] &
R_\lambda \ar[r] \ar[u] &
R \ar[u]
}.
$$
转移映射是清楚的。其他断言的证明留给读者。
\end{proof}

\begin{lemma}
\label{algebra-lemma-limit-essentially-finite-type}
假设 $R \to S$ 是局部环的局部同态。
假设 $S$ 在 $R$ 上本质有限型。
则存在有向集 $(\Lambda, \leq)$ 以及局部环的局部同态
$R_\lambda \to S_\lambda$ 所组成的系统，使得
\begin{enumerate}
\item 系统 $R_\lambda \to S_\lambda$ 的余极限等于 $R \to S$。
\item 每个 $R_\lambda$ 在 $\mathbf{Z}$ 上本质有限型。
\item 每个 $S_\lambda$ 在 $R_\lambda$ 上本质有限型。
\item 对每个 $\lambda \leq \mu$，映射
$S_\lambda \otimes_{R_\lambda} R_\mu \to S_\mu$
把 $S_\mu$ 表示为 $S_\lambda \otimes_{R_\lambda} R_\mu$ 的一个商的局部化。
\end{enumerate}
\end{lemma}

\begin{proof}
以 $\varphi : R  \to S$ 表示该环映射。
设 $\mathfrak m \subset R$ 是 $R$ 的极大理想，设
$\mathfrak n \subset S$ 是 $S$ 的极大理想。
设 $x_1, \ldots, x_n \in S$ 是这样的元：$S$ 是 $S$ 的由
$x_1, \ldots, x_n$ 生成的 $R$-子代数的一个局部化。
换言之，$S$ 是多项式环 $R[x_1, \ldots, x_n]$ 的一个局部化的商。

\medskip\noindent
设 $\Lambda = \{ A \subset R \mid \# A < \infty\}$ 是 $R$ 的有限子集所组成的集合。
我们取包含关系作为偏序。对每个
$\lambda = A \in \Lambda$，设 $R'_\lambda$ 是由 $a \in A$ 生成的
$\mathbf{Z}$-子代数，并设 $S'_\lambda$ 是由 $\varphi(a)$、$a \in A$
和元 $x_1, \ldots, x_n$ 生成的 $\mathbf{Z}$-子代数。
设 $R_\lambda$ 是 $R'_\lambda$ 在素理想 $R'_\lambda \cap \mathfrak m$ 处的局部化，
并设 $S_\lambda$ 是 $S'_\lambda$ 在素理想
$S'_\lambda \cap \mathfrak n$ 处的局部化。用图表示为
$$
\xymatrix{
\varphi(A) \amalg \{x_i\} \ar[r] &
S'_\lambda \ar[r] &
S_\lambda \ar[r] &
S \\
A \ar[r] \ar[u] &
R'_\lambda \ar[r] \ar[u] &
R_\lambda \ar[r] \ar[u] &
R \ar[u]
}
$$
显然，若 $A \subset B$ 对应于 $\Lambda$ 中的 $\lambda \leq \mu$，
则存在典范映射 $R_\lambda \to R_\mu$ 和 $S_\lambda \to S_\mu$，
从而在有向集 $\Lambda$ 上得到一个系统。

\medskip\noindent
断言 $R = \colim R_\lambda$ 是清楚的，因为所有映射
$R_\lambda \to R$ 都是单射，而 $R$ 的任意元最终都在像中。
同样的论证适用于 $S = \colim S_\lambda$。
断言 (2)、(3) 按构造成立。最后一个断言成立，因为映射
$S'_\lambda \otimes_{R'_\lambda} R'_\mu \to S'_\mu$ 显然是满射。
\end{proof}

\begin{lemma}
\label{algebra-lemma-limit-essentially-finite-presentation}
假设 $R \to S$ 是局部环的局部同态。
假设 $S$ 在 $R$ 上本质有限表示。
% P01-E0634: the frozen statement uses singular “a system of local homomorphism”.
则存在有向集 $(\Lambda, \leq)$ 以及局部环的局部同态
$R_\lambda \to S_\lambda$ 所组成的系统，使得
\begin{enumerate}
\item 系统 $R_\lambda \to S_\lambda$ 的余极限等于 $R \to S$。
\item 每个 $R_\lambda$ 在 $\mathbf{Z}$ 上本质有限型。
\item 每个 $S_\lambda$ 在 $R_\lambda$ 上本质有限型。
\item 对每个 $\lambda \leq \mu$，映射
$S_\lambda \otimes_{R_\lambda} R_\mu \to S_\mu$
把 $S_\mu$ 表示为 $S_\lambda \otimes_{R_\lambda} R_\mu$
在一个素理想处的局部化。
\end{enumerate}
\end{lemma}

\begin{proof}
由假设，可以选取同构
$\Phi : (R[x_1, \ldots, x_n]/I)_{\mathfrak q} \to S$，
其中 $I \subset R[x_1, \ldots, x_n]$ 是有限生成理想，
$\mathfrak q \subset R[x_1, \ldots, x_n]/I$ 是素理想。
（注意，$R \cap \mathfrak q$ 等于 $R$ 的极大理想 $\mathfrak m$。）
我们还选取理想 $I$ 的生成元 $f_1, \ldots, f_m \in I$。
以任意方式把 $R$ 写成有向集 $(\Lambda, \leq )$ 上的余极限
$R = \colim R_\lambda$，其中每个 $R_\lambda$ 都是局部环，
并在 $\mathbf{Z}$ 上本质有限型。
存在某个 $\lambda_0 \in \Lambda$，使得 $f_j$ 是某个
$f_{j, \lambda_0} \in R_{\lambda_0}[x_1, \ldots, x_n]$ 的像。
对所有 $\lambda \geq \lambda_0$，以
$f_{j, \lambda} \in R_{\lambda}[x_1, \ldots, x_n]$ 表示 $f_{j, \lambda_0}$ 的像。
从而得到环映射系统
$$
R_\lambda[x_1, \ldots, x_n]/(f_{1, \lambda}, \ldots, f_{m, \lambda})
\to
R[x_1, \ldots, x_n]/(f_1, \ldots, f_m) \to S
$$
% P01-E0635: the frozen phrase “Set q_lambda the inverse image” omits “to be”.
设 $\mathfrak q_\lambda$ 为 $\mathfrak q$ 的逆像。置
$S_\lambda = (R_\lambda[x_1, \ldots, x_n]/
(f_{1, \lambda}, \ldots, f_{m, \lambda}))_{\mathfrak q_\lambda}$。
这一构造确实符合要求，其验证留给读者。
\end{proof}

\begin{remark}
\label{algebra-remark-suitable-systems-limits}
假设 $R \to S$ 是局部环的局部同态，并且它本质上是有限表示的。
取任意系统 $(\Lambda, \leq)$、$R_\lambda \to S_\lambda$，使它具有引理
\ref{algebra-lemma-limit-essentially-finite-type} 中列出的性质。
可能出现的情况是，这是“错误的”系统；也就是说，
引理 \ref{algebra-lemma-limit-essentially-finite-presentation} 的性质 (4) 可能不成立。
下面是一个例子。设 $k = \mathbf{F}_2$。考虑环
$$
R = \text{localization of }
k[z, y_1, y_2, \ldots]/(y_i^2 - zy_{i + 1})
\text{ 在 }(z, y_1, y_2, \ldots)
$$
置 $S = R/zR$。取 $\Lambda = \mathbf{N}$ 作为系统，并令
$$
R_n = \text{localization of }
k[z, y_1, \ldots, y_n]/(\{y_i^2 - zy_{i + 1}\}_{i \leq n-1})
\text{ 在 }(z, y_1, \ldots, y_n)
$$
以及 $S_n = R_n/(z, y_n^2)$。
% P01-E0636: “All the maps ... are not” has ambiguous quantifier scope; the proof means none of them is a localization.
所有映射 $S_n \otimes_{R_n} R_{n + 1} \to S_{n + 1}$
都不是局部化（在本情形中也就是同构），因为
$1 \otimes y_{n + 1}^2$ 映到零。
若改取 $S_n' = R_n/zR_n$，则映射
$S'_n \otimes_{R_n} R_{n + 1} \to S'_{n + 1}$ 是同构。
本备注的要点是：选择这些系统时确实必须稍加小心。
\end{remark}

\begin{lemma}
\label{algebra-lemma-limit-module-essentially-finite-presentation}
假设 $R \to S$ 是局部环的局部同态。
假设 $S$ 在 $R$ 上本质有限表示。
设 $M$ 是有限表示 $S$-模。
则存在有向集 $(\Lambda, \leq)$，以及局部环的局部同态
$R_\lambda \to S_\lambda$ 所组成的系统和 $S_\lambda$-模 $M_\lambda$，使得
\begin{enumerate}
\item 系统 $R_\lambda \to S_\lambda$ 的余极限等于 $R \to S$。
系统 $M_\lambda$ 的余极限是 $M$。
\item 每个 $R_\lambda$ 在 $\mathbf{Z}$ 上本质有限型。
\item 每个 $S_\lambda$ 在 $R_\lambda$ 上本质有限型。
\item 每个 $M_\lambda$ 在 $S_\lambda$ 上都是有限的。
\item 对每个 $\lambda \leq \mu$，映射
$S_\lambda \otimes_{R_\lambda} R_\mu \to S_\mu$
把 $S_\mu$ 表示为 $S_\lambda \otimes_{R_\lambda} R_\mu$
在一个素理想处的局部化。
\item 对每个 $\lambda \leq \mu$，映射
$M_\lambda \otimes_{S_\lambda} S_\mu \to M_\mu$ 是同构。
\end{enumerate}
\end{lemma}

\begin{proof}
与引理 \ref{algebra-lemma-limit-essentially-finite-presentation} 的证明一样，
首先可以把 $R = \colim R_\lambda$ 写成局部 $\mathbf{Z}$-代数的有向余极限，
其中各代数本质上是有限型的。
接着，可以假设对某个 $\lambda_1 \in \Lambda$，存在
$f_{j, \lambda_1} \in R_{\lambda_1}[x_1, \ldots, x_n]$，使得
$$
S =
\colim_{\lambda \geq \lambda_1} S_\lambda, \text{ 其中 }
S_\lambda =
(R_\lambda[x_1, \ldots, x_n]/
(f_{1, \lambda}, \ldots, f_{m, \lambda}))_{\mathfrak q_\lambda}
$$
选取 $M$ 在 $S$ 上的一个表示
$$
S^{\oplus s} \to S^{\oplus t} \to M \to 0
$$
设 $A \in \text{Mat}(t \times s, S)$ 是该表示的矩阵。
对某个 $\lambda_2 \in \Lambda$、$\lambda_2 \geq \lambda_1$，
可以找到映到 $A$ 的矩阵
$A_{\lambda_2} \in \text{Mat}(t \times s, S_{\lambda_2})$。
对所有 $\lambda \geq \lambda_2$，令
$M_\lambda = \Coker(S_\lambda^{\oplus s} \xrightarrow{A_\lambda}
S_\lambda^{\oplus t})$。这一构造确实符合要求，其验证留给读者。
\end{proof}

\begin{lemma}
\label{algebra-lemma-limit-no-condition}
假设 $R \to S$ 是环映射。
则存在有向集 $(\Lambda, \leq)$ 以及环映射
$R_\lambda \to S_\lambda$ 所组成的系统，使得
\begin{enumerate}
\item 系统 $R_\lambda \to S_\lambda$ 的余极限等于 $R \to S$。
\item 每个 $R_\lambda$ 在 $\mathbf{Z}$ 上为有限型的。
\item 每个 $S_\lambda$ 在 $R_\lambda$ 上为有限型的。
\end{enumerate}
\end{lemma}

\begin{proof}
这是引理 \ref{algebra-lemma-limit-no-condition-local} 的非局部版本。
证明类似，留给读者。
\end{proof}

\begin{lemma}
\label{algebra-lemma-limit-integral}
假设 $R \to S$ 是环映射。
假设 $S$ 在 $R$ 上是整的。
则存在有向集 $(\Lambda, \leq)$ 以及环映射
$R_\lambda \to S_\lambda$ 所组成的系统，使得
\begin{enumerate}
\item 系统 $R_\lambda \to S_\lambda$ 的余极限等于 $R \to S$。
\item 每个 $R_\lambda$ 在 $\mathbf{Z}$ 上为有限型的。
\item 每个 $S_\lambda$ 在 $R_\lambda$ 上都是有限的。
\end{enumerate}
\end{lemma}

\begin{proof}
考虑由对 $(E, F)$ 组成的集合 $\Lambda$，其中 $E \subset R$ 是有限子集，
$F \subset S$ 是有限子集，并且每个元 $f \in F$ 都是某个首一多项式
$P(X) \in R[X]$ 的根，而该多项式的系数属于 $E$。
定义 $(E, F) \leq (E', F')$ 当且仅当 $E \subset E'$ 且 $F \subset F'$。
给定 $\lambda = (E, F) \in \Lambda$，设 $R_\lambda \subset R$ 等于由 $E$ 生成的
$R$ 的 $\mathbf{Z}$-子代数，并设 $S_\lambda \subset S$ 等于由 $F$ 和 $E$
在 $S$ 中的像生成的 $\mathbf{Z}$-子代数。
显然，$R = \colim R_\lambda$。
% P01-E0637: the frozen proof says every element of S is integral over S, which is tautological and does not justify the construction; the hypothesis is integral over R.
我们有 $S = \colim S_\lambda$，因为 $S$ 的每个元在 $S$ 上都是整的。
由引理 \ref{algebra-lemma-characterize-finite-in-terms-of-integral} 以及下述事实，
环映射 $R_\lambda \to S_\lambda$ 是有限的：$S_\lambda$ 在 $R_\lambda$
上由 $F$ 中的元生成，而按我们对对 $(E, F)$ 施加的条件，
这些元在 $R_\lambda$ 上是整的。引理得证。
\end{proof}

\begin{lemma}
\label{algebra-lemma-limit-finite-type}
假设 $R \to S$ 是环映射。
假设 $S$ 在 $R$ 上为有限型的。
则存在有向集 $(\Lambda, \leq)$ 以及环映射
$R_\lambda \to S_\lambda$ 所组成的系统，使得
\begin{enumerate}
\item 系统 $R_\lambda \to S_\lambda$ 的余极限等于 $R \to S$。
\item 每个 $R_\lambda$ 在 $\mathbf{Z}$ 上为有限型的。
\item 每个 $S_\lambda$ 在 $R_\lambda$ 上为有限型的。
\item 对每个 $\lambda \leq \mu$，映射
$S_\lambda \otimes_{R_\lambda} R_\mu \to S_\mu$
把 $S_\mu$ 表示为 $S_\lambda \otimes_{R_\lambda} R_\mu$ 的商。
\end{enumerate}
\end{lemma}

\begin{proof}
这是引理 \ref{algebra-lemma-limit-essentially-finite-type} 的非局部版本。
证明类似，留给读者。
\end{proof}

\begin{lemma}
\label{algebra-lemma-limit-finite-presentation}
假设 $R \to S$ 是环映射。
假设 $S$ 在 $R$ 上为有限表示的。
则存在有向集 $(\Lambda, \leq)$ 以及环映射
$R_\lambda \to S_\lambda$ 所组成的系统，使得
\begin{enumerate}
\item 系统 $R_\lambda \to S_\lambda$ 的余极限等于 $R \to S$。
\item 每个 $R_\lambda$ 在 $\mathbf{Z}$ 上为有限型的。
\item 每个 $S_\lambda$ 在 $R_\lambda$ 上为有限型的。
\item 对每个 $\lambda \leq \mu$，映射
$S_\lambda \otimes_{R_\lambda} R_\mu \to S_\mu$ 是同构。
\end{enumerate}
\end{lemma}

\begin{proof}
这是引理 \ref{algebra-lemma-limit-essentially-finite-presentation} 的非局部版本。
证明类似，留给读者。
\end{proof}

\begin{lemma}
\label{algebra-lemma-limit-module-finite-presentation}
假设 $R \to S$ 是环映射。
假设 $S$ 在 $R$ 上为有限表示的。
设 $M$ 是有限表示 $S$-模。
则存在有向集 $(\Lambda, \leq)$，以及环映射
$R_\lambda \to S_\lambda$ 所组成的系统和 $S_\lambda$-模 $M_\lambda$，使得
\begin{enumerate}
\item 系统 $R_\lambda \to S_\lambda$ 的余极限等于 $R \to S$。
系统 $M_\lambda$ 的余极限是 $M$。
\item 每个 $R_\lambda$ 在 $\mathbf{Z}$ 上为有限型的。
\item 每个 $S_\lambda$ 在 $R_\lambda$ 上为有限型的。
\item 每个 $M_\lambda$ 在 $S_\lambda$ 上都是有限的。
\item 对每个 $\lambda \leq \mu$，映射
$S_\lambda \otimes_{R_\lambda} R_\mu \to S_\mu$ 是同构。
\item 对每个 $\lambda \leq \mu$，映射
$M_\lambda \otimes_{S_\lambda} S_\mu \to M_\mu$ 是同构。
\end{enumerate}
特别地，对每个 $\lambda \in \Lambda$，我们有
$$
M = M_\lambda \otimes_{S_\lambda} S
= M_\lambda \otimes_{R_\lambda} R.
$$
\end{lemma}

\begin{proof}
这是引理 \ref{algebra-lemma-limit-module-essentially-finite-presentation} 的非局部版本。
证明类似，留给读者。
\end{proof}

\section{更多平坦性准则}
\label{algebra-section-more-flatness-criteria}

\noindent
下面的引理在代数几何中常用于证明：从正规曲面到光滑曲面的有限态射是平坦的。
它是引理 \ref{algebra-lemma-finite-flat-over-regular-CM} 的一个部分逆命题，
因为单射有限局部环映射当然满足条件 (3)。

\begin{lemma}
\label{algebra-lemma-CM-over-regular-flat}
\begin{slogan}
奇迹平坦性
\end{slogan}
设 $R \to S$ 是 Noether 局部环的局部同态。假设
\begin{enumerate}
\item $R$ 是正则的，
% P01-E0638: the frozen second assumption omits “is”.
\item $S$ Cohen--Macaulay，
\item $\dim(S) = \dim(R) + \dim(S/\mathfrak m_R S)$。
\end{enumerate}
则 $R \to S$ 是平坦的。
\end{lemma}

\begin{proof}
对 $\dim(R)$ 作归纳。情形 $\dim(R) = 0$ 是显然的，因为此时
$R$ 是域。假设 $\dim(R) > 0$。由 (3)，这蕴含 $\dim(S) > 0$。
设 $\mathfrak q_1, \ldots, \mathfrak q_r$ 是 $S$ 的极小素理想。
注意 $\mathfrak q_i \not \supset \mathfrak m_R S$，因为
$$
\dim(S/\mathfrak q_i) = \dim(S) > \dim(S/\mathfrak m_R S)
$$
% P01-E0639: the frozen sentence lacks punctuation after the display before its explanation continues.
其中第一个等式由引理 \ref{algebra-lemma-maximal-chain-CM} 得出，
不等式由 (3) 得出。因此
$\mathfrak p_i = R \cap \mathfrak q_i$ 不等于 $\mathfrak m_R$。
选取 $x \in \mathfrak m_R$，使得 $x \not \in \mathfrak m_R^2$ 且
$x \not \in \mathfrak p_i$；见引理 \ref{algebra-lemma-silly}。
从而可知 $x$ 不包含于 $S$ 的任一极小素理想中。
因此由 (2)，$x$ 是 $S$ 上的非零因子；见引理
\ref{algebra-lemma-reformulate-CM}，并且 $S/xS$ 是 Cohen--Macaulay 的，而且
$\dim(S/xS) = \dim(S) - 1$。
由 (1) 和引理 \ref{algebra-lemma-regular-ring-CM}，环 $R/xR$ 是正则的，且
$\dim(R/xR) = \dim(R) - 1$。
由归纳可知 $R/xR \to S/xS$ 是平坦的。
因此，由引理 \ref{algebra-lemma-variant-local-criterion-flatness} 及其后的备注，得出结论。
\end{proof}

\begin{lemma}
\label{algebra-lemma-flat-over-regular}
设 $R \to S$ 是 Noether 局部环的同态。
假设 $R$ 是正则局部环，并且一个正则参数系映到 $S$ 中的一个正则列。
则 $R \to S$ 是平坦的。
\end{lemma}

\begin{proof}
假设 $x_1, \ldots, x_d$ 是 $R$ 的一个参数系，它在 $S$ 中映到一个正则列。
注意，$S/(x_1, \ldots, x_d)S$ 在 $R/(x_1, \ldots, x_d)$ 上平坦，
因为后者是域。此时 $x_d$ 是 $S/(x_1, \ldots, x_{d - 1})S$ 中的非零因子，
因此由平坦性局部准则，$S/(x_1, \ldots, x_{d - 1})S$ 在
$R/(x_1, \ldots, x_{d - 1})$ 上平坦（见引理
\ref{algebra-lemma-variant-local-criterion-flatness} 及其后的备注）。
接着，$x_{d - 1}$ 是 $S/(x_1, \ldots, x_{d - 2})S$ 中的非零因子，
因此由平坦性局部准则，$S/(x_1, \ldots, x_{d - 2})S$ 在
$R/(x_1, \ldots, x_{d - 2})$ 上平坦（见引理
\ref{algebra-lemma-variant-local-criterion-flatness} 及其后的备注）。
继续如此，直到得出 $S$ 在 $R$ 上平坦的结论。
\end{proof}

\noindent
下面的引理是证明如下事实的关键：基环上有限表示环上的有限表示模的结果，
可以从 Noether 情形中有限模的相应结果推出。

\begin{lemma}
\label{algebra-lemma-colimit-eventually-flat}
设 $R \to S$、$M$、$\Lambda$、$R_\lambda \to S_\lambda$、$M_\lambda$
如引理 \ref{algebra-lemma-limit-module-essentially-finite-presentation} 中所述。
假设 $M$ 在 $R$ 上平坦。则对某个 $\lambda \in \Lambda$，模
$M_\lambda$ 在 $R_\lambda$ 上平坦。
\end{lemma}

\begin{proof}
选取某个 $\lambda \in \Lambda$，并考虑
$$
\text{Tor}_1^{R_\lambda}(M_\lambda, R_\lambda/\mathfrak m_\lambda)
=
\Ker(\mathfrak m_\lambda \otimes_{R_\lambda} M_\lambda
\to M_\lambda).
$$
见备注 \ref{algebra-remark-Tor-ring-mod-ideal}。右端表明这是有限生成
$S_\lambda$-模（因为 $S_\lambda$ 是 Noether 的，而所涉模都是有限的）。
设 $\xi_1, \ldots, \xi_n$ 是生成元。
由于 $M$ 在 $R$ 上平坦，我们有
$0 = \Ker(\mathfrak m_\lambda R \otimes_R M \to M)$。
由于 $\otimes$ 与余极限可换，存在 $\lambda' \geq \lambda$，
使得每个 $\xi_i$ 都在
$\mathfrak m_{\lambda}R_{\lambda'} \otimes_{R_{\lambda'}} M_{\lambda'}$ 中映到零。
因此映射
$$
\text{Tor}_1^{R_\lambda}(M_\lambda, R_\lambda/\mathfrak m_\lambda)
\longrightarrow
\text{Tor}_1^{R_{\lambda'}}(M_{\lambda'},
R_{\lambda'}/\mathfrak m_{\lambda}R_{\lambda'})
$$
为零。注意，$M_\lambda \otimes_{R_\lambda} R_\lambda/\mathfrak m_\lambda$
在 $R_\lambda/\mathfrak m_\lambda$ 上平坦，因为后一个环是域。
因此可以应用引理 \ref{algebra-lemma-another-variant-local-criterion-flatness}，得到
$M_{\lambda'}$ 在 $R_{\lambda'}$ 上平坦。
\end{proof}

\noindent
使用上面的引理，我们可以开始在非 Noether 情形中重新证明第
\ref{algebra-section-criteria-flatness} 节的各个结果。

\begin{lemma}
\label{algebra-lemma-mod-injective-general}
假设 $R \to S$ 是局部环的局部同态。
% P01-E0640: three nearby statements use “Denote m the maximal ideal” without “by/as”.
以 $\mathfrak m$ 表示 $R$ 的极大理想。
设 $u : M \to N$ 是 $S$-模的映射。假设
\begin{enumerate}
\item $S$ 在 $R$ 上本质有限表示，
\item $M$、$N$ 在 $S$ 上为有限表示的，
\item $N$ 在 $R$ 上平坦，以及
\item $\overline{u} : M/\mathfrak mM \to N/\mathfrak mN$ 是单射。
\end{enumerate}
则 $u$ 是单射，且 $N/u(M)$ 在 $R$ 上平坦。
\end{lemma}

\begin{proof}
由引理 \ref{algebra-lemma-limit-module-essentially-finite-presentation} 及其证明，
可以找到局部环映射系统 $R_\lambda \to S_\lambda$ 以及
$S_\lambda$-模映射 $u_\lambda : M_\lambda \to N_\lambda$，它们对该引理中的
% P01-E0641: the frozen phrase “conclusions (1)--(6) for both N and M of that lemma” has malformed attachment.
$N$ 和 $M$ 都满足结论 (1) -- (6)，并且这些映射 $u_\lambda$ 的余极限是 $u$。
由引理 \ref{algebra-lemma-colimit-eventually-flat}，可以假设对所有充分大的
$\lambda$，$N_\lambda$ 在 $R_\lambda$ 上平坦。
以 $\mathfrak m_\lambda \subset R_\lambda$ 表示极大理想，并以
$\kappa_\lambda = R_\lambda / \mathfrak m_\lambda$，分别以
$\kappa = R/\mathfrak m$ 表示剩余域。

\medskip\noindent
考虑映射
$$
\Psi_\lambda :
M_\lambda/\mathfrak m_\lambda M_\lambda \otimes_{\kappa_\lambda} \kappa
\longrightarrow
M/\mathfrak m M.
$$
由于 $S_\lambda/\mathfrak m_\lambda S_\lambda$ 在域 $\kappa_\lambda$ 上本质有限型，
张量积 $S_\lambda/\mathfrak m_\lambda S_\lambda \otimes_{\kappa_\lambda} \kappa$
在 $\kappa$ 上本质有限型。因此它是 Noether 环，从而得出
$\Psi_\lambda$ 的核是有限生成的。
由于 $M/\mathfrak m M$ 是系统 $M_\lambda/\mathfrak m_\lambda M_\lambda$ 的余极限，
而 $\kappa$ 是这些域 $\kappa_\lambda$ 的余极限，
% P01-E0642: a directed preorder guarantees a later index >= lambda, not necessarily a strict index > lambda.
存在 $\lambda' > \lambda$，使得 $\Psi_\lambda$ 的核由以下映射的核生成：
$$
\Psi_{\lambda, \lambda'} :
M_\lambda/\mathfrak m_\lambda M_\lambda
\otimes_{\kappa_\lambda}
\kappa_{\lambda'}
\longrightarrow
M_{\lambda'}/\mathfrak m_{\lambda'} M_{\lambda'}.
$$
按构造，存在乘法子集
$W \subset S_\lambda \otimes_{R_\lambda} R_{\lambda'}$，使得
$S_{\lambda'} = W^{-1}(S_\lambda \otimes_{R_\lambda} R_{\lambda'})$，且
$$
W^{-1}(M_\lambda/\mathfrak m_\lambda M_\lambda
\otimes_{\kappa_\lambda}
\kappa_{\lambda'})
=
M_{\lambda'}/\mathfrak m_{\lambda'} M_{\lambda'}.
$$
现在，假设 $x$ 是以下映射的核中的元：
$$
\Psi_{\lambda'} :
M_{\lambda'}/\mathfrak m_{\lambda'} M_{\lambda'}
\otimes_{\kappa_{\lambda'}} \kappa
\longrightarrow
M/\mathfrak m M.
$$
% P01-E0643: the frozen membership for wx omits the tensor base and places a localized element back in the earlier module without specifying a representative.
则对某个 $w \in W$，有
$wx \in M_\lambda/\mathfrak m_\lambda M_\lambda \otimes \kappa$。
因此 $wx \in \Ker(\Psi_\lambda)$。所以 $wx$ 是
$\Psi_{\lambda, \lambda'}$ 的核中各元的线性组合。
从而 $wx = 0$ 在 $M_{\lambda'}/\mathfrak m_{\lambda'} M_{\lambda'}
\otimes_{\kappa_{\lambda'}} \kappa$ 中成立；又因为 $w$ 在 $S_{\lambda'}$ 中可逆，有 $x = 0$。
我们得出，对所有充分大的 $\lambda'$，$\Psi_{\lambda'}$ 的核都是零！

\medskip\noindent
由前一段的结果，可以假设对所有充分大的 $\lambda$，
$\Psi_\lambda$ 的核都是零，这蕴含映射
$M_\lambda/\mathfrak m_\lambda M_\lambda \to M/\mathfrak m M$ 是单射。
结合 $\overline{u}$ 为单射这一事实，形式地推出
$\overline{u_\lambda} : M_\lambda/\mathfrak m_\lambda M_\lambda
\to N_\lambda/\mathfrak m_\lambda N_\lambda$ 也是单射。
由引理 \ref{algebra-lemma-mod-injective}，得出（对所有充分大的 $\lambda$）
映射 $u_\lambda$ 是单射，且 $N_\lambda/u_\lambda(M_\lambda)$ 在
$R_\lambda$ 上平坦。引理得证。
\end{proof}

\begin{lemma}
\label{algebra-lemma-grothendieck-general}
假设 $R \to S$ 是局部环的局部环同态。
以 $\mathfrak m$ 表示 $R$ 的极大理想。假设
\begin{enumerate}
\item $S$ 在 $R$ 上本质有限表示，
\item $S$ 在 $R$ 上平坦，以及
\item $f \in S$ 是 $S/{\mathfrak m}S$ 中的非零因子。
\end{enumerate}
则 $S/fS$ 在 $R$ 上平坦，而且 $f$ 是 $S$ 中的非零因子。
\end{lemma}

\begin{proof}
% P01-E0645: the frozen one-sentence proof omits its subject.
这直接由引理 \ref{algebra-lemma-mod-injective-general} 得出。
\end{proof}

\begin{lemma}
\label{algebra-lemma-grothendieck-regular-sequence-general}
假设 $R \to S$ 是局部环的局部环同态。
以 $\mathfrak m$ 表示 $R$ 的极大理想。假设
\begin{enumerate}
\item $R \to S$ 本质上是有限表示的，
\item $R \to S$ 是平坦的，以及
\item $f_1, \ldots, f_c$ 是 $S$ 中的元列，且这些像
$\overline{f}_1, \ldots, \overline{f}_c$ 在 $S/{\mathfrak m}S$ 中构成正则列。
\end{enumerate}
则 $f_1, \ldots, f_c$ 是 $S$ 中的正则列，而且每个商
$S/(f_1, \ldots, f_i)$ 在 $R$ 上平坦。
\end{lemma}

\begin{proof}
% P01-E0646: the frozen proof is the noun phrase “Induction and Lemma ...”.
对列长作归纳，并应用引理 \ref{algebra-lemma-grothendieck-general}。
\end{proof}

\noindent
下面给出平坦性局部准则在局部有限表示的局部环映射情形的版本。

\begin{lemma}
\label{algebra-lemma-variant-local-criterion-flatness-general}
设 $R \to S$ 是局部环的局部同态。
设 $I \not = R$ 是 $R$ 中的理想，并设 $M$ 是 $S$-模。假设
\begin{enumerate}
\item $S$ 在 $R$ 上本质有限表示，
\item $M$ 在 $S$ 上有限表示，
\item $\text{Tor}_1^R(M, R/I) = 0$，以及
\item $M/IM$ 在 $R/I$ 上平坦。
\end{enumerate}
则 $M$ 在 $R$ 上平坦。
\end{lemma}

\begin{proof}
设 $\Lambda$、$R_\lambda \to S_\lambda$、$M_\lambda$ 如引理
\ref{algebra-lemma-limit-module-essentially-finite-presentation} 中所述。
% P01-E0647: the frozen declaration “Denote I_lambda ... the inverse image” omits “by/as”.
以 $I_\lambda \subset R_\lambda$ 表示 $I$ 的逆像。
% P01-E0648: the frozen approximation data uses R_lambda rather than R_lambda/I_lambda as the source of the quotient map.
此时，下列目标与系统数据
$R/I \to S/IS$、$M/IM$、$R_\lambda \to S_\lambda/I_\lambda S_\lambda$
以及 $M_\lambda/I_\lambda M_\lambda$ 也满足引理
\ref{algebra-lemma-limit-module-essentially-finite-presentation} 的各项结论。
因此由引理 \ref{algebra-lemma-colimit-eventually-flat}，可以假设（缩小索引集
$\Lambda$ 之后）对所有 $\lambda$，$M_\lambda/I_\lambda M_\lambda$ 都是平坦的。
% P01-E0649: the frozen statement does not name the base over which this quotient module is flat.
选取某个 $\lambda$，并考虑
$$
\text{Tor}_1^{R_\lambda}(M_\lambda, R_\lambda/I_\lambda)
=
\Ker(I_\lambda \otimes_{R_\lambda} M_\lambda
\to M_\lambda).
$$
见备注 \ref{algebra-remark-Tor-ring-mod-ideal}。右端表明这是有限生成
$S_\lambda$-模（因为 $S_\lambda$ 是 Noether 的，而所涉模都是有限的）。
设 $\xi_1, \ldots, \xi_n$ 是生成元。
由于 $\text{Tor}_1^R(M, R/I) = 0$，并且 $\otimes$ 与余极限可换，
可知存在 $\lambda' \geq \lambda$，使得每个 $\xi_i$ 都在
$\text{Tor}_1^{R_{\lambda'}}(M_{\lambda'}, R_{\lambda'}/I_{\lambda'})$
中映到零。以下各映射的复合
$$
\xymatrix{
R_{\lambda'} \otimes_{R_\lambda}
\text{Tor}_1^{R_\lambda}(M_\lambda, R_\lambda/I_\lambda)
\ar[d]^{\text{surjective by Lemma \ref{algebra-lemma-surjective-on-tor-one}}} \\
\text{Tor}_1^{R_\lambda}(M_\lambda, R_{\lambda'}/I_\lambda R_{\lambda'})
\ar[d]^{\text{surjective up to localization by
Lemma \ref{algebra-lemma-surjective-on-tor-one-trivial}}} \\
\text{Tor}_1^{R_{\lambda'}}(M_{\lambda'}, R_{\lambda'}/I_\lambda R_{\lambda'})
\ar[d]^{\text{surjective by Lemma \ref{algebra-lemma-surjective-on-tor-one}}} \\
\text{Tor}_1^{R_{\lambda'}}(M_{\lambda'}, R_{\lambda'}/I_{\lambda'}).
}
$$
由所示理由，在局部化之后是满射。
这里需要局部化，是因为 $M_{\lambda'}$ 不等于
$M_\lambda \otimes_{R_\lambda} R_{\lambda'}$；事实上，它等于
$M_\lambda \otimes_{S_\lambda} S_{\lambda'}$，且 $S_{\lambda'}$
是 $S_{\lambda} \otimes_{R_\lambda} R_{\lambda'}$ 的局部化，
% P01-E0650: the frozen clause “whence the statement up to a localization” is grammatically malformed.
因而得到局部化之后的断言（或等价地，与 $S_{\lambda'}$ 作张量积之后的断言）。
注意，引理 \ref{algebra-lemma-surjective-on-tor-one} 可用于第一支和第三支箭头，
因为 $M_\lambda/I_\lambda M_\lambda$ 在 $R_\lambda/I_\lambda$ 上平坦，
而 $M_{\lambda'}/I_\lambda M_{\lambda'}$ 在
$R_{\lambda'}/I_\lambda R_{\lambda'}$ 上平坦，后者是平坦模
$M_\lambda/I_\lambda M_\lambda$ 的基变换。
如上所述，该复合把生成元 $\xi_i$ 映到零。
最后得出
$\text{Tor}_1^{R_{\lambda'}}(M_{\lambda'}, R_{\lambda'}/I_{\lambda'})$
为零。由引理 \ref{algebra-lemma-variant-local-criterion-flatness}，这蕴含
$M_{\lambda'}$ 在 $R_{\lambda'}$ 上平坦。
\end{proof}

\noindent
请将下面的引理与引理 \ref{algebra-lemma-criterion-flatness-fibre-Noetherian}
（Noether 局部环的情形）及引理
\ref{algebra-lemma-criterion-flatness-fibre-nilpotent}
（基环中幂零理想的情形）比较。

\begin{lemma}[纤维平坦性准则]
\label{algebra-lemma-criterion-flatness-fibre}
设 $R$、$S$、$S'$ 是局部环，并设 $R \to S \to S'$ 是局部环同态。
设 $M$ 是 $S'$-模。设 $\mathfrak m \subset R$ 是极大理想。假设
\begin{enumerate}
\item 环映射 $R \to S$ 和 $R \to S'$ 本质上都是有限表示的。
\item 模 $M$ 在 $S'$ 上有限表示。
\item 模 $M$ 非零。
\item 模 $M/\mathfrak mM$ 是平坦的 $S/\mathfrak mS$-模。
\item 模 $M$ 是平坦的 $R$-模。
\end{enumerate}
则 $S$ 在 $R$ 上平坦，且 $M$ 是平坦的 $S$-模。
\end{lemma}

\begin{proof}
如引理 \ref{algebra-lemma-limit-essentially-finite-presentation} 的证明中那样，
首先可将 $R = \colim R_\lambda$ 写成局部有限型
$\mathbf{Z}$-代数的有向余极限。
% P01-E0651: the frozen declaration “Denote p_lambda the maximal ideal” omits “by/as”.
以 $\mathfrak p_\lambda$ 表示 $R_\lambda$ 的极大理想。
其次，可以假设对某个 $\lambda_1 \in \Lambda$，存在
$f_{j, \lambda_1} \in R_{\lambda_1}[x_1, \ldots, x_n]$，使得
$$
S =
\colim_{\lambda \geq \lambda_1} S_\lambda, \text{ 其中 }
S_\lambda =
(R_\lambda[x_1, \ldots, x_n]/
(f_{1, \lambda}, \ldots, f_{u, \lambda}))_{\mathfrak q_\lambda}
$$
对某个 $\lambda_2 \in \Lambda$，$\lambda_2 \geq \lambda_1$，存在
$g_{j, \lambda_2} \in R_{\lambda_2}[x_1, \ldots, x_n, y_1, \ldots, y_m]$，
其像为
$\overline{g}_{j, \lambda_2} \in S_{\lambda_2}[y_1, \ldots, y_m]$，
并且
$$
S' =
\colim_{\lambda \geq \lambda_2} S'_\lambda, \text{ 其中 }
S'_\lambda =
(S_\lambda[y_1, \ldots, y_m]/
(\overline{g}_{1, \lambda}, \ldots,
\overline{g}_{v, \lambda}))_{\overline{\mathfrak q}'_\lambda}
$$
注意，这也蕴含
% P01-E0652: the frozen R_lambda-presentation omits the defining f-relations of S_lambda.
$$
S'_\lambda =
(R_\lambda[x_1, \ldots, x_n, y_1, \ldots, y_m]/
(g_{1, \lambda}, \ldots, g_{v, \lambda}))_{\mathfrak q'_\lambda}
$$
选取 $M$ 在 $S'$ 上的一个表示
$$
(S')^{\oplus s} \to (S')^{\oplus t} \to M \to 0
$$
设 $A \in \text{Mat}(t \times s, S')$ 是该表示的矩阵。
对某个 $\lambda_3 \in \Lambda$，$\lambda_3 \geq \lambda_2$，可以找到映到
$A$ 的矩阵
% P01-E0653: the frozen approximating presentation matrix is placed over S_lambda rather than S'_lambda.
$A_{\lambda_3} \in \text{Mat}(t \times s, S_{\lambda_3})$。
对所有 $\lambda \geq \lambda_3$，令
$M_\lambda = \Coker((S'_\lambda)^{\oplus s} \xrightarrow{A_\lambda}
(S'_\lambda)^{\oplus t})$。

\medskip\noindent
按这些选取，对每个 $\lambda_3 \leq \lambda \leq \mu$，
$S_\lambda \otimes_{R_{\lambda}} R_\mu \to S_\mu$ 是局部化，
$S'_\lambda \otimes_{S_{\lambda}} S_\mu \to S'_\mu$ 是局部化，且映射
$M_\lambda \otimes_{S'_\lambda} S'_\mu \to M_\mu$ 是同构。
这还蕴含 $S'_\lambda \otimes_{R_{\lambda}} R_\mu \to S'_\mu$ 是局部化。
因此，由于 $M$ 在 $R$ 上平坦，由引理
\ref{algebra-lemma-colimit-eventually-flat} 可知，对所有充分大的 $\lambda$，
模 $M_\lambda$ 在 $R_\lambda$ 上平坦。
此外，注意
$
\mathfrak m = \colim \mathfrak p_\lambda
$、
$
S/\mathfrak mS = \colim S_\lambda/\mathfrak p_\lambda S_\lambda
$、
$
S'/\mathfrak mS' = \colim S'_\lambda/\mathfrak p_\lambda S'_\lambda
$，并且
$
M/\mathfrak mM = \colim M_\lambda/\mathfrak p_\lambda M_\lambda
$。另外，对每个 $\lambda_3 \leq \lambda \leq \mu$，由上列性质可知
$$
S'_\lambda/\mathfrak p_\lambda S'_\lambda
\otimes_{S_{\lambda}/\mathfrak p_\lambda S_\lambda}
S_\mu/\mathfrak p_\mu S_\mu
\longrightarrow
S'_\mu/\mathfrak p_\mu S'_\mu
$$
是局部化，而映射
$$
M_\lambda / \mathfrak p_\lambda M_\lambda
\otimes_{S'_\lambda/\mathfrak p_\lambda S'_\lambda}
S'_\mu /\mathfrak p_\mu S'_\mu
\longrightarrow
M_\mu/\mathfrak p_\mu M_\mu
$$
是同构。因此，系统
$(S_\lambda/\mathfrak p_\lambda S_\lambda \to
S'_\lambda/\mathfrak p_\lambda S'_\lambda,
M_\lambda/\mathfrak p_\lambda M_\lambda)$
也如引理 \ref{algebra-lemma-limit-module-essentially-finite-presentation} 中所述。
由于假设 $M/\mathfrak m M$ 在 $S/\mathfrak mS$ 上平坦，
可以再次应用引理 \ref{algebra-lemma-colimit-eventually-flat}，得到对所有充分大的
$\lambda$，$M_\lambda/\mathfrak p_\lambda M_\lambda$ 在
$S_\lambda/\mathfrak p_\lambda S_\lambda$ 上平坦。
因此，对充分大的 $\lambda$，数据
$R_\lambda \to S_\lambda \to S'_\lambda, M_\lambda$ 满足引理
\ref{algebra-lemma-criterion-flatness-fibre-Noetherian} 的假设。
选取这样的 $\lambda$。于是
% P01-E0654: the frozen proof asserts equality with the unlocalized base change although the constructed transition map is only a localization.
$S = S_\lambda \otimes_{R_\lambda} R$ 在 $R$ 上平坦，并且
$M = M_\lambda \otimes_{S_\lambda} S$ 在 $S$ 上平坦
（因为平坦模的基变换仍平坦）。
\end{proof}

\noindent
下面的结论容易从“纤维平坦性准则”引理
\ref{algebra-lemma-criterion-flatness-fibre} 推出。关于此类结论的更多内容，
见“更多平坦性”第 \ref{flat-section-introduction} 节。

\begin{lemma}
\label{algebra-lemma-criterion-flatness-fibre-fp-over-ft}
设 $R$、$S$、$S'$ 是局部环，并设 $R \to S \to S'$ 是局部环同态。
设 $M$ 是 $S'$-模。设 $\mathfrak m \subset R$ 是极大理想。假设
\begin{enumerate}
\item $R \to S'$ 本质上是有限表示的，
\item $R \to S$ 本质上是有限型的，
\item $M$ 在 $S'$ 上有限表示，
\item $M$ 非零，
\item $M/\mathfrak mM$ 是平坦的 $S/\mathfrak mS$-模，以及
\item $M$ 是平坦的 $R$-模。
\end{enumerate}
则 $S$ 在 $R$ 上本质有限表示且平坦，并且 $M$ 是平坦的 $S$-模。
\end{lemma}

\begin{proof}
% P01-E0655: the frozen proof assumes the finite-presentation conclusion instead of the finite-type hypothesis.
由于 $S$ 在 $R$ 上本质有限表示，可对某个有限型 $R$-代数 $C$ 写成
$S = C_{\overline{\mathfrak q}}$。
写 $C = R[x_1, \ldots, x_n]/I$。
% P01-E0656: the frozen declaration “Denote q ... the prime ideal” omits “by/as”.
以 $\mathfrak q \subset R[x_1, \ldots, x_n]$ 表示对应于
$\overline{\mathfrak q}$ 的素理想。于是 $S = B/J$，其中
$B = R[x_1, \ldots, x_n]_{\mathfrak q}$ 在 $R$ 上本质有限表示，且
$J = IB$。可以找到 $f_1, \ldots, f_k \in J$，使得像
$\overline{f}_i \in B/\mathfrak mB$ 生成 $J$ 在 Noether 环
$B/\mathfrak mB$ 中的像 $\overline{J}$。因此存在有限生成理想
% P01-E0657: the frozen plural existential is paired with the single symbol J'.
$J' \subset J$，使得 $B/J' \to B/J$ 诱导同构
$$
(B/J') \otimes_R R/\mathfrak m \longrightarrow
B/J \otimes_R R/\mathfrak m = S/\mathfrak mS.
$$
对任意上述 $J'$，引理 \ref{algebra-lemma-criterion-flatness-fibre} 可用于环映射
$$
R \longrightarrow B/J' \longrightarrow S'
$$
和模 $M$。因此，对上述任意选取的 $J'$，$B/J'$ 在 $R$ 上平坦。
现在，若
% P01-E0658: the frozen chain repeats J' where the middle ideal must be J''.
$J' \subset J' \subset J$ 是两个如上的有限生成理想，则可知
$B/J' \to B/J''$ 是本质有限表示的平坦 $R$-代数之间的满射，
% P01-E0659: the frozen relative-clause chain has two “which” clauses with conflicting antecedents.
而且模掉 $\mathfrak m$ 后它是同构。因此，引理
\ref{algebra-lemma-mod-injective-general} 蕴含 $B/J' = B/J''$，即
$J' = J''$。这显然意味着 $J$ 是有限生成的，也就是说，
$S$ 在 $R$ 上本质有限表示。因此可以将引理
\ref{algebra-lemma-criterion-flatness-fibre} 用于 $R \to S \to S'$，结论得证。
\end{proof}

\begin{lemma}[纤维平坦性准则：局部幂零情形]
\label{algebra-lemma-criterion-flatness-fibre-locally-nilpotent}
设
$$
\xymatrix{
S \ar[rr] & & S' \\
& R \ar[lu] \ar[ru]
}
$$
是环范畴中的交换图。设 $I \subset R$ 是局部幂零理想，
$M$ 是 $S'$-模。假设
\begin{enumerate}
\item $R \to S$ 是有限型的，
\item $R \to S'$ 是有限表示的，
\item $M$ 是有限表示的 $S'$-模，
\item $M/IM$ 作为 $S/IS$-模是平坦的，以及
\item $M$ 作为 $R$-模是平坦的。
\end{enumerate}
则 $M$ 是平坦的 $S$-模，并且对每个满足
$M \otimes_S \kappa(\mathfrak q)$ 非零的 $\mathfrak q \subset S$，
$S_\mathfrak q$ 在 $R$ 上平坦且本质有限表示。
\end{lemma}

\begin{proof}
如果 $M \otimes_S \kappa(\mathfrak q)$ 非零，那么
$S' \otimes_S \kappa(\mathfrak q)$ 非零，因而存在位于
$\mathfrak q$ 上方的素理想 $\mathfrak q' \subset S'$（引理
\ref{algebra-lemma-in-image}）。设 $\mathfrak p \subset R$ 是
$\mathfrak q$ 在 $\Spec(R)$ 中的像。由于 $I$ 局部幂零，
$I \subset \mathfrak p$，故 $M/\mathfrak p M$ 在
$S/\mathfrak pS$ 上平坦。因此，可以把引理
\ref{algebra-lemma-criterion-flatness-fibre-fp-over-ft} 用于
$R_\mathfrak p \to S_\mathfrak q \to S'_{\mathfrak q'}$
和 $M_{\mathfrak q'}$。得出 $M_{\mathfrak q'}$ 在 $S$ 上平坦，且
$S_\mathfrak q$ 在 $R$ 上平坦并本质有限表示。
由于 $\mathfrak q'$ 是 $S'$ 的任意素理想，还可知 $M$ 在 $S$ 上平坦
（引理 \ref{algebra-lemma-flat-localization}）。
\end{proof}

\section{平坦轨迹的开性}
\label{algebra-section-open-flat}

\noindent
我们使用引理 \ref{algebra-lemma-colimit-eventually-flat} 化约到 Noether 情形。
Noether 情形则用第 \ref{algebra-section-complex-exact} 节给出的正合复形刻画来处理。

\begin{lemma}
\label{algebra-lemma-CM-dim-finite-type}
设 $k$ 是域，$S$ 是有限型 $k$-代数。设 $f_1, \ldots, f_i$ 是
$S$ 中的元。假设 $S$ 是 Cohen--Macaulay 的、等维且维数为 $d$，并且
$\dim V(f_1, \ldots, f_i) \leq d - i$。则等号成立，而且
% P01-E0660: the frozen plural subject “f_1, ..., f_i” takes singular “forms”.
$f_1, \ldots, f_i$ 对 $V(f_1, \ldots, f_i)$ 的每个素理想
$\mathfrak q$，都在 $S_{\mathfrak q}$ 中构成正则列。
% P01-E0661: the frozen phrase “prime of V(...)” uses the wrong point/set preposition.
\end{lemma}

\begin{proof}
如果 $S$ 是维数为 $d$ 的 Cohen--Macaulay 等维环，则对 $S$ 的所有极大理想
$\mathfrak m$ 都有 $\dim(S_{\mathfrak m}) = d$；见引理
\ref{algebra-lemma-disjoint-decomposition-CM-algebra}。
由命题 \ref{algebra-proposition-CM-module} 可知，对所有极大理想
$\mathfrak m \in V(f_1, \ldots, f_i)$，该序列在 $S_{\mathfrak m}$ 中是正则列，
而局部环 $S_{\mathfrak m}/(f_1, \ldots, f_i)$ 是维数为 $d - i$ 的
Cohen--Macaulay 环。这实际上意味着 $S/(f_1, \ldots, f_i)$ 是维数为
$d - i$ 的 Cohen--Macaulay 等维环。
\end{proof}

\begin{lemma}
\label{algebra-lemma-open-regular-sequence}
设 $R \to S$ 是有限型环映射。设 $d$ 是整数，使得所有纤维
$S \otimes_R \kappa(\mathfrak p)$ 都是维数为 $d$ 的 Cohen--Macaulay 等维环。
设 $f_1, \ldots, f_i$ 是 $S$ 中的元。集合
$$
\{ \mathfrak q \in V(f_1, \ldots, f_i)
\mid f_1, \ldots, f_i
\text{ are a regular sequence in }
S_{\mathfrak q}/\mathfrak p S_{\mathfrak q}
\text{ 其中 }\mathfrak p = R \cap \mathfrak q
\}
$$
在 $V(f_1, \ldots, f_i)$ 中是开的。
\end{lemma}

\begin{proof}
写 $\overline{S} = S/(f_1, \ldots, f_i)$。
假设 $\mathfrak q$ 是引理中所定义集合的元，$\mathfrak p$ 是 $R$ 中相应的素理想。
我们将使用定义 \ref{algebra-definition-relative-dimension} 中定义的相对维数。
首先，由引理 \ref{algebra-lemma-dimension-at-a-point-finite-type-field}，注意到
$d = \dim_{\mathfrak q}(S/R) =
\dim(S_{\mathfrak q}/\mathfrak pS_{\mathfrak q}) +
\text{trdeg}_{\kappa(\mathfrak p)}\ \kappa(\mathfrak q)$。
由于 $f_1, \ldots, f_i$ 在 Noether 局部环
$S_{\mathfrak q}/\mathfrak pS_{\mathfrak q}$ 中构成正则列，
引理 \ref{algebra-lemma-one-equation} 告诉我们
$\dim(\overline{S}_{\mathfrak q}/\mathfrak p\overline{S}_{\mathfrak q})
= \dim(S_{\mathfrak q}/\mathfrak pS_{\mathfrak q}) - i$。
由引理 \ref{algebra-lemma-dimension-at-a-point-finite-type-field} 得出
$\dim_{\mathfrak q}(\overline{S}/R)
= \dim(\overline{S}_{\mathfrak q}/\mathfrak p\overline{S}_{\mathfrak q}) +
\text{trdeg}_{\kappa(\mathfrak p)}\ \kappa(\mathfrak q)
= d - i$。由引理 \ref{algebra-lemma-dimension-fibres-bounded-open-upstairs}，
在 $\mathfrak q$ 的一个邻域中，对所有
$\mathfrak q' \in V(f_1, \ldots, f_i) = \Spec(\overline{S})$，都有
$\dim_{\mathfrak q'}(\overline{S}/R) \leq d - i$。
因此，用某个 $g \in S$、$g \not \in \mathfrak q$ 的 $S_g$ 代替 $S$ 后，
可以假设该不等式对所有 $\mathfrak q'$ 都成立。
结论由引理 \ref{algebra-lemma-CM-dim-finite-type} 得出。
\end{proof}

\begin{lemma}
\label{algebra-lemma-exact-on-fibres-open}
设 $R \to S$ 是环映射。考虑由有限自由 $S$-模组成的有限同调复形：
$$
F_{\bullet} :
0
\to
S^{n_e}
\xrightarrow{\varphi_e}
S^{n_{e-1}}
\xrightarrow{\varphi_{e-1}}
\ldots
\xrightarrow{\varphi_{i + 1}}
S^{n_i}
\xrightarrow{\varphi_i}
S^{n_{i-1}}
\xrightarrow{\varphi_{i-1}}
\ldots
\xrightarrow{\varphi_1}
S^{n_0}
$$
对 $S$ 的每个素理想 $\mathfrak q$，考虑复形
$\overline{F}_{\bullet, \mathfrak q} =
F_{\bullet, \mathfrak q} \otimes_R \kappa(\mathfrak p)$，其中
$\mathfrak p$ 是 $\mathfrak q$ 在 $R$ 中的逆像。
% P01-E0662: the frozen hypothesis omits coordinated finite verbs around “finite type, flat with fibres ... Cohen-Macaulay”.
假设 $R$ 是 Noether 的，并且存在整数 $d$，使得 $R \to S$ 是有限型、平坦的，
其纤维 $S \otimes_R \kappa(\mathfrak p)$ 是维数为 $d$ 的 Cohen--Macaulay 环。
集合
$$
\{\mathfrak q \in \Spec(S) \mid
\overline{F}_{\bullet, \mathfrak q}\text{ is exact}\}
$$
在 $\Spec(S)$ 中是开的。
\end{lemma}

\begin{proof}
设 $\mathfrak q$ 是引理所定义集合的元。
我们将使用命题 \ref{algebra-proposition-what-exact} 来证明：存在
$g \in S$、$g \not \in \mathfrak q$，使得 $D(g)$ 包含于引理所定义的集合中。
换言之，我们将证明：用 $S_g$ 代替 $S$ 后，引理中的集合是整个
$\Spec(S)$。因此，在证明过程中，我们会有限多次地用这样的局部化代替 $S$。
回顾命题 \ref{algebra-proposition-what-exact}：在环为局部环时，它用映射
$\varphi_i$ 的秩和理想 $I(\varphi_i)$ 刻画复形的正合性。
我们先处理秩条件。令
$r_i = n_i - n_{i + 1} + \ldots + (-1)^{e - i} n_e$。
注意 $r_i + r_{i + 1} = n_i$，并且在正合情形下，$r_i$ 是
$\varphi_i$ 的预期秩。

\medskip\noindent
由引理 \ref{algebra-lemma-complex-exact-mod} 可知，如果
$\overline{F}_{\bullet, \mathfrak q}$ 正合，则局部化
$F_{\bullet, \mathfrak q}$ 正合。特别地，由某个
$g \in S$、$g \not \in \mathfrak q$ 局部化后，复形 $F_\bullet$ 变为正合。
在这种情况下，把命题 \ref{algebra-proposition-what-exact} 用于 $S$ 在所有素理想处的局部化，
可知 $\varphi_i$ 的所有 $(r_i + 1) \times (r_i + 1)$-子式均为零。
因此 $\varphi_i$ 的秩至多为 $r_i$。

\medskip\noindent
设 $I_i \subset S$ 表示由 $\varphi_i$ 的矩阵的
$r_i \times r_i$-子式生成的理想。由命题
\ref{algebra-proposition-what-exact}，复形
$\overline{F}_{\bullet, \mathfrak q}$ 正合，当且仅当对每个
$1 \leq i \leq e$，或者 $(I_i)_{\mathfrak q} = S_{\mathfrak q}$，
或者 $(I_i)_{\mathfrak q}$ 包含一个长度为 $i$ 的
$S_{\mathfrak q}/\mathfrak p S_{\mathfrak q}$-正则列。
事实上，由上面对 $r_i$ 的选取以及对 $\varphi_i$ 秩的界，
这是命题 \ref{algebra-proposition-what-exact} 的条件得到满足的唯一方式。

\medskip\noindent
如果 $(I_i)_{\mathfrak q} = S_{\mathfrak q}$，则在某个
$g \not\in \mathfrak q$ 处局部化 $S$ 后，可以假设 $I_i = S$。
显然，这是一个开条件。

\medskip\noindent
如果 $(I_i)_{\mathfrak q} \not = S_{\mathfrak q}$，则有序列
% P01-E0663: the frozen proof chooses f_j in the localized ideal but then treats them as global elements at every prime without explicitly clearing denominators.
$f_1, \ldots, f_i \in (I_i)_{\mathfrak q}$，它们在
$S_{\mathfrak q}/\mathfrak pS_{\mathfrak q}$ 中构成正则列。
注意，对任意满足 $(f_1, \ldots, f_i) \not \subset \mathfrak q'$ 的素理想
$\mathfrak q' \subset S$，都有 $(I_i)_{\mathfrak q'} = S_{\mathfrak q'}$。
因此结论由引理 \ref{algebra-lemma-open-regular-sequence} 得出。
\end{proof}


\begin{theorem}
\label{algebra-theorem-openness-flatness}
设 $R$ 是环，$R \to S$ 是有限表示的环映射，$M$ 是有限表示的 $S$-模。
集合
$$
\{ \mathfrak q \in \Spec(S) \mid
M_{\mathfrak q}\text{ is flat over }R\}
$$
在 $\Spec(S)$ 中是开的。
\end{theorem}

\begin{proof}
设 $\mathfrak q \in \Spec(S)$ 是素理想，$\mathfrak p \subset R$ 是
$\mathfrak q$ 在 $R$ 中的逆像。注意，$M_{\mathfrak q}$ 在 $R$ 上平坦，
当且仅当它在 $R_{\mathfrak p}$ 上平坦。
假设 $M_{\mathfrak q}$ 在 $R$ 上平坦。我们断言：存在
$g \in S$、$g \not \in \mathfrak q$，使得 $M_g$ 在 $R$ 上平坦。

\medskip\noindent
首先化约到 $R$ 和 $S$ 都是在 $\mathbf{Z}$ 上有限型的情形。
选取有向集 $\Lambda$ 以及引理
\ref{algebra-lemma-limit-module-finite-presentation} 中的系统
$(R_\lambda \to S_\lambda, M_\lambda)$。
令 $\mathfrak p_\lambda$ 等于 $\mathfrak p$ 在 $R_\lambda$ 中的逆像，
并令 $\mathfrak q_\lambda$ 等于 $\mathfrak q$ 在 $S_\lambda$ 中的逆像。
则系统
$$
((R_\lambda)_{\mathfrak p_\lambda},
(S_\lambda)_{\mathfrak q_\lambda},
(M_\lambda)_{\mathfrak q_{\lambda}})
$$
如引理 \ref{algebra-lemma-limit-module-essentially-finite-presentation} 中所述。
因此由引理 \ref{algebra-lemma-colimit-eventually-flat} 可知，对某个 $\lambda$，
模 $M_\lambda$ 在素理想 $\mathfrak q_{\lambda}$ 处在 $R_\lambda$ 上平坦。
假设可以对系统
$(R_\lambda \to S_\lambda, M_\lambda, \mathfrak q_{\lambda})$
证明我们的断言。换言之，假设可以找到
$g \in S_\lambda$、$g \not\in \mathfrak q_\lambda$，使得
$(M_\lambda)_g$ 在 $R_\lambda$ 上平坦。
由引理 \ref{algebra-lemma-limit-module-finite-presentation}，有
$M = M_\lambda \otimes_{R_\lambda} R$，因而还有
$M_g = (M_\lambda)_g \otimes_{R_\lambda} R$。
于是由引理 \ref{algebra-lemma-flat-base-change}，得出系统
$(R \to S, M, \mathfrak q)$ 的断言。

\medskip\noindent
至此，可以假设 $R$ 和 $S$ 都在 $\mathbf{Z}$ 上有限型。
可以把 $S$ 写成多项式环 $R[x_1, \ldots, x_n]$ 的商。
当然，可以用 $R[x_1, \ldots, x_n]$ 代替 $S$，并假设 $S$ 是
$R$ 上的多项式环。特别地，$R \to S$ 是平坦的，而且所有纤维环
% P01-E0664: the frozen phrase “all fibres rings” has the wrong plural modifier.
$S \otimes_R \kappa(\mathfrak p)$ 的整体维数都是 $n$。

\medskip\noindent
选取 $M$ 在 $S$ 上的分解 $F_\bullet$，其中每个 $F_i$ 都有限自由；
见引理 \ref{algebra-lemma-resolution-by-finite-free}。
令 $K_n = \Ker(F_{n-1} \to F_{n-2})$。
注意，$(K_n)_{\mathfrak q}$ 在 $R$ 上平坦，因为每个 $F_i$ 在 $R$ 上平坦，
并且对 $M$ 的假设成立；见引理 \ref{algebra-lemma-flat-ses}。
此外，序列
$$
0 \to
K_n/\mathfrak p K_n \to
F_{n-1}/ \mathfrak p F_{n-1} \to
\ldots \to
F_0 / \mathfrak p F_0 \to
M/\mathfrak p M \to
0
$$
在 $\mathfrak q$ 处局部化后是正合的，这是因为
$\text{Tor}_i^{R_\mathfrak p}(\kappa(\mathfrak p), M_{\mathfrak q})$ 消失。
% P01-E0665: the frozen proof says the localized polynomial fibre has global dimension exactly n; it is only bounded above by n in general.
由于 $S_\mathfrak q/\mathfrak p S_{\mathfrak q}$ 的整体维数是 $n$，
可知 $K_n / \mathfrak p K_n$ 在 $\mathfrak q$ 处的局部化是
$S_\mathfrak q/\mathfrak p S_{\mathfrak q}$ 上的有限自由模。
由引理 \ref{algebra-lemma-free-fibre-flat-free}，$(K_n)_{\mathfrak q}$ 在
$S_{\mathfrak q}$ 上自由。特别地，存在
$g \in S$、$g \not \in \mathfrak q$，使得 $(K_n)_g$ 在 $S_g$ 上有限自由。

\medskip\noindent
由引理 \ref{algebra-lemma-exact-on-fibres-open}，存在进一步的局部化 $S_g$，
使得复形
$$
0 \to K_n \to F_{n-1} \to \ldots \to F_0
$$
在 $R \to S$ 的{\it 所有纤维}上正合。
由引理 \ref{algebra-lemma-complex-exact-mod}，这蕴含 $F_1 \to F_0$ 的余核平坦。
这就在 Noether 情形中证明了该定理。
\end{proof}

\section{Cohen--Macaulay 轨迹的开性}
\label{algebra-section-CM-open}

\noindent
本节用平坦性刻画有限型代数的 Cohen--Macaulay 性质。
% P01-E0666: the frozen introduction omits “that” before the object clause “the set ... is open”.
然后利用这一刻画证明：这样的代数为 Cohen--Macaulay 的点集是开的。

\begin{lemma}
\label{algebra-lemma-where-CM}
设 $S$ 是域 $k$ 上的有限型代数。设
$\varphi : k[y_1, \ldots, y_d] \to S$ 是拟有限环映射。
作为 $\Spec(S)$ 的子集，有
% P01-E0667: the frozen left-hand predicate omits “is” before “flat”.
% P01-E0668: the frozen right-hand predicate omits “is” before “CM”.
$$
\{ \mathfrak q \mid
S_{\mathfrak q} \text{ flat over }k[y_1, \ldots, y_d]\}
=
\{ \mathfrak q \mid
S_{\mathfrak q} \text{ CM and }\dim_{\mathfrak q}(S/k) = d\}
$$
记号见定义 \ref{algebra-definition-relative-dimension}。
\end{lemma}

\begin{proof}
设 $\mathfrak q \subset S$ 是素理想。
% P01-E0669: the frozen “Denote p = ...” construction lacks a grammatical complement.
记 $\mathfrak p = k[y_1, \ldots, y_d] \cap \mathfrak q$。
注意，例如由引理 \ref{algebra-lemma-dimension-inequality-quasi-finite}，总有
$\dim(S_{\mathfrak q}) \leq \dim(k[y_1, \ldots, y_d]_{\mathfrak p})$。
此外，域扩张 $\kappa(\mathfrak q)/\kappa(\mathfrak p)$ 是有限的，因而
$\text{trdeg}_k(\kappa(\mathfrak p)) = \text{trdeg}_k(\kappa(\mathfrak q))$。

\medskip\noindent
设 $\mathfrak q$ 是左端集合的元。此时引理
\ref{algebra-lemma-finite-flat-over-regular-CM} 适用，得到
$S_{\mathfrak q}$ 是 Cohen--Macaulay 的，并且
$\dim(S_{\mathfrak q}) = \dim(k[y_1, \ldots, y_d]_{\mathfrak p})$。
结合上面超越次数的等式及引理
\ref{algebra-lemma-dimension-at-a-point-finite-type-field}，这蕴含
$\dim_{\mathfrak q}(S/k) = d$。因此 $\mathfrak q$ 是右端集合的元。

\medskip\noindent
设 $\mathfrak q$ 是右端集合的元。由上面超越次数的等式、假设
$\dim_{\mathfrak q}(S/k) = d$ 及引理
\ref{algebra-lemma-dimension-at-a-point-finite-type-field}，得出
$\dim(S_{\mathfrak q}) = \dim(k[y_1, \ldots, y_d]_{\mathfrak p})$。
因此引理 \ref{algebra-lemma-CM-over-regular-flat} 适用，从而
$\mathfrak q$ 是左端集合的元。
\end{proof}

\begin{lemma}
\label{algebra-lemma-finite-type-over-field-CM-open}
设 $S$ 是域 $k$ 上的有限型代数。使得 $S_{\mathfrak q}$ 为
Cohen--Macaulay 的素理想 $\mathfrak q$ 的集合在 $S$ 中是开的。
% P01-E0670: the frozen statement says “open in S” although the locus is a subset of Spec(S).
\end{lemma}

\noindent
该引理是下方引理
\ref{algebra-lemma-finite-presentation-flat-CM-locus-open} 的特殊情形，
所以若愿意，可以直接跳到该引理的证明。

\begin{proof}
设 $\mathfrak q \subset S$ 是使得 $S_{\mathfrak q}$ 为
Cohen--Macaulay 的素理想。必须证明存在
$g \in S$、$g \not \in \mathfrak q$，使得环 $S_g$ 是 Cohen--Macaulay 的。
对任意 $g \in S$、$g \not \in \mathfrak q$，可以用 $S_g$ 代替 $S$，
并用 $\mathfrak qS_g$ 代替 $\mathfrak q$。结合引理
\ref{algebra-lemma-Noether-normalization-at-point} 和
\ref{algebra-lemma-dimension-at-a-point-finite-type-field}，可以假设存在有限单射环映射
$k[y_1, \ldots, y_d] \to S$，其中
$d = \dim(S_{\mathfrak q}) + \text{trdeg}_k(\kappa(\mathfrak q))$。
令 $\mathfrak p = k[y_1, \ldots, y_d] \cap \mathfrak q$。
按构造，$\mathfrak q$ 是引理 \ref{algebra-lemma-where-CM} 所示等式右端集合的元，
因而也是左端集合的元。

\medskip\noindent
由定理 \ref{algebra-theorem-openness-flatness}，可知对某个
$g \in S$、$g \not \in \mathfrak q$，环 $S_g$ 在
$k[y_1, \ldots, y_d]$ 上平坦。因此再次由引理
\ref{algebra-lemma-where-CM} 的等式，得出 $S_g$ 的所有局部环都如所需那样是
Cohen--Macaulay 的。
\end{proof}

\begin{lemma}
\label{algebra-lemma-generic-CM}
设 $k$ 是域，$S$ 是有限型 $k$ 代数。
% P01-E0671: the frozen compound “finite type k algebra” lacks standard hyphenation.
Cohen--Macaulay 素理想的集合构成稠密开集 $U \subset \Spec(S)$。
\end{lemma}

\begin{proof}
由引理 \ref{algebra-lemma-finite-type-over-field-CM-open}，该集合是开的。
它包含所有极小素理想 $\mathfrak q \subset S$，因为极小素理想处的局部环
$S_{\mathfrak q}$ 维数为零，因而是 Cohen--Macaulay 的。
\end{proof}

\begin{lemma}
\label{algebra-lemma-finite-presentation-flat-CM-locus-open}
设 $R$ 是环。设 $R \to S$ 是有限表示且平坦的。对任意 $d \geq 0$，集合
$$
\left\{
\begin{matrix}
\mathfrak q \in \Spec(S)
\text{ such that setting }\mathfrak p = R \cap \mathfrak q
\text{ the fibre ring}\\
S_{\mathfrak q}/\mathfrak pS_{\mathfrak q}
\text{ 为 Cohen--Macaulay}
\text{ 且 } \dim_{\mathfrak q}(S/R) = d
\end{matrix}
\right\}
$$
在 $\Spec(S)$ 中是开的。
\end{lemma}

\begin{proof}
设 $\mathfrak q$ 是所示集合的元，$\mathfrak p$ 是 $R$ 中相应的素理想。
% P01-E0672: the frozen proof goal mentions Cohen-Macaulay fibres but omits the fixed relative-dimension condition defining the displayed locus.
必须找到 $g \in S$、$g \not \in \mathfrak q$，使得
$R \to S_g$ 的所有纤维环都是 Cohen--Macaulay 的。
在证明过程中，可以（有限多次地）用某个
$g \in S$、$g \not \in \mathfrak q$ 的 $S_g$ 代替 $S$。
因此，由引理 \ref{algebra-lemma-quasi-finite-over-polynomial-algebra}，
可以假设存在拟有限环映射 $R[t_1, \ldots, t_d] \to S$，其中
$d = \dim_{\mathfrak q}(S/R)$。令
$\mathfrak q' = R[t_1, \ldots, t_d] \cap \mathfrak q$。
由引理 \ref{algebra-lemma-where-CM}，可知环映射
$$
R[t_1, \ldots, t_d]_{\mathfrak q'} /
\mathfrak p R[t_1, \ldots, t_d]_{\mathfrak q'}
\longrightarrow
S_{\mathfrak q}/\mathfrak p S_{\mathfrak q}
$$
是平坦的。因此，由纤维平坦性准则引理
\ref{algebra-lemma-criterion-flatness-fibre}，可知
$R[t_1, \ldots, t_d]_{\mathfrak q'} \to S_{\mathfrak q}$ 是平坦的。
由定理 \ref{algebra-theorem-openness-flatness}，可知对某个
% P01-E0673: the frozen sentence lacks a comma after the localized-point condition before “the ring map”.
$g \in S$、$g \not \in \mathfrak q$，环映射
$R[t_1, \ldots, t_d] \to S_g$ 是平坦的。用 $S_g$ 代替 $S$ 后，
可知对每个素理想 $\mathfrak r \subset S$，令
$\mathfrak r' = R[t_1, \ldots, t_d] \cap \mathfrak r$ 及
$\mathfrak p' = R \cap \mathfrak r$，局部环映射
$R[t_1, \ldots, t_d]_{\mathfrak r'} \to S_{\mathfrak r}$ 是平坦的。
因此其基变换
$$
R[t_1, \ldots, t_d]_{\mathfrak r'} /
\mathfrak p' R[t_1, \ldots, t_d]_{\mathfrak r'}
\longrightarrow
S_{\mathfrak r}/\mathfrak p' S_{\mathfrak r}
$$
也是平坦的。因此，把引理 \ref{algebra-lemma-where-CM} 用于
$k = \kappa(\mathfrak p')$，可知 $\mathfrak r$ 如所需那样属于引理中的集合。
\end{proof}

\begin{lemma}
\label{algebra-lemma-generic-CM-flat-finite-presentation}
设 $R$ 是环，$R \to S$ 是有限表示的平坦映射。
使得纤维环 $S_{\mathfrak q} \otimes_R \kappa(\mathfrak p)$
为 Cohen--Macaulay 的素理想 $\mathfrak q$ 的集合，其中
$\mathfrak p = R \cap \mathfrak q$，在 $\Spec(S) \to \Spec(R)$ 的每条纤维中
都是开且稠密的。
\end{lemma}

\begin{proof}
由引理 \ref{algebra-lemma-finite-presentation-flat-CM-locus-open}，该集合（记作 $W$）是开的。
它在各纤维中稠密，因为 $W$ 与一条纤维的交正是该纤维中可应用引理
\ref{algebra-lemma-generic-CM} 的相应集合。
\end{proof}

\begin{lemma}
\label{algebra-lemma-extend-field-CM-locus}
设 $k$ 是域，$S$ 是有限型 $k$-代数。设 $K/k$ 是域扩张，并令
$S_K = K \otimes_k S$。设 $\mathfrak q \subset S$ 是 $S$ 的素理想，
并设 $\mathfrak q_K \subset S_K$ 是位于 $\mathfrak q$ 上方的 $S_K$ 的素理想。
则 $S_{\mathfrak q}$ 是 Cohen--Macaulay 的，当且仅当
$(S_K)_{\mathfrak q_K}$ 是 Cohen--Macaulay 的。
\end{lemma}

\begin{proof}
在证明过程中，可以（有限多次地）对任意
$g \in S$、$g \not \in \mathfrak q$ 用 $S_g$ 代替 $S$。
因此，利用引理 \ref{algebra-lemma-Noether-normalization-at-point}，
可以假设 $\dim(S) = \dim_{\mathfrak q}(S/k) =: d$，并找到有限单射
$k[x_1, \ldots, x_d] \to S$。注意，基变换还诱导有限单射
$K[x_1, \ldots, x_d] \to S_K$。
由引理 \ref{algebra-lemma-dimension-at-a-point-preserved-field-extension}，有
$\dim_{\mathfrak q_K}(S_K/K) = d$。令
$\mathfrak p = k[x_1, \ldots, x_d] \cap \mathfrak q$ 且
$\mathfrak p_K = K[x_1, \ldots, x_d] \cap \mathfrak q_K$。
考虑下列 Noether 局部环交换图：
$$
\xymatrix{
S_{\mathfrak q} \ar[r] &
(S_K)_{\mathfrak q_K} \\
k[x_1, \ldots, x_d]_{\mathfrak p} \ar[r] \ar[u] &
K[x_1, \ldots, x_d]_{\mathfrak p_K} \ar[u]
}
$$
由引理 \ref{algebra-lemma-where-CM}，必须证明左侧竖直箭头平坦，当且仅当右侧竖直箭头平坦。
由于底部箭头平坦，该等价性由引理 \ref{algebra-lemma-base-change-flat-up-down} 成立。
\end{proof}

\begin{lemma}
\label{algebra-lemma-CM-locus-commutes-base-change}
设 $R$ 是环，$R \to S$ 是有限型的。设 $R \to R'$ 是任意环映射，令
$S' = R' \otimes_R S$。
% P01-E0674: the frozen declaration “Denote f ... the map” omits “by/as”.
以 $f : \Spec(S') \to \Spec(S)$ 表示与环映射 $S \to S'$ 相关联的映射。
令 $W$ 等于满足如下条件的素理想 $\mathfrak q$ 的集合：纤维环
$S_{\mathfrak q} \otimes_R \kappa(\mathfrak p)$ 是 Cohen--Macaulay 的，
其中 $\mathfrak p = R \cap \mathfrak q$；并令 $W'$ 表示 $S'/R'$ 的类似集合。
则 $W' = f^{-1}(W)$。
\end{lemma}

\begin{proof}
% P01-E0675: the frozen one-sentence proof omits its subject.
这直接由引理 \ref{algebra-lemma-extend-field-CM-locus} 和定义得出。
\end{proof}

\begin{lemma}
\label{algebra-lemma-relative-dimension-CM}
设 $R$ 是环。设 $R \to S$ 是满足如下条件的环映射：(a) 平坦；
(b) 有限表示；(c) 具有 Cohen--Macaulay 纤维。
% P01-E0676: the frozen list is grammatically nonparallel: “is ... (c) has ...”.
则可以把 $S = S_0 \times \ldots \times S_n$ 写成 $R$-代数 $S_d$ 的乘积，
使得每个 $S_d$ 都满足 (a)、(b)、(c)，并且所有纤维都是维数为 $d$ 的等维环。
\end{lemma}

\begin{proof}
% P01-E0677: the frozen declaration “denote W_d ... the set” omits “by/as”.
对每个整数 $d$，以 $W_d \subset \Spec(S)$ 表示引理
\ref{algebra-lemma-finite-presentation-flat-CM-locus-open} 中定义的集合。
显然有 $\Spec(S) = \coprod W_d$，且由刚引用的引理，每个 $W_d$ 都是开的。
因此结论由引理 \ref{algebra-lemma-disjoint-implies-product} 得出。
\end{proof}

\section{微分}
\label{algebra-section-differentials}

\noindent
本节定义环映射的微分模。

\begin{definition}
\label{algebra-definition-derivation}
设 $\varphi : R \to S$ 是环映射，$M$ 是 $S$-模。
一个{\it 导子}，更准确地说，一个取值于 $M$ 的
{\it $R$-导子}，是映射 $D : S \to M$，它具有可加性，
在 $\varphi(R)$ 的元上为零，并且满足{\it Leibniz 法则}：
$D(ab) = aD(b) + bD(a)$。
\end{definition}

\noindent
注意，如果 $r \in R$ 且 $a \in S$，则 $D(ra) = rD(a)$。
等价地，$R$-导子是满足 Leibniz 法则的 $R$-线性映射 $D : S \to M$。
所有 $R$-导子的集合构成一个 $S$-模：给定两个 $R$-导子 $D, D'$，
其和 $D + D' : S \to M$，$a \mapsto D(a)+D'(a)$ 是 $R$-导子；
给定 $R$-导子 $D$ 和元 $c\in S$，其标量倍
$cD : S \to M$，$a \mapsto cD(a)$ 也是 $R$-导子。这个 $S$-模记作
$$
\text{Der}_R(S, M).
$$
另外，如果 $\alpha : M \to N$ 是 $S$-模映射，则复合
$\alpha \circ D$ 是到 $N$ 的 $R$-导子。这样，赋值
$M \mapsto \text{Der}_R(S, M)$ 是协变函子。

\medskip\noindent
考虑下列自由 $S$-模映射
$$
\bigoplus\nolimits_{(a, b)\in S^2} S[(a, b)] \oplus
\bigoplus\nolimits_{(f, g)\in S^2} S[(f, g)] \oplus
\bigoplus\nolimits_{r\in R} S[r]
\longrightarrow
\bigoplus\nolimits_{a\in S} S[a]
$$
它由如下规则定义：
$$
[(a, b)] \longmapsto [a + b] - [a] - [b],\quad
[(f, g)] \longmapsto [fg] -f[g] - g[f],\quad
[r]      \longmapsto [\varphi(r)]
$$
记号的含义是显然的。令 $\Omega_{S/R}$ 为该映射的余核。
有映射 $\text{d} : S \to \Omega_{S/R}$，它把 $a$ 映到
$[a]$ 在余核中的类 $\text{d}a$。由 $\Omega_{S/R}$ 上所施加的关系，
这是一个 $R$-导子；换言之，
$$
\text{d}(a + b) = \text{d}a + \text{d}b, \quad
\text{d}(fg) = f\text{d}g + g\text{d}f, \quad
\text{d}\varphi(r) = 0
$$
其中 $a,b,f,g \in S$ 且 $r \in R$。

\begin{definition}
\label{algebra-definition-differentials}
二元组 $(\Omega_{S/R}, \text{d})$ 称为 $S$ 在 $R$ 上的
{\it K"ahler 微分模}，或简称{\it 微分模}。
\end{definition}

\begin{lemma}
\label{algebra-lemma-universal-omega}
\begin{slogan}
从微分模出发的映射与导子相同。
\end{slogan}
$S$ 在 $R$ 上的微分模具有如下泛性质。映射
$$
\Hom_S(\Omega_{S/R}, M)
\longrightarrow
\text{Der}_R(S, M), \quad
\alpha
\longmapsto
\alpha \circ \text{d}
$$
是函子的同构。
\end{lemma}

\begin{proof}
按定义，$R$-导子是一条规则，它为每个 $a \in S$ 配置一个元
$D(a) \in M$。因此 $D$ 给出映射 $[D] : \bigoplus S[a] \to M$。
而成为 $R$-导子的各项条件恰好意味着，$[D]$ 在上面所示
$\Omega_{S/R}$ 的表示映射的像上为零。
\end{proof}

\begin{lemma}
\label{algebra-lemma-trivial-differential-surjective}
假设 $R \to S$ 是满射。则 $\Omega_{S/R} = 0$。
\end{lemma}

\begin{proof}
这或者可由所有 $R$-导子显然必须为零看出，或者可由
$\Omega_{S/R}$ 的表示中的映射是满射看出。
\end{proof}

\noindent
假设
\begin{equation}
\label{algebra-equation-functorial-omega}
\vcenter{
\xymatrix{
S \ar[r]_\varphi
&
S'
\\
R \ar[r]^\psi \ar[u]^\alpha
&
R' \ar[u]_\beta
}
}
\end{equation}
是环的交换图。此时有微分模之间的自然映射，它嵌入交换图
$$
\xymatrix{
\Omega_{S/R} \ar[r] &
\Omega_{S'/R'}
\\
S \ar[u]^{\text{d}} \ar[r]^{\varphi}
&
S' \ar[u]_{\text{d}}
}
$$
要构造该映射，只须使用 $\Omega_{S/R}$ 与 $\Omega_{S'/R'}$
的表示之间显然的映射。具体地说，
\begin{equation}
\label{algebra-equation-map-presentations}
\vcenter{
\xymatrix{
\bigoplus S'[(a', b')]
\oplus
\bigoplus S'[(f', g')]
\oplus
\bigoplus S'[r'] \ar[r]
&
\bigoplus S' [a'] \\
 \\
\bigoplus S[(a, b)]
\oplus
\bigoplus S[(f, g)]
\oplus
\bigoplus S[r] \ar[r]
\ar[uu]^{
\begin{matrix}
[(a, b)] \mapsto [(\varphi(a), \varphi(b))] \\
[(f, g)] \mapsto [(\varphi(f), \varphi(g))] \\
[r]\mapsto [\psi(r)]
\end{matrix}
} &
\bigoplus S[a] \ar[uu]_{[a] \mapsto [\varphi(a)]}
}
}
\end{equation}
所得结果就是：$f\text{d}g \in \Omega_{S/R}$ 映到
$\varphi(f)\text{d}\varphi(g)$。


\begin{lemma}
\label{algebra-lemma-colimit-differentials}
设 $I$ 是有向集。设 $(R_i \to S_i, \varphi_{ii'})$ 是 $I$ 上的环映射系统；
见“范畴”第 \ref{categories-section-posets-limits} 节。则
% P01-E0678: the frozen display ends with a period before the following “where” clause.
$$
\Omega_{S/R} =
\colim_i \Omega_{S_i/R_i}.
$$
其中 $R \to S = \colim (R_i \to S_i)$。
\end{lemma}

\begin{proof}
这由 $\Omega_{S/R}$ 的定义表示以及上面所述的这一构造的函子性立即得出。
% P01-E0679: the frozen phrase “the functoriality of this described above” lacks a noun modified by “of this”.
\end{proof}

\begin{lemma}
\label{algebra-lemma-differential-surjective}
在图 (\ref{algebra-equation-functorial-omega}) 中，假设 $S \to S'$ 是满射，
其核为 $I \subset S$。则 $\Omega_{S/R} \to \Omega_{S'/R'}$ 是满射，
其核作为 $S$-模由如下元生成：$\text{d}a$，其中 $a \in S$ 满足
$\varphi(a) \in \beta(R')$。（特别地，这包括元 $\text{d}(i)$，$i \in I$。）
\end{lemma}

\begin{proof}[第一种证明]
考虑表示映射 (\ref{algebra-equation-map-presentations})。
显然，自由模之间右侧的竖直映射是满射。因此该映射是满射。
假设 $\Omega_{S/R}$ 的某个元 $\eta$ 在 $\Omega_{S'/R'}$ 中映到零。
把 $\eta$ 写成某个 $s_i, a_i \in S$ 的 $\sum s_i[a_i]$ 的像。
那么 $\sum \varphi(s_i)[\varphi(a_i)]$ 是上方左角中元
$$
\theta = \sum s_j'[a_j', b_j'] + \sum s_k'[f_k', g_k'] + \sum s_l'[r_l']
$$
的像。由于 $\varphi$ 是满射，各项 $s_j'[a_j', b_j']$ 和
$s_k'[f_k', g_k']$ 都是下方右角中各元的像。
% P01-E0680: the frozen proof names the lower right corner, but relation generators lift from the lower left corner.
因此，从 $\eta$ 和 $\theta$ 中减去这些元的像来修改它们后，可以假设
$\theta = \sum s_l'[r_l']$。换言之，
$\sum \varphi(s_i)[\varphi(a_i)]$ 具有形式
$\sum s'_l [\beta(r'_l)]$。
接着，可以假设存在某个 $a' \in S'$，使得对所有 $i$ 都有
$a' = \varphi(a_i)$，并且对所有 $l$ 都有 $a' = \beta(r_l')$。
% P01-E0681: the frozen proof suppresses the needed componentwise decomposition before assuming one common basis element a'.
这由图中右上角的直和分解立即得出。选取满足
$\varphi(a) = a'$ 的 $a \in S$。于是可以对某些 $x_i \in I$ 写成
$a_i = a + x_i$。因此利用容许的关系，可以假设所有 $a_i$ 都等于 $a$。
但此时可假设我们的元具有形式 $s[a]$。仍然知道
% P01-E0682: the frozen formula applies phi to coefficients s'_l that already lie in S'.
$\varphi(s)[a'] = \sum \varphi(s_l')[\beta(r_l')]$。
所以，或者 $\varphi(s) = 0$，此时已得结论；或者
$a' = \varphi(a)$ 在 $\beta$ 的像中，此时同样已得结论。
\end{proof}

\begin{proof}[第二种证明]
我们将使用引理 \ref{algebra-lemma-universal-omega} 给出的微分模泛性质，
不再另行说明。

\medskip\noindent
在 (\ref{algebra-equation-functorial-omega}) 中令 $R'' = S \times_{S'} R'$。
则有下图：
% P01-E0683: the frozen phrase “we have following diagram” omits “the”.
$$
\xymatrix{
S \ar[r] & S \ar[r] & S' \\
R \ar[r] \ar[u] & R'' \ar[r] \ar[u] & R' \ar[u]
}
$$
设 $M$ 是 $S$-模。由定义立即可知，$R$-导子 $D : S \to M$ 是
$R''$-导子，当且仅当它在 $R'' \to S$ 的像中的元上为零。
泛性质把这转化为如下断言：自然映射
$\Omega_{S/R} \to \Omega_{S/R''}$ 是满射，其核作为 $S$-模由
$R''$ 的像生成。
% P01-E0684: the frozen kernel description names elements of R'' rather than their differentials in Omega.

\medskip\noindent
由前一段可知，只须证明：当 $S \to S'$ 是满射且
$R = S \times_{S'} R'$ 时，$\Omega_{S/R} \to \Omega_{S'/R'}$ 是同构。
设 $M'$ 是 $S'$-模。注意，任意 $R'$-导子 $D' : S' \to M'$ 与
$S \to S'$ 前复合后，都给出一个 $R$-导子。
反过来，假设 $M$ 是 $S$-模，$D : S \to M$ 是 $R$-导子。
如果 $i \in I$，则存在 $a \in R$ 使得 $\alpha(a) = i$
（因为 $R = S \times_{S'} R'$）。于是 $D(i) = 0$，从而对所有
$s \in S$，有 $0 = D(is) = iD(s)$。因此，$D$ 的像包含于
由 $I$ 零化的子模 $M' \subset M$ 中，而且诱导映射
$S \to M'$ 分解为一个 $R'$-导子 $S' \to M'$。
利用泛性质看出这意味着 $\Omega_{S/R} \to \Omega_{S'/R'}$ 是同构，
这是一个练习；细节从略。
\end{proof}

\begin{lemma}
\label{algebra-lemma-exact-sequence-differentials}
设 $A \to B \to C$ 是环映射。则有典范的 $C$-模正合列
$$
C \otimes_B \Omega_{B/A} \to
\Omega_{C/A} \to
\Omega_{C/B} \to 0
$$
\end{lemma}

\begin{proof}
在图 (\ref{algebra-equation-functorial-omega}) 中置
$R = A$、$S = C$、$R' = B$ 及 $S' = C$，得到一个图。
由引理 \ref{algebra-lemma-differential-surjective}，映射
$\Omega_{C/A} \to \Omega_{C/B}$ 是满射，其核由元
$\text{d}(c)$ 生成，其中 $c \in C$ 属于 $B \to C$ 的像。引理得证。
\end{proof}

\begin{lemma}
\label{algebra-lemma-differentials-localize}
设 $\varphi : A \to B$ 是环映射。
\begin{enumerate}
\item 若 $S \subset A$ 是一个映到 $B$ 的可逆元中的乘法子集，则
$\Omega_{B/A} = \Omega_{B/S^{-1}A}$。
\item 若 $S \subset B$ 是乘法子集，则
$S^{-1}\Omega_{B/A} = \Omega_{S^{-1}B/A}$。
\end{enumerate}
\end{lemma}

\begin{proof}
要证明 (1) 中的等式，只须证明任意 $A$-导子
$D : B \to M$ 都在元 $\varphi(s)^{-1}$ 上为零。
这由把 Leibniz 法则应用于 $1 = \varphi(s) \varphi(s)^{-1}$ 立即可得。
% P01-E0697: the frozen introductory phrase “To show (2) note” lacks a comma after “(2)”.
要证明 (2)，注意有显然的映射
$S^{-1}\Omega_{B/A} \to \Omega_{S^{-1}B/A}$。
要证明它是同构，只须证明存在一个取值于 $S^{-1}\Omega_{B/A}$ 的
$S^{-1}B$ 的 $A$-导子 $\text{d}'$。
% P01-E0685: the frozen source says “a A-derivation” instead of “an A-derivation”.
为定义它，只须置
$\text{d}'(b/s) = (1/s)\text{d}b - (1/s^2)b\text{d}s$。
细节从略。
\end{proof}

\begin{lemma}
\label{algebra-lemma-differential-seq}
在图 (\ref{algebra-equation-functorial-omega}) 中，假设 $S \to S'$ 是满射，
其核为 $I \subset S$，并且假设 $R' = R$。
则有典范的 $S'$-模正合列
$$
I/I^2
\longrightarrow
\Omega_{S/R} \otimes_S S'
\longrightarrow
\Omega_{S'/R}
\longrightarrow
0
$$
% P01-E0692: the frozen source omits terminal punctuation after this displayed exact sequence.
最左侧映射由如下规则刻画：$f \in I$ 映到
$\text{d}f \otimes 1$。
\end{lemma}

\begin{proof}
中间项是 $\Omega_{S/R} \otimes_S S/I$。
对 $f \in I$，用 $\overline{f}$ 表示 $f$ 在 $I/I^2$ 中的像。
% P01-E0686: the frozen phrase “denote \bar f the image” omits “by/as”.
要证明映射 $\overline{f} \mapsto \text{d}f \otimes 1$ 是良定义的，
只须验证当 $f_1, f_2 \in I$ 时
$\text{d} f_1f_2 \otimes 1 = 0$。
这由 Leibniz 法则立即可得：
$\text{d} f_1f_2 \otimes 1 = (f_1 \text{d}f_2 + f_2 \text{d} f_1 )\otimes 1 =
\text{d}f_2 \otimes f_1 + \text{d}f_1 \otimes f_2 = 0$。
类似的计算表明此映射是 $S' = S/I$-线性的。
% P01-E0687: the frozen source has subject-verb disagreement in “A similar computation show”.

\medskip\noindent
映射 $\Omega_{S/R} \otimes_S S' \to \Omega_{S'/R}$ 是与
$S$-线性映射 $\Omega_{S/R} \to \Omega_{S'/R}$ 相伴的典范
$S'$-线性映射。它是满射，因为由引理
\ref{algebra-lemma-differential-surjective}，映射
$\Omega_{S/R} \to \Omega_{S'/R}$ 是满射。

\medskip\noindent
两个映射的复合为零，因为当 $f \in I$ 时，
$\text{d}f$ 在 $\Omega_{S'/R}$ 中映到零。
注意，正合性恰好是说 $\Omega_{S/R} \to \Omega_{S'/R}$ 的核
作为 $S$-子模由子模 $I\Omega_{S/R}$ 与元 $\text{d}f$（$f \in I$）
共同生成。由引理 \ref{algebra-lemma-differential-surjective}，我们知道
此核由元 $\text{d}(a)$ 生成，其中对某个 $r \in R$ 有
$\varphi(a) = \beta(r)$。但这时
$a = \alpha(r) + a - \alpha(r)$，所以
$\text{d}(a) = \text{d}(a - \alpha(r))$。并且
$a - \alpha(r) \in I$，因为 $\varphi(a - \alpha(r)) =
\varphi(a) - \varphi(\alpha(r)) = \beta(r) - \beta(r) = 0$。
由此可知，元 $\text{d}f$（$f \in I$）已经作为 $S$-模生成此核，
正如所需。
% P01-E0696: the frozen sentence “We conclude the elements...” omits “that”.
\end{proof}

\begin{lemma}
\label{algebra-lemma-differential-seq-split}
在图 (\ref{algebra-equation-functorial-omega}) 中，假设 $S \to S'$ 是满射，
其核为 $I \subset S$，并假设 $R' = R$。进一步，假设存在一个
$R$-代数映射 $S' \to S$，它是 $S \to S'$ 的右逆。
则引理 \ref{algebra-lemma-differential-seq} 的 $S'$-模正合列变为短正合列
$$
0 \longrightarrow
I/I^2
\longrightarrow
\Omega_{S/R} \otimes_S S'
\longrightarrow
\Omega_{S'/R}
\longrightarrow
0
$$
而且它甚至是分裂短正合列。
\end{lemma}

\begin{proof}
令 $\beta : S' \to S$ 是满射 $\alpha : S \to S'$ 的右逆。
考虑映射
$$
D : S \longrightarrow I/I^2,
\quad
x \longmapsto x - \beta(\alpha(x))
$$
% P01-E0693: the frozen source omits terminal punctuation after the displayed map.
容易证明 $D$ 是 $R$-导子（从略）。
此外，若 $x \in I, s \in S$，则 $x D(s) = 0$。
因此由泛性质，$D$ 诱导映射
$\tau : \Omega_{S/R} \otimes_S S' \to I/I^2$。
它是 $\text{d} : I/I^2  \to \Omega_{S/R} \otimes_S S'$ 的左逆；
此验证从略。因此结论成立。
\end{proof}

\begin{lemma}
\label{algebra-lemma-differential-mod-power-ideal}
设 $R \to S$ 是环映射，$I \subset S$ 是理想，$n \geq 1$ 是整数。
置 $S' = S/I^{n + 1}$。映射
$\Omega_{S/R} \to \Omega_{S'/R}$ 诱导同构
$$
\Omega_{S/R} \otimes_S S/I^n
\longrightarrow
\Omega_{S'/R} \otimes_{S'} S/I^n.
$$
\end{lemma}

\begin{proof}
这由引理 \ref{algebra-lemma-differential-seq} 以及
$\text{d}$ 的 Leibniz 法则所给出的
$\text{d}(I^{n + 1}) \subset I^n\Omega_{S/R}$ 得出。
\end{proof}

\begin{lemma}
\label{algebra-lemma-differentials-base-change}
假设有环映射 $R \to R'$ 和 $R \to S$。
置 $S' = S \otimes_R R'$，于是得到图
(\ref{algebra-equation-functorial-omega})。则上面定义的典范映射诱导同构
$\Omega_{S/R} \otimes_R R' = \Omega_{S'/R'}$。
\end{lemma}

\begin{proof}
令 $\text{d}' : S' = S \otimes_R R' \to \Omega_{S/R} \otimes_R R'$
表示映射
$\text{d}'( \sum a_i \otimes x_i ) = \sum \text{d}(a_i) \otimes x_i$。
它存在，因为映射
$S \times R' \to \Omega_{S/R} \otimes_R R'$，
$(a, x)\mapsto \text{d}a \otimes_R x$ 是 $R$-双线性的。
这是一个 $R'$-导子；直接计算即可验证。
我们将证明 $(\Omega_{S/R} \otimes_R R', \text{d}')$ 满足泛性质。
设 $D : S' \to M'$ 是一个取值于某个 $S'$-模的 $R'$-导子。
复合 $S \to S' \to M'$ 是 $R$-导子，所以得到一个
$S$-线性映射 $\varphi_D : \Omega_{S/R} \to M'$。
将它与 $R'$ 作张量积，得到映射
$\varphi'_D : \Omega_{S/R} \otimes_R R' \to M'$，
$\varphi'_D(\eta \otimes x) = x\varphi_D(\eta)$。
显然 $D = \varphi'_D \circ \text{d}'$。
\end{proof}

\noindent
乘法映射 $S \otimes_R S \to S$ 是把 $a \otimes b$ 映到
$S$ 中 $ab$ 的 $R$-代数映射。若通过两个映射
$S \to S \otimes_R S$ 中的任一个，把 $S \otimes_R S$
视为 $S$-代数，则它也是 $S$-代数映射。

\begin{lemma}
\label{algebra-lemma-differentials-diagonal}
设 $R \to S$ 是环映射，令 $J = \Ker(S \otimes_R S \to S)$
为乘法映射的核。存在典范的 $S$-模同构
$\Omega_{S/R} \to J/J^2$，
$a \text{d} b \mapsto a \otimes b - ab \otimes 1$。
\end{lemma}

\begin{proof}[第一种证明]
把引理 \ref{algebra-lemma-differential-seq-split} 应用于交换图
$$
\xymatrix{
S \otimes_R S \ar[r] & S \\
S \ar[r] \ar[u] & S \ar[u]
}
$$
其中左侧竖直箭头是 $a \mapsto a \otimes 1$。得到正合列
$0 \to J/J^2 \to
\Omega_{S \otimes_R S/S} \otimes_{S \otimes_R S} S \to \Omega_{S/S} \to 0$。
由引理 \ref{algebra-lemma-trivial-differential-surjective}，
项 $\Omega_{S/S}$ 是 $0$，因而得到另两项之间的同构。
由引理 \ref{algebra-lemma-differentials-base-change}，
由于 $S \to S \otimes_R S$ 是 $R \to S$ 的基变换，有
$\Omega_{S \otimes_R S/S} = \Omega_{S/R} \otimes_S (S \otimes_R S)$，
所以
$$
\Omega_{S \otimes_R S/S} \otimes_{S \otimes_R S} S =
\Omega_{S/R} \otimes_S (S \otimes_R S) \otimes_{S \otimes_R S} S =
\Omega_{S/R}
$$
% P01-E0694: the frozen source omits terminal punctuation after this displayed chain of equalities.
映射由引理中的规则给出；此验证从略。
\end{proof}

\begin{proof}[第二种证明]
先证明规则 $a \text{d} b \mapsto a \otimes b - ab \otimes 1$
是良定义的。为此必须证明 $\text{d}r$ 和
$a\text{d}b + b \text{d}a - d(ab)$ 映到零。
% P01-E0688: the frozen relation uses d(ab) instead of the established \text{d}(ab) notation.
第一项为零，因为由张量积的定义，
$r \otimes 1 - 1 \otimes r = 0$；第二项为零，因为
% P01-E0689: the frozen “The first because...” clause is a sentence fragment.
$$
(a \otimes b - ab \otimes 1) +
(b \otimes a - ba \otimes 1) -
(1 \otimes ab - ab \otimes 1)
=
(a \otimes 1 - 1\otimes a)(1\otimes b - b \otimes 1)
$$
属于 $J^2$。

\medskip\noindent
我们构造反方向的映射。可以把
$S \to S \otimes_R S$，$a \mapsto a \otimes 1$
视为 $R \to S$ 的基变换。因此由引理
\ref{algebra-lemma-differentials-base-change} 有
$\Omega_{S \otimes_R S/S} = \Omega_{S/R} \otimes_S (S \otimes_R S)$。
此时，引理 \ref{algebra-lemma-differential-seq} 的序列给出映射
$$
J/J^2  \to \Omega_{S \otimes_R S/ S} \otimes_{S \otimes_R S} S
= (\Omega_{S/R} \otimes_S (S \otimes_R S))\otimes_{S \otimes_R S} S
= \Omega_{S/R}.
$$
其为上述映射的逆；留给读者验证。
% P01-E0690: the frozen phrase “see it is the inverse” omits “that”.
\end{proof}

\begin{lemma}
\label{algebra-lemma-differentials-polynomial-ring}
若 $S = R[x_1, \ldots, x_n]$，则 $\Omega_{S/R}$ 是一个有限自由
$S$-模，其基为 $\text{d}x_1, \ldots, \text{d}x_n$。
\end{lemma}

\begin{proof}
先证明 $\text{d}x_1, \ldots, \text{d}x_n$ 作为 $S$-模生成
$\Omega_{S/R}$。为此，对任意 $g \in R[x_1, \ldots, x_n]$，
我们证明 $\text{d}g$ 可表示为和
$\sum g_i \text{d}x_i$。对 $g$ 的（全）次数作归纳。
若 $g$ 的次数为 $0$，则结论显然，因为此时 $\text{d}g = 0$。
若 $g$ 的次数为 $> 0$，则可把 $g$ 写成
$c + \sum g_i x_i$，其中 $c\in R$ 且
$\deg(g_i) < \deg(g)$。由 Leibniz 法则，
$\text{d}g = \sum g_i \text{d} x_i + \sum x_i \text{d}g_i$，
因而由归纳法得到结论。

\medskip\noindent
考虑 $R$-导子
$\partial / \partial x_i :
R[x_1, \ldots, x_n] \to R[x_1, \ldots, x_n]$。
（它的定义留给读者；其刻画性质是
$\partial / \partial x_i (x_j) = \delta_{ij}$。）
由泛性质，它对应于一个 $S$-模映射
$l_i : \Omega_{S/R} \to R[x_1, \ldots, x_n]$，
后者把 $\text{d}x_i$ 映到 $1$，并在 $j \not = i$ 时把
$\text{d}x_j$ 映到 $0$。因此，元
$\text{d}x_1, \ldots, \text{d}x_n$ 之间显然不存在
$S$-线性关系。
\end{proof}

\begin{lemma}
\label{algebra-lemma-differentials-finitely-presented}
假设 $R \to S$ 是有限表现的。则 $\Omega_{S/R}$ 是有限表现
$S$-模。
\end{lemma}

\begin{proof}
写成 $S = R[x_1, \ldots, x_n]/(f_1, \ldots, f_m)$，
并置 $I = (f_1, \ldots, f_m)$。由引理
\ref{algebra-lemma-differential-seq}，有 $S$-模正合列
$$
I/I^2
\to
\Omega_{R[x_1, \ldots, x_n]/R} \otimes_{R[x_1, \ldots, x_n]} S
\to
\Omega_{S/R}
\to
0
$$
% P01-E0695: the frozen source omits terminal punctuation after this displayed exact sequence.
结论由以下事实得出：$I/I^2$ 是有限 $S$-模（由 $f_i$ 的像生成），
而由引理 \ref{algebra-lemma-differentials-polynomial-ring}，中间项是有限自由的。
\end{proof}

\begin{lemma}
\label{algebra-lemma-differentials-finitely-generated}
假设 $R \to S$ 是有限型的。则 $\Omega_{S/R}$ 是有限生成
$S$-模。
% P01-E0691: the frozen source omits the article in “is finitely generated S-module”.
\end{lemma}

\begin{proof}
这与引理 \ref{algebra-lemma-differentials-finitely-presented} 的证明非常相似，
但更为容易。
\end{proof}

\section{de Rham 复形}
\label{algebra-section-de-rham-complex}

\noindent
设 $A \to B$ 是环映射。用 $\text{d} : B \to \Omega_{B/A}$
表示第 \ref{algebra-section-differentials} 节构造的微分模及其泛 $A$-导子。
% P01-E0698: the frozen designation grammatically identifies the map d as “the module of differentials” rather than as its universal derivation.
对 $i \geq 0$，令
$\Omega_{B/A}^i = \wedge^i_B(\Omega_{B/A})$ 为第
\ref{algebra-section-tensor-algebra} 节所述的第 $i$ 个外幂。
{\it $B$ 在 $A$ 上的 de Rham 复形}是复形
$$
\Omega_{B/A}^0 \to \Omega_{B/A}^1 \to \Omega_{B/A}^2 \to \ldots
$$
其 $A$-线性微分将在下面构造并描述。

\medskip\noindent
映射 $\text{d} : \Omega^0_{B/A} \to \Omega^1_{B/A}$
就是泛导子 $\text{d} : B \to \Omega_{B/A}$。
注意它确实是 $A$-线性的。

\medskip\noindent
对 $p \geq 1$，我们断言存在唯一的 $A$-线性映射
$\text{d} : \Omega_{B/A}^p \to \Omega_{B/A}^{p + 1}$，使得
\begin{equation}
\label{algebra-equation-rule}
\text{d}\left(b_0\text{d}b_1 \wedge \ldots \wedge \text{d}b_p\right) =
\text{d}b_0 \wedge \text{d}b_1 \wedge \ldots \wedge \text{d}b_p
\end{equation}
% P01-E0699: the frozen numbered equation lacks terminal punctuation before the next sentence.
回忆 $\Omega_{B/A}$ 作为 $B$-模由元 $\text{d}b$ 生成。
因此，$\Omega^p_{B/A}$ 作为 $A$-模由元
$b_0\text{d}b_1 \wedge \ldots \wedge \text{d}b_p$ 生成，
从而映射 $\text{d} : \Omega^p_{B/A} \to \Omega^{p + 1}_{B/A}$
若存在便是唯一的。

\medskip\noindent
构造 $\text{d} : \Omega_{B/A}^1 \to \Omega_{B/A}^2$。
由定义 \ref{algebra-definition-differentials}，元 $\text{d}b$
作为 $B$-模自由生成 $\Omega_{B/A}$，但须满足关系：
当 $a \in A$ 时 $\text{d}a = 0$，以及当 $b', b'' \in B$ 时
$\text{d}(b' + b'') = \text{d}b' + \text{d}b''$ 和
$\text{d}(b'b'') = b'\text{d}b'' + b''\text{d}b'$。
因此，要证明规则
$$
\sum b'_i \text{d}b_i \longmapsto \sum \text{d}b'_i \wedge \text{d}b_i
$$
是良定义的，必须证明诸元
$$
b\text{d}a,
\quad\text{且}\quad
b\text{d}(b' + b'') - b\text{d}b' - b\text{d}b''
\quad\text{且}\quad
b\text{d}(b'b'') - bb'\text{d}b'' - bb''\text{d}b'
$$
在 $a \in A$ 且 $b, b', b'' \in B$ 时映到零。
使用 $\text{d}$ 的 Leibniz 法则直接计算即可看出这一点。

\medskip\noindent
注意，复合
$\Omega^0_{B/A} \to \Omega^1_{B/A} \to \Omega^2_{B/A}$
为零，因为 $\text{d}(\text{d}(b)) = \text{d}(1 \text{d}b) =
\text{d}(1) \wedge \text{d}(b) = 0 \wedge \text{d}b = 0$。
这里 $\text{d}(1) = 0$，因为 $1 \in B$ 属于 $A \to B$ 的像。
下面将使用这一点。

\medskip\noindent
对 $p \geq 2$ 构造
$\text{d} : \Omega_{B/A}^p \to \Omega_{B/A}^{p + 1}$。
我们将证明由公式
$$
\gamma : \Omega^1_{B/A} \otimes_A \ldots \otimes_A \Omega^1_{B/A}
\longrightarrow
\Omega_{B/A}^{p + 1}
$$
定义的 $A$-线性映射
$$
\omega_1 \otimes \ldots \otimes \omega_p
\longmapsto
\sum (-1)^{i + 1}
\omega_1 \wedge \ldots \wedge \text{d}(\omega_i) \wedge \ldots \wedge \omega_p
$$
% P01-E0700: the frozen alternating-sum formula omits the summation bounds i = 1,...,p.
通过自然满射
$\Omega^1_{B/A} \otimes_A \ldots \otimes_A \Omega^1_{B/A} \to \Omega^p_{B/A}$
分解，从而给出所需映射
$\text{d} : \Omega^p_{B/A} \to \Omega^{p + 1}_{B/A}$。
由引理 \ref{algebra-lemma-present-wedge}，
$\Omega^1_{B/A} \otimes_A \ldots \otimes_A \Omega^1_{B/A} \to \Omega^p_{B/A}$
的核作为 $A$-模由以下元生成：某些 $i \not = j$ 满足
$\omega_i = \omega_j$ 的元
$\omega_1 \otimes \ldots \otimes \omega_p$，以及某个 $f \in B$ 对应的元
$\omega_1 \otimes \ldots \otimes f\omega_i \otimes \ldots \otimes \omega_p -
\omega_1 \otimes \ldots \otimes f\omega_j \otimes \ldots \otimes \omega_p$。
直接计算表明，第一类元被 $\gamma$ 映到 $0$；换言之，
$\gamma$ 是交替的。为完成证明，还须证明
$$
\gamma(
\omega_1 \otimes \ldots \otimes f\omega_i \otimes \ldots \otimes \omega_p)
=
\gamma(
\omega_1 \otimes \ldots \otimes f\omega_j \otimes \ldots \otimes \omega_p)
$$
对 $f \in B$ 成立。由 $A$-线性性和交替性，只须对
$p = 2$、$i = 1$、$j = 2$、
$\omega_1 = b \text{d}b'$ 及 $\omega_2 = c \text{d} c'$
证明，其中 $b, b', c, c' \in B$。因此须证明
\begin{align*}
& \text{d}(fb) \wedge \text{d}b' \wedge c \text{d}c'
- fb \text{d}b' \wedge \text{d}c \wedge \text{d}c' \\
& =
\text{d}b \wedge \text{d}b' \wedge fc\text{d}c'
- b \text{d}b' \wedge \text{d}(fc) \wedge \text{d}c'
\end{align*}
% P01-E0701: the frozen align display lacks punctuation before the continuing “in other words” clause.
换言之，须证明
$$
(c \text{d}(fb) + fb \text{d}c - fc \text{d}b - b \text{d}(fc))
\wedge \text{d}b' \wedge \text{d}c' = 0.
$$
这由 Leibniz 法则得出。注意，$\gamma$ 在元
$b_0\text{d}b_1 \otimes \text{d}b_2 \otimes \ldots \otimes \text{d}b_p$
上的值是
$\text{d}b_0 \wedge \text{d}b_1 \wedge \ldots \wedge \text{d}b_p$，
所以如此得到的映射
$\text{d} : \Omega^p_{B/A} \to \Omega^{p + 1}_{B/A}$
满足 (\ref{algebra-equation-rule})。

\medskip\noindent
最后，由于 $\Omega^p_{B/A}$ 作为加法群由元
$b_0\text{d}b_1 \wedge \ldots \wedge \text{d}b_p$ 生成，
且
$\text{d}(b_0\text{d}b_1 \wedge \ldots \wedge \text{d}b_p) =
\text{d}b_0 \wedge \ldots \wedge \text{d}b_p$，
完全相同的论证表明，对 $p \geq 1$，复合
$
\Omega^p_{B/A} \to
\Omega^{p + 1}_{B/A} \to
\Omega^{p + 2}_{B/A}
$
为零。因此，de Rham 复形确实是一个复形。

\medskip\noindent
给定一个环 $R$，置 $\Omega_R = \Omega_{R/\mathbf{Z}}$。
它有时称为 $R$ 的绝对微分模；此称谓是合理的：若
$\Omega_R$ 是只要求 Leibniz 法则而不要求 $\text{d}1$ 为零的微分模，
则 Leibniz 法则给出
$\text{d}1 = \text{d}(1 \cdot 1) =
1 \text{d}1 + 1 \text{d}1 = 2 \text{d}1$，从而
$\Omega_R$ 中 $\text{d}1 = 0$。在这种情形下，
{\it $R$ 的绝对 de Rham 复形}是相应的复形
$$
\Omega_R^0 \to \Omega_R^1 \to \Omega_R^2 \to \ldots
$$
% P01-E0702: the frozen display lacks a comma before the following “where” clause.
其中置 $\Omega^i_R = \Omega^i_{R/\mathbf{Z}}$，其余类推。

\medskip\noindent
假设有环的交换图
$$
\xymatrix{
B \ar[r] & B' \\
A \ar[r] \ar[u] & A' \ar[u]
}
$$
% P01-E0703: the frozen diagram lacks terminal punctuation before the next sentence.
存在自然的 de Rham 复形映射
$$
\Omega^\bullet_{B/A} \longrightarrow \Omega^\bullet_{B'/A'}
$$
% P01-E0704: the frozen displayed map lacks terminal punctuation before “Namely”.
具体而言，在次数 $0$ 它是映射 $B \to B'$；在次数 $1$，
它是第 \ref{algebra-section-differentials} 节构造的映射
$\Omega_{B/A} \to \Omega_{B'/A'}$；对 $p \geq 2$，它是诱导映射
$\Omega^p_{B/A} = \wedge^p_B(\Omega_{B/A}) \to
\wedge^p_{B'}(\Omega_{B'/A'}) = \Omega^p_{B'/A'}$。
与微分的相容性由公式 (\ref{algebra-equation-rule}) 对微分的刻画得出。

\begin{lemma}
\label{algebra-lemma-de-rham-complex-base-change}
假设有环映射 $A \to A'$ 和 $A \to B$。
置 $B' = B \otimes_A A'$，于是得到如上的图。
则上面定义的典范映射诱导复形同构
$\Omega^\bullet_{B/A} \otimes_A A' = \Omega^\bullet_{B'/A'}$。
\end{lemma}

\begin{proof}
这由引理 \ref{algebra-lemma-differentials-base-change} 以及取外幂与基变换可交换
这一事实得出。
\end{proof}

\begin{lemma}
\label{algebra-lemma-de-rham-complex}
设 $A \to B$ 是环映射，$\pi : \Omega_{B/A} \to \Omega$
是满的 $B$-模映射。用 $\text{d} : B \to \Omega$ 表示 $\pi$
与泛导子 $\text{d}_{B/A} : B \to \Omega_{B/A}$ 的复合。
% P01-E0705: the frozen designation “Denote d ... the composition” omits “by/as”.
置 $\Omega^i = \wedge_B^i(\Omega)$。
假设 $\pi$ 的核作为 $B$-模由如下元
$\omega \in \Omega_{B/A}$ 生成：$\text{d}_{B/A}(\omega) \in
\Omega_{B/A}^2$ 在 $\Omega^2$ 中映到零。
则存在 de Rham 复形
$$
\Omega^0 \to \Omega^1 \to \Omega^2 \to \ldots
$$
其微分由规则
$$
\text{d} : \Omega^p \to \Omega^{p + 1},\quad
\text{d}\left(f_0\text{d}f_1 \wedge \ldots \wedge \text{d}f_p\right) =
\text{d}f_0 \wedge \text{d}f_1 \wedge \ldots \wedge \text{d}f_p
$$
% P01-E0706: the frozen defining display lacks terminal punctuation at the end of the lemma.
\end{lemma}

\begin{proof}
我们将证明存在交换图
$$
\xymatrix{
\Omega_{B/A}^0 \ar[d] \ar[r]_{\text{d}_{B/A}} &
\Omega_{B/A}^1 \ar[d]_\pi \ar[r]_{\text{d}_{B/A}} &
\Omega_{B/A}^2 \ar[d]_{\wedge^2\pi} \ar[r]_{\text{d}_{B/A}} &
\ldots \\
\Omega^0 \ar[r]^{\text{d}} &
\Omega^1 \ar[r]^{\text{d}} &
\Omega^2 \ar[r]^{\text{d}} &
\ldots
}
$$
% P01-E0707: the frozen source runs the diagram directly into a lowercase independent clause without punctuation or a conjunction.
映射 $\text{d}$ 的描述将由上述
$B$ 在 $A$ 上的 de Rham 复形的微分
$\text{d}_{B/A} : \Omega^p_{B/A} \to \Omega^{p + 1}_{B/A}$
的构造得出。由于最左侧竖直箭头是同构，第一个方块成立。
% P01-E0708: the frozen source writes the closed compound adjective “left most” instead of “leftmost”.
由于 $\pi$ 是满射，要得到第二个方块，只须证明
$\text{d}_{B/A}$ 把 $\pi$ 的核映到 $\wedge^2\pi$ 的核中。
已知 $\pi$ 的核的任意元都具有形式
$\sum b_i\omega_i$，其中 $\pi(\omega_i) = 0$ 且
$\wedge^2\pi(\text{d}_{B/A}(\omega_i)) = 0$。
由 $\text{d}_{B/A}$ 的 Leibniz 法则，
$\text{d}_{B/A}(\sum b_i\omega_i) = \sum b_i \text{d}_{B/A}(\omega_i) +
\sum \text{d}_{B/A}(b_i) \wedge \omega_i$。因此它在
$\wedge^2\pi$ 下映到零。

\medskip\noindent
对 $i > 1$，注意 $\wedge^i \pi$ 是满射，其核为映射
$\Ker(\pi) \wedge \Omega^{i - 1}_{B/A}
\to \Omega_{B/A}^i$ 的像。
对 $\omega_1 \in \Ker(\pi)$ 和
$\omega_2 \in \Omega^{i - 1}_{B/A}$，有
$$
\text{d}_{B/A}(\omega_1 \wedge \omega_2) =
\text{d}_{B/A}(\omega_1) \wedge \omega_2 -
\omega_1 \wedge \text{d}_{B/A}(\omega_2)
$$
% P01-E0709: the frozen display lacks a comma before the following “which” clause.
由上面刚刚证明的结论，它属于 $\wedge^{i + 1}\pi$ 的核。
因此得到上图中的第 $(i + 1)$ 个方块。证明完毕。
\end{proof}

\section{有限阶微分算子}
\label{algebra-section-differential-operators}

\noindent
本节引入有限阶微分算子。

\begin{definition}
\label{algebra-definition-differential-operators}
设 $R \to S$ 是环映射，$M$、$N$ 是 $S$-模，$k \geq 0$ 是整数。
归纳地定义一个{\it $k$ 阶微分算子 $D : M \to N$}：
它是一个 $R$-线性映射，使得对所有 $g \in S$，映射
$m \mapsto D(gm) - gD(m)$ 是 $k - 1$ 阶微分算子。
对于基始情形 $k = 0$，定义 $0$ 阶微分算子为 $S$-线性映射。
\end{definition}

\noindent
若 $D : M \to N$ 是 $k$ 阶微分算子，则对所有 $g \in S$，
映射 $gD$ 是 $k$ 阶微分算子。两个 $k$ 阶微分算子的和仍是
这类算子。因此，全体这样的算子所成的集合
$$
\text{Diff}^k(M, N) = \text{Diff}^k_{S/R}(M, N)
$$
是 $S$-模。我们有
$$
\text{Diff}^0(M, N) \subset
\text{Diff}^1(M, N) \subset
\text{Diff}^2(M, N) \subset \ldots
$$
% P01-E0710: the frozen inclusion display lacks terminal punctuation at the end of the paragraph.

\begin{lemma}
\label{algebra-lemma-composition-differential-operators}
设 $R \to S$ 是环映射，$L, M, N$ 是 $S$-模。
若 $D : L \to M$ 和 $D' : M \to N$ 分别是 $k$ 阶和 $k'$ 阶
微分算子，则 $D' \circ D$ 是 $k + k'$ 阶微分算子。
\end{lemma}

\begin{proof}
设 $g \in S$。则把 $x \in L$ 映到
$$
D'(D(gx)) - gD'(D(x)) = D'(D(gx)) - D'(gD(x)) + D'(gD(x)) - gD'(D(x))
$$
的映射是两个较低阶微分算子复合的和。因此对 $k + k'$ 作归纳即可。
\end{proof}

\begin{lemma}
\label{algebra-lemma-module-principal-parts}
设 $R \to S$ 是环映射，$M$ 是 $S$-模，并且 $k \geq 0$。
存在一个 $S$-模 $P^k_{S/R}(M)$ 以及典范同构
$$
\text{Diff}^k_{S/R}(M, N) = \Hom_S(P^k_{S/R}(M), N)
$$
它关于 $S$-模 $N$ 是函子的。
\end{lemma}

\noindent
通过 Yoneda 引理，这说明存在一个 $k$ 阶的“泛”微分算子
$D_{univ} : M \to P^k_{S/R}(M)$，使得对每个 $k$ 阶微分算子
$D : M \to N$，存在唯一的 $S$-线性映射
$\alpha : P^k_{S/R}(M) \to N$ 满足
$D = D_{univ} \circ \alpha$。
% P01-E0711: the frozen universal factorization reverses the composition order; the typed composite is alpha after D_univ.
偶对 $(P^k_{S/R}(M), D_{univ})$ 在唯一同构的意义下唯一。

\begin{proof}
$P^k_{S/R}(M)$ 的存在性可由一般范畴论论证推出
（在此插入未来的参考文献），但我们也给出一个构造。
% P01-E0712: the frozen source contains the editorial placeholder “insert future reference here”.
置 $F = \bigoplus_{m \in M} S[m]$，其中 $[m]$ 是表示与 $m$
对应的直和项中基元的符号。给定任意微分算子
$D : M \to N$，得到 $S$-线性映射 $L_D : F \to N$，
它把 $[m]$ 映到 $D(m)$。若 $D$ 的阶为 $0$，则 $L_D$ 零化诸元
$$
[m + m'] - [m] - [m'],\quad
g_0[m] - [g_0m]
$$
其中 $g_0 \in S$ 且 $m, m' \in M$。
若 $D$ 的阶为 $1$，则 $L_D$ 零化诸元
$$
[m + m'] - [m] - [m'],\quad
f[m] - [fm], \quad
g_0g_1[m] - g_0[g_1m] - g_1[g_0m] + [g_1g_0m]
$$
其中 $f \in R$、$g_0, g_1 \in S$ 且 $m \in M$。
若 $D$ 的阶为 $k$，则 $L_D$ 零化诸元
$[m + m'] - [m] - [m']$、$f[m] - [fm]$ 以及诸元
$$
g_0g_1\ldots g_k[m] - \sum g_0 \ldots \hat g_i \ldots g_k[g_im] + \ldots
+(-1)^{k + 1}[g_0\ldots g_km]
$$
% P01-E0713: the frozen general relation display lacks terminal punctuation before “Conversely”.
反之，若 $S$-线性映射 $L : F \to N$ 零化前一句所列的所有元，
则 $m \mapsto L([m])$ 是 $k$ 阶微分算子。
因此，$P^k_{S/R}(M)$ 是 $F$ 关于这些元所生成子模的商。
\end{proof}

\begin{definition}
\label{algebra-definition-module-principal-parts}
设 $R \to S$ 是环映射，$M$ 是 $S$-模。
引理 \ref{algebra-lemma-module-principal-parts} 中构造的模
$P^k_{S/R}(M)$ 称为 $M$ 的{\it $k$ 阶主部模}。
\end{definition}

\noindent
注意，包含关系
$$
\text{Diff}^0(M, N) \subset
\text{Diff}^1(M, N) \subset
\text{Diff}^2(M, N) \subset \ldots
$$
通过 Yoneda 引理（范畴，引理 \ref{categories-lemma-yoneda}）
对应于满射
$$
\ldots \to P^2_{S/R}(M) \to P^1_{S/R}(M) \to P^0_{S/R}(M) = M
$$
% P01-E0714: the frozen surjection display lacks terminal punctuation at the end of the paragraph.

\begin{example}
\label{algebra-example-derivations-and-differential-operators}
设 $R \to S$ 是环映射，$N$ 是 $S$-模。
注意 $\text{Diff}^1(S, N) = \text{Der}_R(S, N) \oplus N$。
事实上，若 $D : S \to N$ 是 $1$ 阶微分算子，则由
$\sigma_D(g) := D(g) - gD(1)$ 定义的
$\sigma_D : S \to N$ 是 $R$-导子，并且
$D = \sigma_D + \lambda_{D(1)}$，其中
$\lambda_x : S \to N$ 是把 $g$ 映到 $gx$ 的线性映射。
由 $\Omega_{S/R}$ 的泛性质可得
$P^1_{S/R} = \Omega_{S/R} \oplus S$。
\end{example}

\begin{lemma}
\label{algebra-lemma-sequence-of-principal-parts}
设 $R \to S$ 是环映射，$M$ 是 $S$-模。存在典范短正合列
$$
0 \to \Omega_{S/R} \otimes_S M \to P^1_{S/R}(M) \to M \to 0
$$
它关于 $M$ 是函子的，称为{\it 主部列}。
\end{lemma}

\begin{proof}
映射 $P^1_{S/R}(M) \to M$ 已在上面给出。
设 $N$ 是 $S$-模，$D : M \to N$ 是 $1$ 阶微分算子。
对 $m \in M$，映射
$$
g \longmapsto D(gm) - gD(m)
$$
依照 $1$ 阶微分算子的公理，是一个 $R$-导子 $S \to N$。
因此，它对应于线性映射 $D_m : \Omega_{S/R} \to N$，
后者由规则 $a\text{d}b \mapsto aD(bm) - abD(m)$ 决定
（见引理 \ref{algebra-lemma-universal-omega}）。映射
$$
\Omega_{S/R} \times M \longrightarrow N,\quad
(\eta, m) \longmapsto D_m(\eta)
$$
是 $S$-双线性的（细节从略），因而决定一个 $S$-线性映射
$$
\sigma_D : \Omega_{S/R} \otimes_S M \to N
$$
% P01-E0715: the frozen sigma_D display lacks terminal punctuation before the next sentence.
这样便得到关于 $N$ 函子的映射
$\text{Diff}^1(M, N) \to \Hom_S(\Omega_{S/R} \otimes_S M, N)$，
$D \mapsto \sigma_D$。由 Yoneda 引理，这对应于一个映射
$\Omega_{S/R} \otimes_S M \to P^1_{S/R}(M)$。
% P01-E0716: the frozen phrase “this corresponds a map” omits “to”.
由构造立即可见，此映射关于 $M$ 是函子的。序列
$$
\Omega_{S/R} \otimes_S M \to P^1_{S/R}(M) \to M \to 0
$$
是正合的，因为对每个模 $N$，序列
$$
0 \to \Hom_S(M, N) \to
\text{Diff}^1(M, N) \to
\Hom_S(\Omega_{S/R} \otimes_S M, N)
$$
由直接检验是正合的。

\medskip\noindent
要看出 $\Omega_{S/R} \otimes_S M \to P^1_{S/R}(M)$ 是单射，
论证如下。选取正合列
$$
0 \to M' \to F \to M \to 0
$$
其中 $F$ 是自由 $S$-模。这对所有 $N$ 诱导正合列
$$
0 \to \text{Diff}^1(M, N) \to \text{Diff}^1(F, N) \to \text{Diff}^1(M', N)
$$
这证明在交换图
$$
\xymatrix{
0 \ar[r] &
\Omega_{S/R} \otimes_S M' \ar[r] \ar[d] &
P^1_{S/R}(M') \ar[r] \ar[d] &
M' \ar[r] \ar[d] &
0 \\
0 \ar[r] &
\Omega_{S/R} \otimes_S F \ar[r] \ar[d] &
P^1_{S/R}(F) \ar[r] \ar[d] &
F \ar[r] \ar[d] &
0 \\
0 \ar[r] &
\Omega_{S/R} \otimes_S M \ar[r] \ar[d] &
P^1_{S/R}(M) \ar[r] \ar[d] &
M \ar[r] \ar[d] &
0 \\
& 0 & 0 & 0
}
$$
中，中间列是正合的。由
$\Omega_{S/R} \otimes_S -$ 的右正合性，左列是正合的。
由蛇形引理（见第 \ref{algebra-section-snake} 节），只须对自由模 $F$
证明左侧的正合性。利用 $P^1_{S/R}(-)$ 与直和可交换，
归结到 $M = S$ 的情形。此情形由例
\ref{algebra-example-derivations-and-differential-operators} 中的讨论得出。
\end{proof}

\begin{remark}
\label{algebra-remark-functoriality-principal-parts}
假设给定环的交换图
$$
\xymatrix{
B \ar[r] & B' \\
A \ar[u] \ar[r] & A' \ar[u]
}
$$
% P01-E0717: the frozen multi-object “Suppose given ...” sentence lacks a comma after the displayed diagram.
一个 $B$-模 $M$、一个 $B'$-模 $M'$，以及一个 $B$-线性映射
$M \to M'$。则得到相容的模映射系统
$$
\xymatrix{
\ldots \ar[r] &
P^2_{B'/A'}(M') \ar[r] &
P^1_{B'/A'}(M') \ar[r] &
P^0_{B'/A'}(M') \\
\ldots \ar[r] &
P^2_{B/A}(M) \ar[r] \ar[u] &
P^1_{B/A}(M) \ar[r] \ar[u] &
P^0_{B/A}(M) \ar[u]
}
$$
% P01-E0718: the frozen system diagram lacks terminal punctuation before the next sentence.
这些映射与这类映射的进一步复合相容。
看出这一点最容易的方法是使用引理
\ref{algebra-lemma-module-principal-parts} 的证明中通过生成元和关系给出的
$P^k_{B/A}(M)$ 的描述；也可以直接从这些模的泛性质看出。
% P01-E0719: the frozen compound sentence lacks a comma before “but”.
此外，这些映射与引理 \ref{algebra-lemma-sequence-of-principal-parts}
的短正合列相容。
\end{remark}

\begin{lemma}
\label{algebra-lemma-principal-parts-base-change}
假设有环映射 $A \to A'$ 和 $A \to B$，并有一个 $B$-模 $M$。
置 $B' = B \otimes_A A'$，并把 $M' = M \otimes_A A'$
视为 $B'$-模。评注 \ref{algebra-remark-functoriality-principal-parts} 的映射
诱导同构 $P^k_{B/A}(M) \otimes_A A' = P^k_{B'/A'}(M')$。
\end{lemma}

\begin{proof}
令 $D_{univ} : M \to P^k_{B/A}(M)$ 是泛微分算子。
诱导的 $A'$-线性映射
$D_{univ} \otimes 1 : M' = M \otimes_A A' \to P^k_{B/A}(M) \otimes_A A'$
容易验证为关于 $M'/B'/A'$ 的 $k$ 阶微分算子。
我们将证明 $D_{univ} \otimes 1$ 满足泛性质。
设 $D' : M' \to N'$ 是一个取值于 $B'$-模 $N'$ 的 $k$ 阶微分算子。
则复合 $D : M \to M' \to N'$ 是微分算子，
其中把其靶模 $N'$ 视为 $B$-模。因此存在 $B$-线性映射
$\gamma : P^k_{B/A}(M) \to N'$ 使得
$D = \gamma \circ D_{univ}$。读者容易验证，诱导的 $B'$-线性映射
$\gamma' : P^k_{B/A}(M) \otimes_A A' \to N'$ 满足
$\gamma' \circ (D_{univ} \otimes 1) = D'$。
$\gamma'$ 的唯一性证明从略。
\end{proof}

\begin{lemma}
\label{algebra-lemma-principal-parts-diagonal}
设 $R \to S$ 是环映射，$M$ 是 $S$-模。令
$J = \Ker(S \otimes_R S \to S)$ 为乘法映射的核。
存在典范的 $S$-模同构
$$
P^k_{S/R}(M) \longrightarrow (S \otimes_R M)/J^{k + 1}(S \otimes_R M)
$$
其中 $s \in S$ 通过乘以 $s \otimes 1$ 作用在靶上，并且
$k$ 阶泛微分算子对应于映射
$m \mapsto \text{下列对象的类： }1 \otimes m$。
% P01-E0720: the frozen statement omits the verb “corresponds” and mismatches singular/plural in its universal-operator clause.
\end{lemma}

\begin{proof}
考虑映射 $T : M \to S \otimes M$，$m \mapsto 1 \otimes m$。
% P01-E0721: the frozen map writes S tensor M without the base-ring subscript R used by every surrounding tensor product.
% P01-E0722: the frozen sentence also omits punctuation after the displayed map before “Since”.
由于 $T$ 是 $R$-线性的，而且
\begin{align*}
g_0g_1\ldots g_k T(m)
- \sum g_0 \ldots \hat g_i \ldots g_k T(g_im) + \ldots
+(-1)^{k + 1}T(g_0\ldots g_km) \\
=
\prod_{i = 0,\ldots,k} (g_i \otimes 1 - 1 \otimes g_i) 1 \otimes m
\end{align*}
对 $m \in M$ 和 $g_0, \ldots, g_k \in S$ 成立，所以可知引理陈述中的
规则 $m \mapsto \text{下列对象的类： }1 \otimes m$ 是 $k$ 阶微分算子
（见引理 \ref{algebra-lemma-module-principal-parts} 的证明）。
由 $P^k_{S/R}(M)$ 的泛性质，得到引理陈述中的箭头。
另一方面，若 $D : M \to N$ 是 $k$ 阶微分算子，则可考虑
$S$-线性映射
$L : S \otimes_R M \to N$，$g \otimes m \mapsto gD(m)$。
读者用与上述类似的计算以及 $J$ 由形如
$g \otimes 1 - 1 \otimes g$ 的元生成这一事实，可以验证
$L$ 零化 $J^{k + 1}(S \otimes_R M)$。
特别地，把它应用于泛微分算子
$D_{univ} : M \to P^k_{S/R}(M)$，便得到 $S$-线性映射
$(S \otimes_R M)/J^{k + 1}(S \otimes_R M) \to P^k_{S/R}(M)$。
此映射是引理陈述中映射的逆；留给读者验证。
\end{proof}

\begin{lemma}
\label{algebra-lemma-differentials-de-rham-complex-order-1}
设 $A \to B$ 是环映射。微分
$\text{d} : \Omega^i_{B/A}  \to \Omega^{i + 1}_{B/A}$
是 $1$ 阶微分算子。
\end{lemma}

\begin{proof}
给定 $b \in B$，必须证明
$\text{d} \circ b - b \circ \text{d}$ 是线性算子。
因此必须证明
$$
\text{d} \circ b \circ b' - b \circ \text{d} \circ b' -
b' \circ \text{d} \circ b + b' \circ b \circ \text{d} = 0
$$
% P01-E0723: the frozen commutator display lacks terminal punctuation before the next sentence.
为此，只须在 $\Omega^i_{B/A}$ 的加法生成元上检验。
因而只须证明
$$
\text{d}(bb'b_0\text{d}b_1 \wedge \ldots \wedge \text{d}b_i) -
b\text{d}(b'b_0\text{d}b_1 \wedge \ldots \wedge \text{d}b_i) -
b'\text{d}(bb_0\text{d}b_1 \wedge \ldots \wedge \text{d}b_i) +
bb'\text{d}(b_0\text{d}b_1 \wedge \ldots \wedge \text{d}b_i)
$$
为零。这是使用 Leibniz 法则的一项愉快计算，留给读者完成。
\end{proof}

\begin{lemma}
\label{algebra-lemma-check-differential-operators}
设 $A \to B$ 是环映射。设 $g_i \in B$（$i \in I$）是
$B$ 作为 $A$-代数的一组生成元。设 $M, N$ 是 $B$-模，
$D : M \to N$ 是 $A$-线性映射。要证明 $D$ 是 $k$ 阶微分算子，
只须证明对 $i \in I$，$D \circ g_i - g_i \circ D$
是 $k - 1$ 阶微分算子。
% P01-E0724: the frozen criterion omits k >= 1, although order k - 1 is undefined by the preceding definition when k = 0.
\end{lemma}

\begin{proof}
也就是说，我们断言满足
$D \circ g - g \circ D$ 是 $k - 1$ 阶微分算子的元
$g \in B$ 构成 $B$ 的一个 $A$-子代数。这由关系
$$
D \circ (g + g') - (g + g') \circ D =
(D \circ g - g \circ D) + (D \circ g' - g' \circ D)
$$
以及
$$
D \circ gg' - gg' \circ D =
(D \circ g - g \circ D) \circ g' + g \circ (D \circ g' - g' \circ D)
$$
% P01-E0725: the frozen product display lacks terminal punctuation before “Strictly speaking”.
得出。严格来说，为了对乘积得出结论，还要使用引理
\ref{algebra-lemma-composition-differential-operators}。
\end{proof}

\begin{lemma}
\label{algebra-lemma-invert-system-differential-operators}
设 $A \to B$ 是环映射，$M, N$ 是 $B$-模，
$S \subset B$ 是乘法子集。任意 $k$ 阶微分算子
$D : M \to N$ 唯一地延拓为 $k$ 阶微分算子
$E : S^{-1}M \to S^{-1}N$。
\end{lemma}

\begin{proof}
对 $k$ 作归纳。若 $k = 0$，则 $D$ 是 $B$-线性的，
所以由局部化的函子性得到延拓。给定 $b \in B$，
算子 $L_b : m \mapsto D(bm) - bD(m)$ 的阶为 $k - 1$。
因此，由归纳假设，它唯一地延拓为 $k - 1$ 阶微分算子
$E_b : S^{-1}M \to S^{-1}N$。
此外，计算表明
$L_{b'b} = L_{b'} \circ b + b' \circ L_b$，所以由唯一性得到
$E_{b'b} = E_{b'} \circ b + b' \circ E_b$。
% P01-E0728: the frozen sentence lacks punctuation before “hence”.
类似地，得到
$E_{b'} \circ b - b \circ E_{b'} =
E_b \circ b' - b' \circ E_b$。
现在，对 $m \in M$ 和 $g \in S$，置
$$
E(m/g) = (1/g)(D(m) - E_g(m/g))
$$
% P01-E0729: the frozen defining display lacks terminal punctuation before the next sentence.
要证明这是良定义的，只须证明：对 $g' \in S$，
若使用代表元 $g'm/g'g$，便得到同一结果。计算如下：
\begin{align*}
(1/g'g)(D(g'm) - E_{g'g}(g'm/gg'))
& =
(1/gg')(g'D(m) + E_{g'}(m) - E_{g'g}(g'm/gg')) \\
& =
(1/g'g)(g'D(m) - g' E_g(m/g))
\end{align*}
% P01-E0730: the frozen align display lacks a comma before the following “which” clause.
这与先前结果相同。显然 $E$ 是 $R$-线性的，因为
$D$ 和 $E_g$ 都是 $R$-线性的。
% P01-E0726: the frozen proof uses an undefined ring R; the lemma's base ring is A.
取 $g = 1$ 并使用 $E_1 = 0$，看出 $E$ 延拓 $D$。
由引理 \ref{algebra-lemma-check-differential-operators}，要证明 $E$ 是
$k$ 阶微分算子，现在只须证明以下两个算子都是 $k - 1$ 阶
微分算子：对 $b \in B$ 的 $E \circ b - b \circ E$，以及对
$g' \in S$ 的 $E \circ 1/g' - 1/g' \circ E$。
对于前者，选取 $S^{-1}M$ 中的元 $m/g$，并注意
\begin{align*}
E(b m/g) - bE(m/g)
& =
(1/g)(D(bm) - bD(m) - E_g(bm/g) + bE_g(m/g)) \\
& =
(1/g)(L_b(m) - E_b(m) + gE_b(m/g)) \\
& =
E_b(m/g)
\end{align*}
% P01-E0731: the frozen align display lacks a comma before the following “which” clause.
这是 $k - 1$ 阶微分算子。最后，
\begin{align*}
E(m/g'g) - (1/g')E(m/g)
& =
(1/g'g)(D(m) - E_{g'g}(m/g'g)) -
(1/g'g)(D(m) - E_g(m/g)) \\
& =
-(1/g')E_{g'}(m/g'g)
\end{align*}
% P01-E0732: the frozen align display lacks a comma before the following “which” clause.
这也是 $k - 1$ 阶微分算子，因为它是线性映射
（乘以 $1/g'$ 和符号）与 $E_{g'}$ 的复合。
唯一性的证明从略。
\end{proof}

\begin{lemma}
\label{algebra-lemma-base-change-differential-operators}
设 $R \to A$ 和 $R \to B$ 是环映射。设 $M$ 和 $M'$ 是 $A$-模，
$D : M \to M'$ 是关于 $R \to A$ 的 $k$ 阶微分算子。
设 $N$ 是任意 $B$-模。则映射
$$
D \otimes \text{id}_N : M \otimes_R N \to M' \otimes_R N
$$
是关于 $B \to A \otimes_R B$ 的 $k$ 阶微分算子。
\end{lemma}

\begin{proof}
显然 $D' = D \otimes \text{id}_N$ 是 $B$-线性的。
由引理 \ref{algebra-lemma-check-differential-operators}，只须证明
$$
D' \circ a \otimes 1 -  a \otimes 1 \circ D' =
(D \circ a - a \circ D) \otimes \text{id}_N
$$
是 $k - 1$ 阶微分算子；这由对 $k$ 作归纳得出。
% P01-E0727: the frozen proof invokes the k - 1 criterion without treating k = 0, where negative order was never defined.
\end{proof}

\section{朴素余切复形}
\label{algebra-section-netherlander}

\noindent
设 $R \to S$ 是环映射。用 $R[S]$ 表示以元 $s \in S$ 为变量的多项式环。
% P01-E0733: the frozen designation “Denote R[S] the polynomial ring” omits “by/as”.
用 $[s] \in R[S]$ 表示与 $s \in S$ 对应的变量。
于是，$R[S]$ 是自由 $R$-模，其基元为
$[s_1] \ldots [s_n]$，其中 $s_1, \ldots, s_n$
遍历 $S$ 的所有无序元列。
% P01-E0734: the frozen source has subject-verb disagreement in “s_1,...,s_n ranges”.
存在典范满射
\begin{equation}
\label{algebra-equation-canonical-presentation}
R[S] \longrightarrow S,\quad [s] \longmapsto s
\end{equation}
其核记为 $I \subset R[S]$。容易观察到，$I$ 由诸元
$[s + s'] - [s] - [s']$、$[s][s'] - [ss']$ 和 $[r] - r$ 生成。
由引理 \ref{algebra-lemma-differential-seq}，存在典范映射
\begin{equation}
\label{algebra-equation-naive-cotangent-complex}
I/I^2 \longrightarrow \Omega_{R[S]/R} \otimes_{R[S]} S
\end{equation}
其余核典范同构于 $\Omega_{S/R}$。
注意，$S$-模 $\Omega_{R[S]/R} \otimes_{R[S]} S$
以诸元 $\text{d}[s]$ 为自由生成元。

\begin{definition}
\label{algebra-definition-naive-cotangent-complex}
设 $R \to S$ 是环映射。{\it 朴素余切复形} $\NL_{S/R}$
是链复形 (\ref{algebra-equation-naive-cotangent-complex})
$$
\NL_{S/R} = \left(I/I^2 \longrightarrow \Omega_{R[S]/R} \otimes_{R[S]} S\right)
$$
其中把 $I/I^2$ 放在（同调）次数 $1$，把
$\Omega_{R[S]/R} \otimes_{R[S]} S$ 放在次数 $0$。
用 $H_1(L_{S/R}) = H_1(\NL_{S/R})$\footnote{文献中有时把此模记为
$\Gamma_{S/R}$。}表示次数 $1$ 的同调。
% P01-E0735: the frozen notation sentence “denote ... the homology” omits “by/as”.
\end{definition}

\noindent
在继续之前，先略述真正的余切复形（余切，第
\ref{cotangent-section-cotangent-ring-map} 节）。
给定环映射 $R \to S$，存在典范的单纯 $R$-代数 $P_\bullet$，
其各项都是多项式代数，并带有典范同伦等价
$$
P_\bullet \longrightarrow S
$$
% P01-E0736: the frozen display lacks terminal punctuation before the next sentence.
$S$ 在 $R$ 上的余切复形 $L_{S/R}$ 定义为与单纯模
$$
\Omega_{P_\bullet/R} \otimes_{P_\bullet} S
$$
% P01-E0737: the frozen display lacks terminal punctuation before the next sentence.
相伴的链复形。上面定义的朴素余切复形典范同构于截断
$\tau_{\leq 1}L_{S/R}$（见同调，第
\ref{homology-section-truncations} 节，以及余切，第
\ref{cotangent-section-surjections} 节）。
特别地，确有 $H_1(\NL_{S/R}) = H_1(L_{S/R})$，所以我们的定义
与使用余切复形的定义相容。此外，如上所见，
$H_0(L_{S/R}) = H_0(\NL_{S/R}) = \Omega_{S/R}$。

\medskip\noindent
设 $R \to S$ 是环映射。$S$ 在 $R$ 上的一个{\it 表现}
是 $R$-代数满射 $\alpha : P \to S$，其中 $P$ 是一个
（以某变量集为变量的）多项式代数。通常，当 $S$ 在 $R$ 上有限型时，
我们会说：“设 $R[x_1, \ldots, x_n] \to S$ 是 $S/R$ 的一个表现”；
若还想指明 $I$ 是表现的核，则说：“设
$0 \to I \to R[x_1, \ldots, x_n] \to S \to 0$ 是 $S/R$ 的一个表现”。
注意，用来定义朴素余切复形的映射 $R[S] \to S$ 是一个表现的例子。

\medskip\noindent
注意，对每个表现 $\alpha$，都得到一个两项 $S$-模链复形
% P01-E0738: the frozen compound modifier “two term” should be hyphenated.
$$
\NL(\alpha) : I/I^2 \longrightarrow \Omega_{P/R} \otimes_P S.
$$
这里把项 $I/I^2$ 放在次数 $1$，把项
$\Omega_{P/R} \otimes S$ 放在次数 $0$。
$I/I^2$ 中 $f \in I$ 的类映到
$\Omega_{P/R} \otimes S$ 中的 $\text{d}f \otimes 1$。
% P01-E0739: the frozen two subsequent tensor products omit the base P shown in the defining display.
此复形的余核典范地为 $\Omega_{S/R}$，见引理
\ref{algebra-lemma-differential-seq}。把复形 $\NL(\alpha)$ 称为
{\it 与 $S/R$ 的表现 $\alpha : P \to S$ 相伴的朴素余切复形}。
注意，若 $P = R[S]$ 且取其到 $S$ 的典范满射，则恢复
$\NL_{S/R}$。若 $P = R[x_1, \ldots, x_n]$，则有时用记号
$I/I^2 \to \bigoplus_{i = 1, \ldots, n} S\text{d}x_i$
表示此复形。
% P01-E0740: the frozen phrase “then will sometimes use” omits the subject “we”.

\medskip\noindent
假设给定环的交换图
\begin{equation}
\label{algebra-equation-functoriality-NL}
\vcenter{
\xymatrix{
S \ar[r]_{\phi} & S' \\
R \ar[r] \ar[u] & R' \ar[u]
}
}
\end{equation}
设 $\alpha : P \to S$ 是 $S$ 在 $R$ 上的表现，
$\alpha' : P' \to S'$ 是 $S'$ 在 $R'$ 上的表现。
从 $\alpha : P \to S$ 到 $\alpha' : P' \to S'$ 的一个
{\it 表现态射}定义为 $R$-代数映射
$$
\varphi : P \to P'
$$
使得 $\phi \circ \alpha = \alpha' \circ \varphi$。
注意，这时 $\varphi(I) \subset I'$，其中
$I = \Ker(\alpha)$ 且 $I' = \Ker(\alpha')$。
因此，$\varphi$ 诱导 $S$-模映射
$I/I^2 \to I'/(I')^2$，并由微分的函子性还诱导 $S$-模映射
$\Omega_{P/R} \otimes S \to \Omega_{P'/R'} \otimes S'$。
% P01-E0741: the frozen functoriality formula omits the tensor bases P and P'.
这些映射与 $\NL(\alpha)$ 和 $\NL(\alpha')$ 的微分相容，
从而得到朴素余切复形的映射
$$
\NL(\alpha) \longrightarrow \NL(\alpha').
$$
考虑诱导映射 $\NL(\alpha) \otimes_S S' \to \NL(\alpha')$
往往很方便。

\medskip\noindent
在特殊情形 $P = R[S]$、$P' = R'[S']$ 下，映射
$\phi : S \to S'$ 通过规则 $[s] \mapsto [\phi(s)]$
诱导典范环映射 $\varphi : P \to P'$。
因此上述构造确定了与图 (\ref{algebra-equation-functoriality-NL}) 相伴的
典范（！）链复形映射
$$
\NL_{S/R} \longrightarrow \NL_{S'/R'},\quad\text{且}\quad
\NL_{S/R} \otimes_S S' \longrightarrow \NL_{S'/R'}
$$
注意，此构造与复合相容：给定交换图
$$
\xymatrix{
S \ar[r]_{\phi} & S' \ar[r]_{\phi'} & S'' \\
R \ar[r] \ar[u] & R' \ar[u] \ar[r] & R'' \ar[u]
}
$$
% P01-E0742: the frozen diagram lacks a comma before the continuing “we see” clause.
可见，复合
$$
\NL_{S/R} \longrightarrow \NL_{S'/R'} \longrightarrow \NL_{S''/R''}
$$
就是外方块给出的映射 $\NL_{S/R} \to \NL_{S''/R''}$。

\medskip\noindent
事实证明，$\NL(\alpha)$ 与 $\NL_{S/R}$ 同伦等价，
且上面构造的映射在同伦意义下良定义（复形映射的同伦见同调，第
\ref{homology-section-complexes} 节；但在下面引理的证明中，
我们也明确说明这些陈述的确切含义）。

\begin{lemma}
\label{algebra-lemma-NL-homotopy}
假设给定图 (\ref{algebra-equation-functoriality-NL})。
设 $\alpha : P \to S$ 和 $\alpha' : P' \to S'$ 是表现。
\begin{enumerate}
\item 存在一个从 $\alpha$ 到 $\alpha'$ 的表现态射。
\item 任意两个表现态射诱导同伦的复形态射
$\NL(\alpha) \to \NL(\alpha')$。
\item 该构造与表现态射的复合相容（确切陈述见证明）。
\item 若 $R \to R'$ 和 $S \to S'$ 是同构，则对任意从
$\alpha$ 到 $\alpha'$ 的表现映射 $\varphi$，诱导映射
$\NL(\alpha) \to \NL(\alpha')$ 是同伦等价和拟同构。
\end{enumerate}
特别地，把 $\alpha$ 与典范表现
(\ref{algebra-equation-canonical-presentation}) 比较，得到一个在同伦意义下
良定义并与所有函子性相容（不计同伦）的拟同构
$\NL(\alpha) \to \NL_{S/R}$。
\end{lemma}

\begin{proof}
由于 $P$ 是 $R$ 上的多项式代数，可对某个集合 $A$ 写成
$P = R[x_a, a \in A]$。由于 $\alpha'$ 是满射，可对每个
$a \in A$ 选取 $f_a \in P'$，使得
$\alpha'(f_a) = \phi(\alpha(x_a))$。令
$\varphi : P = R[x_a, a \in A] \to P'$ 为满足
$\varphi(x_a) = f_a$ 的唯一 $R$-代数映射。
这给出 (1) 中的态射。

\medskip\noindent
设 $\varphi$ 和 $\varphi'$ 是从 $\alpha$ 到 $\alpha'$ 的表现态射。
% P01-E0743: the frozen source omits “be” in “Let phi and phi' morphisms”.
令 $I = \Ker(\alpha)$ 且 $I' = \Ker(\alpha')$。
必须构造图中的对角映射 $h$：
$$
\xymatrix{
I/I^2 \ar[r]^-{\text{d}}
\ar@<1ex>[d]^{\varphi'_1} \ar@<-1ex>[d]_{\varphi_1}
&
\Omega_{P/R} \otimes_P S
\ar@<1ex>[d]^{\varphi'_0} \ar@<-1ex>[d]_{\varphi_0}
\ar[ld]_h
\\
I'/(I')^2 \ar[r]^-{\text{d}}
&
\Omega_{P'/R'} \otimes_{P'} S'
}
$$
其中竖直映射由 $\varphi$、$\varphi'$ 诱导，并且该对角映射满足
% P01-E0744: the frozen “where ... induced by phi, phi' such that” clause omits a conjunction and the subject h.
$$
\varphi_1 - \varphi'_1 = h \circ \text{d}
\quad\text{且}\quad
\varphi_0 - \varphi'_0 = \text{d} \circ h
$$
考虑映射 $\varphi - \varphi' : P \to P'$。
由于 $\varphi$ 和 $\varphi'$ 都与 $\alpha$、$\alpha'$ 相容，
得到 $\varphi - \varphi' : P \to I'$。
这说明 $\varphi, \varphi' : P \to P'$ 在 $I'/(I')^2$ 上诱导同一
$P$-模结构，因为
$\varphi(p)i' - \varphi'(p)i' = (\varphi - \varphi')(p)i' \in (I')^2$。
此外，$\varphi - \varphi'$ 是 $R$-线性的，并且
$$
(\varphi - \varphi')(fg) =
\varphi(f)(\varphi - \varphi')(g) + (\varphi - \varphi')(f)\varphi'(g)
$$
所以诱导映射 $D : P \to I'/(I')^2$ 是一个 $R$-导子。
% P01-E0745: the frozen source says “a R-derivation” instead of “an R-derivation”.
从而得到典范映射
$h : \Omega_{P/R} \otimes_P S \to I'/(I')^2$，
使得 $D = h \circ \text{d}$。计算（从略）表明 $h$ 正是所需同伦。

\medskip\noindent
假设有交换图
$$
\xymatrix{
S \ar[r]_{\phi} & S' \ar[r]_{\phi'} & S'' \\
R \ar[r] \ar[u] & R' \ar[u] \ar[r] & R'' \ar[u]
}
$$
并且
\begin{enumerate}
\item $\alpha : P \to S$，
\item $\alpha' : P' \to S'$，以及
\item $\alpha'' : P'' \to S''$
\end{enumerate}
是表现。假设
\begin{enumerate}
\item $\varphi : P \to P'$ 是从 $\alpha$ 到 $\alpha'$ 的表现态射，并且
\item $\varphi' : P' \to P''$ 是从 $\alpha'$ 到 $\alpha''$ 的表现态射。
\end{enumerate}
则立即可见，$\varphi' \circ \varphi : P \to P''$
是从 $\alpha$ 到 $\alpha''$ 的表现态射，并且由
$\varphi$ 和 $\varphi'$ 诱导的映射
$\NL(\alpha) \to \NL(\alpha')$ 与
$\NL(\alpha') \to \NL(\alpha'')$ 的复合，正是诱导的朴素余切复形映射
$\NL(\alpha) \to \NL(\alpha'')$。

\medskip\noindent
在两项复形的简单情形下，拟同构就是在两项间映射的核和余核上都诱导
同构的映射。
% P01-E0746: the frozen compound modifier “2 term” should be hyphenated.
注意，同伦的两项复形映射（如上所述）在核和余核上定义相同的映射。
因此，若 $\varphi$ 是从 $S$ 在 $R$ 上的表现 $\alpha$
到自身的映射，则诱导映射
$\NL(\alpha) \to \NL(\alpha)$ 是拟同构，因为由 (2) 它同伦于恒等映射。
% P01-E0747: the frozen phrase “is a quasi-isomorphism being homotopic” uses an awkward/dangling participial construction.
为了在一般情形下证明 (4)，考虑由 (1) 存在的从
$\alpha'$ 到 $\alpha$ 的态射 $\varphi'$。复合
$\NL(\alpha) \to \NL(\alpha') \to \NL(\alpha)$ 和
$\NL(\alpha') \to \NL(\alpha) \to \NL(\alpha')$
由 (3) 同伦于恒等映射，因此按照定义，这些映射是同伦等价。
% P01-E0748: compatibility in (3) alone does not give homotopy to the identity; the argument also needs part (2).
形式上立即推出，两个映射
$\NL(\alpha) \to \NL(\alpha')$ 和
$\NL(\alpha') \to \NL(\alpha)$ 都是拟同构。略去一些细节。
\end{proof}

\begin{lemma}
\label{algebra-lemma-NL-polynomial-algebra}
设 $A \to B$ 是多项式代数。则 $\NL_{B/A}$ 与链复形
$(0 \to \Omega_{B/A})$ 同伦等价，其中 $\Omega_{B/A}$ 位于次数 $0$。
\end{lemma}

\begin{proof}
这由引理 \ref{algebra-lemma-NL-homotopy} 以及如下事实得出：
$\text{id}_B : B \to B$ 是核为零的 $B$ 在 $A$ 上的表现。
% P01-E0749: the frozen proof is the sentence fragment “Follows from ...”.
\end{proof}

\noindent
下面的引理是引入朴素余切复形的部分动机。
余切复形把它扩张为真正的长正合上同调列。
若 $B \to C$ 是局部完全交，则可在序列左侧补上一个零；
见代数进阶，引理 \ref{more-algebra-lemma-transitive-lci-at-end}。

\begin{lemma}[Jacobi--Zariski 序列]
\label{algebra-lemma-exact-sequence-NL}
设 $A \to B \to C$ 是环映射。选取核为 $I$ 的表现
$\alpha : A[x_s, s \in S] \to B$，并选取核为 $J$ 的表现
$\beta : B[y_t, t \in T] \to C$。令
$\gamma : A[x_s, y_t] \to C$ 为所诱导的 $C$ 的表现，其核为 $K$。
则得到典范交换图
$$
\xymatrix{
0 \ar[r] &
\Omega_{A[x_s]/A} \otimes C \ar[r] &
\Omega_{A[x_s, y_t]/A} \otimes C \ar[r] &
\Omega_{B[y_t]/B} \otimes C \ar[r] &
0 \\
&
I/I^2 \otimes C \ar[r] \ar[u] &
K/K^2 \ar[r] \ar[u] &
J/J^2 \ar[r] \ar[u] &
0
}
$$
其两行正合。得到如下同调群正合列
$$
H_1(\NL_{B/A} \otimes_B C) \to
H_1(L_{C/A}) \to
H_1(L_{C/B}) \to
C \otimes_B \Omega_{B/A} \to
\Omega_{C/A} \to
\Omega_{C/B} \to 0
$$
它是扩张引理 \ref{algebra-lemma-exact-sequence-differentials} 中序列的
$C$-模序列。若
$\text{Tor}_1^B(\Omega_{B/A}, C) = 0$ 且
$\text{Tor}_2^B(\Omega_{B/A}, C) = 0$，则
$H_1(\NL_{B/A} \otimes_B C) = H_1(L_{B/A}) \otimes_B C$。
\end{lemma}

\begin{proof}
映射的精确定义从略。上行的正合性来自如下事实：
$\text{d}x_s$、$\text{d}y_t$ 构成中间模的一组基。
映射 $\gamma$ 分解为
$$
A[x_s, y_t] \to B[y_t] \to C
$$
其中第一支箭头是满射，第二支箭头等于 $\beta$。
因此看出 $K \to J$ 是满射。此外，第一个显示箭头的核是
$IA[x_s, y_t]$。故 $I/I^2 \otimes C$ 满射到
$K/K^2 \to J/J^2$ 的核上。最后，可用引理
\ref{algebra-lemma-NL-homotopy} 把各项认作朴素余切复形的同调群。

\medskip\noindent
最后一个断言是同调代数中的陈述。回忆
$\NL_{B/A} = (N^{-1} \to N^0)$ 是一个两项 $B$-模复形，
其中 $N^0$ 自由，而上同调模为
$H^0 = \Omega_{B/A}$ 和 $H^{-1} = H_1(L_{B/A})$。
记 $M \subset N^0$ 为微分的像。若
$\text{Tor}_1^B(H^0, C) = 0$，则有正合列
$$
0 \to M \otimes_B C \to N^0 \otimes_B C \to H^0 \otimes_B C \to 0
$$
% P01-E0750: the frozen display lacks terminal punctuation before the next sentence.
由于 $N^0$ 自由，还看出
$\text{Tor}_2^B(H^0, C) = \text{Tor}_1^B(M, C)$。
因此，若 $\text{Tor}_2^B(H^0, C) = 0$，则还有正合列
$$
0 \to H^{-1} \otimes_B C \to N^{-1} \otimes_B C \to M \otimes_B C \to 0
$$
% P01-E0751: the frozen display lacks terminal punctuation before the next sentence.
综合起来可见，若 $\text{Tor}_1^B(H^0, C) = 0$ 且
$\text{Tor}_2^B(H^0, C) = 0$，则
$H^{-1} \otimes_B C$ 正是
$N^{-1} \otimes_B C \to N^0 \otimes_B C$ 的核，正如所需。
\end{proof}

\begin{remark}
\label{algebra-remark-composition-homotopy-equivalent-to-zero}
设 $A \to B$ 和 $\phi : B \to C$ 是环映射。
则复合 $\NL_{B/A} \to \NL_{C/A} \to \NL_{C/B}$ 同伦于零。
具体而言，此复合是朴素余切复形关于方块
$$
\xymatrix{
B \ar[r]_\phi & C \\
A \ar[r] \ar[u] & B \ar[u]
}
$$
% P01-E0752: the frozen diagram lacks terminal punctuation before the next sentence.
的函子性。写 $J = \Ker(B[C] \to C)$。
一个显式同伦由映射
$\Omega_{A[B]/A} \otimes_{A[B]} B \to J/J^2$ 给出，
它把基元 $\text{d}[b]$ 映到 $J/J^2$ 中
$[\phi(b)] - b$ 的类。
\end{remark}

\begin{lemma}
\label{algebra-lemma-NL-surjection}
设 $A \to B$ 是核为 $I$ 的满环映射。
则 $\NL_{B/A}$ 与链复形 $(I/I^2 \to 0)$ 同伦等价，
其中 $I/I^2$ 位于次数 $1$。特别地，
$H_1(L_{B/A}) = I/I^2$。
\end{lemma}

\begin{proof}
这由引理 \ref{algebra-lemma-NL-homotopy} 以及如下事实得出：
$A \to B$ 是 $B$ 在 $A$ 上的一个表现。
% P01-E0753: the frozen proof is the sentence fragment “Follows from ...”.
\end{proof}

\begin{lemma}
\label{algebra-lemma-application-NL}
设 $A \to B \to C$ 是环映射。假设 $A \to C$ 是满射
（所以 $B \to C$ 也是满射）。记
$I = \Ker(A \to C)$ 且 $J = \Ker(B \to C)$。
则序列
$$
I/I^2 \to J/J^2 \to \Omega_{B/A} \otimes_B B/J \to 0
$$
是正合的。
\end{lemma}

\begin{proof}
这由引理 \ref{algebra-lemma-exact-sequence-NL} 以及引理
\ref{algebra-lemma-NL-surjection} 中对朴素余切复形
$\NL_{C/B}$ 和 $\NL_{C/A}$ 的描述得出。
% P01-E0754: the frozen proof is the sentence fragment “Follows from ...”.
\end{proof}

\begin{lemma}[平坦基变换]
\label{algebra-lemma-change-base-NL}
设 $R \to S$ 是环映射，$\alpha : P \to S$ 是一个表现，
$R \to R'$ 是平坦环映射。令
$\alpha' : P \otimes_R R' \to S' = S \otimes_R R'$
为诱导表现。则
$\NL(\alpha) \otimes_R R' = \NL(\alpha) \otimes_S S' = \NL(\alpha')$。
特别地，若 $R \to R'$ 平坦，则典范映射
$$
\NL_{S/R} \otimes_S S' \longrightarrow \NL_{S \otimes_R R'/R'}
$$
是同伦等价。
\end{lemma}

\begin{proof}
这是因为 $R \to R'$ 平坦，故
$\Ker(\alpha') = R' \otimes_R \Ker(\alpha)$。
\end{proof}

\begin{lemma}
\label{algebra-lemma-colimits-NL}
设 $R_\lambda \to S_\lambda$ 是有向集 $\Lambda$ 上的环映射系统。
置 $R = \colim R_\lambda$ 且 $S = \colim S_\lambda$。
则 $\NL_{S/R} = \colim \NL_{S_\lambda/R_\lambda}$。
\end{lemma}

\begin{proof}
回忆 $\NL_{S/R}$ 是复形
$I/I^2 \to \bigoplus_{s \in S} S\text{d}[s]$，其中
$I \subset R[S]$ 是典范表现 $R[S] \to S$ 的核。
现在显然 $R[S] = \colim R_\lambda[S_\lambda]$；类似地，
$I = \colim I_\lambda$，其中
$I_\lambda = \Ker(R_\lambda[S_\lambda] \to S_\lambda)$。
引理即得。
\end{proof}

\begin{lemma}
\label{algebra-lemma-NL-of-localization}
若 $S \subset A$ 是 $A$ 的一个乘法子集，则
$\NL_{S^{-1}A/A}$ 与零复形同伦等价。
\end{lemma}

\begin{proof}
由于 $A \to S^{-1}A$ 平坦，由平坦基变换
（引理 \ref{algebra-lemma-change-base-NL}）可知
$\NL_{S^{-1}A/A} \otimes_A S^{-1}A \to \NL_{S^{-1}A/S^{-1}A}$
是同伦等价。箭头的源同构于 $\NL_{S^{-1}A/A}$，
而箭头的目标为零（由引理 \ref{algebra-lemma-NL-surjection}），故得证。
\end{proof}

\begin{lemma}
\label{algebra-lemma-NL-localize-bottom}
设 $S \subset A$ 是 $A$ 的一个乘法子集。
设 $S^{-1}A \to B$ 是环映射。
则 $\NL_{B/A} \to \NL_{B/S^{-1}A}$ 是同伦等价。
% P01-E0755: the frozen statement has “Let S ... is” instead of “be”.
\end{lemma}

\begin{proof}
选取 $B$ 在 $A$ 上的一个表现 $\alpha : P \to B$。
则 $\beta : S^{-1}P \to B$ 是 $B$ 在 $S^{-1}A$ 上的一个表现。
直接计算表明 $\NL(\alpha) = \NL(\beta)$；又由引理
\ref{algebra-lemma-NL-homotopy}，朴素余切复形在同伦意义下良定义，
故引理成立。
% P01-E0756: the frozen concluding clause is weakly joined by an unpunctuated “as”.
\end{proof}

\begin{lemma}
\label{algebra-lemma-principal-localization-NL}
\begin{slogan}
朴素余切复形的构造与在一个元素处的局部化交换。
\end{slogan}
设 $A \to B$ 是环映射。设 $g \in B$。假设 $\alpha : P \to B$
是核为 $I$ 的一个表现。则 $B_g$ 在 $A$ 上的一个表现是映射
$$
\beta : P[x] \longrightarrow B_g
$$
它延拓 $\alpha$ 并把 $x$ 映到 $1/g$。
$\beta$ 的核 $J$ 由 $I$ 和元素 $f x - 1$ 生成，
其中 $f \in P$ 是由 $\alpha$ 映到 $g \in B$ 的一个元素。在此情形下，
\begin{enumerate}
\item $J/J^2 = (I/I^2)_g \oplus B_g (f x - 1)$,
\item $\Omega_{P[x]/A} \otimes_{P[x]} B_g =
\Omega_{P/A} \otimes_P B_g \oplus B_g \text{d}x$,
\item $\NL(\beta) \cong
\NL(\alpha) \otimes_B B_g \oplus (B_g \xrightarrow{g} B_g)$
\end{enumerate}
% P01-E0757: the frozen enumeration has no terminal punctuation.
因此典范映射 $\NL_{B/A} \otimes_B B_g \to \NL_{B_g/A}$
是同伦等价。
\end{lemma}

\begin{proof}
由于 $P[x]/(I, fx - 1) = B[x]/(gx - 1) = B_g$，便得到关于
$I$ 和 $fx - 1$ 生成 $J$ 的断言。考虑交换图
$$
\xymatrix{
0 \ar[r] &
\Omega_{P/A} \otimes B_g \ar[r] &
\Omega_{P[x]/A} \otimes B_g \ar[r] &
\Omega_{B[x]/B} \otimes B_g \ar[r] &
0 \\
&
(I/I^2)_g \ar[r] \ar[u] &
J/J^2 \ar[r] \ar[u] &
(gx - 1)/(gx - 1)^2 \ar[r] \ar[u] &
0
}
$$
其两行由引理 \ref{algebra-lemma-exact-sequence-NL} 知为正合。
% P01-E0758: several tensor products in the frozen diagram and prose omit their inferable base rings.
$B_g$-模 $\Omega_{B[x]/B} \otimes B_g$ 是以
$\text{d}x$ 为基的秩 $1$ 自由模。
$B_g$-模 $\Omega_{P[x]/A} \otimes B_g$ 中的元素 $\text{d}x$
给出上面一行的一个分裂。元素 $gx - 1 \in (gx - 1)/(gx - 1)^2$
被映到 $\Omega_{B[x]/B} \otimes B_g$ 中的 $g\text{d}x$，
因而 $(gx - 1)/(gx - 1)^2$ 是秩为 $1$ 的自由 $B_g$-模。
（也可如下说明：$gx - 1$ 在 $B[x]$ 中是非零因子，
因为它是常数项可逆的多项式；而由引理
\ref{algebra-lemma-regular-quasi-regular}，任何非零因子均给出长度为 $1$
的拟正则序列。）

\medskip\noindent
我们来证明 $(I/I^2)_g \to J/J^2$ 是单射。考虑 $P$-代数映射
$$
\pi : P[x] \to (P/I^2)_f = P_f/I_f^2
$$
它把 $x$ 映到 $1/f$。由于 $J$ 由 $I$ 和 $fx - 1$ 生成，
可知 $\pi(J) \subset (I/I^2)_f = (I/I^2)_g$。由于后者是平方为零的理想，
可知 $\pi(J^2) = 0$。
% P01-E0759: the frozen sentence “Let us prove ... injective” omits “is”.
若 $a \in I$ 在 $J$ 中映到 $J^2$ 的一个元素，则
$\pi(a) = 0$，这蕴含 $a$ 在 $I_f/I_f^2$ 中的像为零。
% P01-E0760: the frozen phrase “maps to an element of J^2 in J” is grammatically malformed.
这证明了所需的单射性。

\medskip\noindent
因此有二项复形的短正合序列
% P01-E0761: the frozen compound adjective “two term” is unhyphenated.
$$
0 \to \NL(\alpha) \otimes_B B_g \to \NL(\beta)
\to (B_g \xrightarrow{g} B_g) \to 0
$$
这样的短正合序列在复形范畴中总可分裂。
% P01-E0762: the frozen blanket splitting claim is false for arbitrary short exact sequences of complexes; the explicit termwise splittings below justify only this case.
在本例中，可取如下分裂：
$$
J/J^2 = (I/I^2)_g \oplus B_g (fx - 1)\quad\text{且}\quad
\Omega_{P[x]/A} \otimes B_g = \Omega_{P/A} \otimes B_g \oplus
B_g (g^{-2}\text{d}f + \text{d}x)
$$
% P01-E0763: the frozen display lacks terminal punctuation before “This works”.
这是因为在 $\Omega_{P[x]/A} \otimes B_g$ 中，
$\text{d}(fx - 1) = x\text{d}f + f \text{d}x =
g(g^{-2}\text{d}f + \text{d}x)$。
\end{proof}

\begin{lemma}
\label{algebra-lemma-localize-NL}
设 $A \to B$ 是环映射。设 $S \subset B$ 是乘法子集。
典范映射 $\NL_{B/A} \otimes_B S^{-1}B \to \NL_{S^{-1}B/A}$
是拟同构。
\end{lemma}

\begin{proof}
我们有 $S^{-1}B = \colim_{g \in S} B_g$，这里把 $S$
视为有向集（以整除关系为序）；见引理
\ref{algebra-lemma-localization-colimit}。
由引理 \ref{algebra-lemma-principal-localization-NL}，每个映射
$\NL_{B/A} \otimes_B B_g \to \NL_{B_g/A}$
都是拟同构。
% P01-E0764: the frozen subject “each” incorrectly takes the plural verb “are”.
本引理由引理 \ref{algebra-lemma-colimits-NL} 得出。
\end{proof}

\begin{lemma}
\label{algebra-lemma-sum-two-terms}
设 $R$ 是环。
设 $A_1 \to A_0$ 和 $B_1 \to B_0$ 是二项复形。
% P01-E0765: the frozen sentence has both a spurious comma after the first complex and an unhyphenated “two term”.
假设存在复形态射 $\varphi : A_\bullet \to B_\bullet$
和 $\psi : B_\bullet \to A_\bullet$，使得
$\varphi \circ \psi$ 和 $\psi \circ \varphi$
均与恒等映射同伦。
则作为 $R$-模，$A_1 \oplus B_0 \cong B_1 \oplus A_0$。
\end{lemma}

\begin{proof}
选取映射 $h : A_0 \to A_1$，使得
$$
\text{id}_{A_1} - \psi_1 \circ \varphi_1 = h \circ d_A
\text{ 且 }
\text{id}_{A_0} - \psi_0 \circ \varphi_0 = d_A \circ h.
$$
类似地，选取映射 $h' : B_0 \to B_1$，使得
$$
\text{id}_{B_1} - \varphi_1 \circ \psi_1 = h' \circ d_B
\text{ 且 }
\text{id}_{B_0} - \varphi_0 \circ \psi_0 = d_B \circ h'.
$$
直接计算表明
$$
\left(
\begin{matrix}
\text{id}_{A_1} & -h' \circ \psi_1 + h \circ \psi_0 \\
0 & \text{id}_{B_0}
\end{matrix}
\right)
=
\left(
\begin{matrix}
\psi_1 & h \\
-d_B & \varphi_0
\end{matrix}
\right)
\left(
\begin{matrix}
\varphi_1 & - h' \\
d_A & \psi_0
\end{matrix}
\right)
$$
% P01-E0766: the frozen matrix display lacks terminal punctuation.
这表明右边的两个矩阵都可逆，从而证明了引理。
% P01-E0767: invertibility of the displayed product alone does not imply both factors are invertible over arbitrary modules; the reverse product or explicit inverses are not supplied.
\end{proof}

\begin{lemma}
\label{algebra-lemma-conormal-module}
设 $R \to S$ 是有限型环映射。
对任意下列两个表现：
$$
\alpha : R[x_1, \ldots, x_n] \to S
$$
以及
$$
\beta : R[y_1, \ldots, y_m] \to S
$$
均有
% P01-E0768: the frozen statement has a spurious comma before “and” between the two presentations.
$$
I/I^2 \oplus S^{\oplus m} \cong J/J^2 \oplus S^{\oplus n}
$$
作为 $S$-模成立，其中 $I = \Ker(\alpha)$ 且 $J = \Ker(\beta)$。
% P01-E0769: the frozen phrase “as S-modules where” lacks a comma before the defining clause.
\end{lemma}

\begin{proof}
见引理 \ref{algebra-lemma-NL-homotopy} 和 \ref{algebra-lemma-sum-two-terms}。
\end{proof}

\begin{lemma}
\label{algebra-lemma-conormal-module-localize}
设 $R \to S$ 是有限型环映射。
设 $g \in S$。对任意表现
$\alpha : R[x_1, \ldots, x_n] \to S$ 和
$\beta : R[y_1, \ldots, y_m] \to S_g$，都有
% P01-E0770: the frozen statement has a spurious comma before “and” between the two presentations.
$$
(I/I^2)_g \oplus S^{\oplus m}_g \cong J/J^2 \oplus S_g^{\oplus n}
$$
作为 $S_g$-模成立，其中
$I = \Ker(\alpha)$ 且 $J = \Ker(\beta)$。
% P01-E0771: the frozen phrase “as S_g-modules where” lacks a comma before the defining clause.
\end{lemma}

\begin{proof}
令 $\beta' : R[x_1, \ldots, x_n, x] \to S_g$ 为从 $\alpha$
出发、按引理 \ref{algebra-lemma-principal-localization-NL} 构造的表现。
于是已知 $\NL(\alpha) \otimes_S S_g$ 与 $\NL(\beta')$ 同伦等价。
由引理 \ref{algebra-lemma-NL-homotopy}，$\NL(\beta)$ 与
$\NL(\beta')$ 同伦等价。因此 $\NL(\alpha) \otimes_S S_g$
与 $\NL(\beta)$ 同伦等价。最后应用引理
\ref{algebra-lemma-conormal-module}。
\end{proof}

\section{局部完全交}
\label{algebra-section-lci}

\noindent
作为局部完全交这一性质是 Noether 局部环的一个内蕴性质。
这将在《除幂代数》的第 \ref{dpa-section-lci} 节讨论。
不过眼下我们只对域上的有限型代数定义这一性质。

\begin{definition}
\label{algebra-definition-lci-field}
设 $k$ 是域。
设 $S$ 是有限型 $k$-代数。
\begin{enumerate}
\item 如果存在一个表现
$S = k[x_1, \ldots, x_n]/(f_1, \ldots, f_c)$，使得
$\dim(S) = n - c$，则称 $S$ 是一个\textit{$k$ 上的全局完全交}。
\item 如果存在覆盖 $\Spec(S) = \bigcup D(g_i)$，使得每个环
$S_{g_i}$ 都是 $k$ 上的全局完全交，则称 $S$ 是一个
\textit{$k$ 上的局部完全交}。
\end{enumerate}
我们还约定零环是 $k$ 上的全局完全交。
\end{definition}

\noindent
假设 $S$ 是如定义 \ref{algebra-definition-lci-field} 所述的全局完全交
$S = k[x_1, \ldots, x_n]/(f_1, \ldots, f_c)$。
对极大理想 $\mathfrak m \subset k[x_1, \ldots, x_n]$，
有 $\dim(k[x_1, \ldots, x_n]_\mathfrak m) = n$
（引理 \ref{algebra-lemma-dim-affine-space}）。
若 $(f_1, \ldots, f_c) \subset \mathfrak m$，则由引理
\ref{algebra-lemma-one-equation} 可知 $\dim(S_\mathfrak m) \geq n - c$。
% P01-E0772: the frozen paragraph indexes S by a maximal ideal of the polynomial ring without first replacing it by the corresponding maximal ideal of S.
由于定义 \ref{algebra-definition-lci-field} 给出 $\dim(S) = n - c$，
便得对 $S$ 的所有极大理想都有 $\dim(S_\mathfrak m) = n - c$，
而且 $\Spec(S)$ 是维数为 $n - c$ 的等维空间
（《拓扑》，定义 \ref{topology-definition-equidimensional}）；见引理
\ref{algebra-lemma-dimension-at-a-point-finite-type-over-field}。
下文会经常使用这一点而不再说明。

\begin{lemma}
\label{algebra-lemma-localize-lci}
设 $k$ 是域。
设 $S$ 是有限型 $k$-代数。
设 $g \in S$。
\begin{enumerate}
\item 若 $S$ 是全局完全交，则 $S_g$ 也是。
\item 若 $S$ 是局部完全交，则 $S_g$ 也是。
\end{enumerate}
\end{lemma}

\begin{proof}
第二个断言立即由第一个断言得出。现证明第一个断言。
% P01-E0773: the frozen “Proof of the first statement.” is a sentence fragment.
若 $S_g$ 是零环，断言成立。假设 $S_g$ 非零。
按定义 \ref{algebra-definition-lci-field}，写成
$S = k[x_1, \ldots, x_n]/(f_1, \ldots, f_c)$，其中
$n - c = \dim(S)$。由定义后的说明，$\dim(S_g) = n - c$。
令 $g' \in k[x_1, \ldots, x_n]$ 是剩余类对应于 $g$ 的一个元素。
则
$S_g =  k[x_1, \ldots, x_n, x_{n + 1}]/(f_1, \ldots, f_c, x_{n + 1}g' - 1)$，
正合所需。
\end{proof}

\begin{lemma}
\label{algebra-lemma-lci-CM}
设 $k$ 是域。设 $S$ 是有限型 $k$-代数。
若 $S$ 是局部完全交，则 $S$ 是 Cohen--Macaulay 环。
\end{lemma}

\begin{proof}
选取 $S$ 的一个极大素理想 $\mathfrak m$。
% P01-E0774: the frozen proof says “maximal prime” where the local criterion uses a maximal ideal.
我们须证明 $S_\mathfrak m$ 是 Cohen--Macaulay 的。
由假设，不妨设 $S = k[x_1, \ldots, x_n]/(f_1, \ldots, f_c)$，
其中 $\dim(S) = n - c$。令
$\mathfrak m' \subset k[x_1, \ldots, x_n]$ 为对应于
$\mathfrak m$ 的极大理想。由命题
\ref{algebra-proposition-finite-gl-dim-polynomial-ring}，局部环
$k[x_1, \ldots, x_n]_{\mathfrak m'}$ 是维数为 $n$ 的正则局部环。
特别地，由引理 \ref{algebra-lemma-regular-ring-CM}，它是 Cohen--Macaulay 的。
对引理 \ref{algebra-lemma-one-equation} 应用 $c$ 次可知，局部环
$S_{\mathfrak m} = k[x_1, \ldots, x_n]_{\mathfrak m'}/(f_1, \ldots, f_c)$
的维数 $\geq n - c$。由假设，$\dim(S_{\mathfrak m}) \leq n - c$。
故等号成立。这蕴含 $f_1, \ldots, f_c$ 在
$k[x_1, \ldots, x_n]_{\mathfrak m'}$ 中是正则序列，且
$S_{\mathfrak m}$ 是 Cohen--Macaulay 的；见命题
\ref{algebra-proposition-CM-module}。
\end{proof}

\noindent
下面的引理是本节其余内容的技术关键。它的一个重要特征是：
我们可以为环 $S$ 任取一个表现，而条件 (1) 不依赖于这一选择。

\begin{lemma}
\label{algebra-lemma-lci}
设 $k$ 是域。
设 $S$ 是有限型 $k$-代数。
设 $\mathfrak q$ 是 $S$ 的一个素理想。
任取一个表现 $S = k[x_1, \ldots, x_n]/I$。
令 $\mathfrak q'$ 为 $k[x_1, \ldots, x_n]$ 中对应于
$\mathfrak q$ 的素理想。置
$c = \text{高度}(\mathfrak q') - \text{高度}(\mathfrak q)$，
换言之，$\dim_{\mathfrak q}(S) = n - c$
（见引理 \ref{algebra-lemma-codimension}）。下列条件等价：
\begin{enumerate}
\item 存在 $g \in S$，$g \not \in \mathfrak q$，使得
$S_g$ 是 $k$ 上的全局完全交。
\item 理想 $I_{\mathfrak q'} \subset k[x_1, \ldots, x_n]_{\mathfrak q'}$
可由 $c$ 个元素生成。
\item 余法模 $(I/I^2)_{\mathfrak q}$ 作为
$S_{\mathfrak q}$-模可由 $c$ 个元素生成。
\item 余法模 $(I/I^2)_{\mathfrak q}$ 是秩为 $c$ 的自由
$S_{\mathfrak q}$-模。
\item 理想 $I_{\mathfrak q'}$ 在正则局部环
$k[x_1, \ldots, x_n]_{\mathfrak q'}$ 中可由一个正则序列生成。
\end{enumerate}
在此情形下，$I_{\mathfrak q'}$ 中任意 $c$ 个生成
$I_{\mathfrak q'}/\mathfrak q'I_{\mathfrak q'}$ 的元素，
都在局部环 $k[x_1, \ldots, x_n]_{\mathfrak q'}$ 中构成正则序列。
\end{lemma}

\begin{proof}
置 $R = k[x_1, \ldots, x_n]_{\mathfrak q'}$。这是维数为
$\text{高度}(\mathfrak q')$ 的 Cohen--Macaulay 局部环；
例如见引理 \ref{algebra-lemma-lci-CM}。此外，
$\overline{R} = R/IR = R/I_{\mathfrak q'} = S_{\mathfrak q}$
是维数为 $\text{高度}(\mathfrak q)$ 的商环。
取 $f_1, \ldots, f_c \in I_{\mathfrak q'}$ 为生成
$(I/I^2)_{\mathfrak q}$ 的元素。由引理 \ref{algebra-lemma-NAK} 可知，
$f_1, \ldots, f_c$ 生成 $I_{\mathfrak q'}$。
由维数恰好吻合，再由命题 \ref{algebra-proposition-CM-module} 可知
$f_1, \ldots, f_c$ 在 $R$ 中是正则序列。
由引理 \ref{algebra-lemma-regular-quasi-regular} 可知
$(I/I^2)_{\mathfrak q}$ 是自由的。
这些论证表明 (2)、(3)、(4) 等价，并且它们蕴含引理的最后一个断言，
从而也蕴含 (5)。

\medskip\noindent
若 (5) 成立，设 $I_{\mathfrak q'}$ 由长度为 $e$ 的一个正则序列生成，
则由维数理论
$$
\text{高度}(\mathfrak q) = \dim(S_{\mathfrak q}) =
\dim(k[x_1, \ldots, x_n]_{\mathfrak q'}) - e =
\text{高度}(\mathfrak q') - e
$$
（见第 \ref{algebra-section-dimension} 节）。因此 $e = c$，
故 (5) 蕴含 (2)。

\medskip\noindent
继续沿用第一段引入的记号。对每个 $f_i$，可找到
$d_i \in k[x_1, \ldots, x_n]$，$d_i \not \in \mathfrak q'$，
使得 $f_i' = d_i f_i \in k[x_1, \ldots, x_n]$。
仍有 $I_{\mathfrak q'} = (f_1', \ldots, f_c')R$。
因此存在 $g' \in k[x_1, \ldots, x_n]$，$g' \not \in \mathfrak q'$，
使得 $I_{g'} = (f_1', \ldots, f_c')$。
此外，选取 $g'' \in k[x_1, \ldots, x_n]$，
$g'' \not \in \mathfrak q'$，使得
$\dim(S_{g''}) = \dim_{\mathfrak q} \Spec(S)$。
由引理 \ref{algebra-lemma-codimension}，这一维数等于 $n - c$。
最后，令 $g$ 为 $g'g''$ 在 $S$ 中的像。于是
$$
S_g \cong k[x_1, \ldots, x_n, x_{n + 1}]
/
(f_1', \ldots, f_c', x_{n + 1}g'g'' - 1)
$$
并且由 $g''$ 的选取，此环的维数为 $n - c$。
故它是全局完全交。
因此 (2)、(3)、(4) 中的每一个都蕴含 (1)。

\medskip\noindent
假设 (1) 成立。令
$S_g \cong k[y_1, \ldots, y_m]/(f_1, \ldots, f_t)$
为 $S_g$ 作为全局完全交的一个表现。写
$J = (f_1, \ldots, f_t)$。令
$\mathfrak q'' \subset k[y_1, \ldots, y_m]$
为对应于 $\mathfrak qS_g$ 的素理想。注意
$t = m - \dim(S_g) =
\text{高度}(\mathfrak q'') - \text{高度}(\mathfrak q)$；
最后一个等号见引理 \ref{algebra-lemma-codimension}。
如引理 \ref{algebra-lemma-lci-CM} 的证明（以及上文）所示，
$f_1, \ldots, f_t$ 在局部环
$k[y_1, \ldots, y_m]_{\mathfrak q''}$ 中构成正则序列。
由引理 \ref{algebra-lemma-regular-quasi-regular} 可知
$(J/J^2)_{\mathfrak q}$ 是秩为 $t$ 的自由模。
由引理 \ref{algebra-lemma-conormal-module-localize}，有
$$
J/J^2 \oplus S_g^n \cong (I/I^2)_g \oplus S_g^m
$$
% P01-E0775: the frozen free-module summands use ordinary powers S_g^n and S_g^m instead of direct-sum notation.
因此 $(I/I^2)_{\mathfrak q}$ 是秩为
$t + n - m = m - \dim(S_g) + n - m = n - \dim(S_g) =
\text{高度}(\mathfrak q') - \text{高度}(\mathfrak q) = c$
的自由模。故得到 (4)。
\end{proof}

\noindent
引理 \ref{algebra-lemma-lci} 的结果提示我们作如下定义。

\begin{definition}
\label{algebra-definition-lci-local-ring}
设 $k$ 是域。设 $S$ 是在 $k$ 上本质有限型的局部 $k$-代数。
若存在一个局部 $k$-代数 $R$ 以及元素
$f_1, \ldots, f_c \in \mathfrak m_R$，使得
\begin{enumerate}
\item $R$ 在 $k$ 上本质有限型；
\item $R$ 是正则局部环；
\item $f_1, \ldots, f_c$ 在 $R$ 中构成正则序列；并且
\item 作为 $k$-代数，$S \cong R/(f_1, \ldots, f_c)$，
\end{enumerate}
则称 $S$ 是一个\textit{（$k$ 上的）完全交}。
\end{definition}

\noindent
由 Cohen 结构定理（见定理 \ref{algebra-theorem-cohen-structure-theorem}），
任意完备 Noether 局部环都可以写成某个正则完备局部环的商。
因此可以利用上述定义，为任意完备 Noether 局部环定义完全交环的概念。
我们将在《除幂代数》的第 \ref{dpa-section-lci} 节讨论这一点。
与此同时，下面的引理表明这样的定义是有意义的。

\begin{lemma}
\label{algebra-lemma-ci-well-defined}
设 $A \to B \to C$ 是满的局部环同态。
假设 $A$ 和 $B$ 是正则局部环。下列条件等价：
% P01-E0776: the frozen sentence introducing the equivalence list lacks a colon.
\begin{enumerate}
\item $\Ker(A \to C)$ 由一个正则序列生成；
\item $\Ker(A \to C)$ 可由 $\dim(A) - \dim(C)$ 个元素生成；
\item $\Ker(B \to C)$ 由一个正则序列生成；并且
\item $\Ker(B \to C)$ 可由 $\dim(B) - \dim(C)$ 个元素生成。
\end{enumerate}
\end{lemma}

\begin{proof}
正则局部环是 Cohen--Macaulay 的；见引理
\ref{algebra-lemma-regular-ring-CM}。因此 (1) $\Leftrightarrow$ (2) 以及
(3) $\Leftrightarrow$ (4)；见命题 \ref{algebra-proposition-CM-module}。
% P01-E0777: the frozen “Hence the equivalences ...” clause has no finite verb.
由引理 \ref{algebra-lemma-regular-quotient-regular}，理想
$\Ker(A \to B)$ 可由 $\dim(A) - \dim(B)$ 个元素生成。
故 (4) 蕴含 (2)。

\medskip\noindent
还须证明 (1) 蕴含 (4)。我们对 $\dim(A) - \dim(B)$ 作归纳。
$\dim(A) - \dim(B) = 0$ 的情形是平凡的。
假设 $\dim(A) > \dim (B)$。
写 $I = \Ker(A \to C)$ 且 $J = \Ker(A \to B)$。
注意 $J \subset I$。我们的假设是 $I$ 的最小生成元数为
$\dim(A) - \dim(C)$。令 $\mathfrak m \subset A$ 为极大理想。
考虑映射
$$
J/ \mathfrak m J \to  I / \mathfrak m I \to \mathfrak m /\mathfrak m^2
$$
% P01-E0778: the frozen display lacks terminal punctuation before a new sentence.
由引理 \ref{algebra-lemma-regular-quotient-regular} 及其证明，上述复合是单射。
任取 $x \in J$，使其在 $J /\mathfrak mJ$ 中非零。
由上述结论和 Nakayama 引理，$x$ 是 $I$ 的某个最小生成集中的元素。
因此可以把 $A$ 替换为 $A/xA$，把 $I$ 替换为 $I/xA$；
这使 $\dim(A)$ 以及 $I$ 的最小生成元数都减少 $1$。故得证。
\end{proof}

\begin{lemma}
\label{algebra-lemma-lci-local}
设 $k$ 是域。设 $S$ 是在 $k$ 上本质有限型的局部 $k$-代数。
下列条件等价：
\begin{enumerate}
\item $S$ 是 $k$ 上的完全交；
\item 对任意满射 $R \to S$，若 $R$ 是在 $k$ 上本质有限表示的
正则局部环，则理想 $\Ker(R \to S)$ 可由一个正则序列生成；
\item 对某个满射 $R \to S$，若 $R$ 是在 $k$ 上本质有限表示的
正则局部环，则理想 $\Ker(R \to S)$ 可由
$\dim(R) - \dim(S)$ 个元素生成；
\item 存在 $k$ 上的全局完全交 $A$ 以及 $A$ 的一个素理想
$\mathfrak a$，使得 $S \cong A_{\mathfrak a}$；并且
\item 存在 $k$ 上的局部完全交 $A$ 以及 $A$ 的一个素理想
$\mathfrak a$，使得 $S \cong A_{\mathfrak a}$。
\end{enumerate}
\end{lemma}

\begin{proof}
显然 (2) 蕴含 (1)，而 (1) 蕴含 (3)。同样显然 (4) 蕴含 (5)。
我们来证明 (3) 蕴含 (4)。因此假设存在满射
$R \to S$，其中 $R$ 是在 $k$ 上本质有限表示的正则局部环，
并且理想 $\Ker(R \to S)$ 可由
$\dim(R) - \dim(S)$ 个元素生成。
可写成 $R = (k[x_1, \ldots, x_n]/J)_{\mathfrak q}$，
其中 $J \subset k[x_1, \ldots, x_n]$，而
$\mathfrak q \subset k[x_1, \ldots, x_n]$ 是满足
$J \subset \mathfrak q$ 的某个素理想。令
$I \subset k[x_1, \ldots, x_n]$ 为映射
$k[x_1, \ldots, x_n] \to S$ 的核，从而
$S \cong (k[x_1, \ldots, x_n]/I)_{\mathfrak q}$。
由假设，$(I/J)_{\mathfrak q}$ 可由
$\dim(R) - \dim(S)$ 个元素生成。引理
\ref{algebra-lemma-ci-well-defined} 蕴含 $I_{\mathfrak q}$ 可由
$\dim(k[x_1, \ldots, x_n]_{\mathfrak q}) - \dim(S)$
个元素生成。由引理 \ref{algebra-lemma-lci} 可知，存在
$g \in k[x_1, \ldots, x_n]$，$g \not \in \mathfrak q$，使得代数
$(k[x_1, \ldots, x_n]/I)_g$ 是全局完全交，而 $S$
同构于它的一个局部环。

\medskip\noindent
为了完成引理的证明，还须证明 (5) 蕴含 (2)。
假设 (5)，并令 $\pi : R \to S$ 为满射，其中 $R$
是在 $k$ 上本质有限型的正则局部 $k$-代数。
由假设，对 $k$ 上的某个局部完全交 $A$，有
$S = A_{\mathfrak a}$。选取表现
$R = (k[y_1, \ldots, y_m]/J)_{\mathfrak q}$，其中
$J \subset \mathfrak q \subset k[y_1, \ldots, y_m]$。
可以且确实假设 $J$ 是映射 $k[y_1, \ldots, y_m] \to R$ 的核。
令 $I \subset k[y_1, \ldots, y_m]$ 为映射
$k[y_1, \ldots, y_m] \to S = A_{\mathfrak a}$ 的核。
则 $J \subset I$，而 $(I/J)_{\mathfrak q}$ 是满射
$\pi : R \to S$ 的核。因此
$S = (k[y_1, \ldots, y_m]/I)_{\mathfrak q}$。

\medskip\noindent
由引理 \ref{algebra-lemma-isomorphic-local-rings} 可知，存在
$g \in A$，$g \not \in \mathfrak a$，以及
$g' \in k[y_1, \ldots, y_m]$，$g' \not \in \mathfrak q$，
使得 $A_g \cong (k[y_1, \ldots, y_m]/I)_{g'}$。
把 $A$ 替换为 $A_g$，把 $k[y_1, \ldots, y_m]$ 替换为
$k[y_1, \ldots, y_{m + 1}]$ 后，可以假设
$A \cong k[y_1, \ldots, y_m]/I$。考虑局部环的满映射
$$
k[y_1, \ldots, y_m]_{\mathfrak q} \to R \to S.
$$
我们须证明 $R \to S$ 的核由一个正则序列生成。
由引理 \ref{algebra-lemma-lci} 可知，
$k[y_1, \ldots, y_m]_{\mathfrak q} \to A_{\mathfrak a} = S$
具有这一性质（因为 $A$ 是 $k$ 上的局部完全交）。
再由引理 \ref{algebra-lemma-ci-well-defined} 即得证。
\end{proof}

\begin{lemma}
\label{algebra-lemma-lci-at-prime}
设 $k$ 是域。设 $S$ 是有限型 $k$-代数。
设 $\mathfrak q$ 是 $S$ 的一个素理想。下列条件等价：
\begin{enumerate}
\item 局部环 $S_{\mathfrak q}$ 是完全交环
（定义 \ref{algebra-definition-lci-local-ring}）。
\item 存在 $g \in S$，$g \not \in \mathfrak q$，使得
$S_g$ 是 $k$ 上的局部完全交。
\item 存在 $g \in S$，$g \not \in \mathfrak q$，使得
$S_g$ 是 $k$ 上的全局完全交。
\item 对任意表现 $S = k[x_1, \ldots, x_n]/I$，若
$\mathfrak q' \subset k[x_1, \ldots, x_n]$ 对应于 $\mathfrak q$，
则引理 \ref{algebra-lemma-lci} 的等价条件 (1) -- (5) 中的任一个成立。
% P01-E0779: the frozen plural phrase “any of the equivalent conditions ... hold” has number disagreement and is logically clearer as “any one ... holds”.
\end{enumerate}
\end{lemma}

\begin{proof}
这是引理 \ref{algebra-lemma-lci}、\ref{algebra-lemma-lci-local} 与各定义的结合。
\end{proof}

\begin{lemma}
\label{algebra-lemma-lci-global}
设 $k$ 是域。设 $S$ 是有限型 $k$-代数。
下列条件等价：
\begin{enumerate}
\item 环 $S$ 是 $k$ 上的局部完全交。
\item $S$ 的所有局部环都是 $k$ 上的完全交环。
\item $S$ 在极大理想处的所有局部化都是 $k$ 上的完全交环。
\end{enumerate}
\end{lemma}

\begin{proof}
这由引理 \ref{algebra-lemma-lci-at-prime}、$\Spec(S)$ 的拟紧性以及各定义得出。
\end{proof}

\noindent
下面的引理说明，完全交这一性质在基域变换下保持（而且是在强意义下）。

\begin{lemma}
\label{algebra-lemma-lci-field-change-local}
设 $K/k$ 是域扩张。
设 $S$ 是 $k$ 上的有限型代数。
令 $\mathfrak q_K$ 为 $S_K = K \otimes_k S$ 的一个素理想，
并令 $\mathfrak q$ 为 $S$ 中对应的素理想。
则 $S_{\mathfrak q}$ 是 $k$ 上的完全交
（定义 \ref{algebra-definition-lci-local-ring}），当且仅当
$(S_K)_{\mathfrak q_K}$ 是 $K$ 上的完全交。
\end{lemma}

\begin{proof}
选取表现 $S = k[x_1, \ldots, x_n]/I$。
这给出表现 $S_K = K[x_1, \ldots, x_n]/I_K$，其中
$I_K = K \otimes_k I$。令
$\mathfrak q_K' \subset K[x_1, \ldots, x_n]$ 与
$\mathfrak q' \subset k[x_1, \ldots, x_n]$ 分别为对应的素理想。
我们将证明：引理 \ref{algebra-lemma-lci} 的等价条件对二元组
$(S = k[x_1, \ldots, x_n]/I, \mathfrak q)$ 成立，当且仅当它们对二元组
$(S_K = K[x_1, \ldots, x_n]/I_K, \mathfrak q_K)$ 成立。
由此即可得出引理（见引理 \ref{algebra-lemma-lci-at-prime}）。

\medskip\noindent
由引理 \ref{algebra-lemma-dimension-at-a-point-preserved-field-extension}，有
$\dim_{\mathfrak q} S = \dim_{\mathfrak q_K} S_K$。
因此，引理 \ref{algebra-lemma-lci} 中出现的整数 $c$ 对二元组
$(S = k[x_1, \ldots, x_n]/I, \mathfrak q)$ 与二元组
$(S_K = K[x_1, \ldots, x_n]/I_K, \mathfrak q_K)$ 相同。
另一方面，
\begin{eqnarray*}
I \otimes_{k[x_1, \ldots, x_n]} \kappa(\mathfrak q')
\otimes_{\kappa(\mathfrak q')} \kappa(\mathfrak q_K')
& = &
I \otimes_{k[x_1, \ldots, x_n]} \kappa(\mathfrak q_K') \\
& = &
I \otimes_{k[x_1, \ldots, x_n]} K[x_1, \ldots, x_n]
\otimes_{K[x_1, \ldots, x_n]} \kappa(\mathfrak q_K') \\
& = &
(K \otimes_k I) \otimes_{K[x_1, \ldots, x_n]} \kappa(\mathfrak q_K') \\
& = &
I_K \otimes_{K[x_1, \ldots, x_n]} \kappa(\mathfrak q'_K).
\end{eqnarray*}
因此，
$\dim_{\kappa(\mathfrak q')}
I \otimes_{k[x_1, \ldots, x_n]} \kappa(\mathfrak q')
=
\dim_{\kappa(\mathfrak q'_K)}
I_K \otimes_{K[x_1, \ldots, x_n]} \kappa(\mathfrak q_K')$。
% P01-E0780: the frozen display alternates between q_K' and q'_K for the same extended prime.
于是由 Nakayama 引理 \ref{algebra-lemma-NAK}，
$I_{\mathfrak q'}$ 的最小生成元数等于
$(I_K)_{\mathfrak q'_K}$ 的最小生成元数。
因此，引理由引理 \ref{algebra-lemma-lci} 的刻画 (2) 得出。
\end{proof}

\begin{lemma}
\label{algebra-lemma-lci-field-change}
设 $k \to K$ 是域扩张。
设 $S$ 是有限型 $k$-代数。
则 $S$ 是 $k$ 上的局部完全交，当且仅当
$S \otimes_k K$ 是 $K$ 上的局部完全交。
\end{lemma}

\begin{proof}
这由引理 \ref{algebra-lemma-lci-global} 与
\ref{algebra-lemma-lci-field-change-local} 结合得出。
不过，我们也给出另一个证明（它基于同样的原理）。

\medskip\noindent
置 $S' = S \otimes_k K$。令 $\alpha : k[x_1, \ldots, x_n] \to S$
是核为 $I$ 的一个表现。令
$\alpha' : K[x_1, \ldots, x_n] \to S'$ 是所诱导的表现，其核为 $I'$。

\medskip\noindent
假设 $S$ 是局部完全交。取素理想 $\mathfrak q \subset S'$。
% P01-E0781: the frozen designation sentence omits “by/as” after “Denote”.
分别用 $\mathfrak q'$、$\mathfrak p$ 和 $\mathfrak p'$ 表示
$K[x_1, \ldots, x_n]$ 中对应的素理想、$S$ 中对应的素理想以及
$k[x_1, \ldots, x_n]$ 中对应的素理想。考虑如下 Noether 局部环的图
$$
\xymatrix{
S'_{\mathfrak q} &  K[x_1, \ldots, x_n]_{\mathfrak q'} \ar[l] \\
S_{\mathfrak p}\ar[u] &  k[x_1, \ldots, x_n]_{\mathfrak p'} \ar[u] \ar[l]
}
$$
% P01-E0782: the frozen diagram lacks terminal punctuation before the next sentence.
由引理 \ref{algebra-lemma-lci}，我们知道 $S_{\mathfrak p}$ 在
$k[x_1, \ldots, x_n]_{\mathfrak p'}$ 中由某个正则序列
$f_1, \ldots, f_c$ 截出。由于右边的竖直箭头是平坦的，
$f_1, \ldots, f_c$ 的像在
$K[x_1, \ldots, x_n]_{\mathfrak q'}$ 中构成正则序列。
由于在 $k$ 上与 $K$ 作张量是正合函子，我们有
$S'_{\mathfrak q} = K[x_1, \ldots, x_n]_{\mathfrak q'}/(f_1, \ldots, f_c)$。
故再次由引理 \ref{algebra-lemma-lci} 可知，$S'$ 在 $\mathfrak q$ 的一个邻域内
是局部完全交。由于 $\mathfrak q$ 是任意的，$S'$ 是 $K$ 上的局部完全交。

\medskip\noindent
反过来，假设 $S'$ 是局部完全交。取 $S$ 的一个极大理想
$\mathfrak m$。令 $\mathfrak m'$ 表示 $k[x_1, \ldots, x_n]$
中对应的极大理想。
% P01-E0783: the frozen residue-field designation omits “by/as”.
记 $\kappa = \kappa(\mathfrak m)$ 为剩余域。
由备注 \ref{algebra-remark-fundamental-diagram}，$S'$ 中位于
$\mathfrak m$ 上方的素理想对应于 $K \otimes_k \kappa$ 中的素理想。
由 Hilbert 零点定理 \ref{algebra-theorem-nullstellensatz}，有
$[\kappa : k] < \infty$。因此 $K \otimes_k \kappa$
是 $K$ 上非零的有限代数。于是 $K \otimes_k \kappa$ 有有限个
（其个数 $> 0$）素理想；它们全是极大理想，并且各自的剩余域都是 $K$ 的有限扩张
（见第 \ref{algebra-section-artinian} 节）。
% P01-E0784: the frozen prose twice says “a finite number > 0” of primes.
因此，位于 $\mathfrak m$ 上方的素理想
$\mathfrak n \subset S'$ 有有限个（其个数 $> 0$）；它们全是极大理想，并且各自的剩余域
都是 $K$ 的有限扩张。取其中一个，记为 $\mathfrak n \subset S'$，
并令 $\mathfrak n' \subset K[x_1, \ldots, x_n]$ 表示
$K[x_1, \ldots, x_n]$ 中对应的素理想。
注意，由于 $V(\mathfrak mS')$ 是有限的，$\mathfrak n$
是其中的一个孤立闭点，从而 $\mathfrak mS'_{\mathfrak n}$
是 $S'_{\mathfrak n}$ 的一个定义理想。例如，由引理
\ref{algebra-lemma-dimension-base-fibre-equals-total}，这蕴含
$\dim(S_{\mathfrak m}) = \dim(S'_{\mathfrak n})$。
（也可以利用引理
\ref{algebra-lemma-dimension-at-a-point-preserved-field-extension} 看出这一点。）
考虑相应的 Noether 局部环之图
$$
\xymatrix{
S'_{\mathfrak n} &  K[x_1, \ldots, x_n]_{\mathfrak n'} \ar[l] \\
S_{\mathfrak m}\ar[u] &  k[x_1, \ldots, x_n]_{\mathfrak m'} \ar[u] \ar[l]
}
$$
% P01-E0785: the frozen diagram lacks terminal punctuation before the next sentence.
由引理 \ref{algebra-lemma-change-base-NL}，有
$\NL(\alpha) \otimes_S S' = \NL(\alpha')$，特别地，
$I'/(I')^2 = I/I^2 \otimes_S S'$。因此
$(I/I^2)_{\mathfrak m} \otimes_{S_{\mathfrak m}} \kappa$
与
$(I'/(I')^2)_{\mathfrak n} \otimes_{S'_{\mathfrak n}} \kappa(\mathfrak n)$
具有相同的维数。由于 $(I'/(I')^2)_{\mathfrak n}$
是秩为 $n - \dim S'_{\mathfrak n}$ 的自由模，我们推出
$(I/I^2)_{\mathfrak m}$ 可由
$n - \dim S'_{\mathfrak n} = n - \dim S_{\mathfrak m}$ 个元素生成。
由引理 \ref{algebra-lemma-lci} 可知，$S$ 在 $\mathfrak m$ 的一个邻域内
是局部完全交。由于 $\mathfrak m$ 是任意极大理想，我们得出
$S$ 是局部完全交。
\end{proof}

\noindent
最后给出一个引理。稍后我们将用它证明：给定环映射
$T \to A \to B$，若 $B$ 在 $T$ 上是 syntomic 的，且
$B$ 在 $A$ 上是 syntomic 的，则 $A$ 在 $T$ 上是 syntomic 的。

\begin{lemma}
\label{algebra-lemma-lci-permanence-initial}
设
$$
\xymatrix{
B & S \ar[l] \\
A \ar[u] & R \ar[l] \ar[u]
}
$$
% P01-E0786: the frozen diagram lacks terminal punctuation before “be a commutative square”.
是局部环的交换方块。假设
\begin{enumerate}
\item $R$ 和 $\overline{S} = S/\mathfrak m_R S$ 是正则局部环，
\item 对某些理想 $I$、$J$，有 $A = R/I$ 和 $B = S/J$，
\item $J \subset S$ 以及
% P01-E0787: the frozen quotient omits parentheses and likely the extension of m_R to S.
$\overline{J} = J/\mathfrak m_R \cap J \subset \overline{S}$
由正则序列生成，并且
\item $A \to B$ 和 $R \to S$ 是平坦的。
\end{enumerate}
则 $I$ 由一个正则序列生成。
\end{lemma}

\begin{proof}
置 $\overline{B} = B/\mathfrak m_RB = B/\mathfrak m_AB$，于是
$\overline{B} = \overline{S}/\overline{J}$。
取元素 $f_1, \ldots, f_{\overline{c}} \in J$，使得
$\overline{f}_1, \ldots, \overline{f}_{\overline{c}} \in \overline{J}$
构成生成 $\overline{J}$ 的正则序列。注意
$\overline{c} = \dim(\overline{S}) - \dim(\overline{B})$；
见引理 \ref{algebra-lemma-ci-well-defined}。
由引理 \ref{algebra-lemma-grothendieck-regular-sequence}，环
$S/(f_1, \ldots, f_{\overline{c}})$ 在 $R$ 上平坦。
% P01-E0788: the frozen quotient by the sum lacks parentheses in both occurrences.
因此 $S/(f_1, \ldots, f_{\overline{c}}) + IS$ 在 $A$ 上平坦。
映射 $S/(f_1, \ldots, f_{\overline{c}}) + IS \to B$
于是是有限 $S/IS$-模之间的满射，这些模在 $A$ 上平坦；它模去
$\mathfrak m_A$ 后是同构，故由引理 \ref{algebra-lemma-mod-injective}
可知它是同构。换言之，
$J = (f_1, \ldots, f_{\overline{c}}) + IS$。

\medskip\noindent
再次由引理 \ref{algebra-lemma-ci-well-defined}，理想 $J$ 由一个长度为
$c = \dim(S) - \dim(B)$ 的正则序列生成。因此
$J/\mathfrak m_SJ$ 是维数为 $c$ 的向量空间。
由上述对 $J$ 的描述，存在
$g_1, \ldots, g_{c - \overline{c}} \in I$，使得 $J$ 由
$f_1, \ldots, f_{\overline{c}}, g_1, \ldots, g_{c - \overline{c}}$
生成（使用 Nakayama 引理 \ref{algebra-lemma-NAK}）。考虑环
$A' = R/(g_1, \ldots, g_{c - \overline{c}})$ 以及满射
$A' \to A$。由上述结论可知，
$B = S/(f_1, \ldots, f_{\overline{c}}, g_1, \ldots, g_{c - \overline{c}})$
在 $A'$ 上平坦（因为 $S/(f_1, \ldots, f_{\overline{c}})$
在 $R$ 上平坦）。因此 $A' \to B$ 是单射（因为它忠实平坦；
见引理 \ref{algebra-lemma-local-flat-ff}）。
由于该映射经由 $A$ 分解，我们得到 $A' = A$。
注意 $\dim(B) = \dim(A) + \dim(\overline{B})$，且
$\dim(S) = \dim(R) + \dim(\overline{S})$；见引理
\ref{algebra-lemma-dimension-base-fibre-equals-total}。
% P01-E0789: the frozen dimension difference omits spacing before the second dim command.
由初等代数可知 $c - \overline{c} = \dim(R) -\dim(A)$。
故由引理 \ref{algebra-lemma-ci-well-defined}，
$I = (g_1, \ldots, g_{c - \overline{c}})$ 由一个正则序列生成。
\end{proof}

\section{syntomic 态射}
\label{algebra-section-syntomic}

\noindent
syntomic 环映射是平坦、有限表示，并且所有纤维都是局部完全交的环映射。
一般的局部完全交环映射将在《代数进阶》的第
\ref{more-algebra-section-lci} 节中讨论。

\begin{definition}
\label{algebra-definition-lci}
若环映射 $R \to S$ 是平坦且有限表示的，并且它的所有纤维环
$S \otimes_R \kappa(\mathfrak p)$ 都是局部完全交（见定义
\ref{algebra-definition-lci-field}），则称它是 {\it syntomic} 的，
也称 $S$ 是 $R$ 上的一个{\it 平坦局部完全交}。
\end{definition}

\noindent
显然，域上的代数在该域上是 syntomic 的，当且仅当它是局部完全交。
这个定义具有如下令人满意的性质。

\begin{lemma}
\label{algebra-lemma-syntomic-descends}
\begin{slogan}
syntomic 性质在基底上是 fpqc 局部的。
\end{slogan}
设 $R \to S$ 是环映射。
设 $R \to R'$ 是忠实平坦环映射。
置 $S' = R'\otimes_R S$。
则 $R \to S$ 是 syntomic 的，当且仅当 $R' \to S'$ 是 syntomic 的。
\end{lemma}

\begin{proof}
由引理 \ref{algebra-lemma-finite-presentation-descends} 和
\ref{algebra-lemma-flatness-descends}，结论对平坦性和有限表示性成立。
映射 $\Spec(R') \to \Spec(R)$ 是满射；见引理
\ref{algebra-lemma-ff-rings}。因此，只须证明：给定位于
$\mathfrak p \subset R$ 上方的素理想 $\mathfrak p' \subset R'$，
% P01-E0790: the frozen “it suffices to show given primes ... that” construction is grammatically malformed.
$S \otimes_R \kappa(\mathfrak p)$ 是局部完全交，当且仅当
$S' \otimes_{R'} \kappa(\mathfrak p')$ 是局部完全交。注意
$S' \otimes_{R'} \kappa(\mathfrak p') =
S \otimes_R \kappa(\mathfrak p)
\otimes_{\kappa(\mathfrak p)} \kappa(\mathfrak p')$。
于是可应用引理 \ref{algebra-lemma-lci-field-change}。
\end{proof}

\begin{lemma}
\label{algebra-lemma-base-change-syntomic}
syntomic 映射的任意基变换仍是 syntomic 映射。
\end{lemma}

\begin{proof}
平坦性、有限表示性以及纤维为局部完全交这三个性质分别由引理
\ref{algebra-lemma-flat-base-change}、\ref{algebra-lemma-compose-finite-type} 和
\ref{algebra-lemma-lci-field-change} 得到。
\end{proof}

\begin{lemma}
\label{algebra-lemma-local-syntomic}
设 $R \to S$ 是环映射。
假设有生成单位理想的 $g_1, \ldots g_m \in S$，并且每个
% P01-E0791: the frozen mathematical list omits the comma after \ldots.
$R \to S_{g_i}$ 都是 syntomic 的。
则 $R \to S$ 是 syntomic 的。
\end{lemma}

\begin{proof}
由引理 \ref{algebra-lemma-flat-localization} 和 \ref{algebra-lemma-cover-upstairs}，
这对平坦性和有限表示性成立。纤维环为局部完全交这一性质按照其定义
在 $S$ 上是局部的；见定义 \ref{algebra-definition-lci-field}。
\end{proof}

\begin{definition}
\label{algebra-definition-relative-global-complete-intersection}
设 $R \to S$ 是环映射。若存在表现
$S = R[x_1, \ldots, x_n]/(f_1, \ldots, f_c)$，并且
$\Spec(S) \to \Spec(R)$ 的每个非空纤维都具有维数 $n - c$，
则称 $R \to S$ 是一个{\it 相对全局完全交}。
我们将用“令
$S = R[x_1, \ldots, x_n]/(f_1, \ldots, f_c)$
为一个相对全局完全交”来表示这种情形。
\end{definition}

\noindent
下面的引理有时可用于寻找全局表现。

\begin{lemma}
\label{algebra-lemma-huber}
设 $S$ 是有限表示 $R$-代数，并且它有一个表现
$S = R[x_1, \ldots, x_n]/I$，使得 $I/I^2$ 是自由 $S$-模。
则 $S$ 有一个表现
$S = R[y_1, \ldots, y_m]/(f_1, \ldots, f_c)$，使得
$(f_1, \ldots, f_c)/(f_1, \ldots, f_c)^2$ 是自由模，且
$f_1, \ldots, f_c$ 的类给出其一组基。
\end{lemma}

\begin{proof}
由引理 \ref{algebra-lemma-finite-presentation-independent}，$I$ 是有限生成理想。
取映到 $I/I^2$ 的一组基的元素 $f_1, \ldots, f_c \in I$。
由 Nakayama 引理（引理 \ref{algebra-lemma-NAK}），存在
$g \in 1 + I$，使得
$$
g \cdot I \subset (f_1, \ldots, f_c)
$$
并且 $I_g \cong (f_1, \ldots, f_c)_g$。因此，如所要求的，
$$
S \cong R[x_1, \ldots, x_n]/(f_1, \ldots, f_c)[1/g]
\cong R[x_1, \ldots, x_n, x_{n + 1}]/(f_1, \ldots, f_c, gx_{n + 1} - 1)
$$
进而可知 $f_1, \ldots, f_c,gx_{n + 1} - 1$ 构成
% P01-E0792: the frozen source omits a space after the comma before gx_{n+1}-1.
$(f_1, \ldots, f_c, gx_{n + 1} - 1)/(f_1, \ldots, f_c, gx_{n + 1} - 1)^2$
的一组基。例如，这可由应用引理
\ref{algebra-lemma-principal-localization-NL} 得出。
\end{proof}



\begin{example}
\label{algebra-example-factor-polynomials}
设 $n , m \geq 1$ 是整数。考虑环映射
\begin{eqnarray*}
R = \mathbf{Z}[a_1, \ldots, a_{n + m}]
& \longrightarrow &
S = \mathbf{Z}[b_1, \ldots, b_n, c_1, \ldots, c_m] \\
a_1 & \longmapsto & b_1 + c_1 \\
a_2 & \longmapsto & b_2 + b_1 c_1 + c_2 \\
\ldots & \ldots & \ldots \\
a_{n + m} & \longmapsto & b_n c_m
\end{eqnarray*}
换言之，这是如上所示的多项式环之间唯一的环映射，使得多项式分解
$$
x^{n + m} + a_1 x^{n + m - 1} + \ldots + a_{n + m}
=
(x^n + b_1 x^{n - 1} + \ldots + b_n)
(x^m + c_1 x^{m - 1} + \ldots + c_m)
$$
成立。注意 $S$ 在 $R$ 上由 $n + m$ 个元素（即 $b_i, c_j$）
生成，并且有 $n + m$ 个方程（即 $a_k = a_k(b_i, c_j)$）。
为了证明 $S$ 是 $R$ 上的相对全局完全交，只须证明所有纤维的维数为 $0$。

\medskip\noindent
为此，设 $R \to k$ 是到域 $k$ 的环映射。
设 $a_i$ 映到 $\alpha_i \in k$。
考虑纤维环 $S_k = k \otimes_R S$。设 $k \to K$ 是域扩张。
一个 $k$-代数映射 $S_k \to K$ 等价于找到
% P01-E0793: the frozen prose says “a map of S_k to K” instead of “from S_k to K”.
$\beta_1, \ldots, \beta_n, \gamma_1, \ldots, \gamma_m \in K$，使得
$$
x^{n + m} + \alpha_1 x^{n + m - 1} + \ldots + \alpha_{n + m}
=
(x^n + \beta_1 x^{n - 1} + \ldots + \beta_n)
(x^m + \gamma_1 x^{m - 1} + \ldots + \gamma_m).
$$
因此，在 $K$ 中这样的 $n + m$-元组只有有限种选择。
% P01-E0794: the frozen “at most finitely many” is redundant.
这证明所有纤维都只有有限个闭点（例如，使用 Hilbert 零点定理可知，
它们都对应于 $\overline{k}$ 中的解），因而 $R \to S$
是相对全局完全交。

\medskip\noindent
另一种论证方式是证明
$\mathbf{Z}[a_1, \ldots, a_{n + m}] \to
\mathbf{Z}[b_1, \ldots, b_n, c_1, \ldots, c_m]$
实际上也是一个{\it 有限}环映射。事实上，由引理
\ref{algebra-lemma-polynomials-divide}，每个 $b_i, c_j$ 在 $R$ 上都是整的，
故由引理 \ref{algebra-lemma-characterize-integral}，$R \to S$ 是有限的。
\end{example}

\begin{example}
\label{algebra-example-roots-universal-polynomial}
考虑环映射
\begin{eqnarray*}
R = \mathbf{Z}[a_1, \ldots, a_n]
& \longrightarrow &
S = \mathbf{Z}[\alpha_1, \ldots, \alpha_n] \\
a_1 & \longmapsto &
\alpha_1 + \ldots + \alpha_n \\
\ldots & \ldots & \ldots \\
a_n & \longmapsto & \alpha_1 \ldots \alpha_n
\end{eqnarray*}
换言之，这是如上所示的多项式环之间唯一的环映射，使得
$$
x^n + a_1 x^{n - 1} + \ldots + a_n
=
\prod\nolimits_{i = 1}^n (x + \alpha_i)
$$
在 $\mathbf{Z}[\alpha_i, x]$ 中成立。也就是说，$a_i$ 映到
$\alpha_1, \ldots, \alpha_n$ 的第 $i$ 个初等对称函数。
% P01-E0795: the frozen English typesets “i” and “th” without a hyphen.
由初等对称多项式的通常理论（略去细节），$S$ 是有限自由 $R$-模，
它的一组基由满足 $0 \leq e_i \leq n - i$ 的元素
$\alpha_1^{e_1} \alpha_2^{e_2} \ldots \alpha_n^{e_n}$ 给出。
尤其，$R \to S$ 的纤维环是有限的，因而具有维数 $0$。
另一方面，$S$ 在 $R$ 上由 $n$ 个元素和 $n$ 个方程生成。
因此 $S$ 是 $R$ 上的相对全局完全交。
由于 $S$ 在 $R$ 上的秩为正，还可看出 $S$ 在 $R$ 上忠实平坦。
\end{example}

\begin{lemma}
\label{algebra-lemma-base-change-relative-global-complete-intersection}
设 $S = R[x_1, \ldots, x_n]/(f_1, \ldots, f_c)$
是一个相对全局完全交（定义
\ref{algebra-definition-relative-global-complete-intersection}）。
% P01-E0796: the frozen statement lacks punctuation before the enumeration.
\begin{enumerate}
\item 对任意 $R \to R'$，基变换
$R' \otimes_R S = R'[x_1, \ldots, x_n]/(f_1, \ldots, f_c)$
是相对全局完全交。
\item 对任意 $g \in S$，若它是 $h \in R[x_1, \ldots, x_n]$ 的像，
则环
$$
S_g = R[x_1, \ldots, x_n, x_{n + 1}]/(f_1, \ldots, f_c, hx_{n + 1} - 1)
$$
是相对全局完全交。
\item 若对某个 $f \in R$，$R \to S$ 分解为 $R \to R_f \to S$，
% P01-E0797: the frozen source incorrectly separates the if-clause and “Then” with a period.
则环 $S = R_f[x_1, \ldots, x_n]/(f_1, \ldots, f_c)$
是 $R_f$ 上的相对全局完全交。
\end{enumerate}
\end{lemma}

\begin{proof}
由引理 \ref{algebra-lemma-dimension-preserved-field-extension}，基变换的纤维
与原映射的纤维具有相同的维数。此外，下式成立：
$$
R' \otimes_R R[x_1, \ldots, x_n]/(f_1, \ldots, f_c)
= R'[x_1, \ldots, x_n]/(f_1, \ldots, f_c)
$$
故 (1) 成立。
(2) 的证明在于，关于一个元素的局部化可以写成
$S_g \cong S[x_{n + 1}]/(gx_{n + 1} - 1)$。
断言 (3) 由 (1) 得出，因为在 (3) 的假设下有
$R_f \otimes_R S \cong S$。
\end{proof}

\begin{lemma}
\label{algebra-lemma-localize-relative-complete-intersection}
设 $R$ 是环。设 $S = R[x_1, \ldots, x_n]/(f_1, \ldots, f_c)$。
在以下每种情形中，我们都将找到
$h \in R[x_1, \ldots, x_n]$，其像为 $g \in S$，使得
$$
S_g = R[x_1, \ldots, x_n, x_{n + 1}]/(f_1, \ldots, f_c, hx_{n + 1} - 1)
$$
是相对全局完全交，并且具有定义
\ref{algebra-definition-relative-global-complete-intersection} 中的表现：
\begin{enumerate}
\item 设 $I \subset R$ 是理想。若
$\Spec(S/IS) \to \Spec(R/I)$ 的纤维具有维数 $n - c$，
则可以找到如上的 $(h, g)$，使得 $g$ 的像为 $1 \in S/IS$。
\item 设 $\mathfrak p \subset R$ 是素理想。若
$\dim(S \otimes_R \kappa(\mathfrak p)) = n - c$，
则可以找到如上的 $(h, g)$，使得 $g$ 在
$S \otimes_R \kappa(\mathfrak p)$ 中的像为单位。
\item 设 $\mathfrak q \subset S$ 是位于
$\mathfrak p \subset R$ 上方的素理想。若
$\dim_{\mathfrak q}(S/R) = n - c$，则可以找到如上的
$(h, g)$，使得 $g \not \in \mathfrak q$。
\end{enumerate}
\end{lemma}

\begin{proof}
对于 (1)，由引理 \ref{algebra-lemma-dimension-fibres-bounded-open-upstairs}，
存在包含 $V(IS)$ 的开子集 $W \subset \Spec(S)$，使得
$W \to \Spec(R)$ 的所有纤维都具有维数 $\leq n - c$。
写成 $W = \Spec(S) \setminus V(J)$。于是
$V(J) \cap V(IS) = \emptyset$，故可以找到 $g \in J$，
其像为 $1 \in S/IS$。
任取 $h \in R[x_1, \ldots, x_n]$ 为 $g$ 的原像。

\medskip\noindent
对于 (2)，由引理 \ref{algebra-lemma-dimension-fibres-bounded-open-upstairs}，
存在包含 $\Spec(S \otimes_R \kappa(\mathfrak p))$ 的开子集
$W \subset \Spec(S)$，使得 $W \to \Spec(R)$ 的所有纤维
都具有维数 $\leq n - c$。
写成 $W = \Spec(S) \setminus V(J)$。于是
$V(J \cdot S \otimes_R \kappa(\mathfrak p)) = \emptyset$。
故可以找到 $g \in J$，其在
$S \otimes_R \kappa(\mathfrak p)$ 中的像为单位（略去细节）。
任取 $h \in R[x_1, \ldots, x_n]$ 为 $g$ 的原像。

\medskip\noindent
对于 (3)，由引理 \ref{algebra-lemma-dimension-fibres-bounded-open-upstairs}，
存在 $g \in S$，$g \not \in \mathfrak q$，使得
$R \to S_g$ 的所有非空纤维都具有维数 $\leq n - c$。
任取 $h \in R[x_1, \ldots, x_n]$ 为映到 $g$ 的元素。
\end{proof}

\noindent
下面的引理说明，可以对相对全局完全交作绝对 Noether 逼近。

\begin{lemma}
\label{algebra-lemma-relative-global-complete-intersection-Noetherian}
设 $R$ 是环。设
$S = R[x_1, \ldots, x_n]/(f_1, \ldots, f_c)$
是相对全局完全交（定义
\ref{algebra-definition-relative-global-complete-intersection}）。
存在一个有限型 $\mathbf{Z}$-子代数 $R_0 \subset R$，使得
% P01-E0798: the frozen “There exist a ... subalgebra” has subject-verb disagreement.
$f_i \in R_0[x_1, \ldots, x_n]$，且
$$
S_0 = R_0[x_1, \ldots, x_n]/(f_1, \ldots, f_c)
$$
是相对全局完全交。
\end{lemma}

\begin{proof}
令 $R_0 \subset R$ 为 $R$ 中由多项式
$f_1, \ldots, f_c$ 的所有系数生成的 $\mathbf{Z}$-代数。定义：
$$
S_0 = R_0[x_1, \ldots, x_n]/(f_1, \ldots, f_c)
$$
显然 $S = R \otimes_{R_0} S_0$。
取素理想 $\mathfrak q \subset S$，并用
$\mathfrak p \subset R$、$\mathfrak q_0 \subset S_0$ 和
$\mathfrak p_0 \subset R_0$ 表示它所位于其上的素理想。
% P01-E0799: the frozen “denote ... the primes” construction omits by/as and compresses a chain of contractions into “it lies over”.
由于 $\dim (S \otimes_R \kappa(\mathfrak p) ) = n - c$，由引理
\ref{algebra-lemma-dimension-preserved-field-extension}，还有
$\dim (S_0 \otimes_{R_0} \kappa(\mathfrak p_0)) = n - c$。
由引理 \ref{algebra-lemma-dimension-fibres-bounded-open-upstairs}，
存在 $g \in S_0$，$g \not \in \mathfrak q_0$，使得
$R_0 \to (S_0)_g$ 的所有非空纤维都具有维数 $\leq n - c$。
由于 $\mathfrak q$ 是任意的且 $\Spec(S)$ 拟紧，可以找到有限多个
$g_1, \ldots, g_m \in S_0$，使得：
(a) 对 $j = 1, \ldots, m$，映射
$R_0 \to (S_0)_{g_j}$ 的非空纤维都具有维数 $\leq n - c$；
(b) $\Spec(S) \to \Spec(S_0)$ 的像包含于
$D(g_1) \cup \ldots \cup D(g_m)$。
换言之，$g_1, \ldots, g_m$ 在
$S = R \otimes_{R_0} S_0$ 中的像生成单位理想。
扩大 $R_0$ 后，可以假设 $g_1, \ldots, g_m$ 在 $S_0$
中生成单位理想。由 (a)，$R_0 \to S_0$ 的非空纤维
都具有维数 $\leq n - c$，结论得证。
\end{proof}

\begin{lemma}
\label{algebra-lemma-relative-global-complete-intersection-conormal}
设 $R$ 是环。设
$S = R[x_1, \ldots, x_n]/(f_1, \ldots, f_c)$
是相对全局完全交（定义
\ref{algebra-definition-relative-global-complete-intersection}）。
对 $S$ 的每个素理想 $\mathfrak q$，令 $\mathfrak q'$ 表示
$R[x_1, \ldots, x_n]$ 中对应的素理想。则
\begin{enumerate}
\item $f_1, \ldots, f_c$ 在局部环
$R[x_1, \ldots, x_n]_{\mathfrak q'}$ 中构成正则序列；
% P01-E0800: the frozen English uses singular “is” with the list f_1,...,f_c.
\item 每个环
$R[x_1, \ldots, x_n]_{\mathfrak q'}/(f_1, \ldots, f_i)$
在 $R$ 上平坦；并且
\item $S$-模 $(f_1, \ldots, f_c)/(f_1, \ldots, f_c)^2$
是自由的，且元素 $f_i \bmod (f_1, \ldots, f_c)^2$ 给出其一组基。
\end{enumerate}
\end{lemma}

\begin{proof}
先假设 $R$ 是 Noether 环。令 $\mathfrak p = R \cap \mathfrak q'$。
例如，由引理 \ref{algebra-lemma-lci} 可知，$f_1, \ldots, f_c$
在局部环
$R[x_1, \ldots, x_n]_{\mathfrak q'} \otimes_R \kappa(\mathfrak p)$
中构成正则序列。此外，局部环
$R[x_1, \ldots, x_n]_{\mathfrak q'}$ 在 $R_{\mathfrak p}$ 上平坦。
由于 $R$ 是 Noether 环，因而
$R[x_1, \ldots, x_n]_{\mathfrak q'}$ 也是 Noether 环，
由引理 \ref{algebra-lemma-grothendieck-regular-sequence} 可知 (1) 和 (2) 成立。

\medskip\noindent
现在设 $R$ 是一般的环。将
$R = \colim_{\lambda \in \Lambda} R_\lambda$ 写成有限型
$\mathbf{Z}$-子代数的滤过余极限（与第
\ref{algebra-section-colimits-flat} 节比较）。可以假设对所有 $\lambda$ 都有
$f_1, \ldots, f_c \in R_\lambda[x_1, \ldots, x_n]$。
令 $R_0 \subset R$ 如引理
\ref{algebra-lemma-relative-global-complete-intersection-Noetherian} 中所述。
于是可以假设对所有 $\lambda$ 都有 $R_0 \subset R_\lambda$。
由此，$S_\lambda = R_\lambda[x_1, \ldots, x_n]/(f_1, \ldots, f_c)$
是相对全局完全交（它是 $S_0$ 经由 $R_0 \to R_\lambda$
作基变换所得；见引理
\ref{algebra-lemma-base-change-relative-global-complete-intersection}）。
分别用 $\mathfrak p_\lambda$、$\mathfrak q_\lambda$、
$\mathfrak q'_\lambda$ 表示由 $\mathfrak p$、$\mathfrak q$、
$\mathfrak q'$ 诱导的 $R_\lambda$、$S_\lambda$、
$R_\lambda[x_1, \ldots, x_n]$ 中的素理想。
% P01-E0801: the frozen designation omits by/as and uses singular “the prime” for three primes.
在此记号下，对每个 $\lambda$，(1) 和 (2) 都成立。由于
$$
R[x_1, \ldots, x_n]_{\mathfrak q'}/(f_1, \ldots, f_i)
=
\colim R_\lambda[x_1, \ldots, x_n]_{\mathfrak q_\lambda'}/(f_1, \ldots, f_i)
$$
由引理 \ref{algebra-lemma-colimit-rings-flat}，可推出 (2) 中在 $R$ 上的平坦性。
又因为
\begin{align*}
R[x_1, \ldots, x_n]_{\mathfrak q'}/(f_1, \ldots, f_i)
\xrightarrow{f_{i + 1}}
R[x_1, \ldots, x_n]_{\mathfrak q'}/(f_1, \ldots, f_i) \\
=
\colim
\left(
R_\lambda[x_1, \ldots, x_n]_{\mathfrak q_\lambda'}/(f_1, \ldots, f_i)
\xrightarrow{f_{i + 1}}
R_\lambda[x_1, \ldots, x_n]_{\mathfrak q_\lambda'}/(f_1, \ldots, f_i)
\right)
\end{align*}
并且滤过余极限是正合的（引理
\ref{algebra-lemma-directed-colimit-exact}），所以 (1) 成立。

\medskip\noindent
最后证明 (3)。记 $N$ 为 $S$-模
$(f_1, \ldots, f_c)/(f_1, \ldots, f_c)^2$，并记
$e_i \in N$ 为 $f_i$ 的像。
% P01-E0802: the frozen “Denote N the S-module” construction omits by/as.
由引理 \ref{algebra-lemma-regular-quasi-regular} 和 (1)，对 $S$
的所有素理想 $\mathfrak q$，$e_1, \ldots, e_c$ 都是
$N_\mathfrak q$ 的一组基。
% P01-E0803: the frozen subject e_1,...,e_c takes singular “is a basis”; “form a basis” is required.
由引理 \ref{algebra-lemma-characterize-zero-local}，得出 (3)。
\end{proof}

\begin{lemma}
\label{algebra-lemma-relative-global-complete-intersection}
相对全局完全交是 syntomic 的，即它是平坦的。
% P01-E0804: the frozen “syntomic, i.e., flat” can misleadingly identify syntomicity with flatness outside the already-constrained relative-global-CI context.
\end{lemma}

\begin{proof}
设 $R \to S$ 是相对全局完全交。
它的纤维都是全局完全交，并且 $S$ 是有限表示 $R$-代数。
因此只须证明 $R \to S$ 平坦。这由引理
\ref{algebra-lemma-relative-global-complete-intersection-conormal} 的 (2) 得出。
\end{proof}

\begin{lemma}
\label{algebra-lemma-adjoin-roots}
设 $A$ 是环，并且
$P(x) = x^n + b_1 x^{n-1} + \ldots + b_n \in A[x]$
是 $A$ 上的首一多项式。则存在一个 syntomic、有限自由且忠实平坦的
环扩张 $A \subset A'$，使得对某些 $\beta_i \in A'$ 有
$P(x) = \prod_{i = 1, \ldots, n} (x - \beta_i)$。
\end{lemma}

\begin{proof}
取 $A' = A \otimes_R S$，其中 $R$ 和 $S$ 如例
\ref{algebra-example-roots-universal-polynomial} 中所述，$R \to A$
把 $a_i$ 映到 $b_i$；并令 $\beta_i = -1 \otimes \alpha_i$。
注意，由引理 \ref{algebra-lemma-relative-global-complete-intersection}，
$R \to S$ 忠实平坦、有限自由且 syntomic。
基变换 $A \to A'$ 继承这些性质；略去一些细节。
\end{proof}

\begin{lemma}
\label{algebra-lemma-syntomic}
设 $R \to S$ 是环映射。
设 $\mathfrak q \subset S$ 是位于 $R$ 的素理想 $\mathfrak p$
上方的素理想。下列条件等价：
\begin{enumerate}
\item 存在元素 $g \in S$，$g \not \in \mathfrak q$，使得
$R \to S_g$ 是 syntomic 的。
\item 存在元素 $g \in S$，$g \not \in \mathfrak q$，使得
$S_g$ 是 $R$ 上的相对全局完全交。
\item 存在元素 $g \in S$，$g \not \in \mathfrak q$，使得
$R \to S_g$ 是有限表示的，局部环映射
$R_{\mathfrak p} \to S_{\mathfrak q}$ 是平坦的，并且局部环
$S_{\mathfrak q}/\mathfrak pS_{\mathfrak q}$ 是
$\kappa(\mathfrak p)$ 上的完全交环（见定义
\ref{algebra-definition-lci-local-ring}）。
\end{enumerate}
\end{lemma}

\begin{proof}
蕴含 (1) $\Rightarrow$ (3) 是引理 \ref{algebra-lemma-lci-at-prime}。
蕴含 (2) $\Rightarrow$ (1) 是引理
\ref{algebra-lemma-relative-global-complete-intersection}。
还须证明 (3) 蕴含 (2)。

\medskip\noindent
假设 (3)。对某个 $g \in S$，$g\not\in \mathfrak q$，
% P01-E0805: the frozen sentence lacks a comma after the localization condition.
把 $S$ 替换为 $S_g$ 后，可以假设 $S$ 是有限表示 $R$-代数。
选取表现 $S = R[x_1, \ldots, x_n]/I$。令
$\mathfrak q' \subset R[x_1, \ldots, x_n]$ 为对应于
$\mathfrak q$ 的素理想。写 $\kappa(\mathfrak p) = k$。
注意 $S \otimes_R k = k[x_1, \ldots, x_n]/\overline{I}$，其中
$\overline{I} \subset k[x_1, \ldots, x_n]$ 是由 $I$ 的像生成的理想。
令 $\overline{\mathfrak q}' \subset k[x_1, \ldots, x_n]$
为由 $\mathfrak q'$ 的像生成的素理想。
由引理 \ref{algebra-lemma-lci-at-prime}，引理 \ref{algebra-lemma-lci}
的等价条件对 $\overline{I}$ 和 $\overline{\mathfrak q}'$ 成立。
设
$\overline{I}_{\overline{\mathfrak q}'}/
\overline{\mathfrak q}'\overline{I}_{\overline{\mathfrak q}'}$
在 $\kappa(\overline{\mathfrak q}')$ 上的维数为 $c$。
取 $f_1, \ldots, f_c \in I$，其像构成该向量空间的一组基。
像 $\overline{f}_j \in \overline{I}$ 生成
$\overline{I}_{\overline{\mathfrak q}'}$（由引理
\ref{algebra-lemma-lci}）。置
$S' = R[x_1, \ldots, x_n]/(f_1, \ldots, f_c)$。
令 $J$ 为满射 $S' \to S$ 的核。
由于 $S$ 是有限表示的，$J$ 是有限生成理想
（引理 \ref{algebra-lemma-compose-finite-type}）。考虑短正合列
$$
0 \to J \to S' \to S \to 0
$$
% P01-E0806: the frozen short exact display lacks terminal punctuation before the next sentence.
由于 $S_\mathfrak q$ 在 $R$ 上平坦，
$J_{\mathfrak q'} \otimes_R k \to S'_{\mathfrak q'} \otimes_R k$
是单射（引理 \ref{algebra-lemma-flat-tor-zero}）。
然而，按构造，$S'_{\mathfrak q'} \otimes_R k$ 同构地映到
$S_\mathfrak q \otimes_R k$。因此
$J_{\mathfrak q'} \otimes_R k =
J_{\mathfrak q'}/\mathfrak pJ_{\mathfrak q'} = 0$。
由 Nakayama 引理（引理 \ref{algebra-lemma-NAK}），存在
$g \in R[x_1, \ldots, x_n]$，$g \not \in \mathfrak q'$，使得
$J_g = 0$。换言之，$S'_g \cong S_g$。
进一步局部化后，由引理
\ref{algebra-lemma-localize-relative-complete-intersection} 可知，
$S'$（因而 $S$）成为相对全局完全交，如所要求。
\end{proof}

\begin{lemma}
\label{algebra-lemma-syntomic-presentation-ideal-mod-squares}
设 $R$ 是环。设 $S = R[x_1, \ldots, x_n]/I$，其中
$I$ 是有限生成理想。若 $g \in S$ 使得 $S_g$ 在 $R$ 上
是 syntomic 的，则 $(I/I^2)_g$ 是有限投射 $S_g$-模。
\end{lemma}

\begin{proof}
由引理 \ref{algebra-lemma-syntomic}，存在有限多个元素
$g_1, \ldots, g_m \in S$，它们在 $S_g$ 中生成单位理想，
并且每个 $S_{gg_j}$ 都是 $R$ 上的相对全局完全交。
由引理 \ref{algebra-lemma-finite-projective}，只须证明
$(I/I^2)_{gg_j}$ 有限投射，故可以假设 $S_g$
是相对全局完全交。在这种情况下，结论由引理
\ref{algebra-lemma-conormal-module-localize} 和
\ref{algebra-lemma-relative-global-complete-intersection-conormal} 得出。
\end{proof}

\begin{lemma}
\label{algebra-lemma-composition-syntomic}
设 $R \to S$、$S \to S'$ 是环映射。
\begin{enumerate}
\item 若 $R \to S$ 和 $S \to S'$ 是 syntomic 的，则
$R \to S'$ 是 syntomic 的。
\item 若 $R \to S$ 和 $S \to S'$ 是相对全局完全交，则
$R \to S'$ 是相对全局完全交。
\end{enumerate}
\end{lemma}

\begin{proof}
先证明 (2)。设 $R \to S$ 和 $S \to S'$ 是相对全局完全交，
% P01-E0807: the frozen “Proof of (2).” is a sentence fragment.
并且有定义 \ref{algebra-definition-relative-global-complete-intersection}
中的表现
$S =  R[x_1, \ldots, x_n]/(f_1, \ldots, f_c)$ 和
$S' = S[y_1, \ldots, y_m]/(h_1, \ldots, h_d)$。
则
$$
S' \cong
R[x_1, \ldots, x_n, y_1, \ldots, y_m]/(f_1, \ldots, f_c, h'_1, \ldots, h'_d)
$$
其中 $h_j' \in R[x_1, \ldots, x_n, y_1, \ldots, y_m]$
是 $h_j$ 的某些提升。
% P01-E0808: the frozen display and prose place the prime on h in inconsistent source order (h'_j versus h_j').
因此，只须对纤维环的维数给出界。于是可以假设 $R = k$ 是域。
在这种情况下，环 $S$ 是有限型 $k$-代数，且等维、维数为
$n - c$；有限型环映射 $S \to S'$ 的每个非空纤维环
都是等维的，且维数为 $m - d$。
例如，把引理 \ref{algebra-lemma-dimension-base-fibre-total}
应用于 $S'$ 在极大理想处的局部化，便可看出
$\dim(S') \leq n - c + m - d$，如所要求。

\medskip\noindent
我们把 (1) 归约到 (2)。假设 $R \to S$ 和 $S \to S'$
是 syntomic 的。设 $\mathfrak q' \subset S$ 是位于
$\mathfrak q \subset S$ 上方的素理想。
% P01-E0809: the frozen source places both q' and q in S; q' must be a prime of S' for the subsequent localization and containment statements.
由引理 \ref{algebra-lemma-syntomic}，存在
$g' \in S'$，$g' \not \in \mathfrak q'$，使得
$S \to S'_{g'}$ 是相对全局完全交。
类似地，可以找到 $g \in S$，$g \not \in \mathfrak q$，使得
$R \to S_g$ 是相对全局完全交。
由引理 \ref{algebra-lemma-base-change-relative-global-complete-intersection}，
环映射 $S_g \to S_{gg'}$ 是相对全局完全交。
% P01-E0810: the frozen base-change target is S_{gg'}, but g' belongs to S'; the target should be S'_{gg'}.
由 (2)，$R \to S_{gg'}$ 是相对全局完全交，并且
$gg' \not \in \mathfrak q'$。
% P01-E0811: the frozen composite repeats S_{gg'} where the target must be the localization of S'.
由于 $\mathfrak q'$ 是任意的，结合引理 \ref{algebra-lemma-syntomic}
和 \ref{algebra-lemma-local-syntomic}，可知 $R \to S'$ 是 syntomic 的
（这里也使用了 $S'$ 的谱拟紧；见引理
\ref{algebra-lemma-quasi-compact}）。
\end{proof}

\noindent
下面的引理稍后会得到加强；见《环同态的平滑化》的命题
\ref{smoothing-proposition-lift-smooth}。

\begin{lemma}
\label{algebra-lemma-lift-syntomic}
设 $R$ 是环，$I \subset R$ 是理想。
设 $R/I \to \overline{S}$ 是 syntomic 映射。
则存在生成 $\overline{S}$ 的单位理想的元素
$\overline{g}_i \in \overline{S}$，使得每个
% P01-E0812: the frozen “there exists elements” has subject-verb disagreement.
$\overline{S}_{g_i} \cong S_i/IS_i$，其中 $S_i$
是 $R$ 上的某个相对全局完全交。
% P01-E0813: the frozen localization uses g_i although only \overline{g}_i is introduced; the same mismatch recurs in the proof.
\end{lemma}

\begin{proof}
由引理 \ref{algebra-lemma-syntomic}，可以找到生成 $\overline{S}$
的单位理想的一族元素 $\overline{g}_i \in \overline{S}$，
使得每个 $\overline{S}_{g_i}$ 都是 $R/I$ 上的相对全局完全交。
因此可以假设 $\overline{S}$ 是相对全局完全交。写成
$\overline{S} =
(R/I)[x_1, \ldots, x_n]/(\overline{f}_1, \ldots, \overline{f}_c)$，
如定义 \ref{algebra-definition-relative-global-complete-intersection} 中所述。
选取 $f_1, \ldots, f_c \in R[x_1, \ldots, x_n]$，
它们分别提升 $\overline{f}_1, \ldots, \overline{f}_c$。
置 $S = R[x_1, \ldots, x_n]/(f_1, \ldots, f_c)$。
注意 $S/IS \cong \overline{S}$。
由引理 \ref{algebra-lemma-localize-relative-complete-intersection}，
可以找到 $g \in S$，它在 $\overline{S}$ 中的像为 $1$，
并且 $S_g$ 是 $R$ 上的相对全局完全交。
由于 $\overline{S} \cong S_g/IS_g$，证明完成。
\end{proof}


\section{光滑环映射}
\label{algebra-section-smooth}

\noindent
先用一个例子说明光滑环映射定义的动机。
设 $R$ 是环，且对某个非零 $f \in R[x, y]$ 有
$S = R[x, y]/(f)$。此时有正合列
$$
S \to
S\text{d}x \oplus S\text{d}y \to
\Omega_{S/R} \to 0
$$
其中第一支箭头把 $1$ 映到
$\frac{\partial f}{\partial x} \text{d}x +
\frac{\partial f}{\partial y} \text{d}y$；见第
% P01-E0814: the frozen sentence lacks punctuation before “see Section”.
\ref{algebra-section-netherlander} 节。
由此可知，若 $f$ 的偏导数在 $S$ 中生成单位理想，则
$\Omega_{S/R}$ 是秩为 $1$ 的局部自由模。
在这种情况下，$S$ 在 $R$ 上是相对维数为 $1$ 的光滑代数。
但 $\Omega_{S/R}$ 也可能是秩为 $2$ 的局部自由模，即 $f$
的两个偏导数都为零时如此。例如，若对某个素数 $p$，
$R$ 中有 $p = 0$，并且 $f = x^p + y^p$，就会出现这种情况。
这里 $R \to S$ 是相对维数为 $1$ 的相对全局完全交，但并不光滑。
因此，要检验环映射是否光滑，仅检验微分模是否自由是不够的。
正确的条件如下。

\begin{definition}
\label{algebra-definition-smooth}
若环映射 $R \to S$ 是有限表示的，并且朴素余切复形
$\NL_{S/R}$ 拟同构于置于次数 $0$ 的某个有限投射 $S$-模，
则称它是{\it 光滑的}；这意味着
$H_1(\NL_{S/R}) = 0$，并且 $\Omega_{S/R}$ 是有限投射 $S$-模。
\end{definition}

\noindent
设 $R \to S$ 是环映射。由引理 \ref{algebra-lemma-NL-homotopy}，
对任意表现 $\alpha : P \to S$，都有
$H_1(\NL(\alpha)) = H_1(\NL_{S/R})$ 以及
$H_0(\NL(\alpha)) = \Omega_{S/R}$。
因此，若 $R \to S$ 光滑，则对任意核为 $I$ 的满射
$\alpha : R[x_1, \ldots, x_n] \to S$，映射
$$
I/I^2
\longrightarrow
\Omega_{R[x_1, \ldots, x_n]/R} \otimes_{R[x_1, \ldots, x_n]} S
$$
都是单射，并且其余核是有限投射 $S$-模。
故所显示的箭头是分裂单射。换言之，作为 $S$-模，有
$\bigoplus_{i = 1}^n S \text{d}x_i \cong I/I^2 \oplus \Omega_{S/R}$，
而 $I/I^2$ 也是有限投射 $S$-模。
反过来，若 $R \to S$ 是有限表示的，并且存在核为 $I$ 的满射
$\alpha : R[x_1, \ldots, x_n] \to S$，使得所显示的箭头
是分裂单射，则 $R \to S$ 光滑。

\begin{lemma}
\label{algebra-lemma-localize-smooth}
设 $R \to S$ 是光滑环映射。
任意局部化 $S_g$ 在 $R$ 上都是光滑的。
若 $f \in R$ 在 $S$ 中的像可逆，则 $R_f \to S$ 光滑。
\end{lemma}

\begin{proof}
由引理 \ref{algebra-lemma-localize-NL}，$S_g$ 在 $R$ 上的朴素余切复形
是 $S$ 在 $R$ 上的朴素余切复形的基变换。
假设是 $S/R$ 的朴素余切复形拟同构于 $\Omega_{S/R}$，
% P01-E0816: the frozen proof says the naive cotangent complex “is” Omega although the definition only gives a quasi-isomorphism.
并且后者是有限投射 $S$-模。因此，它的基变换也具有同样的性质。
故 $S_g$ 在 $R$ 上光滑。

\medskip\noindent
第二个断言以同样方式由引理 \ref{algebra-lemma-NL-localize-bottom} 得出。
\end{proof}

\begin{lemma}
\label{algebra-lemma-base-change-smooth}
\begin{slogan}
光滑性在基变换下保持。
% P01-E0815: the frozen slogan lacks terminal punctuation.
\end{slogan}
设 $R \to S$ 是光滑环映射。
设 $R \to R'$ 是任意环映射。
则基变换 $R' \to S' = R' \otimes_R S$ 光滑。
\end{lemma}

\begin{proof}
设 $\alpha : R[x_1, \ldots, x_n] \to S$ 是核为 $I$ 的一个表现。
设 $\alpha' : R'[x_1, \ldots, x_n] \to R' \otimes_R S$
为所诱导的表现。令 $I' = \Ker(\alpha')$。
由于 $0 \to I \to R[x_1, \ldots, x_n] \to S \to 0$ 正合，序列
$R' \otimes_R I \to R'[x_1, \ldots, x_n] \to R' \otimes_R S \to 0$
正合。因此 $R' \otimes_R I \to I'$ 是满射。
由定义 \ref{algebra-definition-smooth}，有短正合列
$$
0 \to I/I^2 \to
\Omega_{R[x_1, \ldots, x_n]/R} \otimes_{R[x_1, \ldots, x_n]} S \to
\Omega_{S/R} \to
0
$$
并且 $S$-模 $\Omega_{S/R}$ 是有限投射的。特别地，$I/I^2$
是 $\Omega_{R[x_1, \ldots, x_n]/R}
\otimes_{R[x_1, \ldots, x_n]} S$ 的直和项。
考虑交换图
$$
\xymatrix{
R' \otimes_R (I/I^2) \ar[r] \ar[d] &
R' \otimes_R (\Omega_{R[x_1, \ldots, x_n]/R} \otimes_{R[x_1, \ldots, x_n]} S)
\ar[d] \\
I'/(I')^2 \ar[r] &
\Omega_{R'[x_1, \ldots, x_n]/R'}
\otimes_{R'[x_1, \ldots, x_n]} (R' \otimes_R S)
}
$$
由于右边的竖直映射是同构，结合上面所述，可知左边的竖直映射
既是单射也是满射。因此，$\NL(\alpha')$ 拟同构于置于次数 $0$ 的
$\Omega_{S'/R'} \cong S' \otimes_S \Omega_{S/R}$。
该模是有限投射的，因为它是有限投射模的基变换。
\end{proof}

\begin{lemma}
\label{algebra-lemma-smooth-over-field}
设 $k$ 是域。
设 $S$ 是光滑 $k$-代数。
则 $S$ 是局部完全交。
\end{lemma}

\begin{proof}
由引理 \ref{algebra-lemma-base-change-smooth} 和
\ref{algebra-lemma-lci-field-change}，只须证明 $k$ 为代数闭域的情形。
选取核为 $I$ 的表现
$\alpha : k[x_1, \ldots, x_n] \to S$。
设 $\mathfrak m$ 是 $S$ 的极大理想，并令
$\mathfrak m' \supset I$ 为 $k[x_1, \ldots, x_n]$
中对应的极大理想。我们将证明引理 \ref{algebra-lemma-lci}
的条件 (5) 成立（以 $\mathfrak m$ 代替 $\mathfrak q$）。
因为 $k$ 是代数闭域，可以对某些 $a_i \in k$ 写成
$\mathfrak m' = (x_1 - a_1, \ldots, x_n - a_n)$；见定理
\ref{algebra-theorem-nullstellensatz}。
由 $k \to S$ 光滑这一假设，$S$-模映射
$\text{d} : I/I^2 \to \bigoplus_{i = 1}^n S \text{d}x_i$
是分裂单射。因此，相应的映射
$I/\mathfrak m' I \to \bigoplus \kappa(\mathfrak m') \text{d}x_i$
是单射。设
$\dim_{\kappa(\mathfrak m')}(I/\mathfrak m' I) = c$，
并取 $f_1, \ldots, f_c \in I$，其像构成
$I/\mathfrak m' I$ 的一组 $\kappa(\mathfrak m')$-基。
由 Nakayama 引理 \ref{algebra-lemma-NAK}，$f_1, \ldots, f_c$
在 $k[x_1, \ldots, x_n]_{\mathfrak m'}$ 上生成
$I_{\mathfrak m'}$。考虑交换图
$$
\xymatrix{
I \ar[r] \ar[d] & I/I^2 \ar[rr] \ar[d] & &
I/\mathfrak m'I \ar[d] \\
\Omega_{k[x_1, \ldots, x_n]/k} \ar[r] &
\bigoplus S\text{d}x_i \ar[rr]^{\text{d}x_i \mapsto x_i - a_i} & &
\mathfrak m'/(\mathfrak m')^2
}
$$
（略去交换性的证明）。中间的竖直映射就是定义
$\alpha$ 的朴素余切复形的映射。注意，右下方的水平箭头诱导同构
$\bigoplus \kappa(\mathfrak m') \text{d}x_i
\to \mathfrak m'/(\mathfrak m')^2$。
因此，$I_{\mathfrak m'}$ 的生成元 $f_1, \ldots, f_c$
映到 $k[x_1, \ldots, x_n]_{\mathfrak m'}$ 中的一族元素，
它们在 $\mathfrak m'/(\mathfrak m')^2$ 中的类在
$\kappa(\mathfrak m')$ 上线性无关。
故由引理 \ref{algebra-lemma-regular-ring-CM}，它们在环
$k[x_1, \ldots, x_n]_{\mathfrak m'}$ 中构成正则序列。
这验证了引理 \ref{algebra-lemma-lci} 的条件 (5)，
% P01-E0817: the frozen source joins the lemma citation directly to “hence” without punctuation.
因而对某个 $g \in S$，$g \not \in \mathfrak m$，$S_g$
是 $k$ 上的全局完全交。由于这对 $S$ 的任意极大理想都成立，
可知 $S$ 是 $k$ 上的局部完全交。
\end{proof}

\begin{definition}
\label{algebra-definition-standard-smooth}
设 $R$ 是环。给定整数 $n \geq c \geq 0$ 以及
$f_1, \ldots, f_c \in R[x_1, \ldots, x_n]$，称
% P01-E0818: the frozen fronted data lack a comma before “we say”.
$$
R[x_1, \ldots, x_n]/(f_1, \ldots, f_c)
$$
为 $R$ 上的一个{\it 标准光滑代数}，如果多项式
$$
g =
\det
\left(
\begin{matrix}
\partial f_1/\partial x_1 &
\partial f_2/\partial x_1 &
\ldots &
\partial f_c/\partial x_1 \\
\partial f_1/\partial x_2 &
\partial f_2/\partial x_2 &
\ldots &
\partial f_c/\partial x_2 \\
\ldots & \ldots & \ldots & \ldots \\
\partial f_1/\partial x_c &
\partial f_2/\partial x_c &
\ldots &
\partial f_c/\partial x_c
\end{matrix}
\right)
$$
在 $R[x_1, \ldots, x_n]/(f_1, \ldots, f_c)$ 中的像可逆。
若存在 $n \geq c \geq 0$ 和
$f_1, \ldots, f_c \in R[x_1, \ldots, x_n]$，使得下列代数：
$$
R[x_1, \ldots, x_n]/(f_1, \ldots, f_c)
$$
是 $R$ 上的标准光滑代数，并且 $S$ 作为 $R$-代数与下列代数同构：
$$
R[x_1, \ldots, x_n]/(f_1, \ldots, f_c)
$$
则称 $R$-代数 $S$ 是{\it 标准光滑的}，也称环映射
$R \to S$ 是{\it 标准光滑的}。
\end{definition}

\begin{lemma}
\label{algebra-lemma-standard-smooth}
设
$S = R[x_1, \ldots, x_n]/(f_1, \ldots, f_c)
= R[x_1, \ldots, x_n]/I$
是标准光滑代数。则
\begin{enumerate}
\item 环映射 $R \to S$ 光滑；
\item $S$-模 $\Omega_{S/R}$ 以
$\text{d}x_{c + 1}, \ldots, \text{d}x_n$ 为一组基；
\item $S$-模 $I/I^2$ 以 $f_1, \ldots, f_c$ 的类为一组基；
\item 对任意 $g \in S$，环映射 $R \to S_g$ 标准光滑；
\item 对任意环映射 $R \to R'$，基变换
$R' \to R'\otimes_R S$ 标准光滑；
\item 若 $f \in R$ 在 $S$ 中的像可逆，则 $R_f \to S$ 标准光滑；并且
\item 环 $S$ 是 $R$ 上的相对全局完全交。
\end{enumerate}
\end{lemma}

\begin{proof}
考虑给定表现的朴素余切复形
$$
(f_1, \ldots, f_c)/(f_1, \ldots, f_c)^2
\longrightarrow
\bigoplus\nolimits_{i = 1}^n S \text{d}x_i
$$
% P01-E0819: the frozen display lacks terminal punctuation before the next sentence.
将此映射与到直和的前 $c$ 个直和项上的投影复合。
按照标准光滑代数的定义，类
$f_i \bmod (f_1, \ldots, f_c)^2$ 映到
$\bigoplus_{i = 1}^c S\text{d}x_i$ 的一组基。
由此得出 $(f_1, \ldots, f_c)/(f_1, \ldots, f_c)^2$
是秩为 $c$ 的自由模，元素
$f_i \bmod (f_1, \ldots, f_c)^2$ 给出其一组基；并且朴素余切复形
在次数 $0$ 的同调，即 $\Omega_{S/R}$，是一个自由 $S$-模，
其一组基为 $\text{d}x_{c + j}$（$j = 1, \ldots, n - c$）的像。
特别地，这证明 $R \to S$ 光滑。

\medskip\noindent
略去 (4) 和 (6) 的证明。不过，可参见下面的例子以及引理
\ref{algebra-lemma-base-change-relative-global-complete-intersection} 的证明。

\medskip\noindent
设 $\varphi : R \to R'$ 是任意环映射。
% P01-E0820: the frozen “Denote S' = ... where ...” construction omits by/as.
置 $S' = R'[x_1, \ldots, x_n]/(f_1^\varphi, \ldots, f_c^\varphi)$，
其中 $f^\varphi$ 是对 $f \in R[x_1, \ldots, x_n]$ 的所有系数
应用 $\varphi$ 后得到的多项式。则 $S' \cong R' \otimes_R S$。
此外，定义 \ref{algebra-definition-standard-smooth} 中 $S'/R'$
的行列式等于 $g^\varphi$。因此，它在 $S'$ 中的像就是 $g$
经由 $R[x_1, \ldots, x_n] \to S \to S'$ 的像，故可逆。
这证明了 (5)。

\medskip\noindent
要证明 (7)，只须证明对每个素理想 $\mathfrak p \subset R$，
$S \otimes_R \kappa(\mathfrak p)$ 的维数都是 $n - c$。
由 (5)，只须证明域 $k$ 上任意标准光滑代数
$k[x_1, \ldots, x_n]/(f_1, \ldots, f_c)$ 的维数为 $n - c$。
由引理 \ref{algebra-lemma-smooth-over-field}，我们已经知道
$k[x_1, \ldots, x_n]/(f_1, \ldots, f_c)$ 是局部完全交。
因此，由于 $I/I^2$ 是秩为 $c$ 的自由模，例如由引理
\ref{algebra-lemma-lci} 可知
$k[x_1, \ldots, x_n]/(f_1, \ldots, f_c)$ 的维数为 $n - c$。
\end{proof}

\begin{example}
\label{algebra-example-make-standard-smooth}
设 $R$ 是环。
设 $f_1, \ldots, f_c \in R[x_1, \ldots, x_n]$。
令
$$
h =
\det
\left(
\begin{matrix}
\partial f_1/\partial x_1 &
\partial f_2/\partial x_1 &
\ldots &
\partial f_c/\partial x_1 \\
\partial f_1/\partial x_2 &
\partial f_2/\partial x_2 &
\ldots &
\partial f_c/\partial x_2 \\
\ldots & \ldots & \ldots & \ldots \\
\partial f_1/\partial x_c &
\partial f_2/\partial x_c &
\ldots &
\partial f_c/\partial x_c
\end{matrix}
\right).
$$
置 $S = R[x_1, \ldots, x_{n + 1}]/(f_1, \ldots, f_c, x_{n + 1}h - 1)$。
这是标准光滑代数的一个例子，只不过该表现中的变量次序不对，
应当依次为
$x_1, \ldots, x_c, x_{n + 1}, x_{c + 1}, \ldots, x_n$。
\end{example}

\begin{lemma}
\label{algebra-lemma-compose-standard-smooth}
标准光滑环映射的合成仍是标准光滑环映射。
\end{lemma}

\begin{proof}
假设 $R \to S$ 和 $S \to S'$ 标准光滑。选取表现
$S =  R[x_1, \ldots, x_n]/(f_1, \ldots, f_c)$
以及
$S' = S[y_1, \ldots, y_m]/(g_1, \ldots, g_d)$。
选取映到 $g_j$ 的元素
$g_j' \in R[x_1, \ldots, x_n, y_1, \ldots, y_m]$。
由此可见
$S' = R[x_1, \ldots, x_n, y_1, \ldots, y_m]/
(f_1, \ldots, f_c, g'_1, \ldots, g'_d)$。
% P01-E0821: the frozen source alternates g_j' and g'_j for the same lifts.
要证明 $S'$ 标准光滑，只须验证行列式
$$
\det
\left(
\begin{matrix}
\partial f_1/\partial x_1 &
\ldots &
\partial f_c/\partial x_1 &
\partial g_1/\partial x_1 &
\ldots &
\partial g_d/\partial x_1 \\
\ldots &
\ldots &
\ldots &
\ldots &
\ldots &
\ldots  \\
\partial f_1/\partial x_c &
\ldots &
\partial f_c/\partial x_c &
\partial g_1/\partial x_c &
\ldots &
\partial g_d/\partial x_c \\
0 &
\ldots &
0 &
\partial g_1/\partial y_1 &
\ldots &
\partial g_d/\partial y_1 \\
\ldots &
\ldots &
\ldots &
\ldots &
\ldots &
\ldots  \\
0 &
\ldots &
0 &
\partial g_1/\partial y_d &
\ldots &
\partial g_d/\partial y_d
\end{matrix}
\right)
$$
在 $S'$ 中可逆。这是显然的，因为它是由假设可逆的两个行列式
的乘积。
\end{proof}

\begin{lemma}
\label{algebra-lemma-smooth-syntomic}
设 $R \to S$ 是光滑环映射。
存在由标准开集 $D(g)$ 构成的 $\Spec(S)$ 的开覆盖，使得每个
$S_g$ 在 $R$ 上标准光滑。特别地，$R \to S$ 是 syntomic 的。
\end{lemma}

\begin{proof}
选取核为 $I = (f_1, \ldots, f_m)$ 的表现
$\alpha : R[x_1, \ldots, x_n] \to S$。
对每个子集 $E \subset \{1, \ldots, m\}$，考虑开子集
$U_E$：在其上，类 $f_e, e\in E$ 自由生成有限投射
$S$-模 $I/I^2$；见引理 \ref{algebra-lemma-cokernel-flat}。
可以用标准开集 $D(g)$ 覆盖 $\Spec(S)$，使每个这样的开集
都完全包含于某个 $U_E$。对这样的 $g$，考虑表现
$$
\beta : R[x_1, \ldots, x_n, x_{n + 1}] \longrightarrow S_g
$$
它把 $x_{n + 1}$ 映到 $1/g$。置 $J = \Ker(\beta)$。
由引理 \ref{algebra-lemma-principal-localization-NL}，
$J/J^2 \cong (I/I^2)_g \oplus S_g$ 是自由模。
我们可以且确实把 $S$ 替换为 $S_g$。于是，利用引理
\ref{algebra-lemma-huber}，可以假设有核为
$I = (f_1, \ldots, f_c)$ 的表现
$\alpha : R[x_1, \ldots, x_n] \to S$，并且 $I/I^2$
以 $f_1, \ldots, f_c$ 的类为一组基。

\medskip\noindent
利用上一段末尾得到的表现 $\alpha$，对
$\Omega_{R[x_1, \ldots, x_n]/R}$ 的基元素
$\text{d}x_1, \ldots, \text{d}x_n$ 作大致相同的论证。
具体地，对任意基数为 $c$ 的子集
$E \subset \{1, \ldots, n\}$，考虑 $\Spec(S)$ 的开子集
$U_E$：在其上，$\NL(\alpha)$ 的微分与投影的复合
$$
S^{\oplus c} \cong I/I^2
\longrightarrow
\Omega_{R[x_1, \ldots, x_n]/R} \otimes_{R[x_1, \ldots, x_n]} S
\longrightarrow
\bigoplus\nolimits_{i \in E} S\text{d}x_i
$$
是同构。再次可以找到由（有限多个）标准开集 $D(g)$ 构成的
$\Spec(S)$ 的覆盖，使每个 $D(g)$ 完全包含于某个 $U_E$。
重编号后，可以假设 $E = \{1, \ldots, c\}$。
对满足 $D(g) \subset U_E$ 的 $g$，考虑表现
$$
\beta : R[x_1, \ldots, x_n, x_{n + 1}] \to S_g
$$
它把 $x_{n + 1}$ 映到 $1/g$。置 $J = \Ker(\beta)$。
由引理 \ref{algebra-lemma-principal-localization-NL} 可知
$J = (f_1, \ldots, f_c, fx_{n + 1} - 1)$，其中
$\alpha(f) = g$，并且复合
$$
J/J^2 \longrightarrow
\Omega_{R[x_1, \ldots, x_{n + 1}]/R}
\otimes_{R[x_1, \ldots, x_{n + 1}]} S_g
\longrightarrow
\bigoplus\nolimits_{i = 1}^c S_g\text{d}x_i \oplus S_g \text{d}x_{n + 1}
$$
是同构。将坐标重排为
$x_1, \ldots, x_c, x_{n + 1}, x_{c + 1}, \ldots, x_n$，
便得 $S_g$ 在 $R$ 上标准光滑，如所要求。

\medskip\noindent
这就完成了证明，因为标准光滑代数是 syntomic 的
（引理 \ref{algebra-lemma-standard-smooth} 和
\ref{algebra-lemma-relative-global-complete-intersection}），
而在 $R$ 上为 syntomic 的性质在 $S$ 上是局部的
（引理 \ref{algebra-lemma-local-syntomic}）。
\end{proof}

\begin{definition}
\label{algebra-definition-smooth-at-prime}
设 $R \to S$ 是环映射。
设 $\mathfrak q$ 是 $S$ 的素理想。
若存在 $g \in S$，$g \not \in \mathfrak q$，使得
$R \to S_g$ 光滑，则称 $R \to S$ {\it 在 $\mathfrak q$ 处光滑}。
\end{definition}

\noindent
对有限表示环映射，可以如下刻画这一性质。

\begin{lemma}
\label{algebra-lemma-smooth-at-point}
设 $R \to S$ 是有限表示的。设 $\mathfrak q$ 是 $S$ 的素理想。
下列条件等价：
% P01-E0822: the frozen statement lacks a colon before the enumeration.
\begin{enumerate}
\item $R \to S$ 在 $\mathfrak q$ 处光滑；
\item $H_1(L_{S/R})_\mathfrak q = 0$，并且
$\Omega_{S/R, \mathfrak q}$ 是有限自由 $S_\mathfrak q$-模；
\item $H_1(L_{S/R})_\mathfrak q = 0$，并且
$\Omega_{S/R, \mathfrak q}$ 是投射 $S_\mathfrak q$-模；并且
\item $H_1(L_{S/R})_\mathfrak q = 0$，并且
$\Omega_{S/R, \mathfrak q}$ 是平坦 $S_\mathfrak q$-模。
% P01-E0823: the frozen criterion switches from the section's \NL_{S/R} notation to an otherwise unexplained L_{S/R}.
\end{enumerate}
\end{lemma}

\begin{proof}
我们将不再另行说明地使用朴素余切复形的形成与局部化可交换这一事实；
见第 \ref{algebra-section-netherlander} 节，特别是引理
\ref{algebra-lemma-localize-NL}。
注意 $\Omega_{S/R}$ 是有限表示 $S$-模；见引理
\ref{algebra-lemma-differentials-finitely-presented}。
因此，由引理 \ref{algebra-lemma-finite-projective}，(2)、(3) 和 (4) 等价。
显然 (1) 蕴含等价条件 (2)、(3) 和 (4)。
假设 (2) 成立。把 $S_\mathfrak q$ 写成主局部化的余极限，
由引理 \ref{algebra-lemma-colimit-category-fp-modules} 可知，存在
$g \in S$，$g \not \in \mathfrak q$，使得
$(\Omega_{S/R})_g$ 有限自由。
选取核为 $I$ 的表现 $\alpha : R[x_1, \ldots, x_n] \to S$。
由引理 \ref{algebra-lemma-NL-homotopy}，可以用 $\NL(\alpha)$
代替 $\NL_{S/R}$。满射
$$
\Omega_{R[x_1, \ldots, x_n]/R} \otimes_{R[x_1, \ldots, x_n]} S
\to \Omega_{S/R} \to 0
$$
在把 $g$ 变为可逆元后有右逆，因为 $(\Omega_{S/R})_g$ 是投射模。
因此
$\text{d} : (I/I^2)_g \to
\Omega_{R[x_1, \ldots, x_n]/R}
\otimes_{R[x_1, \ldots, x_n]} S_g$
的像是直和项，而且该映射也有右逆。
% P01-E0824: a split cokernel makes Im(d) a direct summand, but d itself need not have a right inverse onto the whole free target; only the surjection onto its image splits.
由此得出 $H_1(L_{S/R})_g$ 是 $(I/I^2)_g$ 的商。
特别地，$H_1(L_{S/R})_g$ 是有限 $S_g$-模。
因此，$H_1(L_{S/R})_{\mathfrak q}$ 的消失蕴含：对某个
$g' \in S$，$g' \not \in \mathfrak q$，
$H_1(L_{S/R})_{gg'}$ 消失。于是按照定义，
$R \to S_{gg'}$ 光滑。
\end{proof}

\begin{lemma}
\label{algebra-lemma-locally-smooth}
\begin{slogan}
环映射光滑，当且仅当它在目标的所有素理想处光滑。
% P01-E0825: the frozen slogan lacks terminal punctuation.
\end{slogan}
设 $R \to S$ 是环映射。
则 $R \to S$ 光滑，当且仅当 $R \to S$ 在 $S$
的每个素理想 $\mathfrak q$ 处光滑。
\end{lemma}

\begin{proof}
正向蕴含是平凡的。假设 $R \to S$ 在 $S$ 的每个素理想
$\mathfrak q$ 处都光滑。由于 $\Spec(S)$ 拟紧（见引理
\ref{algebra-lemma-quasi-compact}），存在有限覆盖
$\Spec(S) = \bigcup D(g_i)$，使得每个 $S_{g_i}$ 都光滑。
由引理 \ref{algebra-lemma-cover-upstairs}，这蕴含 $S$ 是有限表示
$R$-代数。由引理 \ref{algebra-lemma-localize-NL}，
$\NL_{S/R} \otimes_S S_{g_i}$ 拟同构于置于次数 $0$
的某个有限投射 $S_{g_i}$-模。
由引理 \ref{algebra-lemma-finite-projective}，这蕴含 $\NL_{S/R}$
拟同构于置于次数 $0$ 的某个有限投射 $S$-模。
\end{proof}

\begin{lemma}
\label{algebra-lemma-compose-smooth}
光滑环映射的合成仍是光滑环映射。
\end{lemma}

\begin{proof}
这可以用许多不同的方法证明。一种方法是使用蛇形引理（引理
\ref{algebra-lemma-snake}）、Jacobi--Zariski 序列（引理
\ref{algebra-lemma-exact-sequence-NL}），并结合投射模是自由模直和项的刻画
（引理 \ref{algebra-lemma-characterize-projective}）。
另一个证明可以通过结合引理 \ref{algebra-lemma-smooth-syntomic}、
\ref{algebra-lemma-compose-standard-smooth} 和
\ref{algebra-lemma-locally-smooth} 得到。
\end{proof}

\begin{lemma}
\label{algebra-lemma-product-smooth}
设 $R$ 是环。设 $S = S' \times S''$ 是 $R$-代数的乘积。
则 $S$ 在 $R$ 上光滑，当且仅当 $S'$ 和 $S''$
都在 $R$ 上光滑。
\end{lemma}

\begin{proof}
略。提示：由引理 \ref{algebra-lemma-locally-smooth}，可以逐个素理想检验
光滑性。由引理 \ref{algebra-lemma-spec-product}，
$\Spec(S)$ 是 $\Spec(S')$ 与 $\Spec(S'')$ 的不交并；
因此，$R \to S$ 在 $\mathfrak q$ 处的光滑性对应于
$R \to S'$ 在相应素理想处的光滑性，或者
$R \to S''$ 在相应素理想处的光滑性。
\end{proof}

\begin{lemma}
\label{algebra-lemma-relative-global-complete-intersection-smooth}
设 $R$ 是环。设
$S = R[x_1, \ldots, x_n]/(f_1, \ldots, f_c)$
是相对全局完全交。
设 $\mathfrak q \subset S$ 是素理想。则 $R \to S$
在 $\mathfrak q$ 处光滑，当且仅当存在基数为 $c$ 的子集
$I \subset \{1, \ldots, n\}$，使得多项式
$$
g_I = \det (\partial f_j/\partial x_i)_{j = 1, \ldots, c, \ i \in I}.
$$
不映到 $\mathfrak q$ 的元素。
% P01-E0826: the frozen display includes a period and is followed by the predicate fragment “does not map”.
\end{lemma}

\begin{proof}
由引理 \ref{algebra-lemma-relative-global-complete-intersection-conormal}，
与 $S$ 的给定表现相伴的朴素余切复形为
$$
\bigoplus\nolimits_{j = 1}^c S \cdot f_j
\longrightarrow
\bigoplus\nolimits_{i = 1}^n S \cdot \text{d}x_i, \quad
f_j \longmapsto \sum \frac{\partial f_j}{\partial x_i} \text{d}x_i.
$$
% P01-E0827: the frozen summation has no index range even though the target direct sum runs i=1,...,n.
给出该映射的矩阵之极大子式正是多项式 $g_I$。

\medskip\noindent
假设 $g_I$ 映到 $g \in S$，其中 $g \not \in \mathfrak q$。
则代数 $S_g$ 在 $R$ 上光滑。事实上，由引理
\ref{algebra-lemma-localize-NL}，它的朴素余切复形拟同构于上面的复形
在 $g$ 处的局部化。而按构造，它拟同构于置于次数 $0$
的秩为 $n - c$ 的自由模。
% P01-E0828: the frozen phrase “a free rank n-c module” is nonidiomatic.

\medskip\noindent
反过来，假设所有 $g_I$ 最终都属于 $\mathfrak q$。
在这种情况下，上面的复形与 $\kappa(\mathfrak q)$ 作张量后
不具有极大秩；因而不存在元素 $g \in S$、$g \not \in \mathfrak q$，
使得在它处局部化后该映射成为分裂单射。
再次由引理 \ref{algebra-lemma-localize-NL}，不存在这样的在 $R$ 上光滑的
局部化。
\end{proof}

\begin{lemma}
\label{algebra-lemma-flat-fibre-smooth}
设 $R \to S$ 是环映射。
设 $\mathfrak q \subset S$ 是位于 $R$ 的素理想
$\mathfrak p$ 上方的素理想。假设
\begin{enumerate}
\item 存在 $g \in S$，$g \not\in \mathfrak q$，使得
$R \to S_g$ 是有限表示的；
\item 局部环同态 $R_{\mathfrak p} \to S_{\mathfrak q}$ 是平坦的；
\item 纤维 $S \otimes_R \kappa(\mathfrak p)$ 在对应于
$\mathfrak q$ 的素理想处是 $\kappa(\mathfrak p)$ 上光滑的。
\end{enumerate}
则 $R \to S$ 在 $\mathfrak q$ 处光滑。
\end{lemma}

\begin{proof}
由引理 \ref{algebra-lemma-syntomic} 和 \ref{algebra-lemma-smooth-over-field}，
可知存在 $g \in S$，使得 $S_g$ 是相对全局完全交。
% P01-E0829: the frozen application omits the required condition g notin q before replacing S by S_g.
把 $S$ 替换为 $S_g$ 后，可以假设
$S = R[x_1, \ldots, x_n]/(f_1, \ldots, f_c)$
是相对全局完全交。
对任意基数为 $c$ 的子集 $I \subset \{1, \ldots, n\}$，
考虑引理 \ref{algebra-lemma-relative-global-complete-intersection-smooth}
中的多项式
$g_I = \det (\partial f_j/\partial x_i)_{j = 1, \ldots, c, i \in I}$。
注意，$g_I$ 在多项式环
$\kappa(\mathfrak p)[x_1, \ldots, x_n]$ 中的像
$\overline{g}_I$，是 $f_j$ 的像
$\overline{f}_j$ 在环
$\kappa(\mathfrak p)[x_1, \ldots, x_n]$ 中的偏导数之行列式。
因此，把引理
\ref{algebra-lemma-relative-global-complete-intersection-smooth}
分别应用于 $R \to S$ 以及
$\kappa(\mathfrak p) \to S \otimes_R \kappa(\mathfrak p)$，
即可得出本引理。
\end{proof}

\noindent
注意，下面引理中的集合 $U, V$ 按定义是开集。

\begin{lemma}
\label{algebra-lemma-flat-base-change-locus-smooth}
设 $R \to S$ 是有限表示环映射。
设 $R \to R'$ 是平坦环映射。
% P01-E0830: the frozen designation “Denote S'=... the base change” omits by/as.
记 $S' = R' \otimes_R S$ 为基变换。
令 $U \subset \Spec(S)$ 为 $R \to S$ 光滑处的素理想集合。
令 $V \subset \Spec(S')$ 为 $R' \to S'$ 光滑处的素理想集合。
% P01-E0831: the frozen “Let V ... the set” omits the verb “be”.
则在映射 $f : \Spec(S') \to \Spec(S)$ 下，$V$ 是 $U$ 的逆像。
\end{lemma}

\begin{proof}
由引理 \ref{algebra-lemma-change-base-NL}，$\NL_{S/R} \otimes_S S'$
同伦等价于 $\NL_{S'/R'}$。这已经蕴含
$f^{-1}(U) \subset V$。

\medskip\noindent
设 $\mathfrak q' \subset S'$ 是位于
$\mathfrak q \subset S$ 上方的素理想。假设 $\mathfrak q' \in V$。
须证明 $\mathfrak q \in U$。
由于 $S \to S'$ 平坦，$S_{\mathfrak q} \to S'_{\mathfrak q'}$
忠实平坦（引理 \ref{algebra-lemma-local-flat-ff}）。
因此，$H_1(L_{S'/R'})_{\mathfrak q'}$ 的消失蕴含
$H_1(L_{S/R})_{\mathfrak q}$ 的消失。
% P01-E0832: the frozen proof again switches from \NL to the unexplained L notation.
把引理 \ref{algebra-lemma-finite-projective-descends} 应用于
$S_{\mathfrak q}$-模 $(\Omega_{S/R})_{\mathfrak q}$
以及映射 $S_{\mathfrak q} \to S'_{\mathfrak q'}$，可知
$(\Omega_{S/R})_{\mathfrak q}$ 是投射模。
故由引理 \ref{algebra-lemma-smooth-at-point}，$R \to S$
在 $\mathfrak q$ 处光滑。
\end{proof}

\begin{lemma}
\label{algebra-lemma-smooth-field-change-local}
设 $K/k$ 是域扩张。
设 $S$ 是有限型 $k$-代数。
设 $\mathfrak q_K$ 是 $S_K = K \otimes_k S$ 的素理想，
并令 $\mathfrak q$ 为 $S$ 中对应的素理想。
则 $S$ 在 $\mathfrak q$ 处是 $k$ 上光滑的，当且仅当
$S_K$ 在 $\mathfrak q_K$ 处是 $K$ 上光滑的。
\end{lemma}

\begin{proof}
这是引理 \ref{algebra-lemma-flat-base-change-locus-smooth} 的特殊情形。
\end{proof}

\begin{lemma}
\label{algebra-lemma-lift-smooth}
设 $R$ 是环，$I \subset R$ 是理想。
设 $R/I \to \overline{S}$ 是光滑环映射。
则存在生成 $\overline{S}$ 的单位理想的元素
$\overline{g}_i \in \overline{S}$，使得每个
% P01-E0833: the frozen “there exists elements” has subject-verb disagreement.
$\overline{S}_{g_i} \cong S_i/IS_i$，其中 $S_i$ 是
$R$ 上的某个（标准）光滑环。
% P01-E0834: the frozen localization drops the bar from the only introduced localizer, and the proof repeats the mismatch.
\end{lemma}

\begin{proof}
由引理 \ref{algebra-lemma-smooth-syntomic}，可以找到生成
$\overline{S}$ 的单位理想的一族元素
$\overline{g}_i \in \overline{S}$，使得每个
$\overline{S}_{g_i}$ 都是 $R/I$ 上的标准光滑代数。
因此可以假设 $\overline{S}$ 是 $R/I$ 上的标准光滑代数。
按照定义 \ref{algebra-definition-standard-smooth}，写成
$\overline{S} =
(R/I)[x_1, \ldots, x_n]/(\overline{f}_1, \ldots, \overline{f}_c)$。
选取 $f_1, \ldots, f_c \in R[x_1, \ldots, x_n]$，
它们分别提升 $\overline{f}_1, \ldots, \overline{f}_c$。置
$S = R[x_1, \ldots, x_n, x_{n + 1}]/(f_1, \ldots, f_c, x_{n + 1}\Delta - 1)$，
其中
$\Delta = \det(\frac{\partial f_j}{\partial x_i})_{i, j = 1, \ldots, c}$，
如例 \ref{algebra-example-make-standard-smooth} 中所述。这证明了引理。
\end{proof}

\section{形式光滑映射}
\label{algebra-section-formally-smooth}

\noindent
本节定义形式光滑环映射。我们将会看到，有限表现的环映射形式光滑，
当且仅当它光滑，见命题 \ref{algebra-proposition-smooth-formally-smooth}。

\begin{definition}
\label{algebra-definition-formally-smooth}
设 $R \to S$ 是环映射。若对每个由实线箭头组成的交换图
$$
\xymatrix{
S \ar[r] \ar@{-->}[rd] & A/I \\
R \ar[r] \ar[u] & A \ar[u]
}
$$
其中 $I \subset A$ 是平方为零的理想，都存在使图交换的虚线箭头，
则称 $S$ 在 $R$ 上\textit{形式光滑}。
\end{definition}

\begin{lemma}
\label{algebra-lemma-base-change-fs}
设 $R \to S$ 是形式光滑环映射，$R \to R'$ 是任意环映射。
则基变换 $S' = R' \otimes_R S$ 在 $R'$ 上形式光滑。
\end{lemma}

\begin{proof}
给定定义 \ref{algebra-definition-formally-smooth} 中的实线图
$$
\xymatrix{
S \ar[r] \ar@{-->}[rrd] & R' \otimes_R S \ar[r] \ar@{-->}[rd] & A/I \\
R  \ar[u] \ar[r] & R' \ar[r] \ar[u] & A \ar[u]
}
$$
依假设，较长的虚线箭头存在。由张量积的泛性质，得到较短的虚线箭头。
\end{proof}

\begin{lemma}
\label{algebra-lemma-compose-formally-smooth}
形式光滑环映射的合成形式光滑。
\end{lemma}

\begin{proof}
略。（提示：这是完全形式的，由考察一个适当的图即可得到。）
\end{proof}

\begin{lemma}
\label{algebra-lemma-polynomial-ring-formally-smooth}
$R$ 上的多项式环在 $R$ 上形式光滑。
\end{lemma}

\begin{proof}
假设有定义 \ref{algebra-definition-formally-smooth} 中的图，且
$S = R[x_j; j \in J]$。只需对上方水平箭头下每个 $x_j$ 在 $A/I$
中的像选取一个提升 $a_j \in A$，便得到一条虚线箭头。
% P01-E0835: the frozen English has the redundant construction “to which ... map to”.
\end{proof}

\begin{lemma}
\label{algebra-lemma-characterize-formally-smooth}
设 $R \to S$ 是环映射。设 $P \to S$ 是从多项式环 $P$ 到 $S$ 的满
$R$-代数映射，并以 $J \subset P$ 表示其核。则 $R \to S$ 形式光滑，
当且仅当存在 $R$-代数映射 $
\sigma : S \to P/J^2$，它是满射 $P/J^2 \to S$ 的右逆。
% P01-E0836: the frozen English “Denote J ... the kernel” omits “by/as”.
\end{lemma}

\begin{proof}
假设 $R \to S$ 形式光滑。考察交换图
$$
\xymatrix{
S \ar[r] \ar@{-->}[rd] & P/J \\
R \ar[r] \ar[u] &  P/J^2\ar[u]
}
$$
依假设，虚线箭头存在。这证明了 $
\sigma$ 存在。

\medskip\noindent
反之，假设有引理中的 $
\sigma$。给定定义 \ref{algebra-definition-formally-smooth} 中的实线图
$$
\xymatrix{
S \ar[r] \ar@{-->}[rd] & A/I \\
R \ar[r] \ar[u] & A \ar[u]
}
$$
由引理 \ref{algebra-lemma-polynomial-ring-formally-smooth}，$P$ 形式光滑，
故存在提升映射 $P \to S \to A/I$ 的 $R$-代数同态
$\psi : P \to A$。显然 $
\psi(J) \subset I$，而 $I^2 = 0$，所以 $
\psi(J^2) = 0$。因而 $
\psi$ 分解为 $
\overline{\psi} : P/J^2 \to A$。所求虚线箭头即合成
$\overline{\psi} \circ \sigma : S \to A$。
\end{proof}

\begin{remark}
\label{algebra-remark-lemma-characterize-formally-smooth}
只要 $P$ 在 $R$ 上形式光滑，引理
\ref{algebra-lemma-characterize-formally-smooth} 便仍然成立。
\end{remark}

\begin{lemma}
\label{algebra-lemma-characterize-formally-smooth-again}
设 $R \to S$ 是环映射。设 $P \to S$ 是从多项式环 $P$ 到 $S$ 的满
$R$-代数映射，并以 $J \subset P$ 表示其核。则 $R \to S$ 形式光滑，
当且仅当引理 \ref{algebra-lemma-differential-seq} 中的序列
$$
0 \to J/J^2 \to \Omega_{P/R} \otimes_P S \to \Omega_{S/R} \to 0
$$
是分裂正合列。
% P01-E0837: the second frozen “Denote J ... the kernel” again omits “by/as”.
\end{lemma}

\begin{proof}
假设 $S$ 在 $R$ 上形式光滑。由引理
\ref{algebra-lemma-characterize-formally-smooth}，这意味着存在 $R$-代数映射
$S \to P/J^2$，它是典范映射 $P/J^2 \to S$ 的右逆。由引理
\ref{algebra-lemma-differential-mod-power-ideal}，有
$\Omega_{P/R} \otimes_P S = \Omega_{(P/J^2)/R} \otimes_{P/J^2} S$。
由引理 \ref{algebra-lemma-differential-seq-split}，该序列分裂。

\medskip\noindent
假设引理中的正合列分裂。选取分裂映射
$\sigma : \Omega_{S/R} \to \Omega_{P/R} \otimes_P S$。
对每个 $
\lambda \in S$，选取映到 $
\lambda$ 的 $x_\lambda \in P$。
然后对每个 $
\lambda \in S$，选取 $f_\lambda \in J$，使得
$$
\text{d}f_\lambda = \text{d}x_\lambda - \sigma(\text{d}\lambda)
$$
在正合列的中项中成立。我们断言，映射
$s : \lambda \mapsto x_\lambda - f_\lambda \mod J^2$ 是 $R$-代数同态
$s : S \to P/J^2$。为证明这一点，我们将反复使用如下事实：若
$h \in J$，且 $
\text{d}h = 0$ 在 $
\Omega_{P/R} \otimes_R S$ 中成立，则 $h \in J^2$。
% P01-E0838: the frozen tensor base is R here, while the exact sequence's middle term is tensored over P.
设 $
\lambda, \mu \in S$。则
$\sigma(\text{d}\lambda + \text{d}\mu - \text{d}(\lambda + \mu)) = 0$。
这推出
$$
\text{d}(x_\lambda + x_\mu - x_{\lambda + \mu}
- f_\lambda - f_\mu + f_{\lambda + \mu}) = 0
$$
因而 $x_\lambda + x_\mu - x_{\lambda + \mu}
- f_\lambda - f_\mu + f_{\lambda + \mu} \in J^2$，进而
$s(\lambda) + s(\mu) = s(\lambda + \mu)$。同理，
$\sigma(\lambda \text{d}\mu + \mu \text{d}\lambda - \text{d}\lambda \mu) = 0$，
所以
% P01-E0839: the frozen final differential reads d lambda mu rather than d(lambda mu).
$$
\mu(\text{d}x_\lambda - \text{d}f_\lambda) +
\lambda(\text{d}x_\mu - \text{d}f_\mu) -
\text{d}x_{\lambda\mu} + \text{d}f_{\lambda\mu} = 0
$$
在正合列的中项中成立。此外，在该中项中还有
$$
\text{d}(x_\lambda x_\mu) =
x_\lambda \text{d}x_\mu + x_\mu \text{d}x_\lambda =
\lambda \text{d}x_\mu + \mu \text{d} x_\lambda
$$
合并这些等式可知
$x_\lambda x_\mu - x_{\lambda\mu}
- \mu f_\lambda - \lambda f_\mu + f_{\lambda\mu} \in J^2$，
从而由于 $f_\lambda f_\mu \in J^2$，有
$(x_\lambda - f_\lambda)(x_\mu - f_\mu) -
(x_{\lambda\mu} - f_{\lambda\mu}) \in J^2$；这就是说
$s(\lambda)s(\mu) = s(\lambda\mu)$。若 $
\lambda \in R$，则 $
\text{d}\lambda = 0$，于是 $
\text{d}f_\lambda = \text{d}x_\lambda$，故
$\lambda - x_\lambda + f_\lambda \in J^2$，进而如所需有
$s(\lambda) = \lambda$。现在应用引理
\ref{algebra-lemma-characterize-formally-smooth}，便得 $S/R$ 形式光滑。
\end{proof}

\begin{proposition}
\label{algebra-proposition-characterize-formally-smooth}
设 $R \to S$ 是环映射。考察形式光滑的 $R$-代数 $P$ 以及核为 $J$ 的满射
$P \to S$。下列条件等价：
\begin{enumerate}
\item $S$ 在 $R$ 上形式光滑；
\item 对某个上述 $P \to S$，映射 $P/J^2 \to S$ 有截面；
\item 对每个上述 $P \to S$，映射 $P/J^2 \to S$ 有截面；
\item 对某个上述 $P \to S$，序列
$0 \to J/J^2 \to \Omega_{P/R} \otimes S \to \Omega_{S/R} \to 0$ 分裂正合；
\item 对每个上述 $P \to S$，序列
$0 \to J/J^2 \to \Omega_{P/R} \otimes S \to \Omega_{S/R} \to 0$ 分裂正合；以及
% P01-E0840: the two frozen tensor products omit the base P used everywhere else in the section.
\item 朴素余切复形 $
\NL_{S/R}$ 拟同构于置于 $0$ 次的投射 $S$-模：
这意味着 $H_1(\NL_{S/R}) = 0$，且 $
\Omega_{S/R}$ 是投射 $S$-模。
\end{enumerate}
\end{proposition}

\begin{proof}
显然 (1) 推出 (3)，(3) 推出 (2)，见引理
\ref{algebra-lemma-characterize-formally-smooth} 证明的第一部分。
又有 (3) 推出 (5)，(5) 推出 (4)，且 (2) 推出 (4)，见引理
\ref{algebra-lemma-characterize-formally-smooth-again} 证明的第一部分。
最后，把引理 \ref{algebra-lemma-characterize-formally-smooth-again}
应用于典范满射 $R[S] \to S$（\ref{algebra-equation-canonical-presentation}），
便知 (1) 推出 (6)。

\medskip\noindent
假设 (4)，来证明 (6)。考察与环映射 $R \to P \to S$ 相关的引理
\ref{algebra-lemma-exact-sequence-NL} 中的序列。由上面已证明的蕴含
(1) $\Rightarrow$ (6)，可知 $
\NL_{P/R} \otimes_R S$ 拟同构于置于 $0$ 次的
$\Omega_{P/R} \otimes_P S$。因此
$H_1(\NL_{P/R} \otimes_P S) = 0$。由于 $P \to S$ 满，
由引理 \ref{algebra-lemma-NL-surjection}，$
\NL_{S/P}$ 同伦等价于置于 $1$ 次的 $J/J^2$。于是得到正合列
$0 \to H_1(L_{S/R}) \to J/J^2 \to \Omega_{P/R} \otimes_P S \to
\Omega_{S/R} \to 0$。由假设可知 $H_1(L_{S/R}) = 0$，且
$\Omega_{S/R}$ 是投射 $S$-模。因此 (6) 成立。
% P01-E0841: the first frozen tensor base is R, while the following occurrence is P.
% P01-E0842: the frozen proof switches from NL to unexplained L in both H_1 terms.

\medskip\noindent
最后证明 (6) 推出 (1)。假设意味着复形
$J/J^2 \to \Omega_{P/R} \otimes S$ 拟同构于置于 $0$ 次的投射
$S$-模，其中 $P = R[S]$，且 $P \to S$ 是典范满射
（\ref{algebra-equation-canonical-presentation}）。因此引理
\ref{algebra-lemma-characterize-formally-smooth-again} 表明 $S$ 在 $R$ 上形式光滑。
% P01-E0843: this frozen tensor product again omits the base P.
\end{proof}

\begin{lemma}
\label{algebra-lemma-ses-formally-smooth}
设 $A \to B \to C$ 是环映射，并假设 $B \to C$ 形式光滑。
则引理 \ref{algebra-lemma-exact-sequence-differentials} 中的序列
$$
0 \to \Omega_{B/A} \otimes_B C \to \Omega_{C/A} \to \Omega_{C/B} \to 0
$$
是分裂短正合列。
\end{lemma}

\begin{proof}
这由命题 \ref{algebra-proposition-characterize-formally-smooth}
和引理 \ref{algebra-lemma-exact-sequence-NL} 得到。
\end{proof}

\begin{lemma}
\label{algebra-lemma-differential-seq-formally-smooth}
设 $A \to B \to C$ 是环映射，其中 $A \to C$ 形式光滑，
$B \to C$ 满且核为 $J \subset B$。则引理
\ref{algebra-lemma-differential-seq} 中的正合列
$$
0 \to J/J^2 \to \Omega_{B/A} \otimes_B C \to \Omega_{C/A} \to 0
$$
分裂正合。
\end{lemma}

\begin{proof}
这由命题 \ref{algebra-proposition-characterize-formally-smooth}、
引理 \ref{algebra-lemma-exact-sequence-NL} 和引理
\ref{algebra-lemma-differential-seq} 得到。
\end{proof}

\begin{lemma}
\label{algebra-lemma-application-NL-formally-smooth}
设 $A \to B \to C$ 是环映射。假设 $A \to C$ 满（因而
$B \to C$ 也满），且 $A \to B$ 形式光滑。记
$I = \Ker(A \to C)$ 和 $J = \Ker(B \to C)$。则引理
\ref{algebra-lemma-application-NL} 中的序列
$$
0 \to I/I^2 \to J/J^2 \to \Omega_{B/A} \otimes_B B/J \to 0
$$
分裂正合。
% P01-E0844: the frozen “Denote I = ... and J = ...” omits “by/as”.
\end{lemma}

\begin{proof}
由于 $A \to B$ 形式光滑，存在环映射
$\sigma : B \to A/I^2$，使它与 $A \to B$ 的合成等于商映射
$A \to A/I^2$。于是 $\sigma$ 诱导映射 $J/J^2 \to I/I^2$，
它是映射 $I/I^2 \to J/J^2$ 的逆。
\end{proof}

\begin{lemma}
\label{algebra-lemma-lift-formal-smoothness}
设 $R \to S$ 是环映射，$I \subset R$ 是理想。假设
\begin{enumerate}
\item $I^2 = 0$；
\item $R \to S$ 平坦；并且
\item $R/I \to S/IS$ 形式光滑。
\end{enumerate}
则 $R \to S$ 形式光滑。
\end{lemma}

\begin{proof}
假设 (1)、(2) 和 (3)。取 $R$-代数满射
$P = R[\{x_t\}_{t \in T}] \to S$，核为 $J$。于是
$0 \to J \to P \to S \to 0$ 是平坦 $R$-模的短正合列。
这推出 $I \otimes_R S = IS$、$I \otimes_R P = IP$ 和
$I \otimes_R J = IJ$，以及 $J \cap IP = IJ$。在整个证明中我们将使用
$$
\Omega_{(S/IS)/(R/I)} = \Omega_{S/R} \otimes_S (S/IS)
= \Omega_{S/R} \otimes_R R/I = \Omega_{S/R} / I\Omega_{S/R}
$$
对 $P$ 也有类似结论（见引理 \ref{algebra-lemma-differentials-base-change}）。
由引理 \ref{algebra-lemma-characterize-formally-smooth-again}，序列
\begin{equation}
\label{algebra-equation-split}
0 \to J/(IJ + J^2) \to
\Omega_{P/R} \otimes_P S/IS \to
\Omega_{S/R} \otimes_S S/IS \to 0
\end{equation}
分裂正合。当然，中项是
$\bigoplus_{t \in T} S/IS \text{d}x_t$。选取分裂映射
$\sigma : \Omega_{P/R} \otimes_P S/IS \to J/(IJ + J^2)$。
对每个 $t \in T$，选取 $f_t \in J$，使它在 $J/(IJ + J^2)$
中的像为 $\sigma(\text{d}x_t)$。这唯一确定了 $S$-模映射
$$
\tilde \sigma : \Omega_{P/R} \otimes_R S
= \bigoplus S\text{d}x_t \longrightarrow J/J^2
$$
并满足 $\tilde\sigma(\text{d}x_t) = f_t$。
% P01-E0845: the frozen tensor base is R although the displayed direct-sum identification requires tensoring over P.
由于 $\sigma$ 是 $\text{d}$ 的截面，差
$$
\Delta = \text{id}_{J/J^2} - \tilde \sigma \circ \text{d}
$$
是自映射 $J/J^2 \to J/J^2$，其像包含在
$(IJ + J^2)/J^2$ 中。特别地，因为 $I^2 = 0$，有
$\Delta((IJ + J^2)/J^2) = 0$。这意味着 $\Delta$ 分解为
$$
J/J^2 \to J/(IJ + J^2) \xrightarrow{\overline{\Delta}}
(IJ + J^2)/J^2 \to J/J^2
$$
其中 $\overline{\Delta}$ 是 $S/IS$-模映射。再次利用序列
（\ref{algebra-equation-split}）分裂，可以找到 $S/IS$-模映射
$\overline{\delta} : \Omega_{P/R} \otimes_P S/IS \to (IJ + J^2)/J^2$，
使得 $\overline{\delta} \circ d$ 等于 $\overline{\Delta}$。
与上面相同，映射 $\overline{\delta}$ 确定一个 $S$-模映射
$\delta : \Omega_{P/R} \otimes_P S \to J/J^2$。
% P01-E0846: the frozen source uses bare d here, unlike the section's \text{d} notation.
把 $\tilde \sigma$ 替换为 $\tilde \sigma + \delta$ 后，简单计算表明
$\Delta = 0$。换言之，$\tilde \sigma$ 是
$J/J^2 \to \Omega_{P/R} \otimes_P S$ 的截面。
由引理 \ref{algebra-lemma-characterize-formally-smooth-again}，
可知 $R \to S$ 形式光滑。
\end{proof}

\begin{proposition}
\label{algebra-proposition-smooth-formally-smooth}
设 $R \to S$ 是环映射。下列条件等价：
% P01-E0847: the frozen proposition omits the colon before its enumeration.
\begin{enumerate}
\item $R \to S$ 是有限表现且形式光滑的；
\item $R \to S$ 光滑。
\end{enumerate}
\end{proposition}

\begin{proof}
这由命题 \ref{algebra-proposition-characterize-formally-smooth}
和定义 \ref{algebra-definition-smooth} 得到。
（注意，若 $R \to S$ 是有限表现的，则 $\Omega_{S/R}$ 是有限表现的
$S$-模，见引理 \ref{algebra-lemma-differentials-finitely-presented}。）
\end{proof}

\begin{lemma}
\label{algebra-lemma-finite-presentation-fs-Noetherian}
设 $R \to S$ 是光滑环映射。则存在在 $\mathbf{Z}$ 上有限型的子环
$R_0 \subset R$ 和光滑环映射 $R_0 \to S_0$，使得
$S \cong R \otimes_{R_0} S_0$。
\end{lemma}

\begin{proof}
我们将使用光滑等价于有限表现且形式光滑这一事实，见命题
\ref{algebra-proposition-smooth-formally-smooth}。写成
$S = R[x_1, \ldots, x_n]/(f_1, \ldots, f_m)$，并记
$I = (f_1, \ldots, f_m)$。按引理
\ref{algebra-lemma-characterize-formally-smooth}，选取投影到 $S$ 的右逆
$\sigma : S \to R[x_1, \ldots, x_n]/I^2$。选取
$h_i \in R[x_1, \ldots, x_n]$，使
$\sigma(x_i \bmod I) = h_i \bmod I^2$。由于 $x_i - h_i \in I$，
存在 $b_{ij} \in R[x_1, \ldots, x_n]$，使得
$$
x_i - h_i = \sum\nolimits_j b_{ij} f_j
$$
$\sigma$ 是 $R$-代数同态
$R[x_1, \ldots, x_n]/I \to R[x_1, \ldots, x_n]/I^2$
这一事实等价于如下条件：对某些
$a_{kl} \in R[x_1, \ldots, x_n]$，有
$$
f_j(h_1, \ldots, h_n) = \sum\nolimits_{j_1 j_2} a_{j_1 j_2} f_{j_1} f_{j_2}
$$
令 $R_0 \subset R$ 为所有多项式 $f_j, h_i, a_{kl}, b_{ij}$ 的系数
在 $\mathbf{Z}$ 上生成的子环。定义：
$$
S_0 = R_0[x_1, \ldots, x_n]/(f_1, \ldots, f_m)
$$
并令：
$$
I_0 = (f_1, \ldots, f_m)
$$
由于第二个陈列等式在
$R_0[x_1, \ldots, x_n]$ 中成立，可以令
$$
\sigma_0 : S_0 \to R_0[x_1, \ldots, x_n]/I_0^2
$$
为由规则 $x_i \mapsto h_i \bmod I_0^2$ 定义的 $R_0$-代数映射。
由于第一个陈列等式在 $R_0[x_1, \ldots, x_n]$ 中成立，
可知 $\sigma_0$ 是投影
$$
R_0[x_1, \ldots, x_n] / I_0^2 \to R_0[x_1, \ldots, x_n] / I_0 = S_0
$$
的右逆。因此由引理 \ref{algebra-lemma-characterize-formally-smooth}，
环 $S_0$ 在 $R_0$ 上形式光滑。
\end{proof}

\begin{lemma}
\label{algebra-lemma-smooth-descends-through-colimit}
设 $A = \colim A_i$ 是环的滤过余极限，$A \to B$ 是光滑环映射。
则存在某个 $i$ 和光滑环映射 $A_i \to B_i$，使得
$B = B_i \otimes_{A_i} A$。
\end{lemma}

\begin{proof}
这由引理 \ref{algebra-lemma-finite-presentation-fs-Noetherian} 得到，
因为由引理 \ref{algebra-lemma-characterize-finite-presentation}，
对某个 $i$，$R_0 \to A$ 将分解经过 $A_i$。
\end{proof}

\begin{lemma}
\label{algebra-lemma-descent-formally-smooth}
设 $R \to S$ 是环映射，$R \to R'$ 是忠实平坦环映射。置
$S' = S \otimes_R R'$。则 $R \to S$ 形式光滑，当且仅当
$R' \to S'$ 形式光滑。
\end{lemma}

\begin{proof}
若 $R \to S$ 形式光滑，则由引理 \ref{algebra-lemma-base-change-fs}，
$R' \to S'$ 形式光滑。为证逆命题，假设 $R' \to S'$ 形式光滑。
注意，对任意 $S$-模 $N$，有
$N \otimes_R R' = N \otimes_S S'$。特别地，$S \to S'$ 也忠实平坦。
选取多项式环 $P = R[\{x_i\}_{i \in I}]$ 以及核为 $J$ 的
$R$-代数满射 $P \to S$。注意，$P' = P \otimes_R R'$ 是 $R'$ 上的
多项式代数。由于 $R \to R'$ 平坦，满射 $P' \to S'$ 的核
$J'$ 是 $J \otimes_R R'$。因此，分裂正合列（见引理
\ref{algebra-lemma-characterize-formally-smooth-again}）
$$
0 \to J'/(J')^2 \to \Omega_{P'/R'} \otimes_{P'} S' \to \Omega_{S'/R'} \to 0
$$
是相应序列
$$
J/J^2 \to \Omega_{P/R} \otimes_P S \to \Omega_{S/R} \to 0
$$
经 $S \to S'$ 的基变换，见引理
\ref{algebra-lemma-differential-seq}。
由于 $S \to S'$ 忠实平坦，得到两件事：
(1) 这个不带 ${}'$ 的序列也正合；(2) $\Omega_{S/R}$ 是投射 $S$-模。
事实上，$\Omega_{S'/R'}$ 作为自由模
$\Omega_{P'/R'} \otimes_{P'} S'$ 的直和因子是投射的，并且由上文有
$\Omega_{S/R} \otimes_S {S'} = \Omega_{S'/R'}$。因此 (2) 由投射性沿忠实平坦环
映射的下降得到，见定理 \ref{algebra-theorem-ffdescent-projectivity}。
所以序列
$0 \to J/J^2 \to \Omega_{P/R} \otimes_P S \to \Omega_{S/R} \to 0$
也正合，再次应用引理
\ref{algebra-lemma-characterize-formally-smooth-again} 即得结论。
\end{proof}

\noindent
事实证明，光滑环映射满足下面这个很强的提升性质。

\begin{lemma}
\label{algebra-lemma-smooth-strong-lift}
设 $R \to S$ 是光滑环映射。给定由实线箭头组成的交换图
$$
\xymatrix{
S \ar[r] \ar@{-->}[rd] & A/I \\
R \ar[r] \ar[u] & A \ar[u]
}
$$
其中 $I \subset A$ 是局部幂零理想，则存在使图交换的虚线箭头。
\end{lemma}

\begin{proof}
由引理 \ref{algebra-lemma-finite-presentation-fs-Noetherian}，可把该图扩张成交换图
$$
\xymatrix{
S_0 \ar[r] & S \ar[r] \ar@{-->}[rd] & A/I \\
R_0 \ar[r] \ar[u] & R \ar[r] \ar[u] & A \ar[u]
}
$$
其中 $R_0 \to S_0$ 光滑，$R_0$ 在 $\mathbf{Z}$ 上有限型，且
$S = S_0 \otimes_{R_0} R$。设 $x_1, \ldots, x_n \in S_0$
是在 $R_0$ 上生成 $S_0$ 的元。设 $a_1, \ldots, a_n$ 是 $A$ 中的元，
它们在 $A/I$ 中的像与 $x_1, \ldots, x_n$ 的像分别相同。
记 $A_0 \subset A$ 为由 $R_0$ 的像和元 $a_1, \ldots, a_n$ 生成的子环。
置 $I_0 = A_0 \cap I$。则 $A_0/I_0 \subset A/I$，且
$S_0 \to A/I$ 映入 $A_0/I_0$。因此，只需在图
% P01-E0848: the frozen “Denote A_0 ... the subring” omits “by/as”.
$$
\xymatrix{
S_0 \ar[r] \ar@{-->}[rd] & A_0/I_0 \\
R_0 \ar[r] \ar[u] & A_0 \ar[u]
}
$$
中找到虚线箭头。按构造，环 $A_0$ 在 $\mathbf{Z}$ 上有限型。
故 $A_0$ 是 Noether 环，从而 $I_0$ 幂零，见引理
\ref{algebra-lemma-Noetherian-power}。设 $I_0^n = 0$。由命题
\ref{algebra-proposition-smooth-formally-smooth}，可以依次把 $R_0$-代数映射
$S_0 \to A_0/I_0$ 提升为
$S_0 \to A_0/I_0^2$、$S_0 \to A_0/I_0^3$、$\ldots$，
最终提升为 $S_0 \to A_0/I_0^n = A_0$。
\end{proof}

\section{光滑性与微分}
\label{algebra-section-smooth-differential}

\noindent
下面给出一些关于微分与光滑环映射的结果。

\begin{lemma}
\label{algebra-lemma-triangle-differentials-smooth}
设有环映射 $A \to B \to C$，且 $B \to C$ 光滑，则引理
\ref{algebra-lemma-exact-sequence-differentials} 中的序列
$$
0 \to C \otimes_B \Omega_{B/A} \to \Omega_{C/A} \to \Omega_{C/B} \to 0
$$
正合。
% P01-E0849: the frozen English combines “Given ...” with an unnecessary “then”.
\end{lemma}

\begin{proof}
这由更一般的引理 \ref{algebra-lemma-ses-formally-smooth} 得到，
因为光滑环映射形式光滑，见命题
\ref{algebra-proposition-smooth-formally-smooth}。它也可以直接由引理
\ref{algebra-lemma-exact-sequence-NL} 得到，因为
$H_1(L_{C/B}) = 0$ 是 $B \to C$ 光滑定义的一部分。
% P01-E0850: the frozen proof switches from the section's NL notation to unexplained L.
\end{proof}

\begin{lemma}
\label{algebra-lemma-differential-seq-smooth}
设 $A \to B \to C$ 是环映射，其中 $A \to C$ 光滑，
$B \to C$ 满且核为 $J \subset B$。则引理
\ref{algebra-lemma-differential-seq} 中的正合列
$$
0 \to J/J^2 \to \Omega_{B/A} \otimes_B C \to \Omega_{C/A} \to 0
$$
分裂正合。
\end{lemma}

\begin{proof}
这由更一般的引理
\ref{algebra-lemma-differential-seq-formally-smooth} 得到，
因为光滑环映射形式光滑，见命题
\ref{algebra-proposition-smooth-formally-smooth}。
\end{proof}

\begin{lemma}
\label{algebra-lemma-application-NL-smooth}
设 $A \to B \to C$ 是环映射。假设 $A \to C$ 满（因而
$B \to C$ 也满），且 $A \to B$ 光滑。记
$I = \Ker(A \to C)$ 和 $J = \Ker(B \to C)$。则引理
\ref{algebra-lemma-application-NL} 中的序列
$$
0 \to I/I^2 \to J/J^2 \to \Omega_{B/A} \otimes_B B/J \to 0
$$
正合。
% P01-E0851: the frozen “Denote I = ... and J = ...” omits “by/as”.
\end{lemma}

\begin{proof}
这由更一般的引理
\ref{algebra-lemma-application-NL-formally-smooth} 得到，
因为光滑环映射形式光滑，见命题
\ref{algebra-proposition-smooth-formally-smooth}。
\end{proof}

\begin{lemma}
\label{algebra-lemma-section-smooth}
\begin{slogan}
若 $R$ 是 $S$ 的直和因子，且 $S$ 在 $R$ 上光滑，则 $S$ 的
$I$-进完备化通常是 $R$ 上的幂级数环，其中 $I$ 是从 $S$ 到 $R$
的投影映射的核。
% P01-E0852: the frozen slogan says “a power series over R” rather than a power series ring.
% P01-E0853: the frozen slogan lacks terminal punctuation.
\end{slogan}
设 $\varphi : R \to S$ 是光滑环映射，$
\sigma : S \to R$ 是 $
\varphi$ 的左逆。置 $I = \Ker(\sigma)$。则
\begin{enumerate}
\item $I/I^2$ 是有限局部自由 $R$-模；并且
\item 若 $I/I^2$ 自由，则作为 $R$-代数有
$S^\wedge \cong R[[t_1, \ldots, t_d]]$，其中 $S^\wedge$ 是 $S$ 的
$I$-进完备化。
\end{enumerate}
\end{lemma}

\begin{proof}
把引理 \ref{algebra-lemma-differential-seq-split} 应用于 $R \to S \to R$，
可知 $I/I^2 = \Omega_{S/R} \otimes_{S, \sigma} R$。
按光滑态射的定义，模 $
\Omega_{S/R}$ 在 $S$ 上有限局部自由，故 (1) 成立。
若 $I/I^2$ 自由，则选取 $f_1, \ldots, f_d \in I$，使它们在
$I/I^2$ 中的像构成一个 $R$-基。考察由下式定义的 $R$-代数映射：
$$
\Psi : R[[x_1, \ldots, x_d]] \longrightarrow S^\wedge, \quad
x_i \longmapsto f_i.
$$
记 $P = R[[x_1, \ldots, x_d]]$ 和 $J = (x_1, \ldots, x_d) \subset P$。
以 $
\Psi_n : P/J^n \to S/I^n$ 表示诱导的商环映射。注意
$S/I^2 = \varphi(R) \oplus I/I^2$，故 $
\Psi_2$ 是同构。记
$\sigma_2 : S/I^2 \to P/J^2$ 为 $
\Psi_2$ 的逆。
% P01-E0854: the frozen “Denote P = ... and J = ...” omits “by/as”.
% P01-E0855: the frozen “Denote sigma_2 ... the inverse” omits “by/as”.
我们将对 $n$ 归纳证明：对每个 $n > 2$，$
\Psi_n$ 都有逆映射
$\sigma_n : S/I^n \to P/J^n$。事实上，由命题
\ref{algebra-proposition-smooth-formally-smooth}，$S$ 在 $R$ 上形式光滑，
所以在 $R$-代数的实线图
$$
\xymatrix{
S \ar@{..>}[r] \ar[rd]_{\sigma_{n - 1}} & P/J^n \ar[d] \\
 & P/J^{n - 1}
}
$$
中，可以填入某个 $R$-代数映射
$\tau : S \to P/J^n$ 作为虚线箭头，使该图交换。它诱导一个
$R$-代数映射 $
\overline{\tau} : S/I^n \to P/J^n$，模 $J^{n - 1}$ 后等于
$\sigma_{n - 1}$。按构造，映射 $
\Psi_n$ 是满射，而现在
$\overline{\tau} \circ \Psi_n$ 是 $P/J^n$ 的 $R$-代数自同态，
它把 $x_i$ 映到 $x_i + \delta_{i, n}$，其中
$\delta_{i, n} \in J^{n - 1}/J^n$。由此可知 $
\Psi_n$ 是同构，
因而它有逆 $
\sigma_n$。这证明了引理。
\end{proof}

\section{域上的光滑代数}
\label{algebra-section-smooth-over-field}

\noindent
警告：一般而言，下面两个引理在非完美域上并不成立。

\begin{lemma}
\label{algebra-lemma-rank-omega}
设 $k$ 是代数闭域，$S$ 是有限型 $k$-代数，
$\mathfrak m \subset S$ 是极大理想。则
$$
\dim_{\kappa(\mathfrak m)} \Omega_{S/k} \otimes_S \kappa(\mathfrak m)
=
\dim_{\kappa(\mathfrak m)} \mathfrak m/\mathfrak m^2.
$$
\end{lemma}

\begin{proof}
考察引理 \ref{algebra-lemma-differential-seq} 中的正合列
$$
\mathfrak m/\mathfrak m^2 \to
\Omega_{S/k} \otimes_S \kappa(\mathfrak m) \to
\Omega_{\kappa(\mathfrak m)/k} \to 0
$$
我们要证明第一个映射是同构。由于 $k$ 代数闭，由定理
\ref{algebra-theorem-nullstellensatz}，合成 $k \to \kappa(\mathfrak m)$
是同构。因此满射 $S \to \kappa(\mathfrak m)$ 作为 $k$-代数映射分裂，
而引理 \ref{algebra-lemma-differential-seq-split} 表明上述序列在左端也正合。
由于 $\Omega_{\kappa(\mathfrak m)/k} = 0$，结论成立。
\end{proof}

\begin{lemma}
\label{algebra-lemma-characterize-smooth-kbar}
设 $k$ 是代数闭域，$S$ 是有限型 $k$-代数，
$\mathfrak m \subset S$ 是极大理想。下列条件等价：
\begin{enumerate}
\item 环 $S_{\mathfrak m}$ 是正则局部环。
\item 有
$\dim_{\kappa(\mathfrak m)} \Omega_{S/k} \otimes_S \kappa(\mathfrak m)
\leq \dim(S_{\mathfrak m})$。
\item 有
$\dim_{\kappa(\mathfrak m)} \Omega_{S/k} \otimes_S \kappa(\mathfrak m)
= \dim(S_{\mathfrak m})$。
\item 存在 $g \in S$、$g \not \in \mathfrak m$，使得 $S_g$ 在 $k$ 上光滑。
换言之，$S/k$ 在 $\mathfrak m$ 处光滑。
\end{enumerate}
\end{lemma}

\begin{proof}
由引理 \ref{algebra-lemma-rank-omega} 和定义 \ref{algebra-definition-regular}，
(1)、(2) 和 (3) 等价。

\medskip\noindent
假设 $S$ 在 $\mathfrak m$ 处光滑。由引理
\ref{algebra-lemma-smooth-syntomic}，对某个适当的
$g \in S$、$g \not \in \mathfrak m$，$S_g$ 在 $k$ 上标准光滑。
因此由引理 \ref{algebra-lemma-standard-smooth}，
$\Omega_{S_g/k}$ 自由且秩为 $\dim(S_g)$。再由引理
\ref{algebra-lemma-rank-omega}，有
$\dim(S_{\mathfrak m}) = \dim (\mathfrak m/\mathfrak m^2)$，
换言之，$S_\mathfrak m$ 正则。

\medskip\noindent
反之，假设 $S_{\mathfrak m}$ 正则。令
$d = \dim(S_{\mathfrak m}) = \dim \mathfrak m/\mathfrak m^2$。
选取表现 $S = k[x_1, \ldots, x_n]/I$，使对所有 $i$，$x_i$ 都映到
$\mathfrak m$ 的一个元。换言之，
$\mathfrak m'' = (x_1, \ldots, x_n)$ 是
$k[x_1, \ldots, x_n]$ 中相应的极大理想。注意，有序列
$$
I/\mathfrak m''I \to \mathfrak m''/(\mathfrak m'')^2
\to \mathfrak m/(\mathfrak m)^2 \to 0
$$
% P01-E0856: the frozen prose calls this a short exact sequence although the display has no initial 0 and the first map need not be injective.
选取 $c = n - d$ 个元 $f_1, \ldots, f_c \in I$，使它们在
$\mathfrak m''/(\mathfrak m'')^2$ 中的像张成到
$\mathfrak m/\mathfrak m^2$ 的映射的核。这显然可以做到。记
$J = (f_1, \ldots, f_c)$，于是 $J \subset I$。置
$S' = k[x_1, \ldots, x_n]/J$，从而有满射 $S' \to S$。记
$\mathfrak m' = \mathfrak m''S'$ 为 $S'$ 中相应的极大理想。于是有
% P01-E0857: three consecutive frozen “Denote ...” clauses omit “by/as”.
$$
\xymatrix{
k[x_1, \ldots, x_n] \ar[r] & S' \ar[r] & S \\
\mathfrak m'' \ar[u] \ar[r] & \mathfrak m' \ar[r] \ar[u] &
\mathfrak m \ar[u]
}
$$
由 $J$ 的选取，正合列
$$
J/\mathfrak m''J \to \mathfrak m''/(\mathfrak m'')^2
\to \mathfrak m'/(\mathfrak m')^2 \to 0
$$
表明 $\dim( \mathfrak m'/(\mathfrak m')^2 ) = d$。
由于 $S'_{\mathfrak m'}$ 满射到 $S_{\mathfrak m}$，可知
$\dim(S_{\mathfrak m'}) \geq d$。因此由定义
\ref{algebra-definition-regular-local} 前的讨论，$S'_{\mathfrak m'}$
也正则且维数为 $d$。
% P01-E0858: the frozen dimension term localizes S at m', although m' is an ideal of S'; the surrounding argument requires S'_{m'}.
由于 $S'$ 由 $c = n - d$ 个方程切出，由引理
\ref{algebra-lemma-lci}，存在 $g' \in S'$、$g' \not \in \mathfrak m'$，
使 $S'_{g'}$ 是 $k$ 上的全局完全交。映射
$S'_{\mathfrak m'} \to S_{\mathfrak m}$ 又是相同维数的 Noether 局部整环
之间的满射，故为同构。因此 $S' \to S$ 是有限生成核的满射，
且在 $\mathfrak m'$ 处局部化后成为同构。于是可以找到
$g' \in S'$、$g \not \in \mathfrak m'$，使
$S'_{g'} \to S_{g'}$ 是同构。综上，把 $S$ 替换为某个主局部化后，
可以假设 $S$ 是全局完全交。
% P01-E0859: the frozen localization sentence introduces g' but states g notin m', and then localizes S at g' although g' belongs to S'.

\medskip\noindent
此时可写
$S = k[x_1, \ldots, x_n]/(f_1, \ldots, f_c)$，且
$\dim S = n - c$。回顾一下，这个代数的朴素余切复形为
$$
\bigoplus S \cdot f_j
\to
\bigoplus S \cdot \text{d}x_i
$$
见引理 \ref{algebra-lemma-relative-global-complete-intersection-conormal}。
由引理 \ref{algebra-lemma-relative-global-complete-intersection-smooth}，
为证明 $S$ 在 $\mathfrak m$ 处光滑，只需证明给出上述映射的矩阵
``$A$'' 的某个 $c \times c$ 子式 $g_I$ 在 $\mathfrak m$ 处不消失。
由引理 \ref{algebra-lemma-rank-omega}，矩阵 $A \bmod \mathfrak m$ 的秩为 $c$。
因此结论成立。
\end{proof}

\begin{lemma}
\label{algebra-lemma-characterize-smooth-over-field}
设 $k$ 是任意域，$S$ 是有限型 $k$-代数，$X = \Spec(S)$。
设素理想 $\mathfrak q \subset S$ 对应于 $x \in X$。下列条件等价：
\begin{enumerate}
\item $k$-代数 $S$ 在 $\mathfrak q$ 处相对于 $k$ 光滑。
\item 有
$\dim_{\kappa(\mathfrak q)} \Omega_{S/k} \otimes_S \kappa(\mathfrak q)
\leq \dim_x X$。
\item 有
$\dim_{\kappa(\mathfrak q)} \Omega_{S/k} \otimes_S \kappa(\mathfrak q)
= \dim_x X$。
\end{enumerate}
此外，在这些条件成立时，局部环 $S_{\mathfrak q}$ 正则。
\end{lemma}

\begin{proof}
若 $S$ 在 $\mathfrak q$ 处相对于 $k$ 光滑，则由引理
\ref{algebra-lemma-smooth-syntomic}，存在
$g \in S$、$g \not \in \mathfrak q$，使 $S_g$ 在 $k$ 上标准光滑。
$k$ 上的标准光滑代数，其微分模自由且秩等于维数，见引理
\ref{algebra-lemma-standard-smooth}（这里使用了域上的相对全局完全交的维数
等于变量数减去方程数这一事实）。因此 (1) 推出 (3)。
为完成引理的证明，只需证明 (2) 推出 (1)，并推出
$S_{\mathfrak q}$ 正则。

\medskip\noindent
假设 (2)。由 Nakayama 引理 \ref{algebra-lemma-NAK}，
$\Omega_{S/k, \mathfrak q}$ 可由 $\leq \dim_x X$ 个元生成。
可以把 $S$ 替换为某个 $S_g$，其中
$g \in S$、$g \not \in \mathfrak q$，使 $\Omega_{S/k}$
由至多 $\dim_x X$ 个元生成。取代数闭域扩张 $K/k$，使得存在
$k$-代数映射 $\psi : \kappa(\mathfrak q) \to K$。
考察 $S_K = K \otimes_k S$。令 $\mathfrak m \subset S_K$
为下列满射对应的极大理想：
$$
\xymatrix{
S_K = K \otimes_k S \ar[r] &
K \otimes_k \kappa(\mathfrak q)
\ar[r]^-{\text{id}_K \otimes \psi} &
K.
}
$$
注意，$\mathfrak m \cap S = \mathfrak q$；换言之，
$\mathfrak m$ 位于 $\mathfrak q$ 之上。由引理
\ref{algebra-lemma-dimension-at-a-point-preserved-field-extension}，
$X_K = \Spec(S_K)$ 在 $\mathfrak m$ 对应点处的维数为 $\dim_x X$。
由引理 \ref{algebra-lemma-dimension-closed-point-finite-type-field}，
它等于 $\dim((S_K)_{\mathfrak m})$。由引理
\ref{algebra-lemma-differentials-base-change}，$S_K$ 在 $K$ 上的微分模
是 $\Omega_{S/k}$ 的基变换，因此也可由至多
$\dim_x X = \dim((S_K)_{\mathfrak m})$ 个元生成。由引理
\ref{algebra-lemma-characterize-smooth-kbar}，$S_K$ 在 $\mathfrak m$ 处相对于
$K$ 光滑。由引理 \ref{algebra-lemma-flat-base-change-locus-smooth}，
这推出 $S$ 在 $\mathfrak q$ 处相对于 $k$ 光滑，从而证明了 (1)。
此外，由引理 \ref{algebra-lemma-characterize-smooth-kbar}，局部环
$(S_K)_{\mathfrak m}$ 正则。由于
$S_{\mathfrak q} \to (S_K)_{\mathfrak m}$ 平坦，由引理
\ref{algebra-lemma-flat-under-regular} 可知 $S_{\mathfrak q}$ 正则。
\end{proof}

\noindent
下面的引理可以显著推广（而且有几种不同的推广方式）。

\begin{lemma}
\label{algebra-lemma-computation-differential}
设 $k$ 是域，$R$ 是包含 $k$ 的 Noether 局部环。假设剩余域
$\kappa = R/\mathfrak m$ 是 $k$ 的有限生成可分扩张。则映射
$$
\text{d} :
\mathfrak m/\mathfrak m^2
\longrightarrow
\Omega_{R/k} \otimes_R \kappa(\mathfrak m)
$$
是单射。
\end{lemma}

\begin{proof}
可以把 $R$ 替换为 $R/\mathfrak m^2$，故可假设
$\mathfrak m^2 = 0$。由假设可写
$\kappa = k(\overline{x}_1, \ldots, \overline{x}_r, \overline{y})$，
其中 $\overline{x}_1, \ldots, \overline{x}_r$ 是 $\kappa$ 在 $k$ 上的
超越基，而 $\overline{y}$ 在
$k(\overline{x}_1, \ldots, \overline{x}_r)$ 上可分代数。令
$P(T)$ in $k(\overline{x}_1, \ldots, \overline{x}_r)[T]$ 为
$\overline{y}$ 在 $k(\overline{x}_1, \ldots, \overline{x}_r)$ 上的
最小多项式。于是 $P(\overline{y}) = 0$，但
$P'(\overline{y}) \not = 0$。任取 $x_i \in R$，提升
$\overline{x}_i \in \kappa$。这给出 $k$-代数的交换图
$$
\xymatrix{
R \ar[r] & \kappa \\
& k(\overline{x}_1, \ldots, \overline{x}_r) \ar[lu]^\varphi \ar[u]
}
$$
我们要把向左上的箭头 $\varphi$ 扩张为从 $\kappa$ 到 $R$ 的
$k$-代数映射。为此，任取提升 $\overline{y}$ 的 $y \in R$。
为说明它定义了以
$\kappa \cong k(\overline{x}_1, \ldots, \overline{x}_r)[T]/(P)$
为定义域的 $k$-代数映射，只需证明可以选取 $y$，使
$P^\varphi(y) = 0$。否则，对 $\delta \in \mathfrak m$ 计算得到
$$
P^\varphi(y + \delta) = P^\varphi(y) + (P')^\varphi(y)\delta
$$
因为 $\mathfrak m^2 = 0$。由于 $(P')^\varphi(y)$ 是 $R$ 的单位，
可以按所需调整我们的选择。这表明作为 $k$-代数有
$R \cong \kappa \oplus \mathfrak m$！现在，直接计算
$\Omega_{\kappa \oplus \mathfrak m/k}$，或者应用引理
\ref{algebra-lemma-differential-seq-split}，都可完成证明。
\end{proof}

\begin{lemma}
\label{algebra-lemma-separable-smooth}
设 $k$ 是域，$S$ 是有限型 $k$-代数，$\mathfrak q \subset S$ 是素理想。
假设 $\kappa(\mathfrak q)$ 在 $k$ 上可分。下列条件等价：
\begin{enumerate}
\item 代数 $S$ 在 $\mathfrak q$ 处相对于 $k$ 光滑。
\item 环 $S_{\mathfrak q}$ 正则。
\end{enumerate}
\end{lemma}

\begin{proof}
记 $R = S_{\mathfrak q}$，以 $\mathfrak m$ 表示其极大理想，
以 $\kappa$ 表示其剩余域。由引理
\ref{algebra-lemma-computation-differential} 和
\ref{algebra-lemma-differential-seq}，有短正合列
% P01-E0860: the frozen designation omits “ideal” after “its maximal” and uses an incomplete “denote ... by” construction.
% P01-E0861: the frozen citation says “By Lemma X and [bare reference Y]”, omitting “Lemma” before the second reference.
$$
0 \to \mathfrak m/\mathfrak m^2 \to
\Omega_{R/k} \otimes_R \kappa \to
\Omega_{\kappa/k} \to 0
$$
注意，$\Omega_{R/k} = \Omega_{S/k, \mathfrak q}$，见引理
\ref{algebra-lemma-differentials-localize}。此外，由于 $\kappa$ 在 $k$ 上可分，
有 $\dim_{\kappa} \Omega_{\kappa/k} = \text{trdeg}_k(\kappa)$。
因而得到
$$
\dim_{\kappa} \Omega_{R/k} \otimes_R \kappa
=
\dim_\kappa \mathfrak m/\mathfrak m^2 + \text{trdeg}_k (\kappa)
\geq
\dim R + \text{trdeg}_k (\kappa)
=
\dim_{\mathfrak q} S
$$
（最后一个等式见引理
\ref{algebra-lemma-dimension-at-a-point-finite-type-field}），
且等号成立当且仅当 $R$ 正则。因此应用引理
\ref{algebra-lemma-characterize-smooth-over-field} 即得结论。
% P01-E0862: the frozen parenthetical citation lacks punctuation before the following “with equality” clause.
\end{proof}

\begin{lemma}
\label{algebra-lemma-characteristic-zero}
设 $R \to S$ 是 $\mathbf{Q}$-代数映射。设 $f \in S$，并且对某个
$S$-子模 $C$ 有 $\Omega_{S/R} = S \text{d}f \oplus C$。则
\begin{enumerate}
\item $f$ 不是幂零元；并且
\item 若 $S$ 是 Noether 局部环，则 $f$ 是 $S$ 中的非零因子。
\end{enumerate}
\end{lemma}

\begin{proof}
对 $a \in S$，写成
$\text{d}(a) = \theta(a)\text{d}f + c(a)$，其中
$\theta(a) \in S$ 且 $c(a) \in C$。考察 $R$-导子
$S \to S$、$a \mapsto \theta(a)$。注意，$\theta(f) = 1$。

\medskip\noindent
若 $f^n = 0$，且 $n > 1$ 最小，则
$0 = \theta(f^n) = n f^{n - 1}$，这与 $n$ 的最小性矛盾。
所以 $f$ 不是幂零元。

\medskip\noindent
假设 $fa = 0$。若 $f$ 是单位，则 $a = 0$，结论成立。
假设 $f$ 不是单位。由 Leibniz 法则，
$0 = \theta(fa) = f\theta(a) + a$，故 $a \in (f)$。
归纳地，假设已经证明 $fa = 0 \Rightarrow a \in (f^n)$。
写 $a = f^nb$，便有
$0 = \theta(f^{n + 1}b) = (n + 1)f^nb + f^{n + 1}\theta(b)$。
因此
$a = f^n b = -f^{n + 1}\theta(b)/(n + 1) \in (f^{n + 1})$。
在 Noether 局部环 $S$ 中有 $\bigcap (f^n) = 0$，见引理
\ref{algebra-lemma-intersect-powers-ideal-module-zero}，所以结论成立。
\end{proof}

\noindent
下面这个结论在应用中大概没有什么用。

\begin{lemma}
\label{algebra-lemma-characteristic-zero-local-smooth}
设 $k$ 是特征 $0$ 的域，$S$ 是有限型 $k$-代数，
$\mathfrak q \subset S$ 是素理想。下列条件等价：
\begin{enumerate}
\item 代数 $S$ 在 $\mathfrak q$ 处相对于 $k$ 光滑。
\item $S_{\mathfrak q}$-模 $\Omega_{S/k, \mathfrak q}$（有限）自由。
\item 环 $S_{\mathfrak q}$ 正则。
\end{enumerate}
\end{lemma}

\begin{proof}
在特征零时，任意域扩张都可分，因此 (1) 与 (3) 的等价性由引理
\ref{algebra-lemma-separable-smooth} 得到。又由光滑代数的定义，(1) 推出 (2)。
假设 $\Omega_{S/k, \mathfrak q}$ 在 $S_{\mathfrak q}$ 上自由。
我们将使用引理 \ref{algebra-lemma-separable-smooth} 的证明中所用的记号和观察。
因此 $R = S_{\mathfrak q}$，其极大理想为 $\mathfrak m$，
剩余域为 $\kappa$。目标是证明 $R$ 正则。

\medskip\noindent
若 $\mathfrak m/\mathfrak m^2 = 0$，则 $\mathfrak m = 0$，且
$R \cong \kappa$。所以 $R$ 正则，结论成立。

\medskip\noindent
若 $\mathfrak m/ \mathfrak m^2 \not = 0$，则任取
$f \in \mathfrak m$，使其在 $\mathfrak m/ \mathfrak m^2$ 中的像非零。
由引理 \ref{algebra-lemma-computation-differential}，可知
$\text{d}f$ 在 $\Omega_{R/k}/\mathfrak m\Omega_{R/k}$ 中的像非零。
按假设，$\Omega_{R/k} = \Omega_{S/k, \mathfrak q}$ 有限自由，
因而由 Nakayama 引理 \ref{algebra-lemma-NAK}，可知
$\text{d}f$ 生成一个直和因子。应用引理
\ref{algebra-lemma-characteristic-zero}，得到 $f$ 是 $R$ 中的非零因子。
再由引理 \ref{algebra-lemma-differential-seq}，得到正合列
$$
(f)/(f^2) \to \Omega_{R/k} \otimes_R R/fR \to \Omega_{(R/fR)/k} \to 0
$$
这推出 $\Omega_{(R/fR)/k}$ 也有限自由。由归纳可知
$R/fR$ 是正则局部环。由于 $f \in \mathfrak m$ 是非零因子，
由引理 \ref{algebra-lemma-regular-mod-x} 可知 $R$ 正则。
\end{proof}

\begin{example}
\label{algebra-example-characteristic-p}
引理 \ref{algebra-lemma-characteristic-zero-local-smooth} 在特征
$p > 0$ 时不成立。标准例子是环映射
$$
\mathbf{F}_p \longrightarrow \mathbf{F}_p[x]/(x^p)
$$
其微分模自由，但显然不光滑；另一个例子是环映射（$p > 2$）
$$
\mathbf{F}_p(t) \to \mathbf{F}_p(t)[x, y]/(x^p + y^2 + \alpha)
$$
它在素理想 $\mathfrak q = (y, x^p + \alpha)$ 处不光滑，但正则。
% P01-E0863: alpha is not defined or constrained in the frozen characteristic-p example.
\end{example}

\noindent
利用上面的材料，可以用域扩张刻画一般点处的光滑性。

\begin{lemma}
\label{algebra-lemma-smooth-at-generic-point}
设 $R \to S$ 是单射的有限型环映射，且 $R$、$S$ 都是整环。
则 $R \to S$ 在 $\mathfrak q = (0)$ 处光滑，当且仅当其分式域的诱导扩张
$L/K$ 可分。
\end{lemma}

\begin{proof}
假设 $R \to S$ 在 $(0)$ 处光滑。可以把 $S$ 替换为某个
$S_g$，其中 $g \in S$ 非零，并假设 $R \to S$ 光滑。
于是 $K \to S \otimes_R K$ 光滑（引理
\ref{algebra-lemma-base-change-smooth}）。此外，对任意域扩张
$K'/K$，环映射 $K' \to S \otimes_R K'$ 也光滑。因此由引理
\ref{algebra-lemma-characterize-smooth-over-field}，
$S \otimes_R K'$ 是正则环，特别地是约化环。这表明
$S \otimes_R K$ 在 $K$ 上几何约化。因此由引理
\ref{algebra-lemma-geometrically-reduced-permanence}，$L$ 在 $K$ 上几何约化。
所以由引理 \ref{algebra-lemma-characterize-separable-field-extensions}，
$L/K$ 可分。
% P01-E0864: the frozen sentence has the extra article “a” in “is a geometrically reduced over K”.

\medskip\noindent
反之，假设 $L/K$ 可分。由引理
\ref{algebra-lemma-generic-finite-presentation}，可以假设 $R \to S$ 是有限表现的。
由引理 \ref{algebra-lemma-flat-base-change-locus-smooth}，只需证明
$K \to S \otimes_R K$ 在 $(0)$ 处光滑。这由引理
\ref{algebra-lemma-separable-smooth}、域是正则环这一事实以及 $L/K$ 可分的假设得到。
\end{proof}

\section{Noether 情形的光滑环映射}
\label{algebra-section-smooth-Noetherian}

\begin{definition}
\label{algebra-definition-small-extension}
设 $\varphi : B' \to B$ 是环映射。若 $B'$ 和 $B$ 是局部 Artin 环，
$\varphi$ 满，且 $I = \Ker(\varphi)$ 作为 $B'$-模长度为 $1$，
则称 $\varphi$ 是\textit{小扩张}。
\end{definition}

\noindent
显然，这意味着 $I^2 = 0$，并且对某个
$x \in B'$ 有 $I = (x)$ 及 $\mathfrak m' x = 0$，其中
$\mathfrak m' \subset B'$ 是极大理想。

\begin{lemma}
\label{algebra-lemma-smooth-test-artinian}
设 $R \to S$ 是环映射。设 $\mathfrak q$ 是 $S$ 的素理想，
位于 $\mathfrak p \subset R$ 之上。假设 $R$ 是 Noether 环，
且 $R \to S$ 是有限型的。下列条件等价：
% P01-E0865: the frozen premise omits “is” before “of finite type”.
\begin{enumerate}
\item $R \to S$ 在 $\mathfrak q$ 处光滑；
\item 对每个局部 $R$-代数满射
$(B', \mathfrak m') \to (B, \mathfrak m)$，若
$\Ker(B' \to B)$ 平方为零，并且对每个实线部分交换图
$$
\xymatrix{
S \ar[r] \ar@{-->}[rd] & B \\
R \ar[r] \ar[u] & B' \ar[u]
}
$$
满足 $\mathfrak q = S \cap \mathfrak m$，则存在使图交换的虚线箭头；
% P01-E0866: the frozen equality uses an intersection although S -> B need not be injective; it should state inverse image of m equals q.
\item 与 (2) 相同，但 $B' \to B$ 只遍历小扩张；并且
\item 与 (2) 相同，但 $B' \to B$ 只遍历还满足如下条件的小扩张：
$S \to B$ 诱导同构
$\kappa(\mathfrak q) \cong \kappa(\mathfrak m)$。
\end{enumerate}
\end{lemma}

\begin{proof}
假设 (1)。这意味着存在
$g \in S$、$g \not \in \mathfrak q$，使 $R \to S_g$ 光滑。
由命题 \ref{algebra-proposition-smooth-formally-smooth}，
$R \to S_g$ 形式光滑。注意，给定 (2) 中的任意图，
映射 $S \to B$ 自动分解经过 $S_{\mathfrak q}$，因而更分解经过 $S_g$。
$S_g$ 在 $R$ 上形式光滑，给出一个态射 $S_g \to B'$，
使其嵌入以 $S_g$ 为左上角的类似图中。与 $S \to S_g$ 合成，
便得到所需箭头。换言之，我们已经证明 (1) 推出 (2)。

\medskip\noindent
显然 (2) 推出 (3)，(3) 推出 (4)。

\medskip\noindent
假设 (4)。我们来证明 (1)，从而完成引理的证明。选取表现
$S = R[x_1, \ldots, x_n]/(f_1, \ldots, f_m)$。
这是可行的，因为 $S$ 在 Noether 环 $R$ 上有限型，因而有限表现
（见引理
\ref{algebra-lemma-Noetherian-finite-type-is-finite-presentation}）。
置 $I = (f_1, \ldots, f_m)$。考察这个表现的朴素余切复形
$$
\text{d} : I/I^2
\longrightarrow
\bigoplus\nolimits_{j = 1}^n S\text{d}x_j
$$
（见第 \ref{algebra-section-netherlander} 节）。只需证明把这个复形在
$\mathfrak q$ 处局部化后，该映射成为分裂单射，见引理
\ref{algebra-lemma-smooth-at-point}。记
$S' = R[x_1, \ldots, x_n]/I^2$。由引理
\ref{algebra-lemma-differential-mod-power-ideal}，有
% P01-E0867: the frozen “Denote S' = ...” omits “by/as”.
$$
S \otimes_{S'} \Omega_{S'/R} =
S \otimes_{R[x_1, \ldots, x_n]} \Omega_{R[x_1, \ldots, x_n]/R} =
\bigoplus\nolimits_{j = 1}^m S\text{d}x_j.
$$
% P01-E0868: the frozen direct sum runs to m (number of equations), but the differential module of the n-variable polynomial ring has basis dx_1,...,dx_n.
因此映射
$$
\text{d} :
I/I^2
\longrightarrow
S \otimes_{S'} \Omega_{S'/R}
$$
与上述朴素余切复形中的映射相同。特别地，我们要证明的断言是否成立，
只依赖三个环 $R \to S' \to S$。设
$\mathfrak q' \subset R[x_1, \ldots, x_n]$ 是对应于 $\mathfrak q$ 的素理想。
由于取微分模与局部化交换（引理
\ref{algebra-lemma-differentials-localize}），只需证明来自
$R \to S'_{\mathfrak q'} \to S_{\mathfrak q}$ 的映射
\begin{equation}
\label{algebra-equation-target-map}
\text{d} :
I_{\mathfrak q'}/I_{\mathfrak q'}^2
\longrightarrow
S_{\mathfrak q} \otimes_{S'_{\mathfrak q'}} \Omega_{S'_{\mathfrak q'}/R}
\end{equation}
是分裂单射。

\medskip\noindent
设 $N \in \mathbf{N}$ 是整数。考察环
$$
B'_N = S'_{\mathfrak q'} / (\mathfrak q')^N S'_{\mathfrak q'}
= (S'/(\mathfrak q')^N S')_{\mathfrak q'}
$$
及其商 $B_N = B'_N/IB'_N$。注意，
$B_N \cong S_{\mathfrak q}/\mathfrak q^NS_{\mathfrak q}$。
观察到 $B'_N$ 是局部 Artin 环，因为它是一个局部 Noether 环
对其极大理想的幂取商。取 $B'_N \to B_N$ 的核 $I_N$ 的一个
$B'_N$-子模滤过
$$
0 \subset J_{N, 1} \subset J_{N, 2} \subset \ldots \subset J_{N, n(N)} = I_N
$$
使每个相继商 $J_{N, i}/J_{N, i - 1}$ 的长度为 $1$。
（由于 $B'_N$ 是 Artin 环，这样的滤过存在。）这给出一列小扩张
$$
B'_N  \to B'_N/J_{N, 1} \to B'_N/J_{N, 2} \to \ldots \to
B'_N/J_{N, n(N)} = B'_N/I_N
= B_N = S_{\mathfrak q}/\mathfrak q^NS_{\mathfrak q}
$$
从映射 $S \to B_N$ 开始，对这些小扩张依次应用条件 (4)，
可知存在交换图
% P01-E0869: the frozen phrase “we see there exists” omits “that”.
$$
\xymatrix{
S \ar[r] \ar[rd] & B_N  \\
R \ar[r] \ar[u] & B'_N \ar[u]
}
$$
显然，环映射 $S \to B'_N$ 分解为
$S \to S_{\mathfrak q} \to B'_N$，其中
$S_{\mathfrak q} \to B'_N$ 是局部环的局部同态。此外，由于
$B'_N$ 的极大理想的 $N$ 次幂为零，
$S_{\mathfrak q} \to B'_N$ 分解经过
$S_{\mathfrak q}/(\mathfrak q)^NS_{\mathfrak q} = B_N$。
换言之，我们已经证明，对每个 $N \in \mathbf{N}$，
$R$-代数满射 $B'_N \to B_N$ 都有分裂映射。

\medskip\noindent
考察由核为 $I_N$ 的满射 $B'_N \to B_N$ 得到的表现
$$
I_N \to B_N \otimes_{B'_N} \Omega_{B'_N/R} \to \Omega_{B_N/R} \to 0
$$
（见引理 \ref{algebra-lemma-differential-seq}）。由上可知，$R$-代数映射
$B'_N \to B_N$ 有右逆。因此由引理
\ref{algebra-lemma-differential-seq-split}，上述序列分裂正合！于是对每个 $N$，
映射
$$
I_N \longrightarrow B_N \otimes_{B'_N} \Omega_{B'_N/R}
$$
是分裂单射。证明的其余部分只需展开这句话的确切含义。注意
$$
I_N = I_{\mathfrak q'}/
(I_{\mathfrak q'}^2 + (\mathfrak q')^N \cap I_{\mathfrak q'})
$$
由 Artin--Rees 引理（引理 \ref{algebra-lemma-Artin-Rees}），
可找到 $c \geq 0$，使得对所有 $N \geq c$ 都有
$$
S_{\mathfrak q}/\mathfrak q^{N - c}S_{\mathfrak q}
\otimes_{S_{\mathfrak q}} I_N =
S_{\mathfrak q}/\mathfrak q^{N - c}S_{\mathfrak q}
\otimes_{S_{\mathfrak q}}
I_{\mathfrak q'}/I_{\mathfrak q'}^2
$$
（这些张量积不过是除以
$\mathfrak q^{N - c}$ 的一种华丽写法）。当然可以假设
$c \geq 1$。
% P01-E0870: the frozen parenthetical says “these tensor product are”, with singular “product” after plural “these”.
由引理 \ref{algebra-lemma-differential-mod-power-ideal}，有
$$
S'_{\mathfrak q'}/(\mathfrak q')^{N - c}S'_{\mathfrak q'}
\otimes_{S'_{\mathfrak q'}} \Omega_{B'_N/R} =
S'_{\mathfrak q'}/(\mathfrak q')^{N - c}S'_{\mathfrak q'}
\otimes_{S'_{\mathfrak q'}} \Omega_{S'_{\mathfrak q'}/R}
$$
进一步与 $B_N = S_{\mathfrak q}/\mathfrak q^N$ 作张量积，得到
% P01-E0871: the frozen source continues with lowercase “we” after the display and has no punctuation ending the preceding sentence.
$$
S_{\mathfrak q}/\mathfrak q^{N - c}S_{\mathfrak q}
\otimes_{S'_{\mathfrak q'}} \Omega_{B'_N/R} =
S_{\mathfrak q}/\mathfrak q^{N - c}S_{\mathfrak q}
\otimes_{S'_{\mathfrak q'}} \Omega_{S'_{\mathfrak q'}/R}.
$$
由于分裂单射与任意对象作张量积后仍是分裂单射，可知
$$
S_{\mathfrak q}/\mathfrak q^{N - c}S_{\mathfrak q}
\otimes_{S_{\mathfrak q}}
(\ref{algebra-equation-target-map}) =
S_{\mathfrak q}/\mathfrak q^{N - c}S_{\mathfrak q}
\otimes_{S_{\mathfrak q}/\mathfrak q^N S_{\mathfrak q}}
(I_N \longrightarrow B_N \otimes_{B'_N} \Omega_{B'_N/R})
$$
对所有 $N \geq c$ 都是分裂单射。由引理
\ref{algebra-lemma-split-injection-after-completion}，可知
（\ref{algebra-equation-target-map}）是分裂单射。这完成了证明。
\end{proof}

\section{光滑环映射结果概览}
\label{algebra-section-smooth-overview}

\noindent
下面列出我们在前几节证明的关于光滑环映射的结果。
更精确的陈述和定义请查阅所给参考文献。
\begin{enumerate}
\item 环映射 $R \to S$ 光滑，是指它有限表现，且 $S/R$ 的朴素余切复形
拟同构于置于 $0$ 次的有限投射 $S$-模，见定义
\ref{algebra-definition-smooth}。
\item 若 $S$ 在 $R$ 上光滑，则 $\Omega_{S/R}$ 是有限投射
$S$-模，见定义 \ref{algebra-definition-smooth} 后的讨论。
\item 光滑性在 $S$ 上是局部的，见引理
\ref{algebra-lemma-locally-smooth}。
\item 光滑性在基变换下稳定，见引理
\ref{algebra-lemma-base-change-smooth}。
\item 光滑性在合成下稳定，见引理
\ref{algebra-lemma-compose-smooth}。
\item 光滑环映射是 syntomic 的，特别地平坦，见引理
\ref{algebra-lemma-smooth-syntomic}。
\item 光滑纤维环的有限表现平坦环映射光滑，见引理
\ref{algebra-lemma-flat-fibre-smooth}。
\item 有限表现环映射 $R \to S$ 光滑，当且仅当它形式光滑，见命题
\ref{algebra-proposition-smooth-formally-smooth}。
\item 若 $R \to S$ 是有限型环映射，且 $R$ 是 Noether 环，
则为检验 $R \to S$ 光滑，只需检验沿局部 Artin 环的小扩张的形式光滑
提升性质，见引理 \ref{algebra-lemma-smooth-test-artinian}。
\item 光滑环映射 $R \to S$ 是某个光滑环映射
$R_0 \to S_0$ 的基变换，其中 $R_0$ 在 $\mathbf{Z}$ 上有限型，见引理
\ref{algebra-lemma-finite-presentation-fs-Noetherian}。
\item 环映射的光滑点集的形成与平坦基变换交换，见引理
\ref{algebra-lemma-flat-base-change-locus-smooth}。
\item 若 $S$ 在代数闭域 $k$ 上有限型，且
$\mathfrak m \subset S$ 是极大理想，则下列条件等价：
% P01-E0872: the frozen premise omits “is” before “a maximal ideal”, and the list introduction omits a colon.
\begin{enumerate}
\item $S$ 在 $\mathfrak m$ 的某个邻域中相对于 $k$ 光滑；
\item $S_{\mathfrak m}$ 是正则局部环；
\item $\dim(S_{\mathfrak m}) =
\dim_{\kappa(m)} \Omega_{S/k} \otimes_S \kappa(\mathfrak m)$。
\end{enumerate}
见引理 \ref{algebra-lemma-characterize-smooth-kbar}。
% P01-E0873: the frozen dimension formula uses kappa(m) once although only the ideal mathfrak m is defined and the same formula ends with kappa(mathfrak m).
\item 若 $S$ 在域 $k$ 上有限型，且
$\mathfrak q \subset S$ 是素理想，则下列条件等价：
\begin{enumerate}
\item $S$ 在 $\mathfrak q$ 的某个邻域中相对于 $k$ 光滑；
\item $\dim_{\mathfrak q}(S/k) =
\dim_{\kappa(\mathfrak q)} \Omega_{S/k} \otimes_S \kappa(\mathfrak q)$。
\end{enumerate}
见引理 \ref{algebra-lemma-characterize-smooth-over-field}。
% P01-E0874: the frozen prime-ideal premise omits “is” and the list introduction omits a colon; the same omitted predicate recurs in the next two items.
\item 若 $S$ 在域上光滑，则它的所有局部环都正则，见引理
\ref{algebra-lemma-characterize-smooth-over-field}。
\item 若 $S$ 在域 $k$ 上有限型，$\mathfrak q \subset S$ 是素理想，
域扩张 $\kappa(\mathfrak q)/k$ 可分，且 $S_{\mathfrak q}$ 正则，
则 $S$ 在 $\mathfrak q$ 处相对于 $k$ 光滑，见引理
\ref{algebra-lemma-separable-smooth}。
\item 若 $S$ 在域 $k$ 上有限型，$k$ 的特征为 $0$，
$\mathfrak q \subset S$ 是素理想，且
$\Omega_{S/k, \mathfrak q}$ 自由，则 $S$ 在 $\mathfrak q$ 处相对于
$k$ 光滑，见引理 \ref{algebra-lemma-characteristic-zero-local-smooth}。
\end{enumerate}
其中一些结果是用标准光滑环映射的概念证明的，见定义
\ref{algebra-definition-standard-smooth}。它对应于 syntomic 态射情形中的
相对全局完全交映射，也是构造例子的最简便方式。

\section{平展环映射}
\label{algebra-section-etale}

\noindent
平展环映射是相对维数等于零的光滑环映射。这与下面稍微更直接的定义
等价。

\begin{definition}
\label{algebra-definition-etale}
设 $R \to S$ 是环映射。若它是有限表现的，并且朴素余切复形
$\NL_{S/R}$ 拟同构于零，即 $H_1(\NL_{S/R}) = 0$
且 $\Omega_{S/R} = 0$，则称 $R \to S$ 是{\it 平展的}。给定 $S$ 的素理想
$\mathfrak q$，若存在 $g \in S$，$g \not \in \mathfrak q$，使得
$R \to S_g$ 平展，则称 $R \to S$ {\it 在 $\mathfrak q$ 处平展}。
\end{definition}

\noindent
特别地，若 $S$ 在 $R$ 上平展，则 $\Omega_{S/R} = 0$。
若 $R \to S$ 光滑，则 $R \to S$ 平展当且仅当 $\Omega_{S/R} = 0$。
由关于光滑环映射的结果，我们自动得到大量关于平展映射的结果。
下面在引理 \ref{algebra-lemma-etale} 中汇总这些结果。但在此之前，我们先证明
{\it 任意}平展环映射都是标准光滑的。

\begin{lemma}
\label{algebra-lemma-etale-standard-smooth}
任意平展环映射都是标准光滑的。更确切地说，若 $R \to S$ 平展，
则存在一个表现
$S = R[x_1, \ldots, x_n]/(f_1, \ldots, f_n)$，使得
$\det(\partial f_j/\partial x_i)$ 在 $S$ 中的像可逆。
\end{lemma}

\begin{proof}
设 $R \to S$ 平展。选取一个表现 $S = R[x_1, \ldots, x_n]/I$。
由 $R \to S$ 平展可知
$$
\text{d} :
I/I^2
\longrightarrow
\bigoplus\nolimits_{i = 1, \ldots, n} S\text{d}x_i
$$
是同构，特别地，$I/I^2$ 是自由 $S$-模。因此由引理
\ref{algebra-lemma-huber}，可以假设（必要时更换表现）
$I = (f_1, \ldots, f_c)$，并且各类 $f_i \bmod I^2$
构成 $I/I^2$ 的一组基。由上面显示的映射是同构，立刻得到
$c = n$，且 $\det(\partial f_j/\partial x_i)$ 在 $S$ 中可逆。
\end{proof}

\begin{lemma}
\label{algebra-lemma-etale}
关于平展环映射的结果如下。
\begin{enumerate}
\item 对任意环 $R$ 和任意 $f \in R$，环映射 $R \to R_f$ 平展。
\item 平展环映射的合成平展。
\item 平展环映射的基变换平展。
\item 平展性是局部性质：给定环映射 $R \to S$ 和生成单位理想的元素
$g_1, \ldots, g_m \in S$，若对 $j = 1, \ldots, m$，
$R \to S_{g_j}$ 平展，则 $R \to S$ 平展。
\item 设 $R \to S$ 是有限表现环映射，$R \to R'$ 是平坦环映射，
并令 $S' = R' \otimes_R S$。$R' \to S'$ 的平展点集，是
$R \to S$ 的平展点集在 $\Spec(S') \to \Spec(S)$ 下的逆像。
\item 平展环映射是 syntomic 的，特别地平坦。
\item 若 $S$ 在域 $k$ 上有限型，则 $S$ 在 $k$ 上平展，当且仅当
$\Omega_{S/k} = 0$。
\item 任意平展环映射 $R \to S$ 都是某个平展环映射
$R_0 \to S_0$ 的基变换，其中 $R_0$ 在 $\mathbf{Z}$ 上有限型。
\item 设 $A = \colim A_i$ 是环的滤过余极限，并设 $A \to B$
是平展环映射。则对某个 $i$，存在平展环映射 $A_i \to B_i$，
使得 $B \cong A \otimes_{A_i} B_i$。
\item 设 $A$ 是环，$S$ 是 $A$ 的乘法子集，并设
$S^{-1}A \to B'$ 平展。则存在平展环映射 $A \to B$，使得
$B' \cong S^{-1}B$。
\item 设 $A$ 是环，并设 $B = B' \times B''$ 是 $A$-代数的乘积。
则 $B$ 在 $A$ 上平展，当且仅当 $B'$ 和 $B''$ 都在 $A$ 上平展。
\end{enumerate}
\end{lemma}

\begin{proof}
在每一种情形中，我们都使用光滑环映射的相应结果，再加上一个说明
$\Omega_{S/R}$ 为零的小论证。

\medskip\noindent
(1) 的证明。环映射 $R \to R_f$ 光滑，且 $\Omega_{R_f/R} = 0$。

\medskip\noindent
(2) 的证明。光滑映射 $A \to B$ 与 $B \to C$ 的合成
$A \to C$ 光滑，见引理 \ref{algebra-lemma-compose-smooth}。由引理
\ref{algebra-lemma-exact-sequence-differentials}，由于 $\Omega_{C/B}$ 和
$\Omega_{B/A}$ 都为零，可知 $\Omega_{C/A}$ 为零。

\medskip\noindent
(3) 的证明。设 $R \to S$ 平展，而 $R \to R'$ 任意。
则 $R' \to S' = R' \otimes_R S$ 光滑，见引理
\ref{algebra-lemma-base-change-smooth}。由引理
\ref{algebra-lemma-differentials-base-change}，
$\Omega_{S'/R'} = S' \otimes_S \Omega_{S/R}$，故
$\Omega_{S'/R'} = 0$。因此 $R' \to S'$ 平展。

\medskip\noindent
(4) 的证明。假设 (4) 的条件成立。由引理
\ref{algebra-lemma-locally-smooth}，$R \to S$ 光滑。又已知对所有 $i$，
$\Omega_{S_{g_i}/R} = (\Omega_{S/R})_{g_i} = 0$。
于是 $\Omega_{S/R} = 0$，见引理 \ref{algebra-lemma-cover}。

\medskip\noindent
(5) 的证明。光滑映射的相应结果是引理
\ref{algebra-lemma-flat-base-change-locus-smooth}。在该引理的证明中，
我们用到了 $\NL_{S/R} \otimes_S S'$ 与 $\NL_{S'/R'}$
同伦等价。这样只需证明：若 $M$ 是有限表现 $S$-模，则 $S'$
的满足 $(M \otimes_S S')_{\mathfrak q'} = 0$ 的素理想
$\mathfrak q'$ 所成集合，是 $S$ 的满足 $M_{\mathfrak q} = 0$
的素理想 $\mathfrak q$ 所成集合的逆像。这由引理
\ref{algebra-lemma-support-base-change} 得到。

\medskip\noindent
(6) 的证明。直接由光滑环映射的相应结果（引理
\ref{algebra-lemma-smooth-syntomic}）得到。

\medskip\noindent
(7) 的证明。由引理 \ref{algebra-lemma-characterize-smooth-over-field}
和定义得到。

\medskip\noindent
(8) 的证明。引理 \ref{algebra-lemma-finite-presentation-fs-Noetherian}
给出光滑环映射的相应结果。所得光滑环映射 $R_0 \to S_0$
满足引理 \ref{algebra-lemma-relative-dimension-CM} 的条件，因而可以把
$S_0$ 换成它在 $R_0$ 上相对维数为 $0$ 的因子。

\medskip\noindent
(9) 的证明。由 (8) 得到，因为由引理
\ref{algebra-lemma-characterize-finite-presentation}，对某个 $i$，
$R_0 \to A$ 会通过 $A_i$ 分解。

\medskip\noindent
(10) 的证明。由 (9)、(1) 和 (2) 得到，因为 $S^{-1}A$
是 $A$ 的主局部化的滤过余极限。

\medskip\noindent
(11) 的证明。先用引理 \ref{algebra-lemma-product-smooth} 得到光滑性的
相应结果，再用 $\Omega_{B/A}$ 为零当且仅当
$\Omega_{B'/A}$ 和 $\Omega_{B''/A}$ 都为零。
\end{proof}

\noindent
下面更详细地说明在域上平展意味着什么。

\begin{lemma}
\label{algebra-lemma-etale-over-field}
设 $k$ 是域。环映射 $k \to S$ 平展，当且仅当 $S$ 作为
$k$-代数同构于 $k$ 的若干有限可分扩张的有限乘积。
\end{lemma}

\begin{proof}
以下事实将不再另行说明：若 $S = S_1 \times \ldots \times S_n$
是 $k$-代数的有限乘积，则 $S$ 在 $k$ 上平展，当且仅当每个
$S_i$ 都在 $k$ 上平展。见引理 \ref{algebra-lemma-etale} 的 (11)。

\medskip\noindent
若 $k'/k$ 是有限可分域扩张，则可写成
$k' = k(\alpha) \cong k[x]/(f)$，其中 $f$ 是元素 $\alpha$
的极小多项式。由于 $k'$ 在 $k$ 上可分，有
$\gcd(f, f') = 1$。这说明
$\text{d} : k'\cdot f \to k' \cdot \text{d}x$ 是同构，故
$k \to k'$ 平展。因此，若 $S$ 是 $k$ 的若干有限可分扩张的
有限乘积，则 $S$ 在 $k$ 上平展。

\medskip\noindent
反过来，假设 $k \to S$ 平展。则 $S$ 在 $k$ 上光滑，且
$\Omega_{S/k} = 0$。由引理
\ref{algebra-lemma-characterize-smooth-over-field}，对 $S$ 的每个极大理想
$\mathfrak m$，有 $\dim_\mathfrak m \Spec(S) = 0$。因此
$\dim(S) = 0$。由命题 \ref{algebra-proposition-dimension-zero-ring}，
$S$ 是有限多个 Artin 局部环的乘积。由已经使用过的引理
\ref{algebra-lemma-characterize-smooth-over-field}，这些局部环都是域。
所以可以假设 $S = k'$ 是域。由 Hilbert 零点定理（定理
\ref{algebra-theorem-nullstellensatz}），扩张 $k'/k$ 有限。环映射
$k \to k'$ 的光滑性由引理 \ref{algebra-lemma-smooth-at-generic-point}
推出 $k'/k$ 是可分扩张，证明完成。
\end{proof}

\begin{lemma}
\label{algebra-lemma-etale-at-prime}
设 $R \to S$ 是环映射，$\mathfrak q \subset S$ 是位于 $R$
中的 $\mathfrak p$ 之上的素理想。若 $S/R$ 在 $\mathfrak q$
处平展，则
\begin{enumerate}
\item $\mathfrak p S_{\mathfrak q} = \mathfrak qS_{\mathfrak q}$，
它是局部环 $S_{\mathfrak q}$ 的极大理想；
\item 域扩张 $\kappa(\mathfrak q)/\kappa(\mathfrak p)$ 有限可分。
\end{enumerate}
\end{lemma}

\begin{proof}
首先可以把 $S$ 换成某个 $S_g$，其中 $g \in S$ 且
$g \not \in \mathfrak q$，从而假设 $R \to S$ 平展。展开
$S \otimes_R \kappa(\mathfrak p)$ 在 $\kappa(\mathfrak p)$
上平展这一事实，再应用引理 \ref{algebra-lemma-etale-over-field}，
即可得到本引理。
\end{proof}

\begin{lemma}
\label{algebra-lemma-etale-quasi-finite}
平展环映射是拟有限的。
\end{lemma}

\begin{proof}
设 $R \to S$ 是平展环映射。由定义，$R \to S$ 是有限型的。
对任意素理想 $\mathfrak p \subset R$，纤维环
$S \otimes_R \kappa(\mathfrak p)$ 在 $\kappa(\mathfrak p)$ 上平展，
因而是 $\kappa(\mathfrak p)$ 的若干有限可分域扩张的有限乘积；
特别地，它在 $\kappa(\mathfrak p)$ 上有限。
% P01-E0875: the frozen source says “a finite products of fields”.
故由引理 \ref{algebra-lemma-quasi-finite}，$R \to S$ 拟有限。
\end{proof}

\begin{lemma}
\label{algebra-lemma-characterize-etale}
设 $R \to S$ 是环映射。设 $\mathfrak q$ 是 $S$ 的素理想，
它位于 $R$ 的素理想 $\mathfrak p$ 之上。若
\begin{enumerate}
\item $R \to S$ 是有限表现的；
\item $R_{\mathfrak p} \to S_{\mathfrak q}$ 平坦；
\item $\mathfrak p S_{\mathfrak q}$ 是局部环
$S_{\mathfrak q}$ 的极大理想；
\item 域扩张 $\kappa(\mathfrak q)/\kappa(\mathfrak p)$ 有限可分，
\end{enumerate}
则 $R \to S$ 在 $\mathfrak q$ 处平展。
\end{lemma}

\begin{proof}
应用引理 \ref{algebra-lemma-isolated-point-fibre}，可找到
$g \in S$，$g \not \in \mathfrak q$，使得 $\mathfrak q$ 是
$S_g$ 中唯一位于 $\mathfrak p$ 之上的素理想。我们可以并且确实
把 $S$ 换成 $S_g$。于是 $S \otimes_R \kappa(\mathfrak p)$
只有一个素理想，因而是局部环，并且等于
$S_{\mathfrak q}/\mathfrak pS_{\mathfrak q}
\cong \kappa(\mathfrak q)$。
由引理 \ref{algebra-lemma-flat-fibre-smooth}，存在
$g \in S$，$g \not \in \mathfrak q$，使得 $R \to S_g$ 光滑。
再次把 $S$ 换成 $S_g$，便可假设 $R \to S$ 光滑。
% P01-E0876: the frozen source has the malformed transition “Replace S by S_g again we may assume”.
由引理 \ref{algebra-lemma-smooth-syntomic}，甚至可以假设
$R \to S$ 标准光滑，比如
$S = R[x_1, \ldots, x_n]/(f_1, \ldots, f_c)$。
由于 $S \otimes_R \kappa(\mathfrak p) = \kappa(\mathfrak q)$
的维数为 $0$，可知 $n = c$，即 $R \to S$ 平展。
\end{proof}

\noindent
下面是一个全新的现象。

\begin{lemma}
\label{algebra-lemma-map-between-etale}
设 $R \to S$ 和 $R \to S'$ 平展。则任意 $R$-代数映射
$S' \to S$ 都平展。
\end{lemma}

\begin{proof}
首先，由引理 \ref{algebra-lemma-compose-finite-type}，$S' \to S$
是有限表现的。设素理想 $\mathfrak q \subset S$ 位于
$\mathfrak q' \subset S'$ 和 $\mathfrak p \subset R$ 之上。
由引理 \ref{algebra-lemma-etale-at-prime}，环映射
$S'_{\mathfrak q'}/\mathfrak p S'_{\mathfrak q'} \to
S_{\mathfrak q}/\mathfrak p S_{\mathfrak q}$
是 $\kappa(\mathfrak p)$ 的有限可分扩张之间的映射。
特别地，它平坦。因此由引理 \ref{algebra-lemma-criterion-flatness-fibre}，
$S'_{\mathfrak q'} \to S_{\mathfrak q}$ 平坦。所以 $S' \to S$
平坦。此外，上述论证还说明 $\mathfrak q'S_{\mathfrak q}$ 是
$S_{\mathfrak q}$ 的极大理想，并且
$S'_{\mathfrak q'} \to S_{\mathfrak q}$ 的剩余域扩张有限可分。
故由引理 \ref{algebra-lemma-characterize-etale}，$S' \to S$ 在
$\mathfrak q$ 处平展。由于平展性是局部的（见引理
\ref{algebra-lemma-etale}），结论成立。
\end{proof}

\begin{lemma}
\label{algebra-lemma-surjective-flat-finitely-presented}
设 $\varphi : R \to S$ 是环映射。若 $R \to S$ 满射、平坦且
有限表现，则存在幂等元 $e \in R$，使得 $S = R_e$。
% P01-E0877: the frozen source says “there exist an idempotent”.
\end{lemma}

\begin{proof}[第一种证明]
设 $I$ 是 $\varphi$ 的核。由于 $\varphi$ 是有限表现的，
由引理 \ref{algebra-lemma-finite-presentation-independent}，$I$ 有限生成。
此外，由于 $S$ 在 $R$ 上平坦，把正合列
$0 \to I \to R \to S \to 0$ 在 $R$ 上与 $S$ 作张量积，得到
$I/I^2 = 0$。现在由引理
\ref{algebra-lemma-ideal-is-squared-union-connected} 即得结论。
\end{proof}

\begin{proof}[第二种证明]
由于 $\Spec(S) \to \Spec(R)$ 是到某个闭子集上的同胚（见引理
\ref{algebra-lemma-spec-closed}），并且是开映射（见命题
\ref{algebra-proposition-fppf-open}），可知其像等于某个幂等元
$e \in R$ 所确定的 $D(e)$（见引理
\ref{algebra-lemma-disjoint-decomposition}）。因此 $R_e \to S$ 在谱上诱导
双射。比如由引理 \ref{algebra-lemma-finite-flat-local} 和
\ref{algebra-lemma-NAK}，该映射在所有局部环上都诱导同构。于是
$R_e \to S$ 也为单射，例如可见引理
\ref{algebra-lemma-characterize-zero-local}。
\end{proof}

\begin{lemma}
\label{algebra-lemma-lift-etale}
\begin{slogan}
平展环映射可以沿环的满射提升
\end{slogan}
设 $R$ 是环，$I \subset R$ 是理想。设
$R/I \to \overline{S}$ 是平展环映射。则存在平展环映射
$R \to S$，使得作为 $R/I$-代数，
$\overline{S} \cong S/IS$。
\end{lemma}

\begin{proof}
由引理 \ref{algebra-lemma-etale-standard-smooth}，可以如定义
\ref{algebra-definition-standard-smooth} 中那样写成
$\overline{S} =
(R/I)[x_1, \ldots, x_n]/(\overline{f}_1, \ldots, \overline{f}_n)$，
其中
$\overline{\Delta} =
\det(\frac{\partial \overline{f}_i}{\partial x_j})_{i, j = 1, \ldots, n}$
在 $\overline{S}$ 中可逆。任取一些提升 $f_i$，
并令
$S = R[x_1, \ldots, x_n, x_{n+1}]/(f_1, \ldots, f_n, x_{n + 1}\Delta - 1)$，
其中
$\Delta = \det(\frac{\partial f_i}{\partial x_j})_{i, j = 1, \ldots, n}$，
如例 \ref{algebra-example-make-standard-smooth} 所示。这就证明了引理。
\end{proof}

\begin{lemma}
\label{algebra-lemma-lift-etale-infinitesimal}
考虑交换图
$$
\xymatrix{
0 \ar[r] &
J \ar[r] &
B' \ar[r] &
B \ar[r] & 0 \\
0 \ar[r] &
I \ar[r] \ar[u] &
A' \ar[r] \ar[u] &
A \ar[r] \ar[u] & 0
}
$$
其各行正合，其中 $B' \to B$ 和 $A' \to A$ 是满环映射，
且其核都是平方为零的理想。若 $A \to B$ 平展，且
$J = I \otimes_A B$，则 $A' \to B'$ 平展。
\end{lemma}

\begin{proof}
由引理 \ref{algebra-lemma-lift-etale}，存在平展环映射 $A' \to C$，
使得 $C/IC = B$。于是 $A' \to C$ 形式光滑（由命题
\ref{algebra-proposition-smooth-formally-smooth}），故得到 $A'$-代数映射
$\varphi : C \to B'$。由于 $A' \to C$ 平坦，有
$I \otimes_A B = I \otimes_A C/IC = IC$。因此，假设
$J = I \otimes_A B$ 蕴含 $\varphi$ 诱导同构 $IC \to J$
和同构 $C/IC \to B'/IB'$，从而 $\varphi$ 是同构。
\end{proof}

\begin{example}
\label{algebra-example-factor-polynomials-etale}
设 $n , m \geq 1$ 是整数。考虑例
\ref{algebra-example-factor-polynomials} 中的环映射
% P01-E0878: the frozen source has an anomalous space before the comma in “n , m”.
\begin{eqnarray*}
R = \mathbf{Z}[a_1, \ldots, a_{n + m}]
& \longrightarrow &
S = \mathbf{Z}[b_1, \ldots, b_n, c_1, \ldots, c_m] \\
a_1 & \longmapsto & b_1 + c_1 \\
a_2 & \longmapsto & b_2 + b_1 c_1 + c_2 \\
\ldots & \ldots & \ldots \\
a_{n + m} & \longmapsto & b_n c_m
\end{eqnarray*}
符号化地写成
$$
S = R[b_1, \ldots, c_m]/(\{a_k(b_i, c_j) - a_k\}_{k = 1, \ldots, n + m})
$$
其中，例如 $a_1(b_i, c_j) = b_1 + c_1$。偏导数矩阵为
$$
\left(
\begin{matrix}
1 & c_1 & \ldots & c_m & 0 & \ldots & \ldots & 0 \\
0 & 1 & c_1 & \ldots & c_m & 0 & \ldots & 0 \\
\ldots & \ldots & \ldots & \ldots & \ldots & \ldots & \ldots & \ldots \\
0 & \ldots & 0 & 1 & c_1 & c_2 & \ldots & c_m \\
1 & b_1 & \ldots & b_{n - 1} & b_n & 0 & \ldots & 0 \\
0 & 1 & b_1 & \ldots & b_{n - 1} & b_n & \ldots & 0 \\
\ldots & \ldots & \ldots & \ldots & \ldots & \ldots & \ldots & \ldots \\
0 & \ldots & \ldots & 0 & 1 & b_1 & \ldots & b_n
\end{matrix}
\right)
$$
这个矩阵的行列式 $\Delta$ 更常称为多项式
$g = x^n + b_1 x^{n - 1} + \ldots + b_n$ 与
$h = x^m + c_1 x^{m - 1} + \ldots + c_m$ 的{\it 结式}，
上面的矩阵称为与 $g, h$ 相伴的 {\it Sylvester 矩阵}。
用公式表示就是 $\Delta = \text{Res}_x(g, h)$。Sylvester 矩阵是
线性映射矩阵的转置：
\begin{eqnarray*}
S[x]_{< m} \oplus S[x]_{< n} & \longrightarrow & S[x]_{< n + m} \\
a \oplus b & \longmapsto & ag + bh
\end{eqnarray*}
设 $\mathfrak q \subset S$ 是任意素理想。由上述讨论，下列条件等价：
\begin{enumerate}
\item $R \to S$ 在 $\mathfrak q$ 处平展；
\item $\Delta = \text{Res}_x(g, h) \not \in \mathfrak q$；
\item 多项式 $g, h$ 在 $\kappa(\mathfrak q)[x]$ 中的像
$\overline{g}, \overline{h} \in \kappa(\mathfrak q)[x]$ 互素。
\end{enumerate}
(2) 与 (3) 等价，是因为 Sylvester 矩阵在
$\text{Mat}(n + m, \kappa(\mathfrak q))$ 中的像有核，
当且仅当多项式 $\overline{g}, \overline{h}$ 有公共因子。
由此可知环映射
$$
R \longrightarrow S[\frac{1}{\Delta}] = S[\frac{1}{\text{Res}_x(g, h)}]
$$
平展。
\end{example}


\begin{lemma}
\label{algebra-lemma-factor-mod-lift-etale}
设 $R$ 是环，$f \in R[x]$ 是首一多项式，$\mathfrak p$ 是
$R$ 的素理想。设
$f \bmod \mathfrak p = \overline{g} \overline{h}$ 是 $f$ 在
$\kappa(\mathfrak p)[x]$ 中的像的一个分解式。若
$\gcd(\overline{g}, \overline{h}) = 1$，则存在
\begin{enumerate}
\item 平展环映射 $R \to R'$；
\item 位于 $\mathfrak p$ 之上的素理想 $\mathfrak p' \subset R'$；
\item $R'[x]$ 中的分解式 $f = g h$，
\end{enumerate}
使得
\begin{enumerate}
\item $\kappa(\mathfrak p) = \kappa(\mathfrak p')$；
\item $\overline{g} = g \bmod \mathfrak p'$，
$\overline{h} = h \bmod \mathfrak p'$；
\item 多项式 $g, h$ 生成 $R'[x]$ 中的单位理想。
\end{enumerate}
\end{lemma}

\begin{proof}
假设
$\overline{g} = \overline{b}_0 x^n + \overline{b}_1 x^{n - 1} + \ldots
+ \overline{b}_n$，且
$\overline{h} = \overline{c}_0 x^m + \overline{c}_1 x^{m - 1} + \ldots
+ \overline{c}_m$，其中
$\overline{b}_0, \overline{c}_0 \in \kappa(\mathfrak p)$ 非零。
在 $R$ 的某个不属于 $\mathfrak p$ 的元素处对 $R$ 局部化后，可以假设
$\overline{b}_0$ 是某个可逆元 $b_0 \in R$ 的像。用
$\overline{g}/b_0$ 替换 $\overline{g}$，用
$b_0\overline{h}$ 替换 $\overline{h}$，便约化到
$\overline{g}$、$\overline{h}$ 都首一的情形（验证从略）。设
$\overline{g} = x^n + \overline{b}_1 x^{n - 1} + \ldots + \overline{b}_n$，
且
$\overline{h} = x^m + \overline{c}_1 x^{m - 1} + \ldots + \overline{c}_m$。
写成
$f = x^{n + m} + a_1 x^{n + m - 1} + \ldots + a_{n + m}$。
考虑纤维积
$$
R' = R \otimes_{\mathbf{Z}[a_1, \ldots, a_{n + m}]}
\mathbf{Z}[b_1, \ldots, b_n, c_1, \ldots, c_m]
$$
其中映射 $\mathbf{Z}[a_k] \to \mathbf{Z}[b_i, c_j]$
如例 \ref{algebra-example-factor-polynomials} 和
\ref{algebra-example-factor-polynomials-etale} 所示。由构造，存在
$R$-代数映射
$$
R' = R \otimes_{\mathbf{Z}[a_1, \ldots, a_{n + m}]}
\mathbf{Z}[b_1, \ldots, b_n, c_1, \ldots, c_m]
\longrightarrow
\kappa(\mathfrak p)
$$
它把 $b_i$ 映到 $\overline{b}_i$，把 $c_j$ 映到
$\overline{c}_j$。记这个映射的核为 $\mathfrak p' \subset R'$。
% P01-E0879: the frozen source says “Denote p' the kernel” and omits “by/as”.
由于按假设多项式 $\overline{g}, \overline{h}$ 互素，可知元素
$\Delta = \text{Res}_x(g, h) \in \mathbf{Z}[b_i, c_j]$
（见例 \ref{algebra-example-factor-polynomials-etale}）在上述映射下
在 $\kappa(\mathfrak p)$ 中的像不为零。因此 $R \to R'$
在 $\mathfrak p'$ 处平展。事实上，本引理所提问题的一个解是
环映射 $R \to R'[1/\Delta]$ 和素理想
$\mathfrak p' R'[1/\Delta]$。由于在此环中
$\text{Res}_x(f, g)$ 可逆，Sylvester 矩阵在 $R'[1/\Delta]$ 上可逆，因而对某些
$a, b \in R'[1/\Delta][x]$ 有 $1 = a g +  b h$，见例
\ref{algebra-example-factor-polynomials-etale}。
% P01-E0880: the frozen source invokes Res_x(f,g), but Delta and the Sylvester matrix are for Res_x(g,h); since f=gh, Res_x(f,g) cannot support the stated invertibility.
\end{proof}

\section{平展环映射的局部结构}
\label{algebra-section-etale-local-structure}

\noindent
引理 \ref{algebra-lemma-etale-standard-smooth} 告诉我们，把标准平展态射定义为
相对维数为 $0$ 的标准光滑态射其实没有多大意义。作为平展态射的模型，
我们取域的有限可分扩张 $k'/k$ 给出的例子。也就是说，总能找到元素
$\alpha \in k'$，使得 $k' = k(\alpha)$，并且 $\alpha$ 的极小多项式
$f(x) \in k[x]$ 的导数 $f'$ 与 $f$ 互素。

\begin{definition}
\label{algebra-definition-standard-etale}
设 $R$ 是环。设 $g , f  \in R[x]$。假设 $f$ 首一，并且导数
$f'$ 在局部化 $R[x]_g/(f)$ 中可逆。此时称环映射
$R \to R[x]_g/(f)$ 是{\it 标准平展的}。
% P01-E0881: the frozen source has anomalous spaces in the mathematical list “g , f  in R[x]”.
\end{definition}

\noindent
在命题 \ref{algebra-proposition-etale-locally-standard} 中，我们将证明每个平展
环映射都局部标准平展。

\begin{lemma}
\label{algebra-lemma-standard-etale}
设 $R \to R[x]_g/(f)$ 标准平展。
\begin{enumerate}
\item 环映射 $R \to R[x]_g/(f)$ 平展。
\item 对任意环映射 $R \to R'$，标准平展环映射
$R \to R[x]_g/(f)$ 的基变换 $R' \to R'[x]_g/(f)$ 标准平展。
\item $R[x]_g/(f)$ 的任意主局部化在 $R$ 上标准平展。
\item 一般而言，标准平展映射的合成{\bf 不}标准平展。
\end{enumerate}
\end{lemma}

\begin{proof}
从略。下面给出 (4) 的一个例子。环映射
$\mathbf{F}_2 \to \mathbf{F}_{2^2}$ 标准平展。环映射
$\mathbf{F}_{2^2} \to \mathbf{F}_{2^2} \times \mathbf{F}_{2^2}
\times \mathbf{F}_{2^2} \times \mathbf{F}_{2^2}$ 标准平展。
但是环映射
$\mathbf{F}_2 \to \mathbf{F}_{2^2} \times \mathbf{F}_{2^2}
\times \mathbf{F}_{2^2} \times \mathbf{F}_{2^2}$ 不标准平展。
\end{proof}

\noindent
标准平展态射是构造平展映射的一种方便方法。下面给出一个例子。

\begin{lemma}
\label{algebra-lemma-make-etale-map-prescribed-residue-field}
设 $R$ 是环，$\mathfrak p$ 是 $R$ 的素理想，并设
$L/\kappa(\mathfrak p)$ 是有限可分域扩张。则存在平展环映射
$R \to R'$ 以及位于 $\mathfrak p$ 之上的素理想 $\mathfrak p'$，
使得域扩张 $\kappa(\mathfrak p')/\kappa(\mathfrak p)$ 同构于
$\kappa(\mathfrak p) \subset L$。
\end{lemma}

\begin{proof}
由本原元定理，可以写成 $L = \kappa(\mathfrak p)[\alpha]$。
令 $\overline{f} \in \kappa(\mathfrak p)[x]$ 表示 $\alpha$
的极小多项式（特别地，它首一）。用某个 $c \in R$，
$c\not \in \mathfrak p$ 的 $c\alpha$ 替换 $\alpha$ 后，可以假设
$\overline{f}$ 的所有系数都在 $R \to \kappa(\mathfrak p)$ 的像中
（验证从略）。所以可以找到首一多项式 $f \in R[x]$，它在
$\kappa(\mathfrak p)[x]$ 中映到 $\overline{f}$。由于
$\kappa(\mathfrak p) \subset L$ 可分，可知
$\gcd(\overline{f}, \overline{f}') = 1$。因此存在
$\gamma \in L$，使得 $\overline{f}'(\alpha) \gamma = 1$。
于是得到 $R$-代数映射
\begin{eqnarray*}
R[x, 1/f']/(f) & \longrightarrow & L \\
x & \longmapsto & \alpha \\
1/f' & \longmapsto & \gamma
\end{eqnarray*}
左端是 $R$ 上的标准平展代数 $R'$，而该环映射的核给出所需的素理想。
\end{proof}

\begin{proposition}
\label{algebra-proposition-etale-locally-standard}
设 $R \to S$ 是环映射，$\mathfrak q \subset S$ 是素理想。
若 $R \to S$ 在 $\mathfrak q$ 处平展，则存在
$g \in S$，$g \not \in \mathfrak q$，使得 $R \to S_g$
标准平展。
\end{proposition}

\begin{proof}
下面的证明稍显迂回，或许有缩短它的方法。

\medskip\noindent
第 1 步。由定义 \ref{algebra-definition-etale}，存在
$g \in S$，$g \not \in \mathfrak q$，使得 $R \to S_g$ 平展。
所以可以假设 $S$ 在 $R$ 上平展。

\medskip\noindent
第 2 步。由引理 \ref{algebra-lemma-etale}，存在平展环映射
$R_0 \to S_0$，其中 $R_0$ 在 $\mathbf{Z}$ 上有限型；还存在环映射
$R_0 \to R$，使得 $R = R \otimes_{R_0} S_0$。记
$\mathfrak q_0$ 为 $S_0$ 中对应于 $\mathfrak q$ 的素理想。
% P01-E0882: the frozen base-change identity has R on the left, although the construction and the next sentence require S = R tensor_{R_0} S_0.
% P01-E0883: the frozen source says “Denote q_0 the prime” and omits “by/as”.
若对 $(R_0 \to S_0, \mathfrak q_0)$ 证明了结论，则由基变换，
结论也对 $(R \to S, \mathfrak q)$ 成立。因此可以假设 $R$
是 Noether 环。

\medskip\noindent
第 3 步。注意，由引理 \ref{algebra-lemma-etale-quasi-finite}，
$R \to S$ 拟有限。由引理
\ref{algebra-lemma-quasi-finite-open-integral-closure}，存在有限环映射
$R \to S'$、$R$-代数映射 $S' \to S$，以及元素
$g' \in S'$，满足 $g' \not \in \mathfrak q$，使得
$S' \to S$ 诱导同构 $S'_{g'} \cong S_{g'}$。
% P01-E0884: the frozen sentence repeats “such that”, and writes g' notin q although g' is introduced in S' while q is an ideal of S without explicitly invoking its image.
（当然，一般而言，$S'$ 在 $R$ 上并不平展。）所以可以假设：
(a) $R$ 是 Noether 环；(b) $R \to S$ 有限；(c) $R \to S$
在 $\mathfrak q$ 处平展（但不再必然在所有素理想处平展）。

\medskip\noindent
第 4 步。设 $\mathfrak p \subset R$ 是对应于 $\mathfrak q$
的素理想。考虑纤维环 $S \otimes_R \kappa(\mathfrak p)$。
这是 $\kappa(\mathfrak p)$ 上的有限代数，因此是 Artin 环
（见引理 \ref{algebra-lemma-finite-dimensional-algebra}），从而是局部环的
有限乘积
$$
S \otimes_R \kappa(\mathfrak p) = \prod\nolimits_{i = 1}^n A_i
$$
（见命题 \ref{algebra-proposition-dimension-zero-ring}）。其中一个因子，比如
$A_1$，是局部环 $S_{\mathfrak q}/\mathfrak pS_{\mathfrak q}$，
它同构于 $\kappa(\mathfrak q)$，见引理
\ref{algebra-lemma-etale-at-prime}。其他因子对应于 $S$ 中位于
$\mathfrak p$ 之上的其他素理想，比如
$\mathfrak q_2, \ldots, \mathfrak q_n$。
% P01-E0885: the frozen display is followed directly by “see Proposition” without sentence punctuation.

\medskip\noindent
第 5 步。可以选择非零元素 $\alpha \in \kappa(\mathfrak q)$，
它生成有限可分域扩张
$\kappa(\mathfrak q)/\kappa(\mathfrak p)$（因此即使域扩张平凡，
也不允许 $\alpha = 0$）。注意，对任意
$\lambda \in \kappa(\mathfrak p)^*$，元素 $\lambda \alpha$ 也在
$\kappa(\mathfrak p)$ 上生成 $\kappa(\mathfrak q)$。考虑元素
$$
\overline{t} =
(\alpha, 0, \ldots, 0) \in
\prod\nolimits_{i = 1}^n A_i =
S \otimes_R \kappa(\mathfrak p).
$$
必要时如上用 $\lambda \alpha$ 替换 $\alpha$，可以假设
$\overline{t}$ 是某个 $t \in S$ 的像。设 $I \subset R[x]$
是 $R$-代数映射 $R[x] \to S$ 的核，其中该映射把 $x$ 映到
$t$。令 $S' = R[x]/I$，于是 $S' \subset S$。有交换图
$$
\xymatrix{
R[x] \ar[r] & S' \ar[r] & S \\
R \ar[u] \ar[ru] \ar[rru] & &
}
$$
由构造，$S$ 的素理想 $\mathfrak q_j$（$j \geq 2$）都位于
$R[x]$ 的素理想 $(\mathfrak p, x)$ 之上，而素理想
$\mathfrak q$ 位于 $R[x]$ 的另一个素理想之上，因为
$\alpha \not = 0$。

\medskip\noindent
第 6 步。记 $\mathfrak q' \subset S'$ 为 $S'$ 中对应于
$\mathfrak q$ 的素理想。由上述讨论，$\mathfrak q$ 是 $S$
中唯一位于 $\mathfrak q'$ 之上的素理想。因此
$S_{\mathfrak q} = S_{\mathfrak q'}$，见引理
\ref{algebra-lemma-unique-prime-over-localize-below}（对 $S' \to S$
有上升性质，见引理 \ref{algebra-lemma-integral-going-up}，因为
$S' \to S$ 有限，这是由于 $R \to S$ 有限）。
% P01-E0886: the frozen source says “Denote q' ... the prime” and omits “by/as”.
% P01-E0887: the frozen parenthesis says “we have going up” rather than “going up holds”.
由此，$S'_{\mathfrak q'} \to S_{\mathfrak q}$ 有限且为单射，
因为它是有限单射环映射 $S' \to S$ 的局部化。考虑局部环映射
$$
R_{\mathfrak p} \to S'_{\mathfrak q'} \to S_{\mathfrak q}
$$
第二个映射有限且为单射。由引理 \ref{algebra-lemma-etale-at-prime}，
$S_{\mathfrak q}/\mathfrak pS_{\mathfrak q} = \kappa(\mathfrak q)$。
从而更有
$S_{\mathfrak q}/\mathfrak q'S_{\mathfrak q} = \kappa(\mathfrak q)$。
由于
$$
\kappa(\mathfrak p) \subset \kappa(\mathfrak q') \subset \kappa(\mathfrak q)
$$
并且 $\alpha$ 属于 $\kappa(\mathfrak q')$ 在
$\kappa(\mathfrak q)$ 中的像，可知
$\kappa(\mathfrak q') = \kappa(\mathfrak q)$。于是，对
$S'_{\mathfrak q'}$-模映射
$S'_{\mathfrak q'} \to S_{\mathfrak q}$ 应用 Nakayama 引理
\ref{algebra-lemma-NAK}，可知映射
$S'_{\mathfrak q'} \to S_{\mathfrak q}$ 满。换言之，
$S'_{\mathfrak q'} \cong S_{\mathfrak q}$。

\medskip\noindent
第 7 步。由引理 \ref{algebra-lemma-isomorphic-local-rings}，存在
$g \in S$，$g \not \in \mathfrak q$ 和
$g' \in S'$，$g' \not \in \mathfrak q'$，使得
$S'_{g'} \cong S_g$。由于 $R$ 是 Noether 环，$S'$ 在 $R$
上有限，因为它是有限 $R$-模 $S$ 的 $R$-子模。因此用 $S'$
替换 $S$ 后，可以假设：(a) $R$ 是 Noether 环；
(b) $S$ 在 $R$ 上有限；(c) $S$ 在 $\mathfrak q$ 处相对于
$R$ 平展；(d) $S = R[x]/I$。
% P01-E0888: item (b) in the frozen source omits the verb “is” before “finite over R”.

\medskip\noindent
第 8 步。考虑环
$S \otimes_R \kappa(\mathfrak p) = \kappa(\mathfrak p)[x]/\overline{I}$，
其中 $\overline{I} = I \cdot \kappa(\mathfrak p)[x]$ 是 $I$ 在
$\kappa(\mathfrak p)[x]$ 中生成的理想。由于
$\kappa(\mathfrak p)[x]$ 是主理想整环，可知对某个首一多项式
$\overline{h} \in \kappa(\mathfrak p)[x]$，有
$\overline{I} = (\overline{h})$。用某个
$\lambda \in \kappa(\mathfrak p)$ 的
$\lambda \cdot \overline{h}$ 替换 $\overline{h}$ 后，可以假设
$\overline{h}$ 是某个 $h \in I \subset R[x]$ 的像。
（问题在于我们不知道能否选择首一的 $h$。）此外，与第 4 步中一样，
有 $S \otimes_R \kappa(\mathfrak p) = A_1 \times \ldots \times A_n$，
其中 $A_1 = \kappa(\mathfrak q)$ 是
$\kappa(\mathfrak p)$ 的有限可分扩张，而
$A_2, \ldots, A_n$ 是局部环。这蕴含
$$
\overline{h} = \overline{h}_1 \overline{h}_2^{e_2} \ldots \overline{h}_n^{e_n}
$$
其中 $\overline{h}_i \in \kappa(\mathfrak p)[x]$ 是两两互素的
不可约首一多项式，$e_2, \ldots, e_n \geq 1$。编号的选取使得
作为 $\kappa(\mathfrak p)[x]$-代数，
$A_i = \kappa(\mathfrak p)[x]/(\overline{h}_i^{e_i})$。
% P01-E0889: the frozen formula uses e_i without defining e_1, while only e_2,...,e_n are introduced and the A_1 factor is separately a field.
注意，$\overline{h}_1$ 是 $\alpha \in \kappa(\mathfrak q)$
的极小多项式，因而是可分多项式（其导数与它本身互素）。

\medskip\noindent
第 9 步。设 $m \in I$ 是首一元素；这样的元素存在，因为环扩张
$R \to R[x]/I$ 有限，因而整。记 $\overline{m}$ 为它在
$\kappa(\mathfrak p)[x]$ 中的像。可以分解为
% P01-E0890: the frozen source says “Denote mbar the image” and omits “by/as”.
$$
\overline{m} = \overline{k}
\overline{h}_1^{d_1} \overline{h}_2^{d_2} \ldots \overline{h}_n^{d_n}
$$
其中 $d_1 \geq 1$，$d_j \geq e_j$，$j = 2, \ldots, n$，
并且 $\overline{k} \in \kappa(\mathfrak p)[x]$ 与所有
$\overline{h}_i$ 互素。令 $f = m^l + h$，其中
$l \deg(m) > \deg(h)$，且 $l \geq 2$。于是 $f$ 是 $R$
上的首一多项式。此外，$f$ 在 $\kappa(\mathfrak p)[x]$
中的像 $\overline{f}$ 分解为
$$
\overline{f} =
\overline{h}_1 \overline{h}_2^{e_2} \ldots \overline{h}_n^{e_n}
+
\overline{k}^l \overline{h}_1^{ld_1} \overline{h}_2^{ld_2}
\ldots \overline{h}_n^{ld_n}
=
\overline{h}_1(\overline{h}_2^{e_2} \ldots \overline{h}_n^{e_n}
+
\overline{k}^l
\overline{h}_1^{ld_1 - 1} \overline{h}_2^{ld_2} \ldots \overline{h}_n^{ld_n})
= \overline{h}_1 \overline{w}
$$
其中 $\overline{w}$ 是与 $\overline{h}_1$ 互素的多项式。
令 $g = f'$（关于 $x$ 的导数）。

\medskip\noindent
第 10 步。环映射 $R[x] \to S = R[x]/I$ 具有以下性质：
(1) 它把 $f$ 映到零；(2) 它把 $g$ 映到
$S \setminus \mathfrak q$ 中的元素。第一条断言是显然的，因为
$f$ 是 $I$ 的元素。要证明第二条断言，只需说明 $g$ 在
$\kappa(\mathfrak q) = \kappa(\mathfrak p)[x]/(\overline{h}_1)$
中的像不为零。$g$ 在 $\kappa(\mathfrak p)[x]$ 中的像是
$\overline{f}$ 的导数。因此 (2) 是显然的，因为
$$
\overline{g} =
\frac{\text{d}\overline{f}}{\text{d}x} =
\overline{w}\frac{\text{d}\overline{h}_1}{\text{d}x} +
\overline{h}_1\frac{\text{d}\overline{w}}{\text{d}x},
$$
$\overline{w}$ 与 $\overline{h}_1$ 互素，而
$\overline{h}_1$ 可分。

\medskip\noindent
第 11 步。由此可知，$\varphi : R[x]/(f) \to S$ 是满环映射，
$R[x]_g/(f)$ 在 $R$ 上平展（因为它标准平展，见引理
\ref{algebra-lemma-standard-etale}），且 $\varphi(g) \not \in \mathfrak q$。
选择元素 $g' \in R[x]/(f)$，使得
$\varphi(g') \not \in \mathfrak q$，且 $S_{\varphi(g')}$
在 $R$ 上也平展（这样的元素存在，因为 $S$ 在 $\mathfrak q$
处相对于 $R$ 平展）。于是环映射
$R[x]_{gg'}/(f) \to S_{\varphi(gg')}$ 是 $R$ 上平展代数之间的
满射。因此由引理 \ref{algebra-lemma-map-between-etale}，它平展。
再由引理 \ref{algebra-lemma-surjective-flat-finitely-presented}，
它是局部化。这样，$S$ 在某个不属于 $\mathfrak q$ 的元素处的
局部化，同构于 $R$ 上某个标准平展代数的局部化，这正是所要证明的。
\end{proof}

\noindent
下面两个引理说明，平展拓扑比由 Zariski 覆盖和有限平坦态射生成的
拓扑更粗。初次阅读时应当跳过它们。

\begin{lemma}
\label{algebra-lemma-standard-etale-finite-flat-Zariski}
设 $R \to S$ 是标准平展态射。则存在具有下列性质的环映射
$R \to S'$：
% P01-E0891: the frozen source omits a colon after “the following properties”.
\begin{enumerate}
\item $R \to S'$ 有限、有限表现且平坦
（换言之，$S'$ 作为 $R$-模有限投射）；
\item $\Spec(S') \to \Spec(R)$ 满射；
\item 对每个位于 $\mathfrak p \subset R$ 之上的素理想
$\mathfrak q \subset S$，以及每个位于 $\mathfrak p$ 之上的素理想
$\mathfrak q' \subset S'$，存在 $g' \in S'$，
$g' \not \in \mathfrak q'$，使得环映射 $R \to S'_{g'}$
通过映射 $\varphi : S \to S'_{g'}$ 分解，并满足
$\varphi^{-1}(\mathfrak q'S'_{g'}) = \mathfrak q$。
\end{enumerate}
\end{lemma}

\begin{proof}
设 $S = R[x]_g/(f)$ 是定义 \ref{algebra-definition-standard-etale}
中那样的 $S$ 的表现。写成
$f = x^n + a_1 x^{n - 1} + \ldots + a_n$，其中 $a_i \in R$。
由引理 \ref{algebra-lemma-adjoin-roots}，存在有限自由且忠实平坦的环映射
$R \to S'$，使得对某些 $\alpha_i \in S'$，有
$f = \prod (x - \alpha_i)$。因此 $R \to S'$ 满足条件 (1)、(2)。
设 $\mathfrak q \subset R[x]/(f)$ 是满足
$g \not \in \mathfrak q$ 的素理想（即它对应于 $S$ 的一个素理想）。
令 $\mathfrak p = R \cap \mathfrak q$，并设
$\mathfrak q' \subset S'$ 是位于 $\mathfrak p$ 之上的素理想。
注意，存在 $n$ 个 $R$-代数映射
\begin{eqnarray*}
\varphi_i : R[x]/(f) & \longrightarrow & S' \\
x & \longmapsto & \alpha_i
\end{eqnarray*}
为完成证明，需要说明对某个 $i$，有
(a) $\varphi_i(g)$ 在 $\kappa(\mathfrak q')$ 中的像不为零；
(b) $\varphi_i^{-1}(\mathfrak q') = \mathfrak q$。
因为这时只需取 $g' = \varphi_i(g)$，并对该 $i$ 取
$\varphi = \varphi_i$。

\medskip\noindent
设 $\overline{f}$ 表示 $f$ 在 $\kappa(\mathfrak p)[x]$ 中的像。
注意，作为 $\Spec(\kappa(\mathfrak p)[x]/(\overline{f}))$
的一个点，素理想 $\mathfrak q$ 对应于 $\overline{f}$ 的某个
不可约因子 $f_1$。此外，$g \not \in \mathfrak q$ 意味着
$f_1$ 不整除 $g$ 在 $\kappa(\mathfrak p)[x]$ 中的像
$\overline{g}$。记
$\overline{\alpha}_1, \ldots, \overline{\alpha}_n$ 为
$\alpha_1, \ldots, \alpha_n$ 在 $\kappa(\mathfrak q')$ 中的像。
% P01-E0892: the frozen source says “Denote alpha-bar_1,... the images” and omits “by/as”.
注意，多项式 $\overline{f}$ 在 $\kappa(\mathfrak q')[x]$
中完全分裂，即
$$
\overline{f} = \prod\nolimits_i (x - \overline{\alpha}_i)
$$
此外，$\varphi_i(g)$ 约化为
$\overline{g}(\overline{\alpha}_i)$。因此可以选择 $i$，使得
$f_1(\overline{\alpha}_i) = 0$ 且
$\overline{g}(\overline{\alpha}_i) \not = 0$。
对这个 $i$，性质 (a)、(b) 成立。略去一些细节。
\end{proof}

\begin{lemma}
\label{algebra-lemma-etale-finite-flat-zariski}
设 $R \to S$ 是环映射。假设
\begin{enumerate}
\item $R \to S$ 平展；
\item $\Spec(S) \to \Spec(R)$ 满射。
\end{enumerate}
则存在环映射 $R \to S'$，使得
\begin{enumerate}
\item $R \to S'$ 有限、有限表现且平坦
（换言之，它作为 $R$-模有限投射）；
\item $\Spec(S') \to \Spec(R)$ 满射；
\item 对每个素理想 $\mathfrak q' \subset S'$，存在
$g' \in S'$，$g' \not \in \mathfrak q'$，使得环映射
$R \to S'_{g'}$ 分解为 $R \to S \to S'_{g'}$。
\end{enumerate}
\end{lemma}

\begin{proof}
由命题 \ref{algebra-proposition-etale-locally-standard} 和
$\Spec(S)$ 的拟紧性（见引理 \ref{algebra-lemma-quasi-compact}），
可以找到生成 $S$ 的单位理想的 $g_1, \ldots, g_n \in S$，
使得每个 $R \to S_{g_i}$ 都标准平展。若对环映射
$R \to \prod_{i = 1, \ldots, n} S_{g_i}$ 证明了引理，
则对环映射 $R \to S$ 也得到引理。因此可以假设
$S = \prod_{i = 1, \ldots, n} S_i$ 是标准平展态射的有限乘积。

\medskip\noindent
对每个 $i$，选取引理
\ref{algebra-lemma-standard-etale-finite-flat-Zariski} 中适用于标准平展态射
$R \to S_i$ 的环映射 $R \to S_i'$。令
$S' = S_1' \otimes_R \ldots \otimes_R S_n'$；下文不再另行说明地
使用 $R$-代数映射 $S_i' \to S'$。我们断言这给出所需构造。
性质 (1)、(2) 是直接的。对性质 (3)，假设
$\mathfrak q' \subset S'$ 是素理想。记 $\mathfrak p$ 为它在
$\Spec(R)$ 中的像。选择 $i \in \{1, \ldots, n\}$，使得
$\mathfrak p$ 属于 $\Spec(S_i) \to \Spec(R)$ 的像；由假设，
这是可能的。令 $\mathfrak q_i' \subset S_i'$ 为
$\mathfrak q'$ 在 $S_i'$ 的谱中的像。
% P01-E0893: the frozen source says “Denote p its image” and omits “by/as”.
% P01-E0894: the frozen source says “Set q_i' ... the image” and omits “to be/as”.
由 $S'_i$ 的构造，存在 $g'_i \in S_i'$，使得
$R \to (S_i')_{g_i'}$ 分解为
$R \to S_i \to (S_i')_{g_i'}$。因此
$R \to S'_{g_i'}$ 也分解为
% P01-E0895: the frozen proof invokes the preceding lemma but omits the inherited condition g'_i notin q'_i, which is needed to match property (3)'s neighborhood statement.
$$
R \to S_i \to (S_i')_{g_i'} \to S'_{g_i'}
$$
如所要求。
\end{proof}

\section{拟有限环映射的平展局部结构}
\label{algebra-section-etale-local-quasi-finite}

\noindent
下面这些引理粗略地说，经过平展扩张后，
拟有限环映射会成为有限环映射。
为帮助理解这些结果，回顾有限型环映射为拟有限的轨迹是开集
（见引理 \ref{algebra-lemma-quasi-finite-open}），而且这一轨迹的形成
与任意基变换交换
（见引理 \ref{algebra-lemma-quasi-finite-base-change}）。

\begin{lemma}
\label{algebra-lemma-produce-finite}
设 $R \to S' \to S$ 是环映射。
设 $\mathfrak p \subset R$ 是素理想。
设 $g \in S'$ 是元素。
假设
\begin{enumerate}
\item $R \to S'$ 是整的；
\item $R \to S$ 是有限型的；
\item $S'_g \cong S_g$；并且
\item $g$ 在 $S' \otimes_R \kappa(\mathfrak p)$ 中可逆。
% P01-E0896: the frozen source says “g invertible” and omits “is”.
\end{enumerate}
则存在 $f \in R$，$f \not \in \mathfrak p$，使得
$R_f \to S_f$ 有限。
% P01-E0897: the frozen source uses the ungrammatical article phrase “a f”.
\end{lemma}

\begin{proof}
由假设，$V(g) \subset \Spec(S')$ 在态射
$\Spec(S') \to \Spec(R)$ 下的像 $T$ 不含
$\mathfrak p$。由第 \ref{algebra-section-going-up} 节，
特别是引理 \ref{algebra-lemma-going-up-closed}，可见 $T$ 是闭集。
% P01-E0898: the frozen source has the malformed citation phrase “By Section ..., especially, Lemma ...”.
取 $f \in R$，$f \not \in \mathfrak p$，使得
$T \cap D(f) = \emptyset$。于是 $g$ 在 $S'_f$ 中变得可逆，
因而 $S'_f \cong S_f$。所以 $S_f$ 在 $R_f$ 上既是有限型的
又是整的，从而是有限的。
\end{proof}

\begin{lemma}
\label{algebra-lemma-etale-makes-quasi-finite-finite-one-prime}
设 $R \to S$ 是环映射。
设 $\mathfrak q \subset S$ 是位于
素理想 $\mathfrak p \subset R$ 之上的素理想。
假设 $R \to S$ 是有限型的，并且在 $\mathfrak q$ 处拟有限。
% P01-E0899: the frozen source says “Assume R -> S finite type” and omits “is”.
则存在：
\begin{enumerate}
\item 平展环映射 $R \to R'$；
\item 位于 $\mathfrak p$ 之上的素理想 $\mathfrak p' \subset R'$；
\item 乘积分解
$$
R' \otimes_R S = A \times B
$$
\end{enumerate}
% P01-E0900: the frozen source omits punctuation after “with the following properties”.
具有下列性质：
\begin{enumerate}
\item $\kappa(\mathfrak p) = \kappa(\mathfrak p')$；
\item $R' \to A$ 有限；
\item $A$ 恰有一个位于 $\mathfrak p'$ 之上的素理想
$\mathfrak r$；
\item $\mathfrak r$ 位于 $\mathfrak q$ 之上；并且
\item $B$ 不含同时位于 $\mathfrak q$ 和 $\mathfrak p'$ 之上的素理想。
\end{enumerate}
\end{lemma}

\begin{proof}
设 $S' \subset S$ 是 $R$ 在 $S$ 中的整闭包。
令 $\mathfrak q' = S' \cap \mathfrak q$。
由 Zariski 主定理 \ref{algebra-theorem-main-theorem}，
存在 $g \in S'$，$g \not \in \mathfrak q'$，使得
$S'_g \cong S_g$。考虑纤维环
$F = S \otimes_R \kappa(\mathfrak p)$ 和
$F' = S' \otimes_R \kappa(\mathfrak p)$。记
$\overline{\mathfrak q}'$ 为 $F'$ 中对应于 $\mathfrak q'$ 的素理想。
% P01-E0901: the frozen source says “Denote q-bar' the prime” and omits “by/as”.
因为 $F'$ 在 $\kappa(\mathfrak p)$ 上整，所以
$\overline{\mathfrak q}'$ 是 $\Spec(F')$ 的闭点，
见引理 \ref{algebra-lemma-integral-over-field}。
注意，$\mathfrak q$ 定义了 $\Spec(F)$ 的一个孤立闭点
$\overline{\mathfrak q}$
（见定义 \ref{algebra-definition-quasi-finite}）。
因为 $S'_g \cong S_g$，所以 $F'_g \cong F_g$，
从而 $\overline{\mathfrak q}$ 和 $\overline{\mathfrak q}'$
在 $\Spec(F)$ 与 $\Spec(F')$ 中有同构的开邻域。
我们断定集合
$\{\overline{\mathfrak q}'\} \subset \Spec(F')$
是开集。再结合 $\mathfrak q'$ 是闭点（已在上面证明），
便断定 $\overline{\mathfrak q}'$ 也定义了
$\Spec(F')$ 的一个孤立闭点。
% P01-E0902: the frozen source says q' is closed although the preceding argument proves q-bar' is closed.

\medskip\noindent
另有一个小观察：在映射
$\Spec(F) \to \Spec(F')$ 下，$\overline{\mathfrak q}$
是映到 $\overline{\mathfrak q}'$ 的唯一一点。
这由上面的讨论推出。

\medskip\noindent
由引理 \ref{algebra-lemma-disjoint-implies-product}，可以写成
$F' = F'_1 \times F'_2$，其中
$\Spec(F'_1) = \{\overline{\mathfrak q}'\}$。
因为 $F' = S' \otimes_R \kappa(\mathfrak p)$，
存在 $s' \in S'$，它映到元素
$(r, 0) \in F'_1 \times F'_2 = F'$，
其中某个 $r \in R$ 满足 $r \not \in \mathfrak p$。
事实上，我们将用到的关于 $s'$ 的事实是：它是 $S'$ 的元素，
不属于 $\mathfrak q'$，而属于任何其他位于
$\mathfrak p$ 之上的素理想。

\medskip\noindent
取首一多项式 $f(x) \in R[x]$，使得 $f(s') = 0$。
记 $\overline{f} \in \kappa(\mathfrak p)[x]$ 为其像。
% P01-E0903: the frozen source says “Denote f-bar the image” and omits “by/as”.
可以将它分解为 $\overline{f} = x^e \overline{h}$，其中
$\overline{h}(0) \not = 0$。必要时将 $f$ 替换为
$x f$，便可假设 $e \geq 1$。
由引理 \ref{algebra-lemma-factor-mod-lift-etale}，
可以找到平展环扩张 $R \to R'$、
位于 $\mathfrak p$ 之上的素理想 $\mathfrak p'$，以及
$R'[x]$ 中的分解 $f = h i$，使得
$\kappa(\mathfrak p) = \kappa(\mathfrak p')$、
$\overline{h} = h \bmod \mathfrak p'$、
$x^e = i \bmod \mathfrak p'$，并且
可以在 $R'[x]$ 中写成 $a h + b i = 1$
（其中 $a, b$ 适当）。

\medskip\noindent
考虑元素 $h(s'), i(s') \in R' \otimes_R S'$。
由构造，$h(s')i(s') = f(s') = 0$。另一方面，
它们生成单位理想，因为 $a(s')h(s') + b(s')i(s') = 1$。
所以 $R' \otimes_R S'$ 是在这些元素处的局部化之积：
$$
R' \otimes_R S'
=
(R' \otimes_R S')_{i(s')}
\times
(R' \otimes_R S')_{h(s')}
=
S'_1 \times S'_2
$$
此外，这一乘积分解与我们为纤维环 $F'$ 得到的乘积分解相容；
这是由我们对 $s', i, h$ 的选择得到的，因为这些选择保证
$\overline{\mathfrak q}'$ 是 $F'$ 中唯一不含
$i(s')$ 在 $F'$ 中之像的素理想。这里使用了如下事实：
$R'\otimes_R S'$ 在 $R'$ 的 $\mathfrak p'$ 处的纤维环
与 $F'$ 相同，因为
$\kappa(\mathfrak p) = \kappa(\mathfrak p')$。
由此，$S'_1$ 恰有一个位于 $\mathfrak p'$ 之上的素理想，
记为 $\mathfrak r'$，并且这个素理想位于
$\mathfrak q'$ 之上。
% P01-E0904: the frozen source contains doubled spacing in “S'_1  has”.
因此元素 $g \in S'$ 在 $S'_1$ 中的像不属于
$\mathfrak r'$。

\medskip\noindent
基变换 $R'\otimes_R S$ 继承类似的乘积分解
$$
R' \otimes_R S
=
(R' \otimes_R S)_{i(s')}
\times
(R' \otimes_R S)_{h(s')}
=
S_1 \times S_2
$$
由上可知，$S_1$ 恰有一个位于 $\mathfrak p'$ 之上的素理想，
记为 $\mathfrak r$（如上考虑纤维环即可），
并且这个素理想位于 $\mathfrak q$ 之上。

\medskip\noindent
现在对环映射 $R' \to S'_1 \to S_1$、
素理想 $\mathfrak p'$ 和元素 $g$ 应用引理
\ref{algebra-lemma-produce-finite}，可知在将 $R'$ 替换为某个主局部化后，
可以假设 $S_1$ 在 $R'$ 上有限，正如所需。
\end{proof}

\begin{lemma}
\label{algebra-lemma-etale-makes-quasi-finite-finite}
设 $R \to S$ 是环映射。
设 $\mathfrak p \subset R$ 是素理想。
假设 $R \to S$ 是有限型的。
则存在：
% P01-E0905: the frozen source introduces the following enumeration without punctuation after “there exists”.
\begin{enumerate}
\item 平展环映射 $R \to R'$；
\item 位于 $\mathfrak p$ 之上的素理想 $\mathfrak p' \subset R'$；
\item 乘积分解
$$
R' \otimes_R S = A_1 \times \ldots \times A_n \times B
$$
\end{enumerate}
具有下列性质：
% P01-E0906: the frozen source omits punctuation after “with the following properties”.
\begin{enumerate}
\item 有 $\kappa(\mathfrak p) = \kappa(\mathfrak p')$；
\item 每个 $A_i$ 在 $R'$ 上有限；
\item 每个 $A_i$ 恰有一个位于 $\mathfrak p'$ 之上的素理想
$\mathfrak r_i$；并且
\item $R' \to B$ 在任何位于 $\mathfrak p'$ 之上的素理想处
都不是拟有限的。
% P01-E0907: the frozen source says “R' -> B not quasi-finite” and omits “is”.
\end{enumerate}
\end{lemma}

\begin{proof}
记 $F = S \otimes_R \kappa(\mathfrak p)$ 为 $S/R$
在素理想 $\mathfrak p$ 处的纤维环。
% P01-E0908: the frozen source says “Denote F ... the fibre ring” and omits “by/as”.
由于 $F$ 在 $\kappa(\mathfrak p)$ 上有限型，所以它是 Noether 环，
因而 $\Spec(F)$ 只有有限多个孤立闭点。若没有孤立闭点，
即不存在 $S$ 中位于 $\mathfrak p$ 之上的素理想
$\mathfrak q$，使得 $S/R$ 在 $\mathfrak q$ 处拟有限，
那么引理成立。
若至少存在一个这样的素理想 $\mathfrak q$，则可以应用引理
\ref{algebra-lemma-etale-makes-quasi-finite-finite-one-prime}。
这给出该引理中的图表
$$
\xymatrix{
S \ar[r] & R'\otimes_R S \ar@{=}[r] & A_1 \times B' \\
R \ar[r] \ar[u] & R' \ar[u] \ar[ru]
}
$$
因为 $\mathfrak p$ 与 $\mathfrak p'$ 处的剩余域相同，
$S/R$ 与 $(A_1 \times B')/R'$ 的纤维环相同。
因此，对纤维中孤立闭点的数目作归纳，可以假设
引理对 $R' \to B'$ 和 $\mathfrak p'$ 成立。
于是得到平展环映射 $R' \to R''$、
素理想 $\mathfrak p'' \subset R''$，以及分解
$$
R'' \otimes_{R'} B' = A_2 \times \ldots \times A_n \times B
$$
% P01-E0909: the frozen source leaves the displayed decomposition without terminal punctuation.
略去如下验证：环映射 $R \to R''$、素理想
$\mathfrak p''$ 以及所得分解
$$
R'' \otimes_R S = (R'' \otimes_{R'} A_1) \times
A_2 \times \ldots \times A_n \times B
$$
给出引理中问题的一个解。
\end{proof}

\begin{lemma}
\label{algebra-lemma-etale-makes-quasi-finite-finite-variant}
设 $R \to S$ 是环映射。
设 $\mathfrak p \subset R$ 是素理想。
假设 $R \to S$ 是有限型的。
% P01-E0910: the frozen source says “Assume R -> S finite type” and omits “is”.
则存在：
% P01-E0911: the frozen source introduces the following enumeration without punctuation after “there exists”.
\begin{enumerate}
\item 平展环映射 $R \to R'$；
\item 位于 $\mathfrak p$ 之上的素理想 $\mathfrak p' \subset R'$；
\item 乘积分解
$$
R' \otimes_R S = A_1 \times \ldots \times A_n \times B
$$
\end{enumerate}
具有下列性质：
% P01-E0912: the frozen source omits punctuation after “with the following properties”.
\begin{enumerate}
\item 每个 $A_i$ 在 $R'$ 上有限；
\item 每个 $A_i$ 恰有一个位于 $\mathfrak p'$ 之上的素理想
$\mathfrak r_i$；
\item 有限域扩张
$\kappa(\mathfrak r_i)/\kappa(\mathfrak p')$
是纯不可分的；并且
\item $R' \to B$ 在任何位于 $\mathfrak p'$ 之上的素理想处
都不是拟有限的。
% P01-E0913: the frozen source says “R' -> B not quasi-finite” and omits “is”.
\end{enumerate}
\end{lemma}

\begin{proof}
证明的策略是作两次平展环扩张：先控制剩余域，
再应用引理 \ref{algebra-lemma-etale-makes-quasi-finite-finite}。

\medskip\noindent
记 $F = S \otimes_R \kappa(\mathfrak p)$ 为 $S/R$
在素理想 $\mathfrak p$ 处的纤维环。
% P01-E0914: the frozen source says “Denote F ... the fibre ring” and omits “by/as”.
与引理 \ref{algebra-lemma-etale-makes-quasi-finite-finite} 的证明中一样，
存在有限多个素理想，记为
$\mathfrak q_1, \ldots, \mathfrak q_n$，它们是 $S$ 中位于
$R$ 之上的素理想，并且环映射 $R \to S$ 在这些素理想处拟有限。
% P01-E0915: the frozen source has the typo “finitely may primes”.
% P01-E0916: the frozen source says these primes of S lie over R, but the fibre argument requires them to lie over p.
设 $\kappa(\mathfrak p) \subset L_i \subset \kappa(\mathfrak q_i)$
是这样的子域：$\kappa(\mathfrak p) \subset L_i$
是可分扩张，而域扩张 $\kappa(\mathfrak q_i)/L_i$
是纯不可分的。设 $L/\kappa(\mathfrak p)$
是一个有限 Galois 扩张，使得对 $i = 1, \ldots, n$，
$L_i$ 可嵌入其中。
由引理 \ref{algebra-lemma-make-etale-map-prescribed-residue-field}，
可以找到平展环扩张
$R \to R'$ 以及位于 $\mathfrak p$ 之上的素理想
$\mathfrak p'$，使得域扩张
$\kappa(\mathfrak p')/\kappa(\mathfrak p)$ 同构于
$\kappa(\mathfrak p) \subset L$。
因此 $R' \otimes_R S$ 在 $\mathfrak p'$ 处的纤维环同构于
$F \otimes_{\kappa(\mathfrak p)} L$。
位于 $\mathfrak q_i$ 之上的素理想对应于
$\kappa(\mathfrak q_i) \otimes_{\kappa(\mathfrak p)} L$
的素理想；由 $L$ 的选择和初等域论，后者是
在 $L$ 上纯不可分的域之积。
由引理 \ref{algebra-lemma-quasi-finite-base-change}，
这些也是位于 $\mathfrak p'$ 之上、使
$R' \to R' \otimes_R S$ 在其处拟有限的全部素理想。
因此，将 $R$ 替换为 $R'$、将 $\mathfrak p$ 替换为
$\mathfrak p'$、并将 $S$ 替换为 $R' \otimes_R S$ 后，
可以假设：对所有位于 $\mathfrak p$ 之上且使
$S/R$ 在其处拟有限的素理想 $\mathfrak q$，
域扩张
$\kappa(\mathfrak q)/\kappa(\mathfrak p)$
都是纯不可分的。

\medskip\noindent
接着应用引理 \ref{algebra-lemma-etale-makes-quasi-finite-finite}。
这就得到所需结果，因为在这次平展环扩张下域扩张不变。
\end{proof}

\section{局部同态}
\label{algebra-section-local-homomorphisms}

\noindent
这里收录若干没有自然所属章节的引理。
第一个引理粗略地说，局部环之间的平展映射，
模掉一个由非单位元生成的主理想的任意次幂后都是同构。

\begin{lemma}
\label{algebra-lemma-lindel}
\begin{reference}
\cite[Lemma on page 321]{Lindel}, \cite[Lemma 4.1.5]{KC}
\end{reference}
设 $(R, \mathfrak m_R) \to (S, \mathfrak m_S)$ 是局部环的
局部同态。假设 $S$ 是 $R$ 的某个平展环扩张的局部化，
并且 $\kappa(\mathfrak m_R) \to \kappa(\mathfrak m_S)$
是同构。则存在 $t \in \mathfrak m_R$，使得对所有
$n \geq 1$，$R/t^nR \to S/t^nS$ 都是同构。
% P01-E0917: the frozen source uses the ungrammatical article phrase “an t”.
% P01-E0918: the frozen source misspells “isomorphism” as “isomorphsm”.
\end{lemma}

\begin{proof}
写成 $S = T_{\mathfrak q}$，其中 $T$ 是某个平展
$R$-代数，$\mathfrak q \subset T$ 是位于
$\mathfrak m_R$ 之上的素理想。
由命题 \ref{algebra-proposition-etale-locally-standard}，
可以假设 $R \to T$ 标准平展。
按照定义 \ref{algebra-definition-standard-etale}，写成
$T = R[x]_g/(f)$。由关于剩余域的假设，可以选择
$a \in R$，使得 $x$ 和 $a$ 在
$\kappa(\mathfrak q) = \kappa(\mathfrak m_S) = \kappa(\mathfrak m_R)$
中的像相同。于是用 $x - a$ 替换 $x$ 后，可以假设
$\mathfrak q$ 在 $T$ 中由 $x$ 和 $\mathfrak m_R$ 生成。
特别地，$t = f(0) \in \mathfrak m_R$。
下面证明 $t = f(0)$ 满足要求。

\medskip\noindent
写成 $f = x^d + \sum_{i = 1, \ldots, d - 1} a_i x^i + t$。
因为 $R \to T$ 标准平展，可知 $a_1$ 是 $R$ 中的单位元：
$f$ 的导数在 $T$ 中可逆，特别地，它不属于
$\mathfrak q$。
% P01-E0919: the frozen source joins the derivative clauses without the needed conjunction or punctuation.
令 $h = a_1 + a_2 x + \ldots + a_{d - 1} x^{d - 2} + x^{d - 1} \in R[x]$，
从而在 $R[x]$ 中有 $f = t + xh$。可见
$h \not \in \mathfrak q$，因此可以用
$R[x]_{hg}/(f)$ 替换 $T$。替换后可见
$$
T/tT = (R/tR)[x]_{hg}/(f) = (R/tR)[x]_{hg}/(xh) = (R/tR)[x]_{hg}/(x)
$$
是 $R/tR$ 的商。由引理
\ref{algebra-lemma-surjective-mod-locally-nilpotent}，可知对所有
$n \geq 1$，$R/t^nR \to T/t^nT$ 都满射。
另一方面，已知平坦局部环映射
$R/t^nR \to S/t^nS$ 对每个 $n$ 都通过
$R/t^nR \to T/t^nT$ 分解，因此这些映射也单射
（局部环之间的平坦局部同态忠实平坦，因而单射，见
引理 \ref{algebra-lemma-local-flat-ff} 和
\ref{algebra-lemma-faithfully-flat-universally-injective}）。
因为 $S$ 是 $T$ 的局部化，所以 $S/t^nS$ 是
$T/t^nT = R/t^nR$ 在位于极大理想之上的某个素理想处的局部化；
但这个环已经是局部环，证明完成。
\end{proof}

\begin{lemma}
\label{algebra-lemma-etale-under-finite-flat}
设 $(R, \mathfrak m_R) \to (S, \mathfrak m_S)$ 是局部环的
局部同态。假设 $S$ 是 $R$ 的某个平展环扩张的局部化。
则存在有限、有限表现且忠实平坦的环映射 $R \to S'$，
使得对 $S'$ 的每个极大理想 $\mathfrak m'$，环映射
$R \to S'_{\mathfrak m'}$ 都有分解
$$
R \to S \to S'_{\mathfrak m'}.
$$
\end{lemma}

\begin{proof}
写成 $S = T_{\mathfrak q}$，其中 $T$ 是某个平展
$R$-代数。由命题 \ref{algebra-proposition-etale-locally-standard}，
可以假设 $T$ 标准平展。
对环映射 $R \to T$ 应用引理
\ref{algebra-lemma-standard-etale-finite-flat-Zariski}，得到
$R \to S'$。特别地，对 $S'$ 的每个极大理想
$\mathfrak m'$，都得到某个 $g' \not \in \mathfrak m'$ 和分解
$\varphi : T \to S'_{g'}$，使得
$\mathfrak q = \varphi^{-1}(\mathfrak m'S'_{g'})$。
因此 $\varphi$ 诱导所需的局部环映射
$S \to S'_{\mathfrak m'}$。
\end{proof}

\section{整闭包与光滑基变换}
\label{algebra-section-integral-closure-smooth-base-change}

\begin{lemma}
\label{algebra-lemma-trick}
设 $R$ 是环。
设 $f \in R[x]$ 是首一多项式。
设 $R \to B$ 是环映射。
若 $h \in B[x]/(f)$ 在 $R$ 上整，则元素
$f' h$ 可以写成 $f'h = \sum_i b_i x^i$，其中
$b_i \in B$ 在 $R$ 上整。
\end{lemma}

\begin{proof}
设在环 $B[x]/(f)$ 中有
$h^e + r_1 h^{e - 1} + \ldots + r_e = 0$，其中
$r_i \in R$。
存在有限自由环扩张 $B \subset B'$，使得对某些
$\alpha_i \in B'$ 有
$f = (x - \alpha_1) \ldots (x - \alpha_d)$，
见引理 \ref{algebra-lemma-adjoin-roots}。
注意，每个 $\alpha_i$ 都在 $R$ 上整。
可以用 $h_i \in B$ 将 $h$ 表成
$h = h_0 + h_1 x + \ldots + h_{d - 1} x^{d - 1}$。
于是有普遍恒等式
$$
f' h =
\sum\nolimits_{i = 1, \ldots, d}
h(\alpha_i)
(x - \alpha_1) \ldots \widehat{(x - \alpha_i)} \ldots (x - \alpha_d)
$$
作为 $B'[x]/(f)$ 的元素。要证明这一点，只需在环
$B_{univ} = \mathbf{Z}[\alpha_i, h_j]$ 上的各点
$\alpha_i$ 处计算等式两边（略去一些细节）。
由假设，h 在环 $B[x]/(f)$ 中满足
$h^e + r_1 h^{e - 1} + \ldots + r_e = 0$，可见
$$
h(\alpha_i)^e + r_1 h(\alpha_i)^{e - 1} + \ldots + r_e = 0
$$
在 $B'$ 中成立。因此 $h(\alpha_i)$ 在 $R$ 上整。
使用上面的公式可知，在 $B'[x]/(f)$ 中有
$f'h \equiv \sum_{j = 0, \ldots, d - 1} b'_j x^j$，
其中 $b'_j \in B'$ 在 $R$ 上整。然而，因为
$f' h \in B[x]/(f)$，而且 $1, x, \ldots, x^{d - 1}$
是 $B'[x]/(f)$ 的一个 $B'$-基，所以如所需，
$b'_j \in B$。
\end{proof}

\begin{lemma}
\label{algebra-lemma-integral-closure-commutes-etale}
设 $R \to S$ 是平展环映射。
设 $R \to B$ 是任意环映射。
设 $A \subset B$ 是 $R$ 在 $B$ 中的整闭包。
设 $A' \subset S \otimes_R B$ 是 $S$ 在
$S \otimes_R B$ 中的整闭包。则典范映射
$S \otimes_R A \to A'$ 是同构。
\end{lemma}

\begin{proof}
映射 $S \otimes_R A \to A'$ 单射，因为 $A \subset B$，
且 $R \to S$ 平坦。我们将反复使用如下事实：
取整闭包与局部化交换，见引理
\ref{algebra-lemma-integral-closure-localize}。
因此，由引理 \ref{algebra-lemma-cover}（判断 $S$-模映射是否为同构
的准则），可以在 $S$ 上局部化。
于是可以假设
$S = R[x]_g/(f) = (R[x]/(f))_g$
在 $R$ 上标准平展，见命题
\ref{algebra-proposition-etale-locally-standard}。
再作一次局部化，可见 $A'$ 是 $(A'')_g$，其中
$A''$ 是 $R[x]/(f)$ 在 $B[x]/(f)$ 中的整闭包。
假设 $a \in A''$。只需证明 $a$ 属于
$S \otimes_R A$。由引理 \ref{algebra-lemma-trick}，
$f' a = \sum a_i x^i$，其中 $a_i \in A$。
因为 $f'$ 在 $S$ 中可逆（由标准平展环映射的定义），
所以如所需，$a \in S \otimes_R A$。
\end{proof}

\begin{example}
\label{algebra-example-fourier}
设 $p$ 是素数。环扩张
$$
R = \mathbf{Z}[1/p] \subset
R' = \mathbf{Z}[1/p][x]/(x^{p - 1} + \ldots + x + 1)
$$
具有如下性质：对 $d < p$，存在元素
$\alpha_0, \ldots, \alpha_{d - 1} \in R'$，使得
$$
\prod\nolimits_{0 \leq i < j < d} (\alpha_i - \alpha_j)
$$
是 $R'$ 中的单位元。具体地，对
$i = 0, \ldots, p - 1$，取 $\alpha_i$ 为
$R'$ 中 $x^i$ 的类。于是，在 $R'[T]$ 中有
$$
T^p - 1 = \prod\nolimits_{i = 0, \ldots, p - 1} (T - \alpha_i)
$$
事实上，环 $\mathbf{Q}[x]/(x^{p - 1} + \ldots + x + 1)$
是域，因为分圆多项式
$x^{p - 1} + \ldots + x + 1$ 在 $\mathbf{Q}$ 上不可约，
而 $\alpha_i$ 是 $T^p - 1$ 的两两不同的根，故有该等式。
对等式两边求导，再代入 $T = \alpha_i$，得到
$$
p \alpha_i^{p - 1}
=
(\alpha_i - \alpha_1) \ldots
\widehat{(\alpha_i - \alpha_i)} \ldots
(\alpha_i - \alpha_1)
$$
% P01-E0920: the frozen derivative-product formula repeats alpha_1 at both endpoints; the full product should run from alpha_0 through alpha_{p-1}, omitting alpha_i.
由此可见，它在 $R'$ 中可逆。
\end{example}

\begin{lemma}
\label{algebra-lemma-integral-closure-commutes-smooth}
设 $R \to S$ 是光滑环映射。
设 $R \to B$ 是任意环映射。
设 $A \subset B$ 是 $R$ 在 $B$ 中的整闭包。
设 $A' \subset S \otimes_R B$ 是 $S$ 在
$S \otimes_R B$ 中的整闭包。则典范映射
$S \otimes_R A \to A'$ 是同构。
\end{lemma}

\begin{proof}
像引理 \ref{algebra-lemma-integral-closure-commutes-etale}
的证明那样论证，可以在 $S$ 上局部化。因此可以假设
$R \to S$ 是标准光滑环映射，见引理
\ref{algebra-lemma-smooth-syntomic}。由标准光滑环映射的定义，
$S$ 在多项式环 $R[x_1, \ldots, x_n]$ 上平展。
由于已经对平展环扩张证明了该结果（引理
\ref{algebra-lemma-integral-closure-commutes-etale}），
问题归结为 $S = R[x]$ 的情形。因此需要证明
$$
f = \sum b_i x^i
\text{ integral over }R[x]
\Leftrightarrow
\text{每个 }b_i\text{ integral over }R.
$$
从右向左的蕴含成立，因为 $B[x]$ 中在 $R[x]$ 上整的元素之集
是一个环（引理 \ref{algebra-lemma-integral-closure-is-ring}），
并且包含 $x$。

\medskip\noindent
假设 $f \in B[x]$ 在 $R[x]$ 上整，并且
$f = \sum_{i < d} b_i x^i$ 的次数 $< d$。
因为整闭包与局部化交换，只需证明存在不同的素数
$p, q$，使得每个 $b_i$ 既在 $R[1/p]$ 上整，
又在 $R[1/q]$ 上整。因此，可以找到有限自由环扩张
$R \subset R'$，使得 $R'$ 含有
$\alpha_1, \ldots, \alpha_d$，并且
$\prod_{i < j} (\alpha_i - \alpha_j)$ 是 $R'$ 中的单位元，
见例 \ref{algebra-example-fourier}。
% P01-E0921: the frozen proof invokes Example Fourier over R without explicitly replacing R by one of the localizations R[1/p] or R[1/q], although invertibility of the chosen integer prime is required.
在这种情形下，有普遍等式
$$
f
=
\sum_i
f(\alpha_i)
\frac{(x - \alpha_1) \ldots \widehat{(x - \alpha_i)} \ldots (x - \alpha_d)}
{(\alpha_i - \alpha_1) \ldots \widehat{(\alpha_i - \alpha_i)} \ldots
(\alpha_i - \alpha_d)}.
$$
此外，元素 $f(\alpha_i)$ 在 $R'$ 上整，因为
$(R' \otimes_R B)[x] \to R' \otimes_R B$，
$h \mapsto h(\alpha_i)$ 是环映射。
% P01-E0922: the frozen source begins this sentence with the editorial artifact “OK,”.
因此可见，$f$ 在 $(R' \otimes_R B)[x]$ 中的系数
在 $R'$ 上整。由于 $R'$ 在 $R$ 上有限（因而在 $R$ 上整），
可知这些系数也在 $R$ 上整，正如所需。
\end{proof}

\begin{lemma}
\label{algebra-lemma-integral-closure-commutes-colim-smooth}
设 $R \to S$ 和 $R \to B$ 是环映射。
设 $A \subset B$ 是 $R$ 在 $B$ 中的整闭包。
设 $A' \subset S \otimes_R B$ 是 $S$ 在
$S \otimes_R B$ 中的整闭包。若 $S$ 是光滑
$R$-代数的滤过余极限，则典范映射
$S \otimes_R A \to A'$ 是同构。
\end{lemma}

\begin{proof}
这由如下直接事实以及引理
\ref{algebra-lemma-integral-closure-commutes-smooth} 推出：
取张量积和取整闭包都与滤过余极限交换。
% P01-E0923: the frozen source has compound-subject agreement “taking ... and taking ... commutes”.
\end{proof}

\section{形式非分歧映射}
\label{algebra-section-formally-unramified}

\noindent
事实表明，在引入形式平展映射的概念之前先定义形式非分歧映射的概念，
在逻辑上更为高效。

\begin{definition}
\label{algebra-definition-formally-unramified}
设 $R \to S$ 是环映射。若对于每个由实线箭头组成的交换图
$$
\xymatrix{
S \ar[r] \ar@{-->}[rd] & A/I \\
R \ar[r] \ar[u] & A \ar[u]
}
$$
其中 $I \subset A$ 是平方为零的理想，使图交换的虚线箭头至多存在一条，
则称 $S$ 在 $R$ 上\textit{形式非分歧}。
\end{definition}

\begin{lemma}
\label{algebra-lemma-base-change-formally-unramified}
设 $R \to S$ 是形式非分歧映射，$R \to R'$ 是任意环映射。
则基变换 $S' = R' \otimes_R S$ 在 $R'$ 上形式非分歧。
\end{lemma}

\begin{proof}
给定定义 \ref{algebra-definition-formally-unramified} 中如下实线图
$$
\xymatrix{
S \ar[r] \ar@{-->}[rrd] & R' \otimes_R S \ar[r] \ar@{-->}[rd] & A/I \\
R  \ar[u] \ar[r] & R' \ar[r] \ar[u] & A \ar[u]
}
$$
依假设，较长的虚线箭头至多存在一条。
由张量积的泛性质，较短的虚线箭头也至多存在一条。
\end{proof}

\begin{lemma}
\label{algebra-lemma-characterize-formally-unramified}
设 $R \to S$ 是环映射。下列条件等价：
\begin{enumerate}
\item $R \to S$ 形式非分歧；
\item 微分模 $\Omega_{S/R}$ 为零。
\end{enumerate}
\end{lemma}

\begin{proof}
设 $J = \Ker(S \otimes_R S \to S)$ 是乘法映射的核，
并设 $A_{univ} = S \otimes_R S/J^2$。回忆
$I_{univ} = J/J^2$ 同构于 $\Omega_{S/R}$，见引理
\ref{algebra-lemma-differentials-diagonal}。此外，两个 $R$-代数映射
$\sigma_1, \sigma_2 : S \to A_{univ}$，其中
$\sigma_1(s) = s \otimes 1 \bmod J^2$ 且
$\sigma_2(s) = 1 \otimes s \bmod J^2$，相差泛导子
$\text{d} : S \to \Omega_{S/R} = I_{univ}$。

\medskip\noindent
假设 $R \to S$ 形式非分歧。于是 $\sigma_1 = \sigma_2$，
从而对所有 $s \in S$ 都有 $\text{d}(s) = 0$。
故 $\Omega_{S/R} = 0$。

\medskip\noindent
假设 $\Omega_{S/R} = 0$。设 $A, I, R \to A, S \to A/I$
构成定义 \ref{algebra-definition-formally-unramified} 中的实线图。
设 $\tau_1, \tau_2 : S \to A$ 是两个使图交换的虚线箭头。
考察由规则 $s_1 \otimes s_2 \mapsto \tau_1(s_1)\tau_2(s_2)$
定义的 $R$-代数映射 $A_{univ} \to A$。
这里省略验证该映射良定义。由于
$I_{univ} = \Omega_{S/R} = 0$，故 $A_{univ} \cong S$，
从而 $\tau_1 = \tau_2$。
\end{proof}

\begin{lemma}
\label{algebra-lemma-formally-unramified-local}
设 $R \to S$ 是环映射。下列条件等价：
\begin{enumerate}
\item $R \to S$ 形式非分歧；
\item 对 $S$ 的所有素理想 $\mathfrak q$，
$R \to S_{\mathfrak q}$ 形式非分歧；
\item 对 $S$ 的所有素理想 $\mathfrak q$，令
$\mathfrak p = R \cap \mathfrak q$，则
$R_{\mathfrak p} \to S_{\mathfrak q}$ 形式非分歧。
\end{enumerate}
\end{lemma}

\begin{proof}
引理 \ref{algebra-lemma-characterize-formally-unramified} 已表明，
(1) 等价于 $\Omega_{S/R} = 0$。类似地，由引理
\ref{algebra-lemma-differentials-localize} 可知，(2) 和 (3) 都等价于：
对所有 $\mathfrak q$ 都有
$(\Omega_{S/R})_{\mathfrak q} = 0$。
因此，该等价性由引理 \ref{algebra-lemma-characterize-zero-local} 推出。
\end{proof}

\begin{lemma}
\label{algebra-lemma-formally-unramified-localize}
设 $A \to B$ 是形式非分歧环映射。
\begin{enumerate}
\item 若 $S \subset A$ 是乘法子集，则
$S^{-1}A \to S^{-1}B$ 形式非分歧。
\item 若 $S \subset B$ 是乘法子集，则
$A \to S^{-1}B$ 形式非分歧。
\end{enumerate}
\end{lemma}

\begin{proof}
这由引理 \ref{algebra-lemma-formally-unramified-local} 推出。
% P01-E0924: the frozen proof begins with the sentence fragment “Follows from”.
（也可以由引理 \ref{algebra-lemma-characterize-formally-unramified}
结合引理 \ref{algebra-lemma-differentials-localize} 推出。）
\end{proof}

\begin{lemma}
\label{algebra-lemma-colimit-formally-unramified}
设 $R$ 是环，$I$ 是有向集，并设
$(S_i, \varphi_{ii'})$ 是 $I$ 上的 $R$-代数系。
若每个 $R \to S_i$ 都形式非分歧，则
$S = \colim_{i \in I} S_i$ 在 $R$ 上形式非分歧。
% P01-E0925: the frozen lemma conclusion lacks a terminal period.
\end{lemma}

\begin{proof}
考察定义 \ref{algebra-definition-formally-unramified} 中的图。
依假设，提升复合映射 $S_i \to S \to A/I$ 的
$R$-代数映射 $S_i \to A$ 至多存在一个。
由于 $S$ 的每个元素都属于某个映射 $S_i \to S$ 的像，
可知嵌入该图的映射 $S \to A$ 至多存在一个。
\end{proof}

\section{余法模与泛加厚}
\label{algebra-section-conormal}

\noindent
事实表明，第一无穷小邻域不仅可对概形的闭浸入定义，
而且已经可对任意形式非分歧态射定义。其基础是如下代数事实。

\begin{lemma}
\label{algebra-lemma-universal-thickening}
设 $R \to S$ 是形式非分歧环映射。存在 $R$-代数满射
$S' \to S$，其核为平方为零的理想，并具有如下泛性质：
给定任意交换图
$$
\xymatrix{
S \ar[r]_a & A/I \\
R \ar[r]^b \ar[u] & A \ar[u]
}
$$
其中 $I \subset A$ 是平方为零的理想，则存在唯一的
$R$-代数映射 $a' : S' \to A$，使得
$S' \to A \to A/I$ 等于 $S' \to S \to A/I$。
\end{lemma}

\begin{proof}
选取 $S$ 作为 $R$-代数的一族生成元 $z_i \in S$，$i \in I$。
% P01-E0926: the frozen proof reuses I first as this generator index set and later as the square-zero ideal in an arbitrary diagram.
以 $P = R[\{x_i\}_{i \in I}]$ 表示以 $x_i$（$i \in I$）
为生成元的多项式环。考察把 $x_i$ 映到 $z_i$ 的
$R$-代数映射 $P \to S$，并令 $J = \Ker(P \to S)$。
考察映射
$$
\text{d} : J/J^2 \longrightarrow \Omega_{P/R} \otimes_P S
$$
见引理 \ref{algebra-lemma-differential-seq}。
由于依假设 $\Omega_{S/R} = 0$，该映射是满射，见引理
\ref{algebra-lemma-characterize-formally-unramified}。
注意 $\Omega_{P/R}$ 是以 $\text{d}x_i$ 为基的自由模，
故 $\Omega_{P/R} \otimes_P S$ 是自由 $S$-模。
因此可选取上述满射的一个分裂，并写成
$$
J/J^2 = K \oplus \Omega_{P/R} \otimes_P S
$$
% P01-E0927: the frozen decomposition display lacks terminal punctuation before the next sentence.
令 $J^2 \subset J' \subset J$ 为 $P$ 的理想，使得
$J'/J^2$ 是上述分解中的第二个直和项。置 $S' = P/J'$。
得到短正合列
$$
0 \to J/J' \to S' \to S \to 0
$$
并且 $J/J' \cong K$ 是 $S'$ 中平方为零的理想。因此
$$
\xymatrix{
S \ar[r]_1 & S \\
R \ar[r] \ar[u] & S' \ar[u]
}
$$
是上述类型的图。事实上，我们断言它是这些图所成范畴中的初始对象。
确切地，设 $(I \subset A, a, b)$ 是任意一个这样的图。
可以选取 $R$-代数映射 $\beta : P \to A$，使图
$$
\xymatrix{
S \ar[r]_1 & S \ar[r]_a & A/I \\
R \ar[r] \ar@/_/[rr]_b \ar[u] & P \ar[u] \ar[r]^\beta & A \ar[u]
}
$$
交换。此时未必有 $\beta(J') = 0$；换言之，$\beta$ 未必通过
$S' = P/J'$ 分解。但显然 $\beta(J') \subset I$，而且由于
$\beta(J) \subset I$ 且 $I^2 = 0$，有 $\beta(J^2) = 0$。
因此，从 $(J/J' \subset S', 1, R \to S')$ 到
$(I \subset A, a, b)$ 的态射存在的“阻碍”，就是相应的
$S$-线性映射 $\overline{\beta} : J'/J^2 \to I$。
选取 $\beta$ 的自由度在于选取 $\beta(x_i)$。
若对 $\beta$ 作另一选择，记作 $\beta'$，则它由
$\beta'(x_i) = \beta(x_i) + \delta_i$ 给出，其中 $\delta_i \in I$。
在这种情形下，对 $g \in J'$ 有
$$
\beta'(g) =
\beta(g) + \sum\nolimits_i \delta_i \beta(\frac{\partial g}{\partial x_i})
$$
% P01-E0928: the frozen display lacks terminal punctuation before the following sentence.
由于按构造，映射
$\text{d}|_{J'/J^2} : J'/J^2 \to \Omega_{P/R} \otimes_P S$
由 $g \mapsto \frac{\partial g}{\partial x_i}\text{d}x_i \otimes 1$
给出且是同构，可知存在唯一的一族 $\delta_i \in I$，使得对所有
$g \in J'$ 都有 $\beta'(g) = 0$。
% P01-E0929: the frozen differential formula omits the summation over i.
（确切地，$\delta_i$ 是 $-\overline{\beta}(g)$，其中
$g \in J'/J^2$ 是在 $\text{d}|_{J'/J^2}$ 下映到
$\text{d}x_i \otimes 1$ 的唯一元素。）
% P01-E0930: the frozen parenthetical uses an unindexed g for the family defining each delta_i; g_i would make the dependence explicit.
解的唯一性蕴含引理所要求的唯一性。
\end{proof}

\noindent
在引理 \ref{algebra-lemma-universal-thickening} 的情形下，
$R$-代数映射 $S' \to S$ 在唯一同构意义下唯一。

\begin{definition}
\label{algebra-definition-universal-thickening}
设 $R \to S$ 是形式非分歧环映射。
\begin{enumerate}
\item $S$ 在 $R$ 上的\textit{泛第一阶加厚}，是引理
\ref{algebra-lemma-universal-thickening} 中的 $R$-代数满射 $S' \to S$。
\item $R \to S$ 的\textit{余法模}，是泛第一阶加厚
$S' \to S$ 的核 $I$，视为 $S$-模。
\end{enumerate}
在这种情形下，我们常以 \textit{$C_{S/R}$} 表示余法模。
\end{definition}

\begin{lemma}
\label{algebra-lemma-universal-thickening-quotient}
设 $I \subset R$ 是环的理想。$R/I$ 在 $R$ 上的泛第一阶加厚是满射
$R/I^2 \to R/I$。$R/I$ 在 $R$ 上的余法模为
$C_{(R/I)/R} = I/I^2$。
\end{lemma}

\begin{proof}
略。
\end{proof}

\begin{lemma}
\label{algebra-lemma-universal-thickening-localize}
设 $A \to B$ 是形式非分歧环映射，并设
$\varphi : B' \to B$ 是 $B$ 在 $A$ 上的泛第一阶加厚。
\begin{enumerate}
\item 设 $S \subset A$ 是乘法子集。则
$S^{-1}B' \to S^{-1}B$ 是 $S^{-1}B$ 在 $S^{-1}A$ 上的
泛第一阶加厚。特别地，
$S^{-1}C_{B/A} = C_{S^{-1}B/S^{-1}A}$。
\item 设 $S \subset B$ 是乘法子集。则
$S' = \varphi^{-1}(S)$ 是 $B'$ 中的乘法子集，并且
$(S')^{-1}B' \to S^{-1}B$ 是 $S^{-1}B$ 在 $A$ 上的
泛第一阶加厚。特别地，$S^{-1}C_{B/A} = C_{S^{-1}B/A}$。
\end{enumerate}
由引理 \ref{algebra-lemma-formally-unramified-localize}，该引理中的说法确有意义。
\end{lemma}

\begin{proof}
采用 (1) 中的记号和假设。设 $(S^{-1}B)' \to S^{-1}B$ 是
$S^{-1}B$ 在 $S^{-1}A$ 上的泛第一阶加厚。
% P01-E0931: both proof parts in the frozen source begin with the sentence fragment “With notation and assumptions as in ...”.
注意 $S^{-1}B' \to S^{-1}B$ 是 $S^{-1}A$-代数的满射，
且其核平方为零。因此按定义得到映射
$(S^{-1}B)' \to S^{-1}B'$，它与到 $S^{-1}B$ 的映射相容。
考察任意交换图
$$
\xymatrix{
B \ar[r] & S^{-1}B \ar[r] & D/I \\
A \ar[r] \ar[u] & S^{-1}A \ar[r] \ar[u] & D \ar[u]
}
$$
其中 $I \subset D$ 是平方为零的理想。由于 $B'$ 是 $B$ 在 $A$ 上的
泛第一阶加厚，得到 $A$-代数映射 $B' \to D$。
但显然，$S$ 在 $D$ 中的像由此映到 $D$ 的可逆元素，
因而得到相容的映射 $S^{-1}B' \to D$。
取 $D = (S^{-1}B)'$，可知得到映射
$S^{-1}B' \to (S^{-1}B)'$。这里省略验证该映射是上述映射的逆。

\medskip\noindent
采用 (2) 中的记号和假设。设 $(S^{-1}B)' \to S^{-1}B$
是 $S^{-1}B$ 在 $A$ 上的泛第一阶加厚。
注意 $(S')^{-1}B' \to S^{-1}B$ 是 $A$-代数的满射，
且其核平方为零。因此按定义得到映射
$(S^{-1}B)' \to (S')^{-1}B'$，它与到 $S^{-1}B$ 的映射相容。
考察任意交换图
$$
\xymatrix{
B \ar[r] & S^{-1}B \ar[r] & D/I \\
A \ar[r] \ar[u] & A \ar[r] \ar[u] & D \ar[u]
}
$$
其中 $I \subset D$ 是平方为零的理想。由于 $B'$ 是 $B$ 在 $A$ 上的
泛第一阶加厚，得到 $A$-代数映射 $B' \to D$。
但显然，$S'$ 在 $D$ 中的像由此映到 $D$ 的可逆元素，
因而得到相容的映射 $(S')^{-1}B' \to D$。
取 $D = (S^{-1}B)'$，可知得到映射
$(S')^{-1}B' \to (S^{-1}B)'$。这里省略验证该映射是上述映射的逆。
\end{proof}

\begin{lemma}
\label{algebra-lemma-differentials-universal-thickening}
设 $R \to A \to B$ 是环映射，并假设 $A \to B$ 形式非分歧。
设 $B' \to B$ 是 $B$ 在 $A$ 上的泛第一阶加厚。
则 $B'$ 在 $A$ 上形式非分歧，并且典范映射
$\Omega_{A/R} \otimes_A B \to \Omega_{B'/R} \otimes_{B'} B$
是同构。
\end{lemma}

\begin{proof}
我们将使用引理 \ref{algebra-lemma-universal-thickening} 的证明中对 $B'$ 的构造，
尽管原则上应该能从定义形式地推出这些结果。确切地，选取表现
$B = P/J$，其中 $P = A[x_i]$ 是 $A$ 上的多项式环。
然后选取元素 $f_i \in J$，使得在
$\Omega_{P/A} \otimes_P B$ 中有
$\text{d}f_i = \text{d}x_i \otimes 1$。作出这些选择后，
$B' = P/J'$，其中 $J' = (f_i) + J^2$，见引理
\ref{algebra-lemma-universal-thickening} 的证明。

\medskip\noindent
考察典范正合列
$$
J'/(J')^2 \to \Omega_{P/A} \otimes_P B' \to \Omega_{B'/A} \to 0
$$
见引理 \ref{algebra-lemma-differential-seq}。
按构造，$f_i \in J'$ 的类在模
$\Omega_{P/A} \otimes_P B'$ 中的像模掉 $J'/J^2$ 后生成该模。
由于 $J'/J^2$ 是幂零理想，可知这些元素生成整个模（由
Nakayama 引理 \ref{algebra-lemma-NAK}）。
% P01-E0932: the frozen source uses J'/J^2 as the modulus and calls it a nilpotent ideal of B'; the relevant square-zero kernel of B' -> B is J/J'.
这证明 $\Omega_{B'/A} = 0$，因而 $B'$ 在 $A$ 上形式非分歧，
见引理 \ref{algebra-lemma-characterize-formally-unramified}。

\medskip\noindent
由于 $P$ 是 $A$ 上的多项式环，有
$\Omega_{P/R} = \Omega_{A/R} \otimes_A P \oplus \bigoplus P\text{d}x_i$。
我们将使用这一分解。考察如下正合列
$$
J'/(J')^2 \to
\Omega_{P/R} \otimes_P B' \to
\Omega_{B'/R} \to 0
$$
见引理 \ref{algebra-lemma-differential-seq}。
可以将其与 $B$ 作张量积，得到正合列
$$
J'/(J')^2 \otimes_{B'} B \to
\Omega_{P/R} \otimes_P B \to
\Omega_{B'/R} \otimes_{B'} B \to 0
$$
% P01-E0933: the frozen exact-sequence display lacks terminal punctuation before the following sentence.
回忆 $J' = (f_i) + J^2$，可知第一条箭头把子模
$J^2/(J')^2$ 化为零。就给出的直和分解
$\Omega_{P/R} \otimes_P B =
\Omega_{A/R} \otimes_A B \oplus \bigoplus B\text{d}x_i $
而言，可知子模 $(f_i)/(J')^2 \otimes_{B'} B$ 同构地映到直和项
$\bigoplus B\text{d}x_i$ 上。
% P01-E0934: the frozen phrase “the direct sum decomposition ... given we see” is missing “above” and punctuation after “given”.
因此，该正合列剩下的部分就是同构
$\Omega_{A/R} \otimes_A B \to \Omega_{B'/R} \otimes_{B'} B$，
正如所需。
\end{proof}

\section{形式平展映射}
\label{algebra-section-formally-etale}

\begin{definition}
\label{algebra-definition-formally-etale}
设 $R \to S$ 是环映射。若对每个由实线箭头组成的交换图
$$
\xymatrix{
S \ar[r] \ar@{-->}[rd] & A/I \\
R \ar[r] \ar[u] & A \ar[u]
}
$$
其中 $I \subset A$ 是平方为零的理想，都存在唯一一条使图交换的虚线箭头，
则称 $S$ \textit{在 $R$ 上形式平展}。
\end{definition}

\noindent
显然，环映射形式平展，当且仅当它既形式光滑又形式非分歧。

\begin{lemma}
\label{algebra-lemma-base-change-formally-etale}
设 $R \to S$ 是形式平展映射，$R \to R'$ 是任意环映射。
则基变换 $S' = R' \otimes_R S$ 在 $R'$ 上形式平展。
\end{lemma}

\begin{proof}
结合引理 \ref{algebra-lemma-base-change-fs} 和
\ref{algebra-lemma-base-change-formally-unramified}。
\end{proof}

\begin{lemma}
\label{algebra-lemma-formally-etale-etale}
设 $R \to S$ 是有限表现环映射。下列条件等价：
\begin{enumerate}
\item $R \to S$ 形式平展；
\item $R \to S$ 平展。
\end{enumerate}
\end{lemma}

\begin{proof}
假设 $R \to S$ 形式平展。由命题
\ref{algebra-proposition-smooth-formally-smooth}，$R \to S$ 光滑。
由引理 \ref{algebra-lemma-characterize-formally-unramified}，有
$\Omega_{S/R} = 0$。故按定义，$R \to S$ 平展。

\medskip\noindent
假设 $R \to S$ 平展。由命题
\ref{algebra-proposition-smooth-formally-smooth}，$R \to S$ 形式光滑。
由引理 \ref{algebra-lemma-characterize-formally-unramified}，它形式非分歧。
因此 $R \to S$ 形式平展。
\end{proof}

\begin{lemma}
\label{algebra-lemma-colimit-formally-etale}
设 $R$ 是环，$I$ 是有向集，并设
$(S_i, \varphi_{ii'})$ 是 $I$ 上的 $R$-代数系。
若每个 $R \to S_i$ 都形式平展，则
$S = \colim_{i \in I} S_i$ 在 $R$ 上形式平展。
% P01-E0935: the frozen lemma conclusion lacks terminal punctuation.
\end{lemma}

\begin{proof}
考察定义 \ref{algebra-definition-formally-etale} 中的图。
依假设，存在唯一的一族 $R$-代数映射 $S_i \to A$
提升复合映射 $S_i \to S \to A/I$。
% P01-E0936: the frozen source says “we get unique ... maps” without a determiner or explicit per-index quantifier.
因此这些映射与转移映射 $\varphi_{ii'}$ 相容，并定义提升
$S \to A$。这证明了存在性。把映射限制到每个 $S_i$ 上，
唯一性便显然成立。
\end{proof}

\begin{lemma}
\label{algebra-lemma-localization-formally-etale}
设 $R$ 是环，$S \subset R$ 是任意乘法子集。
则环映射 $R \to S^{-1}R$ 形式平展。
\end{lemma}

\begin{proof}
设 $I \subset A$ 是平方为零的理想。这里所说的是：
给定环映射 $\varphi : R \to A$，若对所有 $f \in S$，
$\varphi(f) \mod I$ 都可逆，则对所有 $f \in S$，
$\varphi(f)$ 在 $A$ 中也可逆。这是因为，在典范映射
$A \to A/I$ 下，$A^*$ 是 $(A/I)^*$ 的逆像。
\end{proof}

\begin{lemma}
\label{algebra-lemma-formally-etale-lift-infinitesimal}
设 $R \to S$ 是环映射。设 $J \subset S$ 是理想，使得
$R \to S/J$ 为满射；令 $I \subset R$ 为其核。
若 $R \to S$ 形式平展，则对所有 $n$，
$R/I^n \to S/J^n$ 都是同构，并且
$\bigoplus I^n/I^{n + 1} \to \bigoplus J^n/J^{n + 1}$
是分次环的同构。
\end{lemma}

\begin{proof}
归纳使用提升性质，得到虚线箭头
$$
\xymatrix{
S \ar[r] \ar@{-->}[rd] & S/J = R/I \\
R \ar[r] \ar[u] & R/I^2 \ar[u]
}
\quad
\xymatrix{
S \ar[r] \ar@{-->}[rd] & R/I^2 \\
R \ar[r] \ar[u] & R/I^3 \ar[u]
}
\quad
\xymatrix{
S \ar[r] \ar@{-->}[rd] & R/I^3 \\
R \ar[r] \ar[u] & R/I^4 \ar[u]
}
$$
% P01-E0937: the frozen three-diagram display lacks terminal punctuation.
相应的映射 $S/J^n \to R/I^n$ 是同构，因为复合映射
$S/J^n \to R/I^n \to S/J^n$ 由形式平展环映射提升性质中的
唯一性可（归纳地）知为恒等映射。
\end{proof}

\begin{lemma}
\label{algebra-lemma-formally-etale-omega}
设 $R \to S \to S'$ 是环映射。令 $J$ 和 $J'$ 分别为乘法映射
$S \otimes_R S \to S$ 和 $S' \otimes_{R'} S' \to S'$ 的核。
% P01-E0938: R' is undefined in the frozen source; the second tensor product should be over R.
若 $S \to S'$ 形式平展，则对所有 $k \geq 0$，映射
$$
S' \otimes_S \left((S \otimes_R S)/J^{k + 1}\right)
\longrightarrow
(S' \otimes_R S')/(J')^{k + 1}
$$
都是同构。特别地，映射
$S' \otimes_S \Omega_{S/R} \to \Omega_{S'/R}$ 是同构。
\end{lemma}

\begin{proof}
注意 $S' \otimes_S (S \otimes_R S) = S' \otimes_R S$，
而 $J$ 在 $S' \otimes_R S$ 中生成满射
$S' \otimes_R S \to S'$ 的核 $I$。因此左边等于
$S' \otimes_R S/I^{k + 1}$。由引理
\ref{algebra-lemma-base-change-formally-etale}，映射
$S' \otimes_R S \to S' \otimes_R S'$ 形式平展。
于是由引理 \ref{algebra-lemma-formally-etale-lift-infinitesimal}，
引理陈述中的显示箭头是同构。最后一个断言由此以及如下事实推出：
由引理 \ref{algebra-lemma-differentials-diagonal}，
$\Omega_{S/R} = J/J^2$ 且 $\Omega_{S'/R} = J'/(J')^2$。
\end{proof}

\begin{lemma}
\label{algebra-lemma-formally-etale-principal-parts}
设 $R \to S \to S'$ 是环映射，且 $S \to S'$ 形式平展
（例如平展）。设 $M$ 是 $S$-模。置 $M' = S' \otimes_R M$。则
% P01-E0939: the frozen statement tensors M over R, while the lemma and its proof require M' = S' tensor_S M.
$$
S' \otimes_S P^k_{S/R}(M) = P^k_{S'/R}(M')
$$
% P01-E0940: the frozen display lacks terminal punctuation.
由此，对任意 $S$-模 $N$ 以及任意有限阶微分算子
$D : M \to N$，都存在唯一的 $D$ 到同阶或更低阶微分算子的延拓
$D' : M' \to S' \otimes_S N$。
\end{lemma}

\begin{proof}
令 $J$ 和 $J'$ 如引理 \ref{algebra-lemma-formally-etale-omega} 的陈述中所设。
于是
\begin{align*}
S' \otimes_S P^k_{S/R}(M)
& =
S' \otimes_S \left((S \otimes_R M)/J^{k + 1}(S \otimes_R M)\right) \\
& =
S' \otimes_S \left((S \otimes_R S)/J^{k + 1}\right) \otimes_S M \\
& =
\left(S' \otimes_R S'/(J')^{k + 1}\right) \otimes_S M \\
& =
\left(S' \otimes_R S'/(J')^{k + 1}\right) \otimes_{S'} M' \\
& =
S' \otimes_R M'/(J')^{k + 1}(S' \otimes_R M') \\
& =
P^k_{S'/R'}(M')
\end{align*}
% P01-E0941: the frozen final term uses undefined R' instead of R.
第一个和最后一个等式来自引理
\ref{algebra-lemma-principal-parts-diagonal}，第三个等式是引理
\ref{algebra-lemma-formally-etale-omega}。最后一个断言成立，是因为若
$D$ 对应于线性映射 $\gamma : P^k_{S/R}(M) \to N$，
则可令 $D' : M' \to S' \otimes_R N$ 对应于线性映射
$1 \otimes \gamma : S' \otimes_R P^k_{S/R}(M) \to S' \otimes_R N$。
% P01-E0942: the frozen proof tensors both N and the principal-parts module over R, contradicting the stated S-base-change targets.
\end{proof}

\begin{remark}
\label{algebra-remark-formally-etale-differential-operators}
设 $R \to S \to S'$ 是环映射，且 $S \to S'$ 形式平展
（例如平展）。设 $M_i$（$i = 1, 2, 3$）是 $S$-模，并设
$D_i : M_i \to M_{i + 1}$（$i = 1, 2$）是有限阶微分算子。
若 $D'_i : M'_i \to M'_{i + 1}$（$i = 1, 2$）是按引理
\ref{algebra-lemma-formally-etale-principal-parts} 将 $D_i$ 延拓到
$M'_i = S' \otimes_S M_i$ 所得的算子，则
$D'_2 \circ D'_1$ 是 $D_2 \circ D_1$ 的延拓。
特别地，若 $M$ 是 $S$-模，则
$M' = S' \otimes_S M$ 是 $S$-代数
$\text{Diff}_{S/R}(M, M)$ 上的模。
\end{remark}

\section{非分歧环映射}
\label{algebra-section-unramified}

\noindent
我们对非分歧环映射所采用的定义来自 \cite{Henselian}。
我们所谓的 G-非分歧环映射，就是 EGA 所称的非分歧映射。
% P01-E0943: the frozen source has the incorrect article “an G-unramified ring map”.

\begin{definition}
\label{algebra-definition-unramified}
设 $R \to S$ 是环映射。
\begin{enumerate}
\item 若 $R \to S$ 为有限型且 $\Omega_{S/R} = 0$，
则称 $R \to S$ \textit{非分歧}。
\item 若 $R \to S$ 为有限表现且 $\Omega_{S/R} = 0$，
则称 $R \to S$ \textit{G-非分歧}。
\item 给定 $S$ 的素理想 $\mathfrak q$，若存在
$g \in S$、$g \not \in \mathfrak q$，使得 $R \to S_g$ 非分歧，
则称 $S$ \textit{在 $\mathfrak q$ 处非分歧}。
\item 给定 $S$ 的素理想 $\mathfrak q$，若存在
$g \in S$、$g \not \in \mathfrak q$，使得 $R \to S_g$ G-非分歧，
则称 $S$ \textit{在 $\mathfrak q$ 处 G-非分歧}。
\end{enumerate}
\end{definition}

\noindent
当然，G-非分歧映射必非分歧。

\begin{lemma}
\label{algebra-lemma-formally-unramified-unramified}
设 $R \to S$ 是环映射。下列条件等价：
\begin{enumerate}
\item $R \to S$ 形式非分歧且为有限型；
\item $R \to S$ 非分歧。
\end{enumerate}
此外，下列条件也等价：
\begin{enumerate}
\item $R \to S$ 形式非分歧且为有限表现；
\item $R \to S$ G-非分歧。
\end{enumerate}
\end{lemma}

\begin{proof}
由引理 \ref{algebra-lemma-characterize-formally-unramified} 和定义立即得到。
\end{proof}

\begin{lemma}
\label{algebra-lemma-unramified}
非分歧与 G-非分歧环映射具有下列性质。
\begin{enumerate}
\item 非分歧环映射的基变换非分歧。
G-非分歧环映射的基变换 G-非分歧。
\item 非分歧环映射的复合非分歧。
G-非分歧环映射的复合 G-非分歧。
\item 任意主局部化 $R \to R_f$ 都既 G-非分歧又非分歧。
\item 若 $I \subset R$ 是理想，则 $R \to R/I$ 非分歧。
若 $I \subset R$ 是有限生成理想，则 $R \to R/I$ G-非分歧。
\item 平展环映射既 G-非分歧又非分歧。
\item 若 $R \to S$ 为有限型（相应地有限表现），
$\mathfrak q \subset S$ 为素理想且 $(\Omega_{S/R})_{\mathfrak q} = 0$，
则 $R \to S$ 在 $\mathfrak q$ 处非分歧（相应地 G-非分歧）。
\item 若 $R \to S$ 为有限型（相应地有限表现），
$\mathfrak q \subset S$ 为素理想且
$\Omega_{S/R} \otimes_S \kappa(\mathfrak q) = 0$，则
$R \to S$ 在 $\mathfrak q$ 处非分歧（相应地 G-非分歧）。
\item 若 $R \to S$ 为有限型（相应地有限表现），
$\mathfrak q \subset S$ 是位于 $\mathfrak p \subset R$ 之上的素理想，且
$(\Omega_{S \otimes_R \kappa(\mathfrak p)/\kappa(\mathfrak p)})_{\mathfrak q}
= 0$，则 $R \to S$ 在 $\mathfrak q$ 处非分歧（相应地 G-非分歧）。
\item 若 $R \to S$ 为有限型（相应地有限表现），
% P01-E0944: the frozen source says “resp. presentation”, omitting “finite”.
$\mathfrak q \subset S$ 是位于 $\mathfrak p \subset R$ 之上的素理想，且
$(\Omega_{S \otimes_R \kappa(\mathfrak p)/\kappa(\mathfrak p)})
\otimes_{S \otimes_R \kappa(\mathfrak p)} \kappa(\mathfrak q) = 0$，
则 $R \to S$ 在 $\mathfrak q$ 处非分歧（相应地 G-非分歧）。
\item 若 $R \to S$ 是环映射，$g_1, \ldots, g_m \in S$ 生成单位理想，
并且对 $j = 1, \ldots, m$，$R \to S_{g_j}$ 都非分歧
（相应地 G-非分歧），则 $R \to S$ 非分歧（相应地 G-非分歧）。
\item 若环映射 $R \to S$ 在 $S$ 的每个素理想处都非分歧
（相应地 G-非分歧），则 $R \to S$ 非分歧（相应地 G-非分歧）。
\item 若 $R \to S$ G-非分歧，则存在有限型
$\mathbf{Z}$-代数 $R_0$、G-非分歧环映射 $R_0 \to S_0$，
以及环映射 $R_0 \to R$，使得 $S = R \otimes_{R_0} S_0$。
\item 若 $R \to S$ 非分歧，则存在有限型
$\mathbf{Z}$-代数 $R_0$、非分歧环映射 $R_0 \to S_0$，
以及环映射 $R_0 \to R$，使得 $S$ 是
$R \otimes_{R_0} S_0$ 的商。
\end{enumerate}
\end{lemma}

\begin{proof}
依次证明各项。

\medskip\noindent
对 (1)。由引理 \ref{algebra-lemma-differentials-base-change}
和 \ref{algebra-lemma-base-change-finiteness} 得到。

\medskip\noindent
对 (2)。由引理 \ref{algebra-lemma-exact-sequence-differentials}
和 \ref{algebra-lemma-base-change-finiteness} 得到。
% P01-E0945: the frozen proof of composition cites the base-change finiteness lemma; finite-type/finite-presentation composition is lemma-compose-finite-type.

\medskip\noindent
对 (3)。直接计算 $\Omega_{R_f/R}$ 即得；计算从略。

\medskip\noindent
对 (4)。由引理 \ref{algebra-lemma-trivial-differential-surjective}，有
$\Omega_{(R/I)/R} = 0$，且环映射 $R \to R/I$ 为有限型。
若 $I$ 是有限生成理想，则 $R \to R/I$ 为有限表现。

\medskip\noindent
对 (5)。见定义 \ref{algebra-definition-etale} 后的讨论。

\medskip\noindent
对 (6)。此时 $\Omega_{S/R}$ 是有限 $S$-模（见引理
\ref{algebra-lemma-differentials-finitely-generated}），故存在
$g \in S$、$g \not \in \mathfrak q$，使得
$(\Omega_{S/R})_g = 0$。由引理 \ref{algebra-lemma-differentials-localize}，
这意味着 $\Omega_{S_g/R} = 0$，因而 $R \to S_g$ 如所需地非分歧。

\medskip\noindent
对 (7)。使用 Nakayama 引理（引理 \ref{algebra-lemma-NAK}），
可知该条件等价于 (6) 的条件。

\medskip\noindent
对 (8) 和 (9)。二者等价的方式与 (6) 和 (7) 相同。此外，由引理
\ref{algebra-lemma-differentials-base-change}，有
$\Omega_{S \otimes_R \kappa(\mathfrak p)/\kappa(\mathfrak p)} =
\Omega_{S/R} \otimes_S (S \otimes_R \kappa(\mathfrak p))$。
因此 (9) 等价于 (7)，因为两项中的 $\kappa(\mathfrak q)$-向量空间
典范同构。

\medskip\noindent
对 (10)。由引理 \ref{algebra-lemma-cover}
和 \ref{algebra-lemma-differentials-localize} 得到。

\medskip\noindent
对 (11)。由 (6)、(7) 以及 $S$ 的谱拟紧这一事实得到。
% P01-E0946: the frozen proof cites (6) and (7), which assume global finite type/presentation; the local-at-every-prime hypothesis plus quasi-compactness instead yields a finite principal cover and invokes (10).

\medskip\noindent
对 (12)。写成 $S = R[x_1, \ldots, x_n]/(g_1, \ldots, g_m)$。
由于 $\Omega_{S/R} = 0$，可在
$\Omega_{R[x_1, \ldots, x_n]/R}$ 中写成
$$
\text{d}x_i = \sum h_{ij}\text{d}g_j + \sum a_{ijk}g_j\text{d}x_k
$$
其中 $h_{ij}, a_{ijk} \in R[x_1, \ldots, x_n]$。
选取有限生成的 $\mathbf{Z}$-子代数 $R_0 \subset R$，
使其包含多项式 $g_i, h_{ij}, a_{ijk}$ 的所有系数。置
$S_0 = R_0[x_1, \ldots, x_n]/(g_1, \ldots, g_m)$。这就满足要求。

\medskip\noindent
对 (13)。写成 $S = R[x_1, \ldots, x_n]/I$。
由于 $\Omega_{S/R} = 0$，可在
$\Omega_{R[x_1, \ldots, x_n]/R}$ 中写成
$$
\text{d}x_i = \sum h_{ij}\text{d}g_{ij} + \sum g'_{ik}\text{d}x_k
$$
其中 $h_{ij} \in R[x_1, \ldots, x_n]$，且 $g_{ij}, g'_{ik} \in I$。
选取有限生成的 $\mathbf{Z}$-子代数 $R_0 \subset R$，
使其包含多项式 $g_{ij}, h_{ij}, g'_{ik}$ 的所有系数。置
$S_0 = R_0[x_1, \ldots, x_n]/(g_{ij}, g'_{ik})$。这就满足要求。
\end{proof}

\begin{lemma}
\label{algebra-lemma-diagonal-unramified}
设 $R \to S$ 是环映射。若 $R \to S$ 非分歧，则存在幂等元
$e \in S \otimes_R S$，使得 $S \otimes_R S \to S$ 同构于
$S \otimes_R S \to (S \otimes_R S)_e$。
\end{lemma}

\begin{proof}
令 $J = \Ker(S \otimes_R S \to S)$。由假设及引理
\ref{algebra-lemma-differentials-diagonal}，有 $J/J^2 = 0$。
由于 $S$ 在 $R$ 上为有限型，$J$ 有限生成；具体而言，
若 $x_i$ 在 $R$ 上生成 $S$，则生成元可以取为
$x_i \otimes 1 - 1 \otimes x_i$ 生成。
应用引理 \ref{algebra-lemma-ideal-is-squared-union-connected} 即得结论。
\end{proof}

\begin{lemma}
\label{algebra-lemma-unramified-at-prime}
设 $R \to S$ 是环映射。设 $\mathfrak q \subset S$ 是位于
$R$ 中 $\mathfrak p$ 之上的素理想。若 $S/R$ 在 $\mathfrak q$ 处非分歧，
则
\begin{enumerate}
\item $\mathfrak p S_{\mathfrak q} = \mathfrak qS_{\mathfrak q}$，
且它是局部环 $S_{\mathfrak q}$ 的极大理想；
% P01-E0947: the frozen source reads “we have pS_q = qS_q is the maximal ideal”, joining two predicates ungrammatically.
\item 域扩张 $\kappa(\mathfrak q)/\kappa(\mathfrak p)$ 有限可分。
\end{enumerate}
\end{lemma}

\begin{proof}
先用 $S_g$ 替换 $S$，其中 $g \in S$、$g \not \in \mathfrak q$，
便可假设 $R \to S$ 非分歧。由引理 \ref{algebra-lemma-unramified}，
基变换 $S \otimes_R \kappa(\mathfrak p)$ 在
$\kappa(\mathfrak p)$ 上非分歧。由引理
\ref{algebra-lemma-characterize-smooth-over-field}，它在
$\kappa(\mathfrak p)$ 上光滑，因而平展。因此
$F = S \otimes_R \kappa(\mathfrak p)$ 是
$\kappa(\mathfrak p)$ 的有限可分域扩张的有限乘积，见引理
\ref{algebra-lemma-etale-over-field}。采用注记 \ref{algebra-remark-local-ring-fibre}
的记号和结论，可得
$F_{\overline{\mathfrak q}} = S_\mathfrak q/\mathfrak pS_\mathfrak q$
等于 $\kappa(\mathfrak q)$。这就推出引理。
\end{proof}

\begin{lemma}
\label{algebra-lemma-unramified-quasi-finite}
设 $R \to S$ 是有限型环映射，$\mathfrak q$ 是 $S$ 的素理想。
若 $R \to S$ 在 $\mathfrak q$ 处非分歧，则
$R \to S$ 在 $\mathfrak q$ 处拟有限。
特别地，非分歧环映射拟有限。
\end{lemma}

\begin{proof}
非分歧环映射是有限型的，因此第二个断言显然由第一个断言推出。
为证明第一个断言，应用引理 \ref{algebra-lemma-isolated-point-fibre}
的 (2) 中的刻画，并使用引理 \ref{algebra-lemma-unramified-at-prime}。
\end{proof}

\begin{lemma}
\label{algebra-lemma-characterize-unramified}
设 $R \to S$ 是环映射，$\mathfrak q$ 是 $S$ 的素理想，
位于 $R$ 的素理想 $\mathfrak p$ 之上。若
\begin{enumerate}
\item $R \to S$ 为有限型；
\item $\mathfrak p S_{\mathfrak q}$ 是局部环
$S_{\mathfrak q}$ 的极大理想；
\item 域扩张 $\kappa(\mathfrak q)/\kappa(\mathfrak p)$ 有限可分，
\end{enumerate}
则 $R \to S$ 在 $\mathfrak q$ 处非分歧。
\end{lemma}

\begin{proof}
由引理 \ref{algebra-lemma-unramified} (8)，只须证明
$\Omega_{S \otimes_R \kappa(\mathfrak p) / \kappa(\mathfrak p)}$
在 $\mathfrak q$ 处局部化后为零。因此可以用
$S \otimes_R \kappa(\mathfrak p)$ 替换 $S$，并用
$\kappa(\mathfrak p)$ 替换 $R$。换言之，可以假设 $R = k$ 是域，
而 $S$ 是有限型 $k$-代数。此时，假设推出
$S_{\mathfrak q} \cong \kappa(\mathfrak q)$。因此
$(\Omega_{S/k})_{\mathfrak q} = \Omega_{S_\mathfrak q/k} =
\Omega_{\kappa(\mathfrak q)/k}$ 如所需地为零（第一个等号是引理
\ref{algebra-lemma-differentials-localize}）。
\end{proof}

\begin{lemma}
\label{algebra-lemma-etale-flat-unramified-finite-presentation}
设 $R \to S$ 是环映射。下列条件等价：
\begin{enumerate}
\item $R \to S$ 平展；
\item $R \to S$ 平坦且 G-非分歧；
\item $R \to S$ 平坦、非分歧且为有限表现。
\end{enumerate}
\end{lemma}

\begin{proof}
按定义，(2) 和 (3) 等价。由平展环映射为有限表现、引理
\ref{algebra-lemma-etale}（平展映射的平坦性）以及引理
\ref{algebra-lemma-unramified}（平展映射非分歧），可得
(1) $\Rightarrow$ (3)。反过来，引理 \ref{algebra-lemma-characterize-etale}
中对平展环映射的刻画，以及引理 \ref{algebra-lemma-unramified-at-prime}
中非分歧环映射的结构，说明 (3) 推出 (1)。
% P01-E0948: the frozen compound subject “the characterization ... and the structure ... shows” has singular agreement.
（这里使用了如下事实：若 $R \to S$ 在每个素理想
$\mathfrak q \subset S$ 处都平展，则 $R \to S$ 平展；见引理
\ref{algebra-lemma-etale}。）
\end{proof}

\begin{lemma}
\label{algebra-lemma-characterize-etale-over-polynomial-ring}
设 $k$ 是域，并设
$$
\varphi : k[x_1, \ldots, x_n] \to A, \quad x_i \longmapsto a_i
$$
是有限型环映射。则 $\varphi$ 平展，当且仅当下列两个条件成立：
(a) $A$ 在各极大理想处的局部环维数为 $n$；
(b) 元素 $\text{d}(a_1), \ldots, \text{d}(a_n)$ 作为 $A$-模生成
$\Omega_{A/k}$。
\end{lemma}

\begin{proof}
假设 (a) 和 (b) 成立。条件 (b) 推出
$\Omega_{A/k[x_1, \ldots, x_n]} = 0$，因而 $\varphi$ 非分歧。
于是由引理 \ref{algebra-lemma-etale-flat-unramified-finite-presentation}，
只须证明 $\varphi$ 平坦。设 $\mathfrak m \subset A$ 是极大理想。
置 $X = \Spec(A)$，并以 $x \in X$ 表示与 $\mathfrak m$ 对应的闭点。
由引理 \ref{algebra-lemma-dimension-closed-point-finite-type-field}，
$\dim(A_\mathfrak m)$ 就是 $\dim_x X$。因此由引理
\ref{algebra-lemma-characterize-smooth-over-field} 可知，若 (a) 和 (b) 成立，
则对每个极大理想 $\mathfrak m$，$A_\mathfrak m$ 都是正则局部环。
于是由引理 \ref{algebra-lemma-CM-over-regular-flat}，
$k[x_1, \ldots, x_n]_{\varphi^{-1}(\mathfrak m)} \to A_\mathfrak m$
平坦（并使用了正则局部环是 CM 环这一事实，见引理
\ref{algebra-lemma-regular-ring-CM}）。故由引理
\ref{algebra-lemma-flat-localization}，$\varphi$ 平坦。

\medskip\noindent
假设 $\varphi$ 平展。则 $\Omega_{A/k[x_1, \ldots, x_n]} = 0$，
故 (b) 成立。另一方面，平展环映射平坦（引理
\ref{algebra-lemma-etale}）且拟有限（引理 \ref{algebra-lemma-etale-quasi-finite}）。
因此，对 $A$ 的每个极大理想 $\mathfrak m$，可以将引理
\ref{algebra-lemma-dimension-base-fibre-equals-total} 应用于
$k[x_1, \ldots, x_n]_{\varphi^{-1}(\mathfrak m)} \to A_\mathfrak m$，
得到 $\dim(A_\mathfrak m) = n$，故 (a) 成立。
% P01-E0949: the frozen source has the typo “we my apply” instead of “we may apply”.
\end{proof}

\section{非分歧环映射的局部结构}
\label{algebra-section-local-structure-unramified}

\noindent
非分歧态射局部地（在适当意义下）是闭浸入与平展态射的复合。
本节讨论这一事实的代数基础。

\begin{proposition}
\label{algebra-proposition-unramified-locally-standard}
设 $R \to S$ 是环映射，$\mathfrak q \subset S$ 是素理想。
若 $R \to S$ 在 $\mathfrak q$ 处非分歧，则存在
\begin{enumerate}
\item $g \in S$、$g \not \in \mathfrak q$；
\item 标准平展环映射 $R \to S'$；
\item 满 $R$-代数映射 $S' \to S_g$。
\end{enumerate}
\end{proposition}

\begin{proof}
这个证明与命题 \ref{algebra-proposition-etale-locally-standard} 的证明“相同”。
证明稍显迂回，也许可以缩短。

\medskip\noindent
步骤 1。由定义 \ref{algebra-definition-unramified}，存在
$g \in S$、$g \not \in \mathfrak q$，使得 $R \to S_g$ 非分歧。
因此可以假设 $S$ 在 $R$ 上非分歧。

\medskip\noindent
步骤 2。由引理 \ref{algebra-lemma-unramified}，存在非分歧环映射
$R_0 \to S_0$，其中 $R_0$ 是有限型 $\mathbf{Z}$-代数；还存在环映射
$R_0 \to R$，使得 $S$ 是 $R \otimes_{R_0} S_0$ 的商。
记 $\mathfrak q_0$ 为对应于 $\mathfrak q$ 的 $S_0$ 的素理想。
若对 $(R_0 \to S_0, \mathfrak q_0)$ 证明结论，则经基变换即可对
$(R \to S, \mathfrak q)$ 得到结论。因此可以假设 $R$ 是 Noether 环。

\medskip\noindent
步骤 3。由引理 \ref{algebra-lemma-unramified-quasi-finite}，
$R \to S$ 拟有限。由引理 \ref{algebra-lemma-quasi-finite-open-integral-closure}，
存在有限环映射 $R \to S'$、$R$-代数映射 $S' \to S$，以及元素
$g' \in S'$，其中 $g' \not \in \mathfrak q$，使得 $S' \to S$ 诱导同构
$S'_{g'} \cong S_{g'}$。
% P01-E0950: the frozen source compares g' in S' directly with q in S and repeats “such that”; it should refer to the image of g' (or the contracted prime) once.
（注意，$S'$ 未必在 $R$ 上非分歧。）
因此可以假设：(a) $R$ 是 Noether 环；(b) $R \to S$ 有限；
(c) $R \to S$ 在 $\mathfrak q$ 处非分歧
（但不再必然在所有素理想处非分歧）。

\medskip\noindent
步骤 4。设 $\mathfrak p \subset R$ 是与 $\mathfrak q$ 对应的素理想。
考察纤维环 $S \otimes_R \kappa(\mathfrak p)$。它是
$\kappa(\mathfrak p)$ 上的有限代数，因而是 Artin 环
（见引理 \ref{algebra-lemma-finite-dimensional-algebra}），所以是局部环的有限乘积
$$
S \otimes_R \kappa(\mathfrak p) = \prod\nolimits_{i = 1}^n A_i
$$
见命题 \ref{algebra-proposition-dimension-zero-ring}。其中一个因子，例如 $A_1$，
是局部环 $S_{\mathfrak q}/\mathfrak pS_{\mathfrak q}$，
并且它同构于 $\kappa(\mathfrak q)$，见引理
\ref{algebra-lemma-unramified-at-prime}。其余因子对应于 $S$ 中位于
$\mathfrak p$ 之上的其他素理想，记为
$\mathfrak q_2, \ldots, \mathfrak q_n$。

\medskip\noindent
步骤 5。可以选取非零元素 $\alpha \in \kappa(\mathfrak q)$，
使其生成有限可分域扩张 $\kappa(\mathfrak q)/\kappa(\mathfrak p)$
（所以即使该域扩张平凡，也不允许 $\alpha = 0$）。注意，对任意
$\lambda \in \kappa(\mathfrak p)^*$，元素 $\lambda \alpha$ 也在
$\kappa(\mathfrak p)$ 上生成 $\kappa(\mathfrak q)$。考察元素
$$
\overline{t} =
(\alpha, 0, \ldots, 0) \in
\prod\nolimits_{i = 1}^n A_i =
S \otimes_R \kappa(\mathfrak p).
$$
必要时如上用 $\lambda \alpha$ 替换 $\alpha$，可以假设
$\overline{t}$ 是某个 $t \in S$ 的像。令 $I \subset R[x]$ 为
把 $x$ 映到 $t$ 的 $R$-代数映射 $R[x] \to S$ 的核。
置 $S' = R[x]/I$，于是 $S' \subset S$。有交换图
$$
\xymatrix{
R[x] \ar[r] & S' \ar[r] & S \\
R \ar[u] \ar[ru] \ar[rru] & &
}
$$
按构造，$S$ 的素理想 $\mathfrak q_j$（$j \geq 2$）都位于
$R[x]$ 的素理想 $(\mathfrak p, x)$ 之上，而 $\mathfrak q$ 位于
$R[x]$ 的另一个素理想之上，因为 $\alpha \not = 0$。

\medskip\noindent
步骤 6。以 $\mathfrak q' \subset S'$ 表示对应于 $\mathfrak q$ 的
$S'$ 的素理想。由上可知，$\mathfrak q$ 是 $S$ 中位于
$\mathfrak q'$ 之上的唯一素理想。因此由引理
\ref{algebra-lemma-unique-prime-over-localize-below}，有
$S_{\mathfrak q} = S_{\mathfrak q'}$（由于 $R \to S$ 有限，
$S' \to S$ 也有限，故由引理 \ref{algebra-lemma-integral-going-up}，
$S' \to S$ 满足上升性质）。由此，
$S'_{\mathfrak q'} \to S_{\mathfrak q}$ 是有限且单射的，
因为它是有限单射环映射 $S' \to S$ 的局部化。考察局部环映射
$$
R_{\mathfrak p} \to S'_{\mathfrak q'} \to S_{\mathfrak q}
$$
第二个映射有限且单射。由引理 \ref{algebra-lemma-unramified-at-prime}，有
$S_{\mathfrak q}/\mathfrak pS_{\mathfrak q} = \kappa(\mathfrak q)$。
因此更有
$S_{\mathfrak q}/\mathfrak q'S_{\mathfrak q} = \kappa(\mathfrak q)$。
由于
$$
\kappa(\mathfrak p) \subset \kappa(\mathfrak q') \subset \kappa(\mathfrak q)
$$
且 $\alpha$ 属于 $\kappa(\mathfrak q')$ 在 $\kappa(\mathfrak q)$
中的像，可得 $\kappa(\mathfrak q') = \kappa(\mathfrak q)$。
于是把 Nakayama 引理 \ref{algebra-lemma-NAK} 应用于
$S'_{\mathfrak q'}$-模映射 $S'_{\mathfrak q'} \to S_{\mathfrak q}$，
可知映射 $S'_{\mathfrak q'} \to S_{\mathfrak q}$ 是满射。换言之，
$S'_{\mathfrak q'} \cong S_{\mathfrak q}$。

\medskip\noindent
步骤 7。由引理 \ref{algebra-lemma-isomorphic-local-rings}，存在
$g \in S$、$g \not \in \mathfrak q$ 和
$g' \in S'$、$g' \not \in \mathfrak q'$，使得 $S'_{g'} \cong S_g$。
由于 $R$ 是 Noether 环，$S'$ 是有限 $R$-模 $S$ 的 $R$-子模，
故它在 $R$ 上有限。因此，用 $S'$ 替换 $S$ 后，可以假设：
(a) $R$ 是 Noether 环；(b) $S$ 在 $R$ 上有限；
% P01-E0951: the frozen source omits “is” in “(b) S finite over R”.
(c) $S$ 在 $R$ 上于 $\mathfrak q$ 处非分歧；
(d) $S = R[x]/I$。

\medskip\noindent
步骤 8。考察环
$S \otimes_R \kappa(\mathfrak p) = \kappa(\mathfrak p)[x]/\overline{I}$，
其中 $\overline{I} = I \cdot \kappa(\mathfrak p)[x]$ 是 $I$ 在
$\kappa(\mathfrak p)[x]$ 中生成的理想。由于
$\kappa(\mathfrak p)[x]$ 是主理想整环，存在首一的
$\overline{h} \in \kappa(\mathfrak p)$，使得
$\overline{I} = (\overline{h})$。
% P01-E0952: the frozen source places the monic polynomial h-bar in kappa(p) instead of kappa(p)[x].
对某个 $\lambda \in \kappa(\mathfrak p)$，以
$\lambda \cdot \overline{h}$ 替换 $\overline{h}$ 后，可以假设
$\overline{h}$ 是某个 $h \in R[x]$ 的像。
% P01-E0953: the scaling factor must be nonzero to preserve the principal generator, but the frozen source omits the star on kappa(p).
（问题在于我们不知道能否选取首一的 $h$。）此外，与步骤 4 相同，
已知 $S \otimes_R \kappa(\mathfrak p) = A_1 \times \ldots \times A_n$，
其中 $A_1 = \kappa(\mathfrak q)$ 是 $\kappa(\mathfrak p)$ 的有限可分扩张，
而 $A_2, \ldots, A_n$ 是局部环。这推出
$$
\overline{h} = \overline{h}_1 \overline{h}_2^{e_2} \ldots \overline{h}_n^{e_n}
$$
% P01-E0954: after the frozen source replaces the monic generator by an arbitrary nonzero scalar multiple, monicity is no longer assured; this factorization into monic factors needs the resulting leading unit as a coefficient.
其中 $\overline{h}_i \in \kappa(\mathfrak p)[x]$ 是两两互素的不可约
首一多项式，且 $e_2, \ldots, e_n \geq 1$。编号使得
$A_i = \kappa(\mathfrak p)[x]/(\overline{h}_i^{e_i})$ 是
$\kappa(\mathfrak p)[x]$-代数。
% P01-E0955: the frozen formula uses e_i also for i = 1, although only e_2,...,e_n are defined; it should set e_1 = 1 or state the i = 1 case separately.
注意，$\overline{h}_1$ 是 $\alpha \in \kappa(\mathfrak q)$ 的极小多项式，
因而是可分多项式（其导数与自身互素）。

\medskip\noindent
步骤 9。取首一元素 $m \in I$；这种元素存在，因为环扩张
$R \to R[x]/I$ 有限，因而整。以 $\overline{m}$ 表示它在
$\kappa(\mathfrak p)[x]$ 中的像。可以分解为
$$
\overline{m} = \overline{k}
\overline{h}_1^{d_1} \overline{h}_2^{d_2} \ldots \overline{h}_n^{d_n}
$$
其中 $d_1 \geq 1$、$d_j \geq e_j$（$j = 2, \ldots, n$），
且 $\overline{k} \in \kappa(\mathfrak p)[x]$ 与所有
$\overline{h}_i$ 互素。置 $f = m^l + h$，其中
$l \deg(m) > \deg(h)$ 且 $l \geq 2$。于是 $f$ 是 $R$ 上的首一多项式。
此外，$f$ 在 $\kappa(\mathfrak p)[x]$ 中的像 $\overline{f}$ 分解为
$$
\overline{f} =
\overline{h}_1 \overline{h}_2^{e_2} \ldots \overline{h}_n^{e_n}
+
\overline{k}^l \overline{h}_1^{ld_1} \overline{h}_2^{ld_2}
\ldots \overline{h}_n^{ld_n}
=
\overline{h}_1(\overline{h}_2^{e_2} \ldots \overline{h}_n^{e_n}
+
\overline{k}^l
\overline{h}_1^{ld_1 - 1} \overline{h}_2^{ld_2} \ldots \overline{h}_n^{ld_n})
= \overline{h}_1 \overline{w}
$$
其中 $\overline{w}$ 是与 $\overline{h}_1$ 互素的多项式。
置 $g = f'$（对 $x$ 的导数）。

\medskip\noindent
步骤 10。环映射 $R[x] \to S = R[x]/I$ 具有下列性质：
(1) 它把 $f$ 映到零；
(2) 它把 $g$ 映到 $S \setminus \mathfrak q$ 中的元素。
第一个断言显然，因为 $f$ 是 $I$ 的元素。
% P01-E0956: the frozen source only says h lies in R[x], but Step 10 needs h in I to conclude f = m^l + h lies in I.
为证明第二个断言，只须证明 $g$ 在
$\kappa(\mathfrak q) = \kappa(\mathfrak p)[x]/(\overline{h}_1)$
中的像不为零。$g$ 在 $\kappa(\mathfrak p)[x]$ 中的像是
$\overline{f}$ 的导数。因此 (2) 显然成立，因为
$$
\overline{g} =
\frac{\text{d}\overline{f}}{\text{d}x} =
\overline{w}\frac{\text{d}\overline{h}_1}{\text{d}x} +
\overline{h}_1\frac{\text{d}\overline{w}}{\text{d}x},
$$
$\overline{w}$ 与 $\overline{h}_1$ 互素，且
$\overline{h}_1$ 可分。

\medskip\noindent
步骤 11。由此可知，$\varphi : R[x]/(f) \to S$ 是满环映射，
$R[x]_g/(f)$ 在 $R$ 上平展（因为它标准平展；见引理
\ref{algebra-lemma-standard-etale}），并且 $\varphi(g) \not \in \mathfrak q$。
因此映射 $(R[x]/(f))_g \to S_{\varphi(g)}$ 就是所需的满射。
\end{proof}



\begin{lemma}
\label{algebra-lemma-etale-makes-unramified-closed-at-prime}
设 $R \to S$ 是环映射。设 $\mathfrak q$ 是 $S$ 的素理想，
位于 $\mathfrak p \subset R$ 之上。假设 $R \to S$ 为有限型，
且在 $\mathfrak q$ 处非分歧。则存在
\begin{enumerate}
\item 平展环映射 $R \to R'$；
\item 位于 $\mathfrak p$ 之上的素理想 $\mathfrak p' \subset R'$；
\item 乘积分解
$$
R' \otimes_R S = A \times B
$$
\end{enumerate}
满足下列性质：
\begin{enumerate}
\item $R' \to A$ 是满射；
\item $\mathfrak p'A$ 是 $A$ 的素理想，位于 $\mathfrak p'$ 之上，
并且位于 $\mathfrak q$ 之上。
\end{enumerate}
\end{lemma}

\begin{proof}
可以用沿平展环映射 $R \to R'$ 的任意基变换
$(R' \to R'\otimes_R S, \mathfrak p', \mathfrak q')$ 替换
$(R \to S, \mathfrak p, \mathfrak q)$；其中 $\mathfrak p'$ 是位于
$\mathfrak p$ 之上的素理想，并选取同时位于 $\mathfrak q$ 和
$\mathfrak p'$ 之上的 $\mathfrak q'$。还要注意，给定
$R \to R'$ 和 $\mathfrak p'$ 后，总能找到合适的 $\mathfrak q'$。

\medskip\noindent
$R \to S$ 为有限型这一假设意味着可以应用引理
\ref{algebra-lemma-etale-makes-quasi-finite-finite-variant}。因此可以假设
$S = A_1 \times \ldots \times A_n \times B$，每个 $R \to A_i$ 都有限，
并恰有一个位于 $\mathfrak p$ 之上的素理想 $\mathfrak r_i$，使得
$\kappa(\mathfrak p) \subset \kappa(\mathfrak r_i)$ 纯不可分；而
$R \to B$ 在位于 $\mathfrak p$ 之上的任一素理想处都不拟有限。
于是 $\mathfrak q = \mathfrak r_i$ 对某个 $i$ 成立，因为非分歧态射拟有限
（见引理 \ref{algebra-lemma-unramified-quasi-finite}）。不妨设
$\mathfrak q = \mathfrak r_1$。由引理 \ref{algebra-lemma-unramified-at-prime}，
$\kappa(\mathfrak r_1)/\kappa(\mathfrak p)$ 可分，因而是平凡域扩张，
且 $\mathfrak p(A_1)_{\mathfrak r_1}$ 是极大理想。此外，由引理
\ref{algebra-lemma-unique-prime-over-localize-below}（它适用于
$R \to A_1$，因为由引理 \ref{algebra-lemma-integral-going-up}，
有限环映射满足上升性质），有
$(A_1)_{\mathfrak r_1} = (A_1)_{\mathfrak p}$。
由 Nakayama 引理 \ref{algebra-lemma-NAK}，局部环映射
$R_{\mathfrak p} \to (A_1)_{\mathfrak p} = (A_1)_{\mathfrak r_1}$
是满射。由于 $A_1$ 在 $R$ 上有限，存在
$f \in R$、$f \not \in \mathfrak p$，使得
$R_f \to (A_1)_f$ 是满射。用 $R_f$ 替换 $R$ 后即得结论。
\end{proof}

\begin{lemma}
\label{algebra-lemma-etale-makes-unramified-closed}
\begin{slogan}
对于非分歧环映射，可以通过取平展邻域分离一个纤维中的点。
\end{slogan}
设 $R \to S$ 是环映射，$\mathfrak p$ 是 $R$ 的素理想。
若 $R \to S$ 非分歧，则存在
\begin{enumerate}
\item 平展环映射 $R \to R'$；
\item 位于 $\mathfrak p$ 之上的素理想 $\mathfrak p' \subset R'$；
\item 乘积分解
$$
R' \otimes_R S = A_1 \times \ldots \times A_n \times B
$$
\end{enumerate}
满足下列性质：
\begin{enumerate}
\item $R' \to A_i$ 是满射；
\item $\mathfrak p'A_i$ 是 $A_i$ 的素理想，位于 $\mathfrak p'$ 之上；
\item $B$ 中没有位于 $\mathfrak p'$ 之上的素理想。
\end{enumerate}
\end{lemma}

\begin{proof}
可以应用引理 \ref{algebra-lemma-etale-makes-quasi-finite-finite-variant}。
因此，经一次平展基变换后，可以假设
$S = A_1 \times \ldots \times A_n \times B$，每个 $R \to A_i$ 都有限，
并恰有一个位于 $\mathfrak p$ 之上的素理想 $\mathfrak r_i$，使得
$\kappa(\mathfrak p) \subset \kappa(\mathfrak r_i)$ 纯不可分；而
$R \to B$ 在位于 $\mathfrak p$ 之上的任一素理想处都不拟有限。
由于 $R \to S$ 拟有限（见引理 \ref{algebra-lemma-unramified-quasi-finite}），
可知 $B$ 中没有位于 $\mathfrak p$ 之上的素理想。
由引理 \ref{algebra-lemma-unramified-at-prime}，
$\kappa(\mathfrak r_i)/\kappa(\mathfrak p)$ 可分，因而是平凡域扩张，
且 $\mathfrak p(A_i)_{\mathfrak r_i}$ 是极大理想。此外，由引理
\ref{algebra-lemma-unique-prime-over-localize-below}（它适用于
$R \to A_i$，因为由引理 \ref{algebra-lemma-integral-going-up}，
有限环映射满足上升性质），有
$(A_i)_{\mathfrak r_i} = (A_i)_{\mathfrak p}$。
由 Nakayama 引理 \ref{algebra-lemma-NAK}，局部环映射
$R_{\mathfrak p} \to (A_i)_{\mathfrak p} = (A_i)_{\mathfrak r_i}$
是满射。由于 $A_i$ 在 $R$ 上有限，存在
$f \in R$、$f \not \in \mathfrak p$，使得
$R_f \to (A_i)_f$ 是满射。用 $R_f$ 替换 $R$ 后即得结论。
\end{proof}

\section{Hensel 局部环}
\label{algebra-section-henselian}

\noindent
本节简要讨论 Hensel 局部环的概念。设
$(R, \mathfrak m, \kappa)$ 是局部环。对 $a \in R$，以
$\overline{a}$ 表示 $a$ 在 $\kappa$ 中的像。对多项式
$f \in R[T]$，常以 $\overline{f}$ 表示 $f$ 在 $\kappa[T]$ 中的像。
给定多项式 $f \in R[T]$，以 $f'$ 表示 $f$ 关于 $T$ 的导数。
注意 $\overline{f}' = \overline{f'}$。

\begin{definition}
\label{algebra-definition-henselian}
设 $(R, \mathfrak m, \kappa)$ 是局部环。
\begin{enumerate}
\item 若对每个首一多项式 $f \in R[T]$ 以及
$\overline{f}$ 的每个根 $a_0 \in \kappa$，只要
$\overline{f'}(a_0) \not = 0$，就存在 $a \in R$，使得
$f(a) = 0$ 且 $a_0 = \overline{a}$，则称 $R$ \textit{Hensel}。
\item 若 $R$ Hensel，且其剩余域可分代数闭，则称 $R$
\textit{严格 Hensel}。
\end{enumerate}
\end{definition}

\noindent
注意，条件 $\overline{f'}(a_0) \not = 0$ 等价于 $a_0$ 是多项式
$\overline{f}$ 的单根。事实上，它还意味着：若提升
$a \in R$ 存在，则该提升唯一。

\begin{lemma}
\label{algebra-lemma-uniqueness}
设 $(R, \mathfrak m, \kappa)$ 是局部环。设 $f \in R[T]$。
设 $a, b \in R$ 满足 $f(a) = f(b) = 0$、
$a = b \bmod \mathfrak m$，并且 $f'(a) \not \in \mathfrak m$。
则 $a = b$。
\end{lemma}

\begin{proof}
在 $R[x, y]$ 中写成
$f(x + y) - f(x) = f'(x)y + g(x, y) y^2$（展开 $f(x + y)$ 即可看出；
细节从略）。于是对某个 $c \in R$，有
$0 = f(b) - f(a) = f(a + (b - a)) - f(a) =
f'(a)(b - a) + c (b - a)^2$。由假设，$f'(a)$ 是 $R$ 中的单位。
故 $(b - a)(1 + f'(a)^{-1}c(b - a)) = 0$。
又由假设，$b - a \in \mathfrak m$，所以
$1 + f'(a)^{-1}c(b - a)$ 是 $R$ 中的单位。因此 $b - a = 0$ 于 $R$。
\end{proof}

\noindent
下面给出 Hensel 局部环的刻画。

\begin{lemma}
\label{algebra-lemma-characterize-henselian}
\begin{slogan}
Hensel 局部环的刻画
\end{slogan}
设 $(R, \mathfrak m, \kappa)$ 是局部环。下列条件等价：
\begin{enumerate}
\item $R$ Hensel；
\item 对每个 $f \in R[T]$ 以及 $\overline{f}$ 的每个根
$a_0 \in \kappa$，只要 $\overline{f'}(a_0) \not = 0$，
就存在 $a \in R$，使得 $f(a) = 0$ 且 $a_0 = \overline{a}$；
\item 对任意首一 $f \in R[T]$ 以及任意分解
$\overline{f} = g_0 h_0$，若 $\gcd(g_0, h_0) = 1$，则存在
$R[T]$ 中的分解 $f = gh$，使得
$g_0 = \overline{g}$ 且 $h_0 = \overline{h}$；
\item 对任意首一 $f \in R[T]$ 以及任意分解
$\overline{f} = g_0 h_0$，若 $\gcd(g_0, h_0) = 1$，则存在
$R[T]$ 中的分解 $f = gh$，使得
$g_0 = \overline{g}$ 且 $h_0 = \overline{h}$，并且
$\deg_T(g) = \deg_T(g_0)$；
\item 对任意 $f \in R[T]$ 以及任意分解
$\overline{f} = g_0 h_0$，若 $\gcd(g_0, h_0) = 1$，则存在
$R[T]$ 中的分解 $f = gh$，使得
$g_0 = \overline{g}$ 且 $h_0 = \overline{h}$；
\item 对任意 $f \in R[T]$ 以及任意分解
$\overline{f} = g_0 h_0$，若 $\gcd(g_0, h_0) = 1$，则存在
$R[T]$ 中的分解 $f = gh$，使得
$g_0 = \overline{g}$ 且 $h_0 = \overline{h}$，并且
$\deg_T(g) = \deg_T(g_0)$；
\item 对任意平展环映射 $R \to S$ 以及 $S$ 中位于
$\mathfrak m$ 之上的素理想 $\mathfrak q$，若
$\kappa = \kappa(\mathfrak q)$，则存在 $R \to S$ 的收缩
$\tau : S \to R$；
\item 对任意平展环映射 $R \to S$ 以及 $S$ 中位于
$\mathfrak m$ 之上的素理想 $\mathfrak q$，若
$\kappa = \kappa(\mathfrak q)$，则存在唯一的 $R \to S$ 的收缩
$\tau : S \to R$，满足 $\mathfrak q = \tau^{-1}(\mathfrak m)$；
\item 任意有限 $R$-代数都是局部环的乘积；
\item 任意有限 $R$-代数都是局部环的有限乘积；
\item 任意有限型 $R$-代数 $S$ 都可以写成 $A \times B$，
其中 $R \to A$ 有限，而 $R \to B$ 在位于
$\mathfrak m$ 之上的任一素理想处都不拟有限；
\item 任意有限型 $R$-代数 $S$ 都可以写成 $A \times B$，
其中 $R \to A$ 有限，并且
$\Spec(B \otimes_R \kappa)$ 的每个不可约分支维数都 $\geq 1$；
\item 任意拟有限 $R$-代数 $S$ 都可以写成 $S = A \times B$，
其中 $R \to A$ 有限，且 $B \otimes_R \kappa = 0$。
\end{enumerate}
\end{lemma}

\begin{proof}
先列出较容易的蕴含：
\begin{enumerate}
\item 2$\Rightarrow$1，因为 (2) 考察所有多项式，
而 (1) 只考察首一多项式；
\item 5$\Rightarrow$3，因为 (5) 考察所有多项式，
而 (3) 只考察首一多项式；
\item 6$\Rightarrow$4，因为 (6) 考察所有多项式，
而 (4) 只考察首一多项式；
\item 4$\Rightarrow$3 显然；
\item 6$\Rightarrow$5 显然；
\item 8$\Rightarrow$7 显然；
\item 10$\Rightarrow$9 显然；
\item 11$\Leftrightarrow$12，由在素理想处拟有限的定义；
\item 11$\Rightarrow$13，由拟有限的定义。
\end{enumerate}

\noindent
证明 1$\Rightarrow$8。假设 (1) 成立。设 $R \to S$ 平展，
并设 $\mathfrak q \subset S$ 是满足
$\kappa(\mathfrak q) \cong \kappa$ 的素理想。由命题
\ref{algebra-proposition-etale-locally-standard}，可以找到
$g \in S$、$g \not \in \mathfrak q$，使得 $R \to S_g$ 标准平展。
用 $S_g$ 替换 $S$ 后，可以假设
$S = R[t]_g/(f)$ 标准平展（细节从略）。由于素理想
$\mathfrak q$ 的剩余域是 $\kappa$，它对应于 $\overline{f}$ 的根
$a_0$，而该根不是 $\overline{g}$ 的根。由标准平展代数的定义，
这还意味着 $\overline{f'}(a_0) \not = 0$。此外，再次由标准平展代数
的定义，$f$ 首一；所以可以利用 $R$ Hensel，得到
$a \in R$，满足 $a_0 = \overline{a}$ 且 $f(a) = 0$。
这意味着 $g(a)$ 是 $R$ 的单位，并可按规则 $t \mapsto a$ 得到所需映射
$\tau : S = R[t]_g/(f) \to R$。按构造，
$\tau^{-1}(\mathfrak m) = \mathfrak q$。由引理
\ref{algebra-lemma-uniqueness}，映射 $\tau$ 唯一。这证明了 (8)。

\medskip\noindent
证明 7$\Rightarrow$8。（这一部分其实不重要，可以跳过。）
假设 (7) 成立，并假设 $R \to S$ 平展。设
$\mathfrak q_1, \ldots, \mathfrak q_r$ 是 $S$ 中位于
$\mathfrak m$ 之上的其他素理想。可以找到
$g \in S$、$g \not \in \mathfrak q$，且对
$i = 1, \ldots, r$ 有 $g \in \mathfrak q_i$。事实上，可以论证
$\bigcap_{i=1}^{r} \mathfrak{q}_{i} \not\subset \mathfrak{q}$；
否则对某个 $i$ 有 $\mathfrak{q}_{i} \subset \mathfrak{q}$，
但这是不可能的，因为平展态射的纤维是离散的（例如使用引理
\ref{algebra-lemma-etale-over-field}）。将 (7) 应用于平展环映射
$R \to S_g$ 和素理想 $\mathfrak qS_g$，得到收缩
$\tau_g : S_g \to R$，使得复合映射 $\tau : S \to S_g \to R$
具有性质 $\tau^{-1}(\mathfrak m) = \mathfrak q$。细节从略。

\medskip\noindent
证明 8$\Rightarrow$11。假设 (8) 成立，并设 $R \to S$ 是有限型环映射。
应用引理 \ref{algebra-lemma-etale-makes-quasi-finite-finite}。得到平展环映射
$R \to R'$ 以及位于 $\mathfrak m$ 之上的素理想
$\mathfrak m' \subset R'$，满足 $\kappa = \kappa(\mathfrak m')$，
并且 $R' \otimes_R S = A' \times B'$，其中 $A'$ 在 $R'$ 上有限，
而 $B'$ 在位于 $\mathfrak m'$ 之上的任一素理想处都不在 $R'$ 上拟有限。
应用 (8)，得到收缩 $\tau : R' \to R$，满足
$\mathfrak m = \tau^{-1}(\mathfrak m')$。
% P01-E0957: the frozen inverse-image equality is ill-typed and reversed; for tau : R' -> R it should be m' = tau^{-1}(m).
然后使用
$$
S = (S \otimes_R R') \otimes_{R', \tau} R
= (A' \times B') \otimes_{R', \tau} R
= (A' \otimes_{R', \tau} R)  \times  (B' \otimes_{R', \tau} R)
$$
这就给出 (11) 中的分解。

\medskip\noindent
证明 8$\Rightarrow$10。假设 (8) 成立，并设 $R \to S$ 是有限环映射。
应用引理 \ref{algebra-lemma-etale-makes-quasi-finite-finite}。得到平展环映射
$R \to R'$ 以及位于 $\mathfrak m$ 之上的素理想
$\mathfrak m' \subset R'$，满足 $\kappa = \kappa(\mathfrak m')$，
并且 $R' \otimes_R S = A'_1 \times \ldots \times A'_n \times B'$，
其中 $A'_i$ 在 $R'$ 上有限且恰有一个位于 $\mathfrak m'$ 之上的素理想，
而 $B'$ 在位于 $\mathfrak m'$ 之上的任一素理想处都不在 $R'$ 上拟有限。
应用 (8)，得到收缩 $\tau : R' \to R$，满足
$\mathfrak m' = \tau^{-1}(\mathfrak m)$。于是
\begin{align*}
S & = (S \otimes_R R') \otimes_{R', \tau} R \\
& = (A'_1 \times \ldots \times A'_n \times B') \otimes_{R', \tau} R \\
& = (A'_1 \otimes_{R', \tau} R)  \times
\ldots \times (A'_1 \otimes_{R', \tau} R) \times
(B' \otimes_{R', \tau} R) \\
& = A_1 \times \ldots \times A_n \times B
\end{align*}
% P01-E0958: the frozen third line repeats A'_1 in the final A-factor; it should end with A'_n.
因子 $B$ 在 $R$ 上有限，但 $R \to B$ 在位于 $\mathfrak m$ 之上的
任一素理想处都不拟有限。因此 $B = 0$。各因子 $A_i$ 都是有限
$R$-代数，且恰有一个位于 $\mathfrak m$ 之上的素理想，因而都是局部环。
这证明 $S$ 是局部环的有限乘积。

\medskip\noindent
证明 9$\Rightarrow$10。这是因为，若 $S$ 在局部环 $R$ 上有限，
则它至多有有限多个极大理想。事实上，由 $R \to S$ 的上升性质，
$S$ 的所有极大理想都位于 $\mathfrak m$ 之上，而
$S/\mathfrak mS$ 是 Artin 环，因而只有有限多个素理想。

\medskip\noindent
证明 10$\Rightarrow$1。假设 (10) 成立。设 $f \in R[T]$ 是首一多项式，
$a_0 \in \kappa$ 是 $\overline{f}$ 的单根。则
$S = R[T]/(f)$ 是有限 $R$-代数。应用 (10)，得到
$S = A_1 \times \ldots \times A_r$，它是局部 $R$-代数的有限乘积。
特别地，$S/\mathfrak mS = \prod A_i/\mathfrak mA_i$ 是
$\kappa[T]/(\overline{f})$ 分解为局部环乘积的分解。这意味着其中一个因子，
例如 $A_1/\mathfrak mA_1$，就是商
$\kappa[T]/(\overline{f}) \to \kappa[T]/(T - a_0)$。
由于 $A_1$ 是有限自由 $R$-模 $S$ 的直和项，它自身也是有限自由
$R$-模。又因 $A_1/\mathfrak mA_1$ 是维数为 1 的
$\kappa$-向量空间，可知作为 $R$-模有 $A_1 \cong R$。
这显然意味着 $R \to A_1$ 是同构。令 $a \in R$ 为 $T$ 在映射
$R[T] \to S \to A_1 \to R$ 下的像。则如所需，
$f(a) = 0$ 且 $\overline{a} = a_0$。

\medskip\noindent
证明 13$\Rightarrow$1。假设 (13) 成立。设 $f \in R[T]$ 是首一多项式，
$a_0 \in \kappa$ 是 $\overline{f}$ 的单根。则
$S_1 = R[T]/(f)$ 是有限 $R$-代数。取任意 $g \in R[T]$，满足
$\overline{g} = \overline{f}/(T - a_0)$。则 $S = (S_1)_g$ 是拟有限
$R$-代数，并且
$S \otimes_R \kappa \cong \kappa[T]_{\overline{g}}/(\overline{f})
\cong \kappa[T]/(T - a_0) \cong \kappa$。
将 (13) 应用于 $S$，得到 $S = A \times B$，其中 $A$ 在 $R$ 上有限，
且 $B \otimes_R \kappa = 0$。特别地，
$\kappa \cong S/\mathfrak mS = A/\mathfrak mA$。由于 $A$ 是平坦
$R$-代数 $S$ 的直和项，可知它有限平坦，因而在 $R$ 上自由。
又因 $A/\mathfrak mA$ 是维数为 1 的 $\kappa$-向量空间，
可知作为 $R$-模有 $A \cong R$。这显然意味着 $R \to A$ 是同构。
令 $a \in R$ 为 $T$ 在映射 $R[T] \to S \to A \to R$ 下的像。
则如所需，$f(a) = 0$ 且 $\overline{a} = a_0$。

\medskip\noindent
证明 8$\Rightarrow$2。假设 (8) 成立。设 $f \in R[T]$ 是任意多项式，
并设 $a_0 \in \kappa$ 是单根。则代数
$S = R[T]_{f'}/(f)$ 在 $R$ 上平展。
% P01-E0959: the frozen proof calls a_0 “a simple root” without specifying that it is a root of f-bar.
设 $\mathfrak q \subset S$ 是由 $\mathfrak m$ 和 $T - b$ 生成的素理想，
其中 $b \in R$ 是满足 $\overline{b} = a_0$ 的任意元素。
将 (8) 应用于 $S$ 和 $\mathfrak q$，得到 $\tau : S \to R$。
于是像 $\tau(T) = a \in R$ 满足 (2)。

\medskip\noindent
至此可知，(1)、(2)、(7)、(8)、(9)、(10)、(11)、(12)、(13)
全都等价。(3)、(4)、(5)、(6) 中最弱的断言是 (3)，最强的是 (6)。
所以还须证明 (3) 推出 (1)，以及 (1) 推出 (6)。

\medskip\noindent
证明 3$\Rightarrow$1。假设 (3) 成立。设 $f \in R[T]$ 首一，
并设 $a_0 \in \kappa$ 是 $\overline{f}$ 的单根。这给出分解
$\overline{f} = (T - a_0)h_0$，其中 $h_0(a_0) \not = 0$，
故 $\gcd(T - a_0, h_0) = 1$。应用 (3)，得到分解
$f = gh$，其中 $\overline{g} = T - a_0$ 且
$\overline{h} = h_0$。置 $S = R[T]/(f)$，它是有限自由 $R$-代数。
我们也用 $g$、$h$ 表示 $g$、$h$ 在 $S$ 中的像。由于模掉
$\mathfrak m$ 后等式成立，由 Nakayama 引理 \ref{algebra-lemma-NAK}，有
$gS + hS = S$。又因为在 $S$ 中 $gh = f = 0$，还可得
$gS \cap hS = 0$。故由中国剩余定理，得到
$S = S/(g) \times S/(h)$。这意味着 $A = S/(g)$ 是有限自由
$R$-模的直和项，因而有限自由。此外，由
$A/\mathfrak mA = \kappa[T]/(T - a_0)$ 可知 $A$ 的秩为 $1$。
因此映射 $R \to A$ 是同构。令 $a \in R$ 等于 $T$ 在映射
$R[T] \to S \to A \to R$ 下的像，即得到满足
$f(a) = 0$ 且 $\overline{a} = a_0$ 的 $R$ 中元素。

\medskip\noindent
证明 1$\Rightarrow$6。假设 (1) 成立，或者等价地，假设
(1)、(2)、(7)、(8)、(9)、(10)、(11)、(12)、(13) 全部成立。
设 $f \in R[T]$ 是多项式。假设
$\overline{f} = g_0h_0$ 是满足 $\gcd(g_0, h_0) = 1$ 的分解。
可以且将要假设 $g_0$ 首一。由于有这一分解，可知 $f$ 的系数
在 $R$ 中生成单位理想。
% P01-E0960: if g_0 = 1 and h_0 = 0, the coprime factorization is valid but f-bar is zero, so the coefficients of f need not generate the unit ideal; the trivial factorization must be handled separately.
考察 $S = R[T]/(f)$。这意味着 $S$ 在 $R$ 上具有有限纤维，
因而在 $R$ 上拟有限。由引理 \ref{algebra-lemma-grothendieck-general}，
还可知 $S$ 在 $R$ 上平坦。结合 (13) 和 (10)，可写成
$S = A_1 \times \ldots \times A_n \times B$，其中每个 $A_i$
都是在 $R$ 上有限的局部环，且 $B \otimes_R \kappa = 0$。
重新排列因子 $A_1, \ldots, A_n$ 后，可以假设
$$
\kappa[T]/(g_0) =
A_1/\mathfrak m A_1 \times \ldots \times A_r/\mathfrak mA_r,
\ \kappa[T]/(h_0) =
A_{r + 1}/\mathfrak mA_{r + 1} \times \ldots \times A_n/\mathfrak mA_n
$$
作为 $\kappa[T]$ 的商。有限平坦 $R$-代数
$A = A_1 \times \ldots \times A_r$ 作为 $R$-模自由，见引理
\ref{algebra-lemma-finite-flat-local}。它的秩是 $\deg_T(g_0)$。
设 $g \in R[T]$ 是 $R$-线性算子 $T : A \to A$ 的特征多项式。
则 $g$ 是次数 $\deg_T(g) = \deg_T(g_0)$ 的首一多项式，
且 $\overline{g} = g_0$。由 Cayley--Hamilton 定理（引理
\ref{algebra-lemma-charpoly}），有 $g(T_A) = 0$，其中 $T_A$ 表示 $T$ 在
$A$ 中的像。因此得到良定义的满射 $R[T]/(g) \to A$，
它由 Nakayama 引理 \ref{algebra-lemma-NAK} 是同构。按构造，
映射 $R[T] \to A$ 经由 $R[T]/(f)$ 分解，因而对某个 $h$，
可以写成 $f = gh$。证明完成。
\end{proof}

\begin{lemma}
\label{algebra-lemma-finite-over-henselian}
设 $(R, \mathfrak m, \kappa)$ 是 Hensel 局部环。
\begin{enumerate}
\item 若 $R \to S$ 是有限环映射，则 $S$ 是 Hensel 局部环的有限乘积，
其中每个因子都在 $R$ 上有限。
\item 若 $R \to S$ 是有限环映射且 $S$ 是局部环，则 $S$ 是
Hensel 局部环，并且 $R \to S$ 是（有限）局部环映射。
\item 若 $R \to S$ 是有限型环映射，$\mathfrak q$ 是 $S$ 中位于
$\mathfrak m$ 之上的素理想，且 $R \to S$ 在该处拟有限，则
$S_{\mathfrak q}$ Hensel 且在 $R$ 上有限。
\item 若 $R \to S$ 拟有限，则对每个位于 $\mathfrak m$ 之上的素理想
$\mathfrak q$，$S_{\mathfrak q}$ 都 Hensel 且在 $R$ 上有限。
\end{enumerate}
\end{lemma}

\begin{proof}
(2) 推出 (1)，因为由引理 \ref{algebra-lemma-characterize-henselian} (10)，
(1) 中的 $S$ 是它在位于 $\mathfrak m$ 之上的素理想处的局部化的有限乘积。
(2) 也由引理 \ref{algebra-lemma-characterize-henselian} (10) 推出，
因为任意有限 $S$-代数也都是有限 $R$-代数（当然，局部环之间的任意
有限环映射都是局部映射）。

\medskip\noindent
设 $R \to S$ 和 $\mathfrak q$ 如 (3) 中所设。写成
$S = A \times B$，其中 $A$ 在 $R$ 上有限，而 $B$ 在位于
$\mathfrak m$ 之上的任一素理想处都不在 $R$ 上拟有限；见引理
\ref{algebra-lemma-characterize-henselian} (11)。因此 $S_\mathfrak q$
是 $A$ 在某个极大理想处的局部化，由 (1) 可得 (3)。
(4) 由 (3) 推出。
\end{proof}

\begin{lemma}
\label{algebra-lemma-mop-up}
设 $(R, \mathfrak m, \kappa)$ 是 Hensel 局部环。
任意有限型 $R$-代数 $S$ 都可写成
$S = A_1 \times \ldots \times A_n \times B$，其中 $A_i$ 是在
$R$ 上有限的局部环，而 $R \to B$ 在 $B$ 中位于
$\mathfrak m$ 之上的任一素理想处都不拟有限。
\end{lemma}

\begin{proof}
这是引理 \ref{algebra-lemma-characterize-henselian} 的 (11) 与 (10) 的结合。
\end{proof}

\begin{lemma}
\label{algebra-lemma-mop-up-strictly-henselian}
设 $(R, \mathfrak m, \kappa)$ 是严格 Hensel 局部环。
任意有限型 $R$-代数 $S$ 都可写成
$S = A_1 \times \ldots \times A_n \times B$，其中 $A_i$ 是在
$R$ 上有限的局部环，$\kappa \subset \kappa(\mathfrak m_{A_i})$
有限纯不可分，而 $R \to B$ 在 $B$ 中位于 $\mathfrak m$ 之上的
任一素理想处都不拟有限。
\end{lemma}

\begin{proof}
先按引理 \ref{algebra-lemma-mop-up} 写成
$S = A_1 \times \ldots \times A_n \times B$。
域扩张 $\kappa(\mathfrak m_{A_i})/\kappa$ 有限，而 $\kappa$
可分代数闭，故该扩张有限纯不可分。
\end{proof}

\begin{lemma}
\label{algebra-lemma-henselian-cat-finite-etale}
设 $(R, \mathfrak m, \kappa)$ 是 Hensel 局部环。有限平展环扩张
$R \to S$ 的范畴，借助函子 $S \mapsto S/\mathfrak mS$，
等价于有限平展代数 $\kappa \to \overline{S}$ 的范畴。
\end{lemma}

\begin{proof}
以 $\mathcal{C} \to \mathcal{D}$ 表示陈述中的范畴函子。
假设 $R \to S$ 有限平展。利用引理 \ref{algebra-lemma-mop-up} 或引理
\ref{algebra-lemma-characterize-henselian} (10)，可以写成
$$
S = A_1 \times \ldots \times A_n
$$
其中 $A_i$ 是在 $S$ 上有限平展的局部环。特别地，
$A_i/\mathfrak mA_i$ 是 $\kappa$ 的有限可分域扩张，见引理
\ref{algebra-lemma-etale-at-prime}。由此可见，$\mathcal{C}$ 和
$\mathcal{D}$ 中的每个对象都典范分解为不可约部分，
这些部分经给定函子相互对应。
其次，假设 $S_1$、$S_2$ 在 $R$ 上有限平展，并且
$\kappa_1 = S_1/\mathfrak mS_1$ 与
$\kappa_2 = S_2/\mathfrak mS_2$ 都是域（即 $\kappa$ 的有限可分扩张）。
则 $S_1 \otimes_R S_2$ 在 $R$ 上有限平展，并且可以像上面一样写成
$$
S_1 \otimes_R S_2 = A_1 \times \ldots \times A_n
$$
于是 $\Hom_R(S_1, S_2)$ 与满足 $S_2 \to A_i$ 为同构的指标
$i \in \{1, \ldots, n\}$ 的集合相等同。为说明这一点，注意给定任意
$R$-代数映射 $\varphi : S_1 \to S_2$，映射
$\varphi \times 1 : S_1 \otimes_R S_2 \to S_2$ 是满射，
因而等于向某个因子 $A_i$ 的投影。完全相同地，
$\Hom_\kappa(\kappa_1, \kappa_2)$ 与满足
$\kappa_2 \to A_i/\mathfrak mA_i$ 为同构的指标
$i \in \{1, \ldots, n\}$ 的集合相等同。由上面的讨论，
这两个指标集合相同，故该函子完全忠实。
最后，设 $\kappa'/\kappa$ 是有限可分域扩张。由引理
\ref{algebra-lemma-make-etale-map-prescribed-residue-field}，存在平展环映射
$R \to S$ 以及 $S$ 中位于 $\mathfrak m$ 之上的素理想
$\mathfrak q$，使得 $\kappa \subset \kappa(\mathfrak q)$
同构于给定扩张。由 (1)，可以写成
$S = A_1 \times \ldots \times A_n \times B$。
% P01-E0961: the frozen proof cites an undefined “part (1)” here; the required global finite-plus-non-quasi-finite decomposition is lemma-mop-up.
由于 $R \to S$ 拟有限，可知 $B$ 中不存在位于 $\mathfrak m$ 之上的素理想。
因此 $S_{\mathfrak q}$ 等于某个指标 $i$ 所对应的 $A_i$。故 $R \to A_i$ 有限平展，
并给出所指定的剩余域扩张。于是该函子本质满，结论得证。
\end{proof}

\begin{lemma}
\label{algebra-lemma-unramified-over-strictly-henselian}
设 $(R, \mathfrak m, \kappa)$ 是严格 Hensel 局部环。
设 $R \to S$ 是非分歧环映射。则
$$
S = A_1 \times \ldots \times A_n \times B
$$
其中每个 $R \to A_i$ 都是满射，且 $B$ 中没有位于
$\mathfrak m$ 之上的素理想。
\end{lemma}

\begin{proof}
先按引理 \ref{algebra-lemma-mop-up} 写成
$S = A_1 \times \ldots \times A_n \times B$。
此时 $R \to A_i$ 有限非分歧，而 $A_i$ 是局部环。故 $A_i$
的极大理想是 $\mathfrak mA_i$，其剩余域
$A_i / \mathfrak m A_i$ 是 $\kappa$ 的有限可分扩张，见引理
\ref{algebra-lemma-unramified-at-prime}。然而，$R$ 严格 Hensel 意味着
$\kappa$ 可分代数闭，所以 $\kappa = A_i / \mathfrak m A_i$。
由 Nakayama 引理 \ref{algebra-lemma-NAK}，可知 $R \to A_i$ 如所需地为满射。
\end{proof}

\begin{lemma}
\label{algebra-lemma-complete-henselian}
\begin{slogan}
完备局部环由 Newton 法可知为 Hensel 环
\end{slogan}
设 $(R, \mathfrak m, \kappa)$ 是完备局部环，见定义
\ref{algebra-definition-complete-local-ring}。则 $R$ Hensel。
\end{lemma}

\begin{proof}
设 $f \in R[T]$ 首一。以
$f_n \in R/\mathfrak m^{n + 1}[T]$ 表示其像，以 $f'_n$ 表示
$f_n$ 关于 $T$ 的导数。设 $a_0 \in \kappa$ 是 $f_0$ 的单根。
如下归纳地把它提升为 $f$ 在 $R$ 上的解：假设已经给定
$a_n \in R/\mathfrak m^{n + 1}$，满足
$a_n \bmod \mathfrak m = a_0$ 且 $f_n(a_n) = 0$。
任取元素 $b \in R/\mathfrak m^{n + 2}$，满足
$a_n = b \bmod \mathfrak m^{n + 1}$。则
$f_{n + 1}(b) \in \mathfrak m^{n + 1}/\mathfrak m^{n + 2}$。置
$$
a_{n + 1} = b - f_{n + 1}(b)/f'_{n + 1}(b)
$$
（Newton 法）。这是有意义的，因为由对 $a_0$ 的条件，
$f'_{n + 1}(b) \in R/\mathfrak m^{n + 2}$ 可逆。于是计算得
$f_{n + 1}(a_{n + 1}) = f_{n + 1}(b) - f_{n + 1}(b) = 0$
于 $R/\mathfrak m^{n + 2}$。如此构造的元素系
$a_n \in R/\mathfrak m^{n + 1}$ 相容，故得到元素
$a \in \lim R/\mathfrak m^n = R$（这里使用了 $R$ 完备）。
此外，$f(a) = 0$，因为它在每个 $R/\mathfrak m^n$ 中的像都为零。
最后，$\overline{a} = a_0$，结论得证。
\end{proof}

\begin{lemma}
\label{algebra-lemma-local-dimension-zero-henselian}
\begin{slogan}
零维局部环是 Hensel 环。
\end{slogan}
设 $(R, \mathfrak m)$ 是维数为 $0$ 的局部环。则 $R$ Hensel。
\end{lemma}

\begin{proof}
设 $R \to S$ 是有限环映射。由引理
\ref{algebra-lemma-characterize-henselian}，只须证明 $S$ 是局部环的乘积。
由引理 \ref{algebra-lemma-finite-finite-fibres}，$S$ 有有限多个素理想
$\mathfrak m_1, \ldots, \mathfrak m_r$，它们全都位于
$\mathfrak m$ 之上。由引理 \ref{algebra-lemma-integral-no-inclusion}，
这些素理想之间不存在包含关系，故它们全都极大。由引理
\ref{algebra-lemma-Zariski-topology}，
$\mathfrak m_1 \cap \ldots \cap \mathfrak m_r$ 的每个元素都幂零。
于是由引理 \ref{algebra-lemma-product-local}，$S$ 是 $S$ 在各素理想
$\mathfrak m_i$ 处局部化的乘积。
\end{proof}

\noindent
下一个引理将是 Hensel 化与严格 Hensel 化的唯一性和函子性质的关键。

\begin{lemma}
\label{algebra-lemma-map-into-henselian}
设 $R \to S$ 是环映射，且 $S$ Hensel 局部。给定
\begin{enumerate}
\item 平展环映射 $R \to A$；
\item $A$ 中位于 $\mathfrak p = R \cap \mathfrak m_S$ 之上的素理想
$\mathfrak q$；
% P01-E0962: for a possibly noninjective map R -> S, the frozen expression R intersect m_S should be the inverse image of m_S in R.
\item $\kappa(\mathfrak p)$-代数映射
$\tau : \kappa(\mathfrak q) \to S/\mathfrak m_S$，
\end{enumerate}
则存在唯一的 $R$-代数同态 $f : A \to S$，满足
$\mathfrak q = f^{-1}(\mathfrak m_S)$，且 $f$ 在剩余域上诱导映射
$\tau$。
\end{lemma}

\begin{proof}
考察 $A \otimes_R S$。由引理 \ref{algebra-lemma-etale}，
它是在 $S$ 上的平展代数。此外，核
$$
\mathfrak q' = \Ker(A \otimes_R S \to
\kappa(\mathfrak q) \otimes_{\kappa(\mathfrak p)} \kappa(\mathfrak m_S)
\xrightarrow{\tau \otimes 1}
\kappa(\mathfrak m_S))
$$
是位于 $\mathfrak m_S$ 之上的素理想，且其剩余域等于 $S$ 的剩余域。
因此，由引理 \ref{algebra-lemma-characterize-henselian}，存在唯一收缩
$\sigma : A \otimes_R S \to S$，满足
$\sigma^{-1}(\mathfrak m_S) = \mathfrak q'$。令 $f$ 等于复合映射
$A \to A \otimes_R S \to S$。略去对 $f$ 的各项性质的验证；
$f$ 的唯一性来自 $\sigma$ 的唯一性（细节从略）。
\end{proof}

\begin{lemma}
\label{algebra-lemma-strictly-henselian-solutions}
设 $\varphi : R \to S$ 是严格 Hensel 局部环的局部同态。
设 $P_1, \ldots, P_n \in R[x_1, \ldots, x_n]$ 是多项式，
使得 $R[x_1, \ldots, x_n]/(P_1, \ldots, P_n)$ 在 $R$ 上平展。
则映射
$$
R^n \longrightarrow S^n, \quad
(h_1, \ldots, h_n) \longmapsto (\varphi(h_1), \ldots, \varphi(h_n))
$$
在下列两个集合之间诱导双射：
$$
\{
(r_1, \ldots, r_n) \in R^n
\mid
P_i(r_1, \ldots, r_n) = 0, \ i = 1, \ldots, n
\}
$$
以及
$$
\{
(s_1, \ldots, s_n) \in S^n
\mid
P^\varphi_i(s_1, \ldots, s_n) = 0, \ i = 1, \ldots, n
\}
$$
其中 $P^\varphi_i \in S[x_1, \ldots, x_n]$ 是 $P_i$ 在
$\varphi$ 下的像。
\end{lemma}

\begin{proof}
第一个解集典范同构于集合
$$
\Hom_R(R[x_1, \ldots, x_n]/(P_1, \ldots, P_n), R).
$$
由于 $R$ Hensel，映射 $R \to R/\mathfrak m_R$ 在此集合与剩余域
$R/\mathfrak m_R$ 中的解集之间诱导双射，见引理
\ref{algebra-lemma-characterize-henselian}。对 $S$ 同样如此。
现在，由于 $R[x_1, \ldots, x_n]/(P_1, \ldots, P_n)$ 在 $R$ 上平展，
且 $R/\mathfrak m_R$ 可分代数闭，可知
$R/\mathfrak m_R[x_1, \ldots, x_n]/(\overline{P}_1, \ldots, \overline{P}_n)$
是 $R/\mathfrak m_R$ 的若干副本的有限乘积；其中
$\overline{P}_i$ 是 $P_i$ 在
$R/\mathfrak m_R[x_1, \ldots, x_n]$ 中的像。
% P01-E0963: the frozen residue-field polynomial-ring expressions omit parentheses around R/m_R (and later S/m_S), making the displayed quotients formally ill-typed or ambiguous.
因此张量积
$$
R/\mathfrak m_R[x_1, \ldots, x_n]/(\overline{P}_1, \ldots, \overline{P}_n)
\otimes_{R/\mathfrak m_R} S/\mathfrak m_S
=
S/\mathfrak m_S[x_1, \ldots, x_n]/
(\overline{P}^\varphi_1, \ldots, \overline{P}^\varphi_n)
$$
也是 $S/\mathfrak m_S$ 的若干副本在同一个指标集上的有限乘积。
这证明了引理。
\end{proof}

\begin{lemma}
\label{algebra-lemma-split-ML-henselian}
设 $R$ 是 Hensel 局部环。任意可数生成的 Mittag--Leffler
$R$-模都是有限表现 $R$-模的直和。
\end{lemma}

\begin{proof}
设 $M$ 是可数生成的 Mittag--Leffler $R$-模。我们断言：
对任意元素 $x \in M$，都存在直和分解 $M = N \oplus K$，
使得 $x \in N$，模 $N$ 有限表现，而 $K$ 是 Mittag--Leffler 模。

\medskip\noindent
假设该断言成立。选取 $M$ 的生成元 $x_1, x_2, x_3, \ldots$。
由该断言，可以归纳地找到直和分解
$$
M = N_1 \oplus N_2 \oplus \ldots \oplus N_n \oplus K_n
$$
其中 $N_i$ 有限表现，
$x_1, \ldots, x_n \in N_1 \oplus \ldots \oplus N_n$，
且 $K_n$ 是 Mittag--Leffler 模。无限重复即可得到
$M = \bigoplus N_i$。

\medskip\noindent
还须证明该断言。设 $x \in M$。由引理 \ref{algebra-lemma-ML-countable}，
存在自同态 $\alpha : M \to M$，使得 $\alpha$ 经由有限表现模分解，
且 $\alpha (x) = x$。设 $\alpha$ 分解为
$$
\xymatrix{
M \ar[r]^\pi & P \ar[r]^i & M
}
$$
置 $a = \pi \circ \alpha \circ i : P \to P$，则
$i \circ a \circ \pi = \alpha^3$。由引理
\ref{algebra-lemma-charpoly-module}，存在首一多项式 $P \in R[T]$，
使得 $P(a) = 0$。
% P01-E0964: the frozen proof reuses P for both the finitely presented module in the diagram and the annihilating monic polynomial, making later references such as “N is a quotient of P” ambiguous.
注意，这形式地推出 $\alpha^2 P(\alpha) = 0$。因此可以把 $M$
看作 $R[T]/(T^2P)$-模。假设 $x \not = 0$。则
$\alpha(x) = x$ 推出 $0 = \alpha^2P(\alpha)x = P(1)x$，
故 $P(1) = 0$ 于 $R/I$，其中
$I = \{r \in R \mid rx = 0\}$ 是 $x$ 的零化子。
由于 $x \not = 0$，有 $I \subset \mathfrak m_R$，所以 $1$ 是
$\overline{P} = P \bmod \mathfrak m_R \in R/\mathfrak m_R[T]$ 的根。
由于 $R$ Hensel，可以找到分解
$$
T^2P = (T^2 Q_1) Q_2
$$
其中 $Q_1, Q_2 \in R[T]$，并且
$Q_2 = (T - 1)^e \bmod \mathfrak m_R R[T]$、
$Q_1(1) \not = 0 \bmod \mathfrak m_R$；见引理
\ref{algebra-lemma-characterize-henselian}。令
$N = \Im(\alpha^2Q_1(\alpha) : M \to M)$，并令
$K = \Im(Q_2(\alpha) : M \to M)$。由于 $T^2Q_1$ 与 $Q_2$
生成 $R[T]$ 的单位理想，得到直和分解 $M = N \oplus K$。
此外，$Q_2$ 在 $N$ 上作用为零，$T^2Q_1$ 在 $K$ 上作用为零。
注意，$N$ 是 $P$ 的商，因而有限生成。又有 $x \in N$，因为
$\alpha^2Q_1(\alpha)x = Q_1(1)x$，而 $Q_1(1)$ 是 $R$ 中的单位。
由引理 \ref{algebra-lemma-direct-sum-ML}，模 $N$ 和 $K$ 都是
Mittag--Leffler 模。最后，有限生成模 $N$ 有限表现，因为有限生成的
Mittag--Leffler 模有限表现，见例 \ref{algebra-example-ML} (1)。
\end{proof}


\section{平展环映射的滤过余极限}
\label{algebra-section-ind-etale}

\noindent
本节是关于 ind-平展环映射一节（《pro-平展上同调》，第
\ref{proetale-section-ind-etale} 节）的先导。
这些材料也将用于证明局部环的 Hensel 化与严格 Hensel 化的唯一性。

\begin{lemma}
\label{algebra-lemma-base-change-colimit-etale}
设 $R \to A$ 与 $R \to R'$ 是环映射。若 $A$ 是平展
$R$-代数的滤过余极限，则 $R' \otimes_R A$ 也是平展
$R'$-代数的滤过余极限。
% P01-E0965: the frozen English statement duplicates its predicate (“then so is ... is a filtered colimit”).
\end{lemma}

\begin{proof}
这是因为余极限与张量积交换，而平展环映射在基变换下保持
（引理 \ref{algebra-lemma-etale}）。
\end{proof}

\begin{lemma}
\label{algebra-lemma-composition-colimit-etale}
设 $A \to B \to C$ 是环映射。若 $B$ 是平展
$A$-代数的滤过余极限，且 $C$ 是平展 $B$-代数的滤过余极限，
则 $C$ 是平展 $A$-代数的滤过余极限。
\end{lemma}

\begin{proof}
我们将使用引理 \ref{algebra-lemma-when-colimit} 的判别准则。
设 $A \to P \to C$ 是 $A \to C$ 的一个分解，其中 $P$
在 $A$ 上有限表现。写
$B = \colim_{i \in I} B_i$，其中 $I$ 是有向集，而 $B_i$
是平展 $A$-代数。写 $C = \colim_{j \in J} C_j$，其中
$J$ 是有向集，而 $C_j$ 是平展 $B$-代数。由引理
\ref{algebra-lemma-characterize-finite-presentation}，对某个 $j$，
可以把 $P \to C$ 分解为 $P \to C_j \to C$。
由引理 \ref{algebra-lemma-etale}，可以找到 $i \in I$ 以及平展环映射
$B_i \to C'_j$，使得 $C_j = B \otimes_{B_i} C'_j$。
于是
$C_j = \colim_{i' \geq i} B_{i'} \otimes_{B_i} C'_j$，
再次可见 $P \to C_j$ 分解为
$P \to B_{i'} \otimes_{B_i} C'_j \to C$。
由于平展环映射的复合与张量积仍为平展，
$A \to C' = B_{i'} \otimes_{B_i} C'_j$ 是平展的。
% P01-E0966: the frozen English source begins this sentence with “As” and then starts a second causal “as”, leaving a sentence fragment.
因此我们已经把 $P \to C$ 分解为 $P \to C' \to C$，
其中 $C'$ 在 $A$ 上平展，故引理 \ref{algebra-lemma-when-colimit}
的判别准则适用。
\end{proof}

\begin{lemma}
\label{algebra-lemma-colimit-colimit-etale}
设 $R$ 是环。设 $A = \colim A_i$ 是 $R$-代数的滤过余极限，
并且每个 $A_i$ 都是平展 $R$-代数的滤过余极限。
则 $A$ 是平展 $R$-代数的滤过余极限。
\end{lemma}

\begin{proof}
写 $A_i = \colim_{j \in J_i} A_j$，其中 $J_i$ 是有向集，
而 $A_j$ 是平展 $R$-代数。对每个 $i \leq i'$ 和
$j \in J_i$，存在 $j' \in J_{i'}$ 以及 $R$-代数映射
$\varphi_{jj'} : A_j \to A_{j'}$，使图表
$$
\xymatrix{
A_i \ar[r] & A_{i'} \\
A_j \ar[u] \ar[r]^{\varphi_{jj'}} & A_{j'} \ar[u]
}
$$
交换。这是因为 $R \to A_j$ 有限表现，所以引理
\ref{algebra-lemma-characterize-finite-presentation} 适用。
令 $\mathcal{J}$ 为如下范畴：其对象为
$\coprod_{i \in I} J_i$，其态为如上的三元组
$(j, j', \varphi_{jj'})$（复合律是显然的）。
则 $\mathcal{J}$ 是滤过范畴，并且
$A = \colim_\mathcal{J} A_j$。细节从略。
\end{proof}

\begin{lemma}
\label{algebra-lemma-colimit-colimit-etale-better}
设 $I$ 是有向集。设 $i \mapsto (R_i \to A_i)$ 是 $I$
上的环箭头系统。置 $R = \colim R_i$、$A = \colim A_i$。
若每个 $A_i$ 都是平展 $R_i$-代数的滤过余极限，
则 $A$ 是平展 $R$-代数的滤过余极限。
\end{lemma}

\begin{proof}
这是因为
$A = A \otimes_R R = \colim A_i \otimes_{R_i} R$，
因而可以应用引理 \ref{algebra-lemma-colimit-colimit-etale}：
由引理 \ref{algebra-lemma-base-change-colimit-etale}，
$R \to A_i \otimes_{R_i} R$ 是平展环映射的滤过余极限。
\end{proof}

\begin{lemma}
\label{algebra-lemma-colimits-of-etale}
设 $R$ 是环。设 $A \to B$ 是 $R$-代数同态。
若 $A$ 与 $B$ 都是平展 $R$-代数的滤过余极限，
则 $B$ 是平展 $A$-代数的滤过余极限。
\end{lemma}

\begin{proof}
把 $A = \colim A_i$ 与 $B = \colim B_j$ 写成滤过余极限，
其中 $A_i$ 和 $B_j$ 在 $R$ 上平展。对每个 $i$，
可以找到 $j$，使 $A_i \to B$ 经 $B_j$ 分解，见引理
\ref{algebra-lemma-characterize-finite-presentation}。
由引理 \ref{algebra-lemma-map-between-etale}，分解映射
$A_i \to B_j$ 是平展的。由于
$A \to A \otimes_{A_i} B_j$ 平展（引理 \ref{algebra-lemma-etale}），
只须证明
$B = \colim A \otimes_{A_i} B_j$，其中余极限取遍
二元组 $(i, j)$ 以及 $A_i \to B$ 的分解
$A_i \to B_j \to B$（这是一个有向系统；细节从略）。
这是清楚的，因为余极限与张量积交换，故
$\colim A \otimes_{A_i} B_j = A \otimes_A B = B$。
\end{proof}

\begin{lemma}
\label{algebra-lemma-map-into-henselian-colimit}
设 $R \to S$ 是环映射，且 $S$ 是 Hensel 局部环。给定
\begin{enumerate}
\item 一个 $R$-代数 $A$，它是平展 $R$-代数的滤过余极限；
\item $A$ 的一个素理想 $\mathfrak q$，它位于
$\mathfrak p = R \cap \mathfrak m_S$ 之上；
% P01-E0967: as in P01-E0962, the frozen source uses an intersection for the contraction along a possibly noninjective ring map.
\item 一个 $\kappa(\mathfrak p)$-代数映射
$\tau : \kappa(\mathfrak q) \to S/\mathfrak m_S$。
\end{enumerate}
则存在唯一的 $R$-代数同态 $f : A \to S$，使得
$\mathfrak q = f^{-1}(\mathfrak m_S)$ 且
$f \bmod \mathfrak q = \tau$。
\end{lemma}

\begin{proof}
把 $A = \colim A_i$ 写成平展 $R$-代数的滤过余极限。
置 $\mathfrak q_i = A_i \cap \mathfrak q$。
% P01-E0968: the frozen source again writes an intersection although the transition map A_i -> A need not be injective; the contraction is the inverse image.
应用引理 \ref{algebra-lemma-map-into-henselian}，得到
$f_i : A_i \to S$。置 $f = \colim f_i$。
\end{proof}

\begin{lemma}
\label{algebra-lemma-uniqueness-henselian}
设 $R$ 是环。给定环映射的交换图表
$$
\xymatrix{
S \ar[r] & K \\
R \ar[u] \ar[r] & S' \ar[u]
}
$$
其中 $S$、$S'$ 是 Hensel 局部环，$S$、$S'$ 是平展
$R$-代数的滤过余极限，$K$ 是域，并且箭头 $S \to K$ 与
$S' \to K$ 都把 $K$ 等同于 $S$、$S'$ 的剩余域。
则存在唯一的 $R$-代数同构 $S \to S'$，它与到 $K$ 的映射相容。
\end{lemma}

\begin{proof}
这立即由引理 \ref{algebra-lemma-map-into-henselian-colimit} 得出。
\end{proof}

\noindent
下面的引理严格说来并不涉及平展环映射的余极限。

\begin{lemma}
\label{algebra-lemma-colimit-henselian}
沿局部同态所取的（严格）Hensel 局部环的滤过余极限仍是
（严格）Hensel 的。
\end{lemma}

\begin{proof}
《范畴》，引理 \ref{categories-lemma-directed-category-system}
说明，实际上只须考虑有向集上的（严格）Hensel 局部环的余极限。
设 $(R_i, \varphi_{ii'})$ 是这样的系统，每个
$\varphi_{ii'}$ 都是局部的。则 $R = \colim_i R_i$ 是局部环，
其剩余域 $\kappa$ 为 $\colim \kappa_i$（论证从略）。
容易看出，若每个 $\kappa_i$ 都可分代数闭，则
$\colim \kappa_i$ 也可分代数闭；因此只须证明：
若每个 $R_i$ 都是 Hensel 的，则 $R$ 是 Hensel 的。
假设 $f \in R[T]$ 首一，并且 $a_0 \in \kappa$ 是
$\overline{f}$ 的单根。则对某个充分大的 $i$，
存在映到 $f$ 的 $f_i \in R_i[T]$，以及映到 $a_0$ 的
$a_{0, i} \in \kappa_i$。由于
$\overline{f_i}(a_{0, i}) \in \kappa_i$ 与
$\overline{f_i'}(a_{0, i}) \in \kappa_i$ 分别映到
$0 = \overline{f}(a_0) \in \kappa$ 与
$0 \not = \overline{f'}(a_0) \in \kappa$，
可知 $a_{0, i}$ 是 $\overline{f_i}$ 的单根。
由于 $R_i$ 是 Hensel 的，可以找到 $a_i \in R_i$，使得
$f_i(a_i) = 0$ 且 $a_{0, i} = \overline{a_i}$。
于是 $a_i$ 的像 $a \in R$ 就是所求解。
因此 $R$ 是 Hensel 的。
\end{proof}


\section{Hensel 化与严格 Hensel 化}
\label{algebra-section-henselization}

\noindent
本节构造 Hensel 化。阅读这些材料时，我们建议读者始终记住引理
\ref{algebra-lemma-uniqueness-henselian} 已经证明的唯一性，以及引理
\ref{algebra-lemma-map-into-henselian-colimit} 指出的函子性。

\begin{lemma}
\label{algebra-lemma-henselization}
设 $(R, \mathfrak m, \kappa)$ 是局部环。存在局部环映射
$R \to R^h$，它具有下列性质：
\begin{enumerate}
\item $R^h$ 是 Hensel 的；
\item $R^h$ 是平展 $R$-代数的滤过余极限；
\item $\mathfrak m R^h$ 是 $R^h$ 的极大理想；且
\item $\kappa = R^h/\mathfrak m R^h$。
\end{enumerate}
\end{lemma}

\begin{proof}
考虑由偶对 $(S, \mathfrak q)$ 构成的范畴，其中 $R \to S$
是平展环映射，而 $\mathfrak q$ 是 $S$ 中位于
$\mathfrak m$ 之上的素理想，并且
$\kappa = \kappa(\mathfrak q)$。偶对之间的态
$(S, \mathfrak q) \to (S', \mathfrak q')$ 由
$R$-代数映射 $\varphi : S \to S'$ 给出，并且满足
$\varphi^{-1}(\mathfrak q') = \mathfrak q$。置
$$
R^h = \colim_{(S, \mathfrak q)} S.
$$
我们来证明这个偶对范畴是滤过的；见《范畴》，定义
\ref{categories-definition-directed}。该范畴含有偶对
$(R, \mathfrak m)$，因而非空；这证明了《范畴》，定义
\ref{categories-definition-directed} 的 (1)。
对任意偶对 $(S, \mathfrak q)$，素理想 $\mathfrak q$
都是极大理想，且其剩余域为 $\kappa$，因为复合
$\kappa \to S/\mathfrak q \to \kappa(\mathfrak q)$ 是同构。
假设 $(S, \mathfrak q)$ 与 $(S', \mathfrak q')$ 是两个对象。
置
$S'' = S \otimes_R S'$ 且
$\mathfrak q'' = \mathfrak qS'' + \mathfrak q'S''$。
由上面的说明，
$S''/\mathfrak q'' = S/\mathfrak q \otimes_R S'/\mathfrak q' = \kappa$。
此外，由引理 \ref{algebra-lemma-etale}，$R \to S''$ 平展。
这证明了《范畴》，定义 \ref{categories-definition-directed} 的 (2)。

其次，假设
$\varphi, \psi : (S, \mathfrak q) \to (S', \mathfrak q')$
是偶对的两个态。由引理 \ref{algebra-lemma-map-between-etale}，
$\varphi$、$\psi$ 以及 $S' \otimes_R S' \to S'$
都是平展环映射。考虑
$$
S'' = (S' \otimes_{\varphi, S, \psi} S')
\otimes_{S' \otimes_R S'} S'
$$
及其中的素理想
$$
\mathfrak q'' =
(\mathfrak q' \otimes S' + S' \otimes \mathfrak q') \otimes S'
+
(S' \otimes_{\varphi, S, \psi} S') \otimes \mathfrak q'
$$
与上面同样论证（平展映射的基变换是平展的，平展映射的复合也是
平展的），可知 $S''$ 在 $R$ 上平展。此外，典范映射
$S' \to S''$（例如使用最右边的因子）使 $\varphi$ 与 $\psi$
相等。这证明了《范畴》，定义
\ref{categories-definition-directed} 的 (3)。
因此可知，$R^h$ 的元由三元组 $(S, \mathfrak q, f)$ 表示，
其中 $f \in S$；两个这样的三元组
$(S, \mathfrak q, f)$、$(S', \mathfrak q', f')$
定义 $R^h$ 的同一个元，当且仅当存在偶对
$(S'', \mathfrak q'')$ 以及偶对的态
$\varphi : (S, \mathfrak q) \to (S'', \mathfrak q'')$ 和
$\varphi' : (S', \mathfrak q') \to (S'', \mathfrak q'')$，
使得 $\varphi(f) = \varphi'(f')$。

\medskip\noindent
假设 $x \in R^h$。用三元组 $(S, \mathfrak q, f)$ 表示 $x$。
设 $\mathfrak q_1, \ldots, \mathfrak q_r$ 是 $S$ 中位于
$\mathfrak m$ 之上的其他素理想。由于上面已经看到
$\mathfrak q$ 极大，故 $\mathfrak q \not \subset \mathfrak q_i$。
于是，因为 $\mathfrak q$ 是素理想，可以找到
$g \in S$、$g \not \in \mathfrak q$，并且对
$i = 1, \ldots, r$ 有 $g \in \mathfrak q_i$。
考虑偶对的态
$(S, \mathfrak q) \to (S_g, \mathfrak qS_g)$。
这样就看出，总可假设 $x$ 由三元组
$(S, \mathfrak q, f)$ 给出，其中 $\mathfrak q$ 是 $S$ 中
位于 $\mathfrak m$ 之上的唯一素理想，即
$\sqrt{\mathfrak mS} = \mathfrak q$。但由于 $R \to S$
平展，由引理 \ref{algebra-lemma-etale-at-prime}，
$\mathfrak mS_{\mathfrak q} = \mathfrak qS_{\mathfrak q}$。
因此实际上有 $\mathfrak mS = \mathfrak q$。

\medskip\noindent
假设 $x \not \in \mathfrak mR^h$。用满足
$\mathfrak mS = \mathfrak q$ 的三元组
$(S, \mathfrak q, f)$ 表示 $x$。
则 $f \not \in \mathfrak mS$，即 $f \not \in \mathfrak q$。
因此 $(S, \mathfrak q) \to (S_f, \mathfrak qS_f)$
是偶对的态，且 $f$ 的像在其中可逆。
所以 $x$ 可逆，其逆由三元组
$(S_f, \mathfrak qS_f, 1/f)$ 表示。
由此得出 $R^h$ 是局部环，其极大理想为 $\mathfrak mR^h$。
剩余域为 $\kappa$，因为可以定义
$R^h/\mathfrak mR^h \to \kappa$：把三元组
$(S, \mathfrak q, f)$ 映到 $f$ 模 $\mathfrak q$ 的剩余类。

\medskip\noindent
还须证明 $R^h$ 是 Hensel 的。事实上，假设
$P \in R^h[T]$ 是首一多项式，而 $a_0 \in \kappa$
是其约化 $\overline{P} \in \kappa[T]$ 的单根。
则可以找到偶对 $(S, \mathfrak q)$，使得 $P$ 是某个首一多项式
$Q \in S[T]$ 的像。置 $S' = S[T]/(Q)$，并令
$\mathfrak q' \subset S'$ 为极大理想
$\mathfrak q' = \mathfrak qS' + (T - a')S'$，
其中 $a' \in S$ 是 $a_0$ 的任意提升。
由构造，$S \to S'$ 在 $\mathfrak q'$ 处平展，并且
$\kappa = \kappa(\mathfrak q')$。选取
$g \in S'$、$g \not \in \mathfrak q'$，使得
$S'' = S'_g$ 在 $S$ 上平展。于是
$(S, \mathfrak q) \to (S'', \mathfrak q'S'')$ 是偶对的态。
三元组 $(S'', \mathfrak q'S'', \text{下列对象的类： }T)$ 现在确定
$a \in R^h$，满足 $P(a) = 0$ 且
$\overline{a} = a_0$，正如所需。
% P01-E0969: the frozen English source says “Now that triple”, a grammatical substitution of “that” for the article “the”.
\end{proof}

\begin{lemma}
\label{algebra-lemma-strict-henselization}
设 $(R, \mathfrak m, \kappa)$ 是局部环。
设 $\kappa \subset \kappa^{sep}$ 是可分代数闭包。
存在交换图表
$$
\xymatrix{
\kappa \ar[r] & \kappa \ar[r] & \kappa^{sep} \\
R \ar[r] \ar[u] & R^h \ar[r] \ar[u] & R^{sh} \ar[u]
}
$$
它具有下列性质：
\begin{enumerate}
\item 映射 $R^h \to R^{sh}$ 是局部的；
\item $R^{sh}$ 是严格 Hensel 的；
\item $R^{sh}$ 是平展 $R$-代数的滤过余极限；
\item $\mathfrak m R^{sh}$ 是 $R^{sh}$ 的极大理想；且
\item $\kappa^{sep} = R^{sh}/\mathfrak m R^{sh}$。
\end{enumerate}
\end{lemma}

\begin{proof}
证明与引理 \ref{algebra-lemma-henselization} 所用的证明完全相同。
唯一差别在于，不使用偶对，而使用三元组
$(S, \mathfrak q, \alpha)$，其中 $R \to S$ 平展，
$\mathfrak q$ 是 $S$ 中位于 $\mathfrak m$ 之上的素理想，
而 $\alpha : \kappa(\mathfrak q) \to \kappa^{sep}$
是 $\kappa$ 的扩张之间的嵌入。
\end{proof}

\begin{definition}
\label{algebra-definition-henselization}
设 $(R, \mathfrak m, \kappa)$ 是局部环。
\begin{enumerate}
\item 引理 \ref{algebra-lemma-henselization} 中构造的局部环映射
$R \to R^h$ 称为 $R$ 的{\it Hensel 化}。
\item 给定可分代数闭包 $\kappa \subset \kappa^{sep}$，
引理 \ref{algebra-lemma-strict-henselization} 中构造的局部环映射
$R \to R^{sh}$ 称为{\it $R$ 关于
$\kappa \subset \kappa^{sep}$ 的严格 Hensel 化}。
\item 若局部环映射 $R \to R^{sh}$ 同构于引理
\ref{algebra-lemma-strict-henselization} 中构造的某个局部环映射，
则称它为 $R$ 的一个{\it 严格 Hensel 化}
% P01-E0970: the frozen third definition item lacks terminal punctuation.
\end{enumerate}
\end{definition}

\noindent
映射 $R \to R^h \to R^{sh}$ 是平坦局部环同态。
由引理 \ref{algebra-lemma-uniqueness-henselian}，条件
“它们是 Hensel 局部环、是平展 $R$-代数的
滤过余极限，且剩余域分别为 $\kappa$ 与 $\kappa^{sep}$”
把 $R$-代数 $R^h$ 与 $R^{sh}$ 在唯一同构意义下确定。
本节余下部分主要讨论（严格）Hensel 化的函子性。
关于 $R$ 与其 Hensel 化之间关系的更细致结果，
将在《代数进阶》，第
\ref{more-algebra-section-permanence-henselization} 节讨论。

\begin{remark}
\label{algebra-remark-construct-sh-from-h}
也可以由 $R^h$ 构造 $R^{sh}$。事实上，对任意有限可分子扩张
$\kappa^{sep}/\kappa'/\kappa$，存在唯一的（在唯一同构意义下）
有限平展局部环扩张 $R^h \subset R^h(\kappa')$，
其剩余域扩张重现给定扩张；见引理
\ref{algebra-lemma-henselian-cat-finite-etale}。因此可以置
$$
R^{sh} =
\bigcup\nolimits_{\kappa \subset \kappa' \subset \kappa^{sep}}
R^h(\kappa')
$$
该系统中的箭头与剩余域层面上的箭头相容，并由引理
\ref{algebra-lemma-henselian-cat-finite-etale} 保证存在。
由引理 \ref{algebra-lemma-colimit-henselian}，这将产生 Hensel 局部环，
因为每个环 $R^h(\kappa')$ 由引理
\ref{algebra-lemma-finite-over-henselian} 都是 Hensel 的。
由构造，$R^h \to R^{sh}$ 诱导的剩余域扩张
就是域扩张 $\kappa^{sep}/\kappa$。所以这样构造的
$R^{sh}$ 是严格 Hensel 的。由引理
\ref{algebra-lemma-composition-colimit-etale}，$R$-代数 $R^{sh}$
是平展 $R$-代数的余极限。因此，引理
\ref{algebra-lemma-uniqueness-henselian} 的唯一性说明
$R^{sh}$ 就是严格 Hensel 化。
\end{remark}

\begin{lemma}
\label{algebra-lemma-henselian-functorial-prepare}
设 $R \to S$ 是局部环之间的局部映射。
设 $S \to S^h$ 是 Hensel 化。设 $R \to A$ 是平展环映射，
并设 $\mathfrak q$ 是 $A$ 中位于 $\mathfrak m_R$ 之上的素理想，
使得 $R/\mathfrak m_R \cong \kappa(\mathfrak q)$。
则存在唯一的环态 $f : A \to S^h$，使交换图表
$$
\xymatrix{
A \ar[r]_f & S^h \\
R \ar[u] \ar[r] & S \ar[u]
}
$$
成立，并且 $f^{-1}(\mathfrak m_{S^h}) = \mathfrak q$。
\end{lemma}

\begin{proof}
这是引理 \ref{algebra-lemma-map-into-henselian} 的特殊情形。
\end{proof}

\begin{lemma}
\label{algebra-lemma-henselian-functorial}
设 $R \to S$ 是局部环之间的局部映射。
设 $R \to R^h$ 与 $S \to S^h$ 是 Hensel 化。
存在唯一的局部环映射 $R^h \to S^h$，使交换图表
$$
\xymatrix{
R^h \ar[r]_f & S^h \\
R \ar[u] \ar[r] & S \ar[u]
}
$$
成立。
\end{lemma}

\begin{proof}
这立即由引理 \ref{algebra-lemma-map-into-henselian-colimit} 得出。
\end{proof}

\noindent
下面给出 Hensel 化的一个略有不同的构造。

\begin{lemma}
\label{algebra-lemma-henselization-different}
设 $R$ 是环，$\mathfrak p \subset R$ 是素理想。
考虑由偶对 $(S, \mathfrak q)$ 构成的范畴，其中
$R \to S$ 平展，而 $\mathfrak q$ 是位于 $\mathfrak p$
之上的素理想，使得
$\kappa(\mathfrak p) = \kappa(\mathfrak q)$。
这个范畴是滤过的，并且典范地有
$$
(R_{\mathfrak p})^h = \colim_{(S, \mathfrak q)} S
= \colim_{(S, \mathfrak q)} S_{\mathfrak q}
$$
\end{lemma}

\begin{proof}
偶对之间的态 $(S, \mathfrak q) \to (S', \mathfrak q')$
由 $R$-代数映射 $\varphi : S \to S'$ 给出，并且满足
$\varphi^{-1}(\mathfrak q') = \mathfrak q$。
我们来证明偶对范畴是滤过的；见《范畴》，定义
\ref{categories-definition-directed}。该范畴含有偶对
$(R, \mathfrak p)$，因而非空；这证明了《范畴》，定义
\ref{categories-definition-directed} 的 (1)。
假设 $(S, \mathfrak q)$ 与 $(S', \mathfrak q')$ 是两个偶对。
注意，$\mathfrak q$ 与 $\mathfrak q'$ 分别对应于纤维环
$S \otimes \kappa(\mathfrak p)$ 与
$S' \otimes \kappa(\mathfrak p)$ 的素理想，其剩余域为
$\kappa(\mathfrak p)$；因此它们分别对应于
$S \otimes \kappa(\mathfrak p)$ 与
$S' \otimes \kappa(\mathfrak p)$ 的极大理想。
置 $S'' = S \otimes_R S'$。由上可知，存在唯一的素理想
$\mathfrak q'' \subset S''$，它位于 $\mathfrak q$ 与
$\mathfrak q'$ 之上，并且剩余域为 $\kappa(\mathfrak p)$。
由引理 \ref{algebra-lemma-etale}，环映射 $R \to S''$ 平展。
这证明了《范畴》，定义 \ref{categories-definition-directed} 的 (2)。

其次，假设
$\varphi, \psi : (S, \mathfrak q) \to (S', \mathfrak q')$
是偶对的两个态。由引理 \ref{algebra-lemma-map-between-etale}，
$\varphi$、$\psi$ 以及 $S' \otimes_R S' \to S'$
都是平展环映射。考虑
$$
S'' = (S' \otimes_{\varphi, S, \psi} S')
\otimes_{S' \otimes_R S'} S'
$$
与 $S''$ 在 $\mathfrak p$ 上的纤维环
$$
F'' = (F' \otimes_{\varphi, F, \psi} F')
\otimes_{F' \otimes_{\kappa(\mathfrak p)} F'} F'
$$
其中 $F', F$ 分别是 $S'$、$S$ 的纤维环。
与上面同样论证（平展映射的基变换是平展的，平展映射的复合也是
平展的），可知 $S''$ 在 $R$ 上平展。
由于 $\varphi$ 与 $\psi$ 是偶对的态，对应于
$\mathfrak p'$ 的映射 $F' \to \kappa(\mathfrak p)$
延拓为映射 $F'' \to \kappa(\mathfrak p)$，进而对应于素理想
$\mathfrak q'' \subset S''$，其剩余域为
$\kappa(\mathfrak p)$。
% P01-E0971: the frozen source refers to an undefined prime p' here; the pair carries q', which is the prime corresponding to the residue map.
典范映射 $S' \to S''$（例如使用最右边的因子）是偶对的态
$(S', \mathfrak q') \to (S'', \mathfrak q'')$，
并使 $\varphi$ 与 $\psi$ 相等。这证明了《范畴》，定义
\ref{categories-definition-directed} 的 (3)。
因此该范畴是滤过的。

\medskip\noindent
回忆在引理 \ref{algebra-lemma-henselization} 的证明中，我们把
$(R_{\mathfrak p})^h$ 构造成相应的余极限，不过是从
$R_{\mathfrak p}$ 及其极大理想
$\mathfrak pR_{\mathfrak p}$ 出发。现在，对
$(R, \mathfrak p)$ 的任意偶对 $(S, \mathfrak q)$，
得到 $(R_{\mathfrak p}, \mathfrak pR_{\mathfrak p})$ 的偶对
$(S_{\mathfrak p}, \mathfrak qS_{\mathfrak p})$。
此外，在这种情况下
$$
S_{\mathfrak p} = \colim_{f \in R, f \not \in \mathfrak p} S_f.
$$
所以，为了证明引理中的等式，只须证明
$(R_{\mathfrak p}, \mathfrak pR_{\mathfrak p})$ 的任意偶对
$(S_{loc}, \mathfrak q_{loc})$ 都具有形式
$(S_{\mathfrak p}, \mathfrak qS_{\mathfrak p})$，
其中 $(S, \mathfrak q)$ 是
$(R, \mathfrak p)$ 上的某个偶对（略去一些细节）。
这由引理 \ref{algebra-lemma-etale} 得出。
\end{proof}

\begin{lemma}
\label{algebra-lemma-henselian-functorial-improve}
设 $R \to S$ 是环映射。设 $\mathfrak q \subset S$
是位于 $\mathfrak p \subset R$ 之上的素理想。
设 $R \to R^h$ 与 $S \to S^h$ 分别是
$R_\mathfrak p$ 与 $S_\mathfrak q$ 的 Hensel 化。
引理 \ref{algebra-lemma-henselian-functorial} 的局部环映射
$R^h \to S^h$ 把 $S^h$ 等同于
$R^h \otimes_R S$ 在位于 $\mathfrak m^h$ 与
$\mathfrak q$ 之上的唯一素理想处的 Hensel 化。
\end{lemma}

\begin{proof}
由引理 \ref{algebra-lemma-henselization-different}，可知
$R^h$ 与 $S^h$ 分别是平展 $R$-代数与平展 $S$-代数的滤过余极限。
所以 $R^h \otimes_R S$ 是平展 $S$-代数 $A_i$ 的滤过余极限
（引理 \ref{algebra-lemma-etale}）。由引理
\ref{algebra-lemma-colimits-of-etale}，可知 $S^h$ 是平展
$R^h \otimes_R S$-代数的滤过余极限。
又因为 $S^h$ 是 Hensel 局部环，且剩余域等于
$\kappa(\mathfrak q)$，所述结论由引理
\ref{algebra-lemma-uniqueness-henselian} 的唯一性结果得出。
\end{proof}

\begin{lemma}
\label{algebra-lemma-strictly-henselian-functorial-prepare}
设 $\varphi : R \to S$ 是局部环之间的局部映射。
设 $S/\mathfrak m_S \subset \kappa^{sep}$ 是可分代数闭包。
设 $S \to S^{sh}$ 是 $S$ 关于
$S/\mathfrak m_S \subset \kappa^{sep}$ 的严格 Hensel 化。
设 $R \to A$ 是平展环映射，并设 $\mathfrak q$
是 $A$ 中位于 $\mathfrak m_R$ 之上的素理想。
给定任意交换图表
$$
\xymatrix{
\kappa(\mathfrak q) \ar[r]_{\phi} & \kappa^{sep} \\
R/\mathfrak m_R \ar[r]^{\varphi} \ar[u] & S/\mathfrak m_S \ar[u]
}
$$
存在唯一的环态 $f : A \to S^{sh}$，使交换图表
$$
\xymatrix{
A \ar[r]_f & S^{sh} \\
R \ar[u] \ar[r]^{\varphi} & S \ar[u]
}
$$
成立，并且 $f^{-1}(\mathfrak m_{S^h}) = \mathfrak q$，
而所诱导的映射 $\kappa(\mathfrak q) \to \kappa^{sep}$
就是给定映射。
% P01-E0972: the frozen strict-Henselization statement takes the inverse image of m_{S^h}; the codomain is S^{sh}, so the ideal should be m_{S^{sh}}.
\end{lemma}

\begin{proof}
这是引理 \ref{algebra-lemma-map-into-henselian} 的特殊情形。
\end{proof}

\begin{lemma}
\label{algebra-lemma-strictly-henselian-functorial}
设 $R \to S$ 是局部环之间的局部映射。选取可分代数闭包
$R/\mathfrak m_R \subset \kappa_1^{sep}$
以及
$S/\mathfrak m_S \subset \kappa_2^{sep}$。
设 $R \to R^{sh}$ 与 $S \to S^{sh}$ 是相应的严格
Hensel 化。给定任意交换图表
$$
\xymatrix{
\kappa_1^{sep} \ar[r]_{\phi} & \kappa_2^{sep} \\
R/\mathfrak m_R \ar[r]^{\varphi} \ar[u] & S/\mathfrak m_S \ar[u]
}
$$
存在唯一的局部环映射 $R^{sh} \to S^{sh}$，使交换图表
$$
\xymatrix{
R^{sh} \ar[r]_f & S^{sh} \\
R \ar[u] \ar[r] & S \ar[u]
}
$$
成立，并在 $R^{sh}$ 与 $S^{sh}$ 的剩余域上诱导 $\phi$。
% P01-E0973: the frozen English sentence beginning “Given any commutative diagram” is followed by capitalized “There exists” without completing or punctuating the introductory clause.
\end{lemma}

\begin{proof}
这立即由引理 \ref{algebra-lemma-map-into-henselian-colimit} 得出。
\end{proof}

\begin{lemma}
\label{algebra-lemma-strict-henselization-different}
设 $R$ 是环，$\mathfrak p \subset R$ 是素理想。
设 $\kappa(\mathfrak p) \subset \kappa^{sep}$ 是可分代数闭包。
考虑由三元组 $(S, \mathfrak q, \phi)$ 构成的范畴，其中
$R \to S$ 平展，$\mathfrak q$ 是位于 $\mathfrak p$
之上的素理想，而
$\phi : \kappa(\mathfrak q) \to \kappa^{sep}$
是 $\kappa(\mathfrak p)$-代数映射。
这个范畴是滤过的，并且典范地有
$$
(R_{\mathfrak p})^{sh} =
\colim_{(S, \mathfrak q, \phi)} S =
\colim_{(S, \mathfrak q, \phi)} S_{\mathfrak q}
$$
\end{lemma}

\begin{proof}
三元组之间的态
$(S, \mathfrak q, \phi) \to (S', \mathfrak q', \phi')$
由 $R$-代数映射 $\varphi : S \to S'$ 给出，并满足
$\varphi^{-1}(\mathfrak q') = \mathfrak q$ 以及
$\phi' \circ \varphi = \phi$。
我们来证明偶对范畴是滤过的；见《范畴》，定义
\ref{categories-definition-directed}。
% P01-E0974: the frozen proof says “category of pairs” although the objects and morphisms in this lemma are triples.
该范畴含有三元组
$(R, \mathfrak p, \kappa(\mathfrak p) \subset \kappa^{sep})$，
因而非空；这证明了《范畴》，定义
\ref{categories-definition-directed} 的 (1)。
假设 $(S, \mathfrak q, \phi)$ 与
$(S', \mathfrak q', \phi')$ 是两个三元组。
注意，$\mathfrak q$ 与 $\mathfrak q'$ 分别对应于纤维环
$S \otimes \kappa(\mathfrak p)$ 与
$S' \otimes \kappa(\mathfrak p)$ 的素理想，其剩余域是
$\kappa(\mathfrak p)$ 的有限可分扩张；而 $\phi$ 与 $\phi'$
分别对应于到 $\kappa^{sep}$ 的映射。因此这些数据对应于
$\kappa(\mathfrak p)$-代数映射
$$
\phi : S \otimes_R \kappa(\mathfrak p) \longrightarrow \kappa^{sep},
\quad
\phi' : S' \otimes_R \kappa(\mathfrak p) \longrightarrow \kappa^{sep}.
$$
置 $S'' = S \otimes_R S'$。组合上述映射，得到唯一的
$\kappa(\mathfrak p)$-代数映射
% P01-E0975: the frozen source says “Combining the maps the above”, with the article and modifier in the wrong order.
$$
\phi'' = \phi \otimes \phi' :
S'' \otimes_R \kappa(\mathfrak p)
\longrightarrow
\kappa^{sep}
$$
其核对应于素理想 $\mathfrak q'' \subset S''$，
它位于 $\mathfrak q$ 与 $\mathfrak q'$ 之上；其剩余域经
$\phi''$ 映到 $\kappa^{sep}$ 中
$\phi(\kappa(\mathfrak q))$ 与
$\phi'(\kappa(\mathfrak q'))$ 的合成域。
由引理 \ref{algebra-lemma-etale}，环映射 $R \to S''$ 平展。
因此 $(S'', \mathfrak q'', \phi'')$ 是同时支配
$(S, \mathfrak q, \phi)$ 与 $(S', \mathfrak q', \phi')$
的三元组。这证明了《范畴》，定义
\ref{categories-definition-directed} 的 (2)。

其次，假设
$\varphi, \psi : (S, \mathfrak q, \phi) \to
(S', \mathfrak q', \phi')$ 是偶对的两个态。
% P01-E0976: the frozen source again says “morphisms of pairs” although these are morphisms of triples.
由引理 \ref{algebra-lemma-map-between-etale}，$\varphi$、$\psi$ 以及
$S' \otimes_R S' \to S'$ 都是平展环映射。考虑
$$
S'' = (S' \otimes_{\varphi, S, \psi} S')
\otimes_{S' \otimes_R S'} S'
$$
与 $S''$ 在 $\mathfrak p$ 上的纤维环
$$
F'' = (F' \otimes_{\varphi, F, \psi} F')
\otimes_{F' \otimes_{\kappa(\mathfrak p)} F'} F'
$$
其中 $F', F$ 是 $S'$、$S$ 的纤维环。
与上面同样论证（平展映射的基变换是平展的，平展映射的复合也是
平展的），可知 $S''$ 在 $R$ 上平展。
由于 $\varphi$ 与 $\psi$ 是三元组的态，映射
$\phi' : F' \to \kappa^{sep}$ 延拓为映射
$\phi'' : F'' \to \kappa^{sep}$，后者进而对应于素理想
$\mathfrak q'' \subset S''$。典范映射 $S' \to S''$
（例如使用最右边的因子）是三元组之间的态
$(S', \mathfrak q', \phi') \to
(S'', \mathfrak q'', \phi'')$，并使 $\varphi$ 与 $\psi$
相等。这证明了《范畴》，定义
\ref{categories-definition-directed} 的 (3)。
因此该范畴是滤过的。

\medskip\noindent
还须证明，该系统的余极限 $R_{colim}$ 等于
$R_{\mathfrak p}$ 关于 $\kappa^{sep}$ 的严格 Hensel 化。
为此，注意三元组 $(S, \mathfrak q, \phi)$ 的系统含有引理
\ref{algebra-lemma-henselization-different} 的偶对
$(S, \mathfrak q)$ 作为子系统。因此，由该引理的结果，
$R_{colim}$ 含有 $R_{\mathfrak p}^h$。
此外，很清楚，$R_{\mathfrak p}^h \subset R_{colim}$
是平展环扩张的有向余极限。由引理
\ref{algebra-lemma-finite-over-henselian} 与
\ref{algebra-lemma-colimit-henselian}，可知 $R_{colim}$ 是 Hensel 的。
最后，由引理 \ref{algebra-lemma-make-etale-map-prescribed-residue-field}，
可知 $R_{colim}$ 的剩余域等于 $\kappa^{sep}$。
因此 $R_{colim}$ 严格 Hensel，从而如所需，等于
$R_{\mathfrak p}$ 的严格 Hensel 化。略去一些细节。
\end{proof}

\begin{lemma}
\label{algebra-lemma-strictly-henselian-functorial-improve}
设 $R \to S$ 是环映射。设 $\mathfrak q \subset S$
是位于 $\mathfrak p \subset R$ 之上的素理想。
选取可分代数闭包
$\kappa(\mathfrak p) \subset \kappa_1^{sep}$
以及
$\kappa(\mathfrak q) \subset \kappa_2^{sep}$。
设 $R^{sh}$ 与 $S^{sh}$ 分别为
$R_\mathfrak p$ 与 $S_\mathfrak q$ 的相应严格 Hensel 化。
给定任意交换图表
$$
\xymatrix{
\kappa_1^{sep} \ar[r]_{\phi} & \kappa_2^{sep} \\
\kappa(\mathfrak p) \ar[r]^{\varphi} \ar[u] & \kappa(\mathfrak q) \ar[u]
}
$$
引理 \ref{algebra-lemma-strictly-henselian-functorial} 的局部环映射
$R^{sh} \to S^{sh}$ 把 $S^{sh}$ 等同于
$R^{sh} \otimes_R S$ 在一个位于 $\mathfrak q$ 以及极大理想
$\mathfrak m^{sh} \subset R^{sh}$ 之上的素理想处的严格
Hensel 化。
\end{lemma}

\begin{proof}
证明与引理 \ref{algebra-lemma-henselian-functorial-improve} 的证明相同，
只是使用引理 \ref{algebra-lemma-strict-henselization-different}，
而不是引理 \ref{algebra-lemma-henselization-different}。
\end{proof}

\begin{lemma}
\label{algebra-lemma-sh-from-h-map}
设 $R \to S$ 是环映射。设 $\mathfrak q \subset S$
是位于 $\mathfrak p \subset R$ 之上的素理想，使得
$\kappa(\mathfrak p) \to \kappa(\mathfrak q)$ 是同构。
选取 $\kappa(\mathfrak p) = \kappa(\mathfrak q)$ 的一个
可分代数闭包 $\kappa^{sep}$。则
$$
(S_\mathfrak q)^{sh} =
(S_\mathfrak q)^h \otimes_{(R_\mathfrak p)^h} (R_\mathfrak p)^{sh}
$$
\end{lemma}

\begin{proof}
这由备注 \ref{algebra-remark-construct-sh-from-h} 中局部环严格
Hensel 化的另一种构造以及剩余域相等这一事实得出。
略去一些细节。
\end{proof}

\section{Hensel 化与拟有限环映射}
\label{algebra-section-henselization-quasi-finite}

\noindent
本节证明关于拟有限环映射的（严格）Hensel 化之间函子性映射的一些结果。

\begin{lemma}
\label{algebra-lemma-quasi-finite-henselization}
设 $R \to S$ 是环映射。
设 $S$ 的素理想 $\mathfrak q$ 位于 $R$ 的素理想 $\mathfrak p$ 之上。
假设 $R \to S$ 在 $\mathfrak q$ 处是拟有限的。
引理 \ref{algebra-lemma-henselian-functorial} 的交换图表
$$
\xymatrix{
R_{\mathfrak p}^h \ar[r] & S_{\mathfrak q}^h \\
R_{\mathfrak p} \ar[u] \ar[r] & S_{\mathfrak q} \ar[u]
}
$$
把 $S_{\mathfrak q}^h$ 等同于
$R_{\mathfrak p}^h \otimes_{R_{\mathfrak p}} S_{\mathfrak q}$
在由 $\mathfrak q$ 生成的素理想处的局部化。此外，环映射
$R_{\mathfrak p}^h \to S_{\mathfrak q}^h$ 是有限的。
\end{lemma}

\begin{proof}
注意，$R_{\mathfrak p}^h \otimes_R S$ 在与 $\mathfrak q$
对应的素理想处是 $R_{\mathfrak p}^h$ 上拟有限的，见引理
\ref{algebra-lemma-four-rings}。因此，
$R_{\mathfrak p}^h \otimes_{R_{\mathfrak p}} S_{\mathfrak q}$
的局部化 $S'$ 是 Hensel 的，并且在 $R_{\mathfrak p}^h$
上有限，见引理 \ref{algebra-lemma-finite-over-henselian}。
作为一个局部化，$S'$ 是一族 'etale
$R_{\mathfrak p}^h \otimes_{R_{\mathfrak p}} S_{\mathfrak q}$-代数的
滤过余极限。由引理 \ref{algebra-lemma-henselian-functorial-improve}，
可知 $S_\mathfrak q^h$ 是
$R_{\mathfrak p}^h \otimes_{R_{\mathfrak p}} S_{\mathfrak q}$
的 Hensel 化。因此，由引理 \ref{algebra-lemma-uniqueness-henselian}
的唯一性结果可得 $S' = S_\mathfrak q^h$。
\end{proof}

\begin{lemma}
\label{algebra-lemma-quotient-henselization}
\begin{slogan}
Hensel 化与商相容。
\end{slogan}
设 $R$ 是局部环，其 Hensel 化为 $R^h$。
设 $I \subset \mathfrak m_R$。则 $R^h/IR^h$ 是 $R/I$ 的
Hensel 化。
\end{lemma}

\begin{proof}
这是引理 \ref{algebra-lemma-quasi-finite-henselization} 的一个特例。
\end{proof}

\begin{lemma}
\label{algebra-lemma-quasi-finite-strict-henselization}
设 $R \to S$ 是环映射。
设 $S$ 的素理想 $\mathfrak q$ 位于 $R$ 的素理想 $\mathfrak p$ 之上。
假设 $R \to S$ 在 $\mathfrak q$ 处是拟有限的。
设 $\kappa_2^{sep}/\kappa(\mathfrak q)$ 是可分代数闭包，并以
$\kappa_1^{sep} \subset \kappa_2^{sep}$ 表示其中在
$\kappa(\mathfrak p)$ 上可分代数的元素所成的子域
（域，引理 \ref{fields-lemma-separable-first}）。
引理 \ref{algebra-lemma-strictly-henselian-functorial} 的交换图表
$$
\xymatrix{
R_{\mathfrak p}^{sh} \ar[r] & S_{\mathfrak q}^{sh} \\
R_{\mathfrak p} \ar[u] \ar[r] & S_{\mathfrak q} \ar[u]
}
$$
把 $S_{\mathfrak q}^{sh}$ 等同于
$R_{\mathfrak p}^{sh} \otimes_{R_{\mathfrak p}} S_{\mathfrak q}$
在下列映射的核所给出的素理想处的局部化：
$$
R_{\mathfrak p}^{sh} \otimes_{R_{\mathfrak p}} S_{\mathfrak q}
\longrightarrow
\kappa_1^{sep} \otimes_{\kappa(\mathfrak p)} \kappa(\mathfrak q)
\longrightarrow
\kappa_2^{sep}
$$
此外，环映射 $R_{\mathfrak p}^{sh} \to S_{\mathfrak q}^{sh}$
是局部环之间的有限局部同态，其剩余域扩张为
$\kappa_2^{sep}/\kappa_1^{sep}$；该扩张既有限又纯不可分。
\end{lemma}

\begin{proof}
由于 $R \to S$ 在 $\mathfrak q$ 处拟有限，由定义
\ref{algebra-definition-quasi-finite} 和引理
\ref{algebra-lemma-isolated-point-fibre} 可知扩张
$\kappa(\mathfrak q)/\kappa(\mathfrak p)$ 是有限的。
因此，$\kappa_1^{sep}$ 是 $\kappa(\mathfrak p)$ 的一个可分代数闭包
（略去一个小细节）。特别地，引理
\ref{algebra-lemma-strictly-henselian-functorial} 确实适用。
其次，$\kappa(\mathfrak q)$ 与 $\kappa_1^{sep}$ 在
$\kappa_2^{sep}$ 中的合成域是可分代数闭的，因而等于
$\kappa_2^{sep}$。由此可知
$\kappa_2^{sep}/\kappa_1^{sep}$ 是有限的。按构造，扩张
$\kappa_2^{sep}/\kappa_1^{sep}$ 是纯不可分的。
环映射 $R_{\mathfrak p}^{sh} \to S_{\mathfrak q}^{sh}$
确为局部的，并且诱导确为有限纯不可分的剩余域扩张
$\kappa_2^{sep}/\kappa_1^{sep}$。

\medskip\noindent
注意，$R_{\mathfrak p}^{sh} \otimes_R S$ 在引理陈述中给出的素理想
$\mathfrak q'$ 处是 $R_{\mathfrak p}^{sh}$ 上拟有限的，见引理
\ref{algebra-lemma-four-rings}。因此，
$R_{\mathfrak p}^{sh} \otimes_{R_{\mathfrak p}} S_{\mathfrak q}$
在 $\mathfrak q'$ 处的局部化 $S'$ 是 Hensel 的，并且在
$R_{\mathfrak p}^{sh}$ 上有限，见引理
\ref{algebra-lemma-finite-over-henselian}。
注意，$S'$ 的剩余域是 $\kappa_2^{sep}$，因为由上一段的讨论，映射
$\kappa_1^{sep} \otimes_{\kappa(\mathfrak p)} \kappa(\mathfrak q)
\to \kappa_2^{sep}$ 是满射。
此外，作为一个局部化，$S'$ 是一族 'etale
$R_{\mathfrak p}^{sh} \otimes_{R_{\mathfrak p}} S_{\mathfrak q}$-代数的
滤过余极限。由引理
\ref{algebra-lemma-strictly-henselian-functorial-improve}，可知
$S_{\mathfrak q}^{sh}$ 是
$R_{\mathfrak p}^{sh} \otimes_{R_{\mathfrak p}} S_{\mathfrak q}$
在 $\mathfrak q'$ 处的严格 Hensel 化。因此，由引理
\ref{algebra-lemma-uniqueness-henselian} 的唯一性结果可得
$S' = S_\mathfrak q^{sh}$。
\end{proof}

\begin{lemma}
\label{algebra-lemma-quotient-strict-henselization}
设 $R$ 是局部环，其严格 Hensel 化为 $R^{sh}$。
设 $I \subset \mathfrak m_R$。则 $R^{sh}/IR^{sh}$ 是 $R/I$
的一个严格 Hensel 化。
\end{lemma}

\begin{proof}
这是引理 \ref{algebra-lemma-quasi-finite-strict-henselization} 的一个特例。
\end{proof}

\begin{lemma}
\label{algebra-lemma-local-tensor-with-integral}
设 $A \to B$ 和 $A \to C$ 是局部环之间的局部同态。
如果 $A \to C$ 是整的，并且
$\kappa(\mathfrak m_C)/\kappa(\mathfrak m_A)$ 或
$\kappa(\mathfrak m_B)/\kappa(\mathfrak m_A)$ 是纯不可分的，
则 $D = B \otimes_A C$ 是局部环，且 $B \to D$ 和
$C \to D$ 都是局部的。
\end{lemma}

\begin{proof}
由整环映射 $B \to D$ 的上升定理（引理
\ref{algebra-lemma-integral-going-up}），$D$ 的任意极大理想均位于
$B$ 的极大理想之上。现在
$D/\mathfrak m_B D = \kappa(\mathfrak m_B) \otimes_A C =
\kappa(\mathfrak m_B) \otimes_{\kappa(\mathfrak m_A)} C/\mathfrak m_A C$。
$C/\mathfrak m_A C$ 的谱只含一个点，即 $\mathfrak m_C$。
因此，$D/\mathfrak m_B D$ 的谱与
$\kappa(\mathfrak m_B) \otimes_{\kappa(\mathfrak m_A)}
\kappa(\mathfrak m_C)$ 的谱相同。根据我们的假设，
$\kappa(\mathfrak m_C)/\kappa(\mathfrak m_A)$ 或
$\kappa(\mathfrak m_B)/\kappa(\mathfrak m_A)$ 是纯不可分的，
故后一个谱只含一个点。这就证明了 $D$ 是局部的，并且环映射
$B \to D$ 和 $C \to D$ 都是局部的。
\end{proof}

\begin{lemma}
\label{algebra-lemma-base-change-strict-henselization-quasi-finite}
设 $A \to B$ 和 $A \to C$ 是环映射。设 $\kappa$ 是可分代数闭域，
并设 $B \otimes_A C \to \kappa$ 是环同态。以下列图表表示严格
Hensel 化之间相应的映射（见证明）：
$$
\xymatrix{
B^{sh} \ar[r] & (B \otimes_A C)^{sh} \\
A^{sh} \ar[u] \ar[r] & C^{sh} \ar[u]
}
$$
如果
\begin{enumerate}
\item $A \to B$ 在素理想
$\mathfrak p_B = \Ker(B \to \kappa)$ 处是拟有限的；或
\item $B$ 是拟有限 $A$-代数的滤过余极限；或
\item $B_{\mathfrak p_B}$ 是 $A_{\mathfrak p_A}$ 上拟有限代数的
滤过余极限；或
\item $B$ 在 $A$ 上整，
\end{enumerate}
则 $B^{sh} \otimes_{A^{sh}} C^{sh} \to (B \otimes_A C)^{sh}$
是同构。
\end{lemma}

\begin{proof}
记 $D = B \otimes_A C$。记
$\mathfrak p_A = \Ker(A \to \kappa)$，并类似地定义
$\mathfrak p_B$、$\mathfrak p_C$ 和 $\mathfrak p_D$。
以 $\kappa_A \subset \kappa$ 表示 $\kappa(\mathfrak p_A)$
在 $\kappa$ 中的可分代数闭包，并类似地定义
$\kappa_B$、$\kappa_C$ 和 $\kappa_D$。
以 $A^{sh}$ 表示利用可分代数闭包
$\kappa_A/\kappa(\mathfrak p_A)$ 构造的
$A_{\mathfrak p_A}$ 的严格 Hensel 化。类似地定义
$B^{sh}$、$C^{sh}$ 和 $D^{sh}$。由引理
\ref{algebra-lemma-strictly-henselian-functorial} 的函子性，我们得到引理中的
交换图表。

\medskip\noindent
考虑由该交换图表得到的映射
$$
c : B^{sh} \otimes_{A^{sh}} C^{sh} \to D^{sh} = (B \otimes_A C)^{sh}
$$
如果 $A \to B$ 在
$\mathfrak p_B = \Ker(B \to \kappa)$ 处拟有限，则由引理
\ref{algebra-lemma-four-rings}，环映射 $C \to D$ 在 $\mathfrak p_D$
处拟有限。因此，由引理
\ref{algebra-lemma-quasi-finite-strict-henselization}
（以及引理 \ref{algebra-lemma-base-change-integral}），环映射 $c$
是有限 $C^{sh}$-代数之间的同态，并且
$$
B^{sh} = (B \otimes_A A^{sh})_{\mathfrak q}
\quad\text{且}\quad
D^{sh} = (D \otimes_C C^{sh})_{\mathfrak r} =
(B \otimes_A C^{sh})_{\mathfrak r}
$$
对某些素理想 $\mathfrak q$ 和 $\mathfrak r$ 成立。由于
$$
B^{sh} \otimes_{A^{sh}} C^{sh} =
(B \otimes_A A^{sh})_{\mathfrak q} \otimes_{A^{sh}} C^{sh} =
\text{a localization of }
B \otimes_A C^{sh}
$$
可知 $c$ 的源和靶都是 $B \otimes_A C^{sh}$ 的局部化
（且与该映射相容）。因此，只需证明
$B^{sh} \otimes_{A^{sh}} C^{sh}$ 是局部的（略去一个小细节）。
这由引理 \ref{algebra-lemma-local-tensor-with-integral} 以及下述事实得出：
根据上面已经使用过的引理
\ref{algebra-lemma-quasi-finite-strict-henselization}，
$A^{sh} \to B^{sh}$ 是有限的，其剩余域扩张纯不可分。
这证明了引理的情形 (1)。

\medskip\noindent
在情形 (2) 中，把 $B = \colim B_i$ 写成拟有限 $A$-代数的
滤过余极限。相应地得到 $D = \colim D_i$，其中
$D_i = B_i \otimes_A C$。注意 $B^{sh} = \colim B_i^{sh}$。
事实上，由引理 \ref{algebra-lemma-colimit-henselian}，环
$\colim B_i^{sh}$ 是严格 Hensel 局部环。
由引理 \ref{algebra-lemma-colimit-colimit-etale-better}，
$\colim B_i^{sh}$ 也是一族 'etale $B$-代数的滤过余极限。
最后，$\colim B_i^{sh}$ 的剩余域是
$\kappa(\mathfrak p_B)$ 的一个可分代数闭包（略去细节）。
因此，由定义 \ref{algebra-definition-henselization} 后的讨论可得
$B^{sh} = \colim B_i^{sh}$。类似地，
$D^{sh} = \colim D_i^{sh}$。于是由情形 (1) 可得
$$
D^{sh} = \colim D_i^{sh} = \colim B_i^{sh} \otimes_{A^{sh}} C^{sh} =
B^{sh} \otimes_{A^{sh}} C^{sh}
$$
% P01-E0977: frozen English source has subject--verb disagreement in “filtered colimit commute”.
因为滤过余极限与张量积可交换。

\medskip\noindent
情形 (3)。可以分别在 $\mathfrak p_A$、$\mathfrak p_B$ 和
$\mathfrak p_C$ 处用 $A$、$B$、$C$ 的局部化替换它们。
因此，(3) 由 (2) 得出。

\medskip\noindent
由于整环映射是有限环映射的滤过余极限，(4) 也由 (2) 得出。
\end{proof}

\section{Serre 正规性判据}
\label{algebra-section-serre-criterion}

\noindent
我们引入诺特环的下列性质。

\begin{definition}
\label{algebra-definition-conditions}
设 $R$ 是诺特环，$k \geq 0$ 是整数。
\begin{enumerate}
\item 如果对每个高度 $\leq k$ 的素理想 $\mathfrak p$，局部环
$R_{\mathfrak p}$ 都是正则的，则称 $R$ 具有性质
{\it $(R_k)$}。我们也称 $R$ 在余维 $\leq k$ 处{\it 正则}。
\item 如果对每个素理想 $\mathfrak p$，局部环
$R_{\mathfrak p}$ 的深度至少为
$\min\{k, \dim(R_{\mathfrak p})\}$，则称 $R$ 具有性质
{\it $(S_k)$}。
\item 设 $M$ 是有限 $R$-模。如果对每个素理想
$\mathfrak p$，模 $M_{\mathfrak p}$ 的深度至少为
$\min\{k, \dim(\text{Supp}(M_{\mathfrak p}))\}$，
则称 $M$ 具有性质 $(S_k)$。
\end{enumerate}
\end{definition}

\noindent
任意诺特环都具有性质 $(S_0)$，其上的任意有限模也如此。
我们的约定是零模的深度为 $\infty$（见第
\ref{algebra-section-depth} 节），而空集的维数为 $-\infty$
（见拓扑，第 \ref{topology-section-krull-dimension} 节）；
这些约定保证零模对所有 $k$ 都具有性质 $(S_k)$。

\begin{lemma}
\label{algebra-lemma-criterion-no-embedded-primes}
设 $R$ 是诺特环，$M$ 是有限 $R$-模。下列条件等价：
\begin{enumerate}
\item $M$ 没有嵌入相伴素理想；
\item $M$ 具有性质 $(S_1)$。
\end{enumerate}
\end{lemma}

\begin{proof}
设 $\mathfrak p$ 是 $M$ 的一个嵌入相伴素理想。
则 $M$ 还有另一个相伴素理想 $\mathfrak q$，使得
$\mathfrak p \supset \mathfrak q$。特别地，这意味着
$\dim(\text{Supp}(M_{\mathfrak p})) \geq 1$
（因为 $\mathfrak q$ 也在支集中）。另一方面，
$\mathfrak pR_{\mathfrak p}$ 是 $M_{\mathfrak p}$ 的相伴素理想
（引理 \ref{algebra-lemma-associated-primes-localize}），因而
$\text{depth}(M_{\mathfrak p}) = 0$
（见引理 \ref{algebra-lemma-ideal-nonzerodivisor}）。
换言之，$(S_1)$ 不成立。
反过来，如果 $(S_1)$ 不成立，则存在素理想 $\mathfrak p$，使得
$\dim(\text{Supp}(M_{\mathfrak p})) \geq 1$ 且
$\text{depth}(M_{\mathfrak p}) = 0$。由于
$\text{depth}(M_{\mathfrak p}) = 0$，由引理
\ref{algebra-lemma-associated-primes-localize} 和
\ref{algebra-lemma-ideal-nonzerodivisor} 可知
$\mathfrak p \in \text{Ass}(M)$。
由于 $\dim(\text{Supp}(M_{\mathfrak p})) \geq 1$，
存在 $\mathfrak q \in \text{Supp}(M)$，使得
$\mathfrak q \subset \mathfrak p$ 且
$\mathfrak q \not = \mathfrak p$。可以选取这样的
$\mathfrak q$，使其在 $\text{Supp}(M)$ 中极小。于是由命题
\ref{algebra-proposition-minimal-primes-associated-primes}，有
$\mathfrak q \in \text{Ass}(M)$，因此 $\mathfrak p$
是嵌入相伴素理想。
\end{proof}

\begin{lemma}
\label{algebra-lemma-criterion-reduced}
\begin{slogan}
约化等价于 R0 加 S1。
\end{slogan}
设 $R$ 是诺特环。下列条件等价：
\begin{enumerate}
\item $R$ 是约化的；
\item $R$ 具有性质 $(R_0)$ 和 $(S_1)$。
\end{enumerate}
\end{lemma}

\begin{proof}
假设 $R$ 是约化的。根据引理
\ref{algebra-lemma-minimal-prime-reduced-ring}，对 $R$ 的每个极小素理想
$\mathfrak p$，$R_{\mathfrak p}$ 都是域。因此，$(R_0)$ 成立。
设 $\mathfrak p$ 是高度 $\geq 1$ 的素理想。则
$A = R_{\mathfrak p}$ 是维数 $\geq 1$ 的约化局部环。
因此，它的极大理想 $\mathfrak m$ 不是相伴素理想，因为否则将存在
零化子为 $\mathfrak m$ 的 $x \in \mathfrak m$，从而
$x^2 = 0$。于是由引理 \ref{algebra-lemma-ass-zero-divisors}，
$A = R_{\mathfrak p}$ 的深度至少为一。这表明 $(S_1)$ 成立。

\medskip\noindent
反过来，假设 $R$ 满足 $(R_0)$ 和 $(S_1)$。
如果 $\mathfrak p$ 是 $R$ 的极小素理想，则由 $(R_0)$，
$R_{\mathfrak p}$ 是域，因而是约化的。
如果 $\mathfrak p$ 不是极小的，则由 $(S_1)$ 可知
$R_{\mathfrak p}$ 的深度 $\geq 1$，因而存在元素
$t \in \mathfrak pR_{\mathfrak p}$，使得
$R_{\mathfrak p} \to R_{\mathfrak p}[1/t]$ 是单射。
由引理 \ref{algebra-lemma-characterize-zero-local}，
$R_\mathfrak p[1/t]$ 包含于它在各素理想处局部化所得环的乘积中。
这意味着 $R_{\mathfrak p}$ 是下列各局部化之乘积的子环：
在满足 $\mathfrak p \supset \mathfrak q$ 且
$t \not \in \mathfrak q$ 的素理想处对 $R$ 作局部化。
这些素理想的高度更小，故对高度作归纳可知 $R$ 是约化的。
\end{proof}

\begin{lemma}[Serre 正规性判据]
\label{algebra-lemma-criterion-normal}
\begin{reference}
\cite[IV, Theorem 5.8.6]{EGA}
\end{reference}
\begin{slogan}
正规等价于 R1 加 S2。
\end{slogan}
设 $R$ 是诺特环。下列条件等价：
\begin{enumerate}
\item $R$ 是正规环；
\item $R$ 具有性质 $(R_1)$ 和 $(S_2)$。
\end{enumerate}
\end{lemma}

\begin{proof}
证明 (1) $\Rightarrow$ (2)。假设 $R$ 正规，即在各素理想处的所有
局部化 $R_{\mathfrak p}$ 都是正规整环。特别地，由引理
\ref{algebra-lemma-criterion-reduced} 可知 $R$ 具有 $(R_0)$ 和
$(S_1)$。因此，只需证明维数为 $d$ 的局部诺特正规整环 $R$
的深度 $\geq \min(2, d)$，并且当 $d = 1$ 时它是正则的。
当 $d = 1$ 时，该断言由引理 \ref{algebra-lemma-characterize-dvr} 得出。

\medskip\noindent
设 $R$ 是极大理想为 $\mathfrak m$、维数为 $d \geq 2$ 的局部
诺特正规整环。对 $R$ 应用引理 \ref{algebra-lemma-hart-serre-loc-thm}。
显然，$R$ 不属于该引理的情形 (1) 或 (2)。
设 $R \to R'$ 如该引理的 (4) 所述。
由于 $R$ 是整环，有 $R \subset R'$。由于 $\mathfrak m$
不是 $R'$ 的相伴素理想，存在 $x \in \mathfrak m$，它在
$R'$ 上是非零因子。于是 $R_x = R'_x$，所以 $R$ 和 $R'$
是具有相同分式域的整环。然而，$R \subset R'$ 的有限性意味着
$R'$ 的每个元素都整于 $R$（引理
\ref{algebra-lemma-finite-is-integral}）；由于 $R$ 正规，故
$R = R'$。这意味着 (4) 不会发生。因而只剩情形 (3)，即
$\text{depth}(R) \geq 2$，如所需。

\medskip\noindent
证明 (2) $\Rightarrow$ (1)。假设 $R$ 满足 $(R_1)$ 和
$(S_2)$。由引理 \ref{algebra-lemma-criterion-reduced} 可知 $R$ 是约化的。
因此，只需证明：如果 $R$ 是维数为 $d$、满足 $(S_2)$ 和
$(R_1)$ 的约化局部诺特环，则 $R$ 是正规整环。
如果 $d = 0$，结论显然。如果 $d = 1$，则结论由引理
\ref{algebra-lemma-characterize-dvr} 得出。

\medskip\noindent
设 $R$ 是极大理想为 $\mathfrak m$、维数为 $d \geq 2$
并且满足 $(R_1)$ 和 $(S_2)$ 的约化局部诺特环。
由引理 \ref{algebra-lemma-characterize-reduced-ring-normal}，只需证明
$R$ 在其全分式环 $Q(R)$ 中整闭。选取整于 $R$ 的
$x \in Q(R)$。则 $R' = R[x]$ 是 $R$ 的有限环扩张
（引理 \ref{algebra-lemma-characterize-finite-in-terms-of-integral}）。
由于对每个非极大素理想 $\mathfrak p \subset R$ 都有
$\dim(R_\mathfrak p) < d$，归纳可得
$R_\mathfrak p = R'_\mathfrak p$。因此，$R'/R$ 的支集为
$\{\mathfrak m\}$。由此可知 $R'/R$ 被 $\mathfrak m$
的某个幂零化（引理
\ref{algebra-lemma-Noetherian-power-ideal-kills-module}）。
由引理 \ref{algebra-lemma-hart-serre-loc-thm}，这与 $R$ 的深度
$\geq 2 = \min(2, d)$ 的假设矛盾，证明完毕。
\end{proof}

\begin{lemma}
\label{algebra-lemma-regular-normal}
正则环是正规的。
\end{lemma}

\begin{proof}
设 $R$ 是正则环。由引理 \ref{algebra-lemma-criterion-normal}，
只需证明 $R$ 满足 $(R_1)$ 和 $(S_2)$。
正则局部环是 Cohen-Macaulay 的，见引理
\ref{algebra-lemma-regular-ring-CM}，故 $R$ 显然满足 $(S_2)$。
性质 $(R_1)$ 立即成立。
\end{proof}

\begin{lemma}
\label{algebra-lemma-normal-domain-intersection-localizations-height-1}
设 $R$ 是分式域为 $K$ 的诺特正规整环。则
\begin{enumerate}
\item 对任意非零 $a \in R$，商 $R/aR$ 没有嵌入素理想，
且它的所有相伴素理想都具有高度 $1$；
% P01-E0978: frozen English item (1) has no terminal punctuation before the next item.
\item
$$
R = \bigcap\nolimits_{\text{高度}(\mathfrak p) = 1} R_{\mathfrak p}
$$
\item 对任意非零 $x \in K$，商 $R/(R \cap xR)$
没有嵌入素理想，且它的所有相伴素理想都具有高度 $1$。
% P01-E0979: frozen English has the malformed phrase “associates primes”.
\end{enumerate}
\end{lemma}

\begin{proof}
由引理 \ref{algebra-lemma-criterion-normal} 可知 $R$ 具有 $(S_2)$。
因此，对任意非零元素 $a \in R$，$R/aR$ 都具有 $(S_1)$
（例如使用引理 \ref{algebra-lemma-depth-in-ses}）。
% P01-E0980: frozen English omits sentence punctuation after the parenthetical citation.
所以 $R/aR$ 没有嵌入素理想
（引理 \ref{algebra-lemma-criterion-no-embedded-primes}）。
由此可知，$R/aR$ 的相伴素理想恰为 $(a)$ 上方的极小素理想
$\mathfrak p$；由于 $a$ 非零，它们的高度为 $1$
（引理 \ref{algebra-lemma-minimal-over-1}）。这证明了 (1)。

\medskip\noindent
因此，给定 $b \in R$，$b \in aR$ 当且仅当对 $(a)$ 上方的
每个极小素理想 $\mathfrak p$ 都有
$b \in aR_{\mathfrak p}$（见引理
\ref{algebra-lemma-zero-at-ass-zero}）。如上所见，这些素理想的高度均为
$1$，故 $b/a \in R$ 当且仅当对所有高度为 1 的素理想都有
$b/a \in R_{\mathfrak p}$。因此 (2) 成立。

\medskip\noindent
为证明 (3)，写 $x = a/b$。设
$\mathfrak p_1, \ldots, \mathfrak p_r$ 是 $(ab)$ 上方的极小素理想。
由上可知，它们都具有高度 1。于是由本引理的 (2) 可知
$R \cap xR = \bigcap_{i = 1, \ldots, r} (R \cap xR_{\mathfrak p_i})$。
因此，$R/(R \cap xR)$ 是
$\bigoplus R/(R \cap xR_{\mathfrak p_i})$ 的子模。
由于 $R_{\mathfrak p_i}$ 是离散赋值环（这是诺特正规整环
$R$ 的性质 $(R_1)$ 所致，见引理
\ref{algebra-lemma-criterion-normal}），对某个 $e_i \in \mathbf{Z}$ 有
$xR_{\mathfrak p_i} = \mathfrak p_i^{e_i}R_{\mathfrak p_i}$。
因此，该直和等于
$\bigoplus_{e_i > 0} R/\mathfrak p_i^{(e_i)}$，见定义
\ref{algebra-definition-symbolic-power}。
由引理 \ref{algebra-lemma-symbolic-power-associated}，模
$R/\mathfrak p^{(n)}$ 的唯一相伴素理想是 $\mathfrak p$。
因此，$R/(R \cap xR)$ 的相伴素理想集合是
$\{\mathfrak p_i\}$ 的子集，且它们之间不存在包含关系。
% P01-E0981: frozen English has “the set of associate primes” instead of “associated primes”.
这证明了 (3)。
\end{proof}

\section{域的形式光滑性}
\label{algebra-section-p-bases}

\noindent
本节将证明，域扩张形式光滑当且仅当它可分。不过，我们先证明：
有限生成域扩张是可分代数扩张，当且仅当它形式非分歧。

\begin{lemma}
\label{algebra-lemma-characterize-separable-algebraic-field-extensions}
设 $K/k$ 是有限生成域扩张。下列条件等价：
\begin{enumerate}
\item $K$ 是 $k$ 的有限可分域扩张；
\item $\Omega_{K/k} = 0$；
\item $K$ 在 $k$ 上形式非分歧；
\item $K$ 在 $k$ 上非分歧；
\item $K$ 在 $k$ 上形式平展；
\item $K$ 在 $k$ 上平展。
\end{enumerate}
\end{lemma}

\begin{proof}
(2) 与 (3) 的等价性是引理
\ref{algebra-lemma-characterize-formally-unramified}。
由引理 \ref{algebra-lemma-etale-over-field} 可知 (1) 等价于 (6)。
性质 (6) 蕴含 (5) 和 (4)，而后二者又都蕴含 (3)
（引理 \ref{algebra-lemma-formally-etale-etale}、
\ref{algebra-lemma-unramified} 和
\ref{algebra-lemma-formally-unramified-unramified}）。
因此，只需证明 (2) 蕴含 (1)。
选取有限生成 $k$-子代数 $A \subset K$，使得 $K$ 是整环
$A$ 的分式域。令 $S = A \setminus \{0\}$。
由于 $0 = \Omega_{K/k} = S^{-1}\Omega_{A/k}$
（引理 \ref{algebra-lemma-differentials-localize}），并且
$\Omega_{A/k}$ 是有限生成的（引理
\ref{algebra-lemma-differentials-finitely-generated}），可以用某个局部化
$A_f$ 替换 $A$，从而归约到 $\Omega_{A/k} = 0$ 的情形
（略去细节）。于是 $A$ 在 $k$ 上非分歧，故 $K/k$
有限可分；例如可对 $\mathfrak q = (0)$ 应用引理
\ref{algebra-lemma-unramified-at-prime}。
\end{proof}

\begin{lemma}
\label{algebra-lemma-derivative-zero-pth-power}
设 $k$ 是特征 $p > 0$ 的完美域，$K/k$ 是扩张，且
$a \in K$。则 $\text{d}a = 0$ 在 $\Omega_{K/k}$ 中成立，
当且仅当 $a$ 是一个 $p$ 次幂。
\end{lemma}

\begin{proof}
由引理 \ref{algebra-lemma-colimit-differentials} 可知，存在子域
$k \subset L \subset K$，使得 $L/k$ 是有限生成域扩张，
且 $\text{d}a$ 在 $\Omega_{L/k}$ 中为零。
因此，可以假设 $K$ 是 $k$ 的有限生成域扩张。

\medskip\noindent
选取超越基 $x_1, \ldots, x_r \in K$，使得 $K$ 是
$k(x_1, \ldots, x_r)$ 的有限可分扩张。根据定义，这是可能的，
见定义 \ref{algebra-definition-perfect} 和
\ref{algebra-definition-separable-field-extension}。
我们指出，该结果对纯超越子域
$k(x_1, \ldots, x_r) \subset K$ 成立。事实上，
$$
\Omega_{k(x_1, \ldots, x_r)/k} =
\bigoplus\nolimits_{i = 1}^r k(x_1, \ldots, x_r) \text{d}x_i
$$
并且偏导数全为零的任意有理函数都是 $p$ 次幂。此外，还有
$$
\Omega_{K/k} =
\bigoplus\nolimits_{i = 1}^r K\text{d}x_i
$$
因为 $k(x_1, \ldots, x_r) \subset K$ 有限可分
（略去计算）。假设 $a \in K$ 且 $\text{d}a = 0$
在微分模中成立。根据 $x_i$ 的选取，$a$ 的最小多项式
$P(T) \in k(x_1, \ldots, x_r)[T]$ 是可分的。写成
$$
P(T) = T^d + \sum\nolimits_{i = 1}^d a_i T^{d - i}
$$
从而
$$
0 = \text{d}P(a) = \sum\nolimits_{i = 1}^d a^{d - i}\text{d}a_i
$$
在 $\Omega_{K/k}$ 中成立。由上述对 $\Omega_{K/k}$ 的描述以及
$P$ 是 $a$ 的最小多项式这一事实，可知这蕴含
$\text{d}a_i = 0$。因此，对每个 $i$ 都有
$a_i = b_i^p$。于是由域，引理
\ref{fields-lemma-pth-root} 可知 $a$ 是 $p$ 次幂。
\end{proof}

\begin{lemma}
\label{algebra-lemma-size-extension-pth-roots}
设 $k$ 是特征 $p > 0$ 的域。设
$a_1, \ldots, a_n \in k$，且
$\text{d}a_1, \ldots, \text{d}a_n$ 在
$\Omega_{k/\mathbf{F}_p}$ 中线性无关。则域扩张
$k(a_1^{1/p}, \ldots, a_n^{1/p})$ 在 $k$ 上的次数为
$p^n$。
\end{lemma}

\begin{proof}
对 $n$ 归纳。当 $n = 1$ 时，结论是引理
\ref{algebra-lemma-derivative-zero-pth-power}。
在归纳步骤中，假设
$k(a_1^{1/p}, \ldots, a_{n - 1}^{1/p})$ 在 $k$ 上的次数为
$p^{n - 1}$。需要证明 $a_n$ 在
$k(a_1^{1/p}, \ldots, a_{n - 1}^{1/p})$ 中的像不是 $p$ 次幂。
如果它是，则可写成
\begin{align*}
a_n & =
\left(\sum\nolimits_{I = (i_1, \ldots, i_{n - 1}),\ 0 \leq i_j \leq p - 1}
\lambda_I a_1^{i_1/p} \ldots a_{n - 1}^{i_{n - 1}/p}\right)^p \\
& = \sum\nolimits_{I = (i_1, \ldots, i_{n - 1}),\ 0 \leq i_j \leq p - 1}
\lambda_I^p a_1^{i_1} \ldots a_{n - 1}^{i_{n - 1}}
\end{align*}
应用 $\text{d}$，可知 $\text{d}a_n$ 线性依赖于
$\text{d}a_i$，$i < n$。这产生矛盾。
\end{proof}

\begin{lemma}
\label{algebra-lemma-separable-differentials}
设 $k$ 是特征 $p > 0$ 的域。下列条件等价：
\begin{enumerate}
\item 域扩张 $K/k$ 可分
（见定义 \ref{algebra-definition-separable-field-extension}）；
\item 映射
$K \otimes_k \Omega_{k/\mathbf{F}_p} \to \Omega_{K/\mathbf{F}_p}$
是单射。
\end{enumerate}
\end{lemma}

\begin{proof}
把 $K$ 写成有限生成域扩张 $K_i/k$ 的有向余极限
$K = \colim_i K_i$。按定义，$K$ 可分当且仅当每个
$K_i$ 都在 $k$ 上可分；而由引理
\ref{algebra-lemma-colimit-differentials} 可知，
$K \otimes_k \Omega_{k/\mathbf{F}_p} \to \Omega_{K/\mathbf{F}_p}$
是单射，当且仅当每个
$K_i \otimes_k \Omega_{k/\mathbf{F}_p} \to
\Omega_{K_i/\mathbf{F}_p}$ 都是单射。因此，可以假设
$K/k$ 是有限生成域扩张。

\medskip\noindent
假设 $K/k$ 是有限生成可分域扩张。按照引理
\ref{algebra-lemma-generating-finitely-generated-separable-field-extensions}
选取 $x_1, \ldots, x_{r + 1} \in K$。
在此情形下，存在不可约多项式
$G(X_1, \ldots, X_{r + 1}) \in k[X_1, \ldots, X_{r + 1}]$，
使得 $G(x_1, \ldots, x_{r + 1}) = 0$，且
$\partial G/\partial X_{r + 1}$ 不恒为零。
此外，$K$ 是该整环的分式域。
% P01-E0982: the frozen English splits an incomplete domain sentence from a presentation with the wrong base field K instead of k.
$S = K[X_1, \ldots, X_{r + 1}]/(G)$。
写成
$$
G = \sum a_I X^I, \quad X^I = X_1^{i_1}\ldots X_{r + 1}^{i_{r + 1}}.
$$
利用上面对 $S$ 的表现，可知
$$
\Omega_{S/\mathbf{F}_p}
=
\frac{
S \otimes_k \Omega_k \oplus
\bigoplus\nolimits_{i = 1, \ldots, r + 1} S\text{d}X_i
}{
\langle
\sum X^I \text{d}a_I + \sum \partial G/\partial X_i \text{d}X_i
\rangle
}
$$
% P01-E0983: the two frozen presentations use the undefined abbreviation Omega_k instead of Omega_{k/F_p}.
由于 $\Omega_{K/\mathbf{F}_p}$ 是 $S$-模
$\Omega_{S/\mathbf{F}_p}$ 的局部化（见引理
\ref{algebra-lemma-differentials-localize}），可得
$$
\Omega_{K/\mathbf{F}_p}
=
\frac{
K \otimes_k \Omega_k \oplus
\bigoplus\nolimits_{i = 1, \ldots, r + 1} K\text{d}X_i
}{
\langle
\sum X^I \text{d}a_I + \sum \partial G/\partial X_i \text{d}X_i
\rangle
}
$$
现在，由于多项式 $\partial G/\partial X_{r + 1}$ 不恒为零，
可知映射
$K \otimes_k \Omega_{k/\mathbf{F}_p} \to \Omega_{S/\mathbf{F}_p}$
是单射，如所需。
% P01-E0984: the frozen conclusion names Omega_{S/F_p}, although the desired and just-presented codomain is Omega_{K/F_p}.

\medskip\noindent
假设 $K/k$ 是有限生成域扩张，并且
$K \otimes_k \Omega_{k/\mathbf{F}_p} \to \Omega_{K/\mathbf{F}_p}$
是单射。（证明的这一部分与证明引理
\ref{algebra-lemma-characterize-separable-field-extensions} 的论证相同。）
设 $x_1, \ldots, x_r$ 是 $K$ 在 $k$ 上的超越基，并且有限扩张
$k(x_1, \ldots, x_r) \subset K$ 的不可分次数最小。
如果 $K$ 在 $k(x_1, \ldots, x_r)$ 上可分，则已经得到结论。
假设不是如此，以导出矛盾。于是存在元素
$\alpha \in K$，它在 $k(x_1, \ldots, x_r)$ 上不可分。
设 $P(T) \in k(x_1, \ldots, x_r)[T]$ 是它的最小多项式。
因为 $\alpha$ 不可分，实际上 $P$ 是 $T^p$ 的多项式。
清除分母，得到不可约多项式
$$
G(X_1, \ldots, X_r, T) = \sum a_{I, i} X^I T^i \in k[X_1, \ldots, X_r, T]
$$
使得 $G(x_1, \ldots, x_r, \alpha) = 0$ 在 $K$ 中成立。
注意，这意味着
$k[X_1, \ldots, X_r, T]/(G) \subset K$。
可以假设对某个二元组 $(I_0, i_0)$ 有
$a_{I_0, i_0} = 1$。
我们断言，至少对一个 $i$，$\text{d}G/\text{d}X_i$
不恒为零。事实上，如果不是这样，则 $G$ 实际上是
$X_1^p, \ldots, X_r^p, T^p$ 的多项式。于是
$$
\sum\nolimits_{(I, i) \not = (I_0, i_0)} x^I\alpha^i \text{d}a_{I, i}
$$
在 $\Omega_{K/\mathbf{F}_p}$ 中为零。注意，$K$ 的元素集合
$$
\{x^I\alpha^i \mid a_{I, i} \not = 0 \text{ 且 } (I, i) \not = (I_0, i_0)\}
$$
中不存在 $k$-线性关系。因此，映射
$K \otimes_k \Omega_{k/\mathbf{F}_p} \to \Omega_{K/\mathbf{F}_p}$
是单射这一假设蕴含：对所有 $(I, i)$，
$\text{d}a_{I, i} = 0$ 在 $\Omega_{k/\mathbf{F}_p}$ 中成立。
由引理 \ref{algebra-lemma-derivative-zero-pth-power} 可知，每个
$a_{I, i}$ 都是 $p$ 次幂，这意味着 $G$ 是 $p$ 次幂，
与 $G$ 的不可约性矛盾。因此，重新编号后可以假设
$\text{d}G/\text{d}X_1$ 非零。
于是可知 $x_1$ 在 $k(x_2, \ldots, x_r, \alpha)$ 上可分代数，
并且 $x_2, \ldots, x_r, \alpha$ 是 $K$ 在 $k$ 上的超越基。
这意味着有限扩张
$k(x_2, \ldots, x_r, \alpha) \subset K$ 的不可分次数小于有限扩张
$k(x_1, \ldots, x_r) \subset K$ 的不可分次数，产生矛盾。
\end{proof}

\begin{lemma}
\label{algebra-lemma-formally-smooth-implies-separable}
设 $K/k$ 是域扩张。如果 $K$ 在 $k$ 上形式光滑，
则 $K$ 是 $k$ 的可分扩张。
\end{lemma}

\begin{proof}
假设 $K$ 在 $k$ 上形式光滑。由引理
\ref{algebra-lemma-ses-formally-smooth} 可知
$K \otimes_k \Omega_{k/\mathbf{Z}} \to \Omega_{K/\mathbf{Z}}$
是单射。于是由引理 \ref{algebra-lemma-separable-differentials}，
$K$ 在 $k$ 上可分。
% P01-E0985: the frozen proof invokes a characteristic-p lemma without separately handling characteristic zero.
\end{proof}

\begin{lemma}
\label{algebra-lemma-characterize-formally-smooth-field-extension}
设 $K/k$ 是域扩张。则 $K$ 在 $k$ 上形式光滑，当且仅当
$H_1(L_{K/k}) = 0$。
\end{lemma}

\begin{proof}
这由命题 \ref{algebra-proposition-characterize-formally-smooth} 以及
向量空间是自由模（因而是投射模）这一事实得出。
% P01-E0986: frozen English has the number disagreement “a vector spaces is free”.
\end{proof}

\begin{lemma}
\label{algebra-lemma-formally-smooth-extensions-easy}
设 $K/k$ 是域扩张。
\begin{enumerate}
\item 如果 $K$ 在 $k$ 上纯超越，则 $K$ 在 $k$ 上形式光滑。
\item 如果 $K$ 在 $k$ 上可分代数，则 $K$ 在 $k$ 上形式光滑。
\item 如果 $K$ 在 $k$ 上可分，则 $K$ 在 $k$ 上形式光滑。
\end{enumerate}
\end{lemma}

\begin{proof}
对 (1)，写 $K = k(x_j; j \in J)$。假设 $A$ 是 $k$-代数，
$I \subset A$ 是平方为零的理想。设
$\varphi : K \to A/I$ 是 $k$-代数映射。
设 $a_j \in A$ 满足 $a_j \mod I = \varphi(x_j)$。
那么容易看出，存在唯一的 $k$-代数映射 $K \to A$，
它把 $x_j$ 映到 $a_j$，且模 $I$ 约化为 $\varphi$。
因此，$k \subset K$ 形式光滑。

\medskip\noindent
在情形 (2) 中，$k \subset K$ 是平展环扩张的余极限。
平展环映射是形式平展的（引理
\ref{algebra-lemma-formally-etale-etale}）。因此，该情形由引理
\ref{algebra-lemma-colimit-formally-etale} 以及形式平展环映射形式光滑这一
显然事实得出。

\medskip\noindent
在情形 (3) 中，把 $K = \colim K_i$ 写成其有限生成
$k$-子扩张的滤过余极限。
% P01-E0987: frozen English says “finitely generated sub k-extensions” instead of “finitely generated k-subextensions”.
由定义 \ref{algebra-definition-separable-field-extension}，每个 $K_i$
都是 $k$ 的某个纯超越扩张上的可分代数扩张。
因此，由情形 (1)、(2) 以及引理
\ref{algebra-lemma-compose-formally-smooth}，$K_i/k$ 形式光滑。
于是由引理
\ref{algebra-lemma-characterize-formally-smooth-field-extension}，
$H_1(L_{K_i/k}) = 0$。由引理 \ref{algebra-lemma-colimits-NL}，
$H_1(L_{K/k}) = 0$。再次应用引理
\ref{algebra-lemma-characterize-formally-smooth-field-extension}，可知
$K/k$ 形式光滑。
\end{proof}

\begin{lemma}
\label{algebra-lemma-fields-are-formally-smooth}
\begin{slogan}
对域扩张，形式光滑等价于可分。
\end{slogan}
设 $k$ 是域。
\begin{enumerate}
\item 如果 $k$ 的特征为零，则 $k$ 的任意扩域都在 $k$
上形式光滑。
\item 如果 $k$ 的特征为 $p > 0$，则 $K/k$ 形式光滑，
当且仅当它是可分域扩张。
\end{enumerate}
\end{lemma}

\begin{proof}
合并引理 \ref{algebra-lemma-formally-smooth-implies-separable} 和
\ref{algebra-lemma-formally-smooth-extensions-easy}。
\end{proof}

\noindent
下面汇集可分域扩张的所有不同刻画。

\begin{proposition}
\label{algebra-proposition-characterize-separable-field-extensions}
设 $K/k$ 是域扩张。如果 $k$ 的特征为零，则
% P01-E0988: the frozen proposition omits the phrase “the following are equivalent” before this list.
\begin{enumerate}
\item $K$ 在 $k$ 上可分；
\item $K$ 在 $k$ 上几何约化；
\item $K$ 在 $k$ 上形式光滑；
\item $H_1(L_{K/k}) = 0$；
\item 映射 $K \otimes_k \Omega_{k/\mathbf{Z}} \to \Omega_{K/\mathbf{Z}}$
是单射。
\end{enumerate}
如果 $k$ 的特征为 $p > 0$，则下列条件等价：
\begin{enumerate}
\item $K$ 在 $k$ 上可分；
\item 环 $K \otimes_k k^{1/p}$ 约化；
\item $K$ 在 $k$ 上几何约化；
\item 映射 $K \otimes_k \Omega_{k/\mathbf{F}_p} \to \Omega_{K/\mathbf{F}_p}$
是单射；
\item $H_1(L_{K/k}) = 0$；
\item $K$ 在 $k$ 上形式光滑。
\end{enumerate}
\end{proposition}

\begin{proof}
这是下列引理的合并：
\ref{algebra-lemma-characterize-separable-field-extensions}、
\ref{algebra-lemma-fields-are-formally-smooth}
\ref{algebra-lemma-formally-smooth-implies-separable} 以及
\ref{algebra-lemma-separable-differentials}。
% P01-E0989: frozen English omits punctuation or a conjunction between the middle two lemma references.
\end{proof}

\noindent
下面给出有限生成可分域扩张的又一个刻画。

\begin{lemma}
\label{algebra-lemma-localization-smooth-separable}
设 $K/k$ 是有限生成域扩张。则 $K$ 在 $k$ 上可分，
当且仅当 $K$ 是某个光滑 $k$-代数的局部化。
\end{lemma}

\begin{proof}
选取有限型 $k$-代数 $R$，使其为分式域等于 $K$ 的整环。
引理 \ref{algebra-lemma-smooth-at-generic-point} 说明，
$k \to R$ 在 $(0)$ 处光滑，当且仅当 $K/k$ 可分。
这证明了该引理。
\end{proof}

\begin{lemma}
\label{algebra-lemma-colimit-syntomic}
设 $K/k$ 是域扩张。则 $K$ 是 $k$ 上全局完全交代数的
滤过余极限。如果 $K/k$ 可分，则 $K$ 是 $k$ 上光滑代数的
滤过余极限。
\end{lemma}

\begin{proof}
假设 $E \subset K$ 是有限子集。只需证明，存在包含 $E$
的 $k$-子代数 $A \subset K$，并且它是 $k$ 上的全局完全交
（相应地，光滑）。可分/光滑情形由引理
\ref{algebra-lemma-localization-smooth-separable} 得出。
一般地，设 $L \subset K$ 是由 $E$ 生成的子域。
选取它在 $k$ 上的超越基
$x_1, \ldots, x_d \in L$。扩张
$L/k(x_1, \ldots, x_d)$ 是有限的。设
$L = k(x_1, \ldots, x_d)[y_1, \ldots, y_r]$。
归纳选取多项式
$P_i \in k(x_1, \ldots, x_d)[Y_1, \ldots, Y_r]$，
使得 $P_i = P_i(Y_1, \ldots, Y_i)$ 是
$k(x_1, \ldots, x_d)[Y_1, \ldots, Y_{i - 1}]$
上关于 $Y_i$ 的首一多项式，并且映到 $y_i$ 在
$k(x_1, \ldots, x_d)[y_1, \ldots, y_{i - 1}][Y_i]$
中的最小多项式。
% P01-E0990: frozen English says “minimum polynomial” instead of “minimal polynomial”.
于是显然，$P_1, \ldots, P_r$ 是
$k(x_1, \ldots, x_r)[Y_1, \ldots, Y_r]$ 中的正则序列，且
$L = k(x_1, \ldots, x_r)[Y_1, \ldots, Y_r]/(P_1, \ldots, P_r)$。
% P01-E0991: both frozen formulas use x_1,...,x_r where the established transcendence basis is x_1,...,x_d.
如果 $h \in k[x_1, \ldots, x_d]$ 是使得
$P_i \in k[x_1, \ldots, x_d, 1/h, Y_1, \ldots, Y_r]$
的多项式，则 $P_1, \ldots, P_r$ 是
$k[x_1, \ldots, x_d, 1/h, Y_1, \ldots, Y_r]$ 中的正则序列，且
$A = k[x_1, \ldots, x_d, 1/h, Y_1, \ldots, Y_r]/(P_1, \ldots, P_r)$
是全局完全交。适当调整 $h$ 的选取后，可以假设
$E \subset A$，从而得到结论。
\end{proof}

\section{构造平坦环映射}
\label{algebra-section-constructing-flat}

\noindent
下面的引理有时很有用。

\begin{lemma}
\label{algebra-lemma-flat-local-given-residue-field}
设 $(R, \mathfrak m, k)$ 是局部环，$K/k$ 是域扩张。
则存在局部环 $(R', \mathfrak m', k')$ 以及平坦局部环映射
$R \to R'$，使得 $\mathfrak m' = \mathfrak mR'$，并且
$k'$ 作为 $k$ 的扩张同构于 $K$。
\end{lemma}

\begin{proof}
假设 $k' = k(\alpha)$ 是 $k$ 的单生成扩张。
则 $k'$ 是引理所述某平坦局部扩张 $R \subset R'$ 的剩余域。
具体而言，如果 $\alpha$ 在 $k$ 上超越，则令 $R'$ 为
$R[x]$ 在素理想 $\mathfrak mR[x]$ 处的局部化。
如果 $\alpha$ 是代数的，其最小多项式为
$T^d + \sum \overline{\lambda}_iT^{d - i}$，则令
$R' = R[T]/(T^d + \sum \lambda_i T^{d - i})$。

\medskip\noindent
考虑所有三元组 $(k', R \to R', \phi)$ 所成的集合，其中
$k \subset k' \subset K$ 是子域，$R \to R'$ 是引理所述的
局部环映射，并且
$\phi : R' \to k'$ 诱导 $k$-扩张的同构
$R'/\mathfrak mR' \cong k'$。这些对象构成一个``大''范畴
$\mathcal{C}$；态射
$(k_1, R_1, \phi_1) \to (k_2, R_2, \phi_2)$ 由环映射
$\psi : R_1 \to R_2$ 给出，并要求交换图表
$$
\xymatrix{
R_1 \ar[d]_\psi \ar[r]_{\phi_1} & k_1 \ar[r] & K \ar@{=}[d] \\
R_2 \ar[r]^{\phi_2} & k_2 \ar[r] & K
}
$$
交换。这蕴含 $k_1 \subset k_2$。

\medskip\noindent
假设 $I$ 是有向集，而
$((R_i, k_i, \phi_i), \psi_{ii'})$ 是 $I$ 上的系统，见范畴，
第 \ref{categories-section-posets-limits} 节。
% P01-E0992: frozen object order changes from (k', R -> R', phi) to (R_i, k_i, phi_i), and the stated upper bound repeats the reversed order.
在此情形下，可以考虑
$$
R' = \colim_{i \in I} R_i
$$
这是极大理想为 $\mathfrak mR'$ 的局部环，剩余域为
$k' = \bigcup_{i \in I} k_i$。此外，环映射 $R \to R'$
是平坦的，因为它是平坦映射的余极限（且张量积与有向余极限交换）。
因此，$(R', k', \phi')$ 是该系统的一个``上界''。

\medskip\noindent
如果 $\mathcal{C}$ 是集合，则几乎平凡地应用 Zorn 引理即可完成证明，
但它并不是集合。
% P01-E0993: frozen English uses “if C was a set” rather than the irrealis “if C were a set”.
（实际上，可以通过找出所出现局部环之基数的合理界来使这一论证成立。）
为绕开这个问题，我们在 $K$ 上选取一个良序。
对 $x \in K$，令 $K(x)$ 为 $K$ 的如下子域：由 $K$ 中所有
$\leq x$ 的元素所生成。对 $x \in K$ 作超限递归，我们将构造引理所述的
环映射 $R \subset R(x)$，其剩余域扩张为 $K(x)/k$。
此外，按构造，对所有 $y \leq x$，$R(x)$ 都将包含 $R(y)$。
具体而言，如果 $x$ 有前驱 $x'$，则 $K(x) = K(x')[x]$，
因而可以令 $R(x') \subset R(x)$ 为证明第一段所构造的局部环扩张。
如果 $x$ 没有前驱，则先按照证明第三段令
$R'(x) = \colim_{x' < x} R(x')$。
$R'(x)$ 的剩余域为
$K'(x) = \bigcup_{x' < x} K(x')$。由于
$K(x) = K'(x)[x]$，可用证明第一段的构造得到
$R'(x) \subset R(x)$。
% P01-E0994: both frozen field-generation formulas use square brackets, which fail when x is transcendental.
这就完成了引理的证明。
\end{proof}

\begin{lemma}
\label{algebra-lemma-colimit-finite-etale-given-residue-field}
设 $(R, \mathfrak m, k)$ 是局部环。如果 $k \subset K$
是可分代数扩张，则存在有向集 $I$ 以及局部环的有限平展扩张系统
$R \subset R_i$，$i \in I$，使得 $R' = \colim R_i$
的剩余域为 $K$（作为 $k$ 的扩张）。
\end{lemma}

\begin{proof}
设 $R \subset R'$ 是引理
\ref{algebra-lemma-flat-local-given-residue-field} 的证明中构造的扩张。
按构造，$R' = \colim_{\alpha \in A} R_\alpha$，其中 $A$
是良序集，过渡映射 $R_\alpha \to R_{\alpha + 1}$ 是有限平展的，
并且当 $\alpha$ 不是后继元时，
$R_\alpha = \colim_{\beta < \alpha} R_\beta$。
我们将用超限归纳证明该结果。

\medskip\noindent
假设结论对 $R_\alpha$ 成立，即
$R_\alpha = \colim R_i$，其中 $R_i$ 在 $R$ 上有限平展。
由于 $R_\alpha \to R_{\alpha + 1}$ 有限平展，存在一个 $i$
以及有限平展扩张 $R_i \to R_{i, 1}$，使得
$R_{\alpha + 1} = R_\alpha \otimes_{R_i} R_{i, 1}$。
因此，
$R_{\alpha + 1} = \colim_{i' \geq i} R_{i'} \otimes_{R_i} R_{i, 1}$，
结论对 $\alpha + 1$ 成立。
假设 $\alpha$ 不是后继元，并且结论对所有
$\beta < \alpha$ 的 $R_\beta$ 成立。
由于每个有限子集 $E \subset R_\alpha$ 都包含于某个
$\beta < \alpha$ 对应的 $R_\beta$ 中，由假设可知 $E$
包含于某个有限平展子扩张中。因此，结论对 $R_\alpha$ 成立。
% P01-E0995: frozen English inserts an ungrammatical “and” between the finite-subset premise and its conclusion.
\end{proof}

\begin{lemma}
\label{algebra-lemma-finite-free-given-residue-field-extension}
设 $R$ 是环，$\mathfrak p \subset R$ 是素理想，并设
$L/\kappa(\mathfrak p)$ 是有限域扩张。
则存在有限自由环映射 $R \to S$，使得
$\mathfrak q = \mathfrak pS$ 是素理想，并且
$\kappa(\mathfrak q)/\kappa(\mathfrak p)$ 同构于给定扩张
$L/\kappa(\mathfrak p)$。
\end{lemma}

\begin{proof}
对 $\kappa(\mathfrak p) \subset L$ 的次数作归纳。
% P01-E0996: frozen English says “induction of the degree” instead of “induction on the degree”.
如果次数为 $1$，则取 $R = S$。
一般地，如果存在子扩张
$\kappa(\mathfrak p) \subset L' \subset L$，则对次数归纳即可得到
结论（先构造与 $L'/\kappa(\mathfrak p)$ 对应的
$R \subset S'$，再构造与 $L/L'$ 对应的 $S' \subset S$）。
% P01-E0997: the frozen conditional is vacuous without requiring a proper intermediate field, and “then construction” is also malformed.
因此，可以假设 $L \supset \kappa(\mathfrak p)$ 由单个元素
$\alpha \in L$ 生成。设
$X^d + \sum_{i < d} a_iX^i$ 是 $\alpha$ 在
$\kappa(\mathfrak p)$ 上的最小多项式，所以
$a_i \in \kappa(\mathfrak p)$。
可以把 $a_i$ 写成某个 $f_i, g \in R$ 的 $f_i/g$ 的像，
其中 $g \not \in \mathfrak p$。
用 $g\alpha$ 替换 $\alpha$（并相应地用
$g^{d - i}a_i$ 替换 $a_i$）后，可以假设 $a_i$ 是某个
$f_i \in R$ 的像。此时只需取
$S = R[x]/(x^d + \sum f_ix^i)$。
\end{proof}

\begin{lemma}
\label{algebra-lemma-cofinal-system-flat}
设 $A$ 是环，令 $\kappa = \max(|A|, \aleph_0)$。
则每个平坦 $A$-代数 $B$ 都是其满足
$B' \subset B$、$|B'| \leq \kappa$ 的平坦 $A$-子代数的
滤过余极限。（注意，如果 $B$ 在 $A$ 上忠实平坦，
则 $B'$ 也在 $A$ 上忠实平坦。）
\end{lemma}

\begin{proof}
如果 $B$ 的基数 $\leq \kappa$，则结论成立。
设 $E \subset B$ 是满足 $|E| \leq \kappa$ 的 $A$-子代数。
我们将证明 $E$ 包含于某个平坦 $A$-子代数 $B'$ 中，且
$|B'| \leq \kappa$。于是引理成立，因为
(a) $B$ 的每个有限子集都包含于某个基数至多为 $\kappa$ 的
$A$-子代数中，并且 (b) $B$ 的任意两个基数至多为
$\kappa$ 的 $A$-子代数都包含于某个基数至多为 $\kappa$
的 $A$-子代数中。略去细节。

\medskip\noindent
我们将归纳构造一列 $A$-子代数
$$
E = E_0 \subset E_1 \subset E_2 \subset \ldots
$$
其中每个子代数的基数都 $\leq \kappa$；最后证明
$B' = \bigcup E_k$ 在 $A$ 上平坦。

\medskip\noindent
构造如下。令 $E_0 = E$。给定 $k \geq 0$ 时的 $E_k$，
考虑 $E_k$ 中元素之间以 $A$ 中元素为系数的关系所成的集合
$S_k$。因此，元素 $s \in S_k$ 由整数 $n \geq 1$、
$a_1, \ldots, a_n \in A$ 以及
$e_1, \ldots, e_n \in E_k$ 给出，并满足
$\sum a_i e_i = 0$ 在 $E_k$ 中成立。
$A \to B$ 的平坦性以及引理 \ref{algebra-lemma-flat-eq} 表明，
对每个
$s = (n, a_1, \ldots, a_n, e_1, \ldots, e_n) \in S_k$，
可以选取
$$
(m_s, b_{s, 1}, \ldots, b_{s, m_s}, a_{s, 11}, \ldots, a_{s, nm_s})
$$
其中 $m_s \geq 0$ 是整数，$b_{s, j} \in B$，
$a_{s, ij} \in A$，并且
$$
e_i = \sum\nolimits_j a_{s, ij} b_{s, j}, \forall i,
\quad\text{且}\quad
0 = \sum\nolimits_i a_i a_{s, ij}, \forall j.
$$
给定这些选取后，令 $E_{k + 1} \subset B$ 为下列元素生成的
$A$-子代数：
% P01-E0998: frozen English misspells “choices” as “choicse”.
\begin{enumerate}
\item $E_k$；
\item 对每个 $s \in S_k$，元素
$b_{s, 1}, \ldots, b_{s, m_s}$。
\end{enumerate}
一些集合论（略去）表明，$E_{k + 1}$ 的基数至多为
$\kappa$（这里用到归纳已知 $|E_k| \leq \kappa$，因而
$S_k$ 的基数也至多为 $\kappa$）。
% P01-E0999: frozen English says “has at most cardinality kappa” instead of “has cardinality at most kappa”.

\medskip\noindent
为了证明 $B' = \bigcup E_k$ 在 $A$ 上平坦，考虑 $B'$ 中
系数属于 $A$ 的关系
$\sum_{i = 1, \ldots, n} a_i b'_i = 0$。
选取充分大的 $k$，使得对 $i = 1, \ldots, n$ 有
$b'_i \in E_k$。则
$(n, a_1, \ldots, a_n, b'_1, \ldots, b'_n) \in S_k$，
所以该关系在 $E_{k + 1}$ 中是平凡的，从而在 $B'$ 中更是如此。
因此，由引理 \ref{algebra-lemma-flat-eq}，$A \to B'$ 是平坦的。
\end{proof}

\section{Cohen 结构定理}
\label{algebra-section-cohen-structure-theorem}

\noindent
下面是交换代数中的一个基本概念。

\begin{definition}
\label{algebra-definition-complete-local-ring}
设 $(R, \mathfrak m)$ 是局部环。如果到 $R$ 关于
$\mathfrak m$ 的完备化的典范映射
$$
R \longrightarrow \lim_n R/\mathfrak m^n
$$
是同构，则称 $R$ 是{\it 完备局部环}\footnote{这包括条件
$\bigcap \mathfrak m^n = (0)$；在一些文献中，这可能表述为
$R$ 完备且分离。警告：局部环的完备化
$\lim_n R/\mathfrak m^n$ 可能不是完备的，见例，引理
\ref{examples-lemma-noncomplete-completion}。
当 $\mathfrak m$ 有限生成时不会发生这种情况，见引理
\ref{algebra-lemma-hathat-finitely-generated}；在这种情况下，完备化是
Noether 环，见引理 \ref{algebra-lemma-completion-Noetherian}。}。
\end{definition}

\noindent
注意，Artin 局部环 $R$ 是完备局部环，因为对某个
$n > 0$ 有 $\mathfrak m_R^n = 0$。本节主要讨论 Noether
完备局部环。

\begin{lemma}
\label{algebra-lemma-quotient-complete-local}
设 $R$ 是 Noether 完备局部环。$R$ 的任意商也是 Noether
完备局部环。给定有限环映射 $R \to S$，则 $S$ 是 Noether
完备局部环的乘积。
% P01-E1000: frozen English uses the doubled construction “Given ..., then ...”.
\end{lemma}

\begin{proof}
由引理 \ref{algebra-lemma-Noetherian-permanence}，环 $S$ 是 Noether 的。
由引理 \ref{algebra-lemma-completion-tensor}，$S$ 作为 $R$-模是完备的。
因此，由引理 \ref{algebra-lemma-completion-finite-extension}，$S$
是它在各极大理想处的完备化的乘积。
\end{proof}

\begin{lemma}
\label{algebra-lemma-complete-local-ring-Noetherian}
设 $(R, \mathfrak m)$ 是完备局部环。如果 $\mathfrak m$
是有限生成理想，则 $R$ 是 Noether 环。
\end{lemma}

\begin{proof}
见引理 \ref{algebra-lemma-completion-Noetherian}。
\end{proof}

\begin{definition}
\label{algebra-definition-coefficient-ring}
设 $(R, \mathfrak m)$ 是完备局部环。如果子环
$\Lambda \subset R$ 满足下列条件，则称它为{\it 系数环}：
\begin{enumerate}
\item $\Lambda$ 是完备局部环，其极大理想为
$\Lambda \cap \mathfrak m$；
\item $\Lambda$ 的剩余域同构地映到 $R$ 的剩余域；
\item $\Lambda \cap \mathfrak m = p\Lambda$，其中 $p$ 是
$R$ 的剩余域的特征。
\end{enumerate}
\end{definition}

\noindent
下面对该定义作几点说明。分成下列情形讨论：
\begin{enumerate}
\item 局部环 $R$ 包含一个域。如果
$\mathbf{Q} \subset R$，或者 $pR = 0$，其中 $p$ 是
$R/\mathfrak m$ 的特征，就会出现这种情况。此时，系数环
$\Lambda$ 是包含于 $R$ 中的域，并同构地映到
$R/\mathfrak m$。
\item $R/\mathfrak m$ 的特征为 $p > 0$，但 $p$ 的任何幂
在 $R$ 中都不为零。此时，$\Lambda$ 是以 $p$ 为一致化元、
剩余域为 $R/\mathfrak m$ 的完备离散赋值环。
\item $R/\mathfrak m$ 的特征为 $p > 0$，并且对某个
$n > 1$，在 $R$ 中有 $p^{n - 1} \not = 0$、$p^n = 0$。
此时，$\Lambda$ 是 Artin 局部环，其极大理想由 $p$ 生成，
且剩余域为 $R/\mathfrak m$。
\end{enumerate}
上述以 $p$ 为一致化元的完备离散赋值环具有特殊作用，
我们作如下命名。

\begin{definition}
\label{algebra-definition-cohen-ring}
{it Cohen 环}是以某个素数 $p$ 为一致化元的完备离散赋值环。
\end{definition}

\begin{lemma}
\label{algebra-lemma-cohen-rings-exist}
设 $p$ 是素数，$k$ 是特征为 $p$ 的域。
则存在 Cohen 环 $\Lambda$，使得 $\Lambda/p\Lambda \cong k$。
\end{lemma}

\begin{proof}
首先注意，$p$ 进整数 $\mathbf{Z}_p$ 构成
$\mathbf{F}_p$ 的一个 Cohen 环。设 $k$ 是特征 $p$ 的任意域。
取平坦局部环映射 $\mathbf{Z}_p \to R$，使得
$\mathfrak m_R = pR$ 且 $R/pR = k$，见引理
\ref{algebra-lemma-flat-local-given-residue-field}。
由引理 \ref{algebra-lemma-completion-Noetherian}，完备化
$\Lambda = R^\wedge$ 是 Noether 的。由于
$\Lambda/p\Lambda = R/pR$ 是域，可知 $\Lambda$ 是极大理想为
$(p)$ 的完备 Noether 局部环（该等式及完备性由引理
\ref{algebra-lemma-hathat-finitely-generated} 得出）。
对每个 $n$，映射
$\mathbf{Z}/p^n\mathbf{Z} \to \Lambda/p^n\Lambda = R/p^nR$
都是平坦的（该等式再次由引理
\ref{algebra-lemma-hathat-finitely-generated} 得出）。
因此，例如由引理 \ref{algebra-lemma-flat-module-powers} 可知
$\mathbf{Z}_p \to \Lambda$ 平坦。所以 $p$ 是 $\Lambda$
中的非零因子。于是 $\Lambda$ 的维数 $\geq 1$
（引理 \ref{algebra-lemma-one-equation}），从而 $\Lambda$ 是维数为
$1$ 的正则环；即由引理 \ref{algebra-lemma-characterize-dvr}，
它是离散赋值环。因此，$\Lambda$ 是 $k$ 的 Cohen 环。
\end{proof}

\begin{lemma}
\label{algebra-lemma-cohen-ring-formally-smooth}
设 $p > 0$ 是素数。设 $\Lambda$ 是剩余域特征为 $p$
的 Cohen 环。对每个 $n \geq 1$，环映射
$$
\mathbf{Z}/p^n\mathbf{Z} \to \Lambda/p^n\Lambda
$$
都是形式光滑的。
\end{lemma}

\begin{proof}
当 $n = 1$ 时，这由命题
\ref{algebra-proposition-characterize-separable-field-extensions} 得出。
对一般的 $n$，对 $n$ 作归纳。具体而言，如果
$\mathbf{Z}/p^n\mathbf{Z} \to \Lambda/p^n\Lambda$ 形式光滑，
则可将引理 \ref{algebra-lemma-lift-formal-smoothness} 应用于环映射
$\mathbf{Z}/p^{n + 1}\mathbf{Z} \to \Lambda/p^{n + 1}\Lambda$
以及理想
$I = (p^n) \subset \mathbf{Z}/p^{n + 1}\mathbf{Z}$。
\end{proof}

\begin{theorem}[Cohen 结构定理]
\label{algebra-theorem-cohen-structure-theorem}
设 $(R, \mathfrak m)$ 是完备局部环。
\begin{enumerate}
\item $R$ 有一个系数环（见定义
\ref{algebra-definition-coefficient-ring}）；
\item 如果 $\mathfrak m$ 是有限生成理想，则 $R$ 同构于商
$$
\Lambda[[x_1, \ldots, x_n]]/I
$$
其中 $\Lambda$ 是域或 Cohen 环。
\end{enumerate}
\end{theorem}

\begin{proof}
先证明系数环存在。首先考虑剩余域 $\kappa$ 的特征为零的情形。
具体而言，在这种情况下，我们对 $n > 0$ 归纳证明存在映射
$$
\varphi_n : \kappa \longrightarrow R/\mathfrak m^n
$$
它是典范映射 $R/\mathfrak m^n \to \kappa = R/\mathfrak m$
的截面。当 $n = 1$ 时，这显然成立。如果 $n > 1$，
设已给定 $\varphi_{n - 1}$。
由命题 \ref{algebra-proposition-characterize-separable-field-extensions}，
域扩张 $\kappa/\mathbf{Q}$ 形式光滑。因此，可以在下图中找到
虚线箭头：
$$
\xymatrix{
R/\mathfrak m^{n - 1} &
R/\mathfrak m^n \ar[l] \\
\kappa \ar[u]^{\varphi_{n - 1}} \ar@{..>}[ru] & \mathbf{Q} \ar[l] \ar[u]
}
$$
这证明了归纳步骤。把这些映射组合起来，得到
$$
\lim_n\ \varphi_n : \kappa \longrightarrow
R = \lim_n\ R/\mathfrak m^n
$$
其像就是所需的系数环。

\medskip\noindent
其次，证明剩余域 $\kappa$ 的特征为 $p > 0$ 时系数环存在。
具体而言，选取满足 $\kappa = \Lambda/p\Lambda$ 的 Cohen 环
$\Lambda$，见引理 \ref{algebra-lemma-cohen-rings-exist}。
在这种情况下，我们对 $n > 0$ 归纳证明存在映射
$$
\varphi_n :
\Lambda/p^n\Lambda
\longrightarrow
R/\mathfrak m^n
$$
它与约化映射 $R/\mathfrak m^n \to \kappa$ 的合成给出
已知同构 $\Lambda/p\Lambda = \kappa$。
当 $n = 1$ 时，这显然成立。如果 $n > 1$，设已给定
$\varphi_{n - 1}$。由引理
\ref{algebra-lemma-cohen-ring-formally-smooth}，环映射
$\mathbf{Z}/p^n\mathbf{Z} \to \Lambda/p^n\Lambda$ 形式光滑。
因此，可以在下图中找到虚线箭头：
$$
\xymatrix{
R/\mathfrak m^{n - 1} &
R/\mathfrak m^n \ar[l] \\
\Lambda/p^n\Lambda \ar[u]^{\varphi_{n - 1}} \ar@{..>}[ru] &
\mathbf{Z}/p^n\mathbf{Z} \ar[l] \ar[u]
}
$$
这证明了归纳步骤。把这些映射组合起来，得到
$$
\lim_n\ \varphi_n :
\Lambda = \lim_n\ \Lambda/p^n\Lambda
\longrightarrow
R = \lim_n\ R/\mathfrak m^n
$$
其像就是所需的系数环。

\medskip\noindent
定理的最后一个断言容易得出。具体而言，如果
$y_1, \ldots, y_n$ 是理想 $\mathfrak m$ 的生成元，则可利用
刚构造的映射 $\Lambda \to R$ 得到映射
$$
\Lambda[[x_1, \ldots, x_n]] \longrightarrow R,
\quad x_i \longmapsto y_i.
$$
由于两边关于 $(x_1, \ldots, x_n)$ 都是完备的，并且按构造该映射
模 $(x_1, \ldots, x_n)$ 是满射，由引理
\ref{algebra-lemma-completion-generalities} 可知该映射是满射。
\end{proof}

\begin{remark}
\label{algebra-remark-Noetherian-complete-local-ring-universally-catenary}
如果 $k$ 是域，则幂级数环 $k[[X_1, \ldots, X_d]]$ 是维数为
$d$ 的 Noether 完备正则局部环。如果 $\Lambda$ 是 Cohen 环，
则 $\Lambda[[X_1, \ldots, X_d]]$ 是维数为 $d + 1$ 的
完备 Noether 正则局部环。因此，Cohen 结构定理蕴含：
任意 Noether 完备局部环都是正则局部环的商。
特别地，Noether 完备局部环是普遍链状的，见引理
\ref{algebra-lemma-CM-ring-catenary} 和引理
\ref{algebra-lemma-regular-ring-CM}。
\end{remark}

\begin{lemma}
\label{algebra-lemma-regular-complete-containing-coefficient-field}
设 $(R, \mathfrak m)$ 是 Noether 完备局部环。假设 $R$ 正则。
\begin{enumerate}
\item 如果 $R$ 包含 $\mathbf{F}_p$ 或 $\mathbf{Q}$，则 $R$
同构于其剩余域上的幂级数环。
\item 如果 $k$ 是域，并且环映射 $k \to R$ 诱导同构
$k \to R/\mathfrak m$，则 $R$ 作为 $k$-代数同构于 $k$
上的幂级数环。
\end{enumerate}
\end{lemma}

\begin{proof}
在情形 (1) 中，由 Cohen 结构定理（定理
\ref{algebra-theorem-cohen-structure-theorem}），存在一个系数环；
它必须是同构地映到剩余域的域。因此，只需证明 (2)。
在情形 (2) 中，选取
$f_1, \ldots, f_d \in \mathfrak m$，使它们映到
$\mathfrak m/\mathfrak m^2$ 的一组基，并考虑把 $x_i$ 映到
$f_i$ 的连续 $k$-代数映射
$k[[x_1, \ldots, x_d]] \to R$。
由于源和靶都关于 $(x_1, \ldots, x_d)$ 完备，由引理
\ref{algebra-lemma-completion-generalities} 可知该映射是满射。
另一方面，它必为单射，因为否则由引理
\ref{algebra-lemma-one-equation}，$R$ 的维数将 $< d$。
\end{proof}

\begin{lemma}
\label{algebra-lemma-complete-local-Noetherian-domain-finite-over-regular}
设 $(R, \mathfrak m)$ 是 Noether 完备局部整环。
则存在 $R_0 \subset R$，具有下列性质：
% P01-E1001: frozen English says “there exists a R_0 subset R”.
\begin{enumerate}
\item $R_0$ 是正则完备局部环；
\item $R_0 \subset R$ 是有限的，并诱导剩余域上的同构；
\item $R_0$ 同构于 $k[[X_1, \ldots, X_d]]$ 或
$\Lambda[[X_1, \ldots, X_d]]$，其中 $k$ 是域，
$\Lambda$ 是 Cohen 环。
\end{enumerate}
\end{lemma}

\begin{proof}
设 $\Lambda$ 是 $R$ 的系数环。由于 $R$ 是整环，
$\Lambda$ 是域，或者 $\Lambda$ 是 Cohen 环。

\medskip\noindent
情形 I：$\Lambda = k$ 是域。令 $d = \dim(R)$。
选取生成定义理想 $I \subset R$ 的元素
$x_1, \ldots, x_d \in \mathfrak m$。
（见第 \ref{algebra-section-dimension} 节。）
由引理 \ref{algebra-lemma-change-ideal-completion} 可知，$R$
关于 $I$ 也是完备的。考虑把 $X_i$ 映到 $x_i$ 的映射
$R_0 = k[[X_1, \ldots, X_d]] \to R$。
注意，$R_0$ 关于理想 $I_0 = (X_1, \ldots, X_d)$ 完备，
并且 $R/I_0R \cong R/IR$ 在 $k = R_0/I_0$ 上有限
（因为 $\dim(R/I) = 0$，见第 \ref{algebra-section-dimension} 节）。
因此，由引理 \ref{algebra-lemma-finite-over-complete-ring} 可知
$R_0 \to R$ 是有限的。由于 $\dim(R) = \dim(R_0)$，
这蕴含 $R_0 \to R$ 是单射（见引理
\ref{algebra-lemma-integral-dim-up}），引理得证。

\medskip\noindent
情形 II：$\Lambda$ 是 Cohen 环。令 $d + 1 = \dim(R)$。
设 $p > 0$ 是剩余域 $k$ 的特征。由于 $R$ 是整环，
$p$ 是 $R$ 中的非零因子。因此，
$\dim(R/pR) = d$，见引理 \ref{algebra-lemma-one-equation}。
选取 $x_1, \ldots, x_d \in R$，使它们在 $R/pR$
中生成一个定义理想。则
$I = (p, x_1, \ldots, x_d)$ 是 $R$ 的定义理想。
由引理 \ref{algebra-lemma-change-ideal-completion} 可知，$R$
关于 $I$ 也是完备的。考虑把 $X_i$ 映到 $x_i$ 的映射
$R_0 = \Lambda[[X_1, \ldots, X_d]] \to R$。
注意，$R_0$ 关于理想
$I_0 = (p, X_1, \ldots, X_d)$ 完备，并且
$R/I_0R \cong R/IR$ 在 $k = R_0/I_0$ 上有限
（因为 $\dim(R/I) = 0$，见第 \ref{algebra-section-dimension} 节）。
因此，由引理 \ref{algebra-lemma-finite-over-complete-ring} 可知
$R_0 \to R$ 是有限的。由于 $\dim(R) = \dim(R_0)$，
这蕴含 $R_0 \to R$ 是单射（见引理
\ref{algebra-lemma-integral-dim-up}），引理得证。
\end{proof}

\section{日本型环}
\label{algebra-section-japanese}

\noindent
本节开始讨论整闭包的有限性。

\begin{definition}
\label{algebra-definition-N}
\begin{reference}
\cite[Chapter 0, Definition 23.1.1]{EGA}
\end{reference}
设 $R$ 是分式域为 $K$ 的整环。
\begin{enumerate}
\item 若 $R$ 在 $K$ 中的整闭包是有限 $R$-模，则称 $R$ 为
{\it N-1}；
\item 若对任意有限域扩张 $L/K$，$R$ 在 $L$ 中的整闭包
在 $R$ 上有限，则称 $R$ 为 {\it N-2} 或{\it 日本型的}。
\end{enumerate}
\end{definition}

\noindent
这些概念主要用于 Noether 环，不过下面给出一个非 Noether 的例子。

\begin{example}
\label{algebra-example-Japanese-not-Noetherian}
设 $k$ 是域。整环 $R = k[x_1, x_2, x_3, \ldots]$ 是 N-2 的，
但不是 Noether 的。理由如下。假设 $R \subset L$，并且域 $L$
是 $R$ 的分式域的有限扩张。则存在整数 $n$，使得 $L$ 可由有限扩张
$L_0/k(x_1, \ldots, x_n)$ 添加（超越）元素
$x_{n + 1}, x_{n + 2}$ 等而得到。令 $S_0$ 为
$k[x_1, \ldots, x_n]$ 在 $L_0$ 中的整闭包。由下面的命题
\ref{algebra-proposition-ubiquity-nagata}，$S_0$ 在
$k[x_1, \ldots, x_n]$ 上有限。此外，$R$ 在 $L$ 中的整闭包是
$S = S_0[x_{n + 1}, x_{n + 2}, \ldots]$（使用引理
\ref{algebra-lemma-polynomial-domain-normal}），因而在 $R$ 上有限。
同一论证也适用于 $R = \mathbf{Z}[x_1, x_2, x_3, \ldots]$。
\end{example}

\begin{lemma}
\label{algebra-lemma-localize-N}
设 $R$ 是整环。若 $R$ 是 N-1 的，则 $R$ 的任何局部化也是 N-1 的。
N-2 的情形亦然。
\end{lemma}

\begin{proof}
这是因为取整闭包与局部化交换，见引理
\ref{algebra-lemma-integral-closure-localize}。
\end{proof}

\begin{lemma}
\label{algebra-lemma-Japanese-local}
设 $R$ 是整环。设 $f_1, \ldots, f_n \in R$ 生成单位理想。
若每个整环 $R_{f_i}$ 都是 N-1 的，则 $R$ 也是 N-1 的。
N-2 的情形亦然。
\end{lemma}

\begin{proof}
假设 $R_{f_i}$ 是 N-2 的（或 N-1 的）。设 $L$ 是 $R$ 的分式域的
有限扩张（在 N-1 情形中就取该分式域本身）。令 $S$ 为 $R$ 在 $L$
中的整闭包。由引理 \ref{algebra-lemma-integral-closure-localize}，
$S_{f_i}$ 是 $R_{f_i}$ 在 $L$ 中的整闭包。因此按假设，
$S_{f_i}$ 在 $R_{f_i}$ 上有限。于是由引理 \ref{algebra-lemma-cover}，
$S$ 在 $R$ 上有限。
\end{proof}

\begin{lemma}
\label{algebra-lemma-quasi-finite-over-Noetherian-japanese}
设 $R$ 是整环，$R \subset S$ 是整环的拟有限扩张（例如有限扩张）。
假设 $R$ 是 N-2 且 Noether 的。则 $S$ 是 N-2 的。
\end{lemma}

\begin{proof}
设 $L/K$ 为所诱导的分式域扩张。注意，这是有限域扩张（例如，对纤维
$S \otimes_R K$ 应用引理 \ref{algebra-lemma-isolated-point-fibre} (2)，
并使用拟有限环映射的定义）。
令 $S'$ 为 $R$ 在 $S$ 中的整闭包。则 $S'$ 包含于 $R$ 在 $L$
中的整闭包；后者按假设在 $R$ 上有限。由于 $R$ 是 Noether 的，
这蕴含 $S'$ 在 $R$ 上有限。由引理
\ref{algebra-lemma-quasi-finite-open-integral-closure}，存在元素
$g_1, \ldots, g_n \in S'$，使得 $S'_{g_i} \cong S_{g_i}$，
并且 $g_1, \ldots, g_n$ 在 $S$ 中生成单位理想。因此，由引理
\ref{algebra-lemma-localize-N} 和引理 \ref{algebra-lemma-Japanese-local}，
只须证明 $S'$ 是 N-2 的。于是已经约化到 $S$ 在 $R$ 上有限的情形。

\medskip\noindent
假设 $R \subset S$ 满足引理的各项假设，并进一步假设 $S$ 在 $R$
上有限。设 $M$ 是 $S$ 的分式域的有限域扩张。则 $M$ 也是 $K$
的有限域扩张，从而 $R$ 在 $M$ 中的整闭包 $T$ 在 $R$ 上有限。
由引理 \ref{algebra-lemma-integral-closure-transitive}，$T$ 也是 $S$ 在 $M$
中的整闭包；再由引理 \ref{algebra-lemma-integral-permanence} 即得结论。
\end{proof}

\begin{lemma}
\label{algebra-lemma-Laurent-ring-N-1}
设 $R$ 是 Noether 整环。若 $R[z, z^{-1}]$ 是 N-1 的，
则 $R$ 也是 N-1 的。
\end{lemma}

\begin{proof}
设 $R'$ 为 $R$ 在其分式域 $K$ 中的整闭包。设 $S'$ 为
$R[z, z^{-1}]$ 在其分式域中的整闭包。显然 $R' \subset S'$。
由于 $K[z, z^{-1}]$ 是正规整环，可知 $S' \subset K[z, z^{-1}]$。
假设 $f_1, \ldots, f_n \in S'$ 作为 $R[z, z^{-1}]$-模生成 $S'$。
写成 $f_i = \sum a_{ij}z^j$（有限和），其中 $a_{ij} \in K$。
对任意 $x \in R'$，可以写成
$$
x = \sum h_i f_i
$$
其中 $h_i \in R[z, z^{-1}]$。因此，$R'$ 包含于有限 $R$-子模
$\sum Ra_{ij} \subset K$。由于 $R$ 是 Noether 的，可知 $R'$
是有限 $R$-模。
\end{proof}

\begin{lemma}
\label{algebra-lemma-finite-extension-N-2}
设 $R$ 是 Noether 整环，$R \subset S$ 是整环的有限扩张。
若 $S$ 是 N-1 的，则 $R$ 也是 N-1 的。
若 $S$ 是 N-2 的，则 $R$ 也是 N-2 的。
\end{lemma}

\begin{proof}
从略。（提示：$R$ 在扩域中的整闭包包含于 $S$ 在扩域中的整闭包。）
\end{proof}

\begin{lemma}
\label{algebra-lemma-Noetherian-normal-domain-finite-separable-extension}
设 $R$ 是分式域为 $K$ 的 Noether 正规整环。
设 $L/K$ 是有限可分域扩张。则 $R$ 在 $L$ 中的整闭包
在 $R$ 上有限。
\end{lemma}

\begin{proof}
考虑迹配对（《域》，定义 \ref{fields-definition-trace-pairing}）
$$
L \times L \longrightarrow K,
\quad (x, y) \longmapsto \langle x, y\rangle := \text{Trace}_{L/K}(xy).
$$
由于 $L/K$ 可分，该配对非退化（《域》，引理
\ref{fields-lemma-separable-trace-pairing}）。此外，若 $x \in L$
整于 $R$，则 $\text{Trace}_{L/K}(x)$ 属于 $R$。这是因为 $x$ 在 $K$
上的最小多项式的系数属于 $R$（引理
\ref{algebra-lemma-minimal-polynomial-normal-domain}），并且
$\text{Trace}_{L/K}(x)$ 是其中某个系数的整数倍（《域》，引理
\ref{fields-lemma-trace-and-norm-from-minimal-polynomial}）。
选取整于 $R$ 且构成 $L$ 的一组 $K$-基的元素
$x_1, \ldots, x_n \in L$。则整闭包 $S \subset L$ 包含于 $R$-模
$$
M = \{y \in L \mid \langle x_i, y\rangle \in R, \ i = 1, \ldots, n\}
$$
中。由线性代数，$M \cong R^{\oplus n}$ 是 $R$-模。因此，由于
$R$ 是 Noether 的，$S \subset R^{\oplus n}$ 是有限生成 $R$-模。
\end{proof}

\begin{example}
\label{algebra-example-bad-invariants}
若环不是 Noether 的，引理
\ref{algebra-lemma-Noetherian-normal-domain-finite-separable-extension}
不再成立。例如，考虑 $G = \{+1, -1\}$ 在
$A = \mathbf{C}[x_1, x_2, x_3, \ldots]$ 上的作用，其中 $-1$
把 $x_i$ 映到 $-x_i$。不变量环 $R = A^G$ 是由所有 $x_ix_j$
生成的 $\mathbf{C}$-代数。因此，$R \subset A$ 不是有限的。
但是 $R$ 是正规整环，其分式域 $K = L^G$ 是 $A$ 的分式域 $L$
中的 $G$-不变量；而且显然，$A$ 是 $R$ 在 $L$ 中的整闭包。
% P01-E1002: frozen English lacks a separator after K = L^G.
\end{example}

\noindent
对于纯不可分扩张，下面的引理有时可以替代引理
\ref{algebra-lemma-Noetherian-normal-domain-finite-separable-extension}。

\begin{lemma}
\label{algebra-lemma-Noetherian-normal-domain-insep-extension}
设 $R$ 是分式域为 $K$、特征为 $p > 0$ 的 Noether 正规整环。
设 $a \in K$，并且存在导子 $D : R \to R$ 使得
$D(a) \not = 0$。则 $R$ 在 $L = K[x]/(x^p - a)$ 中的整闭包
在 $R$ 上有限。
\end{lemma}

\begin{proof}
对某个 $f \in R$，把 $x$ 替换为 $fx$、把 $a$ 替换为 $f^pa$ 后，
可以假设 $a \in R$。于是也有 $D(a) \in R$。我们对
$i \leq p - 1$ 归纳证明：如果
$$
y = a_0 + a_1x + \ldots + a_i x^i,\quad a_j \in K
$$
整于 $R$，则 $D(a)^i a_j \in R$。因此，整闭包包含于以
$D(a)^{-p + 1}x^j$（$j = 0, \ldots, p - 1$）为基的有限
$R$-模中。由于 $R$ 是 Noether 的，引理得证。

\medskip\noindent
若 $i = 0$，则 $y = a_0$ 整于 $R$ 当且仅当 $a_0 \in R$，
所以断言成立。假设断言对某个 $i < p - 1$ 成立，并假设
$$
y = a_0 + a_1x + \ldots + a_{i + 1} x^{i + 1},\quad a_j \in K
$$
整于 $R$。则
$$
y^p = a_0^p + a_1^p a + \ldots + a_{i + 1}^pa^{i + 1}
$$
属于 $R$（因为它属于 $K$ 且整于 $R$）。应用 $D$，得到
$$
(a_1^p + 2a_2^p a + \ldots + (i + 1)a_{i + 1}^p a^i)D(a)
$$
属于 $R$。因此
$$
D(a)a_1 + 2D(a) a_2 x + \ldots + (i + 1)D(a) a_{i + 1} x^i
$$
整于 $R$。由归纳假设，对 $j = 1, \ldots, i + 1$，有
$D(a)^{i + 1}a_j \in R$。（这里用到 $1, \ldots, i + 1$
均可逆。）进而 $D(a)^{i + 1}a_0$ 也属于 $R$，因为它是
$D(a)^{i + 1}y$ 与 $\sum_{j > 0} D(a)^{i + 1}a_jx^j$ 之差，
而二者都整于 $R$（这是因为 $a \in R$，所以 $x$ 整于 $R$）。
\end{proof}

\begin{lemma}
\label{algebra-lemma-domain-char-zero-N-1-2}
分式域特征为零的 Noether 整环是 N-1 的，当且仅当它是 N-2 的
（即日本型的）。
\end{lemma}

\begin{proof}
这由引理 \ref{algebra-lemma-Noetherian-normal-domain-finite-separable-extension}
立即得出，因为特征为零的每个域扩张都是可分的。
\end{proof}

\begin{lemma}
\label{algebra-lemma-domain-char-p-N-1-2}
设 $R$ 是分式域为 $K$、特征为 $p > 0$ 的 Noether 整环。
则 $R$ 是 N-2 的，当且仅当对每个有限纯不可分扩张 $L/K$，
$R$ 在 $L$ 中的整闭包在 $R$ 上有限。
\end{lemma}

\begin{proof}
假设在 $K$ 的每个有限纯不可分域扩张中，$R$ 的整闭包都是有限的。
设 $L/K$ 是任意有限扩张。我们必须证明 $R$ 在 $L$ 中的整闭包
在 $R$ 上有限。选取包含 $L$ 的有限正规域扩张 $M/K$。
由于 $R$ 是 Noether 的，只须证明 $R$ 在 $M$ 中的整闭包在 $R$
上有限。由《域》，引理 \ref{fields-lemma-normal-case}，存在子扩张
$M/M_{insep}/K$，使得 $M_{insep}/K$ 纯不可分，而
$M/M_{insep}$ 可分。按假设，$R$ 在 $M_{insep}$ 中的整闭包
$R'$ 在 $R$ 上有限。由引理
\ref{algebra-lemma-Noetherian-normal-domain-finite-separable-extension}，
$R'$ 在 $M$ 中的整闭包 $R''$ 在 $R'$ 上有限。再由引理
\ref{algebra-lemma-finite-transitive}，$R''$ 在 $R$ 上有限。
由于 $R''$ 也就是 $R$ 在 $M$ 中的整闭包（见引理
\ref{algebra-lemma-integral-closure-transitive}），结论得证。
\end{proof}

\begin{lemma}
\label{algebra-lemma-polynomial-ring-N-2}
设 $R$ 是 Noether 整环。
若 $R$ 是 N-1 的，则 $R[x]$ 是 N-1 的。
若 $R$ 是 N-2 的，则 $R[x]$ 是 N-2 的。
\end{lemma}

\begin{proof}
假设 $R$ 是 N-1 的。设 $R'$ 为 $R$ 的整闭包；它在 $R$
上有限。因此，$R'[x]$ 也在 $R[x]$ 上有限。环 $R'[x]$
是正规的（见引理 \ref{algebra-lemma-polynomial-domain-normal}），
因而是 N-1 的。这证明了第一个断言。

\medskip\noindent
对于第二个断言，由引理 \ref{algebra-lemma-finite-extension-N-2}，
只须证明 $R'[x]$ 是 N-2 的。换言之，可以并且确实假设 $R$
是正规 N-2 整环。在特征零情形中，由引理
\ref{algebra-lemma-domain-char-zero-N-1-2} 即得结论。
在特征 $p > 0$ 情形中，我们必须证明：对任意有限纯不可分域扩张
$L/K(x)$，$R[x]$ 在该扩域中的整闭包都是有限的，其中 $K$ 是 $R$
的分式域。
% P01-E1003: frozen English says “extension of L/K(x)”.
存在有限纯不可分域扩张 $L'/K$ 和 $q = p^e$，使得
$L \subset L'(x^{1/q})$；一些细节从略。由于 $R[x]$ 是 Noether 的，
只须证明 $R[x]$ 在 $L'(x^{1/q})$ 中的整闭包在 $R[x]$ 上有限。
该整闭包等于 $R'[x^{1/q}]$，其中 $R \subset R' \subset L'$
是 $R$ 在 $L'$ 中的整闭包。
% P01-E1004: frozen English omits “where” before the R' apposition.
由于 $R$ 是 N-2 的，可知 $R'$ 在 $R$ 上有限，故
$R'[x^{1/q}]$ 在 $R[x]$ 上有限。
\end{proof}

\begin{lemma}
\label{algebra-lemma-openness-normal-locus}
设 $R$ 是 Noether 整环。若存在非零 $f \in R$ 使得 $R_f$
正规，则
$$
U = \{\mathfrak p \in \Spec(R) \mid R_{\mathfrak p} \text{ is normal}\}
$$
在 $\Spec(R)$ 中是开集。
\end{lemma}

\begin{proof}
显然，标准开集 $D(f)$ 包含于 $U$。由 Serre 判据，引理
\ref{algebra-lemma-criterion-normal}，$\mathfrak p \not \in U$ 蕴含：
对某个 $\mathfrak q \subset \mathfrak p$，有下列情形之一：
\begin{enumerate}
\item 情形 I：$\text{depth}(R_{\mathfrak q}) < 2$ 且
$\dim(R_{\mathfrak q}) \geq 2$；
\item 情形 II：$R_{\mathfrak q}$ 不正则且
$\dim(R_{\mathfrak q}) = 1$。
\end{enumerate}
特别地，这也意味着 $R_{\mathfrak q}$ 不正规，因而
$f \in \mathfrak q$。在情形 I 中，有
$\text{depth}(R_{\mathfrak q}) =
\text{depth}(R_{\mathfrak q}/fR_{\mathfrak q}) + 1$。
因此，这样的素理想 $\mathfrak q$ 恰是 $R/fR$ 的嵌入相伴素理想。
在情形 II 中，$\mathfrak q$ 是 $R/fR$ 的高度为 1 的相伴素理想。
所以这样的素理想 $\mathfrak q$ 构成有限集合 $E$（见引理
\ref{algebra-lemma-finite-ass}），并且
$$
\Spec(R) \setminus U
=
\bigcup\nolimits_{\mathfrak q \in E} V(\mathfrak q)
$$
如所需。
\end{proof}

\begin{lemma}
\label{algebra-lemma-characterize-N-1}
设 $R$ 是 Noether 整环。则 $R$ 是 N-1 的，当且仅当下列两个条件成立：
\begin{enumerate}
\item 存在非零 $f \in R$，使得 $R_f$ 正规；
\item 对每个极大理想 $\mathfrak m \subset R$，局部环
$R_{\mathfrak m}$ 都是 N-1 的。
\end{enumerate}
\end{lemma}

\begin{proof}
先假设 $R$ 是 N-1 的。设 $R'$ 为 $R$ 在其分式域 $K$ 中的整闭包。
按假设，可以找到 $R'$ 中的 $x_1, \ldots, x_n$，使它们作为
$R$-模生成 $R'$。由于 $R' \subset K$，可以找到非零
$f_i \in R$，使得 $f_i x_i \in R$。于是 $R_f \cong R'_f$，其中
$f = f_1 \ldots f_n$。所以 $R_f$ 正规，从而得到 (1)。
第 (2) 项由引理 \ref{algebra-lemma-localize-N} 得出。

\medskip\noindent
假设 (1) 和 (2)。设 $K$ 是 $R$ 的分式域。假设
$R \subset R' \subset K$ 是包含于 $K$ 的有限 $R$-扩张。
注意，$R_f = R'_f$，因为 $R_f$ 已经正规。因此，由引理
\ref{algebra-lemma-openness-normal-locus}，使 $R'_{\mathfrak p'}$ 不正规的
素理想 $\mathfrak p' \in \Spec(R')$ 所构成的集合在
$\Spec(R')$ 中闭。由于 $\Spec(R') \to \Spec(R)$ 是闭映射，
该集合在 $\Spec(R)$ 中的像是闭集。对这样的环 $R'$，把这个像记作
$Z_{R'} \subset \Spec(R)$。
% P01-E1005: frozen English has malformed word order in the Z_{R'} definition.

\medskip\noindent
选取极大理想 $\mathfrak m \subset R$。设
$R_{\mathfrak m} \subset R_{\mathfrak m}'$ 为该局部环在 $K$
中的整闭包。按假设，这是有限环扩张。由引理
\ref{algebra-lemma-integral-closure-localize}，可以找到有限多个整于 $R$
的元素 $x_1, \ldots, x_n \in K$，使得 $x_1, \ldots, x_n$
在 $R_{\mathfrak m}$ 上生成 $R_{\mathfrak m}'$。
令 $R' = R[x_1, \ldots, x_n] \subset K$。按此选取，显然有
$\mathfrak m \not \in Z_{R'}$。

\medskip\noindent
由于 $\Spec(R)$ 拟紧，上述论证表明，可以找到有限多个
$R \subset R'_i \subset K$，使得 $\bigcap Z_{R'_i} = \emptyset$。
令 $R'$ 为 $K$ 中由所有这些环生成的子环。它在 $R$ 上有限。
而且 $Z_{R'} = \emptyset$。事实上，每个素理想 $\mathfrak p'$
都位于某个素理想 $\mathfrak p'_i$ 之上，而
$(R'_i)_{\mathfrak p'_i}$ 正规。这蕴含
$R'_{\mathfrak p'} = (R'_i)_{\mathfrak p'_i}$ 也正规。
因此 $R'$ 正规；换言之，$R'$ 是 $R$ 在 $K$ 中的整闭包。
\end{proof}

\begin{lemma}[Tate]
\label{algebra-lemma-tate-japanese}
\begin{reference}
\cite[Theorem 23.1.3]{EGA}
\end{reference}
设 $R$ 是环，$x \in R$。假设
\begin{enumerate}
\item $R$ 是 Noether 正规整环；
\item $R/xR$ 是整环且为 N-2；
\item $R \cong \lim_n R/x^nR$ 关于 $x$ 完备。
\end{enumerate}
则 $R$ 是 N-2 的。
\end{lemma}

\begin{proof}
可以假设 $x \not = 0$，否则引理显然。设 $K$ 为 $R$ 的分式域。
若 $K$ 的特征为零，则由 (1) 和引理
\ref{algebra-lemma-domain-char-zero-N-1-2} 即得结论。因此，可以假设
$K$ 的特征为 $p > 0$，并应用引理
\ref{algebra-lemma-domain-char-p-N-1-2}。于是，给定有限纯不可分域扩张
$L/K$，必须证明 $R$ 在 $L$ 中的整闭包 $S$ 在 $R$ 上有限。

\medskip\noindent
设 $q$ 是 $p$ 的幂，并使得 $L^q \subset K$。必要时扩大 $L$，
可以假设存在 $y \in L$ 使得 $y^q = x$。由于 $R \to S$
诱导谱上的同胚（见引理 \ref{algebra-lemma-p-ring-map}），存在唯一素理想
$\mathfrak q \subset S$ 位于素理想 $\mathfrak p = xR$ 之上。
显然
$$
\mathfrak q = \{f \in S \mid f^q \in \mathfrak p\} = yS
$$
因为 $y^q = x$。由引理 \ref{algebra-lemma-characterize-dvr}，
$R_{\mathfrak p}$ 是离散赋值环。于是由 Krull--Akizuki
（引理 \ref{algebra-lemma-krull-akizuki}），$S_{\mathfrak q}$ 是
Noether 的。再次应用引理 \ref{algebra-lemma-characterize-dvr}，
可知 $S_{\mathfrak q}$ 是离散赋值环。由引理
\ref{algebra-lemma-finite-extension-residue-fields-dimension-1}，
$\kappa(\mathfrak q)/\kappa(\mathfrak p)$ 是有限域扩张。
因此，按假设 (2)，$R/xR$ 的整闭包 $S' \subset \kappa(\mathfrak q)$
在 $R/xR$ 上有限。
由于 $S/yS \subset S'$，这蕴含 $S/yS$ 在 $R$ 上有限。
注意，$S/y^nS$ 有有限滤过，其子商为模
$y^iS/y^{i + 1}S \cong S/yS$。因此，每个 $S/y^nS$ 都在 $R$
上有限。特别地，$S/xS$ 在 $R$ 上有限。此外，显然
$\bigcap x^nS = (0)$，因为交中元素的 $q$ 次幂属于
$\bigcap x^nR = (0)$（引理
\ref{algebra-lemma-intersect-powers-ideal-module-zero}）。
所以可以应用引理 \ref{algebra-lemma-finite-over-complete-ring}，得出
$S$ 在 $R$ 上有限，结论得证。
\end{proof}

\begin{lemma}
\label{algebra-lemma-power-series-over-N-2}
设 $R$ 是环。若 $R$ 是 Noether 的整环且为 N-2，
则 $R[[x]]$ 也是如此。
\end{lemma}

\begin{proof}
由引理 \ref{algebra-lemma-Noetherian-power-series}，$R[[x]]$ 是 Noether 的。
设 $R' \supset R$ 为 $R$ 在其分式域中的整闭包。由于 $R$ 是 N-2 的，
这个整闭包在 $R$ 上有限。因此，$R'[[x]]$ 在 $R[[x]]$ 上有限。
由引理 \ref{algebra-lemma-power-series-over-Noetherian-normal-domain}，
$R'[[x]]$ 是正规整环。对元素 $x \in R'[[x]]$ 应用引理
\ref{algebra-lemma-tate-japanese}，可知 $R'[[x]]$ 是 N-2 的。
再由引理 \ref{algebra-lemma-finite-extension-N-2}，$R[[x]]$ 是 N-2 的。
\end{proof}

\section{Nagata 环}
\label{algebra-section-nagata}

\noindent
定义如下。

\begin{definition}
\label{algebra-definition-nagata}
设 $R$ 是环。
\begin{enumerate}
\item 如果对任意有限型环映射 $R \to S$，只要 $S$ 是整环，
就有 $S$ 为 N-2（即日本型的），则称 $R$ 是{\it 广泛日本型的}。
\item 如果 $R$ 是 Noether 的，并且对每个素理想 $\mathfrak p$，
环 $R/\mathfrak p$ 都是 N-2 的，则称 $R$ 是{\it Nagata 环}。
\end{enumerate}
\end{definition}

\noindent
显然，Noether 的广泛日本型环是 Nagata 环。我们的目标是证明
Nagata 环是广泛日本型的。这一点丝毫不显然，需要做一些工作。
不过先给出一个有用的引理。

\begin{lemma}
\label{algebra-lemma-nagata-in-reduced-finite-type-finite-integral-closure}
设 $R$ 是 Nagata 环。设 $R \to S$ 本质上为有限型，且 $S$ 约化。
则 $R$ 在 $S$ 中的整闭包 $A$ 在 $R$ 上有限。
\end{lemma}

\begin{proof}
由于 $S$ 本质上在 $R$ 上有限型，它是 Noether 的，并且只有有限多个
极小素理想 $\mathfrak q_1, \ldots, \mathfrak q_m$；见引理
\ref{algebra-lemma-Noetherian-irreducible-components}。由于 $S$ 约化，
有 $S \subset \prod S_{\mathfrak q_i}$，且每个
$S_{\mathfrak q_i} = K_i$ 都是域；见引理
\ref{algebra-lemma-total-ring-fractions-no-embedded-points}
和 \ref{algebra-lemma-minimal-prime-reduced-ring}。只需证明 $R$ 在每个 $K_i$
中的整闭包 $A_i'$ 在 $R$ 上有限。之所以如此，是因为 $R$ 是 Noether 的，
且 $A \subset \prod A_i'$。设 $\mathfrak p_i \subset R$ 是与
$\mathfrak q_i$ 对应的 $R$ 的素理想。由于 $S$ 本质上在 $R$ 上有限型，
可知 $K_i = S_{\mathfrak q_i} = \kappa(\mathfrak q_i)$ 是
$\kappa(\mathfrak p_i)$ 的有限生成域扩张。因此，
$\kappa(\mathfrak p_i)$ 在 $K_i$ 中的代数闭包 $L_i$ 是
$\kappa(\mathfrak p_i)$ 的有限扩张；见《域》，引理
\ref{fields-lemma-algebraic-closure-in-finitely-generated}。
显然，$A_i'$ 是 $R/\mathfrak p_i$ 在 $L_i$ 中的整闭包，
所以由 Nagata 环的定义即得结论。
\end{proof}

\begin{lemma}
\label{algebra-lemma-check-universally-japanese}
设 $R$ 是环。为检验 $R$ 是广泛日本型的，只需证明如下命题：
若 $R \to S$ 为有限型且 $S$ 是整环，则 $S$ 是 N-1 的。
\end{lemma}

\begin{proof}
假设引理中的条件成立。设 $R \to S$ 是有限型环映射，且 $S$ 是整环。
设 $L$ 是 $S$ 的分式域的有限扩张。则存在有限环扩张
$S \subset S' \subset L$，使 $L$ 为 $S'$ 的分式域。
由假设，$S'$ 是 N-1 的，因此 $S'$ 在 $L$ 中的整闭包 $S''$
在 $S'$ 上有限。于是 $S''$ 在 $S$ 上有限（引理
\ref{algebra-lemma-finite-transitive}），并且 $S''$ 是 $S$ 在 $L$ 中的整闭包
（引理 \ref{algebra-lemma-integral-closure-transitive}）。所以 $R$ 是广泛日本型的。
\end{proof}

\begin{lemma}
\label{algebra-lemma-universally-japanese}
若 $R$ 是广泛日本型的，则任何本质上在 $R$ 上有限型的代数也都是
广泛日本型的。
\end{lemma}

\begin{proof}
在代数在 $R$ 上有限型的情形，结论直接来自定义。一般情形由应用引理
\ref{algebra-lemma-localize-N} 得到。
\end{proof}

\begin{lemma}
\label{algebra-lemma-quasi-finite-over-nagata}
设 $R$ 是 Nagata 环。若 $R \to S$ 是拟有限环映射（例如有限环映射），
则 $S$ 也是 Nagata 环。
\end{lemma}

\begin{proof}
先注意到 $S$ 是 Noether 的，因为 $R$ 是 Noether 的，而拟有限环映射
是有限型的。设 $\mathfrak q \subset S$ 是素理想，并令
$\mathfrak p = R \cap \mathfrak q$。则
$R/\mathfrak p \subset S/\mathfrak q$ 是拟有限的，因而由引理
\ref{algebra-lemma-quasi-finite-over-Noetherian-japanese} 可知
$S/\mathfrak q$ 是 N-2 的，如所求。
\end{proof}

\begin{lemma}
\label{algebra-lemma-nagata-localize}
Nagata 环的局部化仍是 Nagata 环。
\end{lemma}

\begin{proof}
直接由引理 \ref{algebra-lemma-localize-N} 得到。
\end{proof}

\begin{lemma}
\label{algebra-lemma-nagata-local}
设 $R$ 是环。设 $f_1, \ldots, f_n \in R$ 生成单位理想。
\begin{enumerate}
\item  如果每个 $R_{f_i}$ 都是广泛日本型的，则 $R$ 也是如此。
\item  如果每个 $R_{f_i}$ 都是 Nagata 环，则 $R$ 也是 Nagata 环。
\end{enumerate}
\end{lemma}

\begin{proof}
设 $\varphi : R \to S$ 是有限型环映射，且 $S$ 是整环。
则 $\varphi(f_1), \ldots, \varphi(f_n)$ 生成 $S$ 的单位理想。
因此，若每个 $S_{f_i} = S_{\varphi(f_i)}$ 都是 N-1 的，
则 $S$ 也是 N-1 的；见引理 \ref{algebra-lemma-Japanese-local}。这证明了 (1)。

\medskip\noindent
若每个 $R_{f_i}$ 都是 Nagata 环，则每个 $R_{f_i}$ 都是 Noether 的，
因而 $R$ 是 Noether 的；见引理 \ref{algebra-lemma-cover}。再设
$\mathfrak p \subset R$ 是素理想，则每个
$R_{f_i}/\mathfrak pR_{f_i} = (R/\mathfrak p)_{f_i}$ 都是 N-2 的，
所以由引理 \ref{algebra-lemma-Japanese-local} 可知 $R/\mathfrak p$ 是 N-2 的。
这证明了 (2)。
\end{proof}

\begin{lemma}
\label{algebra-lemma-Noetherian-complete-local-Nagata}
Noether 完备局部环是 Nagata 环。
\end{lemma}

\begin{proof}
设 $R$ 是完备的 Noether 局部环。设 $\mathfrak p \subset R$ 是素理想。
则 $R/\mathfrak p$ 也是完备的 Noether 局部环；见引理
\ref{algebra-lemma-quotient-complete-local}。因此，只需证明完备的 Noether
局部整环 $R$ 是 N-2 的。由引理
\ref{algebra-lemma-quasi-finite-over-Noetherian-japanese}
和 \ref{algebra-lemma-complete-local-Noetherian-domain-finite-over-regular}，
可化归到 $R = k[[X_1, \ldots, X_d]]$，其中 $k$ 是域；或者化归到
$R = \Lambda[[X_1, \ldots, X_d]]$，其中 $\Lambda$ 是 Cohen 环。

\medskip\noindent
在 $k[[X_1, \ldots, X_d]]$ 的情形，由引理
\ref{algebra-lemma-power-series-over-N-2} 可化归到域为 N-2 的命题，
这是显然的。在 $\Lambda[[X_1, \ldots, X_d]]$ 的情形，
可化归到 Cohen 环 $\Lambda$ 是 N-2 的命题。再次对
$x = p \in \Lambda$ 应用引理 \ref{algebra-lemma-tate-japanese}，
又一次化归到域的情形。因此结论成立。
\end{proof}

\begin{definition}
\label{algebra-definition-analytically-unramified}
设 $(R, \mathfrak m)$ 是 Noether 局部环。如果它的完备化
$R^\wedge = \lim_n R/\mathfrak m^n$ 是约化的，则称 $R$
是{\it 解析非分歧的}。若素理想 $\mathfrak p \subset R$ 使
$R/\mathfrak p$ 解析非分歧，则称该素理想是
{\it 解析非分歧的}。
\end{definition}

\noindent
至此，对任意 Noether 局部环 $R$，我们已经知道以下事实成立：
映射 $R \to R^\wedge$ 是忠实平坦局部环同态（引理
\ref{algebra-lemma-completion-faithfully-flat}）。完备化 $R^\wedge$
是 Noether 的（引理 \ref{algebra-lemma-completion-Noetherian}），
并且是完备的（引理 \ref{algebra-lemma-completion-complete}）。因此，
完备化 $R^\wedge$ 是 Nagata 环（引理
\ref{algebra-lemma-Noetherian-complete-local-Nagata}）。此外，我们已经在节
\ref{algebra-section-cohen-structure-theorem} 中看到，$R^\wedge$ 是正则局部环的商
（定理 \ref{algebra-theorem-cohen-structure-theorem}），从而是普遍链状的
（评注 \ref{algebra-remark-Noetherian-complete-local-ring-universally-catenary}）。

\begin{lemma}
\label{algebra-lemma-analytically-unramified-easy}
设 $(R, \mathfrak m)$ 是 Noether 局部环。
\begin{enumerate}
\item 若 $R$ 解析非分歧，则 $R$ 约化。
\item 若 $R$ 解析非分歧，则 $R$ 的每个极小素理想都解析非分歧。
\item 若 $R$ 约化，其极小素理想为
$\mathfrak q_1, \ldots, \mathfrak q_t$，并且每个 $\mathfrak q_i$
都解析非分歧，则 $R$ 解析非分歧。
\item 若 $R$ 解析非分歧，则 $R$ 在其全分式环 $Q(R)$ 中的整闭包
在 $R$ 上有限。
\item 若 $R$ 是整环且解析非分歧，则 $R$ 是 N-1 的。
\end{enumerate}
\end{lemma}

\begin{proof}
在本证明中，我们将使用紧接在定义
\ref{algebra-definition-analytically-unramified} 后的评注。由于
$R \to R^\wedge$ 是忠实平坦局部环同态，它是单射，因此 (1) 成立。

\medskip\noindent
设 $\mathfrak q$ 是 $R$ 的极小素理想，并假设 $R$ 解析非分歧。
则 $\mathfrak q$ 是 $R$ 的相伴素理想（见命题
\ref{algebra-proposition-minimal-primes-associated-primes}）。因此存在
$f \in R$，使得 $\{x \in R \mid fx = 0\} = \mathfrak q$。
注意 $(R/\mathfrak q)^\wedge = R^\wedge/\mathfrak q^\wedge$，
并且 $\{x \in R^\wedge \mid fx = 0\} = \mathfrak q^\wedge$，
因为完备化是正合的（引理 \ref{algebra-lemma-completion-flat}）。若
$x \in R^\wedge$ 满足 $x^2 \in \mathfrak q^\wedge$，则
$fx^2 = 0$，从而 $(fx)^2 = 0$，从而 $fx = 0$，从而
$x \in \mathfrak q^\wedge$。所以 $\mathfrak q$ 解析非分歧，
(2) 成立。

\medskip\noindent
假设 $R$ 约化，其极小素理想为
$\mathfrak q_1, \ldots, \mathfrak q_t$，且每个 $\mathfrak q_i$
都解析非分歧。则映射
$R \to R/\mathfrak q_1 \times \ldots \times R/\mathfrak q_t$ 是单射。
由于完备化正合（见引理 \ref{algebra-lemma-completion-flat}），可知
$R^\wedge \subset (R/\mathfrak q_1)^\wedge \times \ldots \times
(R/\mathfrak q_t)^\wedge$。因此 (3) 显然。

\medskip\noindent
假设 $R$ 解析非分歧。设 $\mathfrak p_1, \ldots, \mathfrak p_s$
为 $R^\wedge$ 的极小素理想。则
$$
Q(R^\wedge) =
R^\wedge_{\mathfrak p_1} \times \ldots \times R^\wedge_{\mathfrak p_s}
$$
且每个 $R^\wedge_{\mathfrak p_i}$ 都是域，因为 $R^\wedge$ 约化
（见引理 \ref{algebra-lemma-total-ring-fractions-no-embedded-points}）。
因此，$R^\wedge$ 在 $Q(R^\wedge)$ 中的整闭包 $S$ 等于
$S = S_1 \times \ldots \times S_s$，其中 $S_i$ 是
$R^\wedge/\mathfrak p_i$ 在其分式域中的整闭包。由引理
\ref{algebra-lemma-Noetherian-complete-local-Nagata}，环 $S_i$ 在
$R^\wedge/\mathfrak p_i$ 上有限。因此，$S$ 在 $R^\wedge$ 上有限。
% P01-E1006
用 $R'$ 表示 $R$ 在 $Q(R)$ 中的整闭包。由于
$R \to R^\wedge$ 平坦，可知
$R' \otimes_R R^\wedge \subset Q(R) \otimes_R R^\wedge \subset Q(R^\wedge)$。
此外，$R' \otimes_R R^\wedge$ 整于 $R^\wedge$
（引理 \ref{algebra-lemma-base-change-integral}）。
% P01-E1007
所以 $R' \otimes_R R^\wedge \subset S$ 是 $R^\wedge$-子模。
由于 $R^\wedge$ 是 Noether 的，它是有限 $R^\wedge$-模。
因此可取 $f_1, \ldots, f_n \in R'$，使
$R' \otimes_R R^\wedge$ 作为 $R^\wedge$-模由元素
$f_i \otimes 1$ 生成。由忠实平坦性可知，$R'$ 作为 $R$-模由
$f_1, \ldots, f_n$ 生成。这证明了 (4)。

\medskip\noindent
(5) 是 (4) 的特例。
\end{proof}

\begin{lemma}
\label{algebra-lemma-codimension-1-analytically-unramified}
设 $R$ 是 Noether 局部环。设 $\mathfrak p \subset R$ 是素理想。假设
\begin{enumerate}
\item $R_{\mathfrak p}$ 是离散赋值环；并且
\item $\mathfrak p$ 解析非分歧。
\end{enumerate}
则对 $R^\wedge/\mathfrak pR^\wedge$ 的任意相伴素理想
$\mathfrak q$，局部环 $(R^\wedge)_{\mathfrak q}$ 是离散赋值环。
\end{lemma}

\begin{proof}
假设 (2) 表明 $R^\wedge/\mathfrak pR^\wedge$ 是约化环。
因此，$R^\wedge/\mathfrak pR^\wedge$ 的相伴素理想
$\mathfrak q \subset R^\wedge$，等价于 $\mathfrak pR^\wedge$
之上的极小素理想。特别地，$(R^\wedge)_{\mathfrak q}$ 的极大理想为
$\mathfrak p(R^\wedge)_{\mathfrak q}$。选取 $x \in R$，使得
$xR_{\mathfrak p} = \mathfrak pR_{\mathfrak p}$。由上可知，
$x \in (R^\wedge)_{\mathfrak q}$ 生成极大理想。由于
$R \to R^\wedge$ 忠实平坦，$x$ 是
$(R^\wedge)_{\mathfrak q}$ 中的非零因子。因此结论成立。
\end{proof}

\begin{lemma}
\label{algebra-lemma-criterion-analytically-unramified}
设 $(R, \mathfrak m)$ 是 Noether 局部整环。设 $x \in \mathfrak m$。假设
\begin{enumerate}
\item $x \not = 0$；
\item $R/xR$ 没有嵌入素理想；并且
\item 对 $R/xR$ 的每个相伴素理想 $\mathfrak p \subset R$，都有
\begin{enumerate}
\item 局部环 $R_{\mathfrak p}$ 是正则的；并且
\item $\mathfrak p$ 解析非分歧。
\end{enumerate}
\end{enumerate}
则 $R$ 解析非分歧。
\end{lemma}

\begin{proof}
设 $\mathfrak p_1, \ldots, \mathfrak p_t$ 为 $R$-模 $R/xR$
的相伴素理想。由于 $R/xR$ 没有嵌入素理想，可知每个
$\mathfrak p_i$ 都具有高度 $1$，并且是 $(x)$ 之上的极小素理想。
对每个 $i$，设
$\mathfrak q_{i1}, \ldots, \mathfrak q_{is_i}$ 为 $R^\wedge$-模
$R^\wedge/\mathfrak p_iR^\wedge$ 的相伴素理想。由引理
\ref{algebra-lemma-codimension-1-analytically-unramified}，
$(R^\wedge)_{\mathfrak q_{ij}}$ 是正则的。由引理
\ref{algebra-lemma-bourbaki}，有
$$
\text{Ass}_{R^\wedge}(R^\wedge/xR^\wedge)
=
\bigcup\nolimits_{\mathfrak p \in \text{Ass}_R(R/xR)}
\text{Ass}_{R^\wedge}(R^\wedge/\mathfrak pR^\wedge)
=
\{\mathfrak q_{ij}\}.
$$
设 $y \in R^\wedge$ 且 $y^2 = 0$。由于
$(R^\wedge)_{\mathfrak q_{ij}}$ 正则，从而是整环
（引理 \ref{algebra-lemma-regular-domain}），可知 $y$ 在
$(R^\wedge)_{\mathfrak q_{ij}}$ 中的像为零。因此，由引理
\ref{algebra-lemma-zero-at-ass-zero}，$y$ 在 $R^\wedge/xR^\wedge$
中的像为零。所以 $y = xy'$。由于 $x$ 是非零因子
（因为 $R \to R^\wedge$ 平坦），可知 $(y')^2 = 0$。
因此
$y \in \bigcap x^nR^\wedge = (0)$
（引理 \ref{algebra-lemma-intersect-powers-ideal-module-zero}）。
\end{proof}

\begin{lemma}
\label{algebra-lemma-local-nagata-domain-analytically-unramified}
设 $(R, \mathfrak m)$ 是局部环。若 $R$ 是 Noether 的 Nagata 整环，
则 $R$ 解析非分歧。
\end{lemma}

\begin{proof}
对 $\dim(R)$ 作归纳。当 $\dim(R) = 0$ 时结论平凡。
因此，假设 $\dim(R) = d$，并假设引理对维数 $< d$ 的所有
Noether Nagata 整环都成立。

\medskip\noindent
设 $R \subset S$ 是 $R$ 在 $R$ 的分式域中的整闭包。由假设，$S$ 是有限
$R$-模。由引理 \ref{algebra-lemma-quasi-finite-over-nagata}，$S$ 是 Nagata 环。
由引理 \ref{algebra-lemma-integral-sub-dim-equal}，有
$\dim(R) = \dim(S)$。设 $\mathfrak m_1, \ldots, \mathfrak m_t$
是 $S$ 的极大理想。它们都位于 $R$ 的极大理想 $\mathfrak m$ 之上。
此外，由于 $S/\mathfrak mS$ 是 Artin 的，对充分大的 $n$ 有
$$
(\mathfrak m_1 \cap \ldots \cap \mathfrak m_t)^n \subset \mathfrak mS
$$
由引理 \ref{algebra-lemma-completion-flat}，$R^\wedge \to S^\wedge$
是单射；再把中国剩余引理 \ref{algebra-lemma-chinese-remainder} 与引理
\ref{algebra-lemma-change-ideal-completion} 结合起来，可得
$S^\wedge = \prod S^\wedge_i$，其中 $S^\wedge_i$ 是 $S$
关于极大理想 $\mathfrak m_i$ 的完备化。因此，只需证明
$S_{\mathfrak m_i}$ 解析非分歧。换言之，我们已经化归到
$R$ 是 Noether 正规 Nagata 整环的情形。

\medskip\noindent
假设 $R$ 是 Noether 正规局部 Nagata 整环。在极大理想中选取
非零元素 $x \in \mathfrak m$。我们将应用引理
\ref{algebra-lemma-criterion-analytically-unramified}。必须检验性质
(1)、(2)、(3)(a) 和 (3)(b)。性质 (1) 显然。由引理
\ref{algebra-lemma-normal-domain-intersection-localizations-height-1}，
$R/xR$ 没有嵌入素理想，所以性质 (2) 成立。同一个引理还说明，
$R/xR$ 的每个相伴素理想 $\mathfrak p$ 都具有高度 $1$。
因此，$R_{\mathfrak p}$ 是 $1$-维正规整环，从而是正则的
（引理 \ref{algebra-lemma-characterize-dvr}）。于是 (3)(a) 成立。
最后，(3)(b) 由归纳假设成立，因为 $R/\mathfrak p$ 是 Nagata 环
（由引理 \ref{algebra-lemma-quasi-finite-over-nagata}，或直接由定义）。
所以 $R$ 解析非分歧。
\end{proof}

\begin{lemma}
\label{algebra-lemma-local-nagata-and-analytically-unramified}
% P01-E1008
设 $(R, \mathfrak m)$ 是 Noether 局部环。以下条件等价
\begin{enumerate}
\item $R$ 是 Nagata 环；
\item 对有限映射 $R \to S$，若 $S$ 是整环且
$\mathfrak m' \subset S$ 是极大理想，则局部环
$S_{\mathfrak m'}$ 解析非分歧；
% P01-E1009
\item 对有限局部同态 $(R, \mathfrak m) \to (S, \mathfrak m')$，
若 $S$ 是整环，则 $S$ 解析非分歧。
\end{enumerate}
\end{lemma}

\begin{proof}
假设 $R$ 是 Nagata 环，并设 $R \to S$ 和 $\mathfrak m' \subset S$
如 (2) 所述。由引理 \ref{algebra-lemma-quasi-finite-over-nagata}，$S$ 是
Nagata 环。因此，局部环 $S_{\mathfrak m'}$ 是 Nagata 环
（引理 \ref{algebra-lemma-nagata-localize}）。再由引理
\ref{algebra-lemma-local-nagata-domain-analytically-unramified}，
它解析非分歧。显然，(2) 蕴含 (3)。

\medskip\noindent
假设 (3) 成立。设 $\mathfrak p \subset R$ 是素理想，并设
$L/\kappa(\mathfrak p)$ 是有限域扩张。
% P01-E1010
为证明 (1)，必须证明 $R/\mathfrak p$ 的整闭包在
$R/\mathfrak p$ 上有限。选取 $x_1, \ldots, x_n \in L$，
使它们生成 $L$ 作为 $\kappa(\mathfrak p)$ 的扩张。对每个 $i$，设
$P_i(T) = T^{d_i} + a_{i, 1} T^{d_i - 1} + \ldots + a_{i, d_i}$
是 $x_i$ 在 $\kappa(\mathfrak p)$ 上的极小多项式。
对适当的 $f_i \in R$，$f_i \not \in \mathfrak p$，
用 $f_i x_i$ 代替 $x_i$ 后，可以假设
$a_{i, j} \in R/\mathfrak p$。事实上，再乘以 $\mathfrak m$
中的元素后，可以假设对所有 $i, j$ 都有
$a_{i, j} \in \mathfrak m/\mathfrak p \subset R/\mathfrak p$。
% P01-E1011
完成这些操作后，令 $S = R/\mathfrak p[x_1, \ldots, x_n] \subset L$。
则 $S$ 在 $R$ 上有限，是整环，并且
% P01-E1012
$S/\mathfrak m S$ 是
$R/\mathfrak m[T_1, \ldots, T_n]/(T_1^{d_1}, \ldots, T_n^{d_n})$
的商。因此 $S$ 是局部环。由 (3)，$S$ 解析非分歧；再由引理
\ref{algebra-lemma-analytically-unramified-easy}，它在 $L$ 中的整闭包
$S'$ 在 $S$ 上有限。由于 $S'$ 也是 $R/\mathfrak p$ 在 $L$
中的整闭包，结论得证。
\end{proof}

\noindent
下面的命题特别说明，Nagata 环上的有限型代数仍是 Nagata 环。

\begin{proposition}[Nagata]
\label{algebra-proposition-nagata-universally-japanese}
设 $R$ 是环。以下条件等价：
\begin{enumerate}
\item $R$ 是 Nagata 环；
\item 任意有限型 $R$-代数都是 Nagata 环；并且
\item $R$ 是广泛日本型的且为 Noether 的。
\end{enumerate}
\end{proposition}

\begin{proof}
显然，Noether 的广泛日本型环是广泛 Nagata 的（即条件 (2) 成立）。
设 $R$ 是 Nagata 环。我们将证明任意有限生成 $R$-代数 $S$
都是 Nagata 环。这将证明命题。

\medskip\noindent
第 1 步。存在环映射列
$R = R_0 \to R_1 \to R_2 \to \ldots \to R_n = S$，
使得每个 $R_i \to R_{i + 1}$ 都由单个元素生成。因此，通过归纳，
只需在 $S \cong R[x]/I$ 时证明 $S$ 是 Nagata 环。

\medskip\noindent
第 2 步。设 $\mathfrak q \subset S$ 是 $S$ 的素理想，并设
$\mathfrak p \subset R$ 是相应的 $R$ 的素理想。必须证明
$S/\mathfrak q$ 是 N-2 的。因此，我们已经化归到证明如下命题：
(*) 给定 Nagata 整环 $R$ 和整环的单生成扩张 $R \subset S$，
则 $S$ 是 N-2 的。

\medskip\noindent
第 3 步：设 $R$ 是 Nagata 整环，并设 $R \subset S$ 是整环的
单生成扩张。假设 $R$ 和 $S$ 的分式域之间所诱导的扩张是纯超越的。
在这种情况下，$S = R[x]$。由引理 \ref{algebra-lemma-polynomial-ring-N-2}，
$S$ 是 N-2 的。因此，我们已经化归到证明如下命题：
(**) 给定 Nagata 整环 $R$ 和整环的单生成扩张 $R \subset S$，
若所诱导的分式域扩张有限，则 $S$ 是 N-2 的。

\medskip\noindent
第 4 步。设 $R$ 是正规 Nagata 整环，并设 $R \subset S$ 是整环的
单生成扩张，所诱导的分式域扩张为有限扩张 $L/K$。
选取生成元 $x \in S$，使 $S$ 作为 $R$-代数由它生成。
设 $M/L$ 是有限域扩张。设 $R'$ 是 $R$ 在 $M$ 中的整闭包。
则 $S$ 在 $M$ 中的整闭包 $S'$ 等于 $R'[x]$ 在 $M$ 中的整闭包。
此外，$R'$ 的分式域是 $M$，且 $R \subset R'$ 有限
（由 $R$ 的 Nagata 性质）。这说明 $R'$ 是 Nagata 环
（引理 \ref{algebra-lemma-quasi-finite-over-nagata}）。证明 $S'$ 在 $S$
上有限，等价于证明 $S'$ 在 $R'[x]$ 上有限。用 $R'$ 代替 $R$，
并用 $R'[x]$ 代替 $S$，即可化归到如下命题：
(***) 给定分式域为 $K$ 的正规 Nagata 整环 $R$，以及 $x \in K$，
则由 $R$ 和 $x$ 生成的子环 $S \subset K$ 是 N-1 的。

\medskip\noindent
第 5 步。设 $R$ 是分式域为 $K$ 的正规 Nagata 整环。
设 $x = b/a \in K$。必须证明由 $R$ 和 $x$ 生成的子环
$S \subset K$ 是 N-1 的。注意，$S_a \cong R_a$ 是正规的。
因此，由引理 \ref{algebra-lemma-characterize-N-1}，只需对 $S$ 的每个极大理想
$\mathfrak m$ 证明 $S_{\mathfrak m}$ 是 N-1 的。

\medskip\noindent
在前一段的假设下，选定这样的极大理想，并令
$\mathfrak n = R \cap \mathfrak m$。剩余域扩张
$\kappa(\mathfrak m)/\kappa(\mathfrak n)$ 是有限的
（定理 \ref{algebra-theorem-nullstellensatz}），并由 $x$ 的像生成。
% P01-E1013
对某个 $a \in R$，$a \not \in \mathfrak n$，用 $R_a$ 代替 $R$，
并用 $S_a$ 代替 $S$ 后，可以假设存在首一多项式
% P01-E1014 P01-E1015
$f(X) = X^d + \sum_{i = 1, \ldots, d} a_iX^{d - i}$ 属于 $R[x]$，
使得 $f(x) \in \mathfrak m$（细节略）。设 $K''/K$ 是有限域扩张，
使多项式 $f(X)$ 在 $K''[X]$ 中完全分裂。设 $R'$ 是 $R$ 在
$K''$ 中的整闭包。设 $S' \subset K''$ 是由 $R'$ 和 $x$ 生成的子环。
由于 $R$ 是 Nagata 环，可知 $R'$ 在 $R$ 上有限且为 Nagata 环
（引理 \ref{algebra-lemma-quasi-finite-over-nagata}）。此外，$S'$ 在 $S$
上有限。如果对 $S'$ 的每个极大理想 $\mathfrak m'$，局部环
$S'_{\mathfrak m'}$ 都是 N-1 的，则由引理
\ref{algebra-lemma-characterize-N-1}，$S'_{\mathfrak m}$ 是 N-1 的；
进而由引理 \ref{algebra-lemma-finite-extension-N-2}，
$S_{\mathfrak m}$ 是 N-1 的。用 $R'$ 代替 $R$，用 $S'$ 代替 $S$，
再用位于 $\mathfrak m$ 之上的任意极大理想 $\mathfrak m'$ 代替
$\mathfrak m$ 后，我们达到上述多项式 $f$ 完全分裂的情形：
$f(X) = \prod_{i = 1, \ldots, d} (X - a_i)$，其中 $a_i \in R$。
由于 $f(x) \in \mathfrak m$，对某个 $i$ 有
$x - a_i \in \mathfrak m$。最后，用 $x - a_i$ 代替 $x$ 后，
可以假设 $x \in \mathfrak m$。

\medskip\noindent
概括起来：$R$ 是分式域为 $K$ 的正规 Nagata 整环，$x \in K$，
$S$ 是 $K$ 中由 $x$ 和 $R$ 生成的子环；最后，
$\mathfrak m \subset S$ 是极大理想，且 $x \in \mathfrak m$。必须证明
$S_{\mathfrak m}$ 是 N-1 的。

\medskip\noindent
我们将证明引理 \ref{algebra-lemma-criterion-analytically-unramified}
可应用于局部环 $S_{\mathfrak m}$ 和元素 $x$。这将说明
$S_{\mathfrak m}$ 解析非分歧；随后由引理
\ref{algebra-lemma-analytically-unramified-easy} 可知它是 N-1 的。

\medskip\noindent
必须检验性质 (1)、(2)、(3)(a) 和 (3)(b)。性质 (1) 平凡。
设 $I = \Ker(R[X] \to S)$，其中 $X \mapsto x$。我们断言，
$I$ 由所有满足 $ax = b$ 于 $K$ 中的线性式 $aX - b$ 生成。
显然，所有这些线性式都属于 $I$。若
$g = a_d X^d + \ldots a_1 X + a_0 \in I$，则 $a_dx$ 整于 $R$
（引理 \ref{algebra-lemma-make-integral-trivial}），从而
$b := a_dx \in R$，因为 $R$ 正规。于是
$g - (a_dX - b)X^{d - 1} \in I$，再对次数作归纳即得断言。
作为推论，
$$
S/xS = R[X]/(X, I) = R/J
$$
其中
$$
J = \{b \in R \mid ax = b \text{ 对某个 }a \in R\} = xR \cap R
$$
由引理 \ref{algebra-lemma-normal-domain-intersection-localizations-height-1}，
$S/xS = R/J$ 作为 $R$-模没有嵌入素理想，因而作为 $R/J$-模没有，
因而作为 $S/xS$-模没有，进而作为 $S$-模也没有。这证明了性质 (2)。
取这样的相伴素理想 $\mathfrak q \subset S$，并要求
$\mathfrak q \subset \mathfrak m$（从而它是
$S_{\mathfrak m}/xS_{\mathfrak m}$ 的相伴素理想——这对论证并不重要）。
则 $\mathfrak q$ 是 $xS$ 之上的极小素理想，因而具有高度 $1$。
由上述一系列等式可知，$\mathfrak p = R \cap \mathfrak q$ 是
$R/J$ 的相伴素理想，因此具有高度 $1$（见引理
\ref{algebra-lemma-normal-domain-intersection-localizations-height-1}）。
所以 $R_{\mathfrak p}$ 是离散赋值环，从而包含
$R_{\mathfrak p} \subset S_{\mathfrak q}$ 实为等式。
这说明 $S_{\mathfrak q}$ 正则，证明了性质 (3)(a)。最后，
$(S/\mathfrak q)_{\mathfrak m}$ 是 $S/\mathfrak q$ 的局部化，
而后者是 $S/xS = R/J$ 的商。因此，
$(S/\mathfrak q)_{\mathfrak m}$ 是 Nagata 环 $R$ 的商的局部化，
所以是 Nagata 环（引理 \ref{algebra-lemma-quasi-finite-over-nagata}
和 \ref{algebra-lemma-nagata-localize}），从而解析非分歧
（引理 \ref{algebra-lemma-local-nagata-domain-analytically-unramified}）。
这说明 (3)(b) 成立，证明完毕。
\end{proof}

\begin{proposition}
\label{algebra-proposition-ubiquity-nagata}
以下各类环都是 Nagata 环，因而尤其是广泛日本型环：
\begin{enumerate}
\item 域；
\item Noether 完备局部环；
\item $\mathbf{Z}$；
\item 分式域具有特征零的 Dedekind 整环；
\item 上述任一种环的有限型环扩张。
\end{enumerate}
\end{proposition}

\begin{proof}
Noether 完备局部环的情形是引理
\ref{algebra-lemma-Noetherian-complete-local-Nagata}。在其他情形，
只需对环的每个素理想 $\mathfrak p$ 检验 $R/\mathfrak p$ 是否为 N-2。
只要 $R/\mathfrak p$ 是域，即 $\mathfrak p$ 是极大理想，这就是显然的。
因此，在 Dedekind 环的情形，只需检验 $\mathfrak p = (0)$。
但由于假设分式域具有特征零，引理
\ref{algebra-lemma-domain-char-zero-N-1-2} 给出结论。
\end{proof}

\begin{example}
\label{algebra-example-nonjapanese-dvr}
离散赋值环是 Nagata 环，当且仅当它是 N-2 的（因为模极大理想的商
是域，从而是 N-2 的）。例 \ref{algebra-example-bad-dvr-char-p} 中的离散赋值环
$A$ 不是 Nagata 环，即不是 N-2 的。事实上，有限扩张
$A \subset R = A[f]$ 不是 N-1 的。为此，写
$f = \sum a_i x^i$。对每个 $n \geq 1$，令
$g_n = \sum_{i < n} a_i x^i \in A$。则
$h_n = (f - g_n)/x^n$ 是 $R$ 的分式域中的元素，并且
$h_n^p \in k^p[[x]] \subset A$。所以 $R$ 的整闭包 $R'$ 包含
$h_1, h_2, h_3, \ldots$。现在，若 $R'$ 在 $R$ 上有限，从而在
$A$ 上有限，则对所有 $n$，$f  = x^n h_n + g_n$ 都属于子模
$A + x^nR'$。由 Artin--Rees，这将蕴含 $f \in A$
（引理 \ref{algebra-lemma-intersect-powers-ideal-module-zero}），矛盾。
\end{example}

\begin{lemma}
\label{algebra-lemma-nagata-pth-roots}
设 $(A, \mathfrak m)$ 是 Noether 的 Nagata 局部整环，其分式域
具有特征 $p$。若 $a \in A$ 在 $A^\wedge$ 中有 $p$ 次根，
则 $a$ 在 $A$ 中也有 $p$ 次根。
\end{lemma}

\begin{proof}
设 $\alpha \in A^\wedge$ 是 $a$ 的 $p$ 次根。反设 $a$ 在 $A$
中没有 $p$ 次根。则 $a$ 在 $A$ 的分式域 $K$ 中也没有 $p$ 次根。
确实，若 $\beta = b/c$，其中 $b, c \in A$ 且 $c \not = 0$，
满足 $\beta^p = a$，则在 $A^\wedge$ 中有 $\alpha = b/c$；
由 $A \to A^\wedge$ 的忠实平坦性，这蕴含 $c$ 在 $A$ 中整除 $b$，
因而 $\beta \in A$，矛盾。所以
$B = A[x]/(x^p - a)$ 是整环，因为 $K[x]/(x^p - a)$ 是域。
然而，由所假设的 $\alpha$ 的存在，$B$ 的完备化不是约化的。
这与前面的结果矛盾：$B$ 是 Nagata 环
（命题 \ref{algebra-proposition-nagata-universally-japanese}），
因而由引理 \ref{algebra-lemma-local-nagata-domain-analytically-unramified}
解析非分歧。
\end{proof}

\section{性质的上升}
\label{algebra-section-ascending-properties}

\noindent
本节开始证明一些关于环的性质“上升”的代数事实。为处理环的深度，
我们要用到下述关于模的深度上升的结果，见
\cite[IV, Proposition 6.3.1]{EGA}。

\begin{lemma}
\label{algebra-lemma-apply-grothendieck-module}
\begin{reference}
\cite[IV, Proposition 6.3.1]{EGA}
\end{reference}
有
$$
\text{depth}(M \otimes_R N)
=
\text{depth}(M) + \text{depth}(N/\mathfrak m_RN)
$$
其中 $R \to S$ 是 Noether 局部环的局部同态，$M$ 是有限
$R$-模，而 $N$ 是在 $R$ 上平坦的有限 $S$-模。
\end{lemma}

\begin{proof}
在命题及下面的证明中，我们把 $M$ 的深度取为其作为 $R$-模的深度，
把 $M \otimes_R N$ 的深度取为其作为 $S$-模的深度，并把
$N/\mathfrak m_RN$ 的深度取为其作为 $S/\mathfrak m_RS$-模的深度。
% P01-E1016
以 $n$ 表示右端。先假设 $n$ 为零。于是
$\text{depth}(M) = 0$ 且
$\text{depth}(N/\mathfrak m_RN) = 0$。
这意味着存在 $z \in M$，其零化理想是 $\mathfrak m_R$；还存在
$\overline{y} \in N/\mathfrak m_RN$，其零化理想是
$\mathfrak m_S/\mathfrak m_RS$。取 $\overline{y}$ 的提升 $y \in N$。
由于 $N$ 在 $R$ 上平坦，由该元素给出的映射
$z : R/\mathfrak m_R \to M$ 诱导单射
$N/\mathfrak m_RN \to M \otimes_R N$。因此，
$z \otimes y$ 的零化理想是 $\mathfrak m_S$。于是
$\text{depth}(M \otimes_R N) = 0$。

\medskip\noindent
假设 $n > 0$。若 $\text{depth}(N/\mathfrak m_RN) > 0$，则可选取
$f \in \mathfrak m_S$，其像
$\overline{f} \in S/\mathfrak m_RS$ 是 $N/\mathfrak m_RN$
上的非零因子。由引理 \ref{algebra-lemma-depth-drops-by-one}，
$$
\text{depth}(N/\mathfrak m_RN) =
\text{depth}(N/(f, \mathfrak m_R)N) + 1
$$
根据引理 \ref{algebra-lemma-mod-injective}，元素 $f \in S$ 是 $N$
上的非零因子，且 $N/fN$ 在 $R$ 上平坦。因此，对 $n$ 作归纳可得
$$
\text{depth}(M \otimes_R N/fN) =
\text{depth}(M) + \text{depth}(N/(f, \mathfrak m_R)N).
$$
由于 $N/fN$ 在 $R$ 上平坦，序列
$$
0 \to M \otimes_R N \to M \otimes_R N \to M \otimes_R N/fN \to 0
$$
正合，其中第一个映射是乘以 $f$（引理
\ref{algebra-lemma-flat-tor-zero}）。因此，由引理
\ref{algebra-lemma-depth-drops-by-one} 可得
$\text{depth}(M \otimes_R N) = \text{depth}(M \otimes_R N/fN) + 1$，
从而引理中的公式成立。

\medskip\noindent
若 $n > 0$，但 $\text{depth}(N/\mathfrak m_RN) = 0$，则可选取
$f \in \mathfrak m_R$，使其为 $M$ 上的非零因子。由于 $N$ 在 $R$
上平坦，$f$ 也是 $M \otimes_R N$ 上的非零因子。再次对 $n$ 作归纳，得
$$
\text{depth}(M/fM \otimes_R N) =
\text{depth}(M/fM) + \text{depth}(N/\mathfrak m_RN).
$$
在此情形下，由引理 \ref{algebra-lemma-depth-drops-by-one}，
$\text{depth}(M \otimes_R N) = \text{depth}(M/fM \otimes_R N) + 1$，
且 $\text{depth}(M) = \text{depth}(M/fM) + 1$，
从而引理中的公式成立。
\end{proof}

\begin{lemma}
\label{algebra-lemma-apply-grothendieck}
设 $R \to S$ 是 Noether 局部环之间的平坦局部环同态。则
$$
\text{depth}(S) = \text{depth}(R) + \text{depth}(S/\mathfrak m_RS).
$$
\end{lemma}

\begin{proof}
这是引理 \ref{algebra-lemma-apply-grothendieck-module} 的特例。
\end{proof}

\begin{lemma}
\label{algebra-lemma-CM-goes-up}
% P01-E1017
设 $R \to S$ 是 Noether 局部环之间的平坦局部同态。以下条件等价：
\begin{enumerate}
\item $S$ 是 Cohen--Macaulay 的；并且
\item $R$ 和 $S/\mathfrak m_RS$ 都是 Cohen--Macaulay 的。
\end{enumerate}
\end{lemma}

\begin{proof}
这由定义以及引理 \ref{algebra-lemma-apply-grothendieck} 和
\ref{algebra-lemma-dimension-base-fibre-equals-total} 得出。
\end{proof}

\begin{lemma}
\label{algebra-lemma-Sk-goes-up}
设 $\varphi : R \to S$ 是环映射。假设
\begin{enumerate}
\item $R$ 是 Noether 的；
\item $S$ 是 Noether 的；
\item $\varphi$ 平坦；
\item 纤维环 $S \otimes_R \kappa(\mathfrak p)$ 满足 $(S_k)$；并且
\item $R$ 具有性质 $(S_k)$。
\end{enumerate}
则 $S$ 具有性质 $(S_k)$。
\end{lemma}

\begin{proof}
设 $\mathfrak q$ 是 $S$ 的素理想，位于 $R$ 的素理想
$\mathfrak p$ 之上。由引理 \ref{algebra-lemma-apply-grothendieck}，有
$$
\text{depth}(S_{\mathfrak q}) =
\text{depth}(S_{\mathfrak q}/\mathfrak pS_{\mathfrak q}) +
\text{depth}(R_{\mathfrak p}).
$$
另一方面，由引理 \ref{algebra-lemma-dimension-base-fibre-total}，有
$$
\dim(R_{\mathfrak p})
+
\dim(S_{\mathfrak q}/\mathfrak pS_{\mathfrak q})
\geq
\dim(S_{\mathfrak q})
$$
（实际上，由引理 \ref{algebra-lemma-dimension-base-fibre-equals-total}，
这里等号成立，但严格地说我们并不需要这一点）。最后，由于假设映射的
纤维环满足 $(S_k)$，可知
$\text{depth}(S_{\mathfrak q}/\mathfrak pS_{\mathfrak q})
\geq \min(k, \dim(S_{\mathfrak q}/\mathfrak pS_{\mathfrak q}))$。
因此，引理由下列不等式链得出：
\begin{eqnarray*}
\text{depth}(S_{\mathfrak q}) & = &
\text{depth}(S_{\mathfrak q}/\mathfrak pS_{\mathfrak q}) +
\text{depth}(R_{\mathfrak p}) \\
& \geq &
\min(k, \dim(S_{\mathfrak q}/\mathfrak pS_{\mathfrak q})) +
\min(k, \dim(R_{\mathfrak p})) \\
& = &
\min(2k, \dim(S_{\mathfrak q}/\mathfrak pS_{\mathfrak q}) + k,
k + \dim(R_\mathfrak p),
\dim(S_{\mathfrak q}/\mathfrak pS_{\mathfrak q}) +
\dim(R_{\mathfrak p})) \\
& \geq &
\min(k, \dim(S_{\mathfrak q}))
\end{eqnarray*}
如所需。
\end{proof}

\begin{lemma}
\label{algebra-lemma-Rk-goes-up}
设 $\varphi : R \to S$ 是环映射。假设
\begin{enumerate}
\item $R$ 是 Noether 的；
% P01-E1018
\item $S$ 是 Noether 的；
\item $\varphi$ 平坦；
\item 纤维环 $S \otimes_R \kappa(\mathfrak p)$ 具有性质 $(R_k)$；并且
\item $R$ 具有性质 $(R_k)$。
\end{enumerate}
则 $S$ 具有性质 $(R_k)$。
\end{lemma}

\begin{proof}
设 $\mathfrak q$ 是 $S$ 的素理想，位于 $R$ 的素理想
$\mathfrak p$ 之上。假设 $\dim(S_{\mathfrak q}) \leq k$。
由于由引理 \ref{algebra-lemma-dimension-base-fibre-equals-total} 有
$\dim(S_{\mathfrak q}) = \dim(R_{\mathfrak p})
+ \dim(S_{\mathfrak q}/\mathfrak pS_{\mathfrak q})$，可知
$\dim(R_{\mathfrak p}) \leq k$ 且
$\dim(S_{\mathfrak q}/\mathfrak pS_{\mathfrak q}) \leq k$。
因此，由假设，$R_{\mathfrak p}$ 和
$S_{\mathfrak q}/\mathfrak pS_{\mathfrak q}$ 都正则。再由引理
\ref{algebra-lemma-flat-over-regular-with-regular-fibre}，
$S_{\mathfrak q}$ 正则。
\end{proof}

\begin{lemma}
\label{algebra-lemma-reduced-goes-up-noetherian}
设 $\varphi : R \to S$ 是环映射。假设
\begin{enumerate}
\item $R$ 是 Noether 的；
% P01-E1019
\item $S$ 是 Noether 的；
\item $\varphi$ 平坦；
\item 纤维环 $S \otimes_R \kappa(\mathfrak p)$ 约化；
\item $R$ 约化。
\end{enumerate}
则 $S$ 约化。
\end{lemma}

\begin{proof}
对 Noether 环而言，约化等价于具有性质 $(S_1)$ 和 $(R_0)$，
见引理 \ref{algebra-lemma-criterion-reduced}。因此，已知 $R$ 和纤维环具有这些性质。
于是可应用引理 \ref{algebra-lemma-Sk-goes-up} 和 \ref{algebra-lemma-Rk-goes-up}，
得到 $S$ 具有 $(S_1)$ 和 $(R_0)$；换言之，再由引理
\ref{algebra-lemma-criterion-reduced}，它约化。
\end{proof}

\begin{lemma}
\label{algebra-lemma-reduced-goes-up}
设 $\varphi : R \to S$ 是环映射。假设
\begin{enumerate}
\item $\varphi$ 光滑；
\item $R$ 约化。
\end{enumerate}
则 $S$ 约化。
\end{lemma}

\begin{proof}
注意，$R \to S$ 平坦且纤维正则（见节
\ref{algebra-section-smooth-overview} 中关于光滑环映射的结果列表）。
特别地，纤维约化。因此，若 $R$ 是 Noether 的，则 $S$ 是 Noether 的，
结论由引理 \ref{algebra-lemma-reduced-goes-up-noetherian} 得出。

\medskip\noindent
一般情形下，可找到有限生成的 $\mathbf{Z}$-子代数 $R_0 \subset R$
和光滑环映射 $R_0 \to S_0$，使得
$S \cong R \otimes_{R_0} S_0$，见节
\ref{algebra-section-smooth-overview} 中的评注 (10)。现在，若 $x \in S$
满足 $x^2 = 0$，则可扩大 $R_0$，并假设 $x$ 来自元素
$x_0 \in S_0$。再次扩大 $R_0$ 后，可假设 $x_0^2 = 0$ 于 $S_0$ 中。
% P01-E1020
然而，由于 $R_0 \subset R$ 且后者约化，前者也约化，故可知
$S_0$ 约化，因而如所需 $x_0 = 0$。
\end{proof}

\begin{lemma}
\label{algebra-lemma-normal-goes-up-noetherian}
设 $\varphi : R \to S$ 是环映射。假设
\begin{enumerate}
\item $R$ 是 Noether 的；
\item $S$ 是 Noether 的；
\item $\varphi$ 平坦；
\item 纤维环 $S \otimes_R \kappa(\mathfrak p)$ 正规；并且
\item $R$ 正规。
\end{enumerate}
则 $S$ 正规。
\end{lemma}

\begin{proof}
对 Noether 环而言，正规等价于具有性质 $(S_2)$ 和 $(R_1)$，
见引理 \ref{algebra-lemma-criterion-normal}。因此，已知 $R$ 和纤维环具有这些性质。
于是可应用引理 \ref{algebra-lemma-Sk-goes-up} 和 \ref{algebra-lemma-Rk-goes-up}，
得到 $S$ 具有 $(S_2)$ 和 $(R_1)$；换言之，再由引理
\ref{algebra-lemma-criterion-normal}，它正规。
\end{proof}

\begin{lemma}
\label{algebra-lemma-normal-goes-up}
设 $\varphi : R \to S$ 是环映射。假设
\begin{enumerate}
\item $\varphi$ 光滑；
\item $R$ 正规。
\end{enumerate}
则 $S$ 正规。
\end{lemma}

\begin{proof}
注意，$R \to S$ 平坦且纤维正则（见节
\ref{algebra-section-smooth-overview} 中关于光滑环映射的结果列表）。
特别地，纤维正规。因此，若 $R$ 是 Noether 的，则 $S$ 是 Noether 的，
结论由引理 \ref{algebra-lemma-normal-goes-up-noetherian} 得出。

\medskip\noindent
现考虑一般情形。先注意 $R$ 约化，因而由引理
\ref{algebra-lemma-reduced-goes-up}，$S$ 约化。设 $\mathfrak q$ 是 $S$
的素理想，并设 $\mathfrak p$ 是 $R$ 的相应素理想。注意，
$R_{\mathfrak p}$ 是正规整环。必须证明 $S_{\mathfrak q}$ 是正规整环。
为此，可以用 $R_{\mathfrak p}$ 代替 $R$，并用 $S_{\mathfrak p}$
代替 $S$。因此，可以假设 $R$ 是正规整环。

\medskip\noindent
假设 $R \to S$ 光滑，且 $R$ 是正规整环。可找到有限生成的
$\mathbf{Z}$-子代数 $R_0 \subset R$ 和光滑环映射
$R_0 \to S_0$，使得 $S \cong R \otimes_{R_0} S_0$，见节
\ref{algebra-section-smooth-overview} 中的评注 (10)。由于 $R_0$ 是 Nagata 整环
（见命题 \ref{algebra-proposition-ubiquity-nagata}），可知其整闭包 $R_0'$
在 $R_0$ 上有限。此外，由于 $R$ 是正规整环，显然有
$R_0' \subset R$。因此，可以用 $R_0'$ 代替 $R_0$，并用
$R_0' \otimes_{R_0} S_0$ 代替 $S_0$，从而假设 $R_0$ 是 Noether
正规整环。由证明的第一段可得，$S_0$ 是正规环（当然，它不一定是整环）。
如此可见，$R = \bigcup R_\lambda$ 是 Noether 正规整环的并，
相应地，$S = \colim R_\lambda \otimes_{R_0} S_0$ 是正规环的余极限。
这蕴含 $S$ 是正规环。略去一些细节。
\end{proof}

\begin{lemma}
\label{algebra-lemma-regular-goes-up}
\begin{slogan}
正则性沿光滑环映射上升。
\end{slogan}
设 $\varphi : R \to S$ 是环映射。假设
\begin{enumerate}
\item $\varphi$ 光滑；
\item $R$ 是正则环。
\end{enumerate}
则 $S$ 正则。
\end{lemma}

\begin{proof}
% P01-E1021
对所有性质 $(R_k)$ 应用引理 \ref{algebra-lemma-Rk-goes-up}，并用引理
\ref{algebra-lemma-characterize-smooth-over-field} 验证假设，即得结论。
\end{proof}

\section{性质的下降}
\label{algebra-section-descending-properties}

\noindent
本节开始证明一些关于环的性质“下降”的代数事实。事实表明，性质的下降
往往比其上升“更容易”。换言之，
% P01-E1022
对环映射 $R \to S$ 的假设常常弱于上一节相应引理中的假设。
不过，我们提醒读者：除非也能证明相应的上升结果，否则下降结果往往毫无用处！
下面是说明这种现象的典型结果。

\begin{lemma}
\label{algebra-lemma-descent-Noetherian}
设 $R \to S$ 是环映射。假设
\begin{enumerate}
\item $R \to S$ 忠实平坦；并且
\item $S$ 是 Noether 的。
\end{enumerate}
则 $R$ 是 Noether 的。
\end{lemma}

\begin{proof}
设 $I_0 \subset I_1 \subset I_2 \subset \ldots$ 是 $R$ 的递增理想列。
由假设，对某个 $n$ 有
$I_nS = I_{n+1}S = I_{n+2}S = \ldots$。由忠实平坦性，扩张后再收缩
得到原理想，即对 $R$ 的每个理想 $I$ 都有
$I = R \cap IS$（引理
\ref{algebra-lemma-faithfully-flat-universally-injective}）。因此，如所需，
$I_n = I_{n+1} = I_{n+2} = \ldots$。
\end{proof}

\begin{lemma}
\label{algebra-lemma-descent-reduced}
设 $R \to S$ 是环映射。假设
\begin{enumerate}
\item $R \to S$ 忠实平坦；并且
\item $S$ 约化。
\end{enumerate}
则 $R$ 约化。
\end{lemma}

\begin{proof}
由引理 \ref{algebra-lemma-faithfully-flat-universally-injective}，
$R \to S$ 是单射，故结论显然。
\end{proof}

\begin{lemma}
\label{algebra-lemma-descent-normal}
设 $R \to S$ 是环映射。假设
\begin{enumerate}
\item $R \to S$ 忠实平坦；并且
\item $S$ 是正规环。
\end{enumerate}
则 $R$ 是正规环。
\end{lemma}

\begin{proof}
由于 $S$ 约化，可知 $R$ 约化。设 $\mathfrak p$ 是 $R$ 的素理想。
必须证明 $R_{\mathfrak p}$ 是正规整环。
% P01-E1023
由于 $S_{\mathfrak p}$ 在 $R_{\mathfrak p}$ 上也忠实平坦，
可以假设 $R$ 是极大理想为 $\mathfrak m$ 的局部环。设
$\mathfrak q$ 是 $S$ 的位于 $\mathfrak m$ 之上的素理想。于是
$R \to S_{\mathfrak q}$ 忠实平坦（引理
\ref{algebra-lemma-local-flat-ff}）。因此，可以进一步假设 $S$ 也是局部环。
特别地，$S$ 是正规整环。由于 $R \to S$ 忠实平坦且 $S$
是正规整环，可知 $R$ 是整环。接着，设 $a/b$ 整于 $R$，其中
$a, b \in R$。由于 $S$ 正规，有 $a/b \in S$，因而 $a \in bS$。
这意味着 $a : R \to R/bR$ 经向 $S$ 基变换后成为零映射。
由忠实平坦性，$a \in bR$，所以 $a/b \in R$。因此 $R$ 正规。
\end{proof}

\begin{lemma}
\label{algebra-lemma-descent-regular}
设 $R \to S$ 是环映射。假设
\begin{enumerate}
\item $R \to S$ 忠实平坦；并且
\item $S$ 是正则环。
\end{enumerate}
则 $R$ 是正则环。
\end{lemma}

\begin{proof}
由引理 \ref{algebra-lemma-descent-Noetherian}，$R$ 是 Noether 的。
设 $\mathfrak p \subset R$ 是素理想。选取位于 $\mathfrak p$ 之上的素理想
$\mathfrak q \subset S$。于是，引理 \ref{algebra-lemma-flat-under-regular}
可应用于 $R_\mathfrak p \to S_\mathfrak q$，从而
$R_\mathfrak p$ 正则。由于 $\mathfrak p$ 任意，可知 $R$ 正则。
\end{proof}

\begin{lemma}
\label{algebra-lemma-descent-Sk}
设 $R \to S$ 是环映射。假设
\begin{enumerate}
\item $R \to S$ 忠实平坦；并且
\item $S$ 是 Noether 的且具有性质 $(S_k)$。
\end{enumerate}
则 $R$ 是 Noether 的且具有性质 $(S_k)$。
\end{lemma}

\begin{proof}
我们已经看到，(1) 和 (2) 蕴含 $R$ 是 Noether 的，见引理
\ref{algebra-lemma-descent-Noetherian}。设 $\mathfrak p \subset R$ 是素理想。
选取位于 $\mathfrak p$ 之上的素理想 $\mathfrak q \subset S$，
使其对应于纤维环 $S \otimes_R \kappa(\mathfrak p)$ 的一个极小素理想。
于是，$A = R_{\mathfrak p} \to S_{\mathfrak q} = B$ 是 Noether
局部环之间的平坦局部环同态，且 $\mathfrak m_AB$ 是 $B$ 的定义理想。
因此，$\dim(A) = \dim(B)$（引理
\ref{algebra-lemma-dimension-base-fibre-equals-total}），且
$\text{depth}(A) = \text{depth}(B)$（引理
\ref{algebra-lemma-apply-grothendieck}）。所以，由 $B$ 具有 $(S_k)$，
可知 $A$ 具有 $(S_k)$。
\end{proof}

\begin{lemma}
\label{algebra-lemma-descent-Rk}
设 $R \to S$ 是环映射。假设
\begin{enumerate}
\item $R \to S$ 忠实平坦；并且
\item $S$ 是 Noether 的且具有性质 $(R_k)$。
\end{enumerate}
则 $R$ 是 Noether 的且具有性质 $(R_k)$。
\end{lemma}

\begin{proof}
我们已经看到，(1) 和 (2) 蕴含 $R$ 是 Noether 的，见引理
\ref{algebra-lemma-descent-Noetherian}。设 $\mathfrak p \subset R$ 是素理想，
并假设 $\dim(R_{\mathfrak p}) \leq k$。选取位于 $\mathfrak p$ 之上的
素理想 $\mathfrak q \subset S$，使其对应于纤维环
$S \otimes_R \kappa(\mathfrak p)$ 的一个极小素理想。于是，
$A = R_{\mathfrak p} \to S_{\mathfrak q} = B$ 是 Noether 局部环之间的
平坦局部环同态，且 $\mathfrak m_AB$ 是 $B$ 的定义理想。因此，
$\dim(A) = \dim(B)$（引理
\ref{algebra-lemma-dimension-base-fibre-equals-total}）。由于 $S$ 具有 $(R_k)$，
可知 $B$ 是正则局部环。由引理 \ref{algebra-lemma-flat-under-regular}，
$A$ 正则。
\end{proof}

\begin{lemma}
\label{algebra-lemma-descent-nagata}
设 $R \to S$ 是环映射。假设
\begin{enumerate}
\item $R \to S$ 光滑且在谱上满；并且
\item $S$ 是 Nagata 环。
\end{enumerate}
则 $R$ 是 Nagata 环。
\end{lemma}

\begin{proof}
回忆，Nagata 环等价于 Noether 的广泛日本型环
（命题 \ref{algebra-proposition-nagata-universally-japanese}）。我们已经在引理
\ref{algebra-lemma-descent-Noetherian} 中看到 $R$ 是 Noether 的。设
$R \to A$ 是到整环的有限型环映射。根据引理
\ref{algebra-lemma-check-universally-japanese}，只需检验 $A$ 是 N-1 的。
显然，$B = A \otimes_R S$ 是有限型 $S$-代数，因而是 Nagata 环
（命题 \ref{algebra-proposition-nagata-universally-japanese}）。由于
$A \to B$ 光滑（引理 \ref{algebra-lemma-base-change-smooth}），可知
$B$ 约化（引理 \ref{algebra-lemma-reduced-goes-up}）。由于 $B$ 是 Noether 的，
它只有有限多个极小素理想 $\mathfrak q_1, \ldots, \mathfrak q_t$
（见引理 \ref{algebra-lemma-Noetherian-irreducible-components}）。由于
$A \to B$ 平坦，它们都位于 $(0) \subset A$ 之上
（由下降定理，见引理 \ref{algebra-lemma-flat-going-down}）。
% P01-E1024
全分式环 $Q(B)$ 是 $L_i = \kappa(\mathfrak q_i)$ 的乘积
（引理 \ref{algebra-lemma-total-ring-fractions-no-embedded-points} 和
\ref{algebra-lemma-minimal-prime-reduced-ring}）。此外，$B$ 在 $Q(B)$ 中的整闭包
$B'$ 是各个 $B/\mathfrak q_i$ 在相应因子 $L_i$ 中的整闭包
$B_i'$ 的乘积（与引理
\ref{algebra-lemma-characterize-reduced-ring-normal} 比较）。由于 $B$
是广泛日本型环，环扩张 $B/\mathfrak q_i \subset B_i'$ 有限，
从而 $B' = \prod B_i'$ 在 $B$ 上有限。由于 $A \to B$ 平坦，
$A$ 上的任意非零因子都映为 $B$ 上的非零因子。相应的映射
$$
Q(A) \otimes_A B = (A \setminus \{0\})^{-1}A \otimes_A B
= (A \setminus \{0\})^{-1}B \to Q(B)
$$
是单射（这里使用了引理 \ref{algebra-lemma-tensor-localization}）。
% P01-E1025
经此映射，$A'$ 映入 $B'$。这诱导映射
$$
A' \otimes_A B \longrightarrow B'
$$
它是单射（由上文以及 $A \to B$ 的平坦性）。由于 $B'$ 是有限
$B$-模且 $B$ 是 Noether 的，可知 $A' \otimes_A B$ 是有限 $B$-模。
因此，存在有限多个元素 $x_i \in A'$，使得 $x_i \otimes 1$
作为 $B$-模生成 $A' \otimes_A B$。最后，由 $A \to B$ 的忠实平坦性，
% P01-E1026
这些 $x_i$ 也作为 $A$-模生成 $A'$，结论得证。
\end{proof}

\begin{remark}
\label{algebra-remark-universally-catenary-does-not-descend}
“普遍链状”这一性质不下降，即使沿平展环映射也不下降。在
《例》的节 \ref{examples-section-non-catenary-Noetherian-local}
中，构造了有限环映射 $A \to B$，其中 $A$ 是非普遍链状的
Noether 局部环，而 $B$ 是具有两个极大理想
$\mathfrak m$, $\mathfrak n$ 的半局部环；$B_{\mathfrak m}$ 和
$B_{\mathfrak n}$ 分别是维数 $2$ 和 $1$ 的正则环，
% P01-E1027
并且其剩余域与 $A$ 的剩余域相同。此外，$\mathfrak m_A$ 在
$B_{\mathfrak m}$ 和 $B_{\mathfrak n}$ 中都生成极大理想
（所以 $A \to B$ 既非分歧又有限）。由引理
\ref{algebra-lemma-etale-makes-unramified-closed}，存在局部平展环映射
$A \to A'$，使得 $B \otimes_A A' = B_1 \times B_2$ 分解，
且 $A' \to B_i$ 满。因此，$A'$ 具有两个极小素理想
$\mathfrak q_i$，且 $A'/\mathfrak q_i \cong B_i$。由于 $B_i$
是正则局部环（因为它平展于 $B_{\mathfrak m}$ 或
$B_{\mathfrak n}$ 之一），可得 $A'$ 普遍链状。
\end{remark}

\section{几何正规代数}
\label{algebra-section-geometrically-normal}

\noindent
本节给出环的性质上升和下降的一些应用。

\begin{lemma}
\label{algebra-lemma-geometrically-normal}
设 $k$ 是域。设 $A$ 是 $k$-代数。$A$ 的以下性质等价：
\begin{enumerate}
\item 对每个域扩张 $k'/k$，$k' \otimes_k A$ 都是正规环；
\item 对每个有限生成域扩张 $k'/k$，$k' \otimes_k A$ 都是正规环；
\item 对每个有限纯不可分扩张 $k'/k$，$k' \otimes_k A$ 都是正规环；
\item $k^{perf} \otimes_k A$ 是正规环。
\end{enumerate}
这里的正规环定义见定义 \ref{algebra-definition-ring-normal}。
\end{lemma}

\begin{proof}
显然，(1) $\Rightarrow$ (2) $\Rightarrow$ (3)，且
(1) $\Rightarrow$ (4)。

\medskip\noindent
若 $k'/k$ 是有限纯不可分扩张，则存在 $k$-扩张的嵌入
$k' \to k^{perf}$。环映射
$k' \otimes_k A \to k^{perf} \otimes_k A$ 忠实平坦，
所以若 $k^{perf} \otimes_k A$ 正规，则由引理
\ref{algebra-lemma-descent-normal}，$k' \otimes_k A$ 正规。
由此可见，(4) $\Rightarrow$ (3)。

\medskip\noindent
假设 (2)，并设 $k'/k$ 是任意域扩张。可将
$k' = \colim_i k_i$ 写成有限生成域扩张的有向余极限。因此，
$k' \otimes_k A = \colim_i k_i \otimes_k A$ 是正规环的有向余极限。
由引理 \ref{algebra-lemma-colimit-normal-ring}，$k' \otimes_k A$ 是正规环。
所以 (1) 成立。

\medskip\noindent
假设 (3)，并设 $K/k$ 是有限生成域扩张。由引理
\ref{algebra-lemma-make-separable}，可找到图
$$
\xymatrix{
K \ar[r] & K' \\
k \ar[u] \ar[r] & k' \ar[u]
}
$$
其中 $k'/k$、$K'/K$ 是有限纯不可分域扩张，并且 $K'/k'$ 可分。
由引理 \ref{algebra-lemma-localization-smooth-separable}，存在光滑
$k'$-代数 $B$，使 $K'$ 是 $B$ 的分式域。现在可如下论证：
第 1 步：$k' \otimes_k A$ 是正规环，因为我们假设了 (3)。
第 2 步：$B \otimes_{k'} k' \otimes_k A$ 是正规环，因为
$k' \otimes_k A \to B \otimes_{k'} k' \otimes_k A$ 光滑
（引理 \ref{algebra-lemma-base-change-smooth}），并且正规性沿光滑映射上升
（引理 \ref{algebra-lemma-normal-goes-up}）。
第 3 步。$K' \otimes_{k'} k' \otimes_k A = K' \otimes_k A$
是正规环，因为它是正规环的局部化
（引理 \ref{algebra-lemma-localization-normal-ring}）。
第 4 步。最后，正规性沿忠实平坦环映射
$K \otimes_k A \to K' \otimes_k A$ 下降，所以由引理
\ref{algebra-lemma-descent-normal}，$K \otimes_k A$ 是正规环。
这就证明了引理。
\end{proof}

\begin{definition}
\label{algebra-definition-geometrically-normal}
设 $k$ 是域。若引理 \ref{algebra-lemma-geometrically-normal} 的等价条件成立，
则称 $k$-代数 $R$ 在 $k$ 上{it 几何正规}。
\end{definition}

\begin{lemma}
\label{algebra-lemma-localization-geometrically-normal-algebra}
\begin{slogan}
局部化保持几何正规性。
\end{slogan}
设 $k$ 是域。几何正规 $k$-代数的局部化仍几何正规。
\end{lemma}

\begin{proof}
这是显然的，因为一个环正规与否可在各个素理想处的局部化上检验。
\end{proof}

\begin{lemma}
\label{algebra-lemma-separable-field-extension-geometrically-normal}
设 $k$ 是域。设 $K/k$ 是可分域扩张。则 $K$ 在 $k$ 上几何正规。
\end{lemma}

\begin{proof}
这是因为 $k^{perf} \otimes_k K$ 是域。事实上，由引理
\ref{algebra-lemma-separable-extension-preserves-reducedness}，它约化。
由引理 \ref{algebra-lemma-perfection}（或定义 \ref{algebra-definition-perfection}），
域扩张 $k^{perf}/k$ 纯不可分。因此，由引理
\ref{algebra-lemma-radicial-integral-bijective}，环
$k^{perf} \otimes_k K$ 具有唯一素理想。具有唯一素理想的约化环是域。
\end{proof}

\begin{lemma}
\label{algebra-lemma-geometrically-normal-tensor-normal}
设 $k$ 是域。设 $A, B$ 是 $k$-代数。假设 $A$ 在 $k$ 上几何正规，
且 $B$ 是正规环。则 $A \otimes_k B$ 是正规环。
\end{lemma}

\begin{proof}
% P01-E1028
设 $\mathfrak r$ 是 $A \otimes_k B$ 的素理想。记
$\mathfrak p$、$\mathfrak q$ 分别为 $A$、$B$ 中的相应素理想。则
$(A \otimes_k B)_{\mathfrak r}$ 是
$A_{\mathfrak p} \otimes_k B_{\mathfrak q}$ 的局部化。因此，由引理
\ref{algebra-lemma-localization-normal-ring} 和
\ref{algebra-lemma-localization-geometrically-normal-algebra}，
只需对环 $A_{\mathfrak p} \otimes_k B_{\mathfrak q}$ 证明结论。
所以可以假设 $A$ 和 $B$ 都是整环。

\medskip\noindent
% P01-E1029
假设 $A$ 和 $B$ 是分式域分别为 $K$ 和 $L$ 的整环。注意，
% P01-E1030
$B$ 是其有限型正规 $k$-子代数的滤过余极限（因为 $k$ 是 Nagata 环，
见命题 \ref{algebra-proposition-ubiquity-nagata}；从而，由命题
\ref{algebra-proposition-nagata-universally-japanese}，有限型
$k$-子代数的整闭包仍是有限型 $k$-子代数）。由引理
\ref{algebra-lemma-colimit-normal-ring}，可化归到 $B$ 在 $k$ 上有限型的情形。

\medskip\noindent
% P01-E1029 P01-E1031
假设 $A$ 和 $B$ 是分式域分别为 $K$ 和 $L$ 的整环，且 $B$
在 $k$ 上有限型。在此情形下，环 $K \otimes_k B$ 在 $K$ 上有限型，
因而是 Noether 的（引理 \ref{algebra-lemma-Noetherian-permanence}）。特别地，
$K \otimes_k B$ 只有有限多个极小素理想
（引理 \ref{algebra-lemma-Noetherian-irreducible-components}）。由于
$A \to A \otimes_k B$ 平坦，这蕴含 $A \otimes_k B$ 只有有限多个
极小素理想（由平坦环映射的下降定理——引理
\ref{algebra-lemma-flat-going-down}——这些素理想都位于 $(0) \subset A$ 之上）。
因此，只需证明 $A \otimes_k B$ 在其全分式环中整闭
（引理 \ref{algebra-lemma-characterize-reduced-ring-normal}）。

\medskip\noindent
我们断言，$K \otimes_k B$ 和 $A \otimes_k L$ 都是正规环。
若此断言成立，则 $Q(A \otimes_k B)$ 中任一整于
$A \otimes_k B$ 的元素 $x$ 都属于
$K \otimes_k B \cap A \otimes_k L = A \otimes_k B$
（由引理 \ref{algebra-lemma-normal-ring-integrally-closed}），结论得证。
% P01-E1032
由于由假设 $A \otimes_K L$ 是正规环，只需证明
$K \otimes_k B$ 正规。

\medskip\noindent
由于 $A$ 在 $k$ 上几何正规，可知 $K$ 在 $k$ 上几何正规
（引理 \ref{algebra-lemma-localization-geometrically-normal-algebra}），
因而 $K$ 在 $k$ 上几何约化。所以，$K = \bigcup K_i$ 是
$k$ 的几何约化有限生成域扩张的并
（引理 \ref{algebra-lemma-subalgebra-separable}）。每个 $K_i$ 都是某个光滑
$k$-代数的局部化（引理
\ref{algebra-lemma-localization-smooth-separable}）。因此，
$K_i \otimes_k B$ 是某个光滑 $B$-代数的局部化，故正规
（引理 \ref{algebra-lemma-normal-goes-up}）。于是，
$K \otimes_k B$ 是正规环（引理
\ref{algebra-lemma-colimit-normal-ring}），结论得证。
\end{proof}

\begin{lemma}
\label{algebra-lemma-geometrically-normal-over-separable-algebraic}
设 $k'/k$ 是可分代数域扩张。设 $A$ 是 $k'$ 上的代数。
则 $A$ 在 $k$ 上几何正规，当且仅当它在 $k'$ 上几何正规。
\end{lemma}

\begin{proof}
设 $L/k$ 是有限纯不可分域扩张。则
$L' = k' \otimes_k L$ 是域（见《域》的节
\ref{fields-section-algebraic} 中的材料），且
$A \otimes_k L = A \otimes_{k'} L'$。因此，若 $A$ 在 $k'$
上几何正规，则 $A$ 在 $k$ 上几何正规。

\medskip\noindent
假设 $A$ 在 $k$ 上几何正规。设 $K/k'$ 是域扩张。则
$$
K \otimes_{k'} A = (K \otimes_k A) \otimes_{(k' \otimes_k k')} k'
$$
由引理 \ref{algebra-lemma-separable-algebraic-diagonal}，
$k' \otimes_k k' \to k'$ 是局部化，所以
$K \otimes_{k'} A$ 是正规环的局部化，因而正规。
\end{proof}

\section{几何正则代数}
\label{algebra-section-geometrically-regular}

\noindent
设 $k$ 是域。
设 $A$ 是 Noether 的 $k$-代数。
设 $K/k$ 是有限生成域扩张。
则环 $K \otimes_k A$ 也是 Noether 的，见
引理 \ref{algebra-lemma-Noetherian-field-extension}。
因此，以下引理是有意义的。

\begin{lemma}
\label{algebra-lemma-geometrically-regular}
设 $k$ 是域。设 $A$ 是 $k$-代数。
假设 $A$ 是 Noether 的。
$A$ 的以下性质等价：
\begin{enumerate}
\item 对每个有限生成域扩张 $k'/k$，$k' \otimes_k A$ 都是正则环；
\item 对每个有限纯不可分域扩张 $k'/k$，$k' \otimes_k A$ 都是正则环。
\end{enumerate}
此处正则环的含义如定义 \ref{algebra-definition-regular} 所述。
\end{lemma}

\begin{proof}
由引理前面的说明，此引理是有意义的。
显然 (1) $\Rightarrow$ (2)。

\medskip\noindent
假设 (2)，并设 $K/k$ 是有限生成域扩张。
由引理 \ref{algebra-lemma-make-separable}，可以找到交换图
$$
\xymatrix{
K \ar[r] & K' \\
k \ar[u] \ar[r] & k' \ar[u]
}
$$
其中 $k'/k$、$K'/K$ 是有限纯不可分域扩张，且 $K'/k'$ 可分。由
引理 \ref{algebra-lemma-localization-smooth-separable}，
存在光滑 $k'$-代数 $B$，使 $K'$ 是 $B$ 的分式域。现在论证如下：
步骤 1：$k' \otimes_k A$ 是正则环，因为我们假设了 (2)。
步骤 2：$B \otimes_{k'} k' \otimes_k A$ 是正则环，这是因为
$k' \otimes_k A \to B \otimes_{k'} k' \otimes_k A$ 光滑
（引理 \ref{algebra-lemma-base-change-smooth}），并且正则性沿光滑映射上升
（引理 \ref{algebra-lemma-regular-goes-up}）。
步骤 3：$K' \otimes_{k'} k' \otimes_k A = K' \otimes_k A$ 是正则环，
因为它是正则环的局部化（直接由定义可得）。
步骤 4：最后，由正则性沿忠实平坦环映射
$K \otimes_k A \to K' \otimes_k A$ 的下降
（引理 \ref{algebra-lemma-descent-regular}），$K \otimes_k A$ 是正则环。
这就证明了引理。
\end{proof}

\begin{definition}
\label{algebra-definition-geometrically-regular}
设 $k$ 是域。设 $R$ 是 Noether 的 $k$-代数。
若引理 \ref{algebra-lemma-geometrically-regular} 的等价条件成立，则称
$k$-代数 $R$ 在 $k$ 上\textit{几何正则}。
\end{definition}

\noindent
由定义显然可知，对于 $k$ 的任意有限生成域扩张 $K$，
$K \otimes_k R$ 是在 $K$ 上几何正则的代数。
我们以后将会看到（《代数进阶》，命题
\ref{more-algebra-proposition-characterization-geometrically-regular}），
只需在 $k \subset k' \subset k^{1/p}$（有限）时检验
$R \otimes_k k'$ 正则即可。

\begin{lemma}
\label{algebra-lemma-geometrically-regular-descent}
\begin{slogan}
几何正则性沿代数的忠实平坦映射下降
\end{slogan}
设 $k$ 是域。设 $A \to B$ 是忠实平坦的 $k$-代数映射。
若 $B$ 在 $k$ 上几何正则，则 $A$ 也如此。
\end{lemma}

\begin{proof}
假设 $B$ 在 $k$ 上几何正则。
设 $k'/k$ 是有限纯不可分扩张。
由于 $A \to B$ 的基变换（由
引理 \ref{algebra-lemma-surjective-spec-radical-ideal} 和
\ref{algebra-lemma-flat-base-change}），
$A \otimes_k k' \to B \otimes_k k'$ 忠实平坦；而由我们对 $B$
在 $k$ 上的假设，$B \otimes_k k'$ 正则。于是由引理
\ref{algebra-lemma-descent-regular}，$A \otimes_k k'$ 正则。
\end{proof}

\begin{lemma}
\label{algebra-lemma-geometrically-regular-goes-up}
设 $k$ 是域。设 $A \to B$ 是 $k$-代数间的光滑环映射。
若 $A$ 在 $k$ 上几何正则，则 $B$ 在 $k$ 上几何正则。
\end{lemma}

\begin{proof}
设 $k'/k$ 是有限生成域扩张。
则 $A \otimes_k k' \to B \otimes_k k'$ 是光滑环映射
（引理 \ref{algebra-lemma-base-change-smooth}），且 $A \otimes_k k'$ 正则。
故由引理 \ref{algebra-lemma-regular-goes-up}，$B \otimes_k k'$ 正则。
\end{proof}

\begin{lemma}
\label{algebra-lemma-geometrically-regular-over-subfields}
设 $k$ 是域。设 $A$ 是 $k$ 上的代数。
设 $k = \colim k_i$ 是由子域组成的有向余极限。
若 $A$ 在每个 $k_i$ 上都几何正则，则 $A$ 在 $k$ 上几何正则。
\end{lemma}

\begin{proof}
设 $k'/k$ 是有限纯不可分域扩张。
可以通过向 $k$ 添入有限多个变量并施加有限多个多项式关系得到 $k'$。
因而可见，存在 $i$ 和有限纯不可分域扩张 $k_i'/k_i$，使得
% P01-E1033
$k_i = k \otimes_{k_i} k_i'$。
于是 $A \otimes_k k' = A \otimes_{k_i} k_i'$，引理即得证。
\end{proof}

\begin{lemma}
\label{algebra-lemma-geometrically-regular-over-separable-algebraic}
设 $k'/k$ 是可分代数域扩张。
设 $A$ 是 $k'$ 上的代数。则 $A$ 在 $k$ 上几何正则，当且仅当
它在 $k'$ 上几何正则。
\end{lemma}

\begin{proof}
设 $L/k$ 是有限纯不可分域扩张。
则 $L' = k' \otimes_k L$ 是域（见《域》的节
\ref{fields-section-algebraic} 中的材料），且
$A \otimes_k L = A \otimes_{k'} L'$。因此，若 $A$ 在 $k'$
上几何正则，则 $A$ 在 $k$ 上几何正则。

\medskip\noindent
假设 $A$ 在 $k$ 上几何正则。由于 $k'$ 是 $k$ 的有限扩张的滤过余极限，
由引理 \ref{algebra-lemma-geometrically-regular-over-subfields}，可以假设
$k'/k$ 是有限可分的。考虑环映射
$$
k' \to A \otimes_k k' \to A.
$$
注意，作为 $A$ 到 $k'$ 的基变换，$A \otimes_k k'$ 在 $k'$ 上几何正则。
注意，$A \otimes_k k' \to A$ 是 $k' \otimes_k k' \to k'$
沿映射 $k' \to A$ 的基变换。由于 $k'/k$ 是环的平展扩张，可知
$k' \otimes_k k' \to k'$ 是平展的
（引理 \ref{algebra-lemma-etale}）。故由引理
\ref{algebra-lemma-geometrically-regular-goes-up}，$A$ 在 $k'$ 上几何正则。
\end{proof}

\section{几何 Cohen--Macaulay 代数}
\label{algebra-section-geometrically-CM}

\noindent
本节的名称略有不当，因为 Cohen--Macaulay 代数自动是几何
Cohen--Macaulay 的。具体见下面的
引理 \ref{algebra-lemma-extend-field-CM-locus}
和
引理 \ref{algebra-lemma-CM-geometrically-CM}。

\begin{lemma}
\label{algebra-lemma-tensor-fields-CM}
设 $k$ 是域，$K/k$ 和 $L/k$ 是两个域扩张，并且其中一个是有限型域扩张。
则 $K \otimes_k L$ 是 Noether 的 Cohen--Macaulay 环。
\end{lemma}

\begin{proof}
由引理 \ref{algebra-lemma-Noetherian-field-extension}，
环 $K \otimes_k L$ 是 Noether 的。
不妨设 $K$ 是纯超越扩张 $k(t_1, \ldots, t_r)$ 的有限扩张。则
$k(t_1, \ldots, t_r) \otimes_k L \to K \otimes_k L$
是有限自由环映射。由
引理 \ref{algebra-lemma-finite-flat-over-regular-CM}，
只需证明 $k(t_1, \ldots, t_r) \otimes_k L$ 是 Cohen--Macaulay 的。
这是显然的，因为它是多项式环 $L[t_1, \ldots, t_r]$ 的局部化。
（关于多项式环是 Cohen--Macaulay 的事实，例如见
引理 \ref{algebra-lemma-CM-polynomial-algebra}。）
\end{proof}

\begin{lemma}
\label{algebra-lemma-CM-geometrically-CM}
设 $k$ 是域。设 $S$ 是 Noether 的 $k$-代数。
设 $K/k$ 是有限生成域扩张，并置 $S_K = K \otimes_k S$。
设 $\mathfrak q \subset S$ 是 $S$ 的素理想。设
$\mathfrak q_K \subset S_K$ 是 $S_K$ 中位于 $\mathfrak q$ 之上的素理想。
则 $S_{\mathfrak q}$ 是 Cohen--Macaulay 的，当且仅当
$(S_K)_{\mathfrak q_K}$ 是 Cohen--Macaulay 的。
\end{lemma}

\begin{proof}
由
引理 \ref{algebra-lemma-Noetherian-field-extension}，
环 $S_K$ 是 Noether 的。因此
$S_{\mathfrak q} \to (S_K)_{\mathfrak q_K}$ 是 Noether 局部环之间的
平坦局部同态。注意，其纤维
$$
(S_K)_{\mathfrak q_K} / \mathfrak q (S_K)_{\mathfrak q_K}
\cong (\kappa(\mathfrak q) \otimes_k K)_{\mathfrak q'}
$$
是 Cohen--Macaulay 环（引理 \ref{algebra-lemma-tensor-fields-CM}）
$\kappa(\mathfrak q) \otimes_k K$ 在某个适当素理想
$\mathfrak q'$ 处的局部化。因此，引理由
引理 \ref{algebra-lemma-CM-goes-up} 得出。
\end{proof}

\section{余极限与有限表示映射（二）}
\label{algebra-section-colimits-finite-presentation}

\noindent
本节是节 \ref{algebra-section-colimits-flat} 的续篇。

\medskip\noindent
我们先给出平坦性开性的一个应用。
它说明可以用平坦模逼近平坦模；这一点很有用。

\begin{lemma}
\label{algebra-lemma-flat-finite-presentation-limit-flat}
设 $R \to S$ 是环映射。
设 $M$ 是 $S$-模。
假设
\begin{enumerate}
\item $R \to S$ 是有限表示的；
\item $M$ 是有限表示 $S$-模；
\item $M$ 在 $R$ 上平坦。
\end{enumerate}
在此情形下有：
\begin{enumerate}
\item 存在有限型 $\mathbf{Z}$-代数 $R_0$、有限型环映射
$R_0 \to S_0$ 和有限 $S_0$-模 $M_0$，使得 $M_0$ 在 $R_0$ 上平坦；
% P01-E1034
还存在环映射 $R_0 \to R$、$S_0 \to S$ 和 $S_0$-模映射 $M_0 \to M$，
使得 $S \cong R \otimes_{R_0} S_0$ 且
$M = S \otimes_{S_0} M_0$。
\item 若将 $R = \colim_{\lambda \in \Lambda} R_\lambda$ 写成有向余极限，
则存在 $\lambda$、有限表示环映射 $R_\lambda \to S_\lambda$ 和
有限表示 $S_\lambda$-模 $M_\lambda$，使得 $M_\lambda$ 在
$R_\lambda$ 上平坦，并且
$S = R \otimes_{R_\lambda} S_\lambda$ 且
$M = S \otimes_{S_{\lambda}} M_\lambda$。
\item 若
$$
(R \to S, M) =
\colim_{\lambda \in \Lambda}
(R_\lambda \to S_\lambda, M_\lambda)
$$
写成满足以下条件的有向余极限：
\begin{enumerate}
\item 对 $\mu \geq \lambda$，
$R_\mu \otimes_{R_\lambda} S_\lambda \to S_\mu$ 和
$S_\mu \otimes_{S_\lambda} M_\lambda \to M_\mu$ 是同构；
\item $R_\lambda \to S_\lambda$ 是有限表示的；
\item $M_\lambda$ 是有限表示 $S_\lambda$-模；
\end{enumerate}
则对于所有充分大的 $\lambda$，模 $M_\lambda$ 在 $R_\lambda$ 上平坦。
\end{enumerate}
\end{lemma}

\begin{proof}
先将 $(R \to S, M)$ 写成引理
\ref{algebra-lemma-limit-module-finite-presentation} 所述系统
$(R_\lambda \to S_\lambda, M_\lambda)$ 的有向余极限。
设 $\mathfrak q \subset S$ 是素理想。
设 $\mathfrak p \subset R$、$\mathfrak q_\lambda \subset S_\lambda$
和 $\mathfrak p_\lambda \subset R_\lambda$ 是相应的素理想。
如定理 \ref{algebra-theorem-openness-flatness} 的证明中所见，
$$
((R_\lambda)_{\mathfrak p_\lambda},
(S_\lambda)_{\mathfrak q_\lambda},
(M_\lambda)_{\mathfrak q_{\lambda}})
$$
是引理 \ref{algebra-lemma-limit-module-essentially-finite-presentation}
所述的系统，因而由引理 \ref{algebra-lemma-colimit-eventually-flat}，
可知存在某个 $\lambda_{\mathfrak q} \in \Lambda$，使得对所有
$\lambda \geq \lambda_{\mathfrak q}$，模 $M_\lambda$ 在素理想
$\mathfrak q_{\lambda}$ 处于 $R_\lambda$ 上平坦。

\medskip\noindent
由定理 \ref{algebra-theorem-openness-flatness}，得到开子集
$U_\lambda \subset \Spec(S_\lambda)$，使得 $M_\lambda$ 在
$U_\lambda$ 的所有素理想处于 $R_\lambda$ 上平坦。
记 $V_\lambda \subset \Spec(S)$ 为 $U_\lambda$ 在映射
$\Spec(S) \to \Spec(S_\lambda)$ 下的逆像。
上述论证说明，对每个 $\mathfrak q \in \Spec(S)$，存在
$\lambda_{\mathfrak q}$，使得对所有
$\lambda \geq \lambda_{\mathfrak q}$，都有
$\mathfrak q \in V_\lambda$。由于 $\Spec(S)$ 拟紧，可知这蕴含
存在单个 $\lambda_0 \in \Lambda$，使得
$V_{\lambda_0} = \Spec(S)$。

\medskip\noindent
补集 $\Spec(S_{\lambda_0}) \setminus U_{\lambda_0}$ 等于某个理想
$I \subset S_{\lambda_0}$ 的 $V(I)$。由于
$V_{\lambda_0} = \Spec(S)$，可知 $IS = S$。
% P01-E1035
选取 $f_1, \ldots, f_r \in I$ 和 $s_1, \ldots, s_n \in S$，
使得 $\sum f_i s_i = 1$。由于 $\colim S_\lambda = S$，增大
$\lambda_0$ 后，可以假设存在 $s_{i, \lambda_0} \in S_{\lambda_0}$，
使得 $\sum f_i s_{i, \lambda_0} = 1$。
因此，对这个 $\lambda_0$ 有
$U_{\lambda_0} = \Spec(S_{\lambda_0})$。
这证明了 (1)。

\medskip\noindent
证明 (2)。设 $(R_0 \to S_0, M_0)$ 如 (1) 所述，并假设
$R = \colim R_\lambda$。由于 $R_0$ 是有限型 $\mathbf{Z}$-代数，
存在 $\lambda$ 和映射 $R_0 \to R_\lambda$，使得
$R_0 \to R_\lambda \to R$ 是给定映射 $R_0 \to R$
（见引理 \ref{algebra-lemma-characterize-finite-presentation}）。
于是，取 $S_\lambda = R_\lambda \otimes_{R_0} S_0$ 和
$M_\lambda = S_\lambda \otimes_{S_0} M_0$，即得 (2)。

\medskip\noindent
最后证明 (3)。设 $(R_\lambda \to S_\lambda, M_\lambda)$ 如 (3) 所述。
按 (1) 选取 $(R_0 \to S_0, M_0)$ 和 $R_0 \to R$。
如 (2) 的证明所述，存在 $\lambda_0$ 和环映射
$R_0 \to R_{\lambda_0}$，使得
$R_0 \to R_{\lambda_0} \to R$ 是给定映射 $R_0 \to R$。
由于 $S_0$ 在 $R_0$ 上有限表示且 $S = \colim S_\lambda$，
可知对某个 $\lambda_1 \geq \lambda_0$，存在 $R_0$-代数映射
$S_0 \to S_{\lambda_1}$，使得复合
$S_0 \to S_{\lambda_1} \to S$ 是给定映射 $S_0 \to S$
（见引理 \ref{algebra-lemma-characterize-finite-presentation}）。
对所有 $\lambda \geq \lambda_1$，这给出映射
$$
\Psi_{\lambda} :
R_\lambda \otimes_{R_0} S_0
\longrightarrow
R_\lambda \otimes_{R_{\lambda_1}} S_{\lambda_1}
\cong
S_\lambda
$$
其中最后一个同构来自假设。由构造，
$\colim_\lambda \Psi_\lambda$ 是同构。因此，由引理
\ref{algebra-lemma-colimit-category-fp-algebras}，对所有充分大的 $\lambda$，
$\Psi_\lambda$ 是同构。
同理，存在 $\lambda_2 \geq \lambda_1$ 和 $S_0$-模映射
$M_0 \to M_{\lambda_2}$，使得
$M_0 \to M_{\lambda_2} \to M$ 是给定映射
$M_0 \to M$（见引理 \ref{algebra-lemma-module-map-property-in-colimit}）。
对 $\lambda \geq \lambda_2$，有诱导映射
$$
S_\lambda \otimes_{S_0} M_0
\longrightarrow
S_\lambda \otimes_{S_{\lambda_2}} M_{\lambda_2}
\cong
M_\lambda
$$
而由引理 \ref{algebra-lemma-colimit-category-fp-modules}，对充分大的
$\lambda$，此映射是同构。由于 $M_0$ 在 $R_0$ 上平坦，这就蕴含 (3)。
\end{proof}

\begin{lemma}
\label{algebra-lemma-descend-faithfully-flat-finite-presentation}
设 $R \to A \to B$ 是环映射。
假设 $A \to B$ 是有限表示的忠实平坦映射。
则存在交换图
$$
\xymatrix{
R \ar[r] \ar@{=}[d] &
A_0 \ar[d] \ar[r] &
B_0 \ar[d] \\
R \ar[r] & A \ar[r] & B
}
$$
其中 $R \to A_0$ 是有限表示的，
$A_0 \to B_0$ 是有限表示的忠实平坦映射，并且
$B = A \otimes_{A_0} B_0$。
\end{lemma}

\begin{proof}
% P01-E1036
先证明以 $\mathbf{Z}$ 替代 $R$ 时的引理。
由引理 \ref{algebra-lemma-flat-finite-presentation-limit-flat}，存在交换图
$$
\xymatrix{
A_0 \ar[r] \ar[d] & A \ar[d] \\
B_0 \ar[r] & B
}
$$
其中 $A_0$ 在 $\mathbf{Z}$ 上有限型，$B_0$ 在 $A_0$ 上平坦且有限表示，
并且 $B = A \otimes_{A_0} B_0$。
由于 $A_0 \to B_0$ 平坦且有限表示，可知映射
$\Spec(B_0) \to \Spec(A_0)$ 的像是开的，见
命题 \ref{algebra-proposition-fppf-open}。因此，像的补集等于某个理想
$I_0 \subset A_0$ 的 $V(I_0)$。
由于 $A \to B$ 忠实平坦，映射
$\Spec(B) \to \Spec(A)$ 满射，见引理 \ref{algebra-lemma-ff-rings}。
现在利用基变换的像等于像的基变换这一事实。
于是 $I_0A = A$。选取关系
$\sum f_i r_i = 1$，其中 $r_i \in A$、$f_i \in I_0$。
于是，在扩大 $A_0$ 使其包含元素 $r_i$（并相应扩大 $B_0$）后，
可知 $A_0 \to B_0$ 在谱上也满射，即忠实平坦。

\medskip\noindent
因此，当 $R = \mathbf{Z}$ 时引理成立。
在一般情形下，取刚得到的解 $A_0' \to B_0'$，并置
$A_0 = A_0' \otimes_{\mathbf{Z}} R$、
$B_0 = B_0' \otimes_{\mathbf{Z}} R$。
\end{proof}

\begin{lemma}
\label{algebra-lemma-colimit-finite}
设 $A = \colim_{i \in I} A_i$ 是环的有向余极限。
% P01-E1037
设 $0 \in I$，且 $\varphi_0 : B_0 \to C_0$ 是 $A_0$-代数的映射。
假设
\begin{enumerate}
\item $A \otimes_{A_0} B_0 \to A \otimes_{A_0} C_0$ 是有限的；
\item $C_0$ 在 $B_0$ 上有限型。
\end{enumerate}
则存在 $i \geq 0$，使得映射
$A_i \otimes_{A_0} B_0 \to A_i \otimes_{A_0} C_0$
是有限的。
\end{lemma}

\begin{proof}
设 $x_1, \ldots, x_m$ 是 $C_0$ 在 $B_0$ 上的一组生成元。
选取首一多项式 $P_j \in A \otimes_{A_0} B_0[T]$，使得
$P_j(1 \otimes x_j) = 0$ 在 $A \otimes_{A_0} C_0$ 中成立。
对某个 $i \geq 0$，可以找到映到 $P_j$ 的
$P_{j, i} \in A_i \otimes_{A_0} B_0[T]$。
由于 $\otimes$ 与余极限交换，必要时增大 $i$ 后，可知
$P_{j, i}(1 \otimes x_j)$ 在 $A_i \otimes_{A_0} C_0$ 中为零。
于是这个 $i$ 符合要求。
\end{proof}

\begin{lemma}
\label{algebra-lemma-colimit-surjective}
设 $A = \colim_{i \in I} A_i$ 是环的有向余极限。
设 $0 \in I$，且 $\varphi_0 : B_0 \to C_0$ 是 $A_0$-代数的映射。
假设
\begin{enumerate}
\item $A \otimes_{A_0} B_0 \to A \otimes_{A_0} C_0$ 是满射；
\item $C_0$ 在 $B_0$ 上有限型。
\end{enumerate}
则对某个 $i \geq 0$，映射
$A_i \otimes_{A_0} B_0 \to A_i \otimes_{A_0} C_0$
是满射。
\end{lemma}

\begin{proof}
设 $x_1, \ldots, x_m$ 是 $C_0$ 在 $B_0$ 上的一组生成元。
选取 $b_j \in A \otimes_{A_0} B_0$，使其在
$A \otimes_{A_0} C_0$ 中映到 $1 \otimes x_j$。
对某个 $i \geq 0$，可以找到映到 $b_j$ 的
$b_{j, i} \in A_i \otimes_{A_0} B_0$。
增大 $i$ 后，可以假设对所有 $j = 1, \ldots, m$，
$b_{j, i}$ 在 $A_i \otimes_{A_0} C_0$ 中映到 $1 \otimes x_j$。
于是这个 $i$ 符合要求。
\end{proof}

\begin{lemma}
\label{algebra-lemma-colimit-unramified}
设 $A = \colim_{i \in I} A_i$ 是环的有向余极限。
设 $0 \in I$，且 $\varphi_0 : B_0 \to C_0$ 是 $A_0$-代数的映射。
假设
\begin{enumerate}
\item $A \otimes_{A_0} B_0 \to A \otimes_{A_0} C_0$ 是非分歧的；
\item $C_0$ 在 $B_0$ 上有限型。
\end{enumerate}
则对某个 $i \geq 0$，映射
$A_i \otimes_{A_0} B_0 \to A_i \otimes_{A_0} C_0$
是非分歧的。
\end{lemma}

\begin{proof}
置 $B_i = A_i \otimes_{A_0} B_0$、$C_i = A_i \otimes_{A_0} C_0$、
$B = A \otimes_{A_0} B_0$ 和 $C = A \otimes_{A_0} C_0$。
设 $x_1, \ldots, x_m$ 是 $C_0$ 在 $B_0$ 上的一组生成元。
则 $\text{d}x_1, \ldots, \text{d}x_m$ 在 $C_0$ 上生成
$\Omega_{C_0/B_0}$，并且它们的像在 $C_i$ 上生成
$\Omega_{C_i/B_i}$（引理 \ref{algebra-lemma-differentials-polynomial-ring} 和
\ref{algebra-lemma-differential-seq}）。
注意
$0 = \Omega_{C/B} = \colim \Omega_{C_i/B_i}$
（引理 \ref{algebra-lemma-colimit-differentials}）。
因而存在 $i$，使得 $\text{d}x_1, \ldots, \text{d}x_m$ 映到零，
从而如所需有 $\Omega_{C_i/B_i} = 0$。
\end{proof}

\begin{lemma}
\label{algebra-lemma-colimit-isomorphism}
设 $A = \colim_{i \in I} A_i$ 是环的有向余极限。
设 $0 \in I$，且 $\varphi_0 : B_0 \to C_0$ 是 $A_0$-代数的映射。
假设
\begin{enumerate}
\item $A \otimes_{A_0} B_0 \to A \otimes_{A_0} C_0$
是同构；
\item $B_0 \to C_0$ 是有限表示的。
\end{enumerate}
则对某个 $i \geq 0$，映射
$A_i \otimes_{A_0} B_0 \to A_i \otimes_{A_0} C_0$
是同构。
\end{lemma}

\begin{proof}
由引理 \ref{algebra-lemma-colimit-surjective}，存在 $i$，使得
$A_i \otimes_{A_0} B_0 \to A_i \otimes_{A_0} C_0$
是满射。由于此映射是有限表示的，其核是有限生成理想。设
$g_1, \ldots, g_r \in A_i \otimes_{A_0} B_0$ 生成该核。
于是可以选取 $i' \geq i$，使得 $g_j$ 在
$A_{i'} \otimes_{A_0} B_0$ 中映到零。这样，
$A_{i'} \otimes_{A_0} B_0 \to A_{i'} \otimes_{A_0} C_0$
就是同构。
\end{proof}

\begin{lemma}
\label{algebra-lemma-colimit-etale}
设 $A = \colim_{i \in I} A_i$ 是环的有向余极限。
设 $0 \in I$，且 $\varphi_0 : B_0 \to C_0$ 是 $A_0$-代数的映射。
假设
\begin{enumerate}
\item $A \otimes_{A_0} B_0 \to A \otimes_{A_0} C_0$ 是平展的；
\item $B_0 \to C_0$ 是有限表示的。
\end{enumerate}
则对某个 $i \geq 0$，映射
$A_i \otimes_{A_0} B_0 \to A_i \otimes_{A_0} C_0$
是平展的。
\end{lemma}

\begin{proof}
写成 $C_0 = B_0[x_1, \ldots, x_n]/(f_{1, 0}, \ldots, f_{m, 0})$。
记 $B_i = A_i \otimes_{A_0} B_0$、$C_i = A_i \otimes_{A_0} C_0$。
注意，
$C_i = B_i[x_1, \ldots, x_n]/(f_{1, i}, \ldots, f_{m, i})$，
其中 $f_{j, i}$ 是 $f_{j, 0}$ 在 $B_i$ 上多项式环中的像。
记 $B = A \otimes_{A_0} B_0$、$C = A \otimes_{A_0} C_0$。
注意，
$C = B[x_1, \ldots, x_n]/(f_1, \ldots, f_m)$，
其中 $f_j$ 是 $f_{j, 0}$ 在 $B$ 上多项式环中的像。假设是映射
$$
\text{d} :
(f_1, \ldots, f_m)/(f_1, \ldots, f_m)^2
\longrightarrow
\bigoplus C \text{d}x_k
$$
为同构。因此，对充分大的 $i$，可以找到元素
$$
\xi_{k, i} \in (f_{1, i}, \ldots, f_{m, i})/(f_{1, i}, \ldots, f_{m, i})^2
$$
使得 $\text{d}\xi_{k, i} = \text{d}x_k$ 在
$\bigoplus C_i \text{d}x_k$ 中成立。
此外，必要时增大 $i$ 后，可知
$\sum (\partial f_{j, i}/\partial x_k) \xi_{k, i} =
f_{j, i} \bmod (f_{1, i}, \ldots, f_{m, i})^2$，
因为这在极限中成立。于是这个 $i$ 符合要求。
\end{proof}


\begin{lemma}
\label{algebra-lemma-colimit-smooth}
设 $A = \colim_{i \in I} A_i$ 是环的有向余极限。
设 $0 \in I$，且 $\varphi_0 : B_0 \to C_0$ 是 $A_0$-代数的映射。
假设
\begin{enumerate}
\item $A \otimes_{A_0} B_0 \to A \otimes_{A_0} C_0$ 是光滑的；
\item $B_0 \to C_0$ 是有限表示的。
\end{enumerate}
则对某个 $i \geq 0$，映射
$A_i \otimes_{A_0} B_0 \to A_i \otimes_{A_0} C_0$
是光滑的。
\end{lemma}

\begin{proof}
写成 $C_0 = B_0[x_1, \ldots, x_n]/(f_{1, 0}, \ldots, f_{m, 0})$。
记 $B_i = A_i \otimes_{A_0} B_0$、$C_i = A_i \otimes_{A_0} C_0$。
注意，
$C_i = B_i[x_1, \ldots, x_n]/(f_{1, i}, \ldots, f_{m, i})$，
其中 $f_{j, i}$ 是 $f_{j, 0}$ 在 $B_i$ 上多项式环中的像。
记 $B = A \otimes_{A_0} B_0$、$C = A \otimes_{A_0} C_0$。
注意，
$C = B[x_1, \ldots, x_n]/(f_1, \ldots, f_m)$，
其中 $f_j$ 是 $f_{j, 0}$ 在 $B$ 上多项式环中的像。假设是映射
$$
\text{d} :
(f_1, \ldots, f_m)/(f_1, \ldots, f_m)^2
\longrightarrow
\bigoplus C \text{d}x_k
$$
为分裂单射。设
$\xi_k \in (f_1, \ldots, f_m)/(f_1, \ldots, f_m)^2$ 是满足
$\sum (\partial f_j/\partial x_k) \xi_k =
f_j \bmod (f_1, \ldots, f_m)^2$ 的元素。则对充分大的 $i$，
可以找到元素
$$
\xi_{k, i} \in (f_{1, i}, \ldots, f_{m, i})/(f_{1, i}, \ldots, f_{m, i})^2
$$
使得
$\sum (\partial f_{j, i}/\partial x_k) \xi_{k, i} =
f_{j, i} \bmod (f_{1, i}, \ldots, f_{m, i})^2$，
因为这在极限中成立。于是这个 $i$ 符合要求。
\end{proof}


\begin{lemma}
\label{algebra-lemma-colimit-lci}
设 $A = \colim_{i \in I} A_i$ 是环的有向余极限。
设 $0 \in I$，且 $\varphi_0 : B_0 \to C_0$ 是 $A_0$-代数的映射。
假设
\begin{enumerate}
\item $A \otimes_{A_0} B_0 \to A \otimes_{A_0} C_0$
是 syntomic 的（相应地，是相对全局完全交）；
\item $C_0$ 在 $B_0$ 上有限表示。
\end{enumerate}
则存在 $i \geq 0$，使得映射
$A_i \otimes_{A_0} B_0 \to A_i \otimes_{A_0} C_0$
是 syntomic 的（相应地，是相对全局完全交）。
\end{lemma}

\begin{proof}
假设 $A \otimes_{A_0} B_0 \to A \otimes_{A_0} C_0$
是相对全局完全交。
由引理 \ref{algebra-lemma-relative-global-complete-intersection-Noetherian}，
存在有限型 $\mathbf{Z}$-代数 $R$、环映射
$R \to A \otimes_{A_0} B_0$、相对全局完全交 $R \to S$ 和同构
$$
(A \otimes_{A_0} B_0) \otimes_R S
\longrightarrow
A \otimes_{A_0} C_0
$$
由于 $R$ 在 $\mathbf{Z}$ 上有限型（从而有限表示），存在
$i$ 和映射 $R \to A_i \otimes_{A_0} B_0$，提升映射
$R \to A \otimes_{A_0} B_0$，见引理
\ref{algebra-lemma-characterize-finite-presentation}。
再次应用同一引理，存在 $i' \geq i$，使得
$(A_i \otimes_{A_0} B_0) \otimes_R S \to A \otimes_{A_0} C_0$
来自映射
$(A_i \otimes_{A_0} B_0) \otimes_R S \to A_{i'} \otimes_{A_0} C_0$。
因此，以 $i'$ 替换 $i$ 后，可以假设上述映射来自
$A_i \otimes_{A_0} B_0$-代数映射
$$
(A_i \otimes_{A_0} B_0) \otimes_R S
\longrightarrow
A_i \otimes_{A_0} C_0
$$
由引理 \ref{algebra-lemma-colimit-isomorphism}，增大 $i$ 后，此映射是同构。
这完成了此种情形的证明，因为由引理
\ref{algebra-lemma-base-change-relative-global-complete-intersection}，
相对全局完全交的基变换仍是相对全局完全交。

\medskip\noindent
假设 $A \otimes_{A_0} B_0 \to A \otimes_{A_0} C_0$
是 syntomic 的。则存在元素 $g_1, \ldots, g_m$，它们在
$A \otimes_{A_0} C_0$ 中生成单位理想，并且
$A \otimes_{A_0} B_0 \to (A \otimes_{A_0} C_0)_{g_j}$
是相对全局完全交，见引理 \ref{algebra-lemma-syntomic}。
可以找到 $i$ 和映到 $g_j$ 的元素
$g_{i, j} \in A_i \otimes_{A_0} C_0$。
增大 $i$ 后，可以假设
$g_{i, 1}, \ldots, g_{i, m}$ 在
$A_i \otimes_{A_0} C_0$ 中生成单位理想。
由前一段的结果，再次增大 $i$ 后，可以假设映射
$A_i \otimes_{A_0} B_0 \to (A_i \otimes_{A_0} C_0)_{g_{i, j}}$
都是相对全局完全交。
于是，由引理 \ref{algebra-lemma-local-syntomic}
（以及已经用过的引理 \ref{algebra-lemma-syntomic}），
$A_i \otimes_{A_0} B_0 \to A_i \otimes_{A_0} C_0$
是 syntomic 的。
\end{proof}













\noindent
下面的引理是上述结果的一个应用，似乎不太适合放在其他地方。

\begin{lemma}
\label{algebra-lemma-fppf-fpqf}
设 $R \to S$ 是有限表示的忠实平坦环映射。
则存在交换图
$$
\xymatrix{
S \ar[rr] & & S' \\
& R \ar[lu] \ar[ru]
}
$$
其中 $R \to S'$ 拟有限、忠实平坦且有限表示。
\end{lemma}

\begin{proof}
第一步，将此引理化归到 $R$ 在 $\mathbf{Z}$ 上有限型的情形。
由引理 \ref{algebra-lemma-descend-faithfully-flat-finite-presentation}，
存在交换图
$$
\xymatrix{
S_0 \ar[r] & S \\
R_0 \ar[u] \ar[r] & R \ar[u]
}
$$
其中 $R_0$ 在 $\mathbf{Z}$ 上有限型，而 $S_0$ 在 $R_0$ 上
忠实平坦且有限表示，并且 $S = R \otimes_{R_0} S_0$。
若对环映射 $R_0 \to S_0$ 证明了引理，则由基变换，对
$R \to S$ 也得到引理，因为拟有限环映射的基变换仍拟有限，见
引理 \ref{algebra-lemma-quasi-finite-base-change}。（当然，还使用了平坦映射的
基变换仍平坦，以及有限表示映射的基变换仍有限表示。）

\medskip\noindent
假设 $R \to S$ 是有限表示的忠实平坦环映射，且 $R$ 是 Noether 的
（由前一段可以如此假设）。设 $W \subset \Spec(S)$ 是引理
\ref{algebra-lemma-finite-presentation-flat-CM-locus-open} 的开集。
由于 $R \to S$ 忠实平坦，映射
$\Spec(S) \to \Spec(R)$ 满射，见引理 \ref{algebra-lemma-ff-rings}。
由引理 \ref{algebra-lemma-generic-CM-flat-finite-presentation}，
映射 $W \to \Spec(R)$ 也满射。
因此，将 $S$ 替换为乘积
$S_{g_1} \times \ldots \times S_{g_m}$ 后，可以假设
$W = \Spec(S)$；这里使用了 $\Spec(R)$ 拟紧
（引理 \ref{algebra-lemma-quasi-compact}），以及映射
$\Spec(S) \to \Spec(R)$ 是开的
（命题 \ref{algebra-proposition-fppf-open}）。
假设 $\mathfrak p \subset R$ 是素理想。选取位于
$\mathfrak p$ 之上的素理想 $\mathfrak q \subset S$，它对应于
纤维环 $S \otimes_R \kappa(\mathfrak p)$ 的一个极大理想。
Noether 局部环
$\overline{S}_{\mathfrak q} = S_{\mathfrak q}/\mathfrak pS_{\mathfrak q}$
是 Cohen--Macaulay 的，设其维数为 $d$。可以在
$S_{\mathfrak q}$ 的极大理想中选取 $f_1, \ldots, f_d$，使它们在
$\overline{S}_{\mathfrak q}$ 中的像构成正则序列。
选取 $f_1, \ldots, f_d$ 的一个公分母
$g \in S$，其中 $g \not \in \mathfrak q$，并考虑 $R$-代数
$$
S' = S_g/(f_1, \ldots, f_d).
$$
由构造，存在位于 $\mathfrak p$ 之上的素理想
$\mathfrak q' \subset S'$，它对应于 $\mathfrak q$（通过
$S_g \to S'_g$）。同样由构造，环映射 $R \to S'$ 在
% P01-E1038
$\mathfrak q$ 处拟有限，因为局部环
% P01-E1039
$$
S'_{\mathfrak q'}/\mathfrak pS'_{\mathfrak q'} =
S_{\mathfrak q}/(f_1, \ldots, f_d) + \mathfrak pS_{\mathfrak q} =
\overline{S}_{\mathfrak q}/(\overline{f}_1, \ldots, \overline{f}_d)
$$
维数为零，见引理 \ref{algebra-lemma-isolated-point-fibre}。
同样由构造，$R \to S'$ 是有限表示的。
最后，由引理 \ref{algebra-lemma-grothendieck-regular-sequence}，局部环映射
$R_{\mathfrak p} \to S'_{\mathfrak q'}$ 平坦（这里使用了 $R$
是 Noether 的）。因此，由平坦性的开性
（定理 \ref{algebra-theorem-openness-flatness}）和拟有限性的开性
（引理 \ref{algebra-lemma-quasi-finite-open}），可以用
$gg'$ 替换 $g$，其中某个合适的 $g' \in S$ 满足
$g' \not \in \mathfrak q$，从而假设 $R \to S'$ 平坦且拟有限。
像 $\Spec(S') \to \Spec(R)$ 是开集并包含 $\mathfrak p$。
换言之，我们已经证明，存在引理陈述所述的环 $S'$
（可能忠实性除外），其像包含任意给定的素理想。
再次使用 $\Spec(R)$ 的拟紧性，可知有限多个这样的环的乘积符合要求。
\end{proof}
